% 第六册。直接编译本文件得到合订本；上、下册见“main - 6A.tex”“main - 6B.tex”。
\documentclass[b5paper, openany]{ctexbook}
% 六册共用的导言区。各册 main.tex 在 \documentclass 之后 \input 本文件，
% 再追加本册专用的宏包与命令。改动这里会影响全部 6 册，改完请逐册编译检查。
%
% 宏包加载顺序会影响结果（如 hyperref 应较晚加载），调整顺序前请先确认。

\usepackage[margin=2.5cm]{geometry}

% 脚注用带圈数字 ①②③……，每页重新编号
\usepackage{pifont}
\usepackage[perpage,symbol*]{footmisc}
\DefineFNsymbols{circled}{{\ding{192}}{\ding{193}}{\ding{194}}
{\ding{195}}{\ding{196}}{\ding{197}}{\ding{198}}{\ding{199}}{\ding{200}}{\ding{201}}}
\setfnsymbol{circled}

\usepackage{ulem}

\usepackage{amsmath,amsfonts,mathrsfs,amssymb}
\usepackage{graphicx}

\usepackage[font=bf,labelfont=bf,labelsep=quad]{caption}

\usepackage{tikz}

\usepackage{ntheorem}
\theoremseparator{\;}

\usepackage{blkarray}
\usepackage{bm}
\usepackage[colorlinks=true, linkcolor=black]{hyperref}

\usepackage{enumerate}

% 选做题：用 \staritem 代替 \item，编号照常递增，题号左侧加“*”（同原书）
\makeatletter
\newcommand{\staritem}{%
  \refstepcounter{\@enumctr}%
  \item[\llap{*}\csname label\@enumctr\endcsname]}
\makeatother

\theoremstyle{plain}
\theoremheaderfont{\normalfont\bfseries}
\theorembodyfont{\normalfont}

\usepackage[framemethod=tikz]{mdframed}

% ---- 例题、解答等环境 ----
\newtheorem{example}{\bf 例}[chapter]
\newenvironment{solution}{\noindent {\bf 解：}}{}
\newenvironment{analyze}{\noindent {\bf 分析：}}{}
\newenvironment{rmk}{\noindent {\bf 注意：}}{}
\newenvironment{note}{\noindent {\bf 说明：}}{}

\renewcommand{\proofname}{\bf 证明：}
\newenvironment{proof}{{\noindent \bf 证明：}}{}%{\hfill $\square$\par}

% ---- 常用数学命令 ----
\newcommand{\E}{\mathbb{E}}
\renewcommand{\Pr}{\mathbb{P}}
\newcommand{\EP}{\mathbb{E}^{\mathbb{P}}}
\newcommand{\EQ}{\mathbb{E}^{\mathbb{Q}}}
\newcommand{\dif}{\,{\rm d}}
\newcommand{\Var}{{\rm Var}}
\newcommand{\Cov}{{\rm Cov}}
\newcommand{\x}{\times}

% ---- “练习”框与带标题的框 ----
\usepackage{tcolorbox}
\tcbuselibrary{breakable}
\tcbuselibrary{most}

\newtcolorbox{ex}[1][]
  {colback = white, colframe = cyan!75!black, fonttitle = \bfseries,
    colbacktitle = cyan!85!black, enhanced,
    attach boxed title to top center={yshift=-2mm},breakable,
    title=练习, #1}

\newtcolorbox{blk}[2][]
  {colback = white, colframe = magenta!75!black, fonttitle = \bfseries,
    colbacktitle = magenta!85!black, enhanced,
    attach boxed title to top left={xshift=5mm, yshift=-2mm},breakable,
    title=#2, #1}

% ---- 章节编号：第一节 / 一、/（一） ----
\setcounter{tocdepth}{2}
\setcounter{secnumdepth}{3}

\ctexset {
section = {
	name = {第,节},
 	number = \chinese{section}},
subsection = {
	name = {,、\hspace{-1em}},
	number = \chinese{subsection}
},
subsubsection = {
	name = {（,）\hspace{-1em}},
	number = \chinese{subsubsection},
}
}

\renewcommand{\contentsname}{目~~录}

\newcommand{\poly}{\polynomial[reciprocal]}

\usepackage{mathtools}

\setlength{\abovecaptionskip}{0.cm}
\setlength{\belowcaptionskip}{-0.cm}

\usetikzlibrary{decorations.pathmorphing, patterns}
\usetikzlibrary{calc, patterns, decorations.markings}
\usetikzlibrary{positioning, snakes}

\usepackage{yhmath}   % \wideparen
\usepackage{longdivision}
\usepackage{polynom}
\usepackage{polynomial}
\usepackage{cancel}

\renewcommand{\frac}{\dfrac}
\newcommand{\oc}{$^{\circ}{\rm C}$}

\usepackage{multicol}
\usepackage{cases}

% ---- 合订本 / 上册 / 下册 ----
% “main - nA.tex”“main - nB.tex”先 \def\BookPart{A} 或 {B}，再 % 第六册。直接编译本文件得到合订本；上、下册见“main - 6A.tex”“main - 6B.tex”。
\documentclass[b5paper, openany]{ctexbook}
% 六册共用的导言区。各册 main.tex 在 \documentclass 之后 \input 本文件，
% 再追加本册专用的宏包与命令。改动这里会影响全部 6 册，改完请逐册编译检查。
%
% 宏包加载顺序会影响结果（如 hyperref 应较晚加载），调整顺序前请先确认。

\usepackage[margin=2.5cm]{geometry}

% 脚注用带圈数字 ①②③……，每页重新编号
\usepackage{pifont}
\usepackage[perpage,symbol*]{footmisc}
\DefineFNsymbols{circled}{{\ding{192}}{\ding{193}}{\ding{194}}
{\ding{195}}{\ding{196}}{\ding{197}}{\ding{198}}{\ding{199}}{\ding{200}}{\ding{201}}}
\setfnsymbol{circled}

\usepackage{ulem}

\usepackage{amsmath,amsfonts,mathrsfs,amssymb}
\usepackage{graphicx}

\usepackage[font=bf,labelfont=bf,labelsep=quad]{caption}

\usepackage{tikz}

\usepackage{ntheorem}
\theoremseparator{\;}

\usepackage{blkarray}
\usepackage{bm}
\usepackage[colorlinks=true, linkcolor=black]{hyperref}

\usepackage{enumerate}

% 选做题：用 \staritem 代替 \item，编号照常递增，题号左侧加“*”（同原书）
\makeatletter
\newcommand{\staritem}{%
  \refstepcounter{\@enumctr}%
  \item[\llap{*}\csname label\@enumctr\endcsname]}
\makeatother

\theoremstyle{plain}
\theoremheaderfont{\normalfont\bfseries}
\theorembodyfont{\normalfont}

\usepackage[framemethod=tikz]{mdframed}

% ---- 例题、解答等环境 ----
\newtheorem{example}{\bf 例}[chapter]
\newenvironment{solution}{\noindent {\bf 解：}}{}
\newenvironment{analyze}{\noindent {\bf 分析：}}{}
\newenvironment{rmk}{\noindent {\bf 注意：}}{}
\newenvironment{note}{\noindent {\bf 说明：}}{}

\renewcommand{\proofname}{\bf 证明：}
\newenvironment{proof}{{\noindent \bf 证明：}}{}%{\hfill $\square$\par}

% ---- 常用数学命令 ----
\newcommand{\E}{\mathbb{E}}
\renewcommand{\Pr}{\mathbb{P}}
\newcommand{\EP}{\mathbb{E}^{\mathbb{P}}}
\newcommand{\EQ}{\mathbb{E}^{\mathbb{Q}}}
\newcommand{\dif}{\,{\rm d}}
\newcommand{\Var}{{\rm Var}}
\newcommand{\Cov}{{\rm Cov}}
\newcommand{\x}{\times}

% ---- “练习”框与带标题的框 ----
\usepackage{tcolorbox}
\tcbuselibrary{breakable}
\tcbuselibrary{most}

\newtcolorbox{ex}[1][]
  {colback = white, colframe = cyan!75!black, fonttitle = \bfseries,
    colbacktitle = cyan!85!black, enhanced,
    attach boxed title to top center={yshift=-2mm},breakable,
    title=练习, #1}

\newtcolorbox{blk}[2][]
  {colback = white, colframe = magenta!75!black, fonttitle = \bfseries,
    colbacktitle = magenta!85!black, enhanced,
    attach boxed title to top left={xshift=5mm, yshift=-2mm},breakable,
    title=#2, #1}

% ---- 章节编号：第一节 / 一、/（一） ----
\setcounter{tocdepth}{2}
\setcounter{secnumdepth}{3}

\ctexset {
section = {
	name = {第,节},
 	number = \chinese{section}},
subsection = {
	name = {,、\hspace{-1em}},
	number = \chinese{subsection}
},
subsubsection = {
	name = {（,）\hspace{-1em}},
	number = \chinese{subsubsection},
}
}

\renewcommand{\contentsname}{目~~录}

\newcommand{\poly}{\polynomial[reciprocal]}

\usepackage{mathtools}

\setlength{\abovecaptionskip}{0.cm}
\setlength{\belowcaptionskip}{-0.cm}

\usetikzlibrary{decorations.pathmorphing, patterns}
\usetikzlibrary{calc, patterns, decorations.markings}
\usetikzlibrary{positioning, snakes}

\usepackage{yhmath}   % \wideparen
\usepackage{longdivision}
\usepackage{polynom}
\usepackage{polynomial}
\usepackage{cancel}

\renewcommand{\frac}{\dfrac}
\newcommand{\oc}{$^{\circ}{\rm C}$}

\usepackage{multicol}
\usepackage{cases}

% ---- 合订本 / 上册 / 下册 ----
% “main - nA.tex”“main - nB.tex”先 \def\BookPart{A} 或 {B}，再 \input{main.tex}；
% 直接编译 main.tex 时 \BookPart 未定义，得到合订本。main.tex 中用：
%   \ifPartA ... \fi      合订本和上册包含的内容
%   \ifPartB ... \fi      合订本和下册包含的内容
%   \ifFullBook ... \fi   只有合订本包含的内容
%   \BookPartSuffix       书名后缀：合订本为空，上、下册为“（上）”“（下）”
%   \PartChoice{合订本}{上册}{下册}   按版本取不同的文字（如出版日期）
\newif\ifFullBook
\newif\ifPartA
\newif\ifPartB
\ifdefined\BookPart\else\def\BookPart{}\fi
\def\BookPartNameA{A}
\def\BookPartNameB{B}
\ifx\BookPart\BookPartNameA
  \PartAtrue \def\BookPartSuffix{（上）}
\else\ifx\BookPart\BookPartNameB
  \PartBtrue \def\BookPartSuffix{（下）}
\else
  \FullBooktrue \PartAtrue \PartBtrue \def\BookPartSuffix{}
\fi\fi
\newcommand{\PartChoice}[3]{\ifFullBook#1\else\ifPartA#2\else#3\fi\fi}


% ---- 本册专用 ----
% 注意：不要加载 physics 宏包。它会把 \div 重定义为散度（上册用 \div 表示除号），
% 下册用到的 \dd、\arccot、\arcsec 已在下方单独定义。
\usepackage{pstricks}
\usepackage{bookmark}
\usepackage{circuitikz}
\usepackage{longtable}
\usepackage{tkz-base, tkz-euclide}

\newcommand{\dd}{\,{\rm d}}
\DeclareMathOperator{\arccot}{arccot}\DeclareMathOperator{\arcsec}{arcsec}
\newcommand{\Lim}{\displaystyle\lim}
\newcommand{\Int}{\displaystyle\int}

\newcommand{\NC}{\text{N}/\text{C}}
\newcommand{\ms}{\text{m}/\text{s}}
\newcommand{\cms}{\text{cm}/\text{s}}
\newcommand{\msq}{\text{m}/\text{s}^2}
\newcommand{\cmsq}{\text{cm}/\text{s}^2}
\newcommand{\kmh}{\text{km}/\text{h}}

\newcommand{\R}{\mathbb{R}}
\newcommand{\N}{\mathbb{N}}
\newcommand{\Z}{\mathbb{Z}}

\newcommand{\VEC}{\overrightarrow}

\newcommand{\comb}[2]{{\rm C}^{#1}_{#2}}
\newcommand{\perm}[2]{{\rm P}^{#1}_{#2}}
\newcommand{\rep}[2]{{\rm H}^{#1}_{#2}}
\newcommand{\pc}[1]{\cos{#1}+i\sin{#1}}
\newcommand{\pcx}[1]{\cos\left(#1\right)+i\sin\left(#1\right)}

% 第六册上、下两本的章号各自从“第一章”编起（正文中的“图2.33”“第一章”等引用均按此编号）。
% 合订本用 \volume 分隔两本；上、下册单独编译时不需要。
% \volume{标题}：开始新的一本，章号归零；同时让 hyperref 锚点按“本-章”区分，避免重复。
% \appendixchapters：之后的 \chapter 不编号（用于上册末尾的附录），且不重置页码。
\ctexset{part = {numbering = false}}
\newcounter{volume}
\renewcommand{\theHchapter}{\arabic{volume}.\arabic{chapter}}
\makeatletter
\newcommand{\volume}[1]{\@mainmattertrue\stepcounter{volume}\setcounter{chapter}{0}\part{#1}}
\newcommand{\appendixchapters}{\@mainmatterfalse}
\makeatother

\begin{document}

\title{\Huge\bfseries 中学数学实验教材\\第六册\BookPartSuffix}
\author{\Large 中学数学实验教材编写组编}
\date{\Large 1986年5月}
\maketitle

\frontmatter
\chapter{前~~言}

这一套中学数学实验教材，内容的选取原则是精简实
用，教材的处理力求深入浅出，顺理成章，尽量作到使人人
能懂，到处有用．

    本教材适用于重点中学，侧重在满足学生将来从事理工
方面学习和工作的需要．

    本教材的教学目的是：使学生切实学好从事现代生产、
特别是学习现代科学技术所必需的数学基础知识；通过对数
学理论、应用、思想和方法的学习，培养学生运算能力，思
维能力，空间想象力，从而逐步培养运用数学的思想和方法
去分析和解决实际问题的能力；通过数学的教学和学习，培
养学生良好的学习习惯，严谨的治学态度和科学的思想方
法，逐步形成辩证唯物主义世界观．

   根据上述教学目的，本教材精选了传统数学那些普遍实
用的最基础的部分，这就是在理论上、应用上和思想方法上
都是基本的、长远起作用的通性、通法．比如，代数中的数
系运算律，式的运算，解代数方程，待定系数法；几何中的
图形的基本概念和主要性质，向量，解析几何；分析中的函
数，极限，连续，微分，积分；概率统计以及逻辑、推理论
证等知识．对于那些理论和应用上虽有一定作用，但发展余
地不大，或没有普遍意义和实用价值，或不必要的重复和过
于繁琐的内容，如立体几何中的空间作图，几何体的体积、
表面积计算，几何难题，因式分解，对数计算等作了较大的
精简或删减．

    全套教材共分六册．第一册是代数．在总结小学所学自
然数、小数、分数基础上，明确提出运算律，把数扩充到有
理数和实数系．灵活运用运算律解一元一次、二次方程，二
元、三元一次方程组，然后进一步系统化，引进多项式运
算，综合除法，辗转相除，余式定理及其推论，学到根式、
分式、部分分式．第二册是几何．由直观几何形象分析归纳
出几何基本概念和基本性质，通过集合术语、简易逻辑转入
欧氏推理几何，处理直线形，圆、基本轨迹与作图，三角比
与解三角形等基本内容．第三册是函数．数形结合引入坐
标，研究多项式函数，指数、对数、三角函数，不等式等．
第四册是代数．把数扩充到复数系，进一步加强多项式论，
方程式论，讲线性方程组理论，概率（离散的）统计的初步
知识．第五册是几何．引进向量，用向量和初等几何方法综
合处理几何问题，坐标化处理直线、圆、锥线，坐标变换与
二次曲线讨论，然后讲立体几何，并引进空间向量研究空间
解析几何初步知识．第六册是微积分初步．突出逼近法，讲
实数完备性，函数，极限，连续，变率与微分，求和与积分．

本教材基本上采取代数、几何、分析分科，初中、高中
循环排列的安排体系．教学可按初一、初二代数、几何双科
并进，初三学分析，高一、高二代数（包括概率统计）、几
何双科并进，高三学微积分的程序来安排．

    本教材的处理力求符合历史发展和认识发展的规律，深
入浅出，顺理成章．突出由算术到代数，由实验几何到论证
几何，由综合几何到解析几何，由常量数学到变量数学等四
个重大转折，着力采取措施引导学生合乎规律地实现这些转
折，为此，强调数系运算律，集合逻辑，向量和逼近法分别
在实现这四个转折中的作用．这样既遵循历史发展的规律，
又突出了几个转折关头，缩短了认识过程，有利于学生掌握
数学思想发展的脉络，提高数学教学的思想性．

这一套中学数学实验教材是教育部委托北京师范大学、
中国科学院数学研究所、人民教育出版社、北京师范学院、
北京景山学校等单位组成的领导小组组织“中学数学实验教
材编写组”，根据美国加州大学伯克利分校数学系项武义教
授的《关于中学实验数学教材的设想》编写的．第一版印出
后，由教育部实验研究组和有关省市实验研究组指导在北
京景山学校、北京师院附中、上海大同中学、天津南开中
学、天津十六中学、广东省实验中学、华南师院附中、长春
市实验中学等校试教过两遍，在这个基础上编写组吸收了实
验学校老师们的经验和意见，修改成这一版《中学数学实验
教材》，正式出版，内部发行，供中学选作实验教材，教师
参考书或学生课外读物．在编写和修订过程中，项武义教授
曾数次详细修改过原稿，提出过许多宝贵意见．

    本教材虽然试用过两遍，但是实验基础仍然很不够，这
次修改出版，目的是通过更大范围的实验研究，逐步形成另
一套现代化而又适合我国国情的中学数学教科书．在实验过
程中，我们热忱希望大家多提意见，以便进一步把它修改好．

\begin{flushright}
    中学数学实验教材编写组\\
    一九八五年三月
\end{flushright}








\tableofcontents

\mainmatter

\ifPartA
\ifFullBook\volume{第六册（上）}\fi
\chapter{复数系}

在数学中，数是最基础、应用最广泛的工具之一，我们已经在自然数的基础上，逐步学习了整数、有理数、实数等概念和它们的运算及性质.这一章，我们将进一步学习复数的概念及其运算、性质.

\section{复数的概念}
\subsection{数的概念的发展}

数的概念是中学数学的主要内容之一，是从实践中逐步形成和发展起来的.早在有文字记载的历史以前，由于计数和排序的需要，人类就建立起自然数（正整数）的概念，处理着各种常见的、有天然单元的可数量. 自然数的全体构成了自然数集$\mathbb{N}$.

随着社会生产和科学技术的发展，数的概念相应地得到了不断发展.

为了处理实践中遇到的测量、分割、分配等可以细分的量，人们引进了正分数；为了解决生产实践中的各种具有相反意义的量的表示和运算问题，人们又引进了零及负数，并把自然数看作正整数. 这样一来，把正整数、零、负整数合并在一起，称为整数集$\mathbb{Z}$；并把整数与分数（正分数、负分数）合称为有理数集$\mathbb{Q}$.

任一有理数，可以表示成分数形式$\frac{m}{n}$，其中$m,n\in\mathbb{N}$; 还可以表示成循环小数（包括循环节为0的小数）.

有理数系的形成，使得许多人曾认为这就足以解决可度量的量，特别是在古希腊，他们曾认为：任意两条线段都是可公度的（即：任意两条线段长度的比都是有理数）. 但实际上这种认识是不正确的. 例如，正方形的边长和对角线长，正五边形的边长和对角线长都是不可公度的. 为了解决这些量与量之比值不能用有理数表示的矛盾，于是又引进了无理数的概念，无理数就是无限不循环小数. 无理数与有理数统称为实数，构成了实数集$\mathbb{R}$.

实数集$\mathbb{R}$对加、减、乘、除（零不作除数）四则运算是封闭的，且满足运算通性，即满足加、乘法的交换律、结合律和分配律，数零与1保持在运算中的特性．

实数集$\mathbb{R}$与数轴上的点集之间，可以建立起一一对应的关系. 因此，从度量的角度来说，实数系是一个完整无缺的数系，它足以满足度量的实际需要.

但是，从代数运算和解方程的角度来看，实数集并不完整无缺，因为，在实数范围内，开方运算并不封闭，实系数方程（如$ax^2+bx+c=0$）在实数系内未必有解.

事实上，数的概念的逐步扩充，与代数运算和解方程的需要也是相适应的，例如：

由于自然数集$\mathbb{N}$对于减法、除法运算不封闭，因而方程$x+7=2$, $7x=5$等在自然数集中无解，但在有理数集$\mathbb{Q}$中就有解$x=-5$, $x=\frac{5}{7}$. 

由于有理数集$\mathbb{Q}$对于开方运算不封闭，因而方程$x^2-3
=0$在$\mathbb{Q}$中无解，但在实数集$\mathbb{R}$中就有两个解$x=\pm\sqrt{3}$.

同样，实数集$\mathbb{R}$对开方运算仍不封闭，因而像方程$x^2+1=0$在$\mathbb{R}$中还是无解，因为，没有一个实数的平方等于$-1$．要解决这一类矛盾，数的概念需要进一步发展，这就需要引进新数，并讨论它的运算.

\begin{ex}
\begin{enumerate}
    \item 试叙述实数集$\mathbb{R}$中有哪些数系运算通性？
\item 试分别在自然数集$\mathbb{N}$、整数集$\mathbb{Z}$、有理数集$\mathbb{Q}$和实数集$\mathbb{R}$中，解方程$(x-2)(3x+1)(x+7)(x^2+5)=0$.
\end{enumerate}
\end{ex}

\subsection{复数及其几何表示}
如上所述，由于解方程的需要，早在十六世纪，人们就引进了一个新数$i$，使得方程$x^2+1=0$有解，并把它叫做\textbf{虚数单位}，并规定：
\begin{enumerate}[(1)]
\item $i^2=-1$;
\item 它与实数进行四则运算时，运算通性及数0与1的运算特性仍然保持. 即，满足加、乘法的交换律、结合律、分配律；满足
\[\begin{split}
    0+i&=i+0=i,\\
0\cdot i&=i\cdot 0=0,\\
1\cdot i&=i\cdot 1=i.
\end{split}\]
\end{enumerate}

在这样的规定下，$i$可以与实数$b$相乘，再与实数$a$相加，从而得到$a+bi$.

这样，数的概念扩大了，我们就把形如$a+bi\; (a,b\in\mathbb{R})$的数，叫做复数.全体复数构成的集合，称为\textbf{复数集}，
一般记作$\mathbb{C}$.

对于复数$a+bi\; (a,b\in\mathbb{R})$，当$b=0$时，就是实数$a$；当$b\ne 0$时，这一复数叫做虚数，特别地，当$a=0$且$b\ne 0$时，叫做纯虚数.

$a$，$b$分别叫做复数$z=a+bi$的\textbf{实部}与\textbf{虚部}，并且记作：${\rm Re}(z)=a$, ${\rm Im}(z)=b$.

显然，实数集$\mathbb{R}$是复数集$\mathbb{C}$的真子集，即$\mathbb{R}\subset \mathbb{C}$.

例如，$3+2i$, $-\frac{1}{2}-\sqrt{3}i$, $-0.25i$都是虚数. 其中$-0.25i$是纯虚数．它们的实部和虚部分别为：
\[\begin{split}
&{\rm Re}(3+2i)=3,\qquad {\rm Im}(3+2i)=2;\\
&{\rm Re}\left(-\frac{1}{2}-\sqrt{3}i\right)=-\frac{1}{2},\qquad {\rm Im}\left(-\frac{1}{2}-\sqrt{3}i\right)=-\sqrt{3};\\
&{\rm Re}(-0.25i)=0,\qquad {\rm Im}(-0.25i)=-0.25.\\
\end{split}\]

我们还规定，当且仅当两个复数的实部和虚部分别相等时，这两个复数相等，即
\[a+bi=c+di\Longleftrightarrow a=c\text{且}b=d\]
或
\[z_1=z_2\Longleftrightarrow {\rm Re}(z_1)={\rm Re}(z_2)\quad \text{且} \quad {\rm Im}(z_1)={\rm Im}(z_2)\]
特别地，我们可以得出
\[a+bi=0\Longleftrightarrow a=b=0\]

\begin{example}
    若$3x+2yi-1=(x-1)i+y$，其中$x,y\in\mathbb{R}$. 求$x$与$y$的值.
\end{example}

\begin{solution}
将已知等式整理得
\[(3x-1)+2yi=y+(x-1)i\]
根据复数相等的定义，可得方程组$\begin{cases}
    3x-1=y\\
    2y=x-1
\end{cases}$，
解得
\[x=\frac{1}{5},\qquad y=-\frac{2}{5}\]
\end{solution}

当两个复数的实部相等、虚部互为相反数时，这两个复数叫做\textbf{互为共轭复数}，即
：复数$a+bi$与$a-bi$叫做互为共轭复数. 一般地，复数$z=a+bi$的共轭复数表示为：
$$\overline{z}=a-bi$$

显然，实数$a$（即虚部为零的复数）的共轭复数仍是它自身$a$.

我们已经知道，实数集与数轴上的点之间可建立一一对应关系. 引进虚数，构成复数集之后，是否有相应的对应关系呢？

首先，虚数单位$i$在实数轴上没有立足之地，它不属于实数集；其次，当$b\ne 0$时，复数$a+bi$也不可能在实数轴上找出对应之点，它仍不属于实数集，因此，我们必须另找途径来建立对应关系.

从复数及其相等的定义中，我们知道，任何一个复数
$z=a+bi$，都由一个有序实数对$(a,b)$唯一确定，而每一个有序实数对，又借助平面的坐标化可以唯一确定平面上一个点，因此，只要在平面上建立起直角坐标系，就可以用平面上的点$Z(a,b)$表示出复数$z=a+bi$，如图1.1所示.

\begin{minipage}{.45\textwidth}
\centering
\begin{tikzpicture}[>=stealth]
\draw[->](-1,0)--(3,0)node[right]{$x$};
\draw[->](0,-1)--(0,3)node[right]{$y$};
\draw[dashed](0,1.5)node[left]{$b$}--(2,1.5)node[above]{$Z(a,b)$}--(2,0)node[below]{$a$};
\node[below right]{$O$};
\draw[fill](2,1.5)circle(1pt);
\end{tikzpicture}
\captionof{figure}{ }
\end{minipage}
\hfill
\begin{minipage}{.45\textwidth}
\centering
    \begin{tikzpicture}[>=stealth]
  \draw[->](-1,0)--(3,0)node[right]{$x$};
\draw[->](0,-2)--(0,2)node[right]{$y$};
\draw[dashed](0,1.25)node[left]{$b$}--(2,1.25)node[above]{$Z(a,b)$}--(2,0)node[below right]{$a$}--(2,-1.25)node[below]{$\overline{Z}(a,-b)$}--(0,-1.25)node[left]{$-b$};
\node[below right]{$O$};  
\draw[fill](2,1.25)circle(1pt);
\draw[fill](2,-1.25)circle(1pt);
    \end{tikzpicture}
\captionof{figure}{ }
\end{minipage}


这个建立了坐标系来表示复数的平面，叫做\textbf{复平面}，其中$X$轴叫做\textbf{实轴}，$Y$轴（不含原点）叫做\textbf{虚轴}.

这样一来，每一个复数，在复平面内都有唯一一点和它对应，反之，复平面内的每一个点，都有唯一的一个复数和它对应，因此，复数集$\mathbb{C}$和复平面上的点集是一一对应的. 特别地，实数集和复平面实轴上的点集一一对应；纯虚数集合和复平面虚轴（不含原点）上的点集一一对应.

显然，互为共轭的两个复数$z=a+bi$与$\bar z=a-bi$在复平面上对应的两点$Z(a,b)$与$\overline{Z}(a,-b)$关于实轴对称（图1.2）．

\begin{ex}
\begin{enumerate}
    \item 将下列给出的各数，分类填入表中：
\[1-\sqrt{2},\quad 0.618,\quad -1+i,\quad \frac{1}{4}i,\quad 0,\quad i,\quad i^2,\quad 3+7\sqrt{2}i\]
\[i-1,\quad (\sqrt{2}+\sqrt{3})i,\quad \sqrt{2}i-\sqrt{2},\quad -\frac{\sqrt{3}}{2}i\]
\begin{tabular}{|c|p{.7\textwidth}|}
\hline
    实数&\\
\hline
虚数&\\\hline
纯虚数&\\    \hline
\end{tabular}

\item 填写下列表格：
\begin{center}
\begin{tabular}{|c|p{.1\textwidth}|p{.1\textwidth}|p{.1\textwidth}|p{.1\textwidth}|p{.1\textwidth}|}
\hline 
复数$z$& ${\rm Re}(z)$& ${\rm Im}(z)$&$\bar z$& ${\rm Re}(\bar z)$& ${\rm Im}(\bar z)$\\
\hline
$-\sqrt{3}+i$ &&&&&\\
$-(1+\sqrt{2})$&&&&&\\
$\frac{i-\sqrt{2}}{4}$&&&&&\\
$0$&&&&&\\
$i$&&&&&\\
$a-bi$&&&&&\\
\hline
\end{tabular}
\end{center}
\item 试求出满足下列等式的实数$x$，$y$：
\begin{enumerate}[(1)]
\item $(3x+2y)+(5x-y)i=17-2i$;
\item $3x+2y-7=(2x+3y-8)i$;
\item $7x+(4y-1)i-1=0$.
\end{enumerate}

\item 试说明对任意复数$z$，都有$\overline{\bar z}=z$.
\item 在复平面上，画出下列复数及它们的共轭复数所对应的点：
\[2-2i;\quad -2+2i;\quad -2i; \quad \sqrt{2};\quad 0.\]
\end{enumerate}
\end{ex}

\subsection{复数的模与幅角}
我们已经知道，复数集与复平面上的点集之间是一一对应的，即每一复数$z=a+bi\; (a,b\in\R)$都对应着复平面
上唯一一点$Z(a,b)$，反之亦然。

如果我们在复平面上连结原点$O$与表示复数$z=a+bi$的点$Z(a,b)$，就得到一个向量（从$O$点指向$Z$点），记作$\VEC{OZ}$，这样就可以在复数同向量（从原点出发的向量）之间
建立起联系. 不难看出，任一复数$z$可唯一确定复平面上一点$Z$，进而唯一确定一个向量$\VEC{OZ}$；反之，复平面上任一向量$\VEC{OZ}$，可唯一确定一点$Z$，进而唯一确定一个复数$z$，因此，复数集$\mathbb{C}$与复平面内所有以原点$O$为起点的向量（位置向量）集合也是一一对应的。

这样一来，我们就可以很方便地把复数$z=a+bi\; (a,b\in\R)$说成点$Z(a,b)$，或说成向量$\VEC{OZ}$. 研究复数的运算、性质就可以和研究向量的运算、性质互通了.

如图1.3所示，我们把向量$\VEC{OZ}$的模（$OZ$的长度）$r$叫做复数$z=a+bi$的模（或绝对值），记作
\[r=|z|\quad \text{或}\quad r=|a+bi|\]
\begin{figure}[htp]
    \centering
\begin{tikzpicture}[>=stealth]
\draw[->](-1,0)--(3,0)node[right]{$x$};
\draw[->](0,-1)--(0,3)node[right]{$y$};
\draw[dashed](0,1.5)node[left]{$b$}--(2,1.5)node[above]{$Z(a,b)$}--(2,0)node[below]{$a$};
\draw[very thick,->](0,0)--node[left]{$r$}(2,1.5);
\draw(.5,0) arc (0:35:.5)node[right]{$\theta$};
\node[below right]{$O$};
\end{tikzpicture}
    \caption{}
\end{figure}

显然
\[r=|z|=|a+bi|=\sqrt{a^2+b^2}\ge 0\]

特别地，当$b=0$时，$z$是一个实数，这时
\[r=|z|=|a|\qquad  \text{（与实数的绝对值概念一致）
}\]

又以实轴$OX$的正半轴为始边，向量$\VEC{OZ}$所在的射线为终边的角$\theta$，叫做复数$z=a+bi$的\textbf{幅角}.

由于任一不等于零的复数，按以上规定可有无限多个幅
角（终边相同的角）的值，因而，我们进一步约定：

满足$0\le\theta<2\pi$的幅角$\theta$的值，叫做幅角的\textbf{主值}，并记作$\arg z$，即$0\le \arg z<2\pi$. 
因此，任一非零复数$z$的幅角$\theta=\arg z+2k\pi\; (k\in\Z)$

这样一来，每一个不等于零的复数，就只有唯一的模与唯一的幅角主值，并可被它的模与幅角主值所唯一确定。

应该指出，由于复数$z=0$对应于复平面上的零向量（一个点），而且零向量的方向是任意的，因此，复数$z=0$的模$|z|=0$，它的幅角却是任意的。

\begin{example}
试比较复数$z_1=1+\sqrt{2}i$，$z_2=-\frac{1}{2}-i$的模的大小.
\end{example}

\begin{solution}
$\because\quad |z_1|=\sqrt{1+2}=\sqrt{3},\quad |z_2|=\sqrt{\frac{1}{4}+1}=\frac{1}{2}\sqrt{5},\qquad \sqrt{3}>\frac{1}{2}\sqrt{5}$

$\therefore\quad |z_1|>|z_2|$.
\end{solution}

一般来说，两个复数中，至少有一个不是实数，就不能比较它们的大小，但任意两个复数的模，总是可以比较大小的.

\begin{example}
    如果$a$，$b$都是正实数，试求：
\begin{multicols}{4}
\begin{enumerate}[(1)]
    \item $\arg a$
    \item $\arg (-a)$
    \item $\arg (bi)$
    \item $\arg (-bi)$
\end{enumerate}
\end{multicols}
\end{example}

\begin{solution}
由于$a$，$b$是正实数，因此：
复数$a$，$-a$分别对应着复平面中正、负半实轴上的点（不包含原点）. 所以
\begin{multicols}{2}
    \begin{enumerate}[(1)]
        \item $\arg (a)=0$
        \item $\arg (-a)=\pi$
    \end{enumerate}
    \end{multicols}
复数$bi$，$-bi$分别对应着复平面中正、负半虚轴上的点（不包含原点）. 所以
\begin{multicols}{2}
    \begin{enumerate}[(1)]\setcounter{enumi}{2}
        \item $\arg (bi)=\frac{\pi}{2}$
        \item $\arg (-bi)=\frac{3\pi}{2}$
    \end{enumerate}
    \end{multicols}
\end{solution}

\begin{example}
    设$z\in\mathbb{C}$，试说明满足下列条件的点Z的集合组成什么样的图形？
\begin{multicols}{2}
\begin{enumerate}[(1)]
    \item $|z|=a$
    \item $2<|z|\le 4$
\end{enumerate}
\end{multicols}
\end{example}

\begin{solution}
\begin{enumerate}[(1)]
    \item 复数$z$的模等于$a\ge 0$，就是说它对应于复平面上的向量$\VEC{OZ}$的长度为$a$，也就是说它对应于复平面上的点$Z$到原点的距离为$a$. 所以，满足条件$|z|=a$的点的集合，当$a>0$时组成一个以$O$为圆心、以$a$为半径的圆；当$a=0$时，就是一个点（原点）。
\item 先将条件$2<|z|\le 4$化为$\begin{cases}
    |z|>2\\ |z|\le 4
\end{cases}$

\begin{minipage}{.45\textwidth}
 再分别考虑：满足$|z|>2$的点集组成圆$|z|=2$的外部；满足$|z|\le 4$的点集组成圆$|z|=4$的内部（包含边界）。同时满足以上两条件的点集，就是已求两集合的交集，也就是所要求的集合.

因此，满足$2<|z|\le 4$的点的集合是以$O$为圆心，以2及4为半径的两圆所夹的圆环，其中包含外边界，但不包括内边界（图1.4）.   
\end{minipage}\hfill
\begin{minipage}{.45\textwidth}
    \centering
\begin{tikzpicture}[>=stealth, scale=.4]

\draw[ thick, pattern=north east lines](0,0) circle (4);
\fill[color=white](0,0) circle (2);    
\draw[dashed,  thick](0,0) circle (2);    
\draw[->](-6,0)--(6,0)node[right]{$x$};
\draw[->](0,-6)--(0,6)node[right]{$y$};
\node [below right]{$O$};
\node at (2,0)[below right]{2};
\node at (4,0)[below right]{4};
\node at (0,2)[below right]{2};
\node at (0,4)[above right]{4};

\end{tikzpicture}
    \captionof{figure}{ }
\end{minipage}

\end{enumerate}    
\end{solution}

\begin{ex}
\begin{enumerate}
    \item 求出下列复数的模，并按其大小用不等号连起来：
\[\sqrt{2}-i,\quad -2+i,\quad 7,\quad 2i,\quad -1-i\]
  \item   设$|z|=r$, $\arg z=\alpha$，试求出$|\bar z|$, $\arg \bar z$.
  \item 求证：复数$z_1=-3i$, $z_2=\sqrt{2}+\sqrt{7}i$, $z_3=\sqrt{7}-\sqrt{2}i$, $z_4=-2+\sqrt{5}i$所对应的四个点共圆，并写出这个圆周上的所有点应满足的条件.
\item   你能指出满足下列条件的复数所对应的点集组成什么图形吗？
\begin{multicols}{2}
\begin{enumerate}[(1)]
    \item $\arg z=\frac{\pi}{4}$
    \item $\frac{\pi}{4}\le \arg z\le \frac{\pi}{2}$
\end{enumerate}
\end{multicols}
\end{enumerate}

\end{ex}

\subsection*{习题1.1}
\begin{enumerate}
    \item 设复数$z= (a^{2}- 2a- 3 )+( a^{2}+ 2a- 3) i$, $a\in \mathbb{R} $, 
试问在什么条件下，这一复数是
\begin{multicols}{4}
\begin{enumerate}[(1)]
    \item 实 数 ？
    \item 零？
    \item 虚数？
    \item 纯虚数？
\end{enumerate}
\end{multicols}

\item 实数$a,b$取何值时，才能使复数$z_1=z_2$:
\begin{enumerate}[(1)]
    \item $z_{1}= \left(\frac{1}{2}a+ b \right)+\left( 5a+\frac{2}{3}b \right) i,\quad z_{2}= - 4+ 16i$ ;
    \item $z_{1}=(a+b)+24i,\quad z_{2}=-(5+abi)$;
    \item $z_{1}=8+2abi,\quad z_{2}=6i+a^{2}-b^{2}$;
    \item $z_{1}=2a^{2}-5a+2,\quad z_{2}=(b^{2}+b-2)i$;
    \item $z_{1}= 0,\quad z_{2}= (a- 1)^{2}+ (b^{2}+ 2b)i$.
\end{enumerate}
\item 求出下列各复数的共轭复数，并在复平面上描出表示这些互为共轭复数的点。
\[1- i, \quad - 2+ \sqrt {2}i, \quad - \sqrt {3}i, \quad 2, \quad i, \quad \sqrt {3}+ i\]
\item 求出上题中各复数的模及幅角的主值（可用反三角函数表示）.
\item 实数$x,y$取何值时，才能使$z_1=\bar z_2$:
\begin{enumerate}[(1)]
    \item $z_{1}=( x-y )+5i,\quad  z_{2}=1-( 2x+y ) i$;
    \item $z_{1}=xy-( x+y)i,\quad  z_{2}=5+24 i$; 
    \item $z_{1}=( x^{2}-3x-2 )+( y^{2}-5y-6 ) i,\quad z_{2}=2$;
    \item $z_1=0,\quad z_2=(2x^2-5x+2)+i(y^2+y-2)$.
\end{enumerate}


\item 求证：两个复数互为共轭的必要条件是这两个复数
的模相等。
\item  在下列各题中，要使等式$| z_1|=2| z_2|$成立，试
求实数$x,y$的值：
\begin{enumerate}[(1)]
    \item $z_{1}=x-1-2i,\quad  z_{2}=1+( x+1 )i$;
    \item $z_{1}=( x-y+1 )+2 ( x+y-3 )i,\quad z_{2}=\frac{x-y+1}{2\sqrt{2}}i $
\end{enumerate}

\item  设$a,b\in\mathbb{R}$, 它们取何值时，表示复数$z=(a-b-1)+(2a-b-5)i$在复平面上的点，才能在$y=x$直线上？
\item 试求满足下列条件的复数$x+yi$在复平面上对应的点
的轨迹：
\begin{multicols}{3}
\begin{enumerate}[(1)]
    \item $|x+yi|=3$
    \item $|x+yi|<3$
    \item $|x+yi|\ge 3$
    \item $2\le |x+yi|\le 3$
    \item $0<|x+yi|<3$
\end{enumerate}
\end{multicols}
\item 设$z=a+bi,\; a,b\in\mathbb{R}$. 问：满足下列条件的点
$Z(a,b)$的集合组成什么图形？
\begin{multicols}{2}
\begin{enumerate}[(1)]
    \item $0\le |a|<1$
    \item $a>0,\; b>0,\quad a^2+b^2>4$
\end{enumerate}
\end{multicols}

\item 想想看：若已知一个非零复 数$z=a+bi$, 你能求
出它的幅角来吗？它的主值又如何求出？举例说明。

\item 满足条件$\frac{\pi}{6}\le \arg z\le \frac{\pi}{2}$的复数，在复平面上对应着什么样的点集？

\end{enumerate}

\section{复数的四则运算}
本节讨论复数的四则运算，我们将根据引进的新数——虚数单位$i$的特征$i^2=-1$和希望它遵从“数系运算通性”的要求，合理地提出运算法则，并进一步应用这些法则，有效地进行运算，同时，根据复数的几何意义，我们将相应地给出复数运算的几何意义.

\subsection{复数的加、减法}
\subsubsection{加法}
复数的加法运算按照以下规定的法则进行：设$z_1=a+bi$, $z_2=c+di$是任意两个复数，则
\[z_1+z_2=(a+bi)+(c+di)=(a+c)+(b+d)i\]

可见，任意两个复数的和仍然是一个复数.

特别地，根据这一规定我们可求出任一个实数$a$和任一纯虚数$bi$的和：
\[(a+0i)+(0+bi)=(a+0)+(0+b)i=a+bi\]

由此可知，复数$a+bi$，既可认定为一个数，又可认定它为实数$a$与纯虚数$bi$的和，今后，对于表达复数$a+bi$中的“$+$”号与表达加法运算的“$+$”号，我们将认为是一样的。

不难验证，复数加法法则的规定满足“运算通性”。
即 : 设 $z_1, z_2, z_3\in \mathbb{C}$, 则
\begin{align}
    z_{1}+ z_2 &= z_{2}+ z_{1}\tag{交换律成立}\\
    z_{1}+\left( z_{2}+z_{3} \right)&=\left( z_{1}+z_{2} \right)+z_{3}=z_{1}+z_{2}+z_{3}\tag{结合律成立}\\
    z_{1}+ 0&= 0+ z_{1}= z_{1}\tag{零的运算特性保持}
\end{align}

\begin{example}
    试求下列复数的和：
\begin{enumerate}[(1)]
    \item $(5-3i)+(-\sqrt{2}+i)$
    \item $( 9+ 6i) + ( 9- 6i)$
    \item $( 1- i) + ( - 1+ 2i) + 7+ \sqrt {2}i$
\end{enumerate}
\end{example}

\begin{solution}
\begin{enumerate}[(1)]
    \item $( 5- 3i) + ( - \sqrt {2}+ i) = (5- \sqrt {2})-2i$
    \item $( 9+ 6i) + ( 9- 6i) = 18+ 0i= 18$
    \item $( 1- i) + ( - 1+ 2i) +7+ \sqrt {2}i=(1-1+7)+(-1+2+\sqrt{2})i=7+(1+\sqrt{2})i$
\end{enumerate}
\end{solution}


由于复数可以用复平面上的“位置”向量来表示，因此，复数加法的几何意义就是：复数的加法可以按照向量的“平行四边形”法则来进行，即，先将求和的两个复数所对应的向量$\overrightarrow{OZ_{1}}$, $\overrightarrow{OZ_{2}}$画出，若这两个向量不共线，再以这两个向量为邻边作平行四边形$OZ_1ZZ_2$, 则这个平行 四 边形的对角线$OZ$所示的向量$\overrightarrow{OZ}$, 就 对 应 着 已 知 两 个 复 数 的 和(如图1.5).

\begin{figure}[htp]
    \centering
\begin{tikzpicture}[>=stealth, scale=.8]
\draw[->](-1,0)--(4,0)node[right]{$x$};
\draw[->](0,-1)--(0,4)node[right]{$y$};
\node [below right]{$O$};
\draw[very thick, ->](0,0)--(.5,2)node[above]{$Z_2$};
\draw[very thick, ->](0,0)--(2.5,1)node[right]{$Z_1$};
\draw[very thick, ->](0,0)--(3,3)node[above]{$Z$};
\draw[dashed](.5,2)--(3,3)--(2.5,1);


\end{tikzpicture}
    \caption{}
\end{figure}

若$\overrightarrow{OZ_{1}}$ 与$\overrightarrow{OZ_{2}}$共线，只要在$\overrightarrow{OZ_{1}}$的方向上紧接着延长一个$\overrightarrow{Z_{1}Z}$，使$\overrightarrow{Z_{1}Z}=\overrightarrow{OZ_{2}}$. 这时同样得到向量$\overrightarrow{OZ}$, 它对应着已知两复数的和（相当于作一个退缩成一条线段的平行四边形）.

显然，两个共轭复数的和是一个实数，它对应着实轴上的一个向量. 

\subsubsection{减法}
复数的减法仍定义为加法
的逆运算，即：满足$(c+di)+(x+yi)=a+bi$的复数$x+yi$，就叫做复数$a+bi$减去$c+di$的差，记作
\[x+yi=(a+bi)-(c+di)\]

根据复数加法法则和复数相等的定义，可以得出复数的减法法则
\[(a+bi)-(c+di)=(a-c)+(b-d)i\]

可见，两复数的差，仍是一个确定的复数，就是说，复数集对于减法也是封闭的.

同样，复数的减法的几何意义就是用向量减法的“三角
形法则”来表示，即：两个复数的差
$z_1-z_2$（向量$\VEC{OZ_1}-\VEC{OZ_2}$）与连结两个向量终点并
指向被减向量终点的向量对应
（图1.6），即
\[\VEC{OZ_1}-\VEC{OZ_2}=\VEC{Z_2Z_1}\]

\begin{figure}[htp]
    \centering
\begin{tikzpicture}[>=stealth, scale=.8]
\draw[->](-1,0)--(4,0)node[right]{$x$};
\draw[->](0,-1)--(0,4)node[right]{$y$};
\node [below right]{$O$};
\draw[ thick, ->](0,0)--(3,3)node[above]{$Z_2$};
\draw[ thick, ->](0,0)--(2.5,1)node[right]{$Z_1$};
\draw[very thick, ->](3,3)--(2.5,1);


\end{tikzpicture}
    \caption{}
\end{figure}

综上所述，复数的加、减法与我们熟悉的多项式加、减法是类似的，就是把复数的实部、虚部分别相加、减，即
\[(a+bi)\pm (c+di)=(a\pm c)+(b\pm d)i\]

\begin{example}
    计算$(3-4i)+(-2+i)-(5+2i)$
\end{example}

\begin{solution}
 $$(3-4i)+(-2+i)-(5+2i)
=(3-2-5)+(-4+1-2)i=-4-5i.$$   
\end{solution}

\begin{example}
    设复数$z_1=a+bi$, $z_2=c+di$, 试求这两复数在复平面上所描述的两点之间的距离.
\end{example}

\begin{solution}
    由于$z_1$与$z_2$在复平面上所描述的两点$Z_1$, $Z_2$之间的距离$d$，就是向量$\VEC{Z_1Z_2}$的模，因而，只要求出向量$\VEC{Z_1Z_2}$所对应的复数，并求出它的模即可。

    $\because\quad \VEC{Z_1Z_2}=\VEC{OZ_2}-\VEC{OZ_1}$

$\therefore\quad \VEC{Z_1Z_2}$对应于复数$z_2-z_1$，于是
\[\begin{split}
d&=|\VEC{Z_1Z_2}|=|z_2-z_1|\\
&=|(c+di)-(a+bi)|=|(c-a)+(d-b)i|\\
&=\sqrt{(c-a)^2+(d-b)^2}    
\end{split}\]
\end{solution}

\begin{example}
    根据复数的几何意义及向量表示，试求复平面内以$P(a,b)$点为圆心，$r$为半径的圆的方程.
\end{example}

\begin{figure}[htp]
    \centering
\begin{tikzpicture}[>=stealth, scale=1]
\draw[->](-1,0)--(4,0)node[right]{$x$};
\draw[->](0,-1)--(0,4)node[right]{$y$};
\node [below right]{$O$};
\draw[ thick, ->](0,0)--(1.5,3)node[above]{$Z$};
\draw[ thick, ->](0,0)--(2,2)node[right]{$P$};
\draw[thick](2,2)circle(1.12);
\draw[->, very thick](2,2)--node[right]{$r$}(1.5,3);

\end{tikzpicture}
    \caption{}
\end{figure}


\begin{solution}
    设圆心$P(a,b)$与复数$p=a+bi$对应，圆上任一点$Z(x,y)$与复数$z=x+yi$对应，由于已知圆的半径为$r$，因
此，向量$\VEC{PZ}$的模就是定长$r$. 即
$|\VEC{PZ}|=r$.

又由于$\VEC{PZ}=\VEC{OZ}-\VEC{OP}$，因而向量$\VEC{PZ}$就对应着复数$z-p=(x+yi)-(a+bi)=(x-a)+(y-b)i$.

所以，圆上任一点$Z$应满足$|\VEC{PZ}|=|\VEC{OZ}-\VEC{OP}|=r$，即
\[|z-p|=r\]
这就是复平面上以$P$为圆心，$r$为半径的圆的方程.

若进一步求出$|z-p|=|(x-a)+(y-b)i|
=\sqrt{(x-a)^2+(y-b)^2}$
代入圆方程，并化简。就得出用实数表示的圆的一般方程：
\[(x-a)^2+(y-b)^2=r^2\]
特别地，当$P$点在原点时，圆的方程就变成为
$$|z|=r$$
或
$$x^2+y^2=r^2$$
\end{solution}

\begin{ex}
\begin{enumerate}
    \item 利用复数加法法则和实数、虚数单位运算通性，验证：复数加法的交换律和结合律都成立.
    \item 计算（分别求和，再求差）：
\begin{multicols}{2}
\begin{enumerate}[(1)]
    \item $(2+3i)\pm (1+i)$
    \item $(2+4i)\pm (1-i)$
    \item $(-3i-2)\pm (-1-i)$
    \item $-2i\pm (-4+i)$
    \item $(2-3i)\pm (-5i)\pm(-1+7i)$
\end{enumerate}
\end{multicols}
\item 将上题中的各小题，用向量法求和.
\item 求证：两个互为共轭的复数之和为实数；两互为共轭的复数之差为纯虚数或零.
\item 想想看：在复平面上，以$P(a,b)$为圆心、以r为半径的圆内所有点的集合，可用什么样的条件给定？
\end{enumerate}
\end{ex}

\subsection{复数的乘、除法}
\subsubsection{乘法}
复数的乘法运算按照以下规定的法则进行：设$z_1=a+bi$, $z_2=c+di$是任意两复数，则
\[z_1\cdot z_2=(a+bi)\cdot (c+di)=(ac-bd)+(ad+bc)i\]
其中，由于我们要求乘法运算仍应保持“运算通性”，因而，这一规定说明复数的乘法与多项式乘法是类似的，只要将结果中的$i^2$换成$-1$，分别合并实部与虚部即可．

可见，任意两个复数的乘积仍然是一个复数，也就是说，复数集对于乘法也是封闭的.

特别地，根据这一规定，对两个互为共轭复数$z$与$\bar z$，我们有
\[z\cdot \bar z=(x+yi)\cdot (x-yi)=x^2+y^2\]

因此，互为共轭的两复数之积是一个实数，它等于每一复数的模的平方，即
\[z\cdot \bar z=|z|^2=|\bar z|^2\]

容易验证，复数乘法运算也具有“数系运算通性”，即对于任意$z_1,z_2,z_3\in\mathbb{C}$，都有
\begin{align}
z_1\cdot z_2  &= z_2\cdot z_1 \tag{乘法交换律成立}\\
z_1\cdot (z_2\cdot z_3)  &= (z_1\cdot z_2)\cdot z_3=z_1\cdot z_2\cdot z_3 \tag{乘法结合律成立}\\
z_1\cdot (z_2+ z_3)  &= z_1\cdot z_2+ z_1\cdot z_3 \tag{乘法对加法的分配律成立}\\
z_1\cdot 0  &= 0\cdot z_1=0 \tag{零的运算特性保持}\\
z_1\cdot 1  &= 1\cdot z_1=z_1 \tag{1的运算特性保持}
\end{align}

还应该指出，复数乘方运算、幂的概念与实数完全相同，而且由于复数的乘法满足交换律、结合律，所以在实数集$\R$中的指数运算律，在复数集$\mathbb{C}$中仍然成立，即

对任意$z_1,z_2,z\in\mathbb{C}$及$m,n\in\N$, 都有
\[\begin{split}
    z^m\cdot z^n&=z^{m+n}\\
    (z^m)^n&=z^{m\cdot n}\\
    (z_1\cdot z_2)^n&=z^{n}_1\cdot z^{n}_2\\
\end{split}\]

这样，我们由定义$i^2=-1$，运用指数运算律可以得出：
\[\begin{split}
i^3&=i\cdot i^2=-i\\
i^4&=i^2\cdot i^2=(-1)\cdot (-1)=1\\
i^5&=i^4\cdot i=i\\
\cdots&\cdots
\end{split}\]

一般地，对任意$n\in\mathbb{N}$，我们可归纳出：
\[\begin{split}
    i^{4n}&=(i^4)^n=1^n=1\\
    i^{4n+1}&=i^{4n}\cdot i=i\\
    i^{4n+2}&=i^{4n}\cdot i^2=-1\\
    i^{4n+3}&=i^{4n}\cdot i^3=-i\\
\end{split}\]
即，对于数$n\in\N$，则
\[i^{4n}=1,\qquad i^{4n+1}=i,\qquad i^{4n+2}=-1,\qquad i^{4n+3}=-i\]

\begin{example}
    计算$(5-4i)(1-i)(-2+3i)$.
\end{example}

\begin{solution}
\[(5-4i)(1-i)(-2+3i)=(1- 9i) (- 2+ 3i) = 25+ 21i\]
\end{solution}

\begin{example}
    计算$\left(-\frac12+\frac{\sqrt{3}}{2}i\right)^6$
\end{example}

\begin{solution}
\[\begin{split}
    \left(-\frac12+\frac{\sqrt{3}}{2}i\right)^6&=\left[\left(-\frac12+\frac{\sqrt{3}}{2}i\right)^2\right]^3\\
&= \left(\frac{1}{4}-\frac{\sqrt{3}}{2}i+\frac{3i^{2}}{4}\right)^{3}=\left(-\frac{1}{2}-\frac{\sqrt{3}}{2}i\right)^3\\
&=\left( - \frac 12\right) ^{3}- 3\left( \frac {- 1}{2}\right)^{2}\cdot \frac {\sqrt {3}}2i+3\left(-\frac{1}{2}\right)\left(\frac{\sqrt{3}}{2}i\right)^{2}-\left(\frac{\sqrt{3}}{2}i\right)^{3}\\
&=-\frac{1}{8}-\frac{3\sqrt{3}}{8}i+\frac{9}{8}+\frac{3\sqrt{3}}{8}i=1
\end{split}\]
\end{solution}


\subsubsection{除法}

复数的除法同样定义为乘法的逆运算，即：满足$(c+di)(x+yi)=a+bi\; (c+di\neq0)$的复数$x+yi$叫做复数$a+bi$除以$c+di$的商，记作
$$x+yi=\frac{a+bi}{c+di}\quad \text{或}\quad( a+bi )\div( c+di ) $$

根据两个共轭复数的乘积是一个实数以及乘法运算法
则，可以得出复数的除法法则：

当$c+ di\neq 0$ 时，
\[\frac{a+bi}{c+di}=\frac{(a+bi)(c-di)}{(c+di)(c-di)}=\frac{ac+bd}{c^2+d^2}+\frac{bc-ad}{c^2+d^2}i\]
其中由于$c+di\ne 0$，因而$c$、$d$不同时为零，所以，$c^2+d^2\ne 0$．可见这样求得的商是唯一确定的一个复数，也就是说，复数集$\mathbb{C}$对于除法（除数不为零）也是封闭的.

\begin{example}
    计算$(2-3i)\div (-3+4i)$
\end{example}

\begin{solution}
\[\begin{split}
    (2-3i)\div (-3+4i)&=\frac{2-3i}{-3+4i}=\frac{(2-3i)(-3-4i)}{(-3+4i)(-3-4i)}\\
    &=\frac{-18+i}{25}=-\frac{18}{25}+\frac{1}{25}i
\end{split}\]
\end{solution}

可以看出，复数乘、除法法则的公式并不需要硬记，乘法类似于多项式乘法，除法可视为分式的分子、分母同乘以除数的共轭数，并注意正确应用运算通性及$i^2=-1$，就可得出结果.

\begin{ex}
\begin{enumerate}
    \item 验证复数乘法对于加法满足分配律.
    \item 计算下列各式：
\begin{multicols}{2}
\begin{enumerate}[(1)]
\item $(4-3i)(-\sqrt{3}i)$
\item $\left(\frac{1}{2}-\frac{\sqrt{3}}{2}i\right)(1+i)$
\item $\left(\frac{1}{2}-\frac{\sqrt{3}}{2}i\right)^3\cdot \left(\frac{1}{2}+\frac{\sqrt{3}}{2}i\right)^3$
\item $(-\sqrt{77}+\sqrt{112}i)\cdot 0$
\end{enumerate}
\end{multicols}

\item  （口答）计算：
\[i^{13},\quad i^{22},\quad i^{36},\quad i^{43},\quad i^{70},\quad i^{101},\quad i^{355},\quad i^{404} \]
\[\frac{1}{i},\quad \frac{1}{i^2},\quad \frac{1}{i^3},\quad \frac{1}{i^4},\quad \frac{1}{i^{11}}\]
\item     计算：
\begin{multicols}{3}
\begin{enumerate}[(1)]
    \item $\frac{1}{1+i}$
    \item $\frac{1}{\sqrt{2}i}$
    \item $\frac{2i}{1-i}$
    \item $\frac{2+i}{7+4i}$
    \item $\frac{1}{(9+2i)^2}$
    \item $\frac{(1-i)^2}{(1+i)^3}$
\end{enumerate}
\end{multicols}
\end{enumerate}
\end{ex}

\subsection{复数集内的多项式与方程}
我们已经学习过实数范围内的多项式，而且明确多项式的元（未知数）只要求它满足“运算通性”，因此，对于复系数的多项式运算或未知数（元）取复数值的多项式求值，都可以与实系数多项式同样处理.

\begin{enumerate}
    \item 复系数多项式的四则运算举例
    
例如$f(x)=x+i$, $g(x)=2x+1-i$，则
\[\begin{split}
f(x)+g(x)&=3x+1\\
f(x)-g(x)&=-x-1+2i\\
f(x)\cdot g(x)&=2x^2+(1+i)x+(1+i)\\
g(x)&=2\cdot f(x)+1-3i    
\end{split}\]
这就是说，$g(x)$除
以$f(x)$可得商2，余$1-3i$（由综合除法得）.

\item 多项式求值举例

例如$f(x)=x^2+x+1$，
则
\[\begin{split}
  f(i)&=i^2+i+1=i\\
f(1+i)&=(1+i)^2+(1+i)+1=2+3i.  
\end{split}\]

\item 复系数多项式的因式分解举例

例如：在复数范围内，$f(x)=x^2+1$可以进一步分解
\[x^2+1=x^2-(-1)=x^2-i^2=(x+i)(x-i)\]

\item 复数范围内的方程与方程组，均可按实数范围时的方法求解

\item 由带余除法导出的多项式的余式定理，对于复系数
的多项式仍然适用. 例如
$f(x)=x^4+2x^3-3x^2+5$除以$x-2i$所得的余式为：
\[R=f(2i)=(2i)^4+2(2i)^3-3(2i)^2+5
=33-16i\]
\end{enumerate}

\begin{example}
    在复数范围内分解因式$x^4-16$
\end{example}

\begin{solution}
\[x^{4}-16=(x^{2}-4)(x^{2}+4)=(x- 2) (x+ 2) (x- 2i) (x+ 2i)\]
\end{solution}

\begin{example}
    解方程$x^2-6x+10=0$
\end{example}

\begin{solution}
    利用求根公式，得
$$x=\frac{6\pm\sqrt{36-40}}{2}=\frac{6\pm\sqrt{-4}}{2}=3\pm i$$
$\therefore\quad x_{1}= 3+ i,\quad x_{2}= 3- i$
\end{solution}

\begin{example}
     解方程
$\begin{cases}
    2x+(1+i)y=6-2i\\
    (3-i)x+2i y=13-i
\end{cases}$
\end{example}

\begin{solution}
   用行列式法解，由于
\[\begin{split}
    \Delta&=\begin{vmatrix}
    2& 1+i\\
    3-i& 2i
\end{vmatrix}=4i-(1+i)(3-i)=-4+2i\\
\Delta_{1}&=\begin{vmatrix}
   6- 2i& 1+i\\
    13-i& 2i
\end{vmatrix}=-10\\
\Delta_{2}&=\begin{vmatrix}
    2& 6-2i\\
     3-i& 13-i
 \end{vmatrix}=10+10i\\
\end{split}\]
$\therefore\quad x= \frac {\Delta_{1}}{\Delta} = \frac {- 10}{- 4+ 2i}= 2+ i,\quad 
y=\frac{\Delta_{2}}{\Delta}=\frac{10+10i}{-4+2i}=-1-3i$

因此，方程组的解集为$\{(x,\; y)\}=\{(2+i,\;  -1-3i)\}$.
\end{solution}

\begin{ex}
\begin{enumerate}
    \item 已知$f(x)=2x^2+(1+i)x+(1-i)$，试求$f(i)$, $f(-i)$, $f(1+i)$和$f(0)$的值.
    \item 在复数范围内分解因式$2x^4-5x^2-3$.
    \item 解方程组$\begin{cases}
        ix-y=2(i-1)\\
        x-iy=0
    \end{cases}$
\end{enumerate}
\end{ex}

\subsection*{习题1.2}

\begin{enumerate}
    \item 计算：
\begin{enumerate}[(1)]
    \item $\left(\frac{2}{3}+i\right)+\left(1-\frac{2}{3}i\right)-\left(\frac{1}{2}+\frac{3}{4}i\right)$
    \item $[(a+b)+(a-b)i]-[(a-b)-(a+b)i]$
\end{enumerate}

\item 设复数$z_1=3+2i$, $z_2=-2+5i$在复平面上分别对应着点$Z_1,Z_2$. 试求向量$\VEC{Z_1Z_2}$与$\VEC{Z_2Z_1}$所表示的复数．
\item 试求复平面上点$P(2,1)$与点$Q(3,-1)$之间的距离.
\item 求证：两个复数和的共轭复数，等于这两个复数的共轭复数之和.
\item 设$z_1,z_2$是非零复数，用几何方法证明：
\begin{enumerate}[(1)]
    \item $\big||z_1|-|z_2|\big|\le |z_1\pm z_2|\le |z_1|+|z_2|$
    \item $|z_1+z_2|^2+|z_1-z_2|^2=2|z_1|^2+2|z_2|^2$
\end{enumerate}
\item  写出复平面内任意两点连线段的垂直平分线的方程。
\item  设复平面内有两定点$F_1(-\sqrt{5},0)$, $F_2(\sqrt{5},0)$，试求到这两定点的距离之和等于6的点的轨迹方程（椭圆方程）.
\item  计算：
\begin{enumerate}[(1)]
    \item $(\sqrt{a}+\sqrt{b}i)(\sqrt{a}-\sqrt{b}i)\quad (a,b\in\mathbb{R})$
    \item $(1-2i)(-0.2+0.3i)(0.5-0.4i)$
    \item $(1+i)+(2-i^{3})+(3-i^{5})+(4-i^{7})$
    \item $(\sqrt{2}-\sqrt{2}i)^{2}\cdot\left(\frac{1}{4}+\frac{1}{4}i\right)$
    \item $(a+bi)^{3}+(a-bi)^{3}$
    \item $\begin{vmatrix}
        1-i&i\\1+i&-i
    \end{vmatrix}+\begin{vmatrix}
        i&1&0\\1&2i&1\\0&1&3i
    \end{vmatrix}$
\end{enumerate}

\item 利用公式$(a+bi)(a-bi)=a^2+b^2$, 把下列各式分解因式：
\begin{multicols}{2}
 \begin{enumerate}[(1)]
    \item $a^{4}-b^{4}$
    \item $x^{2}+2x+3$
    \item $a^{2}+2ab+b^{2}+c^{2}$
    \item $x^{3}-2x+4$
\end{enumerate}   
\end{multicols}


\item 设$\omega = - \frac 12+ \frac {\sqrt {3}}2i$, 求证：
\begin{enumerate}[(1)]
    \item $1+\omega+\omega^{2}=0$
    \item $\omega^{3}=1$
    \item $a^{3}+b^{3}+c^{3}-3abc=(a+b+c)(a+\omega b+\omega^{2}c)(a+\omega^{2}b+\omega c)$
\end{enumerate}
\item 计算：
\begin{multicols}{2}
\begin{enumerate}[(1)]
    \item $\frac{1}{2-4i}$
    \item $\frac{3-2i}{1+i}$
    \item $\frac{1+i}{(1-i)^{2}}$
    \item $\frac{1+2i}{2-i^3}$
    \item $\frac{(1-2i)^2}{3-4i}-\frac{(2+i)^2}{4-3i}$
    \item $\frac{\sqrt{3}-\sqrt{2}i}{\sqrt{3}+\sqrt{2}i}-\frac{\sqrt{2}+\sqrt{3}i}{\sqrt{2}-\sqrt{3}i}$
\end{enumerate}
\end{multicols}

\item 若$z_1,z_2\in\mathbb{C}$, 且$z_1\cdot z_2=0$, 则$z_1,z_2$中至少有一个等于0.

\item 已知$\frac1z=\frac{1}{z_{1}}+\frac{1}{z_{2}}$, 且$z_1=2+3i$, $z_{2}=1-i$,
求$z$.

\item 已知$z^2=5-12i$, 求$z$.

\item 证明下列各命题：
\begin{enumerate}[(1)]
    \item 任两个复数的和、差、积、商(除数不为零)的
共轭复数，分别等于这两个复数的共轭复数的和、差、积、商.
\item 若$f(x)$是一个一元多项式，则对于任意复数
$z$, 有$\overline{f(z)}=f(\bar{z})$.
\end{enumerate}

\item 已知$f(z)=\frac{z^2-z+1}{z^2+z+1}$, 求$f(2+3i)$, $f(1-i)$的值.

\staritem  如果规定$i^0=1$, $i^{-m}=\frac{1}{i^m}\; (m\in\mathbb{N})$, 证明：对一切整数$n$，以下等式成立：
$$i^{4n}={1},\quad i^{4n+1}=i,\quad  i^{4n+2}=-1,\quad  i^{4n+3}=-i.$$
\staritem 根据$i^2=-1$及要求复数运算满足“通性”,试说
明复 数 加 法 、 乘 法 法 则 规 定 是 合 理 的。

\item 怎样的复数$z$, 才能满足条件：
\begin{multicols}{2}
\begin{enumerate}[(1)]
    \item $\overline {z}= z$
    \item $\overline {z}= -z$
    \item $z\cdot \bar z=[{\rm Re}(z)]^2$
    \item $z-\bar z=2[{\rm Im}(z)]$
\end{enumerate}    
\end{multicols}

\item 应用余式定理，求下列各余式；
\begin{enumerate}[(1)]
    \item $f( x) = x^{4}- 3x^{3}+ 2x^{2}- x+ i$
除以$g(x)=x+i$
\item $f( x) = x^{4}- 3x^{3}+ 2x^{2}- x+ i$
除以$h(x)=(1+i)x-1+i$
\end{enumerate}
\item 应用综合除法，将多项式
$f\left(x\right)=x^{3}+3\left({1+i}\right)x^{2}-2(5-3i)x-4\left(1+2i\right)$
展开成$(x+i)$的幂的代数和形式.

\item 解方程组$\begin{cases}
    x_1+2x_2=1+i\\
    ix_2+x_3=2i\\
    ix_1-ix_3=1-i
\end{cases}$
\end{enumerate}

\section{复数的三角形式}

我们已经建立了复数集
\[\mathbb{C}=\{x+yi\mid x,y\in\R,\; i^2=-1\}\]
与复平面上的点集以及复平面上的全体向量（由原点出发）集合间的一一对应关系，并且明确了复数的模与幅角、幅角的主值等概念，这一节我们将学习复数的另一种表达形式. 就是由它的模与幅角所确定的三角形式，这将有助于复数乘法、除法、乘方、开方运算的进一步简便。

\subsection{复数的三角形式}

设$z=a+bi,\; (a,b\in\R)$，对应着复平面上的向量$\VEC{OZ}$（图1.8），并设它的模$|\VEC{OZ}|=r$，幅角为$\theta$. 若$z=a+bi\ne 0$，则可得
\[a=r\cdot \cos\theta,\qquad b=r\cdot \sin\theta\]
\begin{figure}[htp]
    \centering
\begin{tikzpicture}[>=stealth, scale=.7]
\draw[->](-1,0)--(4,0)node[right]{$x$};
\draw[->](0,-1)--(0,3)node[right]{$y$};
\node[below left]{$O$};
\draw[->, very thick](0,0)--node[above left]{$r$}(3,2.5)node[above]{$Z$};
\draw[dashed](3,0)--node[right]{$b$}(3,2.5);
\node at (1.5,0)[below]{$a$};
\draw[->](1,0)node[above right]{$\theta$} arc (0:41:1);
\end{tikzpicture}
    \caption{}
\end{figure}


$\therefore\quad z=a+bi=r\left(\pc{\theta}\right)$

因此，对于任何一个复数$z=a+bi$，都可以表示成
$r(\cos\theta +i\sin\theta)$
的形式（其中$r=\sqrt{a^2+b^2}$），并叫做这个复数的三角形式（也叫做复数的极坐标形式）；为了方便，我们将$a+bi$叫做复数的代数形式.

复数的代数形式$z=a+bi$与三角形式$z=r(\cos\theta +i\sin\theta)$之间，由图1.8可知有下列互化的关系：
\[\begin{cases}
 a=r\cos\theta\\
b=r\sin\theta   
\end{cases},\qquad \begin{cases}
    r=\sqrt{a^2+b^2}\\ \tan\theta=b/a
\end{cases}\]
（$\theta$只要取主值，且要考虑$a$，$b$的正负）.

\begin{example}
 将下列复数的代数形式化为三角形式： 
\begin{multicols}{2}
\begin{enumerate}[(1)]
    \item $\sqrt{3}+i$
    \item $-\frac{1}{2}+\frac{\sqrt{3}}{2}i$
    \item $-1$
    \item $\cos\theta-i\sin\theta$
\end{enumerate}
\end{multicols}  
\end{example}

\begin{solution}
\begin{enumerate}[(1)]
    \item $r=\sqrt{3+1}=2$, $\tan\theta=\frac{1}{\sqrt{3}}=\frac{\sqrt{3}}{3}$，且$a,b$均为正值，$\theta$在I象限，因而$\arg(\sqrt{3}
    +i)=\frac{\pi}{6}$

    $\therefore\quad \sqrt{3}+i=2\left(\pc{\frac{\pi}{6}}\right)$
\item $r=\sqrt{\frac{1}{4}+\frac{3}{4}}=1$, $\tan\theta=-\sqrt{3}$，且$a<0, b>0$，$\theta$在II象限，因而$\theta=\frac{2\pi}{3}$（主值）

$\therefore\quad -\frac{1}{2}+\frac{\sqrt{3}}{2}i=\pc{\frac{2\pi}{3}}$

\item $r=1$, $\tan\theta =0$，且$a<0$，因此$\arg(-1)=\pi$

$\therefore\quad -1=\pc{\pi}$

\item $r=\sqrt{\cos^2\theta+\sin^2\theta}=1$

$\therefore\quad \cos\theta-i\sin\theta=\pcx{-\theta}=\pcx{2\pi-\theta}$

\end{enumerate}
\end{solution}


\begin{example}
    将下列复数的三角形式化为代数形式：
\begin{multicols}{2}
\begin{enumerate}[(1)]
    \item $2\left(\pc{\frac{\pi}{4}}\right)$
    \item $\sqrt{2}\left(\pc{\frac{5\pi}{6}}\right)$
    \item $\pc{101\pi}$
    \item $5\left(\pc{\frac{\pi}{2}}\right)$
\end{enumerate}    
\end{multicols}
\end{example}

\begin{solution}
\begin{enumerate}[(1)]
    \item $2\left(\pc{\frac{\pi}{4}}\right)=\sqrt{2}+\sqrt{2}i$
    \item $\sqrt{2}\left(\pc{\frac{5\pi}{6}}\right)=-\frac{\sqrt{6}}{2}+\frac{\sqrt{2}}{2}i$
    \item $\pc{101\pi}=\pc{\pi}=-1$
    \item $5\left(\pc{\frac{\pi}{2}}\right)=5i$
\end{enumerate}  
\end{solution}

\begin{rmk}
复数的三角形式$r(\pc{\theta})$中，必须满足
    $r>0$、式中保持“$+$”号两特征，至于幅角$\theta$则不一定非要取主值，除常用主值表示外，也可以在幅角中任取一个。例如，
\[\sqrt{2}-\sqrt{2}i=2\left[\pcx{-\frac{\pi}{4}}\right]\]
也是复数$\sqrt{2}-\sqrt{2}i$的三角形式.

此外，若$z=a+bi=0$，即$z=0+0i$，这时$r=0$，其幅角是任意实数，则$0=0(\pc{\theta})$.
\end{rmk}

\begin{ex}
\begin{enumerate}
    \item 判断下列给出的是不是复数的三角形式？若不是，请把它们化为三角形式：
\begin{multicols}{2}
    \begin{enumerate}[(1)]
        \item $2\left(\cos\frac{\pi}{4}-i\sin\frac{\pi}{4}\right)$
        \item $-2\left(\pc{\frac{\pi}{3}}\right)$
        \item $\frac{1}{2}\left(\pc{\frac{3\pi}{4}}\right)$
        \item $\cos\frac{7\pi}{5}+i\sin\left(-\frac{3\pi}{5}\right)$
    \end{enumerate}
\end{multicols}
    \item 将复数表示成三角形式：
\begin{multicols}{3}
    \begin{enumerate}[(1)]
        \item $2$
        \item $-3$
        \item $4i$
        \item $-2i$
        \item $1+i$
        \item $-4+3i$
        \item $-1-\sqrt{3}i$
        \item $\frac{\sqrt{3}}{2}-\frac{1}{2}i$
    \end{enumerate}
\end{multicols}
\item    写出下列复数的代数形式：
\begin{multicols}{2}
    \begin{enumerate}[(1)]
        \item $3\left(\pc{\frac{\pi}{3}}\right)$
        \item $\sqrt{2}\left(\pc{\frac{3\pi}{4}}\right)$
        \item $\frac{1}{\sqrt{3}}\left[\pcx{-\frac{\pi}{6}}\right]$
        \item $5\left(\pc{\frac{7\pi}{2}}\right)$
    \end{enumerate}
\end{multicols}

\item 想一想：“两个复数若相等，则它们的模与幅角必定相等”对吗？为什么？
\end{enumerate}
\end{ex}

\subsection{再谈复数的乘、除法和乘方运算}
通过第二节对复数四则运算的学习，我们会感到：复数的代数形式在进行加、减法运算时，步骤十分简明易学，但进行乘、除法和乘方运算时，却显得繁杂.

我们将会看到，复数的三角形式在进行乘、除、乘方甚至开方运算时，比代数形式要简捷得多.

\begin{blk}
 {定理1} 两个复数的乘积是一个复数，它的模等于两乘数的模的乘积；它的幅角等于两乘数的幅角的和，即   

设$z_1=r_1(\pc{\theta_1})$, $z_2=r_2(\pc{\theta_2})$，则
\[\begin{split}
z_1\cdot z_2 &=r_1(\pc{\theta_1})\cdot r_2(\pc{\theta_2})\\
&=r_1\cdot r_2 \left[\pcx{\theta_1+\theta_2}\right]
\end{split}\]
\end{blk}

\begin{proof}
\[\begin{split}
    \because\quad z_1\cdot z_2&=r_1\cdot r_2(\pc{\theta_1})(\pc{\theta_2})\\
&=r_1\cdot r_2 \left[\left(\cos\theta_1\cdot \cos\theta_2-\sin\theta_1\cdot \sin\theta_2\right)\right.\\
&\qquad \qquad +i\left.\left(\cos\theta_1\cdot \sin\theta_2+\sin\theta_1\cdot \cos\theta_2\right)\right]\\
&=r_1\cdot r_2 \left[\pcx{\theta_1+\theta_2}\right]
\end{split}\]

$\therefore\quad $定理成立.
\end{proof}

根据这一定理，联系复数的向量表示，我们可以给出两复数相乘的几何意义：

两复数$z_1=r_1(\pc{\theta_1})$与$z_2=r_2(\pc{\theta_2})$相乘时，可以先在复平
面上画出它们对应的向量$\VEC{OZ_1}$, $\VEC{OZ_2}$，然后将$\VEC{OZ_1}$旋转一个角$|\theta_2|$（若$\theta_2>0$时，就按逆时针方向旋转$\VEC{OZ_1}$；若$\theta_2<0$时，就按顺时针方向旋转$\VEC{OZ_1}$），再把它的模变为原来的$r_2$倍，就可得到向量$\VEC{OZ}$（图1.9所示），这一向量就对应着所求两复数的积.

\begin{figure}[htp]
    \centering
\begin{tikzpicture}[>=stealth, scale=.7]
\draw[->](-1,0)--(5,0)node[right]{$x$};
\draw[->](0,-1)--(0,7)node[right]{$y$};
\node[below left]{$O$};
\coordinate(Z1) at (1.5,1.5);
\coordinate(Z2) at (2.5,1.8);
\coordinate(Z) at (1.05,6.45);
\draw[->, very thick](0,0)--node[above]{$r_1$}(Z1)node[above]{$Z_1$};
\draw[->, very thick](0,0)--node[below right]{$r_2$}(Z2)node[right]{$Z_2$};
\draw[->, very thick](0,0)--node[above right]{$r_1\cdot r_2$}(Z)node[above]{$Z$};
\draw[->](.75,0)node[above right]{$\theta_2$}arc (0:36:.75);
\draw[->](1.75,0)node[above right]{$\theta_1$}arc (0:45:1.75);
\draw[->](2.75,0)arc (0:80:2.75)node[right]{$\theta_1+\theta_2$};
\end{tikzpicture}
    \caption{}
\end{figure}



\begin{blk}
{推论} $n$个复数的积仍是一个复数，它的模等于$n$个乘数
的模的积，它的幅角等于$n$个乘数的幅角的和。（可用数学
归纳法证明）即：
\[\begin{split}
 z_{1}\cdot z_{2}\cdots z_{n}&=r_{1}\left(\cos\theta_{1}+i\sin\theta_{1}\right)\cdot r_{2}\left(\cos\theta_{2}+   i\sin\theta_{2}\right)\\
 &\qquad \qquad \cdots r_n(\pc{\theta_n})\\
 &=r_{1}\cdot r_{2}\cdots  r_{n}[\cos(\theta_1+\theta_2+\cdots+\theta_n)\\
 &\qquad \qquad \qquad +i\sin (\theta_1+\theta_2+\cdots+\theta_n)] 
\end{split}\]
\end{blk}

这样，在计算复数的乘积时，只要先化为三角形式。再求出模之积与幅角之和，就可写出所求乘积的三角形式。

在上边的推论中，如果$z_1=z_2=\cdots=z_n=z$, 即$r_1=r_{2}= \cdots = r_{n}= r$ 且 $\theta _{1}= \theta _{2}= \cdots = \theta _{n}= \theta $. 也就是说有$n$个相同的复数相乘时，那么就得到：
\[\left[r(\pc{\theta})\right]^n=r^n(\pc{n\theta})\quad (n\in\N)\]
这就是著名的棣美佛定理.

\begin{blk}
{定理2（棣美佛定理）}复数的$n$（$n\in\N$）次方幂仍是一个复数，它的模等于这个复数的模的$n$次幂，它的幅角等于这个复数的幅角的$n$倍.    
\end{blk}

\begin{example}
    计算：$\left(\pc{\frac{\pi}{6}}\right)\cdot \sqrt{2}\left(\pc{\frac{\pi}{12}}\right)$
\end{example}

\begin{solution}
由乘法定理可知：
\[\begin{split}
\text{原式}&=\sqrt{2}\left[\pcx{\frac{\pi}{6}+\frac{\pi}{12}}\right]\\
&=\sqrt{2}\left(\pc{\frac{\pi}{4}}\right)\\
&=\sqrt{2}\left(\frac{\sqrt{2}}{2}+\frac{\sqrt{2}}{2}i\right)=1+i
\end{split}\]
\end{solution}

\begin{example}
    计算$(1-i)^{10}$
\end{example}

\begin{solution}
    先将$1-i$化为三角形式：
\[1-i=\sqrt{2}\left[\pcx{-\frac{\pi}{4}}\right]\]
由棣美佛定理，得
\[\begin{split}
    (1-i)^{10}&=\left[\sqrt{2}\left(\pc{\frac{-\pi}{4}}\right)\right]^{10}\\
    &=\left(\sqrt{2}\right)^{10}\cdot \left[\pcx{-\frac{5\pi}{2}}\right]\\
    &=2^5\left(\cos\frac{\pi}{2}-i\sin\frac{\pi}{2}\right)=-32i
\end{split}\]
\end{solution}

\begin{example}
    如果将复数$z=2+3i$在复平面上所对应的向量$\VEC{OZ}$旋转$150^{\circ}$，那么，在复平面上所得新向量$\VEC{OZ_1}$将对应哪一个复数（用代数形式表示）.
\end{example}

\begin{solution}
    由复数乘法的几何意义知道，一个向量（对应一个复数）旋转一个角度$\theta<0$，就相当于这个复数与模长为1，幅角主值为$\theta$的另一复数（对应于幅角为$\theta$的单位向量）作乘积.

因此，所求的复数就等于复数$2+3i$与复数$1\cdot(\cos150^{\circ}+i\sin 150^{\circ})$的乘积，即
\[\begin{split}
    (2+3i)\cdot (\cos150^{\circ}+i\sin 150^{\circ})
    &=(2+3i)\left(-\frac{\sqrt{3}}{2}+\frac{1}{2}i\right)\\
&=-\frac{2\sqrt{3}+3}{2}+\frac{2-3\sqrt{3}}{2}i
\end{split}\]
\end{solution}

\begin{blk}
{定理3} 两个复数的商，在除数不为零的条件下仍是一个复数，它的模等于两个复数的模的商，它的幅角等于两复数的幅角的差（被除数的幅角减去除数的幅角）.即：当$z_2\ne 0$时，有
\[\begin{split}
    \frac{z_1}{z_2}&=\frac{r_1(\pc{\theta_1})}{r_2(\pc{\theta_2})}\\
    &=\frac{r_1}{r_2}\left[\pcx{\theta_1-\theta_2}\right]
\end{split}\]

\end{blk}

\begin{proof}
不妨设$z=\frac{r_1}{r_2}\left[\pcx{\theta_1-\theta_2}\right]$，则由乘法，
\[\begin{split}
    z\cdot z_2 &=\frac{r_1}{r_2}\cdot r_2\left[\pcx{\theta_1-\theta_2+\theta_2}\right]\\
    &=r_1(\pc{\theta_1})=z_1
\end{split} \]

又由除法的定义，可得
\[\frac{z_1}{z_2}=z\quad (z_2\ne 0)\]
\end{proof}

\begin{example}
计算$\left(\sqrt{3}-i\right)\div \left(\sqrt{3}+i\right)$
\end{example}

\begin{solution}
$\because\quad \sqrt{3}-i=2\left(\frac{\sqrt{3}}{2}-\frac{1}{2}i\right)=2\left[\pcx{-\frac{\pi}{6}}\right]$

$\sqrt{3}+i=2\left(\pc{\frac{\pi}{6}}\right)$

\[\begin{split}
    \therefore\quad \left(\sqrt{3}-i\right)\div \left(\sqrt{3}+i\right)&=(2\div 2)\left[\pcx{-\frac{\pi}{6}-\frac{\pi}{6}}\right]\\
&=1\cdot \left[\pcx{-\frac{\pi}{3}}\right]\\
&=\frac{1}{2}-\frac{\sqrt{3}}{2}i
\end{split}\]
\end{solution}

\begin{example}
    已知平面内并列的三个边长为1的正方形（如图1.10所示），利用复数证明：
$\angle 1+\angle 2-\angle 3=0$.
\end{example}

\begin{figure}[htp]
    \centering
\begin{tikzpicture}[>=stealth, scale=1.8]
\draw[->](-.5,0)--(4,0)node[right]{$x$};    
\draw[->](0,-.5)--(0,2)node[right]{$y$};    
\foreach \x/\y in {1/3,2/2,3/1}
{
    \node at (\x,0)[below]{$\x$};
    \draw(\x,0)--(\x,1)node[above]{$Z_{\y}$}--(0,0);
}
\draw(0,1)node[left]{1}--(3,1);
\node[below left]{$O$};

\draw(1-.25,1) arc (180:180+45:.25)node[left]{$3$};
\draw(2-.25,1) arc (180:180+26:.25)node[left]{$2$};
\draw(3-.5,1) arc (180:180+19:.5)node[left]{$1$};

\draw(.25,0)arc (0:45:.25);
\draw(.35,0)arc (0:26:.35);
\draw(.5,0)arc (0:19:.5);


\end{tikzpicture}
    \caption{}
\end{figure}

\begin{solution}
 如图建立坐标系以确定复平面，不在轴上的三个正
方形顶点设为$Z_1,Z_2,Z_3$．显然，它们分别对应着复数$3+i$, $2+i$, $1+i$; 由于平行线的内错角相等，因而$\angle 1$、$\angle 2$、$\angle 3$ 分别等于复数$3+i$, $2+i$，$1+i$的幅角的主值。

又由复数的乘、除法可知，$\angle 1+\angle 2-\angle 3$就是复数$\frac{(3+i)(2+i)}{1+i}$的幅角，而
\[\frac{(3+i)(2+i)}{1+i}=\frac{5+5i}{1+i}=5\]
它的幅角的主值是0，所以$\angle 1+\angle 2-\angle 3=0$.
\end{solution}

\begin{ex}
\begin{enumerate}
    \item 计算
\begin{enumerate}[(1)]
    \item $\sqrt{2}\left(\pc{\frac{\pi}{6}}\right)\cdot \frac{\sqrt{2}}{2}\left(\pc{\frac{\pi}{4}}\right)$
    \item $2\left(\pc{\frac{4\pi}{3}}\right)\cdot 3\left(\pc{\frac{5\pi}{6}}\right)$
    \item $\sqrt{3}(\pc{240^{\circ}})\cdot \frac{\sqrt{3}}{3}(\pc{150^{\circ}})$
    \item $4(\pc{18^{\circ}})\cdot [\pcx{-306^{\circ}}]\cdot 5(\pc{108^{\circ}})$
\end{enumerate}
\item 用棣美佛定理计算：
\begin{enumerate}[(1)]
    \item $\left[2\sqrt{2}(\pc{18^{\circ}})\right]^5$
    \item $(-1-i)^6$
    \item $\left[\sqrt{3}\left(\pc{\frac{3\pi}{4}}\right)\right]^4$
    \item $(2-i)\cdot \left(-\frac{1}{2}+\frac{\sqrt{3}}{2}i\right)^7$
\end{enumerate}
\item 计算：
\begin{enumerate}[(1)]
    \item $12\left(\pc{\frac{7\pi}{4}}\right)\div 6\left(\pc{\frac{2\pi}{3}}\right)$
    \item $2\div (\pc{45^{\circ}})$
    \item $-i\div 2\left(\pc{\frac{4\pi}{3}}\right)$
    \item $(-1+i)\cdot \sqrt{2}(\pc{120^{\circ}})\div 3(\pc{55^{\circ}})$
\end{enumerate}

\item 你能说明复数除法的几何意义吗？
\item 利用复数证明：在${\rm Rt}\triangle ABC$中，
$\angle C=90^{\circ}$, $AC=3BC$，$E$为$AC$上一点，且$CE=2EA$，则：
\[\angle CBE+\angle CBA-\frac{\pi}{4}=\frac{\pi}{2}\]
\end{enumerate}
\end{ex}

\subsection{复数的开方运算}
开方是乘方的逆运算，求一个复数$a+bi$的$n$次方根，就是求一个数$z$，使得$z^n=a+bi$. 凡满足这一等式的数$z$，就叫做复数$a+bi$的$n$次方根.

设$a+bi=r(\cos\theta+i\sin\theta)$, $z=\rho(\cos\varphi+i\sin\varphi)$，
则对于任意自然数$n$，都有
\[r(\cos\theta+i\sin\theta)=[\rho(\cos\varphi+i\sin\varphi)]^n=\rho^n(\cos n\varphi+i\sin n\varphi)\]

由于两复数相等，它们的模相等，而它们的幅角可以相差$2\pi$的整数倍，所以
\[\begin{cases}
    \rho^n=r\\
    n\varphi=\theta+2k\pi\quad (k\in\Z)
\end{cases}\]

因而可得
\[\rho=\sqrt[n]{r},\qquad \varphi=\frac{\theta+2k\pi}{n}\]

因此，复数$r(\pc{\theta})$的$n$次方根是
\[\sqrt[n]{r}\left(\pc{\frac{2k\pi+\theta}{n}}\right)\]
其中，当$k=0,1,\ldots,n-1$各值时，可以得到上式的$n$个不同值；当$k$继续取$n,n+1,\ldots$以及其他的值时，由于正弦、余弦函数都是以$2\pi$为周期，上式的值又重复出现$k$取$0,1,\ldots,n-1$时的同样结果．所以，复数的$n$次（$n\in\N$）方根是$n$个复数，它们的模都等于这个复数的模的$n$次算术根，它们的幅角分别等于这个复数的幅角与$2\pi$的$0,1,\ldots,n-1$倍的和的$n$分之一. 即
\[\sqrt[n]{r(\pc{\theta})}=\sqrt[n]{r}\left(\pc{\frac{\theta+2k\pi}{n}}\right)\quad (k=0,1,\ldots,n-1)\]

\begin{rmk}
在复数范围内，$\sqrt[n]{z}$应表示$n$个复数，这$n$个复数都是$z$的$n$次方根，为了区别于在实数范围内，$\sqrt[n]{a}$仅表示算术根，我们约定：被开方数是实数时，$\sqrt[n]{a}\; (a\in\R)$仍表示算术根（单值），$\sqrt[n]{a+0i}$或$\sqrt[n]{a(\pc{0})}\; (a>0)$或$\sqrt[n]{|a|(\pc{\pi})}$都表示$n$个$n$次方根（多值）. 至于被开方数是虚数时，$\sqrt[n]{a+bi}\; (b\ne 0)$或$\sqrt[n]{r(\pc{\theta})}$都表示$n$个方根.
\end{rmk}

\begin{example}
    求$i$的平方根.
\end{example}

\begin{solution}
$\because\quad i=\pc{\frac{\pi}{2}}$

$\therefore\quad \sqrt{i}=\pc{\frac{\frac{\pi}{2}+2k\pi}{2}}=\pcx{\frac{\pi}{4}+k\pi}\quad (k=0,1)$

因此，$i$的平方根是
\[\begin{split}
z_1=\pc{\frac{\pi}{4}}=\frac{\sqrt{2}}{2}+\frac{\sqrt{2}}{2}i\qquad &\text{（$k=0$时）}\\
z_2=\pc{\frac{5\pi}{4}}=-\frac{\sqrt{2}}{2}-\frac{\sqrt{2}}{2}i\qquad &\text{（$k=1$时）}
\end{split}\]
\end{solution}

\begin{example}
    求$-1+i$的四次方根.
\end{example}

\begin{solution}
$\because\quad -1+i=\sqrt{2}\left(\pc{\frac{3\pi}{4}}\right)$

$\therefore\quad \sqrt[4]{-1+i}=\sqrt[8]{2}\left(\pc{\frac{\frac{3\pi}{4}+2k\pi}{4}}\right)\quad (k=0,1,2,3)$

因此，$-1+i$的四次方根是：
\[\begin{split}
z_1=\sqrt[8]{2}\left(\pc{\frac{3\pi}{16}}\right)\qquad & (k=0)\\
z_2=\sqrt[8]{2}\left(\pc{\frac{11\pi}{16}}\right)\qquad & (k=1)\\
z_3=\sqrt[8]{2}\left(\pc{\frac{19\pi}{16}}\right)\qquad & (k=2)\\
z_4=\sqrt[8]{2}\left(\pc{\frac{27\pi}{16}}\right)\qquad & (k=3)\\
\end{split}\]
\end{solution}

\begin{example}
    求$-3$的平方根.
\end{example}

\begin{solution}
$\because\quad -3=3(\pc{\pi})$

$\therefore\quad \sqrt{-3+0i}=\sqrt{3}\left(\pc{\frac{\pi+2k\pi}{2}}\right)\quad (k=0,1)$

因此，$-3$的平方根有两个复数：
\[\begin{split}
z_1&=\sqrt{3}\left(\pc{\frac{\pi}{2}}\right)=\sqrt{3}i\\
z_2&=\sqrt{3}\left(\pc{\frac{3\pi}{2}}\right)=-\sqrt{3}i\\
\end{split}\]
\end{solution}

一般来说，在复数范围内，任一个负实数$-a$有两个平方根：$\pm\sqrt{a}i$.

这样一来，以前在解实系数一元二次方程$ax^2+bx+c=0$（$a\ne 0$）时，如果$b^2-4ac<0$，在实数范围内无解. 现在，在复数范围内，方程就有两个解：
\[x=\frac{-b\pm\sqrt{-(b^2-4ac)}i}{2a}\qquad (b^2-4ac<0)\]
显然，这两个解是一对共轭复数.

\begin{example}
    解方程$x^2-2x+6=0$.
\end{example}

\begin{solution}
$\because\quad \Delta=b^2-4ac=4-24=-20<0$

$\therefore\quad x=\frac{2\pm\sqrt{20}i}{2}=1\pm\sqrt{5}i$
\end{solution}

由以上可知，复数集$\mathbb{C}$对于开方运算是封闭的，这是复数系不同于实数系的一个重要性质.

利用复数开方的方法，我们可以解一些特殊的高次方程，例如，

形如$a_nx^n+a_0=0\; (a_n\ne 0)$的方程，我们叫做二项方程，就可以完整地解出，即，首先化成$x^n=b\quad \left(b=-\frac{a_0}{a_n}\right)$的形式，从而解出$x$的$n$个值.

\begin{example}
    解方程$x^5+243=0$.
\end{example}

\begin{solution}
$\because\quad x^5=-243=243(-1+0i)=243(\pc{\pi})$

$\therefore\quad x=\sqrt[5]{243}\left(\pc{\frac{\pi+2k\pi}{5}}\right)\quad (k=0,1,2,3,4)$

即
\[\begin{split}
x_1&=3\left(\pc{\frac{\pi}{5}}\right)\qquad\quad (k=0)\\
x_2&=3\left(\pc{\frac{3\pi}{5}}\right)\qquad (k=1)\\
x_3&=3\left(\pc{\pi}\right)\qquad\qquad (k=2)\\
x_4&=3\left(\pc{\frac{7\pi}{5}}\right)\qquad (k=3)\\
x_5&=3\left(\pc{\frac{9\pi}{5}}\right)\qquad (k=4)\\
\end{split}\]
\end{solution}

仔细观察这五个解可以发现：它们的模都相同，它们的幅角依次相差$\frac{2\pi}{5}$：因此，这个方程的五个解所对应的复平面内的五个点是均匀地分布在以原点为圆心，以3为半径的圆周上（图1.11）．这就是此方程的解的几何意义。
\begin{figure}[htp]
    \centering
\begin{tikzpicture}[>=stealth]
\draw[->](-2.5,0)--(2.7,0)node[right]{$x$};
\draw[->](0,-2)--(0,2.7)node[right]{$y$};
\draw(0,0)node[below left]{$O$} circle(1.5);
\coordinate (Z1) at (36:1.5);
\coordinate (Z2) at (36*3:1.5);
\coordinate (Z3) at (36*5:1.5);
\coordinate (Z4) at (36*7:1.5);
\coordinate (Z5) at (36*9:1.5);
\draw[very thick](0,0)--(Z1)node[right]{$Z_1$};
\draw[very thick](0,0)--(Z2)node[above]{$Z_2$};
\draw[very thick](0,0)--(Z3)node[above left]{$Z_3$};
\draw[very thick](0,0)--(Z4)node[below]{$Z_4$};
\draw[very thick](0,0)--(Z5)node[right]{$Z_5$};
\node at (-1.5,0)[below left]{$-3$};
\foreach \x in {1,2,3,4,5}
{
    \draw[fill](Z\x) circle(1.5pt);
}


\end{tikzpicture}
    \caption{}
\end{figure}

一般地，二项方程$x^n=b\; (b\in\mathbb{C})$有$n$个复数根，这些根的几何意义是复平面内的$n$个点（或向量），这些点（或向量的终点）均匀地分布在以原点为圆心，以$\sqrt[n]{|b|}$为半径的圆上.

还要进一步指出：对一般的实系数$n$次方程在复数集内也是恰有$n$个根，但未必都是实数根，可能会有虚根（如例1.25).

可以证明，实系数方程的虚根总是成对地出现的，即

\begin{blk}{定理4（虚根成对定理）}
如果$z=a+bi$是实系数方程
$f(x)=a_nx^n+a_{n-1}x^{n-1}+\cdots+a_1x+a_0=0$
的一个根，那么$\bar z=a-bi$也是方程$f(x)=0$的一个根.
\end{blk}

\begin{proof}
    由于$z=a+bi$是方程的根，所以
\[f(z)=f(a+bi)=0\]
但根据习题1.2第15题（2）已证明的结论，应有
\[f(\bar z)=f(a-bi)=f(\overline{a+bi})=\bar 0=0\]
因此，$\bar z=a-bi$也是方程$f(x)=0$的根.
\end{proof}

\begin{example}
    已知方程
$2x^4-x^3+5x^2+13x+5=0$
的一个根为$1-2i$，求这个方程的其余根.
\end{example}

\begin{solution}
    这是一个实系数一元四次方程，由定理4及已知一根$1-2i$可知，这一方程必有根$1+2i$.

又由因式定理知：$(x-1+2i)(x-1-2i)=x^2-2x＋5$是方程左边的因式，利用长除法可得
\[2x^4-x^3+5x^2+13x+5=(2x^2+3x+1)(x^2-2x+5)\]

所以只要解方程$2x^2+3x+1=0$的根：$x_1=-1$, $x_2=-\frac{1}{2}$，就可知：

原方程的其余三根为$1+2i,\; -1,\; -\frac{1}{2}$.
\end{solution}

\begin{ex}
\begin{enumerate}
    \item 直接说出以下各数的平方根：
\[-4,\quad -2.25,\quad -3,\quad -t\; (t\in\R^+),\quad -m^2\; (m\in\R)\]
\[a\; (a\in\R),\quad a-b\; (a,b\in\R)\]
    \item 求下列各数的方根：
\begin{enumerate}[(1)]
    \item $-i,\; \sqrt{2}-\sqrt{2}i,\; -\frac{1}{2}+\frac{\sqrt{3}}{2}i$的平方根；
    \item $27,\; 1+i$的立方根；
    \item $-i$的五次方根．
\end{enumerate}
    \item 解方程（在复数范围内）：
\begin{multicols}{2}
\begin{enumerate}[(1)]
    \item $x^2+x+16=0$
    \item $x^3-1=0$
    \item $x^4+16=0$
    \item $x^5-i=0$
\end{enumerate}
\end{multicols}
    \item 已知方程$4x^4-2x^3-x-1=0$有一个根为$-\frac{i}{\sqrt{2}}$，试求这个方程的其余三个根.
\end{enumerate}
\end{ex}


\subsection{单位根}
我们把方程$x^n=1$在复数范围内的$n$个根，叫做$n$次单位根．也就是在复数范围内，1的$n$个不同的$n$次方根，都叫$n$次单位根. 即

$\because\quad x^n=1=\pc{0}$

$\therefore\quad x=1\cdot \left(\pc{\frac{2k\pi}{n}}\right)\quad (k=0,1,\ldots,n-1)$

这样可得出$n$个单位根：
\[\begin{split}
x_1&=\pc{0}=1\qquad \qquad \text{（当$k=0$时）}\\
x_2&=\pc{\frac{2\pi}{n}}\qquad \qquad \text{（当$k=1$时）}\\
x_3&=\pc{\frac{4\pi}{n}}\qquad \qquad \text{（当$k=2$时）}\\
\cdots&\cdots\cdots\\ 
x_n&=\pc{\frac{2(n-1)\pi}{n}}\qquad \text{（当$k=n-1$时）}\\
\end{split}\]

由二项方程解的几何意义可以知道：$n$个不同的单位根可用复平面内以原点为圆心的单位圆的一个内接正$n$边形顶点来表示.

\begin{example}
    求三次单位根．
\end{example}

\begin{solution}
$\because\quad x^3=1=\pc{0}$

$\therefore\quad x=\pc{\frac{2k\pi}{3}}\quad (k=0,1,2)$

因此，三个三次单位根为(如图1.12)：
\[\begin{split}
    x_1&=\pc{0}=1\\
    x_2&=\pc{\frac{2\pi}{3}}=-\frac{1}{2}+\frac{\sqrt{3}}{2}i\\
    x_3&=\pc{\frac{4\pi}{3}}=-\frac{1}{2}-\frac{\sqrt{3}}{2}i\\
\end{split}\]
\end{solution}

\begin{minipage}{.45\textwidth}
\centering
\begin{tikzpicture}[>=stealth]
\draw[->](-2.5,0)--(2.7,0)node[right]{$x$};
\draw[->](0,-2)--(0,2.7)node[right]{$y$};
\draw(0,0)node[below left]{$O$} circle(1.5);
\draw[very thick](1.5,0)node[above right]{$x_1$}--(120:1.5)node[above]{$x_2$}--(240:1.5)node[below]{$x_3$}--cycle;

\end{tikzpicture}
\captionof{figure}{ }
\end{minipage}\hfill
\begin{minipage}{.45\textwidth}
\centering
\begin{tikzpicture}[>=stealth]
    \draw[->](-2.5,0)--(2.7,0)node[right]{$x$};
    \draw[->](0,-2)--(0,2.7)node[right]{$y$};
    \draw(0,0)node[below left]{$O$} circle(1.5);
\draw[very thick](1.5,0)node[above right]{$x_1$}--(0,1.5)node[above right]{$x_2$}--(-1.5,0)node[above left]{$x_3$}--(0,-1.5)node[below left]{$x_4$}--cycle;

\end{tikzpicture}
\captionof{figure}{ }
\end{minipage}

\begin{example}
    求四次单位根
\end{example}

\begin{solution}
$\because\quad x^4=1$

$\therefore\quad x=\pc{\frac{2k\pi}{4}}\quad (k=0,1,2,3)$

因此，四个四次单位根为(如图1.13)：
\begin{center}
\begin{tabular}{p{.2\textwidth}p{.15\textwidth}p{.2\textwidth}p{.15\textwidth}}
    $x_1=1$& $(k=0)$ & $x_2=i$ & $(k=1)$\\
    $x_3=-1$ &$ (k=2)$& $x_4=-i$ & $(k=3)$\\
\end{tabular} 
\end{center}
\end{solution}


如果用复数的向量表示法
来看这$n$个不同的$n$次单位根的话，我们会发现单位根之间还有一定的关系.

\begin{figure}[htp]
    \centering
\begin{tikzpicture}[>=stealth]
    \draw[->](-2.5,0)--(2.7,0)node[right]{$x$};
    \draw[->](0,-2)node[below]{$(n=12)$}--(0,2.7)node[right]{$y$};
    \draw(0,0)node[below left]{$O$} circle(1.5);
\foreach \x in {0,1,2,...,11}
{
    \draw[fill](\x*30:1.5) circle(1.5pt);
}
\node at (1.5,0)[below right]{1};
\node at (0:1.5)[above right]{$x_1$};
\node at (30:1.5)[above right]{$x_2$};
\node at (60:1.5)[above right]{$x_3$};
\node at (-30:1.5)[below right]{$x_n$};


\end{tikzpicture}
    \caption{}
\end{figure}




\begin{enumerate}
\item 如图1.14所示，在复平面上$n$个$n$次单位根正好将单位圆$n$等分，其中：
$x_1=1$，单位向量$\VEC{Ox_1}$旋转$\frac{2\pi}{n}$
得：$x_2=\pc{\frac{2\pi}{n}}$
记为$\omega$；再旋转$\frac{2\pi}{n}$
得：$x_3=\pc{\frac{4\pi}{n}}=\omega^2$；再旋转$\frac{2\pi}{n}$可一样地
得：$x_4=\pc{\frac{6\pi}{n}}=\omega^3$，……

同理可得：
\[x_n=\pc{\frac{2(n-1)\pi}{n}}=\omega^{n-1}\]

因此，设$\omega=\pc{\frac{2\pi}{n}}$，则$n$个$n$次单位根又可写成：
\[1,\; \omega,\; \omega^2,\; \ldots, \omega^{n-1}\]

\item 进一步考虑这$n$个单位根的关系，可以发现：

\[\begin{split}
    \because\quad \omega^{n-1}=\pc{\frac{2(n-1)\pi}{n}}&=\pcx{\frac{-2\pi}{n}}\\
    &=\cos\frac{2\pi}{n}-i\sin\frac{2\pi}{n}=\overline{\omega}
\end{split}\]

$\therefore\quad \omega^{n-1}=\overline{\omega}=\frac{1}{\omega}=\omega^{-1}$

同理可得：\[\omega^{n-2}=\overline{\omega^2}=\omega^{-2},\ldots\]

\[\begin{cases}
\text{当$n$为偶数时}& \omega^{\tfrac{n}{2}+1}=\overline{\omega^{\tfrac{n}{2}-1}}=\omega^{-\left(\tfrac{n}{2}-1\right)}=\omega^{-\tfrac{n}{2}+1},\quad \omega^{\tfrac{n}{2}}=-1\\
\text{当$n$为奇数时}& \omega^{\tfrac{n+1}{2}}=\overline{\omega^{\tfrac{n-1}{2}}}=\omega^{-\tfrac{n-1}{2}}
\end{cases}\]
所以，这$n$个$n$次单位根还可以写成：
\begin{itemize}
    \item 当$n$为偶数时（如例1.29）：
\[1,\; \omega,\; \omega^2,\;\ldots,\;  \omega^{\tfrac{n}{2}-1},\;-1,\;  \omega^{-\left(\tfrac{n}{2}-1\right)},\; \ldots, \;  \omega^{-2},\; \omega^{-1}\]
或
\[1,\; \omega,\; \omega^2,\;\ldots,\;  \omega^{\tfrac{n}{2}-1},\;-1,\;  \overline{\omega^{\tfrac{n}{2}-1}},\; \ldots, \;  \overline{\omega^{2}},\; \overline{\omega}\]

\item 当$n$为奇数时（如例1.28）：
\[1,\; \omega,\; \omega^2,\;\ldots,\;  \omega^{\tfrac{n-1}{2}},\;  \omega^{-\tfrac{n-1}{2}},\; \ldots, \;  \omega^{-2},\; \omega^{-1}\]
或
\[1,\; \omega,\; \omega^2,\;\ldots,\;  \omega^{\tfrac{n-1}{2}},\;  \overline{\omega^{\tfrac{n-1}{2}}},\; \ldots, \;  \overline{\omega^{2}},\; \overline{\omega}\]
其中，$\omega=\pc{\frac{2\pi}{n}}$.
\end{itemize}

例如，三次单位根为$1,\; \omega,\; \omega^{-1}$（或$\bar\omega$），其中
\[\omega=\pc{\frac{2\pi}{3}}=-\frac{1}{2}+\frac{\sqrt{3}}{2}i\]

又如，四次单位根为$1,\; \omega,\; -1,\; \omega^{-1}$（或$\bar\omega$），其中
\[\omega=\pc{\frac{2\pi}{4}}=i\]

\item 利用单位根，可以方便地求出任一个正实数$a$（$a>0$且$a\ne 1$）的$n$个$n$次方根．

这就是：若$a>0$且$a\ne 1$，则$a$的$n$个$n$次方根即方程$x^n=a$的解为：
\[\begin{cases}
x_1=\sqrt[n]{a}\\
x_2=\sqrt[n]{a}\cdot \omega\\
x_3=\sqrt[n]{a}\cdot \omega^2\\
\cdots\cdots\cdots\\
x_n=\sqrt[n]{a}\cdot \omega^{n-1}\\
\end{cases}\qquad \text{这里$\omega=\pc{\frac{2\pi}{n}}$是单位根之一}\]
\end{enumerate}

\begin{ex}
试求出六次和七次单位根，并分别检验它们之间的关系.
\end{ex}

\subsection*{习题1.3}
\begin{enumerate}
    \item 把下列复数化成三角形式表示：
\begin{multicols}{3}
\begin{enumerate}[(1)]
    \item 6
    \item $6i$
    \item $1+i$
    \item $\frac{1}{2}-\frac{\sqrt{3}}{2}i$
    \item $-\sqrt{2}-\sqrt{2}i$
\end{enumerate}
\end{multicols}

\item 把下列复数化为代数形式表示，并画出复平面上相应的向量来：
\begin{multicols}{2}
\begin{enumerate}[(1)]
    \item $3\sqrt{2}\left(\pc{\frac{\pi}{4}}\right)$
    \item $2\left(\pc{\frac{11\pi}{6}}\right)$
    \item $\pc{\pi}$
    \item $\sqrt{3}\left[\pcx{-\frac{2\pi}{3}}\right]$
    \item $\pcx{k\cdot \frac{\pi}{4}}$
    
    $(k=0,1,2,\ldots,6,7)$
\end{enumerate}
\end{multicols}

\item 计算
\begin{enumerate}[(1)]
    \item $\sqrt{3}\left(\pc{\frac{\pi}{3}}\right)\cdot 2\sqrt{3}\left(\pc{\frac{\pi}{6}}\right)$
    \item $\sqrt{10}\left(\pc{\frac{\pi}{2}}\right)\cdot \sqrt{2}\left(\pc{\frac{\pi}{4}}\right)$
\end{enumerate}

\item 下列复数是不是三角形式？如果不是，请化成三角形式表示：
\begin{multicols}{2}
\begin{enumerate}[(1)]
    \item $\cos\frac{\pi}{6}-i\sin\frac{\pi}{6}$
    \item $\sqrt{2}\left(\cos\frac{4\pi}{3}+i\sin\frac{2\pi}{3}\right)$
    \item $7\left(\sin\frac{3\pi}{4}+i\cos\frac{3\pi}{4}\right)$
    \item $3\left(\cos\frac{5\pi}{6}+i\sin\frac{-7\pi}{6}\right)$
\end{enumerate}
\end{multicols}

\item 求证：
\begin{enumerate}[(1)]
    \item $\sqrt{2}(\pc{75^{\circ}})\cdot \sqrt{2}\left(\pc{15^{\circ}}\right)=2i$
    \item $(\cos3\theta-i\sin3\theta)(\cos2\theta-i\sin2\theta)=\cos5\theta-i\sin5\theta$
\end{enumerate}

\item 如图1.15所示，菱形$OABC$的一个内角$\angle AOC=\frac{\pi}{4}$，且$A$点所对应的复数为$2+i$, 试求，$B$、$C$两点各对应的复数.

\begin{figure}[htp]
    \centering
\begin{tikzpicture}[>=stealth]
    \draw[->](-.75,0)--(3.4,0)node[right]{$x$};
    \draw[->](0,-.75)--(0,3.4)node[right]{$y$};
    \node[below left]{$O$} ;
    \coordinate (A) at (30:2);
    \coordinate (C) at (75:2);
    \coordinate (B) at (105/2:3.7);
    \draw[very thick](0,0)--(A)node[right]{$A$}--(B)node[above]{$B$}--(C)node[above]{$C$}--cycle;
\end{tikzpicture}
    \caption{}
\end{figure}


\item 已知下列各复数$z$，求$\left|\frac{1}{z}\right|$及$\arg\left(\frac{1}{z}\right)$
\begin{multicols}{2}
\begin{enumerate}[(1)]
    \item $z=2\left(\pc{\frac{\pi}{12}}\right)$
    \item $z=\cos\frac{\pi}{6}-i\sin\frac{\pi}{6}$
    \item $z=\frac{1-i}{\sqrt{2}}$
\end{enumerate}    
\end{multicols}

\item 把复数$3-\sqrt{3}i$对应的向量旋转$\frac{\pi}{3}$，所得向量相对应的复数是什么？

\item 用棣美佛定理计算：
\begin{multicols}{2}
\begin{enumerate}[(1)]
    \item $[2(\pc{15^{\circ}})]^6$
    \item $\left[\frac{1}{\sqrt{2}}(\pc{225^{\circ}})\right]^4$
    \item $(1+\sqrt{3}i)^4$
    \item $(2-2\sqrt{3}i)^4$
    \item $\frac{(\sqrt{3}+i)^5}{-1+\sqrt{3}i}$
    \item $\left(\frac{2+2i}{1-\sqrt{3}i}\right)^8$
\end{enumerate}
\end{multicols}

\item 用棣美佛定理证明：
\begin{enumerate}[(1)]
    \item $\cos2\theta=\cos^2\theta-\sin^2\theta,\qquad \sin2\theta=2\sin\theta\cdot \cos\theta$
    \item $\cos3\theta=4\cos^3\theta-3\cos\theta,\qquad \sin3\theta=3\sin\theta-4\sin^3\theta$
\end{enumerate}

\item 设$n\in\N$，$n$取何值时，$(1+\sqrt{3}i)^n$是一个实数？

\item 计算：
\begin{enumerate}[(1)]
    \item $10\left(\pc{\frac{2\pi}{3}}\right)\div 2\left(\pc{\frac{\pi}{6}}\right)$
    \item $-10i\div 6\left(\pc{\frac{\pi}{6}}\right)$
\end{enumerate}

\item 求证：
\begin{enumerate}[(1)]
    \item $\frac{1}{\cos\theta+i\sin\theta}=\cos\theta-i\sin\theta$
    \item $(\pc{\theta})^{-n}=\pcx{-n\theta}\quad (n\in\mathbb{N})$
    \item $(\cos\theta-i\sin\theta)^n=\cos n\theta-i\sin n \theta\quad (n\in\N)$
\end{enumerate}

\item 利用复数证明余弦定理.

\item 化简：
\begin{enumerate}[(1)]
    \item $\frac{(\pc{7\theta})(\pc{2\theta})}{(\pc{5\theta})(\pc{3\theta})}$
    \item $\frac{\cos\varphi-i\sin\varphi}{\pc{\varphi}}$
    \item $\frac{\pc{\alpha}}{\sin\alpha+i\cos\alpha}$
\end{enumerate}
\item 在复数范围内解下列方程：
\begin{multicols}{2}
\begin{enumerate}[(1)]
    \item $4x^2+9=0$
    \item $2(x^2+4)=5x$
    \item $(x-3)(x-5)=-2$
    \item $\frac{1}{x+3}-\frac{1}{x}=1$
    \item $x^4+3x^2+1=0$
    \item $\frac{x}{x^2+1}+\frac{x^2+1}{x}-\frac{5}{2}=0$
\end{enumerate}
\end{multicols}
\item 解下列方程组：
\begin{multicols}{2}
\begin{enumerate}[(1)]
    \item $\begin{cases}
a+b=2\\ab=2        
    \end{cases}$
    \item $\begin{cases}
        x^2+y^2=0\\
        xy=1
    \end{cases}$
\end{enumerate}
\end{multicols}

\item 求：
\begin{enumerate}[(1)]
    \item $8\left(\pc{\frac{\pi}{3}}\right)$的六次方根;
    \item $-2i$的五次方根.
\end{enumerate}
\item 求：六次单位根、七次单位根，并在复平面上用向量表示出来.

\item 解下列方程：
\begin{multicols}{2}
\begin{enumerate}[(1)]
    \item $x^3+1=i$
    \item $x^4+16=0$
    \item $x^{10}-32x^5+1024=0$
    \item $x^{12}+63x^6-64=0$
\end{enumerate}
\end{multicols}

\staritem 求证：
\begin{enumerate}[(1)]
    \item 不经计算证明$\left|\frac{a+bi}{a-bi}\right|\equiv 1$
    \item 有相同幅角的两个复数之比必为实数；
    \item 复平面上四个点$Z_1$, $Z_2$, $Z_3$, $Z_4$共圆或共直线的充分必要条件是：
\[\left.\frac{Z_3-Z_1}{Z_3-Z_2}\right/\frac{Z_4-Z_1}{Z_4-Z_2} \text{ 为实数}\]
\end{enumerate}
\end{enumerate}

\section{复数的指数形式}
我们已经学习了复数的代数形式、三角形式以及它们的运算.在科学技术，特别是在电工和无线电计算中，常要涉及复数或三角函数的变换和计算，为方便还采用复数的另一种形式——复数的指数形式，在这一节将学习复数的指数形式及其运算、应用.

\subsection{复数的指数形式及其运算}
我们首先引进一个著名的公式——欧拉公式：
\[e^{i\theta}=\pc{\theta}\]
其中$e=2.71828\cdots$是一个重要的常数，也就是自然对数的底数.

欧拉公式在以后的“复变函数论”中可以证明。这个公式表明了：模为1、幅角为$\theta$的复数$\pc{\theta}$与以$e$为底的复指数函数的关系，从这个公式出发，对任一个复数
\[z=a+bi=r(\pc{\theta})\]
就可以表示成$z=r\cdot e^{i\theta}$的形式，我们就把这一表达式叫做复数的指数形式.

例如：
\begin{itemize}
    \item 复数$\pc{\frac{\pi}{2}}$的指数形式为$e^{i\tfrac{\pi}{2}}$
    \item 复数$3\left(\pc{\frac{\pi}{4}}\right)$的指数形式为$3e^{i\tfrac{\pi}{4}}$
\end{itemize}

这样，同一个复数，就有三种不同形式的表达方法——代数形式、三角形式、指数形式.
如：
\[\begin{split}
    \frac{1}{2}+\frac{\sqrt{3}}{2}i&=\pc{\frac{\pi}{3}}=e^{i\tfrac{\pi}{3}}\\
    -2\sqrt{3}+2i&=4\left(\pc{\frac{5\pi}{6}}\right)=4e^{i\tfrac{5\pi}{6}}
\end{split}\]

复数的指数形式便利于乘、除、乘方及开方运算.因为可以证明原有在实数范围内的指数运算律它都仍然满足. 这就是：

设$z_1=r_1e^{i\theta_1}$，$z_2=r_2e^{i\theta_2}$，则有
\[\begin{split}
    z_1\cdot z_2&=\left(r_1e^{i\theta_1}\right)\cdot \left(r_2e^{i\theta_2}\right)\\
    &=r_1(\pc{\theta_1})\cdot r_2(\pc{\theta_2})\\
    &=r_1r_2[\pcx{\theta_1+\theta_2}]=r_1r_2\cdot e^{i(\theta_1+\theta_2)}
\end{split}\]
即
\[r_1e^{i\theta_1}\cdot r_2e^{i\theta_2}=r_1r_2 e^{i(\theta_1+\theta_2)}\]

同理可证：
\[\begin{split}
    \frac{z_1}{z_2}&=\frac{r_1e^{i\theta_1}}{r_2e^{i\theta_2}}=\frac{r_1}{r_2}\cdot e^{i(\theta_1-\theta_2)}\quad (z_2\ne 0)\\
z^n_1&=(r_1e^{i\theta_1})^n=r^n_1 \cdot e^{i n\theta_1}\quad (n\in\N)
\end{split} \]
$z_1$的$n$次方根：$\sqrt[n]{z_1}=\sqrt[n]{r_1e^{i\theta_1}}=\sqrt[n]{r_1}e^{i\tfrac{\theta_1+2k\pi}{n}}\quad (k=0,1,\ldots,n-1)$

\begin{example}
把复数$z=2i$化成指数形式.
\end{example}

\begin{solution}
$\because\quad z=2i=2\left(\pc{\frac{\pi}{2}}\right)$

$\therefore\quad z=2\cdot e^{i\tfrac{\pi}{2}}$
\end{solution}

\begin{example}
已知$z_1=\sqrt{2}\cdot e^{i\left(-\tfrac{\pi}{4}\right)}$，$z_2=\sqrt{5}\cdot e^{i\tfrac{2\pi}{3}}$. 试求：
\begin{multicols}{2}
\begin{enumerate}[(1)]
    \item $z_1\cdot z_2$
    \item $z_1$的平方根
\end{enumerate}
\end{multicols}
要把结果表示成复数的三角形式.
\end{example}

\begin{solution}
\begin{enumerate}[(1)]
    \item \[\begin{split}
z_1\cdot z_2&=\left(\sqrt{2} e^{i\left(-\tfrac{\pi}{4}\right)}\right)\cdot \left(\sqrt{5} e^{i\tfrac{2\pi}{3}}\right)\\
&=\sqrt{10}e^{i\tfrac{5\pi}{12}}=\sqrt{10}\left(\pc{\frac{5\pi}{12}}\right)
    \end{split}\]
\item $\sqrt{\sqrt{2}\cdot e^{i\left(-\tfrac{\pi}{4}\right)}}=\sqrt[4]{2}\cdot e^{i\tfrac{-\tfrac{\pi}{4}+2k\pi}{2}}\quad (k=0,1)$

$\therefore\quad z_1$的两个平方根为
\[\begin{split}
x_1&=\sqrt[4]{2}\cdot e^{i\left(-\tfrac{\pi}{8}\right)}=\sqrt[4]{2}\left(\pc{\frac{15\pi}{8}}\right)\\
x_2&=\sqrt[4]{2}\cdot e^{i\tfrac{7\pi}{8}}=\sqrt[4]{2}\left(\pc{\frac{7\pi}{8}}\right)
\end{split}\]
\end{enumerate}
\end{solution}



\begin{example}
试用$e^{i\theta}$与$e^{-i\theta}$表示$\cos\theta$与$\sin\theta$.    
\end{example}

\begin{solution}
由于
\begin{align}
e^{i\theta}&= \pc{\theta} \tag{1}\\
e^{-i\theta}&= \pcx{-\theta}=\cos\theta-i\sin\theta \tag{2}
\end{align}
(1)(2)联立，可解出
\[\begin{split}
    \cos\theta&=\frac{1}{2}\left(e^{i\theta}+e^{-i\theta}\right)\\
    \sin\theta&=\frac{1}{2i}\left(e^{i\theta}-e^{-i\theta}\right)\\
\end{split}\]
\end{solution}

\begin{ex}
\begin{enumerate}
    \item 将复数$-1$，$2+2i$，$\pc{15^{\circ}}$，$4i$分别表示成指数形式.
    \item 将复数$e^{-i\tfrac{\pi}{2}}$，$3e^{2i}$，$\sqrt{2}e^{i\tfrac{3\pi}{2}}$分别化为三角形式和代数形式.
    \item 用指数形式计算：
\begin{enumerate}[(1)]
    \item $8\left(\pc{\frac{7\pi}{6}}\right)\cdot 2\left(\pc{\frac{\pi}{4}}\right)$
    \item $2\div e^{i\tfrac{\pi}{4}}$
    \item $(1+\sqrt{3}i)^{10}$
    \item 求81的四次方根
\end{enumerate}
\end{enumerate}
\end{ex}

\subsection{证明三角恒等式}

在我们学习过的三角公式和三角恒等式当中，除了公式$\sin(\alpha+\beta)$与$\cos(\alpha+\beta)$是证明棣美佛定理的依据外，其余公式都可以应用复数的指数形式来加以证明.其要点是：首先将公式中的正切、余切、正割、余割诸函数都化为正弦、余弦函数表达；其次再用由复数的指数形式而导出来的公式：（例1.32）
\[\begin{split}
\sin\theta&=\frac{1}{2i}\left(e^{i\theta}-e^{-i\theta}\right)\\
    \cos\theta&=\frac{1}{2}\left(e^{i\theta}+e^{-i\theta}\right)\\
    \end{split}\]
将三角函数转化为指数形式的代数运算；从而达到证明的目的.

\begin{example}
求证：$-2\sin\alpha\cdot \sin\beta=\cos(\alpha+\beta)-\cos(\alpha-\beta)$    
\end{example}

\begin{proof}
$\because\quad \sin\alpha=\frac{1}{2i}\left(e^{i\alpha}-e^{-i\alpha}\right),\quad \sin\beta=\frac{1}{2i}\left(e^{i\beta}-e^{-i\beta}\right)$

\[\begin{split}
\therefore\quad \text{左}&=-2\cdot \frac{1}{4i^2}\left[e^{i(\alpha+\beta)}+e^{-i(\alpha+\beta)}-e^{i(\alpha-\beta)}-e^{i(\beta-\alpha)}\right]\\
&=\frac{1}{2}\left[e^{i(\alpha+\beta)}+e^{-i(\alpha+\beta)}\right]-\frac{1}{2}\left[e^{i(\alpha-\beta)}+e^{-i(\alpha-\beta)}\right]\\
&=\cos(\alpha+\beta)-\cos(\alpha-\beta)=\text{右}
\end{split}\]

$\therefore\quad $原等式成立.
\end{proof}

\begin{example}
求证：$\sin(\alpha+\beta)\cdot \cos(\alpha-\beta)=\sin\alpha\cdot \cos\alpha+\sin\beta\cdot \cos\beta$
\end{example}

\begin{proof}
\[\begin{split}
\text{左式}&=\frac{1}{2i}\left[e^{i(\alpha+\beta)}-e^{-i(\alpha+\beta)}\right]\cdot \frac{1}{2}\left[e^{i(\alpha-\beta)}+e^{-i(\alpha-\beta)}\right]\\
&=\frac{1}{4i}\left[e^{2\alpha i}+e^{2\beta i}-e^{-2\beta i}-e^{-2\alpha i}\right]
\end{split}\]
\[\begin{split}
    \text{右式}&=\frac{1}{2i}\left(e^{i\alpha}-e^{-i\alpha}\right)\cdot\frac{1}{2}\left(e^{i\alpha}+e^{-i\alpha}\right)+\frac{1}{2i}\left(e^{i\beta}-e^{-i\beta}\right)\cdot \frac{1}{2}\left(e^{i\beta}+e^{-i\beta}\right)\\
    &=\frac{1}{4i}\left[e^{2\alpha i}-e^{-2\alpha i}+e^{2\beta i}-e^{-2\beta i}\right]
\end{split}\]

$\because\quad $左式$=$右式

$\therefore\quad $原等式成立.
\end{proof}

对于具有条件$\alpha+\beta+\gamma=\pi$的三角恒等式，应用复数的指数形式加以证明更为有利.

\begin{example}
在$\triangle ABC$中，求证：
\[1+\cos 2A+\cos 2B+\cos 2C+4\cos A\cdot \cos B\cdot \cos C=0\]
\end{example}

\begin{proof}
\[\begin{split}
 \because\quad    4\cos A\cdot \cos B\cdot \cos C &=4\cdot\frac{e^{iA}+e^{-iA}}{2}\cdot\frac{e^{iB}+e^{-iB}}{2}\cdot\frac{e^{iC}+e^{-iC}}{2}\\
 &=\frac{1}{2}\left[e^{i(A+B+C)}+e^{i(B+C-A)}+e^{i(A+C-B)}\right.\\
 &\qquad   +e^{i(A+B-C)}+e^{-i(B+C-A)}+e^{-i(A+C-B)}\\
&\qquad  \left.+e^{-i(A+B-C)}+e^{-i(A+B+C)}\right]\qquad  \qquad  (*)
\end{split}\]

但由于$A+B+C=\pi$，进而有$B+C-A=\pi-2A$, 
$A+C-B=\pi-2B$, $A+B-C=\pi-2C$; 因此可得：
\[\begin{split}
    e^{i(A+B+C)}&=e^{i\pi}=-1\\
    e^{i(B+C-A)}&=e^{i(\pi-2A)}=e^{i\pi}\cdot e^{-2Ai}=-e^{-2Ai}\\
\end{split}\]
同理：$e^{i(A+C-B)}=-e^{-2Bi},\qquad e^{i(A+B-C)}=-e^{-2Ci}$
\[\begin{split}
    e^{-i(A+B+C)}=e^{-i\pi}=-1,\qquad 
    e^{-i(B+C-A)}=-e^{2Ai}\\
    e^{-i(A+C-B)}=-e^{2Bi},\qquad 
    e^{-i(A+B-C)}=-e^{2Ci}
\end{split}\]

将以上结果代入(*)式，得
\[\begin{split}
    4\cos A\cdot \cos B\cdot \cos C &=\frac{1}{2}\left[-2-e^{-2Ai}-e^{2Ai}-e^{-2Bi}-e^{2Bi}-e^{-2Ci}-e^{2Ci}\right]\\
    &=-1-\cos2A-\cos2B-\cos2C
\end{split}\]

$\therefore\quad 1+\cos 2A+\cos 2B+\cos 2C+4\cos A  \cos B  \cos C=0$
\end{proof}

\begin{ex}
利用复数的指数形式，证明以下等式：
\begin{enumerate}
    \item $\sin2\alpha=2\sin\alpha\cdot \cos\alpha$
    \item $\cos^2\beta-\sin^2\alpha=\cos(\alpha+\beta)\cdot \cos(\alpha-\beta)$
\end{enumerate}
\end{ex}

\subsection*{习题1.4}
\begin{enumerate}
    \item 把下列复数化为指数形式：
\begin{multicols}{3}
\begin{enumerate}[(1)]
    \item $5$
    \item $2+2i$
    \item $-2i$
    \item $1-\sqrt{3}i$
    \item $\pc{\frac{\pi}{4}}$
    \item $\pc{3}$
    \item $1+\pc{\frac{\pi}{3}}$
\end{enumerate}
\end{multicols}
\item 设复数$a+bi=r\cdot e^{i\theta}$，写出复数$a-bi$, $-a+bi$, $-a-bi$的指数形式.
\item 把下列复数化为三角形式及代数形式：
\begin{multicols}{3}
\begin{enumerate}[(1)]
    \item $4e^{i\tfrac{\pi}{6}}$
    \item $\sqrt{3}e^{i\tfrac{2\pi}{3}}$
    \item $e^{-i\tfrac{3\pi}{2}}$
\end{enumerate}
\end{multicols}

\item 试求：复数$e^{i\tfrac{4\pi}{5}}$, $e^{i\tfrac{2\pi}{3}}$ 所对应的向量间的夹角$\alpha\; (0\le \alpha\le \pi)$.
\item 用复数的指数形式计算：
\begin{enumerate}[(1)]
    \item $\sqrt{2}\left(\pc{\frac{4\pi}{3}}\right)\cdot\frac{\sqrt{2}}{4}\left(\pc{\frac{5\pi}{6}}\right)$
    \item $\sqrt{3}(\pc{150^{\circ}})\div \sqrt{2}(\pc{225^{\circ}})$
\end{enumerate}

\item 根据欧拉公式求证：$e^{-i\theta}=(\pc{\theta})^{-1}$.

\item 用复数的指数形式计算：
\begin{multicols}{3}
\begin{enumerate}[(1)]
    \item $(1+\sqrt{3}i)^{10}$
    \item $(\sqrt{3}-i)^8$
    \item $\sqrt[4]{64+0i}$
\end{enumerate}
\end{multicols}

\item 用复数的指数形式证明恒等式：
\begin{enumerate}[(1)]
    \item $\sin\alpha-\sin\beta=2\cos\frac{\alpha+\beta}{2}\cdot \sin\frac{\alpha-\beta}{2}$
    \item $(\sin\alpha+\sin\beta)^2+(\cos\alpha+\cos\beta)^2=4\cos^2\frac{\alpha-\beta}{2}$
\end{enumerate}
\end{enumerate}

\section*{本章内容要点}

一、本章里，我们引进了虚数单位($i^2=-1$)，把数的概念扩展到了复数集. 复数的分类如下：
\begin{center}
\begin{tikzpicture}
\node (A) at (0,0)[text width=2cm, align=center]{复数$a+bi$\\$(a,b\in\R)$};
\node (B1) at (2.5,1)[text width=2cm, align=center]{实数\\($b=0$时)}; 
\node (B2) at (2.5,-1)[text width=2cm, align=center]{虚数\\($b\ne 0$时)}; 
\node (C1) at (5,0)[text width=2cm, align=center]{纯虚数\\($a=0$时)}; 
\node (C2) at (5,-2)[text width=2cm, align=center]{非纯虚数\\($a\ne 0$时)}; 
\draw[decorate, decoration={brace, amplitude=5pt}](1.5,-1.3)--(1.5,1.3);
\draw[decorate, decoration={brace, amplitude=5pt}](3.8,-2.3)--(3.8,.3);


\end{tikzpicture}
\end{center}

且复数集$\supset$实数集$\supset$有理数集$\supset$整数集$\supset$自然数集，即：
\[\mathbb{C} \supset \R \supset \mathbb{Q} \supset \Z\supset\N\]

二、复数有四个基本概念：模，幅角，共轭复数，复数的相等；任一个复数有三种表示形式：代数形式$a+bi\; (a,b\in\R)$，三角形式$r(\pc{\theta})$，其中$r=\sqrt{a^2+b^2}$, $\tan\theta =\frac{b}{a}$，指数形式$re^{i\theta}$；复数$z=a+bi$与复平
面内的点$Z(a,b)$对应，并且和从原点出发的向量$\VEC{OZ}$相对应，因此，复数集$\mathbb{C}$与复平面内的点集，与从原点出发的向量集合之间成一一对应关系，即
\[\mathbb{C}:\; \{a+bi\mid a,b\in\R,\; i^2=-1\} \longleftrightarrow \{Z(a,b)\}\longleftrightarrow \{\VEC{OZ}\}\]

据此，一些平面几何，平面解析几何以及平面向量的问题，也可以通过复数来解决.

三、复数集对于四则运算及乘方，开方运算都是封闭的；复数集保持了数系运算的“通性”，即加法、乘法的交换律、结合律以及乘法对于加法的分配律、数0与1的运算特性等；但复数集不再具有顺序性，对任意两个不全为实数的复数，就不能比较大小.

四、复数的运算法则及其几何意义：
\begin{enumerate}
    \item 加、减法运算用复数的代数形式表达比较方便，即
    \[(a+bi)\pm (c+di)=(a\pm c)+(b\pm d)i\]
    其几何意义是对应的向量相加、减，即“平行四边形”法则和“三角形”法则，特别是两复数的差的模表示复平面内两点间的距离.
    \item 乘、除法，乘方、开方的运算用复数的三角形式或指数形式表达比较方便，即
\[\begin{split}
    r_1(\pc{\theta_1})\cdot r_2(\pc{\theta_2})&=r_1\cdot r_2\left[\pcx{\theta_1+\theta_2}\right]\\
\frac{r_1(\pc{\theta_1})}{r_2(\pc{\theta_2})}&=\frac{r_1}{r_2}\left[\pcx{\theta_1-\theta_2}\right]
\end{split}\]
棣美佛定理：
\[[r(\pc{\theta})]^n=r^n(\pc{n\theta})\]
$n$个$n$次方根：
\[\begin{split}
\sqrt[n]{r(\pc{\theta})}&=\sqrt[n]{r}\left(\pc{\frac{\theta+2k\pi}{n}}\right)\\
&\qquad (k=0,1,2,\ldots,n-1)\quad (n\in\N)    
\end{split} \]
\end{enumerate}

复数的这些运算，都有其几何意义.特别是：复数的相乘对应着平面向量的旋转；$n$个$n$次方根对应着圆内接正$n$边形的顶点，即圆周的$n$等分点，这些结论在应用上很重要.

五、在复数集中，本章学习了在实数集中所不能解的有关问题：
\begin{enumerate}
    \item 当$b^2-4ac<0$时，一元二次方程$ax^2+bx+c=0\; (a\ne 0)$有两个共轭的虚根，即
\[x=\frac{-b\pm i\sqrt{4ac-b^2}}{2a}\]
\item 二项方程$x^n=b\; (b\in\mathbb{C})$有$n$个不同的复数根，即设$x^n=b=r(\pc{\theta})$，则
\[x=\sqrt[n]{r}\left(\pc{\frac{\theta+2k\pi}{n}}\right)\]
当$k=0,1,2,\ldots,n-1$时，可得出$n$个不同的解.

\item 特别地，方程$x^n=1$的$n$个复数解，叫做$n$次单位根，若设$\omega=\pc{\frac{2\pi}{n}}$，则当$n$为偶数时，$n$个单位根为
\[1,\; \omega,\; \omega^2,\;\ldots,\;  \omega^{\tfrac{n}{2}-1},\;-1,\;  \overline{\omega^{\tfrac{n}{2}-1}},\; \ldots, \;  \overline{\omega^{2}},\; \overline{\omega}\]
当$n$为奇数时，$n$个单位根为
\[1,\; \omega,\; \omega^2,\;\ldots,\;  \omega^{\tfrac{n-1}{2}},\;  \overline{\omega^{\tfrac{n-1}{2}}},\; \ldots, \;  \overline{\omega^{2}},\; \overline{\omega}\]

\item 实系数一元方程，若有虚根存在，则一定是成对出现，即实系数方程若有根$a+bi\;(b\ne 0)$则必有根$a-bi$.
\end{enumerate}

\section*{复习题一}
\begin{enumerate}
    \item 在复平面上任选一点$Z$（不在原点）表示复数$z$，然后用几何法在平面找出下列各复数的点：
\begin{multicols}{4}
\begin{enumerate}[(1)]
    \item $-z$
    \item $\bar z$
    \item $-\bar z$
    \item $z+\bar z$
    \item $z-\bar z$
    \item $z+|z|$
    \item $z-|z|$
    \item $iz$
    \item $-iz$
    \item $i\bar z$
    \item $-i\bar z$
\end{enumerate}
\end{multicols}

\item 已知$z=a+bi\; (a,b\in\R)$，求下列各复数的实部与虚部：
\begin{multicols}{2}
\begin{enumerate}[(1)]
    \item $z^2$
    \item $z^3$
    \item $\frac{1}{z}$
    \item $v_0 z+\frac{M}{2\pi}\cdot \frac{1}{z}\quad (v_0,M\in\R)$
\end{enumerate}
\end{multicols}

\item 求复数$z=\frac{a^2-b^2+2abi}{ab\sqrt{2}+\sqrt{a^4+b^4}i}$的模.

\item 求证：表示复数$z_0$的点关于直线${\rm Re}(z)={\rm Im}(z)$的对称点表示的是复数$i\bar z_0$.

\item 下列方程（$t$为实参数）给出了怎样的曲线？
\begin{multicols}{2}
\begin{enumerate}[(1)]
    \item $z=t(1+i)$
    \item $z=a\cos t+ib\sin t$
    \item $z=t+\frac{i}{t}$
    \item $z=t^2+\frac{i}{t}$
\end{enumerate}
\end{multicols}

\item 已知$(x+yi)^3=a+bi$，$a,b,x,y\in\R$，求证：
$\frac{a}{x}+\frac{b}{y}=4(x^2-y^2)$

\item 若$a$，$b$是共轭复数，且$(a+b)^2-3abi=4-6i$, 试求$a$及$b$.
\item 若$x,y\in\R$, $z=x+y-30-xyi$且$\bar z=60i-|x+yi|$，试求$x$和$y$的值.
\item 若$z=x+yi$，求证：$|z|\le |x|+|y|\le \sqrt{2}|z|$
\item 试判断以下复数，哪些是实数？哪些是纯虚数？($z,z_1,z_2\in\mathbb{C}$)
\begin{multicols}{2}
\begin{enumerate}[(1)]
    \item $z+\bar z$
    \item $z-\bar z$
    \item $z\cdot \bar z$
    \item $z^2-\bar z^2$
    \item $\frac{\left(\frac{1}{z}+\frac{1}{\bar z}\right)(z+\bar z)}{z-\bar z}$
    \item $z_1\cdot \bar z_2-\bar z_1\cdot z_2$
    \item $z_1\cdot \bar z_2+\bar z_1\cdot z_2$
    \item $\frac{z_1\cdot \bar z_2+\bar z_1\cdot z_2}{z_1\cdot \bar z_1-1}$
    \item $\frac{z_1\cdot \bar z_2-\bar z_1\cdot z_2}{i(z_1\cdot \bar z_1+z_2\cdot \bar z_2)}$
    \item $\frac{i(z_1\cdot \bar z_2+z_2\cdot \bar z_1)}{z_1\cdot \bar z_2-\bar z_1\cdot z_2}$
\end{enumerate}
\end{multicols}

\item 设复数$z_1,z_2$分别对应于点$Z_1,Z_2$，求证：
\begin{enumerate}[(1)]
    \item $OZ_1\bot OZ_2$的充要条件是$z_1\cdot \bar z_2+\bar z_1\cdot z_2=0$
    \item $OZ_1$与$OZ_2$共线的充要条件是$z_1\cdot \bar z_2-\bar z_1\cdot z_2=0$
\end{enumerate}

\item 设$\alpha$是复常数，$k$是实常数，$z\in\mathbb{C}$，试说明下述方程在复平面上的图形是什么？
\begin{multicols}{2}
\begin{enumerate}[(1)]
    \item $\alpha\bar z+\bar\alpha\cdot z=k$;
    \item $z\bar z+z\bar\alpha+\alpha\bar z+k=0$.
\end{enumerate}
\end{multicols}
\item 求证：如果$\frac{1}{z}+z\in\R$，那么${\rm Im}(z)=0$或$|z|=1$.
\item 化简：$\left(\frac{1-i}{1+i}\right)^{4\alpha}$，其中$\alpha=\csc10^{\circ}-\sqrt{3}\sec10^{\circ}$.
\item 解方程组
\begin{multicols}{2}
\begin{enumerate}[(1)]
    \item $\begin{cases}
        x+iy=2+4i\\ ix+y=0
    \end{cases}$
    \item $\begin{cases}
        z_1+2z_2=1+i\\  3z_1+iz_2=2-3i
    \end{cases}$
    \item $\begin{cases}
x^2+y^2=6\\        
y^2+z^2=0\\
x^2+z^2=8i\\       
    \end{cases}$
\end{enumerate}
\end{multicols}

\item 化简
\[\frac{(\cos2\theta-i\sin2\theta)(\pc{\varphi})^2}{\pcx{\theta+\varphi}}\cdot \frac{(\pc{2\theta})^2(\cos2\varphi-i\sin2\varphi)}{\pcx{\theta-\varphi}}\]

\item 将下列复数化为三角形式：
\begin{multicols}{2}
\begin{enumerate}[(1)]
    \item $\frac{(\pc{5\varphi})^2}{(\cos3\varphi-i\sin3\varphi)^3}$
    \item $1-\cos\theta+i\sin\theta,\quad (0\le\theta\le\pi)$
\end{enumerate}
\end{multicols}

\item 设$x=e^{i\alpha}$, $y=e^{i\beta}$, $z=e^{i\gamma}$，求证：
\[(x+y)(y+z)(z+x)=8xyz\cos\frac{\beta-\gamma}{2}\cdot \cos\frac{\gamma-\alpha}{2}\cdot \cos\frac{\alpha-\beta}{2}\]

\item 设$\sin A+\sin B+\sin C=\cos A+\cos B+\cos C=0$，求证：
\begin{enumerate}[(1)]
    \item $\sin3A+\sin3B+\sin3C=3\sin(A+B+C)$
\item $\cos3A+\cos3B+\cos3C=3\cos(A+B+C)$
\end{enumerate}
\item 在复平面上，$z_1=1$, $z_2=2+i$与$z_3$组成一个正三角形，求$z_3$.
\item 若$z_1,z_2,z_3,z_4$是一个正方形的四个顶点，且$z_1=1+2i$, $z_2=3-5i$, 求$z_3,z_4$.
\item 求证：$\triangle z_1z_2z_3$是等边三角形的充要条件是：
\[z_1^2+z_2^2+z_3^2=z_1\cdot z_2+z_2\cdot z_3+z_3\cdot z_1\]
\item 下列条件表示什么样图形？
\begin{enumerate}[(1)]
    \item ${\rm Im}(z^2)=2$
    \item $|z|=1$且$\arg z\in(0,\pi)$
    \item $0<a\le {\rm Im}(z)\le b$且$\frac{\pi}{4}\le \arg z\le \frac{3\pi}{4}$
\end{enumerate}
\item 你能用复数表示出图中的半圆来吗？（包括边界）

\begin{minipage}{.45\textwidth}
    \centering
\begin{tikzpicture}[>=stealth]
\draw[->](-1.5,0)--(1.5,0)node[right]{$x$};
\draw[->](0,-1)--(0,2)node[right]{$y$};
\draw[pattern=north east lines](1,0)node[below]{2} arc (0:180:1)node[below]{$-2$} --cycle;
\node[below left]{$O$};
\node at (0,1)[above right]{2};    
\end{tikzpicture}
\captionof{figure}{ }
\end{minipage}\hfill
\begin{minipage}{.45\textwidth}
    \centering
\begin{tikzpicture}[>=stealth]
    \draw[->](-1.5,0)--(1.5,0)node[right]{$x$};
\draw[->](0,-1.5)--(0,1.5)node[right]{$y$};
\draw[pattern=north east lines](0,1)node[right]{1} arc (90:180+90:1)node[right]{$-1$}--cycle; 
\node[below right]{$O$};   
\node at (-1,0)[below left]{$-1$};   
\end{tikzpicture}
\captionof{figure}{ }
\end{minipage} 

\item 设$z\in\mathbb{C}$，解方程：
\begin{multicols}{2}
\begin{enumerate}[(1)]
    \item $\frac{1}{2}(z-1)=\frac{\sqrt{3}}{2}(1+z)i$
    \item $(z+1)^8=(1+i)^8$
\end{enumerate}
\end{multicols}

\item 设$z$是虚数，解方程：
\begin{multicols}{2}
\begin{enumerate}[(1)]
    \item $z+|\bar z|=2+i$
    \item $z^2=\bar z$
\end{enumerate}
\end{multicols}

\item 求证：
\begin{enumerate}[(1)]
    \item $|z|=1\; (z\in\mathbb{C})$的充要条件是$z^{-1}=\bar z$;
    \item 共轭虚数的$n\; (n\in\N)$次幂仍是共轭虚数；
    \item 虚数的平方根仍是虚数．
\end{enumerate}

\item 在复数集中解下列方程组：
\begin{enumerate}[(1)]
    \item $\begin{cases}
        x^2+y^2-2xy+3(x+y)-2=0\\
        2(x^2+y^2)-2xy-2(x+y)+15=0
    \end{cases}$
    \item $\begin{cases}
        \frac{x-y}{1+xy}=\frac{1}{3}\\  \frac{x+y}{1-xy}=3
    \end{cases}$
\end{enumerate}    

\item 设三次单位根为$1,\omega,\omega^2$，其中
$\omega=-\frac{1}{2}+\frac{\sqrt{3}}{2}i$，求证：
\begin{enumerate}[(1)]
    \item $1+\omega+\omega^2=0$
    \item $1\cdot \omega\cdot \omega^2=1$
    \item $(1+\omega^2)^4=\omega$
    \item $\begin{vmatrix}
        1&1&\omega\\
        1&1&\omega^2\\
        \omega^2&\omega&1
    \end{vmatrix}=-3$
    \item $1^n+\omega^n+(\omega^2)^n=\begin{cases}
        3,& n=3k\\
        0,& n\ne 3k
    \end{cases}(k\in\Z)$

    \item $a^3+b^3=(a+b)(a\omega+b\omega^2)(a\omega^2+b\omega);$
    \item $a^{3}+b^{3}+c^{3}-3abc =(a+b+c)(a+b\omega+c\omega^{2})(a+b\omega^{2}+c\omega)$
\end{enumerate}

\item 利用上题中$\omega$的性质，计算：
\begin{multicols}{2}
\begin{enumerate}[(1)]
    \item $(-1+\sqrt{3}i)^{6}-(-1-\sqrt{3}i)^{6}$;
    \item $\left(\frac{\sqrt{3}+i}{2}\right)^{6}-\left(\frac{\sqrt{3}-i}{2}\right)^{6}$;
    \item $\left(-\frac{1}{2}i-\frac{\sqrt{3}}{2}\right)^{12}$.
\end{enumerate}    
\end{multicols}

\item 已知$f(x)=x^4+3x^3-30x^2+366x-340$的一个根为$3+5i$，求其余的三个根。
\item 若实系数方程$x^4+px^3+qx^2+rx+s=0$有两个根为$2+i$, $-2+i$. 求$p$、$q$、$r$、$s$的值。
\item 若实系数多项式$x^4+px^3+qx^2+rx+s$有一个根是纯虚数，求证：$r^2+p^2s=pqr$.

\item 已知实系数方程$x^2+mx+n=0$的一个根是另一个根的$i$倍，试求$m$、$n$的关系.
\item 已知$f(x)=x^4+(2+i)x^3+(3+2i)x^2-(4-3i)x-4i$, 

$g(x)=x^4-(1-i)x^3-(1+i)x^2+(4+i)x-4i$，求$f(x)$与$g(x)$的公根.
\end{enumerate}

\chapter{排列与组合~~二项式定理}
\section{排列与组合}
\subsection{加法原理和乘法原理}
假如我们要从甲地到乙地，可以乘火车、轮船或公共汽车，火车每天有两班，轮船有两班，公共汽车有六班，问从甲地到乙地有多少种不同的方法？

因为，从甲地到乙地，乘火车有两种不同的方法，乘轮船有两种不同的方法，乘公共汽车有六种不同的方法.而每种方法都可以由甲地到达乙地，因此，从甲地到乙地共有
$2+2+6=10$
种不同的方法.

一般地说，有如下原理：
\begin{blk}{加法原理}
如果完成某件事有$n$种方式，在第一种方式中有$m_1$种方法，在第二种方式中有$m_2$种方法……，在第$n$种方式中有$m_n$种方法. 那么完成这件事共有$m_1+m_2+\cdots+m_n$种不同的方法.
\end{blk}

又假如，要从甲地到丙地，必须经过乙地，现在已知由甲地到乙地有三条道路，由乙地到丙地有两条道路（图2.1），问从甲地到丙地共有多少种不同的走法？
\begin{figure}[htp]
    \centering
\begin{tikzpicture}
\node[draw, circle] (A) at (0,0){甲};    
\node[draw, circle] (B) at (4,0){乙};    
\node[draw, circle] (C) at (8,0){丙};
\draw(A)--(B);
\draw(A)to [bend left=45](B);
\draw(A)to [bend left=-45](B);
\draw(C)to [bend left=15](B);
\draw(C)to [bend left=-15](B);

\end{tikzpicture}
    \caption{}
\end{figure}

因为从甲地到乙地有3种不同的走法，到乙地后又各有两种不同的走法到丙地，因此，从甲地到丙地共有
$3\x2=6$
种不同的走法.

一般地说，有如下原理：

\begin{blk}
    {乘法原理}如果完成某种事需要分成$n$个步骤，第一步骤有$m_1$种方法，第二步骤有$m_2$种方法……，第$n$步骤有$m_n$种方法，那么完成这件事共有$m_1\cdot m_2\cdots m_n$种不同的方法.
\end{blk}

总的来说，如果完成一种事有几种不同的方法，这些方法又是互不相关的，任选一种方法都能完成这件事的，那么完成这件事的方法的总数，应该用加法计算. 如果完成一件事，必须分成若干步骤，每个步骤又有不同的方法，而且每一步骤中的一种方法完成之后，都可以开始下一步骤的工作，依次完成全部步骤，才能达到完成这一件事的目的，那么完成这件事的方法的总数应该用乘法计算.

\begin{ex}
\begin{enumerate}
    \item 完成某件工作，甲有两种不同的方法，乙有3种不同方法，丙有4种不同方法．从三人中选一人完成这件工作，共有几种不同选法？
    \item 一种车床用两个手柄联合控制转速，一个手柄有两档，另一个手柄有3档，问能控制几种不同的转速？
    \item 乘积$(a+b)\cdot (c+d+e)\cdot (m+n+p+q)$展开后共有多少项？
    \item 从甲地到乙地有一条道路，从乙地到丁地有两条道路；从甲地到丙地有三条道路，从丙地到丁地有四条道路.若从甲地到丁地必须经过乙地或丙地，问从甲地到丁地共有多少种不同的走法？
\end{enumerate}
\end{ex}


\subsection{排列}
一般地说，从$n$个不同的元素中任意选取$m\; (m\le n)$个元素，按照一定的顺序排成一列，叫做从$n$个元素中取出$m$个元素的一个排列.

例如：由$1,2,3$三个数字中选取两个数字写成两位数，那么$1,2,3$都被看作是排列的元素。而12与21显然都是由$1,2$两个元素得到的不同的排列．同样，12，21，13，31，23，32等都被认为是不同的排列．只有用相同的元素，又是按相同的顺序排列而成的排列，才叫做相同的排列，例如：123与123就是相同的排列．

\begin{example}
    由数字1,2,3,4可以组成多少个没有重复数字的两位数？
\end{example}

\begin{solution}
    要排出两位数，第一步先要排出十位上的数字（当然也可以先排个位上的数字），第二步是确定个位上的数字.

显然，第一步在十位的位置上，$1,2,3,4$都可以放，故有四种不同的方法.

第二步是放个位上的数字，由于在第一步结束时，剩下的只有3个数字，因此第二步只能有3种不同的选择．因此个位与十位上的数的排列工作都结束后，这个两位数才算排出来，所以，用乘法原理计算得：$4\x3=12$．即以$1,2,3,4$可以排出12个不同的两位数．
\end{solution}

\begin{center}
\begin{tikzpicture}
\begin{scope}
\draw(-.5,-.25) rectangle (.5,.25);
\draw(0,-.25)--(0,.25);
\node at (-0.25,0){1};
\foreach \x/\y in {-1/4,0/3,1/2}
{
    \draw(1, \x-.25) rectangle (2, \x+.25);
    \draw(1.5,\x-.25)--(1.5,\x+.25);
    \node at (1.25,\x){1};
    \node at (1.75,\x){\y};
}
\draw[decorate, decoration={brace, amplitude=5}](.9,-1.25)--(.9,1.25);
\end{scope}
\begin{scope}[xshift=3cm]
    \draw(-.5,-.25) rectangle (.5,.25);
    \draw(0,-.25)--(0,.25);
    \node at (-0.25,0){2};
    \foreach \x/\y in {-1/4,0/3,1/1}
    {
        \draw(1, \x-.25) rectangle (2, \x+.25);
        \draw(1.5,\x-.25)--(1.5,\x+.25);
        \node at (1.25,\x){2};
        \node at (1.75,\x){\y};
    }
    \draw[decorate, decoration={brace, amplitude=5}](.9,-1.25)--(.9,1.25);
    \end{scope}
    \begin{scope}[xshift=6cm]
        \draw(-.5,-.25) rectangle (.5,.25);
        \draw(0,-.25)--(0,.25);
        \node at (-0.25,0){3};
        \foreach \x/\y in {-1/4,0/2,1/1}
        {
            \draw(1, \x-.25) rectangle (2, \x+.25);
            \draw(1.5,\x-.25)--(1.5,\x+.25);
            \node at (1.25,\x){3};
            \node at (1.75,\x){\y};
        }
        \draw[decorate, decoration={brace, amplitude=5}](.9,-1.25)--(.9,1.25);
        \end{scope}
\begin{scope}[xshift=9cm]
\draw(-.5,-.25) rectangle (.5,.25);
\draw(0,-.25)--(0,.25);
\node at (-0.25,0){4};
\foreach \x/\y in {-1/3,0/2,1/1}
{
    \draw(1, \x-.25) rectangle (2, \x+.25);
    \draw(1.5,\x-.25)--(1.5,\x+.25);
    \node at (1.25,\x){4};
    \node at (1.75,\x){\y};
}
\draw[decorate, decoration={brace, amplitude=5}](.9,-1.25)--(.9,1.25);
\end{scope}
\end{tikzpicture}
\end{center}




\begin{example}
$a,b,c,d$四个元素，每次取出三个元素的不同排列有多少种，并写出所有的排列.（在同一排列里，元素不能重复出现.）
\end{example}

\begin{solution}
要把取出的三个元素排成一列，要分三步完成. 第一步先在第一个位置上放一个元素，应该有四种不同的方法；第二步要在第二个位置上放一个元素，由于第一步的工作完成后剩下的是三个元素，从三个元素中选一个放在第二个位置上应有3种不同的方法；第三步要在第三个位置上放一个元素，由于前两步工作完成后只剩下两个元素，所以只有两种不同的方法. 三步工作都完成后才完成排列工作，因此应该用乘法原理计算排列的总数，所以共有：$4\x3\x2=24$种不同的排列．即有如下的24种不同的排列：
\begin{center}
\begin{tabular}{cccccc}
$abc$& $abd$& $acb$ &$acd$ &$adb$& $adc$\\
$bac$& $bad$ &$bca$& $bcd$ &$bda$& $bdc$\\
$cab$& $cad$ &$cba$& $cbd$ &$cda$& $cdb$\\
$dab$& $dac$& $dba$ &$dbc$& $dca$ &$dcb$    
\end{tabular}
\end{center}
\end{solution}

一般地说，从$n$个不同元素中，任取$m$个$(m\le n)$元素（不许重复），按照一定顺序排成一列的个数，叫做从$n$个不同元素中取出$m$个元素的排列数，用符号表示$\perm{m}{n}$.

\begin{blk}
  {定理1} 从$n$个不同元素中，任取$m$个$(m\le n)$没有重复的元素的排列数为$\perm{m}{n}=n(n-1)\cdots (n-m+1)$
\end{blk}

\begin{proof}
    在取出$m$个元素排成一列时，首先在第一个位置上可以选$n$个元素中的任何一个，故有$n$种方法. 在第二个位置上可以选余下的$n-1$个元素中的任一个．如此类推，在第$m$个位置上可以选余下的$(n-m+1)$个元素中的任一个. 根据乘法原理得到总的不同排列数应该是
\[\perm{m}{n}=n(n-1)\cdots(n-m+1)\]

这里的$m$是自然数，并且$m\le n$, 这个公式就叫做排列数公式.
\end{proof}

特别是当$m=n$时 有
\[\perm{n}{n}=n(n-1)\cdots3\cdot2\cdot1\]
这表示，$n$个不同的元素全部取出参加排列的排列数，叫做$n$个不同元素的全排列，自然数1到$n$的连乘积叫做$n$的阶乘，用$n!$表示，所以全排列数公式又可以写成
\[\perm{n}{n}=n!\]
有了阶乘，排列数公式又可以写成
\[\perm{m}{n}=\frac{n!}{(n-m)!}\]
这里，当$m=n$时，分母出现$(n-m)!$就变成$0!$，为了使公式仍能成立，我们特别规定$0!=1$.

例如$\perm{5}{8}=8\cdot  7\cdot 6\cdot 5\cdot 4=6720$，
$\perm{5}{5}=5!=5\cdot 4\cdot 3\cdot 2\cdot 1=120$

\begin{example}
    试证明下列等式成立．
\begin{multicols}{2}
\begin{enumerate}[(1)]
    \item $\perm{m}{n}+m\cdot \perm{m-1}{n}=\perm{m}{n+1}\quad (m\le n)$
\item $\perm{n+1}{n+1}-\perm{n}{n}=n^2\cdot \perm{n-1}{n-1}$
\end{enumerate}
\end{multicols}
\end{example}

\begin{proof}
\begin{enumerate}[(1)]
    \item \[\begin{split}
        \perm{m}{n}+m\cdot \perm{m-1}{n}&=\frac{n!}{(n-m)!}+m\cdot\frac{n!}{(n-m+1)!}\\
        &=\frac{(n-m+1)\cdot n!+m\cdot n!}{(n-m+1)!}=\frac{(n+1)\cdot n!}{(n-m+1)!}\\
        &=\frac{(n+1)!}{[(n+1)-m]!}=\perm{m}{n+1}
    \end{split}\]
所以原等式成立.

\item \[\begin{split}
    \perm{n+1}{n+1}-\perm{n}{n}&=(n+1)!-n!=(n+1)\cdot n!-n!\\
    &=n\cdot n!=n^2\cdot (n-1)!\\
    &=n^2\cdot \perm{n-1}{n-1}
\end{split}\]
所以原等式成立.
\end{enumerate}
\end{proof}

\begin{example}
    用0到9十个数字，可以写出多少个没有重复数的三位数？其中有多少个是偶数？
\end{example}

\begin{solution}
\begin{enumerate}[(1)]
    \item 第一步，因为0不能放在百位的位置上，所以百位上的数有9种不同的选法，
    
第二步是从余下的九个数字中任选两个数字放在十位个位上，这有$\perm{2}{9}$种方法.所以
\[9\cdot \perm{2}{9}=9\cdot 9\cdot 8=648\]
\item 要得到三位偶数，必须个位上的数是0, 2, 4, 6或8．但是0又不能排在百位上，这说明0怎样放在适合的位置上是比其他问题要求更高的. 所以我们还是把0元素的题先解决.

第一类，0放在个位上这一选法确定后，其余百位、十位可以从余下的九个数字中任选两个数字放上去，这有$\perm{2}{9}$种方法.

第二类，个位上可放2,4,6,8中的任一个，在这类方法中，第一步在个位上有4种选法，第二步，在百位上可以放除0以外所余的八个数字的任一个，故有八种选择方法．第三步，余下的八个数字均可以放在十位上，所以
\[\perm{2}{9}+4\x8\x\perm{1}{8}=72+256=328\]
\end{enumerate}
答：有648个没有重复数字的三位数，其中有328个是偶数.

\end{solution}

关于例2.4中的第一个问题，我们还可以有如下的解法：

从0到9十个数字中任选三个数字排列总数为$\perm{3}{10}$．但是其中0在百位上的数有$\perm{2}{9}$个. 所以
\[\perm{3}{10}-\perm{2}{9}=10\cdot 9\cdot 8-9\cdot 8=648\]

\begin{ex}
\begin{enumerate}
    \item 写出从$a,b,c,d,e$中任取两个不同元素的所有排列.
    \item 计算：
\begin{multicols}{3}
\begin{enumerate}[(1)]
    \item $\perm{3}{7}$
    \item $\perm{3}{15}$
    \item $\perm{4}{4}$
    \item $\perm{4}{8}-3\cdot\perm{2}{8}$
    \item $\perm{3}{7}-5\perm{2}{7}$
    \item $\frac{\perm{7}{12}}{\perm{6}{11}}$
\end{enumerate}
\end{multicols}

\item 求证：
\begin{enumerate}[(1)]
    \item $n!=\frac{(n+1)!}{n+1}$
    \item $\perm{m}{n}=\frac{\perm{n+2}{n+2}}{\perm{2}{n+2}\cdot \perm{n-m}{n-m}}\quad (m\le n)$
\end{enumerate}

\item 某一段铁路共有8个车站，需要为这段铁路印制多少种不同的客车票？（都是普通座位票）
\item 6个同学排成一列，有多少种不同的排法？
\item 从12人的1个学习小组中，选1个正组长和1个副组长，问有多少种不同的选法？
\item 用0至5共六个数字，可以组成多少不同的三位数？
\item 用0到9这十个数字，其中有多少个没有重复数字的二位数，其中有多少个二位奇数？
\item 某信号兵用红、黄、蓝三面旗子从上到下挂在竖直的旗杆上表示信号，每次可以任挂一面、二面或三面，并且不同的顺序表示不同的信号，问一共可以表示多少种不同的信号？
\item 有五道不同的工序分配给五个工人工作，其中有一个工人不会做其中两道工序，问有多少种不同的分配法？
\item 由2, 3, 4, 5所组成的没有重复数字的四位数，
\begin{enumerate}[(1)]
\item 求所有这些四位数里各个数字的和；
\item 求所有这些四位数的和。
\end{enumerate}
\end{enumerate}
\end{ex}

\subsection{组合}
我们知道，从一个学习小组中选出正、副组长各1人，这种选法是排列问题.如果改变为从一个学习小组中选出代表2人参加学校的学代会，有多少种不同的选法？”这时的代表无所谓正、副之分，就无先后顺序的意义了，也就不成为排列的问题了.

又例如从1, 2, 3, 4, 5五个数选出两个数字相乘，问有多少个不同的乘积？那么只要所选出的数字是3和5；它们的积只有15，至于是$5\x3$还是$3\x5$，那是不要求加以区别的，也是没有排列次序的区别的.

一般地说，从$n$个不同的元素中，任取$m\; (m\le n)$个元素并成一组，叫做从$n$个不同元素中取出$m$个元素的一个组合.

这里要注意，元素$a$，$b$，$c$的组合与$a,c,b$; $c,a,b$; $b,a,c$; $b,c,a$或$c,b,a$都被看作是相同的组合了. 对于排列来说，上列的是六个不同的排列，而对于组合来说，只算是一个组合，那么组合数是怎样进行计算的呢？

从$n$个不同元素中取出$m\; (m\le n)$个元素的所有组合的个数，叫做从$n$个不同元素中取出$m$个元素的组合数，用符号$\comb{m}{n}$表示。

例如，从四个不同元素中取出3个元素进行排列，有多少种不同的排法？

我们已经知道应该有$\perm{3}{4}=24$种不同的排法，现在我们
想从另外的方法入手解决这个问题.

第一步，从$a,b,c,d$中取三个元素作一组，这样可以得到$a,b,c$; $a,b,d$; $a,c,d$; $b,c,d$四个组，这就是说$\comb{3}{4}=4$.

第二步，把每组的三个元素作全排列，应有$\perm{3}{3}$个不同排列.

根据乘法原理得$\perm{3}{4}=4\cdot \perm{3}{3}$，即$\perm{3}{4}=\comb{3}{4}\cdot \perm{3}{3}$.

又如，从五个不同的元素中取出2个元素的排列数是$\perm{2}{5}$. 如果我们从五个不同元素中取出两个分成一组，我们可以组成如下的小组：

$a,b$; $a,c$; $a,d$; $a,e$; $b,c$; $b,d$; $b,e$; $c,d$, $c,e$和$d,e$等10个组，这就是$\comb{2}{5}$；第二步把每组元素进行全排列应是$\perm{2}{2}$，这样应该有$\perm{2}{5}=\comb{2}{5}\cdot \perm{2}{2}$. 因此，一般地说，从$n$个不同元素中取出$m\; (m\le n)$个元素的排列数$\perm{m}{n}$可以等于$\comb{m}{n}\cdot \perm{m}{m}$.


\begin{blk}
  {定理2} 从$n$个不同元素中任取$m$个（$m\le n$）没有重复元素组合数为$\comb{m}{n}=\frac{n!}{m!(n-m)!}$.  
\end{blk}

\begin{proof}
    由$\perm{m}{n}=\comb{m}{n}\cdot \perm{m}{m}$，所以
\[\comb{m}{n}=\frac{\perm{m}{n}}{\perm{m}{m}}=\frac{\frac{n!}{(n-m)!}}{m!}=\frac{n!}{m!(n-m)!}\]
\end{proof}
注意：$\comb{m}{n}=\frac{n(n-1)\cdots(n-m+1)}{m!}$

\begin{example}
计算$\comb{3}{8}$, $\comb{5}{8}$和$\comb{2}{9}$
\end{example}

\begin{solution}
    根据组合数公式有：
\[\begin{split}
    \comb{3}{8}&=\frac{8\x 7\x 6}{3\x 2\x 1}=56\\
    \comb{5}{8}&=\frac{8\x 7\x 6\x 5\x 4}{5\x 4\x 3\x 2\x 1}=56\\
    \comb{2}{9}&=\frac{9\x 8}{2\x 1}=36
\end{split}\]
\end{solution}

组合数有以下性质：
\begin{blk}
{性质1}$$\comb{m}{n}=\comb{n-m}{n}$$
\end{blk}

\begin{proof}
根据组合数公式：
\[\begin{split}
    \comb{m}{n}&=\frac{n!}{m!(n-m)!}\\
    \comb{n-m}{n}&=\frac{n!}{(n-m)![n-(n-m)]!}=\frac{n!}{m!(n-m)!}
\end{split}\]
所以$\comb{m}{n}=\comb{n-m}{n}$.
\end{proof}

\begin{blk}{性质2}
\[\comb{m}{n+1}=\comb{m}{n}+\comb{m-1}{n}\]
\end{blk}

\begin{proof}
根据组合数公式：
\[\begin{split}
    \comb{m}{n+1}&=\frac{(n+1)!}{m!(n-m+1)!}\\
    \comb{m}{n}+\comb{m-1}{n}&=\frac{n!}{m!(n-m)!}+\frac{n!}{(m-1)!(n-m+1)!}\\
    &=\frac{(n-m+1)\cdot n! +m\cdot n!}{m!(n-m+1)!}=\frac{n!\cdot (n-m+1+m)}{m!(n-m+1)!}\\
    &=\frac{(n+1)!}{m!(n-m+1)!}\\
\end{split}\]
所以$\comb{m}{n+1}=\comb{m}{n}+\comb{m-1}{n}$.
\end{proof}

为了使性质1在$n=m$时也能成立，我们规定$\comb{0}{n}=1$.

\begin{example}
计算$\comb{98}{100}$和$\comb{2}{10}+\comb{3}{10}+\comb{2}{11}+\comb{2}{12}$
\end{example}

\begin{solution}
根据组合数性质1
\[\comb{98}{100}=\comb{100-98}{100}=\comb{2}{100}=\frac{100\x 99}{2\x 1}=4950\]
根据组合数性质2
\[\begin{split}
    \comb{2}{10}+\comb{3}{10}+\comb{2}{11}+\comb{2}{12}&=\comb{3}{11}+\comb{2}{11}+\comb{2}{12}\\
    &=\comb{3}{12}+\comb{2}{12}=\comb{3}{13}\\
    &=\frac{13\x 12\x 11}{3\x 2\x 1}=286
\end{split}\]
\end{solution}

\begin{example}
    解不等式 $\comb{3}{n}>\comb{5}{n}$
\end{example}

\begin{solution}
\[\frac{n(n-1)(n-2)}{3\x 2\x 1}>\frac{n(n-1)(n-2)(n-3)(n-4)}{5\x 4\x 3\x 2\x 1}\]
因为$5\le n$，所以$n,\; n-1,\; n-2,\; n-3,\; n-4$均为正数，所以
\[1>\frac{(n-3)(n-4)}{5\x 4}\]
整理得：$n^2-7n-8<0$

解不等式得：$-1<n<8$，根据题意，得：$n=5,6,7$.
\end{solution}



\begin{example}
高中三个年级，每个年级有六个班。
\begin{enumerate}[(1)]
\item 每级都进行单循环（每班都与级内其他各班比赛一场）篮球预赛，问预赛需要比赛多少场？
\item 从各级预赛中选出每级的冠、亚军，再进行单循环决赛，规定在预赛中已经比赛过的队不再比赛，问决赛需要比赛多少场？
\end{enumerate}
\end{example}

\begin{solution}
\begin{enumerate}[(1)]
    \item 根据题意每个级单循环比赛的场次都是$\comb{2}{6}$，所以
\[3\comb{2}{6}=3\times 15=45\]
\item 决赛共6个班，如果单循环要比赛$\comb{2}{6}$场，但预赛中比赛过的有3场，所以
\[\comb{2}{6}-3=15-3=12\]
\end{enumerate}
答：预赛要比赛45场，决赛要比赛12场.
\end{solution}



\begin{example}
    平面有10个不同的点，其中除4个点在同一直线上外，不再有3个点共线，以每三个点为顶点可以作多少个不同的三角形？
\end{example}


\begin{solution}
    如果10点中的任三点都能连成三角形，那么有$\comb{3}{10}$个三角形. 但是如果在四个共线点中任选三点就不可能构成三角形. 所以
\[\comb{3}{10}-\comb{3}{4}=120-4=116\]
答：可以作116个不同的三角形。
\end{solution}



\begin{example}
    在产品检验时，常从产品中抽出一部分进行检查，假设100件产品中有2件次品，
    \begin{enumerate}[(1)]
        \item 从100件产品中抽出三件检查，其中有1件次品的抽法有多少种？
        \item 从100件产品中抽出4件检查，其中至少有一件次品的抽法有多少种？
    \end{enumerate}
\end{example}

\begin{solution}
\begin{enumerate}[(1)]
    \item 在100件产品中，次品有两件，从2件次品中抽出1件的抽法有$\comb{1}{2}$，从98件正品中抽出2件的方法有$\comb{2}{98}$. 所以
\[\comb{1}{2}\cdot \comb{2}{98}=2\cdot \frac{98\cdot 97}{2}=9506\]
    \item 至少有1件是次品就是包括有1件次品和2件次品两类抽法的和，所以 
\[\comb{1}{2}\cdot \comb{3}{98}+\comb{2}{2}\cdot \comb{2}{98}=2\times \frac{98\cdot 97\cdot 96}{3\cdot 2}+\frac{98\cdot 97}{2}=304192+4753=308945\]
\end{enumerate}
答：3件中有1件是次品的抽法是9506, 4件中至少有1件次品的抽法是308945。
\end{solution}

\begin{example}
    从0, 3, 5, 7, 11五个数中可以得到多少个不同的乘积？
\end{example}

\begin{solution}
    因为零乘以任何数其积都是零，所以不论零与多少个数相乘都只算一个乘积就是0；

在非零的乘积中有$\comb{2}{4}$, $\comb{3}{4}$, $\comb{4}{4}$三类，所以
\[\comb{2}{4}+\comb{3}{4}+\comb{4}{4}+1=6+4+1+1=12\]

答：可以得到12个不同的乘积.
\end{solution}

\begin{ex}
\begin{enumerate}
    \item 写出：
    \begin{enumerate}[(1)]
    \item 从4个相异元素$a$，$b$，$c$，$d$中取出2个元素的所有组合；
    \item 从5个相异元素$a$，$b$，$c$，$d$，$e$中取出3个元素的所有组合.
    \end{enumerate}

    \item 计算：
\begin{multicols}{2}
\begin{enumerate}[(1)]
    \item  $\comb{2}{7}$
    \item $\comb{48}{50}$
    \item $\comb{3}{100}-\comb{3}{90}$
    \item $\comb{1}{5}+\comb{2}{5}+\comb{3}{5}+\comb{4}{5}+\comb{5}{5}$
\end{enumerate}
\end{multicols}
    
\item   平面内有12个点，其中任意3点不在同一直线上，以每3点为顶点作一三角形，一共可以作多少个不同的三角形？
\item 圆上$n$个不同的点，
\begin{enumerate}[(1)]
\item 可以引多少条不同的弦？
\item 可以连接多少个内接三角形？
\end{enumerate}

\item 凸四边形可以连多少条对角线？凸五边形可以连多少条对角线？凸$n$边形可以连多少条对角线？
\item 某班有50名学生，其中正副班长各1人，现选5名学生参加某项活动
\begin{enumerate}[(1)]
\item 如果班长和副班长必须在内，有多少种不同的选法？
\item 如果要有一个班长（正、副都可以）在内，有多少种不同的选法？
\item 如果班长都不参加，有多少种不同的选法？
\item 如果至少要有一个班长（正、副都可以）在内，有多少种不同的选法？
\end{enumerate}
\end{enumerate}
\end{ex}

\section*{习题2.1}
\begin{enumerate}
    \item 计算：
\begin{multicols}{2}
\begin{enumerate}[(1)]
    \item $5\cdot \perm{3}{5}+4\cdot \perm{2}{4}$
    \item $\frac{\perm{5}{7}-\perm{6}{6}}{7!+6!}$
\end{enumerate}
\end{multicols}

\item 求证：
\begin{enumerate}[(1)]
    \item $\frac{(n+1)!}{k!}-\frac{n!}{(k-1)!}=\frac{(n-k+1)\cdot n!}{k!}$
    \item $\frac{(2n)!}{2^n \cdot n!}=1\cdot 3\cdot 5\cdots (2n-1)$
\end{enumerate}

\item 解方程：
\begin{multicols}{2}
\begin{enumerate}[(1)]
    \item $\perm{3}{2n}=10\perm{3}{n}$
    \item $\frac{\perm{5}{n}+\perm{4}{n}}{\perm{3}{n}}=4$
    \item $\perm{4}{2n+1}=140\perm{3}{n}$
\end{enumerate}
\end{multicols}

\item 解不等式：
\begin{multicols}{2}
\begin{enumerate}[(1)]
    \item $\perm{2}{m-2}+m>2$
    \item $\perm{x}{9}>6\perm{x-2}{9}$
\end{enumerate}
\end{multicols}

\item \begin{enumerate}[(1)]
    \item 从多少个不同的元素中取出2个元素的排列数是56？
    \item 已知从$n$个不同的元素中取出2个元素的排列数等于从$n-4$个不同元素中取2个元素的排列数的7倍，求$n$．
\end{enumerate}
\item 有5本不同的书分借给3个同学，每人1本，共有多少种不同的借法？
\item 用0, 1, 2, 3, 4, 5可以组成多少个没有重复数字的五位数？其中有多少个偶数？
\item 用数字1, 2, 3, 4, 5可以组成多少个没有重复数字的自然数？其中有多少个比13000大的自然数？
\item 7部不同型号的机床分配给7个工人使用：
\begin{enumerate}[(1)]
\item 甲必须使用2号机床；
\item 甲必须使用2号机床，乙必须使用4号机床；
\item 甲不适宜使用1号机床，乙不适宜使用7号机床；问三种情况下各有多少种不同的分配方案？
\end{enumerate}

\item 6个儿童站成一排，其中一人不站在排头，也不站在排尾，一共有多少种不同的排法？
\item 五部不同型号的半导体收音机和四部不同型号的电视机陈列成一排，任两部电视机不靠在一起，有多少种不同的陈列方法？
\item 在3000和7000间有多少个没有重复数字的5的倍数？

\item 计算：
\begin{multicols}{2}
\begin{enumerate}[(1)]
    \item $\comb{3}{15}$
    \item $\comb{98}{100}$
    \item $\comb{3}{6}+\comb{6}{8}$
    \item $\comb{n}{n+1}\cdot \comb{n-2}{n}$
\end{enumerate}
\end{multicols}

\item 求证：
\begin{enumerate}[(1)]
    \item $\comb{m}{n+1}=\comb{m-1}{n}+\comb{m}{n-1}+\comb{m-1}{n-1}$
    \item $\comb{k+1}{n}\div \comb{k}{n}=\frac{n-k}{k+1}$
    \item $\comb{0}{m}+\comb{1}{m+1}+\comb{2}{m+2}+\cdots+\comb{n-1}{m+n-1}=\comb{n-1}{m+n}$
    \item $\comb{m+1}{n}+\comb{m-1}{n}+2\comb{m}{n}=\comb{m+1}{n+2}$
\end{enumerate}

\item 解不等式：
\begin{multicols}{2}
\begin{enumerate}[(1)]
    \item $\comb{4}{n}>\comb{6}{n}$
    \item $\comb{n-4}{21}<\comb{n-2}{21}<\comb{n-1}{21}$
\end{enumerate}
\end{multicols}

\item 在200件产品中有2件次品，现在抽5件进行检查：
\begin{enumerate}[(1)]
\item 其中有一件次品的不同抽法有多少种？
\item 没有次品的不同抽法有多少种？
\item 至少有一件次品的不同抽法有多少种？
\end{enumerate}

\item 问2310能被多少个偶数整除？
\item 有8个划船运动员，其中3人能划左舷，2人能划右舷，3人能划左右舷，要选出6人分配在两边，问有多少种不同的选法？

\item  某车间有9名技术熟练的工人，其中2人能当车、钳工，有3人只能当车工，其余4人只能当钳工，现在要抽调三个车工，三个钳工成立一个新车间，问有多少种不同的抽调方法？
\item  空间有12个不同的点，其中有4点在同一平面内，此外不再有4点共面，
\begin{enumerate}[(1)]
\item  这些点可以确定多少个不同的平面？
\item  这些点可以连成多少个不同的四面体？
\end{enumerate}

\item  有8部机器分给陈、李、张三位师傅管理。
\begin{enumerate}[(1)]
\item 陈管4部，李管3部，张管1部，有几种不同分配方法？
\item 陈管4部，其他机器由1人管3部，1人管1部，有多少种不同的分配方法？
\item 1人管4部，一人管3部，一人管1部，有几种不同的分配方法？
\end{enumerate}

\item  从1, 5, 7, 9中任取3个数字，2, 4, 6, 8中任取两个数字，组成没有重复数字的五位数，一共有多少个不同的五位数？
\end{enumerate}

\section{环形排列、重复排列和组合}
\subsection{环形排列}
我们已经掌握了从$n$个相异的元素中取出$m\; (m\le n)$个元素的排列的知识了，那么怎样解决环形排列的问题呢？

我们先从最简单的情况入手，就是先找出环形排列与直
线排列的相互关系，然后利用直线排列的知识解决环形排列的问题.

首先在环形排列中，下面的排列都是当作相同的排列的（图2.2）．
\begin{figure}[htp]
    \centering
\begin{tikzpicture}[>=stealth]
\begin{scope}
\draw circle(.9);
\node at (90:1.2){甲};
\node at (90+180:1.2){乙};
\draw[->](90-20:1.2) arc (90-20:20:1.2);
\end{scope}
\begin{scope}[xshift=3cm]
\draw circle(.9);
\node at (0:1.2){甲};
\node at (0+180:1.2){乙};
\draw[->](0-20:1.2) arc (0-20:-70:1.2);
\end{scope}
\begin{scope}[xshift=6cm]
\draw circle(.9);
\node at (-90:1.2){甲};
\node at (-90+180:1.2){乙};
\draw[->](-90-20:1.2) arc (-90-20:-160:1.2);
\end{scope}
\begin{scope}[xshift=9cm]
\draw circle(.9);
\node at (45:1.2){甲};
\node at (45+180:1.2){乙};
\draw[->](45-20:1.2) arc (45-20:-20:1.2);
\end{scope}

\end{tikzpicture}
    \caption{}
\end{figure}


因为从甲开始，按同一方向，例如顺时针方向看过去都是乙，所以尽管位置不同，但仍然是同一样的元素的排列顺序.
\begin{figure}[htp]
    \centering
\begin{tikzpicture}[>=stealth]
\begin{scope}
\draw circle(.9);
\node at (90:1.2){甲};
\node at (90-120:1.2){乙};
\node at (90+120:1.2){丙};
\draw[->](90-20:1.2) arc (90-20:20:1.2);
\end{scope}
\begin{scope}[xshift=3.5cm]
\draw circle(.9);
\node at (0:1.2){甲};
\node at (0-120:1.2){乙};
\node at (0+120:1.2){丙};
\draw[->](0-20:1.2) arc (0-20:-70:1.2);
\end{scope}
\begin{scope}[xshift=7cm]
\draw circle(.9);
\node at (-90:1.2){甲};
\node at (-90-120:1.2){乙};
\node at (-90+120:1.2){丙};
\draw[->](-90-20:1.2) arc (-90-20:-160:1.2);
\end{scope}
  
\end{tikzpicture}
    \caption{}
\end{figure}


又例如三个不同元素甲、乙、丙如图2.3的三个环形排列也都算是相同的排列，而且与从谁开始无关，即甲$\to $乙$\to $丙$\to $甲、乙$\to $丙$\to $甲$\to $乙与丙$\to $甲$\to $乙$\to $丙，也都算作是相同的环形排列.

从上述的例子，可以得到三个不同元素的环形排列数应该是2．具体的排列是：甲—乙—丙—甲与甲—丙—乙—甲．

现在，我们再来分析一下环形排列与直线排列不同的特点. 我们已经知道下列三个环形排列都是相同的排列（图2.4）
\begin{figure}[htp]
    \centering
\begin{tikzpicture}[>=stealth]
\begin{scope}
\draw circle(.9);
\node at (90:1.2){甲};
\node at (90-120:1.2){乙};
\node at (90+120:1.2){丙};
\draw[->](-.35,1.5)--(-.35,1);
\end{scope}
\begin{scope}[xshift=3.5cm]
    \draw circle(.9);
    \node at (90:1.2){丙};
    \node at (90-120:1.2){甲};
    \node at (90+120:1.2){乙};
    \draw[->](-.35,1.5)--(-.35,1);

    \end{scope}
\begin{scope}[xshift=7cm]
    \draw circle(.9);
    \node at (90:1.2){乙};
    \node at (90-120:1.2){丙};
    \node at (90+120:1.2){甲};
    \draw[->](-.35,1.5)--(-.35,1);

    \end{scope}
\end{tikzpicture}
    \caption{}
\end{figure}

如果我们从图上所示的地方把环切断，并把环拉成直线后，我们却可以得到三个不同的直线的排列，就是：
\[\text{甲、乙、丙；丙、甲、乙；乙、丙、甲}\]

而三个不同元素的全排列数是$\perm{3}{3}$，所以三个不同元素的环形排列数是$\frac{\perm{3}{3}}{3}$

\begin{blk}
    {定理3} $m$个不同元素的环形排列数就是：
\[\frac{\perm{m}{m}}{m}=\frac{m!}{m}=(m-1)!\]
\end{blk}

\begin{proof}
设有$m$个不同元素排成一个环形，如图（图2.5）
\begin{figure}[htp]
    \centering
\begin{tikzpicture}
\draw(0,0) circle(1);
\draw[fill](15:1)node[right]{$a_4$} circle(1pt);
\draw[fill](15+30:1)node[above right]{$a_3$} circle(1pt);
\draw[fill](15+60:1)node[above]{$a_2$} circle(1pt);
\draw[fill](15+90:1)node[above]{$a_1$} circle(1pt);
\draw[fill](15+120:1)node[above left]{$a_m$} circle(1pt);
\draw[fill](15+150:1)node[left]{$a_{m-1}$} circle(1pt);
\end{tikzpicture}
    \caption{}
\end{figure}
就是按顺时针方向的一个排列．其次，$a_1$不动，改变$a_2,a_3,\ldots,a_m$的顺时针排列顺序，就得到第二个不同的环形排列．因此，如果$a_1$不动，要使$a_2,a_3,\ldots,a_m$有不同的排
列，其不同的排列数为$(m-1)!$，故$m$个不同元素的环形排列数为$\perm{m-1}{m-1}=(m-1)!$
\end{proof}


\begin{blk}
  {定理4} 从$n$个不同的元素中，任取$m$个（$m\le n$）没有重复的元素的环形排列数$\comb{m}{n}\cdot (m-1)!$.  
\end{blk}

\begin{proof}
    从$n$个不同元素中取出$m\; (m\le n)$个元素进行环形排列，我们第一步可以先从$n$个元素中先取出$m$个组成一组，第二步是把$m$个元素排成一圈，根据乘法原理，其排列数应为：
\[\comb{m}{n}\cdot \perm{m-1}{m-1}=\comb{m}{n}(m-1)!\]
\end{proof}

\begin{example}
    从10个不同元素中取出6个元素排成一圈，有多少种不同的排法？
\end{example}

\begin{solution}
    从10个不同元素中取出6个的环形排列数为：
\[\comb{6}{10}\cdot\perm{6-1}{6-1}=\comb{4}{10}\cdot  5!=210\cdot 120=25200\]
答：有25200种不同的排法．
\end{solution}

\begin{example}
    有五个男同学与五个女同学围成圆圈跳集体舞，如果男女相间，问有几种不同的排列方法？
\end{example}

\begin{solution}
第一步，五个男同学先作环形排列，其排列数为$4!$；

第二步，五个女同学分别站在两个男同学之间的位置上作全排列.

$\therefore\quad 4! \cdot 5! =2880$

答：有2880种不同的排列方法.    
\end{solution}

\begin{ex}
\begin{enumerate}
\item 把10个人平均分成两组，每组都围坐一个圆桌，问有多少种不同的坐法？
\item 今有5颗不同颜色的珠子，用线把它们串成珠圈，问有多少种串法？
\end{enumerate}
\end{ex}

\subsection{重复排列}

前面我们所研究的都是不允许元素在排列和组合中重复出现的，假如我们允许元素在同一排列中重复出现的话，其排列数应该是多少呢？

\begin{blk}
{定理5} 从$n$个不同的元素中，取出$m$个元素排成一列（允许元素重复出现）的不同排列数为$n^m$.
\end{blk}

\begin{proof}
    选在第一个位置上的元素有$n$个不同的方法，选在第二个位置上的元素仍然有$n$个不同的方法，……，选在第$m$个位置上的元素仍然有$n$个选法，所以根据乘法原理，从$n$个不同的元素中取出$m$个元素的允许重复的排列数应为$n^m$.
\end{proof}

这种排列又称为重复排列，这里应该注意的是，由于元素可以重复选取，因此选取的次数$m$就不会再受$n$个元素的限制. 即$m$不再有$m\le n$的限制了.

\begin{example}
    电报通讯的明码是用0到9十个数字中的每4个数字的重复排列代表一个汉字，问用这种方法制成的电报通讯明码可以表示多少个不同的汉字？
\end{example}

\begin{solution}
    从10个不同数字中选取4个数字的重复排列数是：
\[10^4=10000\]
答：可以表示10000个不同的汉字.
\end{solution}

\begin{example}
    由1, 2, 3, 4四个数字能够组成多少个大于2300的四位数？
\end{example}

\begin{solution}
    这里的四位数显然是包括重复数字的四位数.

第一类：3在千位上的四位数一定大于2300，那么，百、十、个位上的数字每次都有4种选取可能，所以应有$4^3$个四位数.

第二类：4在千位上的四位数，同理也有$4^3$个四位数。

第三类：2在千位上，3在百位上的四位数，由于十位和个位上的数都可以从1, 2, 3, 4中选取，都大于2300，所以，这样的数有$4^2$个．

第四类，2在千位上，4在百位上的四位数有$4^2$个．所以
\[2\cdot 4^3+2\cdot 4^2=2\cdot 4^2(4+1)=160\]

答：这样的四位数有160个．
\end{solution}

\begin{ex}
    \begin{enumerate}
\item 有三封不同的信，投入4个信箱里，一共有多少种投信的方法？
\item 有数学、天文、诗歌三个课外小组，8个同学准备报名，若每人限报一个组，问一共可有多少种报名的方法？
    \end{enumerate}
\end{ex}

\subsection{重复组合}
从$n$个不同元素中取出$m$个元素组成一组，如果允许元素在组合中重复出现，这样的组合简称为重复组合，这种重复组合用符号$\rep{m}{n}$表示.

怎样计算重复组合的组合数呢？我们还是先从具体例子开始考虑，假如有三个不同的元素$a_1,a_2,a_3$，从其中取出两个元素组成重复组合，那么，可以得到如下的六个不同组合：
\begin{equation}
a_1a_1,\; a_1a_2,\; a_1a_3,\; a_2a_2,\; a_2a_3,\; a_3a_3 \tag{1}
\end{equation}

如果是不允许元素重复的话，要选两个元素得到不同的组合数是6的话，应是$\comb{2}{4}$，设这四个元素为$a'_1,a'_2,a'_3,a'_4$,那么得到的组合为：
\begin{equation}
a'_1a'_2,\; a'_1a'_3,\; a'_1a'_4,\; a'_2a'_3,\; a'_2a'_4,\; a'_3a'_4 \tag{2}
\end{equation}
组合（1）与组合（2）不难建立它们之间的一一对应关系．在组合（1）的每个组合的元素下标上分别加上0和1，就得到了下标与（2）完全一致的组合（3）．
\begin{equation}
a_{1+0}a_{1+1},\quad a_{1+0}a_{2+1},\quad a_{1+0}a_{3+1},\quad a_{2+0}a_{2+1},\quad a_{2+0}a_{3+1},\quad a_{3+0}a_{3+1}
\tag{3}
\end{equation}

又例如从4个不同元素中取出3个元素的重复组合有如下的不同组合:
\begin{center}
\begin{tabular}{ccccc}
    $a_{1}a_{1}a_{1}$&$a_{1}a_{1}a_{2}$&$a_{1}a_1a_{3}$&$a_{1}a_{1}a_{4}$\\
    &$a_{1}a_{2}a_{2}$&$a_{1}a_{2}a_{3}$&$a_{1}a_{2}a_{4}$\\
    &$a_{1}a_{3}a_{3}$&$a_{1}a_{3}a_{4}$&$a_{1}a_{4}a_{4}$\\
    &$a_{2}a_{2}a_{2}$&$a_{2}a_{2}a_{3}$&$a_{2}a_{2}a_{4}$\\
    &$a_{2}a_{3}a_{3}$&$a_{2}a_{3}a_{4}$&$a_{2}a_{4}a_{4}$\\
 &$a_{3}a_{3}a_{3}$&$a_{3}a_{3}a_{4}$&$a_{3}a_{4}a_{4}$&$a_{4}a_{4}a_{4}$\\
\end{tabular}
\end{center}

只要我们在每组元素的下标上分别加0, 1, 2, 就得到
\begin{center}
\begin{tabular}{ccccc}
    $a_{1}a_{2}a_{3}$&$a_{1}a_{2}a_{4}$&$a_{1}a_{2}a_{5}$&$a_{1}a_{2}a_{6}$\\
 & $a_{1}a_{3}a_{4}$&$a_{1}a_{3}a_{5}$&$a_{1}a_{3}a_{6}$\\
 &$a_{1}a_{4}a_{5}$&$a_{1}a_{4}a_{6}$&$a_{1}a_{5}a_{6}$\\
 &$a_{2}a_{3}a_{4}$&$a_{2}a_{3}a_{5}$&$a_{2}a_{3}a_{6}$\\
 &$a_{2}a_{4}a_{5}$&$a_{2}a_{4}a_{6}$&$a_{2}a_{5}a_{6}$\\
 &$a_{3}a_{4}a_{5}$&$a_{3}a_{4}a_{6}$&$a_{3}a_{5}a_{6}$&$a_{4}a_{5}a_{6}$
\end{tabular}
\end{center}

这又正是从六个不同元素中取出3个元素的组合的全部。
从上面的例子说明了
\[{\rm H}_{3}^{2}= {\rm C}_{3+ ( 2- 1) }^{2}, \qquad {\rm H}_{4}^{3}={\rm C}_{4+ ( 3- 1) }^3\]

按同样的方法我们可以得出：从$n$个不同元素中取出$m$
个的重复组合数有如下的关系：
$${\rm H}_{n}^{m}= {\rm C}^m_{n+ ( m- 1) }$$

现在我们来证明这个结论.

\begin{blk}
 {定理6} 从$n$个不同元素中，任取$m$个元素 的重复组合为：
 $${\rm H}_{n}^{m}={\rm C}_{n+ ( m-1) }^{m}$$ 
\end{blk}

\begin{proof}
    为了研究问题方便，我们用$1,2,3,\ldots,n$表示这$n$个不同的元素。在其中任取$m$个允许重复的数字作为一个组合，并按大小顺序排成
$$1\le i_1\le i_2\le\cdots\le i_m\le n$$

    又记$j_1=i_1+0,\; j_{2}=i_{2}+1,\;\ldots, j_{m}=i_{m}+(m-1)$
    
    自然，$1\le j_1<j_2<j_3<\cdots<j_m\leq n+(m-1)$，所以$\{j_1,j_2,\ldots,j_m\}$ 是$n+ m- 1$个 数，并 且 是$1, 2, \ldots , n+m-1$中的$m$个不同数的组合.
\end{proof}

下面来证明：在$n$个不同元素中任取$m$个允许重复的元
素的组合，和$n+m-1$个不同元素中任取$m$个不同元素 的组合之间有一个到上的一一对应关系。如果这点证明了，那么$n$个不同元素的$m$个允许重复元素的组合数就与$n+m-1$个不同元素的$m$个不同元素的组合数相等了.

\begin{proof}
    
为了证明这个一一对应，要先证明当允许重复的组合
$\{i_1,i_2,\ldots,i_m\}$与$\{i_1^{\prime},i_2^{\prime},\ldots,i'_m\}$不同 的时候，
（这里取$1\le i_1\le i_2\le\cdots\le i_m\le n$, $1\le i'_1\le i'_2\le\cdots\le i^{\prime}_m\le n$）那么，不同元素不允许重复的组合$\{j_1,j_2,\ldots, j_m\}$与$\{j_1^{\prime},j_2^{\prime},\ldots,j'_m\}$也不同，这里$\{j_1=i_1+0,\; j_2=i_{2}+1,\ldots,j_{m}=i_{m}+m-1,\; j_{1}^{\prime}=i_{1}^{\prime}+0,\; j_{2}^{\prime}=i_{2}^{\prime}+1,\ldots, j_m^{\prime}=i_m^{\prime}+m-1\}$.

事实上，如果后者相同，即有$j_{k}= j'_{k}\; (k= 1, 2, \ldots, m)$ 即
$$i_k+k-1=i_k^{\prime}+k-1\; (k=1,2,\ldots,m)$$ 
所以$i_k= i_k^{\prime }\; ( k= 1, 2, \ldots , m)$, 也就是说$\{i_1, i_2,\ldots,i_m\}$也与$\{i_1^{\prime},i_2^{\prime},\ldots,i_m^{\prime}\}$相同了。这就证明了上述的对应为一一对应.

再证明这是到上的一一对应，即当$\{i_1,i_2,\ldots,i_m\}$取遍$n$个数$1,2,\ldots,n$的一切$m$个数的允许重复组合时，$\{j_1, j_2,\ldots,j_m\}$ 也取遍$n+m-1$个 数 $1,2,\ldots ,n+m-1$一切的 $m$ 个数的不允许重复的组合.

事实上，对任一个组合$\{j_1,j_2,\ldots,j_m\}$, 其中$1\le j_{1}<j_{2}<\cdots<j_{m}\le n+m-1$. 有$1\leq j_{1},\; 2\le j_{2},\; 3\le j_{3},\ldots, m\le j_{m}$. 故$1\le i_1=j_1-0,\; 1\le i_2=j_2-1,\; 1\le i_3=j_3-2,\ldots, 1\leq i_m=j_m-(m-1)$. 即$i_1,i_2,\ldots,i_m$都是自然数，且由$j_{k}>j_{k-1}$, 可知$j_k-j_{k-1}\ge 1$, 故有
\[i_k-i_{k-1}=[j_k-(k-1)]-[j_{k-1}-(k-2)]=j_k-j_{k-1}-1\ge 0\]
再者$i_{m}=j_{m}-(m-1)\le (n+m-1)-(m-1)=n$. 总之有$1\le i_1\le i_2\le \cdots \le i_m\le n$. 所以，$\{i_1,i_2,\ldots,i_m\}$ 是$n$个数$1，2,\ldots,n$中取$m$个数的允许重复的组合，这证明了对应是到上的一一对应.

\end{proof}


\begin{example}
  从3, 5, 7, 9, 11中取出3个数（可以重复选取）相乘，问有多少个不同的乘积？  
\end{example}

\begin{solution}
    因为是乘积，与元素取出的先后顺序无关，同时是可以重复选取的，即是重复组合问题. 所以
\[{\rm H}^3_5=\comb{3}{5+(3-1)}=\comb{3}{7}=35\]
答：有35个不同的乘积．
\end{solution}

\begin{example}
    从三个不同的质数$a_1,a_2,a_3$中取出4个数相乘（可以重复选取），问有多少个不同的乘积？
\end{example}

\begin{solution}
    同例2.16是重复组合问题．所以
    \[{\rm H}^4_3=\comb{4}{3+(4-1)}=\comb{4}{6}=\comb{2}{6}=15\]
答：有15个不同的乘积。
\end{solution}

\begin{example}
    用七种不同的颜色给一个四面体的四个面涂上颜色，规定每个面上只能涂一种颜色，如果四个面的颜色容许相同，问有几种不同的着色方法？
\end{example}

\begin{solution}
    因为只研究对四个面所着的颜色，对那个面先着色并无先后顺序之分，故是组合问题.又因为对四个面所着的色是允许相同的，所以又是重复组合问题. 所以
    \[{\rm H}^4_7=\comb{4}{7+(4-1)}=\comb{4}{10}=210\]
    答：共有210种不同的着色方法。
\end{solution}


\begin{ex}
\begin{enumerate}
\item 展开$(a+b+c+d)^3$，合并同类项后，一共有多少项？
\item 若把六种不同的溶液注满四个烧杯，但不使它们混合，问一共有多少种方法？
\end{enumerate}
\end{ex}

\section*{习题2.2}
\begin{enumerate}
\item 从8个不同的元素中取出5个元素排成圆圈，问有几种不同的排列方法？
\item 两个教师和五个学生代表围坐在圆桌周围开座谈会：
\begin{enumerate}[(1)]
\item 如果不加条件限制，有多少种不同的坐法？
\item 如果两个教师靠在一起，有几种坐法？
\item 如果两个教师不靠在一起，有几种坐法？
\item 学生代表甲坐在教师的中间，有几种坐法？
\end{enumerate}
\item 从三个不同的元素中取出4个元素的重复组合有多少种不同的方法？
\item 用七种颜色涂在正方体的六个面上，如果每个面只能涂一种颜色，如果六个面上的颜色允许相同，问共有几种不同的着色方法？
\item 从2, 3, 5, 7, 11, 13中至少取出两个，至多可以取出五个数相乘（允许因数重复选取），问有多少个不同的乘积？
\item 展开$(a+b+c+d+e+f)^4$，再合并同类项，一共可以得到多少项？
\end{enumerate}

\section{二项式定理}

\subsection{和号$\Sigma$}

我们经常要计算下面各种和
$$1+ 2+ 3+ \cdots + n;\quad 1^{2}+ 2^{2}+ 3^{2}+ \cdots + n^{2},\quad \frac{1\times2}{2}+\frac{2\times3}{2^{2}}+\cdots+\frac{n(n+1)}{2^{n}}$$
等等。 一 般 给 出 一 个 有 顺 序 的 数 列 $a_1, a_2, \ldots , a_n$, 要计算
$a_{1}+a_{2}+\cdots+a_{n}$，
为了方便起见引进和符号$\Sigma$，于是可记
$$1+2+3+\cdots+n=\sum_{k=1}^{n}k$$
\[1^{2}+ 2^{2}+ 3^{2}+ \cdots + n^{2}=\sum^n_{k=1}k^2\]
\[\frac{1\times2}{2}+\frac{2\times3}{2^{2}}+\cdots+\frac{n(n+1)}{2^{n}}=\sum^n_{k=1}\frac{k(k+1)}{2^k}\]

一般可记
$a_1+a_2+\cdots+a_n=\sum_{k=1}^{n}a_k$
符号
$$\sum_{k=1}^{n}a_k$$

表示$k$取$1,2,\ldots,n-1,n$时，得到的$a_k$全部加起来.
例如：
$$\sum_{k=1}^{n}\sin kx=\sin x+\sin2x+\sin3x+\cdots+\sin nx$$
$$\sum _{k= 0}^{n}a_{k} x^{n- k}= a_{0} x^{n}+ a_{1} x^{n- 1}+ \cdots + a_{n- 1} x+ a_{n}$$
又例如，对给定的函数$f_1(x),f_2(x),\ldots,f_n(x)$
有
\[\begin{split}
\sum^n_{k=1}f_k(x)&=f_1(x)+f_2(x)+\cdots +f_n(x)
\\
\sum^n_{k=0}\comb{k}{n} &=\comb{0}{n}+\comb{1}{n}+\comb{2}{n}+\cdots +\comb{n}{n}
\end{split}\]
等等.

关于求和符号还有两点说明，首先，凡是指示指标从1到$n$，且仅仅指示从1到$n$．所以可以用其它的字母来代替$k$而和号仍然是表示了同一个意思. 例如
\[\sum^n_{k=1}a_k=a_1+a_2+\cdots+a_n,\qquad \sum^n_{\ell=1}a_{\ell}=a_1+a_2+\cdots+a_n\]

一般$\sum\limits^n_{*=1}a_*=a_1+a_2+\cdots+a_n$
这里$*$可以用任何一种指标来表示，其次，求和也不限于从1到$n$，也可以从0到$n+2$等等. 例如
\[\begin{split}
    \sum^{n+2}_{k=0}a_k&=a_0+a_1+\cdots+a_{n+1}+a_{n+2}\\
    \sum^{n+1}_{k=3}a_k&=a_3+a_4+\cdots+a_{n+1}
\end{split}\]

\begin{ex}
\begin{enumerate}
    \item 把下列各式用和号表示：
\begin{enumerate}[(1)]
\item $1+3+5+7\cdots+(2n-1)$ 
\item $1\cdot2+2\cdot3+3\cdot4+\cdots+n(n+1)$ 
\item $\cos x+\cos2x+\cos3x+\cdots+\cos nx$
\item $1^{2}+2^{2}+3^{2}+\cdots+(n+2)^{2}$
\end{enumerate}

\item 把下列和号式改写为一般的和式
\begin{multicols}{2}
\begin{enumerate}[(1)]
\item $\sum\limits _{k= 1}^{n}\frac 1{k( 2k+ 1) }$ 
\item $\sum\limits _{k= 1}^{n+ 2}\frac {k^{2}}{k! }$
\item $\sum\limits^n_{k=2}\frac{\sin kx}{k}$
\item $\sum\limits^n_{k=2}\frac{(-1)^k}{k^2+k-2}$
\end{enumerate}
\end{multicols}
\end{enumerate}
\end{ex}

\subsection{二项式定理}
下面我们用数学归纳法来证明二项式定理。我们已经知道：
\[\begin{split}
(x+y)^1&=x+y\\
(x+y)^2&=x^2+2xy+y^2\\
(x+y)^3&=x^3+3x^2y+3xy^2+y^3    
\end{split}\]

现在考虑二项式$(x+y)^n$的展开式，这里$n$是任意自然数，首先研究$(x+y)^4$
\[(x+y)^4=(x+y)(x+y)(x+y)(x+y)\]
等号右边的积应有下面形式的各项：
\[x^4,\; x^3y,\; x^2y^2,\; xy^3,\; y^4\]
显然$x^4$, $y^4$的系数是1，其他三项的系数，例如$x^2y^2$的系数是从四个括弧中两个里取$y$，从其他两个里取$x$作出的积的个
数，从四个括弧中任取两个$y$的取法 有$\comb{2}{4}=\frac{4\cdot 3}{1\cdot 2}=6$种，因
此$x^2y^2$的系数是6，同样，可以求得$x^3y,xy^3$的系数分别
是$\comb{1}{4}=4$，因为规定$\comb{0}{n}=1$，又知$\comb{4}{4}=1$，所以
\[(x+y)^4=\comb{0}{4}x^4+\comb{1}{4}x^3y+\comb{2}{4}x^2y^2+\comb{3}{4}xy^3+\comb{4}{4}y^4\]

一般地，可以归纳出以下的定理：
\begin{blk}
    {定理} 对于两个不定元$x$，$y$的多项式$(x+y)^n$，有下列展开式：
\[\begin{split}
    (x+y)^n&=\sum^n_{k=0}\comb{k}{n}x^{n-k}y^k\\
    &=\comb{0}{n}x^n+\comb{1}{n}x^{n-1}y+\comb{2}{n}x^{n-2}y^2+\cdots+\comb{n}{n}y^n
\end{split}\]
\end{blk}

\begin{proof}
    用数学归纳法来证明定理.
\begin{enumerate}[(1)]
    \item 当$n=1$时，$(x+y)^1=x+y$
\[\sum^1_{k=0}\comb{k}{1}x^{1-k}y^k=\comb{0}{1}x^1y^0+\comb{1}{1}x^{0}y^1=x+y\]
故$n=1$时原命题成立.
\item 假设$n=\ell$时有公式 $(x+y)^{\ell}=\sum\limits^\ell_{k=0}\comb{k}{\ell}x^{\ell-k}y^k$

现在来计算$(x+y)^{\ell+1}$，令
\[\begin{split}
    (x+y)^{\ell+1}&=(x+y)\cdot (x+y)^{\ell}=(x+y)\cdot \sum^\ell_{k=0}\comb{k}{\ell}x^{\ell-k}y^k\\
&=\sum^\ell_{k=0}\comb{k}{\ell}x^{\ell-k+1}y^k +\sum^\ell_{k=0}\comb{k}{\ell}x^{\ell-k}y^{k+1}\\
&=\comb{0}{\ell}x^{\ell+1}+\sum^\ell_{k=1}\comb{k}{\ell}x^{1+\ell-k}y^k+\sum^{\ell-1}_{k=0}\comb{k}{\ell}x^{\ell-k}y^{k+1}+\comb{\ell}{\ell}y^{\ell+1}\\
&=\comb{0}{\ell+1}x^{\ell+1}+\sum^\ell_{k=1}\comb{k}{\ell}x^{1+\ell-k}y^k+\sum^{\ell}_{k=1}\comb{k-1}{\ell}x^{\ell+1-k}y^{k}+\comb{\ell+1}{\ell+1}y^{\ell+1}\\
&=\comb{0}{\ell+1}x^{\ell+1}+\sum^\ell_{k=1}\left(\comb{k}{\ell}+\comb{k-1}{\ell}\right)x^{\ell+1-k}y^k+\comb{\ell+1}{\ell+1}y^{\ell+1}\\
&=\comb{0}{\ell+1}x^{\ell+1}+\sum^\ell_{k=1}\comb{k}{\ell+1}x^{\ell+1-k}y^k +\comb{\ell+1}{\ell+1}y^{\ell+1}\\
&=\sum^{\ell+1}_{k=0}\comb{k}{\ell+1}x^{\ell+1-k}y^k 
\end{split} \]
所以$n=\ell+1$时命题成立.
\end{enumerate}
根据（1）和（2），对于一切自然数$n$都有：$(x+y)^n=\sum\limits^n_{k=0}\comb{k}{n}x^{n-k}y^k$成立.
\end{proof}

这个定理叫做二项式定理，右边的多项式$\sum\limits^n_{k=0}\comb{k}{n}x^{n-k}y^k$
叫做$(x+y)^n$的二项展开式. 式中的$\comb{k}{n}x^{n-k}y^k$叫做二项展开式的通项，用${\rm T}_{k+1}$表示，即通项为展开式的第$k+1$项.
\[{\rm T}_{k+1}=\comb{k}{n}x^{n-k}y^k\]

由二项式定理，取$y=1$，有
\[(x+1)^n=\sum^n_{k=0}\comb{k}{n}x^{n-k}\]
取$y=-1$，有$$(x-1)^n=\sum\limits^n_{k=0}(-1)^k\comb{k}{n}x^{n-k}$$

在二项式定理中，如果遇到$n$是较小的正整数时展开式的系数，也可以直接用下表计算，
\begin{center}
\begin{tikzpicture}
\foreach \x in {0,1,2,...,6}
{
    \node at (0,-\x){$(x+y)^{\x}$};
}
\foreach \x/\y in {3/1, 4/6, 5/15, 6/20, 7/15, 8/6, 9/1}
{
    \node (A\x) at (\x,-6){\y};
}
\foreach \x/\y in {3/1, 4/5, 5/10, 6/10, 7/5, 8/1}
{
    \node (B\x) at (\x+.5,-5){\y};
}
\foreach \x/\y in {3/1, 4/4, 5/6, 6/4, 7/1}
{
    \node (C\x) at (\x+1,-4){\y};
}
\foreach \x/\y in {3/1, 4/3, 5/3, 6/1}
{
    \node (D\x) at (\x+1.5,-3){\y};
}
\foreach \x/\y in {3/1, 4/2, 5/1}
{
    \node (E\x) at (\x+2,-2){\y};
}
\foreach \x/\y in {3/1, 4/1}
{
    \node (F\x) at (\x+2.5,-1){\y};
}
\node (G) at (6,0){1};

\draw(F3)--(E4)--(F4);
\draw(E3)--(D4)--(E4)--(D5)--(E5);
\draw(D3)--(C4)--(D4)--(C5)--(D5)--(C6)--(D6);
\draw(C3)--(B4)--(C4)--(B5)--(C5)--(B6)--(C6)--(B7)--(C7);
\draw(B3)--(A4)--(B4)--(A5)--(B5)--(A6)--(B6)--(A7)--(B7)--(A8)--(B8);

\draw[dotted](1,0)--(G);
\draw[dotted](1,-1)--(F3);
\draw[dotted](1,-2)--(E3);
\draw[dotted](1,-3)--(D3);
\draw[dotted](1,-4)--(C3);
\draw[dotted](1,-5)--(B3);
\draw[dotted](1,-6)--(A3);


\end{tikzpicture}
\end{center}
表中每行两端都是1，而且除1以外的每一个数都等于它肩上两个数的和，这个表叫做杨辉三角，它首先载于我国宋朝数学家杨辉1261年所著的《详解九章算法》一书．

\begin{example}
展开$\left(1{+\frac{1}{x}}\right)^5$
\end{example}


\begin{solution}
\[\begin{split}
    \left(1{+\frac{1}{x}}\right)^5 &=1+5\left(\frac{1}{x}\right)+10\left(\frac{1}{x}\right)^2+10\left(\frac{1}{x}\right)^3+5\left(\frac{1}{x}\right)^4+\left(\frac{1}{x}\right)^5\\
    &=1+\frac{5}{x}+\frac{10}{x^2}+\frac{10}{x^3}+\frac{5}{x^4}+\frac{1}{x^5}
\end{split}\]
\end{solution}

\begin{example}
    展开$\left(2\sqrt{x}-\frac{1}{\sqrt{x}}\right)^6$
\end{example}

\begin{solution}
\[\begin{split}
    \left(2\sqrt{x}-\frac{1}{\sqrt{x}}\right)^6&=\left(\frac{2x-1}{\sqrt{x}}\right)^6 =\frac{1}{x^3}(2x-1)^6\\
&=\frac{1}{x^3}\left[(2x)^6-\comb{1}{6}(2x)^5+\comb{2}{6}(2x)^4-\comb{3}{6}(2x)^3+\comb{4}{6}(2x)^2\right.\\
&\qquad \qquad \left.-\comb{5}{6}(2x)+\comb{6}{6}\right]\\
&=\frac{1}{x^3}\left[64x^6-6\cdot 32x^5+15\cdot 16x^4-20\cdot 8 x^3+15\cdot 4x^2\right.\\
&\qquad \qquad \left.-6\cdot 2x+1\right]\\
&=64x^3-192x^2+240x-160+\frac{60}{x}-\frac{12}{x^2}+\frac{1}{x^3}
\end{split}\]
\end{solution}

\begin{example}
    求$(x+a)^{12}$的展开式中的第四项和倒数第四项。
\end{example}

\begin{solution}
    展开式的通项公式$\mathrm{T}_{k+1}=\mathrm{C}_n^k x^{n-k}y^k$。所以
$$\mathrm{T}_{4}=\mathrm{C}_{12}^{3} x^{12-3} a^{3}=\frac{12\times11\times 10}{3\times2}x^{9} a^{3}=220x^{9} a^{3}$$
又展开式共$12+1=13$, 故倒数第四项是
$$\mathrm{T}_{10}=\mathrm{C}_{12}^{9} x^{12-9} a^{9}=\mathrm{C}_{12}^{3}x^{3} a^{9}=220x^{3} a^{9}$$
\end{solution}

\begin{example}
  求$(x-\frac{1}{x})^{9}$展开式中$x^3$的系数，有没有$x^{0}$项，如
果有，试求出该项来。  
\end{example}

\begin{solution}
    展开式的通项是：
$$T_{k+1} = ( -1 )^{k} C_{9}^{k} x^{9-k} (\frac{1}{x})^{k}=( -1)^{k} C_{9}^{k} x^{9-k-k}$$
依题意有，$9-2k=3$, 所以$k=3$. 因此$x^3$的项系数是
$$( -1 )^{3}{\rm C}_{9}^{3}=- \frac{9\times8\times7}{3\times2}=- 84$$

又若$9-2k=0$, 即$2k=9$, 方程没有整数解，故展开式
里没有$x^0$项.
\end{solution}

\begin{example}
   在$\left(a+\frac1a\right)^{2n}$的展开式里，已知第四项和第六项
系数相等，求展开式里不含有$a$的项. 
\end{example}

\begin{solution}
因为
\[\begin{split}
    {\rm T}_4&=\comb{3}{2n}a^{2n-3}\frac{1}{a^3}=\comb{3}{2n}a^{2n-6}\\
    {\rm T}_6&=\comb{5}{2n}a^{2n-5}\frac{1}{a^5}=\comb{5}{2n}a^{2n-10}
\end{split}\]
依题意得$\comb{3}{2n}=\comb{5}{2n}$，即$\comb{2n-3}{2n}=\comb{5}{2n}$，所以
\[2n-3=5,\quad 2n=8,\quad n=4\]
故所给的二项式是$\left(a+\frac{1}{a}\right)^8$

因为${\rm T}_{k+1}=\comb{k}{8}a^{8-k}\frac{1}{a^k}=\comb{k}{8}a^{8-2k}$，
若${\rm T}_{k+1}$为不含$a$项，则$8-2k=0$，所以$k=4$.
\[{\rm T}_5=\comb{4}{8}\cdot a^0=\frac{8\x 7\x 6\x 5}{4\x 3\x 2}=70\]
\end{solution}

\begin{ex}
\begin{enumerate}
    \item 展开$(x^{3}-2x)^{7}$
    \item 求$(a+\sqrt{b})^{12}$的展开式里的第4项和第9项
    \item \begin{enumerate}[(1)]
        \item 求$(x+a)^{12}$的展开式里的中间一项；
        \item 求$\left(\frac x2+\frac2x\right)^7$的展开式里的中间两项.
    \end{enumerate}  
    \item 在$(x+a)^{10}$的展开式里，求
\begin{multicols}{2}
    \begin{enumerate}[(1)]
        \item 含有$a^{8}$的项；
        \item 含有$x^8$的项. 
    \end{enumerate}    
\end{multicols}

    \item 在$\left(x^{2}+\frac{1}{x^3}\right)^{15}$的展开式里，有没有
\begin{multicols}{2}
    \begin{enumerate}
        \item 不含$x$的项？
        \item 含有$x^{8}$的项？
    \end{enumerate}
\end{multicols}
如果有，把这项求出来.
    
\item 求$\left(3x+\frac{1}{\sqrt[3]{x}}\right)^{8}$的展开式里的常数项. 
\item 已知$\left(\sqrt[8]{a^{2}}+\frac1a\right)^{n}$的展开式里的第3项含有$a^2$,    
    求$n$.
\end{enumerate}
\end{ex}

\subsection{二项式的性质}

我们研究二项式$(x+y)^n$的展开式中各项系数的性质。二项展开式各项系数是：
\[\comb{0}{n},\;\comb{1}{n},\;\comb{2}{n},\;\ldots,\; \comb{n-1}{n},\;\comb{n}{n}\]

展开式各项系数有下列性质
\begin{enumerate}
\item 在二项展开式中，与首末两端“等距离”的两项系数相等. 我们已知$\comb{m}{n}=\comb{n-m}{n}$，所以:
\[\comb{0}{n}=\comb{n}{n},\; \comb{1}{n}=\comb{n-1}{n},\; \comb{2}{n}=\comb{n-2}{n},\; \ldots,\; \comb{k}{n}=\comb{n-k}{n},\ldots\]
\item 如果二项式的幂指数是偶数，中间一项的系数大，如果二项式的幂指数是奇数，中间两项系数相同并且大.
\end{enumerate}

由于展开式各项系数顺次是
\[\comb{0}{n}=1,\quad \comb{1}{n}=n,\quad \comb{2}{n}=\frac{n(n-1)}{1\cdot 2},\ldots\]
\[\comb{k}{n}=\frac{n(n-1)(n-2)\cdots(n-k+1)}{1\cdot 2\cdot 3\cdots k},\ldots,\quad \comb{n}{n}=1\]
其中，后一个系数的分子是前一个系数的分子乘以逐次减小1的数，(如$n,\; n-1,\; n-2,\; \ldots$), 分母是乘以逐次增大1的数(如$1,2,3,\ldots$)得到的，因而，各项系数从开始起总是逐渐增大，又因为与首末两端“等距离”的两项系数相等，所以系数增大到某一项时就逐渐减小，而最大项必在中间.

当$n$是偶数时，$n+1$是奇数，展开式共有$n+1$项，所以
展开式有中间一项，即第($\frac{n}{2}+1$)项，并且系数最大。

当$n$是奇数时，$n+1$是偶数，展开式共有$n+1$项，所以有中间两项，即第($\frac{n+1}2$)和第($\frac{n+1}2+1$)两项，系数相同并且最大.

\begin{example}
    求$(x+y)^{10}$和$(x+y)^{11}$的最大系数的项.
\end{example}

\begin{solution}
    $(x+y)^{10}$最大系数项的项数是
$\frac{10}{2}+1=6$，所以
$$\mathrm{T}_{6}=\mathrm{C}_{10}^{5} x^{10-5}y^{5}=252x^{5}y^{5}$$

$(x+y)^{11}$最大系数的项 数是$\frac{11+1}{2}=6$或$6+1=7$, 所以
\[\begin{split}
    \mathrm{T}_{6}&=\mathrm{C}_{11}^{5}x^{11-5}y^{5}=462x^{6}y^{5}\\
    \mathrm{T}_{7}&=\mathrm{C}_{11}^{6}x^{11-6}y^{6}=462x^{5}y^{6}
\end{split}\]
\end{solution}

\begin{example}
    在$(5x+2y)^{29}$展开式中第几项的系数最大，并
求它的系数。
\end{example}

\begin{solution}
    设$T_{k+1}$系数最大，则它必须满足
\[\begin{cases}
\frac{{\rm T}_{k+1}\text{系数}}{{\rm T}_{k}\text{系数}}=\frac{\comb{k}{29}5^{29-k}2^k}{\comb{k-1}{29}5^{29-k+1}2^{k-1}}=\frac{2(29-k+1)}{5k}\ge 1\\
\frac{{\rm T}_{k+1}\text{系数}}{{\rm T}_{k+2}\text{系数}}=\frac{\comb{k}{29}5^{29-k}2^k}{\comb{k+1}{29}5^{29-k-1}2^{k+1}}=\frac{5(k+1)}{2(29-k)}\ge 1\\
\end{cases}\]
即
$\begin{cases}
    7k\le 60\\
    \frac{53}{7}\le k\le 29
\end{cases}$

解得$7\frac{4}{7}\le k\le 8\frac{4}{7}$. 因为$k$为正整数，故$k=8$，即系数最大项系数是
\[{\rm T}_{8+1}=\comb{8}{29}5^{29-8}\cdot 2^8=\comb{8}{29}5^{21}\cdot 2^8\]
\end{solution}

\begin{example}
    求证：$\sum\limits^n_{k=0}\comb{k}{n}=\comb{0}{n}+\comb{1}{n}+\comb{2}{n}+\cdots+\comb{n-1}{n}+\comb{n}{n}=2^n$
\end{example}

\begin{proof}
在二项式公式中取$x=y=1$，有：
\[(1+1)^n=\sum^n_{k=0}\comb{k}{n}1^{n-k}1^k=\sum^n_{k=0}\comb{k}{n}\]
所以$\sum\limits^n_{k=0}\comb{k}{n}=2^n$
\end{proof}

\begin{example}
证明在$(x+y)^n$展开式中，奇数项系数的和等于
偶数项系数的和.
\end{example}

\begin{proof}
    在二项式公式中取$x=1$, $y=-1$，有
    \[(1-1)^n=\sum^n_{k=0}(-1)^k\comb{k}{n}1^{n-k}1^k=\sum^n_{k=0}(-1)^k\comb{k}{n}\]
即
\[0=\left(\comb{0}{n}+\comb{2}{n}+\cdots\right)-\left(\comb{1}{n}+\comb{3}{n}+\cdots\right)\]
所以
\[\comb{0}{n}+\comb{2}{n}+\cdots=\comb{1}{n}+\comb{3}{n}+\cdots\]
\end{proof}

\begin{example}
求证：$\sum\limits^n_{k=1}2k\comb{k}{n}=n\cdot 2^n$
\end{example}

\begin{proof}
因为
\[\begin{split}
    k\comb{k}{n}&=k\cdot \frac{n(n-1)\cdots (n-k+1)}{k!}\\
    &=n\cdot \frac{(n-1)(n-2)\cdots (n-k+1)}{(k-1)!}=n\comb{k-1}{n-1}
\end{split}\]
所以
\[\sum^n_{k=1}2k\comb{k}{n}=2\sum^n_{k=1}k\comb{k}{n}=2n\sum^{n-1}_{k=0}\comb{k}{n-1}=2n\cdot 2^{n-1}=n\cdot 2^n\]
\end{proof}


\begin{example}
    计算$(1.009)^{20}$和$(0.991)^{5}$的近似值（精确到0.001)
\end{example}

\begin{solution}
\[\begin{split}
    (1.009)^{20}&=(1+0.009)^{20}\\
&=1+20\x0.009+190\x(0.009)^2+1140\x(0.009)^3+\cdots\\
&=1+0.18+0.01539+0.00083106+\cdots
\end{split}\]
可以看出，要使所求的近似值的误差不超过0.001，需取展开式中前三项的和，就是
\[\begin{split}
(1.009)^{20}&\approx 1+0.18+0.01539\approx 1.195\\
(0.991)^5&=(1-0.009)^5=1-5\x0.009+10\x(0.009)^2+\cdots    
\end{split}\]
根据精确度要求，第三项以后各项都可以删去，所以，
\[(0.991)^5\approx 1-0.045=0.955\]
\end{solution}

\begin{ex}
\begin{enumerate}
    \item 求二项展开式里系数最大的项，（系数只需用组合数符号表示，不必算出）：
\begin{multicols}{2}
\begin{enumerate}[(1)]
    \item $(a+b)^{10}$
    \item $(1+x)^{17}$
\end{enumerate}
\end{multicols}
    \item  应用二项式定理展开$(x^2+2x+2)^5$
    \item  $(x+y)^{2n}$的展开式里如果第5项的系数与第13项的系数相等，求展开式里系数最大的项（系数用组合符号表示，不必算出）.
    \item  求$(1-x)^{13}$的奇次项系数的和.
    \item 证明：$\comb{0}{n}+\comb{2}{n}+\comb{4}{n}+\cdots+\comb{n}{n}=2^{n-1}$（$n$为偶数）
    \item 计算$(3.002)^0$和$(3.998)^5$误差不超过0.001的近似值.
\end{enumerate}
\end{ex}

\section*{习题2.3}
\begin{enumerate}
    \item 用二项式定理展开。
\begin{multicols}{2}
\begin{enumerate}[(1)]
    \item $\left( \frac {\sqrt {x}}2- \frac 2{\sqrt {x}}\right) ^{7}$
    \item $\left( \frac {x^{2}}2+ \frac {3}{x}\right) ^{8}$
\end{enumerate}    
\end{multicols}
\item 化简：
\begin{enumerate}[(1)]
    \item $(1+\sqrt{x})^{5}+(1-\sqrt{x})^{5}$
    \item $\left( 2x^{\tfrac 12}+ 3x^{- \tfrac 12}\right) ^{4}- \left( 2x^{\tfrac 12}- 3x^{- \tfrac 12}\right) ^{4}$
\end{enumerate}
\item \begin{enumerate}[(1)]
    \item 求$(1-2x)^{15}$的展开式中前四项；
    \item 求$(2a^{3}-2b^{2})^{10}$的展开式中第11项;
    \item 求$\left(\frac{\sqrt{x}}{3}+\frac{3}{\sqrt{x}}\right)^{12}$的展开式的中间一项;
\item 求$(x\sqrt{y}-y\sqrt{x})^{15}$的展开式的中间两项.
\end{enumerate}
\item 求下列各式的展开式中指定各项的系数。
\begin{enumerate}[(1)]
    \item $\left( 1- \frac 1{2x}\right)^{10}$含$\frac1{x^5}$的项；
    \item $\left(2x^{3}-\frac{1}{2x^{2}}\right)^{10}$的常数项.
\end{enumerate}
 
\item \begin{enumerate}[(1)]
    \item $\left({\sqrt {x}} + \frac 1{\sqrt [3]{x^{2}}}\right) ^{n}$的展开式里第5项系数与
第3项系数的比是7:2, 求展开式里含有$x$的项.
\item 已知$(1+x)^n$的展开式里第四项与第八项的系数相等，求这两项的系数.
\end{enumerate}

\item \begin{enumerate}[(1)]
    \item 求$\left(\frac{x}{y}-\frac{y}{x}\right)^{13}$展开式里系数最大的项.
    \item 写出$\left(x^2-\frac{1}{\sqrt{x}}\right)^{10}$的展开式里系数绝对值最大的项.
    \item 求$(1+x+x^2)(1-x)^5$的展开式里的中间项.
\end{enumerate}

\item 用二项式定理证明：
\begin{enumerate}[(1)]
    \item $(n+1)^{n}-1$能被$n^2$整除；
    \item $99^{10}-1$能被1000整除.
\end{enumerate}

\item 证明：
\begin{enumerate}[(1)]
    \item $\left(x-\frac{1}{x}\right)^{2n}$的展开式中常数项是$(-2)^n\frac{1\cdot 3\cdot 5\cdots (2n-1)}{n!}$
    \item $(1+x)^{2n}$的展开式的中间一项是$\frac{1\cdot 3\cdot 5\cdots (2n-1)}{n!}(2x)^n$
\end{enumerate}

\item 求证：
\begin{enumerate}[(1)]
    \item $\sum\limits^m_{k=0}\comb{n}{n+k}=\comb{n+1}{n+m+1}$
    \item $2^n-\comb{1}{n}2^{n-1}+\comb{2}{n}2^{n-2}+\cdots+(-1)^{n-1}\comb{n-1}{n}\cdot 2+(-1)^n=1$
    \item $\sum\limits^n_{k=0}(2k+1)\comb{k}{n}=(n+1)2^n$
    \item $1+2\comb{1}{n}+\comb{2}{n}+2\comb{3}{n}+\comb{4}{n}+2\comb{5}{n}+\cdots +2\comb{n-1}{n}+\comb{n}{n}=3\cdot 2^{n-1}$（其中$n$是偶数）
\end{enumerate}
\item 求$(1+x_1)(1+x_2)(1+x_3)\cdots(1+x_n)$的展开式里各项系数之和.
\end{enumerate}


\section*{本章内容要点}
一、本章主要内容是排列与组合，二项式定理；重点是不许重复的排列与组合，二项式定理及其性质.

二、两个基本原理——加法原理与乘法原理，是推导排列与组合数公式的基础，也是直接解决某些计数问题的方法. 应用中，要注意区别两个原理：

加法原理与分类有关，即如果完成一件事有$n$类相互独立的办法，那么求完成这件事情的方法种数，就要用加法原理；

乘法原理与分步有关，即如果完成一件事需要分$n$步完成，缺一不可，而每一步又各有若干方法，那么，求完成这件事的方法种数，就要用乘法原理。

三、排列与组合的区别是看研究的问题是否与元素选取的顺序有关，与顺序有关就属于排列；与顺序无关就属于组合.

排列又分为不许重复选取的全排列与选排列、可重复的排列、环形排列；

组合又分为不许重复的组合与可重复的组合。

四、排列与组合的计数公式有：
\begin{enumerate}
    \item 从$n$个不同元素中，任取$m$个没有重复的元素：

排列数公式
\begin{itemize}
    \item 选排列：$\perm{m}{n}=n(n-1)\cdots(n-m+1)=\frac{n!}{(n-m)!}\quad (m<n)$
    \item 全排列：$\perm{n}{n}=n!$
\end{itemize}

组合数公式：$\comb{m}{n}=\frac{\perm{m}{n}}{\perm{m}{m}}\quad (m\le n)$

组合数的性质：$\comb{m}{n}=\comb{n-m}{n},\qquad \comb{m}{n+1}=\comb{m}{n}+\comb{m-1}{n}\quad (m\le n)$

环形排列数公式：$\comb{m}{n}\cdot \perm{m-1}{m-1}=(m-1)! \comb{m}{n}\quad (m\le n)$

\item 可重复的排列数与组合数公式：
\begin{itemize}
    \item 重复排列数$n^m$
    \item 重复组合数${\rm H}^m_n=\comb{m}{n+m-1}$
\end{itemize}
\end{enumerate}

五、二项式定理指出了二项代数和的幂的展开公式。
\[(a+ b) ^{n}=\mathrm{C}_{n}^{0}a^{n}+{\rm C}_n^{1}a^{n- 1}b+ \cdots + \mathrm{C}_{n}^{r}a^{n- r}b^{r}+ \cdots + \mathrm{C} _{n}^{n}b^{n}\]
其中，展开式的通项(第$r+1$项)为：
$$\mathrm{T}_{r+1}=\mathrm{C}_{n}^{r}a^{n-r}b^{r}$$

二项式的展开式中，系数、幂指数、项数等都有很好的
性质，其中主要性质有：
\begin{enumerate}
    \item 
展开式中，与首末两端“等距”的两项系数相等，
二项式$n$次幂的展开式共有$n+1$项.
\item 
当 幂 指 数 $n$为 偶 数 时 , 中 间 一 项 的 系 数${\rm C}_n^{\tfrac{n}{2}}$最大(第$\frac n2+1$项); 当$n$为奇数时，中间有两项系数最大，即第$\frac{n+1}2$项的系数$C_{n}^{\tfrac{n-1}2}$与第$\frac{n+1}2+1$项的系数${\rm C}_{n}^{\tfrac{n+1}2}$相等且最大.
\end{enumerate}

\section*{复习题二}
\begin{enumerate}
    \item 解不等式：
\begin{multicols}{2}
\begin{enumerate}[(1)]
    \item $2<\frac{\perm{m+1}{m+1}}{\perm{m-1}{m-1}}\le 42$
    \item $x\perm{3}{8}<\perm{3}{x}$
\end{enumerate}
\end{multicols}

\item 求证：
\begin{enumerate}[(1)]
    \item $1!+2\cdot 2!+3\cdot 3!+\cdots+n\cdot n!=(n+1)!-1$
    
    （提示：$n\cdot n!=(n+1)!-n!$）

\item $\comb{m}{n-1}+\comb{m}{n-2}+\comb{m}{n-3}+\cdots+\comb{m}{m+1}+\comb{m}{m}=\comb{m+1}{n}$
\end{enumerate}

\item \begin{enumerate}[(1)]
    \item 已知$\frac{1}{\comb{m}{5}}-\frac{1}{\comb{m}{6}}=\frac{7}{10\cdot\comb{m}{7}}$，求$\comb{m}{8}$
    \item 已知$\frac{\comb{m-1}{n}}{2}=\frac{\comb{m}{n}}{3}=\frac{\comb{m+1}{n}}{4}$，求$n$与$m$.
\end{enumerate}

\item 用数字0, 1, 2, 3, 4, 5组成没有重复数字的数：
\begin{enumerate}[(1)]
\item 能够组成多少个六位奇数？
\item 能够组成多少个大于201345的自然数？ 
\end{enumerate}

\item 8个不同的元素排成一行：
\begin{enumerate}[(1)]
 \item 其中某2个元素要排在一起，有多少种排法？
\item 其中某2个元素不排在一起，有多少种排法？
\item 其中某4个元素要排在一起，另外4个元素也要排在一起，有多少种排法？  
\end{enumerate}

\item 书架上有4本不同的数学书，5本不同的物理书，3本不同的化学书，全部竖起排成一排，如果不使同类的书分开，一共有多少种排法？
\item 从1到9这九个数字每次取出5个数字：
\begin{enumerate}[(1)]
\item 可以组成多少个能被25整除的五位数？
\item 可以组成多少个小于98000的五位数？
\item 可以组成多少个五位偶数？
\item 可以组成多少个奇数位是奇数，偶数位是偶数的五位数？   
\end{enumerate}

\item \begin{enumerate}[(1)]
    \item 一个集合由8个不同的元素组成，这个集合中3个元素的子集有几个？
    \item 一个集合由5个不同的元素组成，其中含1个，2个，3个，4个元素的子集共有几个？
\end{enumerate}

\item \begin{enumerate}[(1)]
\item 平面内有$n$条直线，其中没有两条互相平行，也没有三条相交于一点，一共有多少个交点？
\item 空间有$n$个平面，其中没有两个互相平行，也没有三个相交于一直线，一共有多少条交线？
\item 平面有两组平行线，一组$m$条，另一组$n$条，这两组平行线相交，可以构成多少个平行四边形？
\item 空间有三组平行平面，第一组有$m$个，第二组有$n$个，第三组有$\ell$个，不同组的平面都相交，可以构成多少个平行六面体？
\end{enumerate}

\item 一个班有48名学生，分成六排，每排八人
\begin{enumerate}[(1)]
\item 问：有多少种不同坐法？
\item 若有2人必须坐在前排，3人必须在最后一排，有多少种不同的坐法？
\end{enumerate}

\item 由0至9这10个数字组成的没有重复数字的九位数字有多少个？其中能被3整除的有多少个？
\item 有$n+1$件不同的奖品，全部赠给$n$个先进学生代表．如果每人至少得到一件，有多少种不同的赠送方法？
\item \begin{enumerate}[(1)]
    \item 求$\left(9x-\frac{1}{\sqrt[3]{x}}\right)^{10}$展开式的常数项;
    \item 求$(1+x+x^2)(1-x)^{10}$展开式中$x^4$的系数;
    \item 求$(1-x)^{11}+x(1-2x)^{10}-x^2(1-3x)^9$展开式中$x^5$的系数;
    \item 求$（1+2x-3x^2)^6$展开式中$x^5$的系数;
    \item 求$(x^{-1}-1+x)^5$展开式里不含$x$的项.
\end{enumerate}
\item \begin{enumerate}[(1)]
    \item 用二项式定理证明$55^{55}+9$能被8整除；
    \item 用二项式定理求$89^{10}$除以88的余数．
\end{enumerate}

\item 证明：$(1+x)^{2n}$展开式中$x^n$的系数等于$(1+x)^{2n-1}$展开式中$x^n$的系数的2倍．
\item 已知$\left(\sqrt{x}+\frac{1}{\sqrt[3]{x}}\right)^n$展开式的系数之和比$(a+b)^{2n}$展开式的系数之和小240，求：
\begin{enumerate}[(1)]
    \item $\left(\sqrt{x}+\frac{1}{\sqrt[3]{x}}\right)^n$展开式的第三项：
    \item $(a+b)^{2n}$展开式的中间项.
\end{enumerate}

\item 已知$(1+x)^n$展开式里连续三项系数的比是3:8:14，求展开式里系数最大的项．
\item 已知$\left(\sqrt[3]{x^{-2}}+x^2\right)^{2n}$的展开式的系数之和比$\left(\sqrt[3]{x^{-2}}+x^2\right)^{n}$的展开式的系数之和大992，求$\left(\sqrt[3]{x^{-2}}-x^2\right)^{n}$的展开式里含$x^{\tfrac{4}{3}}$项的系数.

\item 证明：
\begin{enumerate}[(1)]
    \item $\sum\limits^n_{k=0}\left(\comb{k}{n}\right)^2=\frac{(2n)!}{n!\cdot n!}$
    \item $\comb{0}{n}\comb{k}{m}+\comb{1}{n}\comb{k-1}{m}+\cdots+\comb{k}{n}\comb{0}{m}=\comb{k}{n+m}$
    \item $\left(\comb{0}{n}-\comb{2}{n}+\comb{4}{n}-\cdots\right)^2+ \left(\comb{1}{n}-\comb{3}{n}+\comb{5}{n}-\cdots\right)^2=\comb{0}{n}+\comb{1}{n}+\comb{2}{n}+\cdots+\comb{n}{n}$
\end{enumerate}

\item 在$(a-3b)^{35}$的展开式中第几项的系数的绝对值最大？
\end{enumerate}


\chapter{概率论初步}

\section{随机事件与概率}

\subsection{随机事件及其概率}
在一定的条件下，可能发生也可能不发生的事件称为随机事件.

\begin{example}
    掷一枚硬币，“正面向上”这一事件可能发生，也可能不发生，因此，它是一随机事件。
\end{example}

\begin{example}
    对准靶发射一枪，“中靶”这一事件可能发生，也可能不发生，因此，它是一随机事件.
\end{example}

\begin{example}
    将1,2,3这三个数字随便排列，就得一个三位数，“这个三位数大于300”可能发生，也可能不发生，因此，它是一随机事件.
\end{example}


在现实世界中，有许多事件在一定的条件下，可能发生也可能不发生，这些都是随机事件. 通常，我们用大写的英文字母$A,B,C,\ldots$表示随机事件，概率就是反映随机事件出现的可能性大小的一个量；在同一条件下，有些事件发生的可能性比另一些更大些，它们相应的概率也就比另一些的大一些，所以随机事件的概率是客观存在的，问题是如何求出随机事件相应的概率. 为了讨论方便起见，习惯上把“必然事件”与“不可能事件”也看作随机事件。“必然事件”就是在给定条件下必然发生的事件，通常用字母$U$表示. “不可能事件”就是在给定条件下必然不发生的事件，通常用字母$V$表示．在例3.1中，“正面向上或反面向上”就是一个必然事件，“正面和反面都向上”就是一个不可能事件；在例3.3中，“三位数小于400”是一个必然事件，“三位数小于110”就是一个不可能事件．

为了讨论反映随机事件的可能性的大小的量——概率。我们进一步考察“投掷一枚硬币”、“从50只三级管中任抽一只检查”这样一类随机现象。例如，投掷一枚硬币，共有“出现正面”和“出现反面”两种试验结果，通常总认为硬币是“均匀的”，所以出现这两种结果的可能性是相同的。类似地，从50只三级管中任抽一只，共有50种不同的结果，而且只要检查员事先不带主观偏见，抽到任何一只的可能性都相同. 今后，我们称每一试验结果为基本事件，在具体问题中，弄清楚所有基本事件是十分重要的事，首先要注意到不同的基本事件是不能同时发生的；其次要注意到事件总是由若干个基本事件组成. 某一事件所包含的基本事件发生了，就说该事件发生了.

从这两个例子可以看到，这一类随机现象有两个共同的特征：
\begin{enumerate}[(1)]
    \item 有限性：上面每次试验只可能有有限种不同的结果，即有限个基本事件；
    \item 等可能性：由于其自然对称性，出现每一基本事件的可能性是相同的.
\end{enumerate}

这是一类最简单但却是常见的随机现象，我们称它为古典概率模型（简称古典概型），如何来描述这类事件的可能性的大小呢？例如我们问：“50只三级管中有2只次品，从
中任抽一只，抽得次品的可能性有多大？”人们常常用“次品率”来描述这种可能性，这时，次品数（2）和全产品数（50）之比为$2/50=4\%$．我们就说这批三级管的次品率是4\%．对于上面的分母和分子，我们也可以这样来看：50是抽一只检查（看作一次试验）时所有可能的基本事件数，2是“抽一只检查得次品”这一事件所包含的基本事件数.

因此对于古典概型，如果试验的基本事件总数为$n$，随机事件$A$所包含的基本事件数为$m$，我们就用数$m/n$来描述每次试验中事件$A$出现的可能性的大小，称它为事件$A$的概率，记作$\Pr(A)$. 即我们定义
\[\Pr(A)=\frac{m}{n}\]

于是对“投掷一枚硬币出现正面”这一随机事件中，由于试验的基本事件总数为2，而“出现正面”只包含其中一个基本事件，所以它的概率等于1/2．对于这样定义的概率显然有

\begin{blk}
    {性质1} \[\Pr(U)=1,\qquad \Pr(V)=0\]
\end{blk}

因为必然事件$U$包含了所有基本事件，所以$\Pr(U)=\frac{n}{n}=1$，而不可能事件$V$不含任何基本事件，所以$\Pr(V)=\frac{0}{n}=0$.

\begin{blk}
 { 性质2}\[0\le \Pr(A)\le 1\]  
\end{blk}

因为事件$A$所含基本事件数$m$满足不等式$0\le m\le n$，所以$0\le \Pr(A)\le 1$.

要算出基本事件总数及A所包含的基本事件数，“排列”、“组合”常常是十分有用的工具.

\begin{example}
    $1,2,3,4,5,6$六个数字中任取两数，试计算它们都是偶数的概率.
\end{example}

\begin{solution}
    六个数字中任取两个，共有$\comb{2}{6}=15$种不同的取法，这就是全体基本事件，即$n=15$. 而事件$A$：“任取两数都是偶数”则要从$2,4,6$三个偶数中取两个，也就是说事件$A$包含$\comb{2}{3}=3$个基本事件，因此它的概率
\[\Pr(A)=\frac{\comb{2}{3}}{\comb{2}{6}}=\frac{3}{15}=\frac{1}{5}\]
\end{solution}


\begin{example}
    在一批零件中有$n$个一等品，$m$个三等品，逐个进行检查. 若已查明前$k\;(k<n)$个均为一等品，求第$k+1$次检查时仍得一等品的概率.
\end{example}

\begin{solution}
    由于已查明前$k$个均为一等品，所求第$k+1$次检查基本事件总数为$n+m-k$，“得一等品”这一事件包含的基本事件数为$n-k$. 因此所求的概率
    \[P=\frac{n-k}{n+m-k}\]
\end{solution}


\begin{example}
    一部四卷文集，按任意次序放到书架上，问各卷自左至右或自右至左的卷号恰为1, 2, 3, 4的顺序的概率是多少？
\end{example}

\begin{solution}
    一部四卷文集，在书架上可有$\perm{4}{4}=4!=24$种不同的排列方法，其中自左至右或自右至左卷号恰为1, 2, 3, 4的顺序的共有2种．因此所求的概率
    \[P=\frac{2}{24}=\frac{1}{12}\]
\end{solution}

\begin{example}
    某停车场有12个位置列成一行．求有8个位置停了车而空着的四个位置连在一起的概率.
\end{example}

\begin{solution}
12个位置中占去8个，共有$\comb{8}{12}$种方法，把四个接连的空位看作一个位置，发生四个接连的位置空着的情形相当于这个假想的位置插入8个停车位置中间或者它们的两端，共有9种不同方法，因此所求的概率为
\[P=\frac{9}{\comb{8}{12}}=\frac{1}{55}\]
\end{solution}

\begin{example}
    鱼池中共有鱼$N$条，从中捕得$t$条，加了标志后立即放回池中，经过一段时间后，再从池中抽出$n$条鱼，问其中有$s$条标志鱼的概率是多少？
\end{example}

\begin{solution}
    从$N$条鱼中捕得$n$条，共有$\comb{n}{N}$种可能结果，它就是基本事件总数. 在捕得的$n$条鱼中恰有$s$条标志鱼，它们是从$t$条标志鱼中捕得的，这种捕法共有$\comb{s}{t}$种. 而对于这样的每一种捕法，其余的$n-s$条未加标志的鱼是从$N-t$条鱼中捕得的，共有$\comb{n-s}{N-t}$种捕法. 因此“捕得$n$条鱼，其中$s$条是标志鱼”的捕法共有$\comb{s}{t}\comb{n-s}{N-t}$种，所求概率
\[P_s(N)=\frac{\comb{s}{t}\comb{n-s}{N-t}}{\comb{n}{N}}\]
\end{solution}

\begin{example}
    某小组有成员3人，每人在某星期7天中参加劳动一天，如果劳动日期可随机安排，求3人在不同的3天参加劳动的概率.
\end{example}

\begin{solution}
    对于某一组员（甲）来说，他可以安排在7天中的任何一天参加劳动，即有7种安排方法，而对于甲的任何一种安排另一组员（乙）又可有7种安排法，因此甲乙两人共有$7^2$种不同的安排法。类似地，对于$7^2$种的任何一种安排，第三名组员（丙）又可有7种安排法．因此所有可能的安排法——总的基本事件数等于$7^3$，而3人安排在7天中的不同的3天有${\rm A}^3_7$种不同的方法，所以要求的概率为
\[P=\frac{{\rm A}^3_7}{7^3}=\frac{30}{49}\]
\end{solution}

上述例3.9颇具典型性，$N$个人在$n$天中的任意安排法相当于把$N$个球随机地放入$n$个盒子里去的方法，共有$n^N$种. 这一类问题称为盒子模型.

\begin{ex}
\begin{enumerate}
    \item 指出下列事物是必然事件，不可能事件，还是随机事件.
\begin{enumerate}[(1)]
\item 如果$a$，$b$都是实数，那么$a+b=b+a$;
\item 从一副扑克牌中任意抽取两张，得到“黑桃”；
\item 某电话总机在一分钟内接到20次呼唤．
\end{enumerate}
\item 某班有50位同学而其中女同学占15名．今碰到这个班的同学，正好碰到一个男同学的概率是多少？正好碰到一个女同学的概率是多少？这两个概率之间有什么关系？为什么？
\item 盒中有100个铁钉，其中有90个是合格的，而10个是坏的，从中任意抽取10个，问其中没有一个为坏的概率是多少？
\item 设袋中有8个球，其中5个白球3个红球，从中任意抽取4个，问恰好抽到3个白球的概率是多少？
\item 两袋分别盛着写有0, 1, 2, 3, 4, 5六个数字的六张卡片，从每袋中各取一张，求所得两数之和等于6的概率．现在有人给出下述两种不同解答：

解1：两数之和共有$0,1,2,\ldots,10$十一种不同结果，因此，所求的概率为$1/11$．

解2：从每袋中各任取一张卡片，共有$6^2$种取法，其中和数为6的情形有5种；
\[(1,5);\quad (2,4);\quad (3,3);\quad (4,2);\quad (5,1)\]
因此所求的概率为$5/36$．

试问哪一种解法正确？为什么？
\item 号码锁上有六个拨盘，每个拨盘上有0—9共10个数字，当这6个拨盘上的数字组成某一个6位数时（第一位可以是0），锁才能打开。如果不知道锁的号码，一次就把锁打开的概率是多少？
\item 设有50张考签，分别加以标号$1,2,\ldots,50$，一学生任意抽一张进行考试，求抽到前10号考签的概率？
\item 一批产品共有100件，其中有次品5件，从中任取10件，求所取10件中至多有2件是次品的概率。
\item 袋中装有$n$个白球和$m$个黑球，从中任取$a+b$个，求所取的球恰含$a$个白球和$b$个黑球的概率？
\end{enumerate}
\end{ex}

\subsection{等可能性事件的概率}
在古典概型里，我们曾规定了$\Pr(A)=m/n$，这是在等可能性的大前提下得出的，这种类型在我们计算事件发生的
概率时是常常碰到的，也是较为重要的，为着重视这种类型，我们再深入研究这类问题.

\begin{example}
    一个均匀材料造的正方体玩具，各个面上分别标以数字$1,2,3,4,5,6$，这种玩具叫骰子，
\begin{enumerate}[(1)]
\item 将骰子抛掷一次，朝上的一面出现奇数的概率是多少？
\item 抛掷2次，朝上的一面的数字之和为7的概率是多少？
\end{enumerate}
\end{example}

\begin{solution}
\begin{enumerate}[(1)]
\item 朝上的一面可为$1,2,3,4,5,6$等六种之一，且每种可能性相等，朝上一面出现奇数，可为$1,3,5$等三种之一，且每种可能性相等，设所问的事件概率为$\Pr(A)$，则$\Pr(A)=\frac{3}{6}=\frac{1}{2}$
\item 抛掷二次，每次朝上一面可有六种情况，两次则共有$6\x6=36$种情况，朝上一面的数字之和为7，则有第一次是1，第二次是6，第一次是6，第二次是1等两种情况。
2与5；3与4，又各有两种，总共有6种，故概率$P=\frac{6}{36}=\frac{1}{6}$.
\end{enumerate}
\end{solution}

\begin{example}
    一次掷三枚骰子，得3点，4点，5点，10点的概率各是多少？
\end{example}

\begin{solution}
\begin{enumerate}[(1)]
    \item 得3点的概率，掷三颗骰子共有216种情况，而得3点只有三颗都为1时的一种，故掷得3点的概率为$\frac{1}{216}$；
    \item 掷三颗骰子共有216种情况，而得4点有$2,1,1$; $1,2,1$; $1,1,2$等三种情况，故掷得4点的概率为$\frac{3}{216}=\frac{1}{72}$；
\item 掷三颗骰子共有216种情况，而得5点有$1,1,3$; $1,3,1$; $3,1,1$; $2,2,1$; $2,1,2$; $1,2,2$，即$\comb{1}{3}+\comb{2}{3}=6$种，故掷得5点的概率为$\frac{6}{216}=\frac{1}{36}$；
\item 掷三颗骰子共有216种情况，而得10点有$1,3,6$; $1,4,5$; $2,2,6$; $2,3,5$; $2,4,4$; $3,3,4$等六种
情况．而$1,3,6$；$1,4,5$；$2,3,5$各有$\perm{3}{3}$种排列．$2,2,6$；$2,4,4$；$3,3,4$有$\comb{1}{3}$种组合. 共有$3\perm{3}{3}+3\comb{1}{3}=27$种. 10点的概率为$\frac{27}{216}=\frac{1}{8}$.
\end{enumerate}  

读者试将掷三颗骰子得3点至18点各点数的概率分别求看看各个概率之间有什么关系，各个概率之和又是什么？
\end{solution}

\begin{example}
    袋中装有$a$个白球和$b$个黑球，从中任意地取出$k+1$（$k+1\le a+b$）个球，每次取一球，取后不放回，求最后取出的一球是白球的概率？
\end{example}

\begin{solution}
    从$a+b$个球中无放回地接连取出$k+1$个球，一共有
$(a+b)(a+b-1)\cdots(a+b-k-1+1)={\rm A}^{k+1}_{a+b}$种取法。

$A$表“取出的$k+1$个球中最后一球是白球”的事件。最后取出的白球可以是$a$个白球中的任何一个，有$a$种；而其余$k$个可以是余下的$a+b-1$中的任意$k$个，有${\rm A}^k_{a+b-1}$种取法，因此事件$A$共含$a\cdot {\rm A}^k_{a+b-1}$个不同的基本事件，故所求概率为
\[\Pr(A)=\frac{a{\rm A}^{k}_{a+b-1}}{{\rm A}^{k+1}_{a+b}}=\frac{a}{a+b}\]

从上式可以看出，$\Pr(A)$的值与$k$无关，它说明了无论在哪个序号上，取出白球的概率总是相同的。
\end{solution}

\begin{example}
    有$n$个可辨的球，随机地放入$N\; (N\ge n)$个盒中，每一个盒可以放任意多个球，试求：
\begin{enumerate}
\item 指定的$n$个盒中各有一球的概率？
\item 某$n$个盒中各恰有一球的概率？
\item 某指定的一个盒中恰有$m$（$m\le n$）个球的概率？
\end{enumerate}
\end{example}

\begin{solution}
    设球以同样的可能性落入$N$个盒中的每一个，$N$个盒中落一个球一共有$N$种方法，$n$个球可以有$N^n$种方法，这就是总的基本事件数.
\begin{enumerate}
\item 以$A$表“指定的$n$个盒中各有一球”的事件．今固定某$n$个盒，第一个球可以落入这$n$个盒中的任何一个，有$n$种方法，第二个球可落在余下的$n-1$个盒中的任何一个，有$n-1$种方法，……，第$n$个球落在最后的一个盒，只能有一种方法，因此事件$A$共含有$n!$个不同的基本事件，故所求概率为
\[\Pr(A)= \frac{n!}{N^n} \]
\item 以$B$表“某$n$个盒中各恰有一球”的事件．因为$n$个盒可从$N$个盒中任意选取，共有$\comb{n}{N}$种选法. 选出这$n$个盒后再按问题1知事件$B$共含$\comb{n}{N}n!$个不同的基本事件数，故所求概率为
\[\Pr(B)=\frac{\comb{n}{N}n!}{N^n}\]
\item 
以$C$表“某指定的一个盒中恰有$m$个球”的事件，因为$m$个球可从$n$个球中任意选出，共有$\comb{m}{n}$种选法，其余$n-m$个球可以任意落入其余的$N-1$个盒中，共有$(N-1)^{n-m}$种落法，因此事件$C$共含$\comb{m}{n}(N-1)^{n-m}$个不同的基本事件，故所求概率为
\[\Pr(C)=\frac{\comb{m}{n}(N-1)^{n-m}}{N^n}\]
\end{enumerate}
\end{solution}

\begin{ex}
\begin{enumerate}
\item 有人说，先后抛掷两枚硬币，只有“两枚都是正面”，“两枚都是反面”，“一枚正面，一枚反面”这3种结果，因此，“两枚都出现正面”这一事件的概率是1/3，这种说法错在哪里？
\item 有100张已编号的卡片（从1号到100号），从中任取1张，计算：
\begin{enumerate}[(1)]
    \item 卡片号是偶数的概率；    \item 卡片号是3的倍数的概率.
\end{enumerate}
\item 掷骰子一枚
\begin{enumerate}[(1)]
\item 掷一次，出现奇数的概率是多少？出现偶数的概率是多少？
\item 抛掷2次，得数字之和为8的概率是多少？
\end{enumerate}
\item 在7张数的卡片中，有4张正数卡片和3张负数卡片．从中任取2张作乘法练习，其积为正数的概率是多少？其积为负数的概率是多少？
\item 某种产品90件，其中甲等品40件，乙等品30件，丙等品20件，在运送这些产品的路上损坏了3件，如果每件产品被
损坏的可能性相同，计算这三等产品中恰好各损坏一件的概率.
\item 在80件产品中，有50件一等品，20件二等品，10件三等品．从中任取3件，计算
\begin{enumerate}[(1)]
\item 三件都是一等品的概率；    \item 2件是一等品，1件是二等品的概率；    \item 一等品，二等品三等品各有一件的概率.
\end{enumerate}

\item 六本不同的书，三本封面是红色，三本封面是蓝色，将它们任意排列在书架的同一层上，问六本书封面颜色相间的概率是多少？
\item 在一副扑克牌（52张）中，有“黑桃，红心，梅花，方块”这四种花色的牌各13张，从中任取4张，这4张牌的花色相同的概率是多少？这四张牌的花色各不相同的概率是多少？
\item 有五根细木棍，它们的长度分别为1, 3, 5, 7, 9厘米，从中任取三根，它们能搭成一个三角形的概率是多少？
\end{enumerate}
\end{ex}

\subsection{互斥事件与加法定理}
在10个乒乓球中，有7个一等品，2个二等品，1个三等品，我们把从中任取一个，取出一等品叫做事件$A$，取得二等品叫做事件$B$，取出三等品叫做事件$C$. 我们看到，如果取出的乒乓球是一等品，即事件$A$发生，那么事件$B$就不发生. 如果取出的是二等品，那么事件$B$发生，那么事件$A$就不发生，也就是说，事件$A$、$B$不可能同时发生. 这种不可能同时发生的两个事件叫做互斥事件. 同理，事件$B$，$C$是互斥事件，事件$A$，$C$是互斥事件. 换句话说，事件$A$，$B$，$C$中，任何两个都是互斥事件。这时我们说事件$A$，$B$，$C$彼此互斥，一般地，如果事件$A_1,A_2,\ldots,A_n$中任何两个都是互斥事件，那么就说事件$A_1,A_2,\ldots,A_n$彼此互斥.

在上面的问题里，因为是任取一个，共有10种等可能的取法，其中得到一等品，二等品，三等品的取法分别有7种，
2种，1种，因此，$\Pr(A)=\frac{7}{10}$, $\Pr(B)=\frac{2}{10}$, $\Pr(C)=\frac{1}{10}$.

现在问：“任取一个乒乓球，取出一等品或二等品”这一事件的概率是多少？这一事件，我们记作“$A+B$”. 因为不论取出一等品还是二等品，都表示这个事件，发生而得到一等品或二等品的取法共有$7+2$种，所以取出一等品或二等品的概率：
\[\Pr(A+B)=\frac{7+2}{10}\]

由$\frac{7+2}{10}=\frac{7}{10}+\frac{2}{10}$，我们看到
\[\Pr(A+B)=\Pr(A)+\Pr(B)\]
它告诉我们：如果事件$A$，$B$互斥，那么事件“$A+B$”发生（即$A,B$中有一个发生）的概率，等于事件$A$，$B$分别发生的概率的和.

一般地，如果事件$A_1,A_2,\ldots,A_n$彼此互斥，那么事件“$A_1+A_2+\cdots+A_n$”发生（即$A_1,A_2,\ldots,A_n$中有一个事件发生，也叫和事件）的概率，等于这$n$个事件分别发生的概率的和，即
\[\Pr(A_1+A_2+\cdots+A_n)=\Pr(A_1)+\Pr(A_2)+\cdots+\Pr(A_n)\]
这叫做互斥事件的加法定理，也就是说互斥事件的和的概率等于各事件概率之和。

\begin{example}
    某地区的年降水量，在100—150毫米范围内的概率是0.12，在150-200毫米范围内的概率是0.25，在200-250毫米范围内的概率是0.16，在250-300毫米范围内的概率是0.14．计算年降水量在100-200毫米范围内的概率与在150-300毫米范围内的概率．
\end{example}

\begin{solution}
    我们把这个地区的年降水量在100-150毫米，150-200毫米，200-250毫米，250-300毫米范围内分别叫做事件$A$，$B$，$C$，$D$. 很明显，这四个事件彼此互斥，根据公式，年降水量在100-200毫米范围内的概率是
\[\Pr(A+B)=\Pr(A)+\Pr(B)=0.12+0.25=0.37\]
年降水量在150-300毫米范围内的概率是
\[\Pr(B+C+D)=\Pr(B)+\Pr(C)+\Pr(D)=0.25+0.16+0.14=0.55\]
\end{solution}


\begin{example}
    在20件产品中，有15件一级品，5件二级品，从中任取3件，其中至少有1件为二级品的概率是多少？
\end{example}

\begin{solution}
    我们把从20件产品中任取3件，其中恰有1件二级品，叫做事件$A_1$；恰有2件二级品，叫做事件$A_2$；3件全是二级品，叫做事件$A_3$，这样，事件$A_1$，$A_2$，$A_3$的概率分别是
\[\begin{split}
    \Pr(A_1)&=\frac{\comb{1}{5}\cdot \comb{2}{15}}{\comb{3}{20}}=\frac{105}{228}\\
    \Pr(A_2)&=\frac{\comb{2}{5}\cdot \comb{1}{15}}{\comb{3}{20}}=\frac{30}{228}\\
    \Pr(A_3)&=\frac{\comb{3}{5}}{\comb{3}{20}}=\frac{2}{228}\\
\end{split}\]

很明显，事件$A_1$，$A_2$，$A_3$彼此互斥，根据公式，3件产品中至少有1件为二级品的概率是
\[\Pr(A_1+A_2+A_3)=\Pr(A_1)+\Pr(A_2)+\Pr(A_3)=\frac{105}{228}+\frac{30}{228}+\frac{2}{228}=\frac{137}{228}\]

从20件产品中任取3件，或者都是一级品，或者不都是一级品（即其中至少有一件是二级品），这两个互斥事件必有一个发生，这种其中必有一个发生的两个互斥事件，叫做对立事件. 一个事件$A$的对立事件通常记作$\overline{A}$，根据对立事件的意义，$A+\overline{A}$是一个必然事件，它的概率等于1．又由于$A$与$\overline{A}$互斥，我们得到
\[\Pr(A)+\Pr(\overline{A})=\Pr(A+\overline{A})=1\]
这就是说，两个对立事件的概率的和等于1．

从上面公式还可得到$\Pr(\overline{A})=1-\Pr(A)$. 
运用这个公式计算事件的概率，有时比较简单．如例3.15还可以这样来解：

从20件产品中任取3件，3件全是一级品（记作事件$A$）的概率：
\[\Pr(A)=\frac{\comb{3}{15}}{\comb{3}{20}}=\frac{91}{228}\]
由于“任取3件，至少有一件为二级品”是事件$A$的对立事件$\overline{A}$，因此，
\[\Pr(\overline{A})=1-\Pr(A)=1-\frac{91}{228}=\frac{137}{228}\]

读者可想想，什么时候用事件的对立事件的概率来计算就比较简单.
\end{solution}

\begin{example}
一个匣子中有红球5个，白球2个，从匣中随便取出两个球，而球的颜色恰好一样的概率是多少？
\end{example}

\begin{solution}
\textbf{解1：} 令$A=$两球颜色一样，$B_1=$两球颜色均为红色，$B_2=$两球颜色均为白色，可见$B_1$和$B_2$互斥，并且$A=B_1+B_2$，因此
\[\Pr(A)=\Pr(B_1)+\Pr(B_2)\]

今从7个球中随便取出两个球，共有$\comb{2}{7}=\frac{7\x 6}{2}=21$种不同的取法，即简单事件总数为21．$B_1$为两球都为红色，即此两球由红球中取得，所以共有$\comb{2}{5}=\frac{5\x4}{2}=10$个简单事件，
因此$\Pr(B_1)=\frac{\comb{2}{5}}{\comb{2}{7}}=\frac{10}{21}$

$B_2$为两球均为白色，即此两球由白球中取得，所以共有
$\comb{2}{2}=1$个简单事件，于是$\Pr(B_2)=\frac{\comb{2}{2}}{\comb{2}{7}}=\frac{1}{21}$，
因此
\[\Pr(A)=\Pr(B_1)+\Pr(B_2)=\frac{10}{21}+\frac{1}{21}=\frac{11}{21}\]

\textbf{解2：}令$A=$“两球颜色一样”，则$\overline{A}=$“两球颜色不一样”，所以
\[\Pr(\overline{A})=\frac{\comb{1}{5}\cdot \comb{1}{2}}{\comb{2}{7}}=\frac{5\x 2}{21}=\frac{10}{21}\]
因此
\[\Pr(A)=1-\Pr(\overline{A})=1-\frac{10}{21}=\frac{11}{21}\]
\end{solution}

一般地说，如果匣中有$m$个红球，$n$个白球，如果从匣中随便抽取两个球，两球颜色相同的概率$P$可以用完全一样的方法，求得其值为
\[P=\frac{\comb{2}{m}}{\comb{2}{m+n}}+\frac{\comb{2}{n}}{\comb{2}{m+n}}=\frac{m(m-1)+n(n-1)}{(m+n)(m+n-1)}=1-\frac{\comb{1}{m}\cdot \comb{1}{n}}{\comb{2}{m+n}}\]

更一般地，如果匣中有$k$种球，第$i$种球共有$n_i$个，从匣中随便取出两个球，这两个球的颜色是一样的概率是多少呢？这也可用类似的方法求出，读者自己去算一下.

\begin{example}
    从52张扑克牌中随便抽取三张．问：
\begin{itemize}
\item $A=$三张花色相同，
\item $B=$三张是“顺子”（即点数相连），
\item $C=$三张是同花顺子（即花色相同，且点数相连）
\end{itemize}
各自出现的概率是多少？
\end{example}

\begin{solution}
从52张牌中随机抽取三张牌，全部基本事件的总个数是
\[\comb{3}{52}=\frac{52\x 51\x 50}{1\x 2\x 3}=52\x 17\x 25\]
注意到花色有四种，令
\[\begin{split}
    A_1=\text{三张都是梅花},&\qquad 
A_2=\text{三张都是方块}\\
A_3=\text{三张都是红心},&\qquad 
A_4=\text{三张都是黑桃}
\end{split}\]
则$A_1,A_2,A_3,A_4$是互斥事件，且$A=A_1+A_2+A_3+A_4$. 又
\[\begin{split}
    \Pr(A_1)&=\Pr(A_2)=\Pr(A_3)=\Pr(A_4)\\
    &=\frac{\comb{3}{13}}{\comb{3}{52}}=\frac{13\x 12\x 11}{52\x 51\x 50}\\
    &=\frac{11}{17\x 50}=\frac{11}{850}
\end{split} \]
因此，$\Pr(A)=4\Pr(A_1)=\frac{22}{425}$.

三张是顺子，从顺子的点数相连这一要求来看，不计花色，共有“$2,3,4$”，“$3,4,5$”，……“$J,Q,K$”，“$Q,K,A$”这十一种不同的情况，依次用$B_1,B_2,\ldots,B_{11}$来表示它们，于是$B_1,B_2,\ldots,B_{11}$互斥，而且是等概率的，$B=B_1+\cdots+B_{11}$，于是
\[\Pr(B)=\Pr(B_1)+\cdots +\Pr(B_{11})=11\Pr(B_1)\]

现在来求$\Pr(B_1)$的值．出现“$2,3,4$”这样的顺子，共有$(\comb{1}{4})^3$这么多种，这是因为这时不考虑花色，同点数的牌每样又有四张，所以出现“$2,3,4$”必须在“2”的四张中取出一张，“3”的四张中取出一张，“4”的四张中取一张，因此，共有$(\comb{1}{4})^3=4^3$这么多种，于是
\[\begin{split}
    \Pr(B_1)&=\frac{4^3}{\comb{3}{52}}=\frac{4^3}{52\x 17\x 25}=\frac{16}{13\x 17\x 25}\\
    \Pr(B)&=\frac{11\x 16}{13\x 17\x 25}=\frac{176}{5525}=\frac{8}{13}\Pr(A)
\end{split}\]
显然$\Pr(B)<\Pr(A)$.

再来计算事件$C$的概率，很明显，从上面的讨论中可以看出来，“$2,3,4$”同花的只可能有四种不同的情形．因此
\[\Pr(C)=\frac{11\x 4}{\comb{3}{52}}=\frac{11}{13\x17\x25} =\frac{1}{16}\Pr(B)\]
显然$\Pr(C)<\Pr(B)$

\end{solution}

\begin{example}
   掷四枚硬币，用“1”表示正面向上，用“0”表示反面向上，于是记$x_i$为第$i$枚硬币的结果，每掷一次，$(x_1,x_2,x_3,x_4)$就表示各枚硬币正、反面的出现情况，例如，
\begin{center}
   （1, 1, 0, 0）就是（正，正，反，反）\\
   （1, 0, 0, 1）就是（正，反，反，正）    
\end{center} 

   求下列各事件的概率：
\[\begin{split}
   A&=\text{恰有两个正面向上}\\
   B&=\text{至少有两个正面向上}\\
   C&=\text{第一枚、第二枚硬币的结果正好相反，第三枚、第四枚硬币的结果正好相反}    
\end{split}\]
\end{example}

\begin{solution}
用$x_1,x_2,x_3,x_4$来表示，就有
\[\begin{split}
   A&=\{x_1+x_2+x_3+x_4=2\}\\
   B&=\{2\le x_1+x_2+x_3+x_4\le 4\}\\
   C&=\{x_1=1-x_2,\; x_3=1-x_4\}    
\end{split}\]

   全部简单事件的个数是一个有重复的排列数，它是$2^4=16$.

\begin{enumerate}[(1)]
    \item   $x_1+x_2+x_3+x_4=2$表示四个数中只有两个是“1"，其余两个一定是“0”．而且只要选定了“1”的位置，则“0”的位置也就随之而定了，从四个位置中任选两个，则共有$\comb{2}{4}$种，即6种，因此
   \[\Pr(A)=\frac{6}{16}=\frac{3}{8}\]
   \item $x_1+x_2+x_3+x_4\ge 2$，这一情形可以分成下面三种等式的情形：
\[\begin{split}
    x_1+x_2+x_3+x_4=2&\qquad (\text{有$\comb{2}{4}=6$种情形})\\
    x_1+x_2+x_3+x_4=3&\qquad (\text{有$\comb{3}{4}=4$种情形})\\
    x_1+x_2+x_3+x_4=4&\qquad (\text{有$\comb{4}{4}=1$种情形})\\
\end{split}\]
因此，
\[\Pr(B)=\frac{6+4+1}{16}=\frac{11}{16}\]
\item 此时$x_1=1-x_2$, $x_3=1-x_4$，于是只要看$x_1,x_3$联合起来有多少种不同的情形. 从重复排列的计算公式可以知道$x_1$，$x_3$联合起来共有$2^2=4$种不同情况. 因此，可得
\[\Pr(C)=\frac{4}{16}=\frac{1}{4}\]

\end{enumerate}
   
\end{solution}

\begin{ex}
\begin{enumerate}
    \item 判别下列每对事件是不是互斥事件，如果是，再判别它们是不是对立事件.
从一堆产品中任取2件，其中：
\begin{enumerate}[(1)]
\item 恰有1件次品和恰有2件次品；
\item 有次品和全是次品；
\item 有正品和有次品；
\item 有次品和全是正品．
\end{enumerate}

\item 从一批乒乓球产品中任取一个，如果其重量小于2.45克的概率是0.22，重量不小于2.50克的概率是0.20，那么重量在2.45—2.50克范围内概率是多少？
\item 某射手在一次射击中射中10环，9环，8环的概率分别为0.24, 0.28, 0.19．计算这个射手在一次射击中：
\begin{enumerate}[(1)]
    \item 射中10环或9环的概率；
    \item 不够8环的概率．
\end{enumerate}

\item 一个匣子内有9张票，其号数分别为$1,2,\ldots,9$．从中任取2张，其号数至少有1个为奇数的概率是多少？
\item 在50件产品中有45件合格品，5件次品，从中任取3件，计算其中有次品的概率.
\item 从52张中随便抽取五张．问
\begin{itemize}
\item $A=$五张花色相同；
\item $B=$五张是“顺子”；
\item $C=$五张同花顺子；
\item $D=$五张中有四张点数相同（J作11点，Q作12点，K作13点）．
\end{itemize}
$A$，$B$，$C$，$D$各自出现的概率是多少？
\item 掷$n$枚硬币，（$n$为偶数）
\begin{itemize}
\item $A=$恰有两个正面向上，
\item $B=$至少有两个正面向上，
\item $C=$第一枚、第二枚硬币的结果正好相反，第三枚、第
四枚的结果正好相反，……第$n-1$枚与第$n$枚的结果正
好相反.
\end{itemize}
求$A$，$B$，$C$各事件的概率.
\item 同时掷三颗骰子，得3点至10点的概率是多少？得11点至18点的概率是多少？
\item 掷三颗骰子，得三颗的点数相同的概率是多少？
\item 从52张扑克中任取五张，问得到三张点数相同，其他两张点数也相同的概率.
\end{enumerate}
\end{ex}

\subsection{事件与集合的对应}
我们现在讨论的一些问题，都离不开基本事件. 所以我们可以把基本事件看作衡量随机事件的“最小单位”，它好像是几何中的“点”，集合中的“元素”.而一切随机事件都是由基本事件组成的，这样，事件之间的各种运算，就和集合之间的各种运算完全对应起来，这样既使我们掌握了事件的运算，还使我们进一步理解了集合运算的意义.

在讨论集合运算时，集合中的一个元素称为一个点，用$a$来表示，则$a$和$\{a\}$是不一样的. 前者是一个点，后者是由$a$这一个点组成的单点集，它是一个集合.

当我们用$E_1,\ldots,E_n$表示基本事件时，它相当于集合中的点，随机事件$A$总是由某些基本事件合起来组成的.如
果把一个基本事件$E_i$也看成是随机事件，那就是指由这个$E_i$组成的集合$\{E_i\}$，这个随机事件$\{E_i\}$相当于一个点组成的单点集. 这一点一定要加以注意，不要一开始就发生混淆.

下面将事件间的运算和事件间的关系的各种定义和集合的运算，以及集合的关系对照着列成一张表，以更进一步看清它们之间的联系，也是为了可以用集合语言来理解概率论.

\begin{longtable}{p{.45\textwidth}p{.45\textwidth}}
\hline
事件& 集合\\
\hline
1. 基本事件  &1. 元素（点）\\
$E_1,E_2,\ldots, E_n$ &$a_1,a_2,\ldots,a_n$\\
\hline
2. 必然事件 & 2. 全集$A$\\
$U=\{ E_1,E_2,\ldots, E_n\}$ & $A=\{a_1,a_2,\ldots,a_n\}$\\
\hline
3. 不可能事件$V$  &  3. 空集$\emptyset$\\
$V$中不含任何$E_i$& $\emptyset$中不含任何$a_i$\\
\hline
4. 随机事件$F$ &  4. 子集合$S$\\
$F=\{E_{i_1},E_{i_2},\ldots,E_{i_m}\}$&$S=\{a_{i_1},a_{i_2},\ldots,a_{i_m}\}$\\
$1\le i_1< i_2<\cdots <{i_m}\le n$&$1\le i_1<i_2<\cdots<i_m\le n$\\
\hline
5. 事件$F_1$包含$F_2$写成$F_1\supset F_2$或$F_2\subset F_1$  &  5. 集合$S_1$包含$S_2$，写成$S_1\supset S_2$或$S_2\subset S_1$\\
即$F_2$中的基本事件均为$F_1$中的基本事件。或$F_2$发生时，$F_1$一定发生。 &即$S_2$中的点均为$S_1$中的点\\
\hline
6. $F_1=\{E_{i1},E_{i2},\ldots,E_{ik}\}$，&   6. $S_1=\{a_{i1},a_{i2},\ldots,a_{ik}\}$，\\
$F_2=\{E_{j1},E_{j2},\ldots,E_{jl}\}$，& $S_2=\{a_{j1},a_{j2},\ldots,a_{jl}\}$，\\
\small $F_1+F_2=\{E_{i1},\ldots,E_{ik}, E_{j1},\ldots,E_{jl}\}$  &$S_1+S_2=\{a_{i1},\ldots,a_{ik},a_{j1},\ldots,a_{jl}\}$\\
则$F_1+F_2$称为$F_1$与$F_2$之和，即$F_1, F_2$中至少有一个事件发生。这个运算为加法运算.
& 这是集合的加法，即把$S_1$和$S_2$内的点合并在一起构成的新集合。\\
\hline
7. 由事件$F_1$和$F_2$中的公共基本事件组成的事件称为$F_1$与$F_2$之积，写成$F_1F_2$. 即“$F_1,F_2$这两个事件都发生”这一事件。这个运算称为乘法运算。&7. 集合$S_1$与$S_2$的公共元素组成的集合，称为$S_1$和$S_2$之积，写成$S_1S_2$。即$S_1$与$S_2$的公共部分。这是集合的乘法运算。\\
\hline
8. $F_1F_2=V$，则称$F_1,F_2$是互斥的，即$F_1$和$F_2$中没有公共的基本事件，即$F_1$与$F_2$同时发生这一事件是不可能的。&
8. $S_1S_2=\emptyset$，则称$S_1,S_2$是互不相交的，即$S_1$与$S_2$没有公共的元素
\\
\hline
9. 逆事件（对立事件）&9. 余集（补集）\\
如果$F=\{E_{i1},\ldots,E_{ik}\}$，则不在$F$中的那些基本事件全体组成的随机事件称为$F$的逆事件，用$\overline{F}$表示。$\overline{F}$就是“$F$不发生”这一事件。
& 
如果$S=\{a_{i1},\ldots,a_{ik}\}$，则由不在$S$中的那些点的全体组成的集，称为$S$的余集，用$\overline{S}$表示。也即在全集$A$中，但不在$S$中的那些点组成的集合。\\
\hline
10. 设$F=\{E_{i1},\ldots,E_{ik}\}$，则由加法的定义可知$F=\{E_{i1}\}+\cdots +\{E_{ik}\}$ & 10. 设$S=\{a_{i1},\ldots,a_{ik}\}$，则由加法的定义可知$S=\{a_{i1}\}+\cdots+\{a_{ik}\}$\\
\hline
\end{longtable}

由这张对照表，可以由集合的关系式得到事件之间的关
系式，符号的含义和上表相同
\begin{enumerate}
\item $V\subset F\subset U$ (相当于$\emptyset\subset S\subset A$) 

不可能事件$\subset$随机事件$\subset$必然事件
\item 若$F_1\subset F_2$, $F_2\subset F_3$, 则$F_1\subset F_3$
(相当于$S_1\subset S_2,\; S_2\subset S_3$时，必有$S_1\subset S_3$)
\item $F\overline {F}= V$, $F+ \overline {F}= U$ (相当于$S\overline{S} = \emptyset$, $S+ \overline{S} = A$) 

注意：此即$F$与$\overline F$之和必然发生，$F$与$\overline{F}$ 之积是不可能事
件。$\overline{F}$是$F$的逆事件，$F$是$\overline{F}$的逆事件，即彼此互为逆事件。
\item 当$F_1\subset F_2$时，$F_1$中所含的简单事件均为$F_2$中的简单事件，由古典概率公式得
$\Pr\left(F_{1}\right)\le  \Pr\left(F_{2}\right)$
\item 由于$F\overline{F}= V$, $F+ \overline {F}= U$, $F$与$\overline{F}$互斥，从加法定理
得到
$$1=\Pr(U)=\Pr(F)+\Pr(\overline{F})$$
因此，
$\Pr(\overline{F})=1-\Pr\left(F\right)$
\end{enumerate}

\section*{习题3.1}
\begin{enumerate}
\item 设有$n$个房间，分给$n$个不同的人，每个人都以同等可能进入每一房间，而且每间房里的人数没有限制，试求不出现空房的概率？
\item 将$n$个可辨的球投入$N$个盒中，每个球都以同等可能落入每一个盒内，求指定某盒是空的概率.
\item 一批产品共$n$件，其中有$k$件废品，求任意抽出的$m$件产品中恰有$\ell$件废品的概率？
\item 某人有$n$把钥匙，其中只有一把能开他的门，他逐把地取钥匙试开，这试开的手续可能需要$1,2,\ldots,n$次试验，试证明在第$i$次，第$j$次打开锁的概率是一样.
\item 袋中装有10个白球5个黑球，第一次从袋中取出一球（不放回），发现是白球，求第二次从袋中取得白球的概率？并比较两次取得白球概率的大小. 如果第一次取得黑球，求第二次取得白球的概率？
\item 在52张扑克牌中任意取13张，问其中至少有一张是10点的概率是多少？
\item 一盒中有5个白球，7个黑球．
\begin{enumerate}[(1)]
    \item 从中任取一球，问取得是白球和取得是黑球的概率各为多少？
    \item 从中任取两球，问取得都是白球的概率、取得一白一黑的概率各为多少？
\end{enumerate}

\item 箱子所装球的总重量为$N$公斤，而每个球为$n$公斤，其中有$m$个是黑球. 求从箱子中任取一球恰巧是黑球的概率.
\item 一表面为红色的正方体被分割成1000个同样大小的小正方体.试求从中任取一小正方体其两面涂有红色的概率.
\item 从写有$a$，$b$，$c$，$d$，$e$的五张卡片中任取两张，求这两张卡片的字母恰好是按字母顺序相邻排着的概率.
\item 在相应地写有2, 4, 6, 7, 8, 11, 12和13数字的八张卡片中任意取两张，求由所取得的两个数字构成的分数为可约的概率.
\item 从1, 2, 3, 4, 5诸数中任取三数，排成三位数，问如此所得三位数是偶数的概率是多少？
\item 求任意取一整数$N$的：（1）平方；（2）四次方；（3）乘任意一整数后，其尾数为1的概率。
（提示：只需讨论个位数.）
\item 在书架中任意放着10本书，求某给定的三本书放在一起的概率.
\item 4个女孩，3个男孩排成一排，求男女相间的概率.
\item 一个匣子里有90只好的、10只次的螺丝钉．如果从中任意取用10只，恰都是好的螺丝钉的概率是多少？
\item 在一个坛子中放有$n_1$个白球，$n_2$个黑球，逐一全部取出，问第一个和最后一个球都为白球的概率是多少？

\item 求出下述两事件的概率：
\begin{enumerate}[(1)]
    \item 10个人的生日在10个不同的月份；
\item 6个人的生日恰巧在两个月中。（提示：“恰巧在两个月中”不包括“集中在某一个月”的情形. 可先考虑6个人的生日在某指定的两个月里有几种可能．再考虑生日集中在任意两个月里有几种可能.）
\end{enumerate}

\item 电话号码由五个数字组成．每个数字可以是$0,1,2,\ldots,9$中的任意一个数，求电话号码是由完全不相同的数字组成的概率.
\item 从一副扑克（52张）中任取四张，求四张牌的花色各不相同的概率.
\item 为了减少比赛场次，把20个球队分成两组（每组10队）进行比赛，求最强的两队被分在不同组内的概率．
\item 掷三枚骰子，得5点以上（包括5点）的概率是多少？
\item 某人的口袋装有1枚五分的硬币，2枚贰分的硬币以及3枚壹分的硬币，他从中任取3枚，取出的总钱数不少于伍分的概率是多少？
\item 把10个运动队先均分成两组预赛，求最强两队被分在：
\begin{multicols}{2}
\begin{enumerate}[(1)]
\item 不同组内；    \item 同一组内的概率．
\end{enumerate}
\end{multicols}

\item 一座楼房有四个大门，两人住在此楼内，问恰好住在同一大门内的概率是多少？
\item 设在一批有1000件一等品，50件二等品的产品中进行质量检查，如果在该批产品中任意抽查10件，结果全是一等品，问在未被抽查的产品中再任意地连续取两件，至少
有一件是一等品的概率是多少？
\item 用火车运载两个工厂生产的同类产品，其中甲厂30件，乙厂20件，有消息证实，在路途中有两件产品损坏，求损坏的是不同厂的产品的概率.
\end{enumerate}

\section{条件概率与独立事件}
\starsubsection{条件概率}
若一袋内有7个白球，5个红球，设事件$A=$“取出一个白球”，事件$B=$“取出两个白球”，显然
\[\Pr(A)=\frac{\comb{1}{7}}{\comb{1}{12}}=\frac{7}{12},\qquad \Pr(B)=\frac{\comb{2}{7}}{\comb{2}{12}}=\frac{7}{22}\]

如果我们取出一个白球后，不再放回，那么这时再取出两个白球（注意，实际上是取出了三个白球，但取法是先一个，再两个地取，而不是一起取出）的概率
\[\Pr(B')=\frac{\comb{2}{6}}{\comb{2}{11}}=\frac{3}{11}\]
同样的，如果我们取出两个白球后，不再放回，那么这时再取出一个白球的概率
\[\Pr(A')=\frac{\comb{1}{5}}{\comb{1}{10}}=\frac{1}{2}\]

这个例子告诉我们，我们所计算的某个事件出现的概率是根据另一事件的出现这一事实来确定的.

我们把“出现事件$B$的条件下，出现事件$A$”的概率称为事件$A$关于事件$B$的条件概率，记作$\Pr(A|B)$. 对于上面例子，则
\[\Pr(A')=\Pr(A|B)=\frac{1}{2},\qquad \Pr(B')=\Pr(B|A)=\frac{3}{11}\]

另一方面，考虑事件$AB=$“取出三个白球”，则
\[\begin{split}
    \Pr(AB)=\frac{\comb{3}{7}}{\comb{3}{12}}=\frac{7}{44}&=\frac{7}{22}\x\frac{1}{2}=\Pr(B)\Pr(A|B)\\
    &=\frac{7}{12}\x\frac{3}{11}=\Pr(A)\Pr(B|A)\\
\end{split} \]

就是说$A$与$B$事件同时出现的概率，可以通过$B$出现的概率$\Pr(B)$与$A$关于$B$的条件概率$\Pr(A|B)$来表示.

一般地，对任两个随机事件$A$与$B$，若$\Pr(B)>0$，则有：
\begin{equation}
  \Pr(AB)=\Pr(A|B)\cdot \Pr(B)  \qquad \Pr(AB)=\Pr(B|A)\cdot \Pr(A)\tag{1}
\end{equation}
即
\begin{equation}
    \Pr(A|B)=\frac{\Pr(AB)}{\Pr(B)}\qquad \Pr(B|A)=\frac{\Pr(AB)}{\Pr(A)}\tag{2}
\end{equation}

\begin{proof}
    设基本事件的总数为$e$，其中事件$A$所包含的基本事件有$m$个，事件$B$所包含的基本事件有$n$个，因为我们并没有假定事件$A$和$B$互斥，因此一般都存在着既属于事件$A$，
也属于事件$B$的基本事件，（即事件$A\cdot B$所包含的基本事件）设这样的事件共有$r$个（如图）于是：
\[\Pr(B)=\frac{n}{e},\qquad \Pr(AB)=\frac{r}{e}\]

\begin{center}
    \begin{tikzpicture}
\draw(0,0) rectangle (5,3.5);
\node at (.5,3.25){$e$};
\draw(2,2)node{$A$} circle(1);
\draw(3.5,2)node{$B$} circle(1);
\draw(1.75,1.75)--(1.5,.6)node[below]{$m$};
\draw(3.75,1.75)--(4,.6)node[below]{$n$};
\draw(2.8,1.75)--(2.7,.6)node[below]{$r$};

    \end{tikzpicture}
\end{center}

再看条件概率$\Pr(A|B)$，因为这是在事件$B$发生的条
件下考虑问题的，这时总的基本事件数就是事件$B$所包含的个数$n$，在$B$发生的条件下$A$再发生所包含的基本事件，就必然是属于事件$AB$的$r$个基本事件.

故$\Pr(A|B)=\frac{r}{n}$

这样就有：
\[\Pr(A|B)\cdot\Pr(B)=\frac{r}{n}\cdot \frac{n}{e}=\frac{r}{e}=\Pr(AB)\]
同样可以证明$\Pr(B|A)\cdot\Pr(A)=\Pr(AB)$.
\end{proof}

\begin{example}
    某人提出一个问题，甲先答，答对的概率是0.4；如果答错，由乙答，答对的概率是0.5，求问题由乙解出的概率．
\end{example}

\begin{solution}
“问题由乙解出”相当于“甲答错”（$A$）与“乙答对”（$B$）两事件一起发生。（即事件$A$，$B$同时发生）.

由题意：“甲答错”的概率：$\Pr(A)=1-0.4=0.6$；而“甲答错的条件下乙答对”的概率$\Pr(B|A)=0.5$，
所以“问题由乙解出”的概率：
\[\Pr(AB)=\Pr(B|A)\cdot \Pr(A)=0.5\x0.6=0.3\]
\end{solution}

\begin{example}
    一只袋中有2只白球和3只红球，从袋中取出一只球，然后在第一只不放回的前提下取出第二只球，那么所取出两只球都是红球的概率是多少？
\end{example}

\begin{solution}
    “取出的两只球都是红球”必定是“第一次取得的是红球”（$A$），且“第二次取得的是红球”（$B$）两事件一起发生（即事件$AB$发生）

    “第一次取得红球”的概率：$\Pr(A)=\frac{3}{5}$，而“第一次取得红球的条件下，第二次取得红球”的概率$\Pr(B|A)=\frac{2}{4}=\frac{1}{2}$，所以“两个球都是红球”的概率：
\[\Pr(AB)=\Pr(B|A)\cdot \Pr(A)=\frac{1}{2}\cdot \frac{3}{5}=\frac{3}{10}\]
\end{solution}

\begin{example}
    某篮球队在己方半场抢得篮板球的概率为0.75而抢得篮板球且反攻投中得分的概率为0.6，问该队已抢得篮板球在手，问这球反攻投中得分的概率有多大？
\end{example}

\begin{solution}
    设事件$A=$“抢得篮板球”；事件$B=$“反攻投中得分”
按题意：$\Pr(A)=0.75$，“抢得篮板球且反攻投中得分”的概率（即积事件$A\cdot B$的概率）：$\Pr(AB)=0.6$

所以“抢得篮板球在手，这球反攻投中得分”的概率。
\[\Pr(B|A)=\frac{\Pr(AB)}{\Pr(A)}=\frac{0.6}{0.75}=0.8\]
\end{solution}

\begin{ex}
\begin{enumerate}
    \item 计算在一分钟的区间内进入某邮局的人数的概率如下表：
\begin{center}
\begin{tabular}{cccccc}
\hline
    1分钟内进入的人数&0人&1人&2人&3人&大于3人\\
\hline
概率&0.05&0.15&0.22&0.22&0.36\\
\hline
\end{tabular}
\end{center}    
考虑下述事件：$A=$至少有一个人到达，$B=$至少有二个人到达，$C=$3个顾客到达.

求$\Pr(C|A)$和$\Pr(C|B)$
\item 一个绿骰子和一个红骰子被转动，令$x$表示在红骰子面上出现点子的数目，$y$表示绿骰子面上出现点子的数目.

令：
\begin{multicols}{2}
    $A=$“$x+y>7$的事件”；
    \\
    $B=$“$x+y<6$的事件”；
    \\
    $C=$“$x+y=7$的事件”；\\$D=$“$y=4$的事件”；\\
$E=$“$y=2$的事件”；\\$F=$“$x$或$y$是5”；\\$G=$“$x$和$y$都是5”；\\$H=$“$x+y=10$”
\end{multicols}
求：
\begin{multicols}{3}
 \begin{enumerate}[(1)]
    \item $\Pr(A|D)$
    \item $\Pr(B|E)$
    \item $\Pr(C|D)$
    \item $\Pr(G|F)$
    \item $\Pr(G|H)$
    \item $\Pr(F|H)$
\end{enumerate}   
\end{multicols}
\end{enumerate}
\end{ex}

\subsection{独立事件与乘法定理}
先看下面的问题：“餐桌上摆着10把外形一样的餐刀，其中7把是不锈钢的；3把是镀铬的，记事件$A=$“甲从中任取一把恰为镀铬的”；事件$B=$“乙从中任取一把恰为镀铬的”.

显然，当只有甲单独取用时：$\Pr(A)=\frac{3}{10}$；当只有乙单独取用时：$\Pr(B)=\frac{3}{10}$.

下面再考虑乙先取用的两种情况：
\begin{enumerate}[(1)]
\item 当乙先取用且取得的是镀铬的，用后归还，这相当于返回抽样.

由于这时甲仍是从10把餐刀中取用，其中镀铬的也仍是3把，所以
“乙取得的是镀铬的条件下（用后归还），甲取得的是镀铬的”这一事件的概率：
\[\Pr(A|B)=\frac{3}{10}=\Pr(A)\]
根据同样的分析：$\Pr(B|A)=\frac{3}{10}=\Pr(B)$
\item 当乙先取用且取得是镀铬的，尚未归还甲就取用，这相当于不返回抽样.

这时由于事件$B$的发生，只剩下9把餐刀，其中镀铬的剩下2把，所以
“乙取得的是镀铬的条件下（用后未归还），甲取得的是镀铬的”这一事件的概率：
\[\Pr(A|B)=\frac{2}{9}\ne \Pr(A)\]
根据同样的分析：$\Pr(B|A)=\frac{2}{9}\ne \Pr(B)$
\end{enumerate}

总结上面的分析得出：当乙（甲）的取用不影响甲（乙）的取用，则有下面两等式成立
\[\Pr(A|B)=\Pr(A),\qquad \Pr(B|A)=\Pr(B)\]

当乙（甲）的取用影响到甲（乙）的取用，则：
\[\Pr(A|B)\ne \Pr(A),\qquad \Pr(B|A)\ne \Pr(B)\]

我们定义：如果对于事件$A$和$B$来说，有
\[\Pr(A|B)= \Pr(A),\qquad \Pr(B|A)= \Pr(B)\]
同时成立，那么称事件$A$和事件$B$为相互独立事件，或称事件$A$和$B$相互独立，否则，就是相关的.

对本例，事实上用后归还的话，不管“乙取用为镀铬的”这一事件发生与否，对“甲取用为镀铬的”概率是没有影响的，反之也一样。

因此，通俗地说，事件$A$与$B$相互独立，就是：事件$A$发生与否对事件$B$发生的概率没有影响，同样，事件$B$发生与否对事件$A$发生的概率也没有影响，例如：同时抛掷两枚硬币，记$A=$“第一枚正面向上”，$B=$“第二枚正面向上”；因第一枚是否正面向上对第二枚正面向上的概率没有影响，反之也一样，所以事件$A$与$B$相互独立.

容易看出：当事件$A$和$B$相互独立时，则事件$A$和$\overline{B}$；$\overline{A}$和$B$；$\overline{A}$和$\overline{B}$都相互独立的。（读者自己证明）

当事件$A$和$B$相互独立时，上节公式(1)
\[\Pr(AB)=\Pr(A|B)\cdot \Pr(B),\qquad \Pr(AB)=\Pr(B|A)\cdot \Pr(A)\]
变为
\begin{equation}
    \Pr(AB)=\Pr(A)\cdot \Pr(B) \tag{3}
\end{equation}

这就是说：两个相互独立的事件同时发生的概率等于每个事件发生的概率的乘积.

公式(1)和(3)通称为概率的乘法定理。

一般地，若$A_1,A_2,\ldots,A_n$为相互独立事件，则有：
\begin{equation}
    \Pr(A_1\cdot A_2\cdots A_n)=\Pr(A_1)\cdot \Pr(A_2)\cdots \Pr(A_n)\tag{4}
\end{equation}
独立性的概念在概率中是十分重要的，它有广泛的应用. 在实际问题中判断事件是否独立是很重要的.

\begin{example}
    甲乙两人进行一次射击，如果两人射中目标的概率甲为0.6，乙为0.7，计算：
\begin{enumerate}[(1)]
\item 两人都击中目标的概率；
\item 其中恰有一人击中的概率；
\item 至少有一人击中目标的概率。
\end{enumerate}
\end{example}

\begin{solution}
\begin{enumerate}[(1)]
    \item “两人都击中目标”（$C$），就是“甲击中目标”（$A$），且“乙击中目标”（$B$），故$C=A\cdot B$. 事件$A$和$B$显然相互独立，所以
  \[  \Pr(C)=\Pr(A\cdot B)=\Pr(A)\cdot \Pr(B)=0.6\x0.7=0.42\]
    \item “其中恰有一人击中”（$D$），就是“甲击中，乙未击中”（$A\cdot \overline{B}$）或“甲未击中，乙击中”（$\overline{A}\cdot B$），故$D=A\cdot \overline{B}+\overline{A}\cdot B$，显然事件$A\cdot \overline{B}$与$\overline{A}\cdot B$在两人各射击一次时是不可能同时发生的，即事件$A\cdot \overline{B}$与$\overline{A}\cdot B$互斥，所以
\[\begin{split}
    \Pr(D)&=\Pr(A\cdot \overline{B}+\overline{A}\cdot B)=\Pr(A\cdot \overline{B})+\Pr(\overline{A}\cdot B)\\
    &=0.6\x(1-0.7)+(1-0.6)\x 0.7=0.18+0.28=0.46
\end{split}\]
    \item “至少有一人击中”（$E$），就是“甲击中，乙未击中”（$A\cdot \overline{B}$）或“甲击中，乙击中”（${A}\cdot B$），或“甲未中乙击中”（$\overline{A}\cdot B$），故$E=A\cdot \overline{B}+A\cdot B+\overline{A}\cdot B$，所以
\[\begin{split}
    \Pr(E) &= \Pr(A\cdot \overline{B}+A\cdot B+\overline{A}\cdot B)\\
    &=\Pr(A\cdot \overline{B})+\Pr(A\cdot B)+\Pr(\overline{A}\cdot B)\\
    &=0.6\x0.3+0.6\x0.7+0.4\x0.7\\
&=0.18+0.42+0.28=0.88
\end{split}\]

此小题还可采用下面解法：
“没有一个击中”（$\overline{E}$），就是“甲未击中”（$\overline{A}$）且“乙未击中”（$\overline{B}$），故$\overline{E}=\overline{A}\cdot \overline{B}$. 所以：
$$\Pr(\overline{E})=\Pr(\overline{A}\cdot \overline{B})=\Pr(\overline{A})\cdot\Pr(\overline{B})=0.4\x0.3=0.12$$
因此：$\Pr(E)=1-\Pr(\overline{E})=1-0.12=0.88$
\end{enumerate}
\end{solution}

\begin{example}
   假设每枚地-空导弹击中飞机的概率均为0.8，如要有99\%的把握击中来犯敌机，问需几枚导弹同时发射？ 
\end{example}


\begin{solution}
设需$n$枚地空-导弹同时发射，

记事件$A=$“敌机被击中”，则“敌机未被击中”（$\overline{A}$），就是“第1枚未击中”（$\overline{A}_1$），且“第2枚未击中”($\overline{A}_2$)，……，“第$n$枚未击中”（$\overline{A}_n$）同时发生。

故$\overline{A}=\overline{A}_1\cdot \overline{A}_2\cdots \overline{A}_n$，而$\overline{A}_1, \overline{A}_2,\ldots,\overline{A}_n$是彼此独立的，所以
\[\Pr(\overline{A})=\Pr(\overline{A}_1)\cdot \Pr(\overline{A}_2)\cdots\Pr(\overline{A}_n)=(1-0.8)^n=0.2^n\]
要求击中概率达99\%，即$\Pr(A)=0.99$，故
\[\Pr(\overline{A})=1-\Pr(A)=0.01\]
所以$n$应满足：$0.2^n=0.01$. 取对数
\[\begin{split}
    n\lg0.2&=\lg0.01\\
n&=\frac{\lg0.01}{\lg0.2}=\frac{-2}{-0.699}\approx 3
\end{split} \]
答：3枚导弹同时发射能以99\%的把握击中来犯敌机．    
\end{solution}

\begin{example}
    一个元件能正常工作的概率$r$称为该元件的可靠性，由元件组成的系统能正常工作的概率称为系统的可靠性。今设所用元件的可靠性都为$r\; (0<r<1)$，且每个元件能否正常工作是相互独立的。试求下面两系统的可靠性，并

\begin{enumerate}[(1)]
    \item 附加通路系统，它由两条$n\; (n\ge 2)$个元件串联成的通路并联而成，如下图：
\begin{center}
    \begin{circuitikz}[european]
\draw(0,0)--(.5,0)--(.5,.5)to [R](2.5,.5) to [R] (4.5,.5)--(5,.5);
\draw(6,.5) to [R](7.5,.5);
\draw(.5,0)--(.5,-.5)to [R](2.5,-.5) to [R](4.5,-.5)--(5,-.5);
\draw[dashed](4.5,-.5)--(6,-.5);
\draw[dashed](4.5,.5)--(6,.5);

\draw(6,-.5) to [R](7.5,-.5);
\draw(7.5,-.5)--(7.5,.5);
\draw(7.5,0)--(8.25,0);


\foreach \x/\y in {3/1,7/2,13.5/$n$}
{
    \node at (\x/2,.5){\y};
    \node at (\x/2,-.5){\y};
}

    \end{circuitikz}
\end{center}
    \item 附加元件系统，它由$n$对并联元件串联而成，如下图：
\begin{center}
    \begin{circuitikz}[european]
\draw(0,0)--(.5,0)--(.5,.5) to [R] (2,.5)--(2,-.5) to [R](.5,-.5)--(.5,0);
\draw(2.5,0)--(3,0)--(3,.5) to [R] (4.5,.5)--(4.5,-.5) to [R](3,-.5)--(3,0);
\draw(2,0)--(2.5,0);
\draw(4.5,0)--(5,0);
\draw(8.5,0)--(8,0);
\draw[dashed](6,0)--(5,0);
\draw(6,0)--(6.5,0)--(6.5,.5) to [R] (8,.5)--(8,-.5) to [R](6.5,-.5)--(6.5,0);
\foreach \x/\y in {2.5/1,7.5/2,14.5/$n$}
{
    \node at (\x/2,.5){\y};
    \node at (\x/2,-.5){\y};
}
    \end{circuitikz}
\end{center}
\end{enumerate}
\end{example}

\begin{solution}
    系统(1)有两条通路. 对每条通路来说，必须当每个元件都正常工作时才正常工作，故其可靠性为：$R_c=r^n$

    所以每条通路发生故障的概率为：$1-r^n$.

    系统发生故障，必定是两通路同时发生故障. 故系统发生故障的概率为：$(1-r^n)^2$.

    因此系统(1)的可靠性为：$$R_s=1-(1-r^n)^2=r^n(2-r^n)$$
    所以$R_s>R_c$，即系统在附加了一条通道后，比一条通路增加了可靠性。

    现在考虑系统(2). 从上述讨论易得每对并联元件的可靠性为：$$R'=1-(1-r)^2=r(2-r)$$
    因系统是由各对并联元件串联而成，故其可靠性为：
\[    R_{s'}=(R')^n=r^n(2-r)^n\]
    显然$R_{s'}>R_c$，即用附加元件的方法也能增加系统的可靠性。

两种系统所用的元件一样，其优劣就看可靠性的大小了. 为比较$R_s$与$R_{s'}$的大小，令
$Q_n=\frac{R_{s'}}{R_s}\; (n\ge 2)$，则
\[Q_n=\frac{(2-r)^n}{2-r^n}\]
当$n=2$时:
\[(2-r)^2=2-r^2+2(1-r)^2>2-r^2\]
所以$Q_2>1$

假设当$n=k$时，$Q_k>1$，即$(2-r)^k>2-r^k$，则
\[\begin{split}
 (2-r)^{k+1}>(2-r^k)\cdot (2-r)&=4-2r^k-2r+r^{k+1}\\
&=2-r^{k+1}+2(1-r)(1-r^k)>2-r^{k+1}
\end{split}\]
即当$n=k+1$时也有$Q_{k+1}>1$

所以对于任意大于等于2的自然数$n$，$Q_n>1$总成立，即$R_{s'}>R_s$，所以系统(2)较系统(1)可靠.
\end{solution}

\begin{ex}
\begin{enumerate}
    \item 一张牌随机地从一副扑克牌中抽取，令：    $A=$“这张牌是梅花”；    $B=$“这张牌是个$K$”；    $C=$“这张牌是高级的”（即$J$、$Q$、$K$、$A$）
    
    问：$A,B$；$A,C$和$B,C$的哪一对是独立事件？
    
    \item 掷三个硬币，令：    $A=$“不多于一个正面”；    $B=$“每面至少有一个”。
    
    问：$A$和$B$是独立的吗？
\end{enumerate}
\end{ex}

\section*{习题3.2}
\begin{enumerate}
    \item 一个盒子中装有5只好灯泡，一个调皮的孩子又放进三只坏灯泡，假如八只灯泡里先取出1只（不放回），再取出第2只，求两只都是好灯泡的概率．
    \item 在$1,2,\ldots,100$中随机选出一个数，求满足下列条件的概率：
\begin{enumerate}[(1)]
\item 已知选出的数是偶数的条件下，能被5整除；
\item 已知它是被4整除的条件下它被10整除．
\end{enumerate}

\item 一个口袋内装有2个白球和2个黑球，把“从中任意摸出一个球，得到白球”叫做事件A，把“从剩下的3个球中任意摸出一个球，得到白球”叫做事件B.事件A与事件B是相互独立的吗？再求在先摸出白球后，再摸出 白球的概率，及先摸出黑球后，再摸出白球的概率，并以求出结果验证你上面的结论.
\item 一学校的学生按照下表分类.
\begin{center}
\begin{tabular}{c|cc}
    \hline
   & 戴眼镜&    不戴眼镜\\
   \hline
    男孩&    120&    280\\
    女孩&    180& 420   \\
    \hline
\end{tabular}
\end{center}
令：事件$A=$“学生是男孩”，$B=$“戴眼镜的学生” 

问：
\begin{multicols}{2}
  \begin{enumerate}[(1)]
    \item $A$和$B$
    \item $A$和$\overline{B}$
    \item $\overline{A}$和$\overline{B}$是相互独立事件吗？
\end{enumerate}  
\end{multicols}

\item 生产一种零件，甲车间的合格率是90\%，乙车间的合格率是97\%，从它们生产的零件中各抽取一件，都抽到合格品的概率是多少？
\item 有一问题，在半小时内，甲能解决它的概率是$\frac{1}{2}$，乙能解决它的概率是$\frac{1}{3}$，如果两人都试图独立地在半小时内解决它，计算：
\begin{multicols}{2}
\begin{enumerate}[(1)]
\item 两人都未解决的概率；    \item 问题得到解
决的概率.
\end{enumerate}
\end{multicols}

\item 将一个硬币连掷5次，5次都出现正面的概率是多少？
\item 制造一种零件，甲机床的废品率是0.04，乙机床的废品率为0.05，从它们制造的产品中各任抽一件，其中恰有一件废品的概率是多少？
\item 电子设备的某一部件由9个元件组成，如果其中任何一个元件坏了，这个部件就不能工作，假定每个元件能使用到3000小时的概率是0.99，计算这个部件能工作3000小时的概率.
\item 甲、乙两人解出习题的概率分别为$\frac{2}{3}$和$\frac{4}{5}$，求下面的概率.
\begin{multicols}{2}
\begin{enumerate}[(1)]
    \item 只有1人解出习题：
    \item 至少1人解出习题．
\end{enumerate}
\end{multicols}


\end{enumerate}


\section{独立重复试验}
某射手射击一次，击中目标的概率是0.9，他射击4次，恰好击中3次的概率是多少？

如果我们把射击一次看作一次试验，那么射击4次就相当于连续依次做4次试验，这四次试验有下面的特点：
\begin{enumerate}[(1)]
    \item 这4次试验的结果彼此无关，是相互独立的。
    \item 每一个试验的结果只有两个：“击中目标”（$A$）或“击不中目标（$\overline{A}$），且每次试验中“击中目标”的概率$\Pr(A)$都是0.9，从而“击不中目标”概率$\Pr(\overline{A})$都是0.1
\end{enumerate}

将上面的例子的特点，推广到一般情形，我们把满足上面两个条件的试验，叫独立重复试验，简称独立试验. 这一类的概率问题称为独立试验概型. 若在试验中，我们只关心出现的事件$A$或$\overline{A}$，这种独立试验概型又称为贝努里（J. Bernoulli）概型. 上面问题（及后面将要讨论的）就是属于贝努里概型.

我们继续讨论其解法：

令${A}_i=$“第$i$次击中”，$i=1,2,3,4$，则$\overline{A}_i=$“第$i$次没有击中”，射击4次击中3次，共有下面几种情况：
\[\overline{A}_1A_2A_3A_4;\quad A_1\overline{A}_2A_3A_4;\quad A_1A_2\overline{A}_3A_4;\quad A_1A_2A_3\overline{A}_4\]
上述每一种情况，都可以看成是在4个位置上取出3个写上$A$，另一个写上$\overline{A}$，所以这些情况的种数等于从4个元素中，取出3个元素的组合数$\comb{3}{4}$，即4种．

由于4次射击，其结果是相互独立的，根据乘法定理，前3次击中，第4次未击中的概率为：
\[\Pr(A_1A_2A_3\overline{A}_4)=\Pr(A_1)\cdot \Pr(A_2)\cdot \Pr(A_3)\cdot \Pr(\overline{A}_4)=0.9^3\x(1-0.9)^{4-3}\]
同理：
$$\Pr(A_1A_2\overline{A}_3A_4)=\Pr(A_1\overline{A}_2A_3A_4)=\Pr(\overline{A}_1A_2A_3A_4)=0.9^3\x(1-0.9)^{4-3}$$
这就是说，上面射击4次，击中3次的4种情况的概率都是$0.9^3\x(1-0.9)^{4-3}$．因为这4种情况彼此互斥，根据加法定理，射击4次，击中3次的概率
\[\begin{split}
    P&=\Pr(A_1A_2A_3\overline{A}_4)+\Pr(A_1A_2\overline{A}_3A_4)+\Pr(A_1\overline{A}_2A_3A_4)+\Pr(\overline{A}_1A_2A_3A_4)\\
    &=\comb{3}{4}\x 0.9^3\x(1-0.9)^{4-3}\\
    &=4\x 0.9^3\x0.1\approx 0.29
\end{split}\]

如果射击4次改为$n$次，击中3次改为$k$次，也可得出类似的结果.

一般地，如果在一次试验中事件$A$发生的概率是$P$，那么在$n$次独立试验中事件$A$恰好发生$k$次的概率：
\[P_n(k)=\comb{k}{n}\cdot P^k(1-P)^{n-k}\]

在应用这个公式的时候，首先要分析实际问题是否满足独立重复试验的两个条件.


\begin{example}
    某气象站天气预报的准确率为80\%，计算：
\begin{enumerate}[(1)]
\item 5次预报恰有4次准确的概率；
\item 5次预报至少有4次准确的概率；
\item 5次预报准确的不多于4次的概率；
\item 5次预报准确的少于4次的概率．
\end{enumerate}

\end{example}

\begin{solution}
\begin{enumerate}[(1)]
    \item 设$A=$“预报一次，结果准确，”预报5次相当于作5次独立试验，知$\Pr(A)=0.8$，所以5次预报恰有4次准确的概率：
\[P_5(4)=\comb{4}{5}\x0.8^4\x(1-0.8)^{5-4}=5\x0.8^4\x0.2\approx 0.41\]
\item 5次预报都准确的概率：
\[P_5(5)=\comb{5}{5}\x0.8^5\x(1-0.8)^{5-5}=0.8^5\approx 0.33\]
因此5次预报至少有4次准确的概率：
\[P=P_5(4)+P_5(5)=0.41+0.33=0.74\]
\item 5次预报准确的不多于4次的概率：
\[\begin{split}
    P&=P_{5}(0)+P_{5}(1)+P_{5}(2)+P_{5}(3)+P_{5}(4)\\
&={\rm C}_{5}^{0}\times 0.8^{0}\times(1-0.8)^{5-0}+{\rm C}_{5}^{1}\times0.8^{1}\times(1-0.8)^{5-1}+{\rm C}_{5}^{2}\times0.8^{2}\times(1-0.8)^{5-2}\\
&\qquad +{\rm C}_{5}^{3}\times0.8^{3}\times(1-0.8)^{5-3}+{\rm C}_{5}^{4}\times 0.8^{4}\times (1-0.8)^{5- 4}\\
&=0.2^{5}+5\times0.8\times0.2^{4}+10\times0.8^{2}\times0.2^{3}+10\times0.8^{3}\times0.2^{2}+5\times0.8^{4}\times0.2\\
&=0.00032+0.0064+0.0512+0.2048+0.4096\approx0.67
\end{split} \]
或
\[P=1-P_{5}(5)\approx 1-0.33=0.67\]

\item 5次预报准确的小于4次的概率：
\[\begin{split}
    P&=P_{5}(0)+P_{5}(1)+P_{5}(2)+P_{5}(3)\\
    &=0.00032+0.0064+0.0512+0.2048\approx 0.26
\end{split}\]
或
\[P=1-P_{5}(4)-P_{5}(5)=1-0.4096-0.33\approx 0.26\]
\end{enumerate}
\end{solution}

\begin{example}
    设有$\ell$升经过紫外线消毒的自来水，其中含有$n$个大肠菌. 今从其中任取一升水检验，问在取出的一升水中含有$k\; (k=0,1,\ldots,n)$个大肠菌的概率是多少？
\end{example}

\begin{solution}
  对于每一个大肠菌来说只有两种可能结果：落入或不落入取出的一升水中，且可以认为每个大肠菌落入与否是彼此独立的. 这样$n$个大肠菌落入所取出的一升水中可以看作做了$n$次独立试验，每次试验结果是某一大肠菌落入该升水（$A$）或不落入该升水（$\overline{A}$）. 而每个大肠菌落入该升水的概率为$P=\Pr(A)=\frac{1}{\ell}$，故在取出的一升水中有$k$个大肠菌
的概率为：
\[P_n(k)=\comb{k}{n}\cdot P^k\cdot (1-P)^{n-k}=\comb{k}{n}\cdot \left(\frac{1}{\ell}\right)^k\cdot \left(1-\frac{1}{\ell}\right)^{n-k}  \]
\end{solution}



\section*{习题3.3}
\begin{enumerate}
    \item 生产一种零件，出现次品的概率是0.04，生产这种零件4件，其中恰有1件次品、恰有2件次品、至多有一件次品的概率各是多少？
    \item 某工厂生产中出现次品的概率为0.01，进行重复抽样检查，共取3个样品，求其中次品数等于0, 1, 2, 3的概率.
    \item 设一射手每次中靶的概率为3/5，求四次射击中：
\begin{multicols}{2}
\begin{enumerate}[(1)]
\item 击中一次，    \item 第三次击中，
\item 击中两次，    \item 第二、三两次击中；
\item 至少击中一次的概率。
\end{enumerate}    
\end{multicols}
    \item 有甲乙丙三批罐头，每批100个，其中各有一个是不合格的. 从三批罐头中各抽出一个，计算：
\begin{enumerate}[(1)]
  \item 3个中恰有一个不合格的概率；
\item 3个中至少有一个不合格的概率． 
\end{enumerate}
    \item 某工厂生产过程中出现次品的概率为0.05，每一百个产品为一批，检查产品质量时，在每批中任取一半来检查，如果发现次品不多于1个，则可以认为这批产品是合格的，求一批产品被认为是合格的概率.
\end{enumerate}

\section*{本章内容要点}
一、本章通过大量实例，初步介绍了随机事件的概率及其计算公式.

二、在一定条件下，可能发生也可能不发生的事件，称为随机事件.反映随机事件出现的可能性的大小的一个量，叫做概率.

三、具有有限性、等可能性的特征的随机现象，我们称为古典概型. 如果试验的基本事件总数为$n$，随机事件$A$所包含的事件数为$m$，那么，$A$的概率就定义为$\frac{m}{n}$，记作
\[\Pr(A)=\frac{m}{n}\]
特别地，必然事件$U$的概率为$\Pr(U)=1$; 不可能事件$V$的概率为$\Pr(V)=0$，显然，$0\le \Pr(A)\le 1$.

四、不可能同时发生的两个事件叫 做互斥事件，当$A$，$B$互斥时，
\begin{equation}
    \Pr(A+B)=\Pr(A)+\Pr(B)\tag{*}
\end{equation} 

五、在事件$B$出现的条件下，事件$A$的出现的概率，我们称为$A$关于$B$的条件概率，记作$\Pr(A|B)$. 且有公式
\begin{equation}
   \Pr(A|B)=\frac{\Pr(AB)}{\Pr(B)}\tag{**} 
\end{equation}

六、对于事件$A$，$B$来说，如果有
\[\Pr(A|B)=\Pr(A),\qquad \Pr(B|A)=\Pr(B)\]
同时成立，那么就称事件$A$，$B$为相互独立事件，两个相互
独立的事件同时发生的概率为：
\begin{equation}
    \Pr(A\cdot B)=\Pr(A)\cdot \Pr(B) \tag{***}
\end{equation}

应该注意区别互斥事件、独立事件这两个概念，注意运用公式(*), (***)的前提条件.

七、其中必有一个发生的两个互斥事件，叫做对立事件，当$A$，$B$是对立事件时，有
\[\Pr(A)=1-\Pr(B)\]

八、具有重复试验的结果是相互独立的，而且每次试验的结果只有$A$与$\overline{A}$的特征的随机现象，我们称为重复独立试验概型（又称贝努里概型）. 如果事件$A$在一次试验中发生的概率是$P$，那么，它在$n$次独立重复试验中恰好发生$k$次的概率是
\[P_n(k)=\comb{k}{n}P^k(1-P)^{n-k}\]


\section*{复习题三}
\begin{enumerate}
    \item 某种集成电路用到2000小时还能正常工作的概率是94\%，使用到3000小时仍能正常工作的概率是78\%，问已经工作了2000小时的集成电路，能继续工作到3000小时的概率是多少？
    \item 若$M$件产品中包含$m$件废品，今在其中任取两件，求：
\begin{enumerate}[(1)]
\item 已知取出两件中有一件为废品的条件下，另一件也是废品概率；
\item 已知两件中有一件不是废品的条件下，另一件是废品的概率.
\end{enumerate}
    \item 一只瓶含有6个黑球，4个红球和2个绿球，两个球是逐个而不再放回地随机取出，求两个球都是：（1）黑的，（2）红的．（3）绿的，（4）相同的颜色的概率.
    \item 甲袋装有白球3个，黑球5个；乙袋中装有白球4个，黑球6个，现在从甲袋中随意地取出一个球放入乙袋，充分掺混后再从乙袋随意地取出一个球放入甲袋，求这时甲袋中的白球数成下面的情况的概率：
\begin{multicols}{2}
\begin{enumerate}[(1)]
\item 增加，    \item 不变．
\end{enumerate}
\end{multicols}

    \item 在$1,2,\ldots,100$中任取一数，求它既能被2整除又能被5整除的概率，又它能被2整除或能被5整除的概率是多少？
\item 假设总人口中男孩和女孩的比例是1:1，在所有恰有三个孩子的家庭中，求：
\begin{enumerate}[(1)]
\item 一个家庭有两个女孩和一个男孩的概率；
\item 已知最大的孩子是女孩的条件下，一个家庭有两个女孩和一个男孩的概率.
\end{enumerate}

\item 在一个恰有三个孩子的家庭里，若已知：已有（1）一个女孩，（2）一个男孩，求恰有两个女孩的概率是多少？
\item 甲、乙、丙三儿童，甲有红球$a$只，白球$b$只；乙全是红球；丙红、白球都有，三人各拿出1球，其中甲是随机地取出的，丙拿红球或白球的概率都是$\frac{1}{2}$，然后让甲从取出的三个球任意取回一个，求甲的球成下列情况的概率：
\begin{multicols}{3}
\begin{enumerate}[(1)]
\item 红球增加；
\item 红球个数不变；
\item 红球减少．
\end{enumerate}
\end{multicols}

\item 下图是某市街道地图的一部分，一人计划从角落$A$处散步到$B$处，这里往北要走过7条街，往东要过4条街，求：
\begin{enumerate}[(1)]
    \item 从$A$到$B$有几种不同的走法？
    \item 他通过$G$或者$H$的概率是多少？
    \item 他通过标有$G$的拐角处的概率是多少？
    \item 已知他通过$G$或$H$. 求他经过拐角$J$的条件概率.
\end{enumerate}

\begin{center}
\begin{tikzpicture}[xscale=.6, yscale=.4]
\foreach \y in {0,1,2,...,9}
{
    \draw[gray](-.5,\y)--(6.5,\y);
}    
\foreach \x in {0,1,2,...,6}
{
    \draw[gray](\x,-.5)--(\x,9.5);
}    
\tkzDefPoints{1/1/A, 3/3/H, 2/6/G, 4/6/J, 5/8/B}
\tkzDrawPoints(A,H,G,J,B)
\tkzLabelPoints[above left](A,G)
\tkzLabelPoints[above right](B,J,H)


\end{tikzpicture}
\end{center}



\item 轰炸机要完成它的使命，必须驾驶员找到了轰炸目标，投弹员又投中了目标，设驾驶员甲和乙找到目标的概率分别为0.9和0.8；又设投弹员丙和丁在驾驶员找到了目标的条件下投中目标的概率分别为0.7和0.6．现在配备两架轰炸机的人员，问甲乙与丙丁应该怎样搭配才能使完成使命有较大的概率？又这个概率比起另一种配合大多少？（只要有一架飞机投中目标即完成使命）.
\item 一个工人看管四台机床，如果在一小时内这些机床不需要人去照顾的概率，第一台是0.79，第二台是0.79，第三台是0.80，第四台是0.81，假设各台机床是否需要照顾相互之间没有影响，计算在这个小时内，这四台机床都不需要照顾的概率.
\item 有甲、乙、丙三批罐头，每批分别有100、120、
150个，其中各有3个是不合格的，从三批罐头中各抽出一个，计算：
\begin{enumerate}[(1)]
\item 3个恰有一个合格的概率；
\item 3个中至少有一个不合格的概率。
\end{enumerate}

\item 甲、乙、丙三步枪手射中某目标的概率各为0.7、0.8及0.85；现向该目标各发一弹，问射中3, 2, 1, 0弹的概率各为多少？
\item 两部车床独立地工作，每部车床不需要工人照顾的概率分别为0.9和0.85，求
\begin{enumerate}[(1)]
\item 都不需要照顾的概率；
\item 恰好一台需要照顾的概率；
\item 同时需要照顾的概率．
\end{enumerate}

\item 有一乘客在电车汽车联合站候车，任何一种车都能使他到达目的地. 如果在五分钟内电车到站的概率是$\frac{1}{2}$，汽车到站的概率是$\frac{1}{3}$. 问这个乘客在五分钟内能够登上任何一种车的概率是多少？
\item 甲袋子内有$m$个白球，$n$个黑球，乙袋子内有$n$个白球，$m$个黑球. 从两个袋子内各任意摸出一个，得到一个白球一个黑球的概率是多少？
\item 一箱螺钉共十万只，次品率为0.004，问事件“从中任意抽验一只，是次品；然后又抽验两只，这两只中至少有一只是次品”的概率有多少？
\item 街道交叉口上交通的畅通程度可用任意给定的一秒钟内有一部车子通过的概率P来描述；假定在不同时刻车子的经过是互不影响的，如果一个步行者只能在三秒钟没有车子通过的情形下才能跨过街道，求他恰巧等待0, 1, 2, 3秒钟的概率.
\item 某仪表有$m$个同样的电子管，其中任一个电子管损坏，这个仪表就不能工作，如果在某段时间内每个电子管损坏的概率是$P$，计算在这段时间内，这个仪表不能工作的概率。
\item 在下图所示的开关电路中，开关$A$、$B$、$C$开或关的概率均独立地为$\frac{1}{2}$，求灯点亮的概率.
\begin{center}
\begin{circuitikz}
\draw(.5,1.5) to [normal open switch, l=$A$] (2.5,1.5) -- (3,1.5);
\draw(3,1.5) to [normal open switch, l=$B$]  (3.5,1.5)--(4,1.5);
\draw(.5,1.5)--(.5,.5) to [normal open switch, l=$C$] (4,.5)--(4,1.5);
\draw(4,1)--(4.5,1)--(4.5,-0.5) to [lamp] (1.5,-0.5) to [battery] (0,-0.5)--(0,1)--(.5,1)--(.5,1.5);
\end{circuitikz}
\end{center}


\item 某种大炮击中目标的概率是0.3，只要以多少门这样的大炮同时射击一次，就可以使击中目标的概率：
\begin{multicols}{2}
\begin{enumerate}[(1)]
\item 超过95\%；    \item 不小于0.99．
\end{enumerate}
\end{multicols}

\item 某售货员在三个柜面上售货，如果在某一小时内柜面需要售货员照顾的概率，第一柜面是0.9，第二柜面为0.8，第三柜面为0.7，假定各个柜面是否需要照顾相互之间不受影响，计算在这个小时内，至少有一个柜面需要售货员的概率.
\item 甲、乙、丙三猎手向同一野猪射击，设甲射中的概率为0.4；乙射中的概率为0.5；丙射中的概率为0.7，若只有一人射中，野猪被击毙的概率为0.2；若其中二人射中，则野猪被击毙的概率为0.6；若三人同时射中，则野猪必定被击毙。求野猪被击毙的概率。
\item 某射手射击一次，击中目标的概率是0.9，他射击了四次。分别写出恰好击中4次，3次，2次，1次，0次的概率，并将它们与$(0.9+0.1)^4$的展开式的各项进行比较，你可得出什么结论？
\item 一次测量中出现正误差和负误差的概率都是$\frac{1}{2}$，在3次测量中，恰好出现2次正误差的概率是多少？恰好出现两次负误差的概率是多少？
\item 一头病牛用某药后被治愈的概率为95\%，计算服用这种药的4头病牛中至少有3头被治愈的概率．
\item 在乒乓球比赛中，甲、乙两人球艺相当，胜负机会相同，试比较下列概率哪一个较大？一方在：
\begin{enumerate}[(1)]
\item 4盘中胜3盘与8盘中胜5盘；
\item 4盘中胜3盘以上与8盘中胜5盘以上；
\item $2n$盘中胜的盘数不多于$n$与$2n$盘中胜$n$盘以上．
\end{enumerate}

\item 每一用户在一小时内向总机拨话的概率等于0.02，该总机有300个分机，试问在一小时内有不少于7个分机向总机拨话的概率是多少？

\item 某一批蚕豆种籽，如果每一粒发芽的概率为90\%，播下5粒种籽，计算：
\begin{enumerate}[(1)]
    \item 其中恰有4粒发芽的概率；
    \item 其中恰有4粒未发芽的概率．
\end{enumerate}

\item 在四次独立试验中，事件$A$至少出现一次的概率为$\frac{80}{81}$，求事件$A$在各次试验中出现的概率.
\item 某小组有10台各为7.5千瓦的机床，如果每台机床的使用情况是互相独立的，且每台机床平均每小时开动12分钟，问全部机床用电超过48千瓦的可能性有多大．
\item 某一学校有730名学生，对每一个人来说，他的生日在某一天的概率设为$\frac{1}{360}$，求：
\begin{enumerate}[(1)]
\item 恰有3名学生的生日为元月1日；
\item 生日在元月1日学生不多于3人；
\item 没有一个学生的生日在元月1日；
\item 至少有一个学生的生日在元月1日的概率．
\end{enumerate}

\item 两个篮球运动员在罚球线投球的命中率分别是0.7与0.6，每人投球3次，计算两人都投进2球的概率．
\item 一个织布工人看管12台布机、在$a$小时内一台布机需要工人照顾的概率等于$\frac{1}{3}$
\begin{enumerate}[(1)]
\item 在$a$小时内，有4台布机需要照顾的概率；
\item 在$a$小时内，需要照顾的布机不多于6台的概率．
\end{enumerate}

\item 在多次试验中，当事件$A$出现不少于3次时，事件B才出现．而事件$A$在每次试验中出现的概率等于0.3，若
\begin{multicols}{2}
    \begin{enumerate}[(1)]
    \item 作了5次试验；
    \item 作了7次试验
\end{enumerate}
\end{multicols}
试求$\Pr(B)$．

\item 一个通讯小组有两套通讯设备，只要其中有一套设备能正常工作，就能进行通讯，每套设备由3个部件组成．只要其中有一个部件出故障，这套设备就不能正常工作，如果在某段时间内每个部件不出故障的概率都是$P$，计算在这段时间内能进行通讯的概率.
\item 甲、乙两个篮球队员各投篮3次，每次投篮时投中得分的概率各为0.6和0.7．求
\begin{multicols}{2}
\begin{enumerate}[(1)]
    \item 得分相同的概率；
    \item 甲得分比乙多的概率．
\end{enumerate}    
\end{multicols}

\end{enumerate}


\appendixchapters
\chapter{附录：统计学简介}

\section*{平均值(均值)}

给了一组数据$x_1,x_2,\ldots,x_n$, 相应地有一个平均数
(也称为均值), 它是如下的算术平均数
\[\frac{1}{n} \sum_{i=1}^nx_i\]
通常用$\bar x$来表示它，平均数是日 常生活中常见的量。例如， 工厂中某一产品的平均日产量，农村中某一作物的平均亩产量，商品的平均销售额，学生考试的平均分数等等。自然会产生这样的问题，为什么要用平均数呢？原因是平均数具有代表性，下面来解释它的意思。

对给定的一组数据$x_1,x_2,\ldots,x_n$, 如果用一个数$C$来代表这些$x_1,x_2,\ldots,x_n$, 则$x_i-C$就表达了$x_i$与$C$的偏差， $x_i-C$的绝对值小，$C$就接近$x_i$。即用$C$来代表$x_i$就代表得好. 由于$C$要代表所有数$x_1,x_2,\ldots,x_n$，因此必须考虑几个偏差$x_{1}-C,x_{2}-C,\ldots,x_{n}-C$. 为了消除正负号的影响， 用$(x_i-C)^2$来衡量$x_i$与$C$ 的差别，于是总的差异就是$\sum\limits_{i= 1}^{n}\left ( x_{i}- C\right ) ^{2}$, 当$x_{1},x_{2},\ldots,x_{n}$确定不变时，应选取$C$使$\sum\limits_{i= 1}^{n}\left ( x_{i}- C\right ) ^{2}$达到最小值，这样的$C$才是代表$x_{1},x_{2},\ldots,x_{n}$这$n$个数据最合适的数。

下面来证明这个数就是$\bar x$
\[\begin{split}
    \sum_{i=1}^{n} ( x_{i}-C)^{2} &=\sum_{i=1}^{n} ( x_{i}-\bar{x}+\bar{x}-{C} )^{2} \\
    &=\sum_{i=1}^{n}\left[\left( x_{i}-\bar{x}\right) ^{2}+2\left( x_{i}-\bar{x}\right) ( \bar{x}-C)+(\bar{x}-C)^{2}\right]\\
    &=\sum_{i=1}^{n}\left( x_{i}-\bar{x}\right)^{2}+2\left[\sum_{i=1}^{n}\left( x_{i}-\bar{x}\right)\right](\bar{x}-C)+n(\bar{x}-C)^{2}
\end{split}\]
由于
$$\sum_{i=1}^{n} ( x_{i}-\bar{x} ) =\sum_{i=1}^{n} x_{i}-n\bar{x}=n\bar{x}-n\bar{x}=0$$
因此
$$\sum_{i=1}^{n}(x_{i}-C)^{2}=\sum_{i=1}^{n}(x_{i}-\bar{x})^{2}+n(\bar{x}-C)^{2}\ge \sum_{i=1}^{n}(x_{i}-\bar{x})^{2}$$
显然，等号成立当且仅当$\bar x=C$.

这就证明了：如果以$\sum\limits_{i=1}^n(x_i-C)^2$来衡量$C$代表$x_1,
x_{2},\ldots,x_{n}$的好坏标准，则$\bar x$是代表性最好的值.

现在来看几个特例：

\begin{example}
    如果$x_1=x_2=\cdots=x_n=a$, 则 $\bar{x}=a$.
\end{example}

\begin{proof}
    $$\bar{x}= \frac{1}{n}\sum_{i=1}^{n}x_{i}= \frac{1}{n}\sum_{i= 1}^{n}a= a$$
\end{proof}

\begin{example}
    如果$x_i=a+id$, $i=0,1,2,\ldots,n-1$,
即$x_0,x_1,\ldots,x_{n-1}$是等差数列，这时
\[\begin{split}
   \bar{x}= \frac {1}{n}\sum _{i= 0}^{n-1}x_{i}&= \frac{1}{n}\sum_{i= 0}^{n- 1}( a+ id) \\
   &= \frac{1}{n}\left [ na+ \frac {(  n- 1 )n}{2}d \right ]=  a+\frac{(n-1)d}{2}
\end{split}\]
\end{example}
例如求$1,2,\ldots, 100$的平均 数，这 时$a=d=1$, $n=100$于是$\bar{x} = 1+ \frac {99}{2}= 50.5$; 又 如求$3, 6, 9, 12, \ldots, 27 $这几 个数的均值，则$a= d= 3$, $n= 9$, 于是$\bar{x}=3+\frac{8}{2}\times 3=15$.


\section*{均方差（方差）}

从上面的讨论可以看出，平均数$\bar x$代表$x_1,\ldots,x_n$的好坏程度是由$\sum\limits^n_{i=1}(x_i-\bar x)^2$来反映的，它是$n$个数据$x_i$对
$\bar x$的偏差平方的总和. 如果用$n$除一下，就得到平均的偏差平方和，它反映了每个数据$x_i$偏离$\bar x$的“平均”状况，因此称它为均方差，简称为方差，通常用$S^2$来表示，即有
\[S^2=\frac{1}{n}\sum^n_{i=1}(x_i-\bar x)^2\]
方差表明了一组数据$x_1,x_2,\ldots,x_n$的分散程度. $S^2$越小，表示$x_i$之间的差异小，$\bar x$的代表性就好；$S^2$越大，表示$x_i$之间的差异也大，数据便很分散，此时$\bar x$代表性就不强.

有时，我们要比较的就是两组数据的分散程度. 这用方差就很方便. 例如有两班学生，甲班有$n$个学生，乙班有$m$个学生，进行教学测验后，甲班的成绩为$x_1,x_2,\ldots,x_n$，乙班的成绩为$y_1,y_2,\ldots,y_m$，平均数$\bar x$与$\bar y$反映了两个班的平均成绩，各自方差$S^2_x$和$S^2_y$反映了两班成绩不齐的程度，方差小的班，学生之间的差异就小.

如果$x_1=x_2=\cdots=x_n=a$，则$\bar x=a$，这时方差$S^2=0$. 如果$x_i=i,\; i=1,2,\ldots,n$，则$\bar x=\frac{n+1}{2}$，而方差
\[\begin{split}
    S^2=\frac{1}{n}\sum^n_{i=1}(x_i-\bar x)^2 &=\frac{1}{n}\sum^n_{i=1}\left(i-\frac{n+1}{2}\right)^2\\
&=\frac{1}{n}\sum^n_{i=1}i^2-\left(\frac{n+1}{2}\right)^2\\
&=\frac{(n+1)(2n+1)}{6}-\frac{(n+1)^2}{4}\\
&=\frac{(n+1)(4n+2-3n-3)}{12}=\frac{n^2-1}{12}
\end{split}\]

平均数$\bar x$与平均方差$S^2$是一组数据$x_1,x_2,\ldots,x_{n}$的两个重要指标，给了一组数$x_1,x_2,\ldots,x_{n}$后，如何具体地算出这两个值，我们在下一小节中来详细讨论这一问题.

\section*{计算公式}
设$x_1,x_2,\ldots,x_{n}$相应的均值和方差为$\bar x$和$S^2_x$，设$y_1,y_2,\ldots,y_{n}$相应的均值和方差为$\bar y$和$S^2_y$，下面先导出几个公式，然后利用它们来具体计算。

\begin{enumerate}
    \item 如果$y_i=x_i+a,\; i=1,2,\cdots,n$。则 $\bar {y}= \bar{x} + a$.
    
\begin{proof}
     $\bar{y} = \frac 1n\sum\limits_{i= 1}^{n}y_{i}= \frac 1n\sum\limits _{i= 1}^{n} ({x_{i}+ a}) = \frac 1n\sum\limits _{i= 1}^{n}x_{i}+ a= \bar x+ a$
\end{proof}
    
\item 如果$y_i=bx_i,i=1,2,...n$, 则$\bar{y}=b\bar{x}$

\begin{proof}
$\bar y= {\frac 1n}\sum\limits _{i= 1}^{n}y_{i}= {\frac 1n}\sum\limits _{i= 1}^{n}bx_{i}= b {\bar {x}}$
\end{proof}

    \item 如果$y_i= bx_{i}+ a,\; i= 1, 2, \ldots, n$, 则$\bar {y}=b\bar{x} + a$
  
\begin{proof}
    由1，2立即可得.
\end{proof}  
    \item $\sum\limits _{i= 1}^{n}\left (  x_{i}- \bar {x} \right ) ^{2}= \sum\limits _{i= 1}^{n}x_{i}^{2}- n\bar {x}^{2}$

\item 如果$y_i=bx_i+a,\; i=1,2,\ldots,n$, 则$S_y^{2}=b^{2}S_{x}^{2}$

\begin{proof}
\[\begin{split}
    S_{y}^{2}= {\frac 1n}\sum\limits  _{i= 1}^{n}(y_{i}- \bar{y})^2&=\frac{1}{n}\sum\limits ^n_{i=1}\left ( bx_{i}+ a- b\bar{x} - a\right ) ^{2}\\
    &=\frac{b^{2}}{n}\sum\limits _{i=1}^{n}(x_{i}-\bar{x})^{2}=b^{2}S_{x}^{2}
\end{split}\]
\end{proof}
\end{enumerate}

公式5告诉我们，当$b=\pm1$时，$S_x^2=S_y^2$. 即一组数据同
加一个常数，或同时改变符号时，它的方差是不变的.

\begin{example}
设某地历年夏季的雨量为：（单位mm）
\[248.7\quad 249.4\quad 133.2\quad 153.5\quad 211.7 \]
求它的平均值$\bar x$及方差$S^2$
\end{example}

取$y_i=x_i-200$，所以$x_i=y_i+200$，于是$\bar x=\bar y+200$，$S^2_x=S^2_y$

列表演算如下：
\begin{center}
\begin{tabular}{ccc}
\hline 
$x_i$ & $y_i=x_i-200$ & $y^2_i$\\
\hline
248.7  & 48.7  & 2371.69\\
249.4& 49.4&2440.36\\
133.2&$-66.8$&4462.24\\
153.5&$-46.5$&2162.25\\
211.7&11.7&136.89\\
\hline
$\Sigma$ &$-3.5$& $11573.43$\\
\hline
\end{tabular}
\end{center}

所以
\[\begin{split}
    \bar x&=\bar y+200=-\frac{3.5}{5}+200=199.3\\
S^2_x&=S^2_y=\frac{1}{n}\left(\sum^n_{i=1} y^2_i-n\bar y^2\right)=\frac{1}{5}(11573.43-2.45)=2314.196
\end{split}
    \]

从演算过程可以看出，$y_i$的平方、求和的计算量比$x_i$的平方、求和的计算量小．这里选200是为了使$x_i-200$数字变小而且减法本身也较易.

\begin{example}
    设$x_i=a+id,\; i=1,2,\ldots,n$，即$x_1,x_2,\ldots,x_n$为等差数列，试求方差$S^2_x$.
\end{example}

为了便于比较，下面用两种方法求解.

\textbf{解一：} 已知$\bar x=a+\frac{n+1}{2}d$，所以$x_i-\bar x=\left(i-\frac{n+1}{2}\right)d$.
\[\sum^n_{i=1}(x_i-\bar x)^2 =\sum^n_{i=1}\left(i-\frac{n+1}{2}\right)^2 d^2=\frac{d^2n(n^2-1)}{12} \]
所以
\[S^2_x=\frac{d^2(n^2-1)}{12} \]

\textbf{解二：} 取$y_i=\frac{x_i-a}{d}=i$, $i=1,2,\ldots,n$，于是$S^2_y=\frac{1}{d^2}S^2_x$，即$S^2_x=d^2S^2_y$. 前面已算过，这时$S^2_y=\frac{n^2-1}{12}$
所以$S^2_x=\frac{d^2(n^2-1)}{12}$.

\begin{example}
    对两组数据$x_1,x_2,\ldots,x_n$及$y_1,y_2,\ldots,y_m$, 设它们的均值及均方差分别为$\bar x,\; S^2_x; \; \bar y,\;  S^2_y$，将这两组数据合并为一组数据，求合并后的数据的均值$\bar z$及均方差$S^2_z$．
\end{example}

\begin{solution}
令$\bar z=\frac{1}{n+m}\left(\sum\limits^n_{i=1}x_i+\sum\limits^m_{i=1}y_i\right)=\frac{1}{n+m}(n\bar x+m\bar y)$

又
\[\begin{split}
    S^2_z&=\frac{1}{n+m}\left[\sum\limits^n_{i=1}(x_i-\bar z)^2+\sum\limits^m_{i=1}(y_i-\bar z)^2\right]\\
    &=\frac{1}{n+m}\left[\sum\limits^n_{i=1}(x_i-\bar x)^2 +n(\bar x-\bar z)^2+\sum\limits^m_{i=1}(y_i-\bar y)^2+m(\bar y-\bar z)^2\right]\\
    &=\frac{1}{n+m}\left[nS^2_x+n(\bar x-\bar z)^2+mS^2_y+m(\bar y-\bar z)^2\right]
\end{split}\]
但是$\bar z=\frac{n}{n+m}\bar x+\frac{m}{n+m}\bar y$，代入有
\[\bar x-\bar z=\frac{m}{n+m}(\bar x-\bar y),\qquad \bar y-\bar z=\frac{n}{n+m}(\bar y-\bar x)\]
\[S_{z}^{2}= \frac n{n+ m} S_{x}^{2}+ \frac m{n+ m} S_{y}^{2}+ \frac n{n+ m}\left(\frac m{n+ m}\right)^{2}(\bar x- \bar{y} ) ^{2}+\frac{m}{n+m}\left(\frac{n}{m+n}\right)^{2}(\bar{y}-\bar x)^{2}\]
所以
\[S_{z}^{2}=\frac{1}{n+m}\left(nS_{x}^{2}+mS_{y}^{2}\right)+\frac{nm}{(n+m)^{2}}(\bar{x}-\bar{y})^2\]
总之有
$$\bar z=\frac{1}{n+m}(n\bar{x}+m\bar{y})$$
\[S_{z}^{2}=\frac{1}{n+m}\left(nS_{x}^{2}+mS_{y}^{2}\right)+\frac{nm}{(n+m)^{2}}(\bar{x}-\bar{y})^2\]
\end{solution}

特别，当$m=1$时，$S^2_y=0$，记$y_1=x_{n+1}$，这就得到增加一个新数据的递推公式
\[\begin{split}
    \bar z&= \frac{1}{n+1}(n\bar x+x_{n+1})\\
    S^2_z&=\frac{n}{n+1}S^2_x+\frac{n}{(n+1)^2}(\bar x-x_{n+1})^2
\end{split}\]

不难看出，如果从$x_1,x_2,\ldots,x_n$中减少一个数据，例如删去$x_n$，那么$x_1,x_2,\ldots,x_{n-1}$的均值$\bar x_{*}$与均方差$S^2_{*}$可以用$\bar x$与$S^2_x$及$x_n$表示，即有递推公式
\[\begin{split}
    \bar x_{*}&= \frac{1}{n-1}(n\bar x-x_{n})\\
    S^2_{*}&=\frac{n}{n-1}S^2_x-\frac{n}{(n-1)^2}(\bar x-x_{n})^2
\end{split}\]

\section*{最小二乘法}
在实际工作中经常需要从实测的数据求出变量之间的关系式. 例如年龄和血压是有关系的，随着年龄的增长，血压是会增高的，调查了几百人的年龄与血压的情况，将资料按年龄分组. 用各组的均值（年龄的均值和血压的均值），得到如下的数据：

\begin{minipage}{.35\textwidth}
\begin{center}
    \begin{tabular}{cc}
\hline
        年龄$x_i$ & 心脏收缩压$y_i$\\
\hline
35&114\\
45&124\\
55&143\\
65&158\\
75&166\\
\hline
    \end{tabular}
\end{center}
\end{minipage}\hfill
\begin{minipage}{.6\textwidth}
\centering
\begin{tikzpicture}[>=stealth, scale =.8]
\draw[->](-1,0)--(6,0)node[above]{$x$（年龄）};
\draw[->](0,-1)--(0,8)node[right]{$y$（血压）};
\foreach \x/\y in {1/35,2/45,3/55,4/65,5/75}
{
    \draw(\x,0)node[below]{$\y$}--(\x,.1);
}
\foreach \x/\y in {1/110,2/120,3/130,4/140,5/150,6/160,7/170}
{
    \draw(0,\x)node[left]{$\y$}--(.1,\x);
}
\tkzDefPoints{1/1.4/A, 2/2.4/B, 3/4.3/C, 4/5.8/D, 5/6.6/E}
\tkzDrawPoints(A,B,C,D,E)
\node at (A) [left]{$P_1$};
\node at (B) [right]{$P_2$};
\node at (C) [left]{$P_3$};
\node at (D) [left]{$P_4$};
\node at (E) [right]{$P_5$};
\draw[domain=.5:5.5, smooth, thick]plot(\x, 1.35*\x+.1);
\node[below left]{$O$};
\end{tikzpicture}

\end{minipage}

从表上的数据或图上的点来看，
血压值$y$与年龄数$x$似乎有直线的关系，但是这五个点又不正好在一条直线上，于是发生了一个问题：如何求一个直线方程，使得这条线与五个点最接近.

用$P_i$表示点$(x_i,y_i)\; i=1,2,3,4,5$．对给定的直线$\ell:\; y=a+bx$. 在直线$\ell$上取五个点，使其横坐标与$P_i$的横坐标$x_i$相同，这五个点应是
\[Q_i(x_i, a+bx_i) \qquad i=1,2,3,4,5\]
很明显，这五个点的纵坐标之差为$y_i-a-bx_i$. 用数值$\sum\limits^5_{i=1}(y_i-a-bx_i)^2$
来衡量直线$\ell$上五个点$Q_1,\ldots,Q_5$和已给的五个点$P_1,\ldots,P_5$的总差距，这个差距愈小，就认为这条直线愈接近这些点. 因此，最接近的直线应使这个差距达到最小值，下
面用和以前类似的方法来求出$a$和$b$.

已知$n$个点有坐标$(x_i,y_i)\; i=1,2,\ldots,n$．求直线$\ell:\; y=a+bx$，使和这$n$个点在上述意义下最接近，也就是求数值$a$和$b$，使
\[Q=\sum^n_{i=1}(y_i-a-bx_i)^2\]
达到最小值.

\begin{blk}{定理（最小二乘法）}
使$Q$达最小值的$a,b$为
\[\hat a=\bar y-\hat b \bar x,\qquad \hat b=\frac{\sum\limits^n_{i=1}(x_i-\bar x)(y_i-\bar y)}{\sum\limits^n_{i=1}(x_i-\bar x)^2}\]
\end{blk}

\begin{proof}
对任意数$a,b,\hat a,\hat b$，有
\[\begin{split}
Q&=\sum^n_{i=1}(y_i-a-bx_i)^2 =\sum^n_{i=1}\left[y_i-\hat a-\hat b x_i+\hat a-a+(\hat b-b)x_i\right]^2\\ 
&=\sum^n_{i=1}\left(y_i-\hat a-\hat b x_i\right)^2+2\sum^n_{i=1}\left(y_i-\hat a-\hat b x_i\right)\cdot \left[\hat a-a+(\hat b-b)x_i\right]\\
&\qquad +\sum^n_{i=1}\left[\hat a-a+(\hat b-b)x_i\right]^2
\end{split}\]
取$\hat a,\hat b$使得
\[\sum^n_{i=1}\left(y_i-\hat a-\hat b x_i\right)=0,\qquad \sum^n_{i=1}x_i\left(y_i-\hat a-\hat b x_i\right)=0\]
那么上面
\[\begin{split}
    Q&=\sum^n_{i=1}\left(y_i-\hat a-\hat b x_i\right)^2+\sum^n_{i=1}\left[\hat a-a+(\hat b-b)x_i\right]^2\ge \sum^n_{i=1}\left(y_i-\hat a-\hat b x_i\right)^2
\end{split}\]
因此若能由$\hat a$，$\hat b$之条件唯一
引出$\hat a$，$\hat b$，那么它就是所求的使$Q$达极小值的$a$，$b$了。

将条件改变为
\[\hat a+\hat b \bar x=\bar y,\quad n\hat a \bar x+\hat b\sum^n_{i=1}x^2_i=\sum^n_{i=1}x_iy_i\]
由$\hat a=\bar y-\hat b\bar x$，代入后式，有
\[n\left(\bar y-\hat b\bar x\right)\bar x+\hat b\sum^n_{i=1}x^2_i=\sum^n_{i=1}x_iy_i\]
即有
\[\hat b\left[\sum^n_{i=1}(x_i-\bar x)^2\right]=\sum^n_{i=1}x_iy_i-n\bar x\bar y=\sum^n_{i=1}(x_i-\bar x)\cdot (y_i-\bar y)\]
所以\[\hat b=\frac{\sum\limits^n_{i=1}(x_i-\bar x)(y_i-\bar y)}{\sum\limits^n_{i=1}(x_i-\bar x)^2}\]
于是$\hat a=\bar y-\hat b\bar x$, 这证明了定理.
\end{proof}

回到原来的例子，下面列表进行计算.
\begin{center}
\begin{tabular}{cccccc}
\hline
$x_i$ & $y_i$&$x_i-\bar x$&$y_i-\bar y$&$(x_i-\bar x)^2$&$(x_i-\bar x)(y_i-\bar y)$\\
\hline
35&114&$-20$&$-27$&400&540\\
45&124&$-10$&$-17$&100&170\\
55&143&0&2&0&0\\
65&158&10&17&100&170\\
75&166&20&25&400&500\\
\hline
$\Sigma 275$&705&&&1000&1380\\
平均55&141\\
\hline
\end{tabular}
\end{center}

\[\hat b=\frac{1380}{1000}=1.38,\qquad \hat a=141-1.38\x 55=65.1\]
于是求得最接近这五点的直线为
\[y=65.1+1.38x\]

现在，进一步来考查一下，由上述方程算得的$x=35$, 45, 55, 65, 75时$y$值是多少呢？它们与实际的观测值$y_i$的差又是多少呢？为了以示区别，由方程算出的$65.1+1.38x_i$记为$\hat y_i$有

\begin{center}
\begin{tabular}{cccc}
\hline
实测值&方程算得的值& 偏差& 偏差平方\\
$y_i$&$\hat y_i$&$y_i-\hat y_i$&$\left(y_i-\hat y_i\right)^2$\\
\hline
114&113.4&0.6&0.36\\
124&127.2&$-3.2$&10.24\\
143&141.0&2.0&4.00\\
158&154.8&3.2&10.24\\
166&168.6&$-2.6$&6.76\\
\hline
$\Sigma$&&0&31.60\\
\hline
\end{tabular}
\end{center}

从上表可以看出
\[\sum^n_{i=1}\left(y_i-\hat y_i\right)=0\]
这不是偶然的，可以证明，在一般情况下，总有$\sum\limits^n_{i=1}\left(y_i-\hat y_i\right)=0$. 事实上，
\[\begin{split}
\sum^n_{i=1}\left(y_i-\hat y_i\right)&=\sum^n_{i=1}y_i-\sum^n_{i=1}\left(\hat a+\hat b x_i\right)    \\
&=n\bar y -n\hat a-n\hat b \bar x=n\left(\bar y-\hat a-\hat b \bar x\right)
\end{split} \]
由于已知$\hat a=\bar y-\hat b\bar x$，所以有$\sum\limits^n_{i=1}\left(y_i-\hat y_i\right)=0$


因此，这一结果常常可以作为验算所计算的结果是否正确，但是，在实际使用时，在对大量数据进行计算时，由于
舍入误差的累积会使$\sum\limits^n_{i=1}\left(y_i-\hat y_i\right)\ne 0$，不过这个$\sum\limits^n_{i=1}\left(y_i-\hat y_i\right)$不会大，而在$\sum\limits^n_{i=1}\left(y_i-\hat y_i\right)$的绝对值很大时，这往往是由于计算上的错误所造成的.


\fi

\ifPartB
\ifFullBook\volume{第六册（下）}\fi
\chapter{函数的极限和连续函数的性质}

\section{函数的极限}
\subsection{函数极限的概念}

在第四册下，我们研究了数列的极限，数列是一种特殊的函数，这里的自变数$n$取自然数列$1, 2, 3,\ldots,n,\ldots$的值，现在我们来研究更一般的情形，即函数$f(x)$随$x$连续变化而变化的情形，下面转到一般函数的极限．

\begin{blk}{定义1}
    如果$x$通过任何一个无限增大的数列$\{x_n\}$, 对应的函数值数列
$f (x_1 ) ,f (x_2) , \ldots,f (x_n ) ,\ldots$
都以定数$\ell$为它的极限，就说函数$f(x)$, 当$x\to +\infty$时，以$\ell$为极限，记作
\begin{equation}
    \lim_{x\to+\infty}f(x)=\ell\qquad \text{或}\qquad f(x)\to \ell\quad (x\to +\infty)
\end{equation}
\end{blk}
 
从几何上看，极限式(1.1)表示随着$x$无限增大，曲线$y=f(x)$以直线$\ell$为渐近线（图1.1）．
\begin{figure}[htp]
    \centering
\begin{tikzpicture}[>=latex]
\draw[->](-3,0)--(5,0)node[right]{$x$};  
\draw[->](0,-1)--(0,3)node[right]{$y$};    
\draw[thick](-3,1)--(5,1)node[right]{$y=\ell$};
\draw[very thick, rounded corners=10pt](-3,.1) --(-1.5,.3) --(-.5,.8) --(1,1.5) --(2.5,.7) -- (3,1.2)--(3.5,.8) --(4.5,1.05);
\node at (0,0)[below right]{$O$};

\end{tikzpicture}
    \caption{}
\end{figure}

类似地，可以定义函数极限$\Lim_{x\to -\infty} f(x)=\ell$, 这时，变量
$x$通过代数值无限地变小，而绝对值无限地增大的任何一个数列$\{x_n\}$．

如果函数$f(x)$当$x\to+\infty$和$x\to-\infty$时，都以定值为极限，就说$f(x)$当$x\to \infty$时，以定值$\ell$为极限，记作$\Lim_{x\to\infty}f(x)=\ell$, 或者$f(x)\to \ell\; (x\to\infty)$．

\begin{example}
    证明$\Lim_{x\to\infty}\frac{1}{x}=0$
\end{example}

\begin{proof}
任何数列$\{x_n\}$的值$x_1,x_2,\ldots,x_n,\ldots$趋向于$+\infty$或$-\infty$时，对应的函数数列
\[\frac{1}{x_1},\frac{1}{x_2},\ldots, \frac{1}{x_n},\ldots\]
的绝对值$\left|\frac{1}{x_n}\right|$便趋向于零，即
    \[\lim_{n\to\infty}\left|\frac{1}{x_n}\right|=0\]
    从而$\Lim_{x\to\infty}\frac{1}{x}=0$
\end{proof}

\begin{example}
说明函数$\sin x$, 当$x\to\infty$时，没有极限．
\end{example}

\begin{solution}
令自变量$x$取数列$x_n=-\frac{\pi}{2}+2n\pi\quad (n=1,2,3,\ldots)$的值趋向$+\infty$，则对应的函数值数列
\[\begin{split}
    \sin x_n&=\sin\left(-\frac{\pi}{2}+2n\pi\right)\\
&=\sin\left(-\frac{\pi}{2}\right) =-1 \qquad (n=1, 2, 3, \ldots) .
\end{split}\]
恒取定值$-1$, 于是
\[\lim_{n\to\infty} \sin x_n=\lim_{n\to\infty} \sin \left(-\frac{\pi}{2}+2n\pi\right)=-1\]

再令自变量$x$取数列
\[x_n=\frac{\pi}{2}+2n\pi\qquad  (n=1, 2, 3, \ldots )\]
的值趋向于$+\infty$, 则对应的函数值数列
\[\sin x_n =\sin\left(\frac{\pi}{2}+2n\pi\right) =1\qquad  (n=1, 2, 3, \ldots)\]
恒取定值1,于是
\[\lim_{n\to\infty} \sin x_n=\lim_{n\to\infty} \sin\left(\frac{\pi}{2}+2n\pi\right) =1\]
由于$x$取趋向$+\infty$的两个不同数列时，$y=\sin x$可以有不同的极限，因此
$\Lim_{x\to\infty} \sin x$不存在．
\end{solution}

\begin{blk}{定义2}
 设函数$f(x)$在点$a$附近有定义（但在$x=a$时，可以没有定义），如果当自变量$x$不论通过怎样一个以$a$为极限但始终不等于$a$的数列$\{x_n\}$, 对应的函数值数列$f (x_1) ,f (x_2) , \ldots,f(x_n ) ,\ldots$
总有极限$\ell$, 就说：

当$x$趋近于$a$时，$f(x)$以$\ell$为极限，记作
\begin{equation}
   \lim_{x\to a}f(x)=\ell,\qquad \text{或}\qquad f(x)\to \ell\quad (x\to a) 
\end{equation}
\end{blk}

极限式(1.2)的几何意义如图1.2所示：当$x$无限地靠近$a$，但总不能等于$a$时，曲线$y=f(x)$上的点$(x,f(x))$无限地靠近$(a,\ell)$点．

\begin{figure}[htp]
    \centering
\begin{tikzpicture}[>=latex]
\draw[->](-3,0)--(5,0)node[right]{$x$};  
\draw[->](0,-1)--(0,3)node[right]{$y$};    
\draw[domain=-2:4, samples=100, very thick]plot(\x, {.5*cos(\x r)+1});
\node at (0,0)[below right]{$O$};
\node at (-1,1.5){$y=f(x)$};
\draw[dashed](pi/2,0)node[below]{$a$}--(pi/2,1)node[above]{$(a,\ell)$};

\tkzDefPoints{1.57/1/A, 1.3/1.133/B, 1.8/0.886/C}
\tkzDrawPoints(A,B,C)
\tkzLabelPoint[below left](B){$(x,f(x))$}
\tkzLabelPoint[above right](C){$(x,f(x))$}
\end{tikzpicture}
    \caption{}
\end{figure}

初学的人常常要问：为什么在定义中谈到$x$趋近于$a$时，要限制$x$始终不等于$a$呢？这是因为我们关心的是函数$f(x)$在$a$附近的变化趋势，它和函数$f(x)$在
$x=a$这一点的值没有什么必然关系，这也就是说，无论$f(x)$在点$a$取什么值甚至没有定义，都不影响在这一点的极限的存在和极限值．

\begin{example}
设$f(x)=\frac{3}{4}\cdot \frac{x^2-1}{x-1}$，$x\in(-\infty,1)\cup(1,+\infty)$，求$\Lim_{x\to 1}f(x)$
\end{example}

\begin{solution}
    $f(x)=\frac{3}{4}\cdot \frac{x^2-1}{x-1}$在$x=1$时无意义，因为那时
    分母就变成零，因此，这里没有函数值$f(1)$, 曲线$y=f(x)$也没有相应于横坐标为1的那个点，但是让$x$任意地趋近于1是完全可以的，若$x\ne 1$, 则有
  \[  f (x) =\frac{3}{4}\cdot  \frac{(x-1)(x+1)}{x-1}=\frac{3}{4}(x+1)\]
  因此，不论$x$通过怎样一个以1为极限的数列$\{x_n\}$, 对于相应的数列$\{f(x_n)\}$, 我们都有
  \[\lim_{x_n\to 1} f (x_n) =\frac{3}{4} (1+1) =\frac{3}{2}\]

    从几何上看，曲线$y=f(x)$除去点$\left(1,1\frac{1}{2}\right)$外是与直线$y=\frac{3}{4}(x+1)$一致的，唯独在那一点，曲线有个空隙，而在$x=1$的邻近的点只要充分接近于点1, 所对应的函数值与$\frac{3}{2}$的差的绝对值可以任意小，如图1.3．
\end{solution}
    
\begin{figure}[htp]
    \centering
\begin{tikzpicture}[>=latex]
    \draw[->](-3,0)--(5,0)node[right]{$x$};
    \draw[->](0,-1.5)--(0,4.5)node[right]{$y$};
\foreach \x in {1,2,3}
{
    \draw(0,\x)node[left]{$\x$}--(.1,\x);
    \draw(\x,0)node[below]{$\x$}--(\x,.1);
}
\node at (.25,-.25){$O$};
\draw[dashed](0,1.5)--(1,1.5)--(1,0);
\draw[domain=-2:4, samples=10, very thick]plot(\x, {0.75*(\x+1)});
\draw(1,1.5)[fill=white]  circle(1.5pt) node[right]{$\left(1,1\tfrac{1}{2}\right)$};
\node at (3.25,3)[right]{$y=\frac{3}{4}\cdot \frac{x^2-1}{x-1}$};
\end{tikzpicture}
    \caption{}
\end{figure}

\begin{example}
证明当$x\to 0$时，函数$f(x)=\sin\frac{1}{x}$没有极限．
\end{example}

\begin{solution}
函数$f(x)=\sin\frac{1}{x}$对于一切$x\ne 0$的值有定义，因此这个函数在点$x=0$的邻域内有定义．

当$x$取数列$\{x_n\}=\left\{\frac{2}{(2n+1)\pi}\Big| n=1,2,3,\ldots\right\}$的值而趋于零时，数列$\left\{\frac{1}{x_n}\right\}$相应的值是
\[\frac{3\pi}{2},\frac{5\pi}{2},\frac{7\pi}{2},\ldots,(2n+1)\frac{\pi}{2},\ldots\]
此时数列$\left\{\sin\frac{1}{x_n}\right\}$便交替地取$-1$和$+1$这两个数值，换言之
\[\sin\frac{1}{x_n}=(-1)^n,\qquad n=1,2,3,\ldots\]
因此，当$n\to\infty$时，数列$\left\{\sin\frac{1}{x_n}\right\}$不趋于任何极限值．这就证明了当$x\to 0$时，函数$f(x)=\sin\frac{1}{x}$的极限不存在．
\end{solution}

函数的图象大致如图1.4所示．曲线关于原点对称，在包含原点的每一个对称邻域$(-\delta, \delta)$内，曲线$y=\sin\frac{1}{x}$
在原点的邻近作无数次振动，且曲线的振幅恒为1, 虽将原点的邻域的长缩小，也不能减少振动的次数．

\begin{figure}[htp]
    \centering
\begin{tikzpicture}[>=latex]
    \draw[->](-4,0)--(4,0)node[right]{$x$};
    \draw[->](0,-2)--(0,2)node[right]{$y$};
    \foreach \x in {1,-1}
    {
        \draw[dashed](-4,\x)--(4,\x)node[right]{$y=\x$};
    }
\draw[domain=-4:-.05, samples=1000, thick]plot(\x, {sin(180/pi/\x)});
\draw[domain=.05:4, samples=1000, thick]plot(\x, {sin(180/pi/\x)});
\foreach \x/\xtext in {1/1, .637/\frac{2}{\pi}}
{
    \draw(\x, 0)node[below]{$\xtext$}--(\x,.1);
}
\end{tikzpicture}    
    \caption{}
\end{figure}

上面是用数列的极限来说明函数的极限，其实我们也可以直接定义函数极限．

\begin{blk}{定义3}
  如果函数$f$在点$a$邻域上有定义（可能去掉点$a$本身），并且对于任给的$\varepsilon>0$，总存在$\delta>0$，使得当$0<|x-a|<\delta$时，就有$|f(x)-\ell|<\varepsilon$, 那么就说$\ell$为当$x$趋近于$a$时，函数$f$在点$a$的极限值．
\end{blk}

我们对这个定义需要说明以下几点：
\begin{enumerate}
    \item 用定义3验证某数$\ell$是函数$f$在点$a$的极限的办法就是对于任给的$\varepsilon>0$, 要找到这样的正数$\delta$使得能够由不等式$|x-a|<\delta$推出不等式$|f(x)-\ell|<\varepsilon$，虽然$\varepsilon$是任意的正数，但是在找$\delta$的过程中，$\varepsilon$是固定不变的，$\delta$依赖于$\varepsilon$．
    \item 对于已给的$\varepsilon$, 只要证明有一个$\delta>0$存在就行．因为如果有一个$\delta$存在，把$\delta$再缩小一些，显然仍满足我们的要求．
    \item 不等式$|x-a|>0$只是说明$x\ne a$, 即把$x$等于$a$的情况去掉，这是因为我们关心的是函数$f$在点$a$附近的变化趋势，而和函数在$x=a$这点的值无关．
    \item 我们指出定义2和定义3是等价的．
\end{enumerate}

\begin{example}
用定义3证明$\Lim_{x\to 1}\frac{x^3-1}{x-1}=3$
\end{example}

\begin{proof}
    任给$\varepsilon>0$, 要找$\delta>0$, 使由$0<|x-1|<\delta$推出
    $\left|\frac{x^3-1}{x-1}-3\right|<\varepsilon$成立．

    当$x\ne 1$时，
\[\begin{split}
    \left|\frac{x^3-1}{x-1}-3\right|&=|(x^2+x+1)-3|\\
    &=|(x-1)(x+2)|=|x-1|\cdot |x+2|
\end{split}\]

要由$|x-1|\cdot |x+2|<\varepsilon$找$\delta$，显然，这里因子$|x+2|$引起了麻烦．为方便起见，先假定$0<|x-1|<1$，即取$\delta_1=1$，这样
\[0<|x-1|<1\quad\to \quad |x+2|=|(x-1)+3|\le |x-1|+3<4\]因此，要使
\[\begin{cases}
    0<|x-1|<1\\
    |x-1|\cdot |x+2|<\varepsilon
\end{cases}\]
只须
\[\begin{cases}
    0<|x-1|<1\\
    4|x-1|<\varepsilon
\end{cases}\to \quad \begin{cases}
    0<|x-1|<1\\
    |x-1|<\frac{\varepsilon}{4}
\end{cases}\]
由此可见，只须取$\delta=\min\left(1,\frac{\varepsilon}{4}\right)$，即取$\delta$为1与$\frac{\varepsilon}{4}$中的较小者．

$\therefore\quad $对于任意$\varepsilon>0$，取$\delta=\min\left(1,\frac{\varepsilon}{4}\right)$，则当$0<|x-1|<\delta$时，即有
\[\left|\frac{x^3-1}{x-1}-3\right|<\varepsilon\]
这也就证明了
\[\lim_{x\to 1}\frac{x^3-1}{x-1}=3\]
\end{proof}

\begin{example}
    证明$\Lim_{x\to a}\sqrt{x}=\sqrt{a}\quad (a>0)$
\end{example}

\begin{proof}
对于任意的$\varepsilon>0$, 我们必须找到一个$\delta>0$, 使得当$|x-a|<\delta$时，$|\sqrt{x}-\sqrt{a}|<\varepsilon$成立，因为
\[|\sqrt{x}-\sqrt{a}|=\frac{|x-a|}{\sqrt{x}+\sqrt{a}}<\frac{|x-a|}{\sqrt{a}}\]
所以要使$\frac{|x-a|}{\sqrt{a}}<\varepsilon$，只须$|x-a|<\sqrt{a}\varepsilon$

如果取$\delta=\sqrt{a}\varepsilon$，则$\frac{|x-a|}{\sqrt{a}}<\frac{\sqrt{a}\varepsilon}{\sqrt{a}}=\varepsilon$
因此，对于任意$\varepsilon>0$，取$\delta=\sqrt{a}\varepsilon$，则当$|x-a|<\delta$时，就有$|\sqrt{x}-\sqrt{a}|<\varepsilon$成立．这也就证明了
\[\lim_{x\to a}\sqrt{x}=\sqrt{a}\quad (a>0)\]
\end{proof}

有时虽然$f(x)$在某点的左边（或右边）没有定义，如图1.5中的$a,b$两点，我们也可以谈论$f$在点$a$或点$b$两点的极限，譬如对于所有小于$a$的数，$f$虽然没有定义，但是我们可以考察，当$x$从$a$的右侧趋近于$a$时，函数$f$的变化趋势，也就是考察$f$的单边极限是否存在．

\begin{figure}[htp]
    \centering
\begin{tikzpicture}[>=latex]
    \draw[->](-1,0)--(3,0)node[right]{$x$};
    \draw[->](0,-1)--(0,2.5)node[right]{$y$};
\draw[very thick](.5,1)to[bend right=-15]node[above]{$y=f(x)$} (2.5,1.5);
\draw(.5,1)--(.5,0)node[below]{$a$};
\draw(2.5,1.5)--(2.5,0)node[below]{$b$};
\node at (-.25,-.25){$O$};
\end{tikzpicture}
    \caption{}
\end{figure}

\begin{blk}{定义}
     设$f(x)$在区间$(a,b)$上有意义，如果任给$\varepsilon>0$, 总存在某个$\delta>0$使得当$x\in (a,a+\delta)$时，总有$|f(x)-\ell|<\varepsilon$，我们就说函数$f$在点$a$以$\ell$为\textbf{右极限}，记作：
\[\lim_{x\to a^+}f(x)=\ell\]
\end{blk}

类似地，可以定义\textbf{左极限}，只需把开区间$(a,a+\delta)$换成$(b-\delta,b)$就行了，并记作
\[\lim_{x\to b^-} f(x)=\ell\]

\begin{example}
    设函数$f(x)=\begin{cases}
        x-1,&x\le 1\\
x+1,&x>1
    \end{cases}\quad $（图1.6）

    求$\Lim_{x\to 1^-}f(x)$和$\Lim_{x\to 1^+}f(x)$
\end{example}

\begin{figure}[htp]
    \centering
    \begin{minipage}[t]{0.48\textwidth}
    \centering
    \begin{tikzpicture}[>=stealth, scale=.6]
        \draw[->](-3,0)--(4,0)node[right]{$x$};
        \draw[->](0,-3)--(0,4.5)node[right]{$y$};
        \draw(0,1)node[left]{1}--(.1,1);
        \draw(0,2)node[left]{2}--(.1,2);
        \draw(1,0)node[below]{1}--(1,.1);
        \node at (-.25,-.25){$O$};
        \draw[domain=-2:1, samples=10, very thick]plot(\x,{\x-1});
        \draw[domain=1:3, samples=10, very thick]plot(\x,{\x+1});
    \node at (-2,-3)[left]{$y=x-1$};
    \node at (3,4)[right]{$y=x+1$};
    \draw[dashed](1,0)--(1,2)--(0,2);
    \draw (1,2)[fill=white] circle(2.5pt);
    
    \end{tikzpicture}
    \caption{}
    \end{minipage}
    \begin{minipage}[t]{0.48\textwidth}
    \centering
    \begin{tikzpicture}[>=stealth, scale=.8]
        \draw[->](-3,0)--(4,0)node[right]{$x$};
    \draw[->](0,-2.5)--(0,3)node[right]{$y$};
    \foreach \x in {-2,-1,1,2,3}
    {
        \draw(\x,0)node[below]{$\x$}--(\x,.1);
    }
\foreach \x in {-2,-1,1,2}
{
    \draw(0,\x)--(.1,\x)node[right]{$\x$};
}
\foreach \x in {-2,-1,...,2}
{
    \draw[very thick](\x,\x)--(\x+1,\x);
    \draw (\x+1,\x)[fill=white]circle(2pt);
}
    \node at (.25,-.25){$O$};
    \end{tikzpicture}
    \caption{}
    \end{minipage}
    \end{figure}


\begin{solution}
\[\begin{split}
    \lim_{x\to 1^-}f(x)&=\lim_{x\to 1^-}(x-1)=0\\
    \lim_{x\to 1^+}f(x)&=\lim_{x\to 1^+}(x+1)=2
\end{split}\]
图1.6显示了上面的结果．
\end{solution}

\begin{example}
    函数$y=[x]$代表不超过$x$的最大整数，即若$n\le x<n+1$，$n\in\mathbb{Z}$，则$y=[x]=n$．（图1.7）

求$\Lim_{x\to 2^+}\frac{[x]}{x},\quad \Lim_{x\to 2^-}\frac{[x]}{x}$
\end{example}

\begin{solution}
    显然，当$x=2$时，$\frac{[x]}{x}=\frac{[2]}{2}=\frac{2}{2}=1$, 又
\[\frac{[x]}{x}=\begin{cases}
    \frac{1}{x},& x\in(1,2)\\
    \frac{2}{x},& x\in(2,3)
\end{cases}\]
所以
\[\begin{split}
    \lim_{x\to 2^+}\frac{[x]}{x}&=\lim_{x\to 2^+}\frac{2}{x}=1\\
     \lim_{x\to 2^-}\frac{[x]}{x}&=\lim_{x\to 2^-}\frac{1}{x}=\frac{1}{2}
\end{split}\]    
\end{solution}


下面的命题说明函数的极限与函数的单边极限的关系：

\begin{blk}{定理}
极限$\Lim_{x\to a} f(x)$
存在的必要和充分条件是左极限$\Lim_{x\to a^-} f(x)$和右极限$\Lim_{x\to a^+} f(x)$都存在，并且二者相等．
\end{blk}

\begin{proof}
    必要性
    
    如果$\Lim_{x\to a} f(x)=\ell$, 就是说任给$\varepsilon>0$, 总存在$\delta>0$,
使得当$0<|x-a|<\delta$, 即当$x\in (a-\delta,a)\cup(a,a+\delta)$时，有$|f(x)-\ell|<\varepsilon$．换言之，当$x\in (a-\delta,a)$和$x\in (a,a+\delta)$时，都有$|f (x) -\ell|<\varepsilon$，因此
\[\Lim_{x\to a^-} f(x)=\ell,\qquad \Lim_{x\to a^+} f(x)=\ell\]

充分性

如果$\Lim_{x\to a^-} f(x)=\ell$且$\Lim_{x\to a^+} f(x)=\ell$, 那么总存在$\delta_1>0$, 使得当$x\in(a-\delta_1,a)$时，有
$|f (x) -\ell|<\varepsilon$．

又存在$\delta_2>0$, 使得当$x\in (a,a+\delta_2)$时，有$|f (x) -\ell|<\varepsilon$．

取$\delta=\min(\delta_1,\delta_2)$，于是当$x\in(a-\delta,a+\delta)$时，有$|f (x) -\ell|<\varepsilon$．
这就是说：
\[\Lim_{x\to a} f(x)=\ell\]
\end{proof}

\begin{example}
说明$\Lim_{x\to 3}\frac{|x-3|}{x-3}$是否存在？
\end{example}

\begin{solution}
\[\begin{split}|x-3|&=\begin{cases}
    x-3, & x>3\\
    3-x, &x<3
\end{cases}\\
    \Lim_{x\to 3^-}\frac{|x-3|}{x-3}&=\Lim_{x\to 3^-}\frac{3-x}{x-3}=\Lim_{x\to 3^-}(-1)=-1\\
    \Lim_{x\to 3^+}\frac{|x-3|}{x-3}&=\Lim_{x\to 3^+}\frac{x-3}{x-3}=\Lim_{x\to 3^+}(1)=1
\end{split}\]

$\because\quad \Lim_{x\to 3^-}\frac{|x-3|}{x-3}\ne \Lim_{x\to 3^+}\frac{|x-3|}{x-3}$
    
$\therefore\quad \Lim_{x\to 3}\frac{|x-3|}{x-3}$
不存在．
\end{solution}

\subsection{函数值趋于无穷大}

如果函数$f$在点$a$的邻域上有定义（可能去掉点$a$本身）对于无论多么大的正数$G$, 总存在一个够小的正数$\delta$, 使得当$0<|x-a|<\delta$时，就有$|f(x)|>G$, 那么就说当$x$趋于$a$时，函数$f(x)$趋于\textbf{无穷大}，记作
\[\lim_{x\to a} f (x) =\infty\]

\begin{example}
    求证$\Lim_{x\to 0}\frac{1}{x}=\infty$
\end{example}
   
\begin{proof}
设$G$是任意给定的正数，我们要求出一个$\delta>0$, 使得当$|x|<\delta$时，$|f(x)|=\left|\frac{1}{x}\right|>G$．

事实上，要使$\frac{1}{|x|}>G$, 
只须$0<|x|<\frac{1}{G} $．取$\delta=\frac{1}{G}$, 于是当$|x|<\delta$时，就有$\frac{1}{|x|}>G$
因此，
$$\Lim_{x\to 0}\frac{1}{x}=\infty$$
\end{proof}

\begin{example}
    证明当$x\to 0$时，函数$f(x)=\frac{1}{x}\sin\frac{1}{x}\quad (x\ne 0)$不趋于无穷大．
\end{example}

\begin{proof}
如果自变量$x$取数列$\{x_n\}=\left\{\frac{2}{(2n+1)\pi}\Big|n=1, 2, 3,\ldots\right\}$的值趋于0时，$\sin\frac{1}{x}$在原点的任意邻域内无限次交替地取$-1, 1$这两个值，对于这些值，$|f(x)|=\frac{1}{x}=\frac{\pi}{2}(2n+1)$趋于无穷大．

但是当$x$取数列$\{x'_n\}=\left\{\frac{1}{n\pi}\Big| n=1, 2, 3,\ldots\right\}$的
值趋于0时，由于
$\sin\frac{1}{x}=\sin (n\pi) =0$, 
故对于这些值，
$\Lim_{x'_n\to 0} f (x) =0$.

可见在原点的邻近不存在这样的$\delta>0$, 使得当$|x|<\delta$时，$|f(x)|>G$, 因此，当$x\to 0$时，$f(x)=\frac{1}{x}\sin\frac{1}{x}\quad (x\ne 0)$不趋于无穷大．

$y=f(x)$的图象位于两条双曲线$xy=\pm 1$之间，且在原点的邻近作无限多次振动，越靠近原点，曲线的振幅越大（图1.8）.

\begin{figure}[htp]
    \centering
\begin{tikzpicture}[>=latex, scale=.8]
\draw[->](-5,0)--(5,0)node[right]{$x$};
\draw[->](0,-4)--(0,4)node[right]{$y$};
\draw[domain=-4.5:-.3, samples=100]plot(\x,{-1/\x}); 
\draw[domain=-4.5:-.3, samples=100]plot(\x,{1/\x}); 
\draw[domain=.3:4.5, samples=100]plot(\x,{-1/\x}); 
\draw[domain=.3:4.5, samples=100]plot(\x,{1/\x}); 
\draw [domain=.5:4.5, samples=200,  thick]plot(\x, {sin(180/\x*pi)/\x});
\draw [domain=.5:4.5, samples=200,  thick]plot(-\x, {sin(180/\x*pi)/\x});
\end{tikzpicture}
    \caption{}
\end{figure}

\end{proof}

如果对于任何$G>0$, 存在$\delta>0$, 当$0<|x-a|<\delta$时，有$f(x)>G$, 就说当$x\to a$时，函数$f(x)$趋于正无穷大，记作
\[\lim_{x\to a}f(x)=+\infty\]

如果对于任何$G>0$, 存在$\delta>0$, 当$0<|x-a|<\delta$时，有$f(x)<-G$, 就说当$x\to a$时，函数$f(x)$趋于负无穷大，记作
\[\lim_{x\to a}f(x)=-\infty\]

例如，我们有：
\[\lim_{x\to 0}\frac{1}{x^2}=+\infty,\qquad \lim_{x\to 0}\frac{(-1)}{x^2}=-\infty\]

类似地，我们可以定义：
\[\lim_{x\to a^-}f(x)=+\infty,\quad \lim_{x\to a^+}f(x)=+\infty,\quad \lim_{x\to a^-}f(x)=-\infty,\quad \lim_{x\to a^+}f(x)=-\infty\]
的含义，这里不再写出．建议读者将这些定义严格地写出来．

例如，我们有：
\[\lim_{\theta\to \tfrac{\pi}{2}^-}\tan\theta =+\infty,\qquad \lim_{\theta\to \tfrac{\pi}{2}^+}\tan\theta =-\infty\]
\[\lim_{x\to 0^+}\log_a x =-\infty\quad (a>1),\qquad \lim_{x\to 0^+}\log_a x =+\infty\quad (0<a<1)\]

\begin{ex}
\begin{enumerate}
    \item 用函数极限定义证明：
    \begin{multicols}{2}
        \begin{enumerate}
            \item $\Lim_{x\to \infty}\frac{1}{2x+1}=0$
            \item $\Lim_{x\to 2}x^2=4$
            \item $\Lim_{x\to -1}\frac{x-3}{x^2-9}=\frac{1}{2}$
            \item $\Lim_{x\to 1}\frac{1}{x^2}=1$
            \item $\Lim_{x\to 1}\frac{x^3-x}{x-1}=2$
        \end{enumerate}
    \end{multicols}
    \item 说明：$\Lim_{x\to 3}\frac{x}{x^2-9}=\infty$
    \item 下面极限是否存在？
\begin{multicols}{2}
    \begin{enumerate}
        \item $\Lim_{x\to 1}\frac{2x|x-1|}{x-1}$
        \item $\Lim_{x\to 3}\frac{[x]^2-9}{x^2-9}$
    \end{enumerate}
\end{multicols}
\end{enumerate}
\end{ex}

\subsection{函数极限算法定理}
函数的极限算法定理与数列的极限算法定理类似，因为所谓$\Lim_{x\to a}u(x)=A$, $\Lim_{x\to a}v(x)=B$的意思就是对于任何一个各项都不同于$a$并且以$a$为极限的数列$x_n\to a$, 便有函数值数列$\{u(x_n)\}$, $\{v(x_n)\}$, 并且$\Lim_{n\to \infty}u(x_n)=A$, 
$\Lim_{n\to \infty} v(x_n)=B$, 因此，根据第四册下第三章的定理就可以直接得到相应的结果．现在给出函数的极限运算定理如下：

\begin{blk}{定理}
设$\Lim_{x\to a}u(x)=A$, $\Lim_{x\to a}v(x)=B$，那么
\begin{enumerate}
    \item $\Lim_{x\to a}[u(x)+v(x)]=\Lim_{x\to a}u(x)+\Lim_{x\to a}v(x)$
    \item $\Lim_{x\to a}[u(x)\cdot v(x)]=\Lim_{x\to a} u(x)\cdot \Lim_{x\to a}v(x)$
    \item $\Lim_{x\to a}[c\cdot u(x)]=c\cdot \Lim_{x\to a}u(x)$
    \item $\Lim_{x\to a}\frac{u(x)}{v(x)}=\frac{\Lim_{x\to a}u(x)}{\Lim_{x\to a}v(x)}$ 
    
    只要$v(x)$恒不为0，而且$\Lim_{x\to a}v(x)\ne 0$
    \item $\Lim_{x\to a}\sqrt[n]{u(x)}=\sqrt[n]{\Lim_{x\to a}u(x)}$
    \item 如$\Lim_{x\to a}u(x)=A=\Lim_{x\to a} v(x)$，而$|x-a|<\delta$时，$u(x)<f(x)<v(x)$，则：
    \[\Lim_{x\to a} f(x)=A\]
    即：如果$u(x)$与$v(x)$趋向同一极限$A$，且$f(x)$在$u(x)$与$v(x)$之间，那么，$f(x)$便也趋向那个极限$A$．
\end{enumerate}
\end{blk}

现在只有5需要补证．

设$\Lim_{x\to a}u(x)=A>0$, 则根据函数极限定义：对于$\varepsilon=\frac{A}{2}>0$, 存在$\delta>0$，使得$0<|x-a|<\delta$时，有
\begin{equation}
    u(x)>\frac{A}{2}>0
\end{equation}

于是，
\[\sqrt[n]{u(x)}-\sqrt[n]{A}=\frac{u(x)-A}{\sum^n_{k=1}[u(x)]^{\tfrac{n-k}{n}}A^{\tfrac{k-1}{n}}}\]
由(1.3)式
\[[u(x)]^{\tfrac{n-k}{n}}>\left(\frac{A}{2}\right)^{\tfrac{n-k}{n}}\]
故
\[\begin{split}
    \sum^n_{k=1}[u(x)]^{\tfrac{n-k}{n}}A^{\tfrac{k-1}{n}}&>\sum^n_{k=1}\left(\frac{A}{2}\right)^{\tfrac{n-k}{n}}A^{\tfrac{k-1}{n}}\\
    &=A^{\tfrac{n-1}{n}}\sum^n_{k=1}\left(\frac{1}{2}\right)^{\tfrac{n-k}{n}}\\
    &=A^{\tfrac{n-1}{n}}\left\{\left(\frac{1}{2}\right)^{\tfrac{n-1}{n}}+\left(\frac{1}{2}\right)^{\tfrac{n-2}{n}}+\cdots+\left(\frac{1}{2}\right)^{\tfrac{1}{n}}+1\right\}\\
    &=A^{\tfrac{n-1}{n}}\left\{\frac{1-\left(\frac{1}{2}\right)^{\tfrac{n-1}{n}}\cdot \left(\frac{1}{2}\right)^{\tfrac{1}{n}}}{1-\left(\frac{1}{2}\right)^{\tfrac{1}{n}}}\right\}\\
    &=A^{\tfrac{n-1}{n}}\cdot\frac{\frac{1}{2}}{1-\left(\frac{1}{2}\right)^{\tfrac{1}{n}}}=\frac{A^{\tfrac{n-1}{n}}}{2-2^{\tfrac{n-1}{n}}}\\
\end{split}\]
$\therefore\quad \left|\sqrt[n]{u(x)}-\sqrt[n]{A}\right|<|u(x)-A|\cdot \frac{2-2^{\tfrac{n-1}{n}}}{A^{\tfrac{n-1}{n}}}$

由题设$\Lim_{x\to a}u(x)=A$，而$\frac{2-2^{\tfrac{n-1}{n}}}{A^{\tfrac{n-1}{n}}}$是一个与$x$无关的常数，所以
\[\lim_{x\to a} \sqrt[n]{u(x)}=\sqrt[n]{A}=\sqrt[n]{\lim_{x\to a}u(x)}\]

下面我们来证明两个重要极限公式，为此先介绍一个引理．

\begin{blk}{引理}
     对于任意实数$\theta$, 都有$|\sin\theta|\le |\theta|$
\end{blk}

\begin{proof}
  由于对于任意实数．有$|\sin\theta|\le 1$, 那么仅考虑$\theta\in (0,\pi/2)$即可．
  
  在单位圆$O$上（图1.9），
 截取弧$\wideparen{P_0M}$使弧长$|\wideparen{P_0M}|$等于$\theta$, 其中$P_0$和$M$分别有坐标$(1, 0)$和$(x,y)$, 于是
 \[\sin\theta =y<\sqrt{y^2+ (1-x)^2} =|P_0M|<|\wideparen{P_0M}|=\theta\]
 
 显然，当$\theta=0$时，有$\sin\theta=\theta$. 因此
\[  |\sin\theta | \le |\theta|,\qquad  \theta\in  (0,\pi/2)\]

  如果$-\pi/2<\theta<0$, 则仍有$|MN|<\wideparen{P_0M}$，于是
\[|\sin\theta|<|\theta|,\qquad \theta\in (-\pi/2, 0) \]

  如果$|\theta|\ge \pi/2$, 则因为$\pi/2>1$与$|\sin\theta|\le 1$, 而同样地也得到$|\sin\theta|<|\theta|$. 
  
  因此，对于一切$\theta\ne 0$的值，有$|\sin\theta|<|\theta|$; 当$\theta=0$时，有$|\sin\theta|=|\theta|$．
\end{proof}



\begin{figure}[htp]\centering
    \begin{minipage}[t]{0.48\textwidth}
    \centering
\begin{tikzpicture}[>=latex]
\draw[->](-2.5,0)--(2.5,0)node[right]{$x$};
\draw[->](0,-2)--(0,2.5)node[right]{$y$};
\draw(0,0) circle(1.5);
\draw[very thick](0,0)--(45:1.5)node[above right]{$M(x,y)$}--(1.5,0)node[below right]{$P_0(1,0)$};
\draw[dashed](1.5/1.414,0)node[below]{$N$}--(1.5/1.414,1.5/1.414);
\node at (0,0)[below left]{$O$};
\end{tikzpicture}    
    \caption{}
    \end{minipage}
    \begin{minipage}[t]{0.48\textwidth}
    \centering
    \begin{tikzpicture}[>=latex, scale=1.3]

\draw[->](-1.5,0)--(2.5,0)node[right]{$x$};
\draw[->](0,-1.5)--(0,2)node[right]{$y$};
\draw(0,0) circle(1);
\draw[very thick](0,0)--(60:2)node[above right]{$C$}--(1,0)node[below right]{$A(1,0)$};
\draw[thick](.5,0)node[below]{$B$}--(60:1)node[above]{$D$};
\node at (0,0)[below left]{$O$};   
\draw(0.2,0) arc (0:60:.2)node[right]{$\theta$};   
    \end{tikzpicture}
    \caption{}
    \end{minipage}
    \end{figure}

\begin{blk}{定理}
\[\lim_{\theta\to 0}\frac{\sin\theta }{\theta}=1,\qquad \lim_{\theta\to 0}\frac{\cos\theta-1}{\theta}=0\]
\end{blk}

\begin{proof}
    让我们先证明$\Lim_{\theta\to 0}\frac{\sin\theta }{\theta}=1$．由图1.10容易看出：

假定$\theta$的单位是弧度，于是：
\[\begin{split}
    \triangle OBD\text{的面积}&=\frac{1}{2}OB\cdot BD=\frac{1}{2}\cos\theta\cdot \sin\theta\\
    \triangle OAC\text{的面积}&=\frac{1}{2}OA\cdot AC=\frac{1}{2}\cdot 1\cdot \tan\theta\\
    \text{扇形$OAD$的面积}&=\frac{1}{2}OA\cdot \wideparen{AD}=\frac{1}{2}\cdot 1^2\cdot \theta
\end{split}\]

假如角是用180等分平角的“度”作为单位，则扇形面积就是$\frac{1}{2}\cdot 1\cdot \frac{\pi}{180}\theta$．

因为上述扇形是夹在$\triangle OBD$和$\triangle OAC$之间，所以
\begin{equation}
    \frac{1}{2}\cos\theta\sin\theta<\frac{1}{2}\theta<\frac{1}{2}\tan\theta=\frac{1}{2}\frac{\sin\theta}{\cos\theta}
\end{equation}
即有，若$0<\theta<\frac{\pi}{2}$，则
\begin{equation}
    \cos\theta<\frac{\sin\theta}{\theta}<\frac{1}{\cos\theta}
\end{equation}
注意到：$\sin(-\theta)=-\sin\theta$和$\cos(-\theta)=\cos\theta$, 不等式(1.5)也蕴含，若$-\frac{\pi}{2}<\theta<0$, 则
\begin{equation}
    \cos(-\theta)<\frac{\sin(-\theta)}{-\theta}<\frac{1}{\cos(-\theta)}
\end{equation}
将不等式(1.5)和(1.6)合并为，若$0<|\theta|<\frac{\pi}{2}$，则
\begin{equation}
    \cos\theta<\frac{\sin\theta}{\theta}<\frac{1}{\cos\theta}
\end{equation}
又因为
\[0\le |1-\cos\theta|=\left|2\sin^2\frac{\theta}{2}\right|\le 2\left|\sin\frac{\theta}{2}\right|<|\theta|\]
从而
\[0\le \lim_{\theta\to 0}|1-\cos\theta|\le \lim_{\theta\to 0}|\theta|=0\]
所以
\[\lim_{\theta\to 0}|1-\cos\theta|=0\]
即：$\Lim_{\theta\to 0}\cos\theta=1$
由(1.7)知被夹逼在两者之间的$\frac{\sin\theta}{\theta}$极限值也就一定是1了！

以$\Lim_{\theta\to 0}\frac{\sin\theta}{\theta}=1$为基础，则：
\[\begin{split}
\lim_{\theta\to 0}\frac{\cos\theta-1}{\theta}&=\lim_{\theta\to 0}\left\{\frac{-2\left(\sin\frac{\theta}{2}\right)^2}{\theta}\right\}\\
&= \lim_{\theta\to 0}\left(-\sin\frac{\theta}{2}\right)\frac{\sin\frac{\theta}{2}}{\frac{\theta}{2}}    \\
&=\lim_{\theta\to 0}\left(-\sin\frac{\theta}{2}\right)\cdot \lim_{\theta\to 0}\frac{\sin\frac{\theta}{2}}{\frac{\theta}{2}}   \\
&=0\cdot 1=0
\end{split}   
\]
\end{proof}

\begin{example}
求：(1) $\Lim_{x\to 2}\frac{x^2+3x-10}{3x^2-5x-2}$；\qquad (2) $\Lim_{x\to 0}\frac{\sqrt[3]{1+x^2}-1}{x^2}$
\end{example}

\begin{solution}
    如果应用1.3中的定理去分别求(1)、(2)中的分子、分母的极限时，其分子、分母的极限值均为零，即我们得到$\frac{0}{0}$的不定形式．我们通常要对原代数式变形，找出分子、分母中具有$x-a$的因子，消去$x-a$的公共因子，一般问题就可解决．

(1)
\[\begin{split}
    \Lim_{x\to 2}\frac{x^2+3x-10}{3x^2-5x-2}&=\Lim_{x\to 2}\frac{(x+5)(x-2)}{(3x+1)(x-2)}\\
    &=\Lim_{x\to 2}\frac{x+5}{3x+1}\\
    &=\frac{2+5}{3\x 2+1}=1
\end{split}\]

(2)
\[\begin{split}
    \Lim_{x\to 0}\frac{\sqrt[3]{1+x^2}-1}{x^2}&=\Lim_{x\to 0}\frac{{(1+x^2)}-1}{x^2\left(\sqrt[3]{(1+x^2)^2}+\sqrt[3]{1+x^2}+1\right)}\\
    &=\Lim_{x\to 0}\frac{1}{\sqrt[3]{(1+x^2)^2}+\sqrt[3]{1+x^2}+1}=\frac{1}{3}
\end{split}\]
\end{solution}

\begin{example}
    求：(1) $\Lim_{x\to 0}\frac{\cos x-1}{x^2}$；\qquad (2) $\Lim_{x\to \pi}\frac{\sin mx}{\sin nx}\;(m,n\in \mathbb{N})$
\end{example}

\begin{solution}
(1)
\[\begin{split}
    \Lim_{x\to 0}\frac{\cos x-1}{x^2}&=\Lim_{x\to 0}\frac{-2\left(\sin\frac{x}{2}\right)^2}{x^2}\\
    &=\frac{-2}{2^2}\Lim_{x\to 0}\frac{\sin\frac{x}{2}}{\frac{x}{2}}\cdot \Lim_{x\to 0}\frac{\sin\frac{x}{2}}{\frac{x}{2}}\\
    &=\left(-\frac{1}{2}\right)\cdot 1\cdot 1=-\frac{1}{2}
\end{split}\]

(2) 由三角函数诱导公式易知
\[\sin m(\pi-x)=\sin(m\pi-mx)=(-1)^{m-1}\sin mx\]
同样得到
\[\sin n(\pi-x)=(-1)^{n-1}\sin nx\]
因此：
\[(-1)^{m-n}\frac{\sin mx}{\sin nx}=\frac{\sin m(\pi-x)}{\sin n(\pi-x)}\]
两边乘以$(-1)^{m-n}$，得
\[\frac{\sin mx}{\sin nx}=(-1)^{m-n}\frac{\sin m(\pi-x)}{\sin n(\pi-x)}\]
设$y=\pi-x$, 则当$x\to \pi$, 有$y\to 0$, $my\to 0$, $ny\to 0$. 因此：
\[\begin{split}
    \Lim_{x\to \pi}\frac{\sin mx}{\sin nx}&=\lim_{y\to 0}(-1)^{m-n}\frac{\sin my}{\sin ny}\\
&=(-1)^{m-n}\lim_{y\to 0}\left(\frac{m}{n}\cdot \frac{\sin my}{my}\cdot \frac{ny}{\sin ny}\right)\\
&=(-1)^{m-n}\cdot \frac{m}{n}\cdot 1\cdot 1=(-1)^{m-n}\frac{m}{n}
\end{split}\]
\end{solution}

\begin{example}
由条件$\Lim_{x\to -\infty}\left(\sqrt[]{x^2-x+1}-ax-b\right)=0$，求出$a,b$的值．
\end{example}

\begin{solution}
$\because\quad \Lim_{x\to -\infty}\left(\sqrt{x^2-x+1}-ax-b\right)=0,\quad \Lim_{x\to -\infty}\frac{1}{x}=0$

$\therefore\quad \Lim_{x\to -\infty}\left(\sqrt{x^2-x+1}-ax-b\right)\cdot \frac{1}{x}=0$

即：
\[\Lim_{x\to -\infty}\left(\frac{\sqrt{x^2-x+1}}{x}-a-\frac{b}{x}\right)=0\]
又因为：当$x<0$时，
\[\sqrt{x^2-x+1}=\sqrt{x^2\left(1-\frac{1}{x}+\frac{1}{x^2}\right)}=-x\sqrt{1-\frac{1}{x}+\frac{1}{x^2}}\]
所以上式可写成
\[\lim_{x\to -\infty}\left(-\sqrt{1-\frac{1}{x}+\frac{1}{x^2}}-a-\frac{b}{x}\right)=0\]
即：$-1-a-0=0\quad \to \quad a=-1$

代入原条件，得
\[\lim_{x\to -\infty}\left(\sqrt{x^2-x+1}+x-b\right)=0\]
即：
\[\begin{split}
    b&=\lim_{x\to-\infty}\left(\sqrt{x^2-x+1}+x\right)\\
    &=\lim_{x\to-\infty} \frac{x^2-x+1-x^2}{\sqrt{x^2-x+1}-x}  \\
    &=\lim_{x\to-\infty}\frac{-x\left(1-\frac{1}{x}\right)}{-x\left(\sqrt{1-\frac{1}{x}+\frac{1}{x^2}}+1\right)}    \\
    &=\lim_{x\to-\infty}\frac{1-\frac{1}{x}}{\sqrt{1-\frac{1}{x}+\frac{1}{x^2}}+1} =\frac{1}{2}   \\
\end{split}\]

$\therefore\quad a=-1,\quad b=\frac{1}{2}$为所求．
\end{solution}

\begin{ex}
\begin{enumerate}
    \item 说明当$x\to \infty$时，函数$\frac{x^5-4}{x^3+x}$的变化趋势．
    \item 求极限
    \begin{enumerate}
\begin{multicols}{2}
    \item  $\Lim _{x \to 0} \frac{(x+2)^{2}}{x^{2}+4}$
    \item  $\Lim _{x \to 2} \frac{x^{2}-4}{x-2}$
    \item  $\Lim _{x \to 2} \frac{5 x^{2}+x-8}{x^{2}-4}$
    \item  $\Lim _{x \to 0} \frac{\left(4 x^{3}-3\right)(1-2 x)}{7 x^{3}-6 x+4}$
    \item  $\Lim _{x \to \infty} \frac{3 x^{5}}{x^{5}-x^{2}+1}$
    \item  $\Lim _{x \to \infty} \frac{\left(x^{2}-5\right)\left(x^{2}+7\right)}{x^{4}+35}$
    \item  $\Lim _{x \to 3} \frac{x^{2}-8 x+15}{x^{2}-7 x+12}$
    \item  $\Lim _{x \to -3} \frac{x^{2}-9}{x^{2}+9 x+18}$
    \item $\Lim _{x \to -1} \frac{x\left(x^{2}+4 x+3\right)}{x^{3}+3 x^{2}+5 x+3}$
    \item  $\Lim _{x \to 1} \frac{x^{3}+x^{2}-2}{x^{4}+2 x^{2}-2 x-1}$
    \item $\Lim _{x \to \infty} \frac{x+\sin x}{2 x+5}$
    \item  $\Lim _{x \to 0} \frac{(1+x)^{5}-(1+5 x)}{x^{2}+x^{5}}$
\end{multicols}
    \item  $\Lim _{x \to 0} \frac{(1+m x)^{n}-(1+n x)^{m}}{x^{2}}\qquad (m, n \in \mathbb{N})$
    \item  $\Lim _{x \to 1} \frac{x^{2}+x^{4}+\cdots+x^{2 n}-n}{x-1}$
\end{enumerate}

    \item 求极限  
\begin{enumerate}
\begin{multicols}{2}
\item $\Lim_{x\to 0^+}\sqrt{3x} $
\item $\Lim_{x\to 1^-}\sqrt{1-x} $
\item $\Lim_{x\to 3^-} 3[x]$
\item $\Lim_{x\to 1^+} [x+3]$
\item $\Lim_{x\to -3^+}\left(1+\sqrt{x+3}\right) $
\item $\Lim_{x\to 4^-} \left(\sqrt{4-x}+[x-1]\right)$
    \item $\Lim_{x\to 1^+}\frac{2x|x-1|}{x-1} $
    \item $\Lim_{x\to 4^-} ([x]-x)$
\end{multicols}
    \item $\Lim_{x\to 0^+}\frac{b}{x}\left[\frac{x}{a}\right] \qquad (a>0,\; b>0)$
    \item $\Lim_{x\to 0^+} \frac{x}{a}\left[\frac{b}{x}\right]\qquad (a>0,\; b>0)$
\end{enumerate}    

    \item 求极限  
\begin{enumerate}
\begin{multicols}{2}
\item $\Lim _{x \to  1} \frac{1-x}{\sqrt{1-x^{2}}}$;
\item  $\Lim _{x \to  1} \frac{1-\sqrt{x}}{1-x}$;
\item  $\Lim _{x \to  1} \frac{(2 x-3)(\sqrt{x}-1)}{2 x^{2}+x-3}$
\item  $\Lim _{x \to  2} \frac{x-2}{\sqrt{x^{2}-2}-\sqrt{2}}$
\item  $\Lim _{t \to +\infty}(\sqrt{1+t}-\sqrt{t})$;
\item  $\Lim _{x \to +\infty} \sqrt{x}(\sqrt{x+a}-\sqrt{x})$
\item  $\Lim _{x \to  5} \frac{\sqrt{x-1}-2}{x-5}$;
\item $\Lim _{x \to  0} \frac{\sqrt[3]{1+x}-\sqrt[3]{1-x}}{x}$
\item  $\Lim _{x \to  1} \frac{4 x+\sqrt{x-1}}{2 x-\sqrt{x+1}}$
\item  $\Lim _{x \to  1} \frac{x-1}{\sqrt{x^{2}-1}+\sqrt{x-1}}$
\item $\Lim _{x \to  a} \frac{\sqrt{x-a}+\sqrt{x}-\sqrt{a}}{\sqrt{x^{2}-a^{2}}}$
\item  $\Lim _{x \to -\infty} \frac{x-2}{\sqrt{x^{2}-4 x+3}}$
\item $\Lim _{x \to -\infty}\left[\frac{2(x-1)^{2}}{\sqrt{4 x^{2}+2 x+1}}+x\right]$
\end{multicols}
\item $\Lim _{x \to +\infty}\left(\sqrt{x+\sqrt{x+\sqrt{x}}}-\sqrt{x}\right)$
\item $\Lim _{x \to +\infty} x\left(\sqrt{x^{2}+2 x}-2 \sqrt{x^{2}+x}+x\right)$
\end{enumerate}   

\item \begin{enumerate}
    \item $f(x)=\begin{cases}
        0,& x>1\\
        1,& x=1\\
        x^2+2, &x<1
    \end{cases}$

    求$f(x)$在$x=1$的左右极限．

    \item $f(x)=\begin{cases}
        x\sin\frac{1}{x},& x>0\\
        1+x^2,& x<0
    \end{cases}$

    求$f(x)$在$x=0$的左右极限．

    \item $\Lim_{x\to 3}\frac{[x]^2-9}{x^2-9}$是否存在？
    \item $\Lim_{x\to 0}\frac{x}{|x|}$是否存在？
\end{enumerate}

\item 若$0<x<\frac{\pi}{2}$，求证：
\begin{multicols}{2}
    \begin{enumerate}
        \item $\tan x>x$
        \item $x-\frac{x^2}{4}<\sin x<x$
    \end{enumerate}
\end{multicols}

\item 求极限  
\begin{enumerate}
\begin{multicols}{2}
    \item $\Lim_{x\to 0} \frac{\sin ax}{x}$
    \item $\Lim_{x\to 0} \frac{\sin 7x}{4x}$
    \item $\Lim_{\theta\to 0} \frac{\sec2\theta\sin\theta}{\theta}$
    \item $\Lim_{\theta \to 0} \frac{\tan\theta}{\theta}$
    \item $\Lim_{t\to 0} \frac{15t}{\tan 6t}$
    \item $\Lim_{x\to 0} \frac{\tan 2x}{\sin 7x}$
    \item $\Lim_{\theta \to 0} \frac{\tan\theta-\sin\theta}{\theta^3}$
    \item $\Lim_{x\to 0} \frac{\csc x-\cot x}{x}$
    \item $\Lim_{x\to 0} \frac{\sin^2\frac{x}{2}}{x^2}$
    \item $\Lim_{x\to \pi} \frac{\sin x}{\pi-x}$
    \item $\Lim_{x\to 1} \frac{1+\cos\pi x}{\tan^2 \pi x}$
    \item $\Lim_{x\to \tfrac{\pi}{4}}\tan 2x\tan\left(\frac{\pi}{4}-x\right) $ 
    \item $\Lim_{x\to 1} (1-x)\tan\frac{\pi x}{2}$
    \item $\Lim_{x\to 0} \frac{\sin mx}{\sin nx}$
    \item $\Lim_{x\to 0} \frac{x+\sin x}{x+2\sin x}$
    \item $\Lim_{x\to 0} \frac{\sqrt{\cos x}-1}{x^2}$
\end{multicols}
\end{enumerate}   

\end{enumerate}
\end{ex}

\subsection{数$e$}
我们在这里要利用数列的极限来定义一个新的数，这一个数不论对于分析本身或者对于它的应用来说都是非常重要的．

考虑数列$a_n=\left(1+\frac{1}{n}\right)^n,\; n=1, 2, 3,\ldots$的极限．

首先我们计算一下数列$\{a_n\}$的数值：
\begin{center}
\begin{tabular}{llll}
   $a_1=2.0$ & $a_2=2.250$ & $a_3=2.370$& $a_4=2.441$\\
$a_5=2.488$ &$a_6=2.522$& $a_7=2.546$& $a_8=2.565$\\
$\cdots$ &$a_{256}=2.712$&$\cdots$&$a_{1024}=2.717$\\ 
\end{tabular}
\end{center}

观察这一串数，我们会猜测这个数列是递增的，可能收敛于一个极限值，现在我们来证明这个猜想成立．

依牛顿二项式定理，即有
\[\begin{split}
    \left(1+\frac{1}{n}\right)^n&=1+n\cdot\frac{1}{n}+\frac{n(n-1)}{1\cdot 2}\cdot \frac{1}{n^2}+\frac{n(n-1)(n-2)}{1\cdot 2\cdot 3}\cdot \frac{1}{n^3}\\
    &\quad +\cdots+\frac{n(n-1)(n-2)\cdots[n-(n-1)]}{1\cdot 2\cdot 3\cdots n}\cdot \frac{1}{n^n}
\end{split}\]
这个等式的右边的第二项等于1，其他各项可以变换如下：
\[\begin{split}
\frac{n(n-1)}{1\cdot 2}\cdot \frac{1}{n^2}&=\frac{n(n-1)}{n^2}\cdot\frac{1}{2!}\\&=\left(1-\frac{1}{n}\right) \cdot\frac{1}{2!}\\
\frac{n(n-1)(n-2)}{1\cdot 2\cdot 3}\cdot \frac{1}{n^3}&=\frac{n(n-1)(n-2)}{n^3}\cdot\frac{1}{3!}\\&=\left(1-\frac{1}{n}\right)\cdot \left(1-\frac{2}{n}\right)\cdot\frac{1}{3!}\\
\cdots & \cdots\cdots\\
\frac{n(n-1)(n-2)\cdots[n-(n-1)]}{1\cdot 2\cdot 3\cdots n}\cdot \frac{1}{n^n}&=\frac{n(n-1)(n-2)\cdots[n-(n-1)]}{n^n}\cdot \frac{1}{n!}\\
&=\left(1-\frac{1}{n}\right) \left(1-\frac{2}{n}\right)\cdots\left(1-\frac{n-1}{n}\right)\cdot\frac{1}{n!}
\end{split}\]
其中$n!=1\cdot 2\cdot 3\cdots n$．因此：
\begin{equation}
\begin{split}
\left(1+\frac{1}{n}\right)^n&=1+1+\left(1-\frac{1}{n}\right)\frac{1}{2!}+ \left(1-\frac{1}{n}\right)\left(1-\frac{2}{n}\right)\cdot\frac{1}{3!}\\
&+\cdots+ \left(1-\frac{1}{n}\right) \left(1-\frac{2}{n}\right)\cdots\left(1-\frac{n-1}{n}\right)\cdot\frac{1}{n!}  
\end{split}
\end{equation}

用$n+1$代替$n$代入此公式，求出：
\begin{equation}
\begin{split}
    \left(1+\frac{1}{n+1}\right)^{n+1}&=1+1+\left(1-\frac{1}{n+1}\right)\frac{1}{2!}+ \left(1-\frac{1}{n+1}\right)\left(1-\frac{2}{n+1}\right)\cdot\frac{1}{3!}\\
    &+\cdots+ \left(1-\frac{1}{n+1}\right) \left(1-\frac{2}{n+1}\right)\cdots\left(1-\frac{n}{n+1}\right)\cdot\frac{1}{(n+1)!}  
\end{split}
\end{equation}
(1.9)的展开式项数比(1.8)的展开式的项数多一项，此外(1.9)的展开式各项由第三项起大于(1.8)的展开式对应项，这是因为
\[1-\frac{1}{n+1}>1-\frac{1}{n},1-\frac{2}{n+1}>1-\frac{2}{n},\ldots, 1-\frac{n-1}{n+1}>1-\frac{n-1}{n}\]
所以
\[\left(1+\frac{1}{n+1}\right)^{n+1}>\left(1+\frac{1}{n}\right)^n\]
这就是说数列$a_n=\left(1+\frac{1}{n}\right)^n,\; n=1,2,3,\ldots$是递增的．

但是由等式(1.8)可知
\[\left(1+\frac{1}{n}\right)^n<1+1+\frac{1}{2!}+\frac{1}{3!}+\cdots+\frac{1}{n!}\]
又因为
\[\frac{1}{1\cdot 2\cdot 3}<\frac{1}{2^2},\frac{1}{1\cdot 2\cdot 3\cdot 4}<\frac{1}{2^3},\ldots,\frac{1}{1\cdot 2\cdot 3\cdots n}<\frac{1}{2^{n-1}}\]
所以：
\[\left(1+\frac{1}{n}\right)^n<2+\frac{1}{2}+\frac{1}{2^2}+\cdots+\frac{1}{2^{n-1}}\]
但是由第二项起的级数总和小于$\frac{\tfrac{1}{2}}{1-\tfrac{1}{2}}=1$，所以：
\[\left(1+\frac{1}{n}\right)^n<2+\frac{\frac{1}{2}}{1-\frac{1}{2}}=3\]
这就是说数列$\left\{\left(1+\frac{1}{n}\right)^n\right\}$是递增有上界的，所以它有极限，设它的极限用$e$表示，因而我们有数
\[e=\Lim_{n\to \infty} \left(1+\frac{1}{n}\right)^n\]
这个数$e$是无理数，取它的十进小数到十五位，其值是
\[e=2.718281828459045\cdots\]

现在要把变数$n$推广到实数$x$, 这是一个以后常要用到的重要极限．

\begin{blk}{定理}
\[\Lim_{x\to \infty} \left(1+\frac{1}{x}\right)^x=e\qquad (x\in\mathbb{R})\] 
\end{blk}

\begin{proof}
先证$\Lim_{x\to +\infty} \left(1+\frac{1}{x}\right)^x=e$

令$[x]=n$，于是$n\le x<n+1$，所以，当$x\to +\infty$时，$n\to +\infty$，而且
\[1+\frac{1}{n+1}<1+\frac{1}{x}\le 1+\frac{1}{n}\]
根据指数函数的单调性和幂函数在$x>0$半直线上的单调性，有：
\[\left(1+\frac{1}{n+1}\right)^n\le \left(1+\frac{1}{n+1}\right)^x<\left(1+\frac{1}{x}\right)^x\le \left(1+\frac{1}{n}\right)^x<\left(1+\frac{1}{n}\right)^{n+1}\]
因为：
\[\begin{split}
    \Lim_{n\to +\infty} \left(1+\frac{1}{n}\right)^{n+1}&=\Lim_{n\to +\infty} \left(1+\frac{1}{n}\right)^{n}\left(1+\frac{1}{n}\right)\\
&=\Lim_{n\to +\infty} \left(1+\frac{1}{n}\right)^{n}\cdot \Lim_{n\to +\infty} \left(1+\frac{1}{n}\right)=e
\end{split}\]
\[\begin{split}
    \Lim_{n\to +\infty} \left(1+\frac{1}{n+1}\right)^{n}&=\Lim_{n\to +\infty} \frac{\left(1+\frac{1}{n+1}\right)^{n+1}}{1+\frac{1}{n+1}}\\
    &=\frac{\Lim_{n\to +\infty}\left(1+\frac{1}{n+1}\right)^{n+1} }{\Lim_{n\to +\infty}\left(1+\frac{1}{n+1}\right) }=e
\end{split}\]
所以
\[e\le\Lim_{x\to +\infty}\left(1+\frac{1}{x}\right)^{x} \le e\]
即：$\Lim_{x\to +\infty}\left(1+\frac{1}{x}\right)^{x}=e$

当$x\to -\infty$时，令$x=-y$，那么$y\to +\infty$，这时
\[\begin{split}
    \left(1+\frac{1}{x}\right)^x&=\left(1-\frac{1}{y}\right)^{-y}=\left(\frac{y-1}{y}\right)^{-y}=\left(\frac{y}{y-1}\right)^y\\
    &=\left(1+\frac{1}{y-1}\right)^y=\left(1+\frac{1}{y-1}\right)^{y-1}\left(1+\frac{1}{y-1}\right)
\end{split}\]
所以
\[\begin{split}
    \Lim_{x\to -\infty}\left(1+\frac{1}{x}\right)^{x}&=\Lim_{y\to +\infty}\left(1+\frac{1}{y-1}\right)^{y-1}\left(1+\frac{1}{y-1}\right)\\
    &=\Lim_{y\to +\infty}\left(1+\frac{1}{y-1}\right)^{y-1}\cdot \Lim_{y\to +\infty}\left(1+\frac{1}{y-1}\right)=e
\end{split}\]
综合上面两个结果，即得
$$\Lim_{x\to \infty}\left(1+\frac{1}{x}\right)^{x}=e$$
\end{proof}

\begin{blk}{推论}
\[\Lim_{y\to 0}\left(1+y\right)^{\tfrac{1}{y}}=e\]
\end{blk}

事实上，令$x=\frac{1}{y}$，当$y\to 0$时，$x\to\infty$，于是：
\[\Lim_{y\to 0}\left(1+y\right)^{\tfrac{1}{y}}=\Lim_{x\to \infty}\left(1+\frac{1}{x}\right)^{x}=e\]

\begin{ex}
\begin{enumerate}
    \item 求下面变量的极限：
\begin{multicols}{2}
\begin{enumerate}
    \item $\Lim_{x\to\infty}\left(1+\frac{8}{x}\right)^x$
    \item $\Lim_{y\to 0}(1+y)^{\tfrac{1}{3y}}$
    \item $\Lim_{t\to\infty}\left(1-\frac{1}{t}\right)^t$
    \item $\Lim_{t\to\infty}\left(\frac{t}{1+t}\right)^t$
    \item $\Lim_{n\to+\infty}\left(1+\frac{x}{n}\right)^n$
    \item $\Lim_{n\to +\infty}\left(1+\frac{1}{n^2}\right)^n$
    \item $\Lim_{x\to+\infty}\left(\frac{2x+3}{2x+1}\right)^{x+1}$
\end{enumerate}
\end{multicols}

\item 证明：
\begin{multicols}{2}
\begin{enumerate}
    \item $\Lim_{x\to\infty}\frac{a^x}{x}=\infty$
    \item $\Lim_{x\to\infty}\frac{a^x}{x^a}=\infty$
\end{enumerate}
\end{multicols}
其中$a$为任意正常数．
\end{enumerate}
\end{ex}

\section{函数的连续性}
\subsection{连续函数的概念}

在第四册下，我们用数列的语言规定了函数$f$在点$a$连续的意义．现在用$\varepsilon-\delta$语言来描述函数$f$在其定义域中的点$a$是连续的含义：

\begin{blk}
   {定义1} 函数$f$在其定义域M中的点$a$是连续的，如果对于每一个正数$\varepsilon$, 有$a$的一个$\delta$邻域$(a-\delta,a+\delta)$, 使得对于定义域中满足$|x-a|<\delta$的一切$x$值，不等式
\[|f (x) -f(a)|<\varepsilon\]
成立． 
\end{blk}

\begin{rmk}
在函数极限定义那里不要求$f(a)$存在，所以特别注意要$x\ne a$, 现在，在函数于某点$a$连续的定义中，只盼$\Lim_{x\to a} f(x)=f(a)$, 因此，保证$|f(x)-f(a)|<\varepsilon$
成立的充分条件$|x-a|<\delta$不再要求$x\ne a$了．
\end{rmk}

对于一个在点$a$处连续的函数，有
\[\lim_{x\to a}f (x) =f \left(\lim_{x\to a} x\right) =f (a)\]  
极限符号同函数符号可以改变次序或者说可以互换．

\begin{blk}
 {定义2 }如果函数$f(x)$在一个区间$(a,b)$的每一点都连续，我们就说，它在这个区间$(a,b)$上连续．   
\end{blk}

函数也可以只在一点的一侧连续，当
\[\lim_{x\to x_0^+} f (x) =f (x_0)\]
时，便说$f(x)$在$x_0$的右侧连续；当
\[\lim_{x\to x_0^-} f (x) =f (x_0)\]
时，便说$f(x)$在$x_0$的左侧连续．显然，在一点的两侧都连续时，必在这点上连续．

所谓函数$f(x)$在闭区间$[a,b]$上连续，除了它在$(a,b)$内每一点都连续外，在$a$点右侧连续，在$b$点左侧连续．

在一般自然现象中连续是常态，不连续则是在特殊情形的突变．例如，在火箭的发射过程中，随着火箭的燃烧，质量逐渐变化，当每一级火箭烧尽时，该级火箭的壳自行脱落，于是质量突然减小，质量变化如图1.11所示．所以，常用的函数一般只有几个特别的点是不连续的：如果$f(x)$在$x=a$处不是连续的，则称$a$为函数的间断点．对于左右极限存在的间断点叫\textbf{第一类间断点}，不属于第一类间断点的叫\textbf{第二类间断点}．

例如$y=[x]$在所有的整数点$n$有第一类间断点；$y=\tan x$在$x=\pi/2$有第二类间断点．

\begin{figure}[htp]
    \centering
    \begin{minipage}[t]{0.48\textwidth}
    \centering
    \begin{tikzpicture}[>=latex, scale=1]
        \draw[->](-.5,0)--(4,0)node[right]{$t$时间};
        \draw[->](0,-.5)--(0,4)node[right]{$M$质量};
    \draw[thick](0,3.5)--(2,2.5);
    \draw[thick](2,1.5)--(4,.5);
    \draw[dashed](2,2.5)--(2,0)node[below]{$t_0$};
    \draw(2,2.5)[fill=white] circle (2pt);
    \node at (0,0)[below right]{$O$};
    \end{tikzpicture}
    \caption{}
    \end{minipage}
    \begin{minipage}[t]{0.48\textwidth}
    \centering
    \begin{tikzpicture}[>=latex, scale=1.8]
        \draw[->](-.5,0)--(2.5,0)node[right]{$x$};
    \draw[->](0,-.5)--(0,2)node[right]{$y$};
    \draw[thick](0,0)node[below left]{$O$}--(1,1)--(2,0)node[below]{2};
    \draw[dashed](1,0)node[below]{1}--(1,1);
    \draw(0,1)node[left]{1}--(.05,1);
    \end{tikzpicture}
    \caption{}
    \end{minipage}
    \end{figure}

\begin{example}
    设函数
$f (x) =\begin{cases}
     x,& 0\le x\le 1 \\ 2-x,& 1\le x\le 2\\
\end{cases}$

问在$x=1$处是否连续？
\end{example}

\begin{solution}
  在$x=1$处的两侧，函数的对应规则由不同的式子表示，所以研究在$x=1$处函数是否连续，就需要从左右极限入手（图1.12）．
\[  \lim_{x\to 1^-} f (x) =\lim_{x\to 1^-} x=1,\qquad  \lim_{x\to 1^+} f (x) =\lim_{x\to 1^+} (2-x)=1\]
$\therefore\quad \Lim_{x\to 1}f(x)=1$

又$f(1)=1$, 故有$\Lim_{x\to 1}f(x)=f(1)$, 即$f(x)$在$x=1$处是连续的．
\end{solution}

\begin{example}
    函数$f(x)=x\sin\frac{1}{x}$在$x=0$处是否连续？如不连续，是第几类间断点？
\end{example}

\begin{solution}
    函数$f$只在$x=0$没有定义，故$f$在$x=0$处间断，它的定义域是$\mathbb{R}-\{0\}$.

因为$\left|\sin\frac{1}{x}\right|\le 1$，从而：    
\[\begin{split}
   & 0\le\left|x\sin\frac{1}{x}\right|\le |x|\\
&0\le \lim_{x\to 0}\left|x\sin\frac{1}{x}\right|\le \lim_{x\to 0}|x|=0\\
&\lim_{x\to 0}x\sin\frac{1}{x}=0
\end{split}\]
所以，$f$在$x=0$处左、右极限存在，这就是说$x=0$是第一类间断点，它的图象大致如图1.13所示．

\begin{figure}[htp]
    \centering
\begin{tikzpicture}[>=latex, scale=1.5]
\draw[->](-1.75,0)--(1.75,0)node[right]{$x$};
\draw[->](0,-1.75)--(0,1.75)node[right]{$y$};
\draw(-1.5,-1.5)--(1.5,1.5)node[right]{$y=x$};
\draw(-1.5,1.5)node[left]{$y=-x$}--(1.5,-1.5);
\draw[domain=0.08:1.5, samples=1000, thick]plot(\x, {\x*sin(1/\x r)});
\draw[domain=0.08:1.5, samples=1000, thick]plot(-\x, {\x*sin(1/\x r)});
\node at (0,0)[below left]{$O$};
\draw[dashed](2/pi,0)node[below]{$\frac{2}{\pi}$}--(2/pi,2/pi)--(-2/pi,2/pi);
\end{tikzpicture}
    \caption{}
\end{figure}


如果把上面的函数$f$加以扩充，即我们定义一个新函数$g$:
\[g(x)=\begin{cases}
    x\sin\frac{1}{x},& x\ne 0\\
    0,& x=0
\end{cases}\]
那么我们便得到一个在$x=0$处的连续函数．
\end{solution}

\begin{blk}
    {定义3} $\Lim_{x\to a}f(x)$存在，但不等于$f(a)$, 或者$f(a)$没有定义，则称$f$在$a$处有\textbf{可去间断点}．
\end{blk}

例1.16中的$x=0$就是$f$的可去间断点．

由函数在某一点连续的概念可以直接导出连续函数的一个重要性质如下：

\begin{blk}
{定理} 若函数$f(x)$在$x=a$连续，且$f(a)>0$, 则存在一个$a$的$\delta$邻域$(a-\delta,a+\delta)$, 使得对于每一个$x\in (a-\delta,a+\delta)$, $f(x)$都是正的．    
\end{blk}

\begin{proof}
    因为$\Lim_{x\to a} f(x)=f(a)>0$, 所以对于
$\varepsilon=\frac{1}{2}f(a)$, 必定存在一个$\delta$使得$|x-a|<\delta$时，有
\[f (x) > f (a) -\frac{1}{2}f (a) =\frac{1}{2}f (a) > 0\]
\end{proof}

\subsection{连续函数的运算}
连续函数的概念是函数概念和极限概念一个自然结合，由本章1.3函数极限算法定理直接得出下面定理．

\begin{blk}
 {定理1} 设$f$和$g$在$a$处连续，则
 \begin{enumerate}
     \item $f\pm g$在$a$处连续；
     \item $f\cdot g$在$a$处连续；
     \item $\frac{1}{g}$在$a$处连续（若$g(a)\ne 0$）；
     \item $\sqrt[n]{f}$在$a$处连续（若$f(a)>0$）．
 \end{enumerate}
\end{blk}

\begin{proof}
\begin{enumerate}
    \item 因为$f$和$g$在$a$处连续，所以
    \[\Lim_{x\to a} f (x) =f (a) ,\qquad  \lim_{x\to a} g (x) =g (a)\]
    根据1.3中的函数极限算法定理，得
   \[\begin{split}
    \Lim_{x\to a} (f\pm g)(x)=\Lim_{x\to a}[f(x)\pm g(x)]&=\Lim_{x\to a} f(x)\pm \Lim_{x\to a}g(x)\\
    &=f(a)\pm g(a)=(f\pm g)(a)
   \end{split} \]
    这正好就是断言$f+g$在$a$处连续．
    \item \[\begin{split}
        \Lim_{x\to a} (f\cdot g) (x)=\Lim_{x\to a}f(x)\cdot g(x)&=\Lim_{x\to a} f(x)\cdot\Lim_{x\to a}g (x)\\
        &=f (a) \cdot g (a) =(f\cdot g) (a) 
    \end{split}\]
    即$f\cdot g$在$a$处连续．
    \item 由于$g(a)\ne 0$, 又$g$在$a$处连续，根据2.1中的定理，知道存在点$a$的$\delta$邻域，使得对于每一个$x\in (a-\delta,a+\delta)$，有$g(x)\cdot g(a)>0$，因此
\[\Lim_{x\to a} \frac{1}{g}(x)=\frac{1}{\Lim_{x\to a} g(x)}=\frac{1}{g(a)}=\frac{1}{g}(a)\]

\item 证明留给读者．
\end{enumerate} 
\end{proof}

在第四册下，我们根据定理1证明了一些基本初等函数在其定义域上到处连续的命题，现在总结如下：

\begin{blk}
   {命题1} 任何多项式函数到处连续． 
\end{blk}

\begin{blk}{命题2}
    若$f$和$g$是两个多项式，$g\ne 0$, 那么有理函数$r=f/g$, 除去$g$的零点集合，函数$r$是有定义的，而且在其定义域上到处连续．
\end{blk}

\begin{blk}
   {命题3} $n$次算术方根函数$\sqrt[n]{x}$在$[0,+\infty)$上到处连续．
\end{blk}

因为每一个无理数可以用两串有理数数列从左、右两方面去逼近它，同样无理指数幂也可以用相应的有理指数幂去逼近它，我们用逼近法定义了无理指数幂，从而把有理指数函数开拓为在实数域上到处连续的实指数函数，于是得到下面的命题． 

\begin{blk}
    {命题4} 指数函数$a^x\; (a>0,\; a\ne 1)$到处连续．
\end{blk}

由反函数连续定理，直接得到

\begin{blk}
    {命题5} 对数函数$\log_a x$在$(0,+\infty)$上到处连续．
\end{blk}

现在将继续证明另外一些基本初等函数连续性的命题．

\begin{blk}
    {命题6} 三角函数$\sin x$和$\cos x$到处连续．
\end{blk}

\begin{proof}
设$x_0$是任意给定的实数，由三角函数的和差化积公式，得到
\[|\sin x-\sin x_0|=2\left|\sin\frac{x-x_0}{2}\right| \cdot \left|\cos \frac{x+x_0}{2}\right|\]
因为$\left|\cos\frac{x+x_0}{2}\right|\le 1$, 并根据1.3中的引理得到
\[|\sin x-\sin x_0|\le 2\left|\frac{x-x_0}{2}\right|\]
即
\[|\sin x-\sin x_0|\le \left|x-x_0\right|\]
当$x\to x_0$时，$|x-x_0|\to 0$, 从而$|\sin x-\sin x_0|\to 0$, 即
\[\lim_{x\to x_0} \sin x=\sin x_0\]
因为$x_0$是任意一点，所以$\sin x$到处连续．

同样证明，得到：当$x\to x_0$时，
\[|\cos x-\cos x_0|\le |x-x_0|\to 0\]
即
\[\lim_{x\to x_0} \cos x=\cos x_0\]
因此，$\cos x$到处连续．
\end{proof}
 
由反函数连续性定理直接得出下面命题．

\begin{blk}
    {命题7} 反三角函数$y=\arcsin x\; \left(-\frac{\pi}{2}\le y\le \frac{\pi}{2}\right)$,\quad $y=\arccos x\; \left(0\le y\le \pi\right)$,\quad $y=\arctan x\; \left(-\frac{\pi}{2}< y< \frac{\pi}{2}\right)$,\quad $y=\arccot x\; \left(0<y<\pi\right)$在它们各自的定义域上到处连续．
\end{blk}

\begin{blk}{命题8}
幂函数$x^{\alpha}$（$\alpha$是任何一个实数）在开区间$(0,+\infty)$上到处连续．
\end{blk}

\begin{proof}
为确定起见，假设$\alpha>0$, $n$是任何一个大于$\alpha$的整数
\[(x+h)^{\alpha} -x^{\alpha}=x^{\alpha} \left[ \left(1+\frac{h}{x}\right)^{\alpha} -1\right] \]

若$h>0$, 则$\left(1+\frac{h}{x}\right)^{\alpha} >1$，（实指数幂性质），并且
\begin{equation}
    1<\left(1+\frac{h}{x}\right)^{\alpha}<\left(1+\frac{h}{x}\right)^{n}\quad \text{(指数函数单调性)}
\end{equation}

若$h<0$, 将上面不等式方向反向，得到
\begin{equation}
    \left(1+\frac{h}{x}\right)^{n}<\left(1+\frac{h}{x}\right)^{\alpha}<1
\end{equation}
因为$\left(1+\frac{h}{x}\right)^{n}$是关于$h$的$n$次多项式，并且显然当$h\to 0$时，
\[\left(1+\frac{h}{x}\right)^{n}\to 1\]

于是，由不等式(1.10)和(1.11), 得到当$h\to 0$时，$\left(1+\frac{h}{x}\right)^{\alpha}\to 1$. 因此
\[\lim_{h\to 0}\left[(x+h)^{\alpha}-x^{\alpha}\right]=0\]
即：
\[\lim_{h\to 0}(x+h)^{\alpha}=x^{\alpha}\]
这就是说$x^{\alpha}$在$(0,+\infty)$上处处连续．
\end{proof}

下面介绍常用的复合函数连续性定理．
\begin{blk}
{定理2} 设$X$、$Y$、$Z$是实数集，函数$g:\,X\mapsto Y$, $f:\, g(X)(\subset Y)\mapsto Z$, 又$h:\, X\mapsto Z$定义为$f$和$g$的复合函数，即$h(x)=f (g (x) )\;  (x\in X)$.

若$y=g(x)$在$x=a$处连续；$z=f(y)$在$y=g(a)=b$处连续，那么复合函数$z=h(x)=f(g(x))$在点$a$处也连续．
\end{blk}

\begin{proof}
由于
\begin{equation}
    \lim_{x\to a} g(x)=g(a)=b,\qquad x\in X
\end{equation}    
又因为：
\begin{equation}
    \lim_{y\to b} f(y)=f(b),\qquad y\in g(X)\subset Y
\end{equation}

$\therefore\quad \Lim_{x\to a}h(x)=\Lim_{x\to a}f(g(x))=\Lim_{y\to b}f(y)=f(b)=f(g(a))=h(a)$

这就说明复合函数在点$a$的连续性．
\end{proof}

把上面的等式写成
\[\lim_{x\to a}f (g (x) ) =f (g (a) ) =f \left(\lim_{x\to a} g(x) \right)\]
这个等式表示可以这样求复合函数的极限$\Lim_{x\to a}f(g(x))$, 
如果内层函数$g(x)$在点$a$的极限存在，即$\Lim_{x\to a}g(x)=b$, 
而且$f$在点$b$连续，就可以把极限符号$\Lim_{x\to a}$与函数符号$f$互换，
也就是把极限运算移到内层函数上去施行．

\begin{example}
    说明$h(x)=\sqrt[4]{\frac{1}{1+x^2}}$到处连续，并求$\Lim_{x\to a}h(x)$
\end{example}

\begin{solution}
    $h(x)=\sqrt[4]{\frac{1}{1+x^2}}$可以看作$y=g(x)=\frac{1}{1+x^2},\; x\in\mathbb{R}$
和$z=f(y)=\sqrt[4]{y},\; y\in [0,+\infty)$的复合函数．

由于$g(x)=\frac{1}{1+x^2}$在$\mathbb{R}$上到处连续并且它的值域
\[g(\mathbb{R})=(0,1]\subset [0,+\infty)=Y\]
又幂函数$f(y)=\sqrt[4]{y}$在$Y=[0,+\infty)$上到处连续，因此$f(y)=\sqrt[4]{y}$在$g(\mathbb{R})=(0,1]$上也到处连续，于是
\[\lim_{x\to a}h(x)=f\left(\lim_{x\to a}g(x)\right)=\sqrt[4]{\frac{1}{1+a^2}}\]
\end{solution}


\begin{example}
    求$\Lim_{x\to 2}\sqrt{\frac{1-\cos\pi x}{4-x^2}}$
\end{example}

\begin{solution}
设$y=g(x)=\frac{1-\cos\pi x}{4-x^2}\; (x\ne \pm 2)$，$z=f(y)=\sqrt{y}\; (y\ge 0)$，且知$f$在$[0,+\infty)$上到处连续，于是
\[f(g(x))=\sqrt{\frac{1-\cos\pi x}{4-x^2}}\]
因此：
\[\begin{split}
    \lim_{x\to 2}g(x)&=\lim_{x\to 2}\frac{1-\cos \pi x}{4-x^2}\\
    &=\lim_{x\to 2}\frac{2\sin^2\frac{\pi}{2}x}{4-x^2}=\lim_{x\to 2}\frac{\pi\sin^2\frac{\pi}{2}(2-x)}{\frac{\pi}{2}(2-x)(2+x)}\\
    &=\lim_{x\to 2}\pi\left[\frac{\sin\frac{\pi}{2}(2-x)}{\frac{\pi}{2}(2-x)}\cdot \frac{1}{2+x}\cdot \sin\frac{\pi}{2}(2-x)  \right]\\
    &=\pi\cdot 1\cdot \frac{1}{4}\cdot 0=0\in [0,+\infty)
\end{split}\]

$\therefore\quad \Lim_{x\to 2}\sqrt{\frac{1-\cos\pi x}{4-x^2}}=\sqrt{\Lim_{x\to 2}\left(\frac{1-\cos\pi x}{4-x^2}\right)}=\sqrt{0}=0$
\end{solution}

\begin{example}
    求$\Lim_{x\to 0}\frac{\arcsin x}{x}$
\end{example}

\begin{solution}
设$\alpha=\arcsin x$，这里$-1\le x\le 1$，$-\frac{\pi}{2}\le\alpha\le \frac{\pi}{2}$，
于是，$\sin\alpha=x$．

根据反正弦函数的连续性，有
\[\lim_{x\to 0} \alpha=\lim_{x\to 0}\arcsin x=\arcsin 0=0\]
进行变量替换，得到
\[\lim_{x\to 0} \frac{\arcsin x}{x}=\lim_{x\to 0}\frac{\alpha}{\sin\alpha}=1\]
\end{solution}

下面根据指数函数与对数函数的连续性，我们要建立几个重要的极限：
\begin{enumerate}
    \item $\Lim_{x\to 0} \frac{\ln(1+x)}{x}=1$,\quad 这里$\ln$表示$\log_e$
    \item $\Lim_{x\to 0}\frac{e^x-1}{x}=1$
    \item $\Lim_{x\to 0}\frac{(1+x)^\alpha -1}{x}=\alpha$
\end{enumerate}

\begin{proof}
\begin{enumerate}
    \item $\Lim_{x\to 0} \frac{\ln(1+x)}{x}=\Lim_{x\to 0} \ln(1+x)^{1/x}$
    因为右端在对数号下的式子当$x\to 0$趋向于$e$, 所以根据对数函数的连续性，得到
\[\Lim_{x\to 0} \frac{\ln(1+x)}{x}=\Lim_{x\to 0} \ln(1+x)^{1/x}=\ln e=1\]
\item 令$e^x-1=y$, 于是当$x\to 0$时，由指数函数连续性有$y\to 0$, 其次，我们有$x=\ln(1+y)$, 因此，有
\[\Lim_{x\to 0}\frac{e^x-1}{x}=\Lim_{y\to 0}\frac{y}{\ln(1+y)}=\Lim_{y\to 0}\frac{1}{\ln(1+y)^{1/y}}=1\]
\item 我们令$(1+x)^{\alpha}-1=y$, 当$x\to 0$时，由幂函数的连续性，也有$y\to 0$, 在等式$(1+x)^{\alpha}=1+y$的两边取对数得到
\[\alpha\ln (1+x) =\ln (1+y) \]
利用这一关系式，我们把所给的表达式变形为
\[\frac{(1+x)^\alpha-1}{x}=\frac{y}{x}=\frac{y}{\ln(1+y)}\cdot \alpha\cdot\frac{\ln(1+x)}{x}\]
因此：
\[\Lim_{x\to 0} \frac{(1+x)^\alpha-1}{x}=   \Lim_{x\to 0} \frac{y}{\ln(1+y)}\cdot \alpha\cdot   \Lim_{x\to 0}  \frac{\ln(1+x)}{x}=\alpha\]
\end{enumerate}
\end{proof}

我们讨论了下面的函数：
\begin{enumerate}
    \item $f(x)=c$\quad （$c$为常数）
    \item $x^n$\quad（$n$是整数）
    \item $\sqrt[n]{x}\quad (x\ge 0)$
    \item $a^x\quad (a>0,\; a\ne 1)$
    \item $\log_a x$
    \item $x^\alpha$\quad （$\alpha$是任意实常数）
    \item $\sin x,\; \cos x, \; \tan x,\; \cot x$
    \item $\arcsin x,\; \arccos x, \; \arctan x,\; \arccot x$
\end{enumerate}
上面这些函数统称为\textbf{基本初等函数}．

凡是由基本初等函数及常数，经过有限次四则运算（加、减、乘、除）及有限次函数复合的步骤而得到的且由一个式子表示的函数，及可以化为这种形式的函数，统称为\textbf{初等函数}．

例如，下列函数都是初等函数：
\begin{enumerate}
    \item $f(x)=x\sin\frac{1}{x}+2^x\log_2 x$
    \item $\varphi(x)=\arctan\frac{1}{x}+\sin(x^2+1)$
\end{enumerate}

连续函数的四则运算定理，复合函数连续性定理及反函数连续性定理，在研究初等函数的连续性时具有重要作用．现在根据这些定理，立即得到一个重要结论：

\begin{blk}
 {定理} 一切初等函数在它有定义的任何区间内到处连续．
\end{blk}

\begin{ex}
\begin{enumerate}
    \item 试证 $y=|x|$是一个到处连续的函数．
    \item 说明下列函数到处连续：
\begin{multicols}{2}
\begin{enumerate}
    \item $f(x)=\frac{x^2}{1+x^2}$
    \item $f(x)=\sin|x|$
    \item $\varphi(x)=|\sin x|$
    \item $k(x)=\left\{\cos x-\sin|x|\right\}^3$
\end{enumerate}
\end{multicols}

\item 求出下列函数的间断点，并指出其类型．
\begin{multicols}{2}
\begin{enumerate}
    \item $f(x)=\frac{1}{x-2}$
    \item $f(x)=\frac{x-1}{x^2-1}$
    \item $h(x)=x\cos^2 \frac{1}{x}$
    \item $y=\frac{1}{\sin\pi x}$
    \item $y=\tan(x-1)$
    \item $y=\frac{\sin x}{|x|}$
    \item $f(x)=\begin{cases}
        \frac{\sin x}{|x|},&x\ne 0\\
        1,& x=0
    \end{cases}$
    \item $f(x)=\cot\frac{1}{x}$
    \item $f(x)=\begin{cases}
        \cos\frac{\pi x}{2},& |x|\le 1\\
        |x-1|,& |x|>1
    \end{cases}$
\end{enumerate}
\end{multicols}
\item 试证明下列函数在其定义域内是连续的：
\begin{enumerate}
    \item $f(x)=\begin{cases}
        x^2\sin\frac{1}{x},&x\ne 0\\
        0,&x=0
    \end{cases}$
    \item $\varphi(x)=\begin{cases}
        1-x,& x<0\\
        2x^2+1,& x=0\\
        2+(x-1)^3,& x>0
    \end{cases}$
\end{enumerate}
\item 求下列变量的极限：
\begin{multicols}{2}
    \begin{enumerate}
\item $\Lim_{x\to+\infty}\left(\sin\sqrt{x+1}-\sin\sqrt{x}\right)$
\item $\Lim_{x\to 1}\frac{1+\cos\pi x}{\tan^2 \pi x}$
\item $\Lim_{x\to 0}\frac{x\sin x}{1-\cos x}$
\item $\Lim_{x\to 0}\frac{x}{\sqrt{1-\cos x}}$
\item $\Lim_{x\to 0}\frac{\cos x-\cos 3x}{x^2}$
\item $\Lim_{x\to \tfrac{\pi}{4}}\tan 2x\tan\left(\frac{\pi}{4}-x\right)$
\item $\Lim_{h\to 0}\frac{\arctan (x+h)-\arctan x}{h}$
\item $\Lim_{x\to \infty}\arcsin \frac{1-x}{1+x}$
\item $\Lim_{x\to 0}\frac{1+\sin x-\cos x}{1+\sin px-\cos px}$
\item $\Lim_{x\to 0}\frac{\log_p(1+x)}{x}$
\item $\Lim_{x\to 0}\frac{a^x-1}{x}$
\item $\Lim_{n\to \infty}n\left(\sqrt[n]{x}-1\right)$
    \end{enumerate}
\end{multicols}
\end{enumerate}
\end{ex}

\subsection{闭区间上的连续函数的性质}

我们在2.1中已经定义了闭区间上连续函数的意义．对于这类函数，我们要介绍它的几个重要性质．

\begin{blk}
    {中间值定理} 设$y=f(x)$是一个在闭区间$[a,b]$上到处连续的函数，设$c$是一个介于$f(a)$和$f(b)$之间的常数，则必须存在一个介于$a$、$b$之间的实数$x_0$, 使得$f(x_0)=c$.
\end{blk}

用几何术语来说，$y=f(x),\; a\le x\le b$的图象是一条连结$P(a,f(a))$点和$Q(b,f(b))$点的连续曲线，而$P$、$Q$分居于直线$y=c$的两侧，则曲线$y=f(x),\; a\le x\le b$至少和直线$y=c$有一个交点（$x_0,\; f(x_0)=c$）．(图1.14).
\begin{figure}[htp]
    \centering
\begin{tikzpicture}[>=latex, xscale=1.8]
\draw[->](-1,-.2)--(3,-.2)node[right]{$x$};
\draw[->](.5,-2.5)--(.5,2.5)node[right]{$y$};
\draw[very thick, domain=-.5:2.4, samples=100]plot(\x, {\x*(\x-1)*(\x-2)})node[below left]{$y=f(x)$};
\draw[dashed](-.5,-.2)node[below left]{$a$}--(-.5,-1.875)node[below]{$P(a,f(a))$};
\draw[dashed](2.4,-.2)node[below]{$b$}--(2.4,1.344)node[above]{$Q(b,f(b))$};
\draw[thick](-.75,.2)node[above]{$y=c$}--(2.75,.2);
\draw[dashed](0.7909,.2)--(0.7909,-.2)node[below]{$x_0$};

\end{tikzpicture}
    \caption{}
\end{figure}


对这个定理，我们在第四册（下）4.3中利用二分逼近法及连续函数的性质给出了它的证明，这里不再重述．

\begin{blk}
 {最大值、最小值} 在闭区间$[a,b]$上连续的函数$f$一定在该区间上取到一个最大值和最小值．  
\end{blk}

我们先利用连续函数的图象来直观地说明这个性质，然
后用证明中间值定理的同样方法来证明这个定理．初学者可略过定理的证明，这不影响对本书后面内容的学习．

从图1.15可以看出，在$[a,b]$中，能够找到两点$x_1,x_2$使$f(x_1)=M$, $f(x_2)=m$, 而对于$[a,b]$中所有的$x$, 都有
\[m=f(x_2)\le f(x)\le f(x_1)=M\]


\begin{figure}[htp]
    \centering
    \begin{minipage}[t]{0.48\textwidth}
    \centering
    \begin{tikzpicture}[>=latex, yscale=1.5]
\draw[->](0,0)--(4,0)node[right]{$x$};
\draw[->](.5,-.5)--(.5,2)node[right]{$y$};
\draw[rounded corners=5pt, very thick](1,.5)--(1.25,1)--(1.5,1.3)--(1.75,1)--(2,0)--(2.25,-.5)--(2.6,0)--(2.8,.3)--(3,0)--(3.2,-.3);
\draw[dashed](1,.5)--(1,0)node[below]{$a$};
\draw[dashed](3.2,-.3)--(3.2,0)node[above]{$b$};
\draw[dashed](1.5,1.2)--(1.5,0)node[below]{$x_1$};
\draw[dashed](2.25,-.4)--(2.25,0)node[above]{$x_2$};
\node at (.5,0)[below left]{$O$};

    \end{tikzpicture}
    \caption{}
    \end{minipage}
    \begin{minipage}[t]{0.48\textwidth}
    \centering
    \begin{tikzpicture}[>=latex, scale=1.5]
\draw[->](-.5,0)--(2,0)node[right]{$x$};
\draw[->](0,-.5)--(0,2)node[right]{$y$};
\draw(1,0)node[below]{1}--(1,.1);
\draw(0,1)node[left]{1}--(.1,1);
\node at (0,0)[below right]{$O$};
\draw[very thick](0,0)--(1,1);
\tkzDefPoints{0/0/O, 1/0/B}
\tkzDrawPoints(O,B)
\draw (1,1)[fill=white] circle(1.5pt);
    \end{tikzpicture}
    \caption{}
    \end{minipage}
  \end{figure}


\begin{rmk}
    这个定理是假定$f$在整个闭区间上连续，如果在一点不连续，则定理的结论就可能不成立．
\end{rmk}

例如，函数$f(x)=x-[x]$, $x\in [0, 1]$在$x=1$这一点不连续，它的图象如图1.16, 我们要说明它有最小值0, 而无最大值．

事实上，$f$是分段函数：
\[f(x)=\begin{cases}
    x,& 0\le x<1\\
    0,& x=1
\end{cases}\]
$f$在$[0, 1)$上递增，故
$f(0)=f(1)=0$最小．又在$[0, 1]$上没有一点的函数值为1, 但是对于任何一个小于1的正数$k$, 恒有
\[f (k) =k- [k] =k<1\]
这就说明了$f$在$[0, 1]$上有上界1, 但始终达不到1, 现在我们进一步说明1是它的最小上界，假设定数$k'<1$是$f$在
$[0, 1]$上的上界，我们在$(k',1)$内任意取一个数$x$, 于是
\[f (x) =x- [x] =x> k'\]
这与假设$k'$是$f$在$[0, 1]$上的上界矛盾，因此1是$f$在$[0, 1]$上的最小上界，但是它不属于$f$的值域，这就说明了$f$无最大值．

从这个例中还可以看出，定理中的闭区间不能用开区间来代替．

现在给出定理的证明如下：

\begin{proof}
令$A_1=a$, $B_1=b$, $S_1=\{f(x)|x\in [a,b]\}$，设$K$是$S_1$的最小上界，记$K=l.u.b\{S_1\}$. 再令
\[S_1'=\left\{f (x) ; \; x\in \left[a,\frac{a+b}{2}\right]\right\},\quad S_1''=\left\{f(x);\; x\in\left[\frac{a+b}{2},b\right]\right\}\]
显然有：
\[S_1=S_1'\cup S_1''\]
由此不难看到：
\[K=l.u.b \{S_1\} =\max \left(l.u.b \{S'_1\} ,l.u.b\{S_1''\}\right)\]
所以$l.u.b\{S_1'\}$和$l.u.b\{S_1''\}$两者之中，至少有一个等于$K$, 换句话说，我们总可以在$[A_1,B_1]$的两个半段之中，选取其一为$[A_2,B_2]$, 使得
\[l.u.b \{S_2\} =K,\qquad  S_2= \{f (x) ,\; x\in [A_2, B_2] \}\] 
如此逐步二分，每次由$[A_{n} ,B_{n}]$归纳地选取其半段作为$[A_{n+1} ,B_{n+1}]$，使得
\[\begin{split}
    l.u.b \{f (x) ;\; x\in  [A_{n+1}, B_{n+1}] \}&=l.u.b \{f (x) ;\; x\in [A_n, B_n] \}\\
&=\cdots\\
&=l.u.b \{f (x) ;\; x\in  [a,b] \}=K.
\end{split}\]
对$\{A_n\},\{B_n \}$, 根据实数的连续性，知存在$c$使得$c\in ［A_n,B_n]$, $n=1, 2, 3,\ldots$, 且$B_n-A_n\to 0$.

现在我们要再利用$f(x)$到处连续这一性质，说明$f(c)＝K$. 假若不然，即$f(c)<K$，则可以取$\varepsilon=\frac{1}{2}(K-f(c))>0$，然后利用$f(x)$在$c$点的连续性，知存在点$c$的$\delta$邻域，使得当$|x-c|<\delta$时，有
\[|f(x)-f(c)|<\varepsilon\]
即：
\[\begin{split}
    f(x)<f(c)+\varepsilon&=f(c)+\frac{1}{2}K-\frac{1}{2}f(c)\\
    &=\frac{1}{2}[K+f(c)]\\
    &=\frac{1}{2}[K+K-2\varepsilon]=K-\varepsilon
\end{split}\]
所以
\[l.u.b \{f (x) ;\; x\in  (c-\delta,c+\delta) \} \le K-\varepsilon\]
但是，由$\Lim_{n\to\infty} A_n=\Lim_{n\to\infty}B_n=c$, 知当$n$足够大时，有
\[[A_n,B_n]\subset (c-\delta,c+\delta)\]
这显然和$[A_n,B_n]$的选取法，亦即保持
\[l.u.b \{f (x) ;\; x\in [A_n,B_n]\}=K\]
相矛盾，因此，$f(c)$必须等于$K$. 

类似地，可以证明存在使函数$f$达到它的最小值的一个点．
\end{proof}

\subsection*{习题1.2}
\begin{enumerate}
    \item 求下列函数的极限：
\begin{enumerate}\begin{multicols}{2}
    \item $\Lim_{x\to 0}(\cos x)^{\sin x}$
    \item $\Lim_{x\to 0}\frac{\frac{x^2}{e^{2x}}}{\ln\left(1+\frac{x^2}{e^{2x}}\right)}$
    \item $\Lim_{x\to 0}\frac{\frac{x^2}{e^{2x}}}{\sin^2 \frac{x}{e^{x}}}$
    \item $\Lim_{n\to \infty}n^2\left(x^{\tfrac{1}{n}}-x^{\tfrac{1}{n+1}}\right)$
    \item $\Lim_{x\to 0}\frac{\left(e^x+e^{-x}\right)\left(e^x-e^{-x}\right)}{2x}$
    \item $\Lim_{x\to 0}(1-2x)^{\tfrac{1}{x}}$
    \item $\Lim_{n\to \infty}10\left(1+\frac{0.03}{n}\right)^n$
    \item $\Lim_{x\to +\infty}\arccos\left(\sqrt{x^2+x}-x\right)$
    \item $\Lim_{x\to -\infty}\arctan\frac{x-4}{(x-2)^2}$
    \item $\Lim_{x\to -\infty}\arctan\frac{x}{\sqrt{1+x^2}}$
    \item $\Lim_{x\to -\infty}\frac{x}{\sqrt{1+x^2}}$
\end{multicols}   
    \item $\Lim_{x\to 0}\frac{\ln\frac{1+x}{1-x}}{\arctan(1+x)-\arctan(1-x)}$
\end{enumerate} 

\item 在一个以弦$AB$为界的弓形中，把点$A$与$B$与弧$\wideparen{AB}$的中点$C$用直线连结起来，于是得到一个等腰三角形$\triangle  ABC$, 再在$A$与$B$作切线，它们交于$D$, 试计算$\Lim_{x\to 0}\frac{\triangle ABC}{\triangle ABD}$，其中$x$是
弦$AB$所对应的中心角．

\item $f(x)=\frac{ax^3+bx^2+cx+d}{x^2-x-2}$满足条件
\[\lim_{x\to \infty} f(x)=2,\qquad \lim_{x\to 2} f(x)=3\]
求$f(x)$.
\end{enumerate}



\chapter{变率和微商}

从本章起，我们开始学习单变量的微积分学的基本概念和基础理论．

微积分学是研究变量的数学，变量之间的关系就是函数，因此，函数是微积分学研究的主要对象．

在函数的基本性质中，有两个最基本、最重要的概念－变率与求和，为了解决求函数的变率与求函数$f$在$[a,b]$上的和的问题就相应地产生微分与积分运算，而这两种运算之间，也存在着一种自然的互逆关系，在本章中，我们由函数的变率问题引出函数的微商（导数）概念，并给出初等函数的一套求导法则，在下一章中，我们揭示微分运算与积分运算的互逆关系，这就是微积分学的基本定理．

总起来说，微分反映了函数的局部性质，或在某个点附近的性质；积分则反映了函数的整体性质，或某个区间的性质；函数的局部性质与整体性质之间的有机联系，恰恰反映了微分运算与积分运算之间的互逆关系．

\section{微商（导数）的定义}
函数关系$y=f(x)$就是确定变量$y$如何随着变量$x$的变动而变动的关系，对于给定的函数$y=f(x)$, 变量$y$在变量$x$的不同点附近的变动情况是不尽相同的，这就是说，在变量$x$的
某个值$x_1$处，当$x_1$略加变动时，相应的$y$的变动可能相当剧烈（急增，或急减）；而在变量$x$的另一个值$x_2$处，$y$的变动就可能较为迟缓，但是，用这样的语言来表达函数在某一点处的变率是不精确的，我们需要用“数量”来确切地表达这个意思，这就是函数在某点（或某时刻）的变率的问题，简称变率．

\subsection{直线函数的变率}
一次函数$f (x) =kx+b$
是一种最简单的函数，它的函数图象是一条斜率等于$k$的直线，如图2.1所示．
\begin{figure}[htp]
    \centering
\begin{tikzpicture}[>=latex]
    \draw[->](-2,0)--(4,0)node[right]{$x$};
    \draw[->] (0,-1)--(0,5)node[right]{$y$};
    \draw[very thick] (-1.5,-.5)--(4,5);
\draw[->](-.2,0) arc (0:45:.8)node[below]{$\theta$};
\draw(2,3)node[left]{$P_1$}--(2,0)node[below]{$x_1$};
\draw(3,4)node[left]{$P$}--(3,0)node[below]{$x$};
\draw[|<->|](3.2,4)--node[right]{$f(x)-f(x_1)$}(3.2,3);
\draw[|<->|](2,2.8)--node[below]{$x-x_1$}(3,2.8);
\draw(2,3)--(3,3);\node at (.25,-.25){$O$};
\end{tikzpicture}
    \caption{}
\end{figure}

设$x_1$, $f(x_1)$和$x$, $f(x)$分别是直线上点$P_1$和$P$的坐标，为了反映函数变化快慢的问题，无论$x<x_1$还是$x>x_1$，自然地考虑在点$x_1$邻近，函数与自变量的相应的改变量的比：
\[\frac{f(x)-f(x_1)}{x-x_1}=\frac{(kx+b)-(kx_1+b)}{x-x_1}=k=\tan\theta\]
上面的表达式称为函数的\textbf{差商}，它表示函数在区间$[x_1,x]$或$[x,x_1]$上对于自变量$x$的\textbf{平均变化率}．由于$k$是不随变量$x$变动而变动的常数，因此，一次函数在自变量的任何一个区间内的平均变化率都是常数．

如果让自变量的变化区间的长度无限地缩短，也就是让$x$无限地接近于$x_1$时，平均变化率所趋向的极限
\[\lim_{x\to x_1} \frac{f (x) -f (x_1)}{x-x_1} =k\]

\subsection{平滑曲线的切线与变率}
一般的函数$y=f(x)$的图象通常不是直线，由于函数
和它的图象的多样性，为了讨论的方便起见，我们先把讨论的范围限制在“平滑”的曲线上，常用的函数$y=f(x)$的图象往往是“平滑的”，平滑性的直观内涵是：用愈高倍的显微镜去观察曲线的微段，就愈象直线段，比较明确的几何说法是：一条曲线在$P$点的平滑性就是存在唯一的一条过$P$点的切线，它无限地逼近曲线在$P$点邻近的微段，于是，当$y=f(x)$的图象$C$在$P$点存在唯一的一条切线时，我们就说曲线$C$在$P$点平滑，而$P$点叫做曲线的平滑点，一条在每点都平滑的曲线叫做平滑的曲线，一个函数的图象是平滑曲线时，我们就称这种函数为平滑函数．

同学可能会问这样一个问题：在曲线的点$P$存在唯一的一条切线的含义是什么？因为迄今我们对于一般的曲线的切
线还未下过定义呢！

我们从图2.2和2.3注意到不能把切线定义为与曲线只有一个交点的直线．这样的定义限制得既太紧同时又太松．因
为，照此定义，图2.2所示的直线就不是过曲线上$P$点的切线了，实际上，尽管图2.2的直线与曲线还有一个交点$Q$, 但它在曲
线$P$点邻近却与曲线密合，故仍应该是过曲线$P$点的切线；又图2.3表明过抛物线上任何一点$P$与$y$轴平行的直线虽然与抛物线只有一个交点，但它的其余部分却远离$P$点邻近的弧，故它不应该是抛物线的切线，定义切线的可行途径是从割线开始，并应用极限的概念．

\begin{figure}[htp]
    \centering
    \begin{minipage}[t]{0.48\textwidth}
    \centering
    \begin{tikzpicture}[>=latex, scale=.7]
\draw[->](-4,0)--(4,0)node[right]{$x$};
\draw[->](0,-2)--(0,3.5)node[right]{$y$};
\draw[domain=-3.2:3, very thick, samples=100]plot(\x, {.06*\x^3+.5})node[above]{$y=f(x)$};
\node at (0,0)[below left]{$O$};
\node at (1.5,.7025)[above]{$P$};
\draw[domain=-4:3.5, samples=10, thick]plot(\x, {.405*\x+.095});
\node at (-3,-1.12)[below right]{$Q$};
    \end{tikzpicture}
    \caption{}
    \end{minipage}
    \begin{minipage}[t]{0.48\textwidth}
    \centering
    \begin{tikzpicture}[>=latex, scale=.8]
  \draw[->](-3,0)--(3,0)node[right]{$x$};
\draw[->](0,-1)--(0,4)node[right]{$y$};
\draw[domain=-2:2, very thick, samples=100]plot(\x, {\x*\x});
\node at (-1.2,1)[left] {$f(x)=x^2$};
\draw(1,0.2)--(1,2);
\node at (1,1)[right]{$P$};
\node at (0,0)[below left]{$O$};
\draw[domain=.2:2, samples=10, thick]plot(\x, {2*\x-1});


    \end{tikzpicture}
    \caption{}
    \end{minipage}
  \end{figure}



如图\ref{fig:PQ_tangent}所示，取曲线$y=f(x)$上点$P$附近的另一点$Q$, 通过这两点画一条直线，这直线叫做过曲线上$P$点的割线，让$Q$点沿曲线向点$P$移动，这条割线将达到极限位置，此极限位置与$Q$点从哪一侧趋向于$P$是无关的，我们称这个割线的极限位置为过曲线上$P$点的切线．

\begin{figure}[htp]
    \centering
\begin{tikzpicture}[>=latex, scale=1.5]
    \draw[->](-.5,0)--(4,0)node[right]{$x$};
    \draw[->](0,-.5)--(0,3)node[right]{$y$};
    \node at (0,0)[below left]{$O$};
    \draw[domain=.5:3.2, samples=100, very thick]plot(\x, {-.5*(\x-2.25)^2+2})node[right]{$y=f(x)$};
    \tkzDefPoints{1/1.22/P, 3/1.72/Q, 3.5/1.22/P'}
    \tkzLabelPoints[above](P,Q)
    \tkzDrawSegments(P,P')
    \draw[domain=.2:2, samples=10, thick]plot(\x, {1.25*\x-0.03});
    \tkzDefPoints{3/3.72/Q'}
    \tkzDrawLines(P,Q)
    \tkzMarkAngles[mark=none, size=.7](P',P,Q)
    \tkzMarkAngles[mark=none, size=.3](P',P,Q')
    \tkzLabelAngle[pos=.5](P',P,Q'){$\alpha_1$}
    \tkzLabelAngle[pos=.9](P',P,Q){$\alpha$}
\end{tikzpicture}
    \caption{}\label{fig:PQ_tangent}
\end{figure}

割线的这种极限位置的存在性这一假设，与曲线在点$P$具有唯一的一条切线或确定的方向的假设是等价的．

现在我们要对曲线$y=f(x)$用解析式子把割线的这种极限位置存在的过程表示出来．
设$\alpha$是割线$PQ$同正$x$轴构成的夹角，$\alpha_1$是过点$P$点的切线同正$x$轴构成的夹角，于是
\[\lim_{Q\to P}\alpha=\alpha_1\]
设$x_1,y_1$和$x,y$分别是点$P$和$Q$的坐标，这时，我们立即得到
\[\tan\alpha=\frac{y-y_1}{x-x_1}=\frac{f(x)-f(x_1)}{x-x_1}\]
因此，上述求极限的过程（不考虑垂直切线$\alpha_1=\frac{\pi}{2}$的情况）可由下式来表示：
\[\lim_{x\to x_1}\frac{f(x)-f(x_1)}{x-x_1}=\lim_{\alpha\to \alpha_1}\tan\alpha=\tan\alpha_1\]
这就是说过曲线$y=f(x)$上$P(x_1,y_1)$点的切线的斜率等于$y=f(x)$的差商当$x\to x_1$时的极限．

\begin{example}
    求抛物线$y=ax^2+bx+c$在$x_0$处的切线的斜率．
\end{example}

\begin{solution}
    解依题意$(x_0,f(x_0)=ax_0^2+bx_0+c)$在抛物线上，并设$(x_0+h,f(x_0+h))$是抛物线上点$(x_0,f(x_0))$的附近的一点，我们有
\begin{align*}
&\qquad \lim_{h\to 0}\frac{f(x_0+h)-f(x_0)}{(x_0+h)-x_0}\\
&=\lim_{h\to 0}\frac{[a(x_0+h)^2+b(x_0+h)+c]-[ax^2_0+bx_0+c]}{h}    \\
&=\lim_{h\to 0}\frac{(2ax_0+b)h+ah^2}{h}\\
&=\lim_{h\to 0}[(2ax_0+b)+ah]=2ax_0+b
\end{align*}
所以抛物线$y=ax^2+bx+c$在$x_0$处的切线的斜率是$2ax_0+b$.
\end{solution}

\begin{example}
    一质点沿一直线在$t$秒内移动的距离是$s=s(t)=t^2+4t$．
    
    求：质点的初速度；在两秒末的速度；前两秒内的平均速度．
\end{example}

\begin{analyze}
    如果质点从起点开始所走的距离$s$是时间$t$的线性函数，则由2.1知道该质点在每一时刻的速度都是常数，它的大小由平均速度来确定，即等于一次函数的斜率，此时我们说该质点作匀速运动，但是，如果运动不再是匀速的，即质点的速度每时每刻都是变的，那么我们将时刻$t$的速度（也叫做瞬时速度）理解成什么呢？为了回答这个问题，我们考察差商
\[\frac{\Delta s}{\Delta t}=\frac{s(t)-s(t_0)}{t-t_0}\]
或者写成\[\frac{s(t_0+\Delta t)-s(t_0)}{\Delta t}\]
这个差商称为在$t_0$和$t_0+\Delta t$之间的这段时间间隔上的质点的平均速度，对照着$s=s(t)=t^2+4t$的图象来看，这个平均速度也就是过曲线上的$P(t_0,s(t_0))$
点及它邻近一点$Q(t_0+\Delta t,s(t_0+\Delta t))$的割线的斜率（图2.5）．当$\Delta t$ 很小时，可以认为，从时刻$t_0$到$t_0+\Delta t$这段时间内，速度来不及有很大变化，可以近似地看成匀速运动，因而这段时间内的平均速度就可以看成时刻$t_0$的瞬时速度的近似值．

\begin{figure}[htp]
    \centering
    \begin{tikzpicture}[>=latex, scale=.5]
\draw[->](-5.5,0)--(3,0)node[right]{$t$};
\draw[->](0,-5)--(0,7)node[right]{$s$};

\draw[domain=-5.25:0, samples=300, dashed]plot(\x, {\x*\x+4*\x});
\draw[domain=0:1.25, samples=100, thick]plot(\x, {\x*\x+4*\x});
\foreach \x in {-4,-3,-2,-1,1,2}
{
    \draw(\x,0)--(\x,.2);
}
\node at (-2,0)[below]{$-2$};
\node at (1,0)[below]{$1$};
\node at (0,0)[below left]{$O$};
    \end{tikzpicture}

    \caption{}
\end{figure}



显然，从时刻$t_0$到时刻$t_0+\Delta t$, 质点走过的路程为
\begin{align*}
    \Delta s&=s(t_0+\Delta t)-s(t_0)\\
    &=(t_0+\Delta t)^2+4(t_0+\Delta t)-(t^2_0+4t_0)\\
    &=(2t_0+4)\Delta t+(\Delta t)^2
\end{align*}
所以这段时间内的平均速度为
\begin{align*}
    \frac{\Delta s}{\Delta t}&=\frac{s(t_0+\Delta t)-s(t_0)}{\Delta t}\\
    &=\frac{(2t_0+4)\Delta t+(\Delta t)^2}{\Delta t}\\
    &=(2t_0+4)+\Delta t
\end{align*}
$\Delta t$越小，这个平均速度就越接近时刻$t_0$的瞬时速度$v_0$, 我们自然令$\Delta t\to 0$, 求差商的极限值，得到
\begin{align*}
    \lim_{\Delta t\to 0}\frac{s(t_0+\Delta t)-s(t_0)}{\Delta t}&= \lim_{\Delta t\to 0}[(2t_0+4)+\Delta t]\\
    &=2t_0+4
\end{align*}

这样平均速度$\frac{s(t_0+\Delta t)-s(t_0)}{\Delta t}$，当$\Delta t\to 0$时的极限值
就表达了质点在时刻$t_0$的瞬时速度，把它记作
\[s'(t_0)=   \lim_{\Delta t\to 0}\frac{s(t_0+\Delta t)-s(t_0)}{\Delta t}\]
\end{analyze}

\begin{solution}
\begin{enumerate}
\item 质点的初速度
\[s'(0)=2\x0+4=4(\ms)\]
\item 质点在两秒末的速度
\[s'(2)=2\x2+4=8(\ms)\]
\item \[\text{质点在前两秒内的平均速度}=\frac{\text{在前两秒内所走距离}}{\text{时间}}=\frac{2^2+4\x2}{2}
=6(\ms)\]
\end{enumerate}
\end{solution}

从这个问题可以看出质点在$t_0$时刻的瞬时速度$s'(t_0)$的几何意义就是曲线$s=s(t)$在$P(t_0,s(t_0))$点的切线的斜率，所以函数在某点的变率有确定值与函数的图象在该点有唯一的一条切线是两个密切相关的概念．

\subsection{微商（导数）的定义}
从上面所举的两个例子来看，问题来自不同的领域：
\begin{enumerate}
    \item 求过曲线上一点的切线，
    \item 求函数在某点的变率.
\end{enumerate}
但解决的方法却完全一样，就是计算函数的差商的极限，这种极限反映了自然界中很多不同现象在量方面的共性，因此有必要从这些具体问题中把它抽象出来加以研究，再反过来去解决这类具体问题．

\begin{blk}{定义}
    设$y=f(x)$是定义在闭区间$[a,b]$上的一个函数，$x_0\in (a,b)$, 如果极限
\[\lim_{\Delta x\to 0}\frac{f (x_0+\Delta x) -f (x_0)}{\Delta x}\]
存在，我们就说$f(x)$在点$x_0$处\textbf{可微}，并称这极限为函数$f(x)$在$x_0$点的\textbf{微商}（或导数），记为$f'(x_0)$或
$\frac{\dd y}{\dd x}\Big|_{x=x_0}$.
\end{blk}

显然，$f'(x_0)$的值与点$x_0$有关，当点$x_0$在开区间$(a,b)$内变化时，$f'(x_0)$也将跟着变化，因此，如果函数$f(x)$在开区间$(a,b)$内每点都可微（即存在有限导数），那么$f'(x)$便是一个新的函数，称为$f(x)$的\textbf{导函数}．

求已知函数的导函数$f'(x)$的运算，称为微商 运算，计算过程如下：
\begin{enumerate}
\item 设$\Delta x$为自变量某个值$x$的改变量：
\item 计算$f(x)$在点$x$的相应改变量
\[\Delta y=f (x+\Delta x) -f (x) \]
\item 计算$f(x)$在点$x$的差商
\[\frac{\Delta y}{\Delta x}=\frac{f (x+\Delta x) -f (x)}{\Delta x}\]
\item 计算
\[\lim_{\Delta x\to 0}\frac{f (x+\Delta x) -f (x)}{\Delta x}=f'(x)=\frac{\dd y}{\dd x}\]
\end{enumerate}

应当注意，这里$\frac{\dd y}{\dd x}$是一个独立的记号，它表示函数$f(x)$
在点$x$的导数，不能把它当成一个分数来看待，必须把它看成一个整体．

从微商的定义可以看出：
\begin{enumerate}
\item 曲线在一点的切线的斜率，就是函数在这一点的变率（微商或导数）．    
\item 微商所涉及的是函数的“局部”性质，也就是说，函数$y=f(x)$在一点$x_0$处是否可微只与函数$y=f(x)$在$x=x_0$处及其近旁的性质有关，而与其它地方无关．    
\item 如果$f(x)$在点$x$可微，按照极限存在的条件，必须且只须
\[\lim_{\Delta x\to 0^+}\frac{f (x+\Delta x) -f (x)}{\Delta x},\qquad \lim_{\Delta x\to 0^-}\frac{f (x+\Delta x) -f (x)}{\Delta x}\]
同时存在而且相等．

上面两式分别称为$f(x)$在点$x$的\textbf{右导数}和\textbf{左导数}，记为$f'_+(x)$和$f'_-(x)$.
\end{enumerate}

\begin{example}
今有一个正在膨胀的肥皂泡，
\begin{enumerate}
\item 求肥皂泡的体积对于半径的增大率，
\item 如果肥皂泡的半径每秒增大0.1cm，问当半径为2cm时，体积的增大率是多少？
\end{enumerate}
\end{example}

\begin{solution}
    肥皂泡的体积$V$与半径$r$的函数关系是．
\[ V(r) =\frac{4}{3}\pi r^3\]
体积$V(r)$对于半径$r$的增大率，依导函数定义，就是$V(r)$对于$r$的导函数，故
\begin{align*}
    V'(r)&=\lim_{\Delta r\to 0}\frac{V(r+\Delta r)-V(r)}{\Delta r}\\
    &=\lim_{\Delta r\to 0}\frac{\frac{4\pi }{3}[(r+\Delta r)^3-r^3]}{\Delta r}\\
    &=\frac{4\pi }{3}\lim_{\Delta r\to 0}[3r^2+3r(\Delta r)+(\Delta r)^2]\\
    &=4\pi r^2
\end{align*}

因此，肥皂泡的体积$V$对于半径$r$的增大率是$4\pi r^2$. 当$r=2$时，\[V'(2)=4\pi \cdot 2^2=16\pi \]
上式表示在$r=2$时，肥皂泡体积对于半径的增大率是$16\pi$, 它
的含义也可以这样理解，体积增大的快慢在$r=2$时，是$16\pi$倍于半径增大的快慢．

由于半径的增大率是每秒0.1cm，故体积在半径等于2时的增大率是
\[16\pi \x0.1=1.6\pi ({\rm cm^3/s})\]
\end{solution}

\begin{ex}
\begin{enumerate}
    \item 设$y=\frac{1}{x},\quad (x\ne 0)$，求
\begin{enumerate}
    \item 当$x$取改变量$\Delta x$后，函数改变量$\Delta y$的表达式；
    \item 当$x=3$, $\Delta x=-1$时，$\Delta y$的值；
    \item 当$x=3$时，$\frac{\Delta y}{\Delta x}$的表达式和$\frac{\dd y}{\dd x}\Big|_{x=3}$的值．
\end{enumerate}

    \item 设$\phi(x)=\frac{2}{x^2},\quad (x\ne0)$, 求
\begin{enumerate}
\item 当$x=1$, $\Delta x=-0.5$时，$\Delta \phi(x)$的值；
\item 对于任何$x\; (x\ne 0)$, 在$\Delta x$的间隔内，$\phi(x)$的平均变率$\frac{\Delta \phi}{\Delta x}$的表达式；
\item $\phi(x)$在任何点$x\; (x\ne 0)$处的瞬时变率$\phi'(x)$.
\end{enumerate}

    \item 平均变化率$\frac{\Delta y}{\Delta x}=\frac{f (x+\Delta x) -f (x)}{\Delta x}$依赖于哪两个变量？在平均变化率取极限求瞬时变化率的过程中，$x$是变量还是常量？$\Delta x$是变量还是常量？
    \item 点$P(2, 8)$在抛物线$y=x^2+2x$上，点$Q$为抛物线上任何一点．
\begin{enumerate}
    \item 求割线$PQ$的斜率的表达式；
  \item 当$Q$点横坐标为$2.1, 1.9, 2.002, 1.998, 2+h$时，求割线斜率的值；
\item 求在$P(2, 8)$点处，抛物线$y=x^2+2x$的切线方程和法线方程．
\end{enumerate}    

\item 求圆面积对于它的半径的变率，对于它的直径的变率．
\item 圆的半径的变率为2cm/s, 求圆面积在半径等于4cm
时的变率．
\end{enumerate}
\end{ex}

\subsection{函数的可微性与连续性的关系}

\begin{blk}
    {定理} 如果函数$f(x)$在点$x_0$可微，那么$f(x)$在点$x_0$处连续．
\end{blk}

\begin{proof}
设$f(x)$在点$x_0$处是可微的，也就是$f(x)$在$x_0$处有导数，则
\[f'(x_0)=\lim_{x\to x_0}\frac{f(x)-f(x_0)}{x-x_0}\]
所以
\begin{align*}
    \lim_{x\to x_0}[f(x)-f(x_0)]&=\lim_{x\to x_0}\frac{f(x)-f(x_0)}{x-x_0}(x-x_0)\\
    &=\lim_{x\to x_0}\frac{f(x)-f(x_0)}{x-x_0}\cdot \lim_{x\to x_0}(x-x_0)\\
&=f'(x_0)\cdot 0=0
\end{align*}
\end{proof}

上面的定理说明函数在其有导数的点一定连续，在其不
连续的点一定没有导数．但是它的逆命题不一定成立，即连续的函数不一定有导数．前面我们曾指出一个函数可微的充分必要条件是它的左导数和右导数都存在而且相等，下面给出在某些点不可微的函数的例子．

\begin{example}
    求函数$f(x)=|x^2-1|$的导函数的定义域．
\end{example}

\begin{figure}[htp]
    \centering
\begin{tikzpicture}[>=latex, scale=1]
    \draw[->](-2.5,0)--(2.5,0)node[right]{$x$};
    \draw[->](0,-.5)--(0,5)node[right]{$y$};
    \node at (0,0)[below left]{$O$};
    \draw[domain=-2.3:-1, samples=100, very thick]plot(\x, {\x*\x-1});
    \draw[domain=1:2.3, samples=100, very thick]plot(\x, {\x*\x-1});
    \draw[domain=-1:1, samples=100, very thick]plot(\x, {-\x*\x+1});
    \node at (0,0)[below left]{$O$};
    \node at (-1,0)[below left]{$-1$};
    \node at (1,0)[below right]{$1$};
    
    \draw[domain=-2:-.5, samples=10, thick]plot(\x, {-2*\x-2});
    \draw[domain=-1.5:.5, samples=10, thick]plot(\x, {2*\x+2});
    \draw[domain=-.5:1.5, samples=10, thick]plot(\x, {-2*\x+2});
    \draw[domain=.5:2, samples=10, thick]plot(\x, {2*\x-2});

\end{tikzpicture}
    \caption{}\label{fig:f(x)=|x^2-1|}
\end{figure}
    
\begin{solution}
函数$f(x)=|x^2-1|$是一个到处连续的函数，它的图象如图\ref{fig:f(x)=|x^2-1|}所示．

去掉函数式中的绝对值的符号，$f(x)$便可以写成分段函数式：
\[f(x)=\begin{cases}
    x^2-1,& x\le -1\quad \text{或}\quad x\ge 1\\
    -(x^2-1), & -1\le x\le 1
\end{cases}\]
所以
\[f'(x)=\begin{cases}
    2x ,&  x< -1\quad \text{或}\quad x> 1\\
    -2x,& -1<x<1
\end{cases}\]
但是当$x=-1$和1时，函数$f(x)$的左导数和右导数不相等，即此时导数不存在．因为当$x=-1$时，函数$f(x)$在$-1$点左邻域的表达式是
\[f (x) =x^2-1\]
而在$-1$点右邻域的表达式是
\[f (x) =- (x^2-1)\]
所以$f(x)$在$x=-1$点的左导数是
\begin{align*}
    f'_-(-1)&=\lim_{\Delta x\to 0^-}\frac{f(-1+\Delta x)-f(-1)}{\Delta x}\\
    &=\lim_{\Delta x\to 0^-}\frac{[(-1+\Delta x)^2-1]-0}{\Delta x}\\
    &=\lim_{\Delta x\to 0^-}(-2+\Delta x)=-2
\end{align*}
$f(x)$在$x=-1$点的右导数是
\begin{align*}
    f'_+(-1)&=\lim_{\Delta x\to 0^+}\frac{f(-1+\Delta x)-f(-1)}{\Delta x}\\
    &=\lim_{\Delta x\to 0^+}\frac{-[(-1+\Delta x)^2-1]-0}{\Delta x}\\
    &=\lim_{\Delta x\to 0^+}(2-\Delta x)=2 
\end{align*}
由于$f'_-(-1)\ne f'_+(-1)$, 所以我们说$f(x)$在$x=-1$处的导数不存在，同样可得
\[-2=f'_- (1)\ne f'_+(1)=2\]
故$f(x)$在$x=1$处也不可微．

因此，$f(x)=|x^2-1|$的导函数的定义域是$x\ne\pm 1$的点的集合，它的函数值如下：
\[f'(x)=\begin{cases}
    2x,& |x|>1\\
    -2x,& |x|<1
\end{cases}\]
\end{solution}

对照图\ref{fig:f(x)=|x^2-1|}来看，尽管函数$f(x)=|x^2-1|$处处连续，但是在$x=-1$和1这两点，曲线的特征是当动点由左侧趋于定
点1(或$-1$)时，割线的极限位置存在；当动点由右侧趋于定点1(或$-1$)时，割线的极限位置也存在，这两个极限位置分别叫1(或$-1$)点的\textbf{左、右切线}，并且它们之间的夹角不为\textbf{平角}，点$-1$和1分别是曲线的一种\textbf{角点}．
    
图\ref{fig:vertical_tangent}的情形是曲线$y=f(x)$在$P(x_0,y_0)$点连续并且在此点有平行于$y$轴的切线．现在我们来讨论图\ref{fig:vertical_tangent}中所示各函数在点$x_0$处的导数：

\begin{figure}[htp]
    \centering
\begin{tikzpicture}[>=latex]
\begin{scope}
\draw[->](-.5,0)--(4,0)node[right]{$x$};
\draw[->](0,-.5)--(0,4)node[right]{$y$};
\node at (2,-1){(a)};
\node at (0,0)[below left]{$O$};
\draw(2,0)node[below]{$x_0$}--(2,3.5);
\node at (2,2)[right]{$P(x_0,y_0)$};
\tkzDefPoints{1.5/.2/A, 2/2/B, 2.5/3.8/C}
\draw[very thick](A) to [bend right=18] (B) to [bend left=18](C);

\end{scope}
\begin{scope}[xshift=6cm]
    \draw[->](-.5,0)--(4,0)node[right]{$x$};
\draw[->](0,-.5)--(0,4)node[right]{$y$};
\node at (2,-1){(b)};
\node at (0,0)[below left]{$O$};
\draw(2,0)node[below]{$x_0$}--(2,3.5);
\node at (2,2)[right]{$P(x_0,y_0)$};
\tkzDefPoints{2.5/.2/A, 2/2/B, 1.5/3.8/C}
\draw[very thick](A) to [bend right=-18] (B) to [bend left=-18](C);

\end{scope}
\begin{scope}[yshift=-6cm]
    \draw[->](-.5,0)--(4,0)node[right]{$x$};
\draw[->](0,-.5)--(0,4)node[right]{$y$};
\node at (2,-1){(c)};
\node at (0,0)[below left]{$O$};
\draw(2,0)node[below]{$x_0$}--(2,3.5);
\node at (2,2)[right]{$P(x_0,y_0)$};
\tkzDefPoints{1.5/.2/A, 2/2/B, 2.5/.2/C}
\draw[very thick](A) to [bend right=18] (B) to [bend left=-18](C);


\end{scope}

\begin{scope}[xshift=6cm, yshift=-6cm]
    \draw[->](-.5,0)--(4,0)node[right]{$x$};
\draw[->](0,-.5)--(0,4)node[right]{$y$};
\node at (2,-1){(d)};
\node at (0,0)[below left]{$O$};
\draw(2,0)node[below]{$x_0$}--(2,3.5);
\node at (2,2)[right]{$P(x_0,y_0)$};
\tkzDefPoints{1.5/3.8/A, 2/2/B, 2.5/3.8/C}
\draw[very thick](A) to [bend right=-18] (B) to [bend left=18](C);



\end{scope}
\end{tikzpicture}
    \caption{}\label{fig:vertical_tangent}
\end{figure}

图(a)表示的函数在点$x_0$的邻近递增，因此过曲线上$P(x_0,y_0)$和$Q(x_0+\Delta x,y_0+\Delta y)$两点的割线，无论$Q$点在$P$点的哪一侧，$\Delta x$与$\Delta y$都有相同正、负号，于是
\[\frac{\Delta y}{\Delta x}=\frac{f(x_0+\Delta x)-f(x_0)}{\Delta x}>0\]
而且当$\Delta x\to 0$时，割线的倾斜角$\varphi$（$0^{\circ}\le \varphi\le 180^{\circ}$）便以$\pi/2$为极限，从而
\[\lim_{\Delta x\to 0}\frac{f(x_0+\Delta x)-f(x_0)}{\Delta x}=\lim_{\varphi\to \tfrac{\pi^-}{2}}\tan\varphi=+\infty\]
可见$f(x)$在点$x_0$不可微（不存在有限导数），但是我们常常把这种情况：当$\Delta x\to 0$时，$\frac{\Delta y}{\Delta x}\to\infty$简略地叙述为函
数$f(x)$在$x_0$的导数为正无穷大，记作$f'(x_0)=+\infty$．

图(b)表示的函数$f$在点$x_0$的邻近递减，因此过曲线上$P(x_0,y_0)$和$Q(x_0+\Delta x,y_0+\Delta y)$两点的割线，无论$Q$点在$P$点的哪一侧，$\Delta x$与$\Delta y$都有相反的正、负号，于是
\[\frac{\Delta y}{\Delta x}=\frac{f(x_0+\Delta x)-f(x_0)}{\Delta x}<0\]
而且当$\Delta x\to 0$时，$\varphi\to \frac{\pi^+}{2}$，从而
\[\lim_{\Delta x\to 0}\frac{f(x_0+\Delta x)-f(x_0)}{\Delta x}=\lim_{\varphi\to \tfrac{\pi^+}{2}}\tan\varphi=-\infty\]
因此$f(x)$在点$x_0$处不可微，但是我们常常把这种情况叙述为函数$f(x)$在$x_0$处的导数为负无穷大，记作$f'(x_0)=-\infty$．

图\ref{fig:vertical_tangent}(c)表示的函数在点$x_0$处连续，让$x$由$x_0$变动到
$x_0+\Delta x$，于是由$\Delta x$所引起的相应的函数的改变量$\Delta y$的情形是：
\begin{itemize}
    \item 当$\Delta x<0$时，有$\Delta y<0$，从而$\frac{\Delta y}{\Delta x}>0$；
    \item 当$\Delta x>0$时，有$\Delta y<  0$，从而$\frac{\Delta y}{\Delta x}<0$．
\end{itemize}
于是：
\begin{itemize}
    \item 当$\Delta x\to 0^-$时，$\varphi\to \frac{\pi^-}{2}$，从而$\frac{\Delta y}{\Delta x}=\tan\varphi\to +\infty$；
    \item 当$\Delta x\to 0^+$时，$\varphi\to \frac{\pi^+}{2}$，从而$\frac{\Delta y}{\Delta x}=\tan\varphi\to -\infty$．
\end{itemize}
因此$f(x)$在点$x_0$不可微，但是我们常把这种情况叙述为$f(x)$在点$x_0$的左导数$f'_-(x_0)=+\infty$, 而它的右导数为$f'_+(x_0)=-\infty$, 曲线在$x_0$处由上升转为下降有一个尖点，它的切线垂直于$x$轴．

图(d)表示的函数在点$x_0$处连续，但是
\begin{itemize}
    \item 当$\Delta x\to 0^-$时，$\frac{\Delta y}{\Delta x}=\tan\varphi\to -\infty$；
    \item 当$\Delta x\to 0^+$时，$\frac{\Delta y}{\Delta x}=\tan\varphi\to +\infty$．
\end{itemize}
因此，$f(x)$在点$x_0$不可微，我们常常把这种情况叙述为$f'_-(x_0)=-\infty$和$f'_+(x_0)=+\infty$, 点$x_0$是曲线的一个尖点，它的切线垂直于$x$轴．

\begin{example}
设$f(x)=\begin{cases}
    0,& x=0\\
    x\sin\frac{1}{x},& x\ne 0
\end{cases}$

讨论它在0点的导数．
\end{example}

\begin{solution}
在第一章，我们已经说明了
\[\lim_{x\to 0}x\sin\frac{1}{x}=0=f(0)\]
因此，$f(x)$在点$x=0$连续，很明显$f(x)$在其它各点处也都连续，所以说$f(x)$是一个到处连续的函数，我们将说明它在原点不存在导数．

因为
\[\frac{f(0+\Delta x)-f(0)}{\Delta x}=\frac{\Delta x\sin\frac{1}{\Delta x}}{\Delta x}=\sin\frac{1}{\Delta x}\]
而当$\Delta x\to 0$时，$\sin\frac{1}{\Delta x}$
没有极限，故$f'(0)$不存在，就
其几何意义来说，当动点沿着曲线$y=x\sin\frac{1}{x}$趋于原点$O$
时，割线$OQ$不断地在$-\frac{\pi}{4}\le \theta\le \frac{\pi}{4}$这个幅度之内摆动，而不趋于任何极限位置，即切线不存在（图\ref{fig:x*sin(1/x)}）．
\end{solution}

\begin{figure}[htp]
    \centering
\begin{tikzpicture}[>=latex, scale=3]
\draw[->](-1.8,0)--(1.8,0)node[right]{$x$};
\draw[->](0,-.5)--(0,1.5)node[right]{$y$};
\draw[domain=.02:1.8, very thick, samples=300]plot(\x, {\x*sin(1/\x r)});
\draw[domain=.02:1.8, very thick, samples=300]plot(-\x, {\x*sin(1/\x r)});
\draw[thick](-.4,-.4)--(1.3,1.3);
\draw[thick](.4,-.4)--(-1.3,1.3);
\draw[thick](15:1.5)--(-165:1);
\draw[dashed](-1.8,1)--(1.8,1);

\foreach \x in {1,2,3,4}
{
   \draw(\x/pi,0)node[below]{$\frac{\x}{\pi}$}--(\x/pi,.04);
}

\node at (0,0)[above]{$O$};
\end{tikzpicture}
    \caption{}\label{fig:x*sin(1/x)}
\end{figure}

通过以上的例子，我们知道函数的连续点未必是它的可以微分点．几种常见的不可微分点如例2.4的角点，平行于$y$轴的切线的切点，特别是曲线上的尖点和在例2.5的曲线上的不存在割线的极限位置的点．

假如区间$[a,b]$上的函数$f(x)$, 在开区间$(a,b)$中存在着导数$f'(x)$, 并且$f'_+(a)$和$f'_-(b)$都存在，则称$f(x)$在闭区间上可以微分（可导）．

\begin{ex}
\begin{enumerate}
    \item 说明下面函数在点$x=0$不可微：
\begin{multicols}{2}
\begin{enumerate}
    \item $f(x)=|x|$
    \item $f(x)=\begin{cases}
        x^2,& x\le 0\\
        x,& x\ge 0
    \end{cases}$
    \item $f(x)=\sqrt[3]{x}$
    \item $f(x)=\sqrt{|x|}$
\end{enumerate}
\end{multicols}
\item 说明函数$f(x)=\begin{cases}
    x^2,& x\ge 0\\
    -x^2,& x<0
\end{cases}$
    处处可微．
\item \begin{enumerate}
\item 假设$g(x)=f(x+c)$. 从定义出发证明$g'(x)=f'(x+c)$, 并绘图说明．    
\item 设$f$可微并有周期$\varphi$, 证明$f'$也是有周期的．
\end{enumerate}

\end{enumerate}
\end{ex}

\section{微商运算的基本法则}
\subsection{几个基本函数的微商（导数）}
\subsubsection{常数函数的导数恒等于0}

\begin{proof}
设$f(x)=c$ （$c$是常数），则根据常数函数的性质，由自变数$x$的改变量$\Delta x$引起的相应的函数的改变量为：
\[\Delta y=f (x+\Delta x) -f (x)=c-c=0\]
显然差商$\frac{\Delta y}{\Delta x}=\frac{0}{\Delta x}=0$, 故
\[f' (x) =\lim_{\Delta x\to 0}\frac{\Delta y}{\Delta x}=
\lim_{\Delta x\to 0} 0=0\]
这个事实是明显的，因为当自变数$x$有增、减时，函数值并无增、减，因此，它的变化率为零．
\end{proof}

\subsubsection{一次函数（直线函数）的导数是常数}

若$f(x)=kx+b$, 则$f'(x)=k$. 

\subsubsection{若$f(x)=x^n$ （$n$是正整数），则$f'(x)=nx^{n-1}$}

\begin{proof}
  由自变量的改变量$\Delta x$引起的函数改变量为
\[\Delta y=f (x+\Delta x) -f (x) = (x+\Delta x)^n-x^n\] 
按牛顿二项式定理展开，得  
\begin{align*}
\Delta y&= x^n+nx^{n-1}\Delta x+\frac{n(n-1)}{1\cdot 2}x^{n-2}(\Delta x)^2+\cdots+nx(\Delta x)^{n-1}+(\Delta x)^n-x^n\\
&=nx^{n-1}\Delta x+\frac{n(n-1)}{1\cdot 2}x^{n-2}(\Delta x)^2+\cdots+nx(\Delta x)^{n-1}+(\Delta x)^n
\end{align*}
其中差商
\[\frac{\Delta y}{\Delta x}=nx^{n-1}+\frac{n(n-1)}{1\cdot 2}x^{n-2}\Delta x+\cdots+nx(\Delta x)^{n-2}+(\Delta x)^{n-1}\]
式中除第一项外，其余各项随$\Delta x\to 0$而趋于0，因此：
\[f'(x)=\lim_{\Delta x\to 0}\frac{\Delta y}{\Delta x}=nx^{n-1}\]
\end{proof}

\subsubsection{若$f(x)=x^{\mu}$ (其中$\mu$是任意实数，函数的定义域依赖于$\mu$), 则$f'(x)=\mu x^{\mu-1}$}

\begin{proof}
    \[\frac{\Delta y}{\Delta x}=\frac{(x+\Delta x)^{\mu}-x^{\mu}}{\Delta x}=x^{\mu-1}\cdot \frac{\left(1+\frac{\Delta x}{x}\right)^{\mu}-1}{\frac{\Delta x}{x}}\]
利用第一章中已经算出的极限
\[\lim_{x\to 0}\frac{(1+x)^{\mu}-1}{x}=\mu\]
就得到
\[f'(x)=\lim_{\Delta x\to 0}\frac{\Delta y}{\Delta x}=\mu x^{\mu-1}\]
\end{proof}

特殊情形：
\begin{enumerate}
    \item 若$f(x)=\frac{1}{x}=x^{-1},\; (x\ne 0)$，则$f'(x)=(-1)x^{-2}=-\frac{1}{x^2}$
    \item 若$f(x)=\sqrt{x}=x^{\tfrac{1}{2}},\; (x\ge 0)$，则$f'(x)=\frac{1}{2}x^{-\tfrac{1}{2}}=\frac{1}{2\sqrt{x}}$
\end{enumerate}

\subsubsection{若$f(x)=\log_a x\; (0<a\ne 1,\; 0<x<+\infty)$，则$f'(x)=\frac{\log_a e}{x}$}

\begin{proof}
\begin{align*}
    \frac{\Delta y}{\Delta x}&=\frac{\log_a(x+\Delta x)-\log_a x}{\Delta x}\\
    &=\frac{1}{x}\cdot \frac{\log_a\left(1+\frac{\Delta x}{x}\right)}{\frac{\Delta x}{x}}\\
    &=\frac{1}{x}\cdot \log_a\left(1+\frac{\Delta x}{x}\right)^{\tfrac{1}{\Delta x/x}}
\end{align*}  

利用极限$\Lim_{t\to 0}(1+t)^{1/t}=e$，以及对数函数是连续函数，得到
\begin{align*}
    f'(x)&=\lim_{\Delta x\to 0}\frac{\Delta y}{\Delta x}=\frac{1}{x}\lim_{\tfrac{\Delta x}{x}\to 0}\log_a\left(1+\frac{\Delta x}{x}\right)^{\tfrac{1}{\Delta x/x}}\\
    &=\frac{1}{x}\log_a\left[\lim_{\tfrac{\Delta x}{x}\to 0}\left(1+\frac{\Delta x}{x}\right)^{\tfrac{1}{\Delta x/x}}\right]\\
    &=\frac{1}{x}\log_a e=\frac{\log_a e}{x}
\end{align*}

特别：当$f(x)=\ln x$时，$f'(x)=\frac{1}{x}$．
\end{proof}

上面的结果表明对数函数（$a>1$时）的增大速度是与自
变数的值成反比的，当自变数无限增大时，增大速度就保持正值而趋向于0. 由于自然对数的导数比较简单，所以在理论研究中常常采用自然对数．

\subsubsection{若$f(x)=a^x\; (0<a\ne 1,\; -\infty<x<+\infty)$，则$f'(x)=a^x\cdot \ln a$}

\begin{proof}
    \[\frac{\Delta y}{\Delta x}=\frac{a^{x+\Delta x}-a^x}{\Delta x}=a^x\cdot \frac{a^{\Delta x}-1}{\Delta x}\]
利用$\Lim_{\Delta x\to 0}a^{\Delta x}=1$，则可设$a^{\Delta x}=1+\alpha$，且当$\Delta x\to 0$时，有$\alpha\to 0$. 在此等式的两边取对数，得：
\[\Delta x=\log_a(1+\alpha)\]
于是
\[\frac{\Delta y}{\Delta x}=a^x\cdot \frac{\alpha}{\log_a(1+\alpha)}=a^x\cdot \frac{1}{\log_a(1+\alpha)^{1/\alpha}}\]
因此
\begin{align*}
    f'(x)&=\lim_{\Delta x\to 0}\frac{\Delta y}{\Delta x}=a^x\cdot \frac{1}{\Lim_{\alpha\to 0}\log_a(1+\alpha)^{1/\alpha}}\\
    &=a^x\cdot \frac{1}{\log_a\left[\Lim_{\alpha\to 0}(1+\alpha)^{1/\alpha}\right]}\\
    &=a^x\cdot \frac{1}{\log_a e}
\end{align*}
又因为$\ln a=\frac{1}{\log_a e}$，所以
\[f'(x)=a^x\cdot \ln a\]

特别，若$f(x)=e^x$，则$f'(x)=e^x$．
\end{proof}

这个事实表明指数函数（当$a>1$时）的增大速度与函数值成正比例，当底数为$e$时，导数的结果特别简单．


\subsubsection{若$f(x)=\sin x$, 则$f'(x)=\cos x$}
\begin{proof}
    \begin{align*}
\frac{\Delta y}{\Delta x}&=\frac{\sin(x+\Delta x)-\sin x}{\Delta x}\\
&=\frac{2\cos\left(x+\frac{\Delta x}{2}\right)\sin\frac{\Delta x}{2}}{\Delta x}\\
&=\cos\left(x+\frac{\Delta x}{2}\right)\cdot \frac{\sin\frac{\Delta x}{2}}{\frac{\Delta x}{2}}
    \end{align*}
因此
\begin{align*}
    \lim_{\Delta x\to 0}\frac{\Delta y}{\Delta x}&=\lim_{\Delta x\to 0}\left[\cos\left(x+\frac{\Delta x}{2}\right)\cdot \frac{\sin\frac{\Delta x}{2}}{\frac{\Delta x}{2}}\right]\\
    &=\lim_{\Delta x\to 0}\cos\left(x+\frac{\Delta x}{2}\right)\cdot \lim_{\Delta x/2\to 0}\frac{\sin\frac{\Delta x}{2}}{\frac{\Delta x}{2}}
\end{align*}
根据$\cos x$的连续性与极限$\Lim_{t\to 0}\frac{\sin t}{t}=1$, 得到
$f' (x) =\cos x$.
\end{proof}

\subsubsection{若$f(x)=\cos x$, 则$f'(x)=-\sin x$}
\begin{proof}
\begin{align*}
\frac{\Delta y}{\Delta x}&=\frac{\cos(x+\Delta x)-\cos x}{\Delta x}\\
&=\frac{-2\sin\left(x+\frac{\Delta x}{2}\right)\cdot \sin\frac{\Delta x}{2}}{\Delta x}\\
&=-\sin\left(x+\frac{\Delta x}{2}\right)\cdot \frac{\sin\frac{\Delta x}{2}}{\frac{\Delta x}{2}}
    \end{align*}
因此
\begin{align*}
    \lim_{\Delta x\to 0}\frac{\Delta y}{\Delta x}&=-\lim_{\Delta x\to 0}\left[\sin\left(x+\frac{\Delta x}{2}\right)\cdot \frac{\sin\frac{\Delta x}{2}}{\frac{\Delta x}{2}}\right]\\
    &=-\lim_{\Delta x\to 0}\sin\left(x+\frac{\Delta x}{2}\right)\cdot \lim_{\Delta x/2\to 0}\frac{\sin\frac{\Delta x}{2}}{\frac{\Delta x}{2}}\\
    &=-\sin  x
\end{align*}
\end{proof}

现在将上面的导数公式列成下表：
\begin{enumerate}
    \item $\frac{\dd c}{\dd x}=0$
    \item $\frac{\dd(kx+b)}{\dd x}=k$
    \item $\frac{\dd x^n}{\dd x}=nx^{n-1}$\quad （$n$为自然数）
    \item $\frac{\dd x^\mu}{\dd x}=\mu x^{\mu-1}$\quad （$\mu$为任意实数）
    \item $\frac{\dd\log_a x}{\dd x}=\frac{\log_a e}{x},\qquad \frac{\dd \ln x}{\dd x}=\frac{1}{x}$
    \item $\frac{\dd a^x}{\dd x}=a^x\cdot \ln a,\qquad \frac{\dd e^x}{\dd x}=e^x$
    \item $\frac{\dd\sin x}{\dd x}=\cos x,\qquad \frac{\dd \cos x}{\dd x}=-\sin x$
\end{enumerate}

\subsection{求导的基本法则}

现在要建立一些求导数的公式，用这些公式求函数的导数．

\begin{blk}
{定理1} 函数和的导数等于各函数的导数的和，即：
\[[f (x) +g (x) ] '=f' (x) +g' (x)\] 
\end{blk}

\begin{proof}
    设$u(x)=f(x)+g(x)$, 任意固定$x$, 任取非零数$\Delta x$，有 
\begin{align*}
    \frac{\Delta u(x)}{\Delta x}&=\frac{[f(x+\Delta x)+g(x+\Delta x)]-[f(x)+g(x)]}{\Delta x}\\
    &=\frac{f(x+\Delta x)-f(x)}{\Delta x}+\frac{g(x+\Delta x)-g(x)}{\Delta x}\\
    &=\frac{\Delta f(x)}{\Delta x}+\frac{\Delta g(x)}{\Delta x}
\end{align*}
令$\Delta x\to 0$，即得：
\begin{align*}
    [f(x)+g(x)]'&=\lim_{\Delta x\to 0}\frac{\Delta u(x)}{\Delta x}\\
    &=\lim_{\Delta x\to 0}\frac{\Delta f(x)}{\Delta x}+\lim_{\Delta x\to 0}\frac{\Delta g(x)}{\Delta x}\\
    &=f'(x)+g'(x)
\end{align*}
\end{proof}

这个法则可以推广到任意有限个函数．

\begin{blk}
{定理2} $[A\varphi (x)]'=A\varphi'(x)$, 其中$A$是常数，即常因数不因微分法而改变．    
\end{blk}

\begin{proof}
设$f(x)=A\varphi(x)$， 
\begin{align*}
    [A\varphi(x)]'&=f'(x)=\lim_{\Delta x\to 0}\left[A\frac{\varphi(x+\Delta x)-\varphi(x)}{\Delta x}\right]\\
    &=A\lim_{\Delta x\to 0}\frac{\varphi(x+\Delta x)-\varphi(x)}{\Delta x}\\
    &=A\varphi'(x)
\end{align*}
\end{proof}

\begin{example}
    求$3x^2-7x-1$的导数．
\end{example}

\begin{solution}
\[  (3x^2-7x-1)'=(3x^2)'+(-7x)'+(-1)'=6x-7\]
\end{solution}    

\begin{blk}{推论}
设$f(x)=a_nx^n+a_{n-1}x^{n-1}+\cdots+a_1x+a_0\; (a_n\ne 0)$，则
\[f'(x)=na_nx^{n-1}+(n-1)a_{n-1}x^{n-2}+\cdots+a_1\]
\end{blk}

\begin{example}
求过抛物线$y=3x^2+6$外一点$P(2,-9)$所引抛物线的两条切线的方程．
\end{example}

\begin{solution}
在点$x$处的抛物线的切线的斜率为
\[k=y'_x=6x\]
故过抛物线上一点$T(x_1,y_1)$的切线方程为
\[y-y_1=6x_1 (x-x_1)\] 

$\because\quad y_1=3x^2_1+6$

$\therefore\quad $切线方程可化简为
\[    y=6x_1x-3x_1^2+6\]
设切线通过$P(2,-9)$点，于是
\[-9=12x_1-3x_1^2+6\]
即：$x_1^2-4x_1-5=0$

$\therefore\quad x_1=-1$或$5$，从而$y_1=9$或81.

故所求切线方程为：
\[y-9=-6(x+1)\quad\text{和}\quad y-81=30(x-5)\]
即：
\[y+6x-3=0\quad\text{和}\quad  y-30x+69=0\]
\end{solution}

\begin{example}
求三次函数$f(x)=ax^3+3bx^2+3cx+d$的图象与$x$轴相切的条件．
\end{example}

    
\begin{solution}
由函数$f(x)=ax^3+3bx^2+3cx+d$求得它的图象任意一点$x=x_0$处的切线斜率为
\[f' (x_0) =3ax^2_0+6bx_0+3c=3 (ax^2_0+2bx_0+c) \]
    
要此曲线与$x$轴相切于$(x_0,0)$点，$x_0$必须满足条件
\begin{numcases}{}
    ax^3_0+3bx^2_0+3cx_0+d=0\\
    ax^2_0+2bx_0+c=0
\end{numcases}
将(2.1)改写成
\begin{equation}
    x_0(ax^2_0+2bx_0+c)+bx_0^2+2cx_0+d=0
\end{equation}
(2.2)代入得
\begin{equation}
    bx_0^2+2cx_0+d=0
\end{equation}
于是解(2.1)和(2.2)即解(2.2)和(2.4)，要方程(2.2)和(2.4)有公共解，必须
\[D=\begin{vmatrix}
    a&2b\\b&2c
\end{vmatrix}=2(ac-b^2)\ne 0\]
于是
\[x^2_0=\frac{c^2-bd}{b^2-ac},\qquad x_0=\frac{ad-bc}{2(b^2-ac)}\; (b^2-ac\ne 0)\]
所以，$a,b,c,d$必须适合条件
\[\begin{cases}
    b^2-ac\ne 0\\
    \frac{c^2-bd}{b^2-ac}=\frac{(ad-bc)^2}{4(b^2-ac)^2}
\end{cases}\]
即
\[\begin{cases}
    b^2-ac\ne 0\\
   (ad-bc)^2 =4(b^2-ac)(c^2-bd)
\end{cases}\]
\end{solution}

\begin{blk}
    {定理3} 二函数之积的导数为两项之和，其中第一项是第一个因子的导数与第二个因子的乘积，而第二项是第二个因子的导数与第一个因子的乘积，即
\[[f (x) \cdot g (x) ] '=f' (x) \cdot g (x) +f (x) \cdot g' (x) \]
\end{blk}

\begin{proof}
    设 $h(x)=f(x)\cdot g(x)$, 任取非零数$\Delta x$, 则
\begin{align*}
    [f(x)\cdot g(x)]'&=\lim_{\Delta x\to 0}\frac{\Delta h(x)}{\Delta x}\\
    &=\lim_{\Delta x\to 0}\frac{f(x+\Delta x)g(x+\Delta x)-f(x)g(x)}{\Delta x}\\
    &=\lim_{\Delta x\to 0}\frac{1}{\Delta x}\Bigl[f(x+\Delta x)g(x+\Delta x)-f(x)g(x+\Delta x)\\
    &\qquad\qquad +f(x)g(x+\Delta x)-f(x)g(x) \Bigr]\\
    &=\lim_{\Delta x\to 0}\frac{f(x+\Delta x)-f(x)}{\Delta x}\cdot g(x+\Delta x)\\
    &\qquad +\lim_{\Delta x\to 0}f(x)\cdot \frac{g(x+\Delta x)-g(x)}{\Delta x}\\
    &=f' (x) \cdot g (x) +f (x) \cdot g' (x) 
\end{align*}
\end{proof}

重复应用这个定理，我们便可以求两个以上的函数之积的导数，例如
\[\frac{\dd }{\dd x}(uvw)=(vw)\frac{\dd u}{\dd x}+u\frac{\dd }{\dd x}(vw)\]
但是
\[\frac{\dd }{\dd x}(vw)=w\frac{\dd v}{\dd x}+v\frac{\dd w}{\dd x}\]
代入上式便得到
\[\frac{\dd }{\dd x}(uvw)=vw\frac{\dd u}{\dd x}+uw\frac{\dd v}{\dd x}+uv\frac{\dd w}{\dd x}\]


\begin{example}
    设$f(x)=\sqrt{x}(x^2-3)$，求$f'(x)$
\end{example}


\begin{solution}
    \begin{align*}
f'(x)&=\left[\sqrt{x}(x^2-3)\right]'\\
&=\left(\sqrt{x}\right)'(x^2-3)+\sqrt{x}(x^2-3)'\\
&=\frac{1}{2\sqrt{x}}(x^2-3)+\sqrt{x}\cdot 2x = \frac{5x^2-3}{2\sqrt{x}}
    \end{align*}
\end{solution}

\begin{example}
    设$f(x)=\frac{1}{2}\sin^2 x+x\sin x+\frac{7}{x^2}$，求$f'(x)$
\end{example}

\begin{solution}
\begin{align*}
    f'(x)&=\left(\frac{1}{2}\sin^2 x\right)'+(x\sin x)'+\left(\frac{7}{x^2}\right)'\\
    &=\frac{1}{2}\left[(\sin x)'\sin x+\sin x(\sin x)'\right]+(\sin x+x\cos x)-\frac{14}{x^3}\\
    &=\sin x\cos x+\sin x+x\cos x-\frac{14}{x^3}
\end{align*}
\end{solution}

\begin{blk}
    {定理4} 两个函数的商的导数，等于分子的导数与分母的积，减去分母的导数与分子的积，再除以分母的平方，即
\[\left(\frac{u(x)}{v(x)}\right)'=\frac{v(x)\cdot u'(x)-u(x)\cdot v'(x)}{[v(x)]^2}\]
\end{blk}

\begin{proof}
    设$f(x)=\frac{u(x)}{v(x)}$，任取非零数$\Delta x$，则
\begin{align*}
    \left(\frac{u(x)}{v(x)}\right)'&=\lim_{\Delta x\to 0}\frac{\Delta f(x)}{\Delta x}=\lim_{\Delta x\to 0}\frac{\frac{u(x+\Delta x)}{v(x+\Delta x)}-\frac{u(x)}{v(x)}}{\Delta x}\\
    &=\lim_{\Delta x\to 0}\frac{1}{\Delta x}\left\{\frac{u(x+\Delta x)\cdot v(x)-u(x)\cdot v(x+\Delta x)}{v(x+\Delta x)\cdot v(x)}\right\}\\
    &=\lim_{\Delta x\to 0}\frac{1}{\Delta x}\left\{\frac{v(x)[u(x+\Delta x)-u(x)]}{v(x)v(x+\Delta x)}-\frac{u(x)[v(x+\Delta x)-v(x)]}{v(x)v(x+\Delta x)}\right\}\\
    &=\frac{v(x)\Lim_{\Delta x\to 0}\frac{u(x+\Delta x)-u(x)}{\Delta x}}{v(x)\Lim_{\Delta x\to 0}v(x+\Delta x)}-\frac{u(x)\Lim_{\Delta x\to 0}\frac{v(x+\Delta x)-v(x)}{\Delta x}}{v(x)\Lim_{\Delta x\to 0}v(x+\Delta x)}\\
    &=\frac{v(x) u'(x)-u(x) v'(x)}{v^2(x)}
\end{align*}
\end{proof}


\begin{example}
求证：$(\tan x)'=\sec^2 x,\qquad (\cot x)'=-csc^2 x$
\end{example}

\begin{proof}    
\begin{align*}
    (\tan x)'&=\left(\frac{\sin x}{\cos x}\right)'\\
    &=\frac{\cos x(\sin x)'-\sin x(\cos x)'}{\cos^2 x}\\
    &=\frac{\cos^2 x+\sin^2 x}{\cos^2x}\\
    &=\frac{1}{\cos^2 x}=\sec^2 x
\end{align*}
\[(\cot x)'=\left(\frac{\cos x}{\sin x}\right)'=\frac{-\sin^2 x-\cos^2 x}{\sin^2 x}=-\csc^2 x\]    
\end{proof}

\begin{example}
    求$\left(\frac{2x-1}{x^2+1}\right)'$
\end{example}

\begin{solution}
\begin{align*}
    \left(\frac{2x-1}{x^2+1}\right)'&=\frac{(x^2+1)(2x-1)'-(2x-1)(x^2+1)'}{(x^2+1)^2}\\
&=\frac{(x^2+1)\cdot 2-(2x-1)(2x)}{(x^2+1)^2}\\
&=\frac{2(x^2+1)-(4x^2-2x)}{(x^2+1)^2}\\
&=\frac{-2(x^2-x-1)}{(x^2+1)^2}
\end{align*}
\end{solution}

\begin{ex}
\begin{enumerate}
    \item  求下面各函数的导函数：
\begin{multicols}{2}
\begin{enumerate}
    \item $3 x^{3}-1$
    \item  $x+2 x^{2}+3 x^{3}$
    \item  $1-\frac{1}{x^{2}}$
    \item $\frac{x^{2}+1}{x}$
    \item  $1+\frac{1}{x}-\frac{1}{x^{2}}$
    \item  $\frac{2 x^{5}-7 x^{3}-3 x^{2}}{6 x^{2}}$
    \item  $(x+2)\left(x^{2}+1\right)$
    \item  $x+\sqrt{x}$
    \item $x-\frac{1}{\sqrt{x}}$
    \item  $\frac{x+1}{\sqrt{x}}$
    \item  $\frac{2 x \sqrt{x}-x^{\tfrac{5}{2}}+3 x^{\tfrac{1}{2}}}{-\sqrt{x}}$ \item  $e^{x}-e^{-x}$
    \item $\frac{x^{n}-x^{-n}}{x} $
    \item $\sin x+\cos x$
    \item $\sin x-\cos x$
\end{enumerate}
\end{multicols}
\item  求下面各函数的导函数：
\begin{multicols}{2}
\begin{enumerate}
    \item $\left(x^{2}+1\right)^{3} \sqrt{x}$
    \item $\frac{x+1}{\sqrt{x}}$
    \item  $x \sin x$
    \item $x \cot  x$
    \item  $\frac{\tan  x}{x}$;
    \item  $x \log _{2} x+x^{2} \tan  x$
    \item  $\frac{x^{2}}{x^{3}+c^{3}}$
    \item  $x^{3} e^{-x}$;
    \item  $a^{-x} \sin x$ 
    \item $a^{x}\{(x+1) \cos x+\log_2 x\}$
    \item $\frac{\sqrt{x}+1}{2 x+1}$ 
    \item $\frac{2 x-1}{x^{2}+1}$
    \item $\frac{(x-3)(x-4)}{(x-5)}$
    \item  $\frac{x^{2}\left(x^{2}-1\right)\left(x^{3}-1\right)}{x+1}$
    \item  $\frac{x}{x+1}-\frac{1}{x-1}$
    \item  $x \sin x \cos x$
    \item $\sec x $ 
    \item $\csc x$
    \item $\frac{\sin x}{1+\tan  x}$
\end{enumerate}
\end{multicols}

\item 已知抛物线$y=x^2-x$上一点的切线平行于直线$y=x$．求此点的坐标．
\item 求曲线$y=x^3-x^2$上这样的点，使过该点的切线与$x$轴平行．
\item 若抛物线$y=ax^2+bx+c$通过原点，且过该点的切线的斜率等于2, 又抛物线过$(1, 1)$点，求$a$、$b$、$c$.
\item 对于函数$f(x)=x^3+ax^2+bx+c$, 方程$f(x)+\frac{1}{2}x=0$
有不相等的三根$2,\alpha,\beta$, 若$f'(\alpha)=f'(\beta)$, $f'(1)=0$, 求$a,b,c$.
\item 求三次函数$y=f(x)$, 使它同时满足下面的条件：
\begin{enumerate}
    \item 用$x+1$去除$f(x)$与用$(x-1)(x-2)$去除，$f(x)$所得余式相同；
    \item 过曲线$y=f(x)$上的点$(1,f(1))$的切线方程是
$y=-2x+1$
\end{enumerate}
\item 
过$(2, 0)$点求与曲线$y=\frac{1}{x}$相切的直线方程．
\item 
\begin{enumerate}
\item 在抛物线$y=x^2$上求这样一点，使过该点的切线与过$(a,a^2)$点的切线垂直；
\item 求抛物线两条正交切线交点的轨迹方程．
\end{enumerate}

\item 设抛物线方程为$y=x^2+ax+b$, 试问点$(x_0,y_0)$位于
何处时，可以从点$(x_0,y_0)$对此抛物线作出两条切线或一条切线，或作不出切线？
\item 求抛物线$y=x^2+ax$, $y=x^2+bx\; (a\ne b)$的公切线的方
程．
\end{enumerate} 
\end{ex}

\subsection{复合函数的求导法则}
\begin{blk}
    {定理} 假设函数$y=g(x)$在点$x$可导，而函数$z=f(y)$在点$g(x)$可导，那么复合函数$z=\varphi(x)=f(g(x))$在点$x$可导，并且
\[[f(g(x))]'=f'(g(x))\cdot g'(x)\]
简写成$[f(g(x))]'=f'_g\cdot g'_x$，或者写成
\[\frac{\dd z}{\dd x}=\frac{\dd z}{\dd y}\cdot \frac{\dd y}{\dd x}\]
\end{blk}

\begin{proof}
    设$D_g$, $D_f$分别是函数$g(x),f(y)$的定义域，令$x\in D_g$, $g(x)\in D_f$, 任取非零数$\Delta x$, 且使$g(x+\Delta x)\in D_f$，于是$\Delta y=g(x+\Delta x)-g(x)$, 又根据$g(x)$在点$x$连续，故当$\Delta x\to 0$时，$\Delta y\to 0$. 因为$g(x)$在点$x$可导，可以设
\begin{equation}
    \alpha=\frac{\Delta y}{\Delta x}-g'(x)
\end{equation}
而且由上式知道
\[\lim_{\Delta x\to 0}\alpha=\lim_{\Delta x\to 0}\left(\frac{\Delta y}{\Delta x}-g'(x)\right)=g'(x)-g'(x)=0\]
把(2.5)改写成
\begin{equation}
    \Delta y=g'(x)\Delta x+\alpha\Delta x
\end{equation}
其中$\alpha$随$\Delta x$一同趋向于零，且$\alpha \Delta x$是比$\Delta x$更小的量．

\begin{enumerate}
    \item 若$g'(x)\ne 0$，只要$\Delta x$够小，由(2.6)知$\Delta y\ne 0$，那么
\begin{align*}
    \frac{\varphi(x+\Delta x)-\varphi(x)}{\Delta x}&=\frac{f(g(x+\Delta x))-f(g(x))}{g(x+\Delta x)-g(x)}\cdot \frac{g(x+\Delta x)-g(x)}{\Delta x}\\
    &=\frac{f(g(x)+\Delta y)-f(g(x))}{\Delta y}\cdot \frac{\Delta y}{\Delta x}
\end{align*}
令$\Delta x\to 0$，由上式得
\[\varphi'(x)=f'(g(x))\cdot g'(x)\]
即
\[[f(g(x))]'=f'(g(x))\cdot g'(x)\]

\item 如果$g'(x)=0$，则当$\Delta x\to 0$时，有两种情形：
\begin{enumerate}
    \item 若$\Delta y=g(x+\Delta x)-g(x)=0$，则
\begin{align*}
    \frac{\varphi(x+\Delta x)-\varphi(x)}{\Delta x}&=\frac{f(g(x+\Delta x))-f(g(x))}{\Delta x}\\
    &=\frac{f(g(x))-f(g(x))}{\Delta x}=0
\end{align*}
\item 若$\Delta y\ne 0$，则
\[\frac{\varphi(x+\Delta x)-\varphi(x)}{\Delta x}=\frac{f(g(x)+\Delta y)-f(g(x))}{\Delta y}\cdot \frac{\Delta y}{\Delta x}\]
令$\Delta x\to 0$，得到
\[\varphi'(x)=f'(g(x))\cdot 0=0\]
\end{enumerate}
合并上述两种情形，都有
\[\varphi'(x)=[f(g(x))]'=0\]
于是定理得证．
\end{enumerate}
\end{proof}

这个定理是说：复合函数对自变量的导数，等于已知函数对中间变量的导数，乘以中间变量对自变量的导数．


\begin{example}
    设$\varphi (x)=2^{\sin x}$, 求$\varphi'(x)$．
\end{example}

    
\begin{solution}
    把$\varphi (x)$看成$u=\sin x$和$f(u)=2^u$的复合函数，于是$\varphi  (x) =2^{\sin x}=f (\sin x)$.

依上述定理，得
\begin{align*}
  \varphi' (x) &= [f (u) ]'\cdot u'_x= (2^u)'\cdot (\sin x)'\\
&=2^u\ln2 \cdot \cos x\\
&=2^{\sin x}\cdot \cos x \cdot \ln2  
\end{align*}
\end{solution}

\begin{example}
    设$y=\sin(x^2+x+1)\cdot \cos(x^3+2x^2)$, 求$y'_x$.
\end{example}

\begin{solution}
    设$u=x^2+x+1$, $v=x^3+2x^2$, 则
\begin{align*}
  y'&= (\sin u)'\cdot u'_x\cdot \cos v+\sin u\cdot (\cos v)'\cdot v'_x\\
&=\cos u\cdot u_x'\cdot \cos v-\sin u\cdot \sin v\cdot  v'_x\\
&= (2x+1) \cos (x^2+x+1) \cos (x^3+2x^2)\\
&\qquad - (3x^2+4x) \sin (x^2+x+1) \sin (x^3+2x^2)  
\end{align*}
\end{solution}
    
\begin{example}
设$f(x)=\frac{x}{\sqrt{9+x^2}}$，求$f'(x)$.
\end{example}


\begin{solution}
\begin{align*}
    f'(x)&=\frac{\sqrt{9+x^2}(x)'-x\left(\sqrt{9+x^2}\right)'}{9+x^2}\\
    &=\frac{\sqrt{9+x^2}-x\cdot \frac{1}{2}(9+x^2)^{-\tfrac{1}{2}}(9+x^2)'}{9+x^2}\\
    &=\frac{\sqrt{9+x^2}-\frac{x^2}{\sqrt{9+x^2}}}{9+x^2}=\frac{9}{\left(9+x^2\right)^{3/2}}
\end{align*}
\end{solution}
    
\begin{blk}
{推论1} 如果函数有两个以上的中间变量，上面的求导法则可以推广使用．
\end{blk}

例如，设$y=f(u)$, $u=g(v)$, $v=h(x)$, 那么对于复合函数．$y=\varphi (x)=f\{g[h(x)]\}$的导数可以这样来求：
\[\frac{\dd y}{\dd x}=\frac{\dd y}{\dd u}\cdot \frac{\dd u}{\dd v}\cdot \frac{\dd v}{\dd x}\]
即：$y'_x=y'_u\cdot u'_v\cdot v'_x$

事实上，依前定理有
\[\frac{\dd y}{\dd x}=\frac{\dd y}{\dd u}\cdot \frac{\dd u}{\dd x}\]
这里$u=g[h(x)]$．再用同一定理
\[\frac{\dd u}{\dd x}=\frac{\dd u}{\dd v}\cdot \frac{\dd v}{\dd x}\]
所以
\[\frac{\dd y}{\dd x}=\frac{\dd y}{\dd u}\cdot \frac{\dd u}{\dd v}\cdot \frac{\dd v}{\dd x}\]


\begin{example}
    设$y=\ln\sqrt{5-2x+3x^4},\quad (5-2x+3x^4>0)$，求$y'_x$.
\end{example}

    
\begin{solution}
    设$y=\ln u$，$u=\sqrt{v}$，$v=5-2x+3x^4$，于是：
\begin{align*}
    \frac{\dd y}{\dd x}&=\frac{\dd y}{\dd u}\cdot \frac{\dd u}{\dd v}\cdot \frac{\dd v}{\dd x}\\
    &=\frac{1}{u}\cdot \left(\frac{1}{2}v^{-\tfrac{1}{2}}\right)\cdot (-2+12x^3)\\
    &=\frac{1}{\sqrt{5-2x+3x^4}}\cdot \frac{1}{2}\frac{1}{\sqrt{5-2x+3x^4}}\cdot (-2+12x^3)\\
    &=\frac{6x^3-1}{5-2x+3x^4}
\end{align*}
\end{solution}

求复合函数的导数，关键在于分析清楚函数的复合关系，适当选定中间变量，根据复合函数求导法则由外向里逐层求导，直到最后一个中间变量对自变量求导为止，每次求导时，必须明确是哪个变量对哪个变量求导，所有这些导数的乘积就是复合函数的导数．

\begin{blk}{推论}
    \[(\ln|g(x)|)'=\frac{g'(x)}{g(x)}\quad (g(x)\ne 0)\]
\end{blk}

\begin{proof}
  \[\ln|g(x)|\text{有意义}\quad \Longleftrightarrow\quad g(x)\ne 0\]  
又
\[|g(x)|'\text{存在}\quad \Longleftrightarrow\quad g(x)\ne 0\]
所以由
\[|g(x)|'=\begin{cases}
    g'(x) ,  & g(x)>0\\
    [-g(x)]'=-g'(x), & g(x)<0
\end{cases}\]
和复合函数求导法则，便得
\[[\ln|g(x)|]'=\frac{1}{|g(x)|}\cdot |g(x)|'=\frac{g'(x)}{g(x)}\quad (g(x)\ne 0)\]
\end{proof}

\begin{example}
    求$(\ln|x^2-5x+4|)'$
\end{example}

\begin{solution}
\[(\ln|x^2-5x+4|)'=\frac{(x^2-5x+4)'}{x^2-5x+4}=\frac{2x-5}{x^2-5x+4}\]
其中：$x^2-5x+4\ne 0$，即$x\ne 1$和$x\ne 4$.
\end{solution}



\begin{example}
求$\left[\frac{(x^2+1)^{\tfrac{1}{2}}(x-1)^2}{(x+1)^{\tfrac{3}{2}}}\right]'$
\end{example}

    
\begin{solution}
设$y=\frac{(x^2+1)^{\tfrac{1}{2}}(x-1)^2}{(x+1)^{\tfrac{3}{2}}}$，则
\begin{align*}
    \ln|y|&=\frac{1}{2}\ln(x^2+1)+2\ln|x-1|-\frac{3}{2}\ln|x+1|\\
    \frac{1}{y}\cdot y'&=\frac{1}{2}\frac{2x}{x^2+1}+\frac{2}{x-1}-\frac{3}{2(x+1)}\\
    &=\frac{3x^3+7x^2-x+7}{2(x^2+1)(x-1)(x+1)}
\end{align*}    
\begin{align*}
    y'&=\frac{3x^3+7x^2-x+7}{2(x^2+1)(x-1)(x+1)}\cdot \frac{(x^2+1)^{\tfrac{1}{2}}(x-1)^2}{(x+1)^{\tfrac{3}{2}}}\\
    &=\frac{(3x^3+7x^2-x+7)(x-1)}{2(x^2+1)^{\tfrac{1}{2}}(x+1)^{\tfrac{5}{2}}}
\end{align*}
    
\end{solution}

\begin{example}
设$f(x)=u(x)^{v(x)}\quad (u(x)>0)$，求$f'(x)$.
\end{example}

\begin{solution}
先求此函数的对数，有
\[\ln f(x)=v(x)\ln u(x)\]
两边对$x$求导，即得
\[\frac{1}{f(x)}\cdot f'(x)=v'(x)\ln u(x)+v(x)\cdot\frac{1}{u(x)}\cdot u'(x)\]
    所以
\[f'(x)=u(x)^{v(x)}\left[v'(x)\ln u(x)+\frac{v(x)u'(x)}{u(x)}\right]\]
\end{solution}

\begin{example}
一人在放风筝，且知风筝在80米的高度飘行，放线速度为2米/秒，求当线长为100米时，风筝和此人之间的水平距离的变率．
\end{example}


\begin{solution}
    设在某时刻$t$, 线长为$\ell(t)$米，水平距离为$x(t)$米，根据勾股定理得到绳长和水平距离的关系式
\[\ell^2 (t) =x^2 (t) +6400\]
两边对时间$t$求导数，得到
\[\frac{\dd \ell^2}{\dd \ell}\cdot \frac{\dd \ell}{\dd t}=\frac{\dd(x^2+6400)}{\dd x}\cdot \frac{\dd x}{\dd t}\]
即
\[\ell\cdot \frac{\dd \ell}{\dd t}=x\cdot \frac{\dd x}{\dd t}\]
但是已知$\frac{\dd\ell}{\dd t}=2({\rm m/s})$，又当$\ell=100$时，
\[x=\sqrt{100^2-6400}=\sqrt{3600}=60\]
于是
\[60\frac{\dd x}{\dd t}=100\times 2\]

$\therefore\quad \frac{\dd x}{\dd t}=\frac{10}{3}=3\frac{1}{3}({\rm m/s})$

答：当$\ell=100$米时，风筝和人之间的水平距离的变率为$3\frac{1}{3}$（米／秒）．
\end{solution}

\begin{ex}
\begin{enumerate}
    \item 求下面函数对于 $x$ 的导数:
\begin{multicols}{2}
\begin{enumerate}
    \item $(x+1)^{2} $ 
\item $(1-3 x)^{3} $
\item  $\sqrt{x+1}$
\item  $\frac{1}{\sqrt{x+1}}$
\item  $\sqrt{x^{2}+1}$
\item  $\frac{1}{\sqrt{x^{2}+1}}$
\item  $\sqrt{2 x^{2}+x-1}$
\item  $\frac{1}{\sqrt{\tan x}}$
\item  $\frac{1}{\sin ^{3} x} $
\item  $\cos (a x+b)$
\item  $\cos ^{3}(a x+b) $
\item  $\cot \sqrt{x} $
\item  $\sec \left(\frac{1}{x}\right)$
\item  $x \sqrt{\sin x}$
\item  $x^{2} \sqrt{\sec 2x}$
\item  $\ln \left(1+x^{2}\right)$
\item $\log _{a}(\sin x)$
\item $\log _{a}(\sec x)$
\item  $\sqrt{\ln x}$
\item $e^{a x^{2}}$
\item $e^{1+x-x^{2}}$
\item  $a^{\sqrt{2x+1}}$
\end{enumerate}
\end{multicols}

 \item    对下面各函数先取对数再求导数:
 \begin{multicols}{2}
    \begin{enumerate}
        \item   $\frac{x-1}{(x+1)^{2}}$
        \item   $\frac{\sqrt{x-1}(x+1)}{x^{2}+1}$
        \item   $\frac{(x-3)(x-4)}{(x-5)}$
        \item   $\sqrt{\frac{x-1}{x+1}}$ 
        \item  $(x+2)^{2}(x+3)^{3}(x+4)^{4}$
\end{enumerate}
\end{multicols}
\item  求下列函数对于 $x$ 的微商:
\begin{enumerate}
\begin{multicols}{2}
  \item $e^{2 x}(2 \cos 3 x+3 \sin 3 x)$
\item $\ln \left(x+\sqrt{x^{2} \pm a^{2}}\right)$
\item $\frac{2 x+1}{\sqrt{3-4 x-4 x^{2}}}$  
\item $\cos[\ln(x^2+9)]$
\item $\ln(1-\tan^2 x)$
\item $x^x+(\ln x)^x$
\end{multicols}
\item $(1-x)^{\tfrac{1}{2}}(1+3x)^{\tfrac{3}{20}}(1-2x)^{-\tfrac{2}{5}}$
\end{enumerate}

\item 若$y=e^{ax}\cos(bx+c)$，求证：
\[\frac{1}{\sqrt{a^2+b^2}}\frac{\dd y}{\dd x}=e^{ax}\cos(bx+c+r)\]
其中：$\tan r=b/a$.

\item 设$f(x)=\begin{cases}
    x^2\sin\frac{1}{x},& x\ne 0\\
    0,& x=0
\end{cases}$，求$f(x)$的导函数．

（提示：对于$x\ne 0$，可用本节的方法得出，而$f'(0)$需直接由定义得出）．

\item 有一个长度为5米的梯子贴靠在铅直的墙上，假设其下端
沿地板以3米/秒的速率离开墙脚而滑动，则
\begin{enumerate}
    \item 当其下端离开墙脚1.4米时，梯子上端下滑之速率为多少？
    \item 何时梯子的上下端能以相同的速率移动？
    \item 何时其上端下滑之速度为4米/秒？
\end{enumerate}

\item 求曲线$y=\cos^2\sqrt{x}$在点$x=0$处的切线方程．
\item 在$\triangle ABC$中，$a=7$cm, $b=5$cm, $c=8$cm, $P$为$AB$边上
一动点，从$A$往$B$以每秒2厘米的速度运动．当$P$点在
$AB$边中点时，$CP$长度的变率是多少？
\item 求曲线$y=2\left(e^x+\frac{1}{5}e^{-3x}\right)$在横坐标$x_0=0$的点的切线方程．
\item 试由公式$1+x+x^2+\cdots+x^n=\frac{x^{n+1}-1}{x-1}$分别导出$1+2x+3x^2+\cdots+nx^{n-1}$与$x+2^2x^2+3^2x^3+\cdots +n^2x^n$的公式．
\end{enumerate}  
\end{ex}


\subsection{反函数求导法则}

\begin{blk}
    {定理} 若函数$y=f(x)$在$x_0$的某一邻域内连续，且严格单调，又函数$f(x)$在点$x_0=f^{-1}(y_0)$的导数$f'(x_0)$存在并且异于零，则反函数$x=f^{-1}(y)$在对应点$y_0=f(x_0)$
的导数存在，并且等于$\frac{1}{f'(x_0)}=\frac{1}{f'(f^{-1}(y_0))}$
，即
\[[f^{-1}(y_0)]'=\frac{1}{f'(x_0)}=\frac{1}{f'[f^{-1}(y_0)]}\]
\end{blk}

\begin{proof}
给数值$y=y_0$以任意增量$\Delta y$, 则反函数$x=f^{-1}(y)$也得到对应增量$\Delta x$, 当$\Delta y\ne 0$时，由于$y=f(x)$的严格单调性，从而$x=f^{-1}(y)$也是严格单调的，也有$\Delta x\ne 0$, 于是有
\begin{align*}
    \frac{f^{-1}(y_0+\Delta y)-f^{-1}(y_0)}{\Delta y}&=\frac{\Delta x}{f(x_0+\Delta x)-f(x_0)}\\
    &=\frac{1}{\frac{f(x_0+\Delta x)-f(x_0)}{\Delta x}}
\end{align*}
现在令$\Delta y\to 0$，则由于反函数$x=f^{-1}(y)$也是连续的，故增量$\Delta x\to 0$，于是
\[\lim_{\Delta y\to 0}\frac{f^{-1}(y_0+\Delta y)-f^{-1}(y_0)}{\Delta y}=\frac{1}{\Lim_{\Delta y\to 0}\frac{f(x_0+\Delta x)-f(x_0)}{\Delta x}}\]
上面等式右边的分母就趋于极限$f'(x_0)\ne 0$，而左边按定义就是$[f^{-1}(y_0)]'$，因此
\[[f^{-1}(y_0)]'=\frac{1}{f'(x_0)}\]
再将$x_0=f^{-1}(y_0)$代入上式，得到
\[[f^{-1}(y_0)]'=\frac{1}{f'[f^{-1}(y_0)]}\]
\end{proof}

\begin{blk}{推论}
\begin{align*}
    (\arcsin x)'&=\frac{1}{\sqrt{1-x^2}},\qquad -1<x<1\\
    (\arccos x)'&=\frac{-1}{\sqrt{1-x^2}},\qquad -1<x<1\\
    (\arctan x)'&=\frac{1}{{1+x^2}},\qquad -\infty<x<\infty\\
    (\arccot x)'&=\frac{-1}{{1+x^2}},\qquad -\infty<x<\infty\\
\end{align*}
\end{blk}

\begin{proof}
如果我们限制正弦函数$x=\sin y$的定义域为$\left(-\frac{\pi}{2},\frac{\pi}{2}\right)$，那么$x=\sin y$的反函数存在，即
\[    y=\arcsin x,\qquad x\in (-1, 1)\]
于是
\[    (\arcsin x)'=\frac{1}{(\sin y)'}=\frac{1}{\cos y}=\frac{1}{\cos(\arcsin x)}\]
由于$\arcsin x\in \left(-\frac{\pi}{2},\frac{\pi}{2}\right)$，
所以$\cos y>0$, 从而
\[    \cos (\arcsin x) =\sqrt{1-\sin^2 (\arcsin x)} =\sqrt{1-x^2}\]
    最后得到
\[  (\arcsin x)'=\frac{1}{\sqrt{1-x^2}},\qquad x\in (-1, 1) \]

我们除去了数值$x=\pm 1$, 因为在它的对应值$y=\pm\frac{\pi}{2}$
处导数$(\sin y)'=\cos y=0$.

在第四册， 我们已经证明$\arcsin x+\arccos x=\frac{\pi}{2}$，所以
\[(\arccos x)'=\left(\frac{\pi}{2}-\arcsin x\right)'=-\frac{1}{\sqrt{1-x^2}},\qquad -1<x<1\]
其余同理可证．
\end{proof}

\begin{example}
求$(\arcsec x)'$.
\end{example}


\begin{solution}
函数$x=\sec y$在区间$\left(0,\frac{\pi}{2}\right)$上，由1递增到$+\infty$，而在区间$\left(\frac{\pi}{2},\pi\right)$上，由$-\infty$递增到$-1$．因此，$x=\sec y$，$y\in\left(0,\frac{\pi}{2}\right)\cup \left(\frac{\pi}{2},\pi\right)$有反函数存在，即：
\[y=\arcsec x,\qquad x\in (-\infty,-1)\cup(1,+\infty)\]

依反函数求导法则，有
\begin{align*}
    (\arcsec x)'&=\frac{1}{(\sec y)'}=\frac{1}{\sec y\cdot \tan y}\\
    &=\frac{1}{\sec(\arcsec x)\cdot \tan(\arcsec x)}\\
    &=\begin{cases}
        \frac{1}{x\sqrt{\sec^2(\arcsec x)-1}}, & 0<\arcsec x<\frac{\pi}{2}\\
        \frac{1}{x\left(-\sqrt{\sec^2(\arcsec x)-1}\right)}, & \frac{\pi}{2}<\arcsec x<\pi \\
    \end{cases}\\
    &=\begin{cases}
        \frac{1}{x\sqrt{x^2-1}}, & x>1\\
       - \frac{1}{x\sqrt{x^2-1}}, & x<-1
    \end{cases}
\end{align*}
\end{solution}

\begin{example}
若已知$(\log_a x)'=\frac{\log_a e}{x},\quad x\in(0,+\infty)$，试利用反函数求导法则求$(a^x)'$    
\end{example}

\begin{solution}
    因为$y=a^x,\quad (a>0,\; a\ne 1,\; -\infty<x<+\infty)$是$x=\log_a y\; (y>0)$的反函数，所以
\begin{align*}
    (a^x)'&=\frac{1}{(\log_a y)'}=\frac{1}{\frac{\log_a e}{y}}=\frac{y}{\log_a e}\\
    &=\ln a\cdot (a^x)=a^x\ln a
\end{align*}
\end{solution}

\begin{example}
求函数$f(x)=2x\arctan 2x-\ln\sqrt{1+4x^2}$的导数．
\end{example}

\begin{solution}
\begin{align*}
    f'(x)&=\left[(2x)'\arctan 2x+2x\cdot \frac{1}{1+(2x)^2}\cdot (2x)'\right]\\
    &\qquad -\frac{1}{\sqrt{1+4x^2}}\cdot \frac{1}{2\sqrt{1+4x^2}}\cdot (1+4x^2)'\\
    &=2\arctan 2x+\frac{4x}{1+4x^2}-\frac{4x}{1+4x^2}\\
    &=2\arctan 2x
\end{align*}
\end{solution}


\begin{example}
求曲线$y=\arcsin x$和$y=\arccos x$的交角．
\end{example}

\begin{solution}
设这两条曲线交于$T(x_0,y_0)$点，则
\[y_0=\arcsin x_0=\arccos x_0,\qquad 0<y_0<\frac{\pi}{2}\]
  由此得
\begin{equation}
    \sin(\arcsin x_0)=\sin(\arccos x_0)
\end{equation}  
即
\[x_0=\sqrt{1-x^2_0}\]
解得
\[x_0=\pm\frac{1}{\sqrt{2}}\]
因为$0<y_0<\frac{\pi}{2}$，所以$x_0>0$. 如此应取$x_0=\frac{1}{\sqrt{2}}$. 由于
\[(\arcsin x)'=\frac{1}{\sqrt{1-x^2}},\qquad (\arccos x)'=\frac{-1}{\sqrt{1-x^2}}\]
所以在点$x=\frac{1}{\sqrt{2}}$处，$y=\arcsin x$的切线的斜率为
\[\tan\alpha\Big|_{x=\tfrac{1}{\sqrt{2}}}=\frac{1}{\sqrt{1-\frac{1}{2}}}=\sqrt{2}\]
又$y=\arccos x$的切线的斜率是
\[\tan\beta\Big|_{x=\tfrac{1}{\sqrt{2}}}=\frac{-1}{\sqrt{1-\frac{1}{2}}}=-\sqrt{2}\]

设$\theta$是这两条曲线在交点$\left(\frac{1}{\sqrt{2}},\frac{\pi}{4}\right)$的夹角，也就是过$T$点的两条切线的交角，于是
\[\tan\theta=\frac{\tan\beta-\tan\alpha}{1+\tan\beta\tan\alpha}=\frac{-\sqrt{2}-\sqrt{2}}{1+\left(-\sqrt{2}\right)\sqrt{2}}=\frac{-2\sqrt{2}}{-1}=2\sqrt{2}\]
因此，这两条曲线的交角$\theta=\arctan 2\sqrt{2}$. 
\end{solution}

\begin{example}
一飞机在离地面2公里的高度，以每小时200公里的速度飞临某目标的上空，以便进行航空摄影，试求飞机飞至该目标正上方时摄影机转动的角速度（图2.9）．
\end{example}

\begin{solution}
我们把该目标取为坐标原点，设飞机和目标的水平距离为$x$, 显然它是时间$t$的函数$x=x(t)$. 现在要求出飞机在目标正上方时，$\frac{\dd \theta}{\dd t}$之值，
为此，先找出$\theta$和$x$的关系，从图2.9中可以看出
\[\tan\theta=\frac{2}{x},\quad  0<\theta\le \frac{\pi}{2}\]

$\therefore\quad \theta=\arctan \frac{2}{x}$
于是
\[\frac{\dd \theta}{\dd t}=\frac{-2}{x^2+4}\frac{\dd x}{\dd t}\]
但根据题意知道$\frac{\dd x}{\dd t}=-200$, 这里负号表示$x$在减小，故得
\[\frac{\dd \theta}{\dd t}=\frac{400}{x^2+4}\]
当飞机在目标正上方时，即$x=0$时，
$\frac{\dd \theta}{\dd t}=100\text{弧度/时}$，将弧度化为角度即为$\frac{\dd \theta}{\dd t}=\frac{5}{\pi}\text{度/秒}$
\end{solution}

\begin{figure}[htp]
    \centering
    \begin{minipage}[t]{0.48\textwidth}
    \centering
\begin{tikzpicture}[>=latex, scale=.9]
\draw[->](-1,0)--(5,0)node[right]{$x$};
\draw[->](0,-1)--(0,4.5)node[right]{$y$};
\node at (0,0)[below left]{$O$};
\draw[dashed](0,3.5)--(4,3.5);
\draw[very thick](4,3.5)--(0,0);
\draw[<->](4,3.5)--node[fill=white]{2}(4,0);
\draw[<->](1,0) arc (0:41.2:1)node[below right]{$\theta$};
\draw[<->](3,3.5) arc (180:180+41.2:1)node[above left]{$\theta$};
\draw[ultra thick, <-](4,3.5)--(4.6,3.5);
\end{tikzpicture}
    \caption{}
    \end{minipage}
    \begin{minipage}[t]{0.48\textwidth}
    \centering
    \begin{tikzpicture}[>=latex, scale=1.5]
        \draw[very thick](-1,2.5)--(-.25,2.5)--(-.25,2)--(-1.5,0)--(1.5,0)--(.25,2)--(.25,2.5)--(1,2.5);
\draw[|<->|](1.75,0)--node[fill=white]{12cm}(1.75,2.5);
\draw[dashed](1,2.5)--(1.75,2.5);
\draw[|<->|](0,-.15)--node[below]{$R=5$cm}(1.5,-.15);
\draw(-1.75/2,1)--(1.75/2,1);
\draw[<->](0,0)--node[fill=white]{$y$}(0,1);
\draw[|<->|](0,.8)--node[fill=white]{$r$}(1.75/2,.8);
    \end{tikzpicture}
    \caption{}
    \end{minipage}
    \end{figure}

\begin{example}
有一底半径为5cm, 高为12cm的正圆锥形容器（图2.10）, 以每秒2$\rm cm^3$的速度往容器中倒水，试求容器水位等于锥高一半时，水面上升速度．
\end{example}


\begin{solution}
    设在某时刻$t$, 容器水高为$y$, 则此时水的体积$V$为
 \begin{align*}
     V&=\frac{\pi R^2h}{3}-\frac{\pi r^2}{3}(h-y)\\
     &=\frac{\pi}{3}\cdot 5^2\cdot 12-\frac{\pi r^2}{3}(12-y)\\
     &=100\pi-\frac{\pi r^2}{3}(12-y)
 \end{align*}
而
\[\frac{r}{R}=\frac{h-y}{h}\]
即
\[r=\frac{R}{h}(h-y)=\frac{5}{12}(12-y)\]
代入上式，得
\begin{align*}
    V&=100\pi-\frac{\pi}{3}\cdot \frac{5^2}{12^2}(12-y)^3\\
    &=100\pi-\frac{25\pi}{432}(12-y)^3
\end{align*}
又由题设知：$V=2t$，所以
\[2t=100\pi-\frac{25\pi}{432}(12-y)^3\]
即
\begin{equation}
    t=50\pi-\frac{25\pi}{864}(12-y)^3
\end{equation}
这就确定了水的高度与时间$t$的关系，水面上升速度就是
$y$关于时间$t$的变化率，也就是$\frac{\dd y}{\dd t}$．根据反函数求导法则，有
\[\frac{\dd y}{\dd t}=\frac{1}{\frac{\dd t}{\dd y}}\]
由(2.8)得
\[\frac{\dd t}{\dd y}=-\frac{25\pi}{864}\cdot 3(12-y)^2(12-y)'=\frac{25\pi}{288}(12-y)^2\]

$\therefore\quad \frac{\dd y}{\dd t}=\frac{288}{25\pi(12-y)^2}$

当$y=\frac{1}{2}h=6$时，
\[\left.\frac{\dd y}{\dd t}\right|_{y=6}=\frac{288}{25\x 36\pi}\approx 0.102({\rm cm/s})\]

答：水面上升速度为每秒0.102厘米．
\end{solution}

\begin{ex}
\begin{enumerate}
    \item 求下列函数的微商：
\begin{multicols}{2}
\begin{enumerate}
  \item $\arcsin \frac{x}{4}$
\item  $\arccos x^{2}$
\item  $\arctan \frac{1}{x}$
\item  $\arcsin \sqrt{x}$
\item  $(\arcsin x)^{2}$
\item  $\arctan (\tan  x)$
\item $\arctan \frac{2 x}{1-x^{2}}$
\item  $\arctan (2 \tan  x)$
\item $x^{2} \arccot \frac{x}{a}$
\item $\ln \sqrt{1+4 x^{2}} \cdot \arccot \frac{1-x}{1+x}$
\item  $\frac{1-\arccos 2 x}{\arctan x}$
\item  ${\rm arccsc}\, x$
\end{enumerate}
\end{multicols}

\item 证明
   \[
    \frac{\dd}{\dd x}\left\{\arctan \left(\frac{\sin x}{\cos 2 x}\right)\right\}=\frac{3 \cos x-\cos 3 x}{2(1-\sin x \sin 3 x)}
   \]
    \item 求下列函数的导数：
\begin{enumerate}
\item $\arcsin \frac{6 x^{2}}{\sqrt{4+81 x^{8}}}$
\item $\arctan \frac{1}{1-3 x^{2}}-\arctan \frac{1}{1+3 x^{2}}$
\item $\arccot \frac{4 \sqrt{x}}{1-4 x}$
\end{enumerate}

\item 一圆锥形容器，深10尺，上顶圆半径为4尺；求
\begin{enumerate}
    \item 灌入水时，水的体积$V$对水面高度$h$的变化率；
    \item 体积$V$对容器截面圆半径$R$的变化率．
\end{enumerate}

\item 一圆锥形容器底面朝上放着，它的顶角为$2\arctan\frac{3}{4}$, 
今向里面倒进某种液体，试求
\begin{enumerate}
\item 当液面半径$r$为3, 半径增加的速度$\frac{\dd r}{\dd t}$为$\frac{1}{4}$时，
体积增加的速度$\frac{\dd V}{\dd t}$；
\item 当液面半径为6, 体积增加的速度为24时，半径增加的速度．
\end{enumerate}

\item 水从高为18厘米，底半径为6厘米的圆锥形漏斗中流入半
径为5厘米的圆柱形筒内．已知此漏斗中水深为12厘米时，漏斗中水面的下降速度为1厘米/分，求此时圆筒中水面的上升速度．
\item 底数$a$为什么值时，直线$y=x$才能与对数曲线$y=\log_a x$相切？在何处相切？
\item 求曲线$y=e^{-x}$上的一点，使过该点的切线与直线$y=-ex$平行，并写出该点的法线方程．
\item 求曲线$y=\frac{1}{2}(1+2x^2\pm \sqrt{1+4x^2})$上横坐标$x=u$的点
处的切线方程．这切线还与曲线交于何处？$y$的反函数为何？
\end{enumerate}    
\end{ex}

\subsection{隐函数的导数}

到现在为止，我们所讨论的函数都是可以由自变量明显地表示出来的函数，例如，方程
$y=x^3+1$
就定义了一个明显的函数$f$, 这里$f(x)=x^3+1$. 这种函数叫做显函数．但是并不是所有的函数都可以用这种明显的表达式来定义的，例如由下面的$x,y$的方程
\[x^3-x=y^3-y^2+24\]
就不容易解出用$x$表示$y$的解，或用$y$表示$x$的解．但是，却有可能存在一个函数$f$, 使得方程
\[x^3-x=f^3(x)-f^2(x)+24\]
对于函数$f$的定义域中的一切$x$值都成立，这样一个函数，如果存在的话，就叫做被这个方程定义的隐函数．最简的情形是二元一次方程：
\[Ax+By+C=0\]
当$B\ne 0$时，就定义一个隐函数$f:\; \mathbb{R}\mapsto \mathbb{R}$, 这里$f(x)=-\frac{A}{B}x-\frac{C}{B}$. 

应该注意：
\begin{enumerate}
\item 方程$F(x,y)=0$的解有可能不存在，例如方程$x^2+y^2+1=0$没有实数解，因此，由它不能定义一个隐函数．
\item 由方程$F(x,y)=0$可能定义几个函数．例如，由方程
$x^2+y^2=16$
就定义了两个函数：
\[f_1 (x) =\sqrt{16-x^2},\qquad x\in [-4, 4] \]
和
\[f_2 (x) =-\sqrt{16-x^2},\qquad x\in  [-4, 4] \]
\end{enumerate}

\begin{blk}
    {定义} 方程
    \begin{equation}
 F (x,y) =0       
    \end{equation}
给出了变元$x,y$之间的一个关系．如果存在一个函数$y=
f(x)$, 或$x=g(y)$使得对于函数的定义域中的一切值，等式$F(x,f(x))=0$或$F(g(y),y)=0$恒成立，我们就说方程(2.9)定义了一个\textbf{隐函数}：
\[y=f(x)\qquad \text{或}\qquad x=g(y)\]
\end{blk}

要对方程(2.9)给出的隐函数求导数，通常的办法是直接对这个方程求导数，而不是先解出显函数，再来求导数．


\begin{example}
设$xy+x^2-1=0$，求$y'_x$.
\end{example}

\begin{solution}
先把$y$看成$x$的函数，即$y=f(x)$, 然后，在等式
$xf (x) +x^2-1=0$
的两边对$x$求导数，得到
\begin{align*}
    (x\cdot f(x))'+(x^2)'&=0\\
xf'(x)+f(x)+2x&=0
\end{align*}

$\therefore\quad f' (x) =-\frac{2x+f (x)}{x}=-\frac{2x+\frac{1-x^2}{x}}{x}=-\frac{1+x^2}{x^2}$    
\end{solution}

\begin{example}
    方程$x^2+y^2=16$定义了两个可微函数：
$y=\sqrt{16-x^2}$和$y=-\sqrt{16-x^2},\; x\in [-4, 4]$. 求每一个函数的导数．
\end{example}


\begin{solution}
    我们用$f(x)$代表上面两个函数中的任何一个，则
\[x^2+f^2 (x) =16,\qquad  x\in  [-4, 4] \]
对上面等式的两边微分，得到
\begin{align*}
    2x+2f (x) \cdot f' (x) &=0\\
    f'(x)&=-\frac{x}{f(x)}
\end{align*}
即
\begin{align*}
    \left(\sqrt{16-x^2}\right)'&=-\frac{x}{\sqrt{16-x^2}}\\
    \left(-\sqrt{16-x^2}\right)'&=\frac{x}{\sqrt{16-x^2}},\qquad x\in[-4,4]
\end{align*}
\end{solution}

\begin{rmk}
应用这个方法求隐函数的导数时，我们作了两点假设：
\begin{enumerate}
    \item 存在一个由方程$F(x,y)=0$定义的函数；
    \item 这样定义的函数是可微的．
\end{enumerate} 
至于在什么条件下，这种假设是合理的问题，已经超出了本书的范围，故不在此讨论．
\end{rmk}

\begin{example}
求证椭圆$\frac{x^2}{a^2}+\frac{y^2}{b^2}=1$在点$(x_0,y_0)$处切线的方程为$\frac{x_0x}{a^2}+\frac{y_0y}{b^2}=1$.
\end{example}

\begin{proof}
\begin{align*}
    \left(\frac{x^2}{a^2}+\frac{y^2}{b^2}\right)'_x&=(1)'_x\\
    \frac{2x}{a^2}+\frac{2y}{b^2}\cdot y'_x&=0
\end{align*}
若$y\ne 0$，则：
\[y'_x=-\frac{b^2x}{a^2y}\]
于是，在点$(x_0,y_0)$处（这里$y_0\ne 0$）的切线方程为
\[y-y_0=-\frac{b^2x_0}{a^2y_0}(x-x_0)\]
即：$b^2x_0x+a^2y_0y=b^2x^2_0+a^2y^2_0$

$\because\quad $点$(x_0,y_0)$在椭圆上，

$\therefore\quad b^2x_0^2+a^2y^2_0=a^2b^2$，因此，所求切线方程是
\[b^2x_0 x+a^2y_0 y=a^2b^2\]
也就是
\[\frac{x_0x}{a^2}+\frac{y_0y}{b^2}=1\]

若$y=0$则$x=\pm a$, 显然在点$(-a,0)$和点$(a,0)$处的切线方程分别是$x=-a$, $x=a$, 它们是方程$\frac{x_0x}{a^2}+\frac{y_0y}{b^2}=1$的特殊形式．
\end{proof}
    
\begin{example}
    如果将抛物线$y^2=2px$绕它的对称轴旋转成旋转抛物面，那么与对称轴平行的光线射到曲面上，经曲面反射后通过焦点，试证明之．
\end{example}


\begin{proof}
从物理学知道，光线经镜面反射后，要服从光的反射定律：
\begin{enumerate}
    \item 在均匀介质中，光沿直线传播．
    \item 当光线碰到平面镜时，光线要反射．入射线与镜面法线夹角（记作$i$）等于反射线与镜面法线的夹角（记作$r$），见图2.11, 角$i$称为入射角，角$r$称为反射角．
\item 光从曲面镜的反射也遵守同样的法则：
  \[  \text{入射角}i=\text{反射角}r\]
    此时镜面的法线定义为与镜面在入射点的切线相垂直的直线．
\end{enumerate}

\begin{figure}[htp]
    \centering
    \begin{tikzpicture}[>=latex, scale=1]
        \fill[pattern=north east lines](-2,-.25) rectangle (2,0);
        \draw(-2,0)node[left]{镜面}--(2,0);
        \draw[->, very thick](140:3)node[left]{入射线}--(140:2);
        \draw[->, very thick](140:2)--(0,0)--(40:3)node[right]{反射线};
        \draw(0,0)--(0,3)node[above]{法线};
        \draw[<->, thick](90:1.8) arc (90:40:1.8);
        \draw[<->, thick](90:1.8) arc (90:140:1.8);
        \node at (60:2){$r$};
        \node at (120:2){$i$};
            \end{tikzpicture}
    \caption{}
\end{figure}

假设光线经位于$xy$平面上的抛物线$y^2=2px$形的镜面反射．当入射线与x轴平行且与抛物线交于$P_1(x_1,y_1)$点时，我们要求反射线的方
程，并证明它通过焦点$F\left(\frac{p}{2},0\right)$. 如果过$P_1(x_1,y_1)$点的切线与$x$轴成$\theta$角，那么平行于$x$轴的入射光线与法线的夹角是$90^{\circ}-\theta$. 由光的反射定律，反射光线与法线的夹角也是$90^{\circ}-\theta$. 于是，从图2.12上看出，反射光线与$x$轴的夹角是$2\theta$.

现在，在方程$y^2=2px$的两边对$x$求导数，得到
\[2yy'=2p\quad \Rightarrow\quad y'=\frac{2p}{2y}=\frac{p}{y}\]
所以，由导数的几何意义得出，在$P_1(x_1,y_1)$点的切线的斜率是
\[\tan\theta=y'_{x_1}=\frac{p}{y_1}\]
于是反射光线的斜率是
\begin{align*}
    \tan2\theta=\frac{2\tan\theta}{1-\tan^2\theta}&=\frac{2\cdot \frac{p}{y_1}}{1-\frac{p^2}{y^2_1}}\\
    &=\frac{2py_1}{y^2_1-p^2}=\frac{2py_1}{2px_1-p^2}=\frac{2y_1}{2x_1-p}
\end{align*}
从而，得出反射线方程
\[y-y_1=\frac{2y_1}{2x_1-p}(x-x_1) \]
令$y=0$, 得到反射光线与$x$轴交点的横坐标，
\[0-y_1=\frac{2y_1}{2x_1-p}(x-x_1)\quad \Rightarrow\quad x=\frac{p}{2}\]
即反射线通过焦点$F\left(\frac{p}{2},0\right)$. 由于$P_1(x_1,y_1)$是抛物线上
任意一点，所以只要入射光线与$x$轴平行，反射光线必通过焦
点．如果把光源放在焦点$F\left(\frac{p}{2},0\right)$上，那么从焦点发出的
光线经抛物线镜面反射后的反射光线是与对称轴（$x$轴）平行的光线．我们常利用这个原理获得平行光．
\end{proof}

\begin{figure}[htp]\centering
\begin{tikzpicture}[>=latex, scale=1.2]
\draw[->](-1.5,0)--(5,0)node[right]{$x$};
\draw[->](0,-3)--(0,3)node[right]{$y$};
\draw[domain=2.5:-2.5, samples=100, very thick]plot({0.5*\x*\x}, \x)node[right]{$y^2=2px$};
\draw[domain=-1:3.5, samples=10, thick]plot(\x, {0.707*(\x+1)});
\draw[thick](.5,0)node[below right]{$F\left(\frac{p}{2},0\right)$}--(1,1.414)--(4,1.414);
\draw[->, thick](4,1.414)--(3,1.414);
\tkzDefPoints{-1/0/A, 1/1.414/P, .5/0/F, 4/1.414/B, 3/2.828/C, 5/0/D}
\tkzMarkAngles[mark=none, size=.45](D,A,P D,F,P A,P,F B,P,C)
\tkzLabelAngle[pos=.6](D,A,P){$\theta$}
\tkzLabelAngle[pos=.6](D,F,P){$2\theta$}
\tkzLabelAngle[pos=.6](A,P,F){$\theta$}
\tkzLabelAngle[pos=.6](B,P,C){$\theta$}

\draw[domain=.5:1.5, samples=10]plot(\x, {1.414*(-\x+2)});
\tkzDefPoints{1.5/0.707/E}
\tkzMarkAngles[mark=none, size=.6, <->](E,P,B)
\tkzLabelAngle[pos=1.2](E,P,B){$90^{\circ}-\theta$}
\node at (P) [left]{$P_1(x_1,y_1)$};
\node at (0,0)[below left]{$O$};
\end{tikzpicture}
\caption{}
\end{figure}

\begin{ex}
\begin{enumerate}
    \item 求下列隐函数在指定点的导数$\frac{\dd y}{\dd x}$:
\begin{enumerate}
    \item $y=\cos x+\frac{1}{2}\sin y$, 点$\left(\frac{\pi}{2},0\right)$
    \item $ye^x+\ln y=1$, 点$\left(0,1\right)$
    \item $x^2+xy+y^2=1$, 点$\left(1,-1\right)$
    \item $x^2y^2+x^2+y^2=9$, 点$\left(1,2\right)$
\end{enumerate}
    \item 求下列隐函数的导数$\frac{\dd y}{\dd x}$:
\begin{multicols}{2}
\begin{enumerate}
    \item $\frac{x^2}{a^2}+\frac{y^2}{b^2}=1$
    \item $y^2=8x$
    \item $x^2=2y$
    \item $\sqrt{x}+\sqrt{y}=1$
    \item $x^3+y^3=a^3$
    \item $x^{\tfrac{2}{3}}+y^{\tfrac{2}{3}}=a^{\tfrac{2}{3}}$
    \item $y-\cos(x+y)=0$
    \item $y=x+\arctan y$
    \item $y=1-\ln(x+y)+e^y$
    \item $\arctan \frac{y}{x}=\ln\sqrt{x^2+y^2}$
\end{enumerate}
\end{multicols}

\item 求下列方程在指定点的导数$\frac{\dd y}{\dd x}$:
\begin{enumerate}
    \item $x-\sqrt{xy}+y=2$，点$(x_1,y_1)$
    \item $(ax)^{\tfrac{2}{3}}-(by)^{\tfrac{2}{3}}=(a^2+b^2)^{\tfrac{2}{3}}$，点$(x_0,y_0)$
    \item $(x+y)^{\tfrac{2}{3}}+(x-y)^{\tfrac{2}{3}}=2a^{\tfrac{2}{3}}$在$x=0$的点；在$x=y$的点；在$x=-y$的点
    \item $x^2+xy+y^2=c^2$，在斜率等于$-1$的点
\end{enumerate}
\item 若$x\sqrt{y}+y\sqrt{x}=c$，求$\frac{\dd y}{\dd x}$和$\frac{\dd x}{\dd y}$.
\item 求椭圆$x^2+4y^2=9$的切线方程，使它和直线$2x+3y=0$
平行．
\item 曲线$y=(x-2)(x-3)(x-4)$和$x$轴交于$P(2, 0)$, $Q(3,0)$和$R(4, 0)$点．
\begin{enumerate}
    \item 求证在$P$点和$R$点处的切线互相平行；
    \item 求曲线在$Q$点处的法线与$y$轴的交点．
\end{enumerate}
\item 求曲线 $2x^3-xy+8=0$上的$P$点，使在该点的切线通过
原点，过$P$点的切线与曲线交于另一点$Q$的坐标．
\item \begin{enumerate}
    \item 求证直线$\ell x +my+n=0$与椭圆$\frac{x^2}{a^2}+\frac{y^2}{b^2}=1$, 双曲线$\frac{x^2}{a^2}-\frac{y^2}{b^2}=1$, 抛物线$y^2=2px$相切的条件分别是$a^2\ell^2+b^2m^2=n^2$, $a^2\ell^2-b^2m^2=n^2$, $pm^2=2n\ell$.
    \item 证明对于任意$k$, 直线$y=kx+\sqrt{a^2k^2+b^2}$是椭圆曲线$\frac{x^2}{a^2}+\frac{y^2}{b^2}=1$的切线．    
\end{enumerate}
\item \begin{enumerate}
    \item 求过$P(5, 3)$点与椭圆$\frac{x^2}{16}+{y^2}=1$相切的直线
    方程.
    \item 求过(a)中的两切点的直线方程．
\end{enumerate}

\item 椭圆$\frac{x^2}{16}+{y^2}=1$的切线与$x$轴成$60^{\circ}$角，求该切线的方程．
\item 通过双曲线$\frac{x^2}{a^2}-\frac{y^2}{b^2}=1$的焦点而与横轴垂直的弦的
两个端点分别记作$A$和$B$, 求证过$A$、$B$两点的切线交于
横轴上的点$Q$, 且$Q$点离中心的距离等于$\frac{a}{e}$.
\item 求证过双曲线上一点$P$的切线与过$P$点的焦点半径成等角．
\end{enumerate}
\end{ex}

\subsection{高阶导数}
对于给定的可微函数$f(x)$, 它的导数$f'(x)$仍是$x$的函数．如果$f'(x)$是可微的，它的导函数$(f'(x))'=f''(x)$或$f^{(2)}(x)$称为$f(x)$的\textbf{二阶导函数}，……如此类推，$f(x)$
的$n$阶导函数$f^{(n)}(x)$, 可以归纳地定义为：
\begin{align*}
  f ^{(1)} (x)&=f'(x)\\
  f ^{(n)} (x)&=\left(f^{(n-1)}(x)\right)',\qquad n=2,3,\ldots  
\end{align*}
因此，求高阶导函数，实际上就是连续地求一阶导函数．

高阶导函数有时也记作$\frac{\dd^n y}{\dd x^n}$或$\frac{\dd^n }{\dd x^n}f(x)$. 例如
\begin{align*}
    f''(x)&=\frac{\dd}{\dd x}f'(x)=\frac{\dd}{\dd x}\left(\frac{\dd y}{\dd x}\right)=\frac{\dd^2 y}{\dd x^2}\\
    f^{(3)}(x)&=\frac{\dd}{\dd x}\left(\frac{\dd^2 y}{\dd x^2}\right)=\frac{\dd^3 y}{\dd x^3}
\end{align*}

二阶导数有它的实际意义，例如动点的速度$v$等于它所
经过的路程$S$对于时间$t$的变率：$v=\frac{\dd S}{\dd t}$；而加速度$a$是速度$v$
对于时间$t$的变率：$a=\frac{\dd v}{\dd t}$，
这就是说加速度是路程$S$对于时间$t$的二阶导数：$a=\frac{\dd^2 S}{\dd t^2}$．
在以后的学习中，我们还会知道二阶导数有它的重要的几何意义．我们已经知道一次函数$f(x)=ax+b\; (a\ne 0)$的图象是一条直线，而它的二阶导数$f''(x)=(f'(x))'=(a)'=0$. 因此，若一函数的二阶导数$f''(x)\ne 0$, 那么它的图象就不是直线．

\begin{example}
设$y=\frac{x}{1-x}$，求$\frac{\dd^2 y}{\dd x^2}$, $\frac{\dd^3 y}{\dd x^3}$.
\end{example}

\begin{solution}
\begin{align*}
    y&=\frac{x}{1-x}=-1+\frac{1}{1-x}\\
\frac{\dd y}{\dd x}&=-\frac{(1-x)'}{(1-x)^2}=\frac{1}{(1-x)^2}
\end{align*}
\begin{align*}
    \frac{\dd^2 y}{\dd x^2}&=\frac{-[(1-x)^2]'}{(1-x)^4}=\frac{-2(1-x)(-1)}{(1-x)^4}=\frac{2}{(1-x)^3}\\
    \frac{\dd^3 y}{\dd x^3}&=\frac{-2[(1-x)^3]'}{(1-x)^6}=\frac{-2\cdot 3(1-x)^2(-1)}{(1-x)^6}=\frac{6}{(1-x)^4}  
\end{align*}
\end{solution}


\begin{example}
若$f(x)=\sqrt{a^2+x^2}$，求$f''(x)$
\end{example}


\begin{solution}
\textbf{解法1: }
\begin{align*}
    f'(x)&=\frac{1}{2}(a^2+x^2)^{-\tfrac{1}{2}}(a^2+x^2)'\\
    &=\frac{1}{2}(a^2+x^2)^{-\tfrac{1}{2}}(2x)=\frac{x}{\sqrt{a^2+x^2}}\\
\end{align*}
\begin{align*}
f''(x)&=\frac{\sqrt{a^2+x^2}\cdot x'-x\left(\sqrt{a^2+x^2}\right)'}{a^2+x^2}\\
&=\frac{\sqrt{a^2+x^2}-\frac{x^2}{\sqrt{a^2+x^2}}}{{a^2+x^2}}\\
&=\frac{a^2+x^2-x^2}{(a^2+x^2)^{\tfrac{3}{2}}}=\frac{a^2}{(a^2+x^2)^{\tfrac{3}{2}}}
\end{align*}

\textbf{解法2: } 将原来的函数变形为隐函数，用隐函数微分法求二阶导数．

设$y=\sqrt{a^2+x^2}$, 则$y^2=a^2+x^2$, 对$x$微分得：
\[y\cdot y'_x=x\]    
再对$x$微分得：
\[y\cdot y_x''+(y_x')^2=1\]
将$y'=\frac{x}{y}$代入，得
\[y\cdot y''_x+\frac{x^2}{y^2}=1\]
因此：
\[y''_x=\frac{1-\frac{x^2}{y^2}}{y}=\frac{1-\frac{x^2}{x^2+a^2}}{\sqrt{a^2+x^2}}=\frac{a^2}{(a^2+x^2)^{\tfrac{3}{2}}}\]
\end{solution}

\begin{example}
    若$y=\frac{x^2}{\sqrt{x^2+1}}$，求证：
\[(x^2+1)\frac{\dd^2 y}{\dd x^2}+3x\frac{\dd y}{\dd x}-y=2\sqrt{x^2+1}\]
\end{example}

\begin{proof}
将原式去分母
\[y\sqrt{x^2+1}=x^2\]
对$x$微分得
\[\sqrt{x^2+1}\cdot\frac{\dd y}{\dd x}+y\cdot \frac{x}{\sqrt{x^2+1}}=2x\]
两边乘以$\sqrt{x^2+1}$，得到
\[(x^2+1)\cdot \frac{\dd y}{\dd x}+y\cdot x=2x\sqrt{x^2+1}\]
再对$x$微分得：
\[(x^2+1)\frac{\dd^2 y}{\dd x^2}+(2x)\cdot \frac{\dd y}{\dd x}+x\frac{\dd y}{\dd x}+y=2\sqrt{x^2+1}+\frac{2x^2}{\sqrt{x^2+1}}\]
即：
\[(x^2+1)\frac{\dd^2 y}{\dd x^2}+3x\frac{\dd y}{\dd x}+y=2\sqrt{x^2+1}+2y\]
因此：
\[(x^2+1)\frac{\dd^2 y}{\dd x^2}+3x\frac{\dd y}{\dd x}-y=2\sqrt{x^2+1}\]
\end{proof}

有时需要求一些函数的$n$阶导数，但是在一般情形下，函数的$n$阶导数是如此的复杂，以致求函数的$n$阶导数的通式是办不到的，然而有一些函数，我们不难写出它的$n$阶导数的通式．

\begin{example}
  设$y=(ax+b)^m$, 求$y^{(n)}\; (m>n)$  
\end{example}

\begin{solution}
\begin{align*}
y'&=ma(ax+b)^{m-1}\\
y''&=m(m-1)a^2(ax+b)^{m-2}\\
y^{(3)}&=m(m-1)(m-2)a^3(ax+b)^{m-3}\\
\cdots\quad &\cdots\quad \cdots
\end{align*}
由此推知：
\[y^{(n)}=m(m-1)\cdots(m-n+1)a^n(ax+b)^{m-n}=\frac{m!}{(m-n)!}a^n(ax+b)^{m-n}\]
这个结果的正确性可用数学归纳法证明．
\end{solution}


\begin{example}
设$y=\frac{1}{x^2-4x+3}$，求$y^{(n)}$.
\end{example}


\begin{solution}
\[y=\frac{1}{x^2-4x+3}=\frac{1}{2}\left(\frac{1}{x-3}-\frac{1}{x-1}\right)\]
\begin{align*}
    y'&=\frac{1}{2}\left[\frac{-1}{(x-3)^2}-\frac{-1}{(x-1)^2}\right]=\frac{(-1)}{2}\left[\frac{1}{(x-3)^2}-\frac{1}{(x-1)^2}\right]\\
    y''&=\frac{(-1)}{2}\left[\frac{-2(x-3)}{(x-3)^4}-\frac{-2(x-1)}{(x-1)^4}\right]=\frac{(-1)^2 2!}{2}\left[\frac{1}{(x-3)^3}-\frac{1}{(x-1)^3}\right]\\
    \cdots\quad &\cdots\quad \cdots
\end{align*}
假设一般地有：
\[y^{(n)}=\frac{(-1)^n n!}{2}\left[\frac{1}{(x-3)^{n+1}}-\frac{1}{(x-1)^{n+1}}\right]\]
那么：
\begin{align*}
    y^{(n+1)}&=\frac{(-1)^n n!}{2}\left[\frac{-(n+1)(x-3)^n}{(x-3)^{2(n+1)}}-\frac{-(n+1)(x-1)^n}{(x-1)^{2(n+1)}}\right]\\
    &=\frac{(-1)^{n+1} (n+1)!}{2}\left[\frac{1}{(x-3)^{n+2}}-\frac{1}{(x-1)^{n+2}}\right]
\end{align*}
即仍有此规律．因此，对于任何自然数$n$，都有
\[y^{(n)}=\frac{(-1)^n n!}{2}\left[\frac{1}{(x-3)^{n+1}}-\frac{1}{(x-1)^{n+1}}\right]\]
成立．
\end{solution}

\begin{example}
    设$y=\sin(ax+b)$，求$y^{(n)}$.
\end{example}


\begin{solution}
    \[y'=a\cos(ax+b)=a\sin\left(ax+b+\frac{\pi}{2}\right)\]
因此，求$\sin(ax+b)$的导数是将函数乘以$a$，同时将幅角加上$\frac{\pi}{2}$.

不难推知：
\[y^{(n)}=a^n\sin\left(ax+b+\frac{n\pi}{2}\right)\]
请读者用数学归纳法证明这个猜想正确．
\end{solution}

\begin{example}
设$f(x)=a_0+a_1x+a_2x^2+\cdots+a_nx^n$，$h$为任取非零实数，求证：
\[f(a+h)=f(a)+f'(a)h+\frac{f''(a)}{2!}h^2+\cdots+\frac{f^{(n)}(a)}{n!}h^n\]
\end{example}


\begin{proof}
已给$f(x)=a_0+a_1x+\cdots+a_nx^n$，令$x=a+h$，并按$h$的升幂展开每一项，得到一个形式如
\begin{equation}
    f(a+h)=C_0+C_1 h+C_2 h^2+\cdots +C_n h^n
\end{equation}
的表达式．

我们将说明，系数$C_0,C_1,\ldots,C_n$由$f(x)$在点$a$的函数值$f(a)$以及直到$n$阶为止的各阶导数在$x=a$的值表示，即
\[C_0=f (a) ,\quad  C_i =\frac{1}{i!}f^{(i)} (a) ,\qquad i=1, 2,\ldots ,n\]

为了证明这个事实，(2.10)中把$h$看作自变量，首先令$h\to 0$, 立刻得到
\[f (a) =C_0\]

然后，应用复合函数求导法则，对方程(2.10), 关于$h$逐次求导，依次得到：
\begin{align*}
f'(a+h)&=C_1+2C_2 h+3C_3h^2+\cdots +nC_nh^{n-1}\\
f''(a+h)&=2C_2+3\x 2C_3h+\cdots+n(n-1)C_nh^{n-2}\\
f^{(3)}(a+h)&=3!C_3+\frac{4!}{1!}C_4h+\frac{5!}{2!}C_5h^2+\cdots+\frac{n!}{(n-3)!}C_nh^{n-3}\\
\cdots\quad &\cdots\quad \cdots\\
f^{(i)}(a+h)&=i! C_i+\frac{(i+1)!}{1!}C_{i+1}h+\cdots+\frac{n!}{(n-i)!}C_nh^{n-i}\\
\cdots\quad &\cdots\quad \cdots\\
f^{(n-1)}(a+h)&=(n-1)!C_{n-1}+n!C_nh\\
f^{(n)}(a+h)&=n!C_n
\end{align*}
把$h=0$代入上面一个方程中，便得到
\[C_1=f^{(1)}(a),C_2=\frac{f^{(2)}(a)}{2!},C_3=\frac{f^{(3)}(a)}{3!},\ldots,C_i=\frac{f^{(i)}(a)}{i!},\ldots,C_n=\frac{f^{(n)}(a)}{n!}\]

把上面$C_0=f(a),\; C_i=\frac{f^{(i)}(a)}{i!},\quad i=1,2,\ldots,n$代入(2.10)中，最后得到：
\begin{equation}
    f(a+h)=f(a)+f'(a)h+\frac{f^{(2)}(a)}{2!}h^2+\cdots+\frac{f^{(n)}(a)}{n!}h^n
\end{equation}
这个公式叫做多项式的泰勒公式．

在(2.11)中，令$a+h=x$, 则(2.11)可以改写成
\[    f(x)=f(a)+f'(a)(x-a)+\frac{f^{(2)}(a)}{2!}(x-a)^2+\cdots+\frac{f^{(n)}(a)}{n!}(x-a)^n\]
\end{proof}

\begin{ex}
\begin{enumerate}
    \item 求下列函数的$n$阶导数：
\begin{multicols}{2}
\begin{enumerate}
    \item $f(x)=\frac{1-x}{1+x}$
    \item $f(x)=\sin^2 x$
\end{enumerate}
\end{multicols}
    \item 若$y=\tan(m\cdot \arctan x)$，验证
\[(1+x^2)\frac{\dd^2 y}{\dd x^2}=2(my-x)\frac{\dd y}{\dd x}\]
\item 若$y=\frac{A\cos mx+B\sin mx}{x^n}$，其中$A,B,m,n$是常数，验证
\[\frac{\dd^2 y}{\dd x^2}+2\frac{n}{x}\frac{\dd y}{\dd x}+\left[\frac{n(n-1)}{x^2}+m^2\right]y=0\]

\item 若$h(x)=f(x)g(x)$，其中$f(x)$, $g(x)$和$h(x)$都是$n$阶可微，求证：
\begin{align*}
    h^{(n)}(x)&=f(x)\cdot g^{(n)}(x)+C^1_n f^{(1)}(x)g^{(n-1)}(x)+\cdots \\
    &\qquad +C^h_n f^{(h)}(x)g^{(n-h)}(x)+\cdots\\
    &\qquad  +C^{n-1}_n f^{(n-1)}(x)g^{(1)}(x)+C^n_n f^{(n)}(x)g(x)
\end{align*}
\item 求下列函数的$n$阶导数：
\begin{multicols}{2}
\begin{enumerate}
\item $f(x)=x\ln x$
\item $f(x)=xe^x$
\item $f(x)=e^{ax}\sin\beta x$
\end{enumerate}
\end{multicols}
\item 证明 $x^3-7x^2+15x-9=0$有一个二重根．
\item 检验方程 $x^3-6x^2+11x-6=0$有无重根？
\item 方程$x^4-6x^2-8x-3=0$有重根，解此方程．
\item 应用泰勒公式将多项式$(x+a)^n$按$x$的升幂展开．
\item 若$f(x)=x^3-3kx+2k+8$在实数范围内能分解成一个二重因式，试求$k$值和它的因式乘积．
\end{enumerate}   
\end{ex}

\subsection{参数方程表示的函数的导数}
用方程$y=f(x)$来表示曲线在几何上受很大的限制：这样表示的曲线与平行于$y$轴的任意直线相交不能多于一点．通常，把曲线分成可以表为$y=f(x)$的若干部分来克服这个限制，例如，以原点为中心，$r$为半径的圆，可以由定义在闭区间$[-r,r]$上的两个函数$y=\sqrt{r^2-x^2}$和$y=-\sqrt{r^2-x^2}$给出．

曲线最直接和最灵活的描述法是参数表示．我们不把直角坐标$x$或$y$中的一个看成另一个的函数，而把两个坐标$x$和$y$都看成为第三个自变量$t$的函数，$t$叫做\textbf{参数}或\textbf{参变量}．当$t$在$t$轴上或$t$轴上的一个区间$[a,b]$变化时，由$t$轴上的原象点，通过参数方程$x=\varphi(t)$, $y=\psi(t)$映射到曲线$C$上的一点$(x,y)=(\varphi(t),\psi(t))$, 使得对应于$t$的某个区间$[a,b]$或$(-\infty, \infty)$, 一对函数值$x=\varphi(t)$, $y=\psi(t)$定义的点集
$$\Big\{\big(\varphi(t), \psi(t)\big)\Big|t\in [a,b]\Big\}$$
就是曲线$C$上的全部点，而且没有其它点．

在解析几何中，我们已经知道：

圆的参数方程是
\[x=a\cos t,\qquad y=a\sin t\]
这里$t$表示圆心角．

椭圆的参数方程是
\[x=a\cos t,\qquad y=b\sin t\quad (a>b>0)\]
这里$t$是偏心角．（图2.13）

\begin{figure}[htp]
    \centering
\begin{tikzpicture}[>=latex]
    \draw[<->|](0,-.25)--node[fill=white]{$a$}(2,-.25);
    \draw[<->|](.25,0)--node[fill=white]{$b$}(.25,1.2);
\draw[->](-2.8,0)--(2.8,0)node[right]{$x$};
\draw[->](0,-2.8)--(0,2.8)node[right]{$y$};
\draw[thick](0,0) circle (2);
\draw[thick](0,0) circle (1.2);
\draw[very thick](0,0) ellipse [x radius=2, y radius=1.2];
\node at (0,0)[below left]{$O$};
\draw(0,0)--(30:3);
\draw(30:1.2)node[above]{$B$}--(1.04,0)node[above right]{$N$};
\draw(30:2)node[above]{$A$}--(1.73,0)node[above right]{$M$};
\draw(.5,0) arc (0:30:.5)node[right]{$t$};
\draw(30:1.2)--(1.9,.6)node[right]{$P$};

\end{tikzpicture}
    \caption{}
\end{figure}

曲线$C$可以有许多参数表示法，就好象同一条路可以有许多不同走法，例如抛物线$y^2=2px$可以用在其上的点$(x,y)$的切线的斜率作参数，由于
$2yy'=2p$, 所以抛物线$y^2=2px$的参数方程是
\[x=\frac{p}{2(y')^2},\qquad y=\frac{p}{y'}\]
其中$y'\in (-\infty,0)\cup (0, \infty)$.

如果用$\frac{1}{y'}=\cot\theta$作参数，那么它的参数方程是
\[x=\frac{p}{2}\cot^2 \theta,\qquad y=p\cot\theta\]

我们也可以用$\frac{y}{2p}=\frac{x}{y}=\lambda,\quad \lambda\in (-\infty, \infty)$作参数，于
是又得一参数方程
\[x=2p\lambda^2,\qquad y=2p\lambda\]


\begin{example}
讨论由参数方程
\[x=\sin t,\qquad y=\cos^2 t,\quad t\in[0,2\pi]\]
所定义的曲线$C$.
\end{example}

\begin{solution}
因为$\sin^2 t+\cos^2 t=1 $，所以曲线$C$在抛物线
\begin{equation}
    x^2+y=1
\end{equation}
上．反过来，抛物线(2.12)并不与曲线$C$一致，我们要说明限制在长方形区域$R=\{-1\le x\le 1,\; 0\le y\le 1\}$内的抛物线(2.12)上的部分是曲线$C$（图2.14）．

\begin{figure}[htp]
    \centering
\begin{tikzpicture}[>=latex, scale=1.5]
    \draw[->](-2,0)--(2,0)node[right]{$x$};
    \draw[->](0,-1.5)--(0,2)node[right]{$y$};
\node at (0,0)[below left]{$O$};
\draw(-1,0) rectangle (1,1);
\draw[domain=-1.5:1.5, samples=100, very thick, dashed]plot(\x, {1-\x*\x});
\draw[domain=-1:1, samples=100, very thick]plot(\x, {1-\x*\x});
\node at (-1,0)[below left]{$(-1,0)$};
\node at (1,0)[below right]{$(1,0)$};
\node at (0,1)[above right]{$(0,1)$};
\end{tikzpicture}
    \caption{}
\end{figure}

事实上，$x=\sin t$在闭区间$[0,2\pi]$上处处连续，且在闭区间$\left[\frac{\pi}{2},\frac{3\pi}{2}\right]$上由1单调递减到$-1$，因此根据反函数存在定理，对于每一个$x\in[-1,1]$，存在唯一的一个$t$，使得$\sin t=x$成立，从而得到
\[y=1-x^2=1-\sin^2 t=\cos^2 t\in [0,1]\]
这就说明了，点$(x,1-x^2)\; (x\in [-1, 1])$在曲线$C$上，而且只有这些点在曲线$C$上，换言之，在(2.12)上的除去这些点的其它任何一点不在曲线$C$上．
\end{solution}

在参数方程$x=\varphi(t)$, $y=\psi(t)$中，由于$t$轴是有序的，我们可以用明显的方式把曲线$C$的点规定一个次序或“指向”，即如果$t_1<t_2$，那么由$t_1$映射的点在由$t_2$映射的点的前面．例如例2.38的曲线上的动点$P$随着$t$增加的方向这样运动：$P$点在$t=0$时，由点$(0,1)$开始移动到在$t=\frac{\pi}{2}$时的$(1,0)$点，然后沿原路线返回移动到$t=\pi$时的$(0,1)$点，再继续移动到在$t=\frac{3\pi}{2}$时的$(-1, 0)$点，最后又返回到在$t=2\pi$时的$(0, 1)$点．

如果函数$\varphi(t)$, $\psi(t)$的定义域是$t$轴上的闭区间$[a,b]$, 而且在这区间上的$t$的不同值对应曲线$C$上不同的点，
那么我们称$C$为\textbf{简单弧}，例2.38中的抛物线，由上面知道不是简单弧，而抛物线$x=t$, $y=t^2,\; t\in [0, 1]$是简单弧的一个例子．

如果曲线$C$上的始点$P(\varphi(a),\psi(a))$与终点$P(\varphi(b),\psi(b))$相同，并且除$t=a$和$t=b$外，在$[a,b]$上没有另外一对不同值对应曲线上同一点，那么曲线$C$叫做\textbf{简单闭曲线}．$x=a\cos t$, $y=a\sin t,\; t\in [-\pi,\pi]$就是简单的闭曲线．

有时我们使用$y=f(x)$表示曲线$C$或$C$的一部分是方便的，如果$C$用参数方程$x=\varphi(t)$, $y=\psi(t)$给出，而对于曲线的一部分$t_2\le t\le t_1$, 函数$\varphi(t),\psi(t)$中的一个，譬如说$x=\varphi(t)$是单调的和连续的，那么这样的表示法总是可能的，办法是先由$x=\varphi(t)$解出它的反函数$t=\varphi^{-1}(x)$, 再代入$y=\psi(t)$中，便得到
\[y=\psi\left(\varphi^{-1}(x)\right)=f(x)\]
即$y=f(x)$是由$y=\psi(t)$和$t=\varphi^{-1}(x)$复合而成的复合函数．

要求曲线$C$的斜率，我们不必重新建立$y$对$x$的函数关系$y=f(x)$, 再去求导，只须由方程组
\begin{equation}
    \begin{cases}
        x=\varphi(t)\\ y=\psi(t)
    \end{cases}
\end{equation}
当$\frac{\dd\varphi}{\dd t}\ne 0$时，按照复合函数求导法则可以直接得到$f'(x)$.

事实上，当$\varphi'(t)\ne 0$时，$t=\varphi^{-1}(x)$的导数存在，并且是$\frac{\dd t}{\dd 
x}=\frac{1}{\varphi'(t)}$.

按照复合函数求导法则，得到
\[f'(x)=\frac{\dd y}{\dd t}\cdot \frac{\dd t}{\dd x}=\frac{\psi'(t)}{\varphi'(t)}\]
于是，当$\frac{\dd\varphi}{\dd t}\ne 0$时，就得到曲线$C$的斜率：
\begin{equation}
    \frac{\dd y}{\dd x}=\frac{\psi'(t)}{\varphi'(t)}
\end{equation}
当$\frac{\dd\psi}{\dd t}\ne 0$时，就得到
\begin{equation}
    \frac{\dd x}{\dd y}=\frac{\varphi'(t)}{\psi'(t)}
\end{equation}
因此，只要$[\varphi'(t)]^2+[\psi'(t)]^2\ne 0$, 切线总是存在的：如果$\psi'(t)=0$, 切线是水平的；如果$\varphi'(t)=0$, 切线是垂直的．

当曲线由参数方程组(2.13)给出，那么曲线上点$(x_0,y_0)$处的切线方程为
\begin{equation}
    (y-y_0)\frac{\dd x}{\dd t}-(x-x_0)\frac{\dd y}{\dd t}=0
\end{equation}

法线的斜率是$-\frac{1}{y'_x}$，因此法线方程为
\begin{equation}
    (x-x_0)\frac{\dd x}{\dd t}+(y-y_0)\frac{\dd y}{\dd t}=0
\end{equation}

\begin{example}
求椭圆$\begin{cases}
    x=a\cos t\\ y=b\sin t
\end{cases}$在点$t$的切线方程和法线方程．
\end{example}



\begin{solution}
椭圆在点t的切线方程，按照(2.16)得到
\[(y-b\sin t) (-a\sin t )-(x-a\cos t)(b\cos t)=0\]
    化简得
 \[   y ( - a\sin t) - x(b\cos t)+ab=0\]
    两边除以$-ab$, 得到在点$t$的切线方程
\[\frac{x}{a}\cos t+\frac{y}{b}\sin t=1\]

在点$t$的法线方程，按照(2.17)得到
\[(x-a\cos t) (-a\sin t)+(y-b\sin t)(b\cos t)=0\]
    化简得
\[   -ax\sin t+by\cos t+ (a^2-b^2) \sin t\cdot \cos t=0\]
    或
    \[xa\sin t-yb\cos t-\frac{c^2}{2}\sin 2t=0\]
\end{solution}    

\begin{example}
    求曲线$\begin{cases}
        x=t^2\\y=t^3
    \end{cases}$
    在点$t$的切线方程，并求此切线与曲线的另一交点的坐标.
\end{example}

\begin{solution}
在点$t$的切线方程是
\[(y-t^3)(2t)-(x-t^2)(3t^2)=0\]

若$t\ne 0$, 则得切线方程
  \[  2y-3tx+t^3=0\]
    设此切线与曲线交于$(T^2,T^3)$点，于是数$T$满足等式
   \[ 2T^3-3tT^2+t^3=0\]
    这是$T$的三次方程，它表示切线与曲线相交于三个点，由于切点的$t$值是它的二重根，根据根和系数的关系，得
\[2t+T=-\frac{-3t}{2}\quad \Rightarrow\quad T=-\frac{1}{2}t\]
故切线与曲线的另一交点的坐标是$\left(\frac{1}{4}t^2,-\frac{1}{8}t^3\right)$

如果再假定$\varphi''(t)$，$\psi''(t)$存在，则$y$对$x$的二阶导数可按照复合函数求导法及商的导数法来求：
\begin{align*}
    y''=\frac{\dd y'}{\dd x}&=\frac{\dd y'}{\dd t}\cdot \frac{\dd t}{\dd x}\\
    &=\frac{\dd }{\dd t}\left(\frac{\psi'(t)}{\varphi'(t)}\right)\cdot \frac{1}{\varphi'(t)}\\
    &=\frac{\varphi'(t)\psi''(t)-\psi'(t)\varphi''(t)}{\big\{\varphi'(t)\big\}^3}
\end{align*}
\end{solution}

\begin{ex}
\begin{enumerate}
    \item 设$x=\frac{1}{2}p\lambda^2$, $y=p\lambda$是抛物线以$\lambda$为参数的方程，求在
    点$\lambda$的切线方程和法线方程．
    \item 过点$(x_0,y_0)$作抛物线$y^2=2px$的两条切线，设切点的参数为$\lambda_1$和$\lambda_2$, 求证
    \[\begin{cases}
        p(\lambda_1+\lambda_2)=2y_0\\
        p\lambda_1\lambda_2=2x_0
    \end{cases}\]
    \item 过抛物线$y^2=2px$的定长为$c$的弦的两端作切线，求切线交点的轨迹方程．
    
    \item 求曲线
    $\begin{cases}
        x=t (1-t^2) \\ y=1+t^2
    \end{cases}$在点$t$的切线方程．

\item 曲线$y^2=x^3$上点$P(t^2,t^3)$的切线交$x$轴于$T$点，过$T$
点与$PT$垂直的直线交$y$轴于$N$点，求证$\frac{ON^2}{OT}$是一个与
$t$无关的常数，这里$O$是原点．
\item 如果曲线$ay^2=x^3$上$P(at^2,at^3)$点的切线交曲线于另一点$Q$, 交$y$轴于$R$点，
\begin{enumerate}
    \item 求$Q$点坐标；
    \item 如果
由点$P$引$x$轴的垂线的垂足为$N$点，求证：$OQ$和$RN$的斜率的绝对值相等．
\end{enumerate}
\item \begin{enumerate}
    \item 求以$a$为参数的抛物线系
\[y+\frac{35}{16}a=\frac{a^3}{3}\left(x+\frac{3}{4a}\right)^2\]
的顶点的轨迹方程．
\item 证明对应于参数$a$的抛物线与双曲线$xy=-\frac{7}{6}$ 
相切于点$\left(\frac{1}{a},-\frac{7a}{6}\right)$，又与双曲线$xy=\frac{10}{3}$, 相切
于点$\left(-\frac{2}{a},-\frac{5}{3}a\right)$.
\end{enumerate} 

\item 试证曲线
$x=a(\cos t+t\sin t)$, $y=a(\sin t-t\cos t)$的每
条法线都与圆$x^2+y^2=a^2$相切

\item 若$x=\sec\theta$ 和$y=2\tan\theta$,
\begin{enumerate}
    \item 求$\frac{\dd y}{\dd x}$和$\frac{\dd^2 y}{\dd x^2}$以$\theta$表示之；
\item 求证：$y^3\frac{\dd^2 y}{\dd x^2}+16=0$
\end{enumerate}

\item 求证过椭圆$x=a\cos\varphi$, $y=b\sin\varphi$上两点$\varphi_1$、$\varphi_2$的切线
的交点的坐标为
$\left(a\frac{\cos\frac{\varphi_1+\varphi_2}{2}}{\cos\frac{\varphi_1-\varphi_2}{2}}, b\frac{\sin\frac{\varphi_1+\varphi_2}{2}}{\cos\frac{\varphi_1-\varphi_2}{2}}\right)$
\end{enumerate}
\end{ex}

\subsection{微分与变率（微商）}

\subsubsection{微分的概念}

当函数$y=f(x)$的自变量$x$在点$x_0$处得到一个微小的改变量$\Delta x=x-x_0$时，函数$y$的值得到相应的改变量，即
\begin{equation}
    \Delta y=f (x_0+\Delta x) -f (x_0) 
\end{equation} 
函数的改变量$\Delta y$最能表现函数的特征，无论函数变化的大小，或者是函数在某一点的递增或递减的变化趋势，都是由它来刻画的．

从数的观点来看，一个变量就是一个可以取很多不同的值的符号，所以我们也可以把$\Delta x$看成一个新变量，这样$\Delta y$就是$x_0$和$\Delta x$这两个变量的函数．如

\begin{example}
    设$y=x^3$，那么
\[\Delta y=(x_0+\Delta x)^3-x^3_0 =3x_0^2\Delta x+3x_0(\Delta x)^2+(\Delta x)^3\]
\end{example}

\begin{example}
    设$f(x)=x^2-\frac{1}{x}$，那么
\begin{align*}
    \Delta y&=\left[(x+\Delta x)^2-\frac{1}{x+\Delta x}\right]-\left[x^2-\frac{1}{x}\right]\\
    &=2x\Delta x+(\Delta x)^2+\frac{\Delta x}{x(x+\Delta x)}
\end{align*} 
\end{example}

从上面的例子看出，通常$\Delta y=f(x_0+\Delta x)-f(x_0)$是一个很复杂的关系，但是如果函数$f(x)$在点$x_0$存在变率$f'(x_0)$, 我们可以把$\Delta y$分解成一个线性部分和一个比$\Delta x$变小得更快的微量部分．事实上
\begin{equation}
 f'(x_0)=\lim_{\Delta x\to 0}\frac{f(x_0+\Delta x)-f(x_0)}{\Delta x}
\end{equation}
我们用$\varepsilon(x_0,\Delta x)$表示$\frac{\Delta 
y}{\Delta x}$与$f'(x_0)$的差，那么，我们可以写出
\begin{equation}
\frac{f(x_0+\Delta x)-f(x_0)}{\Delta x}=f'(x_0)+\varepsilon(x_0,\Delta x)
\end{equation}
这里因为(2.19)，而有$\Lim_{\Delta x\to 0}\Delta x \cdot \varepsilon(x_0,\Delta x)=0$.

将(2.20)的两边乘以$\Delta x$，我们得到
\begin{equation}
    f(x_0+\Delta x)-f(x_0)=f'(x_0)\Delta x+\Delta x\cdot \varepsilon(x_0,\Delta x)
\end{equation}
这个等式右边的第一项是与$\Delta x$成比例的量，其比例常数是$f(x)$在点$x_0$的变率$f'(x_0)$, 而第二项$\varepsilon(x_0,\Delta x)\Delta x$和$\Delta x$比较起来，当$|\Delta x|$很小时，$|\varepsilon(x_0,\Delta x)\cdot \Delta x|$就更小，从极限的观点来看，就是
\[\lim_{\Delta x\to 0} \frac{\varepsilon(x_0,\Delta x)\cdot \Delta x}{\Delta x} =\lim_{\Delta x\to 0}\varepsilon(x_0,\Delta x)   =0\]
这也就是说，$\varepsilon(x_0,\Delta x)\cdot \Delta x$比$\Delta x$变小得更快．因此，若$f'(x_0)\ne 0$, (2.21)中的线性部分$f'(x_0)\Delta x$, 当$|\Delta x|$很小时就是函数改变量$f(x_0+\Delta x)-f(x_0)$的很好的近似值，它的误差等于$\varepsilon(x_0,\Delta x)\cdot \Delta x$, 由于这个与$\Delta x$成比例的部分是$\Delta y$的\textbf{主要部分}，我们把$f'(x_0)\Delta x$叫做函数$f(x)$在$x=x_0$处\textbf{微分}．

\begin{blk}
    {定义} 设函数$y=f(x)$是可微的，我们称$f'(x)\Delta x$为函数$y=f(x)$在点$x$处与自变量的改变量$\Delta x$相对应的微分，记作$\dd y$, 于是
\begin{equation}
   \dd y=f' (x) \Delta x 
\end{equation}
\end{blk}

因为函数$f(x)=x$的导数等于1, 所以对于任何$x$来说，微分$\dd f(x)$便等于$\dd x=(x')\Delta x=\Delta x$. 因此，习惯上，常以$\dd x$代替$\Delta x$, 于是(2.22)就可以写成
\begin{equation}
    \dd y=f'(x)\dd x
\end{equation}
从而
\[f' (x) =\frac{\dd y}{\dd x}\]
即变率等于微分的商，这也就是称变率（导数）为微商的根据．

下面我们用前面的例2.41和例2.42将$\dd y$和$\Delta y$的值作个比较．


\begin{example}
\[y=x^3,\qquad \dd y=3x^2\dd x,\qquad \Delta y=3x^2\dd x+(3x+\dd x)(\dd x)^2\]

当$x=1$时，$\dd y=3\dd x$, $\Delta y=3\dd x+(3+\dd x)(\dd x)^2$.

\begin{center}
\begin{tabular}{cccccc}
    \hline
$\dd x$ & 10& 1& 0.1 & 0.01& $\cdots$\\
    \hline
$\dd y$& 30&3&0.3& 0.03&$\cdots$\\
\hline
$\Delta y$& 1330& 7&0.331&0.030301&$\cdots$\\
    \hline
\end{tabular}
\end{center}
我们注意到$|\dd x|$愈小，$\dd y$与$\Delta y$的值愈靠近，而且 $\Delta y$与$\dd y$的差的绝对值比$\dd x$的绝对值变小得更快．
\end{example}

\begin{example}
\[f(x)=x^2-\frac{1}{x},\qquad \dd f(x)=f'(x)\dd x=\left(2x+\frac{1}{x^2}\right)\dd x\]
\[\Delta y=2x\dd x+(\dd x)^2+\frac{\dd x}{x(x+\dd x)}\]
令$x=2$，则$\dd y=4.25\dd x$, $\Delta y=\left[4+\dd x+\frac{1}{2(2+\dd x)}\right]\dd x$
\begin{center}
    \begin{tabular}{cccccc}
        \hline
    $\dd x$ & 1& 0.5& 0.1 & 0.01& $\cdots$\\
        \hline
    $\dd y$& 4.25&2.13&0.425& 0.0425&$\cdots$\\
    \hline
    $\Delta y$& 5.17& 2.35&0.434&0.0426&$\cdots$\\
        \hline
    \end{tabular}
    \end{center}
上面的例子验证了我们在前面讨论的结论：当自变量的改变量$\Delta x$很小时，可以用微分$\dd y$来近似地代替函数改变量
$\Delta y$.
\end{example}

\begin{example}
    半径为$r$的金属圆片加热膨胀，求当半径从$r$变到$r+\Delta r$时，圆面积$S$变化大小和此时圆面积$S$与$\Delta r$相应的微分$\dd S$, 并试用几何图形解释所得的结果．
\end{example}

\begin{solution}
    已知圆面积$S=\pi r^2$, 设$\Delta S$是半径从$r$变到$r+\Delta r$时圆面积改变量，即圆环的面积（图2.15），则
\begin{align*}
    \Delta S=\pi(r+\Delta r)^2-\pi r^2&=\pi\left[2r\Delta r+(\Delta r)^2\right]\\
    &=\frac{1}{2}\left[2\pi r+2\pi(r+\Delta r)\right]\Delta r
\end{align*}
此时，$\dd S=2\pi r\Delta r$.

\begin{figure}[htp]\centering
	\begin{minipage}[t]{0.48\textwidth}
		\centering
        \begin{tikzpicture}[>=latex, scale=.8]
\draw[thick](0,0) circle(2);
\draw[thick](0,0) circle(2.8);
\draw(0,0)--(2.8,0);
\draw[|<->|](0,-.2)--node[below]{$r$}(2,-.2);
\draw[|<->|](2.8,.2)--node[above]{$\Delta r$}(2,.2);
    \end{tikzpicture}
    \caption{}
\end{minipage}
\begin{minipage}[t]{0.48\textwidth}
    \centering
\begin{tikzpicture}[>=latex, scale=.8]
\tkzDefPoints{-1.5/0/A, 1.5/0/B, 2.7/1.2/C, -2.7/1.2/D}
\tkzDrawPolygon(A,B,C,D)
\tkzDefPoints{1.5/1.2/E, -1.5/1.2/F}
\tkzDrawPolygon[pattern=north east lines](A,D,F)
\tkzDrawPolygon[pattern=north west lines](B,C,E)
\node at (0,1.2)[above]{$2\pi(r+\Delta r)$};
\node at (0,0)[below]{$2\pi r$};
\draw[|<->|](3,0)--node[fill=white]{$\Delta r$}(3,1.2);
\end{tikzpicture}
\caption{}
\end{minipage}
\end{figure}

把圆沿半径剖开，把圆周拉直，圆环就“变成”一个梯形，它的两个底边长分别是$2\pi r$, $2\pi(r+\Delta r)$, 高为$\Delta r$（图2.16），其中的矩形面积就是微分$\dd S=2\pi r\cdot \Delta r$, 它表示圆面积$S$在从$r$到$r+\Delta r$的整段变化过程中都以同一瞬时变率$(\pi r^2)'=2\pi r$作匀速变化所产生的面积增量．图2.16中的两个小直角三角形的面积表示$\Delta S$与$\dd S$的误差，即
\[\Delta S-\dd S=\pi  (\Delta r)^2\]
它比$|\Delta r|$更快地变小，即
\[\lim_{\Delta r\to 0}\frac{\Delta S-\dd S}{\Delta r}=\lim_{\Delta r\to 0}\pi(\Delta r)=0\]
\end{solution}

\subsubsection{微分的几何意义}

从几何的观点来说，变率的根本想法就是以直代曲，即以在点$x_0$处的切线$y=f(x_0)＋f'(x_0)(x-x_0)$来逼近曲线$y=f(x)$的微段．因此，$f'(x_0)\Delta x$的几何意义可以从图2.17看出．

\begin{figure}[htp]
    \centering
    \begin{tikzpicture}[>=latex, xscale=3]
\draw[->](-.25,0)--(2,0)node[right]{$x$};
\draw[->](0,-1)--(0,5)node[right]{$y$};
\node at (0,0)[below left]{$O$};
\draw[domain=.1:1.6, samples=100, very thick]plot(\x, {exp(\x)});
\draw(.5,0)node[below]{$M(x_0)$}--(.5,1.65)node[above]{$P$};
\draw(1.4,0)node[below]{$M'(x_0+\Delta x)$}--(1.4,4.055)node[above]{$P'$};
\draw(.1,1.65)--(1.8,1.65);
\draw[thick](1.4,4.055)--(.5,1.65);
\draw[domain=.1:1.8, samples=10, thick]plot(\x, {1.65*\x+0.825});
\draw[<->](.5,1.3)--node[fill=white]{$\Delta x$}(1.4,1.3);
\draw[<->|](1.45,1.65)node[below]{$N$}--node[right]{$f'(x_0)\Delta x$}(1.45,3.135)node[above]{$T$};
\tkzDefPoints{.5/1.65/P, 1.4/1.65/N, 1.4/3.135/P'}
%\tkzMarkAngle[mark=none, size=.3](N,P,P')
\tkzLabelAngle[pos=.2](N,P,P'){$\alpha$}
    \end{tikzpicture}
    \caption{}
\end{figure}

在直角三角形$TPN$中，$PN=\Delta x$, $\tan\alpha=f'(x_0)$, 故
\[NT=PN\cdot \tan\alpha=\Delta x\cdot f'(x_0)=f'(x_0)\Delta x\]
也就
是切线函数$y=f (x_0) +f' (x_0) (x-x_0)$的相应改变量．至于函数的改变量则是曲线$y=f(x)$本身
的纵坐标在这一段的改变量，在图2.17中用线段$NP'$来代表．所以$f(x)$相应于这一段的微分就是以曲线的切线改变量代替函数的改变量．从图2.17可以看出，当$|\Delta x|$减小时，$TP'$跟着减小，但$TP'$比$|\Delta x|$要变小得更快．

回顾上面的讨论，微分的意义是，当$|\Delta x|$充分小时，曲线$y=f(x)$在点$x_0$的切线
\[y=f (x_0) +f' (x_0) (x-x_0)\]
的改变量$\dd y=f'(x_0)\Delta x=f'(x_0)(x-x_0)$, 就是相应于函数改变量$\Delta y=f(x_0+\Delta x)-f(x_0)$的微分，它可以近似地表达函数改变量，而差是一个比$|\Delta x|=|x-x_0|$变小得更快的微量．因此我们得到
\[\Delta y\approx \dd y\qquad \text{或}\qquad f(x_0+\Delta x)-f(x_0)\approx f'(x_0)\Delta x\]
上式可改写为
\[f (x_0+\Delta x) \approx f (x_0) +f' (x_0) \Delta x\]
可见$f(x)$在$x_0+\Delta x$点的函数值，可以通过它在$x_0$的函数值$f(x_0)$加上一个与$\Delta x$成正比例的量$f'(x_0)\Delta x$来表达，比例系数就是$f$在点$x_0$的微商．上面的公式常用来作近似计算．

\begin{example}
    计算$\sqrt{101}$
\end{example}


\begin{solution}
设$y=\sqrt{x}$，令$x_0=100$, $\Delta x=1$，则
\[\dd y=\frac{\Delta x}{2\sqrt{x_0}}=\frac{1}{2\sqrt{100}}=\frac{1}{20}=0.05\]
因此：
\[\sqrt{101}\approx \sqrt{100}+\dd y=10+0.05=10.05\]
\end{solution}



\begin{example}
一个立方体形状的盒子，它的设计边长等于4(dm), 但允许误差为0.05(dm), 求盒子体积的可能误差．
\end{example}


\begin{solution}
设盒子的边长由$x$改为$x+\dd x$, 那么盒子体积的改变量是
\[\Delta V= (x+\dd x)^3-x^3\]
但是
\[\Delta V\approx \dd V=3x^2\dd x\]
令$x=4$, $\dd x=\pm 0.05$, 那么
\[\dd V=3\x16\x (\pm 0. 05) =\pm 2.4\]
因此，体积的可能误差近似地等于$\pm 2.4({\rm dm})^3$.
\end{solution}

\begin{example}
    证明当$|x|$很小时，$\ln(1+x)\approx x$.
\end{example}

\begin{solution}
    设$f(x)=\ln(1+x)$, 令 $x_0=0$, $\Delta x=x$，则
\[\dd f (x) =f' (x_0) \Delta x=\frac{1}{1+x_0}\Delta x=x\]
因此：
\[f (x_0+\Delta x) =\ln (1+x) \approx \ln 1+x=x\]
即：$\ln(1+x)\approx x$
\end{solution}


\begin{example}
    当$|x|$很小时，导出近似公式：$\tan x\approx x$.
\end{example}


\begin{solution}
    设$f(x)=\tan x$, 令 $x_0=0$, $\Delta x=x$, 则
\[\dd f (x) =\sec^2x_0 \dd x=\sec^2 0\cdot x=1\cdot x=x\]
因此：
\[f (x_0+\Delta x) =\tan x\approx \tan 0+x=x\]
即：$\tan x\approx x$.
\end{solution}

\begin{example}
如图2.18, 加工锥形工件时，已知工件两头直径分别为$d_1,d_2$, 长度为$\ell$, 当斜角$\alpha$很小时，导出近似公式：
\[\alpha=28.6^{\circ}\x\frac{d_1-d_2}{\ell}\]
\end{example}

\begin{figure}[htp]
    \centering
\begin{tikzpicture}[>=latex]
    \draw[|<->|](-0.5,.9)--node[fill=white]{$d_2$}(-.5,-.9);
\draw[|<->|](4.8,1.5)--node[fill=white]{$d_1$}(4.8,-1.5);
\draw[very thick](4,0) ellipse [x radius=.5, y radius=1.5];
\draw[dashed](0,0) ellipse [x radius=.3, y radius=.9];
\draw[dashed](0,-.9) rectangle (4,.9);
\draw[very thick](0,-.9)--(4,-1.5);\draw[very thick](0,.9)--(4,1.5);
\draw[dashdotted](0,0)--(4,0);
\draw[dashed](4,-1.5)--(4,1.5);
\draw[|<->|](0,1.8)--node[fill=white]{$\ell$}(4,1.8);
\draw[very thick] (0,.9) arc [start angle =90, end angle =270, x radius=.3, y radius=.9] ;
\draw (1.5,.9) arc (0:9:1.5)node[below right]{$\alpha$};
\end{tikzpicture}  
    \caption{}
\end{figure}

\begin{solution}
由图：$\tan\alpha=\frac{d_1-d_2}{2\ell}$. 当$\alpha$很小时，根据$\tan\alpha\approx \alpha$，得到：
\[\alpha\approx \frac{d_1-d_2}{2\ell}\]
又：$1\text{弧度}=\left(\frac{180}{\pi}\right)^{\circ}\approx 57.3^{\circ}$，因此：
\[\alpha\approx 57.3^{\circ}\x \frac{d_1-d_2}{2\ell}\approx 28.6^{\circ}\x \frac{d_1-d_2}{\ell}\]
\end{solution}

\subsubsection{微分的运算}

如果函数$f(x)$可微，即导数$f'(x)$存在，求函数的微分只要求出导数再乘以自变量的改变量$\Delta x$就行了，求已知函数的微分或求导数的方法都叫做微分法．

我们已经知道当$x$是自变量时，$\dd x=\Delta x$, 原来的微分表达式$\dd y=f'(x)\Delta x$就可以改写为$\dd y=f'(x)\dd x$, 例如$\dd(\sin x)=\cos x \dd x$. 所以我们可以根据函数的和、积、商的求导数法则得到函数的和、积、商的求微分法则，即
\begin{enumerate}
    \item $\dd(cu)=c\dd u$
    \item $\dd(u\pm v)=\dd u\pm \dd v$
    \item $\dd(uv)=v\dd u+u\dd v$
    \item $\dd\left(\frac{u}{v}\right)=\frac{v\dd u-u\dd v}{v^2}$
\end{enumerate}
这里只证明最后一个法则，其余请读者自证．

\begin{proof}
\begin{align*}
\dd\left(\frac{u}{v}\right)=\left(\frac{u}{v}\right)'\dd x&=\frac{u'v-uv'}{v^2}\dd x\\
&=\frac{v(u'\dd x)-u(v'\dd x)}{v^2}\\
&=\frac{v\dd u-u\dd v}{v^2}
\end{align*}
\end{proof}


现在我们来研究求复合函数的微分法则．

如果$y=f(x)$, $x=\varphi (t)$, 那么$y$是$t$的复合函数$y=f(\varphi (t))$, 这时$x$和$y$的微分根据定义都要以$t$为自变量来确定．在这情况下，须注意$\dd x\ne \Delta x$, 比如若$x=t^2$, 则$\dd x=2t\dd t$, 而$\Delta x=(t+\Delta t)^2-t^2=2t\Delta t+(\Delta t)^2$, 显然$\dd x\ne \Delta x$, 对于复合函数$y=f(\varphi (t))$的微分，我们有下面重要定理．

\begin{blk}
    {定理} 设$y=f(x)$, $x=\varphi (t)$都是可微函数，$x$是中间变量，那么关系式
$\dd y=f' (x) \dd x$
仍成立．
\end{blk}

\begin{proof}
   \[ \dd y=[f(\varphi (t))]'\dd t\]
根据复合函数求导数法则有
\[[f (\varphi  (t) ) ] '=f' (\varphi  (t) ) \cdot \varphi ' (t)\]
因此
\[\dd y=f' (\varphi  (t) ) \cdot \varphi' (t) \dd t\]

因为$x=\varphi  (t)$, $\dd x=\varphi ' (t) \dd t$，最后得到
\[\dd y=f' (x) \dd x\]
\end{proof}

这个定理是说不论$x$是自变量或是中间变量，函数$f(x)$
的微分形式不变，只须注意，若$x$是中间变量，$\dd x$不是任意增量$\Delta x$, 而是$x$的微分，$\dd y$不能写成$f'(x)\Delta x$即$\dd y\ne f'(x)\Delta x$, 这个性质叫做\textbf{微分不变性}．有了微分不变性，我们在作微分运算时，就无需指明对哪一个变量的微分，而在求函数的变率时，总要指明对中间变量求导数还是对自变量求导数．另外，求复合函数的导数法则用微分符号写出，就是
\[\frac{\dd y}{\dd t}=\frac{\dd y}{\dd x}\cdot \frac{\dd x}{\dd t}\]
它就成为简单的代数恒等式了．

\begin{example}
    求由方程$y^3=x^2+xy+y^2$确定的隐函数的导数$\frac{\dd y}{\dd x}$.
\end{example}

\begin{solution}
    对等式两边求微分得到
\[3y^2\dd y=2x\dd x+ (x\dd y+y\dd x) +2y\dd y\]
移项
\[(3y^2-x-2y)\dd y=(2x+y)\dd x\]
最后得到
\[\frac{\dd y}{\dd x}=\frac{2x+y}{3y^2-x-2y}\]
\end{solution}

\begin{example}
    求$\dd\sin [\tan(x^2-x+1)]$.
\end{example}


\begin{solution}
\begin{align*}
    \dd\sin \left[\tan(x^2-x+1)\right]&=\cos\left[\tan(x^2-x+1)\right]\dd \tan(x^2-x+1)\\
    &=\cos\left[\tan(x^2-x+1)\right]\cdot \sec^2 (x^2-x+1)\cdot \dd (x^2-x+1)\\
    &=\cos\left[\tan(x^2-x+1)\right]\sec^2(x^2-x+1)\cdot (2x-1)\dd x 
\end{align*}
\end{solution}
（例2.52的计算中，将微分符号“$\dd$”逐步顺着复合函数
的层次从$\sin$到$\tan$到$x^2-x+1$, 由外到里，犹如剥笋．）


\begin{ex}
\begin{enumerate}
    \item 边长4cm的正方形铁皮，在加热中边长增加了0.001cm, 
求此时刻面积的可能误差$\Delta S$, 和可能误差的近似值$\dd s$, 又正方形的面积近似值是多少？
\item 半径为10cm的球，在冷却中$R$缩短了0.001cm, 求此时
刻体积的微分$\dd V$和体积的近似值．
\item 当$|x|$很小时，导出下列近似公式：
\begin{multicols}{2}
\begin{enumerate}
    \item  $\sqrt{1+x} \approx 1+\frac{x}{2}$
\item $e^{x} \approx 1+x$
\item  $(1+x)^{n} \approx 1+n x$
\item  $\sqrt[n]{1+x} \approx 1+\frac{1}{n} x$
\item  $\sin x \approx x$
\item  $\arcsin x \approx x$
\item  $\ln (1+\sin x) \approx x$
\end{enumerate}
\end{multicols}

\item  计算下列各式的值：
\begin{multicols}{2}
\begin{enumerate}
\item  $\sqrt[3]{1.02}$
\item  $\sqrt[3]{0.998}$
\item   $e^{-0.1}$
\item   $\tan 0.01$ 
\item   $\sin 0.1^{\circ}$
\item   $\ln 1.0021$   
\end{enumerate}
\end{multicols}

\item 求下列函数的微分:
\begin{multicols}{2}
\begin{enumerate}
    \item $y=\frac{x+5}{x^{2}-2}$
    \item  $y=\left(\sqrt{1+x^{2}}\right)^{n}$
    \item  $y=\ln (\sin \sqrt{x})$
    \item  $y=\sin ^{3}(\sqrt{x})$
    \item  $y=\arcsin \sqrt{2 x}$
    \item  $y=\arctan (\cos x)$ 
    \item $y=\sin (\ln x)$
    \item $y=x^{n} \ln (n x)$
    \item $y=e^{x} \ln x$
    \item $y=\left(e^{x}+e^{-x}\right)^{2}$
\end{enumerate}
\end{multicols}

\item 一直圆锥内接于半径等于$R$的球中，球心在直圆锥的内部
且到圆锥底面的距离等于$p$, 若$p$增加一个微量$x$而球半径$R$不变，求证圆锥体积的增量约等于$\frac{\pi}{3}(R+p)(R-3p) x$.
\item 直圆锥的底面半径为$r$ cm, 高为$h$ cm, 若半径$r$和高$h$分别
有一个增量$\rho$ cm和$\lambda$ cm, 求证圆锥体积的增量近似地等于$\left(\frac{2\rho}{r}+\frac{\lambda}{h}\right)$.
\end{enumerate}
\end{ex}

\subsection*{习题2.2}
\begin{enumerate}
    \item 求下列函数的导函数：
	\begin{enumerate}
    	\begin{multicols}{2}
			\item $f(x)=2 x^{3}-3 x^{3}+4 x-5 $
			\item $f(x)=x^{5}-6 x^{2}+13$
			\item $f(x)=\frac{x^{3}-4 x}{x^{2}+1} $
			\item $f(x)=\cos m x \sin n x$       
    	\end{multicols}

		\item $f(x)=\left[\sin x+\cos x+\left(x^{2}+x+1\right)\right]^{2}$
		\item $f(x)=\tan\left[\sin x+\cos x+\left(x^{2}+x+1\right)\right]$
		\begin{multicols}{2}
			\item $f(x)=\cot \left(x^{2}+1\right)$
			\item $f(x)=\sqrt{a^{2}+b^{2} x^{2}}$
			\item $f(x)=\sqrt[3]{x^{4}+x^{2}+1}$
			\item $f(x)=\cos \left(\sqrt[3]{3 x^{4}+5 x^{2}+6}\right)$
			\item $f(x)=(\sqrt{x}+1)\left(\frac{1}{\sqrt{x}}-1\right)$
			\item $f(x)=x \sin x+\cos x$
			\item $f(x)=2^{\sin ^{2}\tfrac{1}{x}} $
			\item $f(x)=x^{\sin x}$       
		\end{multicols}
		\begin{multicols}{2}
		\item $f(x)=e^{-x^{2}} \cos \left(e^{-x^{2}}\right) $
		\item  $f(x)=\arctan\sqrt{x^{2}-2 x}$
		\item  $f(x)=x e^{1-\cos x}$
		\item  $f(x)=x \arccos x-\sqrt{1-x^{2}}$
		\item  $f(x)=(\sin x)^{\cos x}$
		\item  $f(x)=\sqrt[5]{\frac{x-5}{\sqrt[5]{x^{2}+2}}}$
		\end{multicols}
		\item  $f(x)=x-\ln \left(2 e^{x}+1+\sqrt{e^{2 x}+4 e^{x}+1}\right)$
		\begin{multicols}{2}
		\item $f(x)=\sqrt{x \sqrt{x \sqrt{x}}}$
		\item  $f(x)=\log _{a} e ^{\sin ^{2} x+1}$
		\end{multicols}
		\item  $f(x)=\ln \tan  \frac{x}{2}-\cot x \cdot \ln (1+\sin x)-x$
	\end{enumerate}

\item 验证函数$f(x)=e^x \sin x$满足方程
\[f'' (x) -2f' (x) +2f (x) =0.\]

\item 验证函数$f(x)=\frac{x-3}{x+4}$满足方程
\[2 [f' (x) ]^2= [f (x) -1] f'' (x).\]

\item $P$点在直线上运动，在时刻$t$, 点$P$与直线上的定点$O$的距离由方程
$s^2=a^2+V^2t^2$
给出，其中$a$和$V$是常量，试用$s$表示$P$点在时刻$t$的速度和加速度．

\item 在$xy$平面上，求通过点$P(-3, 6)$且与曲线$y=x^3-5x^2+x+5$
相切，并且切点的横坐标满足$x\ge 0$的切线方程．

\item 求证方程$16x^4-24x^2+16x-3=0$在$(0, 1)$内有
一个三重根，又在$(-2,-1)$内有一个单根．

\item 设三次函数$y=x^3-3ax^2+bx+c$ （其中$a,b,c$是
常数）．
\begin{enumerate}
    \item 该曲线与直线$y=2x-1$切于点$(2, 3)$, 试用$a$的式子表示$b,c$;
    \item 另外，该曲线还与$y=2x+3$相切，试确定$a,b,c$的值．
\end{enumerate}
\item 在中午十二点整，甲船以6公里/小时的速率向东行，乙船在甲船之北16公里，以8公里/小时的速率向南行，在下午一点整两船相近之速率为多少？在下午两点整两船相离之速率是多少？
\item 已知圆心为$O$, 半径为1的圆与直线$\ell$相切于点$A$, 一动点$P$自切点$A$沿直线$\ell$向右移动时，取弧$\wideparen{AC}$的长为$\frac{2}{3}AP$，
直线$PC$与直线$AO$交于点$M$(图2.19), 又知当$AP=\frac{3\pi}{4}$时，点$P$的速率为$v$, 求这时$M$点的速率．

\begin{figure}[htp]
    \centering
\begin{tikzpicture}
\draw(-2,0)--(2,0)node[right]{$\ell$};
\draw[very thick](0,1) circle (1);
\draw(0,0)node[below]{$A$}--(0,2.5);
\draw(0,1.7)node[left]{$M$}--(1.7,0)node[below]{$P$};
\node at (30:1)[right]{$C$};
\tkzDefPoints{0/1/O}
\tkzLabelPoints[left](O)
\tkzDrawPoints(O)

\end{tikzpicture}    
    \caption{}
\end{figure}

\item 溶液自深18cm, 顶直
径12cm的正圆锥形漏斗中漏入一直径为10cm的圆柱形筒中，开始时漏斗中盛满了水，已知当溶液在漏斗中深为12cm时，其水平面下落速率为1cm/s, 问此时圆柱形筒中水平面上升之速率为多少？
\item 求曲线$y=\arctan\ln x$与$x$轴的交角．
\item 设质点$P$在时刻$t$的位移方程是
\[x=\frac{58}{240}+A\cos\left(\sqrt{24}gt+\alpha\right)\]
其中$A,\alpha$是待定常数，$g\approx 9.8{\rm m/s}$，且知当$t=0$时，$x=0$, $\frac{\dd x}{\dd t}=0$，求$A$和$\alpha$.


\begin{figure}[htp]
    \centering
\begin{tikzpicture}[>=latex, scale=1.7]
    \draw[dashdotted] (-3.5,0)--(1.2,0);
    \draw (-3.5,.5) rectangle (-1.5,-.5);
    \draw (0,0) circle (.75);
    \node at (-.75,0)[below left]{$K_0$};
    \node at (.75,0)[above right]{$K_1$};
\draw [fill=gray!20](-2.5,.5) rectangle (-2.6,-.5);
\draw[very thick] (0,0)node[below]{$O$}--node[above]{$r$}(90+30:.75)node[above]{$K$}--node[above right]{$\ell$}(-2.5,0)node[below]{$M$};
\draw [dashed](-3.3,.5)node[below]{$M_0$}--(-3.3,-.5);
\draw [dashed](-1.7,.5)node[below]{$M_1$}--(-1.7,-.5);
\draw (-.25,0) arc (180:180-60:.25)node[below left]{$\varphi$};
\draw[dashed] (90+30:.75)--(-.75/2,0);
\draw[|<->|] (-3.3,-.7)--node[fill=white]{$x$}(-2.5,-.7);

\draw (-2,0) arc (0:15:.5)node[below right]{$\alpha$};


\end{tikzpicture}
    \caption{}
\end{figure}



\item 图2.20表示曲柄机构装置的图样，连接杆$MK$
的长$\ell=125$cm, 曲柄$OK$的长$r=25$cm, 连杆与汽缸的轴$OM$成$\alpha$角，曲柄与同一轴成$\varphi$角．
\begin{enumerate}
    \item 求用$\varphi$的式子表示活塞$M$的位移$x$的公式；
    \item 如果$\frac{\dd \varphi}{\dd t}=\omega$ （$\omega$
    是常量），求活塞$M$的速率．
\end{enumerate}
\item 在抛物线$y=ax^2+bx+c\; (a\ne 0)$上两点作切线，
它们在何处相交？
\item 求抛物线$y^2=4ax$在点$(a\lambda^2, 2a\lambda)$的切线方程，并求它的两条互相垂直的切线的交点的轨迹方程．
\item 
\begin{enumerate}
    \item 求曲线$4y^3=27x^2$ 在点$(2t^3, 3t^2)$的切
线方程；
\item 求互相垂直的两条切线的交点的轨迹方程．
\end{enumerate}

\item 计算$\sqrt{\frac{2.037^2-1}{2.037^2+1}}$的近似值．
\item 计算$\cos151^{\circ}$的近似值．
\item 设$f(x)=e^{0.1x(1-x)}$, 试计算$f(1.05)$的近
似值．
\item 若球体积的相对误差$\frac{\Delta V}{V}=0.01$, 试计算球半径的
相对误差$\frac{\Delta R}{R}$约等于多少？
\end{enumerate}


\section{微商运算的初步应用{}---{}函数的增减与极值}

\subsection{局部极值和中值定理}
我们在第一章中曾说明在闭区间$[a,b]$上连续的函数必有最大值和最小值．函数$f$在闭区间$M=[a,b]$上的最大值的大小涉及到两个因素，既依赖于$f$，也依赖于$M$.

如图2.21，若$N=[c,d]\subset M=[a,b]$，则
\[
\max_N f \le \max_M f.
\]
相应地，
\[
\min_N f \ge \min_M f.
\]
在这两种情形里都可能出现严格不等式（即不带等号的不等式）．

\begin{figure}[htp]
    \centering
    \begin{tikzpicture}[>=latex]
        \draw[->](-.5,0)--(5.5,0)node[right]{$x$};
        \draw[->](0,-.5)--(0,4.5)node[right]{$y$};
        \node at (0,0)[below left]{$O$};
        \draw[rounded corners=10pt, very thick](.5,1)--(1.2,2)--(2,1.6)--(2.3,2.2)--(2.9,3.6)--(3.6,3.4)--(4.3,2.5)--(4.8,4);
        \node at (3,4){$y=f(x)$};
        \draw(.5,1)--(.5,0)node[below]{$a$};
        \draw(1.2,1.85)--(1.2,0)node[below]{$c$};
        \draw(4.3,2.65)--(4.3,0)node[below]{$d$};
        \draw(4.8,4)--(4.8,0)node[below]{$b$};
        \end{tikzpicture}
    \caption{}
\end{figure}

现在我们将阐明一种只决定于$f$的性质的局部极大与极小的概念．

\begin{blk}
    {定义} 设函数$f$在点$x$的某个$\delta>0$的邻域，即$(x-\delta,x+\delta)$内有定义，如果对于充分小的$|\Delta x|<\delta$, 能使$x+\Delta x\in (x-\delta,x+\delta)$, 而且对于一切$x'=x+\Delta x\in (x-\delta,x+\delta)$, 都有
\[f(x+\Delta x)\le f(x)\quad \text{或}\quad f(x+\Delta x)\ge f(x)\]
那么函数$f$在点$x$达到局部极大值（或局部极小值）．
\end{blk}

\begin{blk}
    {费马定理} 如果函数$f$在区间$(a,b)$上有定义，并且在这个区间内的一点$c$取局部极大值（或极小值），又若$f$在$c$点可微，那么
    \[f' (c) =0\]
\end{blk}

\begin{proof}
    为确定起见，设$f(x)$在点$c\in (a,b)$取最大值、于是对于一切点$x\in (a,b)$, 就有
\[
f (x) \le f (c). 
\]
由导数定义：
\[\lim_{x\to c}\frac{f (x) -f (c)}{x-c} =f' (c) \]
并且极限$f'(c)$的存在与$x$是从右边还是从左边趋向$c$无关．

\begin{enumerate}
    \item 如果$x>c$，我们有$\frac{f (x) -f (c)}{x-c}\le 0$，从而
\[\lim_{x\to c^+}\frac{f (x) -f (c)}{x-c}\le 0\]
\item 如果$x<c$，我们有$\frac{f (x) -f (c)}{x-c}\ge 0$，从而
\[\lim_{x\to c^-}\frac{f (x) -f (c)}{x-c}\ge 0\]
\end{enumerate}
根据假设$f$在$c$点可微，因此这两个极限必相等，且应等于$f'(c)$, 这表示$f'(c)\ge 0$和$f'(c)\le 0$, 由此得：$f' (c) =0$.

$f$在$c$点有局部极小值的情形留给读者自证．
\end{proof}

\begin{blk}
  {罗尔定理} 如果函数$f$在闭区间$[a,b]$上连续，在开区间$(a,b)$内可微，且在两端点的函数值相等，即$f(a)=f(b)$, 那么至少存在一个点$c$, $c\in (a,b)$, 使$f'(c)=0$.  
\end{blk}

\begin{proof}
如果$f$在$[a,b]$上恒为常数，无疑对于任何
$c\in (a,b)$有$f'(c)=0$. 因此，可设$f$在$[a,b]$中不为常数，此时$f$在$[a,b]$上的最大值$M$和最小值$m$至少有一个不等于$f(a)=f(b)$, 为确定起见，不妨设$M\ne f(a)=f(b)$．故$f$在$[a,b]$上的最大值$M$不会在$a$点或$b$点达到，这时根据闭区间上连续函数的性质，在$(a,b)$内至少有一点$c$, 使$f(c)=M$, 于是，由费马定理，即得$f'(c)=0$. 

本定理的几何意义如图2.22, 曲线$y=f(x)$在$a,b$处有相同高度，如果在点$c$处、曲线达到最高度，并有切线存在，那么，曲线在该处的切线与$x$轴平行．同样，在曲线的最低点处的切线也与$x$轴平行．
\end{proof}

\begin{figure}[htp]
    \centering
    \begin{minipage}[t]{0.48\textwidth}
    \centering
\begin{tikzpicture}[>=latex, scale=1]
    \draw[->](-.5,0)--(5,0)node[right]{$x$};
    \draw[->](0,-.5)--(0,2.8)node[right]{$y$};
    \node at (0,0)[below left]{$O$};
    \draw[domain=.5:4.5, samples=100, thick]plot(\x, {sin((\x-.5)*pi/2 r)+.5});
    \draw[dashed](.5,0)node[below]{$a$}--(.5,.5)--(4.5,.5)--(4.5,0)node[below]{$b$};
    \draw(1.5,0)node[below]{$c$}--(1.5,1.5);
\end{tikzpicture}
    \caption{}
    \end{minipage}
    \begin{minipage}[t]{0.48\textwidth}
    \centering
    \begin{tikzpicture}[>=latex, scale=1.5]
        \draw[->](-1.5,0)--(1.5,0)node[right]{$x$};
        \draw[->](0,-.25)--(0,2)node[right]{$y$};
        \draw[thick](-1,0)node[below]{$(-1,0)$}--(0,1)node[above right]{$(0,1)$}--node[right]{$g(x)=1-|x|$}(1,0)node[below]{$(1,0)$};
        \node at (0,0)[below left]{$O$};
    \end{tikzpicture}
    \caption{}
    \end{minipage}
    \end{figure}


\begin{rmk}
    注意 如果罗尔定理中的条件不被满足，结论就不成立．例如$g(x)=1-|x|$在$[-1, 1]$上连续，且$g(-1)=g(1)=0$, 但在$x=0$处，$g'(x)$不存在，对于这个$g$, 在$(-1,1)$中就没有$c$能使$g'(c)=0$. 它的图象如图2.23所示，
\end{rmk}

\begin{blk}
    {拉格朗日中值定理} 如果函数$f$在闭区间$[a,b]$上连续，在开区间$(a,b)$内可导，那么在$(a,b)$内至少有一点$c$, 使得
\[f'(c)=\frac{f(b)-f(a)}{b-a}\]
\end{blk}

罗尔定理是拉格朗日定理的特殊情形，我们借助于罗尔定理来证明这个重要的定理．

\begin{proof}
 我们作一个辅助函数$\Phi$使它具有函数f的性质：即它在$[a,b]$上连续且在$(a,b)$内可微，同时满足$\Phi (a)=\Phi (b)$. 为此，设$\Phi  (x) =f (x) +k (x-a)$, 并由条件$\Phi (a)=\Phi (b)$来决定参数$k$. 

因为当$x=a$时，有$\Phi (x)=f(a)$; 当$x=b$时，有$\Phi (b)=f(b)+k(b-a)$, 所以
\[k=-\frac{f (b) -f (a)}{b-a}\]
这样得到的函数
\[\Phi  (x) =f (x) -\frac{f (b) -f (a) }{b-a}(x-a)\]
就满足罗尔定理中的所有的条件，由罗尔定理，至少存在一点$c\in (a,b)$使得$\Phi '(c)=0$,
但
\[\Phi ' (x) =f' (x) -\frac{f (b) -f (a)}{b-a}\]
所以
\[\Phi ' (c) =f' (c) -\frac{f (b) -f (a)}{b-a} =0\]
即
\[f' (c) =\frac{f (b) -f (a)}{b-a}\]
\end{proof}

为了便于应用，拉格朗日定理的结论，通常写成如下形式：
\begin{equation}
   f (b) -f (a) =f' (c) (b-a),\qquad a<c<b 
\end{equation}
在上式中，令$b=a+h$，$c=a+\theta h$，$0<\theta<1$，得到中值定理的另一形式：
\begin{equation}
    f (a+h) -f (a) =hf' (a+\theta h) ,\qquad 0<\theta<1
\end{equation}

中值定理的几何意义如图2.24所示．在闭区间$[a,b]$上有一条连续的曲线，且过曲线$y=f(x)$上每一点（端点除外），都可以作出一条切线，那么在曲线上至少有一点$M(c, f( c)),\; c\in (a,b)$, 使得过点$M$的切线与弦$AB$平行．

\begin{figure}[htp]
    \centering
\begin{tikzpicture}[>=latex, scale=1.5]
\draw[->](-.5,1)--(4,1)node[right]{$x$};
\node at (0,1)[below left]{$O$};
\draw[->](0,.5)--(0,3)node[right]{$y$};
\draw[domain=.8:3.2, samples=100, very thick]plot(\x, {(\x-1)*(\x-2)*(\x-3)+2});
\draw(.8,1)node[below]{$a$}--(.8,1.472)node[left]{$A$};
\draw(3.2,1)node[below]{$b$}--(3.2,2.528)node[right]{$B$};
\draw[thick](.8,1.472)--(3.2,2.528);
\draw[dashed](.8,1.472)--(3.2,1.472)node[right]{$C$};
\draw[dashed](1.31,2.3603)node[above left]{$M(c,f(c))$}--(1.31,1)node[below]{$c$};

\draw[domain=.5:2.5, samples=10, thick]plot(\x, {0.44*\x+1.784});

\end{tikzpicture}
    \caption{}
\end{figure}



\begin{example}
求$f(x)=1+\frac{4}{x}$在$(1, 4)$内满足中值定理$f(b)-f(a)=f'(c)\cdot (b-a)$的点$c$的值．
\end{example}


\begin{solution}
    显然$f$在$[1, 4]$上连续，又$f'(x)=-\frac{4}{x^2}$, 故$f$在$(1, 4)$内可微，根据中值定理，存在点$c\in (1, 4)$, 使得
\[    f (4) -f (1) = (4-1) f' (c)\]
成立．

$\because\quad f (4) =2,\quad f (1) =5$,

$\therefore\quad 2-5=3 \left(-\frac{4}{c^2}\right),$
即：$c^2=4$，得$c=\pm 2$.

又$\because\quad 2\in (1, 4),\quad -2\notin (1, 4)$,

$\therefore\quad c=2$. 于是
\[    f (4) -f (1) = (4-1) f' (2) \]
\end{solution}

\begin{example}
    设$f(x)=a_0+a_1x+a_2x^2+\cdots+a_nx^n$, 求证在$f'(x)=0$的相邻二根之间，至多有$f(x)=0$的一个根．   
\end{example}

\begin{proof}
因为多项式函数处处连续和可微，所以若在$f'(x)=0$的两相邻实根$x_1,x_2$之间有$f(x)=0$的二实根$\alpha,\beta\; (\alpha<\beta)$, 那么根据罗尔定理，至少存在$f'(x)=0$的一个根$c$, 使得$\alpha<c<\beta$．但是$(\alpha,\beta)\subset (x_1,x_2)$，于是$x_1<c<x_2$，由此得出矛盾，所以在$f'(x)=0$的相邻两根之间，至多有$f(x)=0$的一个根．
\end{proof}

\begin{example}
    求证不等式$\frac{h}{1+h}<\ln(1+h)<h$，
    对于$h\in (-1, 0)\cup (0,+\infty)$成立．
\end{example}


\begin{proof}
    设$f(x)=\ln x$, 于是$(\ln x)'=\frac{1}{x}$. 由于$\ln x$在区间$[a,b],\; (0<a<b)$上连续且在$(a,b)$内可微，根据中值定理$(a,b)$内存在$c$，使得：
\begin{equation}
    \ln b-\ln a=\frac{1}{c}(b-a),\quad a<c<b
\end{equation}

若$h>0$, 令$a=1$, $b=1+h$, 于是$c=1+\theta h$, 这样(2.26)可写成
\[\ln (1+h) =\frac{h}{1+\theta h},\quad 0<\theta<1\]
若把上式中的$\theta$看作变数，由0变到1, 那么$\ln(1+h)=\frac{h}{1+\theta h}=\frac{1}{\frac{1}{h}+\theta}$
便由$h$递减到$\frac{h}{1+h}$，所以
\[h>\ln(1+h)>\frac{h}{1+h},\quad h>0\]

若$-1<h<0$，令$a=1+h$，$b=1$，于是$c=1+\theta h$，这样
(2.26)可以写成
\[\ln1-\ln (1+h) =\frac{-h}{1+\theta h}\]
仍得
\[\ln (1+h) =\frac{h}{1+\theta h}\]
同样可证
\[h> \ln (1+h) >\frac{h}{1+h}\]
因此，当$h>-1$且$h\ne 0$时，不等式
\[\frac{h}{1+h}<\ln(1+h)<h\]
成立．
\end{proof}



\begin{ex}
\begin{enumerate}
    \item 曲线$y=x^2$上的某点的切线平行于连接横坐标为$a,b$两点的弦，求出此点．
    \item 曲线$y=\frac{1}{x}$上的某点的切线平行于连接横坐标为$a,b\;(0<a<b)$两点的弦，求出此点．
    \item 若$\varphi(x)=(x-a)^m(x-b)^n$，验证罗尔定理成立．
    \item 用罗尔定理证明，无论$m$为何值，多项式函数$f_m(x)=x^3-3x+m$在$[0, 1]$之间至多有一个根．
    \item 证明下列不等式：
\begin{enumerate}
    \item $|\sin x-\sin y|\le |x-y|$
    \item $\frac{a-b}{a}<\ln \frac{a}{b}<\frac{a-b}{b},\quad 0<b<a$
    \item $\frac{h}{1+h^2}<\arctan h<h,\quad h>0$
    \item 若$0<m<1$，且$x>1$，则：$\frac{m(x-1)}{x^{1-m}}<x^m-1<m(x-1)$
    
    （提示：利用中值定理形式1，取$b=x,\; a=1$）
    \item 若$h>-1$，$h\ne 0$，$n>1$，则：$(1+h)^{\tfrac{1}{n}}<1+\frac{h}{n}$
    \item 若$h>-1$，$h\ne 0$，则：$1-\frac{h}{2}<\frac{1}{\sqrt{1+h}}<1-\frac{h}{2\sqrt{(1+h)^3}}$
\end{enumerate}
\item 说明方程$3x^5-25x^3+60x-20=0$有三个实数根，给出根所在的区间，使每个区间只含一个根．
\end{enumerate}
\end{ex}

\subsection{函数的增减与极值——中值定理的应用}

\subsubsection{函数的增减性}

我们已经学过增函数与减函数的概念，现在要阐明导数的正负与函数的增减的关系．

读者已经知道常量的导数等于零，现在反过来问：如果
函数$f$在给定区间上的导数永远是零，那么这个函数在这个区间上就是一个常量？由拉格朗日定理可得到肯定的回答．

\begin{blk}
  {定理1} 如果函数$f$在开区间$(a,b)$上的每一点的导数等于零，即$f'(x)=0$, 那么函数$f$在$(a,b)$上是一个常数．  
\end{blk}

\begin{proof}
  我们从区间$(a,b)$中固定某一点$x_0$, 再取其它的任何一点$x\in (a,b)$. 对于区间$[x_0,x]$或$[x,x_0]$, 函数$f$满足拉格朗日定理条件，因此，我们可以写出
  \[f (x) -f (x_0) =f' (c) (x-x_0),\qquad c\in (x_0,x)\text{ 或 }(x,x_0)\]
  也就是$c\in (a,b)$, 依已知条件$f'(c)=0$, 因而对任何一点$x\in (a,b)$, $f(x)=f(x_0)=$常数．  
\end{proof}



\begin{example}
    求证：如果$x>1$，那么$2\arctan x+\arcsin\frac{2x}{1+x^2}=\pi$成立．
\end{example}

\begin{proof}
设$f(x)=2\arctan x+\arcsin\frac{2x}{1+x^2}$，因为当$x>1$时，$\left|\frac{2x}{1+x^2}\right|<1$，故对于所指定的$x$值的范围，$f(x)$有定义，又因为：
\begin{align*}
    f'(x)&=2(\arctan x)'+\left(\arcsin\frac{2x}{1+x^2}\right)'\\
    &=\frac{2}{1+x^2}+\frac{1}{\sqrt{1-\left(\frac{2x}{1+x^2}\right)^2}}\cdot \left(\frac{2x}{1+x^2}\right)'\\
    &=\frac{2}{1+x^2}+\frac{1}{\sqrt{1-\left(\frac{2x}{1+x^2}\right)^2}}\cdot \frac{(1+x^2)\cdot 2-2x(2x)}{(1+x^2)^2}\\
    &=\frac{2}{1+x^2}+\frac{2(1-x^2)}{(x^2-1)(1+x^2)}\\
    &=2\left(\frac{1}{1+x^2}-\frac{1}{1+x^2}\right)=0
\end{align*}
所以：
\[2\arctan x+\arcsin\frac{2x}{1+x^2}=c\; \text{（常数）}\]
另一方面，因为
\[f\left(\sqrt{3}\right)=2\arctan\sqrt{3}+\arcsin\frac{\sqrt{3}}{2}=\frac{2\pi}{3}+\frac{\pi}{3}=\pi\]
因此：
\[2\arctan x+\arcsin\frac{2x}{1+x^2}=\pi,\quad (x>1)\]
\end{proof}


\begin{blk}
    {定理2} 若$f(x)$在$[a,b]$连续，在$(a,b)$内$f'(x)＞0$ (或$f'(x)<0$)，那么$f(x)$在$[a,b]$上严格递增（或严格递减）．
\end{blk}


\begin{proof}
    设$a\le x_1<x_2\le b$，由于在闭区间$[x_1,x_2]$上满足拉格朗日定理的条件，因此，在开区间$(x_1,x_2)$上可以找到点$c$，使得
\[f(x_2)-f(x_1)=(x_2-x_1)f'(c)\quad (x_1<c<x_2)\]
根据条件，如果在$(a,b)$内，$f'(x)>0$ (或$f'(x)<0$), 那么$f'(c)>0$ (或$f'(c)<0$), 从而
\[f(x_2)-f(x_1)>0\quad (\text{或}<0)\]
因为对于任意$x_1$和$x_2$ (其中$a\le x_1<x_2\le b$)，上面不等式成立，所以$f(x)$在$[a,b]$上严格递增（或严格递减）．
\end{proof}

定理2的几何意义是明显的，$f'(x)$是函数图象上点$x$处的切线的斜率$\tan\theta$, 这里$0\le\theta \le \pi$. 这斜率的符号指出切线是向上或向下倾斜的，而曲线本身也就随着它向上升或向下降．

如果$f'(x)$在个别的$x$值等于零，其他条件不变，则定理
2仍然成立，这一点很容易证实，就是把区间$[a,b]$用使$f'(x)=0$的点来分成若干个小区间，再在每个小区间上分别应用定理2. 图2.25和2.26表示定理2的几何意义．

\begin{figure}[htp]
    \centering
    \begin{minipage}[t]{0.48\textwidth}
    \centering
    \begin{tikzpicture}[>=latex, scale=.7]
  \draw[->](-2,0)--(5,0)node[right]{$x$};
\draw[->](0,-1)--(0,5)node[right]{$y$};
\node at (0,0)[below left]{$O$};
\draw[domain=-.2:4, samples=100, very thick]plot(\x, {.2*(\x-2)^3+2.5});
\draw[->, domain=-1.5:3, thick, samples=10]plot(\x, {1.2*\x+1.22});
\draw[->, domain=.5:5, thick, samples=10]plot(\x, {1.2*\x-1.03});
\draw(.2,2.5)--(3.5,2.5);
\draw[->, thick](.86+.5,0) arc (0:50.2:.5)node[right]{$\theta'$};
\draw[->, thick](-1.02+.5,0) arc (0:50.2:.5)node[right]{$\theta$};
    \end{tikzpicture}
    \caption{}
    \end{minipage}
    \begin{minipage}[t]{0.48\textwidth}
    \centering
    \begin{tikzpicture}[>=latex, scale=.7]
  \draw[->](-1,0)--(7.5,0)node[right]{$x$};
\draw[->](0,-1)--(0,5)node[right]{$y$};
\node at (0,0)[below left]{$O$};
\draw[domain=.7:5, samples=100, very thick]plot(\x, {-.2*(\x-3)^3+2.5});
\draw[->, domain=2.5:6.5, thick, samples=10]plot(\x, {-1.2*\x+7.23});
\draw[->, domain=.5:5, thick, samples=10]plot(\x, {-1.2*\x+5});
\draw(.2,2.5)--(5.5,2.5);
\draw[->, thick](6.025+.5,0) arc (0:129.8:.5)node[above right]{$\theta'$};
\draw[->, thick](4.15+.5,0) arc (0:129.8:.5)node[above right]{$\theta$};


    \end{tikzpicture}
    \caption{}
    \end{minipage}
  \end{figure}



\begin{example}
证明$\tan x>x,\quad 0<x<\frac{\pi}{2}$
\end{example}


\begin{proof}
我们只须证明$\tan x-x>0,\; 0<x<\frac{\pi}{2}$．为此，设$f(x)=\tan x-x,\; x\in\left(0,\frac{\pi}{2}\right)$，因为$f(x)$在$\left[0,\frac{\pi}{2}\right)$
上连续，且在$\left(0,\frac{\pi}{2}\right)$内$f'(x)=\sec^2x-1=\tan^2 x>0$，依定理2, $f(x)$在$\left[0,\frac{\pi}{2}\right)$上递增，所以有
\[f (x) > f (0) =\tan 0-0=0,\quad x\in \left(0,\frac{\pi}{2}\right)\]
即：$\tan x-x>0$，
也就是$\tan x>x,\quad 0<x<\frac{\pi}{2}$.
\end{proof}

\begin{example}
        求证当$x>0$时，有
$e^x> 1+\ln (1+x) $.
\end{example}


\begin{proof}
设$f(x)=e^x-1-\ln(1+x)\; (x>0)$, 则
\[f' (x) =e^x-\frac{1}{1+x},\qquad 
f''(x) =e^x +\frac{1}{(1+x)^2}\] 
    显然，当$x>0$时，$f''(x)>0$, 又$f'(x)$在$[0,\infty)$上连续，由定理2, $f'(x)$在$[0,\infty)$上递增，所以
\[f' (x) > f' (0) =e^0-\frac{1}{1+0}=0\quad  (x> 0) \]
又$f(x)$在$[0, \infty)$连续，同理$f(x)$在$[0, \infty)$上递增，所以当$x>0$时，有
\[f (x) > f (0) =e^0-1-\ln (1+0) =0\]
即：
\[e^x-1-\ln(1+x)>0\quad \text{或}\quad e^x>1+\ln(1+x)\]
\end{proof}


\begin{example}
    确定$f(x)=2x^3-9x^2+12x+5$的递增区间和递减区间，并画出这个函数图象的草图．
\end{example}


\begin{solution}
  \[  f'(x)=6x^2-18x+12=6(x-1)(x-2)\]

  当$x<1$或$x>2$时，$f'(x)>0$, 所以$f(x)$在区间$(-\infty,1)$和$(2,+\infty)$内递增；

  当$1<x<2$时，$f'(x)<0$，所以$f(x)$在区间$(1,2)$内递减．
  
  图象的$y$截距$f(0)=5$，为确定图象的$x$截距，根据例2.54和连续函数中间值定理，由计算的结果
  \[f(-1)=-18,\quad f(1)=10,\quad f(2)=9,\quad f(+\infty)>0\]
  得知$f(x)$在$(1, 2)$和$(2,+\infty)$内没有零点，即图象在$x>1$的范围内与$x$轴不相交：由于$f(x)$在$x<1$的区间内递增，并且$f(-1)\cdot f(0)<0$, 故$f(x)$的图象在
  $x<1$的范围内与$x$轴只有一个交点，它的横坐标在区间$(-1,0)$内．

  把$f(x)$的递增、递减情况列成下表，就更加醒目．

  \begin{center}
\begin{tabular}{c|ccccc}
\hline
$x$ & $(-\infty,1)$& 1& $(1,2)$ & 2& $(2,+\infty)$\\
\hline
$f'(x)$ & $+$ &0&$-$&0&$+$\\
$f(x)$ & $\nearrow$ &10&$\searrow$ &9&$\nearrow$ \\
\hline
\end{tabular}
  \end{center}
图象的草图如图2.27.

\begin{figure}[htp]
    \centering
\begin{tikzpicture}[>=latex, yscale=.3]
\draw[->](-2,0)--(4,0)node[right]{$x$};
\draw[->](0,-2)--(0,15)node[right]{$y$};
\draw[domain=-1:3, samples=300, very thick]plot(\x, {2*\x^3-9*\x^2+12*\x+5});
\draw[dashed](1,0)node[below]{1}--(1,10);
\draw[dashed](2,0)node[below]{2}--(2,9);
\draw(-1,0)node[below]{$-1$}--(-1,.2);
\end{tikzpicture}
    \caption{}
\end{figure}
\end{solution}

\begin{ex}
\begin{enumerate}
    \item     若$|f(x)-f(y)|\le (x-y)^n,\; n>1$，
   求证$f$是一常数．
    \item 若$f'(x)=g'(x)$, $x\in(-\infty,+\infty)$, 求证$f(x)=g (x) +c$.
    \item 证明在任何区间内，函数$x-\sin x$是增函数，又$\alpha$为何
    值，函数$\alpha x-\sin x$是增函数，$\alpha x-\sin x$为减函数？
    \item 证明$\frac{x}{\sin x}$在$\left(0,\frac{\pi}{2}\right]$内是增函数．
    \item 若$x>0$, 求证
\begin{multicols}{2}
  \begin{enumerate}
\item  $\cos x-1+\frac{1}{2}x^2> 0$
\item $\sin x-x+\frac{1}{6}x^3> 0$
\item  $x-\frac{x^2}{2}<\ln (1+x)<x$
\end{enumerate}  
\end{multicols}

\item 证明下列不等式当$0<x<\frac{\pi}{2}$时成立：
\begin{multicols}{2}
    \begin{enumerate}
        \item $\tan x> x+\frac{x^3}{3}$
        \item $\frac{\sin x}{x}> \frac{2}{\pi}$
    \end{enumerate}
\end{multicols}

\item 已知$x+e^x=y+e^y$, 问$\sin x=\sin y$是否为真？
\item 若$a>b$, 求证$a^3-3a^2+3a>b^3-3b^2+3b$.
\item 设$f(0)=0$, 且$f'$在$[0,+\infty)$内递增，证明
$ g(x) =\frac{f (x)}{x}$在$(0,+\infty)$内递增．
\item 
\begin{enumerate}
    \item 已知$a,b$为实数，并且$e<a<b$，其中$e$是自然对数的底，证明：$a^b>b^a$.
    \item 如果正实数$a,b$满足$a^b=b^a$, 且$a<1$, 则$a=b$.
\end{enumerate}

\item 证明：对于所有$x\in\left(0,\frac{\pi}{2}\right]$, 有不等式
$(\sin x)^{-2}\le x^{-2}+1-\frac{4}{\pi^2}$
成立．
\end{enumerate}
\end{ex}


\subsubsection{函数的极大值与极小值}

从例2.59我们看到$x=1$和$x=2$分别是函数$f(x)=2x^3-9x^2+12x+5$取到极大值$f(1)=10$和极小值$f(2)=9$的点，我们称$x=1$为极大点，$x=2$为极小点，以后，极大值与极小值统称为\textbf{极值}，极大点与极小点统称为\textbf{极值点}．

下面我们来讨论如何确定一个函数的极值．
\begin{enumerate}
    \item 若$x_0$是$f(x)$的极值点，且$f(x)$在点$x_0$可微，那么总存在$x_0$的一个邻域，使$f(x)$在点$x_0$的邻域中满足费马条件，因而必有$f'(x_0)=0$. 例如例2.59, $f'(1)=f' (2) =0$.
    \item 若$f(x)$在$x_0$不可微，这时$x_0$也可能是极值点，例如，$f(x)=|x|$, 它在$x_0=0$这一点不可微，但从图2.28即可看出$x_0=0$是它的极小点．
\end{enumerate}

\begin{figure}[htp]
    \centering
    \begin{minipage}[t]{0.48\textwidth}
    \centering
\begin{tikzpicture}[>=latex, scale=1.3]
    \draw[->](-2,0)--(2,0)node[right]{$x$};
\draw[->](0,-.5)--(0,2)node[right]{$y$};
\draw[very thick](-1.8,1.8)--(0,0)node[below left]{$O$}--(1.8,1.8)node[right]{$y=|x|$};
\end{tikzpicture}
    \caption{}
    \end{minipage}
    \begin{minipage}[t]{0.48\textwidth}
    \centering
    \begin{tikzpicture}[>=latex, scale=.35]
        \draw[->](-5,0)--(5,0)node[right]{$x$};
\draw[->](0,-8)--(0,9)node[right]{$y$};
\draw[domain=-2:2, samples=100, very thick]plot(\x, {\x^3});
\node at (0,0)[below right]{$O$};
\foreach \x in {2,4,6,8}
{
    \draw(0,\x)node[left]{$\x$}--(.3,\x);
}
\draw(2,0)node[below]{2}--(2,.3);
    \end{tikzpicture}
    \caption{}
    \end{minipage}
    \end{figure}


这就告诉我们连续函数的极值点应该从它的导函数$f'(x)$的零点和$f(x)$的不可微点中去找，但是，$f'(x)$的零点和$f(x)$的不可微点只是$f(x)$可能达到极值的点，并不一定是极值点，例如$f(x)=x^3$, $x=0$是$f'(x)=3x^2$的零点，然而在整个区间$(-\infty,+\infty)$上，$f'(x)\ge 0$, 而且只在$x=0$时，$f'(x)=0$, 即函数$f(x)=x^3$在$(-\infty,+\infty)$上严格递增，因而$x=0$不是极值点，它的图象如图2.29, 在$y$轴的左侧，图象位于切线$y=0$的下方；而在$y$轴的右侧，图象位于切线$y=0$的上方，我们把方程$f'(x)=0$的根称为函数$f(x)$的\textbf{驻点}．

又如函数
\[g(x)=\begin{cases}
    \frac{1}{2}(x-1), & x\in[1,2]\\
    \frac{3}{2}(x-2)+ \frac{1}{2}(2-1), & x\in[2,3]\\
 (x-3) +\frac{1}{2}(2-1) + \frac{3}{2}(3-2), & x\in[3,4]
\end{cases}\]
的图象是折线，如图2.30, 易知
\[g' (x) > 0,\quad x\in (1, 2) \cup (2, 3) \cup (3, 4) \]
但$g'(2)$和$g'(3)$不存在，由于$g(x)$在$[1, 4]$上严格递增，所以$x=2$, $x=3$都不是极值点．

\begin{figure}[htp]
    \centering
\begin{tikzpicture}[>=latex]
    \draw[->](-.5,0)--(5,0)node[right]{$x$};
    \draw[->](0,-.5)--(0,4)node[right]{$y$};
\draw[domain=1:2, samples=5, very thick]plot(\x, {.5*(\x-1)});
\draw[domain=2:3, samples=5, very thick]plot(\x, {1.5*(\x-2)+.5});
\draw[domain=3:4, samples=5, very thick]plot(\x, {(\x-3)+2});
\foreach \x in {1,2,3,4}
{
    \node at (\x,0)[below]{$\x$};
}
\draw[dashed](2,0)--(2,.5);
\draw[dashed](3,0)--(3,2)node[above left]{$y=g(x)$};
\draw[dashed](4,0)--(4,3);
\node at (0,0)[below left]{$O$};
\end{tikzpicture}
    \caption{}
\end{figure}

综上讨论，我们得到函数$f$取极值的必要条件：

\begin{blk}
    {定理1（极值的必要条件）}若$f(x)$在$[a,b]$上连续，$x_0$是$f(x)$的极值点，那么$x_0$只可能是$f'(x)$的零点或$f(x)$的不可微点．
\end{blk}

下面给出极值点的两个充分性判别法．

\begin{blk}
  {定理2} 如果函数$f$在$x_0$的邻域内连续，并且在$(x_0-\delta,x_0)$和$(x_0,x_0+\delta)\; (\delta>0)$可微：
\begin{itemize}
    \item 在$(x_0-\delta,x_0)$内，$f'(x)<0\; (>0)$;
    \item 在$(x_0,x_0+\delta)$内，$f'(x)>0\; (<0)$;
\end{itemize}
那么，$x_0$是$f$的局部极小点（极大点）．
\end{blk}

\begin{proof}
按照前一节中定理2, $f$在$(x_0-\delta,x_0)$内严格递减（递增），而在$(x_0,x_0+\delta)$内严格递增（递减），故$f(x_0)$必为极小值（极大值）．
\end{proof}

\begin{rmk}
    这个定理并不要求函数$f$在点$x_0$处存在导数，如果导数存在，那么根据费马定理，它等于零．
\end{rmk}


\begin{example}
    讨论$y=\sqrt[3]{(x-1)(x+2)^2}$的极值，并绘出图象的草图．
\end{example}

\begin{solution}
将$y=\sqrt[3]{(x-1)(x+2)^2}$的两边立方，得
\begin{align*}
    y^3&= (x-1) (x+2)^2\\
    3y^2\cdot y'&= (x+2)^2+ (x-1)\cdot 2\cdot  (x+2)\\
    &= (x+2) [x+2+2(x-1)]\\
    &=3x (x+2) 
\end{align*}

$\therefore\quad 
y'=\frac{x (x+2)}{y^2}=\frac{x (x+2)}{\sqrt[3]{(x-1)^2(x+2)^4}}$
即
\[y'=\frac{x}{\sqrt[3]{(x-1)^2(x+2)}}\quad 
(x\ne 1,\; x\ne -2)\]
\begin{itemize}
    \item 当$x=0$时，$y'=0$, 函数图象在点$x=0$处有水平切线；
    \item 当$x=-2$和1时，$y'$不存在，但是由$\Lim_{x\to 1}y'=+\infty$和$\Lim_{x\to -2^+}y'=-\infty$, 而$\Lim_{x\to -2^-}y'=+\infty$得知函数图象在点$x=1$和$x=-2$处有平行$y$轴的切线，所以函数的可能的极值点是$x=-2, 0, 1$这三个点．
\end{itemize}

为确定曲线的$y$截距，令$x=0$, 得$y=-\sqrt[3]{4}\approx -1.59$; 为确定它的$x$截距，令$y=0$, 得$x=-2$和$x=1$, 现列表讨论它的增减性与极值如下：
\begin{center}
\begin{tabular}{c|ccccccc}
\hline
$x$ &$(-\infty,-2)$ & $-2$ & $(-2,0)$ & 0&$(0,1)$& 1& $(1,+\infty)$\\
\hline
$y'$ & $+$ & 不存在 & $-$ &0&$+$&不存在&$+$\\
$y$ & $\nearrow$ &极大值&$\searrow$ &极小值&$\nearrow$ &0&$\nearrow$\\ 
&&0&&$-\sqrt[3]{4}\approx -1.59$\\
\hline
\end{tabular}
\end{center}

函数$y=\sqrt[3]{(x-1)(x+2)^2}$的图象如图2.31.
\begin{figure}[htp]
    \centering
\begin{tikzpicture}[>=latex]
    \draw[->](-4,0)--(3,0)node[right]{$x$};
    \draw[->](0,-3)--(0,3)node[right]{$y$};
\draw[domain=-4:-2.001, samples=100, very thick]plot(\x, {-exp(1/3*ln(1-\x)+2/3*ln(-\x-2))});

\draw[domain=-1.999:.999, samples=100, very thick]plot(\x, {-exp(1/3*ln(1-\x)+2/3*ln(\x+2))});
\draw[domain=1.001:2.25, samples=100, very thick]plot(\x, {exp(1/3*ln(\x-1)+2/3*ln(\x+2))});
\draw[very thick](1,.4)--(1,-.4);

\node at (-2,0) [above]{$-2$};
\node at (1,0) [below right]{$1$};
\draw(2,0)node[below]{$2$}--(2,.1);
\draw(-3,0)node[below]{$-3$}--(-3,.1);
\draw(-1,0)node[below]{$-1$}--(-1,.1);
\draw[dashed](-2,-3)--(-2,3);
\draw[dashed](1,-3)--(1,3);
\end{tikzpicture}
    \caption{}
\end{figure}
\end{solution}

\begin{blk}
{定理3} 设函数$f(x)$在区间$(a,b)$中有二阶连续的导函数，又$c\in (a,b)$, 使得$f' (c) =0$,
\begin{enumerate}
    \item 如果$f''(c)<0$, 则$c$是$f(x)$的一个极大点．
    \item 如果$f''(c)>0$, 则$c$是$f(x)$的一个极小点．
\end{enumerate}
\end{blk}

\begin{proof}
设$f''(c)<0$, 由于$f''(x)$的连续性，故存在点$c$的一个$\delta>0$的邻域$(c-\delta,c+\delta)\subset (a,b)$, 使得$f''(x)＜0$, $x\in (c-\delta,c+\delta)$. 按前节定理2, $f'(x)$在$(c-\delta,c+\delta)$中是递减的，又由于$f'(c)=0$, 因此，
$f'(x)$在点$c$处由正变为负，由前面的定理2, $c$是$f(x)$的极大点．

当$f''(c)>0$时，同理可证$c$是$f(x)$的一个极小点．

为说明定理3的应用，我们求例2.60的二阶导函数：
\begin{align*}
   y''&=\frac{y^2 (2x+2) -x (x+2) \cdot 2y\cdot y'}{y^4}\\
&=\frac{2 [ (x+1) y-x (x+2) y']}{y^3} 
\end{align*}

从例2.60知，当$x=0$时，$y'=0$. 又$f(0)=-\sqrt[3]{4}$. 故在$x=0$处
\[y''_{x=0}=\frac{2\left[(0+1)\left(-\sqrt[3]{4}\right)-0\right]}{-4}=\frac{\sqrt[3]{4}}{2}>0\]
由定理3, $x=0$是函数$y=\sqrt[3]{(x-1)(x+2)^2}$的极小点．
\end{proof}




\begin{example}
    讨论$y=x^3+ax^2+bx+c$的极值．
\end{example}


\begin{solution}
\[y=x^3+ax^2+bx+c,\qquad 
    y'=3x^2+2ax+b\]

$ \Delta=a^2-3b$是导函数的判别式．
    我们可以分$\Delta>0$, $\Delta=0$和$\Delta<0$的三种情形讨论$y'$的正负区间．
\begin{enumerate}
    \item 若$\Delta>0$, 则
\[    y'=3x^2+2ax+b=3 (x-x_1) (x-x_2)\]
    这里$x_1=\frac{1}{3}\left(-a-\sqrt{a^2-3b}\right)$和$x_2=\frac{1}{3}\left(-a+\sqrt{a^2-3b}\right)$
    是导函数$y'$的二个实根．
\begin{enumerate}
    \item 当$x\in (-\infty,x_1)\cup (x_2,+\infty)$时，$y'>0$;
    \item 当$x\in (x_1,x_2)$时，$y'<0$
\end{enumerate}
    所以，$x_1$是$y$的一个极大点，$x_2$是$y$的一个极小点．
    
\item 若$\Delta=0$, 则
    \[y'=3x^2+2ax+\frac{a^2}{3}=3\left(x+\frac{a}{3}\right)^2\]
这里，$y'$虽然在$-\frac{a}{3}$处等于0, 但$y'$在$-\frac{a}{3}$的左、右两侧都取正值，即$y=x^3+ax^2+bx+c$是一个递增函数．因此，$-\frac{a}{3}$不是$y$的极值点．
\item 若$\Delta<0$, 则对于任意$x\in\mathbb{R}$, 恒有\[y'=3x^2+2ax+b=3 \left(x+\frac{a}{3}\right)^2+\frac{3b-a^2}{3}> 0\]
因此，$y=x^3+ax^2+bx+c$是一个增函数，也是没有极值点．
\end{enumerate}    
\end{solution}

\begin{ex}
\begin{enumerate}
    \item 先求出$f'(x)$, 然后由$f'(x)$的正负去讨论$f(x)$的增减与极值，并画出函数图象的草图：
\begin{multicols}{2}
    \begin{enumerate}
        \item $f (x) =2x^3+9x^2-24x+5$
        \item $f (x) =x^4-6x^2-8x+15$
        \item $f (x) =-x^5-5x^4$
        \item $f (x) = (x-1)^3$
     \item  $f (x) =|x^2-1|$
     \item $f (x) = (x^2-1)^3+1$
    \end{enumerate}
\end{multicols}
\item 讨论$f(x)=\frac{x^2-2x-3}{x^2+1}$的增减和极值，并画函数图象的草图．
    \item 求$f(\theta)=\sin\theta-\frac{1}{2}\sin2\theta$在$(-\pi,\pi)$内的极大值和极
    小值．
    \item 
\begin{enumerate}
    \item 求两条曲线$y=\frac{4e^x+2e^{-x}}{3}$和$y=e^{-x}$的交点的坐标；
    \item 求函数$y=\frac{4e^x+2e^{-x}}{3}$的极值点，又它在此点取极大值，还是极小值？
\end{enumerate}

\item 讨论函数$y=x^2e^{-x^2}$的增减与极值．
\item 求$f(x)=x^x\; (x>0)$的极值，并说明此极值是极大值
还是极小值．
\item 求函数$\frac{\ln x}{x}$的极大值，

\item 函数$f(x)=\frac{ax+b}{(x-1)(x-4)}$在点$(2,-1)$有一极值，
求$a,b$并证明此极值是极大值．
\item 证明$f(x)=\left(\alpha-\frac{1}{\alpha}-x\right)(4-3x^2)$仅有一个极大值
和一个极小值，且它们的差是$\frac{4}{9}\left(\alpha+\frac{1}{\alpha}\right)^3$, 又此差值的绝对值的最小值是多少？
\item 函数$f(x)=x(x-1)(x-a)$有绝对值相等，符号相反的
极大值与极小值，试确定常数$a$的值．
\item 求曲线$y=e^{-x}\sin x$上的极值点，并证明极值点上的纵坐
标成等比数列，又该等比数列的公比是多少？  
\end{enumerate} 
\end{ex}


\subsubsection{函数的最大值和最小值}
在实践中有一类“最大”、“最小”、“最省”的问题，例如生产搪瓷杯时，就要考虑在一定容积下，杯子的高和直径取多大时，用料最省．又如把一根圆木锯成矩形截面横梁时，怎样选取矩形的长和宽才能使横梁强度最大？这类“最大”、“最小”、“最省”等问题，在数学上叫做最大值、最小值问题．

我们曾在第一章中指出过“在闭区间$[a,b]$上连续的函数$f(x)$总是具有最大值和最小值”，但是，微积分并未提供确定函数$f(x)$的最大值、最小值的直接方法．如果我们讨论的函数$f(x)$在$(a,b)$上有连续的一阶导数，那么只要求出$f(x)$在$(a,b)$内的驻点，即$f'(x)=0$的根，就可以求出$f(x)$的局部极值．显然，$f(x)$在$[a,b]$上的最大值和最小值一定出现在局部极大值、局部极小值和端点值之中，从几何上来说，最大值、最小值如果不是位于定义区间的端点上，则分别是曲线的波峰和波谷．从图2.33可以看到局部极大值可能小于局部极小值．

\begin{figure}[htp]
    \centering
\begin{tikzpicture}[>=latex]
        \draw[->](-.5,0)--(5,0)node[right]{$x$};
        \draw[->](0,-.5)--(0,4.5)node[right]{$y$};
        
        \node at (0,0) {};
        \draw plot[smooth] coordinates{(.5,4)(1,2.5)(1.9, 3.5)(3.2, .8)(4,1.5)(4.4,1)};
%        \pscurve[linewidth=2pt](.5,4)(1,2.5)(1.9, 3.5)(3.2, .8)(4,1.5)(4.4,1)
%        };
        \draw[dashed](.5,4)--(.5,0)node[below]{$a$};
        \draw[dashed](4.4,1)--(4.4,0)node[below]{$b$};
\end{tikzpicture}
    \caption{}
\end{figure}

下面给出求函数$f(x)$在$[a,b]$上最大值和最小值的法则，这个法则适用于在$[a,b]$上连续，在$(a,b)$内可
微的函数$f(x)$,并且$f'(x)$最多在有限个点上等于零．
\begin{enumerate}
    \item 求$f(x)$在$(a,b)$内的驻点，即解$f'(x)=0$;
    \item 计算$f(x)$在驻点和端点的函数值，并把这些值加以比较，其中最大的一个为最大值，最小的一个为最小值．
\end{enumerate}

假如$f'(x)$在$(a,b)$内有一个驻点$c$, 且在整个$(a,b)$
上，$f''(x)<0\; (>0)$, 那么点$c$是最大值（最小值）点．

事实上，如果$f''(x)<0,\; x\in (a,b)$, 那么函数$f'(x)$在$[a,b]$上递减，因而$c$是其唯一的驻点，由中值定理知：

当$a\le x<c$时，$f'(x)-f'(c)=(x-c)\cdot f''(\xi)$，这里$x<\xi<c$．

$\because\quad f' (c) =0,\quad f'' (\xi)<0,\quad x-c<0$,

$\therefore\quad f'(x)>0$

同理说明当$c<x\le b$时，$f'(x)<0$．再根据$f(x)$在$[a,c]$上递增而在$[c,b]$上递减．故当$x\ne c$时有$f(x)<f(c)$．由此可知$c$是最大值点．当$f''(x)>0$时，也可以进行同样的论证．

如果$f(x)$在$[a,b]$内有不可微的点，而且这种点只有有限个，那么求最大值和最小值的法则改为求出三种点的对应函数值再加以比较，即$f(x)$的驻点的函数值，$f(x)$的不可微点的函数值，最后还有端点的函数值$f(a)$和$f(b)$, 这些值中最大的就是最大值，最小的就是最小值．

在实际问题中，如果在$(a,b)$内部，$f(x)$只有一个驻点$c$, 而且从实际问题本身又可以知道在$(a,b)$内必定有最大值或最小值，那么$f(c)$就是所要求的最大值或最小值，不需要算出$f(a)$, $f(b)$了．

\begin{example}
求函数$f(x)=x^{\tfrac{2}{3}}-(x^2-1)^{\tfrac{1}{3}}$在区间$[-2,2]$上的最大值和最小值．
\end{example}


\begin{solution}
\begin{align*}
f'(x)&=\frac{2}{3}x^{-\tfrac{1}{3}}-\frac{1}{3}(x^2-1)^{-\tfrac{2}{3}}\cdot 2x\\
&=\frac{2}{3} \frac{(x^2-1)^{\tfrac{2}{3}}-x^{\tfrac{4}{3}}}{x^{\tfrac{1}{3}}(x^2-1)^{\tfrac{2}{3}}}    
\end{align*}
    故$f(x)$的驻点是$\pm\frac{1}{\sqrt{2}}$, 不可微点是$x=0$, $x=\pm 1$, 且这些点都在$[-2, 2]$内，因此$f(x)$的极值点可能是
\[-1,\quad \frac{-1}{\sqrt{2}},\quad  0,\quad \frac{1}{\sqrt{2}},\quad  1\]
    由于$f(x)$是偶函数，故仅需计算
\[    f (0) =1,\qquad f\left(\frac{1}{\sqrt{2}}\right) =\sqrt[3]{4},\qquad f (1) =1\]
    又在端点$x=2$处的函数值$f(2)=\sqrt[3]{4}-\sqrt[3]{3}$, 把上面这些函数值相互比较，可见$f(x)$在$[-2, 2]$中的最大值为$\sqrt[3]{4}\approx 1.59$, 最小值为$\sqrt[3]{4}-\sqrt[3]{3}\approx 1.59-1. 44\approx 0.15$.
\end{solution}

\begin{example}
    设$a>0$, 证明
$    f (x) =\frac{1}{1+|x|}+\frac{1}{1+|x-a|}$
    的最大值是$\frac{2+a}{1+a}$.
\end{example}    

\begin{proof}
把$f(x)$写成分段函数
\[f(x)=\begin{cases}
    \frac{1}{1-x}+\frac{1}{1+a-x}, & x<0\\
    \frac{1}{1+x}+\frac{1}{1+a-x}, & 0\le x\le a\\
    \frac{1}{1+x}+\frac{1}{1+x-a}, & x>a
\end{cases}\]
于是
\[f'(x)=\begin{cases}
    \frac{1}{(1-x)^2}+\frac{1}{(1+a-x)^2}, & x<0\\
    \frac{-1}{(1+x)^2}+\frac{1}{(1+a-x)^2}, & 0\le x\le a\\
    \frac{-1}{(1+x)^2}-\frac{1}{(1+x-a)^2}, & x>a
\end{cases}\]
因而$f(x)$在$(-\infty, 0]$上递增，而在$[a,+\infty)$上递减，由此得知$f(x)$在$[0,a]$上的最大值也就是在$\mathbb{R}$上的最大值．

设对于$(0,a)$内的$x$有$f'(x)=0$, 则
\[(1+x)^2- (1+a-x)^2=0\]
上式唯一解是$x=\frac{a}{2}$. 
因为$f\left(\frac{a}{2}\right)=\frac{4}{2+a}$, $f(0)=f(a)=\frac{2+a}{1+a}$，
并且
\[\frac{2+a}{1+a}-\frac{4}{2+a}= \frac{a^2}{(1+a)(2+a)}>0\quad  (a>0)\]
所以$f(x)$的最大值是$\frac{2+a}{1+a}$.
\end{proof}

\begin{example}
    用铁皮做体积为$V$的圆柱形封闭容器，如何使材料最省？
\end{example}


\begin{solution}
\textbf{解法1:} 设圆柱形底面半径为$r$, 高为$h$, 于是需用材料即圆柱表面积
\[S=2\pi rh+2\pi r^2\]
但$V=\pi r^2h$为已知，因此，$h=\frac{V}{\pi r^2}$，
从而问题变为求函数$S(r)=\frac{2V}{r}+2\pi r^2\; (0<r<+\infty)$的最小值．

根据经验，当$r$很小时，$S$很大，又当$r$很大时，$S$也很大，因此在$(0,+\infty)$内必有$r$的一个值使$S$最小，即它是$S'(r)=0$的一个根，由
\[S' (r) =-\frac{2V}{r^2}+4\pi r=0\]
得唯一正根$r=\sqrt[3]{\frac{V}{2\pi}}$．
从这个实际问题知道，当$r=\sqrt[3]{\frac{V}{2\pi}}$
时，$S$为最小，这时相应的高为
\[h=\frac{V}{\pi r^2}=\frac{2\pi r^3}{\pi r^2}=2r=2\sqrt[3]{\frac{V}{2\pi}}\]
这说明圆柱的高和直径相等时，用料最省．

\textbf{解法2:}  我们可以不先将下面两个式子中的$h$消去，再求$\frac{\dd S}{\dd r}=0$的根
\begin{align}
    S&=2\pi r h+2\pi r^2\\
    \pi r^2 h&=V
\end{align}
而是保留辅助变量$h$，将
(2.27), (2.28)对$r$求导数，只须记着$V$是常量，$h$是$r$的函数，于是由$\frac{\dd S}{\dd r}=0$，得
\begin{equation}
    \frac{\dd S}{\dd r}=2\pi\left(h+r\frac{\dd h}{\dd r}+2r\right)=0
\end{equation}
由(2.28)得
\begin{equation}
    \pi\left(2rh+r^2\frac{\dd h}{\dd r}\right)=0
\end{equation}
由此得：
\[\frac{\dd h}{\dd r}=-\frac{2h}{r}\]
将它代入(2.29)，得
\[2\pi\left[h+r\left(-\frac{2h}{r}\right)+2r\right]=0\]
所以：$2r=h$. 

同样地说明了圆柱的高和直径相等时，用料最省．
\end{solution}


\begin{example}
由半径为$\ell$的圆铁皮，剪去一个扇形，把剩下的部分围成一个圆锥形容器，为了使所得圆锥容器有最大的容积，问剪去扇形的中心角应该多大？
\end{example}

\begin{figure}[htp]
    \centering
\begin{tikzpicture}[>=latex]
\begin{scope}[scale=1.5]
\draw[pattern=north east lines](0,0)--(-135:1.5) arc (-135:135:1.5)--(0,0);
\draw[<->](-135:.2) arc (-135:135:.2);
\node at (.5,0)[fill=white]{$x$};
\node at (1,1)[above right]{$x\ell=2\pi r$};
\tkzDefPoint(-30:1.5){A}
\tkzDefPoint(-30:2){B}
\tkzDrawCircle(B,A)
\draw[->](B)--node[above]{$r$}+(10:.5);
\draw[decorate,decoration={brace,mirror,raise=2pt}](135:1.5)--node[below left]{$\ell$}(0,0);


\end{scope}
\begin{scope}[xshift=7cm, scale=.7]
\draw[dashed](0,-2) ellipse [x radius=1.5cm, y radius=.5cm];
\draw[thick](-1.5,-2)--(0,3)--node[above right]{$\ell$}(1.5,-2);
\draw[dashed](0,3)--node[left]{$h$}(0,-2)--node[below]{$r$}(1.5,-2);
\draw[thick](1.5,-2) arc [start angle =0, end angle=-180, x radius=1.5cm, y radius=.5cm];
\end{scope}


\end{tikzpicture}
    \caption{}
\end{figure}


\begin{solution}
    设剪去扇形后，所剩材料的圆心角为$x$（弧度），扇形的弧长为$x\ell$, 由它围成一个圆锥，其底面半径为$r$, 母线长为$\ell$, 高为$h$, 如图2.34. 于是
\[r=\frac{\ell x}{2\pi}\]
圆锥高为
\[h=\sqrt{\ell^2-r^2}=\sqrt{\ell^2-\left(\frac{\ell x}{2\pi}\right)^2}=\frac{\ell }{2\pi}\sqrt{4\pi^2-x^2}\]
因此圆锥的体积为
\begin{align*}
    V&=\frac{1}{3}\pi r^2 h\\
    &=\frac{1}{3}\pi\left(\frac{\ell x}{2\pi}\right)^2\cdot \frac{\ell}{2\pi}\sqrt{4\pi^2-x^2}\\
    &=\frac{\ell^3 x^2}{24\pi^2}\sqrt{4\pi^2-x^2}\quad (0\le x\le 2\pi)
\end{align*}
由于$V(0)=0$, $V(2\pi)=0$, 故在$(0, 2\pi)$内必有一个值$x_0$使$V$的值最大，因而使$V'(x_0)=0$. 由于
\begin{align*}
    V'(x)&=\frac{\ell^3}{24\pi^2}\left[2x\sqrt{4\pi^2-x^2}+\frac{1}{2}x^2(4\pi^2-x^2)^{-\tfrac{1}{2}}\cdot (-2x)\right]\\
    &=\frac{\ell^3}{24\pi^2}\left(\frac{8\pi^2 x-3x^3}{\sqrt{4\pi^2-x^2}}\right)
\end{align*}
令$V'(x)=0$, 即$8\pi^2 x-3x^3=0$, 解得
\[x_1=0,\quad x_2=-2\pi\sqrt{\frac{2}{3}}, \quad x_3=2\pi\sqrt{\frac{2}{3}}\]
这里只有$x_3=2\pi\sqrt{\frac{2}{3}}\in (0, 2\pi)$, 故知$x=2\pi\sqrt{\frac{2}{3}}$是$V(x)$在$[0, 2\pi]$内的唯一的极大点，因而也是最大点．于是最后得到结论，剪去扇形的圆心角为
\[2\pi-2\pi\sqrt{\frac{2}{3}}=\frac{2}{3}\pi\left(3-\sqrt{6}\right)\; \text{（弧度）}\]
时，才能使所围成的容器有最大的容积．
\end{solution}

\begin{example}
    求证：若$p>q>1$，又$x\ge 0$且$x\ne 1$，则$\frac{x^p-1}{p}>\frac{x^q-1}{q}$.
\end{example}

\begin{proof}
设$f(x)=\frac{x^p-1}{p}-\frac{x^q-1}{q}\; (x\ge 0)$，则：
\[f'(x)=x^{p-1}-x^{q-1}=x^{q-1}\left(x^{p-q}-1\right)\quad (p>q>1)\]
    函数的增、减如下表所示：
\begin{center}
\begin{tabular}{c|cccc}
\hline
$x$&0&$(0,1)$&1& $(1,+\infty)$\\
\hline
$f'(x)$&0&$-$&0&$+$\\
$f(x)$&$\frac{1}{q}-\frac{1}{p}$&$\searrow$&极小值0& $\nearrow$\\ 
\hline
    \end{tabular}
\end{center}

$\because\quad f(0)=\frac{1}{q}-\frac{1}{p}>0=f(1)$

$\therefore\quad f(1)=0$不仅是极小值也是最小值．于是：
\[f(x)=\frac{x^p-1}{p}-\frac{x^q-1}{q}\ge 0\quad (x\ge 0)\]
因为$x\ne 1$，所以
\[\frac{x^p-1}{p}>\frac{x^q-1}{q}\quad (x\ge 0,\; x\ne 1)\]
\end{proof}

\begin{example}
 （光的折射定律）设有两种介质，$X'X$是它们的交界线，光由一种介质中的某点$A$走向另一种介质中的某点$B$. 又设当光走过交界线$X'X$时，它的速度由$v_1$变成$v_2$. 依据物理学中的一个基本原理“光由$A$到达$B$时，所采取的途径要使用的时间为最小”，如果两种介质是均匀的（例如一为空气，一为水），那么$v_1,v_2$都是常数，从而光在这两种介质中所走的路径是两段直线$AP$, $PB$, $P$表示光达到交界线时与这条线的交点，如果$AB$不与$X'X$垂直，一般$APB$不是直线，我们来定出$AP$, $BP$与$X'X$的法线的夹角$\theta_1$和$\theta_2$的关系．
\end{example}

    \begin{figure}[htp]
    \centering
\begin{tikzpicture}[>=latex, scale=.7]
\draw[->](-1,0)node[left]{$X'$}--(6,0)node[right]{$X$};
\draw[->](0,-4)node[right]{$S'$}--(0,4)node[right]{$S$};
\node at (0,0)[below left]{$O$};
\draw[->, very thick](0,3)node[left]{$A(0,a)$}--(2,0)node[above left]{$P$};
\draw[->, very thick](2,0)--(4.2,-4)node[right]{$B(c,b)$};
\draw[dashed](2,-2)--(2,2);
\draw[<->](0,-.4)--node[fill=white]{$x$}(2,-.4);
\draw[|<->](4.2,.4)--node[fill=white]{$c-x$}(2,.4);
\draw[dashed](4.2,-4)--(4.2,0);
\draw(2,.75) arc (90:90+33.7:.75)node[above]{$\theta_1$};
\draw(2,-.75) arc (-90:-90+28.8:.75)node[below]{$\theta_2$};


\end{tikzpicture}
    \caption{}
\end{figure}

\begin{solution}
如图2.35建立坐标系，$A$点坐标是$(0,a)$, $B$点的坐标是$(c,b)$, 设$P$点的坐标是$(x,0)$, 光由$A$到$P$所需时间是$\frac{\sqrt{a^2+x^2}}{v_1}$，从$P$到$B$所需时间是$\frac{\sqrt{b^2+(c-x)^2}}{v_2}$, 
则整个时间$t(x)$是
\[t (x) =\frac{\sqrt{a^2+x^2}}{v_1}+\frac{\sqrt{b^2+(c-x)^2}}{v_2}\]

为了使$t(x)$极小，应使$t'(x)=0,\; x\in (0,c)$. 由$t(x)$得
\begin{align*}
    t' (x) &=\frac{x}{v_1\sqrt{a^2+x^2}}-\frac{c-x}{v_2\sqrt{b^2+(c-x)^2}}\\
    t'' (x) &=\frac{-x^2}{v_1(a^2+x^2)^{\tfrac{3}{2}}}+\frac{1}{v_1\sqrt{a^2+x^2}}+\frac{1}{v_2\sqrt{b^2+(c-x)^2}}-\frac{(c-x)^2}{v_2\left[b^2+(c-x)^2\right]^{\tfrac{3}{2}}}\\
    &=\frac{a^2}{v_1(a^2+x^2)^{\tfrac{3}{2}}}+\frac{b^2}{v_2\left[b^2+(c-x)^2\right]^{\tfrac{3}{2}}}>0.
\end{align*}
既然$t'(0)<0$, $t'(c)>0$, 及$t'$是严格递增的，故在$(0,c)$内有唯一的$x_0$, 使得
\[t'(x_0)=0\]
从而当$x=x_0$时，$t(x_0)$是极小值也是最小值，于是由$t'(x_0)=0$得到
\[\frac{x_0}{v_1\sqrt{a^2+x^2_0}}=\frac{c-x_0}{v_2\sqrt{b^2+(c-x_0)^2}}\]
即
\[\frac{\sin\theta_1}{v_1}=\frac{\sin\theta_2}{v_2}\]
这就是光的折射定律．
\end{solution}

\begin{ex}
\begin{enumerate}
    \item 由$y=0$, $x=8$, $y=x^2$围成一曲边三角形$OAB$, 在曲边$\wideparen{OB}$上，求一点使得过此点所作$y=x^2$的切线与$OA$、$AB$所围成的三角形面积最大．
    \item 椭圆$\frac{x^2}{a^2}+\frac{y^2}{b^2}=1$的内接矩形中怎样的矩形面积最大？又此最大的矩形面积比椭圆的内接正方形的面积大多少？
    \item 设$\triangle ABC$是定圆的内接三角形，其中$A$、$B$两点固定，$C$点在圆周上变动，试求其面积最大时的位置．
    \item 要做一个底面为长方形的带盖箱子，其体积为72立方厘米，其底边成$1:2$的关系，问各边的长怎样，才能使表面积最小？
    \item 设有底为等边三角形的直棱柱，体积为$V$, 要使其总面积为最小，问底边长应为多少？
    \item 设正三棱锥的侧面积为$S$, 问当侧面与底面所成角$\alpha$为何值时，其体积最大？
    \item 平面上通过一个已知点$P(1, 4)$引一条直线，要使它在两个坐标轴上的截距都为正，且它们的和为最小，求这条直线的方程．
    \item 直圆柱内接于半径为$R$的球中，当圆柱的高为何值时，其
    体积最大．
    \item 三个点$A$, $B$和$C$不在同一直线上，$\angle ABC=60^{\circ}$, 汽车以80公里/小时的速度由$A$向$B$行驶，同时火车以50公里/小时的速度由$B$向$C$行驶．若$AB=200$公里，问运动开始几小时后汽车与火车的距离最小？
    \item 船航行一昼夜的耗费由两部分组成：一为固定部分，设为$a$元；另一为变动部分，设它与速度立方成比例，试问船应以怎样的速度$v$行驶为最经济？
    \item 在某种产品的制造过程中，次品率$y$依赖于日产量$x$, 即$y=y(x)$, 已知
\[y(x)=\begin{cases}
    \frac{1}{101-x},& 0\le x\le 100,\\
    1,& x>100
\end{cases}\]
其中$x$为正整数．又该厂每生产一件产品可盈利$A$元，但每生产出一件次品就要损失$\frac{A}{3}$元．问为了获得最大盈利，该厂的日产量应定为多少？
\item 抛物线$y=x^2$上的哪一点到直线$x-y-2=0$的距离最短？\item 一弹在空中飞行（不计空气阻力），其弹道方程为
\[y=mx-\frac{(m^2+1)x^2}{800}\]
这里原点取在此弹发射之点，其中$m$为曲线在原点处的切线斜率．
\begin{enumerate}
    \item 若要此弹击中同一水平面上最远距离的目标；
    \item 若要此弹击中300米远处一直立墙壁上的最大高度，问$m$之值各为多少？
\end{enumerate}
\item 若直角三角形的周长$2p$一定，求它的面积的最大值．
\end{enumerate}  
\end{ex}

\subsection{函数的凸性与拐点}
上面已经讨论了函数的增减性与极值，这对于我们了解函数的性质与图象形状有很大的帮助．现在我们还要进一步考察函数在什么区间是向上凸的，什么区间是向下凸的，以及上凸函数和下凸函数的一些性质．

\begin{blk}
    {定义1} 函数$f$在区间$(a,b)$内，称为向下（上）凸的，指对于任意$c\in (a,b)$, 曲线位于在点$(c,f(c))$处的切线上（下）面（见图2.37和2.38）.
\end{blk}    


\begin{figure}[htp]
    \centering
    \begin{minipage}[t]{0.48\textwidth}
    \centering
\begin{tikzpicture}[>=latex, scale=.9]
\draw[->](-.5,0)--(5,0)node[right]{$x$};
\draw[->](0,-.5)--(0,4)node[right]{$y$};
\draw[domain=1:4, samples=100, very thick]plot(\x,{0.5*(\x-2)^2+1})node[right]{$y=f(x)$};
\draw[dashed](1,1.5)--(1,0)node[below]{$a$};
\draw[dashed](4,3)--(4,0)node[below]{$b$};
\draw[domain=.5:5, samples=10,thick]plot(\x,{0.5*\x-0.125});
\node at (0,0)[below left]{$O$};
\node at (2.5,1.125)[below]{$(c,f(c))$};
\end{tikzpicture}
    \caption{}
    \end{minipage}
    \begin{minipage}[t]{0.48\textwidth}
    \centering
    \begin{tikzpicture}[>=latex, scale=.9]
\draw[->](-.5,0)--(5,0)node[right]{$x$};
\draw[->](0,-.5)--(0,4)node[right]{$y$};
\draw[domain=1.5:4, samples=100, very thick]plot(\x,{-0.5*(\x-3)^2+2})node[right]{$y=f(x)$};
\draw[dashed](1.5,.875)--(1.5,0)node[below]{$a$};
\draw[dashed](4,1.5)--(4,0)node[below]{$b$};
\node at (2.5,1.875)[above]{$(c,f(c))$};
\draw[domain=.5:5, samples=10,thick]plot(\x,{0.5*\x+0.625});
\node at (0,0)[below left]{$O$};
    \end{tikzpicture}
    \caption{}
    \end{minipage}
    \end{figure}


\begin{blk}
    {定义2} 函数$f$的图象若经过某点$x=c$时，它的图象
$y=f(x)$由向上凸变为向下凸，或者由向下凸变为向上凸，则点$x=c$称为函数$f$的\textbf{拐点}，如图2.39和2.40.
\end{blk}    
    

\begin{figure}[htp]
    \centering
    \begin{minipage}[t]{0.48\textwidth}
    \centering
\begin{tikzpicture}[>=latex, yscale=.8]
    \draw[->](0,.5)--(5,.5)node[right]{$x$};
    \draw[->](1,0)--(1,4)node[right]{$y$};
    \draw[domain=.5*pi:1.5*pi, samples=100, very thick]plot(\x,{sin(\x r)+2});
    \node at (1,.5)[below left]{$O$};
\node at (3.1416, 2)[above right]{$(c,f(c))$};
\draw[domain=2:4.4, samples=10, thick]plot(\x,{-\x+5.1416});
\draw[dashed](3.1416,2)--(3.1416,.5)node[below]{$c$};
\end{tikzpicture}
    \caption{}
    \end{minipage}
    \begin{minipage}[t]{0.48\textwidth}
    \centering
    \begin{tikzpicture}[>=latex, yscale=.8]
   \draw[->](-3,.5)--(2,.5)node[right]{$x$};
    \draw[->](-2,0)--(-2,4)node[right]{$y$};
    \draw[domain=-.5*pi:.5*pi, samples=100, very thick]plot(\x,{sin(\x r)+2});
    \node at (-2,.5)[below left]{$O$};
\node at (0,2)[above left]{$(c,f(c))$};
\draw[domain=-1.25:1.25, samples=10, thick]plot(\x,{\x+2});   
\draw[dashed](0,2)--(0,.5)node[below]{$c$};

    \end{tikzpicture}
    \caption{}
    \end{minipage}
    \end{figure}


这时过$x=c$的点的切线，把曲线分成两部分，换言之，存在充分小的邻域$\delta>0$, 使得对于一切$x\in (c-\delta,c)$, 曲线位于在点$c$处的切线的一方，而对于$x\in (c,c+\delta)$, 曲线在该切线的另一方．

\begin{blk}
  {凸性判定定理} 设函数$f$在区间$(a,b)$内有二阶导数，则
\begin{enumerate}
    \item 函数f在区间$(a,b)$内为向下凸的充分必要条件是，对于每个$x\in (a,b)$, 有$f''(x)\ge 0$;
    \item 函数f在区间$(a,b)$内为向上凸的充分必要条件是，对于每个$x\in (a,b)$, 有$f''(x)\le 0$.
\end{enumerate}
\end{blk}

\begin{proof}
    假设曲线$y=f(x)$在$(a,b)$内是向下凸的，因为函数$f$在$(a,b)$内有二阶导数，则对于任意一点$c\in(a,b)$, $f'(c)$存在，曲线$y=f(x)$在点$x=c$处的切线方程是
\begin{equation}
    y=f (c) +f' (c) (x-c)
\end{equation}
    对于每个点$x\in (a,b)$, 曲线$y=f(x)$的对应点的纵坐
标与切线上的对应点的纵坐标之差用$g(x)$表示，如图2.41, 于是
\begin{equation}
   g (x) =f (x) -[f (c) +f' (c) (x-c)]=[f(x)-f (c) ] - (x-c) f' (c)  
\end{equation}
根据中值定理，得到
\[g (x) = (x-c) f' (d) - (x-c) f'(c)\]
这里$d\in (x,c)$或$(c,x)$. 又$g(x)=(x-c)\cdot [f'(d)-f'(c)]$, 再根据中值定理，得到
\begin{equation}
    g (x) = (x-c) (d-c)f''(e)
\end{equation}
这里$e\in (d,c)$或$(c,d)$.

于是对于一切$x\in (a,b)$, 点$(x,f(x))$不在切线(2.31)的下方的充要条件是$g(x)\ge 0$, 显然只有点$c$使$g(c)=0$. 根据等式(2.33), 并且注意到
\[(x-c)(d-c)>0\]
便知道$g(x)\ge 0$的充分必要条件是$f''(e)\ge 0$. 这里$e\in(a,b)$且随着$c$和$x$的不同值而改变．因此，依曲线向下凸的定义，曲线向下凸的充分必要条件是对于任意$x\in (a,b)$有$f''(x)\ge 0$.

用同样的方法便可说明曲线$y=f(x)$向上凸的充分必要条件是$f''(x)\le 0$.
\end{proof}

\begin{figure}[htp]
    \centering
    \begin{minipage}[t]{0.48\textwidth}
    \centering
\begin{tikzpicture}[>=latex, scale=1]
    \draw[->](-.5,0)--(5,0)node[right]{$x$};
\draw[->](0,-.5)--(0,4)node[right]{$y$};
\draw[domain=.75:3.1, samples=100, very thick]plot(\x, {-1/(\x-3.5)+1})node[right]{$y=f(x)$};
\draw[domain=.75:3, samples=10, thick]plot(\x,{\x-0.5});
\draw[dashed] (2.5,2)node[left]{$(c,f(c))$}--(2.5,0)node[below]{$c$};
\draw[dashed] (1,1.4)--(1,0)node[below]{$x$};
\draw[|<->|](.8,1.4)--node[left]{$g(x)$}(.8,.5);
\node at (0,0) [below left]{$O$};
\end{tikzpicture}
    \caption{}
    \end{minipage}
    \begin{minipage}[t]{0.48\textwidth}
    \centering
    \begin{tikzpicture}[>=latex, scale=1]
       \draw[->](-.5,0)--(5,0)node[right]{$x$};
\draw[->](0,-.5)--(0,4)node[right]{$y$};
\draw[domain=.25:4, samples=100, very thick]plot(\x,{0.4*(\x-1.5)^2+1})node[right]{$y=f(x)$};
\draw[domain=.25:4, samples=10, thick]plot(\x,{.6*\x+.5})node[right]{$y=\ell(x)$};

\draw[dashed] (1,1.1)--(1,0)node[below]{$a$};
\draw[dashed] (3.5,2.6)--(3.5,0)node[below]{$b$};
\draw(2,0)node[below]{$x$}--(2,1.7)node[above left]{$(x,\ell(x))$};
\draw[<->|](2.2,0)--node[right]{$(x,f(x))$}(2.2,1.1);
\node at (0,0)[below left]{$O$};
    \end{tikzpicture}
    \caption{}
    \end{minipage}
    \end{figure}

向下凸或向上凸的函数的重要性质如下面的定理所述：

\begin{blk}
    {定理} 设$f$是$[a,b]$上二次可微函数，并且$f''(x)>0$, $x\in (a,b)$, 那么对于每一个满足$a<x<b$的$x$，有
\[f(x)<f(b)\cdot \frac{x-a}{b-a}+f(a)\cdot \frac{b-x}{b-a}\]
\end{blk}

这一定理有一个明显的几何意义（图2.42），用$\ell(x)$表示右端的一次函数：
\[\ell(x)=f(b)\cdot \frac{x-a}{b-a}+f(a)\cdot \frac{b-x}{b-a}\]
    它在$x=a$和$x=b$的值是$\ell(a)=f(a)$和$\ell(b)=f(b)$, 即与$f(x)$在这些点的值一致．因为$y=\ell(x)$的图象是$y=f(x)$在$[a,b]$上的割线，因此$[a,b]$上的向下凸函数$f$的图象$y=f(x)$位于该区间的割线的下面．





\begin{proof}
设$F(x)=f(x)-f(b)\frac{x-a}{b-a}-f(a)\frac{b-x}{b-a}$，
这里$x\in [a,b]$. 我们需要证明在$[a,b]$上，$F(x)\le 0$. 因为$F(x)$在$[a,b]$上连续，所以$F(x)$在$[a,b]$的某点达到它的最大值，用$c$表示最大值点，我们证明$c$是两端点之一．

假设$a<c<b$，那么$c$也是$F(x)$的极大值点．因此$F'(c)=0$. 又
\begin{align*}
    F'(x)&=f'(x)-\frac{f(b)}{b-a}+\frac{f(a)}{b-a}\\
    F''(x)&=f''(x)>0,\qquad x\in(a,b)
\end{align*}
因此我们得出结论，$F(c)$是$[a,b]$上的最小值，即对于
任意$x\in [a,b]$且$x\ne c$, 有
\[F (x) > F (c) =f (c) -\ell (c)\]
这与$c$是极大值点矛盾，因此，唯一的可能是$c$是一个端点$a$或$b$, 但是$f(a)=\ell(a)$, $f(b)=\ell(b)$, 这表明$F(x)=f(x)-\ell(x)$的最大值是0, 因此在$[a,b]$上除端点以外的所有其它点$x$, 有
\[F (x) =f (x) -f (b) \frac{x-a}{b-a}-f (a) \frac{b-x}{b-a}<0\]
即
\[f (x)<f (b) \frac{x-a}{b-a}+f (a) \frac{b-x}{b-a}\]
这就完成了向下凸函数定理的证明．
\end{proof}

对于向上凸函数有类似的定理，即：区间$[a,b]$上向上凸的函数的图象位于$[a,b]$的割线上面，用式子表示就是
\[f (x)>f (b) \frac{x-a}{b-a}+f (a) \frac{b-x}{b-a},\quad x\in(a,b)\]

\begin{example}
    求证 $\cos\pi x>1-2x,\quad x\in\left(0, \frac{1}{2}\right)$.
\end{example}


\begin{proof}
    设$f(x)=\cos\pi x$, $\ell (x)=1-2x$, 因为
\[f (0) =1=\ell (0) ,\qquad f \left(\frac{1}{2}\right) =0=\ell\left(\frac{1}{2}\right)\]
所以：$y=1-2x$是曲线$y=\cos\pi x$在$\left[0,\frac{1}{2}\right]$上的割线．

又$f'(x)=-\pi \sin\pi x$, $f''(x)=-\pi^2\cos\pi x$, 且当$x\in \left(0,\frac{1}{2}\right)$时，有$f''(x)<0$, 即$f(x)$在$\left(0,\frac{1}{2}\right)$内是向上凸
的．因此，$y=\cos\pi x$在$\left[0,\frac{1}{2}\right]$上的图象是在该区间的割
线$y=1-2x$的上面，从而
\[\cos\pi x>1-2x,\qquad x\in \left(0,\frac{1}{2}\right)\]
\end{proof}

下面给出拐点满足的必要条件．

\begin{blk}
    {定理} 如果$x=c$是曲线$y=f(x)$的拐点，并且$f''(c)$存在，则$f''(c)=0$.
\end{blk}

\begin{proof}
    因为$x=c$是曲线的拐点，即它是曲线$y=f(x)$向上凸和向下凸的分界点，依凸性判定定理，我们可以取$a<c<b$, 使得$f''(x)$在$(a,c)$内的函数值的符号与在$(c,b)$内的函数值的符号相反．

设$g(x)=f'(x),\; x\in (a,b)$, 于是$g'(x)=f''(x)$经过$c$点由负变正或由正变负，因此点$c$是$g(x)$在$(a,b)$内的唯一的极值点，又$g'(c)=f''(c)$存在，根据费马定理
\[g' (c) =f'' (c) =0\]
\end{proof}


\begin{example}
    讨论函数$f(x)=x^3$和$f(x)=x^4$有无拐点．
\end{example}


\begin{solution}
    $f(x)=x^3$有导数$f'(x)=3x^2$, $f''(x)=6x$, 显然$f'(0)=f''(0)=0$, 而且：
当$x>0$时，$f''(x)>0$; 当$x<0$时，$f''(x)<0$, 即曲线$y=x^3$在$(-\infty, 0)$内是向上凸的，而在$(0,+\infty)$区间内是向下凸的．因此，$x=0$是曲线的拐点而不是极值点．

又$f(x)=x^4$有导数$f'(x)=4x^3$, $f''(x)=12x^2$, 显然$f'(0)=f''(0)=0$, 并且对于每一个$x\ne 0$都有$f''(x)>0$, 因此这个曲线$y=x^4$是处处向下凸的，$x=0$不是曲线$y=x^4$的拐点而是一个极小点．
\end{solution}


下面给出判定拐点的充分条件．
    
\begin{blk}
   {定理（拐点的充分条件）} 设$f$在点$c$处连续，在点$c$附近（可以不包括$c$点）具有二阶导数，如果自变量$x$通过$c$时，二阶导数变号，则$x=c$是曲线$y=f(x)$的拐点．
\end{blk}

\begin{proof}
    我们可以取$c$点的一个邻域$(a,b)$使$f''(x)$在
$(a,c)$上的函数值的符号和在$(c,b)$上的相反，依凸性判定定理，$f$经过$c$时将由向上凸变为向下凸或者由向下凸变为向上凸，即$x=c$是曲线$y=f(x)$的拐点．
\end{proof}


\begin{example}
说明$x=0$是曲线$y=x^{\tfrac{1}{3}}$的拐点．
\end{example}

\begin{proof}
    设$f(x)=x^{\tfrac{1}{3}}$，则
\[f'(x)=\frac{1}{3x^{\tfrac{2}{3}}},\qquad f''(x)=\frac{-2}{9x^{\tfrac{5}{3}}}\]
显然$f'(0)$和$f''(0)$都不存在，但是当$x<0$时，$f''(x)>0$, 即曲线$y=x^{\tfrac{1}{3}}$向下凸；当$x>0$时，$f''(x)<0$, 即曲线$y=x^{\tfrac{1}{3}}$向上凸，因此点$x=0$是曲线的拐点．

由$\Lim_{x\to 0}f'(x)=\Lim_{x\to 0}\frac{1}{3x^{\tfrac{2}{3}}}=+\infty$知，曲线在点$x=0$处的切线就是$y$轴本身（见图2.43）．
\end{proof}

\begin{figure}[htp]
    \centering
    \begin{minipage}[t]{0.48\textwidth}
    \centering
\begin{tikzpicture}[>=latex, scale=1]
    \draw[->](-2.5,0)--(2.5,0)node[right]{$x$};
\draw[->](0,-2)--(0,2)node[right]{$y$};
\draw[domain=-1.3:1.3, samples=100, very thick]plot({\x^3},\x)node[above]{$y=x^{\tfrac{1}{3}}$};
\node at (0,0) [above left]{$O$};
\end{tikzpicture}
    \caption{}
    \end{minipage}
    \begin{minipage}[t]{0.48\textwidth}
    \centering
    \begin{tikzpicture}[>=latex, xscale=.6]
       \draw[->](-4.5,0)--(5.5,0)node[right]{$x$};
\draw[->](0,-4)--(0,2)node[right]{$y$};
\draw[->](-4.5,-3)--(5.5,-3)node[right]{$x$};
\node at (0,0)[below left]{$O$};
\draw[domain=-4:5, samples=100, very thick]plot(\x, {sin((\x+pi/3)*180/pi)});
\draw[domain=-4:5, samples=100, very thick]plot(\x, {cos((\x+pi/3)*180/pi)-3});

\foreach \x/\xtext in {1.5*pi/m_3, -.5*pi/m_1, 0/I_1, pi/I_2}
{
   \draw[dashed](\x-pi/3, 0)node[above]{$\xtext$}--(\x-pi/3, -3);
}
\draw[dashed](.5*pi-pi/3, 1)--(.5*pi-pi/3, -3);
\node at (.5*pi-pi/3,0)[above right]{$m_2$};

\draw[domain=-2.5:.5, samples=10, thick]plot(\x, {\x+pi/3});
\draw[domain=1:3.2, samples=10, thick]plot(\x, {-\x+2*pi/3});
\node at (2,1)[right]{$y=f(x)$};
\node at (3,-4)[right]{$y=f'(x)$};
    \end{tikzpicture}
    \caption{}
    \end{minipage}
    \end{figure}

在图2.44中，函数$y=f(x)$的图象上的极值点$m_1$, $m_2$和$m_3$, 对应于函数$y=f'(x)$图象上的零点；$y=f(x)$的图象上的拐点对应于$y=f'(x)$的图象上的极值点．

\begin{ex}
\begin{enumerate}
    \item 讨论下列函数的极值、凸性及拐点，并出函数图象的
    草图；
\begin{multicols}{2}
\begin{enumerate}
    \item $f(x)=\frac{1}{1+x^2}$
    \item $y=(x-1)\sqrt[3]{x^2}$
\end{enumerate}
\end{multicols}
\item 如果函数$f(x)$在区间$(a,b)$内有二阶导数，并且向上凸，求证：对于满足条件$a\le x_1<x_2\le b$的任意两点$x_1,x_2$，恒有
\[f (\lambda x_1+ (1-\lambda) x_2) \ge \lambda f (x_1) + (1-\lambda)f (x_2) ,\quad  0\le\lambda\le 1\]
如果上面的函数在区间$(a,b)$内向下凸，则应有怎样的结论？
\item 求证函数
\begin{enumerate}
    \item $y=\log_ax \; (a>1)$在区间$(0,+\infty)$内是向上凸的；
    \item $f(x)=\sin x$在$[0,\pi]$内是向上凸的；
    \item $\log_a\frac{x_1+x_2}{2}\ge \frac{\log_a x_1+\log_a x_2}{2}\quad (x_1> 0,\; x_2>0)$
    \item $\sin\frac{x_1+x_2}{2}\ge \frac{\sin x_1+\sin x_2}{2}\quad  (0\le x_1,x_2\le \pi)$
\end{enumerate}

\item 如果函数$f(x)$在$(a,b)$内是向上凸的，求证：对于任意$a\le x_1<x_2<x_3\le b$，$\sum^3_{i=1}\lambda_i=1$，有
\[f(\lambda_1x_1+\lambda_2x_2+\lambda_3x_3)\ge \lambda_1f(x_1)+\lambda_2f(x_2)+\lambda_3f(x_3)\]

\item  求证对于任意正数$x_1,x_2,x_3$有
\[\ln\left(\frac{1}{6}x_1+\frac{1}{3}x_2+\frac{1}{2}x_3\right) \ge \ln \left(x_1^{\tfrac{1}{6}}x_2^{\tfrac{1}{3}}x_3^{\tfrac{1}{2}} \right)\]
\item  若$A,B,C$是任意三角形的三个内角，求证
\[\sin A+\sin B+\sin C\le \frac{3\sqrt{3}}{2}\]
又当$A,B,C$为何值时，上面等式成立？
\item  对于$a>0$, $b>0$, $0<\alpha<1$, $0<\beta<1$，$\alpha+\beta=1$，求证：
$a^{\alpha}b^{\beta}\le \alpha a+\beta b$成立．
（提示：利用对数函数是上凸函数）
\item  $f(x)=a_4x^4+a_3x^3+a_2x^2+a_1x+a_0$为具有正系数
的非零多项式，若为偶函数，求证$f(x)$在$(-\infty,+\infty)$内向下凸，且仅有一个极值．
\end{enumerate}
\end{ex}

\subsection{作函数的图象}
给出函数$y=f(x)$, 常常希望画出它的图象来，这对于许多问题是有用的，在初中我们所介绍的用列表法来描绘函数的图象，有很大的不精确性和盲目性，由于所取的点是有穷多个，很可能在所要画出的某两点之间的曲线有比较剧烈的变化，如果我们用比较平滑的曲线弧去连结，于是所得
曲线与实际相差很大．以前各节所介绍的函数性态可以帮助我们断定函数图象的主要特征，因此在取点之前，应先对函数的性态作一番讨论，兹将描绘曲线的一般步骤列出如下：

\begin{enumerate}
\item 确定函数的定义域;
\item 确定曲线关于坐标轴的对称性；
\item 确定曲线与坐标轴的交点（如果不易确定，不必强求）；
\item 确定函数的增减性，极大值与极小值；
\item 确定函数的凸性与拐点；
\item 确定曲线的渐近线；
\item 需要时，还得由曲线方程计算出一些适当的点的坐标；
\item 把上面所得的结果，按自变量大小顺序列入一个表格内，以观察图形的大概形态，然后描绘成图．
\end{enumerate}

在用例子说明如何绘图之前，我们需要补充说明如何求曲线的渐近线．

\begin{blk}
    {定义}假设曲线$y=f(x)$不限制在有限平面内，即其上的点$P$可以趋于无穷远，如果曲线的点$P$到直线$\ell:\; ax+by+c=0$的距离$d$, 当$P$沿曲线趋于无穷远时，
    $d\to 0$，
    那么我们称直线$\ell$为所给曲线的\textbf{渐近线}．
\end{blk}

\subsubsection{水平渐近线及其求法}

设$\ell_1:\; y=\alpha$是曲线$y=f(x)$的一条渐近线，曲线上的$P(x,f(x))$点到$\ell_1$的距离为
\[d=|\alpha-f(x)|\]
因$\ell_1$为渐近线，所以，当$x\to +\infty$或$x\to -\infty$时，应有$d\to 0$, 即
\[\lim_{x\to+\infty}f(x)=\alpha,\qquad \text{或}\qquad \lim_{x\to -\infty}f(x)=\alpha\]
反之，若上式成立，则$y=\alpha$必为$y=f(x)$的一条渐近线．

\subsubsection{铅直渐近线及其求法}

设$\ell_2:\; x=\beta$是曲线$y=f(x)$的一条渐近线，又$P(x,f(x))$为曲线上的任意一点，它到$x=\beta$的距离是
\[d=|\beta-x|\]
因为$\ell_2$是$y=f(x)$的渐近线，那么
\[\lim_{x\to\pm\beta} f (x) =\pm\infty\]
反之，若上式成立，则$\ell_2$为$y=f(x)$的一条渐近线．

例如，对于$y=\frac{1}{x}$, 当$x\to 0$, $y\to \pm\infty$，故$x=0$是$y=\frac{1}{x}$的一条渐近线．

\subsubsection{任意渐近线及其求法}

设直线$\ell_3:\; y=kx+b$是曲线$y=f(x)$的一条渐近线，曲线上任意一点$P(x,f(x))$到直线$\ell_3$的距离是
\[d=\left|\frac{kx+b-f(x)}{\sqrt{1+k^2}}\right|\]
因为$\ell_3$是渐近线，所以当$x\to\pm\infty$时，$d\to 0$, 即
\[kx+b-f (x) \to 0\]
我们必须用上式求出$k,b$, 才能确定$\ell_3$. 由于当$x\to\pm\infty$时，$\frac{1}{x}\to 0$, 所以
\begin{align*}
\lim_{x\to\pm\infty}\left[k+\frac{b}{x}-\frac{f(x)}{x}\right]&=\lim_{x\to\pm\infty}\frac{1}{x}[kx+b-f(x)]\\
&=\lim_{x\to\pm\infty}\frac{1}{x}\cdot \lim_{x\to\pm\infty}[kx+b-f(x)]=0
\end{align*}
由于$b$是常数，故有
\begin{align}
    k&=\lim_{x\to\pm\infty}\frac{f(x)}{x}\qquad \text{(渐近线方向)}\\
b&=\lim_{x\to\pm\infty}[f(x)-kx]
\end{align}

\begin{example}
求函数$f(x)=-3x^5+5x^3$的极值、拐点，并画出它的图象的草图．
\end{example}

\begin{solution}
\begin{enumerate}
\item 曲线的范围与截距：这函数$f$的定义域与值域都是无限的，它的曲线伸展在整个平面内，与$y$轴的交点是原点，而与$x$轴有三个交点，其零点是$0,\pm\sqrt{\frac{5}{3}}$.
\item 对称性：因为$f(-x)=-f(x)$, 所以曲线关于原点对称．
\item 函数的增、减与极值：
\[f' (x) =-15x^4+15x^2=-15x^2 (x-1) (x+1)\]
由此知$f(x)$的驻点是$-1, 0, 1$. $f'(x)$的符号变化的情
况是：
\begin{itemize}
    \item 当$x<-1$时，$f'(x)<0$, 这时$f(x)$递减；
    \item 当$-1<x<1$时, $f'(x)>0$，这时$f(x)$递增；
    \item 当$x>1$时，$f'(x)<0$, 这时$f(x)$递减．
\end{itemize}
因此，$x=-1$是极小点，极小值$f(-1)=-2$; $x=1$是极大点，极大值$f(1)=2$.
\item 函数的凸性与拐点：
\[f''(x)=-60x^3+30x=-60x\left(x+\frac{\sqrt{2}}{2}\right)\left(x-\frac{\sqrt{2}}{2}\right)\]
拐点的可能值是$-\frac{\sqrt{2}}{2},0,\frac{\sqrt{2}}{2}$. 又$f''(x)$的符号变化情况是
\begin{itemize}
    \item 当$x<-\frac{\sqrt{2}}{2}$时，$f''(x)>0$, 此时曲线向下凸；
    \item 当$-\frac{\sqrt{2}}{2}<x<0$时，$f''(x)<0$, 此时曲线向上凸；
    \item 当$0<x<\frac{\sqrt{2}}{2}$时，$f''(x)>0$, 此时曲线向下凸；
    \item 当$x>\frac{\sqrt{2}}{2}$时，$f''(x)<0$, 此时曲线向上凸；
\end{itemize}
因此：$x=-\frac{\sqrt{2}}{2},0,\frac{\sqrt{2}}{2}$的点都是拐点．

\item 列表：今将讨论情况列成下表，并作图（如图2.45）
\end{enumerate}    
\begin{center}\small
    \begin{tabular}{cccccccccccc}
    \hline
    $x$ & &$-1$&&$-\frac{\sqrt{2}}{2}$&&0&&$\frac{\sqrt{2}}{2}$&&1\\
    \hline
    $f'(x)$& $-$ & 0&$+$&$+$&$+$&0&$+$&$+$&$+$&0&$-$\\
    $f''(x)$&$+$&$+$&$+$&0&$-$&0&$+$&0&$-$&$-$&$-$\\
    $f(x)$&下凸&极小值& $\nearrow$  & $-1.24$&上凸 & 0& 下凸 &1.24& 上凸 & 极大值 & $\searrow$  \\
    & $\searrow $ &$-2$&&&$\nearrow $&&$\nearrow $&&$\nearrow $&2\\
    \hline
    \end{tabular}
    \end{center}
    
\begin{figure}[htp]
    \centering
\begin{tikzpicture}[>=latex, xscale=1.5]
\draw[->](-2.5,0)--(2.5,0)node[right]{$x$};
\draw[->](0,-3)--(0,3)node[right]{$y$}; 
\draw[domain=-1.4:1.4, samples=200, very thick]plot(\x, {-3*\x^5+5*\x^3});
\foreach \x in {-1,-2,1,2}
{
    \draw(\x,.1)--(\x,0)node[below]{$\x$};
}
\node at (0,0)[below right]{$O$};
\draw[dashed](-.707,-1.24)--(-.707,0)node[above]{$-\tfrac{\sqrt{2}}{2}$};
\draw[dashed](.707,1.24)--(.707,0)node[below]{$\tfrac{\sqrt{2}}{2}$};
\node at (1.29,0)[above]{$\sqrt{\tfrac{5}{3}}$};

\end{tikzpicture}
    \caption{}
\end{figure}
\end{solution}


\begin{example}    
描绘函数$y=\frac{(x-3)^2}{4(x-1)}$的图象．
\end{example}

\begin{solution}
    这个函数在$x=1$这点没有定义，并且，当$x\to 1^-$时，$y\to -\infty$; 当$x\to 1^+$时，$y\to +\infty$, 所以$x=1$是曲线的铅直渐近线，曲线在$x=1$这点断开．
    
求导数，得
\[y'= \frac{(x+1) (x-3)}{4 (x-1)^2}\]

当$x<-1$时，$y'>0$；当$-1<x<3$（$x\ne1$）时，$y'<0$；当$x>3$时，$y'>0$．因此，$x=-1$是极大点；$x=3$是极小点．又极大值$f(-1)=-2$; 极小值$f(3)=0$

求二阶导数，得
\[y''=\frac{2}{(x-1)^3}\]
\begin{itemize}
    \item 当$x<1$时，$y''<0$, 曲线向上凸；
    \item 当$x>1$时，$y''>0$, 曲线向下凸，但是曲线在$x=1$这一
    点断开，因此$x=1$不是曲线的拐点．
\end{itemize}

又因为
\[y=\frac{(x-3)^2}{4(x-1)}=\frac{1}{4}(x-5)+\frac{1}{x-1}\]
\[\lim_{x\to \infty}\left[\frac{(x-3)^2}{4(x-1)}-\frac{1}{4}(x-5)\right]=\lim_{x\to\infty}\frac{1}{x-1}=0\]
这就是说曲线以$y=\frac{1}{4}x-\frac{5}{4}$为渐近线．

综合起来得到下面结果：
\begin{center}
    \begin{tabular}{cccccccccccc}
    \hline
    $x$ & &$-1$&&0&&1&&3&\\
    \hline
    $y'$& $+$ & 0&$-$&$-$&$-$&$\x$&$-$&0&$+$\\
    $y''$&$-$&$-$&$-$&$-$&$-$&$\x$&$+$&$+$&$+$\\
    $y$&$\nearrow$ &极大值& $\searrow$  & $\searrow$&$\searrow$ & $\x$& $\searrow$ & 极小值 & $\nearrow$  \\
    &  &$-2$&&$-\frac{9}{4}$&$-\infty $&&$+\infty$&0 &\\
    \hline
    \end{tabular}
    \end{center}

画出具体图形如图2.46.
\begin{figure}[htp]
    \centering
\begin{tikzpicture}[>=latex]
    \draw[->](-3.5,0)--(6.5,0)node[right]{$x$};
    \draw[->](0,-5)--(0,3)node[right]{$y$}; 
\draw[domain=-3.5:6.5, samples=10, thick]plot(\x, {0.25*\x -1.25});
\draw[domain=-3:.7, samples=100, very thick]plot(\x, {(\x-3)^2/(4*(\x-1))});
\draw[domain=1.3:6, samples=100, very thick]plot(\x, {(\x-3)^2/(4*(\x-1))});
\node at (0,0)[below left]{$O$};
\draw[dashed](1,-5)--(1,3);
\foreach \x in {1,2,3,4}
{
    \draw(\x,0)node[below]{$\x$}--(\x,.1);
}
\foreach \x in {-1,-2}
{
    \draw(0,\x)node[left]{$\x$}--(.1,\x);
}

\node at (2,-1)[ right]{$y=\frac{1}{4}x-\frac{5}{4}$};
\node at (3,1)[ right]{$y=\frac{(x-3)^2}{4(x-1)}$};

\end{tikzpicture} 
    \caption{}
\end{figure}
\end{solution}



\begin{ex}
\begin{enumerate}
    \item 讨论下列函数的极值、凸性和拐点, 并画出函数图象.
\begin{multicols}{2}
\begin{enumerate}
\item  $y=\frac{6(x-1)}{x^{2}+3}$
\item  $y=(x-1)^{3}-3(x-1)$
\item  $y=\frac{x^{4}}{2}-x^{2}$
\item  $y=1+4 x^{3}-3 x^{4}$
\item  $y=\sqrt{(x+2)^{2}\left(16-x^{2}\right)}$
\item  $y=x^{x}\quad (x>0.1)$
\item  $y=x e^{-x}$
\end{enumerate}
\end{multicols}

    \item 作下列函数的图象, 并画出渐近线及在拐点处的切线:
    \begin{multicols}{2}
        \begin{enumerate}
            \item $y=x+\frac{1}{x}$
            \item $y=\frac{x^{2}}{1+x^{3}}$
            \item $y=\frac{2 x}{1-x^{2}}$
            \item $y=\frac{x^{3}}{x^{2}-4}$
            \item $y=\frac{(x-5)(x+3)}{(x+2)^{2}}$
\end{enumerate}
\end{multicols}
\end{enumerate}
\end{ex}

\subsection*{习题2.3}
\begin{enumerate}
    \item 求函数$f(x)=x^3+ax^2-3bx+c$的图象关于原点对称时，$f(x)$在区间$-2\le x\le 2$上的最小值，其中$b>0$.
    
    \item 函数$f(x)=3x^2-ax^3$在区间$0\le x\le 2$上的最小值
    是$-4$, 求
\begin{enumerate}
    \item $a$的值
    \item $f(x)$在区间$0\le x\le 2$的最大值$M$
\end{enumerate}
\item 设$a,b$为实数，求函数$y=x^3-3ax^2+bx+1$在区间$0\le x\le 1$上单调增加时，点$(a,b)$存在范围，并用图表示出来．

\item 设$f(x)=x^3-5x^2+3x+a$, 在怎样的情况下对于
$x$的正值，总有$f(x)>0$成立？
\item 设$a$是常数，对于函数$f(x)=4x^3-3x+a$
\begin{enumerate}
    \item 求$y=f(x)$的极大值和极小值；
    \item 确定当方程$f(x)=0$有三个不相等的实根时，$a$值的范围．
\end{enumerate}

\item 讨论函数$y=\frac{x^3}{2 (x+2)}$的极值，凸向，拐点，渐近线，
并作出它的图象．

\item 求证$\pi<\frac{\sin\pi x}{x(1-x)}\le 4$.
\item 对于已知半径为$R$的球，求具有最小体积的外切圆锥
的体积．
\item 一长方体的对角线长$\sqrt{11}$cm，全面积为$14{\rm cm}^2$, 求它的体积的最大值，并求体积为最大时三边的长．
\item 一浮标由三部分组成，一个圆筒与两个相等的圆
锥，其中每一个圆锥的高等于圆筒的高，问当表面积一定时，什么样的形状会有最大的体积．
\end{enumerate}     



\chapter{求函数从$a$到$b$的和与积分}
函数有两个基本性质——变率与和，在前一章，我们研究了函数的瞬时变率的概念，即
\[f'(x_0) =\lim_{h\to 0} \frac{f (x_0+h) -f (x_0)}{h}\]
以及它的应用．在本章我们将研究对于一个给定函数$f(x)$“求从$a$到$b$的和”这个概念．在给出定义之前，我们先举几个实例来看一看．

\paragraph{速率与距离} 
一列行驶中的火车，它的行进速率$v$是时间$t$的函数，即
$v=f(t)$, 如图3.1, 我们可以从速率表读得当时的速率，很自然地，我们想知道火车在从$t=a$到$t=b$这一段时间间隔内一共走了多少路程．从函数的观点看，所谓求在时间$[a,b]$内所走过的距离就是求速率函数$v=f(t)$由$t=a$到$t=b$的和．

\begin{figure}[htp]
    \centering
\begin{tikzpicture}[>=latex]
\draw[->](-.5,0)--(6,0)node[right]{$t$};
\draw[->](0,-.5)--(0,3)node[right]{$v=f(t)$};
\draw[very thick](0,0)--(1,1.5)--(1.6,1.5);
\draw[dashed](1,1.5)--(1,0)node[below]{$a$};
\node at (0,0)[below left]{$O$};
\draw plot[smooth] coordinates{
(1.6,1.5)(2.5,2.5)(3.5,2.2)(4,1.2)(4.3,.8)(4.8,.6)(5,0)
};
\draw[dashed](3.5,2.2)--(3.5,0)node[below]{$b$};

\end{tikzpicture}
    \caption{}
\end{figure}

让我们用$D(f,[a,b])$
表示所走过的距离，这个记号强调$D$依赖于$f$和区间$[a,b]$．

\paragraph{变力所作的功}

假定某物体在一个平行于$OX$轴的力$P$的作 用下沿直线
$OX$运动，力$P$的方向与物体运动的方向一致，并且力的大小随离开$O$点的距离而改变，即变力$P$是所在点的横坐标$x$的函数$P=P(x)$, 如图3.2．假定物体在这个变力$P$作用之下，从直线$OX$的一点$a$移到另一点$b$, 那么力$P$所作的功就是变力函数$P(x)$由$x=a$到$x=b$的和，我们用
$W (P, [a,b])$
表示力$P$所作的功，它表明$W$依赖于$P(x)$和$[a,b]$．

\begin{figure}[htp]
    \centering
    \begin{minipage}[t]{0.48\textwidth}
    \centering
    \begin{tikzpicture}[>=latex, scale=1]
\draw[->](-.5,0)--(4,0)node[right]{$x$};
\draw[->](0,-.5)--(0,3)node[right]{$P(x)$};
\draw[domain=0:3.5, samples=100, very thick]plot(\x, {3/(\x+1)-.5});
\draw(1,0)node[below]{$a$}--(1,1);
\draw(3,0)node[below]{$b$}--(3,.25);
\node at (0,0)[below left]{$O$};
    \end{tikzpicture}
    \caption{}
    \end{minipage}
    \begin{minipage}[t]{0.48\textwidth}
    \centering
    \begin{tikzpicture}[>=latex, scale=1]
\draw[->](-.5,0)--(4,0)node[right]{$x$};
\draw[->](0,-.5)--(0,3)node[right]{$y$};
\draw[domain=0:3.5, samples=100, very thick]plot(\x, {.3*(\x-2.2)^2+.5})node[above right]{$y=f(x)$};
\node at (0,0)[below left]{$O$};
\draw(.5,0)node[below]{$a$}--(.5,1.367)node[above right]{$M_0$};
\draw(3,0)node[below]{$b$}--(3,.692)node[above]{$M$};
\fill[pattern=north east lines, domain=.5:3, samples=100, very thick]plot(\x, {.3*(\x-2.2)^2+.5})--(3,.692)--(3,0)--(.5,0)--(.5,1.367);


    \end{tikzpicture}
    \caption{}
    \end{minipage}
\end{figure}

\paragraph{曲边梯形的面积}

令$y=f(x)$是函数$f$的图象，表示一条曲线，如图3.3. 我们要求曲线上的一段弧$M_0M$ 与其两端的纵坐标线及$x$轴上的线段$[a,b]$所围成的图形的面积．这样的图形（它有三条边是直线，其中两条互相平行，第三条与前两条
互相垂直，而第四条边是曲线）叫做曲边梯形．显然曲边梯形的面积$A$依赖于$y=f(x)$和$x$轴上的线段$[a,b]$, 记这一面积为$A (f, [a,b])$．

在上一章，我们利用函数$f(x)$和它的图象之间的对应关系，就可以把函数的变率和它的图象的切线的斜率相对应地一并分析讨论，这样，一方面可以把函数的变率这种“数量”的概念用切线来形象化，便于想象；而另一方面又可以把“切线”这种几何概念数量化，便于计算，在这一章，我们也要利用函数$f(x)$和它的图象之间的对应关系，把函数$f(x)$由$x=a$到$x=b$的“和”与曲线$y=f(x)$的曲边梯形的“面积”相对应地一并分析讨论，并且说明函数的“求从$a$到$b$的和”恰好对应于求曲边梯形的面积．这也就是为什么把函数“求从$a$到$b$的和”这种基本运算叫做积分的道理．

\section{{``}和{''}与{``}面积{''}}

对于任意曲线围成的图形，我们还没有规定它的“面积”的意义，和它密切相关的“函数从$a$到$b$的和”的概念至今也没有明确的解析的定义．在这一节我们要把这两个概念由“直观的定性理解”推进到“数量化的定量定义”．唯有确立了它们的“解析的定义”，它们才真正地成为能算好用的量．

\subsection{{``}和{''}与{``}面积{''}的基本性质}

现在让我们先从“函数的和”与“曲线形的面积”的直观内涵来分析一下它们分别所应有的基本性质．

\subsubsection{{``}曲线形的面积{''}的基本性质}

从面积的直观内涵容易看出下列两点：
\begin{enumerate}
\item (单调性）设区域$R_1$包含在$R_2$之内，即$R_1\subseteq R_2$, 则：
$R_1$的面积$\le R_2$的面积．（图3.4）
\item (可加性）设区域$R$可以用一条曲线分割成两块区
域$R_1+R_2$, 则有：
$R$的面积$=R_1$的面积$+R_2$的面积．（图3.5）
\end{enumerate}

\begin{figure}[htp]
    \centering
    \begin{minipage}[t]{0.48\textwidth}
    \centering
    \begin{tikzpicture}[>=latex, scale=1]
\draw[very thick](0,0) ellipse [x radius=2, y radius =1];
\draw[very thick](-.5,0) ellipse [x radius=1, y radius =.6];
\node at (0,1.3){$R$};
\node at (-.5,0){$R_1$};
\node at (1.3,0){$R_2$};
    \end{tikzpicture}
    \caption{}
    \end{minipage}
    \begin{minipage}[t]{0.48\textwidth}
    \centering
    \begin{tikzpicture}[>=latex, scale=1]
\draw[very thick](0,0) ellipse [x radius=2, y radius =1];
\node at (0,1.3){$R$};
\node at (-1,0){$R_1$};
\node at (1,0){$R_2$};
\draw[very thick](0,1) to [bend left](.3,.6)to [bend left](0,-.3)to [bend right](0,-1);
    \end{tikzpicture}
    \caption{}
    \end{minipage}
\end{figure}

\subsubsection{{``}函数的和{''}的基本性质}

同样地，从“函数从$a$到$b$的和”的直观内涵容易看出下列两个基本性质，即

\begin{blk}{性质1：单调性}
设函数$f(t)$和$g(t)$在$a\le t\le b$上，$f(t)\le g(t)$恒成立，那么：

“$f(t)$由$t=a$到$t=b$的和”$\le $“$g(t)$由$t=a$到$t=b$
的和”．
\end{blk}

例如，两部车子在$t=a$到$t=b$的时间内，甲车的速率$f(t)\le $乙车的速率$g(t)$, 则甲车在上述时间内所经过的里程$\le $乙车在上述时间内所经过的里程．

\begin{blk}{性质2：可加性}
设$a<b<c$，那么，有：

“$f(t)$由$t=a$到$t=b$的和”$+$“$f(t)$由$t=b$到$t=c$的和”$=$“$f(t)$由$t=a$到$t=c$的和”．
\end{blk}

上述基本性质是一目了然的，让我们先用来说明一些简单的基本事实．

\begin{example}
    设函数$y=f(t)=k$（常数），则由和的直观内涵，显然应有：
\[\text{常数函数从$a$到$b$的和}=k(b-a)\]

例如，以等速率每小时$k$公里行驶的火车从$t=a$小时到$t=b$小时内所走的总里程应该是$k(b-a)$公里．

现在让我们利用函数的图象来观察上述数值$k(b-a)$的几何意义．

\begin{figure}[htp]
    \centering
    \begin{minipage}[t]{0.48\textwidth}
    \centering
    \begin{tikzpicture}[>=latex, scale=.9]
\draw[->](-.5,0)--(5,0)node[right]{$t$};
\draw[->](0,-.5)--(0,3)node[right]{$y$};
\fill[pattern=north east lines](1,0) rectangle (4,1.5);
\draw[dashed](0,1.5)--(1,1.5)--(1,0)node[below]{$t=a$};
\draw[dashed](4,1.5)--(4,0)node[below]{$t=b$};
\node at (0,0)[below left]{$O$};
\draw[very thick](1,1.5)--(4,1.5);
\node at (0,.75)[left]{$k>0$};
\node at (2.5,.75){\Large $+$};

    \end{tikzpicture}
    \caption{}
    \end{minipage}
    \begin{minipage}[t]{0.48\textwidth}
    \centering
    \begin{tikzpicture}[>=latex, scale=.9]
\draw[->](-.5,0)--(5,0)node[right]{$t$};
\draw[->](0,-2)--(0,1.5)node[right]{$y$};
\fill[pattern=north east lines](1,0) rectangle (4,-1.5);
\draw[dashed](0,-1.5)--(1,-1.5)--(1,0)node[above]{$t=a$};
\draw[dashed](4,-1.5)--(4,0)node[above]{$t=b$};
\node at (0,0)[below left]{$O$};
\draw[very thick](1,-1.5)--(4,-1.5);
\node at (0,-.9)[left]{$k<0$};
\node at (2.5,-.75){\Large $-$};
    \end{tikzpicture}
    \caption{}
    \end{minipage}
\end{figure}

\begin{enumerate}
    \item 当$k>0$时，函数$f(t)=k$的图象在$t$轴上方，$k(b-a)$对应于一个矩形的面积，如图3.6中阴影部分的面积．
    \item 当$k<0$时，函数$f(t)=k$的图象在$t$轴下方，$k(b-a)$对应于一个矩形的面积的负值，如图3.7中阴影部分的面积的负值．
\end{enumerate}

假如我们把$t$轴之上的面积定义为正的，而把$t$轴之下的面积定义为负的，则当$y=k$（常数）时，常数函数从$a$到$b$的和$=k(b-a)=$上述的有号面积．
\end{example}

\begin{example}
设$y=f(t)$是一个阶梯函数，即分段地是常数函数：
设
\[a=t_0<t_1<t_2<\cdots<t_{i-1}<t_i<\cdots <t_n=b\]
当$t_{i-1}\le t<t_i$时，$f(t)=k_i,\; i=1,2,\ldots,n$．由性质2，我们可以分段地用常数函数求和，就可以得到阶梯函数$f(t)$从$a$到$b$的和．

阶梯函数$f(t)$从$a$到$b$的和
\[=k_1 (t_1-t_0) +k_2 (t_2-t_1) +\cdots +k_n (t_n-t_{n-1})\]

作出阶梯函数的图象（图3.8），上述函数和就等于下列逐段由矩形并起来的区域的“有号面积”．
\begin{figure}[htp]
    \centering
\begin{tikzpicture}[>=latex, yscale=.7]
\draw[->](-5,0)--(7,0)node[right]{$t$};
\draw[->](0,-2)--(0,4)node[right]{$y=f(t)$};
\node at (0,0)[above right]{$O$};

\foreach \x/\xtext in {-3/t_1, -1/t_2, 1.5/t_3, 2.5/t_4, 3.5/t_5, 4.3/t_6, 5.5/t_{n-1},-4.5/{t_0=a},6.5/{t_n=b}}
{
    \node at (\x,0)[below]{$\xtext$};
}

\draw[pattern=north east lines, thick](-4.5,0) rectangle(-3,1.5);
\draw[pattern=north east lines, thick](-3,1) rectangle(-1,0);
\draw[pattern=north east lines, thick](-1,0) rectangle(1.5,-1.5);
\draw[pattern=north east lines, thick](1.5,0) rectangle(2.5,1);
\draw[pattern=north east lines, thick](2.5,0) rectangle(3.5,2);
\draw[pattern=north east lines, thick](3.5,0) rectangle(4.3,-1);
\draw[pattern=north east lines, thick](5.5,0) rectangle(6.5,1);
\foreach \x in {-3.75, -2, 2, 3, 6}
{
    \node at (\x,.5){$+$};
}
\node at (.75,-.9){$-$};
\node at (3.9,-.65){$-$};
\node at (4.8,.5){…};\node at (4.8,-.5){…};

\end{tikzpicture}
    \caption{}
\end{figure}
\end{example}

\subsection{逼近法求和}

对于一个比较一般的函数，例如$y=kt+c$, $y=t^2$, $y=\sin t$等等，我们必须给函数的和下一个适当的定义，同时要提供计算这个和的方法，这里我们将要再一次地用逼近法的观点去解答上述问题．回顾上一章讨论变率时，我们的基本想法是用折线函数去逼近一般的平滑函数，所以折线函数是讨论函数变率的简单好用的基本函数，因为它们的变率十分简单，分段地是个常数而且无限逼近函数在某一点的变率．由上一段的讨论；我们知道阶梯函数的“从$a$到$b$的和”是十分简明的，它们在“函数求从$a$到$b$的和”这方面是不是也扮演着
这种既简单又基本的角色呢？这就得看一看能否用阶梯函数的从$a$到$b$的和去无限逼近一般性的“函数从$a$到$b$的和”了！

现在让我们使用几个简单的实例来试试看．

\begin{example}
假设物体以初速度$u$, 加速度$a>0$在直线上运动，于是物体在任何时刻$t$的速度是$v=f (t) =u+at$, 求物体从$t=0$到$t=T$时所经过的距离，即求速度函数$f(t)$在时间间隔$[0,T]$上的和．
\end{example}

\begin{analyze}
    一个自然的作法是用两个阶梯函数$g_n(t)$和$G_n(t)$把上述函数$f(t)=u+at$夹逼在中间，即
    \[g_n(t)\le f(t)\le G_n(t)\]
对于任何$0\le t\le T$都成立，那么由性质1，$g_n(t)$从0到$T$的和$\le f(t)$从$0$到$T$的和$\le G_n(t)$从0到$T$的和．

如果能使$g_n(t)$从0到$T$的和与$G_n(t)$从0到$T$的和，当$n\to\infty$时的极限相同，则由上述夹逼关系就可以看出$f(t)$从0到$T$的和必须等于这个共同极限．
\end{analyze}

\begin{figure}[htp]
    \centering
\begin{tikzpicture}[>=latex]
\draw[->](-.5,1)--(8,1)node[right]{$t$};
\draw[|<->|](-.2,1)--node[fill=white]{$u$}(-.2,2);
\draw[->](0,0)--(0,6)node[right]{$v=f(t)=u+at,\; (0\le t\le T)$};
\node at (0,1)[below left]{$O$};
\draw[ultra thick](0,2)node[above left]{$A$}--(7,5.5)node[above]{$B$};
\draw[<->](7.3,1)node[below]{$C$}--node[fill=white]{$u$}(7.3,2);
\draw[|<->|](7.3,5.5)--node[fill=white]{$aT$}(7.3,2);
\foreach \x in {1,2,3,...,6,7}
{
    \ifnum\x<7 \node at (\x,1)[below]{$t_{\x}$};\fi
    \draw[dashed](\x,1)--(\x,2+\x*.5);
    \draw[very thick](\x-1, 1.5+\x*0.5)--(\x, 1.5+ \x*0.5)--(\x, 2+\x*0.5);
    \draw[thick](\x-1, 1.5+\x*0.5)--(\x-1, 2+\x*0.5)--(\x, 2+ \x*0.5);
}
\draw[dashed](0,2)--(7,2);
\draw(7,0)--(7,1);
\draw[<->](0,.2)--node[fill=white]{$T$}(7,.2);
\end{tikzpicture}
    \caption{}\label{fig:chapter3_sec1_line_segmentation}
\end{figure}

下面就是把上述想法付诸实践的具体做法之一．

\begin{solution}
\begin{enumerate}
    \item 把时间间隔$n$等分，取分点$t_i=\frac{i}{n}T,\; i=0,1, 2,\ldots,n$, 于是在每个分点$t_i=\frac{i}{n}T$处的速度分别是
\[f (t_0) =f (0) =u,\;  f(t_1) =u+at_1,\; f (t_2) =u+at_2,\ldots,f (t_n) =f (T) =u+aT\]
\item 用阶梯函数近似代替$v=f(t)=u+at$. 

假设物体在时间的每个小区间$[t_i,t_{i+1}],\; i=0,1, 2,\ldots,(n-1)$内，以小区间起点处的速度作匀速运动，我们得到一个速度的阶梯函数
\[g_n (t)=f(t_i)=u+at_i,\quad t_i\le t<t_{i+1},\quad i=0,1,2,\ldots,(n-1)\]
因为$f(t)=u+at$是严格递增的，所以$g_n(t)$有以下性质：
\[g_n (t)=f(t_i)\le f(t),\quad t_i\le t<t_{i+1},\quad i=0,1,2,\ldots,(n-1)\]
从而，对于任何$0\le t\le T$，有$g_n(t)\le f(t)$.

假设物体在时间的每个小区间$[t_i,t_{i+1}],\; i=0,1, 2,\ldots,(n-1)$内，以小区间终点处的速度作匀速运动，我们得到另一个速度的阶梯函数
\[G_n(t)=f(t_{i+1})=u+at_{i+1},\quad t_i< t\le t_{i+1},\quad i=0,1,2,\ldots,(n-1)\]

由于$f(t)=u+at$是递增的，所以$G_n(t)$有以下性质
\[G_n (t)=f(t_{i+1})\ge f(t),\quad t_i< t\le t_{i+1},\quad i=0,1,2,\ldots,(n-1)\]
从而对于任何$0\le t\le T$，有
\[G_n(t)\ge f(t)\]
这样，我们就得到了$v=f(t)$的夹逼阶梯函数：
\[ g_n(t) \le f (t) \le G_n (t),\quad 0\le t\le T\]

（在图\ref{fig:chapter3_sec1_line_segmentation}中，绘出了$f(t)=u+at$的图象和当$n=7$时阶梯函数$g_7(t)$和$G_7(t)$的图象.）

\item 求阶梯函数的总和

应用基本性质2, 我们得到$g_n(t)$从0到$T$的和
\[f(t_0)(t_1-t_0)+f(t_1)(t_2-t_1)+\cdots+f (t_{n-1}) (t_n-t_{n-1})\]

$\because\quad t_1-t_0=t_2-t_1=\cdots=t_n-t_{n-1}=\frac{T}{n}$

因此：
\begin{align*}
\text{$g_n(t)$从0到$T$的和}&=\left[f(t_0)+f(t_1)+\cdots+f(t_{n-1})\right]\cdot \frac{T}{n}\\
&=\left[u+(u+at_1)+(u+at_2)+\cdots +(u+at_{n-1})\right]\cdot \frac{T}{n}\\
&=\left[nu+a(t_1+t_2+\cdots+t_{n-1})\right]\cdot \frac{T}{n}\\
&=\left[nu+a\left(\frac{T}{n}+\frac{2T}{n}+\cdots+\frac{(n-1)T}{n}\right)\right]\cdot \frac{T}{n}\\
&=uT+a\cdot [1+2+\cdots+(n-1)]\left(\frac{T}{n}\right)^2\\
&=uT+a\cdot \frac{n(n-1)}{2}\cdot \frac{T^2}{n^2}\\
&=uT+\frac{1}{2}aT^2\left(1-\frac{1}{n}\right)
\end{align*}

\begin{align*}
    \text{$G_n(t)$从0到$T$的和}&=f(t_1)(t_1-t_0)+f(t_2)(t_2-t_1)+\cdots+f(t_n)(t_n-t_{n-1})  \\
   &= \left[f(t_1)+f(t_2)+\cdots+f(t_{n})\right]\cdot \frac{T}{n}\\
    &=\left[(u+at_1)+(u+at_2)+\cdots +(u+at_{n})\right]\cdot \frac{T}{n}\\
    &=\left[nu+a\left(\frac{T}{n}+\frac{2T}{n}+\cdots+\frac{nT}{n}\right)\right]\cdot \frac{T}{n}\\
    &=uT+a\cdot (1+2+\cdots+n)\cdot \left(\frac{T}{n}\right)^2\\
    &=uT+\frac{1}{2}aT^2\left(1+\frac{1}{n}\right)
    \end{align*}
综合上述计算和基本性质1, 即得：$g_n(t)$从0到$T$的和$\le f(t)$从0到$T$的和$\le G_n(t)$从0到$T$的
和，即
\[uT+\frac{1}{2}aT^2\left(1-\frac{1}{n}\right)\le f(t)
\text{从0到$T$的和}\le uT+\frac{1}{2}aT^2\left(1+\frac{1}{n}\right)\]

\item 求阶梯函数和式的极限．

因为
\begin{align*}
    \lim_{n\to\infty}\left\{uT+\frac{1}{2}aT^2\left(1-\frac{1}{n}\right)\right\}&=\lim_{n\to\infty}\left\{uT+\frac{1}{2}aT^2\left(1+\frac{1}{n}\right)\right\}\\
    &=uT+\frac{1}{2}aT^2
\end{align*}
所以，这个共同的极限值$uT+\frac{1}{2}aT^2$就是物体在变速$v=f(t)=u+at$运动下，从$t=0$到$t=T$所走的距离，也就是函数$v=f(t)=u+at$从$t=0$到$t=T$的和．
\end{enumerate}

从几何上看，上述极限值$uT+\frac{1}{2}aT^2=\frac{[u+(u+aT)]T}{2}$就是在图\ref{fig:chapter3_sec1_line_segmentation}中的梯形$OABC$的面积．
\end{solution}    

\begin{example}
    设函数$f(x)=x^2$, $a=0$, $b>0$, 求$f(x)$由0到$b$的和，相应地，求抛物线$y=x^2$在线段$[0,b]$上所盖的曲边三角形$OBC$的面积（图\ref{fig:chapter3_sec1_y=x^2}）.
\end{example}

\begin{figure}[htp]
    \centering
\begin{tikzpicture}[>=latex, scale=.5]
\draw[->](-5,0)--(5,0)node[right]{$x$};
\draw[->](0,-1.5)--(0,11)node[right]{$y$};
\draw[domain=-3.25:3.25, samples=100, very thick]plot(\x, {\x*\x});
\node at (0,0)[below left]{$O$};
\foreach \x in {.5,1,...,3}
{
    \draw[thick](\x-.5,0) rectangle (\x, \x*\x);
}
\foreach \x in {.5,1,...,2.5}
{
    \draw(\x,\x*\x) -- (\x+.5,\x*\x);
}
\node at (3,9)[right]{$C$};
\node at (3,0)[below]{$B$};
\node at (-2.5,6)[left]{$y=x^2$};
\end{tikzpicture}
    \caption{}\label{fig:chapter3_sec1_y=x^2}
\end{figure}

\begin{solution}
\begin{enumerate}
    \item 把线段$[0,b]$ $n$等分，等分点的坐标是
   $ x_i=\frac{b}{n}i\;  (i=0, 1, 2,\ldots,n)$ 
   \item 作阶梯函数$g_n(x)$和
    $G_n (x)$.

    根据$f(x)=x^2$在$x>0$时递增，定义
\begin{align*}
g_n(x)&=f(x_i)=x^2_i=\left(\frac{b}{n}\right)^2i^2,\quad x_i\le x<x_{i+1}\\
G_n(x)&=f(x_{i+1})=x^2_{i+1}=\left(\frac{b}{n}\right)^2(i+1)^2,\quad x_i<x\le x_{i+1}\\
\end{align*}
其中，$i=0,1,2,\ldots,(n-1)$．于是，$g_n(x)\le f(x)\le G_n(x),\quad 0\le x\le b$

\item 求阶梯函数的和
\begin{align*}
s_n&=g_n(x)\text{从0到$b$的和}\\
&=f\left(x_{0}\right) \cdot \frac{b}{n}+f\left(x_{1}\right) \cdot \frac{b}{n}+\cdots+f\left(x_{n-1}\right) \cdot \frac{b}{n}\\
&=\left\{f\left(x_{0}\right)+f\left(x_{1}\right)+\cdots+f\left(x_{n-1}\right)\right\} \cdot \frac{b}{n} \\
&=\left\{0^{2}+1^{2}+2^{2}+\cdots+(n-1)^{2}\right\} \cdot\left(\frac{b}{n}\right)^3\\
&=\frac{1}{6}n(n-1)(2 n-1) \cdot\left(\frac{b}{n}\right)^{3} =\frac{b^3}{6}\left(1-\frac{1}{n}\right)\left(2-\frac{1}{n}\right) 
\end{align*}
\begin{align*}
   S_{n}= G_{n}(x) \text { 从0到$b$的和}
   &= \left\{f\left(x_{1}\right)+f\left(x_{2}\right)+\cdots+f\left(x_{n}\right)\right\} \cdot \frac{b}{n} \\
   &=\left\{1^2+2^{2}+\cdots+n^{2}\right\} \cdot\left(\frac{b}{n}\right)^{3} \\
   &=\frac{1}{6}n(n+1)(2 n+1) \cdot\left(\frac{b}{n}\right)^{3} \\
   &=\frac{b^3}{6}\left(1+\frac{1}{n}\right)\left(2+\frac{1}{n}\right) 
\end{align*}

综合上述计算和性质1，有：
\begin{align*}
\frac{b^3}{6}\left(1-\frac{1}{n}\right)\left(2-\frac{1}{n}\right)&=s_n\le \begin{bmatrix}
\text{函数$f(x)=x^2$从0到$b$的和，}\\
\text{相应的曲边三角形$OBC$的面积}
\end{bmatrix}\\
&\le S_n=\frac{b^3}{6}\left(1+\frac{1}{n}\right)\left(2+\frac{1}{n}\right)
\end{align*}

\item 令$n\to\infty$，求阶梯函数和式的极限．因为
\begin{align*}
    \frac{b^3}{6} \left(1-\frac{1}{n}\right)\left(2-\frac{1}{n}\right)\to \frac{b^3}{3}\qquad (n\to\infty)\\
    \frac{b^3}{6} \left(1+\frac{1}{n}\right)\left(2+\frac{1}{n}\right)\to \frac{b^3}{3}\qquad (n\to\infty)\\
\end{align*}
所以$\frac{1}{3}b^3$是能被$s_n$、$S_n$左、右夹逼的唯一实数，而所求的“$f(x)=x^2$从0到$b$的和”以及“曲边三角形$OBC$的面积，”这两个值都是夹逼在$s_n$和$S_n$中间的实数，故它们都等于$\frac{1}{3}b^3$.
\end{enumerate}
\end{solution}

\begin{blk}
    {推论} $f(x)=x^2$从$a$到$b$的和$=\frac{1}{3}b^3-\frac{1}{3}a^3\quad  (b> a> 0)$.
\end{blk}


\begin{example}
设$y=f(x)$是一个定义在$[a,b]$上的递增连续函数，试说明它从$a$到$b$的和是可以确定的．
\end{example}

\begin{figure}[htp]
    \centering
\begin{tikzpicture}[>=latex]
\draw[->](-.5,0)--(5,0)node[right]{$x$};
\draw[->](0,-.5)--(0,4)node[right]{$y$};
\draw[domain=.5:4, samples=50, very thick]plot(\x, {.2*\x*\x+.5});
\draw[dashed](1,.7)--(1,0)node[below]{$a$};
\draw[dashed](4,3.7)--(4,0)node[below]{$b$};
\draw[dashed](2,1.3)--(2,0);
\draw[dashed](3,2.3)node[left]{$y=f(x)$}--(3,0);
\node at (0,0)[below left]{$O$};
\end{tikzpicture}
    \caption{}\label{fig:chapter3_sec1_monotone f}
\end{figure}

\begin{solution}
\begin{enumerate}
    \item     取$[a,b]$之间的$n$等分点，将它分成$n$段，即有
\[x_0=a<x_1<x_2<\cdots<x_n=b\]
它的第$i$段是$[x_{i-1},x_i]$，且$x_i-x_{i-1}=\frac{b-a}{n}$，$x_i=a+\frac{i}{n}(b-a)$

\item 定义$f(x)$的上、下夹逼阶梯函数如下：
\[
\begin{split}
    g_n(x)&=f(x_{i-1}),\qquad x_{i-1}\le x\le x_i \\
    G_n(x)&=f(x_i),\qquad x_{i-1}\le x\le x_i 
\end{split} \quad (i=1,2,3,\ldots,n)
\]
由$f(x)$的递增性，对于任何一个$x\in[a,b]$，有
\[g_1(x)\le g_2(x)\le\cdots \le g_n(x)\le\cdots \le f(x)\le \cdots\le G_n(x)\le\cdots\le G_2(x)\le G_1(x)\]
简写成
\[g_n(x)\le f(x)\le G_n(x),\qquad a\le x\le b\]
\item 求相应的阶梯函数从$a$到$b$的和
\begin{equation}
\begin{split}
    s_n=g_n(x)\text{从$a$到$b$的和}
    &=\frac{b-a}{n}\left[f(x_0)+f(x_1)+\cdots+f(x_{n-1})\right]\\
    &=\frac{b-a}{n}\sum^n_{i=1}f(x_{i-1})
\end{split}
\end{equation}
\begin{equation}
    \begin{split}
    S_n=G_n(x)\text{从$a$到$b$的和}
    &=\frac{b-a}{n}\left[f(x_1)+f(x_2)+\cdots+f(x_{n})\right]\\
    &=\frac{b-a}{n}\sum^n_{i=1}f(x_{i})   
    \end{split}
    \end{equation}
\end{enumerate}
由性质1, 有：
\[s_1<s_2<\cdots<s_n<\cdots<f(x)\text{从$a$到$b$的和}<\cdots<S_n<\cdots<S_2<S_1\]
简写成
\begin{equation}
    s_n<f(x)\text{从$a$到$b$的和}<S_n
\end{equation}
在例3.3, 例3.4中，$f(x)$有明确的解析式，我们可以用求和公式直接求出$s_n$和$S_n$的表达式，从而可以得出它们的共同极限值，在这里，我们只知道$f(x)$的递增性，当然无法将和式$s_n$和$S_n$进一步化简，但是我们可以说明，当$n\to\infty$时，$S_n-s_n\to 0$.

事实上，
\begin{align*}
    S_n-s_n&=\frac{b-a}{n}\sum^n_{i=1}f(x_{i})   -\frac{b-a}{n}\sum^n_{i=1}f(x_{i-1}) \\
    &=\frac{b-a}{n}\sum^n_{i=1}\left[f(x_{i})-f(x_{i-1})\right]   \\
    &=\frac{b-a}{n}\left[f(x_{n})-f(x_{0})\right]  =\frac{b-a}{n}\left[f(b)-f(a)\right]   
\end{align*}
因此，当$n\to\infty$时，$S_n-s_n\to 0$.

又$f(x)$从$a$到$b$的和是介于$s_n$和$S_n$中间的唯一实数，
因此
\[\lim_{n\to\infty} s_n=\lim_{n\to\infty} S_n=f(x)\text{ 从$a$到$b$的和}\]

以上简单明了的分析，说明了下面两点互相关联的事实：
\begin{enumerate}
    \item 当函数$f(x)$在$a\le x\le b$上递增时，则它从$a$到$b$的和可以用上述两个夹逼阶梯函数从$a$到$b$的和去无限逼近.
    \item 由于$\Lim_{n\to\infty}s_n$与$\Lim_{n\to\infty}S_n$存在，且$\Lim_{n\to\infty}s_n=\Lim_{n\to\infty}S_n$, 我们可以用$f(x)$的上、下夹逼阶梯函数序列，即：$g_1 (x)\le g_2(x)\le \cdots\le g_n(x)\le \cdots\le f(x)\le \cdots\le G_n(x)\le \cdots\le G_2(x)\le G_1(x),\; (a\le x\le b)$的从$a$到$b$的和的极限$\Lim_{n\to\infty}s_n=\Lim_{n\to\infty}S_n$作为上述$f(x)$从$a$到$b$的和的数量化定义．
\end{enumerate}
\end{solution}

\begin{blk}
    {定义1}若$f(x)$是$[a,b]$上的递增连续函数，将$[a,b]$等分成$n$段，分点坐标$x_i=a+\frac{i}{n}(b-a),\; i=0,1, 2,\ldots,n$，则存在两个阶梯函数：
\[    g_n (x) =f (x_{i-1}) ,\quad x_{i-1}\le  x< x_i\]
    和
\[    G_n (x)=f(x_i),\quad x_{i-1}<x\le x_i,\quad (i=1,2,\ldots,n)\]
满足下面的性质：
\begin{enumerate}
    \item 对于任何$x\in[a,b]$，$g_n(x)\le f(x)\le G_n(x)$
    \item 相应的阶梯函数从$a$到$b$的和
\[s_n=\frac{b-a}{n}\sum^n_{i=1}f(x_{i-1}),\qquad S_n=\frac{b-a}{n}\sum^n_{i=1}f(x_{i})  \]
\end{enumerate}
当$n\to\infty$时，$S_n-s_n\to 0$，那么$f(x)$从$a$到$b$的和是$s_n$与$S_n$的共同极限，即：
\begin{align*}
f(x)\text{从$a$到$b$的和}&=\lim_{n\to\infty}\frac{b-a}{n}\sum^n_{i=1}f(x_{i-1})\\
&=\lim_{n\to\infty}\frac{b-a}{n}\sum^n_{i=1}f(x_{i})
\end{align*}

同样地，如果$f(x)$是$[a,b]$上的递减连续函数．那么将$[a,b]$ $n$等分，分点坐标$x_i=a+\frac{i}{n}(b-a)$, 这时有两个阶梯函数
\[    g_n (x) =f (x_{i}) ,\quad x_{i-1}<  x\le x_i\]
    和
\[    G_n (x)=f(x_{i-1}),\quad x_{i-1}\le x< x_i,\quad (i=1,2,\ldots,n)\]
满足下面的性质：
\begin{enumerate}
    \item 对于任何$x\in[a,b]$，$g_n(x)\le f(x)\le G_n(x)$
    \item 相应的阶梯函数$g_n(x)$与$G_n(x)$相应的和
\begin{align*}
    s_n&=g_n\text{从$a$到$b$的和}=\frac{b-a}{n}\sum^n_{i=1}f(x_{i})\\
    S_n&=G_n\text{从$a$到$b$的和}=\frac{b-a}{n}\sum^n_{i=1}f(x_{i-1})  
\end{align*}
\end{enumerate}
当$n\to\infty$时，$S_n-s_n=\frac{b-a}{n}\left[f(a)-f(b)\right]\to 0$，即：
\[\Lim_{n\to\infty}s_n=\Lim_{n\to\infty}S_n\]
\end{blk}

\begin{blk}
    {定义2} 递减连续函数从$a$到$b$的和为上述夹逼阶梯函数$g_n(x)$和$G_n(x)$从$a$到$b$的和的共同极限，即  
\begin{align*}
    f(x)\text{从$a$到$b$的和}&=\lim_{n\to\infty} \frac{b-a}{n}\sum^n_{i=1}f(x_i)\\
    &=\lim_{n\to\infty} \frac{b-a}{n}\sum^n_{i=1}f(x_{i-1})\\
\end{align*}
\end{blk}

通常我们把递增或递减的函数合称为单调函数，常见的
函数$f(x)$都是分段单调连续的，例如，$y=\sin x$本身虽然不是单调的，但是它在$\left[-\frac{\pi}{2}+2k\pi,\frac{\pi}{2}+2k\pi\right]\; (k\in\mathbb{Z})$这些区间的每一段是递增的；而在$\left[\frac{\pi}{2}+2k\pi,\frac{3\pi}{2}+2k\pi\right]\; (k\in\mathbb{Z})$这些区间的每一段是递减的．因此，对于一般在$[a,b]$上连续的函数，如果存在有限个分点使得$f(x)$在每个分段上都是单调的，我们可以逐段取上、下夹逼阶梯函数，合起来作为定义在$[a,b]$上的函数$f(x)$的阶梯函数，于是得出
\[g_n (x)\le f(x)\le G_n(x),\qquad x\in [a,b]\]
又因为有限个在分段上趋于0的量的和仍趋于0, 所以
\[\text{“$G_n(x)$从$a$到$b$的和”}-\text{“$g_n(x)$从$a$到$b$的和”}\to 0\]

总结以上讨论，我们叙述为存在定理和定义如下：

\begin{blk}
    {定理} 设$y=f(x)$是一个定义在$[a,b]$上分段单调函数，则存在满足下列性质的两系列阶梯函数：
    \[g_n(x)\le f(x)\le G_n(x),\qquad x\in[a,b]\]
而且当$n\to\infty$时，“$G_n(x)$从$a$到$b$的和，”与“$g_n(x)$从$a$到$b$的和”趋于共同极限．
\end{blk}


\begin{blk}
    {定义3} 设$f(x)$是$[a,b]$上的连续函数，如果 存在有限个分点：
\[a=a_0<a_1<\cdots<a_k<\cdots<a_\ell=b\]
使得$f(x)$在每个分段$[a_{k-1},a_k]$都是单调的，再将每个分段
    $[a_{k-1},a_k]$都$n$等分，则当$n\to\infty$时，有
\begin{align*}
f(x)\text{从$a$到$b$的和}&=\lim_{n\to\infty}\frac{a_1-a_0}{n}\sum^n_{i=1}f(x_i)+\lim_{n\to\infty}\frac{a_2-a_1}{n}\sum^{2n}_{i=n+1}f(x_i)+\\
&\cdots +\lim_{n\to\infty}\frac{a_\ell-a_{\ell-1}}{n}\sum^{\ell n}_{i=(\ell-1)n+1}f(x_i)
\end{align*}
\end{blk}

为了说明上述定义是合理的，我们就得证明上述$f(x)$从$a$到$b$的和与夹逼的阶梯函数列$\{G_n(x)\}$和$\{g_n(x)\}$的选取无关，其证明如下：

\begin{proof}
    设$\{\bar G_m(x)\}$和$\{\bar g_m(x)\}$是另外一组满足存在定理的上、下夹逼函数列，则由下述不等式
\[g_n(x)\le f(x)\le G_n(x),\quad \bar g_m(x)\le f(x)\le \bar G_m(x),\qquad a\le x\le b\]
即有
\[g_n(x)\le \bar G_m(x),\qquad \bar g_m(x)\le G_n(x)\]
所以由基本性质1，有
\begin{align*}
    s_n&=g_n(x)\text{从$a$到$b$的和}\le \bar G_m(x)\text{从$a$到$b$的和}=\bar S_m\\
    \bar s_m&=\bar g_m(x)\text{从$a$到$b$的和}\le G_n(x)\text{从$a$到$b$的和}=S_n\\
\end{align*}
所以
\[\lim_{n\to\infty}s_n\le\lim_{m\to\infty}\bar S_m,\qquad \lim_{m\to\infty}\bar s_m\le \lim_{n\to\infty}S_n\]
但是
\[\lim_{n\to\infty}s_n\le \lim_{m\to\infty}\bar S_m=\lim_{m\to\infty}\bar s_m\le \lim_{n\to\infty} S_n=\lim_{n\to\infty}s_n\]
所以上述极限必须相等，即：
\[\lim_{n\to\infty}s_n= \lim_{m\to\infty}\bar S_m=\lim_{m\to\infty}\bar s_m=\lim_{n\to\infty} S_n\]
\end{proof}

\begin{example}
求$f(x)=2x-x^2$从0到2的和．
\end{example}

\begin{solution}
    由$f'(x)=2-2x$, $f''(x)=-2$知$f(1)=2-1=1$是
$f(x)$的极大值，并且$y=2x-x^2$在$[0, 1]$上递增，在$[1,2]$上递减，于是，我们把区间$[0, 1]$和$[1, 2]$都$n$等分，设分点坐标$x_i=\frac{i}{n},\; i=0, 1, 2,\ldots,2n$, 即有
\[0=x_0<x_1<x_2<\cdots <x_n=1<x_{n+1}<x_{n+2}<\cdots <x_{2n}=2\]
且$\Delta x_i=x_i-x_{i-1}=\frac{1}{n}$，由定义3得到：
\begin{align*}
    f(x)&=2x-x^2\text{从0到2的和}\\
    &=\lim_{n\to\infty}\frac{1}{n}\sum^n_{i=1}\left[2\left(\frac{i}{n}\right)-\left(\frac{i}{n}\right)^2\right]+\lim_{n\to\infty}\frac{1}{n}\sum^{2n}_{i=n+1}\left[2\left(\frac{i}{n}\right)-\left(\frac{i}{n}\right)^2\right] \\
&=\lim_{n\to\infty}\frac{1}{n}\sum^{2n}_{i=1}2\left(\frac{i}{n}\right)-\lim_{n\to\infty}\frac{1}{n}\sum^{2n}_{i=1}\left(\frac{i}{n}\right)^2\\
&=2\lim_{n\to\infty}\frac{1}{n^2}\sum^{2n}_{i=1}i-\lim_{n\to\infty}\frac{1}{n^3}\sum^{2n}_{i=1}i^2\\
&=2\lim_{n\to\infty}\frac{1}{n^2}\cdot \frac{(2n+1)2n}{2}-\lim_{n\to\infty}\frac{1}{n^3}\cdot \frac{2n(2n+1)(4n+1)}{6}\\
&=2\lim_{n\to\infty}\left(2+\frac{1}{n}\right)-\lim_{n\to\infty}\frac{1}{3}\left(2+\frac{1}{n}\right)\left(4+\frac{1}{n}\right)\\
&=4-\frac{8}{3}=\frac{4}{3}
\end{align*}
\end{solution}

\begin{ex}
\begin{enumerate}
    \item 已知质点的运动速度$v=t+4$, 试求质点在前10秒内所走的路程．
    \item 求$f(x)=x^3$在$1\le x\le 2$上的和．
\end{enumerate}
\end{ex}


\section{定积分的定义和基本性质}
\subsection{定积分定义}

设函数$f(x)$在$[a,b]$上连续，如果存在有限个分点
\[a=a_0<a_1<a_2<\cdots<a_{\ell-1}<a_{\ell}=b\]
使得$f(x)$在每个分段$[a_{k-1},a_k]\; (k=1, 2,\ldots,\ell)$上都是单调的，我们把$f(x)$从$a$到$b$的和叫做$f(x)$从$a$到$b$的\textbf{定积分}，并记作
\[\int^b_a f(x) \dd x\]

这里，积分符号所使用的是长$S$形的求和号的变形，而从部分区间长$\Delta x_i=\frac{a_k-a_{k-1}}{n}$过渡到极限，则通过字母$\dd$来表示．

我们把数$a$与$b$称为\textbf{积分限}（$a$称为\textbf{下限}，$b$称为\textbf{上限}）．区间$[a,b]$称为\textbf{积分区间}，函数$f(x)$称为\textbf{被积函数}，乘积$f(x)\dd x$称为\textbf{被积表达式}，定积分符号下出现的字母$x$叫做\textbf{积分变量}．

上述定积分定义用定积分符号表示就是：
\[\int^b_a f(x) \dd x=\sum^{\ell}_{k=1}\int^{a_k}_{a_{k-1}}f(x)\dd x\]
其中$\Int^{a_k}_{a_{k-1}}f(x)\dd x$是$f(x)$为单调的第$k$个分段$[a_{k-1},a_k]$上的夹逼阶梯函数$g_n(x)$和$G_n(x)$从$a_{k-1}$到$a_k$的和的共同极限，即：
\begin{align*}
    \int^{a_k}_{a_{k-1}}f(x)\dd x&=\lim_{n\to \infty}\frac{a_k-a_{k-1}}{n}\sum^n_{i=1}f(x_{i-1})\\
    &=\lim_{n\to \infty}\frac{a_k-a_{k-1}}{n}\sum^n_{i=1}f(x_{i})
\end{align*}

定积分定义的另一种表述形式是：设函数$f(x)$在$a\le x\le b$上分段单调连续，如果存在两系列上、下夹逼阶梯函数$\{g_n(x)\}$, $\{G_n(x)\}$, 使得
\[g_n (x)\le f(x)\le G_n(x),\qquad a\le x\le b\]
并且$g_n(x)$从$a$到$b$的和$s_n$与$G_n(x)$从$a$到$b$的和$S_n$具有相同的
极限，这个极限叫做$f(x)$从$a$到$b$的定积分$\Int^b_a f(x) \dd x$.

综上所述，在积分符号中，我们只对$a<b$, 即积分下限小于积分
上限的情形给出了定义，若$a>b$，我们定义
\[\int^a_b f(x) \dd x=-\int^b_a f(x) \dd x\]
此外，由于定积分可以解释为曲边梯形的面积，自然可以规定
\[\int^a_a f(x) \dd x=0\]
作了这样的规定之后，不论$a<b$, $a>b$或$a=b$, 定积分都有意义了．


\begin{ex}
\begin{enumerate}
    \item 求证 $\Int_{a}^{b} k x\dd x=\frac{k b^{2}}{2}-\frac{k a^{2}}{2} \quad(a<b)$
    \item  求 $\Int_{8}^{0}\left(x^{2}-4 x\right) \dd x$
    \item  求 $\Int_{-1}^{2}\left(x^{3}-3 x\right) \dd x$
\end{enumerate}
\end{ex}

\subsection{逼近法求曲线形的面积}

任意一条曲线围成的图形（图\ref{fig:chapter3_sec2_arbitrary_shape}）常常可以用两组互相垂直的直线把它分成若干部分，每一部分都是一个曲边梯形（图\ref{fig:chapter3_sec2_good_piece}）, 在这里并不排除下述情形（图\ref{fig:chapter3_sec2_degenerated_piece}）：两条平行的边中有一条缩成了一点，因而曲边梯形变成了曲边三角形，这样一来，我们的问题就化成了求曲边梯形面积的问题．

\begin{figure}[htp]
    \centering
    \begin{minipage}[t]{0.45\textwidth}
        \centering
        \begin{tikzpicture}[>=latex, scale=1]
    \draw(-2,3)--(2,3);
    \draw(-2,2)--(2,2); 
    \draw(-.5,0)--(-.5,4);
    \draw(.5,0)--(.5,4); 
    \node at (0,0){};
    \draw[thick] plot[smooth cycle] coordinates{
 (-0.4,3.75)(0.4,3.85)(1.0,3.6)(1.15,3.0)(0.95,2.1)(0.55,1.1)(0.0,0.62)(-0.8,0.62)(-1.35,0.95)(-1.3,1.7)(-1.0,2.6)
};
        \end{tikzpicture}
        \caption{}\label{fig:chapter3_sec2_arbitrary_shape}
        \end{minipage}
    \begin{minipage}[t]{0.25\textwidth}
    \centering
    \begin{tikzpicture}[>=latex, scale=1]
    \draw(0,1.3)--(0,0)--(2.5,0)--(2.5,1.8);
    \node at (0,0) {};
    \draw plot[smooth] coordinates{(0,1.3)(1.4,2)(2.5,1.8)};
    \end{tikzpicture}
    \caption{}\label{fig:chapter3_sec2_good_piece}
    \end{minipage}
    \begin{minipage}[t]{0.25\textwidth}
    \centering
    \begin{tikzpicture}[>=latex, scale=1]
    \draw(0,2)--(0,0)--(1.5,0);
    \node at (0,0){};
    \draw plot[smooth] coordinates{(0,2)(0.5,1.8)(1,1.3)(1.5,0)};
    \end{tikzpicture}
    \caption{}\label{fig:chapter3_sec2_degenerated_piece}
    \end{minipage}
  \end{figure}

\begin{blk}
  {命题} 设$y=f(x)$是定义在$a\le x\le b$的分段单调连续函数，而且$f(x)\ge 0$, 则区域$R=\{a\le x\le b,\; 0\le y\le f(x)\}$的
面积等于
\[\int^b_a f (x) \dd x \]  
\end{blk}

\begin{figure}[htp]
    \centering
\begin{tikzpicture}[>=latex]
    \draw[->, thick](-.5,0)--(7,0)node[right]{$x$};
    \draw[->, thick](0,-.5)--(0,5)node[right]{$y$};
\draw[domain=.5:6, samples=100, very thick]plot(\x, {1.5*sin(1.2*(\x-1) r)+3});
\draw(.5,0)node[below]{$a$} rectangle  (1.1, .18+3);
\draw(1.1,0) rectangle  (1.8, 1.23+3);
\draw(1.8,0) rectangle  (2.5, 1.46+3);
\draw[pattern=north east lines](2.5,0) rectangle  (3.5, .212+3);
\draw[pattern=north east lines](3.5,0) rectangle  (4.2, -.965+3);
\draw[pattern=north east lines](4.2,0) rectangle  (5, -1.49+3);
\draw(5,0) rectangle  (5.5, -1.16+3);
\draw(5.5,0) rectangle (6, -.419+3);

\draw[pattern=north east lines](.5,0) rectangle  (1.1, -.847+3);
\draw[pattern=north east lines](1.1,0) rectangle  (1.8, .18+3);
\draw[pattern=north east lines](1.8,0) rectangle  (2.5, 1.23+3);
\draw(2.5,0) rectangle  (3.5, 1.46+3);
\draw(3.5,0) rectangle  (4.2, .212+3);
\draw(4.2,0) rectangle  (5, -.965+3);
\draw[pattern=north east lines](5,0) rectangle  (5.5,-1.49+3);
\draw[pattern=north east lines](5.5,0) rectangle (6, -1.16+3);

\node at (6,0)[below]{$b$};
\node at (0,0)[below left]{$O$};
\node at (4.5,3.2)[above]{$y=f(x)$};

\end{tikzpicture}
    \caption{}\label{fig:chapter3_sec2_riemann_sum}
\end{figure}


\begin{proof}
    由$\Int^b_a f (x) \dd x$的定义得知，存在两系列阶梯函数
$\{g_n(x)\}$和$\{G_n(x)\}$, 满足下面的性质：
\[g_n(x)\le f (x) \le G_n(x)\]
而且$G_n(x)$从$a$到$b$的和与$g_n(x)$从$a$到$b$的和趋于共同的极限$\Int^b_a f (x) \dd x$. 令
\begin{align*}
    R_n&=\{a\le x\le b,\; 0\le y\le g_n(x)\}\\
    R'_n&=\{a\le x\le b,\; 0\le y\le G_n(x)\}
\end{align*}
它们都是由高高低低的狭长方形所组成的区域（图\ref{fig:chapter3_sec2_riemann_sum}），则：
\[R_n\subset \text{曲边梯形区域 }R \subset R'_n\]
而且
\[R_n\text{的面积}=g_n(x)\text{从$a$到$b$的和}s_n,\qquad R'_n\text{的面积}=G_n(x)\text{从$a$到$b$的和}S_n\]
即有
\[s_n<\text{曲边梯形区域$R$的面积}A<S_n\]
由定积分定义有：
\[\lim_{n\to\infty}s_n=\lim_{n\to\infty}S_n=\Int^b_a f (x) \dd x\]
又曲边梯形面积$A(f,\; [a,b])$也是夹逼在$s_n$和$S_n$中间的实数，所以
\[A(f,\; [a,b])=\Int^b_a f (x) \dd x\]
\end{proof}

\begin{ex}
\begin{enumerate}
    \item 求$\Int^1_{-2}e^x\dd x$. 

提示：$\Lim_{\Delta x\to 0}\frac{e^{\Delta x}-1}{\Delta x}=1$.    

\item 求在$x=0$与$x=\pi$间正弦曲线$y=\sin x$与$Ox$轴所包的面积．
\end{enumerate}
\end{ex}    

\subsection{定积分的基本性质}

为简单起见，我们约定以下被积函数在积分区间上连续并且可以分段单调，即$f(x)$是在$[a,b]$上连续并且只有有限个极大值和极小值的函数．

\begin{blk}{性质1}
若$a<b<c$，那么，
\[\int^b_a f(x)\dd x+\int^c_b f(x)\dd x=\int^c_a f(x)\dd x\]
\end{blk}

\begin{proof}
因为$\Int^b_a f(x)\dd x$存在，故对于在$[a,b]$上的
$f(x)$, 存在一组上、下夹逼函数列$\{g_n(x)\}$和$\{G_n(x)\}$, 使得
\[g_n (x) \le f (x) \le G_n (x),\qquad x\in [a,b]\]
且相应的阶梯函数在$[a,b]$上的和$s_n$与$S_n$适合
\begin{equation}
    s_n<\int^b_a f(x)\dd x<S_n,\qquad S_n-s_n\to 0
\end{equation}
对于在$[b,c]$上的$f(x)$, 同样得到
\[\bar g_n (x) \le f (x) \le \bar G_n (x),\qquad x\in [b,c]\]
且
\begin{equation}
    \bar  s_n<\int^c_b f(x)\dd x<\bar S_n,\qquad \bar S_n-\bar s_n\to 0
\end{equation}
$(3.4)+(3.5)$得到
\begin{equation}
    s_n+ \bar  s_n<\int^b_a f(x)\dd x+\int^c_b f(x)\dd x<S_n+\bar S_n
\end{equation}

另一方面，对于函数$f(x)\; (a\le x\le c)$存在下面一组阶梯函数列：
\[\tilde g_n(x)=\begin{cases}
    g_n(x),  &  a\le x<b\\
    \max\left(g_n(b), \bar g_n(b)\right),  &  x=b\\
    \bar g_n(x), & b<x\le c
\end{cases}\]
\[\tilde G_n(x)=\begin{cases}
    G_n(x),  &  a\le x<b\\
    \min\left(G_n(b), \bar G_n(b)\right),  &  x=b\\
    \bar G_n(x), & b<x\le c
\end{cases}\]

由性质2，阶梯函数列$\tilde g_n(x)$与$\tilde G_n(x)$的从$a$到$c$的和分别是$s_n+\bar s_n$和$S_n+\bar S_n$．因为$\tilde g_n(x)$与$\tilde G_n(x)$满足条件：
\begin{enumerate}
    \item $\tilde g_n(x)\le f(x)\le \tilde G_n(x),\quad a\le x\le c$
    \item 
    \begin{align*}
    \lim_{n\to\infty}\left[(S_n+\bar S_n)-(s_n+\bar s_n)\right]&=\lim_{n\to\infty}\left[(S_n-s_n)+(\bar S_n-\bar s_n)\right]\\
    &=\lim_{n\to\infty}(S_n-s_n)+\lim_{n\to\infty}(\bar S_n-\bar s_n)=0
    \end{align*}
\end{enumerate}
所以根据定积分定义得到
\[\lim_{n\to\infty}(S_n+\bar S_n)=\lim_{n\to\infty}(s_n+\bar s_n)=\int^c_a f(x)\dd x \]

因为$\Int^c_a f(x)\dd x $与$\Int^b_a f(x)\dd x +\Int^c_b f(x)\dd x $都是被$s_n+\bar s_n$与$S_n+\bar S_n$所夹逼的唯一实数，所以
\[\Int^c_a f(x)\dd x=\Int^b_a f(x)\dd x +\Int^c_b f(x)\dd x\]
\end{proof}

\begin{blk}{性质2}
若$f(x)=f_1(x)+f_2(x)$，则：
\[\Int^b_a f(x)\dd x=\Int^b_a f_1(x)\dd x +\Int^b_a f_2(x)\dd x\]
\end{blk}

\begin{proof}
设$\{g_n(x)\}$与$\{G_n(x)\}$是$f_1(x)$在$[a,b]$上的上、下夹逼阶梯函数列，又$\{\bar g_n(x)\}$与$\{\bar G_n(x)\}$是$f_2(x)$在$[a,b]$上的上、下夹逼阶梯函数列．$s_n, S_n, \bar s_n, \bar S_n$是相应的阶梯函数从$a$到$b$的和．于是由
$\Int^b_a f_1(x)\dd x$和$\Int^b_a f_2(x)\dd x$的存在，得：
\begin{align}
g_n(x)\le f_1(x)\le G_n(x),\qquad a\le x\le b\\
\bar g_n(x)\le f_2(x)\le \bar G_n(x),\qquad a\le x\le b
\end{align}
并且
\begin{align}
s_n<\Int^b_a f_1(x)\dd x<S_n,\quad \text{且 } S_n-s_n\to 0\\
\bar s_n<\Int^b_a f_2(x)\dd x<\bar S_n,\quad \text{且 } \bar S_n-\bar s_n\to 0
\end{align}
由$(3.9)+(3.10)$，得到
\begin{equation}
   s_n+ \bar s_n <\Int^b_a f_1(x)\dd x+\Int^b_a f_2(x)\dd x<S_n+\bar S_n
\end{equation}
另一方面，由$(3.7)+(3.8)$，得到
\[g_n(x)+\bar g_n(x)\le f_1(x)+ f_2(x)\le G_n(x)+\bar G_n(x),\qquad a\le x\le b\]
而且，$g_n(x)+\bar g_n(x)$与$G_n(x)+\bar G_n(x)$也是在$[a,b]$上的阶梯函数，我们要说明它们是$f_1(x)+ f_2(x)$在$[a,b]$上的一组夹逼阶梯函数列．

因为：
\begin{align*}
    g_n(x)+\bar g_n(x)\text{从$a$到$b$的和}&=\left[ g_n(x)\text{从$a$到$b$的和}\right]+\left[\bar g_n(x)\text{从$a$到$b$的和}\right]\\
    &=s_n+\bar s_n\\
    G_n(x)+\bar G_n(x)\text{从$a$到$b$的和}&=\left[ G_n(x)\text{从$a$到$b$的和}\right]+\left[\bar G_n(x)\text{从$a$到$b$的和}\right]\\
    &=S_n+\bar S_n
\end{align*}
而且
\begin{equation}
    \left(S_n+\bar S_n\right)-\left(s_n+\bar s_n\right)=\left(S_n-s_n\right)+\left(\bar S_n-\bar s_n\right)\to 0
\end{equation}
所以，由$\Int^b_a \bigl(f_1(x)+ f_2(x)\bigr)\dd x$的定义，得到
\begin{equation}
    \lim_{n\to\infty}\left(S_n+\bar S_n\right)=\lim_{n\to\infty}\left(s_n+\bar s_n\right)=\Int^b_a \bigl(f_1(x)+ f_2(x)\bigr)\dd x
\end{equation}
由(3.11)和(3.13)，$\Int^b_a \bigl(f_1(x)+ f_2(x)\bigr)\dd x$与$\Int^b_a f_1(x)\dd x + \Int^b_a  f_2(x)\dd x$
是被$s_n+\bar s_n$与$S_n+\bar S_n$所夹逼的唯一实数，所以
\[\Int^b_a \bigl(f_1(x)+ f_2(x)\bigr)\dd x=\Int^b_a f_1(x)\dd x + \Int^b_a  f_2(x)\dd x\]
\end{proof}

\begin{blk}{性质3}
\[\Int^b_a kf(x)\dd x =k\Int^b_a f(x)\dd x \]
\end{blk}

此法则的证明大致和性质2的证明相同，即当$\{g_n(x)\}$和$\{G_n(x)\}$上、下夹逼$f(x)$时，那么在$k>0$的情形，$\{kg_n(x)\}$和$\{kG_n(x)\}$上、下夹逼$kf(x)$; 在$k<0$的情形，$\{kG_n(x)\}$和$\{kg_n(x)\}$上、下夹逼$kf(x)$. 我们把证明的过程留给读者去写．

\begin{blk}{性质4}
    若$m\le f(x)\le M$, $a\le x\le b$，那么
    \[m(b-a)\le \Int^b_a f(x)\dd x \le M(b-a) \]
    \end{blk}

\begin{proof}
    设$\{g_n(x)\}$与$\{G_n(x)\}$是$f(x)$在$[a,b]$上的夹逼阶梯函数列，$s_n$与$S_n$是相应的阶梯函数在$[a,b]$上的和，根据积分的定义，有
\[\lim_{n\to\infty}s_n=\lim_{n\to\infty}S_n=\int^b_a f(x)\dd x\]
换言之，任给一个正数$\varepsilon$，存在自然数$N$，使得当$n>N$时，有
\begin{equation}
    S_n<\int^b_a f(x)\dd x+\varepsilon
\end{equation}
成立．

又由于$G_n(x)\ge f(x)\ge m,\quad x\in[a,b]$，根据性质1，得到：
\begin{equation}
    S_n\ge m(b-a)
\end{equation}
所以，当$n>N$时，由(3.14)和(3.15)，有
\[m(b-a)\le S_n<\int^b_a f(x)\dd x+\varepsilon\]
即
\begin{equation}
    m(b-a)<\int^b_a f(x)\dd x+\varepsilon
\end{equation}
成立．

因为这个不等式(3.16)对于每个正数$\varepsilon$都成立，所以
\[m(b-a)\le \int^b_a f(x)\dd x\]
同样证明，得到
\[M(b-a)\ge \int^b_a f(x)\dd x\]
因此
\[m(b-a)\le \int^b_a f(x)\dd x\le M(b-a)\]
\end{proof}

\begin{figure}[htp]
    \centering
\begin{tikzpicture}[>=latex]
\draw[->, thick](-1.5,0)--(3,0)node[right]{$x$};
\draw[->, thick](-1,-.5)--(-1,4)node[right]{$y$};
\node at (-1,0)[below left]{$O$};
\draw[domain=-.1:2.3, samples=50, very thick]plot(\x, {\x*(\x-1)*(\x-2)+2});
\draw(-.1,0)node[below]{$a$} rectangle (2.3,2.9);
\node at (2.3,0)[below]{$b$};
\draw(1.577,0)node[below]{$x_0$}--(1.577,1.615);
\draw(-.1,1.615)--(2.3,1.615);
\draw[dashed](-1,1.615)node[left]{$m$}--(-.1,1.615);
\draw[dashed](-1,2.9)node[left]{$M$}--(-.1,2.9);

\node at (1.1,2.5){$y=f(x)$};
\end{tikzpicture}
    \caption{}\label{fig:chapter3_sec2_y=f(x)_bounded}
\end{figure}

从几何图形来看，以曲线$y=f(x)$为一边而以线段$[a,b]$为底边的曲边梯形界于以$[a,b]$为底边，高分别等于$m$和$M$的两个矩形之内，
故曲边梯形的面积$\Int^b_a f(x)\dd x$
在两个矩形面积$m(b-a)$和$M(b-a)$之间（图\ref{fig:chapter3_sec2_y=f(x)_bounded}）．

这个性质是说分段单调的连续函数$f(x)$在区间$[a,b]$上的定积分是有界的．

\begin{blk}{性质5}
如果$f(x)\le \varphi(x),\; a\le x\le b$，那么：
\[\int^b_a f(x)\dd x\le \int^b_a \varphi(x)\dd x\]
\end{blk}

\begin{proof}
    设$\{g_n(x)\}$与$\{G_n(x)\}$是$f(x)$在$[a,b]$上的夹逼阶梯函数列，又$\{\bar g_n(x)\}$与$\{\bar G_n(x)\}$是$\varphi(x)$在$[a,b]$上的夹逼阶梯函数列，
又$s_n, S_n, \bar s_n, \bar S_n$是相应的阶梯函数从$a$到$b$的和．于是根据定积分的定义，有：
\[s_n<\int^b_a f(x)\dd x<S_n,\qquad \bar s_n<\int^b_a \varphi(x)\dd x<\bar S_n\]
当$n\to\infty$时，得到：
\begin{align*}
\lim_{n\to\infty}S_n =\lim_{n\to\infty}s_n =\int^b_a f(x)\dd x\\
\lim_{n\to\infty}\bar S_n =\lim_{n\to\infty}\bar s_n =\int^b_a \varphi(x)\dd x\\
\end{align*}
换言之，任给正数$\varepsilon$，存在$N_1$，使得当$n>N_1$时，有
\[s_n>\int^b_a f(x)\dd x-\frac{\varepsilon}{2}\]

当$n>N_2$时，有
\[\bar S_n<\int^b_a \varphi(x)\dd x+\frac{\varepsilon}{2}\]

因为$\bar G_n(x)>\varphi(x)>f(x)>g_n(x),\quad x\in[a,b]$，根据性质1，得到：$\bar S_n\ge s_n$.

所以，当$n>\max(N_1,N_2)$时，有
\[\int^b_a f(x)\dd x<s_n+\frac{\varepsilon}{2}<\bar S_n+\frac{\varepsilon}{2}<\int^b_a \varphi(x)\dd x+\varepsilon\]
因为这个不等式对于每个正数$\varepsilon$都成立，所以必然有
\[\int^b_a f(x)\dd x\le \int^b_a \varphi(x)\dd x\]
\end{proof}

\begin{blk}{性质6：定积分中值定理}
   设函数$f(x)$在闭区间$[a,b]$上分段单调连续，又
$m=f(c)$, $M=f(d)$分别是$f(x)$在$[a,b]$上的最小值和最大值，则在$[a,b]$上至少存在一点$\xi$, 使得下面的等式成立：
\[\int^b_a f (x) \dd x=f(\xi) (b-a)\] 
\end{blk}

\begin{proof}
    因为$m\le f (x) \le M,\quad x\in [a,b]$, 
由性质4, 即得
\[m (b-a)\le \int^b_a f (x) \dd x\le M(b-a)\]
上面不等式的两端各除以$(b-a)$, 得
\[f(c)=m\le \frac{\int^b_a f (x) \dd x}{b-a}\le M=f(d)\]
因为$f(x)$在$[a,b]$上连续，再由连续函数的中间值定理，必存在一个$\xi\in(c,d)$, 使得
\[f(\xi)=\frac{1}{b-a}\int^b_a f (x) \dd x\]
两边再乘以$(b-a)$，得
\[\int^b_a f (x) \dd x=f(\xi)(b-a)\]
这就是我们所要证明的．
\end{proof}

这个公式的几何意义是：以线段$[a,b]$为底边，以曲线$y=f(x)$为曲边的曲边梯形，它的面积等于同一底边
而高为$f(\xi)$的一个矩形的面积（图\ref{fig:chapter3_sec2_integral_mean_value}）．因此，$f(\xi)$称为曲边梯形的平均高度．

我们也称
\[f(\xi)=\frac{1}{b-a}\int^b_a f (x) \dd x\]
为$f(x)$在$[a,b]$上的平均值．

\begin{figure}[htp]
    \centering
    \begin{minipage}[t]{0.48\textwidth}
    \centering
    \begin{tikzpicture}[>=latex, xscale=.8]
\draw[->, thick](-.5,0)--(5,0)node[right]{$x$};
\draw[->, thick](0,-.5)--(0,4)node[right]{$y$};
\node at (0,0)[below left]{$O$};
\draw[domain=.7:4, very thick, samples=100]plot(\x, {ln(\x)+1.5})node[above]{$y=f(x)$};
\draw[dashed](.7,0)node[below]{$a$} rectangle (4,2.3);
\draw[dashed](2.23,0)node[below]{$\xi$}--(2.23,2.3)node[below right]{$f(\xi)$};
\draw[dashed](4,2.89)--(4,2.3);
\node at (4,0)[below]{$b$};

    \end{tikzpicture}
    \caption{}\label{fig:chapter3_sec2_integral_mean_value}
    \end{minipage}
    \begin{minipage}[t]{0.48\textwidth}
    \centering
    \begin{tikzpicture}[>=latex, scale=.5]
  \draw[->, thick](-2,0)--(6,0)node[right]{$x$};
\draw[->, thick](0,-8.5)--(0,2)node[right]{$y$};
\node at (0,0)[below left]{$O$};
\draw[domain=-1:4, very thick, samples=100]plot(\x, {2*\x-\x*\x});
\fill[pattern=north east lines, domain=2:4, very thick, samples=100]plot(\x, {2*\x-\x*\x})--(4,0)node[above]{4}
--(2,0)node[above]{2};
\draw(1,0)node[below]{1}--(1,.2);
    \end{tikzpicture}
    \caption{}\label{fig:chapter3_sec2_example3.7}
    \end{minipage}
  \end{figure}

\begin{example}
求$\Int^4_2 (2x-x^2)\dd x$.
\end{example}


\begin{solution}
    根据本节的定积分计算法则得：
\begin{align*}
    \Int^4_2 (2x-x^2)\dd x&=\Int^4_2 2x\dd x+\Int^4_2 (-x^2)\dd x\\
    &=2\Int^4_2 x\dd x-\Int^4_2 x^2\dd x
\end{align*}
利用例3.4的计算结果，得
\[\Int^4_2 x\dd x=\frac{4^2-2^2}{2}=6,\qquad \Int^4_2 x^2\dd x=\frac{4^3-2^3}{3}=18\frac{2}{3}\]
因此
\[\Int^4_2 (2x-x^2)\dd x=2\x 6-18\frac{2}{3}=-6\frac{2}{3}\]
它的几何意义是图\ref{fig:chapter3_sec2_example3.7}中阴影区域的有号面积．
\end{solution}

\begin{example}
  设$f(x)=\begin{cases}
      x^2,& 0\le x\le 1\\
      1,& 1<x\le 2
  \end{cases}$,  求$f(x)$在$[0,2]$上的平均值.
\end{example}

\begin{solution}
由平均值定义，再利用性质1，有    
\begin{align*}
    f(\xi)&=\frac{1}{b-a}\int^b_a f(x)\dd x=\frac{1}{2}\int^2_0 f(x)\dd x\\
    &=\frac{1}{2}\left[\int^1_0 f(x)\dd x+\int^2_1 f(x)\dd x\right]\\
    &=\frac{1}{2}\left[\int^1_0 x^2\dd x+\int^2_1 1\dd x\right]\\
    &=\frac{1}{2}\left[\frac{1}{3}+1\right]=\frac{2}{3}
\end{align*}
\end{solution}

\begin{ex}
\begin{enumerate}
    \item 利用以前定积分的结果和定积分的性质，计算下列定积分：
\begin{enumerate}
    \item $\Int^3_{-1}(2x^2-4x)\dd x$
    \item 若$f(x)=\begin{cases}
        x,& 0\le x<1\\
        x-2,& 1\le x\le 2
    \end{cases}$，则$\Int^2_0 f(x)\dd x=?$
    \item 若$f(x)=\begin{cases}
        x,& 0\le x\le 1\\
        x-2,& 1< x\le 2
    \end{cases}$，则$\Int^2_0 f(x)\dd x=?$
\item 若$g(t)=2t^2+|t|-1,\; -1\le t\le 1$，则$\Int^1_{-1} g(t)\dd t=?$
\end{enumerate}

\item 将图中阴影部分的面积用定积分表示．
\item 设$\ell(t)=mt+b$, （$m,b$是常数），$t\in [c,d]$，

求证：$\Int^d_c (mt+b)\dd t=\frac{\ell(c)+\ell(d)}{2}(d-c)$
\item 证明：
\begin{enumerate}
    \item 若$0<x<10$，则$\frac{1}{1016}\le\frac{1}{x^3+16}\le \frac{1}{16}$
    \item $\frac{5}{508}\le \Int^{10}_0\frac{1}{x^3+16}\dd x\le \frac{5}{8}$
\end{enumerate}
\end{enumerate}
\end{ex}

\begin{figure}[htp]
    \centering
\begin{tikzpicture}[>=latex]
    
\begin{scope}[scale=.6]
\draw[->](-2,0)--(4,0)node[right]{$x$};
\draw[->](0,-3)--(0,7)node[right]{$y$};
\fill[domain=-1:3, samples=50, pattern=north east lines]plot(\x,{2*(\x-1)^2-2})--(3,0)--(-1,0)--(-1,6);
\draw[domain=-1:3, samples=50, very thick]plot(\x,{2*(\x-1)^2-2})node[above]{$y=2x^2-4x$};
\node at (0,0)[above right]{$O$};
\node at (-1,0)[below]{$-1$};
\node at (3,0)[below]{$3$};
\node at (-1,3)[left]{$x=-1$};
\node at (3,3)[right]{$x=3$};
\node at (1,-2)[below]{$(1,-2)$};
\node at (1,-3.8){(i)};
\end{scope}
\begin{scope}[xshift=5cm, yshift=-.8cm, scale=.9]
\draw[->](-1,0)--(4,0)node[right]{$x$};
\draw[->](0,-1)--(0,5.5)node[right]{$y$};
\node at (0,0)[below left]{$O$};
\draw[domain=0:2.2, samples=50, very thick]plot(\x,{\x^2})node[above]{$y=x^2$};
\draw[thick](0,0)--(3.5,3.5)node[right]{$y=x$};
\node at (1,0)[below]{$1$};
\node at (2,0)[below]{$2$};
\draw[dashed](0,1)node[left]{1}--(1,1)--(1,0);
\fill[domain=0:1, samples=50, pattern=north east lines]
plot(\x,{\x^2})--(2,2)--(2,0)--(0,0);
\node at (2,-1.5){(ii)};
\end{scope}

\end{tikzpicture}
    \caption*{第2题}
\end{figure}



\subsection*{习题3.2}
\begin{enumerate}
    \item 计算下列定积分：
\begin{multicols}{3}
\begin{enumerate}
    \item $\Int^1_4 |x|\dd x$
    \item $\Int^2_4 |x|\dd x$
    \item $\Int^4_1 |x|\dd x$
\end{enumerate}
\end{multicols}
\item 求$\Int^b_a |x|\dd x$.

提示：分$a<b<0$, $a<0<b$, $0<a<b$三种情况讨论．
\item 证明：
\begin{enumerate}
    \item 若$\frac{\pi}{4}\le x\le\frac{\pi}{2}$，则$\frac{2}{\pi}\le \frac{\sin x}{x}\le \frac{2\sqrt{2}}{\pi}$
    \item $\frac{1}{2}\le \Int^{\tfrac{\pi}{2}}_{\tfrac{\pi}{4}}\frac{\sin x}{x}\dd x\le \frac{\sqrt{2}}{2}$
\end{enumerate}
\item 已知作用在作直线运动的质点上的力是$F=s^2+1$, 试求从距离1到10之间所作的功．
\item 求$y=\begin{cases}
    A\sin\frac{2\pi t}{T}, & 0\le t\le \frac{T}{2}\\
    0,& \frac{T}{2}\le t\le T
\end{cases}$ 在$[0,T]$上的平均值．
\end{enumerate}
\chapter{微积分学基本定理}
在前两章中，我们分别引入了函数的变率（导数），函数的和（定积分）的基本概念，本章将研究函数的导函数与函数的求和函数这两者之间的互逆关系，并说明我们可以用求导函数的逆运算方法来计算定积分．

\section{微积分学基本定理}
\subsection{导函数与求和函数}
函数$f(x)$在点$x_0$处的变化率（导数）的定义是
\[f' (x_0) =\lim_{h\to 0}\frac{f (x_0+h) -f (x_0)}{h}\]
显然， $f'(x_0)$的值与$f(x)$在点$x_0$的值以及在点$x_0$的邻近的函数值有关，当点$x_0$在$(a,b)$内变化时，$f'(x_0)$也跟着变化，那么$f'(x_0)$便是一个新函数称为$f(x)$的导函数．计算一个函数的导函数是一件比较简便的事情．

定积分$\int^b_a f(x)\dif x$的定义是把区间$[a,b]$无限细分
而得到上下夹逼阶梯函数的和的共同极限，其几何意义是曲线$y=f(x)$和直线$x=a$, $x=b$, 及$y=0$所围成的区域的有号面积．

假如我们考虑$f(x)$在一个变动的区间$[a,x]$上的和，
即让区间的左端点固定，右端点变动，则
\[S_f(x)=\text{函数$f$从$a$到$x$的和}=\int^x_a f (x) \dif x\]
可以看作上限变量$x$的函数，在这里，积分符号中的$x$既表示积分变量，又表示积分上限，容易混淆，因此，为了区别起见，我们用字母$t$来代表积分变量，这样上式就写成
\[S_f (x)=\int^x_a f(t)\dif t\]

和函数$S_f(x)$在$x=x_0$处的值$S_f(x_0)$的几何意义就是曲线$y=f(t)$, 直线$t=a$, $t=x_0$, $y=0$所围成的区域的有号面积，它是随区域的变动界线$t=x_0$的变动而变动的（图4.1）．

\begin{figure}[htp]
    \centering
    \begin{tikzpicture}[>=latex]
        \draw[->](-3,0)--(5,0)node[right]{$t$};
        \draw[->](0,-1.5)--(0,2)node[right]{$y$};
    \node at (0,0) [below left]{$O$};
    \node at (0,0){};
        \draw[line width=1.5pt] plot[smooth] coordinates{ (-2.75,-.5)(-2,-1)(-.8,0)(1,1.25)(2.5,0)(3.5,-1.1)(4,-1.3) 
    };
    \draw(-2,-1)node[below]{$M'$}--node[left]{$t=x_0'$}(-2,0)node[above]{$P'$};
    \draw(1,1.25)node[above]{$B$}--node[right]{$t=a$}(1,0)node[below]{$A$};
    \draw(3.5,-1.1)node[below]{$M$}--node[right]{$t=x_0$}(3.5,0)node[above]{$P$};
    \node at (-.8,0)[above]{$D$};
    \node at (2.5,0)[above]{$C$};
    \node at (2,1)[right]{$y=f(t)$};
    
    \end{tikzpicture}
    \caption{}
\end{figure}

例如，当变动界限（积分上限）在图中的$PM$位置时，则
\[S_f(x_0)=\int^{x_0}_a f(t)\dif t=ACB\text{的面积}-CPM\text{的面积}\]

当变动界限在图中的$P'M'$位置时，则
\[\begin{split}
    S_f(x'_0)&=\int^{x_0'}_a f(t)\dif t=-\int^a_{x_0'} f(t)\dif t\\
    &=-\left\{-P'DM'\text{的面积}+DAB\text{的面积}  \right\}\\
    &=P'DM'\text{的面积}-DAB\text{的面积}
\end{split}\]

例如，折线函数
\[g(x)=\begin{cases}
    \frac{1}{2}(x-1) & x\in[1,2]\\
    \frac{3}{2}(x-2)+\frac{1}{2}(2-1) & x\in [2,3]\\
    (x-3)+\frac{1}{2}(2-1)+\frac{3}{2}(3-2) & x\in [3,4]
\end{cases}\]
的导函数（除去在折线段的那些交接点处不作定义外）是阶梯函数
\[g'(x)=\begin{cases}
    \frac{1}{2} & x\in [1,2)\\
    \frac{3}{2} & x\in (2,3)\\
    1& x\in (3,4]
\end{cases}\]
反过来，该阶梯函数的和函数，是上述折线函数$g(x)$, 我们有如下的图解关系：
\begin{center}
    \begin{tikzpicture}[>=latex]
        \node (A) at (0,0) {\{折线函数\}};
        \node (B) at (5,0) {\{阶梯函数\}};
\draw[->](A) to [bend left=30]node[above]{求导} (B);
\draw[->](B) to [bend left=30]node[above]{求和} (A);

    \end{tikzpicture}
\end{center}

上述简明的例子表明“微分”与“积分”（或求函数由$a$到$b$的和）之间的运算关系应该是互逆的．

\subsection{微积分学基本定理}

\begin{blk}
  {定理1} 设$f(t)$是在$[a,b]$上的分段单调连续函数，又它的和函数是
\[S_f(x)=\int^x_a f (t) \dif t\]
那么
\[\frac{\dif }{\dif x}S_f(x)=\frac{\dif }{\dif x}\int^x_a f (t) \dif t=f(x),\qquad a\le x\le b\]
\end{blk}

\begin{figure}[htp]
    \centering
    \begin{tikzpicture}[>=latex, xscale=3]
        \draw[->](.4,0)--(3,0)node[right]{$t$};
        \draw[->](.6,-.5)--(.6,4.5)node[right]{$y$};
    \draw[domain=.9:2.5, very thick, samples=100]plot(\x, {2.5*(\x-1)*(\x-2)*(\x-3)+2})node[right]{$y=f(t)$};
    \draw(.9,0)node[below]{$a$}--(.9,1.4225);
    \draw[pattern=north east lines, domain=1.4:1.7, samples=50]plot(\x, {2.5*(\x-1)*(\x-2)*(\x-3)+2})--(1.7,2.6825)--(1.7,0)node[below]{$x$}--(1.4,0)node[below]{$x+\bar h$}--(1.4,2.96);
    \draw[pattern=north west lines, domain=1.7:2, samples=50]plot(\x, {2.5*(\x-1)*(\x-2)*(\x-3)+2})--(2,2)--(2,0)node[below]{$x+h$}--(1.7,0)--(1.7,2.6825);
    \draw[<-](1.55,2)--(1.55,3.5)node[right]{$S_f(x)-S_f(x+\bar h),\; (\bar h<0)$};
    \draw[<-](1.85,.5)--(2.1,.5)node[right]{$S_f(x+h)-S_f(x),\; (h>0)$};
    \node at (.6,0)[below left]{$O$};
    \end{tikzpicture}
    \caption{}
\end{figure}

\begin{proof}
  如图4.2，设$h>2$，则
\[\begin{split}
    \frac{S_f(x+h)-S_f(x)}{h}&=\frac{\Int^{x+h}_a f(t)\dd t-\Int^{x}_a f(t)\dd t}{h}\\
    &=\frac{\Int^{x+h}_x f(t)\dd t}{h}\qquad \text{（定积分性质1）}\\
    &=\frac{hf(\xi)}{h}\quad (x<\xi<x+h)\\
    &=f(\xi)\qquad \text{（定积分中值定理）}
\end{split}\]
于是
\begin{equation}
    \lim_{h\to 0^+}\frac{S_f(x+h)-S_f(x)}{h}=\lim_{\xi\to x}f(\xi)=f(x)
\end{equation}

设$\bar h<0$，则
\[\begin{split}
    \frac{S_f(x)-S_f(x+\bar h)}{\bar h}&=\frac{\Int^{x}_a f(t)\dd t-\Int^{x+\bar h}_a f(t)\dd t}{\bar h}\\
    &=\frac{\Int^{x}_{x+\bar h} f(t)\dd t}{\bar h}\qquad \text{（定积分性质1）}\\
    &=\frac{-\bar h f(\bar\xi)}{\bar h}=-f(\bar\xi)\quad (x+\bar h<\bar \xi<x)\\
\end{split}\]
即：
\[\frac{S_f(x+\bar h)-S_f(x)}{\bar h}=f(\bar\xi),\qquad (x+\bar h<\bar \xi<x)\]
于是
\begin{equation}
    \lim_{h\to 0^-}\frac{S_f(x+\bar h)-S_f(x)}{\bar h}=\lim_{\bar\xi\to x}f(\bar\xi)=f(x)
\end{equation}
由(4.1)和(4.2)，得
\[\frac{\dd }{\dd x}S_f(x)=\frac{\dd }{\dd x}\int^x_a f(t)\dd t=f(x)\]
\end{proof}

\begin{blk}
    {引理} 如果两个函数$F(x)$和$G(x)$满足条件
\[F' (x) =G' (x) ,\qquad a\le x\le b\]
那么
\[G(x)=F(x)+C,\qquad (C\text{是常数})\]
\end{blk}

\begin{proof}
    由$G'(x)-F'(x)=0$, 得
\[[G (x) -F (x) ]'=0\]
根据中值定理（参看第三章）知，对于每一个$x\in (a,b)$, 恒有
\[G(x)-F(x)=C,\qquad (C\text{是常数})\]
所以
\[G (x) =F (x) +C\]
\end{proof}

\begin{blk}
 {定义} 设$f(x)$是在区间$[a,b]$上给定的函数，如果可微函数$G(x)$满足条件
\[G'(x)=f(x)\qquad (\text{或}\quad \dd G(x)=f(x)\dd x)\]
那么就称$G(x)$为$f(x)$在$[a,b]$上的一个\textbf{原函数}．
\end{blk}

例如，设$G'(x)=\cos x$, 那么$G(x)=\sin x$就称为$\cos x$的一个原函数．

\begin{blk}
    {定理2} 设$f(x)$是给定的分段单调连续函数，如果一个可微函数$G(x)$满足条件
\begin{equation}
    G'(x)=f(x)\qquad \text{（或$\dd G(x)=f(x)\dd x$）},\quad a\le x\le b
\end{equation}    
那么
\[S_f (x)=\int^x_a  f (t) \dd t=G (x) -G (a) \]
\end{blk}

\begin{proof}
    由定理1知，$f(x)$的一个原函数存在，就是它的和函数，即
\[\frac{\dd }{\dd x}\int^x_a f(t)\dd t=f(x)\]
题设$G(x)$是$f(x)$的另一个原函数，故由引理得
\begin{equation}
    \int^x_a    f (t) dt=G (x) +C 
\end{equation}
这里的常数要由条件$\Int^a_a f(t)\dd t=0$来确定．

在(4.4)中，令$x=a$，得
\[0=G(a)+C\]
所以：$C=-G(a)$，代入(4.4)，得
\[\int^x_a f(t)\dd t=G(x)-G(a)\]
\end{proof}

在上面等式中，令$x=b$, 便得到下面著名的莱布尼兹-牛顿公式：
\[\int^b_a f (t) \dd t= G (b) -G (a) \]
这个公式告诉我们，求分段单调连续函数$f(x)$在$[a,b]$上的定积分可以化简为去求$f(x)$的一个原函数$G(x)$, 而函数值的差$G(b)-G(a)$就是定积分的值．

\begin{example}
    设$f(x)=k$（常数函数），则$G(x)=kx$是一个满足$G'(x)=k=f(x)$的原函数，所以
\[\int^b_a k\dd x=kb-ka=k(b-a)\]
以后我们引用一个符号：
\[G(x)\Bigg|^b_a=G(b)-G(a)\]
\end{example}

\begin{example}
    设$f(x)=kx^n,\; (n\in\mathbb{N})$，则
\[G(x)=k\frac{x^{n+1}}{n+1}\]
是满足$G'(x)=kx^n$的一个原函数，所以
\[\int^b_akx^n\dd x=k\frac{x^{n+1}}{n+1}\Bigg|^b_a=\frac{k}{n+1}(b^{n+1}-a^{n+1})\]
\end{example}

\begin{example}
设$f(x)=x^{-n}=\frac{1}{x^n}\; (x\ne 0,\; n\in\mathbb{N})$，则
\[G(x)=\frac{x^{-n+1}}{-n+1}\]
是满足$G'(x)=x^{-n}$的一个原函数．要特别注意定理2中$f(x)$是连续的条件，在任何包含0的区间上，函数$f(x)=\frac{1}{x^n}$无界，即有第二类间断点，如果$a,b$异号，则定积分$\Int^b_a x^{-n}\dd x$无意义，但若$a,b$同时为正或同时为负，则
\[\int^b_ax^{-n}\dd x=\frac{x^{-n+1}}{-n+1}\Bigg|^b_a=\frac{b^{-n+1}}{-n+1}-\frac{a^{-n+1}}{-n+1}\]
当然，又有$n\ne -1$时，上式才能成立．

如果$n=-1$，依复合函数求导法则，得
\[(\ln|x|)'=\frac{1}{|x|}\cdot |x|'\]
若$x>0$，则
\[(\ln|x|)'=\frac{1}{x}\cdot 1=\frac{1}{x}\]
若$x<0$，则
\[(\ln|x|)'=\frac{1}{-x}\cdot (-1)=\frac{1}{x}\]
总之
\[(\ln|x|)'=\frac{1}{x},\quad (x\ne 0)\]
因此
\[\begin{split}
\int^b_a \frac{1}{x}\dd x=\ln |x|\Bigg|^b_a &=\ln|b|-\ln |a|\\
&=\ln\left|\frac{b}{a}\right|=\ln\frac{b}{a}\quad (x\ne 0,\; ab>0)
\end{split}\]
\end{example}

\begin{example}
    求$\Int^3_1 |x-2|\dd x$
\end{example}

\begin{solution}
\[\begin{split}
    \Int^3_1 |x-2|\dd x&=\Int^2_1 -(x-2)\dd x+\Int^3_2 (x-2)\dd x\\
&=\left(-\frac{x^2}{2}+2x\right)\Bigg|^2_1+\left(\frac{x^2}{2}-2x\right)\Bigg|^3_2\\
&=\frac{1}{2}+\frac{1}{2}=1
\end{split}\]
\end{solution}

\begin{example}
求$\Lim_{n\to\infty}\frac{\sum^n_{i=1}i^p}{n^{p+1}}\quad (p\in \mathbb{N})$
\end{example}


\begin{solution}
\[\begin{split}
    \Lim_{n\to\infty}\frac{\sum^n_{i=1}i^p}{n^{p+1}}&=\Lim_{n\to\infty}\sum^n_{i=1}\left(\frac{i}{n}\right)^p\cdot \frac{1}{n}\\
    &=\int^1_0 x^p\dd x=\left.\frac{x^{p+1}}{p+1}\right|^1_0=\frac{1}{p+1}
\end{split}\]
\end{solution}

\begin{ex}
\begin{enumerate}
\item 用求原函数的方法, 计算下列定积分:
\begin{multicols}{2}
\begin{enumerate}
 \item $\Int_{2}^{3}\left(x^{3}+6 x^{2}+4 x+1\right) \dd x$
\item $\Int_{1}^{5} \frac{v^{2}+2}{v^{4}} \dd v$
\item $\Int_{-3}^{-2}\left(x+\frac{2}{x}\right)^{2} \dd x$
\item $\Int_{1}^{10}\left(x+\frac{1}{x}\right) \dd x$
\item $\Int_{a}^{b}|x-1| \dd x$
\end{enumerate}
\end{multicols}
\item 求曲线 $y=2 x-x^{2}$ 与 $x$ 轴之间的面积.
\item 设 $y=x(x-1)(x-2)$，
求 $\Int_{0}^{1} y \dd x,\; \Int_{1}^{2} y\dd x,\; \Int_{0}^{2} y \dd x$,
并说明所得结果的几何意义．

\item 求直线$y=2x$和曲线$y=x^2$之间的面积．
\item 求由曲线$y=\frac{1}{x^2}$和直线$x=1$, $x=4$所围成的曲边梯形的
面积．
\item 求抛物线方程使它具有以下的性质：
\begin{enumerate}
\item 抛物线通过$(0, 0)$, $(1, 2)$两点，
\item 它的对称轴平行$y$轴，且向上凸，
\item 它与$x$轴所围的面积最小．
\end{enumerate}

\item 求正数$a$, 使$\Int^1_0 |x^2-a^2|\dd x$最小．
\item 求下列函数的导数：
\begin{multicols}{2}
    \begin{enumerate}
     \item $F(x)=\Int_{a}^{x}\sin^2 t \dd t$
    \item $F(x)=\Int_{a}^{x^2} \sin^2 t \dd t$
    \end{enumerate}
    \end{multicols}
\end{enumerate}

\end{ex}

\section{不定积分}
由上节的微积分基本定理得知，在$[a,b]$上的两个可
微函数$G_1(x)$和$G_2(x)$均以$f(x)$为其导数，则$G_1(x)$与$G_2(x)$在$[a,b]$上只差一个常数$C$. 由此可见，如果$G'(x)=f(x)$, 且$C$为任意常数，那么$[G(x)+C]'=G'(x)=f(x)$. 这就是说：如果 函数$f(x)$有一个原函数$G(x)$, 那么就有无穷多个原函数，并且所有原函数刚好组成函数族
\[\{G(x)+C\;|\; C\text{是任意常数}\}\]

\begin{blk}
   {定义} 函数$f(x)$在区间$[a,b]$上的原函数全体叫做$f(x)$的不定积分，记作
   \[\int f(x)\dd x\]
   设$G(x)$是$f(x)$的一个原函数，那么由不定积分的定义得到
\[\int f (x) \dd x=G (x) +C\]
   其中$C$是任意常数．
\end{blk}   

关于不定积分，易见有如下性质：

\begin{blk}{性质1}
   积分与微分互为逆运算，即：
\[\begin{cases}
    \dd\Int f(x)\dd x=f(x)\dd x\\
    \Int \dd f(x)=f(x)+C
\end{cases}\]
\end{blk}

事实上，因为$\frac{\dd }{\dd x}\Int^x_a f(t)\dd t=f(x)$，所以
\[\int f(x)\dd x=\int^x_a f(t)\dd t+C\qquad \text{（基本定理1和不定积分定义）}\]
于是
\[\left(\int f(x)\dd x\right)'=\left(\int^x_a f(x)\dd x\right)'=f(x)\]
从而
\[\dd\int f(x)\dd x=f(x)\dd x\]
又
\[\begin{split}
    \int \dd f(x)&=\int f'(x)\dd x=\int^x_a f'(t)\dd t+C_1\\
    &=f(x)-f(a)+C_1\quad \text{（基本定理）}\\
    &=f(x)+C
\end{split}\]
这里$C=-f(a)+C_1$. 

从性质1看出，若不要不定积分等式中的任意常数项，则当符号$\Int$与$\dd$紧接着时，无论哪个在前，哪个在后，都可以消掉．

因为微分，积分是两种互逆的运算，所以它们的运算法则是互相对应的，现将它们相对应的法则叙述如下：

\begin{blk}{性质2}
 \[\begin{cases}
     {\rm d} [f(x)+g(x)]=\dd f(x)+\dd g(x)\\
     \Int [f(x)+g(x)] \dd x=\Int f(x)\dd x+\Int g(x)\dd x+C
 \end{cases}\]
 \end{blk}

 上面不定积分公式的正确性不难证明，只要求出两边的微商，由于所得到的微商恒等，就可以知道它们只差一个常数项，因为，由性质1
\[\begin{split}
    \left(\int[f(x)+g(x)]\dd x\right)'&=f(x)+g(x)\\
    \left(\int f(x)\dd x+\int g(x)\dd x\right)'&=\left(\int f(x)\dd x\right)'+\left(\int g(x)\dd x\right)'=f(x)+g(x)
\end{split}\]
所以$f(x)+g(x)$的两个原函数相差一个常数项，即：
\[\int[f(x)+g(x)]\dd x=\int f(x)\dd x+\int g(x)\dd x+C\]
同理可得：

\begin{blk}{性质3}
\[\begin{cases}
    {\rm d} [a f(x)]=a\dd f(x)\\
    \Int af(x)\dd x=a\Int f(x)\dd x
\end{cases}\qquad (a\ne 0)\]
\end{blk}

\subsection{基本积分表}

由基本微商公式表，倒转顺序，就可以得到下面这个表：
\begin{itemize}
    \item $\frac{\dd}{\dd x}C=0,\qquad \Int 0\dd x=C$
    \item $\frac{\dd }{\dd x}x^{\mu}=\mu x^{\mu-1},\qquad \Int x^{\mu}\dd x=\frac{x^{\mu+1}}{\mu+1}+C\quad \text{（$\mu$是不等于$-1$的实数）}$
    \item $\frac{\dd }{\dd x}\ln|x|=\frac{1}{x},\qquad \Int\frac{\dd x}{x}=\ln|x|+C\quad (x\ne 0)$
    \item $\frac{\dd }{\dd x}e^x=e^x,\qquad \Int e^x\dd x=e^x+C$
\end{itemize}
还因为$a^x=e^{x\ln a}$，由复合函数求导法则，得：
\[(a^x)'=e^{x\ln a}(x\ln a)'=\ln a\cdot e^{x\ln a}=\ln a\cdot a^x\]
所以：
\begin{itemize}
    \item $\frac{\dd}{\dd x}a^x=\ln a\cdot a^x,\qquad \Int a^x\dd x=\frac{a^x}{\ln a}+C$
    \item $\frac{\dd}{\dd x} \sin x=\cos x, \qquad \Int \cos x \dd x=\sin x+C$.
    \item $\frac{\dd}{\dd x} \cos x=-\sin x, \qquad \Int \sin x \dd x=-\cos x+C$.
    \item  $\frac{\dd}{\dd x}\tan  x=\sec ^{2} x, \qquad \Int \sec^{2} x \dd x=\tan x+C$
    \item $\frac{\dd}{\dd x}\cot x=-\csc ^{2} x, \qquad \Int \csc^{2} x \dd x=-\cot x+C$
    \item $\frac{\dd}{\dd x} \arcsin x=\frac{1}{\sqrt{1-x^{2}}} ,\qquad \Int \frac{\dd x}{\sqrt{1-x^{2}}}=\arcsin x+C$ 
    
    $\frac{\dd}{\dd x} \arccos x=\frac{-1}{\sqrt{1-x^{2}}},\qquad \Int\frac{-1}{\sqrt{1-x^{2}}}\dd x=-\arccos x+C'$
    \item  $\frac{\dd}{\dd x}\arctan x=\frac{1}{x^{2}+1}, \qquad \Int \frac{\dd x}{x^{2}+1}=\arctan x+C$ 
    
    $\frac{\dd}{\dd x} \arccot x=\frac{-1}{x^{2}+1}, \qquad \Int \frac{-1}{x^{2}+1}\dd x=-\arccot  x+C'$
\end{itemize}

现在来看不定积分的几何意义．

由不定积分的定义知道，求函数$f(x)=3x^2$的原函数全体，就是求不定积分$\Int 3x^2\dd x$, 也就是求满足微分方程
\begin{equation}
    \frac{\dd y}{\dd x}=3x^2
\end{equation}
的所有的原函数．

显然，$y=x^3$就是微分方程(4.5)的一个原函数，它的图象上各点切线的斜率由(4.5)确定．我们称曲线$y=x^3$为
微分方程(4.5)的积分曲线，把它沿$y$轴上下平移任意一段距离，便得到曲线族
\begin{equation}
    y=x^3+C
\end{equation}
并且(4.6)中任意一条曲线与曲线$y=x^3$在横坐标相同的点的切线平行，其斜率都等于$3x^2$, 因此，曲线族(4.6)是满足微分方程(4.5)的所有积分曲线，也就是不定积分$\Int 3x^2\dd x=x^3+C$的几何表示（图4.3）．

为要确定(4.5)的某条积分曲线的位置，应当给出这条积分曲线必须经过的一点，例如它经过点$P(1, 3)$. 我们说这条件也是微分方程(4.5)的一个相应的原函数要满足的初值条件
$y\big|_{x=1}=3$
把它们代入(4.6), 得到
\[3=1^3+C\quad \Rightarrow\quad C=2\]
因此，满足微分方程$\frac{\dd y}{\dd x}=3x^2$和初值条$y\big|_{x=1}=3$的原函
数或相应的积分曲线是
\[y=x^3+2\]

通常把满足微分方程
\begin{equation}
\frac{\dd y}{\dd x}=f(x)
\end{equation}
的所有解叫做给定函数$f(x)$的不定积分，即
\[y=\int f(x)\dd x\]
把函数$f(x)$的一个原函数$y=G(x)$的图象叫做微分方程(4.7)的一条\textbf{积分曲线}，这样，不定积分
\[y=\int f(x)\dd x=G(x)+C \quad \text{（$C$为任意常数）}\]
的几何意义就表示与$y=G(x)$在各点有相同斜率$\frac{\dd y}{\dd x}=f(x)$的积分曲线族（图4.3）.
\begin{figure}[htp]
    \centering
\begin{tikzpicture}[>=latex, xscale=3, yscale=.8]
\draw[->](-1,0)--(1.25,0)node[right]{$x$};
\draw[->](0,-4)--(0,4)node[right]{$y$};
\node at (0,0)[below left]{$O$};
\foreach \y in {1,2,3}
{
    \draw[domain=-.8:.8, samples=100, very thick]plot(\x, {\x*\x*\x+\y})node[right]{$y=x^3+$\y};
    \draw[domain=-.8:.8, samples=100, very thick]plot(\x, {\x*\x*\x-\y})node[right]{$y=x^3-$\y};
}
\draw[domain=-.8:.8, samples=100, very thick]plot(\x, {\x*\x*\x})node[right]{$y=x^3$};
\end{tikzpicture}
    \caption{}
\end{figure}

由微积分基本定理2知道，如果$y=G(x)$满足下面两个条件：
\[G'(x)=f(x)\quad \text{且}\quad G(a)=0\]
那么，$y=G(x)$就是$f(x)$的求和函数$G(x)=\Int^x_a f(x)\dd x$.

\begin{example}
求$\Int \sin (kx)\dd x,\quad \Int \cos(kx)\dd x,\quad \Int e^{kx}\dd x,\quad \Int \frac{\dd x}{x^2+a^2}$
\end{example}

\begin{solution}
    因为
\[\begin{split}
    (\cos kx)'&=-k\sin kx\\
    (\sin kx)'&=k\cos kx\\
    (e^{kx})'&=ke^{kx}\\
    \left(\arctan\frac{x}{a}\right)'&=\frac{1}{\left(\frac{x}{a}\right)^2+1}\cdot\left(\frac{x}{a}\right)'=\frac{a}{x^2+a^2}
\end{split}\]
所以
\[\begin{split}
    \int \sin (kx)\dd x&=-\frac{1}{k}\cos kx+C\\
    \Int \cos(kx)\dd x&=\frac{1}{k}\sin kx+C\\
    \Int e^{kx}\dd x&= \frac{1}{k}e^{kx}+C\\
    \Int \frac{\dd x}{x^2+a^2}&=\frac{1}{a}\arctan\frac{x}{a}+C \\
\end{split}  \]
\end{solution}

\begin{example}
    证明：若$\Int f(t)\dd t=F(t)+C$，则
\[\int f(kx)\dd x=\frac{1}{k}F(kx)+C\]
\end{example}

\begin{proof}
只须证明$f(kx)$的一个原函数是$\frac{1}{k}F(kx)$即可．
我们求$\frac{1}{k}F(kx)$对$x$的导数，于是
\[\frac{\dd}{\dd x}\left[\frac{1}{k}F(kx)\right]=\frac{1}{k}F'(kx)\cdot (kx)'=F'(kx)\]
由题设知$F' (t) =f (t)$, 所以
\[\frac{\dd}{\dd x}\left[\frac{1}{k}F(kx)\right]=F'(kx)=f(kx)\]
因此
\[\int f (kx) \dd x=\frac{1}{k}F (kx) +C\]
\end{proof}

\begin{example}
求$\Int \sin kx\sin\ell x\dd x$
\end{example}

\begin{solution}
\[\begin{split}
    \Int \sin kx\sin\ell x\dd x&=\int \frac{1}{2}[\cos(k-\ell)x-\cos(k+\ell)x]\dd x\\
    &=\frac{1}{2}\int \cos (k-\ell)x\dd x-\frac{1}{2}\int \cos (k+\ell)x\dd x\\
&=\frac{\sin(k-\ell)x}{2(k-\ell)}-\frac{\sin(k+\ell)x}{2(k+\ell)}+C
\end{split}\]
    这里$k\ne \pm \ell$．（如果$k=\pm\ell$，则另当考虑）
\end{solution}


\begin{example}
求$\Int \left(\frac{2a}{\sqrt{x}}-\frac{b}{x^2}+2xc^{\tfrac{2}{3}}\right)\dd x$
\end{example}


\begin{solution}
    \[\begin{split}
        \Int \left(\frac{2a}{\sqrt{x}}-\frac{b}{x^2}+2xc^{\tfrac{2}{3}}\right)\dd x&=2a\int x^{-\tfrac{1}{2}}\dd x-b\int x^{-2}\dd x+3c\int x^{\tfrac{2}{3}}\dd x\\
        &=4a\sqrt{x}+\frac{b}{x}+\frac{9}{5}cx^{\tfrac{5}{3}}+C
    \end{split}\]
\end{solution}

\begin{example}
    求$\Int \frac{3x^2}{1+x^2}\dd x$
\end{example}

\begin{solution}
\[\begin{split}
    \Int \frac{3x^2}{1+x^2}\dd x&=3\int \left(1-\frac{1}{x^2+1}\right)\dd x\\
    &=3\int \dd x-3\int \frac{\dd x}{x^2+1}\\
    &=3x-3\arctan x+C
\end{split}\]
\end{solution}

以上各例都是利用不定积分的性质2、3, 同时配合使用三角恒等式或设法将被积函数拆成几项，从而把一个比较复杂的积分化成若干个可以查基本积分表的积分．

\begin{ex}
\begin{enumerate}
    \item 计算下列不定积分:
\begin{multicols}{2}
\begin{enumerate}
\item  $\Int\left(x^{3}+3 x+5\right) \dd x$
\item  $\Int(x+2)^{4} \dd x$
\item  $\Int\left(x^{-2}+x^{-1}+x^{\tfrac{1}{3}}+x^{\tfrac{1}{2}}\right) \dd x$
\item  $\Int\left(a^{\tfrac{2}{3}}-x^{\tfrac{2}{3}}\right)^{3} \dd x$
\item  $\Int\left(\frac{1}{x}+e^{x}\right) \dd x$
\item  $\Int \frac{2 \cdot 3^{x}-5 \cdot 2^{x}}{3^{x}} \dd x$
\item  $\Int \left(\frac{3}{1+x^{2}}-\frac{2}{\sqrt{1-x^{2}}}\right) \dd x$
\item  $\Int \frac{2+x^{2}}{1+x^{2}} \dd x$
\item  $\Int \frac{\dd x}{x^{2}\left(1+x^{2}\right)}$
\item  $\Int \frac{\dd x}{x^{2}\left(a^{2}+x^{2}\right)}$
\item  $\Int 3^{x} e^{x} \dd x$
\end{enumerate}
\end{multicols}

\item  利用三角恒等式, 计算下列不定积分:
    \begin{multicols}{2}
        \begin{enumerate}
    \item  $\Int \tan^{2} x \dd x$
    \item  $\Int \frac{\cos 2 x}{\cos x-\sin x} \dd x$
    \item  $\Int \frac{2+\sin ^{2} x}{\cos ^{2} x} \dd x$
    \item  $\Int \sin ^{2} x \dd x$
    \item  $\Int \cos ^{2} x \dd x$
    \item  $\Int \cos k x \cos \ell x \dd x$
    \item  $\Int \sin 2 x \cos 5 x \dd x$
    \item  $\Int \frac{\dd t}{\sin ^{2} t \cos ^{2} t}$
    \item  $\Int \sin ^{3} x \dd x$
    \item  $\Int \cos ^{3} x \dd x$
\end{enumerate}
\end{multicols}    

\item 若$\Int f(t)\dd t=F(t)+C$, 求证
\[\int f (ax+b) \dd x=\frac{1}{a}F (ax+b) +C\]
\item  若$\frac{\dd y}{\dd x}=4x-6$, 且当$x=0$时，$y=9$, 求$y$的极小值．
\item  若$\frac{\dd y}{\dd x}=ax+b$, 求原函数$y=f(x)$使其满足下列条件：
\begin{enumerate}
    \item 当$x=0$时，$y=5$;
    \item 积分曲线向下凸且曲线$y=f(x)$在$x$轴上截得的弦长等于4;
    \item 曲线$y=f(x)$上纵坐标的最小值等于$-4$.
\end{enumerate}
\item 试求一个次数最低的多项式，使之当$x=1$时，有最大值6，而当$x=3$时，有最小值2.
\item  求下列各式的极限：
\begin{enumerate}
    \item $\Lim_{n\to\infty}\frac{\sqrt[n]{e}+\sqrt[n]{e^2}+\cdots+\sqrt[n]{e^n}}{n}$
    \item $\Lim_{n\to\infty}\frac{\sqrt[n]{e}+\sqrt[n]{e^4}+\cdots+\sqrt[n]{e^{2n}}}{n}$
\end{enumerate}
\end{enumerate}

\end{ex}


\section{不定积分的计算方法}
从上一节看到，虽然利用积分运算法则及基本积分表可以求出不少函数的原函数，但是实际遇到的积分仅凭这一些
方法还不能完全解决，本节将介绍几种典型的方法，利用这些方法，我们就可以计算更多的不定积分，此外，我们还制作了较为详细的积分表，可供使用时查阅．

\subsection{第一换元法}

有一些不定积分，将积分变量进行某种变映后就能由基本积分公式求出，例如求$\Int e^{2x}\dd x$, 在基本积分公式中只有$\Int e^x\dd x$, 没有直接求出它的公式，我们看到所求积分的被积函数$e^{2x}$是$y=e^u$和$u=2x$的复合函数，如果在积分表达式$e^{2x}\dd x$中，把被积函数看成中间变量$u$的函数，即设$u=2x$则$e^{2x}=e^u$, 从而$\dd u=2\dd x$, 由此可见只要再给微分$\dd x$凑上一个常数因子2, 就得到中间变量$u$的微分$\dd u$, 于是
\[\int e^{2x}\dd x=\int e^{2x}\cdot \frac{1}{2}\dd (2x) =\frac{1}{2}\int e^u\dd u=\frac{1}{2}e^u+C\]
然后再代回原来的积分变量，得
\[\int e^{2x}\dd x=\frac{1}{2}e^{2x}+C\]

\begin{example}
    求$\Int \frac{\dd x}{\sqrt{a^2-x^2}}\quad (a>0)$
\end{example}

\begin{solution}
    因为被积函数$\frac{1}{\sqrt{a^2-x^2}}=\frac{1}{a\sqrt{1-\left(\frac{x}{a}\right)^2}}$
是$y=\frac{1}{a\sqrt{1-u^2}}$和$u=\frac{x}{a}$的复合函数，且$\dd u=\frac{1}{a}\dd x$，于是
\[\begin{split}
    \int \frac{\dd x}{\sqrt{a^2-x^2}}=\int \frac{\dd x}{a\sqrt{1-\left(\frac{x}{a}\right)^2}}&=\int \frac{\frac{1}{a}\dd x}{\sqrt{1-\left(\frac{x}{a}\right)^2}}\\
    \text{令}u=\frac{x}{a}\to &= \int \frac{\dd u}{\sqrt{1-u^2}}\\
    \text{查表}\to &=\arcsin u+C\\
    u=\frac{x}{a} \to &=\arcsin \frac{x}{a}+C
\end{split}\]
\end{solution}

\begin{example}
    求$\Int \tan x\dd x$
\end{example}


\begin{solution}
\[\begin{split}
    \Int \tan x\dd x=\Int  \frac{\sin x}{\cos x}\dd x&=-\int \frac{\dd \cos x}{\cos x}\\
    u=\cos x\to &=-\int \frac{\dd u}{u}\\
    &=-\ln|u|+C=-\ln |\cos x|+C=\ln |\sec x|+C
\end{split}\]
\end{solution}

从上述两例看到，在求不定积分时，首先要与已知的基本积分公式相对比，并利用简单的变量代换，把要求的积分凑成公式所具有的形式，求出以后，再把原来的变量代回．以上变量代换中关键的一步是把原来的被积函数$f(x)$的某部分$\varphi(x)$换成中间变量$u$, 那么原来的被积表达式$f(x)\dd x$就拆成了$g[\varphi(x)]=g(u)$和$\varphi'(x)\dd x=\dd u$这两部分的乘积，即
\[f (x) \dd x=g [\varphi (x) ] \varphi' (x) \dd x=g (u) \dd u\]
于是：若$\Int g(x)\dd x=F(x)+C$，那么
\[\int g[\varphi(x)]\varphi'(x)\dd x=\int g(u)\dd u=F(u)+C=F[\varphi(x)]+C\]
上述换元法叫做第一换元法．

\begin{example}
求$\Int \sin^{3}x\dd x$
\end{example}

\begin{solution}
设$u=\cos x$，$\dd u=-\sin x\dd x$，于是
\[\begin{split}
    \Int \sin^{3}x\dd x&=\int (1-\cos^2 x)\cdot \sin x\dd x\\
  u=\cos x\to  &=-\int (1-u^2)\dd u\\
  &=-u+\frac{u^3}{3}+C=\frac{\cos^3 x}{3}-\cos x+C
\end{split}\]
\end{solution}

对换元积分法较熟练以后，所设换元变量$u$可以不写出．

\begin{example}
求$\Int \frac{\dd x}{\sqrt{ax+b}}$
\end{example}

\begin{solution}
    \[\Int \frac{\dd x}{\sqrt{ax+b}}=\frac{1}{a}\int \frac{\dd(ax+b)}{\sqrt{ax+b}}=\frac{2}{a}\sqrt{ax+b}+C\]
\end{solution}

\begin{example}
    求$\Int \frac{\dd x}{a^2-x^2}$
\end{example}

\begin{solution}
\[\begin{split}
    \Int \frac{\dd x}{a^2-x^2}&=\frac{1}{2a}\int\left(\frac{1}{a-x}+\frac{1}{a+x}\right)\dd x\\
    &=\frac{1}{2a}\left[-\int\frac{\dd(a-x)}{a-x}+\int\frac{\dd(a+x)}{a+x}\right]\\
    &=\frac{1}{2a}\left[-\ln|a-x|+\ln|a+x|\right]=\frac{1}{2a}\ln\left|\frac{a+x}{a-x}\right|
\end{split}\]
\end{solution}

\begin{example}
    求$\Int \sec x\dd x$
\end{example}

\begin{solution}
    \[\Int \sec x\dd x=\int \frac{\cos x}{\cos^2 x}\dd x=\int \frac{\dd(\sin x)}{1-\sin^2 x}\]
利用上题结果，得：
\[\begin{split}
    \int \sec x\dd x&=\frac{1}{2}\ln\left|\frac{1+\sin x}{1-\sin x}\right|+C\\
    &=\frac{1}{2}\ln\frac{1+\sin x}{1-\sin x}+C\\
    &=\frac{1}{2}\ln\frac{(1+\sin x)^2}{\cos^2 x}+C\\
    &=\ln|\sec x+\tan x|+C
\end{split}\]

\end{solution}

\begin{example}
    求$\Int \frac{1}{\sin^2x +3\cos^2x}\dd x$
\end{example}

\begin{solution}
    \[\begin{split}
        \Int \frac{1}{\sin^2x +3\cos^2x}\dd x&=\Int \frac{1}{(\tan^2x +3)\cos^2x}\dd x\\
        &=\int \frac{\dd (\tan x)}{\tan^2 x+\left(\sqrt{3}\right)^2}=\frac{1}{\sqrt{3}}\int \frac{\dd\left(\frac{\tan x}{\sqrt{3}}\right)}{\left(\frac{\tan x}{\sqrt{3}}\right)^2+1}\\
        &=\frac{1}{\sqrt{3}}\arctan\left(\frac{\tan x}{\sqrt{3}}\right)+C      
    \end{split}\]
\end{solution}

\begin{ex}
求下列不定积分：
\begin{enumerate}
    \begin{multicols}{2}
\item  $\Int(3 x+5)^{100} \dd x$
\item  $\Int\pi e^{-\tfrac{\pi}{4} x} \dd x$
\item  $\Int \frac{\ln x}{x} \dd x$
\item  $\Int\frac{1}{x^{2}} \cdot e^{\tfrac{1}{x}} \dd x$
\item  $\Int\frac{\cos \sqrt{x}}{\sqrt{x}} \dd x$
\item  $\Int \frac{1}{(\arcsin x)^{2} \sqrt{1-x^{2}}} \dd x$
\item  $\Int\frac{\dd x}{x^{2}+a^{2}}$
\item  $\Int \frac{\dd x}{x^{2}-a^{2}}$
\item  $\Int \frac{1}{\sqrt{x}} e^{\sqrt{x}} \dd x$
\item  $\Int \frac{2 a x+b}{2 \sqrt{a x^{2}+b x+c}} \dd x$
\item  $\Int\frac{1}{1+e^{x}} \dd x$        
    \end{multicols}
    \begin{multicols}{2}
\item  $\Int \frac{e^{x}}{e^{2 x}+1} \dd x$    
\item $\Int \frac{\cos x}{e^{\sin x}} \dd x$
\item $\Int \frac{\dd x}{\cos ^{2} x \sqrt{\tan x}}$
\item $\Int  \frac{\cos a x}{b-\sin a x} \dd x\quad (a \neq 0) $
\item $\Int x \sqrt{a x^{2}+c}\; \dd x$
\item $\Int \frac{\sin ^{2} x \cos x \dd x}{1+\sin x}$
\item $\Int \frac{\sqrt{\arctan x}}{1+x^{2}} \dd x$
\item $\Int \frac{1}{1+\sin x} \dd x$        
    \end{multicols}
    \begin{multicols}{2}
\item $\Int \frac{\sqrt{1+\cos x}}{\sin x} \dd x$
\item $\Int \frac{\dd x}{x^{2}+x-6}$
\item $\Int \frac{x \dd x}{3 x^{2}+x+2}$
\item $\Int  \frac{\dd x}{3 x^{2}+x+2}$
\item $\Int  \cot x \dd x$
\item $\Int  \frac{\dd x}{\sin x}$
\item $\Int \frac{\arcsin x}{\sqrt{1-x^{2}}} \dd x$        
    \end{multicols}

\end{enumerate}
\end{ex}

\subsection{第二换元积分法}

有些积分不能很容易地凑出微分，而是一开始就要作代换，把要求的积分化简，然后再求积分．

\begin{example}
    求$\Int \frac{1}{1+\sqrt{x}}\dd x$.
\end{example}

\begin{solution}
    这个积分的表达式不容易拆成中间变量$u=\varphi(x)$
的函数$g(\varphi(x))=g(u)$和中间变量的微分$\varphi'(x)\dd x=\dd u$的乘积，因此前面的第一换元法对此题用起来不便．这时我们设法作一个代换，以便把被积函数中的根号去掉，化成新变量的有理式．

令$x=u^2$, 即作代换$\sqrt{x}=u$, 于是
\[\begin{split}
    \frac{1}{1+\sqrt{x}}&=\frac{1}{1+u}\\
    \dd x&=2u\dd u
\end{split}\]
从而
\[\begin{split}
    \Int \frac{1}{1+\sqrt{x}}\dd x&=\int \frac{1}{1+u}2u\dd u\\
    &=2\int \frac{(1+u)-1}{1+u}\dd u\\
    &=2\left[\int \dd u-\int \frac{1}{1+u}\dd u\right]\\
    &=2\left[u-\ln |1+u|\right]+C\\
    &=2\left[\sqrt{x}-\ln\left(1+\sqrt{x}\right)\right]+C
\end{split}\]
\end{solution}

这里由于把积分变量$x$换成了$u$的函数$x=u^2$, 因此在最后结果中，用变量$u$表出的函数，还要还原成$x$的函数，这就要求从所作的代换$x=u^2$中解出反函数$u=\sqrt{x}$来．

我们把第二换元积分法用定理形式叙述如下：

\begin{blk}{定理}
设$x=\varphi(u)$是可微函数，且$\varphi'(u)\ne 0$，若
\begin{equation}
    \int f[\varphi(u)]\varphi'(u)\dd u=F(u)+C
\end{equation}
那么
\begin{equation}
    \int f(x)\dd x=F\left[\varphi^{-1}(x)\right]+C
\end{equation}
\end{blk}

\begin{proof}
只须证明$\frac{\dd}{\dd x}F[\varphi^{-1}(x)]=f(x)$即可．

因为$\varphi'(u)\ne 0$，故$x=\varphi(u)$的反函数$u=\varphi^{-1}(x)$存在，于是
\[\begin{split}
    \frac{\dd}{\dd x}F[\varphi^{-1}(x)]&=\frac{\dd}{\dd u}F[\varphi^{-1}(x)]\cdot \frac{\dd }{\dd x}\varphi^{-1}(x)\\
    &=\frac{\dd}{\dd u}F(u)\cdot \frac{1}{\frac{\dd}{\dd u}\varphi(u)}\\
    &=F'(u)\cdot \frac{1}{\varphi'(u)}\\
    &=\left\{f[\varphi(u)]\varphi'(u)\right\}\cdot \frac{1}{\varphi'(u)}\\
    &=f[\varphi(u)]\\
    x=\varphi(u)\to &=f(x)
\end{split}\]
\end{proof}


\begin{example}
    求$\Int\frac{\dd x}{\sqrt{x}+\sqrt[3]{x}}$.
\end{example}

\begin{solution}
    \[\Int\frac{\dd x}{\sqrt{x}+\sqrt[3]{x}}=\int \frac{\dd x}{\left(\sqrt[6]{x}\right)^3+\left(\sqrt[6]{x}\right)^2}\]
设$x=u^6$，$u>0$，则可化去被积函数中的根号，于是
\[\frac{1}{\sqrt{x}+\sqrt[3]{x}}=\frac{1}{u^3+u^2}\quad \Rightarrow\quad \dd x=6u^5\dd u\]
从而
\[\begin{split}
    \Int\frac{\dd x}{\sqrt{x}+\sqrt[3]{x}}&=\int \frac{6u^5}{u^3+u^2}\dd u=\int\frac{6u^3}{u+1}\dd u\\
    &=6\int \left(u^2-u+1-\frac{1}{u+1}\right)\dd u\\
    &=6\left\{\frac{u^3}{3}-\frac{u^2}{2}+u-\ln|u+1|\right\}+C\\
    &=2\sqrt{x}-3\sqrt[3]{x}+6\sqrt[6]{x}-\ln\left(\sqrt[6]{x}+1\right)+C
\end{split}
    \]
\end{solution}

\begin{example}
    求$\Int \frac{x+1}{\sqrt[3]{3x+1}}\dd x$.
\end{example}

\begin{solution}
    设$3x+1=t^3$，即$x=\frac{1}{2}(t^3-1)$，于是：$\dd x=t^2\dd t$
\[\begin{split}
    \Int \frac{x+1}{\sqrt[3]{3x+1}}\dd x&=\int \frac{\frac{1}{3}(t^3-1)+1}{t}t^2\dd t\\
&=\frac{1}{3}\int (t^4+2t)\dd t\\
&=\frac{1}{3}\left(\frac{1}{5}t^5+t^2\right)+C=\frac{1}{15}(3x+1)^{\tfrac{5}{3}}+\frac{1}{3}(3x+1)^{\tfrac{2}{3}}+C\\
&=\frac{1}{5}(x+2)(3x+1)^{\tfrac{2}{3}}+C
\end{split}\]
\end{solution}

上面所介绍的两种变量代换法的基本思想是一致的．第一种是把被积函数中的一个小部分看作一个变量，即“化繁为简”；第二种则把积分变量代以一个新变量的函数，表面上看来是“化简为繁”，但实际上克服了求积分的困难．应当注意，在使用第二种变量代换法时，要保证变换是可微的和一对一的，否则会出现错误.

下面介绍当被积函数含有二次根式$\sqrt{a^2-x^2}$, $\sqrt{a^2+x^2}$, $\sqrt{x^2-a^2}$时，怎样作代换．

\begin{example}
    求$\Int \sqrt{a^2-x^2}\dd x\quad (a>0)$.
\end{example}

\begin{solution}
因为\[\sqrt{a^2-x^2}=\sqrt{a^2\left(1-\frac{x^2}{a^2}\right)}=a\sqrt{1-\left(\frac{x}{a}\right)^2}\]
所以可作正弦代换
\[\frac{x}{a}=\sin t\quad \left(-\frac{\pi}{2}\le t\le \frac{\pi}{2}\right)\]
$\sin t$在$\left[-\frac{\pi}{2},\frac{\pi}{2}\right]$上连续且递增，有反函数$t=\arcsin\frac{x}{a},\; -1\le x\le 1$. 于是
\[\begin{split}
    \dd x&=\dd (a\sin t)=a\cos t\dd t\\
    \sqrt{a^2-x^2}&=a\sqrt{1-\left(\frac{x}{a}\right)^2}=a\sqrt{1-\sin^2 t}=a\cos t
\end{split}\]
从而：
\[\begin{split}
    \Int \sqrt{a^2-x^2}\dd x&=\int a\cos t\cdot a\cos t\dd t=a^2\int \cos^2 t\dd t\\ 
&=a^2\int \frac{1+\cos 2t}{2}\dd t=\frac{a^2}{2}\left[\int \dd t+\frac{1}{2}\int \cos 2t\dd (2t)\right]\\
&=\frac{a^2}{2}\left(t+\frac{1}{2}\sin 2t\right)+C=\frac{a^2}{2}\left(t+\sin t\cos t\right)+C\\
&=\frac{a^2}{2}\left[\arcsin\frac{x}{a}+\frac{x}{a}\sqrt{1-\left(\frac{x}{a}\right)^2}+C\right]\\
&=\frac{a^2}{2}\arcsin\frac{x}{a}+\frac{x}{2}\sqrt{a^2-x^2}+C
\end{split}\]
\end{solution}

\begin{example}
    求$\Int \frac{\dd x}{\sqrt{a^2+x^2}}\quad (a>0)$.
\end{example}

\begin{solution}
因为$\sqrt{a^2+x^2}=a\sqrt{1+\left(\frac{x}{a}\right)^2}$，设$\frac{x}{a}=\tan t\; \left(-\frac{\pi}{2}<t<\frac{\pi}{2}\right)$，则存在反函数
\[t=\arctan\frac{x}{a}\quad \Rightarrow\quad \dd x=a\sec^2 t\dd t\]
又$\sqrt{a^2+x^2}=a\sqrt{1+\tan^2 t}=a\sec t$
，从而：
\[\begin{split}
    \Int \frac{\dd x}{\sqrt{a^2+x^2}}&=\int\frac{a\sec^2 t\dd t}{a\sec t}=\int\sec t\dd t\\
    &=\ln|\sec t+\tan t|+C_1\\
    &=\ln\left|\frac{1}{a}\sqrt{a^2+x^2}+\frac{x}{a}\right|+C_1\\
    &=\ln\left|x+\sqrt{a^2+x^2}\right|-\ln a+C_1\\
    &=\ln\left|x+\sqrt{a^2+x^2}\right|+C\qquad (C=C_1-\ln a)
\end{split}\]
\end{solution}

\begin{example}
    求$\Int \frac{\dd x}{\sqrt{x^2-a^2}}\quad (a>0)$
\end{example}

\begin{solution}

    \begin{center}
\begin{tikzpicture}[scale=.7]
\draw[very thick](0,0)--node[below]{$a$}(3,0)--node[right]{$\sqrt{x^2-a^2}$}(3,4)--node[left]{$x$}(0,0);
\draw(.75,0) arc (0:53.13:.75)node[right]{$t$};
\end{tikzpicture}
    \end{center}

    作正割代换．设$x=a\sec t\; \left(0<t<\frac{\pi}{2}\right)$，此函数在$t=\frac{\pi}{2}$点间断，在区间$\left[0,\frac{\pi}{2}\right)$上由$a$递增到$+\infty$，而在区间$\left(\frac{\pi}{2},\pi\right]$上由$-\infty$递增到$-a$．因此，$x=a\sec t$在$\left(0,\frac{\pi}{2}\right]$上的反函数是 
\[t=\arcsec\frac{x}{a}\quad \left(1<\frac{x}{a}<\infty,\; \text{即}x>a\right)\]
于是，当$x>a$时，有
\[\begin{split}
    \sqrt{x^2-a^2}&=\sqrt{a^2\sec^2 t-a^2}=\sqrt{a^2(\sec^2 t-1)}\\
    &=a\sqrt{\tan^2 t}=a|\tan t|\\
    &=a\tan t \quad \left(0\le t<\frac{\pi}{2}\right)
\end{split}\]
\[\dd x=\dd(a\sec t)=a\frac{\sin t}{\cos^2 t}\dd t\]
从而：
\[\begin{split}
    \Int \frac{1}{\sqrt{x^2-a^2}}\dd x&=\int\frac{1}{a\tan t}\cdot \frac{a\sin t}{\cos^2 t}\dd t\\
    &=\int \frac{1}{\cos t}\dd t=\ln|\sec t+\tan t|+C_1\\
    &=\ln\left|\frac{x}{a}+\frac{\sqrt{x^2-a^2}}{a}\right|+C_1\\
    &=\ln\left|x+\sqrt{x^2-a^2}\right|+C_1-\ln a=\ln\left|x+\sqrt{x^2-a^2}\right|+C
\end{split}\]
其中$C=C_1-\ln a$.

又函数$x=a\sec t$在$\left(\frac{\pi}{2},\pi\right]$上的反函数是
\[t=\arcsec\frac{x}{a}\quad \left(-\infty<\frac{x}{a}<-1,\; \text{即}x<-a\right)\]
于是，当$x<-a$时，有
\[\sqrt{x^2-a^2}=a\sqrt{\tan^2 t}=a|\tan t|=-a\tan t\quad \left(\frac{\pi}{2}<t\le \pi\right)\]
\[\dd x=\dd(a\sec t)=a\frac{\sin t}{\cos^2 t}\cdot \dd t\]
\[\begin{split}
    \Int \frac{1}{\sqrt{x^2-a^2}}\dd x&=\int-\frac{1}{a\tan t}\cdot \frac{a\sin t}{\cos^2 t}\dd t\\
    &=-\int \frac{1}{\cos t}\dd t=-\ln|\sec t+\tan t|+C_2\\
    &=-\ln\left|\frac{x}{a}-\frac{\sqrt{x^2-a^2}}{a}\right|+C_2\\
    &=-\ln\left|x-\sqrt{x^2-a^2}\right|+C_2+\ln a\\
    &=\ln\left|\frac{1}{x-\sqrt{x^2-a^2}}\right|+C'\qquad (C'=C+\ln a)\\
    &=\ln\left|\frac{x+\sqrt{x^2-a^2}}{a^2}\right|+C'=\ln\left|x+\sqrt{x^2-a^2}\right|+C'-\ln a^2\\
    &=\ln\left|x+\sqrt{x^2-a^2}\right|+C''\qquad (C''=C'-\ln a^2)
\end{split}\]
因此，不论$x>a$或$x<-a$，都有
\[\int\frac{\dd x}{\sqrt{x^2-a^2}}=\ln\left|x+\sqrt{x^2-a^2}\right|+C\]
\end{solution}

到现在为止，我们又得到四个重要的不定积分公式：

\begin{blk}{}
\begin{align}
\Int \frac{\dd x}{a^2+x^2}&=\frac{1}{a}\arctan\frac{x}{a}+C\\
\Int\frac{\dd x}{a^2-x^2}&=\frac{1}{2a}\ln\left|\frac{a+x}{a-x}\right|+C\\
 \Int\frac{\dd x}{\sqrt{a^2-x^2}}&=\arcsin\frac{x}{a}+C\\
\Int \frac{\dd x}{\sqrt{x^2\pm a^2}}&=\ln\left|x+\sqrt{x^2\pm a^2}\right|+C
\end{align}
\end{blk}

利用这些公式可以进一步求出许多不定积分．

\begin{blk}
 {类型A}
\[\int \frac{\dd x}{ax^2+bx+c}\]
\end{blk}

如果被积函数的分母可以分解因式，则把被积函数分解成部分分式，就很容易分项求出不定积分．

如分母不能分解因式，那么将它配方得
\[ax^2+bx+c=a\left(x+\frac{b}{2a}\right)^2+\frac{4ac-b^2}{4a}\]
于是
\begin{itemize}
    \item 当$4ac>b^2$时，上面不定积分归结到(4.10);
    \item 当$4ac<b^2$时，上面不定积分归结到(4.11).
\end{itemize}

\begin{example}
    求$\Int \frac{\dd x}{3x^2+x-1}$.
\end{example}

\begin{solution}
\[\begin{split}
    \Int \frac{\dd x}{3x^2+x-1}&=\frac{1}{3}\int \frac{\dd x}{\left(x+\frac{1}{6}\right)^2-\frac{13}{36}}\\
&=\frac{1}{3}\frac{1}{2\sqrt{\frac{13}{36}}}\ln\left|\frac{x+\frac{1}{6}-\frac{\sqrt{13}}{6}}{x+\frac{1}{6}+\frac{\sqrt{13}}{6}}\right|+C\\
&=\frac{1}{\sqrt{13}}\ln\left|\frac{6x+1-\sqrt{13}}{6x+1+\sqrt{13}}\right|+C
\end{split}\]    
\end{solution}

\begin{blk}
    {类型B}
   \[\int \frac{px+q}{ax^2+bx+c}\dd x\]
   \end{blk}

注意到$\frac{\dd}{\dd x}(ax^2+bx+c)=2ax+b$，被积函数的分子可以写成
\[px+q=\frac{p}{2a}(2ax+b)+q-\frac{pb}{2a}\]
于是
\[\int\frac{px+q}{ax^2+bx+c}\dd x=\frac{p}{2a}\int\frac{2ax+b}{ax^2+bx+c}\dd x+\left(q-\frac{bp}{2a}\right)\int\frac{\dd x}{ax^2+bx+c}\]
等式右端第一项是$\frac{p}{2a}\ln(ax^2+bx+c)$，第二项可归结到类型A.

\begin{example}
    求$\Int \frac{5x-1}{3x^2+x+2}\dd x$.
\end{example}

\begin{solution}
\[\begin{split}
    \Int \frac{5x-1}{3x^2+x+2}\dd x&=\frac{5}{6}\int \frac{6x+1}{3x^2+x+2}\dd x-\frac{11}{6}\int\frac{\dd x}{3x^2+x+2}\\
&=\frac{5}{6}\ln(3x^2+x+2)-\frac{11}{6}\cdot \frac{1}{3}\cdot \int \frac{\dd x}{\left(x+\frac{1}{6}\right)^2+\frac{23}{36}}\\
&=\frac{5}{6}\ln(3x^2+x+2)-\frac{11}{6}\cdot\frac{1}{3}\cdot\frac{6}{\sqrt{23}}\arctan\frac{6x+1}{\sqrt{23}}+C\\
&=\frac{5}{6}\ln(3x^2+x+2)-\frac{11}{3\sqrt{23}}\arctan\frac{6x+1}{\sqrt{23}}+C
\end{split}\]    
\end{solution}

\begin{blk}{类型C}
\[\int \frac{\dd x}{\sqrt{ax^2+bx+c}}\]
\end{blk}

\begin{itemize}
    \item 若$a>0$, 则被积函数的根底式经配方后，所求积分可归结到积分公式(4.13);
    \item  若$a<0$, 则此积分可归结到积分公式(4.12).
\end{itemize}

\begin{example}
    求$\Int \frac{\dd x}{\sqrt{2+x-3x^2}}$
\end{example}

\begin{solution}
    \[\begin{split}
        \Int \frac{\dd x}{\sqrt{2+x-3x^2}}&=\frac{1}{\sqrt{3}}\int \frac{\dd x}{\sqrt{\frac{2}{3}+\frac{1}{3}x-x^2}}\\
        &=\frac{1}{\sqrt{3}}\int \frac{\dd x}{\sqrt{\frac{25}{36}-\left(x-\frac{1}{6}\right)^2}}\\
        &=\frac{1}{\sqrt{3}}\arcsin\frac{x-\frac{1}{6}}{\frac{5}{6}}+C\\
        &=\frac{1}{\sqrt{3}}\arcsin\frac{6x-1}{5}+C       
    \end{split}\]
\end{solution}

\begin{blk}{类型D}
\[\int \frac{px+q}{\sqrt{ax^2+bx+c}}\dd x\]
\end{blk}

被积函数的分子可写成
\[px+q=\frac{p}{2a}(2ax+b)+q-\frac{pb}{2a}\]
所以
\[
    \int \frac{px+q}{\sqrt{ax^2+bx+c}}\dd x=\frac{p}{2a}\int \frac{2ax+b}{\sqrt{ax^2+bx+c}}\dd x+\left(q-\frac{pb}{2a}\right)\int\frac{\dd x}{\sqrt{ax^2+bx+c}}
\]
等式右端第一个积分等于被积函数的分母的2倍，即$2\sqrt{ax^2+bx+c}$, 第二个是类型C的积分．

\begin{example}
    求$\Int \frac{5x-1}{\sqrt{3x^2+x+2}}\dd x$.
\end{example}

\begin{solution}
    \[\begin{split}
        \Int \frac{5x-1}{\sqrt{3x^2+x+2}}\dd x&=\frac{5}{6}\int \frac{6x+1}{\sqrt{3x^2+x+2}}\dd x-\frac{11}{6}\int \frac{\dd x}{\sqrt{3x^2+x+2}}\\
&=\frac{5}{6}\cdot 2\sqrt{3x^2+x+2}\\
&\qquad -\frac{11}{6}\cdot\frac{1}{\sqrt{3}}\ln\left[x+\frac{1}{6}+\sqrt{\left(x+\frac{1}{6}\right)^2+\frac{23}{36}}\right]+C\\
&=\frac{5}{3}\sqrt{3x^2+x+2}-\frac{11}{6\sqrt{3}}\ln\left[x+\frac{1}{6}+\sqrt{\left(x+\frac{1}{6}\right)^2+\frac{23}{36}}\right]+C        
    \end{split}\]
\end{solution}

\begin{blk}{类型E}
    \[\int \frac{\dd x}{x\sqrt{ax^2+bx+c}}\]
\end{blk}

作变量代换$x=\frac{1}{y}$，所求积分可化为类型C．因为$\dd x=-\frac{1}{y^2}\dd y$，并且
\[\int \frac{\dd x}{x\sqrt{ax^2+bx+c}}=\int\frac{-\frac{1}{y^2}\dd y}{\frac{1}{y}\sqrt{\frac{a}{y^2}+\frac{b}{y}+c}}=-\int\frac{\dd y}{\sqrt{a+by+cy^2}}\]

\begin{example}
    求$\Int\frac{\dd x}{x\sqrt{x^2+x+1}}$
\end{example}

\begin{solution}
\[\begin{split}
    \Int\frac{\dd x}{x\sqrt{x^2+x+1}}&=-\int \frac{-\frac{1}{y^2}\dd y}{\frac{1}{y}\sqrt{\frac{1}{y^2}+\frac{1}{y}+1}}=-\int \frac{\dd y}{\sqrt{y^2+y+1}}\\
    &=-\int \frac{\dd y }{\sqrt{\left(y+\frac{1}{2}\right)^2+\frac{3}{4}}}  \\
    &=-\ln \left[y+\frac{1}{2}+\sqrt{\left(y+\frac{1}{2}\right)^2+\frac{3}{4}}\right]+C \\
    &=-\ln \left[y+\frac{1}{2}+\sqrt{y^2+y+1}\right]+C\\
    &=-\ln\left[\frac{1}{x}+\frac{1}{2}+\sqrt{\frac{1}{x^2}+\frac{1}{x}+1}\right]+C\\
    &=-\ln\frac{2+x+2\sqrt{x^2+x+1}}{2x}+C
\end{split}\]
\end{solution}
    
\begin{ex}
\begin{enumerate}
    \item 求下列不定积分：
\begin{multicols}{2}
\begin{enumerate}
    \item $\Int \frac{1}{1+\sqrt[3]{x}}\dd x$
    \item $\Int \frac{1}{(x+2) \sqrt{x+1}}\dd x$
    \item $\Int  \frac{x}{\sqrt[3]{1-x}}\dd x $
    \item $\Int \frac{\sqrt{x}}{\sqrt[3]{x^{2}}-\sqrt[4] {x}}\dd x$
    \item $\Int  \frac{\sqrt{x^{2}-a^{2}}}{x}\dd x $
    \item $\Int  \frac{\sqrt{x^{2}-a^{2}}}{x^{2}}\dd x$
    \item $\Int  \frac{\sqrt{a^{2}-x^{2}}}{x}\dd x$
    \item $\Int \frac{\sqrt{a^{2}-x^{2}}}{x^{2}}\dd x$
    \item $\Int \frac{x^{2}}{\sqrt{1-x^{2}}}\dd x$
    \item $\Int \frac{1}{x^{4} \sqrt{1+x^{2}}}\dd x$
\end{enumerate}
\end{multicols}

    \item 求下列不定积分：
\begin{multicols}{2}
\begin{enumerate}
\item $\Int \frac{1}{\sqrt{4 x^{2}+6 x}}\dd x $
\item $\Int \frac{x}{\sqrt{4 x^{2}+6 x}}\dd x$
\item $\Int  \frac{1}{\sqrt{1-x-4 x^{2}}}\dd x$
\item $\Int \frac{x+1}{\sqrt{1-x-4 x^{2}}}\dd x$
\item $\Int \frac{1}{x \sqrt{x^{2}+1}}\dd x$
\item $\Int \frac{1}{x \sqrt{1-x^{2}}}\dd x$
\item $\Int \frac{1}{x \sqrt{x^{2}-4 x}}\dd x$
\item $\Int \sqrt{\frac{x-1}{x+1}}\dd x$  
\end{enumerate}
\end{multicols}

\end{enumerate}    
\end{ex}

\subsection{有理函数的积分}
如果$P(x)$, $Q(x)$是两个多项式，则$\Int \frac{P(x)}{Q(x)}\dd x$可以化
归为下列类型的积分：
\begin{center}
\begin{tabular}{p{.3\textwidth}p{.3\textwidth}}
        $\Int \frac{A \dd x}{x-a}$ & $\Int \frac{A \dd x}{(x-a)^{n}}$ \\
        $\Int \frac{p x+q}{x^{2}+b x+c} \dd x$ & $\Int \frac{p x+q}{\left(x^{2}+b x+c\right)^{n}} \dd x$
\end{tabular}
\end{center}

对这些积分，我们是有办法的．下面我们通过举例来说明求这些积分的大概的方法．

如果有理分式的分子次数高于分母的次数时，则只要作一次除法就可以把它写成一个多项式与一个真分式的和，而前者是容易求积的．所以我们只须注意真分式的积分如何求．

\begin{example}
    求$\Int \frac{x^3+x+1}{x^2+1}\dd x$
\end{example}

\begin{solution}
    \[\begin{split}
        \Int \frac{x^3+x+1}{x^2+1}\dd x&=\int\left(x+\frac{1}{x^2+1}\right)\dd x\\
&=\int x\dd x+\int\frac{\dd x}{1+x^2}\\
&=\frac{x^2}{2}+\arctan x+C        
    \end{split}\]
\end{solution}

\begin{example}
    求$\Int \frac{\dd x}{x^4-1}$
\end{example}

\begin{solution}
$\because\quad x^4-1=(x-1)(x+1)(x^2+1)$

$\therefore\quad \frac{1}{x^4-1}=\frac{A}{x-1}+\frac{B}{x+1}+\frac{Cx+D}{x^2+1}$

把等式右端通分，再比较两端系数，于是等式
\[A (x^2+1) (x+1)+B(x^2+1)(x-1)+(Cx+D) (x^2-1)=1\]
的左边多项式的$x^3,x^2,x$的系数应为零，常数项应为1, 从而得到
\begin{numcases}{}
    A+B+C=0\\
    A-B+D=0\\
    A+B-C=0\\
    A-B-D=1
\end{numcases}
由$(4.14)+(4.16)$和$(4.15)+(4.17)$得到：
\[A+B=0,\qquad A-B=\frac{1}{2}\]
解得：
\[A=\frac{1}{4},\quad B=-\frac{1}{4},\quad C=0,\quad D=-\frac{1}{2}\]
即：
\[\frac{1}{x^4-1}=\frac{1}{4(x-1)}-\frac{1}{4(x+1)}-\frac{1}{2(x^2+1)}\]
从而
\[\begin{split}
    \int\frac{\dd x}{x^4-1}&=\frac{1}{4}\int \frac{\dd x}{x-1}-\frac{1}{4}\int\frac{\dd x}{x+1}-\frac{1}{2}\int\frac{\dd x}{1+x^2}\\
    &=\frac{1}{4}\ln\left|\frac{x-1}{x+1}\right|-\frac{1}{2}\arctan x+C
\end{split}\]
\end{solution}

\begin{example}
求$\Int \frac{3x^2+x-2}{(x-2)^2(1-2x)}\dd x$
\end{example}

\begin{solution}
设$\frac{3x^2+x-2}{(x-2)^2(1-2x)}=\frac{A}{1-2x}+\frac{B}{x-2}+\frac{C}{(x-2)^2}$，则
\[3x^2+x-2=A(x-2)^2+B(1-2x)(x-2)+C(1-2x)\]

令$1-2x=0$，则$A=-\frac{1}{3}$；令$x-2=0$，则$C=-4$. 比较等式两端$x^2$的系数，得
\[3=A-2B\quad\Rightarrow\quad B=-\frac{5}{3}\]
所以
\[\frac{3x^2+x-2}{(x-2)^2(1-2x)}=-\frac{1}{3(1-2x)}-\frac{5}{3(x-2)}-\frac{4}{(x-2)^2}\]
从而
\[\begin{split}
\int \frac{3x^2+x-2}{(x-2)^2(1-2x)}\dd x &=\frac{1}{6}\int \frac{\dd x}{x-\frac{1}{2}}-\frac{5}{3}\int \frac{\dd x}{x-2}-4\int \frac{\dd x}{(x-2)^2}\\
&=\frac{1}{6}\ln\left|x-\frac{1}{2}\right|-\frac{5}{3}\ln|x-2|+\frac{4}{x-2}+C
\end{split}\]
\end{solution}

三角函数有理式的积分
\[\int R (\sin x,\cos x) \dd x\]
恒可用“半角变换”法，即设$t=\tan\frac{x}{2}$或$x=2\arctan t$化为$t$的有理函数的积分．

事实上，由$t=\tan\frac{x}{2}$, 知
\[\begin{split}
\sin x&=\frac{2 \sin \frac{x}{2} \cos \frac{x}{2}}{\sin ^{2} \frac{x}{2}+\cos ^{2} \frac{x}{2}}=\frac{2 \tan \frac{x}{2}}{1+\tan^{2} \frac{x}{2}}=\frac{2 t}{1+t^{2}}\\
\cos x &=\frac{\cos ^{2}-\frac{x}{2}-\sin ^{2} \frac{x}{2}}{\sin ^{2} \frac{x}{2}+\cos ^{2} \frac{x}{2}}=\frac{1-\tan^{2} \frac{x}{2}}{1+\tan^{2} \frac{x}{2}} 
=\frac{1-t^{2}}{1+t^{2}}
\end{split}\]
又由 $x=2 \arctan t$, 知
\[\dd x=\frac{2 \dd t}{1+t^{2}}\]
于是
\[
\int R(\sin x,  \cos x) \dd x=\int R\left(\frac{2 t}{1+t^{2}}, \frac{1-t^{2}}{1+t^{2}}\right) \cdot \frac{2}{1+t^{2}} \dd x 
\]
上式右端是 $t$ 的有理函数.

\begin{example}
    求$\Int \frac{\dd\theta}{\sin\theta}$.
\end{example}

\begin{solution}
令$t=\tan\frac{\theta}{2}$，则：$\theta=2\arctan t,\quad \dd\theta=\frac{2\dd t}{1+t^2}$，从而：
\[\begin{split}
    \int\frac{\dd\theta}{\sin\theta}=\int\frac{1+t^2}{2t}\cdot \frac{2\dd t}{1+t^2}&=\int\frac{\dd t}{t}\\
    &=\ln|t|+C=\ln\left|\tan\frac{\theta}{2}\right|+C
\end{split}\]
又$\tan\frac{\theta}{2}=\frac{\sin\frac{\theta}{2}}{\cos\frac{\theta}{2}}=\frac{2\sin^2 \frac{\theta}{2}}{2\sin\frac{\theta}{2}\cos \frac{\theta}{2}}=\frac{1-\cos\theta}{\sin\theta}=\csc\theta-\cot\theta$，
所以
\[\int\frac{\dd\theta}{\sin\theta}=\ln\left|\tan\frac{\theta}{2}\right|+C=\ln|\csc\theta-\cot\theta|+C\]    
\end{solution}

\begin{example}
    求$\Int \frac{\dd\theta}{3+2\cos\theta}$
\end{example}

\begin{solution}
    令$t=\tan\frac{\theta}{2}$，则：$\dd\theta=\frac{2\dd t}{1+t^2}$，于是
\[\begin{split}
    \Int \frac{\dd\theta}{3+2\cos\theta}=\int\frac{\frac{2\dd t}{1+t^2}}{3+2\frac{1-t^2}{1+t^2}}&=\int \frac{2\dd t}{5+t^2}\\
    &=\frac{2}{\sqrt{5}} \arctan\frac{t}{\sqrt{5}}+C\\
    &=\frac{2}{\sqrt{5}}\arctan\left(\frac{1}{\sqrt{5}}\tan\frac{\theta}{2}\right)+C
\end{split}\]
\end{solution}

\begin{ex}
\begin{enumerate}
\item 求下列不定积分:
\begin{multicols}{2}
\begin{enumerate}
\item  $\Int \frac{1}{x^{2}-3 x+2}\dd x$
\item  $\Int  \frac{x}{x^{2}-3 x+2}\dd x$
\item  $\Int  \frac{x^{2}}{x^{2}-3 x+2}\dd x$
\item  $\Int \frac{x^{2}+1}{x+3}\dd x$
\item  $\Int \frac{1}{x^{2}-3 x+1}\dd x$
\item  $\Int \frac{x}{x^{2}-3 x+1}\dd x$
\item  $\Int \frac{x^{2}}{x^{2}-3 x+1}\dd x$
\end{enumerate}
\end{multicols}

\item 求下列不定积分:
\begin{multicols}{2}
    \begin{enumerate}
        \item  $\Int  \frac{\cos \theta \dd \theta}{1+\sin ^{2} \theta}$
        \item  $\Int \frac{\sin x+\cos x}{\cos ^{2} x}\dd x$
        \item  $\Int  \frac{\dd x}{4 \cos x+3 \sin x+5}$
        \item  $\Int \frac{\sin x}{\cos x+\sin x}\dd x$
    \end{enumerate}
\end{multicols}
\end{enumerate}
\end{ex}

\subsection{分部积分法}

分部积分法是从两个函数乘积的微分公式得到的求积分方法．

设$u=u(x)$, $v=v(x)$是两个可微函数，乘积$u\cdot v$的微分法则是$$\dd (u\cdot v) =u\dd v+v\dd u$$
移项，得$u\dd v=\dd (w\cdot v) -v\dd u$，
两边积分得
\begin{equation}
    \int u\dd v=u\cdot v-\int v\dd u
\end{equation}
公式(4.18)表示如果积分$\Int u\dd v=\Int u(x)v'(x)\dd x$不可能直接求出，就改为求积分$\Int v\dd u=\Int v(x)u'(x)\dd x$, 设$F(x)$是$v(x)\cdot u'(x)$的一个原函数，那么$u(x)\cdot v(x)-F(x)$是$u(x)v'(x)$的一个原函数，用公式(4.18)求不定积分的方法叫做\textbf{分部积分法}，公式(4.18)叫做\textbf{分部积分公式}．

\begin{example}
    求$\Int x\sin 3x\dd x$.
\end{example}

\begin{solution}
被积函数$x\sin3x$不能由前面的基本积分表求出它的原函数，我们把$x\sin3x$分成两部分，设$u(x)=x$, $v'(x)=\sin3x$, 于是
\[\int v' (x) \dd x= \int \sin3x\dd x=-\frac{\cos3x}{3}+C\]
即$v(x)=-\frac{\cos3x}{3}$是$v'(x)=\sin3x$的一个原函数，根据公
式(4.18)，有
\[\begin{split}
    \int x\sin 3x\dd x&=\int x\dd\left(-\frac{\cos 3x}{3}\right)\\
    &=x\left(-\frac{\cos 3x}{3}\right)-\int\left(-\frac{\cos3x}{3}\right)\dd x \\
&=-\frac{x\cos 3x}{3}+\frac{1}{3}\int \cos 3x\dd x\\
&=-\frac{x\cos 3x}{3}+\frac{\sin 3x}{9}+C
\end{split}\]
\end{solution}

运用公式(4.18)时，应首先将被积表达式分成函数$u(x)$和微分$\dd v=v'(x)\dd x$两部分，选取$u(x)$和$v'(x)\dd x$的原则是：
\begin{itemize}
    \item 充当微分部分的$v'(x)\dd x$, 其$v'(x)$的原函数能够比较容易地求出；
\item 
充当$u(x)$的部分经求微商后要化简，以便使积分$\Int v\dd u$较$\Int u\dd v$简单，容易求出结果．
\end{itemize}

在应用分部积分法求不定积分时，往往要连续使用这个
方法多次才能求出最后结果．

我们把能够应用分部积分法求积分的例题分为以下几类：

\begin{blk}
 {类型A} 被积函数是$x^ne^{ax}$, $x^n\sin ax$, $x^n\cos ax$, 其中$x$的方幂因子经求微商后降次，而非$x$的方幂因子的原函数是容易求得的．   
\end{blk}

\begin{example}
求$\Int x^2 e^{2x}\dd x$.
\end{example}

\begin{solution}
\[\begin{split}
    \Int x^2 e^{2x}\dd x=\int x^2 \dd\left(\frac{e^{2x}}{2}\right)&=\frac{x^2e^{2x}}{2}-\int \frac{e^{2x}}{2}\dd(x^2)\\
&=\frac{x^2e^{2x}}{2}-\int xe^{2x}\dd x=\frac{x^2e^{2x}}{2}-\left(\frac{xe^{2x}}{2}-\int \frac{e^{2x}}{2}\dd x\right)\\
&=\frac{x^2}{2}e^{2x}-\frac{x}{2}e^{2x}+\frac{e^{2x}}{4}+C\\    
&=e^{2x}\left(\frac{x^2}{2}-\frac{x}{2}+\frac{1}{4}\right)+C
\end{split}\]
\end{solution}

\begin{blk}
    {类型B} 被积函数是$x^n\sin^{-1}x$, $x^n\tan^{-1}x$, $x^n\ln x$, 其中乘积中的因子$\tan^{-1} x$, $\ln x$的导数是$x$的有理函数，$\sin^{-1}x$的
导数是$\frac{1}{\sqrt{1-x^2}}$，
它们的积分可求，而$x^n$的原函数是$\frac{x^{n+1}}{n+1}$.
\end{blk}


\begin{example}
求$\Int x\arctan x\dd x$.    
\end{example}

\begin{solution}
  \[\begin{split}
    \Int x\arctan x\dd x=\int \arctan x \dd\left(\frac{x^2}{2}\right)&=\frac{x^2\arctan x}{2}-\int \frac{x^2}{2}\dd(\arctan x)\\
    &=\frac{x^2\arctan x}{2}-\frac{1}{2}\int \frac{x^2}{1+x^2}\dd x\\
    &=\frac{x^2\arctan x}{2}-\frac{1}{2}\int \left(1-\frac{1}{1+x^2}\right)\dd x\\
    &=\frac{x^2\arctan x}{2}-\frac{1}{2}x+\frac{1}{2}\arctan x+C
  \end{split}\]  
\end{solution}


\begin{example}
    求$\Int x^3\ln x\dd x$.
\end{example}

\begin{solution}
\[\begin{split}
    \Int x^3\ln x\dd x=\int \ln x\dd\left(\frac{x^4}{4}\right)&=\frac{x^4}{4}\ln x-\int \frac{x^4}{4}\dd(\ln x)\\
    &= \frac{x^4}{4}\ln x-\int  \frac{x^4}{4}\cdot \frac{1}{x}\dd x\\
    &= \frac{x^4}{4}\ln x- \frac{x^4}{16}+C
\end{split}\]
\end{solution}

\begin{blk}
    {类型C} 被积函数是$\sin^{-1}x$, $\tan^{-1}x$, $\ln x$, 这里没有乘积的形式，不妨设$v'(x)=1$.
\end{blk}

\begin{example}
求$\Int \arcsin x\dd x$.
\end{example}

\begin{solution}
\[\begin{split}
    \Int \arcsin x\dd x=x\arcsin x-\int x\dd(\arcsin x)&=x\arcsin x-\int\frac{x}{\sqrt{1-x^2}}\dd x  \\
    &=x\arcsin x+\frac{1}{2}\int \frac{\dd(1-x^2)}{\sqrt{1-x^2}}\\
    &=x\arcsin x+\frac{1}{2}\cdot 2\sqrt{1-x^2}+C\\
    &=x\arcsin x+\sqrt{1-x^2}+C
\end{split}\]

另解：设$y=\arcsin x,\; -1\le x\le 1$，则$\sin y=x,\; -\frac{\pi}{2}\le y\le \frac{\pi}{2}$
\[\begin{split}
    \Int \arcsin x\dd x=x\arcsin x-\int x\dd(\arcsin x)&=y\sin y-\int \sin y \dd y\\
    &=y\sin y+\cos y+C\\
    &=x\arcsin x+\sqrt{1-\sin^2 y}+C\\
    &=x\arcsin x+\sqrt{1-x^2}+C
\end{split}\]
\end{solution}

\begin{example}
    求$\Int \ln^2 x\dd x$.
\end{example}

\begin{solution}
    \[\begin{split}
        \Int \ln^2 x\dd x=x(\ln x)^2-\int x\dd (\ln^2 x)&=x(\ln x)^2-\int x\cdot 2\ln x\cdot \frac{1}{x}\dd x\\
        &=x(\ln x)^2-2\int \ln x\dd x\\
        &=x(\ln x)^2-2\left[x\ln x-\int x\d(\ln x)\right]\\
        &=x(\ln x)^2-2x\ln x+2\int \dd x\\
        &=x(\ln x)^2-2x\ln x+2x+C   
    \end{split}\]
\end{solution}

\begin{blk}
    {类型D} 被积函数是$e^{ax}\sin bx$, $\sec^3 x$, $\sqrt{a^2+x^2}$. 对于这类例题，两次使用这个方法可以得到含有这个积分作未知元的一次方程．
\end{blk}

\begin{example}
    求$\Int e^{ax}\sin bx\dd x$
\end{example}

\begin{solution}
\[\begin{split}
    I=\Int e^{ax}\sin bx\dd x
    &=\int \sin bx \dd\left(\frac{e^{ax}}{a}\right)\\
&=\sin bx\cdot \frac{e^{ax}}{a}-\int \frac{e^{ax}}{a}\dd(\sin bx)\\
&=\sin bx\cdot \frac{e^{ax}}{a}-\int \frac{e^{ax}}{a}b\cos bx\dd x\\
&=\sin bx\cdot \frac{e^{ax}}{a}-\frac{b}{a}\int \cos bx\dd\left(\frac{e^{ax}}{a}\right)\\
&=\sin bx\cdot \frac{e^{ax}}{a}-\frac{b}{a}\left[\cos bx\cdot \frac{e^{ax}}{a}-\int \frac{e^{ax}}{a}\dd\left(\cos bx\right)\right]\\
&=\sin bx\cdot \frac{e^{ax}}{a}-\frac{b}{a^2}\cos bx\cdot e^{ax}-\frac{b^2}{a^2}\int e^{ax}\sin bx\dd x\\
&=\sin bx\cdot \frac{e^{ax}}{a}-\frac{b}{a^2}\cos bx\cdot e^{ax}-\frac{b^2}{a^2}I
\end{split}\]
所以：
\[\begin{split}
    I\left(1+\frac{b^2}{a^2}\right)&=\frac{1}{a}\sin bx\cdot e^{ax}-\frac{b}{a^2}\cos bx\cdot e^{ax}\\
    I&=\frac{e^{ax}(a\sin bx-b\cos bx)}{a^2+b^2}+C
\end{split}\]    
\end{solution}

\begin{example}
    求$\Int \sqrt{a^2+x^2}\dd x$.
\end{example}

\begin{solution}
\[    \begin{split}
    I =\int \sqrt{a^{2}+x^{2}} \dd x 
    &=x \sqrt{a^{2}+x^{2}}-\int x \dd\left(\sqrt{a^{2}+x^{2}}\right) \\
    &=x \sqrt{a^{2}+x^{2}}-\int \frac{x^{2}}{\sqrt{a^{2}+x^{2}}} \dd x \\
    &=x \sqrt{a^{2}+x^{2}}-\int \frac{x^{2}+a^{2}}{\sqrt{a^{2}+x^{2}}} \dd x+\int \frac{a^{2}}{\sqrt{a^{2}+x^{2}}} \dd x \\
    &=x \sqrt{a^{2}+x^{2}}-I+a^{2} \ln \left(x+\sqrt{a^{2}+x^{2}}\right)
    \end{split}\]
    所以
\[    I=\frac{1}{2} x \sqrt{a^{2}+x^{2}}+\frac{1}{2} a^{2} \ln \left(x+\sqrt{a^{2}+x^{2}}\right)\]
\end{solution}

应用分部积分公式(4.18)还能导出一些求积分的递归公式．所谓求积分的递归公式就是建立一个被积函数含有正整数$n$的积分与被积函数含有非负整数$n-1$或$n-2$的积分的关系式．每次要求$n=k$时的积分，重复利用这个关系式总能回到当$n=0, 1, 2$时的起始积分，于是所求积分总能用$n$的起始值$n=0, 1, 2$中的一个已知积分表示出来．这也就是递归的
意思．
    
\begin{example}
    求$\Int \sin^n x\dd x$.
\end{example}

\begin{solution}
当$n=0$时，$I_0=\Int (\sin x)^0\dd x=\Int \dd x=x+C$

又：
\[
\begin{split}
I_{n} =\int \sin ^{n} x \dd x&=\int \sin ^{n-1} x \dd(-\cos x) \\
&=-\cos x \sin ^{n-1} x+\int \cos x \dd\left(\sin ^{n-1} x\right) \\
&=-\cos x \sin ^{n-1} x+(n-1) \int \cos x \sin ^{n-2} x \cdot \cos x \dd x \\
&=-\cos x \sin ^{n-1} x+(n-1) \int \sin ^{n-2} x\left(1-\sin ^{2} x\right) \dd x
\end{split}
\]
所以
$$
I_{n}=-\cos x \sin ^{n-1} x+(n-1) I_{n-2}-(n-1) I_{n}
$$
即
\begin{equation}
   n I_{n}=-\cos x \sin ^{n-1} x+(n-1) I_{n-2} 
\end{equation}

现在我们用递归关系(4.19)来求 $\Int \sin ^{4} x \dd x$. 
设 $I_{4}=\Int \sin ^{4} x \dd x$, 于是
\begin{align}
4 I_{4}&=-\cos x \sin ^{3} x+3 I_{2} \\
2 I_{2}&=-\cos x \sin x+I_{0} 
\end{align}
将(4.21)代入(4.20)得：
\[\begin{split}
    4I_4&=-\cos x\sin^3 x+3\left(-\frac{\sin x\cos x}{2}+\frac{1}{2}I_0\right)\\
    &=-\cos x\sin^3 x-\frac{3}{2}\cos x\sin x+\frac{3}{2}x+C
\end{split}\]
所以
\[I_4=\int \sin^4 x\dd x=-\frac{1}{4}\cos x\sin^3 x-\frac{3}{8}\cos x\sin x+\frac{3}{8}x+C\]


\begin{ex}
\begin{enumerate}
    \item 求下列不定积分：
\begin{multicols}{2}
\begin{enumerate}
    \item $\Int x \sin x\dd x$
    \item $\Int x \cos 3 x\dd x$
    \item $\Int x \sin ^{2} \frac{x}{2}\dd x$
    \item $\Int x e^{-x}\dd x$
    \item $\Int (x-1) e^{x}\dd x$
    \item $\Int \left(x^{2}+x\right) e^{-x}\dd x$
    \item $\Int  x \sec ^{2} x\dd x$
    \item $\Int  x \arcsin x\dd x$
    \item $\Int  \frac{\arcsin x}{x^{2}}\dd x$
    \item $\Int  \frac{\arctan x}{x^{2}}\dd x$
    \item $\Int  x^{2} e^{3 x}\dd x$
    \item $\Int  e^{a x} \sin ^{2} b x\dd x$
    \item $\Int  \sin x \cdot \ln (\tan x)\dd x$
    \item $\Int  x^{2} \arctan x\dd x$
    \item $\Int  \frac{x^{2} e^{x}\dd x}{(x+2)^{2}}$
    \item $\Int \frac{x e^{x}}{(x+1)^{2}}\dd x$
\end{enumerate}
\end{multicols}
 \item   求下列不定积分:    
   \begin{multicols}{2}
        \begin{enumerate}
    \item $\Int \arctan \sqrt{x}\dd x$
    \item $\Int \ln ^{2} x \dd {x}$
    \item $\Int \ln \left(x+\sqrt{1+x^{2}}\right)\dd x$
    \item $\Int (\arcsin x)^{2}\dd x$
    \item $\Int \sin(\ln x)\dd x$
    \item $\Int \sqrt{4+x^2}\dd x$
\end{enumerate}
\end{multicols}
\end{enumerate}
\end{ex}
\end{solution}
    
\section{定积分的计算与应用}

\subsection{定积分的计算}

由微积分的基本定理得知，定积分$\Int^b_a f(x)\dd x$的计算归结为不定积分$\Int f(x)\dd x$的计算，由后者计算得到$f(x)$的原函数$F(x)$, 则
\[\Int^b_a f(x)\dd x=F (x)\Bigg|^b_a=F (b) -F (a)\] 
于是计算不定积分的各种方法也可移植过来求定积分．

\subsubsection{定积分的换元积分法}

第一换元法：设$\Int f(u)\dd u=F(u)$, $u=\varphi(x)$可微，
那么
\[\begin{split}
    \int^b_a f(\varphi(x))\varphi'(x)\dd x&=\int^{\varphi(b)}_{\varphi(a)}f(u)\dd u\\
    &=F(u)\Big|^{\varphi(b)}_{\varphi(a)}=F[{\varphi(b)}]-F[{\varphi(a)}]
\end{split}\]

\begin{example}
    求$\Int^{\tfrac{\pi}{2}}_0 \sin^3x \cos x\dd x$.
\end{example}

\begin{solution}
    因为
\[\int \sin^3x\cos x\dd x= \int \sin^3x\dd (\sin x) =\frac{1}{4}\sin^4x+C\]
所以
\[\begin{split}
    \Int^{\tfrac{\pi}{2}}_0 \sin^3x \cos x\dd x&=\frac{1}{4}\sin^4 x\Big|^{\tfrac{\pi}{2}}_0 \\
    &=\frac{1}{4}\left[\sin^4\frac{\pi}{4}-0\right]=\frac{1}{4}
\end{split}\]
\end{solution}

\begin{blk}
 {定理1（第二种变量代换法）}设$f(x)$在$[a,b]$上连续，作代换$x=\varphi (t)$, 其中$\varphi (t)$在某一闭区间$[\alpha,\beta]$上有连续导数$\varphi '(t)$, 当$\alpha\le t\le \beta$时，$a\le \varphi (t)\le b$, 且$\varphi (\alpha)=a$, $\varphi (\beta)=b$, 那么
\[\int^b_a f(x) \dd x= \int^{\beta}_{\alpha} f [\varphi  (t) ] \varphi ' (t) \dd t\]   
\end{blk}

\begin{proof}
    设$f(x)$的一个原函数是$G(x)$, 即有
\[\frac{\dd}{\dd x}G (x) =f (x)\]
于是由微积分基本定理得
\begin{equation}
    \int^b_a f (x)\dd x=G(x)\Big|^b_a=G(b) -G (a)
\end{equation}
根据复合函数求导法则，知
\[\begin{split}
    \frac{\dd}{ \dd t}G[\varphi  (t) ] &=\frac{\dd}{ \dd x}G (x)\cdot  \frac{\dd}{ \dd t}\varphi  (t)\\
    &=f(x)\cdot \varphi'(t)=f[\varphi(t)]\varphi'(t)
\end{split}\]
所以$G[\varphi (t)]$是$f[\varphi (t)]\varphi '(t)$的一个原函数，从而
\begin{equation}
    \begin{split}
\int^{\beta}_{\alpha} f [\varphi  (t) ] \varphi ' (t) \dd t&=G [\varphi  (t) ]\Big|^{\beta}_{\alpha}\\
&=       G [\varphi  (\beta) ] -G [\varphi  (\alpha) ]=G (b) -G (a)
    \end{split}
\end{equation}
因为(4.22)(4.23)两个积分相等，所以
\[\int^b_a f(x) \dd x= \int^{\beta}_{\alpha} f [\varphi  (t) ] \varphi ' (t) \dd t\]
\end{proof}

这个公式与不定积分的换元公式很类似，所不同的是，后者最后需要将变量还原，而现在只须把积分限作相应的改变．

\begin{example}
    求$\Int^a_0 \sqrt{a^2-x^2}\dd x$
\end{example}

\begin{solution}
作代换$x=a\sin t$，当$x$从0变到$a$时，相应地$t$自0变到$\frac{\pi}{2}$，于是
\[\begin{split}
    \Int^a_0 \sqrt{a^2-x^2}\dd x=a^2\int^{\tfrac{\pi}{2}}_0 \cos ^2 t\dd t&=a^2\int^{\tfrac{\pi}{2}}_0 \frac{1+\cos 2t}{2}\dd t\\
    &=\frac{a^2}{2}\left(t+\frac{1}{2}\sin 2t\right)\Bigg|^{\tfrac{\pi}{2}}_0=\frac{\pi a^2}{4} 
\end{split}\]
\end{solution}

\begin{example}
    求$\Int^4_0 \frac{1}{1+\sqrt{x}}\dd x$
\end{example}

\begin{solution}
    设$x=u^2$，于是
\[\frac{1}{1+\sqrt{x}}=\frac{1}{1+u},\qquad \dd x=2u\dd u\]
定限：当$x=0$时，从变换$\sqrt{x}=u$中解出$u=0$；当$x=4$时，解出$u=2$. 于是
\[\begin{split}
    \Int^4_0 \frac{1}{1+\sqrt{x}}\dd x=\int^2_0 \frac{1}{1+u}\cdot 2u\dd u&=2\int^2_0 \left(1-\frac{1}{1+u}\right)\dd u\\
    &=2\left[u-\ln(1+u)\right]\Big|^2_0=2(2-\ln 3)
\end{split}\]
\end{solution}

\begin{example}
求证：$\Int^a_0 f(x)\dd x=\Int^a_0 f(a-x)\dd x$
\end{example}

\begin{proof}
    设$x=a-t$，则$\dd x=-\dd t$，且当$x=0$时，$t=a$；当$x=a$时，$t=0$. 于是：
\[\begin{split}
    \Int^a_0 f(x)\dd x=\int^0_a f(a-t)(-\dd t)&=-\int^0_a f(a-t)\dd t\\
    &=\int^a_0 f(a-t)\dd t=\Int^a_0 f(x)\dd x
\end{split}\]
\end{proof}

我们在第三节曾指出，使用第二种变量代换时，要注意条件是否满足，否则会出现矛盾，这里举一反例．

求$\Int^{1}_{-1}\frac{\dd x}{1+x^2}$时，如作变换$x=\frac{1}{t}$，则
\[\Int^{1}_{-1}\frac{\dd x}{1+x^2}=-\Int^{1}_{-1}\frac{\dd t}{1+t^2}\]
于是$\Int^{1}_{-1}\frac{\dd x}{1+x^2}=0$，这当然不对，问题在于，代换$x=\frac{1}{t}$
在$[-1,1]$上并不是连续变化的．


\subsubsection{定积分的分部积分法}
\begin{blk}
    {定理2} 设$u=u(x)$, $v=v(x)$在$[a,b]$上有连续的导数$u'(x)$, $v'(x)$, 则有分部积分公式
\[\int^b_a u\dd v=u\cdot v\Big|^b_a- \int^b_a v\dd u\]
\end{blk}

\begin{proof}
    由牛顿-莱布尼兹公式，显然有
\[\begin{split}
    u (x) v (x)\Big|^b_a &=\int^b_a (u (x) v (x) ) '\dd x\\
&=\int^b_a u' (x) v (x) \dd x+\int^b_a u (x) v' (x) \dd x
\end{split}\]
移项，得
\[\int^b_a u (x) \dd [v (x) ] =u (b) v (b) -u (a) v (a) -\int^b_a v (x) \dd [u (x) ] \]
\end{proof}

这个公式与不定积分的分部积分公式很相似，但是，这里每一项都带着积分限．

\begin{example}
求$\Int^{\sqrt{3}}_{-\sqrt{3}}|\arctan x|\dd x$
\end{example}

\begin{solution}
因为$|\arctan x|$是偶函数，所以
\[I=\Int^{\sqrt{3}}_{-\sqrt{3}}|\arctan x|\dd x=2\Int^{\sqrt{3}}_{0}|\arctan x|\dd x=2\Int^{\sqrt{3}}_{0}\arctan x\dd x\]
设$u(x)=\arctan x$, $v'(x)=1$, 则$v(x)=x$，利用分部积分公式，得到：
\[\begin{split}
    I=2\Int^{\sqrt{3}}_{0}\arctan x\dd x &=2\left[x\arctan x\Big|^{\sqrt{3}}_0-\int^{\sqrt{3}}_0 x\dd (\arctan x)\right]\\
    &=2\left[\sqrt{3}\arctan\sqrt{3}-\int^{\sqrt{3}}_0 x\cdot \frac{1}{1+x^2}\dd x \right]\\
    &=2\sqrt{3}\cdot \frac{\pi}{3}-\int^{\sqrt{3}}_0 \frac{\dd(1+x^2)}{1+x^2}\\
    &=\frac{2\sqrt{3}}{3}\pi -\ln(1+x^2)\Big|^{\sqrt{3}}_0= \frac{2\sqrt{3}}{3}\pi -\ln 4
\end{split}\]
\end{solution}


\begin{example}
    求$\Int^{\tfrac{\pi}{2}}_0 \sin^n \theta\dd \theta$.
\end{example}

\begin{solution}
    设$I_n=\Int^{\tfrac{\pi}{2}}_0 \sin^n \theta\dd \theta$，利用例4.42的结果，得到：
\[nI_n=\left[-\cos\theta\sin^{n-1}\theta\right]\Big|^{\tfrac{\pi}{2}_0} +(n-1)I_{n-2}\]
因为方括号的式子对于积分上，下限的值等于零，所以上面等式化简为
\[I_n=\frac{n-1}{n}I_{n-2}\]
分别考虑$n$是偶数和奇数两种情况：
\begin{enumerate}
    \item 若$n$是偶数，那么
\[\begin{split}
    I_n=\frac{n-1}{n}I_{n-2}&=\frac{n-1}{n}\cdot \frac{n-3}{n-2}I_{n-4}\\
    &=\cdots=\frac{n-1}{n}\cdot \frac{n-3}{n-2}\cdot \frac{n-5}{n-4}\cdots \frac{1}{2}I_{0}
\end{split}
    \]
此处$I_0=\Int^{\tfrac{\pi}{2}}_0 1\cdot \dd \theta=\frac{\pi}{2}$，所以
\[I_n=\frac{(n-1)(n-3)\cdots 3\cdot 1}{n(n-2)\cdots 4\cdot 2}\cdot \frac{\pi}{2}\]

\item 若$n$是奇数，由于
\[I_1=\Int^{\tfrac{\pi}{2}}_0 \sin\theta\dd\theta =(-\cos\theta)\Big|^{\tfrac{\pi}{2}}_0=1\]
所以
\[I_n=\frac{(n-1)(n-3)\cdots 4\cdot 2}{n(n-2)\cdots 5\cdot 3}\]
\end{enumerate}
\end{solution}

\begin{ex}
\begin{enumerate}
    \item 求下列定积分：
\begin{multicols}{2}
\begin{enumerate}
    \item $\Int_{0}^{\tfrac{\pi}{4}}(1+\sin \theta)^{2} \dd \theta $ 
    \item $\Int_{0}^{\tfrac{\pi}{4}} \frac{\cos \theta \dd \theta}{1+\sin ^{2} \theta}$
    \item $\Int_{0}^{1} \frac{x^{2}}{\left(1+x^{2}\right)^{2}} \dd x$
    \item $\Int_{0}^{\tfrac{1}{\sqrt{2}}} \frac{x^{2}}{\sqrt{1-x^{2}}} \dd x$
    \item $\Int_{\tfrac{1}{\sqrt{2}}}^{1} \sqrt{1-x^{2}} \dd x$
    \item $\Int_{0}^{x} \frac{\cos ^{2} x \dd x}{a^{2} \cos ^{2} x+b^{2} \sin x}$
    \item $\Int_{0}^{2 n}[\sin \omega t -\sin (\omega t+\varphi)]^{2} \dd t$
    \item $\Int_{2}^{4} \sqrt{(4-x)(x-2)} \dd x$
    \item $\Int_{1}^{5} \frac{1}{\sqrt{2 x-1}} \dd x$ 
    \item $\Int_{0}^{1} \frac{x+1}{\sqrt{1 - x^{2}}} \dd x$
    \item $\Int_{\tfrac{\pi}{6}}^{\tfrac{\pi}{2}} \cot \theta \dd \theta$
    \item $\Int_{1}^{0} \frac{x^{2}}{ 1+x^{3}} \dd x$
    \item $\Int^{\tfrac{\pi}{2}}_0 \frac{\sin x}{\cos^2 x}\dd x$
    \item $\Int^{2a}_0 x^2\sqrt{2ax-x^2}\dd x$
\end{enumerate}
\end{multicols}

\item 证明：$\Int^{\tfrac{\pi}{2}}_0 \frac{\cos x}{\cos x+\sin x}\dd x=\Int^{\tfrac{\pi}{2}}_0 \frac{\sin x}{\cos x+\sin x}\dd x$

（提示：用变量代换$x=\frac{\pi}{2}-y$.）
\item \begin{enumerate}
    \item 求常数$A,B,C,D$使等式$$\frac{1+2x^2}{(1-x)(1+x)^2}=A+\frac{B}{1-x}+\frac{C}{(1+x)^2}+\frac{D}{1+x}$$成立，并计算
    \[
    \Int^{\tfrac{1}{2}}_0 \frac{1+2x^3}{(1-x)(1+x)^2}\dd x.
    \]
    \item 求常数$a$，使$\Int^1_0 \frac{x-a}{(x+1)(3x+1)}\dd x=0$.
\end{enumerate}
\item 求极限：$\Lim_{n\to\infty}\left(\frac{1}{n+1}+\frac{1}{n+2}+\cdots +\frac{1}{n+n}\right)$
\item 求一三次多项式$f(x)=ax^3+bx^2+cx+d$, 使它满足下
列条件：
\begin{enumerate}
    \item $f (-1) =f (1) =0$;
    \item $\Int^1_{-1} f(x)\dd x=6$.
\end{enumerate}
\item 
\begin{enumerate}
\item 在同一个坐标系内，作$y=(x-1)(x-2)$和$y=\Int^x_0 (t-1)(t-2)\dd t$的函数图象．
\item 求函数$\Int^x_0 (t-1)(t-2)\dd t$在闭区间$[0, 3]$上的最大值和最小值；
\item 求在曲线$y=(x-1)(x-2)$, $x$轴和纵坐标$x=0$, $x=3$之间的区域的面积．
\end{enumerate}
\item 
\begin{enumerate}
    \item 若$f(2a-x)=-f(x)$, 求证$\Int^{2a}_0 f (x) \dd x=0$；
\item 
若$f(2a-x)=f(x)$, 求证$\Int^{2a}_0 f (x) \dd x=2\int^a_0 f (x) \dd x$. 
\item 试计算$\int^{\pi}_0 \sin^6 x\cos^3x\dd x$和$\Int^{\pi}_0 \sin^3x \cos^2 x\dd x$.
\end{enumerate}

\item 求$\Int^{\tfrac{\pi}{2}}_0 \cos^n \theta\dd \theta$.
\item 求$\Int^{\tfrac{\pi}{4}}_0 x^3\cos2x\dd x$.
\item 求$\Int^{\tfrac{3\pi}{2}}_0 \sin^3\theta \cos^2 \theta\dd \theta$.
\end{enumerate} 
\end{ex}

\subsection{平面图形的面积}
我们已经证明过在曲线$y=f(x)$与$x$轴，$x=a$, $x=b$之间的曲边梯形的有号面积等于
\[\int^b_a y\dd x=\int^b_a f(x)\dd x\]
对于任意的平面图形．我们总可以取适当的直角坐标系，使得它的边界为两条曲线
\[y=f (x) ,\qquad  y=g (x) ,\quad x\in [a,b]\]
所描述，现在我们来求这样的区域$R$的面积（图4.4）．

\begin{blk}
   {定理} 若$f$和$g$是在区间$[a,b]$上的连续函数，而且$f(x)\ge g(x)$, 那么在曲线$y=f(x)$, $y=g(x)$和$x=a$, $x=b$之间的区域$R$的面积等于
\[A=\int^b_a [f (x) -g (x) ] \dd x\] 
\end{blk}

\begin{figure}[htp]\centering
	\begin{minipage}[t]{0.48\textwidth}
		\centering
        \begin{tikzpicture}[>=latex, scale=.8]
\draw[->](-2.5,0)--(3,0)node[right]{$x$};
\draw[->](0,-3)--(0,3)node[right]{$y$};
\draw[domain=-2:2, samples=100, very thick]plot(\x, {.15*(-\x*\x-4*\x +5) })node[right]{$y=f(x)$};
\draw[domain=-2:2, samples=100, very thick]plot(\x, {.1*(\x*\x-4*\x -15)})node[right]{$y=g(x)$};
\draw[domain=-2:2, samples=100, pattern=north east lines, dashed]plot(\x, {.15*(-\x*\x-4*\x +5) })--plot(-\x, {.1*(\x*\x+4*\x -15)});
\node at (0,0)[below left]{$O$};
\draw[dashed](-1.75,-2)--(-1.75,2)node[above]{$x=a$};
\draw[dashed](1.75,-2)--(1.75,2)node[above]{$x=b$};



    \end{tikzpicture}
    \caption{}
\end{minipage}
\begin{minipage}[t]{0.48\textwidth}
    \centering
\begin{tikzpicture}[>=latex, scale=.8]
\draw[->](-2.5,0)--(3,0)node[right]{$x$};
\draw[->](0,-3)--(0,3)node[right]{$y$};
\draw[domain=-2:2, samples=100, very thick]plot(\x, {.15*(-\x*\x-4*\x +5) })node[right]{$y=f(x)-k$};
\draw[domain=-2:2, samples=100, very thick]plot(\x, {.1*(\x*\x-4*\x -15)})node[right]{$y=g(x)-k$};
\draw[domain=-2:2, samples=100, pattern=north east lines, dashed]plot(\x, {.15*(-\x*\x-4*\x +5) })--plot(-\x, {.1*(\x*\x+4*\x -15)});
\node at (0,0)[below left]{$O$};

\node at (0,3)[left]{$y'$};
\draw[thick, ->](-2.5,-2.5)--(3,-2.5)node[right]{$x'$};
\draw[dashed](-1.75,-2.5)node[below]{$a$}--(-1.75,2);
\draw[dashed](1.75,-2.5)node[below]{$b$}--(1.75,2);
\node at (0,-2.5)[below right]{$O'$};
\end{tikzpicture}
\caption{}
\end{minipage}
\end{figure}

\begin{proof}
让我们选择一个常数$k$使其小于函数$g$在区间$[a,b]$上的最小值，则对于每一个$x\in [a,b]$, $g(x)
-k\ge 0$.

又因为$f(x)\ge g(x)$, 因此对于每一个$x\in [a,b]$, 
也有$f(x)-k\ge 0$. 将坐标系$x,y$平行于$y$轴平移到新坐标系$x',y'$, 使新原点在原来坐标系中的坐标为$(0,k)$ (图4.5), 这个坐标系的平移由变换
\[\begin{cases}
    x=x'\\
    y=y'+k
\end{cases}\]
实现，于是原来的函数图象在新坐标的方程是
\[f_1 (x) =f (x) -k,\qquad g_1 (x) =g (x) -k\]
函数$f_1$和$g_1$在区间$[a,b]$上仍连续且取非负值，显然在函数$f_1$和$g_1$之间的区域$R=\{a\le x\le b,\; g_1(x)\le y\le f_1(x)\}$就是原来的区域$R=\{a\le x\le b,\; g(x)\le y\le f(x)\}$, 但是，$R'$的区域的面积等于
\[A=\int^b_a f_1(x)\dd x-\int^b_a g_1(x)\dd x\]
根据定积分的性质进一步得到
\[\begin{split}
    A&=\int^b_a [f_1(x)-g_1(x)]\dd x\\
    &=\int^b_a \Big\{\big[f(x)-k\big]-\big[g(x)-k\big]\Big\}\dd x\\
    &=\int^b_a [f(x)-g(x)]\dd x\\
\end{split}\]
\end{proof}

\begin{example}
求直线$y=x-2$和抛物线$y=2x-x^2$所围成的图形的面积$A$.
\end{example}

\begin{solution}
直线与抛物线交于$(2, 0)$, $(-1,-3)$两点，作出它的草图（图4.6）. 依本节定理，得到
\[\begin{split}
    A=\int^2_1 [(2x-x^2)-(x-2)]\dd x
    &=\int^2_1 (-x^2+x+2)\dd x\\
    &=\left[-\frac{1}{3}x^3+\frac{1}{2}x^2+2x\right]\Bigg|^2_1=\frac{9}{2}
\end{split}\]
\end{solution}


\begin{figure}[htp]\centering
	\begin{minipage}[t]{0.48\textwidth}
		\centering
        \begin{tikzpicture}[>=latex, scale=.7]
\draw[->](-2,0)--(5,0)node[right]{$x$};
\draw[->](0,-5)--(0,2.5)node[right]{$y$};
\draw[domain=-1.25:3.25, samples=100, very thick]plot(\x, {2*\x-\x*\x});
\draw[domain=-2:4, samples=10, very thick]plot(\x, {\x-2});
\node at (1,1)[above]{$(1,1)$};
\node at (2,0)[below right]{$(2,0)$};
\node at (-1,-3)[left]{$(-1,-3)$};
\node at (0,0)[above left]{$O$};
\draw[domain=-1:2, samples=100, pattern=north west lines]plot(\x, {2*\x-\x*\x})--(-1,-3);

    \end{tikzpicture}
    \caption{}
\end{minipage}
\begin{minipage}[t]{0.48\textwidth}
    \centering
\begin{tikzpicture}[>=latex, scale=.8]
    \draw[->](-2,0)--(3.5,0)node[right]{$x$};
\draw[->](0,-3.5)--(0,3.5)node[right]{$y$};
\node at (0,0)[below left]{$O$};
\draw[domain=-2.5:2.5, samples=100, very thick]plot({2-0.5*\x*\x}, \x);
\draw[domain=-2.75:2.75, samples=100, very thick]plot({1-0.25*\x*\x}, \x);
\foreach \x in {2,-2}
{
    \node at (0,\x) [left]{$\x$};
}
\node at (1,0)[below left]{1};
\node at (2,0)[below right]{2};
\draw[domain=-2:2, samples=100, pattern=north east lines]plot({2-0.5*\x*\x}, \x)--plot({1-0.25*\x*\x}, -\x);


\end{tikzpicture}
\caption{}
\end{minipage}
\end{figure}


\begin{example}
求抛物线$y^2=-4(x-1)$, $y^2=-2(x-2)$所围成的图形的面积．
\end{example}

\begin{solution}
    先作出它们的草图（图4.7）, 且解得它的两个交点$(0, 2)$和$(0,-2)$. 于是
\[\begin{split}
    A=\int^2_{-2} \left[\left(2-\frac{1}{2}y^2\right)-\left(1-\frac{1}{4}y^2\right)\right]\dd y
    &=\frac{1}{4}\int^2_{-2} (4-y^2)\dd y\\
    &=\frac{1}{4}\left[4y-\frac{1}{3}y^3\right]\Bigg|^2_{-2}=\frac{8}{3}.
\end{split}\]
\end{solution}

\begin{example}
求半径等于$r$, 中心角等于$\alpha \; \left(0<\alpha\le \frac{\pi}{2}\right)$的扇形的面积．
\end{example}

\begin{figure}[htp]
    \centering
\begin{tikzpicture}[>=latex, scale=1.5]
\draw[->](-1.5,0)--(2,0)node[right]{$x$};
\draw[->](0,-1.5)--(0,1.5)node[right]{$y$};
\draw[very thick] (0,0) circle (1);
\draw[pattern=north west lines](0,0)--(1,0)node[below right]{$B$} arc (0:60:1)node[above right]{$P(r\cos\alpha, r\sin\alpha)$}--(0,0);
\node at (0,0)[below left]{$O$};
\draw(.5,0)node[below]{$C$}--(60:1);
\draw[<->](.2,0) arc (0:60:.2)node[right]{$\alpha$};

\end{tikzpicture}
    \caption{}
\end{figure}

\begin{solution}
    如图4.8, 扇形$OBP$的区域划分为$\triangle OCP$的区域与曲边三角形$CBP$的区域的和，因此
\[A (OBP) =A (OCP) +A (CBP) \]
因为$P$点坐标是$(r\cos\alpha, r\sin\alpha)$，所以
\[A(OCP)=\frac{1}{2}(r\cos\alpha)\cdot (r\sin\alpha)=\frac{r^2}{4}\sin 2\alpha\]
又扇形的圆方程是$x^2+y^2=r^2$，所以
\[A(CBP)=\int^r_{r\cos \alpha}\sqrt{r^2-x^2}\dd x\]
从而
\[A(OBP)=\frac{r^2}{4}\sin2\alpha+\int^r_{r\cos \alpha}\sqrt{r^2-x^2}\dd x\]

令$x=r\sin\theta,\; \theta\in\left[-\frac{\pi}{2},\frac{\pi}{2}\right]$，则
\[\theta=\arcsin\frac{x}{r},\qquad \dd x=r\cos\theta \dd\theta\]
且当$x=r\cos \alpha$时，
\[\theta=\arcsin(\cos\alpha)=\arcsin\left[\sin\left(\frac{\pi}{2}-\alpha\right)\right]=\frac{\pi}{2}-\alpha\]
当$x=r$时，$\theta=\arcsin 1=\frac{\pi}{2}$

所以
\[\begin{split}
    A(OBP)&=\frac{r^2}{4}\sin2\alpha+\int^{\tfrac{\pi}{2}}_{\tfrac{\pi}{2}-\alpha}(r\cos\theta)(r\cos\theta)\dd \theta\\
&=\frac{r^2}{4}\sin2\alpha+r^2\int^{\tfrac{\pi}{2}}_{\tfrac{\pi}{2}-\alpha}\cos^2\theta\dd\theta\\ 
&=\frac{r^2}{4}\sin2\alpha+r^2\left.\left[\frac{\theta}{2}+\frac{\sin 2\theta}{4}\right]\right|^{\tfrac{\pi}{2}}_{\tfrac{\pi}{2}-\alpha}\\
&=\frac{r^2}{2}\alpha
\end{split}\]

当$\alpha=2\pi$时，由上面公式就得到圆面积$\pi r^2$.
\end{solution}

\begin{example}
$P$点是位于曲线$y=x^2$上的动点，从原点$O$向点$A(1, 1)$移动，过二点$O,P$引直线，它与曲线$y=x^2$及直线$x=1$所围图形的面积为图4.9中阴影部分所示，求使这部分面积最小时$P$点的坐标．
\end{example}

\begin{figure}[htp]\centering
        \begin{tikzpicture}[>=latex, scale=3]
\draw[->](-.2,0)--(1.5,0)node[right]{$x$};
\draw[->](0,0)node[below]{$O$}--(0,1.5)node[right]{$y$};
\draw(0,1)--(1,1)node[right]{$A(1,1)$}--(1,0);
\draw[domain=0:1.15, samples=100, very thick]plot(\x, {\x*\x})node[above right]{$y=x^2$};
\draw[domain=0:1, samples=100, thick, pattern=north west lines]plot(\x, {\x*\x})--(1,.8)--(0,0);
\node at (.8,.64)[left]{$P(t,t^2)$};
        \end{tikzpicture}
        \caption{}
    \end{figure}

\begin{solution}
    设$P$点的坐标是$(t,t^2),\; (0\le t\le 1)$. 直线$OP$的方程是$y=tx$，设阴影部分的面积为$S$，则
\[\begin{split}
S&=\int^t_0 (tx-x^2)\dd x+\int^1_t (x^2-tx)\dd x\\
&=\left[\frac{t}{2}x^2-\frac{x^3}{3}\right]\Bigg|^t_0+\left[\frac{x^3}{3}-\frac{t}{2}x^2\right]\Bigg|^1_t\\
&=\frac{t^3}{3}-\frac{t}{2}+\frac{1}{3}
\end{split}\]
\[S'=t^2-\frac{1}{2}=\left(t+\frac{\sqrt{2}}{2}\right)\left(t-\frac{\sqrt{2}}{2}\right)\]
由$S'$符号的变化，当$t=\frac{\sqrt{2}}{2}$时，$S$极小也就是最小，故所求点的坐标是$\left(\frac{\sqrt{2}}{2},\frac{1}{2}\right)$.
\end{solution}

\begin{example}
    求椭圆$\frac{x^2}{a^2}+\frac{y^2}{b^2}=1\; (a>b>0)$的面积$A$（图4.10）.
\end{example}

\begin{figure}[htp]
    \centering
\begin{tikzpicture}[>=latex]
    \draw[->](-3,0)--(3,0)node[right]{$x$};
\draw[->](0,-2)--(0,2)node[right]{$y$};
\node at (0,0) [below left]{$O$};
\draw[very thick](0,0) ellipse [x radius=2.2, y radius=1.3];

\end{tikzpicture}
    \caption{}
\end{figure}

\begin{solution}
由于图形的对称性，只要求出椭圆在第一象限内的面积，再乘以4就是该椭圆的面积，所以
\[\frac{A}{4}=\int^a_0 y\dd x\]
这里将椭圆用参数方程
$\begin{cases}
    x=a\cos t\\ y=b\sin t
\end{cases}$
表出，当$x=a$时，$t=\arccos 1=0$; 当$x=0$时，$t=\arccos 0=\frac{\pi}{2}$, 于是
\[\begin{split}
    \frac{A}{4}=\int^0_{\tfrac{\pi}{2}}b\sin t \dd(a\cos t)&=\int^0_{\tfrac{\pi}{2}}(b\sin t )(-a\sin t)\dd t\\
&=ab \int^{\tfrac{\pi}{2}}_0 \sin^2 t\dd t=\frac{1}{4}ab\pi 
\end{split}\]
所以：$A=\pi ab$.
\end{solution}

\begin{ex}
\begin{enumerate}
    \item 求在曲线$y=\sin2x$, $y=0$, $x\in[0, 2\pi]$之间的面积．
    \item 求下列曲线所围图形的面积：
\begin{enumerate}
    \item $y=8-\frac{1}{2}x^2$和$y=3.2$ 
    \item $y=x^2-5x+7$和$y=-2x^2+10x-5$
    \item $y=\sin x$和$y=\frac{2}{\pi} x,\; x\ge 0$
\end{enumerate}
\item 求以曲线$y=\sin x\; \left(0\le x\le \frac{\pi}{4}\right)$；$y=\cos x\; \left(\frac{\pi}{4}\le x\le \frac{\pi}{2}\right)$和$y=0$为边界的区域的面积．

\item \begin{enumerate}
    \item 证明点$x=t^2-1$, $y=t(t^2-1)$在曲线$y^2=x^2(x+1)$上；
    \item 当$t$由$t=-2$变化到$t=2$时，作出曲线$\begin{cases}
        x=t^2-1\\x=t (t^2-1)
    \end{cases}$的图象；
\item 求由曲线的环路所围区域的面积．
\end{enumerate}

\item 在曲线$y=x^2-2x+2$与过点$(2, 3)$的直线所围的图形
中，求图形面积最小时的直线方程．
\item 求抛物线$y=x^2$及过此抛物线上点$(2, 4)$的切线和$x$轴所围图形的面积．
\item 已知曲线$C:\; y=x^2-7x+10\; (x>0)$与点$A(0, 1)$, 
回答下列问题：
\begin{enumerate}
    \item 求过点$A$所引曲线$C$的切线方程以及切点$T$的坐标．
    \item 设$S$是直线$x=1$与曲线$C$的交点．若曲线$C$上的动点$P$从$S$点运动到$T$点时，将线段$AP$所在的区域用阴影表示出来，并求出这个区域的面积．
\end{enumerate}

\item 求通过$(0, 0)$, $(1, 2)$两点的抛物线，要求它具有以
下的性质：
\begin{enumerate}
    \item 它的对称轴平行于$y$轴，且凸向上，
    \item 它与$x$轴所围区域的面积最小．
\end{enumerate}

\item 对称轴平行$y$轴的抛物线与圆$x^2+y^2=4$切于$A(0, 2)$点并通过$B(-2, 0)$点，计算抛物线与圆所围区域的面积．
\end{enumerate}
\end{ex}

\subsection{利用横断面算体积法}
对于空间的一个立体，我们取一条直线作为坐标轴，取直线上一点作为原点（图4.11），假定这立体紧夹于在$x=a$与$x=b$两点所作的$x$轴的垂直平面之间，又如果在离原点$x$处作这轴的一个垂直平面，并且该垂直平面截得立体的截面的面积是$x$的连续函数$A(x),\; a\le x\le b$, 那么立体的体积可由定积分来计算：
\begin{figure}[htp]
    \centering
\begin{tikzpicture}[>=latex]
    \draw[->](-.5,0)--(6,0)node[right]{$x$};
\tkzDefPoints{.8/0/a, 2/0/A, 3.3/0/x_i, 5/0/b, 0/0/O}
\tkzDrawPoints(a,A,x_i,b,O)
\tkzLabelPoints[below](a,x_i,b,O)
\tkzLabelPoint[below](A){$x_{i-1}$}
\node at (0,0){};
\draw[line width = 1pt] plot[smooth] coordinates{(.8,0)(1.2,.8)(2,1.3)(3.3,1.15)(5,0)(3.3,-1.15)(2,-1.3)(1.2,-.8)(.8,0)
};
\draw[dashed](A) ellipse [x radius=.35, y radius=1.3]; 
\draw[dashed](x_i) ellipse [x radius=.3, y radius=1.15]; 
\draw[thick](2,1.3) arc [start angle =90, end angle =270, x radius=.35, y radius=1.3];
\draw[thick](3.3,1.15) arc [start angle =90, end angle =270, x radius=.3, y radius=1.15];

\end{tikzpicture}
    \caption{}
\end{figure}

\begin{enumerate}
    \item 分$[a,b]$为$n$份，其分点
\[a=x_0<x_1<x_2<\cdots<x_n=b\]
    记$\|P_n\|=\max(x_i-x_{i-1})\quad i=1, 2,\ldots,n$.

    以平面$x=x_i,\; i=0, 1, 2,\ldots,n$截此立体为$n$个小单元，于是整个体积被看作这$n$个小单元的和．
\item  求阶梯函数的总和作为所求立体体积的近似和．考虑这样一个单元，它介于两个截面$x=x_{i-1},x=x_i$之间，于是以$(x_i-x_{i-1})$为高，以$A(x_{i-1})$作为底面积的直棱柱的体积是
  \[  A (x_{i-1}) (x_i-x_{i-1})\]
    于是，我们得到所求的体积$V$的近似表达式
    \[\sum^n_{i=1}A (x_{i-1}) (x_i-x_{i-1})\] 
    \item 对和取极限：当断面无限增加，$\|P_n\|\to 0$时，有
\[V=\lim_{n\to\infty}\sum^n_{i=1}A (x_{i-1}) (x_i-x_{i-1})=\int^b_a A(x)\dd x\]
\end{enumerate}

现在将上面的讨论总结为下面的定理．

\begin{blk}
   {定理} 如果已知一个给定的立体的垂直于一定方向所有的横断面的面积，取这些横断面的垂直方向作为$x$轴的方向，则这立体的体积由公式
\[V=\int^b_a A (x) \dd x\]
表达．其中$A(x)$是横坐标为$x$的横断面的面积，$a$、$b$为这立体的两端断面的横坐标． 
\end{blk}

\begin{example}
    求底边长为$a$, 高为$h$的正四棱锥的体积．
\end{example}

\begin{figure}[htp]
    \centering
\begin{tikzpicture}[>=latex]
    \tkzDefPoints{0/0/A, 3/0/B, 4/1/C}
    \tkzDefPointsBy[translation= from B to C](A){D}
    \tkzInterLL(A,C)(B,D)  \tkzGetPoint{O}
    \tkzDefPoints{0/3.5/E, 5/3.5/F}
    \tkzDefPointBy[projection = onto E--F](O) \tkzGetPoint{OO}
    \draw[|<->|](.7,1.4)--node[fill=white]{$y$}(2.7,1.4);
    \tkzDrawSegments(OO,A OO,B OO,C A,B B,C)
    \tkzDrawSegments[dashed](OO,D OO,O C,D D,A)
    
    \foreach \x in {A,B,C,D,O}
    {
        \tkzDefPointWith[linear, K=.66](OO,\x)  \tkzGetPoint{\x'}
        \tkzDefPointWith[linear, K=.33](OO,\x)  \tkzGetPoint{\x''}
    }
    \tkzDrawSegments(A',B' B',C' A'',B'' B'',C'')
    \tkzDrawSegments[dashed](C',D' D',A' C'',D'' D'',A'')
\draw[->](OO)--(2,4)node[right]{$x$};
\draw(-1,3.5)--(5,3.5);
\draw[dashed](.2,.5)--(O);
\draw(-1,0.5)--(.2,.5);
\draw[<->](-.8,0.5)--node[fill=white]{$h$}(-.8,3.5);
\draw[<->](4.8,0.5)--node[fill=white]{$x$}(4.8,1.5);
\draw[<->](4.8,1.5)--node[fill=white]{$h-x$}(4.8,3.5);
\draw[<->](3.8,1.5)--node[fill=white]{$\Delta x$}(3.8,2.5);
\tkzDefMidPoint(B,C) \tkzGetPoint{G}
\tkzDefMidPoint(B',C') \tkzGetPoint{G'}
\tkzDefMidPoint(B'',C'') \tkzGetPoint{G''}
\tkzDrawSegments[dashed](O,G O',G' O'',G'')
\tkzDefPointWith[linear, K=2](O,G) \tkzGetPoint{H}
\tkzDefPointWith[linear, K=3](O',G') \tkzGetPoint{H'}
\tkzDefPointWith[linear, K=4.5](O'',G'') \tkzGetPoint{H''}
\tkzDrawSegments(G,H G',H' G'',H'')
\tkzLabelPoints[below](O)

\draw[|<->|](0,-0.25)--node[fill=white]{$a$}(3,-.25);


\end{tikzpicture}
    \caption{}
\end{figure}

\begin{solution}
    设$x$为从底面到截面的距离，$y$为截面正方形的边长（图4.12），于是
\[\frac{h-x}{h}=\frac{y}{a}\]
即：$y=\frac{a}{h}(h-x)$.

截面面积
\[A(x)=\frac{a^2}{h^2}(h-x)^2\]
所以
\[\begin{split}
    V&=\int^h_0 \left(\frac{a}{h}\right)^2(h-x)^2 \dd x\\
    &=-\left(\frac{a}{h}\right)^2\frac{(h-x)^3}{3}\Bigg|^h_0=\frac{1}{3}a^2 h
\end{split}\]
\end{solution}

如果立体是一个旋转体，即先给出一条平面曲线
$y=f(x)$, 绕$x$轴旋转所得出的立体，我们的横断面就是以$y$为半径的圆，所以这立体的体积就等于
\[V=\pi \int^b_a y^2\dd x\]
如果$c\le y\le d$，$x=g(y)$绕$y$轴旋转，则
\[V=\pi \int^d_c x^2\dd y\]


\begin{example}
若球的半径是$a$, 球缺的高是$h$, 求球缺的体积．
\end{example}

\begin{figure}[htp]
    \centering
\begin{tikzpicture}[>=latex, scale=1.5]
    \draw[->](-1.5,0)--(2,0)node[right]{$x$};
    \draw[->](0,-1.5)--(0,1.5)node[right]{$y$};
    \draw[very thick](0,0)circle(1);
\draw(.4,1)--(.4,-1.5);
\draw(1,0)--(1,-1.5);
\draw[<->](0,-1.3)--node[below]{$a-h$}(.4,-1.3);
\draw[<->](0.4,-1.1)--node[fill=white]{$h$}(1,-1.1);
\draw[<->](-.2,0)--node[fill=white]{$a$}(-.2,1);

\draw[pattern=north east lines](45:1) arc (45:50:1)--(.64,0)--(1/1.414,0)--(45:1);
\draw[<->](0,1.3)--node[below]{$x$}(.64,1.3);
\draw(.64,1.5)--(.64,0);
\end{tikzpicture}
    \caption{}
\end{figure}


\begin{solution}
球是由圆$x^2+y^2=a^2$绕$x$轴旋转所得到的立体，且球心到球缺底面的距离为$a-h$, 设球心到截面的距离为$x$, 则截面圆的面积$A (x) =\pi y^2=\pi  (a^2-x^2)$.（图4.13）

所以
\[\begin{split}
    V=\pi\int^a_{a-h}(a^2-x^2)\dd x&=\pi\left[a^2 x-\frac{1}{3}x^3\right]\Bigg|^a_{a-h}\\
    &=\frac{\pi h^2}{3}(3a-h)
\end{split} \]
\end{solution}

\begin{example}
    计算底面半径分别为$r$和$R$, 高为$h$的圆台体积$V$.
\end{example}

\begin{solution}
此圆台可看作$xy$平面上的直线段
\[y=\frac{R-r}{h}x+r,\; (0\le x\le h),\quad x=0,\quad x=h,\quad y=0\]
所围成的直角梯形绕$x$轴旋转所得到的立体，从而
\[\begin{split}
    V&=\pi\int^h_0 \left(\frac{R-r}{h}x+r\right)^2 \dd x\\
    &=\pi\left[\left(\frac{R-r}{h}\right)^2\frac{x^3}{3}+2r\cdot \frac{R-r}{h}\cdot \frac{x^2}{2}+r^2 x\right]\Bigg|^h_0\\
    &=\frac{\pi h}{3}(R^2+Rr+r^2)
\end{split}\]
\end{solution}

\begin{example}
在第一象限中，由$y=\frac{x}{1-x}\; (x\ne 1)$, $y=1$, $y=0$和$x=1$所围成的图形绕着直线$x=1$旋转，求所生成的旋转体的体积（图4.14）．
\end{example}

\begin{figure}[htp]
    \centering
\begin{tikzpicture}[>=latex, scale=2.5]
    \draw[->](-.25,0)--(1.8,0)node[right]{$x$};
    \draw[->](0,-.25)--(0,1.8)node[right]{$y$};
\draw[domain=0:.6, samples=100, very thick]plot(\x, {\x/(1-\x)});
\draw(1,0)node[below]{1}--(1,1.5)node[above]{$x=1$};
\draw[dashed](0,1)node[left]{1}--(1,1);
\draw[domain=0:.5, samples=100, pattern=north east lines]plot(\x, {\x/(1-\x)})--(1,1)--(1,0)--(0,0);

\draw[<->|](1.1,0)--node[fill=white]{$y$}(1.1,.4);
\draw[domain=0.286:.333, samples=10, pattern=north west lines]plot(\x, {\x/(1-\x)})--(1,.5)--(1,0.4)--(0.286,0.4);
\draw[|<->|](0.286,0.25)--node[fill=white]{$1-x$}(1,.25);

\end{tikzpicture}
    \caption{}
\end{figure}

\begin{solution}
$y$轴垂直于旋转体的横断面，设原点到横断面的距离为$y$, 则该横断面的面积
\[A (y) =\pi  (1-x)^2\]
由于函数$y=\frac{x}{1-x}$是单射的，故反函数存在，由它解出
\[x=\frac{y}{1+y},\qquad 1-x=\frac{1}{1+y}\]
于是
\[A(y)=\frac{\pi}{(1+y)^2}\]
\[\begin{split}
    V=\int^1_0 A(y)\dd y&=\pi\int^1_0 \frac{\dd y}{(1+y)^2}\\
&=-\pi\left(\frac{1}{1+y}\right)\Bigg|^1_0=-\pi\left(\frac{1}{2}-1\right)=\frac{\pi}{2}
\end{split}\]
\end{solution}

\begin{figure}[htp]
    \centering
    \begin{tikzpicture}[>=latex, scale=1.5]
\draw[dashed](0,0) ellipse [x radius=1.5, y radius=.5];
\draw[dashed](0,3) ellipse [x radius=1.5, y radius=.5];
\draw[dashed](-1.5,0)--(-1.5,3);
\draw[dashed](1.5,0)--(1.5,3);
\draw(.5,-.45)node[below right]{$B$}--(-.5,.45)node[above left]{$A$};
\draw[thick] (-.5,.45) .. controls (1.8,2.8) and (2.1,1.6) .. (.5,-.45);
\draw(0,0)--(1.4,1.7);
\draw[dashed](0,0)--(1.4,.2)--(1.4,1.7);
\tkzDefPoints{0/0/O, 1.4/.2/B1, 1.4/1.7/C}
\tkzMarkAngle[mark=none, size=.4, ->](B1,O,C)
\tkzLabelAngle[pos=.6](B1,O,C){$60^{\circ}$}
\tkzLabelPoints[left](O)
\draw(1.5,0) arc [start angle = 0, end angle =-70, x radius=1.5, y radius=.5];
\draw(1.5,0)--(1.5,1.5);
    \end{tikzpicture}    
    \caption{}
\end{figure}

\begin{ex}
\begin{enumerate}
    \item 底面半径为$a$的直圆柱，用通过底面一条直径并与底面成$60^{\circ}$角的平面去
    截，得到图4.15的立体，求其体积．
    \item 由抛物线$y^2=4ax$和直线$x=h$所围成
    的区域绕$x$轴旋转一周，求所生成旋转体的体积．
    \item 求底面半径为$r$, 高为$h$的圆锥体的
    体积．
    \item 设由椭圆$\frac{x^2}{a^2}+\frac{y^2}{b^2}=1$所围成的区域绕$x$轴旋转，求此旋转
    椭圆体的体积．
    \item 由曲线$y=\frac{\ln x}{x}$, $x$轴和通过曲线的最高点的纵坐标所围
    成的区域绕$x$轴旋转，求该旋转体的体积．
    \item 两条抛物线有公共的顶点和公共的对称轴，但位于两个
    互相垂直的平面内，一个运动着的椭圆，其中心在两抛物线的公共轴上，它的平面与公共轴垂直，它的顶点位于两条抛物线上，如果椭圆从公共顶点开始，移动一段距离$h$, 求由此椭圆所生成的体积．
    \item 求$y=4-x^2$与$x$轴之间的区域绕直线$y=-1$一周所生成的体积．
    \item 求圆$x^2+(y-b)^2=a^2\; (0<a<b)$绕$x$轴旋转所生成的旋转体的体积．
\end{enumerate}
\end{ex}

\section{简单初值问题{}---{}不定积分的简单应用}

在本章第2节，我们曾说过，一个函数的原函数有无穷多个，它们彼此之间相差一个常数．但是，在许多问题中，我们感兴趣的往往并不是这无穷多个原函数的全体（不定积分），而是某一个特定的原函数，在几何上，就是一条特定的积分曲线．为了求得这条积分曲线，必须知道它所经过的一个点，也就是说，要知道初始条件，在实际问题中，求满足一定初始条件的原函数问题是很多的，下面举一些简单的实例．

\begin{example}
    已知一质点$m$作具有匀加速度$a$的直线运动，设它的初速度为$u$, 求质点在$t$秒后所经过的距离．
\end{example}

\begin{solution}
由于$\frac{\dd v}{\dd t}=a$，所以
\begin{equation}
    v=at+C
\end{equation}
因为，当$t=0$时，$v=u$，代入(4.24)得：$C=u$，所以
\begin{equation}
    v=at+u
\end{equation}
由于$\frac{\dd s}{\dd t}=u+at$，分离变量，把它写成
\[\dd s=(u+at)\dd t\]
两边求不定积分，得
\[\int \dd s=\int (u+at)\dd t\]
\begin{equation}
    s=ut+\frac{1}{2}at^2+C'
\end{equation}
由初始条件：当$t=0$时，$s=0$，代入(4.26)，得：$C'=0$，所以
\begin{equation}
    s=ut+\frac{1}{2}at^2
\end{equation}
为所求．(4.25)和(4.27)就是我们常见的运动学公式.
\end{solution}


\begin{example}
假设在原点处斜抛一质点$m$, 它的初速度$v_0$与水平$x$轴成$\alpha$角，又作用在质点$m$上的力只有向下的重力，求质点$m$的轨迹方程．
\end{example}

\begin{solution}
    设质点$m$的坐标是$(x,y)$，它的质量为$M$克，于是沿着$x$轴和$y$轴的速度分量分别是
    \[v_x=\frac{\dd x}{\dd t},\qquad v_y=\frac{\dd y}{\dd t}.\]
    而沿着$x$轴和$y$轴的加速度分量分别是
    \[\frac{\dd v_x}{\dd t}=\frac{\dd^2 x}{\dd t^2},\qquad \frac{\dd v_y}{\dd t}=\frac{\dd^2 y}{\dd t^2}.\]
    根据牛顿第二定律，得到
\[
\begin{cases}
    M\frac{\dd v_x}{\dd t}=0\\
    M\frac{\dd v_y}{\dd t}=-Mg
\end{cases}
\]
即：
\[\begin{cases}
    \dd v_x=0\\ 
    \dd v_y=-g\dd t
\end{cases}\]
所以
%\begin{align}
\[
\begin{cases}
    v_x=C_1\\ 
    v_y=-gt+C_2
\end{cases}
\] %\end{align}
常数$C_1$和$C_2$由初始条件决定，所以
\[ %\begin{align}
\begin{cases}
v_x=\frac{\dd x}{\dd t}=v_0\cos\alpha\\
v_y=\frac{\dd y}{\dd t}=-gt+v_0\sin\alpha
\end{cases}
\] 
即：
\[\begin{cases}
    \dd x=v_0\cos\alpha \dd t\\ \dd y=(-g t+v_0\sin\alpha)\dd t
\end{cases}\]
再求不定积分得
\[\begin{cases}
    x=(v_0\cos\alpha)t+C_3\\
    y=-\frac{gt^2}{2}+(v_0\sin\alpha)t+C_4
\end{cases}\]
由于，当$t=0$时，$x=y=0$，因此：
\[C_3=C_4=0\]
于是得到质点$m$以时间$t$为参变数的轨迹方程
\begin{align}
\begin{cases}
    x=(v_0\cos\alpha)t\\
    y=-\frac{gt^2}{2}+(v_0\sin\alpha)t
\end{cases}
\end{align}
消去$t$，得：
\[y=\frac{-gx^2}{2v^2_0 \cos^2\alpha}+x\tan\alpha\]
化为标准式，得
\[\left(x-\frac{v^2_0\sin2\alpha}{2g}\right)^2=-\frac{2v^2_0\cos^2\alpha}{g}\left(y-\frac{v^2_0\sin^2\alpha}{2g}\right)\]
这是顶点在$\left(\frac{v^2_0\sin2\alpha}{2g},\frac{v^2_0\sin^2\alpha}{2g}\right)$处的抛物线方程（如图4.16），且从顶点到焦点的方向是$y$轴的负方向，又焦点到准线的距离是$\frac{v^2_0 \cos^2\alpha}{g}$，因此，从顶点到它的上方的准线距离等于$\frac{v^2_0 \cos^2\alpha}{2g}$，而准线到$x$轴的距离就等于
\[\frac{v^2_0 \cos^2\alpha}{2g}+\frac{v^2_0 \sin^2\alpha}{2g}=\frac{v^2_0}{2g}\]
故准线的位置仅依赖于初速度的大小，而与抛射角的大小无关．

\begin{figure}[htp]
    \centering
\begin{tikzpicture}[>=latex]
\draw[->](-.5,0)--(5,0)node[right]{$x$};
\draw[->](0,-.5)--(0,4)node[right]{$y$};
\draw[domain=0:4.2, samples=100, very thick]plot(\x, {0.5*(-\x*\x+4*\x)});
\draw(-.5,2.75)--(4.5,2.75)node[right]{$y=\frac{v^2_0}{2g}$};
\node at (3.5,1)[right]{$y=\frac{-gx^2}{2v^2_0 \cos^2\alpha}+x\tan\alpha$};
\node at (0,0)[below left]{$O$};
\draw(2,2.75)--(2,0);
\tkzDefPoints{2/2/A, 2/1.25/S}
\tkzDrawPoints(A,S)
\tkzLabelPoints[above right](A,S)
\draw[->, thick](0,0)--(1,2)node[right]{$v_0$};


\end{tikzpicture}
    \caption{}
\end{figure}
\end{solution}

下面我们来讨论人口的增长与放射性元素蜕变的问题．

在很多单纯的自然现象中，一个随着时间而变化的量$y(t)$的变化率常常和$y(t)$当时的值成正比例，用数学式表达，即存在一个随该现象而定的常数$k$（与$t$无关），使得
\[y' (t) =ky (t) \]
或
\begin{equation}
    \frac{\dd }{\dd t}y(t)=ky(t)
\end{equation}
人口的增长与放射性元素的蜕变就是这样的两个例子．

假如人口按某一固定的“增长率”$k$增长，譬如时间单位是“年”，增长率2\%, 则$k=2\%$. 设人口随时间变化的函数是$y(t)$, 本来$y(t)$是整数，但由于数量很大，我们可以把它看作一个平滑函数，$y(t)$满足下面的条件：
\[y(t+\Delta t)-y(t)=ky(t)\cdot \Delta t\]
换言之，人口增加数量与这段时间开始的人口数，以及这段时间间隔成比例，也即
\[\frac{y(t+\Delta t)-y(t)}{\Delta t}=ky(t)\]
令$\Delta t\to 0$, 由上式得到
\[\frac{\dd }{\dd t}y(t)=ky(t)\]
也就是说，人口在时刻$t$的瞬时变率与当时的值$y(x)$成正比例．为求出函数$y(t)$, 我们来解具有初始条件的微分方程
\begin{numcases}{}
    \frac{\dd }{\dd t}y(t)=ky(t)\\
    y(0)=A
\end{numcases}
将(4.33)改写成
\[\frac{\dd y(t)}{y(t)}=k\dd t\]
即：$\ln[y(t)]=kt+C$，由初始条件确定常数$C$的值，得：
\[\ln A=C\]
所以：$\ln[y(t)]=kt+\ln A$，
移项，得：
\[\ln\frac{y(t)}{A}=kt\]
由此得：
\[\frac{y(t)}{A}=e^{kt}\]
所以：$y(t)=Ae^{kt}$为所求．



总结上面的讨论，得到以下命题：

\begin{blk}{命题}
    设$y'(t)=ky(t)$，$y(0)=A$，则：
    \[y(t)=Ae^{kt}\]
\end{blk}




\begin{example}
    若某城市现有人口10万，按照每年3\%的比率增长，问几年后人口是现有人口的2倍？
\end{example}

\begin{solution}
    设$t_1$年后人口加倍，于是依题意有
\[10e^{3\%t_1}=2\x10\]
即：$e^{0.03t_1}=2$, 两边取常用对数，得
\[\begin{split}
    0.03t_1\lg e&=\lg2\\
t_1&= \frac{100\lg2}{3\lg e}\approx \frac{100}{3}\cdot\frac{0.3010}{0. 4343}\\
&\approx \frac{100}{3}\cdot 0.6931\approx 23.
\end{split} \]
答：约23年后，人口会加倍．
\end{solution}

从问题的解法看出，人口的加倍期仅与增长率有关，即$t_1=\frac{\ln2}{k}$, 下面是二者之间的对照表．

\begin{center}
\begin{tabular}{c|ccccc}
\hline
    增长率$k$ & 0.05&0.04&0.03&0.02&0.01\\
\hline
加倍期$t_1$ & 约14年& 约17年& 约23年& 约34.5年& 约39年\\
\hline
        \end{tabular}
\end{center}

同样地可把上述命题应用于放射性元素蜕变问题．

\begin{example}
    放射性元素的数量随时间$t$变化的函数是$N(t)$, $A=N(0)$是起算时的数量，$k$是蜕变率，且知放射性元素经过$\Delta t$这段时间的蜕变，它减少的数量$\Delta N(t)=N(t+\Delta t)-N(t)$与在时刻$t$的数量$N(t)$及这段蜕变时间$\Delta t$成比
例，即
\[\Delta N (t) =-kN (t) \cdot \Delta t\]
因为放射性元素不断减少，所以$\Delta N$是负数，问经过多少年放射性元素剩下了一半？
\end{example}

\begin{solution}
    依题意，有
\[\frac{\dd N}{\dd t}=-kN (t) \]
根据本节命题，得到
\[N (t) =Ae^{-kt}\]

设经过$t_1$年，放射性元素减少了一半，通常称$t_1$为放射性元素的半衰期，于是
\[t_1=\frac{\ln 2}{k}\approx \frac{0. 6931}{k}\]
\end{solution}

下面是一些放射性元素的半衰期的简表．

\begin{center}
    \begin{tabular}{c|ccc}
    \hline
        元素 & 氩Ar$^{39}$ &碳C$^{14}$&氢H$^3$  \\
    半衰期 & 269年& 5730年& 12.3年 \\
    \hline
    元素 & 碘I$^{129}$ &铀U$^{238}$& 铀U$^{235}$ \\
    半衰期 &  $16\x10^6$年& $4.5\x 10^9$年& $0.71\x 10^9$年\\
    \hline
            \end{tabular}
    \end{center}

我们可以把岩石或化石中所含的各种微量的放射性元素，当作天然的计时仪表，因为各种放射性元素的微量都可以测量十分准确，所以我们可以通过对残留微量的测定，
用半衰期去推算岩石或化石的“年龄”，甚至于我们所居住的地球的“年龄”．

\begin{ex}
\begin{enumerate}
    \item 求下列积分曲线：
\begin{multicols}{2}
\begin{enumerate}
    \item $\begin{cases}
        y'=2x\\
        y|_{x=x_0}=y_0
    \end{cases}$
    \item $\begin{cases}
        y'=2\sqrt{y}\\
        y|_{x=0}=0
    \end{cases}$
    \item $\begin{cases}
        y'=\frac{2}{3\sqrt[3]{x}}\\
        y|_{x=0}=4
    \end{cases}$
\end{enumerate}
\end{multicols}
\item 在$20.2^{\circ}{\rm C}$, 细菌受到5\%的消毒溶液消毒，每小时细菌死
    亡的比率为11\%, 假如开始时有1000个细菌，那么在24小时后还剩下多少？
    \item 化学反应速度$v$和温度$t$的增加是按照有机体生长的规律
    进行的，当$t=0$时，$v=30$; 当$t=9.90$时，$v=60$, 求反应率以及当$t=20$时，$v$的值．
\item 若函数$y=f(x)$的变率$\frac{\dd y}{\dd x}$与当时的函数值的平方成
    正比例，且当$x=0$时，$y=f(0)=2$; 当$x=1$时，$y=f(1)=4$, 求函数$y=f(x)$.
\end{enumerate} 
\end{ex}

\subsection*{习题4.5}
\begin{enumerate}
    \item 计算下列不定积分：
\begin{multicols}{2}
\begin{enumerate}
    \item $\displaystyle\int \frac{\dd x}{1-3x^2}$
    \item $\displaystyle\int \frac{1-\sin^{-3}x}{1-\cos^2 x}\dd x$
    \item $\displaystyle\int \frac{\dd x}{\sin x\cos x}$
    \item $\displaystyle\int x\sec^2 x\dd x$
    \item $\displaystyle\int xe^{2x+1}\dd x$
    \item $\displaystyle\int \frac{x^2+2}{x(x-1)}\dd x$
\end{enumerate}
\end{multicols}

\item 求下列定积分：
\begin{multicols}{2}
    \begin{enumerate}
        \item  $\displaystyle\int_{-1}^{3} \frac{\dd x}{x^{2}+7 x+10} $
        \item  $\displaystyle\int_{\sqrt{3}}^{0} \frac{\dd x}{2+x^{2}}$
        \item  $\displaystyle\int_{0}^{1} x \sqrt{\left(2-x^{2}\right)^{3}} \dd x$
        \item  $\displaystyle\int_{0}^{\ln 3} \frac{\dd x}{e^{x}+5+6 e^{-x}}$
        \item  $\displaystyle\int_{1}^{\tfrac{7}{5}} \frac{x^{\tfrac{3}{4}}+x^{\tfrac{1}{4}}}{x^{\tfrac{9}{4}}-x^{\tfrac{5}{4}}+x^{\tfrac{1}{4}}} \dd x$
        \item  $\displaystyle\int_{1}^{2} \frac{2-x+x^{2}}{(1+x)(1-x)^{2}} \dd x$
        \item  $\displaystyle\int_{1}^{2} x \ln x \dd x$
        \item  $\displaystyle\int_{-\pi}^{\pi} \sin m x \sin n x \dd x$
        (此处 $m, n$ 是整数且 $m \pm n\ne 0$)
\item  $\displaystyle\int_{\tfrac{a}{\sqrt{3}}}^{a} \frac{\dd x}{\left(a^{2}+x^{2}\right)^{3 / 2}}$
\end{enumerate}
\end{multicols}
\item  求下列极限：
    \begin{enumerate}
\item  $\displaystyle\lim _{n \to  \infty} \frac{1}{n}\left\{\left(1+
\frac{1}{n}\right)^2+\left(1+\frac{2}{n}\right)^2+\cdots+
\left(1+\frac{n}{n}\right)^{2}\right\}$
\item $\displaystyle\lim _{n \to  \infty} \frac{1^{4}+2^{4}+\cdots+n^{4}}{n^{5}}$
\item $\Lim_{n\to\infty}\frac{1}{n}\left(\sin\frac{\pi}{n}+\sin\frac{2\pi}{n}+\cdots+\sin\frac{n\pi}{n}\right)$
\item $\Lim_{n\to\infty}\left(\frac{1}{\sqrt{n+1}}+\frac{1}{\sqrt{n+2}}+\cdots+\frac{1}{\sqrt{n+n}}\right)\cdot \frac{1}{\sqrt{n}}$
\end{enumerate}

\item 利用求$y=\frac{1}{x}$的定积分，证明
\[\ln(n+1)<1+\frac{1}{2}+\frac{1}{3}+\cdots+\frac{1}{n}<1+\ln n.\]
\item 设$0<a\le 1$，求证：
\[\int^a_0\frac{\dd x}{x^3+1}<\int^a_0\frac{\dd x}{x^4+1}. \]
\item 设$0\le x\le \frac{\pi}{4}$，利用不等式$0\le \sin x\le x$证明：
\[\frac{\pi}{4}<\int^{\pi/4}_0\frac{\dd x}{\sqrt{1-\sin x}}<2-\sqrt{4-\pi}.\]

\item 若$m,n$是正整数，求证：
\[\int^1_0 x^m (1-x)^n\dd x=\frac{n}{m+1}\int^1_0x^{m+1}(1-x)^{n-1}\dd x. \]
\item 求证：$\displaystyle\int^\pi_0 xf(\sin x)\dd x=\frac{\pi}{2}\int^\pi_0f(\sin x)\dd x$.
\item 若$I_n=\displaystyle\int^{\pi/4}_0 \tan^n \theta \dd \theta\quad (n\ge 2)$. 试求$I_n$和$I_{n-2}$的关系式．
\item 若$I_n=\displaystyle\int \frac{\dd x}{(x^2+a^2)^n}$，求证：
\[2na^2I_{n+1}=\frac{x}{(x^2+a^2)^n}+(2n-1)I_n\]
\item 设$S_n=\frac{1}{\sqrt{1}}+\frac{1}{\sqrt{2}}+\cdots+\frac{1}{\sqrt{n}}$，$n$为正整数．
\[T_n=\frac{1}{\sqrt{1+\frac{1}{2}}}+\frac{1}{\sqrt{2+\frac{1}{2}}}+\cdots+\frac{1}{\sqrt{n+\frac{1}{2}}}.\]
求证：
\begin{multicols}{2}
\begin{enumerate}
    \item $S_n\ge \displaystyle\int^{n+1}_1 \frac{\dd x}{\sqrt{x}}$
    \item $\Lim_{n\to\infty}\frac{1}{S_n}=0$
    \item $S_{n+1}-1\le T_n\le S_n$
    \item $\Lim_{n\to\infty}\frac{T_n}{S_n}=1$
\end{enumerate}
\end{multicols}

\item 设$f(x)=(ax+b)e^{x/2}$, 求$a,b$的值，使得等式
\[f (x) =e^{x/2}-1+\frac{1}{2}\int^x_0 f (t) \dd t\]
对于任何$x$的值都成立．
\item 求由曲线$y=\frac{4}{x^2}$, $y=x-1$, $x=1$所围区域的面积．
\item 求在曲线$x^2+y^2=9$和$y=\frac{1}{3}(-x^2+1)$之间位于
第二象限那部分的面积．
\item 求积分$I(a)=\displaystyle\int^1_{-1}|x-a|e^x\dd x$在$|a|\le 1$时的最大值．
\item 设$0<t<\pi$，考察$x$的函数$f(x)=2\sin x\sin(t-x)$
\begin{enumerate}
    \item 画出$f(x)$在区间$0\le x\le \pi$上的图象，再设
    \[F (t) =\int^\pi_0 |f (x)|\dd x.\]
根据所作的图象说明$F(t)$所表示的是什么？
\item 用关于$t$的式子写出$F(t)$, 并求出使$F(t)$取最小值时$t$的值．
\end{enumerate}
\item  设$T_1$是由抛物线$y=4x^2$和直线$x=a$, $x=1$, $y=0$所
围成的区域，$T_2$是抛物线$y=4x^2$和二直线$x=a$, $y=0$围成的区域，但是$0<a<1$，试求：
\begin{enumerate}
\item $T_1$绕$x$轴旋转而成的旋转体体积$V_1$;    \item $T_2$绕$y$轴旋转而成的旋转体体积$V_2$;    \item 使$V_1+V_2$为最大时$a$的值．
\end{enumerate}
\item 设$0\le t\le \frac{\pi}{2}$，曲线$y=\sin x$与三条直线$x=t$, $x=2t$, 
$y=0$围成部分绕$x$轴旋转所成的旋转体的体积为$V(t)$．当$t=\alpha$时$V(t)$最大，试求$\cos\alpha$的值．
\item 设曲线$y=f(x)$通过点$P(3, 4)$, 并满足微分方程
$\frac{\dd y}{\dd x}=\frac{y}{x}$，
求此积分曲线．

\item 若$\frac{\dd y}{\dd x}=\cos(x+y)$，求不定积分．（提示：设$z=x+y$）
\end{enumerate}

\fi

\end{document}
；
% 直接编译 main.tex 时 \BookPart 未定义，得到合订本。main.tex 中用：
%   \ifPartA ... \fi      合订本和上册包含的内容
%   \ifPartB ... \fi      合订本和下册包含的内容
%   \ifFullBook ... \fi   只有合订本包含的内容
%   \BookPartSuffix       书名后缀：合订本为空，上、下册为“（上）”“（下）”
%   \PartChoice{合订本}{上册}{下册}   按版本取不同的文字（如出版日期）
\newif\ifFullBook
\newif\ifPartA
\newif\ifPartB
\ifdefined\BookPart\else\def\BookPart{}\fi
\def\BookPartNameA{A}
\def\BookPartNameB{B}
\ifx\BookPart\BookPartNameA
  \PartAtrue \def\BookPartSuffix{（上）}
\else\ifx\BookPart\BookPartNameB
  \PartBtrue \def\BookPartSuffix{（下）}
\else
  \FullBooktrue \PartAtrue \PartBtrue \def\BookPartSuffix{}
\fi\fi
\newcommand{\PartChoice}[3]{\ifFullBook#1\else\ifPartA#2\else#3\fi\fi}


% ---- 本册专用 ----
% 注意：不要加载 physics 宏包。它会把 \div 重定义为散度（上册用 \div 表示除号），
% 下册用到的 \dd、\arccot、\arcsec 已在下方单独定义。
\usepackage{pstricks}
\usepackage{bookmark}
\usepackage{circuitikz}
\usepackage{longtable}
\usepackage{tkz-base, tkz-euclide}

\newcommand{\dd}{\,{\rm d}}
\DeclareMathOperator{\arccot}{arccot}\DeclareMathOperator{\arcsec}{arcsec}
\newcommand{\Lim}{\displaystyle\lim}
\newcommand{\Int}{\displaystyle\int}

\newcommand{\NC}{\text{N}/\text{C}}
\newcommand{\ms}{\text{m}/\text{s}}
\newcommand{\cms}{\text{cm}/\text{s}}
\newcommand{\msq}{\text{m}/\text{s}^2}
\newcommand{\cmsq}{\text{cm}/\text{s}^2}
\newcommand{\kmh}{\text{km}/\text{h}}

\newcommand{\R}{\mathbb{R}}
\newcommand{\N}{\mathbb{N}}
\newcommand{\Z}{\mathbb{Z}}

\newcommand{\VEC}{\overrightarrow}

\newcommand{\comb}[2]{{\rm C}^{#1}_{#2}}
\newcommand{\perm}[2]{{\rm P}^{#1}_{#2}}
\newcommand{\rep}[2]{{\rm H}^{#1}_{#2}}
\newcommand{\pc}[1]{\cos{#1}+i\sin{#1}}
\newcommand{\pcx}[1]{\cos\left(#1\right)+i\sin\left(#1\right)}

% 第六册上、下两本的章号各自从“第一章”编起（正文中的“图2.33”“第一章”等引用均按此编号）。
% 合订本用 \volume 分隔两本；上、下册单独编译时不需要。
% \volume{标题}：开始新的一本，章号归零；同时让 hyperref 锚点按“本-章”区分，避免重复。
% \appendixchapters：之后的 \chapter 不编号（用于上册末尾的附录），且不重置页码。
\ctexset{part = {numbering = false}}
\newcounter{volume}
\renewcommand{\theHchapter}{\arabic{volume}.\arabic{chapter}}
\makeatletter
\newcommand{\volume}[1]{\@mainmattertrue\stepcounter{volume}\setcounter{chapter}{0}\part{#1}}
\newcommand{\appendixchapters}{\@mainmatterfalse}
\makeatother

\begin{document}

\title{\Huge\bfseries 中学数学实验教材\\第六册\BookPartSuffix}
\author{\Large 中学数学实验教材编写组编}
\date{\Large 1986年5月}
\maketitle

\frontmatter
\chapter{前~~言}

这一套中学数学实验教材，内容的选取原则是精简实
用，教材的处理力求深入浅出，顺理成章，尽量作到使人人
能懂，到处有用．

    本教材适用于重点中学，侧重在满足学生将来从事理工
方面学习和工作的需要．

    本教材的教学目的是：使学生切实学好从事现代生产、
特别是学习现代科学技术所必需的数学基础知识；通过对数
学理论、应用、思想和方法的学习，培养学生运算能力，思
维能力，空间想象力，从而逐步培养运用数学的思想和方法
去分析和解决实际问题的能力；通过数学的教学和学习，培
养学生良好的学习习惯，严谨的治学态度和科学的思想方
法，逐步形成辩证唯物主义世界观．

   根据上述教学目的，本教材精选了传统数学那些普遍实
用的最基础的部分，这就是在理论上、应用上和思想方法上
都是基本的、长远起作用的通性、通法．比如，代数中的数
系运算律，式的运算，解代数方程，待定系数法；几何中的
图形的基本概念和主要性质，向量，解析几何；分析中的函
数，极限，连续，微分，积分；概率统计以及逻辑、推理论
证等知识．对于那些理论和应用上虽有一定作用，但发展余
地不大，或没有普遍意义和实用价值，或不必要的重复和过
于繁琐的内容，如立体几何中的空间作图，几何体的体积、
表面积计算，几何难题，因式分解，对数计算等作了较大的
精简或删减．

    全套教材共分六册．第一册是代数．在总结小学所学自
然数、小数、分数基础上，明确提出运算律，把数扩充到有
理数和实数系．灵活运用运算律解一元一次、二次方程，二
元、三元一次方程组，然后进一步系统化，引进多项式运
算，综合除法，辗转相除，余式定理及其推论，学到根式、
分式、部分分式．第二册是几何．由直观几何形象分析归纳
出几何基本概念和基本性质，通过集合术语、简易逻辑转入
欧氏推理几何，处理直线形，圆、基本轨迹与作图，三角比
与解三角形等基本内容．第三册是函数．数形结合引入坐
标，研究多项式函数，指数、对数、三角函数，不等式等．
第四册是代数．把数扩充到复数系，进一步加强多项式论，
方程式论，讲线性方程组理论，概率（离散的）统计的初步
知识．第五册是几何．引进向量，用向量和初等几何方法综
合处理几何问题，坐标化处理直线、圆、锥线，坐标变换与
二次曲线讨论，然后讲立体几何，并引进空间向量研究空间
解析几何初步知识．第六册是微积分初步．突出逼近法，讲
实数完备性，函数，极限，连续，变率与微分，求和与积分．

本教材基本上采取代数、几何、分析分科，初中、高中
循环排列的安排体系．教学可按初一、初二代数、几何双科
并进，初三学分析，高一、高二代数（包括概率统计）、几
何双科并进，高三学微积分的程序来安排．

    本教材的处理力求符合历史发展和认识发展的规律，深
入浅出，顺理成章．突出由算术到代数，由实验几何到论证
几何，由综合几何到解析几何，由常量数学到变量数学等四
个重大转折，着力采取措施引导学生合乎规律地实现这些转
折，为此，强调数系运算律，集合逻辑，向量和逼近法分别
在实现这四个转折中的作用．这样既遵循历史发展的规律，
又突出了几个转折关头，缩短了认识过程，有利于学生掌握
数学思想发展的脉络，提高数学教学的思想性．

这一套中学数学实验教材是教育部委托北京师范大学、
中国科学院数学研究所、人民教育出版社、北京师范学院、
北京景山学校等单位组成的领导小组组织“中学数学实验教
材编写组”，根据美国加州大学伯克利分校数学系项武义教
授的《关于中学实验数学教材的设想》编写的．第一版印出
后，由教育部实验研究组和有关省市实验研究组指导在北
京景山学校、北京师院附中、上海大同中学、天津南开中
学、天津十六中学、广东省实验中学、华南师院附中、长春
市实验中学等校试教过两遍，在这个基础上编写组吸收了实
验学校老师们的经验和意见，修改成这一版《中学数学实验
教材》，正式出版，内部发行，供中学选作实验教材，教师
参考书或学生课外读物．在编写和修订过程中，项武义教授
曾数次详细修改过原稿，提出过许多宝贵意见．

    本教材虽然试用过两遍，但是实验基础仍然很不够，这
次修改出版，目的是通过更大范围的实验研究，逐步形成另
一套现代化而又适合我国国情的中学数学教科书．在实验过
程中，我们热忱希望大家多提意见，以便进一步把它修改好．

\begin{flushright}
    中学数学实验教材编写组\\
    一九八五年三月
\end{flushright}








\tableofcontents

\mainmatter

\ifPartA
\ifFullBook\volume{第六册（上）}\fi
\chapter{复数系}

在数学中，数是最基础、应用最广泛的工具之一，我们已经在自然数的基础上，逐步学习了整数、有理数、实数等概念和它们的运算及性质.这一章，我们将进一步学习复数的概念及其运算、性质.

\section{复数的概念}
\subsection{数的概念的发展}

数的概念是中学数学的主要内容之一，是从实践中逐步形成和发展起来的.早在有文字记载的历史以前，由于计数和排序的需要，人类就建立起自然数（正整数）的概念，处理着各种常见的、有天然单元的可数量. 自然数的全体构成了自然数集$\mathbb{N}$.

随着社会生产和科学技术的发展，数的概念相应地得到了不断发展.

为了处理实践中遇到的测量、分割、分配等可以细分的量，人们引进了正分数；为了解决生产实践中的各种具有相反意义的量的表示和运算问题，人们又引进了零及负数，并把自然数看作正整数. 这样一来，把正整数、零、负整数合并在一起，称为整数集$\mathbb{Z}$；并把整数与分数（正分数、负分数）合称为有理数集$\mathbb{Q}$.

任一有理数，可以表示成分数形式$\frac{m}{n}$，其中$m,n\in\mathbb{N}$; 还可以表示成循环小数（包括循环节为0的小数）.

有理数系的形成，使得许多人曾认为这就足以解决可度量的量，特别是在古希腊，他们曾认为：任意两条线段都是可公度的（即：任意两条线段长度的比都是有理数）. 但实际上这种认识是不正确的. 例如，正方形的边长和对角线长，正五边形的边长和对角线长都是不可公度的. 为了解决这些量与量之比值不能用有理数表示的矛盾，于是又引进了无理数的概念，无理数就是无限不循环小数. 无理数与有理数统称为实数，构成了实数集$\mathbb{R}$.

实数集$\mathbb{R}$对加、减、乘、除（零不作除数）四则运算是封闭的，且满足运算通性，即满足加、乘法的交换律、结合律和分配律，数零与1保持在运算中的特性．

实数集$\mathbb{R}$与数轴上的点集之间，可以建立起一一对应的关系. 因此，从度量的角度来说，实数系是一个完整无缺的数系，它足以满足度量的实际需要.

但是，从代数运算和解方程的角度来看，实数集并不完整无缺，因为，在实数范围内，开方运算并不封闭，实系数方程（如$ax^2+bx+c=0$）在实数系内未必有解.

事实上，数的概念的逐步扩充，与代数运算和解方程的需要也是相适应的，例如：

由于自然数集$\mathbb{N}$对于减法、除法运算不封闭，因而方程$x+7=2$, $7x=5$等在自然数集中无解，但在有理数集$\mathbb{Q}$中就有解$x=-5$, $x=\frac{5}{7}$. 

由于有理数集$\mathbb{Q}$对于开方运算不封闭，因而方程$x^2-3
=0$在$\mathbb{Q}$中无解，但在实数集$\mathbb{R}$中就有两个解$x=\pm\sqrt{3}$.

同样，实数集$\mathbb{R}$对开方运算仍不封闭，因而像方程$x^2+1=0$在$\mathbb{R}$中还是无解，因为，没有一个实数的平方等于$-1$．要解决这一类矛盾，数的概念需要进一步发展，这就需要引进新数，并讨论它的运算.

\begin{ex}
\begin{enumerate}
    \item 试叙述实数集$\mathbb{R}$中有哪些数系运算通性？
\item 试分别在自然数集$\mathbb{N}$、整数集$\mathbb{Z}$、有理数集$\mathbb{Q}$和实数集$\mathbb{R}$中，解方程$(x-2)(3x+1)(x+7)(x^2+5)=0$.
\end{enumerate}
\end{ex}

\subsection{复数及其几何表示}
如上所述，由于解方程的需要，早在十六世纪，人们就引进了一个新数$i$，使得方程$x^2+1=0$有解，并把它叫做\textbf{虚数单位}，并规定：
\begin{enumerate}[(1)]
\item $i^2=-1$;
\item 它与实数进行四则运算时，运算通性及数0与1的运算特性仍然保持. 即，满足加、乘法的交换律、结合律、分配律；满足
\[\begin{split}
    0+i&=i+0=i,\\
0\cdot i&=i\cdot 0=0,\\
1\cdot i&=i\cdot 1=i.
\end{split}\]
\end{enumerate}

在这样的规定下，$i$可以与实数$b$相乘，再与实数$a$相加，从而得到$a+bi$.

这样，数的概念扩大了，我们就把形如$a+bi\; (a,b\in\mathbb{R})$的数，叫做复数.全体复数构成的集合，称为\textbf{复数集}，
一般记作$\mathbb{C}$.

对于复数$a+bi\; (a,b\in\mathbb{R})$，当$b=0$时，就是实数$a$；当$b\ne 0$时，这一复数叫做虚数，特别地，当$a=0$且$b\ne 0$时，叫做纯虚数.

$a$，$b$分别叫做复数$z=a+bi$的\textbf{实部}与\textbf{虚部}，并且记作：${\rm Re}(z)=a$, ${\rm Im}(z)=b$.

显然，实数集$\mathbb{R}$是复数集$\mathbb{C}$的真子集，即$\mathbb{R}\subset \mathbb{C}$.

例如，$3+2i$, $-\frac{1}{2}-\sqrt{3}i$, $-0.25i$都是虚数. 其中$-0.25i$是纯虚数．它们的实部和虚部分别为：
\[\begin{split}
&{\rm Re}(3+2i)=3,\qquad {\rm Im}(3+2i)=2;\\
&{\rm Re}\left(-\frac{1}{2}-\sqrt{3}i\right)=-\frac{1}{2},\qquad {\rm Im}\left(-\frac{1}{2}-\sqrt{3}i\right)=-\sqrt{3};\\
&{\rm Re}(-0.25i)=0,\qquad {\rm Im}(-0.25i)=-0.25.\\
\end{split}\]

我们还规定，当且仅当两个复数的实部和虚部分别相等时，这两个复数相等，即
\[a+bi=c+di\Longleftrightarrow a=c\text{且}b=d\]
或
\[z_1=z_2\Longleftrightarrow {\rm Re}(z_1)={\rm Re}(z_2)\quad \text{且} \quad {\rm Im}(z_1)={\rm Im}(z_2)\]
特别地，我们可以得出
\[a+bi=0\Longleftrightarrow a=b=0\]

\begin{example}
    若$3x+2yi-1=(x-1)i+y$，其中$x,y\in\mathbb{R}$. 求$x$与$y$的值.
\end{example}

\begin{solution}
将已知等式整理得
\[(3x-1)+2yi=y+(x-1)i\]
根据复数相等的定义，可得方程组$\begin{cases}
    3x-1=y\\
    2y=x-1
\end{cases}$，
解得
\[x=\frac{1}{5},\qquad y=-\frac{2}{5}\]
\end{solution}

当两个复数的实部相等、虚部互为相反数时，这两个复数叫做\textbf{互为共轭复数}，即
：复数$a+bi$与$a-bi$叫做互为共轭复数. 一般地，复数$z=a+bi$的共轭复数表示为：
$$\overline{z}=a-bi$$

显然，实数$a$（即虚部为零的复数）的共轭复数仍是它自身$a$.

我们已经知道，实数集与数轴上的点之间可建立一一对应关系. 引进虚数，构成复数集之后，是否有相应的对应关系呢？

首先，虚数单位$i$在实数轴上没有立足之地，它不属于实数集；其次，当$b\ne 0$时，复数$a+bi$也不可能在实数轴上找出对应之点，它仍不属于实数集，因此，我们必须另找途径来建立对应关系.

从复数及其相等的定义中，我们知道，任何一个复数
$z=a+bi$，都由一个有序实数对$(a,b)$唯一确定，而每一个有序实数对，又借助平面的坐标化可以唯一确定平面上一个点，因此，只要在平面上建立起直角坐标系，就可以用平面上的点$Z(a,b)$表示出复数$z=a+bi$，如图1.1所示.

\begin{minipage}{.45\textwidth}
\centering
\begin{tikzpicture}[>=stealth]
\draw[->](-1,0)--(3,0)node[right]{$x$};
\draw[->](0,-1)--(0,3)node[right]{$y$};
\draw[dashed](0,1.5)node[left]{$b$}--(2,1.5)node[above]{$Z(a,b)$}--(2,0)node[below]{$a$};
\node[below right]{$O$};
\draw[fill](2,1.5)circle(1pt);
\end{tikzpicture}
\captionof{figure}{ }
\end{minipage}
\hfill
\begin{minipage}{.45\textwidth}
\centering
    \begin{tikzpicture}[>=stealth]
  \draw[->](-1,0)--(3,0)node[right]{$x$};
\draw[->](0,-2)--(0,2)node[right]{$y$};
\draw[dashed](0,1.25)node[left]{$b$}--(2,1.25)node[above]{$Z(a,b)$}--(2,0)node[below right]{$a$}--(2,-1.25)node[below]{$\overline{Z}(a,-b)$}--(0,-1.25)node[left]{$-b$};
\node[below right]{$O$};  
\draw[fill](2,1.25)circle(1pt);
\draw[fill](2,-1.25)circle(1pt);
    \end{tikzpicture}
\captionof{figure}{ }
\end{minipage}


这个建立了坐标系来表示复数的平面，叫做\textbf{复平面}，其中$X$轴叫做\textbf{实轴}，$Y$轴（不含原点）叫做\textbf{虚轴}.

这样一来，每一个复数，在复平面内都有唯一一点和它对应，反之，复平面内的每一个点，都有唯一的一个复数和它对应，因此，复数集$\mathbb{C}$和复平面上的点集是一一对应的. 特别地，实数集和复平面实轴上的点集一一对应；纯虚数集合和复平面虚轴（不含原点）上的点集一一对应.

显然，互为共轭的两个复数$z=a+bi$与$\bar z=a-bi$在复平面上对应的两点$Z(a,b)$与$\overline{Z}(a,-b)$关于实轴对称（图1.2）．

\begin{ex}
\begin{enumerate}
    \item 将下列给出的各数，分类填入表中：
\[1-\sqrt{2},\quad 0.618,\quad -1+i,\quad \frac{1}{4}i,\quad 0,\quad i,\quad i^2,\quad 3+7\sqrt{2}i\]
\[i-1,\quad (\sqrt{2}+\sqrt{3})i,\quad \sqrt{2}i-\sqrt{2},\quad -\frac{\sqrt{3}}{2}i\]
\begin{tabular}{|c|p{.7\textwidth}|}
\hline
    实数&\\
\hline
虚数&\\\hline
纯虚数&\\    \hline
\end{tabular}

\item 填写下列表格：
\begin{center}
\begin{tabular}{|c|p{.1\textwidth}|p{.1\textwidth}|p{.1\textwidth}|p{.1\textwidth}|p{.1\textwidth}|}
\hline 
复数$z$& ${\rm Re}(z)$& ${\rm Im}(z)$&$\bar z$& ${\rm Re}(\bar z)$& ${\rm Im}(\bar z)$\\
\hline
$-\sqrt{3}+i$ &&&&&\\
$-(1+\sqrt{2})$&&&&&\\
$\frac{i-\sqrt{2}}{4}$&&&&&\\
$0$&&&&&\\
$i$&&&&&\\
$a-bi$&&&&&\\
\hline
\end{tabular}
\end{center}
\item 试求出满足下列等式的实数$x$，$y$：
\begin{enumerate}[(1)]
\item $(3x+2y)+(5x-y)i=17-2i$;
\item $3x+2y-7=(2x+3y-8)i$;
\item $7x+(4y-1)i-1=0$.
\end{enumerate}

\item 试说明对任意复数$z$，都有$\overline{\bar z}=z$.
\item 在复平面上，画出下列复数及它们的共轭复数所对应的点：
\[2-2i;\quad -2+2i;\quad -2i; \quad \sqrt{2};\quad 0.\]
\end{enumerate}
\end{ex}

\subsection{复数的模与幅角}
我们已经知道，复数集与复平面上的点集之间是一一对应的，即每一复数$z=a+bi\; (a,b\in\R)$都对应着复平面
上唯一一点$Z(a,b)$，反之亦然。

如果我们在复平面上连结原点$O$与表示复数$z=a+bi$的点$Z(a,b)$，就得到一个向量（从$O$点指向$Z$点），记作$\VEC{OZ}$，这样就可以在复数同向量（从原点出发的向量）之间
建立起联系. 不难看出，任一复数$z$可唯一确定复平面上一点$Z$，进而唯一确定一个向量$\VEC{OZ}$；反之，复平面上任一向量$\VEC{OZ}$，可唯一确定一点$Z$，进而唯一确定一个复数$z$，因此，复数集$\mathbb{C}$与复平面内所有以原点$O$为起点的向量（位置向量）集合也是一一对应的。

这样一来，我们就可以很方便地把复数$z=a+bi\; (a,b\in\R)$说成点$Z(a,b)$，或说成向量$\VEC{OZ}$. 研究复数的运算、性质就可以和研究向量的运算、性质互通了.

如图1.3所示，我们把向量$\VEC{OZ}$的模（$OZ$的长度）$r$叫做复数$z=a+bi$的模（或绝对值），记作
\[r=|z|\quad \text{或}\quad r=|a+bi|\]
\begin{figure}[htp]
    \centering
\begin{tikzpicture}[>=stealth]
\draw[->](-1,0)--(3,0)node[right]{$x$};
\draw[->](0,-1)--(0,3)node[right]{$y$};
\draw[dashed](0,1.5)node[left]{$b$}--(2,1.5)node[above]{$Z(a,b)$}--(2,0)node[below]{$a$};
\draw[very thick,->](0,0)--node[left]{$r$}(2,1.5);
\draw(.5,0) arc (0:35:.5)node[right]{$\theta$};
\node[below right]{$O$};
\end{tikzpicture}
    \caption{}
\end{figure}

显然
\[r=|z|=|a+bi|=\sqrt{a^2+b^2}\ge 0\]

特别地，当$b=0$时，$z$是一个实数，这时
\[r=|z|=|a|\qquad  \text{（与实数的绝对值概念一致）
}\]

又以实轴$OX$的正半轴为始边，向量$\VEC{OZ}$所在的射线为终边的角$\theta$，叫做复数$z=a+bi$的\textbf{幅角}.

由于任一不等于零的复数，按以上规定可有无限多个幅
角（终边相同的角）的值，因而，我们进一步约定：

满足$0\le\theta<2\pi$的幅角$\theta$的值，叫做幅角的\textbf{主值}，并记作$\arg z$，即$0\le \arg z<2\pi$. 
因此，任一非零复数$z$的幅角$\theta=\arg z+2k\pi\; (k\in\Z)$

这样一来，每一个不等于零的复数，就只有唯一的模与唯一的幅角主值，并可被它的模与幅角主值所唯一确定。

应该指出，由于复数$z=0$对应于复平面上的零向量（一个点），而且零向量的方向是任意的，因此，复数$z=0$的模$|z|=0$，它的幅角却是任意的。

\begin{example}
试比较复数$z_1=1+\sqrt{2}i$，$z_2=-\frac{1}{2}-i$的模的大小.
\end{example}

\begin{solution}
$\because\quad |z_1|=\sqrt{1+2}=\sqrt{3},\quad |z_2|=\sqrt{\frac{1}{4}+1}=\frac{1}{2}\sqrt{5},\qquad \sqrt{3}>\frac{1}{2}\sqrt{5}$

$\therefore\quad |z_1|>|z_2|$.
\end{solution}

一般来说，两个复数中，至少有一个不是实数，就不能比较它们的大小，但任意两个复数的模，总是可以比较大小的.

\begin{example}
    如果$a$，$b$都是正实数，试求：
\begin{multicols}{4}
\begin{enumerate}[(1)]
    \item $\arg a$
    \item $\arg (-a)$
    \item $\arg (bi)$
    \item $\arg (-bi)$
\end{enumerate}
\end{multicols}
\end{example}

\begin{solution}
由于$a$，$b$是正实数，因此：
复数$a$，$-a$分别对应着复平面中正、负半实轴上的点（不包含原点）. 所以
\begin{multicols}{2}
    \begin{enumerate}[(1)]
        \item $\arg (a)=0$
        \item $\arg (-a)=\pi$
    \end{enumerate}
    \end{multicols}
复数$bi$，$-bi$分别对应着复平面中正、负半虚轴上的点（不包含原点）. 所以
\begin{multicols}{2}
    \begin{enumerate}[(1)]\setcounter{enumi}{2}
        \item $\arg (bi)=\frac{\pi}{2}$
        \item $\arg (-bi)=\frac{3\pi}{2}$
    \end{enumerate}
    \end{multicols}
\end{solution}

\begin{example}
    设$z\in\mathbb{C}$，试说明满足下列条件的点Z的集合组成什么样的图形？
\begin{multicols}{2}
\begin{enumerate}[(1)]
    \item $|z|=a$
    \item $2<|z|\le 4$
\end{enumerate}
\end{multicols}
\end{example}

\begin{solution}
\begin{enumerate}[(1)]
    \item 复数$z$的模等于$a\ge 0$，就是说它对应于复平面上的向量$\VEC{OZ}$的长度为$a$，也就是说它对应于复平面上的点$Z$到原点的距离为$a$. 所以，满足条件$|z|=a$的点的集合，当$a>0$时组成一个以$O$为圆心、以$a$为半径的圆；当$a=0$时，就是一个点（原点）。
\item 先将条件$2<|z|\le 4$化为$\begin{cases}
    |z|>2\\ |z|\le 4
\end{cases}$

\begin{minipage}{.45\textwidth}
 再分别考虑：满足$|z|>2$的点集组成圆$|z|=2$的外部；满足$|z|\le 4$的点集组成圆$|z|=4$的内部（包含边界）。同时满足以上两条件的点集，就是已求两集合的交集，也就是所要求的集合.

因此，满足$2<|z|\le 4$的点的集合是以$O$为圆心，以2及4为半径的两圆所夹的圆环，其中包含外边界，但不包括内边界（图1.4）.   
\end{minipage}\hfill
\begin{minipage}{.45\textwidth}
    \centering
\begin{tikzpicture}[>=stealth, scale=.4]

\draw[ thick, pattern=north east lines](0,0) circle (4);
\fill[color=white](0,0) circle (2);    
\draw[dashed,  thick](0,0) circle (2);    
\draw[->](-6,0)--(6,0)node[right]{$x$};
\draw[->](0,-6)--(0,6)node[right]{$y$};
\node [below right]{$O$};
\node at (2,0)[below right]{2};
\node at (4,0)[below right]{4};
\node at (0,2)[below right]{2};
\node at (0,4)[above right]{4};

\end{tikzpicture}
    \captionof{figure}{ }
\end{minipage}

\end{enumerate}    
\end{solution}

\begin{ex}
\begin{enumerate}
    \item 求出下列复数的模，并按其大小用不等号连起来：
\[\sqrt{2}-i,\quad -2+i,\quad 7,\quad 2i,\quad -1-i\]
  \item   设$|z|=r$, $\arg z=\alpha$，试求出$|\bar z|$, $\arg \bar z$.
  \item 求证：复数$z_1=-3i$, $z_2=\sqrt{2}+\sqrt{7}i$, $z_3=\sqrt{7}-\sqrt{2}i$, $z_4=-2+\sqrt{5}i$所对应的四个点共圆，并写出这个圆周上的所有点应满足的条件.
\item   你能指出满足下列条件的复数所对应的点集组成什么图形吗？
\begin{multicols}{2}
\begin{enumerate}[(1)]
    \item $\arg z=\frac{\pi}{4}$
    \item $\frac{\pi}{4}\le \arg z\le \frac{\pi}{2}$
\end{enumerate}
\end{multicols}
\end{enumerate}

\end{ex}

\subsection*{习题1.1}
\begin{enumerate}
    \item 设复数$z= (a^{2}- 2a- 3 )+( a^{2}+ 2a- 3) i$, $a\in \mathbb{R} $, 
试问在什么条件下，这一复数是
\begin{multicols}{4}
\begin{enumerate}[(1)]
    \item 实 数 ？
    \item 零？
    \item 虚数？
    \item 纯虚数？
\end{enumerate}
\end{multicols}

\item 实数$a,b$取何值时，才能使复数$z_1=z_2$:
\begin{enumerate}[(1)]
    \item $z_{1}= \left(\frac{1}{2}a+ b \right)+\left( 5a+\frac{2}{3}b \right) i,\quad z_{2}= - 4+ 16i$ ;
    \item $z_{1}=(a+b)+24i,\quad z_{2}=-(5+abi)$;
    \item $z_{1}=8+2abi,\quad z_{2}=6i+a^{2}-b^{2}$;
    \item $z_{1}=2a^{2}-5a+2,\quad z_{2}=(b^{2}+b-2)i$;
    \item $z_{1}= 0,\quad z_{2}= (a- 1)^{2}+ (b^{2}+ 2b)i$.
\end{enumerate}
\item 求出下列各复数的共轭复数，并在复平面上描出表示这些互为共轭复数的点。
\[1- i, \quad - 2+ \sqrt {2}i, \quad - \sqrt {3}i, \quad 2, \quad i, \quad \sqrt {3}+ i\]
\item 求出上题中各复数的模及幅角的主值（可用反三角函数表示）.
\item 实数$x,y$取何值时，才能使$z_1=\bar z_2$:
\begin{enumerate}[(1)]
    \item $z_{1}=( x-y )+5i,\quad  z_{2}=1-( 2x+y ) i$;
    \item $z_{1}=xy-( x+y)i,\quad  z_{2}=5+24 i$; 
    \item $z_{1}=( x^{2}-3x-2 )+( y^{2}-5y-6 ) i,\quad z_{2}=2$;
    \item $z_1=0,\quad z_2=(2x^2-5x+2)+i(y^2+y-2)$.
\end{enumerate}


\item 求证：两个复数互为共轭的必要条件是这两个复数
的模相等。
\item  在下列各题中，要使等式$| z_1|=2| z_2|$成立，试
求实数$x,y$的值：
\begin{enumerate}[(1)]
    \item $z_{1}=x-1-2i,\quad  z_{2}=1+( x+1 )i$;
    \item $z_{1}=( x-y+1 )+2 ( x+y-3 )i,\quad z_{2}=\frac{x-y+1}{2\sqrt{2}}i $
\end{enumerate}

\item  设$a,b\in\mathbb{R}$, 它们取何值时，表示复数$z=(a-b-1)+(2a-b-5)i$在复平面上的点，才能在$y=x$直线上？
\item 试求满足下列条件的复数$x+yi$在复平面上对应的点
的轨迹：
\begin{multicols}{3}
\begin{enumerate}[(1)]
    \item $|x+yi|=3$
    \item $|x+yi|<3$
    \item $|x+yi|\ge 3$
    \item $2\le |x+yi|\le 3$
    \item $0<|x+yi|<3$
\end{enumerate}
\end{multicols}
\item 设$z=a+bi,\; a,b\in\mathbb{R}$. 问：满足下列条件的点
$Z(a,b)$的集合组成什么图形？
\begin{multicols}{2}
\begin{enumerate}[(1)]
    \item $0\le |a|<1$
    \item $a>0,\; b>0,\quad a^2+b^2>4$
\end{enumerate}
\end{multicols}

\item 想想看：若已知一个非零复 数$z=a+bi$, 你能求
出它的幅角来吗？它的主值又如何求出？举例说明。

\item 满足条件$\frac{\pi}{6}\le \arg z\le \frac{\pi}{2}$的复数，在复平面上对应着什么样的点集？

\end{enumerate}

\section{复数的四则运算}
本节讨论复数的四则运算，我们将根据引进的新数——虚数单位$i$的特征$i^2=-1$和希望它遵从“数系运算通性”的要求，合理地提出运算法则，并进一步应用这些法则，有效地进行运算，同时，根据复数的几何意义，我们将相应地给出复数运算的几何意义.

\subsection{复数的加、减法}
\subsubsection{加法}
复数的加法运算按照以下规定的法则进行：设$z_1=a+bi$, $z_2=c+di$是任意两个复数，则
\[z_1+z_2=(a+bi)+(c+di)=(a+c)+(b+d)i\]

可见，任意两个复数的和仍然是一个复数.

特别地，根据这一规定我们可求出任一个实数$a$和任一纯虚数$bi$的和：
\[(a+0i)+(0+bi)=(a+0)+(0+b)i=a+bi\]

由此可知，复数$a+bi$，既可认定为一个数，又可认定它为实数$a$与纯虚数$bi$的和，今后，对于表达复数$a+bi$中的“$+$”号与表达加法运算的“$+$”号，我们将认为是一样的。

不难验证，复数加法法则的规定满足“运算通性”。
即 : 设 $z_1, z_2, z_3\in \mathbb{C}$, 则
\begin{align}
    z_{1}+ z_2 &= z_{2}+ z_{1}\tag{交换律成立}\\
    z_{1}+\left( z_{2}+z_{3} \right)&=\left( z_{1}+z_{2} \right)+z_{3}=z_{1}+z_{2}+z_{3}\tag{结合律成立}\\
    z_{1}+ 0&= 0+ z_{1}= z_{1}\tag{零的运算特性保持}
\end{align}

\begin{example}
    试求下列复数的和：
\begin{enumerate}[(1)]
    \item $(5-3i)+(-\sqrt{2}+i)$
    \item $( 9+ 6i) + ( 9- 6i)$
    \item $( 1- i) + ( - 1+ 2i) + 7+ \sqrt {2}i$
\end{enumerate}
\end{example}

\begin{solution}
\begin{enumerate}[(1)]
    \item $( 5- 3i) + ( - \sqrt {2}+ i) = (5- \sqrt {2})-2i$
    \item $( 9+ 6i) + ( 9- 6i) = 18+ 0i= 18$
    \item $( 1- i) + ( - 1+ 2i) +7+ \sqrt {2}i=(1-1+7)+(-1+2+\sqrt{2})i=7+(1+\sqrt{2})i$
\end{enumerate}
\end{solution}


由于复数可以用复平面上的“位置”向量来表示，因此，复数加法的几何意义就是：复数的加法可以按照向量的“平行四边形”法则来进行，即，先将求和的两个复数所对应的向量$\overrightarrow{OZ_{1}}$, $\overrightarrow{OZ_{2}}$画出，若这两个向量不共线，再以这两个向量为邻边作平行四边形$OZ_1ZZ_2$, 则这个平行 四 边形的对角线$OZ$所示的向量$\overrightarrow{OZ}$, 就 对 应 着 已 知 两 个 复 数 的 和(如图1.5).

\begin{figure}[htp]
    \centering
\begin{tikzpicture}[>=stealth, scale=.8]
\draw[->](-1,0)--(4,0)node[right]{$x$};
\draw[->](0,-1)--(0,4)node[right]{$y$};
\node [below right]{$O$};
\draw[very thick, ->](0,0)--(.5,2)node[above]{$Z_2$};
\draw[very thick, ->](0,0)--(2.5,1)node[right]{$Z_1$};
\draw[very thick, ->](0,0)--(3,3)node[above]{$Z$};
\draw[dashed](.5,2)--(3,3)--(2.5,1);


\end{tikzpicture}
    \caption{}
\end{figure}

若$\overrightarrow{OZ_{1}}$ 与$\overrightarrow{OZ_{2}}$共线，只要在$\overrightarrow{OZ_{1}}$的方向上紧接着延长一个$\overrightarrow{Z_{1}Z}$，使$\overrightarrow{Z_{1}Z}=\overrightarrow{OZ_{2}}$. 这时同样得到向量$\overrightarrow{OZ}$, 它对应着已知两复数的和（相当于作一个退缩成一条线段的平行四边形）.

显然，两个共轭复数的和是一个实数，它对应着实轴上的一个向量. 

\subsubsection{减法}
复数的减法仍定义为加法
的逆运算，即：满足$(c+di)+(x+yi)=a+bi$的复数$x+yi$，就叫做复数$a+bi$减去$c+di$的差，记作
\[x+yi=(a+bi)-(c+di)\]

根据复数加法法则和复数相等的定义，可以得出复数的减法法则
\[(a+bi)-(c+di)=(a-c)+(b-d)i\]

可见，两复数的差，仍是一个确定的复数，就是说，复数集对于减法也是封闭的.

同样，复数的减法的几何意义就是用向量减法的“三角
形法则”来表示，即：两个复数的差
$z_1-z_2$（向量$\VEC{OZ_1}-\VEC{OZ_2}$）与连结两个向量终点并
指向被减向量终点的向量对应
（图1.6），即
\[\VEC{OZ_1}-\VEC{OZ_2}=\VEC{Z_2Z_1}\]

\begin{figure}[htp]
    \centering
\begin{tikzpicture}[>=stealth, scale=.8]
\draw[->](-1,0)--(4,0)node[right]{$x$};
\draw[->](0,-1)--(0,4)node[right]{$y$};
\node [below right]{$O$};
\draw[ thick, ->](0,0)--(3,3)node[above]{$Z_2$};
\draw[ thick, ->](0,0)--(2.5,1)node[right]{$Z_1$};
\draw[very thick, ->](3,3)--(2.5,1);


\end{tikzpicture}
    \caption{}
\end{figure}

综上所述，复数的加、减法与我们熟悉的多项式加、减法是类似的，就是把复数的实部、虚部分别相加、减，即
\[(a+bi)\pm (c+di)=(a\pm c)+(b\pm d)i\]

\begin{example}
    计算$(3-4i)+(-2+i)-(5+2i)$
\end{example}

\begin{solution}
 $$(3-4i)+(-2+i)-(5+2i)
=(3-2-5)+(-4+1-2)i=-4-5i.$$   
\end{solution}

\begin{example}
    设复数$z_1=a+bi$, $z_2=c+di$, 试求这两复数在复平面上所描述的两点之间的距离.
\end{example}

\begin{solution}
    由于$z_1$与$z_2$在复平面上所描述的两点$Z_1$, $Z_2$之间的距离$d$，就是向量$\VEC{Z_1Z_2}$的模，因而，只要求出向量$\VEC{Z_1Z_2}$所对应的复数，并求出它的模即可。

    $\because\quad \VEC{Z_1Z_2}=\VEC{OZ_2}-\VEC{OZ_1}$

$\therefore\quad \VEC{Z_1Z_2}$对应于复数$z_2-z_1$，于是
\[\begin{split}
d&=|\VEC{Z_1Z_2}|=|z_2-z_1|\\
&=|(c+di)-(a+bi)|=|(c-a)+(d-b)i|\\
&=\sqrt{(c-a)^2+(d-b)^2}    
\end{split}\]
\end{solution}

\begin{example}
    根据复数的几何意义及向量表示，试求复平面内以$P(a,b)$点为圆心，$r$为半径的圆的方程.
\end{example}

\begin{figure}[htp]
    \centering
\begin{tikzpicture}[>=stealth, scale=1]
\draw[->](-1,0)--(4,0)node[right]{$x$};
\draw[->](0,-1)--(0,4)node[right]{$y$};
\node [below right]{$O$};
\draw[ thick, ->](0,0)--(1.5,3)node[above]{$Z$};
\draw[ thick, ->](0,0)--(2,2)node[right]{$P$};
\draw[thick](2,2)circle(1.12);
\draw[->, very thick](2,2)--node[right]{$r$}(1.5,3);

\end{tikzpicture}
    \caption{}
\end{figure}


\begin{solution}
    设圆心$P(a,b)$与复数$p=a+bi$对应，圆上任一点$Z(x,y)$与复数$z=x+yi$对应，由于已知圆的半径为$r$，因
此，向量$\VEC{PZ}$的模就是定长$r$. 即
$|\VEC{PZ}|=r$.

又由于$\VEC{PZ}=\VEC{OZ}-\VEC{OP}$，因而向量$\VEC{PZ}$就对应着复数$z-p=(x+yi)-(a+bi)=(x-a)+(y-b)i$.

所以，圆上任一点$Z$应满足$|\VEC{PZ}|=|\VEC{OZ}-\VEC{OP}|=r$，即
\[|z-p|=r\]
这就是复平面上以$P$为圆心，$r$为半径的圆的方程.

若进一步求出$|z-p|=|(x-a)+(y-b)i|
=\sqrt{(x-a)^2+(y-b)^2}$
代入圆方程，并化简。就得出用实数表示的圆的一般方程：
\[(x-a)^2+(y-b)^2=r^2\]
特别地，当$P$点在原点时，圆的方程就变成为
$$|z|=r$$
或
$$x^2+y^2=r^2$$
\end{solution}

\begin{ex}
\begin{enumerate}
    \item 利用复数加法法则和实数、虚数单位运算通性，验证：复数加法的交换律和结合律都成立.
    \item 计算（分别求和，再求差）：
\begin{multicols}{2}
\begin{enumerate}[(1)]
    \item $(2+3i)\pm (1+i)$
    \item $(2+4i)\pm (1-i)$
    \item $(-3i-2)\pm (-1-i)$
    \item $-2i\pm (-4+i)$
    \item $(2-3i)\pm (-5i)\pm(-1+7i)$
\end{enumerate}
\end{multicols}
\item 将上题中的各小题，用向量法求和.
\item 求证：两个互为共轭的复数之和为实数；两互为共轭的复数之差为纯虚数或零.
\item 想想看：在复平面上，以$P(a,b)$为圆心、以r为半径的圆内所有点的集合，可用什么样的条件给定？
\end{enumerate}
\end{ex}

\subsection{复数的乘、除法}
\subsubsection{乘法}
复数的乘法运算按照以下规定的法则进行：设$z_1=a+bi$, $z_2=c+di$是任意两复数，则
\[z_1\cdot z_2=(a+bi)\cdot (c+di)=(ac-bd)+(ad+bc)i\]
其中，由于我们要求乘法运算仍应保持“运算通性”，因而，这一规定说明复数的乘法与多项式乘法是类似的，只要将结果中的$i^2$换成$-1$，分别合并实部与虚部即可．

可见，任意两个复数的乘积仍然是一个复数，也就是说，复数集对于乘法也是封闭的.

特别地，根据这一规定，对两个互为共轭复数$z$与$\bar z$，我们有
\[z\cdot \bar z=(x+yi)\cdot (x-yi)=x^2+y^2\]

因此，互为共轭的两复数之积是一个实数，它等于每一复数的模的平方，即
\[z\cdot \bar z=|z|^2=|\bar z|^2\]

容易验证，复数乘法运算也具有“数系运算通性”，即对于任意$z_1,z_2,z_3\in\mathbb{C}$，都有
\begin{align}
z_1\cdot z_2  &= z_2\cdot z_1 \tag{乘法交换律成立}\\
z_1\cdot (z_2\cdot z_3)  &= (z_1\cdot z_2)\cdot z_3=z_1\cdot z_2\cdot z_3 \tag{乘法结合律成立}\\
z_1\cdot (z_2+ z_3)  &= z_1\cdot z_2+ z_1\cdot z_3 \tag{乘法对加法的分配律成立}\\
z_1\cdot 0  &= 0\cdot z_1=0 \tag{零的运算特性保持}\\
z_1\cdot 1  &= 1\cdot z_1=z_1 \tag{1的运算特性保持}
\end{align}

还应该指出，复数乘方运算、幂的概念与实数完全相同，而且由于复数的乘法满足交换律、结合律，所以在实数集$\R$中的指数运算律，在复数集$\mathbb{C}$中仍然成立，即

对任意$z_1,z_2,z\in\mathbb{C}$及$m,n\in\N$, 都有
\[\begin{split}
    z^m\cdot z^n&=z^{m+n}\\
    (z^m)^n&=z^{m\cdot n}\\
    (z_1\cdot z_2)^n&=z^{n}_1\cdot z^{n}_2\\
\end{split}\]

这样，我们由定义$i^2=-1$，运用指数运算律可以得出：
\[\begin{split}
i^3&=i\cdot i^2=-i\\
i^4&=i^2\cdot i^2=(-1)\cdot (-1)=1\\
i^5&=i^4\cdot i=i\\
\cdots&\cdots
\end{split}\]

一般地，对任意$n\in\mathbb{N}$，我们可归纳出：
\[\begin{split}
    i^{4n}&=(i^4)^n=1^n=1\\
    i^{4n+1}&=i^{4n}\cdot i=i\\
    i^{4n+2}&=i^{4n}\cdot i^2=-1\\
    i^{4n+3}&=i^{4n}\cdot i^3=-i\\
\end{split}\]
即，对于数$n\in\N$，则
\[i^{4n}=1,\qquad i^{4n+1}=i,\qquad i^{4n+2}=-1,\qquad i^{4n+3}=-i\]

\begin{example}
    计算$(5-4i)(1-i)(-2+3i)$.
\end{example}

\begin{solution}
\[(5-4i)(1-i)(-2+3i)=(1- 9i) (- 2+ 3i) = 25+ 21i\]
\end{solution}

\begin{example}
    计算$\left(-\frac12+\frac{\sqrt{3}}{2}i\right)^6$
\end{example}

\begin{solution}
\[\begin{split}
    \left(-\frac12+\frac{\sqrt{3}}{2}i\right)^6&=\left[\left(-\frac12+\frac{\sqrt{3}}{2}i\right)^2\right]^3\\
&= \left(\frac{1}{4}-\frac{\sqrt{3}}{2}i+\frac{3i^{2}}{4}\right)^{3}=\left(-\frac{1}{2}-\frac{\sqrt{3}}{2}i\right)^3\\
&=\left( - \frac 12\right) ^{3}- 3\left( \frac {- 1}{2}\right)^{2}\cdot \frac {\sqrt {3}}2i+3\left(-\frac{1}{2}\right)\left(\frac{\sqrt{3}}{2}i\right)^{2}-\left(\frac{\sqrt{3}}{2}i\right)^{3}\\
&=-\frac{1}{8}-\frac{3\sqrt{3}}{8}i+\frac{9}{8}+\frac{3\sqrt{3}}{8}i=1
\end{split}\]
\end{solution}


\subsubsection{除法}

复数的除法同样定义为乘法的逆运算，即：满足$(c+di)(x+yi)=a+bi\; (c+di\neq0)$的复数$x+yi$叫做复数$a+bi$除以$c+di$的商，记作
$$x+yi=\frac{a+bi}{c+di}\quad \text{或}\quad( a+bi )\div( c+di ) $$

根据两个共轭复数的乘积是一个实数以及乘法运算法
则，可以得出复数的除法法则：

当$c+ di\neq 0$ 时，
\[\frac{a+bi}{c+di}=\frac{(a+bi)(c-di)}{(c+di)(c-di)}=\frac{ac+bd}{c^2+d^2}+\frac{bc-ad}{c^2+d^2}i\]
其中由于$c+di\ne 0$，因而$c$、$d$不同时为零，所以，$c^2+d^2\ne 0$．可见这样求得的商是唯一确定的一个复数，也就是说，复数集$\mathbb{C}$对于除法（除数不为零）也是封闭的.

\begin{example}
    计算$(2-3i)\div (-3+4i)$
\end{example}

\begin{solution}
\[\begin{split}
    (2-3i)\div (-3+4i)&=\frac{2-3i}{-3+4i}=\frac{(2-3i)(-3-4i)}{(-3+4i)(-3-4i)}\\
    &=\frac{-18+i}{25}=-\frac{18}{25}+\frac{1}{25}i
\end{split}\]
\end{solution}

可以看出，复数乘、除法法则的公式并不需要硬记，乘法类似于多项式乘法，除法可视为分式的分子、分母同乘以除数的共轭数，并注意正确应用运算通性及$i^2=-1$，就可得出结果.

\begin{ex}
\begin{enumerate}
    \item 验证复数乘法对于加法满足分配律.
    \item 计算下列各式：
\begin{multicols}{2}
\begin{enumerate}[(1)]
\item $(4-3i)(-\sqrt{3}i)$
\item $\left(\frac{1}{2}-\frac{\sqrt{3}}{2}i\right)(1+i)$
\item $\left(\frac{1}{2}-\frac{\sqrt{3}}{2}i\right)^3\cdot \left(\frac{1}{2}+\frac{\sqrt{3}}{2}i\right)^3$
\item $(-\sqrt{77}+\sqrt{112}i)\cdot 0$
\end{enumerate}
\end{multicols}

\item  （口答）计算：
\[i^{13},\quad i^{22},\quad i^{36},\quad i^{43},\quad i^{70},\quad i^{101},\quad i^{355},\quad i^{404} \]
\[\frac{1}{i},\quad \frac{1}{i^2},\quad \frac{1}{i^3},\quad \frac{1}{i^4},\quad \frac{1}{i^{11}}\]
\item     计算：
\begin{multicols}{3}
\begin{enumerate}[(1)]
    \item $\frac{1}{1+i}$
    \item $\frac{1}{\sqrt{2}i}$
    \item $\frac{2i}{1-i}$
    \item $\frac{2+i}{7+4i}$
    \item $\frac{1}{(9+2i)^2}$
    \item $\frac{(1-i)^2}{(1+i)^3}$
\end{enumerate}
\end{multicols}
\end{enumerate}
\end{ex}

\subsection{复数集内的多项式与方程}
我们已经学习过实数范围内的多项式，而且明确多项式的元（未知数）只要求它满足“运算通性”，因此，对于复系数的多项式运算或未知数（元）取复数值的多项式求值，都可以与实系数多项式同样处理.

\begin{enumerate}
    \item 复系数多项式的四则运算举例
    
例如$f(x)=x+i$, $g(x)=2x+1-i$，则
\[\begin{split}
f(x)+g(x)&=3x+1\\
f(x)-g(x)&=-x-1+2i\\
f(x)\cdot g(x)&=2x^2+(1+i)x+(1+i)\\
g(x)&=2\cdot f(x)+1-3i    
\end{split}\]
这就是说，$g(x)$除
以$f(x)$可得商2，余$1-3i$（由综合除法得）.

\item 多项式求值举例

例如$f(x)=x^2+x+1$，
则
\[\begin{split}
  f(i)&=i^2+i+1=i\\
f(1+i)&=(1+i)^2+(1+i)+1=2+3i.  
\end{split}\]

\item 复系数多项式的因式分解举例

例如：在复数范围内，$f(x)=x^2+1$可以进一步分解
\[x^2+1=x^2-(-1)=x^2-i^2=(x+i)(x-i)\]

\item 复数范围内的方程与方程组，均可按实数范围时的方法求解

\item 由带余除法导出的多项式的余式定理，对于复系数
的多项式仍然适用. 例如
$f(x)=x^4+2x^3-3x^2+5$除以$x-2i$所得的余式为：
\[R=f(2i)=(2i)^4+2(2i)^3-3(2i)^2+5
=33-16i\]
\end{enumerate}

\begin{example}
    在复数范围内分解因式$x^4-16$
\end{example}

\begin{solution}
\[x^{4}-16=(x^{2}-4)(x^{2}+4)=(x- 2) (x+ 2) (x- 2i) (x+ 2i)\]
\end{solution}

\begin{example}
    解方程$x^2-6x+10=0$
\end{example}

\begin{solution}
    利用求根公式，得
$$x=\frac{6\pm\sqrt{36-40}}{2}=\frac{6\pm\sqrt{-4}}{2}=3\pm i$$
$\therefore\quad x_{1}= 3+ i,\quad x_{2}= 3- i$
\end{solution}

\begin{example}
     解方程
$\begin{cases}
    2x+(1+i)y=6-2i\\
    (3-i)x+2i y=13-i
\end{cases}$
\end{example}

\begin{solution}
   用行列式法解，由于
\[\begin{split}
    \Delta&=\begin{vmatrix}
    2& 1+i\\
    3-i& 2i
\end{vmatrix}=4i-(1+i)(3-i)=-4+2i\\
\Delta_{1}&=\begin{vmatrix}
   6- 2i& 1+i\\
    13-i& 2i
\end{vmatrix}=-10\\
\Delta_{2}&=\begin{vmatrix}
    2& 6-2i\\
     3-i& 13-i
 \end{vmatrix}=10+10i\\
\end{split}\]
$\therefore\quad x= \frac {\Delta_{1}}{\Delta} = \frac {- 10}{- 4+ 2i}= 2+ i,\quad 
y=\frac{\Delta_{2}}{\Delta}=\frac{10+10i}{-4+2i}=-1-3i$

因此，方程组的解集为$\{(x,\; y)\}=\{(2+i,\;  -1-3i)\}$.
\end{solution}

\begin{ex}
\begin{enumerate}
    \item 已知$f(x)=2x^2+(1+i)x+(1-i)$，试求$f(i)$, $f(-i)$, $f(1+i)$和$f(0)$的值.
    \item 在复数范围内分解因式$2x^4-5x^2-3$.
    \item 解方程组$\begin{cases}
        ix-y=2(i-1)\\
        x-iy=0
    \end{cases}$
\end{enumerate}
\end{ex}

\subsection*{习题1.2}

\begin{enumerate}
    \item 计算：
\begin{enumerate}[(1)]
    \item $\left(\frac{2}{3}+i\right)+\left(1-\frac{2}{3}i\right)-\left(\frac{1}{2}+\frac{3}{4}i\right)$
    \item $[(a+b)+(a-b)i]-[(a-b)-(a+b)i]$
\end{enumerate}

\item 设复数$z_1=3+2i$, $z_2=-2+5i$在复平面上分别对应着点$Z_1,Z_2$. 试求向量$\VEC{Z_1Z_2}$与$\VEC{Z_2Z_1}$所表示的复数．
\item 试求复平面上点$P(2,1)$与点$Q(3,-1)$之间的距离.
\item 求证：两个复数和的共轭复数，等于这两个复数的共轭复数之和.
\item 设$z_1,z_2$是非零复数，用几何方法证明：
\begin{enumerate}[(1)]
    \item $\big||z_1|-|z_2|\big|\le |z_1\pm z_2|\le |z_1|+|z_2|$
    \item $|z_1+z_2|^2+|z_1-z_2|^2=2|z_1|^2+2|z_2|^2$
\end{enumerate}
\item  写出复平面内任意两点连线段的垂直平分线的方程。
\item  设复平面内有两定点$F_1(-\sqrt{5},0)$, $F_2(\sqrt{5},0)$，试求到这两定点的距离之和等于6的点的轨迹方程（椭圆方程）.
\item  计算：
\begin{enumerate}[(1)]
    \item $(\sqrt{a}+\sqrt{b}i)(\sqrt{a}-\sqrt{b}i)\quad (a,b\in\mathbb{R})$
    \item $(1-2i)(-0.2+0.3i)(0.5-0.4i)$
    \item $(1+i)+(2-i^{3})+(3-i^{5})+(4-i^{7})$
    \item $(\sqrt{2}-\sqrt{2}i)^{2}\cdot\left(\frac{1}{4}+\frac{1}{4}i\right)$
    \item $(a+bi)^{3}+(a-bi)^{3}$
    \item $\begin{vmatrix}
        1-i&i\\1+i&-i
    \end{vmatrix}+\begin{vmatrix}
        i&1&0\\1&2i&1\\0&1&3i
    \end{vmatrix}$
\end{enumerate}

\item 利用公式$(a+bi)(a-bi)=a^2+b^2$, 把下列各式分解因式：
\begin{multicols}{2}
 \begin{enumerate}[(1)]
    \item $a^{4}-b^{4}$
    \item $x^{2}+2x+3$
    \item $a^{2}+2ab+b^{2}+c^{2}$
    \item $x^{3}-2x+4$
\end{enumerate}   
\end{multicols}


\item 设$\omega = - \frac 12+ \frac {\sqrt {3}}2i$, 求证：
\begin{enumerate}[(1)]
    \item $1+\omega+\omega^{2}=0$
    \item $\omega^{3}=1$
    \item $a^{3}+b^{3}+c^{3}-3abc=(a+b+c)(a+\omega b+\omega^{2}c)(a+\omega^{2}b+\omega c)$
\end{enumerate}
\item 计算：
\begin{multicols}{2}
\begin{enumerate}[(1)]
    \item $\frac{1}{2-4i}$
    \item $\frac{3-2i}{1+i}$
    \item $\frac{1+i}{(1-i)^{2}}$
    \item $\frac{1+2i}{2-i^3}$
    \item $\frac{(1-2i)^2}{3-4i}-\frac{(2+i)^2}{4-3i}$
    \item $\frac{\sqrt{3}-\sqrt{2}i}{\sqrt{3}+\sqrt{2}i}-\frac{\sqrt{2}+\sqrt{3}i}{\sqrt{2}-\sqrt{3}i}$
\end{enumerate}
\end{multicols}

\item 若$z_1,z_2\in\mathbb{C}$, 且$z_1\cdot z_2=0$, 则$z_1,z_2$中至少有一个等于0.

\item 已知$\frac1z=\frac{1}{z_{1}}+\frac{1}{z_{2}}$, 且$z_1=2+3i$, $z_{2}=1-i$,
求$z$.

\item 已知$z^2=5-12i$, 求$z$.

\item 证明下列各命题：
\begin{enumerate}[(1)]
    \item 任两个复数的和、差、积、商(除数不为零)的
共轭复数，分别等于这两个复数的共轭复数的和、差、积、商.
\item 若$f(x)$是一个一元多项式，则对于任意复数
$z$, 有$\overline{f(z)}=f(\bar{z})$.
\end{enumerate}

\item 已知$f(z)=\frac{z^2-z+1}{z^2+z+1}$, 求$f(2+3i)$, $f(1-i)$的值.

\staritem  如果规定$i^0=1$, $i^{-m}=\frac{1}{i^m}\; (m\in\mathbb{N})$, 证明：对一切整数$n$，以下等式成立：
$$i^{4n}={1},\quad i^{4n+1}=i,\quad  i^{4n+2}=-1,\quad  i^{4n+3}=-i.$$
\staritem 根据$i^2=-1$及要求复数运算满足“通性”,试说
明复 数 加 法 、 乘 法 法 则 规 定 是 合 理 的。

\item 怎样的复数$z$, 才能满足条件：
\begin{multicols}{2}
\begin{enumerate}[(1)]
    \item $\overline {z}= z$
    \item $\overline {z}= -z$
    \item $z\cdot \bar z=[{\rm Re}(z)]^2$
    \item $z-\bar z=2[{\rm Im}(z)]$
\end{enumerate}    
\end{multicols}

\item 应用余式定理，求下列各余式；
\begin{enumerate}[(1)]
    \item $f( x) = x^{4}- 3x^{3}+ 2x^{2}- x+ i$
除以$g(x)=x+i$
\item $f( x) = x^{4}- 3x^{3}+ 2x^{2}- x+ i$
除以$h(x)=(1+i)x-1+i$
\end{enumerate}
\item 应用综合除法，将多项式
$f\left(x\right)=x^{3}+3\left({1+i}\right)x^{2}-2(5-3i)x-4\left(1+2i\right)$
展开成$(x+i)$的幂的代数和形式.

\item 解方程组$\begin{cases}
    x_1+2x_2=1+i\\
    ix_2+x_3=2i\\
    ix_1-ix_3=1-i
\end{cases}$
\end{enumerate}

\section{复数的三角形式}

我们已经建立了复数集
\[\mathbb{C}=\{x+yi\mid x,y\in\R,\; i^2=-1\}\]
与复平面上的点集以及复平面上的全体向量（由原点出发）集合间的一一对应关系，并且明确了复数的模与幅角、幅角的主值等概念，这一节我们将学习复数的另一种表达形式. 就是由它的模与幅角所确定的三角形式，这将有助于复数乘法、除法、乘方、开方运算的进一步简便。

\subsection{复数的三角形式}

设$z=a+bi,\; (a,b\in\R)$，对应着复平面上的向量$\VEC{OZ}$（图1.8），并设它的模$|\VEC{OZ}|=r$，幅角为$\theta$. 若$z=a+bi\ne 0$，则可得
\[a=r\cdot \cos\theta,\qquad b=r\cdot \sin\theta\]
\begin{figure}[htp]
    \centering
\begin{tikzpicture}[>=stealth, scale=.7]
\draw[->](-1,0)--(4,0)node[right]{$x$};
\draw[->](0,-1)--(0,3)node[right]{$y$};
\node[below left]{$O$};
\draw[->, very thick](0,0)--node[above left]{$r$}(3,2.5)node[above]{$Z$};
\draw[dashed](3,0)--node[right]{$b$}(3,2.5);
\node at (1.5,0)[below]{$a$};
\draw[->](1,0)node[above right]{$\theta$} arc (0:41:1);
\end{tikzpicture}
    \caption{}
\end{figure}


$\therefore\quad z=a+bi=r\left(\pc{\theta}\right)$

因此，对于任何一个复数$z=a+bi$，都可以表示成
$r(\cos\theta +i\sin\theta)$
的形式（其中$r=\sqrt{a^2+b^2}$），并叫做这个复数的三角形式（也叫做复数的极坐标形式）；为了方便，我们将$a+bi$叫做复数的代数形式.

复数的代数形式$z=a+bi$与三角形式$z=r(\cos\theta +i\sin\theta)$之间，由图1.8可知有下列互化的关系：
\[\begin{cases}
 a=r\cos\theta\\
b=r\sin\theta   
\end{cases},\qquad \begin{cases}
    r=\sqrt{a^2+b^2}\\ \tan\theta=b/a
\end{cases}\]
（$\theta$只要取主值，且要考虑$a$，$b$的正负）.

\begin{example}
 将下列复数的代数形式化为三角形式： 
\begin{multicols}{2}
\begin{enumerate}[(1)]
    \item $\sqrt{3}+i$
    \item $-\frac{1}{2}+\frac{\sqrt{3}}{2}i$
    \item $-1$
    \item $\cos\theta-i\sin\theta$
\end{enumerate}
\end{multicols}  
\end{example}

\begin{solution}
\begin{enumerate}[(1)]
    \item $r=\sqrt{3+1}=2$, $\tan\theta=\frac{1}{\sqrt{3}}=\frac{\sqrt{3}}{3}$，且$a,b$均为正值，$\theta$在I象限，因而$\arg(\sqrt{3}
    +i)=\frac{\pi}{6}$

    $\therefore\quad \sqrt{3}+i=2\left(\pc{\frac{\pi}{6}}\right)$
\item $r=\sqrt{\frac{1}{4}+\frac{3}{4}}=1$, $\tan\theta=-\sqrt{3}$，且$a<0, b>0$，$\theta$在II象限，因而$\theta=\frac{2\pi}{3}$（主值）

$\therefore\quad -\frac{1}{2}+\frac{\sqrt{3}}{2}i=\pc{\frac{2\pi}{3}}$

\item $r=1$, $\tan\theta =0$，且$a<0$，因此$\arg(-1)=\pi$

$\therefore\quad -1=\pc{\pi}$

\item $r=\sqrt{\cos^2\theta+\sin^2\theta}=1$

$\therefore\quad \cos\theta-i\sin\theta=\pcx{-\theta}=\pcx{2\pi-\theta}$

\end{enumerate}
\end{solution}


\begin{example}
    将下列复数的三角形式化为代数形式：
\begin{multicols}{2}
\begin{enumerate}[(1)]
    \item $2\left(\pc{\frac{\pi}{4}}\right)$
    \item $\sqrt{2}\left(\pc{\frac{5\pi}{6}}\right)$
    \item $\pc{101\pi}$
    \item $5\left(\pc{\frac{\pi}{2}}\right)$
\end{enumerate}    
\end{multicols}
\end{example}

\begin{solution}
\begin{enumerate}[(1)]
    \item $2\left(\pc{\frac{\pi}{4}}\right)=\sqrt{2}+\sqrt{2}i$
    \item $\sqrt{2}\left(\pc{\frac{5\pi}{6}}\right)=-\frac{\sqrt{6}}{2}+\frac{\sqrt{2}}{2}i$
    \item $\pc{101\pi}=\pc{\pi}=-1$
    \item $5\left(\pc{\frac{\pi}{2}}\right)=5i$
\end{enumerate}  
\end{solution}

\begin{rmk}
复数的三角形式$r(\pc{\theta})$中，必须满足
    $r>0$、式中保持“$+$”号两特征，至于幅角$\theta$则不一定非要取主值，除常用主值表示外，也可以在幅角中任取一个。例如，
\[\sqrt{2}-\sqrt{2}i=2\left[\pcx{-\frac{\pi}{4}}\right]\]
也是复数$\sqrt{2}-\sqrt{2}i$的三角形式.

此外，若$z=a+bi=0$，即$z=0+0i$，这时$r=0$，其幅角是任意实数，则$0=0(\pc{\theta})$.
\end{rmk}

\begin{ex}
\begin{enumerate}
    \item 判断下列给出的是不是复数的三角形式？若不是，请把它们化为三角形式：
\begin{multicols}{2}
    \begin{enumerate}[(1)]
        \item $2\left(\cos\frac{\pi}{4}-i\sin\frac{\pi}{4}\right)$
        \item $-2\left(\pc{\frac{\pi}{3}}\right)$
        \item $\frac{1}{2}\left(\pc{\frac{3\pi}{4}}\right)$
        \item $\cos\frac{7\pi}{5}+i\sin\left(-\frac{3\pi}{5}\right)$
    \end{enumerate}
\end{multicols}
    \item 将复数表示成三角形式：
\begin{multicols}{3}
    \begin{enumerate}[(1)]
        \item $2$
        \item $-3$
        \item $4i$
        \item $-2i$
        \item $1+i$
        \item $-4+3i$
        \item $-1-\sqrt{3}i$
        \item $\frac{\sqrt{3}}{2}-\frac{1}{2}i$
    \end{enumerate}
\end{multicols}
\item    写出下列复数的代数形式：
\begin{multicols}{2}
    \begin{enumerate}[(1)]
        \item $3\left(\pc{\frac{\pi}{3}}\right)$
        \item $\sqrt{2}\left(\pc{\frac{3\pi}{4}}\right)$
        \item $\frac{1}{\sqrt{3}}\left[\pcx{-\frac{\pi}{6}}\right]$
        \item $5\left(\pc{\frac{7\pi}{2}}\right)$
    \end{enumerate}
\end{multicols}

\item 想一想：“两个复数若相等，则它们的模与幅角必定相等”对吗？为什么？
\end{enumerate}
\end{ex}

\subsection{再谈复数的乘、除法和乘方运算}
通过第二节对复数四则运算的学习，我们会感到：复数的代数形式在进行加、减法运算时，步骤十分简明易学，但进行乘、除法和乘方运算时，却显得繁杂.

我们将会看到，复数的三角形式在进行乘、除、乘方甚至开方运算时，比代数形式要简捷得多.

\begin{blk}
 {定理1} 两个复数的乘积是一个复数，它的模等于两乘数的模的乘积；它的幅角等于两乘数的幅角的和，即   

设$z_1=r_1(\pc{\theta_1})$, $z_2=r_2(\pc{\theta_2})$，则
\[\begin{split}
z_1\cdot z_2 &=r_1(\pc{\theta_1})\cdot r_2(\pc{\theta_2})\\
&=r_1\cdot r_2 \left[\pcx{\theta_1+\theta_2}\right]
\end{split}\]
\end{blk}

\begin{proof}
\[\begin{split}
    \because\quad z_1\cdot z_2&=r_1\cdot r_2(\pc{\theta_1})(\pc{\theta_2})\\
&=r_1\cdot r_2 \left[\left(\cos\theta_1\cdot \cos\theta_2-\sin\theta_1\cdot \sin\theta_2\right)\right.\\
&\qquad \qquad +i\left.\left(\cos\theta_1\cdot \sin\theta_2+\sin\theta_1\cdot \cos\theta_2\right)\right]\\
&=r_1\cdot r_2 \left[\pcx{\theta_1+\theta_2}\right]
\end{split}\]

$\therefore\quad $定理成立.
\end{proof}

根据这一定理，联系复数的向量表示，我们可以给出两复数相乘的几何意义：

两复数$z_1=r_1(\pc{\theta_1})$与$z_2=r_2(\pc{\theta_2})$相乘时，可以先在复平
面上画出它们对应的向量$\VEC{OZ_1}$, $\VEC{OZ_2}$，然后将$\VEC{OZ_1}$旋转一个角$|\theta_2|$（若$\theta_2>0$时，就按逆时针方向旋转$\VEC{OZ_1}$；若$\theta_2<0$时，就按顺时针方向旋转$\VEC{OZ_1}$），再把它的模变为原来的$r_2$倍，就可得到向量$\VEC{OZ}$（图1.9所示），这一向量就对应着所求两复数的积.

\begin{figure}[htp]
    \centering
\begin{tikzpicture}[>=stealth, scale=.7]
\draw[->](-1,0)--(5,0)node[right]{$x$};
\draw[->](0,-1)--(0,7)node[right]{$y$};
\node[below left]{$O$};
\coordinate(Z1) at (1.5,1.5);
\coordinate(Z2) at (2.5,1.8);
\coordinate(Z) at (1.05,6.45);
\draw[->, very thick](0,0)--node[above]{$r_1$}(Z1)node[above]{$Z_1$};
\draw[->, very thick](0,0)--node[below right]{$r_2$}(Z2)node[right]{$Z_2$};
\draw[->, very thick](0,0)--node[above right]{$r_1\cdot r_2$}(Z)node[above]{$Z$};
\draw[->](.75,0)node[above right]{$\theta_2$}arc (0:36:.75);
\draw[->](1.75,0)node[above right]{$\theta_1$}arc (0:45:1.75);
\draw[->](2.75,0)arc (0:80:2.75)node[right]{$\theta_1+\theta_2$};
\end{tikzpicture}
    \caption{}
\end{figure}



\begin{blk}
{推论} $n$个复数的积仍是一个复数，它的模等于$n$个乘数
的模的积，它的幅角等于$n$个乘数的幅角的和。（可用数学
归纳法证明）即：
\[\begin{split}
 z_{1}\cdot z_{2}\cdots z_{n}&=r_{1}\left(\cos\theta_{1}+i\sin\theta_{1}\right)\cdot r_{2}\left(\cos\theta_{2}+   i\sin\theta_{2}\right)\\
 &\qquad \qquad \cdots r_n(\pc{\theta_n})\\
 &=r_{1}\cdot r_{2}\cdots  r_{n}[\cos(\theta_1+\theta_2+\cdots+\theta_n)\\
 &\qquad \qquad \qquad +i\sin (\theta_1+\theta_2+\cdots+\theta_n)] 
\end{split}\]
\end{blk}

这样，在计算复数的乘积时，只要先化为三角形式。再求出模之积与幅角之和，就可写出所求乘积的三角形式。

在上边的推论中，如果$z_1=z_2=\cdots=z_n=z$, 即$r_1=r_{2}= \cdots = r_{n}= r$ 且 $\theta _{1}= \theta _{2}= \cdots = \theta _{n}= \theta $. 也就是说有$n$个相同的复数相乘时，那么就得到：
\[\left[r(\pc{\theta})\right]^n=r^n(\pc{n\theta})\quad (n\in\N)\]
这就是著名的棣美佛定理.

\begin{blk}
{定理2（棣美佛定理）}复数的$n$（$n\in\N$）次方幂仍是一个复数，它的模等于这个复数的模的$n$次幂，它的幅角等于这个复数的幅角的$n$倍.    
\end{blk}

\begin{example}
    计算：$\left(\pc{\frac{\pi}{6}}\right)\cdot \sqrt{2}\left(\pc{\frac{\pi}{12}}\right)$
\end{example}

\begin{solution}
由乘法定理可知：
\[\begin{split}
\text{原式}&=\sqrt{2}\left[\pcx{\frac{\pi}{6}+\frac{\pi}{12}}\right]\\
&=\sqrt{2}\left(\pc{\frac{\pi}{4}}\right)\\
&=\sqrt{2}\left(\frac{\sqrt{2}}{2}+\frac{\sqrt{2}}{2}i\right)=1+i
\end{split}\]
\end{solution}

\begin{example}
    计算$(1-i)^{10}$
\end{example}

\begin{solution}
    先将$1-i$化为三角形式：
\[1-i=\sqrt{2}\left[\pcx{-\frac{\pi}{4}}\right]\]
由棣美佛定理，得
\[\begin{split}
    (1-i)^{10}&=\left[\sqrt{2}\left(\pc{\frac{-\pi}{4}}\right)\right]^{10}\\
    &=\left(\sqrt{2}\right)^{10}\cdot \left[\pcx{-\frac{5\pi}{2}}\right]\\
    &=2^5\left(\cos\frac{\pi}{2}-i\sin\frac{\pi}{2}\right)=-32i
\end{split}\]
\end{solution}

\begin{example}
    如果将复数$z=2+3i$在复平面上所对应的向量$\VEC{OZ}$旋转$150^{\circ}$，那么，在复平面上所得新向量$\VEC{OZ_1}$将对应哪一个复数（用代数形式表示）.
\end{example}

\begin{solution}
    由复数乘法的几何意义知道，一个向量（对应一个复数）旋转一个角度$\theta<0$，就相当于这个复数与模长为1，幅角主值为$\theta$的另一复数（对应于幅角为$\theta$的单位向量）作乘积.

因此，所求的复数就等于复数$2+3i$与复数$1\cdot(\cos150^{\circ}+i\sin 150^{\circ})$的乘积，即
\[\begin{split}
    (2+3i)\cdot (\cos150^{\circ}+i\sin 150^{\circ})
    &=(2+3i)\left(-\frac{\sqrt{3}}{2}+\frac{1}{2}i\right)\\
&=-\frac{2\sqrt{3}+3}{2}+\frac{2-3\sqrt{3}}{2}i
\end{split}\]
\end{solution}

\begin{blk}
{定理3} 两个复数的商，在除数不为零的条件下仍是一个复数，它的模等于两个复数的模的商，它的幅角等于两复数的幅角的差（被除数的幅角减去除数的幅角）.即：当$z_2\ne 0$时，有
\[\begin{split}
    \frac{z_1}{z_2}&=\frac{r_1(\pc{\theta_1})}{r_2(\pc{\theta_2})}\\
    &=\frac{r_1}{r_2}\left[\pcx{\theta_1-\theta_2}\right]
\end{split}\]

\end{blk}

\begin{proof}
不妨设$z=\frac{r_1}{r_2}\left[\pcx{\theta_1-\theta_2}\right]$，则由乘法，
\[\begin{split}
    z\cdot z_2 &=\frac{r_1}{r_2}\cdot r_2\left[\pcx{\theta_1-\theta_2+\theta_2}\right]\\
    &=r_1(\pc{\theta_1})=z_1
\end{split} \]

又由除法的定义，可得
\[\frac{z_1}{z_2}=z\quad (z_2\ne 0)\]
\end{proof}

\begin{example}
计算$\left(\sqrt{3}-i\right)\div \left(\sqrt{3}+i\right)$
\end{example}

\begin{solution}
$\because\quad \sqrt{3}-i=2\left(\frac{\sqrt{3}}{2}-\frac{1}{2}i\right)=2\left[\pcx{-\frac{\pi}{6}}\right]$

$\sqrt{3}+i=2\left(\pc{\frac{\pi}{6}}\right)$

\[\begin{split}
    \therefore\quad \left(\sqrt{3}-i\right)\div \left(\sqrt{3}+i\right)&=(2\div 2)\left[\pcx{-\frac{\pi}{6}-\frac{\pi}{6}}\right]\\
&=1\cdot \left[\pcx{-\frac{\pi}{3}}\right]\\
&=\frac{1}{2}-\frac{\sqrt{3}}{2}i
\end{split}\]
\end{solution}

\begin{example}
    已知平面内并列的三个边长为1的正方形（如图1.10所示），利用复数证明：
$\angle 1+\angle 2-\angle 3=0$.
\end{example}

\begin{figure}[htp]
    \centering
\begin{tikzpicture}[>=stealth, scale=1.8]
\draw[->](-.5,0)--(4,0)node[right]{$x$};    
\draw[->](0,-.5)--(0,2)node[right]{$y$};    
\foreach \x/\y in {1/3,2/2,3/1}
{
    \node at (\x,0)[below]{$\x$};
    \draw(\x,0)--(\x,1)node[above]{$Z_{\y}$}--(0,0);
}
\draw(0,1)node[left]{1}--(3,1);
\node[below left]{$O$};

\draw(1-.25,1) arc (180:180+45:.25)node[left]{$3$};
\draw(2-.25,1) arc (180:180+26:.25)node[left]{$2$};
\draw(3-.5,1) arc (180:180+19:.5)node[left]{$1$};

\draw(.25,0)arc (0:45:.25);
\draw(.35,0)arc (0:26:.35);
\draw(.5,0)arc (0:19:.5);


\end{tikzpicture}
    \caption{}
\end{figure}

\begin{solution}
 如图建立坐标系以确定复平面，不在轴上的三个正
方形顶点设为$Z_1,Z_2,Z_3$．显然，它们分别对应着复数$3+i$, $2+i$, $1+i$; 由于平行线的内错角相等，因而$\angle 1$、$\angle 2$、$\angle 3$ 分别等于复数$3+i$, $2+i$，$1+i$的幅角的主值。

又由复数的乘、除法可知，$\angle 1+\angle 2-\angle 3$就是复数$\frac{(3+i)(2+i)}{1+i}$的幅角，而
\[\frac{(3+i)(2+i)}{1+i}=\frac{5+5i}{1+i}=5\]
它的幅角的主值是0，所以$\angle 1+\angle 2-\angle 3=0$.
\end{solution}

\begin{ex}
\begin{enumerate}
    \item 计算
\begin{enumerate}[(1)]
    \item $\sqrt{2}\left(\pc{\frac{\pi}{6}}\right)\cdot \frac{\sqrt{2}}{2}\left(\pc{\frac{\pi}{4}}\right)$
    \item $2\left(\pc{\frac{4\pi}{3}}\right)\cdot 3\left(\pc{\frac{5\pi}{6}}\right)$
    \item $\sqrt{3}(\pc{240^{\circ}})\cdot \frac{\sqrt{3}}{3}(\pc{150^{\circ}})$
    \item $4(\pc{18^{\circ}})\cdot [\pcx{-306^{\circ}}]\cdot 5(\pc{108^{\circ}})$
\end{enumerate}
\item 用棣美佛定理计算：
\begin{enumerate}[(1)]
    \item $\left[2\sqrt{2}(\pc{18^{\circ}})\right]^5$
    \item $(-1-i)^6$
    \item $\left[\sqrt{3}\left(\pc{\frac{3\pi}{4}}\right)\right]^4$
    \item $(2-i)\cdot \left(-\frac{1}{2}+\frac{\sqrt{3}}{2}i\right)^7$
\end{enumerate}
\item 计算：
\begin{enumerate}[(1)]
    \item $12\left(\pc{\frac{7\pi}{4}}\right)\div 6\left(\pc{\frac{2\pi}{3}}\right)$
    \item $2\div (\pc{45^{\circ}})$
    \item $-i\div 2\left(\pc{\frac{4\pi}{3}}\right)$
    \item $(-1+i)\cdot \sqrt{2}(\pc{120^{\circ}})\div 3(\pc{55^{\circ}})$
\end{enumerate}

\item 你能说明复数除法的几何意义吗？
\item 利用复数证明：在${\rm Rt}\triangle ABC$中，
$\angle C=90^{\circ}$, $AC=3BC$，$E$为$AC$上一点，且$CE=2EA$，则：
\[\angle CBE+\angle CBA-\frac{\pi}{4}=\frac{\pi}{2}\]
\end{enumerate}
\end{ex}

\subsection{复数的开方运算}
开方是乘方的逆运算，求一个复数$a+bi$的$n$次方根，就是求一个数$z$，使得$z^n=a+bi$. 凡满足这一等式的数$z$，就叫做复数$a+bi$的$n$次方根.

设$a+bi=r(\cos\theta+i\sin\theta)$, $z=\rho(\cos\varphi+i\sin\varphi)$，
则对于任意自然数$n$，都有
\[r(\cos\theta+i\sin\theta)=[\rho(\cos\varphi+i\sin\varphi)]^n=\rho^n(\cos n\varphi+i\sin n\varphi)\]

由于两复数相等，它们的模相等，而它们的幅角可以相差$2\pi$的整数倍，所以
\[\begin{cases}
    \rho^n=r\\
    n\varphi=\theta+2k\pi\quad (k\in\Z)
\end{cases}\]

因而可得
\[\rho=\sqrt[n]{r},\qquad \varphi=\frac{\theta+2k\pi}{n}\]

因此，复数$r(\pc{\theta})$的$n$次方根是
\[\sqrt[n]{r}\left(\pc{\frac{2k\pi+\theta}{n}}\right)\]
其中，当$k=0,1,\ldots,n-1$各值时，可以得到上式的$n$个不同值；当$k$继续取$n,n+1,\ldots$以及其他的值时，由于正弦、余弦函数都是以$2\pi$为周期，上式的值又重复出现$k$取$0,1,\ldots,n-1$时的同样结果．所以，复数的$n$次（$n\in\N$）方根是$n$个复数，它们的模都等于这个复数的模的$n$次算术根，它们的幅角分别等于这个复数的幅角与$2\pi$的$0,1,\ldots,n-1$倍的和的$n$分之一. 即
\[\sqrt[n]{r(\pc{\theta})}=\sqrt[n]{r}\left(\pc{\frac{\theta+2k\pi}{n}}\right)\quad (k=0,1,\ldots,n-1)\]

\begin{rmk}
在复数范围内，$\sqrt[n]{z}$应表示$n$个复数，这$n$个复数都是$z$的$n$次方根，为了区别于在实数范围内，$\sqrt[n]{a}$仅表示算术根，我们约定：被开方数是实数时，$\sqrt[n]{a}\; (a\in\R)$仍表示算术根（单值），$\sqrt[n]{a+0i}$或$\sqrt[n]{a(\pc{0})}\; (a>0)$或$\sqrt[n]{|a|(\pc{\pi})}$都表示$n$个$n$次方根（多值）. 至于被开方数是虚数时，$\sqrt[n]{a+bi}\; (b\ne 0)$或$\sqrt[n]{r(\pc{\theta})}$都表示$n$个方根.
\end{rmk}

\begin{example}
    求$i$的平方根.
\end{example}

\begin{solution}
$\because\quad i=\pc{\frac{\pi}{2}}$

$\therefore\quad \sqrt{i}=\pc{\frac{\frac{\pi}{2}+2k\pi}{2}}=\pcx{\frac{\pi}{4}+k\pi}\quad (k=0,1)$

因此，$i$的平方根是
\[\begin{split}
z_1=\pc{\frac{\pi}{4}}=\frac{\sqrt{2}}{2}+\frac{\sqrt{2}}{2}i\qquad &\text{（$k=0$时）}\\
z_2=\pc{\frac{5\pi}{4}}=-\frac{\sqrt{2}}{2}-\frac{\sqrt{2}}{2}i\qquad &\text{（$k=1$时）}
\end{split}\]
\end{solution}

\begin{example}
    求$-1+i$的四次方根.
\end{example}

\begin{solution}
$\because\quad -1+i=\sqrt{2}\left(\pc{\frac{3\pi}{4}}\right)$

$\therefore\quad \sqrt[4]{-1+i}=\sqrt[8]{2}\left(\pc{\frac{\frac{3\pi}{4}+2k\pi}{4}}\right)\quad (k=0,1,2,3)$

因此，$-1+i$的四次方根是：
\[\begin{split}
z_1=\sqrt[8]{2}\left(\pc{\frac{3\pi}{16}}\right)\qquad & (k=0)\\
z_2=\sqrt[8]{2}\left(\pc{\frac{11\pi}{16}}\right)\qquad & (k=1)\\
z_3=\sqrt[8]{2}\left(\pc{\frac{19\pi}{16}}\right)\qquad & (k=2)\\
z_4=\sqrt[8]{2}\left(\pc{\frac{27\pi}{16}}\right)\qquad & (k=3)\\
\end{split}\]
\end{solution}

\begin{example}
    求$-3$的平方根.
\end{example}

\begin{solution}
$\because\quad -3=3(\pc{\pi})$

$\therefore\quad \sqrt{-3+0i}=\sqrt{3}\left(\pc{\frac{\pi+2k\pi}{2}}\right)\quad (k=0,1)$

因此，$-3$的平方根有两个复数：
\[\begin{split}
z_1&=\sqrt{3}\left(\pc{\frac{\pi}{2}}\right)=\sqrt{3}i\\
z_2&=\sqrt{3}\left(\pc{\frac{3\pi}{2}}\right)=-\sqrt{3}i\\
\end{split}\]
\end{solution}

一般来说，在复数范围内，任一个负实数$-a$有两个平方根：$\pm\sqrt{a}i$.

这样一来，以前在解实系数一元二次方程$ax^2+bx+c=0$（$a\ne 0$）时，如果$b^2-4ac<0$，在实数范围内无解. 现在，在复数范围内，方程就有两个解：
\[x=\frac{-b\pm\sqrt{-(b^2-4ac)}i}{2a}\qquad (b^2-4ac<0)\]
显然，这两个解是一对共轭复数.

\begin{example}
    解方程$x^2-2x+6=0$.
\end{example}

\begin{solution}
$\because\quad \Delta=b^2-4ac=4-24=-20<0$

$\therefore\quad x=\frac{2\pm\sqrt{20}i}{2}=1\pm\sqrt{5}i$
\end{solution}

由以上可知，复数集$\mathbb{C}$对于开方运算是封闭的，这是复数系不同于实数系的一个重要性质.

利用复数开方的方法，我们可以解一些特殊的高次方程，例如，

形如$a_nx^n+a_0=0\; (a_n\ne 0)$的方程，我们叫做二项方程，就可以完整地解出，即，首先化成$x^n=b\quad \left(b=-\frac{a_0}{a_n}\right)$的形式，从而解出$x$的$n$个值.

\begin{example}
    解方程$x^5+243=0$.
\end{example}

\begin{solution}
$\because\quad x^5=-243=243(-1+0i)=243(\pc{\pi})$

$\therefore\quad x=\sqrt[5]{243}\left(\pc{\frac{\pi+2k\pi}{5}}\right)\quad (k=0,1,2,3,4)$

即
\[\begin{split}
x_1&=3\left(\pc{\frac{\pi}{5}}\right)\qquad\quad (k=0)\\
x_2&=3\left(\pc{\frac{3\pi}{5}}\right)\qquad (k=1)\\
x_3&=3\left(\pc{\pi}\right)\qquad\qquad (k=2)\\
x_4&=3\left(\pc{\frac{7\pi}{5}}\right)\qquad (k=3)\\
x_5&=3\left(\pc{\frac{9\pi}{5}}\right)\qquad (k=4)\\
\end{split}\]
\end{solution}

仔细观察这五个解可以发现：它们的模都相同，它们的幅角依次相差$\frac{2\pi}{5}$：因此，这个方程的五个解所对应的复平面内的五个点是均匀地分布在以原点为圆心，以3为半径的圆周上（图1.11）．这就是此方程的解的几何意义。
\begin{figure}[htp]
    \centering
\begin{tikzpicture}[>=stealth]
\draw[->](-2.5,0)--(2.7,0)node[right]{$x$};
\draw[->](0,-2)--(0,2.7)node[right]{$y$};
\draw(0,0)node[below left]{$O$} circle(1.5);
\coordinate (Z1) at (36:1.5);
\coordinate (Z2) at (36*3:1.5);
\coordinate (Z3) at (36*5:1.5);
\coordinate (Z4) at (36*7:1.5);
\coordinate (Z5) at (36*9:1.5);
\draw[very thick](0,0)--(Z1)node[right]{$Z_1$};
\draw[very thick](0,0)--(Z2)node[above]{$Z_2$};
\draw[very thick](0,0)--(Z3)node[above left]{$Z_3$};
\draw[very thick](0,0)--(Z4)node[below]{$Z_4$};
\draw[very thick](0,0)--(Z5)node[right]{$Z_5$};
\node at (-1.5,0)[below left]{$-3$};
\foreach \x in {1,2,3,4,5}
{
    \draw[fill](Z\x) circle(1.5pt);
}


\end{tikzpicture}
    \caption{}
\end{figure}

一般地，二项方程$x^n=b\; (b\in\mathbb{C})$有$n$个复数根，这些根的几何意义是复平面内的$n$个点（或向量），这些点（或向量的终点）均匀地分布在以原点为圆心，以$\sqrt[n]{|b|}$为半径的圆上.

还要进一步指出：对一般的实系数$n$次方程在复数集内也是恰有$n$个根，但未必都是实数根，可能会有虚根（如例1.25).

可以证明，实系数方程的虚根总是成对地出现的，即

\begin{blk}{定理4（虚根成对定理）}
如果$z=a+bi$是实系数方程
$f(x)=a_nx^n+a_{n-1}x^{n-1}+\cdots+a_1x+a_0=0$
的一个根，那么$\bar z=a-bi$也是方程$f(x)=0$的一个根.
\end{blk}

\begin{proof}
    由于$z=a+bi$是方程的根，所以
\[f(z)=f(a+bi)=0\]
但根据习题1.2第15题（2）已证明的结论，应有
\[f(\bar z)=f(a-bi)=f(\overline{a+bi})=\bar 0=0\]
因此，$\bar z=a-bi$也是方程$f(x)=0$的根.
\end{proof}

\begin{example}
    已知方程
$2x^4-x^3+5x^2+13x+5=0$
的一个根为$1-2i$，求这个方程的其余根.
\end{example}

\begin{solution}
    这是一个实系数一元四次方程，由定理4及已知一根$1-2i$可知，这一方程必有根$1+2i$.

又由因式定理知：$(x-1+2i)(x-1-2i)=x^2-2x＋5$是方程左边的因式，利用长除法可得
\[2x^4-x^3+5x^2+13x+5=(2x^2+3x+1)(x^2-2x+5)\]

所以只要解方程$2x^2+3x+1=0$的根：$x_1=-1$, $x_2=-\frac{1}{2}$，就可知：

原方程的其余三根为$1+2i,\; -1,\; -\frac{1}{2}$.
\end{solution}

\begin{ex}
\begin{enumerate}
    \item 直接说出以下各数的平方根：
\[-4,\quad -2.25,\quad -3,\quad -t\; (t\in\R^+),\quad -m^2\; (m\in\R)\]
\[a\; (a\in\R),\quad a-b\; (a,b\in\R)\]
    \item 求下列各数的方根：
\begin{enumerate}[(1)]
    \item $-i,\; \sqrt{2}-\sqrt{2}i,\; -\frac{1}{2}+\frac{\sqrt{3}}{2}i$的平方根；
    \item $27,\; 1+i$的立方根；
    \item $-i$的五次方根．
\end{enumerate}
    \item 解方程（在复数范围内）：
\begin{multicols}{2}
\begin{enumerate}[(1)]
    \item $x^2+x+16=0$
    \item $x^3-1=0$
    \item $x^4+16=0$
    \item $x^5-i=0$
\end{enumerate}
\end{multicols}
    \item 已知方程$4x^4-2x^3-x-1=0$有一个根为$-\frac{i}{\sqrt{2}}$，试求这个方程的其余三个根.
\end{enumerate}
\end{ex}


\subsection{单位根}
我们把方程$x^n=1$在复数范围内的$n$个根，叫做$n$次单位根．也就是在复数范围内，1的$n$个不同的$n$次方根，都叫$n$次单位根. 即

$\because\quad x^n=1=\pc{0}$

$\therefore\quad x=1\cdot \left(\pc{\frac{2k\pi}{n}}\right)\quad (k=0,1,\ldots,n-1)$

这样可得出$n$个单位根：
\[\begin{split}
x_1&=\pc{0}=1\qquad \qquad \text{（当$k=0$时）}\\
x_2&=\pc{\frac{2\pi}{n}}\qquad \qquad \text{（当$k=1$时）}\\
x_3&=\pc{\frac{4\pi}{n}}\qquad \qquad \text{（当$k=2$时）}\\
\cdots&\cdots\cdots\\ 
x_n&=\pc{\frac{2(n-1)\pi}{n}}\qquad \text{（当$k=n-1$时）}\\
\end{split}\]

由二项方程解的几何意义可以知道：$n$个不同的单位根可用复平面内以原点为圆心的单位圆的一个内接正$n$边形顶点来表示.

\begin{example}
    求三次单位根．
\end{example}

\begin{solution}
$\because\quad x^3=1=\pc{0}$

$\therefore\quad x=\pc{\frac{2k\pi}{3}}\quad (k=0,1,2)$

因此，三个三次单位根为(如图1.12)：
\[\begin{split}
    x_1&=\pc{0}=1\\
    x_2&=\pc{\frac{2\pi}{3}}=-\frac{1}{2}+\frac{\sqrt{3}}{2}i\\
    x_3&=\pc{\frac{4\pi}{3}}=-\frac{1}{2}-\frac{\sqrt{3}}{2}i\\
\end{split}\]
\end{solution}

\begin{minipage}{.45\textwidth}
\centering
\begin{tikzpicture}[>=stealth]
\draw[->](-2.5,0)--(2.7,0)node[right]{$x$};
\draw[->](0,-2)--(0,2.7)node[right]{$y$};
\draw(0,0)node[below left]{$O$} circle(1.5);
\draw[very thick](1.5,0)node[above right]{$x_1$}--(120:1.5)node[above]{$x_2$}--(240:1.5)node[below]{$x_3$}--cycle;

\end{tikzpicture}
\captionof{figure}{ }
\end{minipage}\hfill
\begin{minipage}{.45\textwidth}
\centering
\begin{tikzpicture}[>=stealth]
    \draw[->](-2.5,0)--(2.7,0)node[right]{$x$};
    \draw[->](0,-2)--(0,2.7)node[right]{$y$};
    \draw(0,0)node[below left]{$O$} circle(1.5);
\draw[very thick](1.5,0)node[above right]{$x_1$}--(0,1.5)node[above right]{$x_2$}--(-1.5,0)node[above left]{$x_3$}--(0,-1.5)node[below left]{$x_4$}--cycle;

\end{tikzpicture}
\captionof{figure}{ }
\end{minipage}

\begin{example}
    求四次单位根
\end{example}

\begin{solution}
$\because\quad x^4=1$

$\therefore\quad x=\pc{\frac{2k\pi}{4}}\quad (k=0,1,2,3)$

因此，四个四次单位根为(如图1.13)：
\begin{center}
\begin{tabular}{p{.2\textwidth}p{.15\textwidth}p{.2\textwidth}p{.15\textwidth}}
    $x_1=1$& $(k=0)$ & $x_2=i$ & $(k=1)$\\
    $x_3=-1$ &$ (k=2)$& $x_4=-i$ & $(k=3)$\\
\end{tabular} 
\end{center}
\end{solution}


如果用复数的向量表示法
来看这$n$个不同的$n$次单位根的话，我们会发现单位根之间还有一定的关系.

\begin{figure}[htp]
    \centering
\begin{tikzpicture}[>=stealth]
    \draw[->](-2.5,0)--(2.7,0)node[right]{$x$};
    \draw[->](0,-2)node[below]{$(n=12)$}--(0,2.7)node[right]{$y$};
    \draw(0,0)node[below left]{$O$} circle(1.5);
\foreach \x in {0,1,2,...,11}
{
    \draw[fill](\x*30:1.5) circle(1.5pt);
}
\node at (1.5,0)[below right]{1};
\node at (0:1.5)[above right]{$x_1$};
\node at (30:1.5)[above right]{$x_2$};
\node at (60:1.5)[above right]{$x_3$};
\node at (-30:1.5)[below right]{$x_n$};


\end{tikzpicture}
    \caption{}
\end{figure}




\begin{enumerate}
\item 如图1.14所示，在复平面上$n$个$n$次单位根正好将单位圆$n$等分，其中：
$x_1=1$，单位向量$\VEC{Ox_1}$旋转$\frac{2\pi}{n}$
得：$x_2=\pc{\frac{2\pi}{n}}$
记为$\omega$；再旋转$\frac{2\pi}{n}$
得：$x_3=\pc{\frac{4\pi}{n}}=\omega^2$；再旋转$\frac{2\pi}{n}$可一样地
得：$x_4=\pc{\frac{6\pi}{n}}=\omega^3$，……

同理可得：
\[x_n=\pc{\frac{2(n-1)\pi}{n}}=\omega^{n-1}\]

因此，设$\omega=\pc{\frac{2\pi}{n}}$，则$n$个$n$次单位根又可写成：
\[1,\; \omega,\; \omega^2,\; \ldots, \omega^{n-1}\]

\item 进一步考虑这$n$个单位根的关系，可以发现：

\[\begin{split}
    \because\quad \omega^{n-1}=\pc{\frac{2(n-1)\pi}{n}}&=\pcx{\frac{-2\pi}{n}}\\
    &=\cos\frac{2\pi}{n}-i\sin\frac{2\pi}{n}=\overline{\omega}
\end{split}\]

$\therefore\quad \omega^{n-1}=\overline{\omega}=\frac{1}{\omega}=\omega^{-1}$

同理可得：\[\omega^{n-2}=\overline{\omega^2}=\omega^{-2},\ldots\]

\[\begin{cases}
\text{当$n$为偶数时}& \omega^{\tfrac{n}{2}+1}=\overline{\omega^{\tfrac{n}{2}-1}}=\omega^{-\left(\tfrac{n}{2}-1\right)}=\omega^{-\tfrac{n}{2}+1},\quad \omega^{\tfrac{n}{2}}=-1\\
\text{当$n$为奇数时}& \omega^{\tfrac{n+1}{2}}=\overline{\omega^{\tfrac{n-1}{2}}}=\omega^{-\tfrac{n-1}{2}}
\end{cases}\]
所以，这$n$个$n$次单位根还可以写成：
\begin{itemize}
    \item 当$n$为偶数时（如例1.29）：
\[1,\; \omega,\; \omega^2,\;\ldots,\;  \omega^{\tfrac{n}{2}-1},\;-1,\;  \omega^{-\left(\tfrac{n}{2}-1\right)},\; \ldots, \;  \omega^{-2},\; \omega^{-1}\]
或
\[1,\; \omega,\; \omega^2,\;\ldots,\;  \omega^{\tfrac{n}{2}-1},\;-1,\;  \overline{\omega^{\tfrac{n}{2}-1}},\; \ldots, \;  \overline{\omega^{2}},\; \overline{\omega}\]

\item 当$n$为奇数时（如例1.28）：
\[1,\; \omega,\; \omega^2,\;\ldots,\;  \omega^{\tfrac{n-1}{2}},\;  \omega^{-\tfrac{n-1}{2}},\; \ldots, \;  \omega^{-2},\; \omega^{-1}\]
或
\[1,\; \omega,\; \omega^2,\;\ldots,\;  \omega^{\tfrac{n-1}{2}},\;  \overline{\omega^{\tfrac{n-1}{2}}},\; \ldots, \;  \overline{\omega^{2}},\; \overline{\omega}\]
其中，$\omega=\pc{\frac{2\pi}{n}}$.
\end{itemize}

例如，三次单位根为$1,\; \omega,\; \omega^{-1}$（或$\bar\omega$），其中
\[\omega=\pc{\frac{2\pi}{3}}=-\frac{1}{2}+\frac{\sqrt{3}}{2}i\]

又如，四次单位根为$1,\; \omega,\; -1,\; \omega^{-1}$（或$\bar\omega$），其中
\[\omega=\pc{\frac{2\pi}{4}}=i\]

\item 利用单位根，可以方便地求出任一个正实数$a$（$a>0$且$a\ne 1$）的$n$个$n$次方根．

这就是：若$a>0$且$a\ne 1$，则$a$的$n$个$n$次方根即方程$x^n=a$的解为：
\[\begin{cases}
x_1=\sqrt[n]{a}\\
x_2=\sqrt[n]{a}\cdot \omega\\
x_3=\sqrt[n]{a}\cdot \omega^2\\
\cdots\cdots\cdots\\
x_n=\sqrt[n]{a}\cdot \omega^{n-1}\\
\end{cases}\qquad \text{这里$\omega=\pc{\frac{2\pi}{n}}$是单位根之一}\]
\end{enumerate}

\begin{ex}
试求出六次和七次单位根，并分别检验它们之间的关系.
\end{ex}

\subsection*{习题1.3}
\begin{enumerate}
    \item 把下列复数化成三角形式表示：
\begin{multicols}{3}
\begin{enumerate}[(1)]
    \item 6
    \item $6i$
    \item $1+i$
    \item $\frac{1}{2}-\frac{\sqrt{3}}{2}i$
    \item $-\sqrt{2}-\sqrt{2}i$
\end{enumerate}
\end{multicols}

\item 把下列复数化为代数形式表示，并画出复平面上相应的向量来：
\begin{multicols}{2}
\begin{enumerate}[(1)]
    \item $3\sqrt{2}\left(\pc{\frac{\pi}{4}}\right)$
    \item $2\left(\pc{\frac{11\pi}{6}}\right)$
    \item $\pc{\pi}$
    \item $\sqrt{3}\left[\pcx{-\frac{2\pi}{3}}\right]$
    \item $\pcx{k\cdot \frac{\pi}{4}}$
    
    $(k=0,1,2,\ldots,6,7)$
\end{enumerate}
\end{multicols}

\item 计算
\begin{enumerate}[(1)]
    \item $\sqrt{3}\left(\pc{\frac{\pi}{3}}\right)\cdot 2\sqrt{3}\left(\pc{\frac{\pi}{6}}\right)$
    \item $\sqrt{10}\left(\pc{\frac{\pi}{2}}\right)\cdot \sqrt{2}\left(\pc{\frac{\pi}{4}}\right)$
\end{enumerate}

\item 下列复数是不是三角形式？如果不是，请化成三角形式表示：
\begin{multicols}{2}
\begin{enumerate}[(1)]
    \item $\cos\frac{\pi}{6}-i\sin\frac{\pi}{6}$
    \item $\sqrt{2}\left(\cos\frac{4\pi}{3}+i\sin\frac{2\pi}{3}\right)$
    \item $7\left(\sin\frac{3\pi}{4}+i\cos\frac{3\pi}{4}\right)$
    \item $3\left(\cos\frac{5\pi}{6}+i\sin\frac{-7\pi}{6}\right)$
\end{enumerate}
\end{multicols}

\item 求证：
\begin{enumerate}[(1)]
    \item $\sqrt{2}(\pc{75^{\circ}})\cdot \sqrt{2}\left(\pc{15^{\circ}}\right)=2i$
    \item $(\cos3\theta-i\sin3\theta)(\cos2\theta-i\sin2\theta)=\cos5\theta-i\sin5\theta$
\end{enumerate}

\item 如图1.15所示，菱形$OABC$的一个内角$\angle AOC=\frac{\pi}{4}$，且$A$点所对应的复数为$2+i$, 试求，$B$、$C$两点各对应的复数.

\begin{figure}[htp]
    \centering
\begin{tikzpicture}[>=stealth]
    \draw[->](-.75,0)--(3.4,0)node[right]{$x$};
    \draw[->](0,-.75)--(0,3.4)node[right]{$y$};
    \node[below left]{$O$} ;
    \coordinate (A) at (30:2);
    \coordinate (C) at (75:2);
    \coordinate (B) at (105/2:3.7);
    \draw[very thick](0,0)--(A)node[right]{$A$}--(B)node[above]{$B$}--(C)node[above]{$C$}--cycle;
\end{tikzpicture}
    \caption{}
\end{figure}


\item 已知下列各复数$z$，求$\left|\frac{1}{z}\right|$及$\arg\left(\frac{1}{z}\right)$
\begin{multicols}{2}
\begin{enumerate}[(1)]
    \item $z=2\left(\pc{\frac{\pi}{12}}\right)$
    \item $z=\cos\frac{\pi}{6}-i\sin\frac{\pi}{6}$
    \item $z=\frac{1-i}{\sqrt{2}}$
\end{enumerate}    
\end{multicols}

\item 把复数$3-\sqrt{3}i$对应的向量旋转$\frac{\pi}{3}$，所得向量相对应的复数是什么？

\item 用棣美佛定理计算：
\begin{multicols}{2}
\begin{enumerate}[(1)]
    \item $[2(\pc{15^{\circ}})]^6$
    \item $\left[\frac{1}{\sqrt{2}}(\pc{225^{\circ}})\right]^4$
    \item $(1+\sqrt{3}i)^4$
    \item $(2-2\sqrt{3}i)^4$
    \item $\frac{(\sqrt{3}+i)^5}{-1+\sqrt{3}i}$
    \item $\left(\frac{2+2i}{1-\sqrt{3}i}\right)^8$
\end{enumerate}
\end{multicols}

\item 用棣美佛定理证明：
\begin{enumerate}[(1)]
    \item $\cos2\theta=\cos^2\theta-\sin^2\theta,\qquad \sin2\theta=2\sin\theta\cdot \cos\theta$
    \item $\cos3\theta=4\cos^3\theta-3\cos\theta,\qquad \sin3\theta=3\sin\theta-4\sin^3\theta$
\end{enumerate}

\item 设$n\in\N$，$n$取何值时，$(1+\sqrt{3}i)^n$是一个实数？

\item 计算：
\begin{enumerate}[(1)]
    \item $10\left(\pc{\frac{2\pi}{3}}\right)\div 2\left(\pc{\frac{\pi}{6}}\right)$
    \item $-10i\div 6\left(\pc{\frac{\pi}{6}}\right)$
\end{enumerate}

\item 求证：
\begin{enumerate}[(1)]
    \item $\frac{1}{\cos\theta+i\sin\theta}=\cos\theta-i\sin\theta$
    \item $(\pc{\theta})^{-n}=\pcx{-n\theta}\quad (n\in\mathbb{N})$
    \item $(\cos\theta-i\sin\theta)^n=\cos n\theta-i\sin n \theta\quad (n\in\N)$
\end{enumerate}

\item 利用复数证明余弦定理.

\item 化简：
\begin{enumerate}[(1)]
    \item $\frac{(\pc{7\theta})(\pc{2\theta})}{(\pc{5\theta})(\pc{3\theta})}$
    \item $\frac{\cos\varphi-i\sin\varphi}{\pc{\varphi}}$
    \item $\frac{\pc{\alpha}}{\sin\alpha+i\cos\alpha}$
\end{enumerate}
\item 在复数范围内解下列方程：
\begin{multicols}{2}
\begin{enumerate}[(1)]
    \item $4x^2+9=0$
    \item $2(x^2+4)=5x$
    \item $(x-3)(x-5)=-2$
    \item $\frac{1}{x+3}-\frac{1}{x}=1$
    \item $x^4+3x^2+1=0$
    \item $\frac{x}{x^2+1}+\frac{x^2+1}{x}-\frac{5}{2}=0$
\end{enumerate}
\end{multicols}
\item 解下列方程组：
\begin{multicols}{2}
\begin{enumerate}[(1)]
    \item $\begin{cases}
a+b=2\\ab=2        
    \end{cases}$
    \item $\begin{cases}
        x^2+y^2=0\\
        xy=1
    \end{cases}$
\end{enumerate}
\end{multicols}

\item 求：
\begin{enumerate}[(1)]
    \item $8\left(\pc{\frac{\pi}{3}}\right)$的六次方根;
    \item $-2i$的五次方根.
\end{enumerate}
\item 求：六次单位根、七次单位根，并在复平面上用向量表示出来.

\item 解下列方程：
\begin{multicols}{2}
\begin{enumerate}[(1)]
    \item $x^3+1=i$
    \item $x^4+16=0$
    \item $x^{10}-32x^5+1024=0$
    \item $x^{12}+63x^6-64=0$
\end{enumerate}
\end{multicols}

\staritem 求证：
\begin{enumerate}[(1)]
    \item 不经计算证明$\left|\frac{a+bi}{a-bi}\right|\equiv 1$
    \item 有相同幅角的两个复数之比必为实数；
    \item 复平面上四个点$Z_1$, $Z_2$, $Z_3$, $Z_4$共圆或共直线的充分必要条件是：
\[\left.\frac{Z_3-Z_1}{Z_3-Z_2}\right/\frac{Z_4-Z_1}{Z_4-Z_2} \text{ 为实数}\]
\end{enumerate}
\end{enumerate}

\section{复数的指数形式}
我们已经学习了复数的代数形式、三角形式以及它们的运算.在科学技术，特别是在电工和无线电计算中，常要涉及复数或三角函数的变换和计算，为方便还采用复数的另一种形式——复数的指数形式，在这一节将学习复数的指数形式及其运算、应用.

\subsection{复数的指数形式及其运算}
我们首先引进一个著名的公式——欧拉公式：
\[e^{i\theta}=\pc{\theta}\]
其中$e=2.71828\cdots$是一个重要的常数，也就是自然对数的底数.

欧拉公式在以后的“复变函数论”中可以证明。这个公式表明了：模为1、幅角为$\theta$的复数$\pc{\theta}$与以$e$为底的复指数函数的关系，从这个公式出发，对任一个复数
\[z=a+bi=r(\pc{\theta})\]
就可以表示成$z=r\cdot e^{i\theta}$的形式，我们就把这一表达式叫做复数的指数形式.

例如：
\begin{itemize}
    \item 复数$\pc{\frac{\pi}{2}}$的指数形式为$e^{i\tfrac{\pi}{2}}$
    \item 复数$3\left(\pc{\frac{\pi}{4}}\right)$的指数形式为$3e^{i\tfrac{\pi}{4}}$
\end{itemize}

这样，同一个复数，就有三种不同形式的表达方法——代数形式、三角形式、指数形式.
如：
\[\begin{split}
    \frac{1}{2}+\frac{\sqrt{3}}{2}i&=\pc{\frac{\pi}{3}}=e^{i\tfrac{\pi}{3}}\\
    -2\sqrt{3}+2i&=4\left(\pc{\frac{5\pi}{6}}\right)=4e^{i\tfrac{5\pi}{6}}
\end{split}\]

复数的指数形式便利于乘、除、乘方及开方运算.因为可以证明原有在实数范围内的指数运算律它都仍然满足. 这就是：

设$z_1=r_1e^{i\theta_1}$，$z_2=r_2e^{i\theta_2}$，则有
\[\begin{split}
    z_1\cdot z_2&=\left(r_1e^{i\theta_1}\right)\cdot \left(r_2e^{i\theta_2}\right)\\
    &=r_1(\pc{\theta_1})\cdot r_2(\pc{\theta_2})\\
    &=r_1r_2[\pcx{\theta_1+\theta_2}]=r_1r_2\cdot e^{i(\theta_1+\theta_2)}
\end{split}\]
即
\[r_1e^{i\theta_1}\cdot r_2e^{i\theta_2}=r_1r_2 e^{i(\theta_1+\theta_2)}\]

同理可证：
\[\begin{split}
    \frac{z_1}{z_2}&=\frac{r_1e^{i\theta_1}}{r_2e^{i\theta_2}}=\frac{r_1}{r_2}\cdot e^{i(\theta_1-\theta_2)}\quad (z_2\ne 0)\\
z^n_1&=(r_1e^{i\theta_1})^n=r^n_1 \cdot e^{i n\theta_1}\quad (n\in\N)
\end{split} \]
$z_1$的$n$次方根：$\sqrt[n]{z_1}=\sqrt[n]{r_1e^{i\theta_1}}=\sqrt[n]{r_1}e^{i\tfrac{\theta_1+2k\pi}{n}}\quad (k=0,1,\ldots,n-1)$

\begin{example}
把复数$z=2i$化成指数形式.
\end{example}

\begin{solution}
$\because\quad z=2i=2\left(\pc{\frac{\pi}{2}}\right)$

$\therefore\quad z=2\cdot e^{i\tfrac{\pi}{2}}$
\end{solution}

\begin{example}
已知$z_1=\sqrt{2}\cdot e^{i\left(-\tfrac{\pi}{4}\right)}$，$z_2=\sqrt{5}\cdot e^{i\tfrac{2\pi}{3}}$. 试求：
\begin{multicols}{2}
\begin{enumerate}[(1)]
    \item $z_1\cdot z_2$
    \item $z_1$的平方根
\end{enumerate}
\end{multicols}
要把结果表示成复数的三角形式.
\end{example}

\begin{solution}
\begin{enumerate}[(1)]
    \item \[\begin{split}
z_1\cdot z_2&=\left(\sqrt{2} e^{i\left(-\tfrac{\pi}{4}\right)}\right)\cdot \left(\sqrt{5} e^{i\tfrac{2\pi}{3}}\right)\\
&=\sqrt{10}e^{i\tfrac{5\pi}{12}}=\sqrt{10}\left(\pc{\frac{5\pi}{12}}\right)
    \end{split}\]
\item $\sqrt{\sqrt{2}\cdot e^{i\left(-\tfrac{\pi}{4}\right)}}=\sqrt[4]{2}\cdot e^{i\tfrac{-\tfrac{\pi}{4}+2k\pi}{2}}\quad (k=0,1)$

$\therefore\quad z_1$的两个平方根为
\[\begin{split}
x_1&=\sqrt[4]{2}\cdot e^{i\left(-\tfrac{\pi}{8}\right)}=\sqrt[4]{2}\left(\pc{\frac{15\pi}{8}}\right)\\
x_2&=\sqrt[4]{2}\cdot e^{i\tfrac{7\pi}{8}}=\sqrt[4]{2}\left(\pc{\frac{7\pi}{8}}\right)
\end{split}\]
\end{enumerate}
\end{solution}



\begin{example}
试用$e^{i\theta}$与$e^{-i\theta}$表示$\cos\theta$与$\sin\theta$.    
\end{example}

\begin{solution}
由于
\begin{align}
e^{i\theta}&= \pc{\theta} \tag{1}\\
e^{-i\theta}&= \pcx{-\theta}=\cos\theta-i\sin\theta \tag{2}
\end{align}
(1)(2)联立，可解出
\[\begin{split}
    \cos\theta&=\frac{1}{2}\left(e^{i\theta}+e^{-i\theta}\right)\\
    \sin\theta&=\frac{1}{2i}\left(e^{i\theta}-e^{-i\theta}\right)\\
\end{split}\]
\end{solution}

\begin{ex}
\begin{enumerate}
    \item 将复数$-1$，$2+2i$，$\pc{15^{\circ}}$，$4i$分别表示成指数形式.
    \item 将复数$e^{-i\tfrac{\pi}{2}}$，$3e^{2i}$，$\sqrt{2}e^{i\tfrac{3\pi}{2}}$分别化为三角形式和代数形式.
    \item 用指数形式计算：
\begin{enumerate}[(1)]
    \item $8\left(\pc{\frac{7\pi}{6}}\right)\cdot 2\left(\pc{\frac{\pi}{4}}\right)$
    \item $2\div e^{i\tfrac{\pi}{4}}$
    \item $(1+\sqrt{3}i)^{10}$
    \item 求81的四次方根
\end{enumerate}
\end{enumerate}
\end{ex}

\subsection{证明三角恒等式}

在我们学习过的三角公式和三角恒等式当中，除了公式$\sin(\alpha+\beta)$与$\cos(\alpha+\beta)$是证明棣美佛定理的依据外，其余公式都可以应用复数的指数形式来加以证明.其要点是：首先将公式中的正切、余切、正割、余割诸函数都化为正弦、余弦函数表达；其次再用由复数的指数形式而导出来的公式：（例1.32）
\[\begin{split}
\sin\theta&=\frac{1}{2i}\left(e^{i\theta}-e^{-i\theta}\right)\\
    \cos\theta&=\frac{1}{2}\left(e^{i\theta}+e^{-i\theta}\right)\\
    \end{split}\]
将三角函数转化为指数形式的代数运算；从而达到证明的目的.

\begin{example}
求证：$-2\sin\alpha\cdot \sin\beta=\cos(\alpha+\beta)-\cos(\alpha-\beta)$    
\end{example}

\begin{proof}
$\because\quad \sin\alpha=\frac{1}{2i}\left(e^{i\alpha}-e^{-i\alpha}\right),\quad \sin\beta=\frac{1}{2i}\left(e^{i\beta}-e^{-i\beta}\right)$

\[\begin{split}
\therefore\quad \text{左}&=-2\cdot \frac{1}{4i^2}\left[e^{i(\alpha+\beta)}+e^{-i(\alpha+\beta)}-e^{i(\alpha-\beta)}-e^{i(\beta-\alpha)}\right]\\
&=\frac{1}{2}\left[e^{i(\alpha+\beta)}+e^{-i(\alpha+\beta)}\right]-\frac{1}{2}\left[e^{i(\alpha-\beta)}+e^{-i(\alpha-\beta)}\right]\\
&=\cos(\alpha+\beta)-\cos(\alpha-\beta)=\text{右}
\end{split}\]

$\therefore\quad $原等式成立.
\end{proof}

\begin{example}
求证：$\sin(\alpha+\beta)\cdot \cos(\alpha-\beta)=\sin\alpha\cdot \cos\alpha+\sin\beta\cdot \cos\beta$
\end{example}

\begin{proof}
\[\begin{split}
\text{左式}&=\frac{1}{2i}\left[e^{i(\alpha+\beta)}-e^{-i(\alpha+\beta)}\right]\cdot \frac{1}{2}\left[e^{i(\alpha-\beta)}+e^{-i(\alpha-\beta)}\right]\\
&=\frac{1}{4i}\left[e^{2\alpha i}+e^{2\beta i}-e^{-2\beta i}-e^{-2\alpha i}\right]
\end{split}\]
\[\begin{split}
    \text{右式}&=\frac{1}{2i}\left(e^{i\alpha}-e^{-i\alpha}\right)\cdot\frac{1}{2}\left(e^{i\alpha}+e^{-i\alpha}\right)+\frac{1}{2i}\left(e^{i\beta}-e^{-i\beta}\right)\cdot \frac{1}{2}\left(e^{i\beta}+e^{-i\beta}\right)\\
    &=\frac{1}{4i}\left[e^{2\alpha i}-e^{-2\alpha i}+e^{2\beta i}-e^{-2\beta i}\right]
\end{split}\]

$\because\quad $左式$=$右式

$\therefore\quad $原等式成立.
\end{proof}

对于具有条件$\alpha+\beta+\gamma=\pi$的三角恒等式，应用复数的指数形式加以证明更为有利.

\begin{example}
在$\triangle ABC$中，求证：
\[1+\cos 2A+\cos 2B+\cos 2C+4\cos A\cdot \cos B\cdot \cos C=0\]
\end{example}

\begin{proof}
\[\begin{split}
 \because\quad    4\cos A\cdot \cos B\cdot \cos C &=4\cdot\frac{e^{iA}+e^{-iA}}{2}\cdot\frac{e^{iB}+e^{-iB}}{2}\cdot\frac{e^{iC}+e^{-iC}}{2}\\
 &=\frac{1}{2}\left[e^{i(A+B+C)}+e^{i(B+C-A)}+e^{i(A+C-B)}\right.\\
 &\qquad   +e^{i(A+B-C)}+e^{-i(B+C-A)}+e^{-i(A+C-B)}\\
&\qquad  \left.+e^{-i(A+B-C)}+e^{-i(A+B+C)}\right]\qquad  \qquad  (*)
\end{split}\]

但由于$A+B+C=\pi$，进而有$B+C-A=\pi-2A$, 
$A+C-B=\pi-2B$, $A+B-C=\pi-2C$; 因此可得：
\[\begin{split}
    e^{i(A+B+C)}&=e^{i\pi}=-1\\
    e^{i(B+C-A)}&=e^{i(\pi-2A)}=e^{i\pi}\cdot e^{-2Ai}=-e^{-2Ai}\\
\end{split}\]
同理：$e^{i(A+C-B)}=-e^{-2Bi},\qquad e^{i(A+B-C)}=-e^{-2Ci}$
\[\begin{split}
    e^{-i(A+B+C)}=e^{-i\pi}=-1,\qquad 
    e^{-i(B+C-A)}=-e^{2Ai}\\
    e^{-i(A+C-B)}=-e^{2Bi},\qquad 
    e^{-i(A+B-C)}=-e^{2Ci}
\end{split}\]

将以上结果代入(*)式，得
\[\begin{split}
    4\cos A\cdot \cos B\cdot \cos C &=\frac{1}{2}\left[-2-e^{-2Ai}-e^{2Ai}-e^{-2Bi}-e^{2Bi}-e^{-2Ci}-e^{2Ci}\right]\\
    &=-1-\cos2A-\cos2B-\cos2C
\end{split}\]

$\therefore\quad 1+\cos 2A+\cos 2B+\cos 2C+4\cos A  \cos B  \cos C=0$
\end{proof}

\begin{ex}
利用复数的指数形式，证明以下等式：
\begin{enumerate}
    \item $\sin2\alpha=2\sin\alpha\cdot \cos\alpha$
    \item $\cos^2\beta-\sin^2\alpha=\cos(\alpha+\beta)\cdot \cos(\alpha-\beta)$
\end{enumerate}
\end{ex}

\subsection*{习题1.4}
\begin{enumerate}
    \item 把下列复数化为指数形式：
\begin{multicols}{3}
\begin{enumerate}[(1)]
    \item $5$
    \item $2+2i$
    \item $-2i$
    \item $1-\sqrt{3}i$
    \item $\pc{\frac{\pi}{4}}$
    \item $\pc{3}$
    \item $1+\pc{\frac{\pi}{3}}$
\end{enumerate}
\end{multicols}
\item 设复数$a+bi=r\cdot e^{i\theta}$，写出复数$a-bi$, $-a+bi$, $-a-bi$的指数形式.
\item 把下列复数化为三角形式及代数形式：
\begin{multicols}{3}
\begin{enumerate}[(1)]
    \item $4e^{i\tfrac{\pi}{6}}$
    \item $\sqrt{3}e^{i\tfrac{2\pi}{3}}$
    \item $e^{-i\tfrac{3\pi}{2}}$
\end{enumerate}
\end{multicols}

\item 试求：复数$e^{i\tfrac{4\pi}{5}}$, $e^{i\tfrac{2\pi}{3}}$ 所对应的向量间的夹角$\alpha\; (0\le \alpha\le \pi)$.
\item 用复数的指数形式计算：
\begin{enumerate}[(1)]
    \item $\sqrt{2}\left(\pc{\frac{4\pi}{3}}\right)\cdot\frac{\sqrt{2}}{4}\left(\pc{\frac{5\pi}{6}}\right)$
    \item $\sqrt{3}(\pc{150^{\circ}})\div \sqrt{2}(\pc{225^{\circ}})$
\end{enumerate}

\item 根据欧拉公式求证：$e^{-i\theta}=(\pc{\theta})^{-1}$.

\item 用复数的指数形式计算：
\begin{multicols}{3}
\begin{enumerate}[(1)]
    \item $(1+\sqrt{3}i)^{10}$
    \item $(\sqrt{3}-i)^8$
    \item $\sqrt[4]{64+0i}$
\end{enumerate}
\end{multicols}

\item 用复数的指数形式证明恒等式：
\begin{enumerate}[(1)]
    \item $\sin\alpha-\sin\beta=2\cos\frac{\alpha+\beta}{2}\cdot \sin\frac{\alpha-\beta}{2}$
    \item $(\sin\alpha+\sin\beta)^2+(\cos\alpha+\cos\beta)^2=4\cos^2\frac{\alpha-\beta}{2}$
\end{enumerate}
\end{enumerate}

\section*{本章内容要点}

一、本章里，我们引进了虚数单位($i^2=-1$)，把数的概念扩展到了复数集. 复数的分类如下：
\begin{center}
\begin{tikzpicture}
\node (A) at (0,0)[text width=2cm, align=center]{复数$a+bi$\\$(a,b\in\R)$};
\node (B1) at (2.5,1)[text width=2cm, align=center]{实数\\($b=0$时)}; 
\node (B2) at (2.5,-1)[text width=2cm, align=center]{虚数\\($b\ne 0$时)}; 
\node (C1) at (5,0)[text width=2cm, align=center]{纯虚数\\($a=0$时)}; 
\node (C2) at (5,-2)[text width=2cm, align=center]{非纯虚数\\($a\ne 0$时)}; 
\draw[decorate, decoration={brace, amplitude=5pt}](1.5,-1.3)--(1.5,1.3);
\draw[decorate, decoration={brace, amplitude=5pt}](3.8,-2.3)--(3.8,.3);


\end{tikzpicture}
\end{center}

且复数集$\supset$实数集$\supset$有理数集$\supset$整数集$\supset$自然数集，即：
\[\mathbb{C} \supset \R \supset \mathbb{Q} \supset \Z\supset\N\]

二、复数有四个基本概念：模，幅角，共轭复数，复数的相等；任一个复数有三种表示形式：代数形式$a+bi\; (a,b\in\R)$，三角形式$r(\pc{\theta})$，其中$r=\sqrt{a^2+b^2}$, $\tan\theta =\frac{b}{a}$，指数形式$re^{i\theta}$；复数$z=a+bi$与复平
面内的点$Z(a,b)$对应，并且和从原点出发的向量$\VEC{OZ}$相对应，因此，复数集$\mathbb{C}$与复平面内的点集，与从原点出发的向量集合之间成一一对应关系，即
\[\mathbb{C}:\; \{a+bi\mid a,b\in\R,\; i^2=-1\} \longleftrightarrow \{Z(a,b)\}\longleftrightarrow \{\VEC{OZ}\}\]

据此，一些平面几何，平面解析几何以及平面向量的问题，也可以通过复数来解决.

三、复数集对于四则运算及乘方，开方运算都是封闭的；复数集保持了数系运算的“通性”，即加法、乘法的交换律、结合律以及乘法对于加法的分配律、数0与1的运算特性等；但复数集不再具有顺序性，对任意两个不全为实数的复数，就不能比较大小.

四、复数的运算法则及其几何意义：
\begin{enumerate}
    \item 加、减法运算用复数的代数形式表达比较方便，即
    \[(a+bi)\pm (c+di)=(a\pm c)+(b\pm d)i\]
    其几何意义是对应的向量相加、减，即“平行四边形”法则和“三角形”法则，特别是两复数的差的模表示复平面内两点间的距离.
    \item 乘、除法，乘方、开方的运算用复数的三角形式或指数形式表达比较方便，即
\[\begin{split}
    r_1(\pc{\theta_1})\cdot r_2(\pc{\theta_2})&=r_1\cdot r_2\left[\pcx{\theta_1+\theta_2}\right]\\
\frac{r_1(\pc{\theta_1})}{r_2(\pc{\theta_2})}&=\frac{r_1}{r_2}\left[\pcx{\theta_1-\theta_2}\right]
\end{split}\]
棣美佛定理：
\[[r(\pc{\theta})]^n=r^n(\pc{n\theta})\]
$n$个$n$次方根：
\[\begin{split}
\sqrt[n]{r(\pc{\theta})}&=\sqrt[n]{r}\left(\pc{\frac{\theta+2k\pi}{n}}\right)\\
&\qquad (k=0,1,2,\ldots,n-1)\quad (n\in\N)    
\end{split} \]
\end{enumerate}

复数的这些运算，都有其几何意义.特别是：复数的相乘对应着平面向量的旋转；$n$个$n$次方根对应着圆内接正$n$边形的顶点，即圆周的$n$等分点，这些结论在应用上很重要.

五、在复数集中，本章学习了在实数集中所不能解的有关问题：
\begin{enumerate}
    \item 当$b^2-4ac<0$时，一元二次方程$ax^2+bx+c=0\; (a\ne 0)$有两个共轭的虚根，即
\[x=\frac{-b\pm i\sqrt{4ac-b^2}}{2a}\]
\item 二项方程$x^n=b\; (b\in\mathbb{C})$有$n$个不同的复数根，即设$x^n=b=r(\pc{\theta})$，则
\[x=\sqrt[n]{r}\left(\pc{\frac{\theta+2k\pi}{n}}\right)\]
当$k=0,1,2,\ldots,n-1$时，可得出$n$个不同的解.

\item 特别地，方程$x^n=1$的$n$个复数解，叫做$n$次单位根，若设$\omega=\pc{\frac{2\pi}{n}}$，则当$n$为偶数时，$n$个单位根为
\[1,\; \omega,\; \omega^2,\;\ldots,\;  \omega^{\tfrac{n}{2}-1},\;-1,\;  \overline{\omega^{\tfrac{n}{2}-1}},\; \ldots, \;  \overline{\omega^{2}},\; \overline{\omega}\]
当$n$为奇数时，$n$个单位根为
\[1,\; \omega,\; \omega^2,\;\ldots,\;  \omega^{\tfrac{n-1}{2}},\;  \overline{\omega^{\tfrac{n-1}{2}}},\; \ldots, \;  \overline{\omega^{2}},\; \overline{\omega}\]

\item 实系数一元方程，若有虚根存在，则一定是成对出现，即实系数方程若有根$a+bi\;(b\ne 0)$则必有根$a-bi$.
\end{enumerate}

\section*{复习题一}
\begin{enumerate}
    \item 在复平面上任选一点$Z$（不在原点）表示复数$z$，然后用几何法在平面找出下列各复数的点：
\begin{multicols}{4}
\begin{enumerate}[(1)]
    \item $-z$
    \item $\bar z$
    \item $-\bar z$
    \item $z+\bar z$
    \item $z-\bar z$
    \item $z+|z|$
    \item $z-|z|$
    \item $iz$
    \item $-iz$
    \item $i\bar z$
    \item $-i\bar z$
\end{enumerate}
\end{multicols}

\item 已知$z=a+bi\; (a,b\in\R)$，求下列各复数的实部与虚部：
\begin{multicols}{2}
\begin{enumerate}[(1)]
    \item $z^2$
    \item $z^3$
    \item $\frac{1}{z}$
    \item $v_0 z+\frac{M}{2\pi}\cdot \frac{1}{z}\quad (v_0,M\in\R)$
\end{enumerate}
\end{multicols}

\item 求复数$z=\frac{a^2-b^2+2abi}{ab\sqrt{2}+\sqrt{a^4+b^4}i}$的模.

\item 求证：表示复数$z_0$的点关于直线${\rm Re}(z)={\rm Im}(z)$的对称点表示的是复数$i\bar z_0$.

\item 下列方程（$t$为实参数）给出了怎样的曲线？
\begin{multicols}{2}
\begin{enumerate}[(1)]
    \item $z=t(1+i)$
    \item $z=a\cos t+ib\sin t$
    \item $z=t+\frac{i}{t}$
    \item $z=t^2+\frac{i}{t}$
\end{enumerate}
\end{multicols}

\item 已知$(x+yi)^3=a+bi$，$a,b,x,y\in\R$，求证：
$\frac{a}{x}+\frac{b}{y}=4(x^2-y^2)$

\item 若$a$，$b$是共轭复数，且$(a+b)^2-3abi=4-6i$, 试求$a$及$b$.
\item 若$x,y\in\R$, $z=x+y-30-xyi$且$\bar z=60i-|x+yi|$，试求$x$和$y$的值.
\item 若$z=x+yi$，求证：$|z|\le |x|+|y|\le \sqrt{2}|z|$
\item 试判断以下复数，哪些是实数？哪些是纯虚数？($z,z_1,z_2\in\mathbb{C}$)
\begin{multicols}{2}
\begin{enumerate}[(1)]
    \item $z+\bar z$
    \item $z-\bar z$
    \item $z\cdot \bar z$
    \item $z^2-\bar z^2$
    \item $\frac{\left(\frac{1}{z}+\frac{1}{\bar z}\right)(z+\bar z)}{z-\bar z}$
    \item $z_1\cdot \bar z_2-\bar z_1\cdot z_2$
    \item $z_1\cdot \bar z_2+\bar z_1\cdot z_2$
    \item $\frac{z_1\cdot \bar z_2+\bar z_1\cdot z_2}{z_1\cdot \bar z_1-1}$
    \item $\frac{z_1\cdot \bar z_2-\bar z_1\cdot z_2}{i(z_1\cdot \bar z_1+z_2\cdot \bar z_2)}$
    \item $\frac{i(z_1\cdot \bar z_2+z_2\cdot \bar z_1)}{z_1\cdot \bar z_2-\bar z_1\cdot z_2}$
\end{enumerate}
\end{multicols}

\item 设复数$z_1,z_2$分别对应于点$Z_1,Z_2$，求证：
\begin{enumerate}[(1)]
    \item $OZ_1\bot OZ_2$的充要条件是$z_1\cdot \bar z_2+\bar z_1\cdot z_2=0$
    \item $OZ_1$与$OZ_2$共线的充要条件是$z_1\cdot \bar z_2-\bar z_1\cdot z_2=0$
\end{enumerate}

\item 设$\alpha$是复常数，$k$是实常数，$z\in\mathbb{C}$，试说明下述方程在复平面上的图形是什么？
\begin{multicols}{2}
\begin{enumerate}[(1)]
    \item $\alpha\bar z+\bar\alpha\cdot z=k$;
    \item $z\bar z+z\bar\alpha+\alpha\bar z+k=0$.
\end{enumerate}
\end{multicols}
\item 求证：如果$\frac{1}{z}+z\in\R$，那么${\rm Im}(z)=0$或$|z|=1$.
\item 化简：$\left(\frac{1-i}{1+i}\right)^{4\alpha}$，其中$\alpha=\csc10^{\circ}-\sqrt{3}\sec10^{\circ}$.
\item 解方程组
\begin{multicols}{2}
\begin{enumerate}[(1)]
    \item $\begin{cases}
        x+iy=2+4i\\ ix+y=0
    \end{cases}$
    \item $\begin{cases}
        z_1+2z_2=1+i\\  3z_1+iz_2=2-3i
    \end{cases}$
    \item $\begin{cases}
x^2+y^2=6\\        
y^2+z^2=0\\
x^2+z^2=8i\\       
    \end{cases}$
\end{enumerate}
\end{multicols}

\item 化简
\[\frac{(\cos2\theta-i\sin2\theta)(\pc{\varphi})^2}{\pcx{\theta+\varphi}}\cdot \frac{(\pc{2\theta})^2(\cos2\varphi-i\sin2\varphi)}{\pcx{\theta-\varphi}}\]

\item 将下列复数化为三角形式：
\begin{multicols}{2}
\begin{enumerate}[(1)]
    \item $\frac{(\pc{5\varphi})^2}{(\cos3\varphi-i\sin3\varphi)^3}$
    \item $1-\cos\theta+i\sin\theta,\quad (0\le\theta\le\pi)$
\end{enumerate}
\end{multicols}

\item 设$x=e^{i\alpha}$, $y=e^{i\beta}$, $z=e^{i\gamma}$，求证：
\[(x+y)(y+z)(z+x)=8xyz\cos\frac{\beta-\gamma}{2}\cdot \cos\frac{\gamma-\alpha}{2}\cdot \cos\frac{\alpha-\beta}{2}\]

\item 设$\sin A+\sin B+\sin C=\cos A+\cos B+\cos C=0$，求证：
\begin{enumerate}[(1)]
    \item $\sin3A+\sin3B+\sin3C=3\sin(A+B+C)$
\item $\cos3A+\cos3B+\cos3C=3\cos(A+B+C)$
\end{enumerate}
\item 在复平面上，$z_1=1$, $z_2=2+i$与$z_3$组成一个正三角形，求$z_3$.
\item 若$z_1,z_2,z_3,z_4$是一个正方形的四个顶点，且$z_1=1+2i$, $z_2=3-5i$, 求$z_3,z_4$.
\item 求证：$\triangle z_1z_2z_3$是等边三角形的充要条件是：
\[z_1^2+z_2^2+z_3^2=z_1\cdot z_2+z_2\cdot z_3+z_3\cdot z_1\]
\item 下列条件表示什么样图形？
\begin{enumerate}[(1)]
    \item ${\rm Im}(z^2)=2$
    \item $|z|=1$且$\arg z\in(0,\pi)$
    \item $0<a\le {\rm Im}(z)\le b$且$\frac{\pi}{4}\le \arg z\le \frac{3\pi}{4}$
\end{enumerate}
\item 你能用复数表示出图中的半圆来吗？（包括边界）

\begin{minipage}{.45\textwidth}
    \centering
\begin{tikzpicture}[>=stealth]
\draw[->](-1.5,0)--(1.5,0)node[right]{$x$};
\draw[->](0,-1)--(0,2)node[right]{$y$};
\draw[pattern=north east lines](1,0)node[below]{2} arc (0:180:1)node[below]{$-2$} --cycle;
\node[below left]{$O$};
\node at (0,1)[above right]{2};    
\end{tikzpicture}
\captionof{figure}{ }
\end{minipage}\hfill
\begin{minipage}{.45\textwidth}
    \centering
\begin{tikzpicture}[>=stealth]
    \draw[->](-1.5,0)--(1.5,0)node[right]{$x$};
\draw[->](0,-1.5)--(0,1.5)node[right]{$y$};
\draw[pattern=north east lines](0,1)node[right]{1} arc (90:180+90:1)node[right]{$-1$}--cycle; 
\node[below right]{$O$};   
\node at (-1,0)[below left]{$-1$};   
\end{tikzpicture}
\captionof{figure}{ }
\end{minipage} 

\item 设$z\in\mathbb{C}$，解方程：
\begin{multicols}{2}
\begin{enumerate}[(1)]
    \item $\frac{1}{2}(z-1)=\frac{\sqrt{3}}{2}(1+z)i$
    \item $(z+1)^8=(1+i)^8$
\end{enumerate}
\end{multicols}

\item 设$z$是虚数，解方程：
\begin{multicols}{2}
\begin{enumerate}[(1)]
    \item $z+|\bar z|=2+i$
    \item $z^2=\bar z$
\end{enumerate}
\end{multicols}

\item 求证：
\begin{enumerate}[(1)]
    \item $|z|=1\; (z\in\mathbb{C})$的充要条件是$z^{-1}=\bar z$;
    \item 共轭虚数的$n\; (n\in\N)$次幂仍是共轭虚数；
    \item 虚数的平方根仍是虚数．
\end{enumerate}

\item 在复数集中解下列方程组：
\begin{enumerate}[(1)]
    \item $\begin{cases}
        x^2+y^2-2xy+3(x+y)-2=0\\
        2(x^2+y^2)-2xy-2(x+y)+15=0
    \end{cases}$
    \item $\begin{cases}
        \frac{x-y}{1+xy}=\frac{1}{3}\\  \frac{x+y}{1-xy}=3
    \end{cases}$
\end{enumerate}    

\item 设三次单位根为$1,\omega,\omega^2$，其中
$\omega=-\frac{1}{2}+\frac{\sqrt{3}}{2}i$，求证：
\begin{enumerate}[(1)]
    \item $1+\omega+\omega^2=0$
    \item $1\cdot \omega\cdot \omega^2=1$
    \item $(1+\omega^2)^4=\omega$
    \item $\begin{vmatrix}
        1&1&\omega\\
        1&1&\omega^2\\
        \omega^2&\omega&1
    \end{vmatrix}=-3$
    \item $1^n+\omega^n+(\omega^2)^n=\begin{cases}
        3,& n=3k\\
        0,& n\ne 3k
    \end{cases}(k\in\Z)$

    \item $a^3+b^3=(a+b)(a\omega+b\omega^2)(a\omega^2+b\omega);$
    \item $a^{3}+b^{3}+c^{3}-3abc =(a+b+c)(a+b\omega+c\omega^{2})(a+b\omega^{2}+c\omega)$
\end{enumerate}

\item 利用上题中$\omega$的性质，计算：
\begin{multicols}{2}
\begin{enumerate}[(1)]
    \item $(-1+\sqrt{3}i)^{6}-(-1-\sqrt{3}i)^{6}$;
    \item $\left(\frac{\sqrt{3}+i}{2}\right)^{6}-\left(\frac{\sqrt{3}-i}{2}\right)^{6}$;
    \item $\left(-\frac{1}{2}i-\frac{\sqrt{3}}{2}\right)^{12}$.
\end{enumerate}    
\end{multicols}

\item 已知$f(x)=x^4+3x^3-30x^2+366x-340$的一个根为$3+5i$，求其余的三个根。
\item 若实系数方程$x^4+px^3+qx^2+rx+s=0$有两个根为$2+i$, $-2+i$. 求$p$、$q$、$r$、$s$的值。
\item 若实系数多项式$x^4+px^3+qx^2+rx+s$有一个根是纯虚数，求证：$r^2+p^2s=pqr$.

\item 已知实系数方程$x^2+mx+n=0$的一个根是另一个根的$i$倍，试求$m$、$n$的关系.
\item 已知$f(x)=x^4+(2+i)x^3+(3+2i)x^2-(4-3i)x-4i$, 

$g(x)=x^4-(1-i)x^3-(1+i)x^2+(4+i)x-4i$，求$f(x)$与$g(x)$的公根.
\end{enumerate}

\chapter{排列与组合~~二项式定理}
\section{排列与组合}
\subsection{加法原理和乘法原理}
假如我们要从甲地到乙地，可以乘火车、轮船或公共汽车，火车每天有两班，轮船有两班，公共汽车有六班，问从甲地到乙地有多少种不同的方法？

因为，从甲地到乙地，乘火车有两种不同的方法，乘轮船有两种不同的方法，乘公共汽车有六种不同的方法.而每种方法都可以由甲地到达乙地，因此，从甲地到乙地共有
$2+2+6=10$
种不同的方法.

一般地说，有如下原理：
\begin{blk}{加法原理}
如果完成某件事有$n$种方式，在第一种方式中有$m_1$种方法，在第二种方式中有$m_2$种方法……，在第$n$种方式中有$m_n$种方法. 那么完成这件事共有$m_1+m_2+\cdots+m_n$种不同的方法.
\end{blk}

又假如，要从甲地到丙地，必须经过乙地，现在已知由甲地到乙地有三条道路，由乙地到丙地有两条道路（图2.1），问从甲地到丙地共有多少种不同的走法？
\begin{figure}[htp]
    \centering
\begin{tikzpicture}
\node[draw, circle] (A) at (0,0){甲};    
\node[draw, circle] (B) at (4,0){乙};    
\node[draw, circle] (C) at (8,0){丙};
\draw(A)--(B);
\draw(A)to [bend left=45](B);
\draw(A)to [bend left=-45](B);
\draw(C)to [bend left=15](B);
\draw(C)to [bend left=-15](B);

\end{tikzpicture}
    \caption{}
\end{figure}

因为从甲地到乙地有3种不同的走法，到乙地后又各有两种不同的走法到丙地，因此，从甲地到丙地共有
$3\x2=6$
种不同的走法.

一般地说，有如下原理：

\begin{blk}
    {乘法原理}如果完成某种事需要分成$n$个步骤，第一步骤有$m_1$种方法，第二步骤有$m_2$种方法……，第$n$步骤有$m_n$种方法，那么完成这件事共有$m_1\cdot m_2\cdots m_n$种不同的方法.
\end{blk}

总的来说，如果完成一种事有几种不同的方法，这些方法又是互不相关的，任选一种方法都能完成这件事的，那么完成这件事的方法的总数，应该用加法计算. 如果完成一件事，必须分成若干步骤，每个步骤又有不同的方法，而且每一步骤中的一种方法完成之后，都可以开始下一步骤的工作，依次完成全部步骤，才能达到完成这一件事的目的，那么完成这件事的方法的总数应该用乘法计算.

\begin{ex}
\begin{enumerate}
    \item 完成某件工作，甲有两种不同的方法，乙有3种不同方法，丙有4种不同方法．从三人中选一人完成这件工作，共有几种不同选法？
    \item 一种车床用两个手柄联合控制转速，一个手柄有两档，另一个手柄有3档，问能控制几种不同的转速？
    \item 乘积$(a+b)\cdot (c+d+e)\cdot (m+n+p+q)$展开后共有多少项？
    \item 从甲地到乙地有一条道路，从乙地到丁地有两条道路；从甲地到丙地有三条道路，从丙地到丁地有四条道路.若从甲地到丁地必须经过乙地或丙地，问从甲地到丁地共有多少种不同的走法？
\end{enumerate}
\end{ex}


\subsection{排列}
一般地说，从$n$个不同的元素中任意选取$m\; (m\le n)$个元素，按照一定的顺序排成一列，叫做从$n$个元素中取出$m$个元素的一个排列.

例如：由$1,2,3$三个数字中选取两个数字写成两位数，那么$1,2,3$都被看作是排列的元素。而12与21显然都是由$1,2$两个元素得到的不同的排列．同样，12，21，13，31，23，32等都被认为是不同的排列．只有用相同的元素，又是按相同的顺序排列而成的排列，才叫做相同的排列，例如：123与123就是相同的排列．

\begin{example}
    由数字1,2,3,4可以组成多少个没有重复数字的两位数？
\end{example}

\begin{solution}
    要排出两位数，第一步先要排出十位上的数字（当然也可以先排个位上的数字），第二步是确定个位上的数字.

显然，第一步在十位的位置上，$1,2,3,4$都可以放，故有四种不同的方法.

第二步是放个位上的数字，由于在第一步结束时，剩下的只有3个数字，因此第二步只能有3种不同的选择．因此个位与十位上的数的排列工作都结束后，这个两位数才算排出来，所以，用乘法原理计算得：$4\x3=12$．即以$1,2,3,4$可以排出12个不同的两位数．
\end{solution}

\begin{center}
\begin{tikzpicture}
\begin{scope}
\draw(-.5,-.25) rectangle (.5,.25);
\draw(0,-.25)--(0,.25);
\node at (-0.25,0){1};
\foreach \x/\y in {-1/4,0/3,1/2}
{
    \draw(1, \x-.25) rectangle (2, \x+.25);
    \draw(1.5,\x-.25)--(1.5,\x+.25);
    \node at (1.25,\x){1};
    \node at (1.75,\x){\y};
}
\draw[decorate, decoration={brace, amplitude=5}](.9,-1.25)--(.9,1.25);
\end{scope}
\begin{scope}[xshift=3cm]
    \draw(-.5,-.25) rectangle (.5,.25);
    \draw(0,-.25)--(0,.25);
    \node at (-0.25,0){2};
    \foreach \x/\y in {-1/4,0/3,1/1}
    {
        \draw(1, \x-.25) rectangle (2, \x+.25);
        \draw(1.5,\x-.25)--(1.5,\x+.25);
        \node at (1.25,\x){2};
        \node at (1.75,\x){\y};
    }
    \draw[decorate, decoration={brace, amplitude=5}](.9,-1.25)--(.9,1.25);
    \end{scope}
    \begin{scope}[xshift=6cm]
        \draw(-.5,-.25) rectangle (.5,.25);
        \draw(0,-.25)--(0,.25);
        \node at (-0.25,0){3};
        \foreach \x/\y in {-1/4,0/2,1/1}
        {
            \draw(1, \x-.25) rectangle (2, \x+.25);
            \draw(1.5,\x-.25)--(1.5,\x+.25);
            \node at (1.25,\x){3};
            \node at (1.75,\x){\y};
        }
        \draw[decorate, decoration={brace, amplitude=5}](.9,-1.25)--(.9,1.25);
        \end{scope}
\begin{scope}[xshift=9cm]
\draw(-.5,-.25) rectangle (.5,.25);
\draw(0,-.25)--(0,.25);
\node at (-0.25,0){4};
\foreach \x/\y in {-1/3,0/2,1/1}
{
    \draw(1, \x-.25) rectangle (2, \x+.25);
    \draw(1.5,\x-.25)--(1.5,\x+.25);
    \node at (1.25,\x){4};
    \node at (1.75,\x){\y};
}
\draw[decorate, decoration={brace, amplitude=5}](.9,-1.25)--(.9,1.25);
\end{scope}
\end{tikzpicture}
\end{center}




\begin{example}
$a,b,c,d$四个元素，每次取出三个元素的不同排列有多少种，并写出所有的排列.（在同一排列里，元素不能重复出现.）
\end{example}

\begin{solution}
要把取出的三个元素排成一列，要分三步完成. 第一步先在第一个位置上放一个元素，应该有四种不同的方法；第二步要在第二个位置上放一个元素，由于第一步的工作完成后剩下的是三个元素，从三个元素中选一个放在第二个位置上应有3种不同的方法；第三步要在第三个位置上放一个元素，由于前两步工作完成后只剩下两个元素，所以只有两种不同的方法. 三步工作都完成后才完成排列工作，因此应该用乘法原理计算排列的总数，所以共有：$4\x3\x2=24$种不同的排列．即有如下的24种不同的排列：
\begin{center}
\begin{tabular}{cccccc}
$abc$& $abd$& $acb$ &$acd$ &$adb$& $adc$\\
$bac$& $bad$ &$bca$& $bcd$ &$bda$& $bdc$\\
$cab$& $cad$ &$cba$& $cbd$ &$cda$& $cdb$\\
$dab$& $dac$& $dba$ &$dbc$& $dca$ &$dcb$    
\end{tabular}
\end{center}
\end{solution}

一般地说，从$n$个不同元素中，任取$m$个$(m\le n)$元素（不许重复），按照一定顺序排成一列的个数，叫做从$n$个不同元素中取出$m$个元素的排列数，用符号表示$\perm{m}{n}$.

\begin{blk}
  {定理1} 从$n$个不同元素中，任取$m$个$(m\le n)$没有重复的元素的排列数为$\perm{m}{n}=n(n-1)\cdots (n-m+1)$
\end{blk}

\begin{proof}
    在取出$m$个元素排成一列时，首先在第一个位置上可以选$n$个元素中的任何一个，故有$n$种方法. 在第二个位置上可以选余下的$n-1$个元素中的任一个．如此类推，在第$m$个位置上可以选余下的$(n-m+1)$个元素中的任一个. 根据乘法原理得到总的不同排列数应该是
\[\perm{m}{n}=n(n-1)\cdots(n-m+1)\]

这里的$m$是自然数，并且$m\le n$, 这个公式就叫做排列数公式.
\end{proof}

特别是当$m=n$时 有
\[\perm{n}{n}=n(n-1)\cdots3\cdot2\cdot1\]
这表示，$n$个不同的元素全部取出参加排列的排列数，叫做$n$个不同元素的全排列，自然数1到$n$的连乘积叫做$n$的阶乘，用$n!$表示，所以全排列数公式又可以写成
\[\perm{n}{n}=n!\]
有了阶乘，排列数公式又可以写成
\[\perm{m}{n}=\frac{n!}{(n-m)!}\]
这里，当$m=n$时，分母出现$(n-m)!$就变成$0!$，为了使公式仍能成立，我们特别规定$0!=1$.

例如$\perm{5}{8}=8\cdot  7\cdot 6\cdot 5\cdot 4=6720$，
$\perm{5}{5}=5!=5\cdot 4\cdot 3\cdot 2\cdot 1=120$

\begin{example}
    试证明下列等式成立．
\begin{multicols}{2}
\begin{enumerate}[(1)]
    \item $\perm{m}{n}+m\cdot \perm{m-1}{n}=\perm{m}{n+1}\quad (m\le n)$
\item $\perm{n+1}{n+1}-\perm{n}{n}=n^2\cdot \perm{n-1}{n-1}$
\end{enumerate}
\end{multicols}
\end{example}

\begin{proof}
\begin{enumerate}[(1)]
    \item \[\begin{split}
        \perm{m}{n}+m\cdot \perm{m-1}{n}&=\frac{n!}{(n-m)!}+m\cdot\frac{n!}{(n-m+1)!}\\
        &=\frac{(n-m+1)\cdot n!+m\cdot n!}{(n-m+1)!}=\frac{(n+1)\cdot n!}{(n-m+1)!}\\
        &=\frac{(n+1)!}{[(n+1)-m]!}=\perm{m}{n+1}
    \end{split}\]
所以原等式成立.

\item \[\begin{split}
    \perm{n+1}{n+1}-\perm{n}{n}&=(n+1)!-n!=(n+1)\cdot n!-n!\\
    &=n\cdot n!=n^2\cdot (n-1)!\\
    &=n^2\cdot \perm{n-1}{n-1}
\end{split}\]
所以原等式成立.
\end{enumerate}
\end{proof}

\begin{example}
    用0到9十个数字，可以写出多少个没有重复数的三位数？其中有多少个是偶数？
\end{example}

\begin{solution}
\begin{enumerate}[(1)]
    \item 第一步，因为0不能放在百位的位置上，所以百位上的数有9种不同的选法，
    
第二步是从余下的九个数字中任选两个数字放在十位个位上，这有$\perm{2}{9}$种方法.所以
\[9\cdot \perm{2}{9}=9\cdot 9\cdot 8=648\]
\item 要得到三位偶数，必须个位上的数是0, 2, 4, 6或8．但是0又不能排在百位上，这说明0怎样放在适合的位置上是比其他问题要求更高的. 所以我们还是把0元素的题先解决.

第一类，0放在个位上这一选法确定后，其余百位、十位可以从余下的九个数字中任选两个数字放上去，这有$\perm{2}{9}$种方法.

第二类，个位上可放2,4,6,8中的任一个，在这类方法中，第一步在个位上有4种选法，第二步，在百位上可以放除0以外所余的八个数字的任一个，故有八种选择方法．第三步，余下的八个数字均可以放在十位上，所以
\[\perm{2}{9}+4\x8\x\perm{1}{8}=72+256=328\]
\end{enumerate}
答：有648个没有重复数字的三位数，其中有328个是偶数.

\end{solution}

关于例2.4中的第一个问题，我们还可以有如下的解法：

从0到9十个数字中任选三个数字排列总数为$\perm{3}{10}$．但是其中0在百位上的数有$\perm{2}{9}$个. 所以
\[\perm{3}{10}-\perm{2}{9}=10\cdot 9\cdot 8-9\cdot 8=648\]

\begin{ex}
\begin{enumerate}
    \item 写出从$a,b,c,d,e$中任取两个不同元素的所有排列.
    \item 计算：
\begin{multicols}{3}
\begin{enumerate}[(1)]
    \item $\perm{3}{7}$
    \item $\perm{3}{15}$
    \item $\perm{4}{4}$
    \item $\perm{4}{8}-3\cdot\perm{2}{8}$
    \item $\perm{3}{7}-5\perm{2}{7}$
    \item $\frac{\perm{7}{12}}{\perm{6}{11}}$
\end{enumerate}
\end{multicols}

\item 求证：
\begin{enumerate}[(1)]
    \item $n!=\frac{(n+1)!}{n+1}$
    \item $\perm{m}{n}=\frac{\perm{n+2}{n+2}}{\perm{2}{n+2}\cdot \perm{n-m}{n-m}}\quad (m\le n)$
\end{enumerate}

\item 某一段铁路共有8个车站，需要为这段铁路印制多少种不同的客车票？（都是普通座位票）
\item 6个同学排成一列，有多少种不同的排法？
\item 从12人的1个学习小组中，选1个正组长和1个副组长，问有多少种不同的选法？
\item 用0至5共六个数字，可以组成多少不同的三位数？
\item 用0到9这十个数字，其中有多少个没有重复数字的二位数，其中有多少个二位奇数？
\item 某信号兵用红、黄、蓝三面旗子从上到下挂在竖直的旗杆上表示信号，每次可以任挂一面、二面或三面，并且不同的顺序表示不同的信号，问一共可以表示多少种不同的信号？
\item 有五道不同的工序分配给五个工人工作，其中有一个工人不会做其中两道工序，问有多少种不同的分配法？
\item 由2, 3, 4, 5所组成的没有重复数字的四位数，
\begin{enumerate}[(1)]
\item 求所有这些四位数里各个数字的和；
\item 求所有这些四位数的和。
\end{enumerate}
\end{enumerate}
\end{ex}

\subsection{组合}
我们知道，从一个学习小组中选出正、副组长各1人，这种选法是排列问题.如果改变为从一个学习小组中选出代表2人参加学校的学代会，有多少种不同的选法？”这时的代表无所谓正、副之分，就无先后顺序的意义了，也就不成为排列的问题了.

又例如从1, 2, 3, 4, 5五个数选出两个数字相乘，问有多少个不同的乘积？那么只要所选出的数字是3和5；它们的积只有15，至于是$5\x3$还是$3\x5$，那是不要求加以区别的，也是没有排列次序的区别的.

一般地说，从$n$个不同的元素中，任取$m\; (m\le n)$个元素并成一组，叫做从$n$个不同元素中取出$m$个元素的一个组合.

这里要注意，元素$a$，$b$，$c$的组合与$a,c,b$; $c,a,b$; $b,a,c$; $b,c,a$或$c,b,a$都被看作是相同的组合了. 对于排列来说，上列的是六个不同的排列，而对于组合来说，只算是一个组合，那么组合数是怎样进行计算的呢？

从$n$个不同元素中取出$m\; (m\le n)$个元素的所有组合的个数，叫做从$n$个不同元素中取出$m$个元素的组合数，用符号$\comb{m}{n}$表示。

例如，从四个不同元素中取出3个元素进行排列，有多少种不同的排法？

我们已经知道应该有$\perm{3}{4}=24$种不同的排法，现在我们
想从另外的方法入手解决这个问题.

第一步，从$a,b,c,d$中取三个元素作一组，这样可以得到$a,b,c$; $a,b,d$; $a,c,d$; $b,c,d$四个组，这就是说$\comb{3}{4}=4$.

第二步，把每组的三个元素作全排列，应有$\perm{3}{3}$个不同排列.

根据乘法原理得$\perm{3}{4}=4\cdot \perm{3}{3}$，即$\perm{3}{4}=\comb{3}{4}\cdot \perm{3}{3}$.

又如，从五个不同的元素中取出2个元素的排列数是$\perm{2}{5}$. 如果我们从五个不同元素中取出两个分成一组，我们可以组成如下的小组：

$a,b$; $a,c$; $a,d$; $a,e$; $b,c$; $b,d$; $b,e$; $c,d$, $c,e$和$d,e$等10个组，这就是$\comb{2}{5}$；第二步把每组元素进行全排列应是$\perm{2}{2}$，这样应该有$\perm{2}{5}=\comb{2}{5}\cdot \perm{2}{2}$. 因此，一般地说，从$n$个不同元素中取出$m\; (m\le n)$个元素的排列数$\perm{m}{n}$可以等于$\comb{m}{n}\cdot \perm{m}{m}$.


\begin{blk}
  {定理2} 从$n$个不同元素中任取$m$个（$m\le n$）没有重复元素组合数为$\comb{m}{n}=\frac{n!}{m!(n-m)!}$.  
\end{blk}

\begin{proof}
    由$\perm{m}{n}=\comb{m}{n}\cdot \perm{m}{m}$，所以
\[\comb{m}{n}=\frac{\perm{m}{n}}{\perm{m}{m}}=\frac{\frac{n!}{(n-m)!}}{m!}=\frac{n!}{m!(n-m)!}\]
\end{proof}
注意：$\comb{m}{n}=\frac{n(n-1)\cdots(n-m+1)}{m!}$

\begin{example}
计算$\comb{3}{8}$, $\comb{5}{8}$和$\comb{2}{9}$
\end{example}

\begin{solution}
    根据组合数公式有：
\[\begin{split}
    \comb{3}{8}&=\frac{8\x 7\x 6}{3\x 2\x 1}=56\\
    \comb{5}{8}&=\frac{8\x 7\x 6\x 5\x 4}{5\x 4\x 3\x 2\x 1}=56\\
    \comb{2}{9}&=\frac{9\x 8}{2\x 1}=36
\end{split}\]
\end{solution}

组合数有以下性质：
\begin{blk}
{性质1}$$\comb{m}{n}=\comb{n-m}{n}$$
\end{blk}

\begin{proof}
根据组合数公式：
\[\begin{split}
    \comb{m}{n}&=\frac{n!}{m!(n-m)!}\\
    \comb{n-m}{n}&=\frac{n!}{(n-m)![n-(n-m)]!}=\frac{n!}{m!(n-m)!}
\end{split}\]
所以$\comb{m}{n}=\comb{n-m}{n}$.
\end{proof}

\begin{blk}{性质2}
\[\comb{m}{n+1}=\comb{m}{n}+\comb{m-1}{n}\]
\end{blk}

\begin{proof}
根据组合数公式：
\[\begin{split}
    \comb{m}{n+1}&=\frac{(n+1)!}{m!(n-m+1)!}\\
    \comb{m}{n}+\comb{m-1}{n}&=\frac{n!}{m!(n-m)!}+\frac{n!}{(m-1)!(n-m+1)!}\\
    &=\frac{(n-m+1)\cdot n! +m\cdot n!}{m!(n-m+1)!}=\frac{n!\cdot (n-m+1+m)}{m!(n-m+1)!}\\
    &=\frac{(n+1)!}{m!(n-m+1)!}\\
\end{split}\]
所以$\comb{m}{n+1}=\comb{m}{n}+\comb{m-1}{n}$.
\end{proof}

为了使性质1在$n=m$时也能成立，我们规定$\comb{0}{n}=1$.

\begin{example}
计算$\comb{98}{100}$和$\comb{2}{10}+\comb{3}{10}+\comb{2}{11}+\comb{2}{12}$
\end{example}

\begin{solution}
根据组合数性质1
\[\comb{98}{100}=\comb{100-98}{100}=\comb{2}{100}=\frac{100\x 99}{2\x 1}=4950\]
根据组合数性质2
\[\begin{split}
    \comb{2}{10}+\comb{3}{10}+\comb{2}{11}+\comb{2}{12}&=\comb{3}{11}+\comb{2}{11}+\comb{2}{12}\\
    &=\comb{3}{12}+\comb{2}{12}=\comb{3}{13}\\
    &=\frac{13\x 12\x 11}{3\x 2\x 1}=286
\end{split}\]
\end{solution}

\begin{example}
    解不等式 $\comb{3}{n}>\comb{5}{n}$
\end{example}

\begin{solution}
\[\frac{n(n-1)(n-2)}{3\x 2\x 1}>\frac{n(n-1)(n-2)(n-3)(n-4)}{5\x 4\x 3\x 2\x 1}\]
因为$5\le n$，所以$n,\; n-1,\; n-2,\; n-3,\; n-4$均为正数，所以
\[1>\frac{(n-3)(n-4)}{5\x 4}\]
整理得：$n^2-7n-8<0$

解不等式得：$-1<n<8$，根据题意，得：$n=5,6,7$.
\end{solution}



\begin{example}
高中三个年级，每个年级有六个班。
\begin{enumerate}[(1)]
\item 每级都进行单循环（每班都与级内其他各班比赛一场）篮球预赛，问预赛需要比赛多少场？
\item 从各级预赛中选出每级的冠、亚军，再进行单循环决赛，规定在预赛中已经比赛过的队不再比赛，问决赛需要比赛多少场？
\end{enumerate}
\end{example}

\begin{solution}
\begin{enumerate}[(1)]
    \item 根据题意每个级单循环比赛的场次都是$\comb{2}{6}$，所以
\[3\comb{2}{6}=3\times 15=45\]
\item 决赛共6个班，如果单循环要比赛$\comb{2}{6}$场，但预赛中比赛过的有3场，所以
\[\comb{2}{6}-3=15-3=12\]
\end{enumerate}
答：预赛要比赛45场，决赛要比赛12场.
\end{solution}



\begin{example}
    平面有10个不同的点，其中除4个点在同一直线上外，不再有3个点共线，以每三个点为顶点可以作多少个不同的三角形？
\end{example}


\begin{solution}
    如果10点中的任三点都能连成三角形，那么有$\comb{3}{10}$个三角形. 但是如果在四个共线点中任选三点就不可能构成三角形. 所以
\[\comb{3}{10}-\comb{3}{4}=120-4=116\]
答：可以作116个不同的三角形。
\end{solution}



\begin{example}
    在产品检验时，常从产品中抽出一部分进行检查，假设100件产品中有2件次品，
    \begin{enumerate}[(1)]
        \item 从100件产品中抽出三件检查，其中有1件次品的抽法有多少种？
        \item 从100件产品中抽出4件检查，其中至少有一件次品的抽法有多少种？
    \end{enumerate}
\end{example}

\begin{solution}
\begin{enumerate}[(1)]
    \item 在100件产品中，次品有两件，从2件次品中抽出1件的抽法有$\comb{1}{2}$，从98件正品中抽出2件的方法有$\comb{2}{98}$. 所以
\[\comb{1}{2}\cdot \comb{2}{98}=2\cdot \frac{98\cdot 97}{2}=9506\]
    \item 至少有1件是次品就是包括有1件次品和2件次品两类抽法的和，所以 
\[\comb{1}{2}\cdot \comb{3}{98}+\comb{2}{2}\cdot \comb{2}{98}=2\times \frac{98\cdot 97\cdot 96}{3\cdot 2}+\frac{98\cdot 97}{2}=304192+4753=308945\]
\end{enumerate}
答：3件中有1件是次品的抽法是9506, 4件中至少有1件次品的抽法是308945。
\end{solution}

\begin{example}
    从0, 3, 5, 7, 11五个数中可以得到多少个不同的乘积？
\end{example}

\begin{solution}
    因为零乘以任何数其积都是零，所以不论零与多少个数相乘都只算一个乘积就是0；

在非零的乘积中有$\comb{2}{4}$, $\comb{3}{4}$, $\comb{4}{4}$三类，所以
\[\comb{2}{4}+\comb{3}{4}+\comb{4}{4}+1=6+4+1+1=12\]

答：可以得到12个不同的乘积.
\end{solution}

\begin{ex}
\begin{enumerate}
    \item 写出：
    \begin{enumerate}[(1)]
    \item 从4个相异元素$a$，$b$，$c$，$d$中取出2个元素的所有组合；
    \item 从5个相异元素$a$，$b$，$c$，$d$，$e$中取出3个元素的所有组合.
    \end{enumerate}

    \item 计算：
\begin{multicols}{2}
\begin{enumerate}[(1)]
    \item  $\comb{2}{7}$
    \item $\comb{48}{50}$
    \item $\comb{3}{100}-\comb{3}{90}$
    \item $\comb{1}{5}+\comb{2}{5}+\comb{3}{5}+\comb{4}{5}+\comb{5}{5}$
\end{enumerate}
\end{multicols}
    
\item   平面内有12个点，其中任意3点不在同一直线上，以每3点为顶点作一三角形，一共可以作多少个不同的三角形？
\item 圆上$n$个不同的点，
\begin{enumerate}[(1)]
\item 可以引多少条不同的弦？
\item 可以连接多少个内接三角形？
\end{enumerate}

\item 凸四边形可以连多少条对角线？凸五边形可以连多少条对角线？凸$n$边形可以连多少条对角线？
\item 某班有50名学生，其中正副班长各1人，现选5名学生参加某项活动
\begin{enumerate}[(1)]
\item 如果班长和副班长必须在内，有多少种不同的选法？
\item 如果要有一个班长（正、副都可以）在内，有多少种不同的选法？
\item 如果班长都不参加，有多少种不同的选法？
\item 如果至少要有一个班长（正、副都可以）在内，有多少种不同的选法？
\end{enumerate}
\end{enumerate}
\end{ex}

\section*{习题2.1}
\begin{enumerate}
    \item 计算：
\begin{multicols}{2}
\begin{enumerate}[(1)]
    \item $5\cdot \perm{3}{5}+4\cdot \perm{2}{4}$
    \item $\frac{\perm{5}{7}-\perm{6}{6}}{7!+6!}$
\end{enumerate}
\end{multicols}

\item 求证：
\begin{enumerate}[(1)]
    \item $\frac{(n+1)!}{k!}-\frac{n!}{(k-1)!}=\frac{(n-k+1)\cdot n!}{k!}$
    \item $\frac{(2n)!}{2^n \cdot n!}=1\cdot 3\cdot 5\cdots (2n-1)$
\end{enumerate}

\item 解方程：
\begin{multicols}{2}
\begin{enumerate}[(1)]
    \item $\perm{3}{2n}=10\perm{3}{n}$
    \item $\frac{\perm{5}{n}+\perm{4}{n}}{\perm{3}{n}}=4$
    \item $\perm{4}{2n+1}=140\perm{3}{n}$
\end{enumerate}
\end{multicols}

\item 解不等式：
\begin{multicols}{2}
\begin{enumerate}[(1)]
    \item $\perm{2}{m-2}+m>2$
    \item $\perm{x}{9}>6\perm{x-2}{9}$
\end{enumerate}
\end{multicols}

\item \begin{enumerate}[(1)]
    \item 从多少个不同的元素中取出2个元素的排列数是56？
    \item 已知从$n$个不同的元素中取出2个元素的排列数等于从$n-4$个不同元素中取2个元素的排列数的7倍，求$n$．
\end{enumerate}
\item 有5本不同的书分借给3个同学，每人1本，共有多少种不同的借法？
\item 用0, 1, 2, 3, 4, 5可以组成多少个没有重复数字的五位数？其中有多少个偶数？
\item 用数字1, 2, 3, 4, 5可以组成多少个没有重复数字的自然数？其中有多少个比13000大的自然数？
\item 7部不同型号的机床分配给7个工人使用：
\begin{enumerate}[(1)]
\item 甲必须使用2号机床；
\item 甲必须使用2号机床，乙必须使用4号机床；
\item 甲不适宜使用1号机床，乙不适宜使用7号机床；问三种情况下各有多少种不同的分配方案？
\end{enumerate}

\item 6个儿童站成一排，其中一人不站在排头，也不站在排尾，一共有多少种不同的排法？
\item 五部不同型号的半导体收音机和四部不同型号的电视机陈列成一排，任两部电视机不靠在一起，有多少种不同的陈列方法？
\item 在3000和7000间有多少个没有重复数字的5的倍数？

\item 计算：
\begin{multicols}{2}
\begin{enumerate}[(1)]
    \item $\comb{3}{15}$
    \item $\comb{98}{100}$
    \item $\comb{3}{6}+\comb{6}{8}$
    \item $\comb{n}{n+1}\cdot \comb{n-2}{n}$
\end{enumerate}
\end{multicols}

\item 求证：
\begin{enumerate}[(1)]
    \item $\comb{m}{n+1}=\comb{m-1}{n}+\comb{m}{n-1}+\comb{m-1}{n-1}$
    \item $\comb{k+1}{n}\div \comb{k}{n}=\frac{n-k}{k+1}$
    \item $\comb{0}{m}+\comb{1}{m+1}+\comb{2}{m+2}+\cdots+\comb{n-1}{m+n-1}=\comb{n-1}{m+n}$
    \item $\comb{m+1}{n}+\comb{m-1}{n}+2\comb{m}{n}=\comb{m+1}{n+2}$
\end{enumerate}

\item 解不等式：
\begin{multicols}{2}
\begin{enumerate}[(1)]
    \item $\comb{4}{n}>\comb{6}{n}$
    \item $\comb{n-4}{21}<\comb{n-2}{21}<\comb{n-1}{21}$
\end{enumerate}
\end{multicols}

\item 在200件产品中有2件次品，现在抽5件进行检查：
\begin{enumerate}[(1)]
\item 其中有一件次品的不同抽法有多少种？
\item 没有次品的不同抽法有多少种？
\item 至少有一件次品的不同抽法有多少种？
\end{enumerate}

\item 问2310能被多少个偶数整除？
\item 有8个划船运动员，其中3人能划左舷，2人能划右舷，3人能划左右舷，要选出6人分配在两边，问有多少种不同的选法？

\item  某车间有9名技术熟练的工人，其中2人能当车、钳工，有3人只能当车工，其余4人只能当钳工，现在要抽调三个车工，三个钳工成立一个新车间，问有多少种不同的抽调方法？
\item  空间有12个不同的点，其中有4点在同一平面内，此外不再有4点共面，
\begin{enumerate}[(1)]
\item  这些点可以确定多少个不同的平面？
\item  这些点可以连成多少个不同的四面体？
\end{enumerate}

\item  有8部机器分给陈、李、张三位师傅管理。
\begin{enumerate}[(1)]
\item 陈管4部，李管3部，张管1部，有几种不同分配方法？
\item 陈管4部，其他机器由1人管3部，1人管1部，有多少种不同的分配方法？
\item 1人管4部，一人管3部，一人管1部，有几种不同的分配方法？
\end{enumerate}

\item  从1, 5, 7, 9中任取3个数字，2, 4, 6, 8中任取两个数字，组成没有重复数字的五位数，一共有多少个不同的五位数？
\end{enumerate}

\section{环形排列、重复排列和组合}
\subsection{环形排列}
我们已经掌握了从$n$个相异的元素中取出$m\; (m\le n)$个元素的排列的知识了，那么怎样解决环形排列的问题呢？

我们先从最简单的情况入手，就是先找出环形排列与直
线排列的相互关系，然后利用直线排列的知识解决环形排列的问题.

首先在环形排列中，下面的排列都是当作相同的排列的（图2.2）．
\begin{figure}[htp]
    \centering
\begin{tikzpicture}[>=stealth]
\begin{scope}
\draw circle(.9);
\node at (90:1.2){甲};
\node at (90+180:1.2){乙};
\draw[->](90-20:1.2) arc (90-20:20:1.2);
\end{scope}
\begin{scope}[xshift=3cm]
\draw circle(.9);
\node at (0:1.2){甲};
\node at (0+180:1.2){乙};
\draw[->](0-20:1.2) arc (0-20:-70:1.2);
\end{scope}
\begin{scope}[xshift=6cm]
\draw circle(.9);
\node at (-90:1.2){甲};
\node at (-90+180:1.2){乙};
\draw[->](-90-20:1.2) arc (-90-20:-160:1.2);
\end{scope}
\begin{scope}[xshift=9cm]
\draw circle(.9);
\node at (45:1.2){甲};
\node at (45+180:1.2){乙};
\draw[->](45-20:1.2) arc (45-20:-20:1.2);
\end{scope}

\end{tikzpicture}
    \caption{}
\end{figure}


因为从甲开始，按同一方向，例如顺时针方向看过去都是乙，所以尽管位置不同，但仍然是同一样的元素的排列顺序.
\begin{figure}[htp]
    \centering
\begin{tikzpicture}[>=stealth]
\begin{scope}
\draw circle(.9);
\node at (90:1.2){甲};
\node at (90-120:1.2){乙};
\node at (90+120:1.2){丙};
\draw[->](90-20:1.2) arc (90-20:20:1.2);
\end{scope}
\begin{scope}[xshift=3.5cm]
\draw circle(.9);
\node at (0:1.2){甲};
\node at (0-120:1.2){乙};
\node at (0+120:1.2){丙};
\draw[->](0-20:1.2) arc (0-20:-70:1.2);
\end{scope}
\begin{scope}[xshift=7cm]
\draw circle(.9);
\node at (-90:1.2){甲};
\node at (-90-120:1.2){乙};
\node at (-90+120:1.2){丙};
\draw[->](-90-20:1.2) arc (-90-20:-160:1.2);
\end{scope}
  
\end{tikzpicture}
    \caption{}
\end{figure}


又例如三个不同元素甲、乙、丙如图2.3的三个环形排列也都算是相同的排列，而且与从谁开始无关，即甲$\to $乙$\to $丙$\to $甲、乙$\to $丙$\to $甲$\to $乙与丙$\to $甲$\to $乙$\to $丙，也都算作是相同的环形排列.

从上述的例子，可以得到三个不同元素的环形排列数应该是2．具体的排列是：甲—乙—丙—甲与甲—丙—乙—甲．

现在，我们再来分析一下环形排列与直线排列不同的特点. 我们已经知道下列三个环形排列都是相同的排列（图2.4）
\begin{figure}[htp]
    \centering
\begin{tikzpicture}[>=stealth]
\begin{scope}
\draw circle(.9);
\node at (90:1.2){甲};
\node at (90-120:1.2){乙};
\node at (90+120:1.2){丙};
\draw[->](-.35,1.5)--(-.35,1);
\end{scope}
\begin{scope}[xshift=3.5cm]
    \draw circle(.9);
    \node at (90:1.2){丙};
    \node at (90-120:1.2){甲};
    \node at (90+120:1.2){乙};
    \draw[->](-.35,1.5)--(-.35,1);

    \end{scope}
\begin{scope}[xshift=7cm]
    \draw circle(.9);
    \node at (90:1.2){乙};
    \node at (90-120:1.2){丙};
    \node at (90+120:1.2){甲};
    \draw[->](-.35,1.5)--(-.35,1);

    \end{scope}
\end{tikzpicture}
    \caption{}
\end{figure}

如果我们从图上所示的地方把环切断，并把环拉成直线后，我们却可以得到三个不同的直线的排列，就是：
\[\text{甲、乙、丙；丙、甲、乙；乙、丙、甲}\]

而三个不同元素的全排列数是$\perm{3}{3}$，所以三个不同元素的环形排列数是$\frac{\perm{3}{3}}{3}$

\begin{blk}
    {定理3} $m$个不同元素的环形排列数就是：
\[\frac{\perm{m}{m}}{m}=\frac{m!}{m}=(m-1)!\]
\end{blk}

\begin{proof}
设有$m$个不同元素排成一个环形，如图（图2.5）
\begin{figure}[htp]
    \centering
\begin{tikzpicture}
\draw(0,0) circle(1);
\draw[fill](15:1)node[right]{$a_4$} circle(1pt);
\draw[fill](15+30:1)node[above right]{$a_3$} circle(1pt);
\draw[fill](15+60:1)node[above]{$a_2$} circle(1pt);
\draw[fill](15+90:1)node[above]{$a_1$} circle(1pt);
\draw[fill](15+120:1)node[above left]{$a_m$} circle(1pt);
\draw[fill](15+150:1)node[left]{$a_{m-1}$} circle(1pt);
\end{tikzpicture}
    \caption{}
\end{figure}
就是按顺时针方向的一个排列．其次，$a_1$不动，改变$a_2,a_3,\ldots,a_m$的顺时针排列顺序，就得到第二个不同的环形排列．因此，如果$a_1$不动，要使$a_2,a_3,\ldots,a_m$有不同的排
列，其不同的排列数为$(m-1)!$，故$m$个不同元素的环形排列数为$\perm{m-1}{m-1}=(m-1)!$
\end{proof}


\begin{blk}
  {定理4} 从$n$个不同的元素中，任取$m$个（$m\le n$）没有重复的元素的环形排列数$\comb{m}{n}\cdot (m-1)!$.  
\end{blk}

\begin{proof}
    从$n$个不同元素中取出$m\; (m\le n)$个元素进行环形排列，我们第一步可以先从$n$个元素中先取出$m$个组成一组，第二步是把$m$个元素排成一圈，根据乘法原理，其排列数应为：
\[\comb{m}{n}\cdot \perm{m-1}{m-1}=\comb{m}{n}(m-1)!\]
\end{proof}

\begin{example}
    从10个不同元素中取出6个元素排成一圈，有多少种不同的排法？
\end{example}

\begin{solution}
    从10个不同元素中取出6个的环形排列数为：
\[\comb{6}{10}\cdot\perm{6-1}{6-1}=\comb{4}{10}\cdot  5!=210\cdot 120=25200\]
答：有25200种不同的排法．
\end{solution}

\begin{example}
    有五个男同学与五个女同学围成圆圈跳集体舞，如果男女相间，问有几种不同的排列方法？
\end{example}

\begin{solution}
第一步，五个男同学先作环形排列，其排列数为$4!$；

第二步，五个女同学分别站在两个男同学之间的位置上作全排列.

$\therefore\quad 4! \cdot 5! =2880$

答：有2880种不同的排列方法.    
\end{solution}

\begin{ex}
\begin{enumerate}
\item 把10个人平均分成两组，每组都围坐一个圆桌，问有多少种不同的坐法？
\item 今有5颗不同颜色的珠子，用线把它们串成珠圈，问有多少种串法？
\end{enumerate}
\end{ex}

\subsection{重复排列}

前面我们所研究的都是不允许元素在排列和组合中重复出现的，假如我们允许元素在同一排列中重复出现的话，其排列数应该是多少呢？

\begin{blk}
{定理5} 从$n$个不同的元素中，取出$m$个元素排成一列（允许元素重复出现）的不同排列数为$n^m$.
\end{blk}

\begin{proof}
    选在第一个位置上的元素有$n$个不同的方法，选在第二个位置上的元素仍然有$n$个不同的方法，……，选在第$m$个位置上的元素仍然有$n$个选法，所以根据乘法原理，从$n$个不同的元素中取出$m$个元素的允许重复的排列数应为$n^m$.
\end{proof}

这种排列又称为重复排列，这里应该注意的是，由于元素可以重复选取，因此选取的次数$m$就不会再受$n$个元素的限制. 即$m$不再有$m\le n$的限制了.

\begin{example}
    电报通讯的明码是用0到9十个数字中的每4个数字的重复排列代表一个汉字，问用这种方法制成的电报通讯明码可以表示多少个不同的汉字？
\end{example}

\begin{solution}
    从10个不同数字中选取4个数字的重复排列数是：
\[10^4=10000\]
答：可以表示10000个不同的汉字.
\end{solution}

\begin{example}
    由1, 2, 3, 4四个数字能够组成多少个大于2300的四位数？
\end{example}

\begin{solution}
    这里的四位数显然是包括重复数字的四位数.

第一类：3在千位上的四位数一定大于2300，那么，百、十、个位上的数字每次都有4种选取可能，所以应有$4^3$个四位数.

第二类：4在千位上的四位数，同理也有$4^3$个四位数。

第三类：2在千位上，3在百位上的四位数，由于十位和个位上的数都可以从1, 2, 3, 4中选取，都大于2300，所以，这样的数有$4^2$个．

第四类，2在千位上，4在百位上的四位数有$4^2$个．所以
\[2\cdot 4^3+2\cdot 4^2=2\cdot 4^2(4+1)=160\]

答：这样的四位数有160个．
\end{solution}

\begin{ex}
    \begin{enumerate}
\item 有三封不同的信，投入4个信箱里，一共有多少种投信的方法？
\item 有数学、天文、诗歌三个课外小组，8个同学准备报名，若每人限报一个组，问一共可有多少种报名的方法？
    \end{enumerate}
\end{ex}

\subsection{重复组合}
从$n$个不同元素中取出$m$个元素组成一组，如果允许元素在组合中重复出现，这样的组合简称为重复组合，这种重复组合用符号$\rep{m}{n}$表示.

怎样计算重复组合的组合数呢？我们还是先从具体例子开始考虑，假如有三个不同的元素$a_1,a_2,a_3$，从其中取出两个元素组成重复组合，那么，可以得到如下的六个不同组合：
\begin{equation}
a_1a_1,\; a_1a_2,\; a_1a_3,\; a_2a_2,\; a_2a_3,\; a_3a_3 \tag{1}
\end{equation}

如果是不允许元素重复的话，要选两个元素得到不同的组合数是6的话，应是$\comb{2}{4}$，设这四个元素为$a'_1,a'_2,a'_3,a'_4$,那么得到的组合为：
\begin{equation}
a'_1a'_2,\; a'_1a'_3,\; a'_1a'_4,\; a'_2a'_3,\; a'_2a'_4,\; a'_3a'_4 \tag{2}
\end{equation}
组合（1）与组合（2）不难建立它们之间的一一对应关系．在组合（1）的每个组合的元素下标上分别加上0和1，就得到了下标与（2）完全一致的组合（3）．
\begin{equation}
a_{1+0}a_{1+1},\quad a_{1+0}a_{2+1},\quad a_{1+0}a_{3+1},\quad a_{2+0}a_{2+1},\quad a_{2+0}a_{3+1},\quad a_{3+0}a_{3+1}
\tag{3}
\end{equation}

又例如从4个不同元素中取出3个元素的重复组合有如下的不同组合:
\begin{center}
\begin{tabular}{ccccc}
    $a_{1}a_{1}a_{1}$&$a_{1}a_{1}a_{2}$&$a_{1}a_1a_{3}$&$a_{1}a_{1}a_{4}$\\
    &$a_{1}a_{2}a_{2}$&$a_{1}a_{2}a_{3}$&$a_{1}a_{2}a_{4}$\\
    &$a_{1}a_{3}a_{3}$&$a_{1}a_{3}a_{4}$&$a_{1}a_{4}a_{4}$\\
    &$a_{2}a_{2}a_{2}$&$a_{2}a_{2}a_{3}$&$a_{2}a_{2}a_{4}$\\
    &$a_{2}a_{3}a_{3}$&$a_{2}a_{3}a_{4}$&$a_{2}a_{4}a_{4}$\\
 &$a_{3}a_{3}a_{3}$&$a_{3}a_{3}a_{4}$&$a_{3}a_{4}a_{4}$&$a_{4}a_{4}a_{4}$\\
\end{tabular}
\end{center}

只要我们在每组元素的下标上分别加0, 1, 2, 就得到
\begin{center}
\begin{tabular}{ccccc}
    $a_{1}a_{2}a_{3}$&$a_{1}a_{2}a_{4}$&$a_{1}a_{2}a_{5}$&$a_{1}a_{2}a_{6}$\\
 & $a_{1}a_{3}a_{4}$&$a_{1}a_{3}a_{5}$&$a_{1}a_{3}a_{6}$\\
 &$a_{1}a_{4}a_{5}$&$a_{1}a_{4}a_{6}$&$a_{1}a_{5}a_{6}$\\
 &$a_{2}a_{3}a_{4}$&$a_{2}a_{3}a_{5}$&$a_{2}a_{3}a_{6}$\\
 &$a_{2}a_{4}a_{5}$&$a_{2}a_{4}a_{6}$&$a_{2}a_{5}a_{6}$\\
 &$a_{3}a_{4}a_{5}$&$a_{3}a_{4}a_{6}$&$a_{3}a_{5}a_{6}$&$a_{4}a_{5}a_{6}$
\end{tabular}
\end{center}

这又正是从六个不同元素中取出3个元素的组合的全部。
从上面的例子说明了
\[{\rm H}_{3}^{2}= {\rm C}_{3+ ( 2- 1) }^{2}, \qquad {\rm H}_{4}^{3}={\rm C}_{4+ ( 3- 1) }^3\]

按同样的方法我们可以得出：从$n$个不同元素中取出$m$
个的重复组合数有如下的关系：
$${\rm H}_{n}^{m}= {\rm C}^m_{n+ ( m- 1) }$$

现在我们来证明这个结论.

\begin{blk}
 {定理6} 从$n$个不同元素中，任取$m$个元素 的重复组合为：
 $${\rm H}_{n}^{m}={\rm C}_{n+ ( m-1) }^{m}$$ 
\end{blk}

\begin{proof}
    为了研究问题方便，我们用$1,2,3,\ldots,n$表示这$n$个不同的元素。在其中任取$m$个允许重复的数字作为一个组合，并按大小顺序排成
$$1\le i_1\le i_2\le\cdots\le i_m\le n$$

    又记$j_1=i_1+0,\; j_{2}=i_{2}+1,\;\ldots, j_{m}=i_{m}+(m-1)$
    
    自然，$1\le j_1<j_2<j_3<\cdots<j_m\leq n+(m-1)$，所以$\{j_1,j_2,\ldots,j_m\}$ 是$n+ m- 1$个 数，并 且 是$1, 2, \ldots , n+m-1$中的$m$个不同数的组合.
\end{proof}

下面来证明：在$n$个不同元素中任取$m$个允许重复的元
素的组合，和$n+m-1$个不同元素中任取$m$个不同元素 的组合之间有一个到上的一一对应关系。如果这点证明了，那么$n$个不同元素的$m$个允许重复元素的组合数就与$n+m-1$个不同元素的$m$个不同元素的组合数相等了.

\begin{proof}
    
为了证明这个一一对应，要先证明当允许重复的组合
$\{i_1,i_2,\ldots,i_m\}$与$\{i_1^{\prime},i_2^{\prime},\ldots,i'_m\}$不同 的时候，
（这里取$1\le i_1\le i_2\le\cdots\le i_m\le n$, $1\le i'_1\le i'_2\le\cdots\le i^{\prime}_m\le n$）那么，不同元素不允许重复的组合$\{j_1,j_2,\ldots, j_m\}$与$\{j_1^{\prime},j_2^{\prime},\ldots,j'_m\}$也不同，这里$\{j_1=i_1+0,\; j_2=i_{2}+1,\ldots,j_{m}=i_{m}+m-1,\; j_{1}^{\prime}=i_{1}^{\prime}+0,\; j_{2}^{\prime}=i_{2}^{\prime}+1,\ldots, j_m^{\prime}=i_m^{\prime}+m-1\}$.

事实上，如果后者相同，即有$j_{k}= j'_{k}\; (k= 1, 2, \ldots, m)$ 即
$$i_k+k-1=i_k^{\prime}+k-1\; (k=1,2,\ldots,m)$$ 
所以$i_k= i_k^{\prime }\; ( k= 1, 2, \ldots , m)$, 也就是说$\{i_1, i_2,\ldots,i_m\}$也与$\{i_1^{\prime},i_2^{\prime},\ldots,i_m^{\prime}\}$相同了。这就证明了上述的对应为一一对应.

再证明这是到上的一一对应，即当$\{i_1,i_2,\ldots,i_m\}$取遍$n$个数$1,2,\ldots,n$的一切$m$个数的允许重复组合时，$\{j_1, j_2,\ldots,j_m\}$ 也取遍$n+m-1$个 数 $1,2,\ldots ,n+m-1$一切的 $m$ 个数的不允许重复的组合.

事实上，对任一个组合$\{j_1,j_2,\ldots,j_m\}$, 其中$1\le j_{1}<j_{2}<\cdots<j_{m}\le n+m-1$. 有$1\leq j_{1},\; 2\le j_{2},\; 3\le j_{3},\ldots, m\le j_{m}$. 故$1\le i_1=j_1-0,\; 1\le i_2=j_2-1,\; 1\le i_3=j_3-2,\ldots, 1\leq i_m=j_m-(m-1)$. 即$i_1,i_2,\ldots,i_m$都是自然数，且由$j_{k}>j_{k-1}$, 可知$j_k-j_{k-1}\ge 1$, 故有
\[i_k-i_{k-1}=[j_k-(k-1)]-[j_{k-1}-(k-2)]=j_k-j_{k-1}-1\ge 0\]
再者$i_{m}=j_{m}-(m-1)\le (n+m-1)-(m-1)=n$. 总之有$1\le i_1\le i_2\le \cdots \le i_m\le n$. 所以，$\{i_1,i_2,\ldots,i_m\}$ 是$n$个数$1，2,\ldots,n$中取$m$个数的允许重复的组合，这证明了对应是到上的一一对应.

\end{proof}


\begin{example}
  从3, 5, 7, 9, 11中取出3个数（可以重复选取）相乘，问有多少个不同的乘积？  
\end{example}

\begin{solution}
    因为是乘积，与元素取出的先后顺序无关，同时是可以重复选取的，即是重复组合问题. 所以
\[{\rm H}^3_5=\comb{3}{5+(3-1)}=\comb{3}{7}=35\]
答：有35个不同的乘积．
\end{solution}

\begin{example}
    从三个不同的质数$a_1,a_2,a_3$中取出4个数相乘（可以重复选取），问有多少个不同的乘积？
\end{example}

\begin{solution}
    同例2.16是重复组合问题．所以
    \[{\rm H}^4_3=\comb{4}{3+(4-1)}=\comb{4}{6}=\comb{2}{6}=15\]
答：有15个不同的乘积。
\end{solution}

\begin{example}
    用七种不同的颜色给一个四面体的四个面涂上颜色，规定每个面上只能涂一种颜色，如果四个面的颜色容许相同，问有几种不同的着色方法？
\end{example}

\begin{solution}
    因为只研究对四个面所着的颜色，对那个面先着色并无先后顺序之分，故是组合问题.又因为对四个面所着的色是允许相同的，所以又是重复组合问题. 所以
    \[{\rm H}^4_7=\comb{4}{7+(4-1)}=\comb{4}{10}=210\]
    答：共有210种不同的着色方法。
\end{solution}


\begin{ex}
\begin{enumerate}
\item 展开$(a+b+c+d)^3$，合并同类项后，一共有多少项？
\item 若把六种不同的溶液注满四个烧杯，但不使它们混合，问一共有多少种方法？
\end{enumerate}
\end{ex}

\section*{习题2.2}
\begin{enumerate}
\item 从8个不同的元素中取出5个元素排成圆圈，问有几种不同的排列方法？
\item 两个教师和五个学生代表围坐在圆桌周围开座谈会：
\begin{enumerate}[(1)]
\item 如果不加条件限制，有多少种不同的坐法？
\item 如果两个教师靠在一起，有几种坐法？
\item 如果两个教师不靠在一起，有几种坐法？
\item 学生代表甲坐在教师的中间，有几种坐法？
\end{enumerate}
\item 从三个不同的元素中取出4个元素的重复组合有多少种不同的方法？
\item 用七种颜色涂在正方体的六个面上，如果每个面只能涂一种颜色，如果六个面上的颜色允许相同，问共有几种不同的着色方法？
\item 从2, 3, 5, 7, 11, 13中至少取出两个，至多可以取出五个数相乘（允许因数重复选取），问有多少个不同的乘积？
\item 展开$(a+b+c+d+e+f)^4$，再合并同类项，一共可以得到多少项？
\end{enumerate}

\section{二项式定理}

\subsection{和号$\Sigma$}

我们经常要计算下面各种和
$$1+ 2+ 3+ \cdots + n;\quad 1^{2}+ 2^{2}+ 3^{2}+ \cdots + n^{2},\quad \frac{1\times2}{2}+\frac{2\times3}{2^{2}}+\cdots+\frac{n(n+1)}{2^{n}}$$
等等。 一 般 给 出 一 个 有 顺 序 的 数 列 $a_1, a_2, \ldots , a_n$, 要计算
$a_{1}+a_{2}+\cdots+a_{n}$，
为了方便起见引进和符号$\Sigma$，于是可记
$$1+2+3+\cdots+n=\sum_{k=1}^{n}k$$
\[1^{2}+ 2^{2}+ 3^{2}+ \cdots + n^{2}=\sum^n_{k=1}k^2\]
\[\frac{1\times2}{2}+\frac{2\times3}{2^{2}}+\cdots+\frac{n(n+1)}{2^{n}}=\sum^n_{k=1}\frac{k(k+1)}{2^k}\]

一般可记
$a_1+a_2+\cdots+a_n=\sum_{k=1}^{n}a_k$
符号
$$\sum_{k=1}^{n}a_k$$

表示$k$取$1,2,\ldots,n-1,n$时，得到的$a_k$全部加起来.
例如：
$$\sum_{k=1}^{n}\sin kx=\sin x+\sin2x+\sin3x+\cdots+\sin nx$$
$$\sum _{k= 0}^{n}a_{k} x^{n- k}= a_{0} x^{n}+ a_{1} x^{n- 1}+ \cdots + a_{n- 1} x+ a_{n}$$
又例如，对给定的函数$f_1(x),f_2(x),\ldots,f_n(x)$
有
\[\begin{split}
\sum^n_{k=1}f_k(x)&=f_1(x)+f_2(x)+\cdots +f_n(x)
\\
\sum^n_{k=0}\comb{k}{n} &=\comb{0}{n}+\comb{1}{n}+\comb{2}{n}+\cdots +\comb{n}{n}
\end{split}\]
等等.

关于求和符号还有两点说明，首先，凡是指示指标从1到$n$，且仅仅指示从1到$n$．所以可以用其它的字母来代替$k$而和号仍然是表示了同一个意思. 例如
\[\sum^n_{k=1}a_k=a_1+a_2+\cdots+a_n,\qquad \sum^n_{\ell=1}a_{\ell}=a_1+a_2+\cdots+a_n\]

一般$\sum\limits^n_{*=1}a_*=a_1+a_2+\cdots+a_n$
这里$*$可以用任何一种指标来表示，其次，求和也不限于从1到$n$，也可以从0到$n+2$等等. 例如
\[\begin{split}
    \sum^{n+2}_{k=0}a_k&=a_0+a_1+\cdots+a_{n+1}+a_{n+2}\\
    \sum^{n+1}_{k=3}a_k&=a_3+a_4+\cdots+a_{n+1}
\end{split}\]

\begin{ex}
\begin{enumerate}
    \item 把下列各式用和号表示：
\begin{enumerate}[(1)]
\item $1+3+5+7\cdots+(2n-1)$ 
\item $1\cdot2+2\cdot3+3\cdot4+\cdots+n(n+1)$ 
\item $\cos x+\cos2x+\cos3x+\cdots+\cos nx$
\item $1^{2}+2^{2}+3^{2}+\cdots+(n+2)^{2}$
\end{enumerate}

\item 把下列和号式改写为一般的和式
\begin{multicols}{2}
\begin{enumerate}[(1)]
\item $\sum\limits _{k= 1}^{n}\frac 1{k( 2k+ 1) }$ 
\item $\sum\limits _{k= 1}^{n+ 2}\frac {k^{2}}{k! }$
\item $\sum\limits^n_{k=2}\frac{\sin kx}{k}$
\item $\sum\limits^n_{k=2}\frac{(-1)^k}{k^2+k-2}$
\end{enumerate}
\end{multicols}
\end{enumerate}
\end{ex}

\subsection{二项式定理}
下面我们用数学归纳法来证明二项式定理。我们已经知道：
\[\begin{split}
(x+y)^1&=x+y\\
(x+y)^2&=x^2+2xy+y^2\\
(x+y)^3&=x^3+3x^2y+3xy^2+y^3    
\end{split}\]

现在考虑二项式$(x+y)^n$的展开式，这里$n$是任意自然数，首先研究$(x+y)^4$
\[(x+y)^4=(x+y)(x+y)(x+y)(x+y)\]
等号右边的积应有下面形式的各项：
\[x^4,\; x^3y,\; x^2y^2,\; xy^3,\; y^4\]
显然$x^4$, $y^4$的系数是1，其他三项的系数，例如$x^2y^2$的系数是从四个括弧中两个里取$y$，从其他两个里取$x$作出的积的个
数，从四个括弧中任取两个$y$的取法 有$\comb{2}{4}=\frac{4\cdot 3}{1\cdot 2}=6$种，因
此$x^2y^2$的系数是6，同样，可以求得$x^3y,xy^3$的系数分别
是$\comb{1}{4}=4$，因为规定$\comb{0}{n}=1$，又知$\comb{4}{4}=1$，所以
\[(x+y)^4=\comb{0}{4}x^4+\comb{1}{4}x^3y+\comb{2}{4}x^2y^2+\comb{3}{4}xy^3+\comb{4}{4}y^4\]

一般地，可以归纳出以下的定理：
\begin{blk}
    {定理} 对于两个不定元$x$，$y$的多项式$(x+y)^n$，有下列展开式：
\[\begin{split}
    (x+y)^n&=\sum^n_{k=0}\comb{k}{n}x^{n-k}y^k\\
    &=\comb{0}{n}x^n+\comb{1}{n}x^{n-1}y+\comb{2}{n}x^{n-2}y^2+\cdots+\comb{n}{n}y^n
\end{split}\]
\end{blk}

\begin{proof}
    用数学归纳法来证明定理.
\begin{enumerate}[(1)]
    \item 当$n=1$时，$(x+y)^1=x+y$
\[\sum^1_{k=0}\comb{k}{1}x^{1-k}y^k=\comb{0}{1}x^1y^0+\comb{1}{1}x^{0}y^1=x+y\]
故$n=1$时原命题成立.
\item 假设$n=\ell$时有公式 $(x+y)^{\ell}=\sum\limits^\ell_{k=0}\comb{k}{\ell}x^{\ell-k}y^k$

现在来计算$(x+y)^{\ell+1}$，令
\[\begin{split}
    (x+y)^{\ell+1}&=(x+y)\cdot (x+y)^{\ell}=(x+y)\cdot \sum^\ell_{k=0}\comb{k}{\ell}x^{\ell-k}y^k\\
&=\sum^\ell_{k=0}\comb{k}{\ell}x^{\ell-k+1}y^k +\sum^\ell_{k=0}\comb{k}{\ell}x^{\ell-k}y^{k+1}\\
&=\comb{0}{\ell}x^{\ell+1}+\sum^\ell_{k=1}\comb{k}{\ell}x^{1+\ell-k}y^k+\sum^{\ell-1}_{k=0}\comb{k}{\ell}x^{\ell-k}y^{k+1}+\comb{\ell}{\ell}y^{\ell+1}\\
&=\comb{0}{\ell+1}x^{\ell+1}+\sum^\ell_{k=1}\comb{k}{\ell}x^{1+\ell-k}y^k+\sum^{\ell}_{k=1}\comb{k-1}{\ell}x^{\ell+1-k}y^{k}+\comb{\ell+1}{\ell+1}y^{\ell+1}\\
&=\comb{0}{\ell+1}x^{\ell+1}+\sum^\ell_{k=1}\left(\comb{k}{\ell}+\comb{k-1}{\ell}\right)x^{\ell+1-k}y^k+\comb{\ell+1}{\ell+1}y^{\ell+1}\\
&=\comb{0}{\ell+1}x^{\ell+1}+\sum^\ell_{k=1}\comb{k}{\ell+1}x^{\ell+1-k}y^k +\comb{\ell+1}{\ell+1}y^{\ell+1}\\
&=\sum^{\ell+1}_{k=0}\comb{k}{\ell+1}x^{\ell+1-k}y^k 
\end{split} \]
所以$n=\ell+1$时命题成立.
\end{enumerate}
根据（1）和（2），对于一切自然数$n$都有：$(x+y)^n=\sum\limits^n_{k=0}\comb{k}{n}x^{n-k}y^k$成立.
\end{proof}

这个定理叫做二项式定理，右边的多项式$\sum\limits^n_{k=0}\comb{k}{n}x^{n-k}y^k$
叫做$(x+y)^n$的二项展开式. 式中的$\comb{k}{n}x^{n-k}y^k$叫做二项展开式的通项，用${\rm T}_{k+1}$表示，即通项为展开式的第$k+1$项.
\[{\rm T}_{k+1}=\comb{k}{n}x^{n-k}y^k\]

由二项式定理，取$y=1$，有
\[(x+1)^n=\sum^n_{k=0}\comb{k}{n}x^{n-k}\]
取$y=-1$，有$$(x-1)^n=\sum\limits^n_{k=0}(-1)^k\comb{k}{n}x^{n-k}$$

在二项式定理中，如果遇到$n$是较小的正整数时展开式的系数，也可以直接用下表计算，
\begin{center}
\begin{tikzpicture}
\foreach \x in {0,1,2,...,6}
{
    \node at (0,-\x){$(x+y)^{\x}$};
}
\foreach \x/\y in {3/1, 4/6, 5/15, 6/20, 7/15, 8/6, 9/1}
{
    \node (A\x) at (\x,-6){\y};
}
\foreach \x/\y in {3/1, 4/5, 5/10, 6/10, 7/5, 8/1}
{
    \node (B\x) at (\x+.5,-5){\y};
}
\foreach \x/\y in {3/1, 4/4, 5/6, 6/4, 7/1}
{
    \node (C\x) at (\x+1,-4){\y};
}
\foreach \x/\y in {3/1, 4/3, 5/3, 6/1}
{
    \node (D\x) at (\x+1.5,-3){\y};
}
\foreach \x/\y in {3/1, 4/2, 5/1}
{
    \node (E\x) at (\x+2,-2){\y};
}
\foreach \x/\y in {3/1, 4/1}
{
    \node (F\x) at (\x+2.5,-1){\y};
}
\node (G) at (6,0){1};

\draw(F3)--(E4)--(F4);
\draw(E3)--(D4)--(E4)--(D5)--(E5);
\draw(D3)--(C4)--(D4)--(C5)--(D5)--(C6)--(D6);
\draw(C3)--(B4)--(C4)--(B5)--(C5)--(B6)--(C6)--(B7)--(C7);
\draw(B3)--(A4)--(B4)--(A5)--(B5)--(A6)--(B6)--(A7)--(B7)--(A8)--(B8);

\draw[dotted](1,0)--(G);
\draw[dotted](1,-1)--(F3);
\draw[dotted](1,-2)--(E3);
\draw[dotted](1,-3)--(D3);
\draw[dotted](1,-4)--(C3);
\draw[dotted](1,-5)--(B3);
\draw[dotted](1,-6)--(A3);


\end{tikzpicture}
\end{center}
表中每行两端都是1，而且除1以外的每一个数都等于它肩上两个数的和，这个表叫做杨辉三角，它首先载于我国宋朝数学家杨辉1261年所著的《详解九章算法》一书．

\begin{example}
展开$\left(1{+\frac{1}{x}}\right)^5$
\end{example}


\begin{solution}
\[\begin{split}
    \left(1{+\frac{1}{x}}\right)^5 &=1+5\left(\frac{1}{x}\right)+10\left(\frac{1}{x}\right)^2+10\left(\frac{1}{x}\right)^3+5\left(\frac{1}{x}\right)^4+\left(\frac{1}{x}\right)^5\\
    &=1+\frac{5}{x}+\frac{10}{x^2}+\frac{10}{x^3}+\frac{5}{x^4}+\frac{1}{x^5}
\end{split}\]
\end{solution}

\begin{example}
    展开$\left(2\sqrt{x}-\frac{1}{\sqrt{x}}\right)^6$
\end{example}

\begin{solution}
\[\begin{split}
    \left(2\sqrt{x}-\frac{1}{\sqrt{x}}\right)^6&=\left(\frac{2x-1}{\sqrt{x}}\right)^6 =\frac{1}{x^3}(2x-1)^6\\
&=\frac{1}{x^3}\left[(2x)^6-\comb{1}{6}(2x)^5+\comb{2}{6}(2x)^4-\comb{3}{6}(2x)^3+\comb{4}{6}(2x)^2\right.\\
&\qquad \qquad \left.-\comb{5}{6}(2x)+\comb{6}{6}\right]\\
&=\frac{1}{x^3}\left[64x^6-6\cdot 32x^5+15\cdot 16x^4-20\cdot 8 x^3+15\cdot 4x^2\right.\\
&\qquad \qquad \left.-6\cdot 2x+1\right]\\
&=64x^3-192x^2+240x-160+\frac{60}{x}-\frac{12}{x^2}+\frac{1}{x^3}
\end{split}\]
\end{solution}

\begin{example}
    求$(x+a)^{12}$的展开式中的第四项和倒数第四项。
\end{example}

\begin{solution}
    展开式的通项公式$\mathrm{T}_{k+1}=\mathrm{C}_n^k x^{n-k}y^k$。所以
$$\mathrm{T}_{4}=\mathrm{C}_{12}^{3} x^{12-3} a^{3}=\frac{12\times11\times 10}{3\times2}x^{9} a^{3}=220x^{9} a^{3}$$
又展开式共$12+1=13$, 故倒数第四项是
$$\mathrm{T}_{10}=\mathrm{C}_{12}^{9} x^{12-9} a^{9}=\mathrm{C}_{12}^{3}x^{3} a^{9}=220x^{3} a^{9}$$
\end{solution}

\begin{example}
  求$(x-\frac{1}{x})^{9}$展开式中$x^3$的系数，有没有$x^{0}$项，如
果有，试求出该项来。  
\end{example}

\begin{solution}
    展开式的通项是：
$$T_{k+1} = ( -1 )^{k} C_{9}^{k} x^{9-k} (\frac{1}{x})^{k}=( -1)^{k} C_{9}^{k} x^{9-k-k}$$
依题意有，$9-2k=3$, 所以$k=3$. 因此$x^3$的项系数是
$$( -1 )^{3}{\rm C}_{9}^{3}=- \frac{9\times8\times7}{3\times2}=- 84$$

又若$9-2k=0$, 即$2k=9$, 方程没有整数解，故展开式
里没有$x^0$项.
\end{solution}

\begin{example}
   在$\left(a+\frac1a\right)^{2n}$的展开式里，已知第四项和第六项
系数相等，求展开式里不含有$a$的项. 
\end{example}

\begin{solution}
因为
\[\begin{split}
    {\rm T}_4&=\comb{3}{2n}a^{2n-3}\frac{1}{a^3}=\comb{3}{2n}a^{2n-6}\\
    {\rm T}_6&=\comb{5}{2n}a^{2n-5}\frac{1}{a^5}=\comb{5}{2n}a^{2n-10}
\end{split}\]
依题意得$\comb{3}{2n}=\comb{5}{2n}$，即$\comb{2n-3}{2n}=\comb{5}{2n}$，所以
\[2n-3=5,\quad 2n=8,\quad n=4\]
故所给的二项式是$\left(a+\frac{1}{a}\right)^8$

因为${\rm T}_{k+1}=\comb{k}{8}a^{8-k}\frac{1}{a^k}=\comb{k}{8}a^{8-2k}$，
若${\rm T}_{k+1}$为不含$a$项，则$8-2k=0$，所以$k=4$.
\[{\rm T}_5=\comb{4}{8}\cdot a^0=\frac{8\x 7\x 6\x 5}{4\x 3\x 2}=70\]
\end{solution}

\begin{ex}
\begin{enumerate}
    \item 展开$(x^{3}-2x)^{7}$
    \item 求$(a+\sqrt{b})^{12}$的展开式里的第4项和第9项
    \item \begin{enumerate}[(1)]
        \item 求$(x+a)^{12}$的展开式里的中间一项；
        \item 求$\left(\frac x2+\frac2x\right)^7$的展开式里的中间两项.
    \end{enumerate}  
    \item 在$(x+a)^{10}$的展开式里，求
\begin{multicols}{2}
    \begin{enumerate}[(1)]
        \item 含有$a^{8}$的项；
        \item 含有$x^8$的项. 
    \end{enumerate}    
\end{multicols}

    \item 在$\left(x^{2}+\frac{1}{x^3}\right)^{15}$的展开式里，有没有
\begin{multicols}{2}
    \begin{enumerate}
        \item 不含$x$的项？
        \item 含有$x^{8}$的项？
    \end{enumerate}
\end{multicols}
如果有，把这项求出来.
    
\item 求$\left(3x+\frac{1}{\sqrt[3]{x}}\right)^{8}$的展开式里的常数项. 
\item 已知$\left(\sqrt[8]{a^{2}}+\frac1a\right)^{n}$的展开式里的第3项含有$a^2$,    
    求$n$.
\end{enumerate}
\end{ex}

\subsection{二项式的性质}

我们研究二项式$(x+y)^n$的展开式中各项系数的性质。二项展开式各项系数是：
\[\comb{0}{n},\;\comb{1}{n},\;\comb{2}{n},\;\ldots,\; \comb{n-1}{n},\;\comb{n}{n}\]

展开式各项系数有下列性质
\begin{enumerate}
\item 在二项展开式中，与首末两端“等距离”的两项系数相等. 我们已知$\comb{m}{n}=\comb{n-m}{n}$，所以:
\[\comb{0}{n}=\comb{n}{n},\; \comb{1}{n}=\comb{n-1}{n},\; \comb{2}{n}=\comb{n-2}{n},\; \ldots,\; \comb{k}{n}=\comb{n-k}{n},\ldots\]
\item 如果二项式的幂指数是偶数，中间一项的系数大，如果二项式的幂指数是奇数，中间两项系数相同并且大.
\end{enumerate}

由于展开式各项系数顺次是
\[\comb{0}{n}=1,\quad \comb{1}{n}=n,\quad \comb{2}{n}=\frac{n(n-1)}{1\cdot 2},\ldots\]
\[\comb{k}{n}=\frac{n(n-1)(n-2)\cdots(n-k+1)}{1\cdot 2\cdot 3\cdots k},\ldots,\quad \comb{n}{n}=1\]
其中，后一个系数的分子是前一个系数的分子乘以逐次减小1的数，(如$n,\; n-1,\; n-2,\; \ldots$), 分母是乘以逐次增大1的数(如$1,2,3,\ldots$)得到的，因而，各项系数从开始起总是逐渐增大，又因为与首末两端“等距离”的两项系数相等，所以系数增大到某一项时就逐渐减小，而最大项必在中间.

当$n$是偶数时，$n+1$是奇数，展开式共有$n+1$项，所以
展开式有中间一项，即第($\frac{n}{2}+1$)项，并且系数最大。

当$n$是奇数时，$n+1$是偶数，展开式共有$n+1$项，所以有中间两项，即第($\frac{n+1}2$)和第($\frac{n+1}2+1$)两项，系数相同并且最大.

\begin{example}
    求$(x+y)^{10}$和$(x+y)^{11}$的最大系数的项.
\end{example}

\begin{solution}
    $(x+y)^{10}$最大系数项的项数是
$\frac{10}{2}+1=6$，所以
$$\mathrm{T}_{6}=\mathrm{C}_{10}^{5} x^{10-5}y^{5}=252x^{5}y^{5}$$

$(x+y)^{11}$最大系数的项 数是$\frac{11+1}{2}=6$或$6+1=7$, 所以
\[\begin{split}
    \mathrm{T}_{6}&=\mathrm{C}_{11}^{5}x^{11-5}y^{5}=462x^{6}y^{5}\\
    \mathrm{T}_{7}&=\mathrm{C}_{11}^{6}x^{11-6}y^{6}=462x^{5}y^{6}
\end{split}\]
\end{solution}

\begin{example}
    在$(5x+2y)^{29}$展开式中第几项的系数最大，并
求它的系数。
\end{example}

\begin{solution}
    设$T_{k+1}$系数最大，则它必须满足
\[\begin{cases}
\frac{{\rm T}_{k+1}\text{系数}}{{\rm T}_{k}\text{系数}}=\frac{\comb{k}{29}5^{29-k}2^k}{\comb{k-1}{29}5^{29-k+1}2^{k-1}}=\frac{2(29-k+1)}{5k}\ge 1\\
\frac{{\rm T}_{k+1}\text{系数}}{{\rm T}_{k+2}\text{系数}}=\frac{\comb{k}{29}5^{29-k}2^k}{\comb{k+1}{29}5^{29-k-1}2^{k+1}}=\frac{5(k+1)}{2(29-k)}\ge 1\\
\end{cases}\]
即
$\begin{cases}
    7k\le 60\\
    \frac{53}{7}\le k\le 29
\end{cases}$

解得$7\frac{4}{7}\le k\le 8\frac{4}{7}$. 因为$k$为正整数，故$k=8$，即系数最大项系数是
\[{\rm T}_{8+1}=\comb{8}{29}5^{29-8}\cdot 2^8=\comb{8}{29}5^{21}\cdot 2^8\]
\end{solution}

\begin{example}
    求证：$\sum\limits^n_{k=0}\comb{k}{n}=\comb{0}{n}+\comb{1}{n}+\comb{2}{n}+\cdots+\comb{n-1}{n}+\comb{n}{n}=2^n$
\end{example}

\begin{proof}
在二项式公式中取$x=y=1$，有：
\[(1+1)^n=\sum^n_{k=0}\comb{k}{n}1^{n-k}1^k=\sum^n_{k=0}\comb{k}{n}\]
所以$\sum\limits^n_{k=0}\comb{k}{n}=2^n$
\end{proof}

\begin{example}
证明在$(x+y)^n$展开式中，奇数项系数的和等于
偶数项系数的和.
\end{example}

\begin{proof}
    在二项式公式中取$x=1$, $y=-1$，有
    \[(1-1)^n=\sum^n_{k=0}(-1)^k\comb{k}{n}1^{n-k}1^k=\sum^n_{k=0}(-1)^k\comb{k}{n}\]
即
\[0=\left(\comb{0}{n}+\comb{2}{n}+\cdots\right)-\left(\comb{1}{n}+\comb{3}{n}+\cdots\right)\]
所以
\[\comb{0}{n}+\comb{2}{n}+\cdots=\comb{1}{n}+\comb{3}{n}+\cdots\]
\end{proof}

\begin{example}
求证：$\sum\limits^n_{k=1}2k\comb{k}{n}=n\cdot 2^n$
\end{example}

\begin{proof}
因为
\[\begin{split}
    k\comb{k}{n}&=k\cdot \frac{n(n-1)\cdots (n-k+1)}{k!}\\
    &=n\cdot \frac{(n-1)(n-2)\cdots (n-k+1)}{(k-1)!}=n\comb{k-1}{n-1}
\end{split}\]
所以
\[\sum^n_{k=1}2k\comb{k}{n}=2\sum^n_{k=1}k\comb{k}{n}=2n\sum^{n-1}_{k=0}\comb{k}{n-1}=2n\cdot 2^{n-1}=n\cdot 2^n\]
\end{proof}


\begin{example}
    计算$(1.009)^{20}$和$(0.991)^{5}$的近似值（精确到0.001)
\end{example}

\begin{solution}
\[\begin{split}
    (1.009)^{20}&=(1+0.009)^{20}\\
&=1+20\x0.009+190\x(0.009)^2+1140\x(0.009)^3+\cdots\\
&=1+0.18+0.01539+0.00083106+\cdots
\end{split}\]
可以看出，要使所求的近似值的误差不超过0.001，需取展开式中前三项的和，就是
\[\begin{split}
(1.009)^{20}&\approx 1+0.18+0.01539\approx 1.195\\
(0.991)^5&=(1-0.009)^5=1-5\x0.009+10\x(0.009)^2+\cdots    
\end{split}\]
根据精确度要求，第三项以后各项都可以删去，所以，
\[(0.991)^5\approx 1-0.045=0.955\]
\end{solution}

\begin{ex}
\begin{enumerate}
    \item 求二项展开式里系数最大的项，（系数只需用组合数符号表示，不必算出）：
\begin{multicols}{2}
\begin{enumerate}[(1)]
    \item $(a+b)^{10}$
    \item $(1+x)^{17}$
\end{enumerate}
\end{multicols}
    \item  应用二项式定理展开$(x^2+2x+2)^5$
    \item  $(x+y)^{2n}$的展开式里如果第5项的系数与第13项的系数相等，求展开式里系数最大的项（系数用组合符号表示，不必算出）.
    \item  求$(1-x)^{13}$的奇次项系数的和.
    \item 证明：$\comb{0}{n}+\comb{2}{n}+\comb{4}{n}+\cdots+\comb{n}{n}=2^{n-1}$（$n$为偶数）
    \item 计算$(3.002)^0$和$(3.998)^5$误差不超过0.001的近似值.
\end{enumerate}
\end{ex}

\section*{习题2.3}
\begin{enumerate}
    \item 用二项式定理展开。
\begin{multicols}{2}
\begin{enumerate}[(1)]
    \item $\left( \frac {\sqrt {x}}2- \frac 2{\sqrt {x}}\right) ^{7}$
    \item $\left( \frac {x^{2}}2+ \frac {3}{x}\right) ^{8}$
\end{enumerate}    
\end{multicols}
\item 化简：
\begin{enumerate}[(1)]
    \item $(1+\sqrt{x})^{5}+(1-\sqrt{x})^{5}$
    \item $\left( 2x^{\tfrac 12}+ 3x^{- \tfrac 12}\right) ^{4}- \left( 2x^{\tfrac 12}- 3x^{- \tfrac 12}\right) ^{4}$
\end{enumerate}
\item \begin{enumerate}[(1)]
    \item 求$(1-2x)^{15}$的展开式中前四项；
    \item 求$(2a^{3}-2b^{2})^{10}$的展开式中第11项;
    \item 求$\left(\frac{\sqrt{x}}{3}+\frac{3}{\sqrt{x}}\right)^{12}$的展开式的中间一项;
\item 求$(x\sqrt{y}-y\sqrt{x})^{15}$的展开式的中间两项.
\end{enumerate}
\item 求下列各式的展开式中指定各项的系数。
\begin{enumerate}[(1)]
    \item $\left( 1- \frac 1{2x}\right)^{10}$含$\frac1{x^5}$的项；
    \item $\left(2x^{3}-\frac{1}{2x^{2}}\right)^{10}$的常数项.
\end{enumerate}
 
\item \begin{enumerate}[(1)]
    \item $\left({\sqrt {x}} + \frac 1{\sqrt [3]{x^{2}}}\right) ^{n}$的展开式里第5项系数与
第3项系数的比是7:2, 求展开式里含有$x$的项.
\item 已知$(1+x)^n$的展开式里第四项与第八项的系数相等，求这两项的系数.
\end{enumerate}

\item \begin{enumerate}[(1)]
    \item 求$\left(\frac{x}{y}-\frac{y}{x}\right)^{13}$展开式里系数最大的项.
    \item 写出$\left(x^2-\frac{1}{\sqrt{x}}\right)^{10}$的展开式里系数绝对值最大的项.
    \item 求$(1+x+x^2)(1-x)^5$的展开式里的中间项.
\end{enumerate}

\item 用二项式定理证明：
\begin{enumerate}[(1)]
    \item $(n+1)^{n}-1$能被$n^2$整除；
    \item $99^{10}-1$能被1000整除.
\end{enumerate}

\item 证明：
\begin{enumerate}[(1)]
    \item $\left(x-\frac{1}{x}\right)^{2n}$的展开式中常数项是$(-2)^n\frac{1\cdot 3\cdot 5\cdots (2n-1)}{n!}$
    \item $(1+x)^{2n}$的展开式的中间一项是$\frac{1\cdot 3\cdot 5\cdots (2n-1)}{n!}(2x)^n$
\end{enumerate}

\item 求证：
\begin{enumerate}[(1)]
    \item $\sum\limits^m_{k=0}\comb{n}{n+k}=\comb{n+1}{n+m+1}$
    \item $2^n-\comb{1}{n}2^{n-1}+\comb{2}{n}2^{n-2}+\cdots+(-1)^{n-1}\comb{n-1}{n}\cdot 2+(-1)^n=1$
    \item $\sum\limits^n_{k=0}(2k+1)\comb{k}{n}=(n+1)2^n$
    \item $1+2\comb{1}{n}+\comb{2}{n}+2\comb{3}{n}+\comb{4}{n}+2\comb{5}{n}+\cdots +2\comb{n-1}{n}+\comb{n}{n}=3\cdot 2^{n-1}$（其中$n$是偶数）
\end{enumerate}
\item 求$(1+x_1)(1+x_2)(1+x_3)\cdots(1+x_n)$的展开式里各项系数之和.
\end{enumerate}


\section*{本章内容要点}
一、本章主要内容是排列与组合，二项式定理；重点是不许重复的排列与组合，二项式定理及其性质.

二、两个基本原理——加法原理与乘法原理，是推导排列与组合数公式的基础，也是直接解决某些计数问题的方法. 应用中，要注意区别两个原理：

加法原理与分类有关，即如果完成一件事有$n$类相互独立的办法，那么求完成这件事情的方法种数，就要用加法原理；

乘法原理与分步有关，即如果完成一件事需要分$n$步完成，缺一不可，而每一步又各有若干方法，那么，求完成这件事的方法种数，就要用乘法原理。

三、排列与组合的区别是看研究的问题是否与元素选取的顺序有关，与顺序有关就属于排列；与顺序无关就属于组合.

排列又分为不许重复选取的全排列与选排列、可重复的排列、环形排列；

组合又分为不许重复的组合与可重复的组合。

四、排列与组合的计数公式有：
\begin{enumerate}
    \item 从$n$个不同元素中，任取$m$个没有重复的元素：

排列数公式
\begin{itemize}
    \item 选排列：$\perm{m}{n}=n(n-1)\cdots(n-m+1)=\frac{n!}{(n-m)!}\quad (m<n)$
    \item 全排列：$\perm{n}{n}=n!$
\end{itemize}

组合数公式：$\comb{m}{n}=\frac{\perm{m}{n}}{\perm{m}{m}}\quad (m\le n)$

组合数的性质：$\comb{m}{n}=\comb{n-m}{n},\qquad \comb{m}{n+1}=\comb{m}{n}+\comb{m-1}{n}\quad (m\le n)$

环形排列数公式：$\comb{m}{n}\cdot \perm{m-1}{m-1}=(m-1)! \comb{m}{n}\quad (m\le n)$

\item 可重复的排列数与组合数公式：
\begin{itemize}
    \item 重复排列数$n^m$
    \item 重复组合数${\rm H}^m_n=\comb{m}{n+m-1}$
\end{itemize}
\end{enumerate}

五、二项式定理指出了二项代数和的幂的展开公式。
\[(a+ b) ^{n}=\mathrm{C}_{n}^{0}a^{n}+{\rm C}_n^{1}a^{n- 1}b+ \cdots + \mathrm{C}_{n}^{r}a^{n- r}b^{r}+ \cdots + \mathrm{C} _{n}^{n}b^{n}\]
其中，展开式的通项(第$r+1$项)为：
$$\mathrm{T}_{r+1}=\mathrm{C}_{n}^{r}a^{n-r}b^{r}$$

二项式的展开式中，系数、幂指数、项数等都有很好的
性质，其中主要性质有：
\begin{enumerate}
    \item 
展开式中，与首末两端“等距”的两项系数相等，
二项式$n$次幂的展开式共有$n+1$项.
\item 
当 幂 指 数 $n$为 偶 数 时 , 中 间 一 项 的 系 数${\rm C}_n^{\tfrac{n}{2}}$最大(第$\frac n2+1$项); 当$n$为奇数时，中间有两项系数最大，即第$\frac{n+1}2$项的系数$C_{n}^{\tfrac{n-1}2}$与第$\frac{n+1}2+1$项的系数${\rm C}_{n}^{\tfrac{n+1}2}$相等且最大.
\end{enumerate}

\section*{复习题二}
\begin{enumerate}
    \item 解不等式：
\begin{multicols}{2}
\begin{enumerate}[(1)]
    \item $2<\frac{\perm{m+1}{m+1}}{\perm{m-1}{m-1}}\le 42$
    \item $x\perm{3}{8}<\perm{3}{x}$
\end{enumerate}
\end{multicols}

\item 求证：
\begin{enumerate}[(1)]
    \item $1!+2\cdot 2!+3\cdot 3!+\cdots+n\cdot n!=(n+1)!-1$
    
    （提示：$n\cdot n!=(n+1)!-n!$）

\item $\comb{m}{n-1}+\comb{m}{n-2}+\comb{m}{n-3}+\cdots+\comb{m}{m+1}+\comb{m}{m}=\comb{m+1}{n}$
\end{enumerate}

\item \begin{enumerate}[(1)]
    \item 已知$\frac{1}{\comb{m}{5}}-\frac{1}{\comb{m}{6}}=\frac{7}{10\cdot\comb{m}{7}}$，求$\comb{m}{8}$
    \item 已知$\frac{\comb{m-1}{n}}{2}=\frac{\comb{m}{n}}{3}=\frac{\comb{m+1}{n}}{4}$，求$n$与$m$.
\end{enumerate}

\item 用数字0, 1, 2, 3, 4, 5组成没有重复数字的数：
\begin{enumerate}[(1)]
\item 能够组成多少个六位奇数？
\item 能够组成多少个大于201345的自然数？ 
\end{enumerate}

\item 8个不同的元素排成一行：
\begin{enumerate}[(1)]
 \item 其中某2个元素要排在一起，有多少种排法？
\item 其中某2个元素不排在一起，有多少种排法？
\item 其中某4个元素要排在一起，另外4个元素也要排在一起，有多少种排法？  
\end{enumerate}

\item 书架上有4本不同的数学书，5本不同的物理书，3本不同的化学书，全部竖起排成一排，如果不使同类的书分开，一共有多少种排法？
\item 从1到9这九个数字每次取出5个数字：
\begin{enumerate}[(1)]
\item 可以组成多少个能被25整除的五位数？
\item 可以组成多少个小于98000的五位数？
\item 可以组成多少个五位偶数？
\item 可以组成多少个奇数位是奇数，偶数位是偶数的五位数？   
\end{enumerate}

\item \begin{enumerate}[(1)]
    \item 一个集合由8个不同的元素组成，这个集合中3个元素的子集有几个？
    \item 一个集合由5个不同的元素组成，其中含1个，2个，3个，4个元素的子集共有几个？
\end{enumerate}

\item \begin{enumerate}[(1)]
\item 平面内有$n$条直线，其中没有两条互相平行，也没有三条相交于一点，一共有多少个交点？
\item 空间有$n$个平面，其中没有两个互相平行，也没有三个相交于一直线，一共有多少条交线？
\item 平面有两组平行线，一组$m$条，另一组$n$条，这两组平行线相交，可以构成多少个平行四边形？
\item 空间有三组平行平面，第一组有$m$个，第二组有$n$个，第三组有$\ell$个，不同组的平面都相交，可以构成多少个平行六面体？
\end{enumerate}

\item 一个班有48名学生，分成六排，每排八人
\begin{enumerate}[(1)]
\item 问：有多少种不同坐法？
\item 若有2人必须坐在前排，3人必须在最后一排，有多少种不同的坐法？
\end{enumerate}

\item 由0至9这10个数字组成的没有重复数字的九位数字有多少个？其中能被3整除的有多少个？
\item 有$n+1$件不同的奖品，全部赠给$n$个先进学生代表．如果每人至少得到一件，有多少种不同的赠送方法？
\item \begin{enumerate}[(1)]
    \item 求$\left(9x-\frac{1}{\sqrt[3]{x}}\right)^{10}$展开式的常数项;
    \item 求$(1+x+x^2)(1-x)^{10}$展开式中$x^4$的系数;
    \item 求$(1-x)^{11}+x(1-2x)^{10}-x^2(1-3x)^9$展开式中$x^5$的系数;
    \item 求$（1+2x-3x^2)^6$展开式中$x^5$的系数;
    \item 求$(x^{-1}-1+x)^5$展开式里不含$x$的项.
\end{enumerate}
\item \begin{enumerate}[(1)]
    \item 用二项式定理证明$55^{55}+9$能被8整除；
    \item 用二项式定理求$89^{10}$除以88的余数．
\end{enumerate}

\item 证明：$(1+x)^{2n}$展开式中$x^n$的系数等于$(1+x)^{2n-1}$展开式中$x^n$的系数的2倍．
\item 已知$\left(\sqrt{x}+\frac{1}{\sqrt[3]{x}}\right)^n$展开式的系数之和比$(a+b)^{2n}$展开式的系数之和小240，求：
\begin{enumerate}[(1)]
    \item $\left(\sqrt{x}+\frac{1}{\sqrt[3]{x}}\right)^n$展开式的第三项：
    \item $(a+b)^{2n}$展开式的中间项.
\end{enumerate}

\item 已知$(1+x)^n$展开式里连续三项系数的比是3:8:14，求展开式里系数最大的项．
\item 已知$\left(\sqrt[3]{x^{-2}}+x^2\right)^{2n}$的展开式的系数之和比$\left(\sqrt[3]{x^{-2}}+x^2\right)^{n}$的展开式的系数之和大992，求$\left(\sqrt[3]{x^{-2}}-x^2\right)^{n}$的展开式里含$x^{\tfrac{4}{3}}$项的系数.

\item 证明：
\begin{enumerate}[(1)]
    \item $\sum\limits^n_{k=0}\left(\comb{k}{n}\right)^2=\frac{(2n)!}{n!\cdot n!}$
    \item $\comb{0}{n}\comb{k}{m}+\comb{1}{n}\comb{k-1}{m}+\cdots+\comb{k}{n}\comb{0}{m}=\comb{k}{n+m}$
    \item $\left(\comb{0}{n}-\comb{2}{n}+\comb{4}{n}-\cdots\right)^2+ \left(\comb{1}{n}-\comb{3}{n}+\comb{5}{n}-\cdots\right)^2=\comb{0}{n}+\comb{1}{n}+\comb{2}{n}+\cdots+\comb{n}{n}$
\end{enumerate}

\item 在$(a-3b)^{35}$的展开式中第几项的系数的绝对值最大？
\end{enumerate}


\chapter{概率论初步}

\section{随机事件与概率}

\subsection{随机事件及其概率}
在一定的条件下，可能发生也可能不发生的事件称为随机事件.

\begin{example}
    掷一枚硬币，“正面向上”这一事件可能发生，也可能不发生，因此，它是一随机事件。
\end{example}

\begin{example}
    对准靶发射一枪，“中靶”这一事件可能发生，也可能不发生，因此，它是一随机事件.
\end{example}

\begin{example}
    将1,2,3这三个数字随便排列，就得一个三位数，“这个三位数大于300”可能发生，也可能不发生，因此，它是一随机事件.
\end{example}


在现实世界中，有许多事件在一定的条件下，可能发生也可能不发生，这些都是随机事件. 通常，我们用大写的英文字母$A,B,C,\ldots$表示随机事件，概率就是反映随机事件出现的可能性大小的一个量；在同一条件下，有些事件发生的可能性比另一些更大些，它们相应的概率也就比另一些的大一些，所以随机事件的概率是客观存在的，问题是如何求出随机事件相应的概率. 为了讨论方便起见，习惯上把“必然事件”与“不可能事件”也看作随机事件。“必然事件”就是在给定条件下必然发生的事件，通常用字母$U$表示. “不可能事件”就是在给定条件下必然不发生的事件，通常用字母$V$表示．在例3.1中，“正面向上或反面向上”就是一个必然事件，“正面和反面都向上”就是一个不可能事件；在例3.3中，“三位数小于400”是一个必然事件，“三位数小于110”就是一个不可能事件．

为了讨论反映随机事件的可能性的大小的量——概率。我们进一步考察“投掷一枚硬币”、“从50只三级管中任抽一只检查”这样一类随机现象。例如，投掷一枚硬币，共有“出现正面”和“出现反面”两种试验结果，通常总认为硬币是“均匀的”，所以出现这两种结果的可能性是相同的。类似地，从50只三级管中任抽一只，共有50种不同的结果，而且只要检查员事先不带主观偏见，抽到任何一只的可能性都相同. 今后，我们称每一试验结果为基本事件，在具体问题中，弄清楚所有基本事件是十分重要的事，首先要注意到不同的基本事件是不能同时发生的；其次要注意到事件总是由若干个基本事件组成. 某一事件所包含的基本事件发生了，就说该事件发生了.

从这两个例子可以看到，这一类随机现象有两个共同的特征：
\begin{enumerate}[(1)]
    \item 有限性：上面每次试验只可能有有限种不同的结果，即有限个基本事件；
    \item 等可能性：由于其自然对称性，出现每一基本事件的可能性是相同的.
\end{enumerate}

这是一类最简单但却是常见的随机现象，我们称它为古典概率模型（简称古典概型），如何来描述这类事件的可能性的大小呢？例如我们问：“50只三级管中有2只次品，从
中任抽一只，抽得次品的可能性有多大？”人们常常用“次品率”来描述这种可能性，这时，次品数（2）和全产品数（50）之比为$2/50=4\%$．我们就说这批三级管的次品率是4\%．对于上面的分母和分子，我们也可以这样来看：50是抽一只检查（看作一次试验）时所有可能的基本事件数，2是“抽一只检查得次品”这一事件所包含的基本事件数.

因此对于古典概型，如果试验的基本事件总数为$n$，随机事件$A$所包含的基本事件数为$m$，我们就用数$m/n$来描述每次试验中事件$A$出现的可能性的大小，称它为事件$A$的概率，记作$\Pr(A)$. 即我们定义
\[\Pr(A)=\frac{m}{n}\]

于是对“投掷一枚硬币出现正面”这一随机事件中，由于试验的基本事件总数为2，而“出现正面”只包含其中一个基本事件，所以它的概率等于1/2．对于这样定义的概率显然有

\begin{blk}
    {性质1} \[\Pr(U)=1,\qquad \Pr(V)=0\]
\end{blk}

因为必然事件$U$包含了所有基本事件，所以$\Pr(U)=\frac{n}{n}=1$，而不可能事件$V$不含任何基本事件，所以$\Pr(V)=\frac{0}{n}=0$.

\begin{blk}
 { 性质2}\[0\le \Pr(A)\le 1\]  
\end{blk}

因为事件$A$所含基本事件数$m$满足不等式$0\le m\le n$，所以$0\le \Pr(A)\le 1$.

要算出基本事件总数及A所包含的基本事件数，“排列”、“组合”常常是十分有用的工具.

\begin{example}
    $1,2,3,4,5,6$六个数字中任取两数，试计算它们都是偶数的概率.
\end{example}

\begin{solution}
    六个数字中任取两个，共有$\comb{2}{6}=15$种不同的取法，这就是全体基本事件，即$n=15$. 而事件$A$：“任取两数都是偶数”则要从$2,4,6$三个偶数中取两个，也就是说事件$A$包含$\comb{2}{3}=3$个基本事件，因此它的概率
\[\Pr(A)=\frac{\comb{2}{3}}{\comb{2}{6}}=\frac{3}{15}=\frac{1}{5}\]
\end{solution}


\begin{example}
    在一批零件中有$n$个一等品，$m$个三等品，逐个进行检查. 若已查明前$k\;(k<n)$个均为一等品，求第$k+1$次检查时仍得一等品的概率.
\end{example}

\begin{solution}
    由于已查明前$k$个均为一等品，所求第$k+1$次检查基本事件总数为$n+m-k$，“得一等品”这一事件包含的基本事件数为$n-k$. 因此所求的概率
    \[P=\frac{n-k}{n+m-k}\]
\end{solution}


\begin{example}
    一部四卷文集，按任意次序放到书架上，问各卷自左至右或自右至左的卷号恰为1, 2, 3, 4的顺序的概率是多少？
\end{example}

\begin{solution}
    一部四卷文集，在书架上可有$\perm{4}{4}=4!=24$种不同的排列方法，其中自左至右或自右至左卷号恰为1, 2, 3, 4的顺序的共有2种．因此所求的概率
    \[P=\frac{2}{24}=\frac{1}{12}\]
\end{solution}

\begin{example}
    某停车场有12个位置列成一行．求有8个位置停了车而空着的四个位置连在一起的概率.
\end{example}

\begin{solution}
12个位置中占去8个，共有$\comb{8}{12}$种方法，把四个接连的空位看作一个位置，发生四个接连的位置空着的情形相当于这个假想的位置插入8个停车位置中间或者它们的两端，共有9种不同方法，因此所求的概率为
\[P=\frac{9}{\comb{8}{12}}=\frac{1}{55}\]
\end{solution}

\begin{example}
    鱼池中共有鱼$N$条，从中捕得$t$条，加了标志后立即放回池中，经过一段时间后，再从池中抽出$n$条鱼，问其中有$s$条标志鱼的概率是多少？
\end{example}

\begin{solution}
    从$N$条鱼中捕得$n$条，共有$\comb{n}{N}$种可能结果，它就是基本事件总数. 在捕得的$n$条鱼中恰有$s$条标志鱼，它们是从$t$条标志鱼中捕得的，这种捕法共有$\comb{s}{t}$种. 而对于这样的每一种捕法，其余的$n-s$条未加标志的鱼是从$N-t$条鱼中捕得的，共有$\comb{n-s}{N-t}$种捕法. 因此“捕得$n$条鱼，其中$s$条是标志鱼”的捕法共有$\comb{s}{t}\comb{n-s}{N-t}$种，所求概率
\[P_s(N)=\frac{\comb{s}{t}\comb{n-s}{N-t}}{\comb{n}{N}}\]
\end{solution}

\begin{example}
    某小组有成员3人，每人在某星期7天中参加劳动一天，如果劳动日期可随机安排，求3人在不同的3天参加劳动的概率.
\end{example}

\begin{solution}
    对于某一组员（甲）来说，他可以安排在7天中的任何一天参加劳动，即有7种安排方法，而对于甲的任何一种安排另一组员（乙）又可有7种安排法，因此甲乙两人共有$7^2$种不同的安排法。类似地，对于$7^2$种的任何一种安排，第三名组员（丙）又可有7种安排法．因此所有可能的安排法——总的基本事件数等于$7^3$，而3人安排在7天中的不同的3天有${\rm A}^3_7$种不同的方法，所以要求的概率为
\[P=\frac{{\rm A}^3_7}{7^3}=\frac{30}{49}\]
\end{solution}

上述例3.9颇具典型性，$N$个人在$n$天中的任意安排法相当于把$N$个球随机地放入$n$个盒子里去的方法，共有$n^N$种. 这一类问题称为盒子模型.

\begin{ex}
\begin{enumerate}
    \item 指出下列事物是必然事件，不可能事件，还是随机事件.
\begin{enumerate}[(1)]
\item 如果$a$，$b$都是实数，那么$a+b=b+a$;
\item 从一副扑克牌中任意抽取两张，得到“黑桃”；
\item 某电话总机在一分钟内接到20次呼唤．
\end{enumerate}
\item 某班有50位同学而其中女同学占15名．今碰到这个班的同学，正好碰到一个男同学的概率是多少？正好碰到一个女同学的概率是多少？这两个概率之间有什么关系？为什么？
\item 盒中有100个铁钉，其中有90个是合格的，而10个是坏的，从中任意抽取10个，问其中没有一个为坏的概率是多少？
\item 设袋中有8个球，其中5个白球3个红球，从中任意抽取4个，问恰好抽到3个白球的概率是多少？
\item 两袋分别盛着写有0, 1, 2, 3, 4, 5六个数字的六张卡片，从每袋中各取一张，求所得两数之和等于6的概率．现在有人给出下述两种不同解答：

解1：两数之和共有$0,1,2,\ldots,10$十一种不同结果，因此，所求的概率为$1/11$．

解2：从每袋中各任取一张卡片，共有$6^2$种取法，其中和数为6的情形有5种；
\[(1,5);\quad (2,4);\quad (3,3);\quad (4,2);\quad (5,1)\]
因此所求的概率为$5/36$．

试问哪一种解法正确？为什么？
\item 号码锁上有六个拨盘，每个拨盘上有0—9共10个数字，当这6个拨盘上的数字组成某一个6位数时（第一位可以是0），锁才能打开。如果不知道锁的号码，一次就把锁打开的概率是多少？
\item 设有50张考签，分别加以标号$1,2,\ldots,50$，一学生任意抽一张进行考试，求抽到前10号考签的概率？
\item 一批产品共有100件，其中有次品5件，从中任取10件，求所取10件中至多有2件是次品的概率。
\item 袋中装有$n$个白球和$m$个黑球，从中任取$a+b$个，求所取的球恰含$a$个白球和$b$个黑球的概率？
\end{enumerate}
\end{ex}

\subsection{等可能性事件的概率}
在古典概型里，我们曾规定了$\Pr(A)=m/n$，这是在等可能性的大前提下得出的，这种类型在我们计算事件发生的
概率时是常常碰到的，也是较为重要的，为着重视这种类型，我们再深入研究这类问题.

\begin{example}
    一个均匀材料造的正方体玩具，各个面上分别标以数字$1,2,3,4,5,6$，这种玩具叫骰子，
\begin{enumerate}[(1)]
\item 将骰子抛掷一次，朝上的一面出现奇数的概率是多少？
\item 抛掷2次，朝上的一面的数字之和为7的概率是多少？
\end{enumerate}
\end{example}

\begin{solution}
\begin{enumerate}[(1)]
\item 朝上的一面可为$1,2,3,4,5,6$等六种之一，且每种可能性相等，朝上一面出现奇数，可为$1,3,5$等三种之一，且每种可能性相等，设所问的事件概率为$\Pr(A)$，则$\Pr(A)=\frac{3}{6}=\frac{1}{2}$
\item 抛掷二次，每次朝上一面可有六种情况，两次则共有$6\x6=36$种情况，朝上一面的数字之和为7，则有第一次是1，第二次是6，第一次是6，第二次是1等两种情况。
2与5；3与4，又各有两种，总共有6种，故概率$P=\frac{6}{36}=\frac{1}{6}$.
\end{enumerate}
\end{solution}

\begin{example}
    一次掷三枚骰子，得3点，4点，5点，10点的概率各是多少？
\end{example}

\begin{solution}
\begin{enumerate}[(1)]
    \item 得3点的概率，掷三颗骰子共有216种情况，而得3点只有三颗都为1时的一种，故掷得3点的概率为$\frac{1}{216}$；
    \item 掷三颗骰子共有216种情况，而得4点有$2,1,1$; $1,2,1$; $1,1,2$等三种情况，故掷得4点的概率为$\frac{3}{216}=\frac{1}{72}$；
\item 掷三颗骰子共有216种情况，而得5点有$1,1,3$; $1,3,1$; $3,1,1$; $2,2,1$; $2,1,2$; $1,2,2$，即$\comb{1}{3}+\comb{2}{3}=6$种，故掷得5点的概率为$\frac{6}{216}=\frac{1}{36}$；
\item 掷三颗骰子共有216种情况，而得10点有$1,3,6$; $1,4,5$; $2,2,6$; $2,3,5$; $2,4,4$; $3,3,4$等六种
情况．而$1,3,6$；$1,4,5$；$2,3,5$各有$\perm{3}{3}$种排列．$2,2,6$；$2,4,4$；$3,3,4$有$\comb{1}{3}$种组合. 共有$3\perm{3}{3}+3\comb{1}{3}=27$种. 10点的概率为$\frac{27}{216}=\frac{1}{8}$.
\end{enumerate}  

读者试将掷三颗骰子得3点至18点各点数的概率分别求看看各个概率之间有什么关系，各个概率之和又是什么？
\end{solution}

\begin{example}
    袋中装有$a$个白球和$b$个黑球，从中任意地取出$k+1$（$k+1\le a+b$）个球，每次取一球，取后不放回，求最后取出的一球是白球的概率？
\end{example}

\begin{solution}
    从$a+b$个球中无放回地接连取出$k+1$个球，一共有
$(a+b)(a+b-1)\cdots(a+b-k-1+1)={\rm A}^{k+1}_{a+b}$种取法。

$A$表“取出的$k+1$个球中最后一球是白球”的事件。最后取出的白球可以是$a$个白球中的任何一个，有$a$种；而其余$k$个可以是余下的$a+b-1$中的任意$k$个，有${\rm A}^k_{a+b-1}$种取法，因此事件$A$共含$a\cdot {\rm A}^k_{a+b-1}$个不同的基本事件，故所求概率为
\[\Pr(A)=\frac{a{\rm A}^{k}_{a+b-1}}{{\rm A}^{k+1}_{a+b}}=\frac{a}{a+b}\]

从上式可以看出，$\Pr(A)$的值与$k$无关，它说明了无论在哪个序号上，取出白球的概率总是相同的。
\end{solution}

\begin{example}
    有$n$个可辨的球，随机地放入$N\; (N\ge n)$个盒中，每一个盒可以放任意多个球，试求：
\begin{enumerate}
\item 指定的$n$个盒中各有一球的概率？
\item 某$n$个盒中各恰有一球的概率？
\item 某指定的一个盒中恰有$m$（$m\le n$）个球的概率？
\end{enumerate}
\end{example}

\begin{solution}
    设球以同样的可能性落入$N$个盒中的每一个，$N$个盒中落一个球一共有$N$种方法，$n$个球可以有$N^n$种方法，这就是总的基本事件数.
\begin{enumerate}
\item 以$A$表“指定的$n$个盒中各有一球”的事件．今固定某$n$个盒，第一个球可以落入这$n$个盒中的任何一个，有$n$种方法，第二个球可落在余下的$n-1$个盒中的任何一个，有$n-1$种方法，……，第$n$个球落在最后的一个盒，只能有一种方法，因此事件$A$共含有$n!$个不同的基本事件，故所求概率为
\[\Pr(A)= \frac{n!}{N^n} \]
\item 以$B$表“某$n$个盒中各恰有一球”的事件．因为$n$个盒可从$N$个盒中任意选取，共有$\comb{n}{N}$种选法. 选出这$n$个盒后再按问题1知事件$B$共含$\comb{n}{N}n!$个不同的基本事件数，故所求概率为
\[\Pr(B)=\frac{\comb{n}{N}n!}{N^n}\]
\item 
以$C$表“某指定的一个盒中恰有$m$个球”的事件，因为$m$个球可从$n$个球中任意选出，共有$\comb{m}{n}$种选法，其余$n-m$个球可以任意落入其余的$N-1$个盒中，共有$(N-1)^{n-m}$种落法，因此事件$C$共含$\comb{m}{n}(N-1)^{n-m}$个不同的基本事件，故所求概率为
\[\Pr(C)=\frac{\comb{m}{n}(N-1)^{n-m}}{N^n}\]
\end{enumerate}
\end{solution}

\begin{ex}
\begin{enumerate}
\item 有人说，先后抛掷两枚硬币，只有“两枚都是正面”，“两枚都是反面”，“一枚正面，一枚反面”这3种结果，因此，“两枚都出现正面”这一事件的概率是1/3，这种说法错在哪里？
\item 有100张已编号的卡片（从1号到100号），从中任取1张，计算：
\begin{enumerate}[(1)]
    \item 卡片号是偶数的概率；    \item 卡片号是3的倍数的概率.
\end{enumerate}
\item 掷骰子一枚
\begin{enumerate}[(1)]
\item 掷一次，出现奇数的概率是多少？出现偶数的概率是多少？
\item 抛掷2次，得数字之和为8的概率是多少？
\end{enumerate}
\item 在7张数的卡片中，有4张正数卡片和3张负数卡片．从中任取2张作乘法练习，其积为正数的概率是多少？其积为负数的概率是多少？
\item 某种产品90件，其中甲等品40件，乙等品30件，丙等品20件，在运送这些产品的路上损坏了3件，如果每件产品被
损坏的可能性相同，计算这三等产品中恰好各损坏一件的概率.
\item 在80件产品中，有50件一等品，20件二等品，10件三等品．从中任取3件，计算
\begin{enumerate}[(1)]
\item 三件都是一等品的概率；    \item 2件是一等品，1件是二等品的概率；    \item 一等品，二等品三等品各有一件的概率.
\end{enumerate}

\item 六本不同的书，三本封面是红色，三本封面是蓝色，将它们任意排列在书架的同一层上，问六本书封面颜色相间的概率是多少？
\item 在一副扑克牌（52张）中，有“黑桃，红心，梅花，方块”这四种花色的牌各13张，从中任取4张，这4张牌的花色相同的概率是多少？这四张牌的花色各不相同的概率是多少？
\item 有五根细木棍，它们的长度分别为1, 3, 5, 7, 9厘米，从中任取三根，它们能搭成一个三角形的概率是多少？
\end{enumerate}
\end{ex}

\subsection{互斥事件与加法定理}
在10个乒乓球中，有7个一等品，2个二等品，1个三等品，我们把从中任取一个，取出一等品叫做事件$A$，取得二等品叫做事件$B$，取出三等品叫做事件$C$. 我们看到，如果取出的乒乓球是一等品，即事件$A$发生，那么事件$B$就不发生. 如果取出的是二等品，那么事件$B$发生，那么事件$A$就不发生，也就是说，事件$A$、$B$不可能同时发生. 这种不可能同时发生的两个事件叫做互斥事件. 同理，事件$B$，$C$是互斥事件，事件$A$，$C$是互斥事件. 换句话说，事件$A$，$B$，$C$中，任何两个都是互斥事件。这时我们说事件$A$，$B$，$C$彼此互斥，一般地，如果事件$A_1,A_2,\ldots,A_n$中任何两个都是互斥事件，那么就说事件$A_1,A_2,\ldots,A_n$彼此互斥.

在上面的问题里，因为是任取一个，共有10种等可能的取法，其中得到一等品，二等品，三等品的取法分别有7种，
2种，1种，因此，$\Pr(A)=\frac{7}{10}$, $\Pr(B)=\frac{2}{10}$, $\Pr(C)=\frac{1}{10}$.

现在问：“任取一个乒乓球，取出一等品或二等品”这一事件的概率是多少？这一事件，我们记作“$A+B$”. 因为不论取出一等品还是二等品，都表示这个事件，发生而得到一等品或二等品的取法共有$7+2$种，所以取出一等品或二等品的概率：
\[\Pr(A+B)=\frac{7+2}{10}\]

由$\frac{7+2}{10}=\frac{7}{10}+\frac{2}{10}$，我们看到
\[\Pr(A+B)=\Pr(A)+\Pr(B)\]
它告诉我们：如果事件$A$，$B$互斥，那么事件“$A+B$”发生（即$A,B$中有一个发生）的概率，等于事件$A$，$B$分别发生的概率的和.

一般地，如果事件$A_1,A_2,\ldots,A_n$彼此互斥，那么事件“$A_1+A_2+\cdots+A_n$”发生（即$A_1,A_2,\ldots,A_n$中有一个事件发生，也叫和事件）的概率，等于这$n$个事件分别发生的概率的和，即
\[\Pr(A_1+A_2+\cdots+A_n)=\Pr(A_1)+\Pr(A_2)+\cdots+\Pr(A_n)\]
这叫做互斥事件的加法定理，也就是说互斥事件的和的概率等于各事件概率之和。

\begin{example}
    某地区的年降水量，在100—150毫米范围内的概率是0.12，在150-200毫米范围内的概率是0.25，在200-250毫米范围内的概率是0.16，在250-300毫米范围内的概率是0.14．计算年降水量在100-200毫米范围内的概率与在150-300毫米范围内的概率．
\end{example}

\begin{solution}
    我们把这个地区的年降水量在100-150毫米，150-200毫米，200-250毫米，250-300毫米范围内分别叫做事件$A$，$B$，$C$，$D$. 很明显，这四个事件彼此互斥，根据公式，年降水量在100-200毫米范围内的概率是
\[\Pr(A+B)=\Pr(A)+\Pr(B)=0.12+0.25=0.37\]
年降水量在150-300毫米范围内的概率是
\[\Pr(B+C+D)=\Pr(B)+\Pr(C)+\Pr(D)=0.25+0.16+0.14=0.55\]
\end{solution}


\begin{example}
    在20件产品中，有15件一级品，5件二级品，从中任取3件，其中至少有1件为二级品的概率是多少？
\end{example}

\begin{solution}
    我们把从20件产品中任取3件，其中恰有1件二级品，叫做事件$A_1$；恰有2件二级品，叫做事件$A_2$；3件全是二级品，叫做事件$A_3$，这样，事件$A_1$，$A_2$，$A_3$的概率分别是
\[\begin{split}
    \Pr(A_1)&=\frac{\comb{1}{5}\cdot \comb{2}{15}}{\comb{3}{20}}=\frac{105}{228}\\
    \Pr(A_2)&=\frac{\comb{2}{5}\cdot \comb{1}{15}}{\comb{3}{20}}=\frac{30}{228}\\
    \Pr(A_3)&=\frac{\comb{3}{5}}{\comb{3}{20}}=\frac{2}{228}\\
\end{split}\]

很明显，事件$A_1$，$A_2$，$A_3$彼此互斥，根据公式，3件产品中至少有1件为二级品的概率是
\[\Pr(A_1+A_2+A_3)=\Pr(A_1)+\Pr(A_2)+\Pr(A_3)=\frac{105}{228}+\frac{30}{228}+\frac{2}{228}=\frac{137}{228}\]

从20件产品中任取3件，或者都是一级品，或者不都是一级品（即其中至少有一件是二级品），这两个互斥事件必有一个发生，这种其中必有一个发生的两个互斥事件，叫做对立事件. 一个事件$A$的对立事件通常记作$\overline{A}$，根据对立事件的意义，$A+\overline{A}$是一个必然事件，它的概率等于1．又由于$A$与$\overline{A}$互斥，我们得到
\[\Pr(A)+\Pr(\overline{A})=\Pr(A+\overline{A})=1\]
这就是说，两个对立事件的概率的和等于1．

从上面公式还可得到$\Pr(\overline{A})=1-\Pr(A)$. 
运用这个公式计算事件的概率，有时比较简单．如例3.15还可以这样来解：

从20件产品中任取3件，3件全是一级品（记作事件$A$）的概率：
\[\Pr(A)=\frac{\comb{3}{15}}{\comb{3}{20}}=\frac{91}{228}\]
由于“任取3件，至少有一件为二级品”是事件$A$的对立事件$\overline{A}$，因此，
\[\Pr(\overline{A})=1-\Pr(A)=1-\frac{91}{228}=\frac{137}{228}\]

读者可想想，什么时候用事件的对立事件的概率来计算就比较简单.
\end{solution}

\begin{example}
一个匣子中有红球5个，白球2个，从匣中随便取出两个球，而球的颜色恰好一样的概率是多少？
\end{example}

\begin{solution}
\textbf{解1：} 令$A=$两球颜色一样，$B_1=$两球颜色均为红色，$B_2=$两球颜色均为白色，可见$B_1$和$B_2$互斥，并且$A=B_1+B_2$，因此
\[\Pr(A)=\Pr(B_1)+\Pr(B_2)\]

今从7个球中随便取出两个球，共有$\comb{2}{7}=\frac{7\x 6}{2}=21$种不同的取法，即简单事件总数为21．$B_1$为两球都为红色，即此两球由红球中取得，所以共有$\comb{2}{5}=\frac{5\x4}{2}=10$个简单事件，
因此$\Pr(B_1)=\frac{\comb{2}{5}}{\comb{2}{7}}=\frac{10}{21}$

$B_2$为两球均为白色，即此两球由白球中取得，所以共有
$\comb{2}{2}=1$个简单事件，于是$\Pr(B_2)=\frac{\comb{2}{2}}{\comb{2}{7}}=\frac{1}{21}$，
因此
\[\Pr(A)=\Pr(B_1)+\Pr(B_2)=\frac{10}{21}+\frac{1}{21}=\frac{11}{21}\]

\textbf{解2：}令$A=$“两球颜色一样”，则$\overline{A}=$“两球颜色不一样”，所以
\[\Pr(\overline{A})=\frac{\comb{1}{5}\cdot \comb{1}{2}}{\comb{2}{7}}=\frac{5\x 2}{21}=\frac{10}{21}\]
因此
\[\Pr(A)=1-\Pr(\overline{A})=1-\frac{10}{21}=\frac{11}{21}\]
\end{solution}

一般地说，如果匣中有$m$个红球，$n$个白球，如果从匣中随便抽取两个球，两球颜色相同的概率$P$可以用完全一样的方法，求得其值为
\[P=\frac{\comb{2}{m}}{\comb{2}{m+n}}+\frac{\comb{2}{n}}{\comb{2}{m+n}}=\frac{m(m-1)+n(n-1)}{(m+n)(m+n-1)}=1-\frac{\comb{1}{m}\cdot \comb{1}{n}}{\comb{2}{m+n}}\]

更一般地，如果匣中有$k$种球，第$i$种球共有$n_i$个，从匣中随便取出两个球，这两个球的颜色是一样的概率是多少呢？这也可用类似的方法求出，读者自己去算一下.

\begin{example}
    从52张扑克牌中随便抽取三张．问：
\begin{itemize}
\item $A=$三张花色相同，
\item $B=$三张是“顺子”（即点数相连），
\item $C=$三张是同花顺子（即花色相同，且点数相连）
\end{itemize}
各自出现的概率是多少？
\end{example}

\begin{solution}
从52张牌中随机抽取三张牌，全部基本事件的总个数是
\[\comb{3}{52}=\frac{52\x 51\x 50}{1\x 2\x 3}=52\x 17\x 25\]
注意到花色有四种，令
\[\begin{split}
    A_1=\text{三张都是梅花},&\qquad 
A_2=\text{三张都是方块}\\
A_3=\text{三张都是红心},&\qquad 
A_4=\text{三张都是黑桃}
\end{split}\]
则$A_1,A_2,A_3,A_4$是互斥事件，且$A=A_1+A_2+A_3+A_4$. 又
\[\begin{split}
    \Pr(A_1)&=\Pr(A_2)=\Pr(A_3)=\Pr(A_4)\\
    &=\frac{\comb{3}{13}}{\comb{3}{52}}=\frac{13\x 12\x 11}{52\x 51\x 50}\\
    &=\frac{11}{17\x 50}=\frac{11}{850}
\end{split} \]
因此，$\Pr(A)=4\Pr(A_1)=\frac{22}{425}$.

三张是顺子，从顺子的点数相连这一要求来看，不计花色，共有“$2,3,4$”，“$3,4,5$”，……“$J,Q,K$”，“$Q,K,A$”这十一种不同的情况，依次用$B_1,B_2,\ldots,B_{11}$来表示它们，于是$B_1,B_2,\ldots,B_{11}$互斥，而且是等概率的，$B=B_1+\cdots+B_{11}$，于是
\[\Pr(B)=\Pr(B_1)+\cdots +\Pr(B_{11})=11\Pr(B_1)\]

现在来求$\Pr(B_1)$的值．出现“$2,3,4$”这样的顺子，共有$(\comb{1}{4})^3$这么多种，这是因为这时不考虑花色，同点数的牌每样又有四张，所以出现“$2,3,4$”必须在“2”的四张中取出一张，“3”的四张中取出一张，“4”的四张中取一张，因此，共有$(\comb{1}{4})^3=4^3$这么多种，于是
\[\begin{split}
    \Pr(B_1)&=\frac{4^3}{\comb{3}{52}}=\frac{4^3}{52\x 17\x 25}=\frac{16}{13\x 17\x 25}\\
    \Pr(B)&=\frac{11\x 16}{13\x 17\x 25}=\frac{176}{5525}=\frac{8}{13}\Pr(A)
\end{split}\]
显然$\Pr(B)<\Pr(A)$.

再来计算事件$C$的概率，很明显，从上面的讨论中可以看出来，“$2,3,4$”同花的只可能有四种不同的情形．因此
\[\Pr(C)=\frac{11\x 4}{\comb{3}{52}}=\frac{11}{13\x17\x25} =\frac{1}{16}\Pr(B)\]
显然$\Pr(C)<\Pr(B)$

\end{solution}

\begin{example}
   掷四枚硬币，用“1”表示正面向上，用“0”表示反面向上，于是记$x_i$为第$i$枚硬币的结果，每掷一次，$(x_1,x_2,x_3,x_4)$就表示各枚硬币正、反面的出现情况，例如，
\begin{center}
   （1, 1, 0, 0）就是（正，正，反，反）\\
   （1, 0, 0, 1）就是（正，反，反，正）    
\end{center} 

   求下列各事件的概率：
\[\begin{split}
   A&=\text{恰有两个正面向上}\\
   B&=\text{至少有两个正面向上}\\
   C&=\text{第一枚、第二枚硬币的结果正好相反，第三枚、第四枚硬币的结果正好相反}    
\end{split}\]
\end{example}

\begin{solution}
用$x_1,x_2,x_3,x_4$来表示，就有
\[\begin{split}
   A&=\{x_1+x_2+x_3+x_4=2\}\\
   B&=\{2\le x_1+x_2+x_3+x_4\le 4\}\\
   C&=\{x_1=1-x_2,\; x_3=1-x_4\}    
\end{split}\]

   全部简单事件的个数是一个有重复的排列数，它是$2^4=16$.

\begin{enumerate}[(1)]
    \item   $x_1+x_2+x_3+x_4=2$表示四个数中只有两个是“1"，其余两个一定是“0”．而且只要选定了“1”的位置，则“0”的位置也就随之而定了，从四个位置中任选两个，则共有$\comb{2}{4}$种，即6种，因此
   \[\Pr(A)=\frac{6}{16}=\frac{3}{8}\]
   \item $x_1+x_2+x_3+x_4\ge 2$，这一情形可以分成下面三种等式的情形：
\[\begin{split}
    x_1+x_2+x_3+x_4=2&\qquad (\text{有$\comb{2}{4}=6$种情形})\\
    x_1+x_2+x_3+x_4=3&\qquad (\text{有$\comb{3}{4}=4$种情形})\\
    x_1+x_2+x_3+x_4=4&\qquad (\text{有$\comb{4}{4}=1$种情形})\\
\end{split}\]
因此，
\[\Pr(B)=\frac{6+4+1}{16}=\frac{11}{16}\]
\item 此时$x_1=1-x_2$, $x_3=1-x_4$，于是只要看$x_1,x_3$联合起来有多少种不同的情形. 从重复排列的计算公式可以知道$x_1$，$x_3$联合起来共有$2^2=4$种不同情况. 因此，可得
\[\Pr(C)=\frac{4}{16}=\frac{1}{4}\]

\end{enumerate}
   
\end{solution}

\begin{ex}
\begin{enumerate}
    \item 判别下列每对事件是不是互斥事件，如果是，再判别它们是不是对立事件.
从一堆产品中任取2件，其中：
\begin{enumerate}[(1)]
\item 恰有1件次品和恰有2件次品；
\item 有次品和全是次品；
\item 有正品和有次品；
\item 有次品和全是正品．
\end{enumerate}

\item 从一批乒乓球产品中任取一个，如果其重量小于2.45克的概率是0.22，重量不小于2.50克的概率是0.20，那么重量在2.45—2.50克范围内概率是多少？
\item 某射手在一次射击中射中10环，9环，8环的概率分别为0.24, 0.28, 0.19．计算这个射手在一次射击中：
\begin{enumerate}[(1)]
    \item 射中10环或9环的概率；
    \item 不够8环的概率．
\end{enumerate}

\item 一个匣子内有9张票，其号数分别为$1,2,\ldots,9$．从中任取2张，其号数至少有1个为奇数的概率是多少？
\item 在50件产品中有45件合格品，5件次品，从中任取3件，计算其中有次品的概率.
\item 从52张中随便抽取五张．问
\begin{itemize}
\item $A=$五张花色相同；
\item $B=$五张是“顺子”；
\item $C=$五张同花顺子；
\item $D=$五张中有四张点数相同（J作11点，Q作12点，K作13点）．
\end{itemize}
$A$，$B$，$C$，$D$各自出现的概率是多少？
\item 掷$n$枚硬币，（$n$为偶数）
\begin{itemize}
\item $A=$恰有两个正面向上，
\item $B=$至少有两个正面向上，
\item $C=$第一枚、第二枚硬币的结果正好相反，第三枚、第
四枚的结果正好相反，……第$n-1$枚与第$n$枚的结果正
好相反.
\end{itemize}
求$A$，$B$，$C$各事件的概率.
\item 同时掷三颗骰子，得3点至10点的概率是多少？得11点至18点的概率是多少？
\item 掷三颗骰子，得三颗的点数相同的概率是多少？
\item 从52张扑克中任取五张，问得到三张点数相同，其他两张点数也相同的概率.
\end{enumerate}
\end{ex}

\subsection{事件与集合的对应}
我们现在讨论的一些问题，都离不开基本事件. 所以我们可以把基本事件看作衡量随机事件的“最小单位”，它好像是几何中的“点”，集合中的“元素”.而一切随机事件都是由基本事件组成的，这样，事件之间的各种运算，就和集合之间的各种运算完全对应起来，这样既使我们掌握了事件的运算，还使我们进一步理解了集合运算的意义.

在讨论集合运算时，集合中的一个元素称为一个点，用$a$来表示，则$a$和$\{a\}$是不一样的. 前者是一个点，后者是由$a$这一个点组成的单点集，它是一个集合.

当我们用$E_1,\ldots,E_n$表示基本事件时，它相当于集合中的点，随机事件$A$总是由某些基本事件合起来组成的.如
果把一个基本事件$E_i$也看成是随机事件，那就是指由这个$E_i$组成的集合$\{E_i\}$，这个随机事件$\{E_i\}$相当于一个点组成的单点集. 这一点一定要加以注意，不要一开始就发生混淆.

下面将事件间的运算和事件间的关系的各种定义和集合的运算，以及集合的关系对照着列成一张表，以更进一步看清它们之间的联系，也是为了可以用集合语言来理解概率论.

\begin{longtable}{p{.45\textwidth}p{.45\textwidth}}
\hline
事件& 集合\\
\hline
1. 基本事件  &1. 元素（点）\\
$E_1,E_2,\ldots, E_n$ &$a_1,a_2,\ldots,a_n$\\
\hline
2. 必然事件 & 2. 全集$A$\\
$U=\{ E_1,E_2,\ldots, E_n\}$ & $A=\{a_1,a_2,\ldots,a_n\}$\\
\hline
3. 不可能事件$V$  &  3. 空集$\emptyset$\\
$V$中不含任何$E_i$& $\emptyset$中不含任何$a_i$\\
\hline
4. 随机事件$F$ &  4. 子集合$S$\\
$F=\{E_{i_1},E_{i_2},\ldots,E_{i_m}\}$&$S=\{a_{i_1},a_{i_2},\ldots,a_{i_m}\}$\\
$1\le i_1< i_2<\cdots <{i_m}\le n$&$1\le i_1<i_2<\cdots<i_m\le n$\\
\hline
5. 事件$F_1$包含$F_2$写成$F_1\supset F_2$或$F_2\subset F_1$  &  5. 集合$S_1$包含$S_2$，写成$S_1\supset S_2$或$S_2\subset S_1$\\
即$F_2$中的基本事件均为$F_1$中的基本事件。或$F_2$发生时，$F_1$一定发生。 &即$S_2$中的点均为$S_1$中的点\\
\hline
6. $F_1=\{E_{i1},E_{i2},\ldots,E_{ik}\}$，&   6. $S_1=\{a_{i1},a_{i2},\ldots,a_{ik}\}$，\\
$F_2=\{E_{j1},E_{j2},\ldots,E_{jl}\}$，& $S_2=\{a_{j1},a_{j2},\ldots,a_{jl}\}$，\\
\small $F_1+F_2=\{E_{i1},\ldots,E_{ik}, E_{j1},\ldots,E_{jl}\}$  &$S_1+S_2=\{a_{i1},\ldots,a_{ik},a_{j1},\ldots,a_{jl}\}$\\
则$F_1+F_2$称为$F_1$与$F_2$之和，即$F_1, F_2$中至少有一个事件发生。这个运算为加法运算.
& 这是集合的加法，即把$S_1$和$S_2$内的点合并在一起构成的新集合。\\
\hline
7. 由事件$F_1$和$F_2$中的公共基本事件组成的事件称为$F_1$与$F_2$之积，写成$F_1F_2$. 即“$F_1,F_2$这两个事件都发生”这一事件。这个运算称为乘法运算。&7. 集合$S_1$与$S_2$的公共元素组成的集合，称为$S_1$和$S_2$之积，写成$S_1S_2$。即$S_1$与$S_2$的公共部分。这是集合的乘法运算。\\
\hline
8. $F_1F_2=V$，则称$F_1,F_2$是互斥的，即$F_1$和$F_2$中没有公共的基本事件，即$F_1$与$F_2$同时发生这一事件是不可能的。&
8. $S_1S_2=\emptyset$，则称$S_1,S_2$是互不相交的，即$S_1$与$S_2$没有公共的元素
\\
\hline
9. 逆事件（对立事件）&9. 余集（补集）\\
如果$F=\{E_{i1},\ldots,E_{ik}\}$，则不在$F$中的那些基本事件全体组成的随机事件称为$F$的逆事件，用$\overline{F}$表示。$\overline{F}$就是“$F$不发生”这一事件。
& 
如果$S=\{a_{i1},\ldots,a_{ik}\}$，则由不在$S$中的那些点的全体组成的集，称为$S$的余集，用$\overline{S}$表示。也即在全集$A$中，但不在$S$中的那些点组成的集合。\\
\hline
10. 设$F=\{E_{i1},\ldots,E_{ik}\}$，则由加法的定义可知$F=\{E_{i1}\}+\cdots +\{E_{ik}\}$ & 10. 设$S=\{a_{i1},\ldots,a_{ik}\}$，则由加法的定义可知$S=\{a_{i1}\}+\cdots+\{a_{ik}\}$\\
\hline
\end{longtable}

由这张对照表，可以由集合的关系式得到事件之间的关
系式，符号的含义和上表相同
\begin{enumerate}
\item $V\subset F\subset U$ (相当于$\emptyset\subset S\subset A$) 

不可能事件$\subset$随机事件$\subset$必然事件
\item 若$F_1\subset F_2$, $F_2\subset F_3$, 则$F_1\subset F_3$
(相当于$S_1\subset S_2,\; S_2\subset S_3$时，必有$S_1\subset S_3$)
\item $F\overline {F}= V$, $F+ \overline {F}= U$ (相当于$S\overline{S} = \emptyset$, $S+ \overline{S} = A$) 

注意：此即$F$与$\overline F$之和必然发生，$F$与$\overline{F}$ 之积是不可能事
件。$\overline{F}$是$F$的逆事件，$F$是$\overline{F}$的逆事件，即彼此互为逆事件。
\item 当$F_1\subset F_2$时，$F_1$中所含的简单事件均为$F_2$中的简单事件，由古典概率公式得
$\Pr\left(F_{1}\right)\le  \Pr\left(F_{2}\right)$
\item 由于$F\overline{F}= V$, $F+ \overline {F}= U$, $F$与$\overline{F}$互斥，从加法定理
得到
$$1=\Pr(U)=\Pr(F)+\Pr(\overline{F})$$
因此，
$\Pr(\overline{F})=1-\Pr\left(F\right)$
\end{enumerate}

\section*{习题3.1}
\begin{enumerate}
\item 设有$n$个房间，分给$n$个不同的人，每个人都以同等可能进入每一房间，而且每间房里的人数没有限制，试求不出现空房的概率？
\item 将$n$个可辨的球投入$N$个盒中，每个球都以同等可能落入每一个盒内，求指定某盒是空的概率.
\item 一批产品共$n$件，其中有$k$件废品，求任意抽出的$m$件产品中恰有$\ell$件废品的概率？
\item 某人有$n$把钥匙，其中只有一把能开他的门，他逐把地取钥匙试开，这试开的手续可能需要$1,2,\ldots,n$次试验，试证明在第$i$次，第$j$次打开锁的概率是一样.
\item 袋中装有10个白球5个黑球，第一次从袋中取出一球（不放回），发现是白球，求第二次从袋中取得白球的概率？并比较两次取得白球概率的大小. 如果第一次取得黑球，求第二次取得白球的概率？
\item 在52张扑克牌中任意取13张，问其中至少有一张是10点的概率是多少？
\item 一盒中有5个白球，7个黑球．
\begin{enumerate}[(1)]
    \item 从中任取一球，问取得是白球和取得是黑球的概率各为多少？
    \item 从中任取两球，问取得都是白球的概率、取得一白一黑的概率各为多少？
\end{enumerate}

\item 箱子所装球的总重量为$N$公斤，而每个球为$n$公斤，其中有$m$个是黑球. 求从箱子中任取一球恰巧是黑球的概率.
\item 一表面为红色的正方体被分割成1000个同样大小的小正方体.试求从中任取一小正方体其两面涂有红色的概率.
\item 从写有$a$，$b$，$c$，$d$，$e$的五张卡片中任取两张，求这两张卡片的字母恰好是按字母顺序相邻排着的概率.
\item 在相应地写有2, 4, 6, 7, 8, 11, 12和13数字的八张卡片中任意取两张，求由所取得的两个数字构成的分数为可约的概率.
\item 从1, 2, 3, 4, 5诸数中任取三数，排成三位数，问如此所得三位数是偶数的概率是多少？
\item 求任意取一整数$N$的：（1）平方；（2）四次方；（3）乘任意一整数后，其尾数为1的概率。
（提示：只需讨论个位数.）
\item 在书架中任意放着10本书，求某给定的三本书放在一起的概率.
\item 4个女孩，3个男孩排成一排，求男女相间的概率.
\item 一个匣子里有90只好的、10只次的螺丝钉．如果从中任意取用10只，恰都是好的螺丝钉的概率是多少？
\item 在一个坛子中放有$n_1$个白球，$n_2$个黑球，逐一全部取出，问第一个和最后一个球都为白球的概率是多少？

\item 求出下述两事件的概率：
\begin{enumerate}[(1)]
    \item 10个人的生日在10个不同的月份；
\item 6个人的生日恰巧在两个月中。（提示：“恰巧在两个月中”不包括“集中在某一个月”的情形. 可先考虑6个人的生日在某指定的两个月里有几种可能．再考虑生日集中在任意两个月里有几种可能.）
\end{enumerate}

\item 电话号码由五个数字组成．每个数字可以是$0,1,2,\ldots,9$中的任意一个数，求电话号码是由完全不相同的数字组成的概率.
\item 从一副扑克（52张）中任取四张，求四张牌的花色各不相同的概率.
\item 为了减少比赛场次，把20个球队分成两组（每组10队）进行比赛，求最强的两队被分在不同组内的概率．
\item 掷三枚骰子，得5点以上（包括5点）的概率是多少？
\item 某人的口袋装有1枚五分的硬币，2枚贰分的硬币以及3枚壹分的硬币，他从中任取3枚，取出的总钱数不少于伍分的概率是多少？
\item 把10个运动队先均分成两组预赛，求最强两队被分在：
\begin{multicols}{2}
\begin{enumerate}[(1)]
\item 不同组内；    \item 同一组内的概率．
\end{enumerate}
\end{multicols}

\item 一座楼房有四个大门，两人住在此楼内，问恰好住在同一大门内的概率是多少？
\item 设在一批有1000件一等品，50件二等品的产品中进行质量检查，如果在该批产品中任意抽查10件，结果全是一等品，问在未被抽查的产品中再任意地连续取两件，至少
有一件是一等品的概率是多少？
\item 用火车运载两个工厂生产的同类产品，其中甲厂30件，乙厂20件，有消息证实，在路途中有两件产品损坏，求损坏的是不同厂的产品的概率.
\end{enumerate}

\section{条件概率与独立事件}
\starsubsection{条件概率}
若一袋内有7个白球，5个红球，设事件$A=$“取出一个白球”，事件$B=$“取出两个白球”，显然
\[\Pr(A)=\frac{\comb{1}{7}}{\comb{1}{12}}=\frac{7}{12},\qquad \Pr(B)=\frac{\comb{2}{7}}{\comb{2}{12}}=\frac{7}{22}\]

如果我们取出一个白球后，不再放回，那么这时再取出两个白球（注意，实际上是取出了三个白球，但取法是先一个，再两个地取，而不是一起取出）的概率
\[\Pr(B')=\frac{\comb{2}{6}}{\comb{2}{11}}=\frac{3}{11}\]
同样的，如果我们取出两个白球后，不再放回，那么这时再取出一个白球的概率
\[\Pr(A')=\frac{\comb{1}{5}}{\comb{1}{10}}=\frac{1}{2}\]

这个例子告诉我们，我们所计算的某个事件出现的概率是根据另一事件的出现这一事实来确定的.

我们把“出现事件$B$的条件下，出现事件$A$”的概率称为事件$A$关于事件$B$的条件概率，记作$\Pr(A|B)$. 对于上面例子，则
\[\Pr(A')=\Pr(A|B)=\frac{1}{2},\qquad \Pr(B')=\Pr(B|A)=\frac{3}{11}\]

另一方面，考虑事件$AB=$“取出三个白球”，则
\[\begin{split}
    \Pr(AB)=\frac{\comb{3}{7}}{\comb{3}{12}}=\frac{7}{44}&=\frac{7}{22}\x\frac{1}{2}=\Pr(B)\Pr(A|B)\\
    &=\frac{7}{12}\x\frac{3}{11}=\Pr(A)\Pr(B|A)\\
\end{split} \]

就是说$A$与$B$事件同时出现的概率，可以通过$B$出现的概率$\Pr(B)$与$A$关于$B$的条件概率$\Pr(A|B)$来表示.

一般地，对任两个随机事件$A$与$B$，若$\Pr(B)>0$，则有：
\begin{equation}
  \Pr(AB)=\Pr(A|B)\cdot \Pr(B)  \qquad \Pr(AB)=\Pr(B|A)\cdot \Pr(A)\tag{1}
\end{equation}
即
\begin{equation}
    \Pr(A|B)=\frac{\Pr(AB)}{\Pr(B)}\qquad \Pr(B|A)=\frac{\Pr(AB)}{\Pr(A)}\tag{2}
\end{equation}

\begin{proof}
    设基本事件的总数为$e$，其中事件$A$所包含的基本事件有$m$个，事件$B$所包含的基本事件有$n$个，因为我们并没有假定事件$A$和$B$互斥，因此一般都存在着既属于事件$A$，
也属于事件$B$的基本事件，（即事件$A\cdot B$所包含的基本事件）设这样的事件共有$r$个（如图）于是：
\[\Pr(B)=\frac{n}{e},\qquad \Pr(AB)=\frac{r}{e}\]

\begin{center}
    \begin{tikzpicture}
\draw(0,0) rectangle (5,3.5);
\node at (.5,3.25){$e$};
\draw(2,2)node{$A$} circle(1);
\draw(3.5,2)node{$B$} circle(1);
\draw(1.75,1.75)--(1.5,.6)node[below]{$m$};
\draw(3.75,1.75)--(4,.6)node[below]{$n$};
\draw(2.8,1.75)--(2.7,.6)node[below]{$r$};

    \end{tikzpicture}
\end{center}

再看条件概率$\Pr(A|B)$，因为这是在事件$B$发生的条
件下考虑问题的，这时总的基本事件数就是事件$B$所包含的个数$n$，在$B$发生的条件下$A$再发生所包含的基本事件，就必然是属于事件$AB$的$r$个基本事件.

故$\Pr(A|B)=\frac{r}{n}$

这样就有：
\[\Pr(A|B)\cdot\Pr(B)=\frac{r}{n}\cdot \frac{n}{e}=\frac{r}{e}=\Pr(AB)\]
同样可以证明$\Pr(B|A)\cdot\Pr(A)=\Pr(AB)$.
\end{proof}

\begin{example}
    某人提出一个问题，甲先答，答对的概率是0.4；如果答错，由乙答，答对的概率是0.5，求问题由乙解出的概率．
\end{example}

\begin{solution}
“问题由乙解出”相当于“甲答错”（$A$）与“乙答对”（$B$）两事件一起发生。（即事件$A$，$B$同时发生）.

由题意：“甲答错”的概率：$\Pr(A)=1-0.4=0.6$；而“甲答错的条件下乙答对”的概率$\Pr(B|A)=0.5$，
所以“问题由乙解出”的概率：
\[\Pr(AB)=\Pr(B|A)\cdot \Pr(A)=0.5\x0.6=0.3\]
\end{solution}

\begin{example}
    一只袋中有2只白球和3只红球，从袋中取出一只球，然后在第一只不放回的前提下取出第二只球，那么所取出两只球都是红球的概率是多少？
\end{example}

\begin{solution}
    “取出的两只球都是红球”必定是“第一次取得的是红球”（$A$），且“第二次取得的是红球”（$B$）两事件一起发生（即事件$AB$发生）

    “第一次取得红球”的概率：$\Pr(A)=\frac{3}{5}$，而“第一次取得红球的条件下，第二次取得红球”的概率$\Pr(B|A)=\frac{2}{4}=\frac{1}{2}$，所以“两个球都是红球”的概率：
\[\Pr(AB)=\Pr(B|A)\cdot \Pr(A)=\frac{1}{2}\cdot \frac{3}{5}=\frac{3}{10}\]
\end{solution}

\begin{example}
    某篮球队在己方半场抢得篮板球的概率为0.75而抢得篮板球且反攻投中得分的概率为0.6，问该队已抢得篮板球在手，问这球反攻投中得分的概率有多大？
\end{example}

\begin{solution}
    设事件$A=$“抢得篮板球”；事件$B=$“反攻投中得分”
按题意：$\Pr(A)=0.75$，“抢得篮板球且反攻投中得分”的概率（即积事件$A\cdot B$的概率）：$\Pr(AB)=0.6$

所以“抢得篮板球在手，这球反攻投中得分”的概率。
\[\Pr(B|A)=\frac{\Pr(AB)}{\Pr(A)}=\frac{0.6}{0.75}=0.8\]
\end{solution}

\begin{ex}
\begin{enumerate}
    \item 计算在一分钟的区间内进入某邮局的人数的概率如下表：
\begin{center}
\begin{tabular}{cccccc}
\hline
    1分钟内进入的人数&0人&1人&2人&3人&大于3人\\
\hline
概率&0.05&0.15&0.22&0.22&0.36\\
\hline
\end{tabular}
\end{center}    
考虑下述事件：$A=$至少有一个人到达，$B=$至少有二个人到达，$C=$3个顾客到达.

求$\Pr(C|A)$和$\Pr(C|B)$
\item 一个绿骰子和一个红骰子被转动，令$x$表示在红骰子面上出现点子的数目，$y$表示绿骰子面上出现点子的数目.

令：
\begin{multicols}{2}
    $A=$“$x+y>7$的事件”；
    \\
    $B=$“$x+y<6$的事件”；
    \\
    $C=$“$x+y=7$的事件”；\\$D=$“$y=4$的事件”；\\
$E=$“$y=2$的事件”；\\$F=$“$x$或$y$是5”；\\$G=$“$x$和$y$都是5”；\\$H=$“$x+y=10$”
\end{multicols}
求：
\begin{multicols}{3}
 \begin{enumerate}[(1)]
    \item $\Pr(A|D)$
    \item $\Pr(B|E)$
    \item $\Pr(C|D)$
    \item $\Pr(G|F)$
    \item $\Pr(G|H)$
    \item $\Pr(F|H)$
\end{enumerate}   
\end{multicols}
\end{enumerate}
\end{ex}

\subsection{独立事件与乘法定理}
先看下面的问题：“餐桌上摆着10把外形一样的餐刀，其中7把是不锈钢的；3把是镀铬的，记事件$A=$“甲从中任取一把恰为镀铬的”；事件$B=$“乙从中任取一把恰为镀铬的”.

显然，当只有甲单独取用时：$\Pr(A)=\frac{3}{10}$；当只有乙单独取用时：$\Pr(B)=\frac{3}{10}$.

下面再考虑乙先取用的两种情况：
\begin{enumerate}[(1)]
\item 当乙先取用且取得的是镀铬的，用后归还，这相当于返回抽样.

由于这时甲仍是从10把餐刀中取用，其中镀铬的也仍是3把，所以
“乙取得的是镀铬的条件下（用后归还），甲取得的是镀铬的”这一事件的概率：
\[\Pr(A|B)=\frac{3}{10}=\Pr(A)\]
根据同样的分析：$\Pr(B|A)=\frac{3}{10}=\Pr(B)$
\item 当乙先取用且取得是镀铬的，尚未归还甲就取用，这相当于不返回抽样.

这时由于事件$B$的发生，只剩下9把餐刀，其中镀铬的剩下2把，所以
“乙取得的是镀铬的条件下（用后未归还），甲取得的是镀铬的”这一事件的概率：
\[\Pr(A|B)=\frac{2}{9}\ne \Pr(A)\]
根据同样的分析：$\Pr(B|A)=\frac{2}{9}\ne \Pr(B)$
\end{enumerate}

总结上面的分析得出：当乙（甲）的取用不影响甲（乙）的取用，则有下面两等式成立
\[\Pr(A|B)=\Pr(A),\qquad \Pr(B|A)=\Pr(B)\]

当乙（甲）的取用影响到甲（乙）的取用，则：
\[\Pr(A|B)\ne \Pr(A),\qquad \Pr(B|A)\ne \Pr(B)\]

我们定义：如果对于事件$A$和$B$来说，有
\[\Pr(A|B)= \Pr(A),\qquad \Pr(B|A)= \Pr(B)\]
同时成立，那么称事件$A$和事件$B$为相互独立事件，或称事件$A$和$B$相互独立，否则，就是相关的.

对本例，事实上用后归还的话，不管“乙取用为镀铬的”这一事件发生与否，对“甲取用为镀铬的”概率是没有影响的，反之也一样。

因此，通俗地说，事件$A$与$B$相互独立，就是：事件$A$发生与否对事件$B$发生的概率没有影响，同样，事件$B$发生与否对事件$A$发生的概率也没有影响，例如：同时抛掷两枚硬币，记$A=$“第一枚正面向上”，$B=$“第二枚正面向上”；因第一枚是否正面向上对第二枚正面向上的概率没有影响，反之也一样，所以事件$A$与$B$相互独立.

容易看出：当事件$A$和$B$相互独立时，则事件$A$和$\overline{B}$；$\overline{A}$和$B$；$\overline{A}$和$\overline{B}$都相互独立的。（读者自己证明）

当事件$A$和$B$相互独立时，上节公式(1)
\[\Pr(AB)=\Pr(A|B)\cdot \Pr(B),\qquad \Pr(AB)=\Pr(B|A)\cdot \Pr(A)\]
变为
\begin{equation}
    \Pr(AB)=\Pr(A)\cdot \Pr(B) \tag{3}
\end{equation}

这就是说：两个相互独立的事件同时发生的概率等于每个事件发生的概率的乘积.

公式(1)和(3)通称为概率的乘法定理。

一般地，若$A_1,A_2,\ldots,A_n$为相互独立事件，则有：
\begin{equation}
    \Pr(A_1\cdot A_2\cdots A_n)=\Pr(A_1)\cdot \Pr(A_2)\cdots \Pr(A_n)\tag{4}
\end{equation}
独立性的概念在概率中是十分重要的，它有广泛的应用. 在实际问题中判断事件是否独立是很重要的.

\begin{example}
    甲乙两人进行一次射击，如果两人射中目标的概率甲为0.6，乙为0.7，计算：
\begin{enumerate}[(1)]
\item 两人都击中目标的概率；
\item 其中恰有一人击中的概率；
\item 至少有一人击中目标的概率。
\end{enumerate}
\end{example}

\begin{solution}
\begin{enumerate}[(1)]
    \item “两人都击中目标”（$C$），就是“甲击中目标”（$A$），且“乙击中目标”（$B$），故$C=A\cdot B$. 事件$A$和$B$显然相互独立，所以
  \[  \Pr(C)=\Pr(A\cdot B)=\Pr(A)\cdot \Pr(B)=0.6\x0.7=0.42\]
    \item “其中恰有一人击中”（$D$），就是“甲击中，乙未击中”（$A\cdot \overline{B}$）或“甲未击中，乙击中”（$\overline{A}\cdot B$），故$D=A\cdot \overline{B}+\overline{A}\cdot B$，显然事件$A\cdot \overline{B}$与$\overline{A}\cdot B$在两人各射击一次时是不可能同时发生的，即事件$A\cdot \overline{B}$与$\overline{A}\cdot B$互斥，所以
\[\begin{split}
    \Pr(D)&=\Pr(A\cdot \overline{B}+\overline{A}\cdot B)=\Pr(A\cdot \overline{B})+\Pr(\overline{A}\cdot B)\\
    &=0.6\x(1-0.7)+(1-0.6)\x 0.7=0.18+0.28=0.46
\end{split}\]
    \item “至少有一人击中”（$E$），就是“甲击中，乙未击中”（$A\cdot \overline{B}$）或“甲击中，乙击中”（${A}\cdot B$），或“甲未中乙击中”（$\overline{A}\cdot B$），故$E=A\cdot \overline{B}+A\cdot B+\overline{A}\cdot B$，所以
\[\begin{split}
    \Pr(E) &= \Pr(A\cdot \overline{B}+A\cdot B+\overline{A}\cdot B)\\
    &=\Pr(A\cdot \overline{B})+\Pr(A\cdot B)+\Pr(\overline{A}\cdot B)\\
    &=0.6\x0.3+0.6\x0.7+0.4\x0.7\\
&=0.18+0.42+0.28=0.88
\end{split}\]

此小题还可采用下面解法：
“没有一个击中”（$\overline{E}$），就是“甲未击中”（$\overline{A}$）且“乙未击中”（$\overline{B}$），故$\overline{E}=\overline{A}\cdot \overline{B}$. 所以：
$$\Pr(\overline{E})=\Pr(\overline{A}\cdot \overline{B})=\Pr(\overline{A})\cdot\Pr(\overline{B})=0.4\x0.3=0.12$$
因此：$\Pr(E)=1-\Pr(\overline{E})=1-0.12=0.88$
\end{enumerate}
\end{solution}

\begin{example}
   假设每枚地-空导弹击中飞机的概率均为0.8，如要有99\%的把握击中来犯敌机，问需几枚导弹同时发射？ 
\end{example}


\begin{solution}
设需$n$枚地空-导弹同时发射，

记事件$A=$“敌机被击中”，则“敌机未被击中”（$\overline{A}$），就是“第1枚未击中”（$\overline{A}_1$），且“第2枚未击中”($\overline{A}_2$)，……，“第$n$枚未击中”（$\overline{A}_n$）同时发生。

故$\overline{A}=\overline{A}_1\cdot \overline{A}_2\cdots \overline{A}_n$，而$\overline{A}_1, \overline{A}_2,\ldots,\overline{A}_n$是彼此独立的，所以
\[\Pr(\overline{A})=\Pr(\overline{A}_1)\cdot \Pr(\overline{A}_2)\cdots\Pr(\overline{A}_n)=(1-0.8)^n=0.2^n\]
要求击中概率达99\%，即$\Pr(A)=0.99$，故
\[\Pr(\overline{A})=1-\Pr(A)=0.01\]
所以$n$应满足：$0.2^n=0.01$. 取对数
\[\begin{split}
    n\lg0.2&=\lg0.01\\
n&=\frac{\lg0.01}{\lg0.2}=\frac{-2}{-0.699}\approx 3
\end{split} \]
答：3枚导弹同时发射能以99\%的把握击中来犯敌机．    
\end{solution}

\begin{example}
    一个元件能正常工作的概率$r$称为该元件的可靠性，由元件组成的系统能正常工作的概率称为系统的可靠性。今设所用元件的可靠性都为$r\; (0<r<1)$，且每个元件能否正常工作是相互独立的。试求下面两系统的可靠性，并

\begin{enumerate}[(1)]
    \item 附加通路系统，它由两条$n\; (n\ge 2)$个元件串联成的通路并联而成，如下图：
\begin{center}
    \begin{circuitikz}[european]
\draw(0,0)--(.5,0)--(.5,.5)to [R](2.5,.5) to [R] (4.5,.5)--(5,.5);
\draw(6,.5) to [R](7.5,.5);
\draw(.5,0)--(.5,-.5)to [R](2.5,-.5) to [R](4.5,-.5)--(5,-.5);
\draw[dashed](4.5,-.5)--(6,-.5);
\draw[dashed](4.5,.5)--(6,.5);

\draw(6,-.5) to [R](7.5,-.5);
\draw(7.5,-.5)--(7.5,.5);
\draw(7.5,0)--(8.25,0);


\foreach \x/\y in {3/1,7/2,13.5/$n$}
{
    \node at (\x/2,.5){\y};
    \node at (\x/2,-.5){\y};
}

    \end{circuitikz}
\end{center}
    \item 附加元件系统，它由$n$对并联元件串联而成，如下图：
\begin{center}
    \begin{circuitikz}[european]
\draw(0,0)--(.5,0)--(.5,.5) to [R] (2,.5)--(2,-.5) to [R](.5,-.5)--(.5,0);
\draw(2.5,0)--(3,0)--(3,.5) to [R] (4.5,.5)--(4.5,-.5) to [R](3,-.5)--(3,0);
\draw(2,0)--(2.5,0);
\draw(4.5,0)--(5,0);
\draw(8.5,0)--(8,0);
\draw[dashed](6,0)--(5,0);
\draw(6,0)--(6.5,0)--(6.5,.5) to [R] (8,.5)--(8,-.5) to [R](6.5,-.5)--(6.5,0);
\foreach \x/\y in {2.5/1,7.5/2,14.5/$n$}
{
    \node at (\x/2,.5){\y};
    \node at (\x/2,-.5){\y};
}
    \end{circuitikz}
\end{center}
\end{enumerate}
\end{example}

\begin{solution}
    系统(1)有两条通路. 对每条通路来说，必须当每个元件都正常工作时才正常工作，故其可靠性为：$R_c=r^n$

    所以每条通路发生故障的概率为：$1-r^n$.

    系统发生故障，必定是两通路同时发生故障. 故系统发生故障的概率为：$(1-r^n)^2$.

    因此系统(1)的可靠性为：$$R_s=1-(1-r^n)^2=r^n(2-r^n)$$
    所以$R_s>R_c$，即系统在附加了一条通道后，比一条通路增加了可靠性。

    现在考虑系统(2). 从上述讨论易得每对并联元件的可靠性为：$$R'=1-(1-r)^2=r(2-r)$$
    因系统是由各对并联元件串联而成，故其可靠性为：
\[    R_{s'}=(R')^n=r^n(2-r)^n\]
    显然$R_{s'}>R_c$，即用附加元件的方法也能增加系统的可靠性。

两种系统所用的元件一样，其优劣就看可靠性的大小了. 为比较$R_s$与$R_{s'}$的大小，令
$Q_n=\frac{R_{s'}}{R_s}\; (n\ge 2)$，则
\[Q_n=\frac{(2-r)^n}{2-r^n}\]
当$n=2$时:
\[(2-r)^2=2-r^2+2(1-r)^2>2-r^2\]
所以$Q_2>1$

假设当$n=k$时，$Q_k>1$，即$(2-r)^k>2-r^k$，则
\[\begin{split}
 (2-r)^{k+1}>(2-r^k)\cdot (2-r)&=4-2r^k-2r+r^{k+1}\\
&=2-r^{k+1}+2(1-r)(1-r^k)>2-r^{k+1}
\end{split}\]
即当$n=k+1$时也有$Q_{k+1}>1$

所以对于任意大于等于2的自然数$n$，$Q_n>1$总成立，即$R_{s'}>R_s$，所以系统(2)较系统(1)可靠.
\end{solution}

\begin{ex}
\begin{enumerate}
    \item 一张牌随机地从一副扑克牌中抽取，令：    $A=$“这张牌是梅花”；    $B=$“这张牌是个$K$”；    $C=$“这张牌是高级的”（即$J$、$Q$、$K$、$A$）
    
    问：$A,B$；$A,C$和$B,C$的哪一对是独立事件？
    
    \item 掷三个硬币，令：    $A=$“不多于一个正面”；    $B=$“每面至少有一个”。
    
    问：$A$和$B$是独立的吗？
\end{enumerate}
\end{ex}

\section*{习题3.2}
\begin{enumerate}
    \item 一个盒子中装有5只好灯泡，一个调皮的孩子又放进三只坏灯泡，假如八只灯泡里先取出1只（不放回），再取出第2只，求两只都是好灯泡的概率．
    \item 在$1,2,\ldots,100$中随机选出一个数，求满足下列条件的概率：
\begin{enumerate}[(1)]
\item 已知选出的数是偶数的条件下，能被5整除；
\item 已知它是被4整除的条件下它被10整除．
\end{enumerate}

\item 一个口袋内装有2个白球和2个黑球，把“从中任意摸出一个球，得到白球”叫做事件A，把“从剩下的3个球中任意摸出一个球，得到白球”叫做事件B.事件A与事件B是相互独立的吗？再求在先摸出白球后，再摸出 白球的概率，及先摸出黑球后，再摸出白球的概率，并以求出结果验证你上面的结论.
\item 一学校的学生按照下表分类.
\begin{center}
\begin{tabular}{c|cc}
    \hline
   & 戴眼镜&    不戴眼镜\\
   \hline
    男孩&    120&    280\\
    女孩&    180& 420   \\
    \hline
\end{tabular}
\end{center}
令：事件$A=$“学生是男孩”，$B=$“戴眼镜的学生” 

问：
\begin{multicols}{2}
  \begin{enumerate}[(1)]
    \item $A$和$B$
    \item $A$和$\overline{B}$
    \item $\overline{A}$和$\overline{B}$是相互独立事件吗？
\end{enumerate}  
\end{multicols}

\item 生产一种零件，甲车间的合格率是90\%，乙车间的合格率是97\%，从它们生产的零件中各抽取一件，都抽到合格品的概率是多少？
\item 有一问题，在半小时内，甲能解决它的概率是$\frac{1}{2}$，乙能解决它的概率是$\frac{1}{3}$，如果两人都试图独立地在半小时内解决它，计算：
\begin{multicols}{2}
\begin{enumerate}[(1)]
\item 两人都未解决的概率；    \item 问题得到解
决的概率.
\end{enumerate}
\end{multicols}

\item 将一个硬币连掷5次，5次都出现正面的概率是多少？
\item 制造一种零件，甲机床的废品率是0.04，乙机床的废品率为0.05，从它们制造的产品中各任抽一件，其中恰有一件废品的概率是多少？
\item 电子设备的某一部件由9个元件组成，如果其中任何一个元件坏了，这个部件就不能工作，假定每个元件能使用到3000小时的概率是0.99，计算这个部件能工作3000小时的概率.
\item 甲、乙两人解出习题的概率分别为$\frac{2}{3}$和$\frac{4}{5}$，求下面的概率.
\begin{multicols}{2}
\begin{enumerate}[(1)]
    \item 只有1人解出习题：
    \item 至少1人解出习题．
\end{enumerate}
\end{multicols}


\end{enumerate}


\section{独立重复试验}
某射手射击一次，击中目标的概率是0.9，他射击4次，恰好击中3次的概率是多少？

如果我们把射击一次看作一次试验，那么射击4次就相当于连续依次做4次试验，这四次试验有下面的特点：
\begin{enumerate}[(1)]
    \item 这4次试验的结果彼此无关，是相互独立的。
    \item 每一个试验的结果只有两个：“击中目标”（$A$）或“击不中目标（$\overline{A}$），且每次试验中“击中目标”的概率$\Pr(A)$都是0.9，从而“击不中目标”概率$\Pr(\overline{A})$都是0.1
\end{enumerate}

将上面的例子的特点，推广到一般情形，我们把满足上面两个条件的试验，叫独立重复试验，简称独立试验. 这一类的概率问题称为独立试验概型. 若在试验中，我们只关心出现的事件$A$或$\overline{A}$，这种独立试验概型又称为贝努里（J. Bernoulli）概型. 上面问题（及后面将要讨论的）就是属于贝努里概型.

我们继续讨论其解法：

令${A}_i=$“第$i$次击中”，$i=1,2,3,4$，则$\overline{A}_i=$“第$i$次没有击中”，射击4次击中3次，共有下面几种情况：
\[\overline{A}_1A_2A_3A_4;\quad A_1\overline{A}_2A_3A_4;\quad A_1A_2\overline{A}_3A_4;\quad A_1A_2A_3\overline{A}_4\]
上述每一种情况，都可以看成是在4个位置上取出3个写上$A$，另一个写上$\overline{A}$，所以这些情况的种数等于从4个元素中，取出3个元素的组合数$\comb{3}{4}$，即4种．

由于4次射击，其结果是相互独立的，根据乘法定理，前3次击中，第4次未击中的概率为：
\[\Pr(A_1A_2A_3\overline{A}_4)=\Pr(A_1)\cdot \Pr(A_2)\cdot \Pr(A_3)\cdot \Pr(\overline{A}_4)=0.9^3\x(1-0.9)^{4-3}\]
同理：
$$\Pr(A_1A_2\overline{A}_3A_4)=\Pr(A_1\overline{A}_2A_3A_4)=\Pr(\overline{A}_1A_2A_3A_4)=0.9^3\x(1-0.9)^{4-3}$$
这就是说，上面射击4次，击中3次的4种情况的概率都是$0.9^3\x(1-0.9)^{4-3}$．因为这4种情况彼此互斥，根据加法定理，射击4次，击中3次的概率
\[\begin{split}
    P&=\Pr(A_1A_2A_3\overline{A}_4)+\Pr(A_1A_2\overline{A}_3A_4)+\Pr(A_1\overline{A}_2A_3A_4)+\Pr(\overline{A}_1A_2A_3A_4)\\
    &=\comb{3}{4}\x 0.9^3\x(1-0.9)^{4-3}\\
    &=4\x 0.9^3\x0.1\approx 0.29
\end{split}\]

如果射击4次改为$n$次，击中3次改为$k$次，也可得出类似的结果.

一般地，如果在一次试验中事件$A$发生的概率是$P$，那么在$n$次独立试验中事件$A$恰好发生$k$次的概率：
\[P_n(k)=\comb{k}{n}\cdot P^k(1-P)^{n-k}\]

在应用这个公式的时候，首先要分析实际问题是否满足独立重复试验的两个条件.


\begin{example}
    某气象站天气预报的准确率为80\%，计算：
\begin{enumerate}[(1)]
\item 5次预报恰有4次准确的概率；
\item 5次预报至少有4次准确的概率；
\item 5次预报准确的不多于4次的概率；
\item 5次预报准确的少于4次的概率．
\end{enumerate}

\end{example}

\begin{solution}
\begin{enumerate}[(1)]
    \item 设$A=$“预报一次，结果准确，”预报5次相当于作5次独立试验，知$\Pr(A)=0.8$，所以5次预报恰有4次准确的概率：
\[P_5(4)=\comb{4}{5}\x0.8^4\x(1-0.8)^{5-4}=5\x0.8^4\x0.2\approx 0.41\]
\item 5次预报都准确的概率：
\[P_5(5)=\comb{5}{5}\x0.8^5\x(1-0.8)^{5-5}=0.8^5\approx 0.33\]
因此5次预报至少有4次准确的概率：
\[P=P_5(4)+P_5(5)=0.41+0.33=0.74\]
\item 5次预报准确的不多于4次的概率：
\[\begin{split}
    P&=P_{5}(0)+P_{5}(1)+P_{5}(2)+P_{5}(3)+P_{5}(4)\\
&={\rm C}_{5}^{0}\times 0.8^{0}\times(1-0.8)^{5-0}+{\rm C}_{5}^{1}\times0.8^{1}\times(1-0.8)^{5-1}+{\rm C}_{5}^{2}\times0.8^{2}\times(1-0.8)^{5-2}\\
&\qquad +{\rm C}_{5}^{3}\times0.8^{3}\times(1-0.8)^{5-3}+{\rm C}_{5}^{4}\times 0.8^{4}\times (1-0.8)^{5- 4}\\
&=0.2^{5}+5\times0.8\times0.2^{4}+10\times0.8^{2}\times0.2^{3}+10\times0.8^{3}\times0.2^{2}+5\times0.8^{4}\times0.2\\
&=0.00032+0.0064+0.0512+0.2048+0.4096\approx0.67
\end{split} \]
或
\[P=1-P_{5}(5)\approx 1-0.33=0.67\]

\item 5次预报准确的小于4次的概率：
\[\begin{split}
    P&=P_{5}(0)+P_{5}(1)+P_{5}(2)+P_{5}(3)\\
    &=0.00032+0.0064+0.0512+0.2048\approx 0.26
\end{split}\]
或
\[P=1-P_{5}(4)-P_{5}(5)=1-0.4096-0.33\approx 0.26\]
\end{enumerate}
\end{solution}

\begin{example}
    设有$\ell$升经过紫外线消毒的自来水，其中含有$n$个大肠菌. 今从其中任取一升水检验，问在取出的一升水中含有$k\; (k=0,1,\ldots,n)$个大肠菌的概率是多少？
\end{example}

\begin{solution}
  对于每一个大肠菌来说只有两种可能结果：落入或不落入取出的一升水中，且可以认为每个大肠菌落入与否是彼此独立的. 这样$n$个大肠菌落入所取出的一升水中可以看作做了$n$次独立试验，每次试验结果是某一大肠菌落入该升水（$A$）或不落入该升水（$\overline{A}$）. 而每个大肠菌落入该升水的概率为$P=\Pr(A)=\frac{1}{\ell}$，故在取出的一升水中有$k$个大肠菌
的概率为：
\[P_n(k)=\comb{k}{n}\cdot P^k\cdot (1-P)^{n-k}=\comb{k}{n}\cdot \left(\frac{1}{\ell}\right)^k\cdot \left(1-\frac{1}{\ell}\right)^{n-k}  \]
\end{solution}



\section*{习题3.3}
\begin{enumerate}
    \item 生产一种零件，出现次品的概率是0.04，生产这种零件4件，其中恰有1件次品、恰有2件次品、至多有一件次品的概率各是多少？
    \item 某工厂生产中出现次品的概率为0.01，进行重复抽样检查，共取3个样品，求其中次品数等于0, 1, 2, 3的概率.
    \item 设一射手每次中靶的概率为3/5，求四次射击中：
\begin{multicols}{2}
\begin{enumerate}[(1)]
\item 击中一次，    \item 第三次击中，
\item 击中两次，    \item 第二、三两次击中；
\item 至少击中一次的概率。
\end{enumerate}    
\end{multicols}
    \item 有甲乙丙三批罐头，每批100个，其中各有一个是不合格的. 从三批罐头中各抽出一个，计算：
\begin{enumerate}[(1)]
  \item 3个中恰有一个不合格的概率；
\item 3个中至少有一个不合格的概率． 
\end{enumerate}
    \item 某工厂生产过程中出现次品的概率为0.05，每一百个产品为一批，检查产品质量时，在每批中任取一半来检查，如果发现次品不多于1个，则可以认为这批产品是合格的，求一批产品被认为是合格的概率.
\end{enumerate}

\section*{本章内容要点}
一、本章通过大量实例，初步介绍了随机事件的概率及其计算公式.

二、在一定条件下，可能发生也可能不发生的事件，称为随机事件.反映随机事件出现的可能性的大小的一个量，叫做概率.

三、具有有限性、等可能性的特征的随机现象，我们称为古典概型. 如果试验的基本事件总数为$n$，随机事件$A$所包含的事件数为$m$，那么，$A$的概率就定义为$\frac{m}{n}$，记作
\[\Pr(A)=\frac{m}{n}\]
特别地，必然事件$U$的概率为$\Pr(U)=1$; 不可能事件$V$的概率为$\Pr(V)=0$，显然，$0\le \Pr(A)\le 1$.

四、不可能同时发生的两个事件叫 做互斥事件，当$A$，$B$互斥时，
\begin{equation}
    \Pr(A+B)=\Pr(A)+\Pr(B)\tag{*}
\end{equation} 

五、在事件$B$出现的条件下，事件$A$的出现的概率，我们称为$A$关于$B$的条件概率，记作$\Pr(A|B)$. 且有公式
\begin{equation}
   \Pr(A|B)=\frac{\Pr(AB)}{\Pr(B)}\tag{**} 
\end{equation}

六、对于事件$A$，$B$来说，如果有
\[\Pr(A|B)=\Pr(A),\qquad \Pr(B|A)=\Pr(B)\]
同时成立，那么就称事件$A$，$B$为相互独立事件，两个相互
独立的事件同时发生的概率为：
\begin{equation}
    \Pr(A\cdot B)=\Pr(A)\cdot \Pr(B) \tag{***}
\end{equation}

应该注意区别互斥事件、独立事件这两个概念，注意运用公式(*), (***)的前提条件.

七、其中必有一个发生的两个互斥事件，叫做对立事件，当$A$，$B$是对立事件时，有
\[\Pr(A)=1-\Pr(B)\]

八、具有重复试验的结果是相互独立的，而且每次试验的结果只有$A$与$\overline{A}$的特征的随机现象，我们称为重复独立试验概型（又称贝努里概型）. 如果事件$A$在一次试验中发生的概率是$P$，那么，它在$n$次独立重复试验中恰好发生$k$次的概率是
\[P_n(k)=\comb{k}{n}P^k(1-P)^{n-k}\]


\section*{复习题三}
\begin{enumerate}
    \item 某种集成电路用到2000小时还能正常工作的概率是94\%，使用到3000小时仍能正常工作的概率是78\%，问已经工作了2000小时的集成电路，能继续工作到3000小时的概率是多少？
    \item 若$M$件产品中包含$m$件废品，今在其中任取两件，求：
\begin{enumerate}[(1)]
\item 已知取出两件中有一件为废品的条件下，另一件也是废品概率；
\item 已知两件中有一件不是废品的条件下，另一件是废品的概率.
\end{enumerate}
    \item 一只瓶含有6个黑球，4个红球和2个绿球，两个球是逐个而不再放回地随机取出，求两个球都是：（1）黑的，（2）红的．（3）绿的，（4）相同的颜色的概率.
    \item 甲袋装有白球3个，黑球5个；乙袋中装有白球4个，黑球6个，现在从甲袋中随意地取出一个球放入乙袋，充分掺混后再从乙袋随意地取出一个球放入甲袋，求这时甲袋中的白球数成下面的情况的概率：
\begin{multicols}{2}
\begin{enumerate}[(1)]
\item 增加，    \item 不变．
\end{enumerate}
\end{multicols}

    \item 在$1,2,\ldots,100$中任取一数，求它既能被2整除又能被5整除的概率，又它能被2整除或能被5整除的概率是多少？
\item 假设总人口中男孩和女孩的比例是1:1，在所有恰有三个孩子的家庭中，求：
\begin{enumerate}[(1)]
\item 一个家庭有两个女孩和一个男孩的概率；
\item 已知最大的孩子是女孩的条件下，一个家庭有两个女孩和一个男孩的概率.
\end{enumerate}

\item 在一个恰有三个孩子的家庭里，若已知：已有（1）一个女孩，（2）一个男孩，求恰有两个女孩的概率是多少？
\item 甲、乙、丙三儿童，甲有红球$a$只，白球$b$只；乙全是红球；丙红、白球都有，三人各拿出1球，其中甲是随机地取出的，丙拿红球或白球的概率都是$\frac{1}{2}$，然后让甲从取出的三个球任意取回一个，求甲的球成下列情况的概率：
\begin{multicols}{3}
\begin{enumerate}[(1)]
\item 红球增加；
\item 红球个数不变；
\item 红球减少．
\end{enumerate}
\end{multicols}

\item 下图是某市街道地图的一部分，一人计划从角落$A$处散步到$B$处，这里往北要走过7条街，往东要过4条街，求：
\begin{enumerate}[(1)]
    \item 从$A$到$B$有几种不同的走法？
    \item 他通过$G$或者$H$的概率是多少？
    \item 他通过标有$G$的拐角处的概率是多少？
    \item 已知他通过$G$或$H$. 求他经过拐角$J$的条件概率.
\end{enumerate}

\begin{center}
\begin{tikzpicture}[xscale=.6, yscale=.4]
\foreach \y in {0,1,2,...,9}
{
    \draw[gray](-.5,\y)--(6.5,\y);
}    
\foreach \x in {0,1,2,...,6}
{
    \draw[gray](\x,-.5)--(\x,9.5);
}    
\tkzDefPoints{1/1/A, 3/3/H, 2/6/G, 4/6/J, 5/8/B}
\tkzDrawPoints(A,H,G,J,B)
\tkzLabelPoints[above left](A,G)
\tkzLabelPoints[above right](B,J,H)


\end{tikzpicture}
\end{center}



\item 轰炸机要完成它的使命，必须驾驶员找到了轰炸目标，投弹员又投中了目标，设驾驶员甲和乙找到目标的概率分别为0.9和0.8；又设投弹员丙和丁在驾驶员找到了目标的条件下投中目标的概率分别为0.7和0.6．现在配备两架轰炸机的人员，问甲乙与丙丁应该怎样搭配才能使完成使命有较大的概率？又这个概率比起另一种配合大多少？（只要有一架飞机投中目标即完成使命）.
\item 一个工人看管四台机床，如果在一小时内这些机床不需要人去照顾的概率，第一台是0.79，第二台是0.79，第三台是0.80，第四台是0.81，假设各台机床是否需要照顾相互之间没有影响，计算在这个小时内，这四台机床都不需要照顾的概率.
\item 有甲、乙、丙三批罐头，每批分别有100、120、
150个，其中各有3个是不合格的，从三批罐头中各抽出一个，计算：
\begin{enumerate}[(1)]
\item 3个恰有一个合格的概率；
\item 3个中至少有一个不合格的概率。
\end{enumerate}

\item 甲、乙、丙三步枪手射中某目标的概率各为0.7、0.8及0.85；现向该目标各发一弹，问射中3, 2, 1, 0弹的概率各为多少？
\item 两部车床独立地工作，每部车床不需要工人照顾的概率分别为0.9和0.85，求
\begin{enumerate}[(1)]
\item 都不需要照顾的概率；
\item 恰好一台需要照顾的概率；
\item 同时需要照顾的概率．
\end{enumerate}

\item 有一乘客在电车汽车联合站候车，任何一种车都能使他到达目的地. 如果在五分钟内电车到站的概率是$\frac{1}{2}$，汽车到站的概率是$\frac{1}{3}$. 问这个乘客在五分钟内能够登上任何一种车的概率是多少？
\item 甲袋子内有$m$个白球，$n$个黑球，乙袋子内有$n$个白球，$m$个黑球. 从两个袋子内各任意摸出一个，得到一个白球一个黑球的概率是多少？
\item 一箱螺钉共十万只，次品率为0.004，问事件“从中任意抽验一只，是次品；然后又抽验两只，这两只中至少有一只是次品”的概率有多少？
\item 街道交叉口上交通的畅通程度可用任意给定的一秒钟内有一部车子通过的概率P来描述；假定在不同时刻车子的经过是互不影响的，如果一个步行者只能在三秒钟没有车子通过的情形下才能跨过街道，求他恰巧等待0, 1, 2, 3秒钟的概率.
\item 某仪表有$m$个同样的电子管，其中任一个电子管损坏，这个仪表就不能工作，如果在某段时间内每个电子管损坏的概率是$P$，计算在这段时间内，这个仪表不能工作的概率。
\item 在下图所示的开关电路中，开关$A$、$B$、$C$开或关的概率均独立地为$\frac{1}{2}$，求灯点亮的概率.
\begin{center}
\begin{circuitikz}
\draw(.5,1.5) to [normal open switch, l=$A$] (2.5,1.5) -- (3,1.5);
\draw(3,1.5) to [normal open switch, l=$B$]  (3.5,1.5)--(4,1.5);
\draw(.5,1.5)--(.5,.5) to [normal open switch, l=$C$] (4,.5)--(4,1.5);
\draw(4,1)--(4.5,1)--(4.5,-0.5) to [lamp] (1.5,-0.5) to [battery] (0,-0.5)--(0,1)--(.5,1)--(.5,1.5);
\end{circuitikz}
\end{center}


\item 某种大炮击中目标的概率是0.3，只要以多少门这样的大炮同时射击一次，就可以使击中目标的概率：
\begin{multicols}{2}
\begin{enumerate}[(1)]
\item 超过95\%；    \item 不小于0.99．
\end{enumerate}
\end{multicols}

\item 某售货员在三个柜面上售货，如果在某一小时内柜面需要售货员照顾的概率，第一柜面是0.9，第二柜面为0.8，第三柜面为0.7，假定各个柜面是否需要照顾相互之间不受影响，计算在这个小时内，至少有一个柜面需要售货员的概率.
\item 甲、乙、丙三猎手向同一野猪射击，设甲射中的概率为0.4；乙射中的概率为0.5；丙射中的概率为0.7，若只有一人射中，野猪被击毙的概率为0.2；若其中二人射中，则野猪被击毙的概率为0.6；若三人同时射中，则野猪必定被击毙。求野猪被击毙的概率。
\item 某射手射击一次，击中目标的概率是0.9，他射击了四次。分别写出恰好击中4次，3次，2次，1次，0次的概率，并将它们与$(0.9+0.1)^4$的展开式的各项进行比较，你可得出什么结论？
\item 一次测量中出现正误差和负误差的概率都是$\frac{1}{2}$，在3次测量中，恰好出现2次正误差的概率是多少？恰好出现两次负误差的概率是多少？
\item 一头病牛用某药后被治愈的概率为95\%，计算服用这种药的4头病牛中至少有3头被治愈的概率．
\item 在乒乓球比赛中，甲、乙两人球艺相当，胜负机会相同，试比较下列概率哪一个较大？一方在：
\begin{enumerate}[(1)]
\item 4盘中胜3盘与8盘中胜5盘；
\item 4盘中胜3盘以上与8盘中胜5盘以上；
\item $2n$盘中胜的盘数不多于$n$与$2n$盘中胜$n$盘以上．
\end{enumerate}

\item 每一用户在一小时内向总机拨话的概率等于0.02，该总机有300个分机，试问在一小时内有不少于7个分机向总机拨话的概率是多少？

\item 某一批蚕豆种籽，如果每一粒发芽的概率为90\%，播下5粒种籽，计算：
\begin{enumerate}[(1)]
    \item 其中恰有4粒发芽的概率；
    \item 其中恰有4粒未发芽的概率．
\end{enumerate}

\item 在四次独立试验中，事件$A$至少出现一次的概率为$\frac{80}{81}$，求事件$A$在各次试验中出现的概率.
\item 某小组有10台各为7.5千瓦的机床，如果每台机床的使用情况是互相独立的，且每台机床平均每小时开动12分钟，问全部机床用电超过48千瓦的可能性有多大．
\item 某一学校有730名学生，对每一个人来说，他的生日在某一天的概率设为$\frac{1}{360}$，求：
\begin{enumerate}[(1)]
\item 恰有3名学生的生日为元月1日；
\item 生日在元月1日学生不多于3人；
\item 没有一个学生的生日在元月1日；
\item 至少有一个学生的生日在元月1日的概率．
\end{enumerate}

\item 两个篮球运动员在罚球线投球的命中率分别是0.7与0.6，每人投球3次，计算两人都投进2球的概率．
\item 一个织布工人看管12台布机、在$a$小时内一台布机需要工人照顾的概率等于$\frac{1}{3}$
\begin{enumerate}[(1)]
\item 在$a$小时内，有4台布机需要照顾的概率；
\item 在$a$小时内，需要照顾的布机不多于6台的概率．
\end{enumerate}

\item 在多次试验中，当事件$A$出现不少于3次时，事件B才出现．而事件$A$在每次试验中出现的概率等于0.3，若
\begin{multicols}{2}
    \begin{enumerate}[(1)]
    \item 作了5次试验；
    \item 作了7次试验
\end{enumerate}
\end{multicols}
试求$\Pr(B)$．

\item 一个通讯小组有两套通讯设备，只要其中有一套设备能正常工作，就能进行通讯，每套设备由3个部件组成．只要其中有一个部件出故障，这套设备就不能正常工作，如果在某段时间内每个部件不出故障的概率都是$P$，计算在这段时间内能进行通讯的概率.
\item 甲、乙两个篮球队员各投篮3次，每次投篮时投中得分的概率各为0.6和0.7．求
\begin{multicols}{2}
\begin{enumerate}[(1)]
    \item 得分相同的概率；
    \item 甲得分比乙多的概率．
\end{enumerate}    
\end{multicols}

\end{enumerate}


\appendixchapters
\chapter{附录：统计学简介}

\section*{平均值(均值)}

给了一组数据$x_1,x_2,\ldots,x_n$, 相应地有一个平均数
(也称为均值), 它是如下的算术平均数
\[\frac{1}{n} \sum_{i=1}^nx_i\]
通常用$\bar x$来表示它，平均数是日 常生活中常见的量。例如， 工厂中某一产品的平均日产量，农村中某一作物的平均亩产量，商品的平均销售额，学生考试的平均分数等等。自然会产生这样的问题，为什么要用平均数呢？原因是平均数具有代表性，下面来解释它的意思。

对给定的一组数据$x_1,x_2,\ldots,x_n$, 如果用一个数$C$来代表这些$x_1,x_2,\ldots,x_n$, 则$x_i-C$就表达了$x_i$与$C$的偏差， $x_i-C$的绝对值小，$C$就接近$x_i$。即用$C$来代表$x_i$就代表得好. 由于$C$要代表所有数$x_1,x_2,\ldots,x_n$，因此必须考虑几个偏差$x_{1}-C,x_{2}-C,\ldots,x_{n}-C$. 为了消除正负号的影响， 用$(x_i-C)^2$来衡量$x_i$与$C$ 的差别，于是总的差异就是$\sum\limits_{i= 1}^{n}\left ( x_{i}- C\right ) ^{2}$, 当$x_{1},x_{2},\ldots,x_{n}$确定不变时，应选取$C$使$\sum\limits_{i= 1}^{n}\left ( x_{i}- C\right ) ^{2}$达到最小值，这样的$C$才是代表$x_{1},x_{2},\ldots,x_{n}$这$n$个数据最合适的数。

下面来证明这个数就是$\bar x$
\[\begin{split}
    \sum_{i=1}^{n} ( x_{i}-C)^{2} &=\sum_{i=1}^{n} ( x_{i}-\bar{x}+\bar{x}-{C} )^{2} \\
    &=\sum_{i=1}^{n}\left[\left( x_{i}-\bar{x}\right) ^{2}+2\left( x_{i}-\bar{x}\right) ( \bar{x}-C)+(\bar{x}-C)^{2}\right]\\
    &=\sum_{i=1}^{n}\left( x_{i}-\bar{x}\right)^{2}+2\left[\sum_{i=1}^{n}\left( x_{i}-\bar{x}\right)\right](\bar{x}-C)+n(\bar{x}-C)^{2}
\end{split}\]
由于
$$\sum_{i=1}^{n} ( x_{i}-\bar{x} ) =\sum_{i=1}^{n} x_{i}-n\bar{x}=n\bar{x}-n\bar{x}=0$$
因此
$$\sum_{i=1}^{n}(x_{i}-C)^{2}=\sum_{i=1}^{n}(x_{i}-\bar{x})^{2}+n(\bar{x}-C)^{2}\ge \sum_{i=1}^{n}(x_{i}-\bar{x})^{2}$$
显然，等号成立当且仅当$\bar x=C$.

这就证明了：如果以$\sum\limits_{i=1}^n(x_i-C)^2$来衡量$C$代表$x_1,
x_{2},\ldots,x_{n}$的好坏标准，则$\bar x$是代表性最好的值.

现在来看几个特例：

\begin{example}
    如果$x_1=x_2=\cdots=x_n=a$, 则 $\bar{x}=a$.
\end{example}

\begin{proof}
    $$\bar{x}= \frac{1}{n}\sum_{i=1}^{n}x_{i}= \frac{1}{n}\sum_{i= 1}^{n}a= a$$
\end{proof}

\begin{example}
    如果$x_i=a+id$, $i=0,1,2,\ldots,n-1$,
即$x_0,x_1,\ldots,x_{n-1}$是等差数列，这时
\[\begin{split}
   \bar{x}= \frac {1}{n}\sum _{i= 0}^{n-1}x_{i}&= \frac{1}{n}\sum_{i= 0}^{n- 1}( a+ id) \\
   &= \frac{1}{n}\left [ na+ \frac {(  n- 1 )n}{2}d \right ]=  a+\frac{(n-1)d}{2}
\end{split}\]
\end{example}
例如求$1,2,\ldots, 100$的平均 数，这 时$a=d=1$, $n=100$于是$\bar{x} = 1+ \frac {99}{2}= 50.5$; 又 如求$3, 6, 9, 12, \ldots, 27 $这几 个数的均值，则$a= d= 3$, $n= 9$, 于是$\bar{x}=3+\frac{8}{2}\times 3=15$.


\section*{均方差（方差）}

从上面的讨论可以看出，平均数$\bar x$代表$x_1,\ldots,x_n$的好坏程度是由$\sum\limits^n_{i=1}(x_i-\bar x)^2$来反映的，它是$n$个数据$x_i$对
$\bar x$的偏差平方的总和. 如果用$n$除一下，就得到平均的偏差平方和，它反映了每个数据$x_i$偏离$\bar x$的“平均”状况，因此称它为均方差，简称为方差，通常用$S^2$来表示，即有
\[S^2=\frac{1}{n}\sum^n_{i=1}(x_i-\bar x)^2\]
方差表明了一组数据$x_1,x_2,\ldots,x_n$的分散程度. $S^2$越小，表示$x_i$之间的差异小，$\bar x$的代表性就好；$S^2$越大，表示$x_i$之间的差异也大，数据便很分散，此时$\bar x$代表性就不强.

有时，我们要比较的就是两组数据的分散程度. 这用方差就很方便. 例如有两班学生，甲班有$n$个学生，乙班有$m$个学生，进行教学测验后，甲班的成绩为$x_1,x_2,\ldots,x_n$，乙班的成绩为$y_1,y_2,\ldots,y_m$，平均数$\bar x$与$\bar y$反映了两个班的平均成绩，各自方差$S^2_x$和$S^2_y$反映了两班成绩不齐的程度，方差小的班，学生之间的差异就小.

如果$x_1=x_2=\cdots=x_n=a$，则$\bar x=a$，这时方差$S^2=0$. 如果$x_i=i,\; i=1,2,\ldots,n$，则$\bar x=\frac{n+1}{2}$，而方差
\[\begin{split}
    S^2=\frac{1}{n}\sum^n_{i=1}(x_i-\bar x)^2 &=\frac{1}{n}\sum^n_{i=1}\left(i-\frac{n+1}{2}\right)^2\\
&=\frac{1}{n}\sum^n_{i=1}i^2-\left(\frac{n+1}{2}\right)^2\\
&=\frac{(n+1)(2n+1)}{6}-\frac{(n+1)^2}{4}\\
&=\frac{(n+1)(4n+2-3n-3)}{12}=\frac{n^2-1}{12}
\end{split}\]

平均数$\bar x$与平均方差$S^2$是一组数据$x_1,x_2,\ldots,x_{n}$的两个重要指标，给了一组数$x_1,x_2,\ldots,x_{n}$后，如何具体地算出这两个值，我们在下一小节中来详细讨论这一问题.

\section*{计算公式}
设$x_1,x_2,\ldots,x_{n}$相应的均值和方差为$\bar x$和$S^2_x$，设$y_1,y_2,\ldots,y_{n}$相应的均值和方差为$\bar y$和$S^2_y$，下面先导出几个公式，然后利用它们来具体计算。

\begin{enumerate}
    \item 如果$y_i=x_i+a,\; i=1,2,\cdots,n$。则 $\bar {y}= \bar{x} + a$.
    
\begin{proof}
     $\bar{y} = \frac 1n\sum\limits_{i= 1}^{n}y_{i}= \frac 1n\sum\limits _{i= 1}^{n} ({x_{i}+ a}) = \frac 1n\sum\limits _{i= 1}^{n}x_{i}+ a= \bar x+ a$
\end{proof}
    
\item 如果$y_i=bx_i,i=1,2,...n$, 则$\bar{y}=b\bar{x}$

\begin{proof}
$\bar y= {\frac 1n}\sum\limits _{i= 1}^{n}y_{i}= {\frac 1n}\sum\limits _{i= 1}^{n}bx_{i}= b {\bar {x}}$
\end{proof}

    \item 如果$y_i= bx_{i}+ a,\; i= 1, 2, \ldots, n$, 则$\bar {y}=b\bar{x} + a$
  
\begin{proof}
    由1，2立即可得.
\end{proof}  
    \item $\sum\limits _{i= 1}^{n}\left (  x_{i}- \bar {x} \right ) ^{2}= \sum\limits _{i= 1}^{n}x_{i}^{2}- n\bar {x}^{2}$

\item 如果$y_i=bx_i+a,\; i=1,2,\ldots,n$, 则$S_y^{2}=b^{2}S_{x}^{2}$

\begin{proof}
\[\begin{split}
    S_{y}^{2}= {\frac 1n}\sum\limits  _{i= 1}^{n}(y_{i}- \bar{y})^2&=\frac{1}{n}\sum\limits ^n_{i=1}\left ( bx_{i}+ a- b\bar{x} - a\right ) ^{2}\\
    &=\frac{b^{2}}{n}\sum\limits _{i=1}^{n}(x_{i}-\bar{x})^{2}=b^{2}S_{x}^{2}
\end{split}\]
\end{proof}
\end{enumerate}

公式5告诉我们，当$b=\pm1$时，$S_x^2=S_y^2$. 即一组数据同
加一个常数，或同时改变符号时，它的方差是不变的.

\begin{example}
设某地历年夏季的雨量为：（单位mm）
\[248.7\quad 249.4\quad 133.2\quad 153.5\quad 211.7 \]
求它的平均值$\bar x$及方差$S^2$
\end{example}

取$y_i=x_i-200$，所以$x_i=y_i+200$，于是$\bar x=\bar y+200$，$S^2_x=S^2_y$

列表演算如下：
\begin{center}
\begin{tabular}{ccc}
\hline 
$x_i$ & $y_i=x_i-200$ & $y^2_i$\\
\hline
248.7  & 48.7  & 2371.69\\
249.4& 49.4&2440.36\\
133.2&$-66.8$&4462.24\\
153.5&$-46.5$&2162.25\\
211.7&11.7&136.89\\
\hline
$\Sigma$ &$-3.5$& $11573.43$\\
\hline
\end{tabular}
\end{center}

所以
\[\begin{split}
    \bar x&=\bar y+200=-\frac{3.5}{5}+200=199.3\\
S^2_x&=S^2_y=\frac{1}{n}\left(\sum^n_{i=1} y^2_i-n\bar y^2\right)=\frac{1}{5}(11573.43-2.45)=2314.196
\end{split}
    \]

从演算过程可以看出，$y_i$的平方、求和的计算量比$x_i$的平方、求和的计算量小．这里选200是为了使$x_i-200$数字变小而且减法本身也较易.

\begin{example}
    设$x_i=a+id,\; i=1,2,\ldots,n$，即$x_1,x_2,\ldots,x_n$为等差数列，试求方差$S^2_x$.
\end{example}

为了便于比较，下面用两种方法求解.

\textbf{解一：} 已知$\bar x=a+\frac{n+1}{2}d$，所以$x_i-\bar x=\left(i-\frac{n+1}{2}\right)d$.
\[\sum^n_{i=1}(x_i-\bar x)^2 =\sum^n_{i=1}\left(i-\frac{n+1}{2}\right)^2 d^2=\frac{d^2n(n^2-1)}{12} \]
所以
\[S^2_x=\frac{d^2(n^2-1)}{12} \]

\textbf{解二：} 取$y_i=\frac{x_i-a}{d}=i$, $i=1,2,\ldots,n$，于是$S^2_y=\frac{1}{d^2}S^2_x$，即$S^2_x=d^2S^2_y$. 前面已算过，这时$S^2_y=\frac{n^2-1}{12}$
所以$S^2_x=\frac{d^2(n^2-1)}{12}$.

\begin{example}
    对两组数据$x_1,x_2,\ldots,x_n$及$y_1,y_2,\ldots,y_m$, 设它们的均值及均方差分别为$\bar x,\; S^2_x; \; \bar y,\;  S^2_y$，将这两组数据合并为一组数据，求合并后的数据的均值$\bar z$及均方差$S^2_z$．
\end{example}

\begin{solution}
令$\bar z=\frac{1}{n+m}\left(\sum\limits^n_{i=1}x_i+\sum\limits^m_{i=1}y_i\right)=\frac{1}{n+m}(n\bar x+m\bar y)$

又
\[\begin{split}
    S^2_z&=\frac{1}{n+m}\left[\sum\limits^n_{i=1}(x_i-\bar z)^2+\sum\limits^m_{i=1}(y_i-\bar z)^2\right]\\
    &=\frac{1}{n+m}\left[\sum\limits^n_{i=1}(x_i-\bar x)^2 +n(\bar x-\bar z)^2+\sum\limits^m_{i=1}(y_i-\bar y)^2+m(\bar y-\bar z)^2\right]\\
    &=\frac{1}{n+m}\left[nS^2_x+n(\bar x-\bar z)^2+mS^2_y+m(\bar y-\bar z)^2\right]
\end{split}\]
但是$\bar z=\frac{n}{n+m}\bar x+\frac{m}{n+m}\bar y$，代入有
\[\bar x-\bar z=\frac{m}{n+m}(\bar x-\bar y),\qquad \bar y-\bar z=\frac{n}{n+m}(\bar y-\bar x)\]
\[S_{z}^{2}= \frac n{n+ m} S_{x}^{2}+ \frac m{n+ m} S_{y}^{2}+ \frac n{n+ m}\left(\frac m{n+ m}\right)^{2}(\bar x- \bar{y} ) ^{2}+\frac{m}{n+m}\left(\frac{n}{m+n}\right)^{2}(\bar{y}-\bar x)^{2}\]
所以
\[S_{z}^{2}=\frac{1}{n+m}\left(nS_{x}^{2}+mS_{y}^{2}\right)+\frac{nm}{(n+m)^{2}}(\bar{x}-\bar{y})^2\]
总之有
$$\bar z=\frac{1}{n+m}(n\bar{x}+m\bar{y})$$
\[S_{z}^{2}=\frac{1}{n+m}\left(nS_{x}^{2}+mS_{y}^{2}\right)+\frac{nm}{(n+m)^{2}}(\bar{x}-\bar{y})^2\]
\end{solution}

特别，当$m=1$时，$S^2_y=0$，记$y_1=x_{n+1}$，这就得到增加一个新数据的递推公式
\[\begin{split}
    \bar z&= \frac{1}{n+1}(n\bar x+x_{n+1})\\
    S^2_z&=\frac{n}{n+1}S^2_x+\frac{n}{(n+1)^2}(\bar x-x_{n+1})^2
\end{split}\]

不难看出，如果从$x_1,x_2,\ldots,x_n$中减少一个数据，例如删去$x_n$，那么$x_1,x_2,\ldots,x_{n-1}$的均值$\bar x_{*}$与均方差$S^2_{*}$可以用$\bar x$与$S^2_x$及$x_n$表示，即有递推公式
\[\begin{split}
    \bar x_{*}&= \frac{1}{n-1}(n\bar x-x_{n})\\
    S^2_{*}&=\frac{n}{n-1}S^2_x-\frac{n}{(n-1)^2}(\bar x-x_{n})^2
\end{split}\]

\section*{最小二乘法}
在实际工作中经常需要从实测的数据求出变量之间的关系式. 例如年龄和血压是有关系的，随着年龄的增长，血压是会增高的，调查了几百人的年龄与血压的情况，将资料按年龄分组. 用各组的均值（年龄的均值和血压的均值），得到如下的数据：

\begin{minipage}{.35\textwidth}
\begin{center}
    \begin{tabular}{cc}
\hline
        年龄$x_i$ & 心脏收缩压$y_i$\\
\hline
35&114\\
45&124\\
55&143\\
65&158\\
75&166\\
\hline
    \end{tabular}
\end{center}
\end{minipage}\hfill
\begin{minipage}{.6\textwidth}
\centering
\begin{tikzpicture}[>=stealth, scale =.8]
\draw[->](-1,0)--(6,0)node[above]{$x$（年龄）};
\draw[->](0,-1)--(0,8)node[right]{$y$（血压）};
\foreach \x/\y in {1/35,2/45,3/55,4/65,5/75}
{
    \draw(\x,0)node[below]{$\y$}--(\x,.1);
}
\foreach \x/\y in {1/110,2/120,3/130,4/140,5/150,6/160,7/170}
{
    \draw(0,\x)node[left]{$\y$}--(.1,\x);
}
\tkzDefPoints{1/1.4/A, 2/2.4/B, 3/4.3/C, 4/5.8/D, 5/6.6/E}
\tkzDrawPoints(A,B,C,D,E)
\node at (A) [left]{$P_1$};
\node at (B) [right]{$P_2$};
\node at (C) [left]{$P_3$};
\node at (D) [left]{$P_4$};
\node at (E) [right]{$P_5$};
\draw[domain=.5:5.5, smooth, thick]plot(\x, 1.35*\x+.1);
\node[below left]{$O$};
\end{tikzpicture}

\end{minipage}

从表上的数据或图上的点来看，
血压值$y$与年龄数$x$似乎有直线的关系，但是这五个点又不正好在一条直线上，于是发生了一个问题：如何求一个直线方程，使得这条线与五个点最接近.

用$P_i$表示点$(x_i,y_i)\; i=1,2,3,4,5$．对给定的直线$\ell:\; y=a+bx$. 在直线$\ell$上取五个点，使其横坐标与$P_i$的横坐标$x_i$相同，这五个点应是
\[Q_i(x_i, a+bx_i) \qquad i=1,2,3,4,5\]
很明显，这五个点的纵坐标之差为$y_i-a-bx_i$. 用数值$\sum\limits^5_{i=1}(y_i-a-bx_i)^2$
来衡量直线$\ell$上五个点$Q_1,\ldots,Q_5$和已给的五个点$P_1,\ldots,P_5$的总差距，这个差距愈小，就认为这条直线愈接近这些点. 因此，最接近的直线应使这个差距达到最小值，下
面用和以前类似的方法来求出$a$和$b$.

已知$n$个点有坐标$(x_i,y_i)\; i=1,2,\ldots,n$．求直线$\ell:\; y=a+bx$，使和这$n$个点在上述意义下最接近，也就是求数值$a$和$b$，使
\[Q=\sum^n_{i=1}(y_i-a-bx_i)^2\]
达到最小值.

\begin{blk}{定理（最小二乘法）}
使$Q$达最小值的$a,b$为
\[\hat a=\bar y-\hat b \bar x,\qquad \hat b=\frac{\sum\limits^n_{i=1}(x_i-\bar x)(y_i-\bar y)}{\sum\limits^n_{i=1}(x_i-\bar x)^2}\]
\end{blk}

\begin{proof}
对任意数$a,b,\hat a,\hat b$，有
\[\begin{split}
Q&=\sum^n_{i=1}(y_i-a-bx_i)^2 =\sum^n_{i=1}\left[y_i-\hat a-\hat b x_i+\hat a-a+(\hat b-b)x_i\right]^2\\ 
&=\sum^n_{i=1}\left(y_i-\hat a-\hat b x_i\right)^2+2\sum^n_{i=1}\left(y_i-\hat a-\hat b x_i\right)\cdot \left[\hat a-a+(\hat b-b)x_i\right]\\
&\qquad +\sum^n_{i=1}\left[\hat a-a+(\hat b-b)x_i\right]^2
\end{split}\]
取$\hat a,\hat b$使得
\[\sum^n_{i=1}\left(y_i-\hat a-\hat b x_i\right)=0,\qquad \sum^n_{i=1}x_i\left(y_i-\hat a-\hat b x_i\right)=0\]
那么上面
\[\begin{split}
    Q&=\sum^n_{i=1}\left(y_i-\hat a-\hat b x_i\right)^2+\sum^n_{i=1}\left[\hat a-a+(\hat b-b)x_i\right]^2\ge \sum^n_{i=1}\left(y_i-\hat a-\hat b x_i\right)^2
\end{split}\]
因此若能由$\hat a$，$\hat b$之条件唯一
引出$\hat a$，$\hat b$，那么它就是所求的使$Q$达极小值的$a$，$b$了。

将条件改变为
\[\hat a+\hat b \bar x=\bar y,\quad n\hat a \bar x+\hat b\sum^n_{i=1}x^2_i=\sum^n_{i=1}x_iy_i\]
由$\hat a=\bar y-\hat b\bar x$，代入后式，有
\[n\left(\bar y-\hat b\bar x\right)\bar x+\hat b\sum^n_{i=1}x^2_i=\sum^n_{i=1}x_iy_i\]
即有
\[\hat b\left[\sum^n_{i=1}(x_i-\bar x)^2\right]=\sum^n_{i=1}x_iy_i-n\bar x\bar y=\sum^n_{i=1}(x_i-\bar x)\cdot (y_i-\bar y)\]
所以\[\hat b=\frac{\sum\limits^n_{i=1}(x_i-\bar x)(y_i-\bar y)}{\sum\limits^n_{i=1}(x_i-\bar x)^2}\]
于是$\hat a=\bar y-\hat b\bar x$, 这证明了定理.
\end{proof}

回到原来的例子，下面列表进行计算.
\begin{center}
\begin{tabular}{cccccc}
\hline
$x_i$ & $y_i$&$x_i-\bar x$&$y_i-\bar y$&$(x_i-\bar x)^2$&$(x_i-\bar x)(y_i-\bar y)$\\
\hline
35&114&$-20$&$-27$&400&540\\
45&124&$-10$&$-17$&100&170\\
55&143&0&2&0&0\\
65&158&10&17&100&170\\
75&166&20&25&400&500\\
\hline
$\Sigma 275$&705&&&1000&1380\\
平均55&141\\
\hline
\end{tabular}
\end{center}

\[\hat b=\frac{1380}{1000}=1.38,\qquad \hat a=141-1.38\x 55=65.1\]
于是求得最接近这五点的直线为
\[y=65.1+1.38x\]

现在，进一步来考查一下，由上述方程算得的$x=35$, 45, 55, 65, 75时$y$值是多少呢？它们与实际的观测值$y_i$的差又是多少呢？为了以示区别，由方程算出的$65.1+1.38x_i$记为$\hat y_i$有

\begin{center}
\begin{tabular}{cccc}
\hline
实测值&方程算得的值& 偏差& 偏差平方\\
$y_i$&$\hat y_i$&$y_i-\hat y_i$&$\left(y_i-\hat y_i\right)^2$\\
\hline
114&113.4&0.6&0.36\\
124&127.2&$-3.2$&10.24\\
143&141.0&2.0&4.00\\
158&154.8&3.2&10.24\\
166&168.6&$-2.6$&6.76\\
\hline
$\Sigma$&&0&31.60\\
\hline
\end{tabular}
\end{center}

从上表可以看出
\[\sum^n_{i=1}\left(y_i-\hat y_i\right)=0\]
这不是偶然的，可以证明，在一般情况下，总有$\sum\limits^n_{i=1}\left(y_i-\hat y_i\right)=0$. 事实上，
\[\begin{split}
\sum^n_{i=1}\left(y_i-\hat y_i\right)&=\sum^n_{i=1}y_i-\sum^n_{i=1}\left(\hat a+\hat b x_i\right)    \\
&=n\bar y -n\hat a-n\hat b \bar x=n\left(\bar y-\hat a-\hat b \bar x\right)
\end{split} \]
由于已知$\hat a=\bar y-\hat b\bar x$，所以有$\sum\limits^n_{i=1}\left(y_i-\hat y_i\right)=0$


因此，这一结果常常可以作为验算所计算的结果是否正确，但是，在实际使用时，在对大量数据进行计算时，由于
舍入误差的累积会使$\sum\limits^n_{i=1}\left(y_i-\hat y_i\right)\ne 0$，不过这个$\sum\limits^n_{i=1}\left(y_i-\hat y_i\right)$不会大，而在$\sum\limits^n_{i=1}\left(y_i-\hat y_i\right)$的绝对值很大时，这往往是由于计算上的错误所造成的.


\fi

\ifPartB
\ifFullBook\volume{第六册（下）}\fi
\chapter{函数的极限和连续函数的性质}

\section{函数的极限}
\subsection{函数极限的概念}

在第四册下，我们研究了数列的极限，数列是一种特殊的函数，这里的自变数$n$取自然数列$1, 2, 3,\ldots,n,\ldots$的值，现在我们来研究更一般的情形，即函数$f(x)$随$x$连续变化而变化的情形，下面转到一般函数的极限．

\begin{blk}{定义1}
    如果$x$通过任何一个无限增大的数列$\{x_n\}$, 对应的函数值数列
$f (x_1 ) ,f (x_2) , \ldots,f (x_n ) ,\ldots$
都以定数$\ell$为它的极限，就说函数$f(x)$, 当$x\to +\infty$时，以$\ell$为极限，记作
\begin{equation}
    \lim_{x\to+\infty}f(x)=\ell\qquad \text{或}\qquad f(x)\to \ell\quad (x\to +\infty)
\end{equation}
\end{blk}
 
从几何上看，极限式(1.1)表示随着$x$无限增大，曲线$y=f(x)$以直线$\ell$为渐近线（图1.1）．
\begin{figure}[htp]
    \centering
\begin{tikzpicture}[>=latex]
\draw[->](-3,0)--(5,0)node[right]{$x$};  
\draw[->](0,-1)--(0,3)node[right]{$y$};    
\draw[thick](-3,1)--(5,1)node[right]{$y=\ell$};
\draw[very thick, rounded corners=10pt](-3,.1) --(-1.5,.3) --(-.5,.8) --(1,1.5) --(2.5,.7) -- (3,1.2)--(3.5,.8) --(4.5,1.05);
\node at (0,0)[below right]{$O$};

\end{tikzpicture}
    \caption{}
\end{figure}

类似地，可以定义函数极限$\Lim_{x\to -\infty} f(x)=\ell$, 这时，变量
$x$通过代数值无限地变小，而绝对值无限地增大的任何一个数列$\{x_n\}$．

如果函数$f(x)$当$x\to+\infty$和$x\to-\infty$时，都以定值为极限，就说$f(x)$当$x\to \infty$时，以定值$\ell$为极限，记作$\Lim_{x\to\infty}f(x)=\ell$, 或者$f(x)\to \ell\; (x\to\infty)$．

\begin{example}
    证明$\Lim_{x\to\infty}\frac{1}{x}=0$
\end{example}

\begin{proof}
任何数列$\{x_n\}$的值$x_1,x_2,\ldots,x_n,\ldots$趋向于$+\infty$或$-\infty$时，对应的函数数列
\[\frac{1}{x_1},\frac{1}{x_2},\ldots, \frac{1}{x_n},\ldots\]
的绝对值$\left|\frac{1}{x_n}\right|$便趋向于零，即
    \[\lim_{n\to\infty}\left|\frac{1}{x_n}\right|=0\]
    从而$\Lim_{x\to\infty}\frac{1}{x}=0$
\end{proof}

\begin{example}
说明函数$\sin x$, 当$x\to\infty$时，没有极限．
\end{example}

\begin{solution}
令自变量$x$取数列$x_n=-\frac{\pi}{2}+2n\pi\quad (n=1,2,3,\ldots)$的值趋向$+\infty$，则对应的函数值数列
\[\begin{split}
    \sin x_n&=\sin\left(-\frac{\pi}{2}+2n\pi\right)\\
&=\sin\left(-\frac{\pi}{2}\right) =-1 \qquad (n=1, 2, 3, \ldots) .
\end{split}\]
恒取定值$-1$, 于是
\[\lim_{n\to\infty} \sin x_n=\lim_{n\to\infty} \sin \left(-\frac{\pi}{2}+2n\pi\right)=-1\]

再令自变量$x$取数列
\[x_n=\frac{\pi}{2}+2n\pi\qquad  (n=1, 2, 3, \ldots )\]
的值趋向于$+\infty$, 则对应的函数值数列
\[\sin x_n =\sin\left(\frac{\pi}{2}+2n\pi\right) =1\qquad  (n=1, 2, 3, \ldots)\]
恒取定值1,于是
\[\lim_{n\to\infty} \sin x_n=\lim_{n\to\infty} \sin\left(\frac{\pi}{2}+2n\pi\right) =1\]
由于$x$取趋向$+\infty$的两个不同数列时，$y=\sin x$可以有不同的极限，因此
$\Lim_{x\to\infty} \sin x$不存在．
\end{solution}

\begin{blk}{定义2}
 设函数$f(x)$在点$a$附近有定义（但在$x=a$时，可以没有定义），如果当自变量$x$不论通过怎样一个以$a$为极限但始终不等于$a$的数列$\{x_n\}$, 对应的函数值数列$f (x_1) ,f (x_2) , \ldots,f(x_n ) ,\ldots$
总有极限$\ell$, 就说：

当$x$趋近于$a$时，$f(x)$以$\ell$为极限，记作
\begin{equation}
   \lim_{x\to a}f(x)=\ell,\qquad \text{或}\qquad f(x)\to \ell\quad (x\to a) 
\end{equation}
\end{blk}

极限式(1.2)的几何意义如图1.2所示：当$x$无限地靠近$a$，但总不能等于$a$时，曲线$y=f(x)$上的点$(x,f(x))$无限地靠近$(a,\ell)$点．

\begin{figure}[htp]
    \centering
\begin{tikzpicture}[>=latex]
\draw[->](-3,0)--(5,0)node[right]{$x$};  
\draw[->](0,-1)--(0,3)node[right]{$y$};    
\draw[domain=-2:4, samples=100, very thick]plot(\x, {.5*cos(\x r)+1});
\node at (0,0)[below right]{$O$};
\node at (-1,1.5){$y=f(x)$};
\draw[dashed](pi/2,0)node[below]{$a$}--(pi/2,1)node[above]{$(a,\ell)$};

\tkzDefPoints{1.57/1/A, 1.3/1.133/B, 1.8/0.886/C}
\tkzDrawPoints(A,B,C)
\tkzLabelPoint[below left](B){$(x,f(x))$}
\tkzLabelPoint[above right](C){$(x,f(x))$}
\end{tikzpicture}
    \caption{}
\end{figure}

初学的人常常要问：为什么在定义中谈到$x$趋近于$a$时，要限制$x$始终不等于$a$呢？这是因为我们关心的是函数$f(x)$在$a$附近的变化趋势，它和函数$f(x)$在
$x=a$这一点的值没有什么必然关系，这也就是说，无论$f(x)$在点$a$取什么值甚至没有定义，都不影响在这一点的极限的存在和极限值．

\begin{example}
设$f(x)=\frac{3}{4}\cdot \frac{x^2-1}{x-1}$，$x\in(-\infty,1)\cup(1,+\infty)$，求$\Lim_{x\to 1}f(x)$
\end{example}

\begin{solution}
    $f(x)=\frac{3}{4}\cdot \frac{x^2-1}{x-1}$在$x=1$时无意义，因为那时
    分母就变成零，因此，这里没有函数值$f(1)$, 曲线$y=f(x)$也没有相应于横坐标为1的那个点，但是让$x$任意地趋近于1是完全可以的，若$x\ne 1$, 则有
  \[  f (x) =\frac{3}{4}\cdot  \frac{(x-1)(x+1)}{x-1}=\frac{3}{4}(x+1)\]
  因此，不论$x$通过怎样一个以1为极限的数列$\{x_n\}$, 对于相应的数列$\{f(x_n)\}$, 我们都有
  \[\lim_{x_n\to 1} f (x_n) =\frac{3}{4} (1+1) =\frac{3}{2}\]

    从几何上看，曲线$y=f(x)$除去点$\left(1,1\frac{1}{2}\right)$外是与直线$y=\frac{3}{4}(x+1)$一致的，唯独在那一点，曲线有个空隙，而在$x=1$的邻近的点只要充分接近于点1, 所对应的函数值与$\frac{3}{2}$的差的绝对值可以任意小，如图1.3．
\end{solution}
    
\begin{figure}[htp]
    \centering
\begin{tikzpicture}[>=latex]
    \draw[->](-3,0)--(5,0)node[right]{$x$};
    \draw[->](0,-1.5)--(0,4.5)node[right]{$y$};
\foreach \x in {1,2,3}
{
    \draw(0,\x)node[left]{$\x$}--(.1,\x);
    \draw(\x,0)node[below]{$\x$}--(\x,.1);
}
\node at (.25,-.25){$O$};
\draw[dashed](0,1.5)--(1,1.5)--(1,0);
\draw[domain=-2:4, samples=10, very thick]plot(\x, {0.75*(\x+1)});
\draw(1,1.5)[fill=white]  circle(1.5pt) node[right]{$\left(1,1\tfrac{1}{2}\right)$};
\node at (3.25,3)[right]{$y=\frac{3}{4}\cdot \frac{x^2-1}{x-1}$};
\end{tikzpicture}
    \caption{}
\end{figure}

\begin{example}
证明当$x\to 0$时，函数$f(x)=\sin\frac{1}{x}$没有极限．
\end{example}

\begin{solution}
函数$f(x)=\sin\frac{1}{x}$对于一切$x\ne 0$的值有定义，因此这个函数在点$x=0$的邻域内有定义．

当$x$取数列$\{x_n\}=\left\{\frac{2}{(2n+1)\pi}\Big| n=1,2,3,\ldots\right\}$的值而趋于零时，数列$\left\{\frac{1}{x_n}\right\}$相应的值是
\[\frac{3\pi}{2},\frac{5\pi}{2},\frac{7\pi}{2},\ldots,(2n+1)\frac{\pi}{2},\ldots\]
此时数列$\left\{\sin\frac{1}{x_n}\right\}$便交替地取$-1$和$+1$这两个数值，换言之
\[\sin\frac{1}{x_n}=(-1)^n,\qquad n=1,2,3,\ldots\]
因此，当$n\to\infty$时，数列$\left\{\sin\frac{1}{x_n}\right\}$不趋于任何极限值．这就证明了当$x\to 0$时，函数$f(x)=\sin\frac{1}{x}$的极限不存在．
\end{solution}

函数的图象大致如图1.4所示．曲线关于原点对称，在包含原点的每一个对称邻域$(-\delta, \delta)$内，曲线$y=\sin\frac{1}{x}$
在原点的邻近作无数次振动，且曲线的振幅恒为1, 虽将原点的邻域的长缩小，也不能减少振动的次数．

\begin{figure}[htp]
    \centering
\begin{tikzpicture}[>=latex]
    \draw[->](-4,0)--(4,0)node[right]{$x$};
    \draw[->](0,-2)--(0,2)node[right]{$y$};
    \foreach \x in {1,-1}
    {
        \draw[dashed](-4,\x)--(4,\x)node[right]{$y=\x$};
    }
\draw[domain=-4:-.05, samples=1000, thick]plot(\x, {sin(180/pi/\x)});
\draw[domain=.05:4, samples=1000, thick]plot(\x, {sin(180/pi/\x)});
\foreach \x/\xtext in {1/1, .637/\frac{2}{\pi}}
{
    \draw(\x, 0)node[below]{$\xtext$}--(\x,.1);
}
\end{tikzpicture}    
    \caption{}
\end{figure}

上面是用数列的极限来说明函数的极限，其实我们也可以直接定义函数极限．

\begin{blk}{定义3}
  如果函数$f$在点$a$邻域上有定义（可能去掉点$a$本身），并且对于任给的$\varepsilon>0$，总存在$\delta>0$，使得当$0<|x-a|<\delta$时，就有$|f(x)-\ell|<\varepsilon$, 那么就说$\ell$为当$x$趋近于$a$时，函数$f$在点$a$的极限值．
\end{blk}

我们对这个定义需要说明以下几点：
\begin{enumerate}
    \item 用定义3验证某数$\ell$是函数$f$在点$a$的极限的办法就是对于任给的$\varepsilon>0$, 要找到这样的正数$\delta$使得能够由不等式$|x-a|<\delta$推出不等式$|f(x)-\ell|<\varepsilon$，虽然$\varepsilon$是任意的正数，但是在找$\delta$的过程中，$\varepsilon$是固定不变的，$\delta$依赖于$\varepsilon$．
    \item 对于已给的$\varepsilon$, 只要证明有一个$\delta>0$存在就行．因为如果有一个$\delta$存在，把$\delta$再缩小一些，显然仍满足我们的要求．
    \item 不等式$|x-a|>0$只是说明$x\ne a$, 即把$x$等于$a$的情况去掉，这是因为我们关心的是函数$f$在点$a$附近的变化趋势，而和函数在$x=a$这点的值无关．
    \item 我们指出定义2和定义3是等价的．
\end{enumerate}

\begin{example}
用定义3证明$\Lim_{x\to 1}\frac{x^3-1}{x-1}=3$
\end{example}

\begin{proof}
    任给$\varepsilon>0$, 要找$\delta>0$, 使由$0<|x-1|<\delta$推出
    $\left|\frac{x^3-1}{x-1}-3\right|<\varepsilon$成立．

    当$x\ne 1$时，
\[\begin{split}
    \left|\frac{x^3-1}{x-1}-3\right|&=|(x^2+x+1)-3|\\
    &=|(x-1)(x+2)|=|x-1|\cdot |x+2|
\end{split}\]

要由$|x-1|\cdot |x+2|<\varepsilon$找$\delta$，显然，这里因子$|x+2|$引起了麻烦．为方便起见，先假定$0<|x-1|<1$，即取$\delta_1=1$，这样
\[0<|x-1|<1\quad\to \quad |x+2|=|(x-1)+3|\le |x-1|+3<4\]因此，要使
\[\begin{cases}
    0<|x-1|<1\\
    |x-1|\cdot |x+2|<\varepsilon
\end{cases}\]
只须
\[\begin{cases}
    0<|x-1|<1\\
    4|x-1|<\varepsilon
\end{cases}\to \quad \begin{cases}
    0<|x-1|<1\\
    |x-1|<\frac{\varepsilon}{4}
\end{cases}\]
由此可见，只须取$\delta=\min\left(1,\frac{\varepsilon}{4}\right)$，即取$\delta$为1与$\frac{\varepsilon}{4}$中的较小者．

$\therefore\quad $对于任意$\varepsilon>0$，取$\delta=\min\left(1,\frac{\varepsilon}{4}\right)$，则当$0<|x-1|<\delta$时，即有
\[\left|\frac{x^3-1}{x-1}-3\right|<\varepsilon\]
这也就证明了
\[\lim_{x\to 1}\frac{x^3-1}{x-1}=3\]
\end{proof}

\begin{example}
    证明$\Lim_{x\to a}\sqrt{x}=\sqrt{a}\quad (a>0)$
\end{example}

\begin{proof}
对于任意的$\varepsilon>0$, 我们必须找到一个$\delta>0$, 使得当$|x-a|<\delta$时，$|\sqrt{x}-\sqrt{a}|<\varepsilon$成立，因为
\[|\sqrt{x}-\sqrt{a}|=\frac{|x-a|}{\sqrt{x}+\sqrt{a}}<\frac{|x-a|}{\sqrt{a}}\]
所以要使$\frac{|x-a|}{\sqrt{a}}<\varepsilon$，只须$|x-a|<\sqrt{a}\varepsilon$

如果取$\delta=\sqrt{a}\varepsilon$，则$\frac{|x-a|}{\sqrt{a}}<\frac{\sqrt{a}\varepsilon}{\sqrt{a}}=\varepsilon$
因此，对于任意$\varepsilon>0$，取$\delta=\sqrt{a}\varepsilon$，则当$|x-a|<\delta$时，就有$|\sqrt{x}-\sqrt{a}|<\varepsilon$成立．这也就证明了
\[\lim_{x\to a}\sqrt{x}=\sqrt{a}\quad (a>0)\]
\end{proof}

有时虽然$f(x)$在某点的左边（或右边）没有定义，如图1.5中的$a,b$两点，我们也可以谈论$f$在点$a$或点$b$两点的极限，譬如对于所有小于$a$的数，$f$虽然没有定义，但是我们可以考察，当$x$从$a$的右侧趋近于$a$时，函数$f$的变化趋势，也就是考察$f$的单边极限是否存在．

\begin{figure}[htp]
    \centering
\begin{tikzpicture}[>=latex]
    \draw[->](-1,0)--(3,0)node[right]{$x$};
    \draw[->](0,-1)--(0,2.5)node[right]{$y$};
\draw[very thick](.5,1)to[bend right=-15]node[above]{$y=f(x)$} (2.5,1.5);
\draw(.5,1)--(.5,0)node[below]{$a$};
\draw(2.5,1.5)--(2.5,0)node[below]{$b$};
\node at (-.25,-.25){$O$};
\end{tikzpicture}
    \caption{}
\end{figure}

\begin{blk}{定义}
     设$f(x)$在区间$(a,b)$上有意义，如果任给$\varepsilon>0$, 总存在某个$\delta>0$使得当$x\in (a,a+\delta)$时，总有$|f(x)-\ell|<\varepsilon$，我们就说函数$f$在点$a$以$\ell$为\textbf{右极限}，记作：
\[\lim_{x\to a^+}f(x)=\ell\]
\end{blk}

类似地，可以定义\textbf{左极限}，只需把开区间$(a,a+\delta)$换成$(b-\delta,b)$就行了，并记作
\[\lim_{x\to b^-} f(x)=\ell\]

\begin{example}
    设函数$f(x)=\begin{cases}
        x-1,&x\le 1\\
x+1,&x>1
    \end{cases}\quad $（图1.6）

    求$\Lim_{x\to 1^-}f(x)$和$\Lim_{x\to 1^+}f(x)$
\end{example}

\begin{figure}[htp]
    \centering
    \begin{minipage}[t]{0.48\textwidth}
    \centering
    \begin{tikzpicture}[>=stealth, scale=.6]
        \draw[->](-3,0)--(4,0)node[right]{$x$};
        \draw[->](0,-3)--(0,4.5)node[right]{$y$};
        \draw(0,1)node[left]{1}--(.1,1);
        \draw(0,2)node[left]{2}--(.1,2);
        \draw(1,0)node[below]{1}--(1,.1);
        \node at (-.25,-.25){$O$};
        \draw[domain=-2:1, samples=10, very thick]plot(\x,{\x-1});
        \draw[domain=1:3, samples=10, very thick]plot(\x,{\x+1});
    \node at (-2,-3)[left]{$y=x-1$};
    \node at (3,4)[right]{$y=x+1$};
    \draw[dashed](1,0)--(1,2)--(0,2);
    \draw (1,2)[fill=white] circle(2.5pt);
    
    \end{tikzpicture}
    \caption{}
    \end{minipage}
    \begin{minipage}[t]{0.48\textwidth}
    \centering
    \begin{tikzpicture}[>=stealth, scale=.8]
        \draw[->](-3,0)--(4,0)node[right]{$x$};
    \draw[->](0,-2.5)--(0,3)node[right]{$y$};
    \foreach \x in {-2,-1,1,2,3}
    {
        \draw(\x,0)node[below]{$\x$}--(\x,.1);
    }
\foreach \x in {-2,-1,1,2}
{
    \draw(0,\x)--(.1,\x)node[right]{$\x$};
}
\foreach \x in {-2,-1,...,2}
{
    \draw[very thick](\x,\x)--(\x+1,\x);
    \draw (\x+1,\x)[fill=white]circle(2pt);
}
    \node at (.25,-.25){$O$};
    \end{tikzpicture}
    \caption{}
    \end{minipage}
    \end{figure}


\begin{solution}
\[\begin{split}
    \lim_{x\to 1^-}f(x)&=\lim_{x\to 1^-}(x-1)=0\\
    \lim_{x\to 1^+}f(x)&=\lim_{x\to 1^+}(x+1)=2
\end{split}\]
图1.6显示了上面的结果．
\end{solution}

\begin{example}
    函数$y=[x]$代表不超过$x$的最大整数，即若$n\le x<n+1$，$n\in\mathbb{Z}$，则$y=[x]=n$．（图1.7）

求$\Lim_{x\to 2^+}\frac{[x]}{x},\quad \Lim_{x\to 2^-}\frac{[x]}{x}$
\end{example}

\begin{solution}
    显然，当$x=2$时，$\frac{[x]}{x}=\frac{[2]}{2}=\frac{2}{2}=1$, 又
\[\frac{[x]}{x}=\begin{cases}
    \frac{1}{x},& x\in(1,2)\\
    \frac{2}{x},& x\in(2,3)
\end{cases}\]
所以
\[\begin{split}
    \lim_{x\to 2^+}\frac{[x]}{x}&=\lim_{x\to 2^+}\frac{2}{x}=1\\
     \lim_{x\to 2^-}\frac{[x]}{x}&=\lim_{x\to 2^-}\frac{1}{x}=\frac{1}{2}
\end{split}\]    
\end{solution}


下面的命题说明函数的极限与函数的单边极限的关系：

\begin{blk}{定理}
极限$\Lim_{x\to a} f(x)$
存在的必要和充分条件是左极限$\Lim_{x\to a^-} f(x)$和右极限$\Lim_{x\to a^+} f(x)$都存在，并且二者相等．
\end{blk}

\begin{proof}
    必要性
    
    如果$\Lim_{x\to a} f(x)=\ell$, 就是说任给$\varepsilon>0$, 总存在$\delta>0$,
使得当$0<|x-a|<\delta$, 即当$x\in (a-\delta,a)\cup(a,a+\delta)$时，有$|f(x)-\ell|<\varepsilon$．换言之，当$x\in (a-\delta,a)$和$x\in (a,a+\delta)$时，都有$|f (x) -\ell|<\varepsilon$，因此
\[\Lim_{x\to a^-} f(x)=\ell,\qquad \Lim_{x\to a^+} f(x)=\ell\]

充分性

如果$\Lim_{x\to a^-} f(x)=\ell$且$\Lim_{x\to a^+} f(x)=\ell$, 那么总存在$\delta_1>0$, 使得当$x\in(a-\delta_1,a)$时，有
$|f (x) -\ell|<\varepsilon$．

又存在$\delta_2>0$, 使得当$x\in (a,a+\delta_2)$时，有$|f (x) -\ell|<\varepsilon$．

取$\delta=\min(\delta_1,\delta_2)$，于是当$x\in(a-\delta,a+\delta)$时，有$|f (x) -\ell|<\varepsilon$．
这就是说：
\[\Lim_{x\to a} f(x)=\ell\]
\end{proof}

\begin{example}
说明$\Lim_{x\to 3}\frac{|x-3|}{x-3}$是否存在？
\end{example}

\begin{solution}
\[\begin{split}|x-3|&=\begin{cases}
    x-3, & x>3\\
    3-x, &x<3
\end{cases}\\
    \Lim_{x\to 3^-}\frac{|x-3|}{x-3}&=\Lim_{x\to 3^-}\frac{3-x}{x-3}=\Lim_{x\to 3^-}(-1)=-1\\
    \Lim_{x\to 3^+}\frac{|x-3|}{x-3}&=\Lim_{x\to 3^+}\frac{x-3}{x-3}=\Lim_{x\to 3^+}(1)=1
\end{split}\]

$\because\quad \Lim_{x\to 3^-}\frac{|x-3|}{x-3}\ne \Lim_{x\to 3^+}\frac{|x-3|}{x-3}$
    
$\therefore\quad \Lim_{x\to 3}\frac{|x-3|}{x-3}$
不存在．
\end{solution}

\subsection{函数值趋于无穷大}

如果函数$f$在点$a$的邻域上有定义（可能去掉点$a$本身）对于无论多么大的正数$G$, 总存在一个够小的正数$\delta$, 使得当$0<|x-a|<\delta$时，就有$|f(x)|>G$, 那么就说当$x$趋于$a$时，函数$f(x)$趋于\textbf{无穷大}，记作
\[\lim_{x\to a} f (x) =\infty\]

\begin{example}
    求证$\Lim_{x\to 0}\frac{1}{x}=\infty$
\end{example}
   
\begin{proof}
设$G$是任意给定的正数，我们要求出一个$\delta>0$, 使得当$|x|<\delta$时，$|f(x)|=\left|\frac{1}{x}\right|>G$．

事实上，要使$\frac{1}{|x|}>G$, 
只须$0<|x|<\frac{1}{G} $．取$\delta=\frac{1}{G}$, 于是当$|x|<\delta$时，就有$\frac{1}{|x|}>G$
因此，
$$\Lim_{x\to 0}\frac{1}{x}=\infty$$
\end{proof}

\begin{example}
    证明当$x\to 0$时，函数$f(x)=\frac{1}{x}\sin\frac{1}{x}\quad (x\ne 0)$不趋于无穷大．
\end{example}

\begin{proof}
如果自变量$x$取数列$\{x_n\}=\left\{\frac{2}{(2n+1)\pi}\Big|n=1, 2, 3,\ldots\right\}$的值趋于0时，$\sin\frac{1}{x}$在原点的任意邻域内无限次交替地取$-1, 1$这两个值，对于这些值，$|f(x)|=\frac{1}{x}=\frac{\pi}{2}(2n+1)$趋于无穷大．

但是当$x$取数列$\{x'_n\}=\left\{\frac{1}{n\pi}\Big| n=1, 2, 3,\ldots\right\}$的
值趋于0时，由于
$\sin\frac{1}{x}=\sin (n\pi) =0$, 
故对于这些值，
$\Lim_{x'_n\to 0} f (x) =0$.

可见在原点的邻近不存在这样的$\delta>0$, 使得当$|x|<\delta$时，$|f(x)|>G$, 因此，当$x\to 0$时，$f(x)=\frac{1}{x}\sin\frac{1}{x}\quad (x\ne 0)$不趋于无穷大．

$y=f(x)$的图象位于两条双曲线$xy=\pm 1$之间，且在原点的邻近作无限多次振动，越靠近原点，曲线的振幅越大（图1.8）.

\begin{figure}[htp]
    \centering
\begin{tikzpicture}[>=latex, scale=.8]
\draw[->](-5,0)--(5,0)node[right]{$x$};
\draw[->](0,-4)--(0,4)node[right]{$y$};
\draw[domain=-4.5:-.3, samples=100]plot(\x,{-1/\x}); 
\draw[domain=-4.5:-.3, samples=100]plot(\x,{1/\x}); 
\draw[domain=.3:4.5, samples=100]plot(\x,{-1/\x}); 
\draw[domain=.3:4.5, samples=100]plot(\x,{1/\x}); 
\draw [domain=.5:4.5, samples=200,  thick]plot(\x, {sin(180/\x*pi)/\x});
\draw [domain=.5:4.5, samples=200,  thick]plot(-\x, {sin(180/\x*pi)/\x});
\end{tikzpicture}
    \caption{}
\end{figure}

\end{proof}

如果对于任何$G>0$, 存在$\delta>0$, 当$0<|x-a|<\delta$时，有$f(x)>G$, 就说当$x\to a$时，函数$f(x)$趋于正无穷大，记作
\[\lim_{x\to a}f(x)=+\infty\]

如果对于任何$G>0$, 存在$\delta>0$, 当$0<|x-a|<\delta$时，有$f(x)<-G$, 就说当$x\to a$时，函数$f(x)$趋于负无穷大，记作
\[\lim_{x\to a}f(x)=-\infty\]

例如，我们有：
\[\lim_{x\to 0}\frac{1}{x^2}=+\infty,\qquad \lim_{x\to 0}\frac{(-1)}{x^2}=-\infty\]

类似地，我们可以定义：
\[\lim_{x\to a^-}f(x)=+\infty,\quad \lim_{x\to a^+}f(x)=+\infty,\quad \lim_{x\to a^-}f(x)=-\infty,\quad \lim_{x\to a^+}f(x)=-\infty\]
的含义，这里不再写出．建议读者将这些定义严格地写出来．

例如，我们有：
\[\lim_{\theta\to \tfrac{\pi}{2}^-}\tan\theta =+\infty,\qquad \lim_{\theta\to \tfrac{\pi}{2}^+}\tan\theta =-\infty\]
\[\lim_{x\to 0^+}\log_a x =-\infty\quad (a>1),\qquad \lim_{x\to 0^+}\log_a x =+\infty\quad (0<a<1)\]

\begin{ex}
\begin{enumerate}
    \item 用函数极限定义证明：
    \begin{multicols}{2}
        \begin{enumerate}
            \item $\Lim_{x\to \infty}\frac{1}{2x+1}=0$
            \item $\Lim_{x\to 2}x^2=4$
            \item $\Lim_{x\to -1}\frac{x-3}{x^2-9}=\frac{1}{2}$
            \item $\Lim_{x\to 1}\frac{1}{x^2}=1$
            \item $\Lim_{x\to 1}\frac{x^3-x}{x-1}=2$
        \end{enumerate}
    \end{multicols}
    \item 说明：$\Lim_{x\to 3}\frac{x}{x^2-9}=\infty$
    \item 下面极限是否存在？
\begin{multicols}{2}
    \begin{enumerate}
        \item $\Lim_{x\to 1}\frac{2x|x-1|}{x-1}$
        \item $\Lim_{x\to 3}\frac{[x]^2-9}{x^2-9}$
    \end{enumerate}
\end{multicols}
\end{enumerate}
\end{ex}

\subsection{函数极限算法定理}
函数的极限算法定理与数列的极限算法定理类似，因为所谓$\Lim_{x\to a}u(x)=A$, $\Lim_{x\to a}v(x)=B$的意思就是对于任何一个各项都不同于$a$并且以$a$为极限的数列$x_n\to a$, 便有函数值数列$\{u(x_n)\}$, $\{v(x_n)\}$, 并且$\Lim_{n\to \infty}u(x_n)=A$, 
$\Lim_{n\to \infty} v(x_n)=B$, 因此，根据第四册下第三章的定理就可以直接得到相应的结果．现在给出函数的极限运算定理如下：

\begin{blk}{定理}
设$\Lim_{x\to a}u(x)=A$, $\Lim_{x\to a}v(x)=B$，那么
\begin{enumerate}
    \item $\Lim_{x\to a}[u(x)+v(x)]=\Lim_{x\to a}u(x)+\Lim_{x\to a}v(x)$
    \item $\Lim_{x\to a}[u(x)\cdot v(x)]=\Lim_{x\to a} u(x)\cdot \Lim_{x\to a}v(x)$
    \item $\Lim_{x\to a}[c\cdot u(x)]=c\cdot \Lim_{x\to a}u(x)$
    \item $\Lim_{x\to a}\frac{u(x)}{v(x)}=\frac{\Lim_{x\to a}u(x)}{\Lim_{x\to a}v(x)}$ 
    
    只要$v(x)$恒不为0，而且$\Lim_{x\to a}v(x)\ne 0$
    \item $\Lim_{x\to a}\sqrt[n]{u(x)}=\sqrt[n]{\Lim_{x\to a}u(x)}$
    \item 如$\Lim_{x\to a}u(x)=A=\Lim_{x\to a} v(x)$，而$|x-a|<\delta$时，$u(x)<f(x)<v(x)$，则：
    \[\Lim_{x\to a} f(x)=A\]
    即：如果$u(x)$与$v(x)$趋向同一极限$A$，且$f(x)$在$u(x)$与$v(x)$之间，那么，$f(x)$便也趋向那个极限$A$．
\end{enumerate}
\end{blk}

现在只有5需要补证．

设$\Lim_{x\to a}u(x)=A>0$, 则根据函数极限定义：对于$\varepsilon=\frac{A}{2}>0$, 存在$\delta>0$，使得$0<|x-a|<\delta$时，有
\begin{equation}
    u(x)>\frac{A}{2}>0
\end{equation}

于是，
\[\sqrt[n]{u(x)}-\sqrt[n]{A}=\frac{u(x)-A}{\sum^n_{k=1}[u(x)]^{\tfrac{n-k}{n}}A^{\tfrac{k-1}{n}}}\]
由(1.3)式
\[[u(x)]^{\tfrac{n-k}{n}}>\left(\frac{A}{2}\right)^{\tfrac{n-k}{n}}\]
故
\[\begin{split}
    \sum^n_{k=1}[u(x)]^{\tfrac{n-k}{n}}A^{\tfrac{k-1}{n}}&>\sum^n_{k=1}\left(\frac{A}{2}\right)^{\tfrac{n-k}{n}}A^{\tfrac{k-1}{n}}\\
    &=A^{\tfrac{n-1}{n}}\sum^n_{k=1}\left(\frac{1}{2}\right)^{\tfrac{n-k}{n}}\\
    &=A^{\tfrac{n-1}{n}}\left\{\left(\frac{1}{2}\right)^{\tfrac{n-1}{n}}+\left(\frac{1}{2}\right)^{\tfrac{n-2}{n}}+\cdots+\left(\frac{1}{2}\right)^{\tfrac{1}{n}}+1\right\}\\
    &=A^{\tfrac{n-1}{n}}\left\{\frac{1-\left(\frac{1}{2}\right)^{\tfrac{n-1}{n}}\cdot \left(\frac{1}{2}\right)^{\tfrac{1}{n}}}{1-\left(\frac{1}{2}\right)^{\tfrac{1}{n}}}\right\}\\
    &=A^{\tfrac{n-1}{n}}\cdot\frac{\frac{1}{2}}{1-\left(\frac{1}{2}\right)^{\tfrac{1}{n}}}=\frac{A^{\tfrac{n-1}{n}}}{2-2^{\tfrac{n-1}{n}}}\\
\end{split}\]
$\therefore\quad \left|\sqrt[n]{u(x)}-\sqrt[n]{A}\right|<|u(x)-A|\cdot \frac{2-2^{\tfrac{n-1}{n}}}{A^{\tfrac{n-1}{n}}}$

由题设$\Lim_{x\to a}u(x)=A$，而$\frac{2-2^{\tfrac{n-1}{n}}}{A^{\tfrac{n-1}{n}}}$是一个与$x$无关的常数，所以
\[\lim_{x\to a} \sqrt[n]{u(x)}=\sqrt[n]{A}=\sqrt[n]{\lim_{x\to a}u(x)}\]

下面我们来证明两个重要极限公式，为此先介绍一个引理．

\begin{blk}{引理}
     对于任意实数$\theta$, 都有$|\sin\theta|\le |\theta|$
\end{blk}

\begin{proof}
  由于对于任意实数．有$|\sin\theta|\le 1$, 那么仅考虑$\theta\in (0,\pi/2)$即可．
  
  在单位圆$O$上（图1.9），
 截取弧$\wideparen{P_0M}$使弧长$|\wideparen{P_0M}|$等于$\theta$, 其中$P_0$和$M$分别有坐标$(1, 0)$和$(x,y)$, 于是
 \[\sin\theta =y<\sqrt{y^2+ (1-x)^2} =|P_0M|<|\wideparen{P_0M}|=\theta\]
 
 显然，当$\theta=0$时，有$\sin\theta=\theta$. 因此
\[  |\sin\theta | \le |\theta|,\qquad  \theta\in  (0,\pi/2)\]

  如果$-\pi/2<\theta<0$, 则仍有$|MN|<\wideparen{P_0M}$，于是
\[|\sin\theta|<|\theta|,\qquad \theta\in (-\pi/2, 0) \]

  如果$|\theta|\ge \pi/2$, 则因为$\pi/2>1$与$|\sin\theta|\le 1$, 而同样地也得到$|\sin\theta|<|\theta|$. 
  
  因此，对于一切$\theta\ne 0$的值，有$|\sin\theta|<|\theta|$; 当$\theta=0$时，有$|\sin\theta|=|\theta|$．
\end{proof}



\begin{figure}[htp]\centering
    \begin{minipage}[t]{0.48\textwidth}
    \centering
\begin{tikzpicture}[>=latex]
\draw[->](-2.5,0)--(2.5,0)node[right]{$x$};
\draw[->](0,-2)--(0,2.5)node[right]{$y$};
\draw(0,0) circle(1.5);
\draw[very thick](0,0)--(45:1.5)node[above right]{$M(x,y)$}--(1.5,0)node[below right]{$P_0(1,0)$};
\draw[dashed](1.5/1.414,0)node[below]{$N$}--(1.5/1.414,1.5/1.414);
\node at (0,0)[below left]{$O$};
\end{tikzpicture}    
    \caption{}
    \end{minipage}
    \begin{minipage}[t]{0.48\textwidth}
    \centering
    \begin{tikzpicture}[>=latex, scale=1.3]

\draw[->](-1.5,0)--(2.5,0)node[right]{$x$};
\draw[->](0,-1.5)--(0,2)node[right]{$y$};
\draw(0,0) circle(1);
\draw[very thick](0,0)--(60:2)node[above right]{$C$}--(1,0)node[below right]{$A(1,0)$};
\draw[thick](.5,0)node[below]{$B$}--(60:1)node[above]{$D$};
\node at (0,0)[below left]{$O$};   
\draw(0.2,0) arc (0:60:.2)node[right]{$\theta$};   
    \end{tikzpicture}
    \caption{}
    \end{minipage}
    \end{figure}

\begin{blk}{定理}
\[\lim_{\theta\to 0}\frac{\sin\theta }{\theta}=1,\qquad \lim_{\theta\to 0}\frac{\cos\theta-1}{\theta}=0\]
\end{blk}

\begin{proof}
    让我们先证明$\Lim_{\theta\to 0}\frac{\sin\theta }{\theta}=1$．由图1.10容易看出：

假定$\theta$的单位是弧度，于是：
\[\begin{split}
    \triangle OBD\text{的面积}&=\frac{1}{2}OB\cdot BD=\frac{1}{2}\cos\theta\cdot \sin\theta\\
    \triangle OAC\text{的面积}&=\frac{1}{2}OA\cdot AC=\frac{1}{2}\cdot 1\cdot \tan\theta\\
    \text{扇形$OAD$的面积}&=\frac{1}{2}OA\cdot \wideparen{AD}=\frac{1}{2}\cdot 1^2\cdot \theta
\end{split}\]

假如角是用180等分平角的“度”作为单位，则扇形面积就是$\frac{1}{2}\cdot 1\cdot \frac{\pi}{180}\theta$．

因为上述扇形是夹在$\triangle OBD$和$\triangle OAC$之间，所以
\begin{equation}
    \frac{1}{2}\cos\theta\sin\theta<\frac{1}{2}\theta<\frac{1}{2}\tan\theta=\frac{1}{2}\frac{\sin\theta}{\cos\theta}
\end{equation}
即有，若$0<\theta<\frac{\pi}{2}$，则
\begin{equation}
    \cos\theta<\frac{\sin\theta}{\theta}<\frac{1}{\cos\theta}
\end{equation}
注意到：$\sin(-\theta)=-\sin\theta$和$\cos(-\theta)=\cos\theta$, 不等式(1.5)也蕴含，若$-\frac{\pi}{2}<\theta<0$, 则
\begin{equation}
    \cos(-\theta)<\frac{\sin(-\theta)}{-\theta}<\frac{1}{\cos(-\theta)}
\end{equation}
将不等式(1.5)和(1.6)合并为，若$0<|\theta|<\frac{\pi}{2}$，则
\begin{equation}
    \cos\theta<\frac{\sin\theta}{\theta}<\frac{1}{\cos\theta}
\end{equation}
又因为
\[0\le |1-\cos\theta|=\left|2\sin^2\frac{\theta}{2}\right|\le 2\left|\sin\frac{\theta}{2}\right|<|\theta|\]
从而
\[0\le \lim_{\theta\to 0}|1-\cos\theta|\le \lim_{\theta\to 0}|\theta|=0\]
所以
\[\lim_{\theta\to 0}|1-\cos\theta|=0\]
即：$\Lim_{\theta\to 0}\cos\theta=1$
由(1.7)知被夹逼在两者之间的$\frac{\sin\theta}{\theta}$极限值也就一定是1了！

以$\Lim_{\theta\to 0}\frac{\sin\theta}{\theta}=1$为基础，则：
\[\begin{split}
\lim_{\theta\to 0}\frac{\cos\theta-1}{\theta}&=\lim_{\theta\to 0}\left\{\frac{-2\left(\sin\frac{\theta}{2}\right)^2}{\theta}\right\}\\
&= \lim_{\theta\to 0}\left(-\sin\frac{\theta}{2}\right)\frac{\sin\frac{\theta}{2}}{\frac{\theta}{2}}    \\
&=\lim_{\theta\to 0}\left(-\sin\frac{\theta}{2}\right)\cdot \lim_{\theta\to 0}\frac{\sin\frac{\theta}{2}}{\frac{\theta}{2}}   \\
&=0\cdot 1=0
\end{split}   
\]
\end{proof}

\begin{example}
求：(1) $\Lim_{x\to 2}\frac{x^2+3x-10}{3x^2-5x-2}$；\qquad (2) $\Lim_{x\to 0}\frac{\sqrt[3]{1+x^2}-1}{x^2}$
\end{example}

\begin{solution}
    如果应用1.3中的定理去分别求(1)、(2)中的分子、分母的极限时，其分子、分母的极限值均为零，即我们得到$\frac{0}{0}$的不定形式．我们通常要对原代数式变形，找出分子、分母中具有$x-a$的因子，消去$x-a$的公共因子，一般问题就可解决．

(1)
\[\begin{split}
    \Lim_{x\to 2}\frac{x^2+3x-10}{3x^2-5x-2}&=\Lim_{x\to 2}\frac{(x+5)(x-2)}{(3x+1)(x-2)}\\
    &=\Lim_{x\to 2}\frac{x+5}{3x+1}\\
    &=\frac{2+5}{3\x 2+1}=1
\end{split}\]

(2)
\[\begin{split}
    \Lim_{x\to 0}\frac{\sqrt[3]{1+x^2}-1}{x^2}&=\Lim_{x\to 0}\frac{{(1+x^2)}-1}{x^2\left(\sqrt[3]{(1+x^2)^2}+\sqrt[3]{1+x^2}+1\right)}\\
    &=\Lim_{x\to 0}\frac{1}{\sqrt[3]{(1+x^2)^2}+\sqrt[3]{1+x^2}+1}=\frac{1}{3}
\end{split}\]
\end{solution}

\begin{example}
    求：(1) $\Lim_{x\to 0}\frac{\cos x-1}{x^2}$；\qquad (2) $\Lim_{x\to \pi}\frac{\sin mx}{\sin nx}\;(m,n\in \mathbb{N})$
\end{example}

\begin{solution}
(1)
\[\begin{split}
    \Lim_{x\to 0}\frac{\cos x-1}{x^2}&=\Lim_{x\to 0}\frac{-2\left(\sin\frac{x}{2}\right)^2}{x^2}\\
    &=\frac{-2}{2^2}\Lim_{x\to 0}\frac{\sin\frac{x}{2}}{\frac{x}{2}}\cdot \Lim_{x\to 0}\frac{\sin\frac{x}{2}}{\frac{x}{2}}\\
    &=\left(-\frac{1}{2}\right)\cdot 1\cdot 1=-\frac{1}{2}
\end{split}\]

(2) 由三角函数诱导公式易知
\[\sin m(\pi-x)=\sin(m\pi-mx)=(-1)^{m-1}\sin mx\]
同样得到
\[\sin n(\pi-x)=(-1)^{n-1}\sin nx\]
因此：
\[(-1)^{m-n}\frac{\sin mx}{\sin nx}=\frac{\sin m(\pi-x)}{\sin n(\pi-x)}\]
两边乘以$(-1)^{m-n}$，得
\[\frac{\sin mx}{\sin nx}=(-1)^{m-n}\frac{\sin m(\pi-x)}{\sin n(\pi-x)}\]
设$y=\pi-x$, 则当$x\to \pi$, 有$y\to 0$, $my\to 0$, $ny\to 0$. 因此：
\[\begin{split}
    \Lim_{x\to \pi}\frac{\sin mx}{\sin nx}&=\lim_{y\to 0}(-1)^{m-n}\frac{\sin my}{\sin ny}\\
&=(-1)^{m-n}\lim_{y\to 0}\left(\frac{m}{n}\cdot \frac{\sin my}{my}\cdot \frac{ny}{\sin ny}\right)\\
&=(-1)^{m-n}\cdot \frac{m}{n}\cdot 1\cdot 1=(-1)^{m-n}\frac{m}{n}
\end{split}\]
\end{solution}

\begin{example}
由条件$\Lim_{x\to -\infty}\left(\sqrt[]{x^2-x+1}-ax-b\right)=0$，求出$a,b$的值．
\end{example}

\begin{solution}
$\because\quad \Lim_{x\to -\infty}\left(\sqrt{x^2-x+1}-ax-b\right)=0,\quad \Lim_{x\to -\infty}\frac{1}{x}=0$

$\therefore\quad \Lim_{x\to -\infty}\left(\sqrt{x^2-x+1}-ax-b\right)\cdot \frac{1}{x}=0$

即：
\[\Lim_{x\to -\infty}\left(\frac{\sqrt{x^2-x+1}}{x}-a-\frac{b}{x}\right)=0\]
又因为：当$x<0$时，
\[\sqrt{x^2-x+1}=\sqrt{x^2\left(1-\frac{1}{x}+\frac{1}{x^2}\right)}=-x\sqrt{1-\frac{1}{x}+\frac{1}{x^2}}\]
所以上式可写成
\[\lim_{x\to -\infty}\left(-\sqrt{1-\frac{1}{x}+\frac{1}{x^2}}-a-\frac{b}{x}\right)=0\]
即：$-1-a-0=0\quad \to \quad a=-1$

代入原条件，得
\[\lim_{x\to -\infty}\left(\sqrt{x^2-x+1}+x-b\right)=0\]
即：
\[\begin{split}
    b&=\lim_{x\to-\infty}\left(\sqrt{x^2-x+1}+x\right)\\
    &=\lim_{x\to-\infty} \frac{x^2-x+1-x^2}{\sqrt{x^2-x+1}-x}  \\
    &=\lim_{x\to-\infty}\frac{-x\left(1-\frac{1}{x}\right)}{-x\left(\sqrt{1-\frac{1}{x}+\frac{1}{x^2}}+1\right)}    \\
    &=\lim_{x\to-\infty}\frac{1-\frac{1}{x}}{\sqrt{1-\frac{1}{x}+\frac{1}{x^2}}+1} =\frac{1}{2}   \\
\end{split}\]

$\therefore\quad a=-1,\quad b=\frac{1}{2}$为所求．
\end{solution}

\begin{ex}
\begin{enumerate}
    \item 说明当$x\to \infty$时，函数$\frac{x^5-4}{x^3+x}$的变化趋势．
    \item 求极限
    \begin{enumerate}
\begin{multicols}{2}
    \item  $\Lim _{x \to 0} \frac{(x+2)^{2}}{x^{2}+4}$
    \item  $\Lim _{x \to 2} \frac{x^{2}-4}{x-2}$
    \item  $\Lim _{x \to 2} \frac{5 x^{2}+x-8}{x^{2}-4}$
    \item  $\Lim _{x \to 0} \frac{\left(4 x^{3}-3\right)(1-2 x)}{7 x^{3}-6 x+4}$
    \item  $\Lim _{x \to \infty} \frac{3 x^{5}}{x^{5}-x^{2}+1}$
    \item  $\Lim _{x \to \infty} \frac{\left(x^{2}-5\right)\left(x^{2}+7\right)}{x^{4}+35}$
    \item  $\Lim _{x \to 3} \frac{x^{2}-8 x+15}{x^{2}-7 x+12}$
    \item  $\Lim _{x \to -3} \frac{x^{2}-9}{x^{2}+9 x+18}$
    \item $\Lim _{x \to -1} \frac{x\left(x^{2}+4 x+3\right)}{x^{3}+3 x^{2}+5 x+3}$
    \item  $\Lim _{x \to 1} \frac{x^{3}+x^{2}-2}{x^{4}+2 x^{2}-2 x-1}$
    \item $\Lim _{x \to \infty} \frac{x+\sin x}{2 x+5}$
    \item  $\Lim _{x \to 0} \frac{(1+x)^{5}-(1+5 x)}{x^{2}+x^{5}}$
\end{multicols}
    \item  $\Lim _{x \to 0} \frac{(1+m x)^{n}-(1+n x)^{m}}{x^{2}}\qquad (m, n \in \mathbb{N})$
    \item  $\Lim _{x \to 1} \frac{x^{2}+x^{4}+\cdots+x^{2 n}-n}{x-1}$
\end{enumerate}

    \item 求极限  
\begin{enumerate}
\begin{multicols}{2}
\item $\Lim_{x\to 0^+}\sqrt{3x} $
\item $\Lim_{x\to 1^-}\sqrt{1-x} $
\item $\Lim_{x\to 3^-} 3[x]$
\item $\Lim_{x\to 1^+} [x+3]$
\item $\Lim_{x\to -3^+}\left(1+\sqrt{x+3}\right) $
\item $\Lim_{x\to 4^-} \left(\sqrt{4-x}+[x-1]\right)$
    \item $\Lim_{x\to 1^+}\frac{2x|x-1|}{x-1} $
    \item $\Lim_{x\to 4^-} ([x]-x)$
\end{multicols}
    \item $\Lim_{x\to 0^+}\frac{b}{x}\left[\frac{x}{a}\right] \qquad (a>0,\; b>0)$
    \item $\Lim_{x\to 0^+} \frac{x}{a}\left[\frac{b}{x}\right]\qquad (a>0,\; b>0)$
\end{enumerate}    

    \item 求极限  
\begin{enumerate}
\begin{multicols}{2}
\item $\Lim _{x \to  1} \frac{1-x}{\sqrt{1-x^{2}}}$;
\item  $\Lim _{x \to  1} \frac{1-\sqrt{x}}{1-x}$;
\item  $\Lim _{x \to  1} \frac{(2 x-3)(\sqrt{x}-1)}{2 x^{2}+x-3}$
\item  $\Lim _{x \to  2} \frac{x-2}{\sqrt{x^{2}-2}-\sqrt{2}}$
\item  $\Lim _{t \to +\infty}(\sqrt{1+t}-\sqrt{t})$;
\item  $\Lim _{x \to +\infty} \sqrt{x}(\sqrt{x+a}-\sqrt{x})$
\item  $\Lim _{x \to  5} \frac{\sqrt{x-1}-2}{x-5}$;
\item $\Lim _{x \to  0} \frac{\sqrt[3]{1+x}-\sqrt[3]{1-x}}{x}$
\item  $\Lim _{x \to  1} \frac{4 x+\sqrt{x-1}}{2 x-\sqrt{x+1}}$
\item  $\Lim _{x \to  1} \frac{x-1}{\sqrt{x^{2}-1}+\sqrt{x-1}}$
\item $\Lim _{x \to  a} \frac{\sqrt{x-a}+\sqrt{x}-\sqrt{a}}{\sqrt{x^{2}-a^{2}}}$
\item  $\Lim _{x \to -\infty} \frac{x-2}{\sqrt{x^{2}-4 x+3}}$
\item $\Lim _{x \to -\infty}\left[\frac{2(x-1)^{2}}{\sqrt{4 x^{2}+2 x+1}}+x\right]$
\end{multicols}
\item $\Lim _{x \to +\infty}\left(\sqrt{x+\sqrt{x+\sqrt{x}}}-\sqrt{x}\right)$
\item $\Lim _{x \to +\infty} x\left(\sqrt{x^{2}+2 x}-2 \sqrt{x^{2}+x}+x\right)$
\end{enumerate}   

\item \begin{enumerate}
    \item $f(x)=\begin{cases}
        0,& x>1\\
        1,& x=1\\
        x^2+2, &x<1
    \end{cases}$

    求$f(x)$在$x=1$的左右极限．

    \item $f(x)=\begin{cases}
        x\sin\frac{1}{x},& x>0\\
        1+x^2,& x<0
    \end{cases}$

    求$f(x)$在$x=0$的左右极限．

    \item $\Lim_{x\to 3}\frac{[x]^2-9}{x^2-9}$是否存在？
    \item $\Lim_{x\to 0}\frac{x}{|x|}$是否存在？
\end{enumerate}

\item 若$0<x<\frac{\pi}{2}$，求证：
\begin{multicols}{2}
    \begin{enumerate}
        \item $\tan x>x$
        \item $x-\frac{x^2}{4}<\sin x<x$
    \end{enumerate}
\end{multicols}

\item 求极限  
\begin{enumerate}
\begin{multicols}{2}
    \item $\Lim_{x\to 0} \frac{\sin ax}{x}$
    \item $\Lim_{x\to 0} \frac{\sin 7x}{4x}$
    \item $\Lim_{\theta\to 0} \frac{\sec2\theta\sin\theta}{\theta}$
    \item $\Lim_{\theta \to 0} \frac{\tan\theta}{\theta}$
    \item $\Lim_{t\to 0} \frac{15t}{\tan 6t}$
    \item $\Lim_{x\to 0} \frac{\tan 2x}{\sin 7x}$
    \item $\Lim_{\theta \to 0} \frac{\tan\theta-\sin\theta}{\theta^3}$
    \item $\Lim_{x\to 0} \frac{\csc x-\cot x}{x}$
    \item $\Lim_{x\to 0} \frac{\sin^2\frac{x}{2}}{x^2}$
    \item $\Lim_{x\to \pi} \frac{\sin x}{\pi-x}$
    \item $\Lim_{x\to 1} \frac{1+\cos\pi x}{\tan^2 \pi x}$
    \item $\Lim_{x\to \tfrac{\pi}{4}}\tan 2x\tan\left(\frac{\pi}{4}-x\right) $ 
    \item $\Lim_{x\to 1} (1-x)\tan\frac{\pi x}{2}$
    \item $\Lim_{x\to 0} \frac{\sin mx}{\sin nx}$
    \item $\Lim_{x\to 0} \frac{x+\sin x}{x+2\sin x}$
    \item $\Lim_{x\to 0} \frac{\sqrt{\cos x}-1}{x^2}$
\end{multicols}
\end{enumerate}   

\end{enumerate}
\end{ex}

\subsection{数$e$}
我们在这里要利用数列的极限来定义一个新的数，这一个数不论对于分析本身或者对于它的应用来说都是非常重要的．

考虑数列$a_n=\left(1+\frac{1}{n}\right)^n,\; n=1, 2, 3,\ldots$的极限．

首先我们计算一下数列$\{a_n\}$的数值：
\begin{center}
\begin{tabular}{llll}
   $a_1=2.0$ & $a_2=2.250$ & $a_3=2.370$& $a_4=2.441$\\
$a_5=2.488$ &$a_6=2.522$& $a_7=2.546$& $a_8=2.565$\\
$\cdots$ &$a_{256}=2.712$&$\cdots$&$a_{1024}=2.717$\\ 
\end{tabular}
\end{center}

观察这一串数，我们会猜测这个数列是递增的，可能收敛于一个极限值，现在我们来证明这个猜想成立．

依牛顿二项式定理，即有
\[\begin{split}
    \left(1+\frac{1}{n}\right)^n&=1+n\cdot\frac{1}{n}+\frac{n(n-1)}{1\cdot 2}\cdot \frac{1}{n^2}+\frac{n(n-1)(n-2)}{1\cdot 2\cdot 3}\cdot \frac{1}{n^3}\\
    &\quad +\cdots+\frac{n(n-1)(n-2)\cdots[n-(n-1)]}{1\cdot 2\cdot 3\cdots n}\cdot \frac{1}{n^n}
\end{split}\]
这个等式的右边的第二项等于1，其他各项可以变换如下：
\[\begin{split}
\frac{n(n-1)}{1\cdot 2}\cdot \frac{1}{n^2}&=\frac{n(n-1)}{n^2}\cdot\frac{1}{2!}\\&=\left(1-\frac{1}{n}\right) \cdot\frac{1}{2!}\\
\frac{n(n-1)(n-2)}{1\cdot 2\cdot 3}\cdot \frac{1}{n^3}&=\frac{n(n-1)(n-2)}{n^3}\cdot\frac{1}{3!}\\&=\left(1-\frac{1}{n}\right)\cdot \left(1-\frac{2}{n}\right)\cdot\frac{1}{3!}\\
\cdots & \cdots\cdots\\
\frac{n(n-1)(n-2)\cdots[n-(n-1)]}{1\cdot 2\cdot 3\cdots n}\cdot \frac{1}{n^n}&=\frac{n(n-1)(n-2)\cdots[n-(n-1)]}{n^n}\cdot \frac{1}{n!}\\
&=\left(1-\frac{1}{n}\right) \left(1-\frac{2}{n}\right)\cdots\left(1-\frac{n-1}{n}\right)\cdot\frac{1}{n!}
\end{split}\]
其中$n!=1\cdot 2\cdot 3\cdots n$．因此：
\begin{equation}
\begin{split}
\left(1+\frac{1}{n}\right)^n&=1+1+\left(1-\frac{1}{n}\right)\frac{1}{2!}+ \left(1-\frac{1}{n}\right)\left(1-\frac{2}{n}\right)\cdot\frac{1}{3!}\\
&+\cdots+ \left(1-\frac{1}{n}\right) \left(1-\frac{2}{n}\right)\cdots\left(1-\frac{n-1}{n}\right)\cdot\frac{1}{n!}  
\end{split}
\end{equation}

用$n+1$代替$n$代入此公式，求出：
\begin{equation}
\begin{split}
    \left(1+\frac{1}{n+1}\right)^{n+1}&=1+1+\left(1-\frac{1}{n+1}\right)\frac{1}{2!}+ \left(1-\frac{1}{n+1}\right)\left(1-\frac{2}{n+1}\right)\cdot\frac{1}{3!}\\
    &+\cdots+ \left(1-\frac{1}{n+1}\right) \left(1-\frac{2}{n+1}\right)\cdots\left(1-\frac{n}{n+1}\right)\cdot\frac{1}{(n+1)!}  
\end{split}
\end{equation}
(1.9)的展开式项数比(1.8)的展开式的项数多一项，此外(1.9)的展开式各项由第三项起大于(1.8)的展开式对应项，这是因为
\[1-\frac{1}{n+1}>1-\frac{1}{n},1-\frac{2}{n+1}>1-\frac{2}{n},\ldots, 1-\frac{n-1}{n+1}>1-\frac{n-1}{n}\]
所以
\[\left(1+\frac{1}{n+1}\right)^{n+1}>\left(1+\frac{1}{n}\right)^n\]
这就是说数列$a_n=\left(1+\frac{1}{n}\right)^n,\; n=1,2,3,\ldots$是递增的．

但是由等式(1.8)可知
\[\left(1+\frac{1}{n}\right)^n<1+1+\frac{1}{2!}+\frac{1}{3!}+\cdots+\frac{1}{n!}\]
又因为
\[\frac{1}{1\cdot 2\cdot 3}<\frac{1}{2^2},\frac{1}{1\cdot 2\cdot 3\cdot 4}<\frac{1}{2^3},\ldots,\frac{1}{1\cdot 2\cdot 3\cdots n}<\frac{1}{2^{n-1}}\]
所以：
\[\left(1+\frac{1}{n}\right)^n<2+\frac{1}{2}+\frac{1}{2^2}+\cdots+\frac{1}{2^{n-1}}\]
但是由第二项起的级数总和小于$\frac{\tfrac{1}{2}}{1-\tfrac{1}{2}}=1$，所以：
\[\left(1+\frac{1}{n}\right)^n<2+\frac{\frac{1}{2}}{1-\frac{1}{2}}=3\]
这就是说数列$\left\{\left(1+\frac{1}{n}\right)^n\right\}$是递增有上界的，所以它有极限，设它的极限用$e$表示，因而我们有数
\[e=\Lim_{n\to \infty} \left(1+\frac{1}{n}\right)^n\]
这个数$e$是无理数，取它的十进小数到十五位，其值是
\[e=2.718281828459045\cdots\]

现在要把变数$n$推广到实数$x$, 这是一个以后常要用到的重要极限．

\begin{blk}{定理}
\[\Lim_{x\to \infty} \left(1+\frac{1}{x}\right)^x=e\qquad (x\in\mathbb{R})\] 
\end{blk}

\begin{proof}
先证$\Lim_{x\to +\infty} \left(1+\frac{1}{x}\right)^x=e$

令$[x]=n$，于是$n\le x<n+1$，所以，当$x\to +\infty$时，$n\to +\infty$，而且
\[1+\frac{1}{n+1}<1+\frac{1}{x}\le 1+\frac{1}{n}\]
根据指数函数的单调性和幂函数在$x>0$半直线上的单调性，有：
\[\left(1+\frac{1}{n+1}\right)^n\le \left(1+\frac{1}{n+1}\right)^x<\left(1+\frac{1}{x}\right)^x\le \left(1+\frac{1}{n}\right)^x<\left(1+\frac{1}{n}\right)^{n+1}\]
因为：
\[\begin{split}
    \Lim_{n\to +\infty} \left(1+\frac{1}{n}\right)^{n+1}&=\Lim_{n\to +\infty} \left(1+\frac{1}{n}\right)^{n}\left(1+\frac{1}{n}\right)\\
&=\Lim_{n\to +\infty} \left(1+\frac{1}{n}\right)^{n}\cdot \Lim_{n\to +\infty} \left(1+\frac{1}{n}\right)=e
\end{split}\]
\[\begin{split}
    \Lim_{n\to +\infty} \left(1+\frac{1}{n+1}\right)^{n}&=\Lim_{n\to +\infty} \frac{\left(1+\frac{1}{n+1}\right)^{n+1}}{1+\frac{1}{n+1}}\\
    &=\frac{\Lim_{n\to +\infty}\left(1+\frac{1}{n+1}\right)^{n+1} }{\Lim_{n\to +\infty}\left(1+\frac{1}{n+1}\right) }=e
\end{split}\]
所以
\[e\le\Lim_{x\to +\infty}\left(1+\frac{1}{x}\right)^{x} \le e\]
即：$\Lim_{x\to +\infty}\left(1+\frac{1}{x}\right)^{x}=e$

当$x\to -\infty$时，令$x=-y$，那么$y\to +\infty$，这时
\[\begin{split}
    \left(1+\frac{1}{x}\right)^x&=\left(1-\frac{1}{y}\right)^{-y}=\left(\frac{y-1}{y}\right)^{-y}=\left(\frac{y}{y-1}\right)^y\\
    &=\left(1+\frac{1}{y-1}\right)^y=\left(1+\frac{1}{y-1}\right)^{y-1}\left(1+\frac{1}{y-1}\right)
\end{split}\]
所以
\[\begin{split}
    \Lim_{x\to -\infty}\left(1+\frac{1}{x}\right)^{x}&=\Lim_{y\to +\infty}\left(1+\frac{1}{y-1}\right)^{y-1}\left(1+\frac{1}{y-1}\right)\\
    &=\Lim_{y\to +\infty}\left(1+\frac{1}{y-1}\right)^{y-1}\cdot \Lim_{y\to +\infty}\left(1+\frac{1}{y-1}\right)=e
\end{split}\]
综合上面两个结果，即得
$$\Lim_{x\to \infty}\left(1+\frac{1}{x}\right)^{x}=e$$
\end{proof}

\begin{blk}{推论}
\[\Lim_{y\to 0}\left(1+y\right)^{\tfrac{1}{y}}=e\]
\end{blk}

事实上，令$x=\frac{1}{y}$，当$y\to 0$时，$x\to\infty$，于是：
\[\Lim_{y\to 0}\left(1+y\right)^{\tfrac{1}{y}}=\Lim_{x\to \infty}\left(1+\frac{1}{x}\right)^{x}=e\]

\begin{ex}
\begin{enumerate}
    \item 求下面变量的极限：
\begin{multicols}{2}
\begin{enumerate}
    \item $\Lim_{x\to\infty}\left(1+\frac{8}{x}\right)^x$
    \item $\Lim_{y\to 0}(1+y)^{\tfrac{1}{3y}}$
    \item $\Lim_{t\to\infty}\left(1-\frac{1}{t}\right)^t$
    \item $\Lim_{t\to\infty}\left(\frac{t}{1+t}\right)^t$
    \item $\Lim_{n\to+\infty}\left(1+\frac{x}{n}\right)^n$
    \item $\Lim_{n\to +\infty}\left(1+\frac{1}{n^2}\right)^n$
    \item $\Lim_{x\to+\infty}\left(\frac{2x+3}{2x+1}\right)^{x+1}$
\end{enumerate}
\end{multicols}

\item 证明：
\begin{multicols}{2}
\begin{enumerate}
    \item $\Lim_{x\to\infty}\frac{a^x}{x}=\infty$
    \item $\Lim_{x\to\infty}\frac{a^x}{x^a}=\infty$
\end{enumerate}
\end{multicols}
其中$a$为任意正常数．
\end{enumerate}
\end{ex}

\section{函数的连续性}
\subsection{连续函数的概念}

在第四册下，我们用数列的语言规定了函数$f$在点$a$连续的意义．现在用$\varepsilon-\delta$语言来描述函数$f$在其定义域中的点$a$是连续的含义：

\begin{blk}
   {定义1} 函数$f$在其定义域M中的点$a$是连续的，如果对于每一个正数$\varepsilon$, 有$a$的一个$\delta$邻域$(a-\delta,a+\delta)$, 使得对于定义域中满足$|x-a|<\delta$的一切$x$值，不等式
\[|f (x) -f(a)|<\varepsilon\]
成立． 
\end{blk}

\begin{rmk}
在函数极限定义那里不要求$f(a)$存在，所以特别注意要$x\ne a$, 现在，在函数于某点$a$连续的定义中，只盼$\Lim_{x\to a} f(x)=f(a)$, 因此，保证$|f(x)-f(a)|<\varepsilon$
成立的充分条件$|x-a|<\delta$不再要求$x\ne a$了．
\end{rmk}

对于一个在点$a$处连续的函数，有
\[\lim_{x\to a}f (x) =f \left(\lim_{x\to a} x\right) =f (a)\]  
极限符号同函数符号可以改变次序或者说可以互换．

\begin{blk}
 {定义2 }如果函数$f(x)$在一个区间$(a,b)$的每一点都连续，我们就说，它在这个区间$(a,b)$上连续．   
\end{blk}

函数也可以只在一点的一侧连续，当
\[\lim_{x\to x_0^+} f (x) =f (x_0)\]
时，便说$f(x)$在$x_0$的右侧连续；当
\[\lim_{x\to x_0^-} f (x) =f (x_0)\]
时，便说$f(x)$在$x_0$的左侧连续．显然，在一点的两侧都连续时，必在这点上连续．

所谓函数$f(x)$在闭区间$[a,b]$上连续，除了它在$(a,b)$内每一点都连续外，在$a$点右侧连续，在$b$点左侧连续．

在一般自然现象中连续是常态，不连续则是在特殊情形的突变．例如，在火箭的发射过程中，随着火箭的燃烧，质量逐渐变化，当每一级火箭烧尽时，该级火箭的壳自行脱落，于是质量突然减小，质量变化如图1.11所示．所以，常用的函数一般只有几个特别的点是不连续的：如果$f(x)$在$x=a$处不是连续的，则称$a$为函数的间断点．对于左右极限存在的间断点叫\textbf{第一类间断点}，不属于第一类间断点的叫\textbf{第二类间断点}．

例如$y=[x]$在所有的整数点$n$有第一类间断点；$y=\tan x$在$x=\pi/2$有第二类间断点．

\begin{figure}[htp]
    \centering
    \begin{minipage}[t]{0.48\textwidth}
    \centering
    \begin{tikzpicture}[>=latex, scale=1]
        \draw[->](-.5,0)--(4,0)node[right]{$t$时间};
        \draw[->](0,-.5)--(0,4)node[right]{$M$质量};
    \draw[thick](0,3.5)--(2,2.5);
    \draw[thick](2,1.5)--(4,.5);
    \draw[dashed](2,2.5)--(2,0)node[below]{$t_0$};
    \draw(2,2.5)[fill=white] circle (2pt);
    \node at (0,0)[below right]{$O$};
    \end{tikzpicture}
    \caption{}
    \end{minipage}
    \begin{minipage}[t]{0.48\textwidth}
    \centering
    \begin{tikzpicture}[>=latex, scale=1.8]
        \draw[->](-.5,0)--(2.5,0)node[right]{$x$};
    \draw[->](0,-.5)--(0,2)node[right]{$y$};
    \draw[thick](0,0)node[below left]{$O$}--(1,1)--(2,0)node[below]{2};
    \draw[dashed](1,0)node[below]{1}--(1,1);
    \draw(0,1)node[left]{1}--(.05,1);
    \end{tikzpicture}
    \caption{}
    \end{minipage}
    \end{figure}

\begin{example}
    设函数
$f (x) =\begin{cases}
     x,& 0\le x\le 1 \\ 2-x,& 1\le x\le 2\\
\end{cases}$

问在$x=1$处是否连续？
\end{example}

\begin{solution}
  在$x=1$处的两侧，函数的对应规则由不同的式子表示，所以研究在$x=1$处函数是否连续，就需要从左右极限入手（图1.12）．
\[  \lim_{x\to 1^-} f (x) =\lim_{x\to 1^-} x=1,\qquad  \lim_{x\to 1^+} f (x) =\lim_{x\to 1^+} (2-x)=1\]
$\therefore\quad \Lim_{x\to 1}f(x)=1$

又$f(1)=1$, 故有$\Lim_{x\to 1}f(x)=f(1)$, 即$f(x)$在$x=1$处是连续的．
\end{solution}

\begin{example}
    函数$f(x)=x\sin\frac{1}{x}$在$x=0$处是否连续？如不连续，是第几类间断点？
\end{example}

\begin{solution}
    函数$f$只在$x=0$没有定义，故$f$在$x=0$处间断，它的定义域是$\mathbb{R}-\{0\}$.

因为$\left|\sin\frac{1}{x}\right|\le 1$，从而：    
\[\begin{split}
   & 0\le\left|x\sin\frac{1}{x}\right|\le |x|\\
&0\le \lim_{x\to 0}\left|x\sin\frac{1}{x}\right|\le \lim_{x\to 0}|x|=0\\
&\lim_{x\to 0}x\sin\frac{1}{x}=0
\end{split}\]
所以，$f$在$x=0$处左、右极限存在，这就是说$x=0$是第一类间断点，它的图象大致如图1.13所示．

\begin{figure}[htp]
    \centering
\begin{tikzpicture}[>=latex, scale=1.5]
\draw[->](-1.75,0)--(1.75,0)node[right]{$x$};
\draw[->](0,-1.75)--(0,1.75)node[right]{$y$};
\draw(-1.5,-1.5)--(1.5,1.5)node[right]{$y=x$};
\draw(-1.5,1.5)node[left]{$y=-x$}--(1.5,-1.5);
\draw[domain=0.08:1.5, samples=1000, thick]plot(\x, {\x*sin(1/\x r)});
\draw[domain=0.08:1.5, samples=1000, thick]plot(-\x, {\x*sin(1/\x r)});
\node at (0,0)[below left]{$O$};
\draw[dashed](2/pi,0)node[below]{$\frac{2}{\pi}$}--(2/pi,2/pi)--(-2/pi,2/pi);
\end{tikzpicture}
    \caption{}
\end{figure}


如果把上面的函数$f$加以扩充，即我们定义一个新函数$g$:
\[g(x)=\begin{cases}
    x\sin\frac{1}{x},& x\ne 0\\
    0,& x=0
\end{cases}\]
那么我们便得到一个在$x=0$处的连续函数．
\end{solution}

\begin{blk}
    {定义3} $\Lim_{x\to a}f(x)$存在，但不等于$f(a)$, 或者$f(a)$没有定义，则称$f$在$a$处有\textbf{可去间断点}．
\end{blk}

例1.16中的$x=0$就是$f$的可去间断点．

由函数在某一点连续的概念可以直接导出连续函数的一个重要性质如下：

\begin{blk}
{定理} 若函数$f(x)$在$x=a$连续，且$f(a)>0$, 则存在一个$a$的$\delta$邻域$(a-\delta,a+\delta)$, 使得对于每一个$x\in (a-\delta,a+\delta)$, $f(x)$都是正的．    
\end{blk}

\begin{proof}
    因为$\Lim_{x\to a} f(x)=f(a)>0$, 所以对于
$\varepsilon=\frac{1}{2}f(a)$, 必定存在一个$\delta$使得$|x-a|<\delta$时，有
\[f (x) > f (a) -\frac{1}{2}f (a) =\frac{1}{2}f (a) > 0\]
\end{proof}

\subsection{连续函数的运算}
连续函数的概念是函数概念和极限概念一个自然结合，由本章1.3函数极限算法定理直接得出下面定理．

\begin{blk}
 {定理1} 设$f$和$g$在$a$处连续，则
 \begin{enumerate}
     \item $f\pm g$在$a$处连续；
     \item $f\cdot g$在$a$处连续；
     \item $\frac{1}{g}$在$a$处连续（若$g(a)\ne 0$）；
     \item $\sqrt[n]{f}$在$a$处连续（若$f(a)>0$）．
 \end{enumerate}
\end{blk}

\begin{proof}
\begin{enumerate}
    \item 因为$f$和$g$在$a$处连续，所以
    \[\Lim_{x\to a} f (x) =f (a) ,\qquad  \lim_{x\to a} g (x) =g (a)\]
    根据1.3中的函数极限算法定理，得
   \[\begin{split}
    \Lim_{x\to a} (f\pm g)(x)=\Lim_{x\to a}[f(x)\pm g(x)]&=\Lim_{x\to a} f(x)\pm \Lim_{x\to a}g(x)\\
    &=f(a)\pm g(a)=(f\pm g)(a)
   \end{split} \]
    这正好就是断言$f+g$在$a$处连续．
    \item \[\begin{split}
        \Lim_{x\to a} (f\cdot g) (x)=\Lim_{x\to a}f(x)\cdot g(x)&=\Lim_{x\to a} f(x)\cdot\Lim_{x\to a}g (x)\\
        &=f (a) \cdot g (a) =(f\cdot g) (a) 
    \end{split}\]
    即$f\cdot g$在$a$处连续．
    \item 由于$g(a)\ne 0$, 又$g$在$a$处连续，根据2.1中的定理，知道存在点$a$的$\delta$邻域，使得对于每一个$x\in (a-\delta,a+\delta)$，有$g(x)\cdot g(a)>0$，因此
\[\Lim_{x\to a} \frac{1}{g}(x)=\frac{1}{\Lim_{x\to a} g(x)}=\frac{1}{g(a)}=\frac{1}{g}(a)\]

\item 证明留给读者．
\end{enumerate} 
\end{proof}

在第四册下，我们根据定理1证明了一些基本初等函数在其定义域上到处连续的命题，现在总结如下：

\begin{blk}
   {命题1} 任何多项式函数到处连续． 
\end{blk}

\begin{blk}{命题2}
    若$f$和$g$是两个多项式，$g\ne 0$, 那么有理函数$r=f/g$, 除去$g$的零点集合，函数$r$是有定义的，而且在其定义域上到处连续．
\end{blk}

\begin{blk}
   {命题3} $n$次算术方根函数$\sqrt[n]{x}$在$[0,+\infty)$上到处连续．
\end{blk}

因为每一个无理数可以用两串有理数数列从左、右两方面去逼近它，同样无理指数幂也可以用相应的有理指数幂去逼近它，我们用逼近法定义了无理指数幂，从而把有理指数函数开拓为在实数域上到处连续的实指数函数，于是得到下面的命题． 

\begin{blk}
    {命题4} 指数函数$a^x\; (a>0,\; a\ne 1)$到处连续．
\end{blk}

由反函数连续定理，直接得到

\begin{blk}
    {命题5} 对数函数$\log_a x$在$(0,+\infty)$上到处连续．
\end{blk}

现在将继续证明另外一些基本初等函数连续性的命题．

\begin{blk}
    {命题6} 三角函数$\sin x$和$\cos x$到处连续．
\end{blk}

\begin{proof}
设$x_0$是任意给定的实数，由三角函数的和差化积公式，得到
\[|\sin x-\sin x_0|=2\left|\sin\frac{x-x_0}{2}\right| \cdot \left|\cos \frac{x+x_0}{2}\right|\]
因为$\left|\cos\frac{x+x_0}{2}\right|\le 1$, 并根据1.3中的引理得到
\[|\sin x-\sin x_0|\le 2\left|\frac{x-x_0}{2}\right|\]
即
\[|\sin x-\sin x_0|\le \left|x-x_0\right|\]
当$x\to x_0$时，$|x-x_0|\to 0$, 从而$|\sin x-\sin x_0|\to 0$, 即
\[\lim_{x\to x_0} \sin x=\sin x_0\]
因为$x_0$是任意一点，所以$\sin x$到处连续．

同样证明，得到：当$x\to x_0$时，
\[|\cos x-\cos x_0|\le |x-x_0|\to 0\]
即
\[\lim_{x\to x_0} \cos x=\cos x_0\]
因此，$\cos x$到处连续．
\end{proof}
 
由反函数连续性定理直接得出下面命题．

\begin{blk}
    {命题7} 反三角函数$y=\arcsin x\; \left(-\frac{\pi}{2}\le y\le \frac{\pi}{2}\right)$,\quad $y=\arccos x\; \left(0\le y\le \pi\right)$,\quad $y=\arctan x\; \left(-\frac{\pi}{2}< y< \frac{\pi}{2}\right)$,\quad $y=\arccot x\; \left(0<y<\pi\right)$在它们各自的定义域上到处连续．
\end{blk}

\begin{blk}{命题8}
幂函数$x^{\alpha}$（$\alpha$是任何一个实数）在开区间$(0,+\infty)$上到处连续．
\end{blk}

\begin{proof}
为确定起见，假设$\alpha>0$, $n$是任何一个大于$\alpha$的整数
\[(x+h)^{\alpha} -x^{\alpha}=x^{\alpha} \left[ \left(1+\frac{h}{x}\right)^{\alpha} -1\right] \]

若$h>0$, 则$\left(1+\frac{h}{x}\right)^{\alpha} >1$，（实指数幂性质），并且
\begin{equation}
    1<\left(1+\frac{h}{x}\right)^{\alpha}<\left(1+\frac{h}{x}\right)^{n}\quad \text{(指数函数单调性)}
\end{equation}

若$h<0$, 将上面不等式方向反向，得到
\begin{equation}
    \left(1+\frac{h}{x}\right)^{n}<\left(1+\frac{h}{x}\right)^{\alpha}<1
\end{equation}
因为$\left(1+\frac{h}{x}\right)^{n}$是关于$h$的$n$次多项式，并且显然当$h\to 0$时，
\[\left(1+\frac{h}{x}\right)^{n}\to 1\]

于是，由不等式(1.10)和(1.11), 得到当$h\to 0$时，$\left(1+\frac{h}{x}\right)^{\alpha}\to 1$. 因此
\[\lim_{h\to 0}\left[(x+h)^{\alpha}-x^{\alpha}\right]=0\]
即：
\[\lim_{h\to 0}(x+h)^{\alpha}=x^{\alpha}\]
这就是说$x^{\alpha}$在$(0,+\infty)$上处处连续．
\end{proof}

下面介绍常用的复合函数连续性定理．
\begin{blk}
{定理2} 设$X$、$Y$、$Z$是实数集，函数$g:\,X\mapsto Y$, $f:\, g(X)(\subset Y)\mapsto Z$, 又$h:\, X\mapsto Z$定义为$f$和$g$的复合函数，即$h(x)=f (g (x) )\;  (x\in X)$.

若$y=g(x)$在$x=a$处连续；$z=f(y)$在$y=g(a)=b$处连续，那么复合函数$z=h(x)=f(g(x))$在点$a$处也连续．
\end{blk}

\begin{proof}
由于
\begin{equation}
    \lim_{x\to a} g(x)=g(a)=b,\qquad x\in X
\end{equation}    
又因为：
\begin{equation}
    \lim_{y\to b} f(y)=f(b),\qquad y\in g(X)\subset Y
\end{equation}

$\therefore\quad \Lim_{x\to a}h(x)=\Lim_{x\to a}f(g(x))=\Lim_{y\to b}f(y)=f(b)=f(g(a))=h(a)$

这就说明复合函数在点$a$的连续性．
\end{proof}

把上面的等式写成
\[\lim_{x\to a}f (g (x) ) =f (g (a) ) =f \left(\lim_{x\to a} g(x) \right)\]
这个等式表示可以这样求复合函数的极限$\Lim_{x\to a}f(g(x))$, 
如果内层函数$g(x)$在点$a$的极限存在，即$\Lim_{x\to a}g(x)=b$, 
而且$f$在点$b$连续，就可以把极限符号$\Lim_{x\to a}$与函数符号$f$互换，
也就是把极限运算移到内层函数上去施行．

\begin{example}
    说明$h(x)=\sqrt[4]{\frac{1}{1+x^2}}$到处连续，并求$\Lim_{x\to a}h(x)$
\end{example}

\begin{solution}
    $h(x)=\sqrt[4]{\frac{1}{1+x^2}}$可以看作$y=g(x)=\frac{1}{1+x^2},\; x\in\mathbb{R}$
和$z=f(y)=\sqrt[4]{y},\; y\in [0,+\infty)$的复合函数．

由于$g(x)=\frac{1}{1+x^2}$在$\mathbb{R}$上到处连续并且它的值域
\[g(\mathbb{R})=(0,1]\subset [0,+\infty)=Y\]
又幂函数$f(y)=\sqrt[4]{y}$在$Y=[0,+\infty)$上到处连续，因此$f(y)=\sqrt[4]{y}$在$g(\mathbb{R})=(0,1]$上也到处连续，于是
\[\lim_{x\to a}h(x)=f\left(\lim_{x\to a}g(x)\right)=\sqrt[4]{\frac{1}{1+a^2}}\]
\end{solution}


\begin{example}
    求$\Lim_{x\to 2}\sqrt{\frac{1-\cos\pi x}{4-x^2}}$
\end{example}

\begin{solution}
设$y=g(x)=\frac{1-\cos\pi x}{4-x^2}\; (x\ne \pm 2)$，$z=f(y)=\sqrt{y}\; (y\ge 0)$，且知$f$在$[0,+\infty)$上到处连续，于是
\[f(g(x))=\sqrt{\frac{1-\cos\pi x}{4-x^2}}\]
因此：
\[\begin{split}
    \lim_{x\to 2}g(x)&=\lim_{x\to 2}\frac{1-\cos \pi x}{4-x^2}\\
    &=\lim_{x\to 2}\frac{2\sin^2\frac{\pi}{2}x}{4-x^2}=\lim_{x\to 2}\frac{\pi\sin^2\frac{\pi}{2}(2-x)}{\frac{\pi}{2}(2-x)(2+x)}\\
    &=\lim_{x\to 2}\pi\left[\frac{\sin\frac{\pi}{2}(2-x)}{\frac{\pi}{2}(2-x)}\cdot \frac{1}{2+x}\cdot \sin\frac{\pi}{2}(2-x)  \right]\\
    &=\pi\cdot 1\cdot \frac{1}{4}\cdot 0=0\in [0,+\infty)
\end{split}\]

$\therefore\quad \Lim_{x\to 2}\sqrt{\frac{1-\cos\pi x}{4-x^2}}=\sqrt{\Lim_{x\to 2}\left(\frac{1-\cos\pi x}{4-x^2}\right)}=\sqrt{0}=0$
\end{solution}

\begin{example}
    求$\Lim_{x\to 0}\frac{\arcsin x}{x}$
\end{example}

\begin{solution}
设$\alpha=\arcsin x$，这里$-1\le x\le 1$，$-\frac{\pi}{2}\le\alpha\le \frac{\pi}{2}$，
于是，$\sin\alpha=x$．

根据反正弦函数的连续性，有
\[\lim_{x\to 0} \alpha=\lim_{x\to 0}\arcsin x=\arcsin 0=0\]
进行变量替换，得到
\[\lim_{x\to 0} \frac{\arcsin x}{x}=\lim_{x\to 0}\frac{\alpha}{\sin\alpha}=1\]
\end{solution}

下面根据指数函数与对数函数的连续性，我们要建立几个重要的极限：
\begin{enumerate}
    \item $\Lim_{x\to 0} \frac{\ln(1+x)}{x}=1$,\quad 这里$\ln$表示$\log_e$
    \item $\Lim_{x\to 0}\frac{e^x-1}{x}=1$
    \item $\Lim_{x\to 0}\frac{(1+x)^\alpha -1}{x}=\alpha$
\end{enumerate}

\begin{proof}
\begin{enumerate}
    \item $\Lim_{x\to 0} \frac{\ln(1+x)}{x}=\Lim_{x\to 0} \ln(1+x)^{1/x}$
    因为右端在对数号下的式子当$x\to 0$趋向于$e$, 所以根据对数函数的连续性，得到
\[\Lim_{x\to 0} \frac{\ln(1+x)}{x}=\Lim_{x\to 0} \ln(1+x)^{1/x}=\ln e=1\]
\item 令$e^x-1=y$, 于是当$x\to 0$时，由指数函数连续性有$y\to 0$, 其次，我们有$x=\ln(1+y)$, 因此，有
\[\Lim_{x\to 0}\frac{e^x-1}{x}=\Lim_{y\to 0}\frac{y}{\ln(1+y)}=\Lim_{y\to 0}\frac{1}{\ln(1+y)^{1/y}}=1\]
\item 我们令$(1+x)^{\alpha}-1=y$, 当$x\to 0$时，由幂函数的连续性，也有$y\to 0$, 在等式$(1+x)^{\alpha}=1+y$的两边取对数得到
\[\alpha\ln (1+x) =\ln (1+y) \]
利用这一关系式，我们把所给的表达式变形为
\[\frac{(1+x)^\alpha-1}{x}=\frac{y}{x}=\frac{y}{\ln(1+y)}\cdot \alpha\cdot\frac{\ln(1+x)}{x}\]
因此：
\[\Lim_{x\to 0} \frac{(1+x)^\alpha-1}{x}=   \Lim_{x\to 0} \frac{y}{\ln(1+y)}\cdot \alpha\cdot   \Lim_{x\to 0}  \frac{\ln(1+x)}{x}=\alpha\]
\end{enumerate}
\end{proof}

我们讨论了下面的函数：
\begin{enumerate}
    \item $f(x)=c$\quad （$c$为常数）
    \item $x^n$\quad（$n$是整数）
    \item $\sqrt[n]{x}\quad (x\ge 0)$
    \item $a^x\quad (a>0,\; a\ne 1)$
    \item $\log_a x$
    \item $x^\alpha$\quad （$\alpha$是任意实常数）
    \item $\sin x,\; \cos x, \; \tan x,\; \cot x$
    \item $\arcsin x,\; \arccos x, \; \arctan x,\; \arccot x$
\end{enumerate}
上面这些函数统称为\textbf{基本初等函数}．

凡是由基本初等函数及常数，经过有限次四则运算（加、减、乘、除）及有限次函数复合的步骤而得到的且由一个式子表示的函数，及可以化为这种形式的函数，统称为\textbf{初等函数}．

例如，下列函数都是初等函数：
\begin{enumerate}
    \item $f(x)=x\sin\frac{1}{x}+2^x\log_2 x$
    \item $\varphi(x)=\arctan\frac{1}{x}+\sin(x^2+1)$
\end{enumerate}

连续函数的四则运算定理，复合函数连续性定理及反函数连续性定理，在研究初等函数的连续性时具有重要作用．现在根据这些定理，立即得到一个重要结论：

\begin{blk}
 {定理} 一切初等函数在它有定义的任何区间内到处连续．
\end{blk}

\begin{ex}
\begin{enumerate}
    \item 试证 $y=|x|$是一个到处连续的函数．
    \item 说明下列函数到处连续：
\begin{multicols}{2}
\begin{enumerate}
    \item $f(x)=\frac{x^2}{1+x^2}$
    \item $f(x)=\sin|x|$
    \item $\varphi(x)=|\sin x|$
    \item $k(x)=\left\{\cos x-\sin|x|\right\}^3$
\end{enumerate}
\end{multicols}

\item 求出下列函数的间断点，并指出其类型．
\begin{multicols}{2}
\begin{enumerate}
    \item $f(x)=\frac{1}{x-2}$
    \item $f(x)=\frac{x-1}{x^2-1}$
    \item $h(x)=x\cos^2 \frac{1}{x}$
    \item $y=\frac{1}{\sin\pi x}$
    \item $y=\tan(x-1)$
    \item $y=\frac{\sin x}{|x|}$
    \item $f(x)=\begin{cases}
        \frac{\sin x}{|x|},&x\ne 0\\
        1,& x=0
    \end{cases}$
    \item $f(x)=\cot\frac{1}{x}$
    \item $f(x)=\begin{cases}
        \cos\frac{\pi x}{2},& |x|\le 1\\
        |x-1|,& |x|>1
    \end{cases}$
\end{enumerate}
\end{multicols}
\item 试证明下列函数在其定义域内是连续的：
\begin{enumerate}
    \item $f(x)=\begin{cases}
        x^2\sin\frac{1}{x},&x\ne 0\\
        0,&x=0
    \end{cases}$
    \item $\varphi(x)=\begin{cases}
        1-x,& x<0\\
        2x^2+1,& x=0\\
        2+(x-1)^3,& x>0
    \end{cases}$
\end{enumerate}
\item 求下列变量的极限：
\begin{multicols}{2}
    \begin{enumerate}
\item $\Lim_{x\to+\infty}\left(\sin\sqrt{x+1}-\sin\sqrt{x}\right)$
\item $\Lim_{x\to 1}\frac{1+\cos\pi x}{\tan^2 \pi x}$
\item $\Lim_{x\to 0}\frac{x\sin x}{1-\cos x}$
\item $\Lim_{x\to 0}\frac{x}{\sqrt{1-\cos x}}$
\item $\Lim_{x\to 0}\frac{\cos x-\cos 3x}{x^2}$
\item $\Lim_{x\to \tfrac{\pi}{4}}\tan 2x\tan\left(\frac{\pi}{4}-x\right)$
\item $\Lim_{h\to 0}\frac{\arctan (x+h)-\arctan x}{h}$
\item $\Lim_{x\to \infty}\arcsin \frac{1-x}{1+x}$
\item $\Lim_{x\to 0}\frac{1+\sin x-\cos x}{1+\sin px-\cos px}$
\item $\Lim_{x\to 0}\frac{\log_p(1+x)}{x}$
\item $\Lim_{x\to 0}\frac{a^x-1}{x}$
\item $\Lim_{n\to \infty}n\left(\sqrt[n]{x}-1\right)$
    \end{enumerate}
\end{multicols}
\end{enumerate}
\end{ex}

\subsection{闭区间上的连续函数的性质}

我们在2.1中已经定义了闭区间上连续函数的意义．对于这类函数，我们要介绍它的几个重要性质．

\begin{blk}
    {中间值定理} 设$y=f(x)$是一个在闭区间$[a,b]$上到处连续的函数，设$c$是一个介于$f(a)$和$f(b)$之间的常数，则必须存在一个介于$a$、$b$之间的实数$x_0$, 使得$f(x_0)=c$.
\end{blk}

用几何术语来说，$y=f(x),\; a\le x\le b$的图象是一条连结$P(a,f(a))$点和$Q(b,f(b))$点的连续曲线，而$P$、$Q$分居于直线$y=c$的两侧，则曲线$y=f(x),\; a\le x\le b$至少和直线$y=c$有一个交点（$x_0,\; f(x_0)=c$）．(图1.14).
\begin{figure}[htp]
    \centering
\begin{tikzpicture}[>=latex, xscale=1.8]
\draw[->](-1,-.2)--(3,-.2)node[right]{$x$};
\draw[->](.5,-2.5)--(.5,2.5)node[right]{$y$};
\draw[very thick, domain=-.5:2.4, samples=100]plot(\x, {\x*(\x-1)*(\x-2)})node[below left]{$y=f(x)$};
\draw[dashed](-.5,-.2)node[below left]{$a$}--(-.5,-1.875)node[below]{$P(a,f(a))$};
\draw[dashed](2.4,-.2)node[below]{$b$}--(2.4,1.344)node[above]{$Q(b,f(b))$};
\draw[thick](-.75,.2)node[above]{$y=c$}--(2.75,.2);
\draw[dashed](0.7909,.2)--(0.7909,-.2)node[below]{$x_0$};

\end{tikzpicture}
    \caption{}
\end{figure}


对这个定理，我们在第四册（下）4.3中利用二分逼近法及连续函数的性质给出了它的证明，这里不再重述．

\begin{blk}
 {最大值、最小值} 在闭区间$[a,b]$上连续的函数$f$一定在该区间上取到一个最大值和最小值．  
\end{blk}

我们先利用连续函数的图象来直观地说明这个性质，然
后用证明中间值定理的同样方法来证明这个定理．初学者可略过定理的证明，这不影响对本书后面内容的学习．

从图1.15可以看出，在$[a,b]$中，能够找到两点$x_1,x_2$使$f(x_1)=M$, $f(x_2)=m$, 而对于$[a,b]$中所有的$x$, 都有
\[m=f(x_2)\le f(x)\le f(x_1)=M\]


\begin{figure}[htp]
    \centering
    \begin{minipage}[t]{0.48\textwidth}
    \centering
    \begin{tikzpicture}[>=latex, yscale=1.5]
\draw[->](0,0)--(4,0)node[right]{$x$};
\draw[->](.5,-.5)--(.5,2)node[right]{$y$};
\draw[rounded corners=5pt, very thick](1,.5)--(1.25,1)--(1.5,1.3)--(1.75,1)--(2,0)--(2.25,-.5)--(2.6,0)--(2.8,.3)--(3,0)--(3.2,-.3);
\draw[dashed](1,.5)--(1,0)node[below]{$a$};
\draw[dashed](3.2,-.3)--(3.2,0)node[above]{$b$};
\draw[dashed](1.5,1.2)--(1.5,0)node[below]{$x_1$};
\draw[dashed](2.25,-.4)--(2.25,0)node[above]{$x_2$};
\node at (.5,0)[below left]{$O$};

    \end{tikzpicture}
    \caption{}
    \end{minipage}
    \begin{minipage}[t]{0.48\textwidth}
    \centering
    \begin{tikzpicture}[>=latex, scale=1.5]
\draw[->](-.5,0)--(2,0)node[right]{$x$};
\draw[->](0,-.5)--(0,2)node[right]{$y$};
\draw(1,0)node[below]{1}--(1,.1);
\draw(0,1)node[left]{1}--(.1,1);
\node at (0,0)[below right]{$O$};
\draw[very thick](0,0)--(1,1);
\tkzDefPoints{0/0/O, 1/0/B}
\tkzDrawPoints(O,B)
\draw (1,1)[fill=white] circle(1.5pt);
    \end{tikzpicture}
    \caption{}
    \end{minipage}
  \end{figure}


\begin{rmk}
    这个定理是假定$f$在整个闭区间上连续，如果在一点不连续，则定理的结论就可能不成立．
\end{rmk}

例如，函数$f(x)=x-[x]$, $x\in [0, 1]$在$x=1$这一点不连续，它的图象如图1.16, 我们要说明它有最小值0, 而无最大值．

事实上，$f$是分段函数：
\[f(x)=\begin{cases}
    x,& 0\le x<1\\
    0,& x=1
\end{cases}\]
$f$在$[0, 1)$上递增，故
$f(0)=f(1)=0$最小．又在$[0, 1]$上没有一点的函数值为1, 但是对于任何一个小于1的正数$k$, 恒有
\[f (k) =k- [k] =k<1\]
这就说明了$f$在$[0, 1]$上有上界1, 但始终达不到1, 现在我们进一步说明1是它的最小上界，假设定数$k'<1$是$f$在
$[0, 1]$上的上界，我们在$(k',1)$内任意取一个数$x$, 于是
\[f (x) =x- [x] =x> k'\]
这与假设$k'$是$f$在$[0, 1]$上的上界矛盾，因此1是$f$在$[0, 1]$上的最小上界，但是它不属于$f$的值域，这就说明了$f$无最大值．

从这个例中还可以看出，定理中的闭区间不能用开区间来代替．

现在给出定理的证明如下：

\begin{proof}
令$A_1=a$, $B_1=b$, $S_1=\{f(x)|x\in [a,b]\}$，设$K$是$S_1$的最小上界，记$K=l.u.b\{S_1\}$. 再令
\[S_1'=\left\{f (x) ; \; x\in \left[a,\frac{a+b}{2}\right]\right\},\quad S_1''=\left\{f(x);\; x\in\left[\frac{a+b}{2},b\right]\right\}\]
显然有：
\[S_1=S_1'\cup S_1''\]
由此不难看到：
\[K=l.u.b \{S_1\} =\max \left(l.u.b \{S'_1\} ,l.u.b\{S_1''\}\right)\]
所以$l.u.b\{S_1'\}$和$l.u.b\{S_1''\}$两者之中，至少有一个等于$K$, 换句话说，我们总可以在$[A_1,B_1]$的两个半段之中，选取其一为$[A_2,B_2]$, 使得
\[l.u.b \{S_2\} =K,\qquad  S_2= \{f (x) ,\; x\in [A_2, B_2] \}\] 
如此逐步二分，每次由$[A_{n} ,B_{n}]$归纳地选取其半段作为$[A_{n+1} ,B_{n+1}]$，使得
\[\begin{split}
    l.u.b \{f (x) ;\; x\in  [A_{n+1}, B_{n+1}] \}&=l.u.b \{f (x) ;\; x\in [A_n, B_n] \}\\
&=\cdots\\
&=l.u.b \{f (x) ;\; x\in  [a,b] \}=K.
\end{split}\]
对$\{A_n\},\{B_n \}$, 根据实数的连续性，知存在$c$使得$c\in ［A_n,B_n]$, $n=1, 2, 3,\ldots$, 且$B_n-A_n\to 0$.

现在我们要再利用$f(x)$到处连续这一性质，说明$f(c)＝K$. 假若不然，即$f(c)<K$，则可以取$\varepsilon=\frac{1}{2}(K-f(c))>0$，然后利用$f(x)$在$c$点的连续性，知存在点$c$的$\delta$邻域，使得当$|x-c|<\delta$时，有
\[|f(x)-f(c)|<\varepsilon\]
即：
\[\begin{split}
    f(x)<f(c)+\varepsilon&=f(c)+\frac{1}{2}K-\frac{1}{2}f(c)\\
    &=\frac{1}{2}[K+f(c)]\\
    &=\frac{1}{2}[K+K-2\varepsilon]=K-\varepsilon
\end{split}\]
所以
\[l.u.b \{f (x) ;\; x\in  (c-\delta,c+\delta) \} \le K-\varepsilon\]
但是，由$\Lim_{n\to\infty} A_n=\Lim_{n\to\infty}B_n=c$, 知当$n$足够大时，有
\[[A_n,B_n]\subset (c-\delta,c+\delta)\]
这显然和$[A_n,B_n]$的选取法，亦即保持
\[l.u.b \{f (x) ;\; x\in [A_n,B_n]\}=K\]
相矛盾，因此，$f(c)$必须等于$K$. 

类似地，可以证明存在使函数$f$达到它的最小值的一个点．
\end{proof}

\subsection*{习题1.2}
\begin{enumerate}
    \item 求下列函数的极限：
\begin{enumerate}\begin{multicols}{2}
    \item $\Lim_{x\to 0}(\cos x)^{\sin x}$
    \item $\Lim_{x\to 0}\frac{\frac{x^2}{e^{2x}}}{\ln\left(1+\frac{x^2}{e^{2x}}\right)}$
    \item $\Lim_{x\to 0}\frac{\frac{x^2}{e^{2x}}}{\sin^2 \frac{x}{e^{x}}}$
    \item $\Lim_{n\to \infty}n^2\left(x^{\tfrac{1}{n}}-x^{\tfrac{1}{n+1}}\right)$
    \item $\Lim_{x\to 0}\frac{\left(e^x+e^{-x}\right)\left(e^x-e^{-x}\right)}{2x}$
    \item $\Lim_{x\to 0}(1-2x)^{\tfrac{1}{x}}$
    \item $\Lim_{n\to \infty}10\left(1+\frac{0.03}{n}\right)^n$
    \item $\Lim_{x\to +\infty}\arccos\left(\sqrt{x^2+x}-x\right)$
    \item $\Lim_{x\to -\infty}\arctan\frac{x-4}{(x-2)^2}$
    \item $\Lim_{x\to -\infty}\arctan\frac{x}{\sqrt{1+x^2}}$
    \item $\Lim_{x\to -\infty}\frac{x}{\sqrt{1+x^2}}$
\end{multicols}   
    \item $\Lim_{x\to 0}\frac{\ln\frac{1+x}{1-x}}{\arctan(1+x)-\arctan(1-x)}$
\end{enumerate} 

\item 在一个以弦$AB$为界的弓形中，把点$A$与$B$与弧$\wideparen{AB}$的中点$C$用直线连结起来，于是得到一个等腰三角形$\triangle  ABC$, 再在$A$与$B$作切线，它们交于$D$, 试计算$\Lim_{x\to 0}\frac{\triangle ABC}{\triangle ABD}$，其中$x$是
弦$AB$所对应的中心角．

\item $f(x)=\frac{ax^3+bx^2+cx+d}{x^2-x-2}$满足条件
\[\lim_{x\to \infty} f(x)=2,\qquad \lim_{x\to 2} f(x)=3\]
求$f(x)$.
\end{enumerate}



\chapter{变率和微商}

从本章起，我们开始学习单变量的微积分学的基本概念和基础理论．

微积分学是研究变量的数学，变量之间的关系就是函数，因此，函数是微积分学研究的主要对象．

在函数的基本性质中，有两个最基本、最重要的概念－变率与求和，为了解决求函数的变率与求函数$f$在$[a,b]$上的和的问题就相应地产生微分与积分运算，而这两种运算之间，也存在着一种自然的互逆关系，在本章中，我们由函数的变率问题引出函数的微商（导数）概念，并给出初等函数的一套求导法则，在下一章中，我们揭示微分运算与积分运算的互逆关系，这就是微积分学的基本定理．

总起来说，微分反映了函数的局部性质，或在某个点附近的性质；积分则反映了函数的整体性质，或某个区间的性质；函数的局部性质与整体性质之间的有机联系，恰恰反映了微分运算与积分运算之间的互逆关系．

\section{微商（导数）的定义}
函数关系$y=f(x)$就是确定变量$y$如何随着变量$x$的变动而变动的关系，对于给定的函数$y=f(x)$, 变量$y$在变量$x$的不同点附近的变动情况是不尽相同的，这就是说，在变量$x$的
某个值$x_1$处，当$x_1$略加变动时，相应的$y$的变动可能相当剧烈（急增，或急减）；而在变量$x$的另一个值$x_2$处，$y$的变动就可能较为迟缓，但是，用这样的语言来表达函数在某一点处的变率是不精确的，我们需要用“数量”来确切地表达这个意思，这就是函数在某点（或某时刻）的变率的问题，简称变率．

\subsection{直线函数的变率}
一次函数$f (x) =kx+b$
是一种最简单的函数，它的函数图象是一条斜率等于$k$的直线，如图2.1所示．
\begin{figure}[htp]
    \centering
\begin{tikzpicture}[>=latex]
    \draw[->](-2,0)--(4,0)node[right]{$x$};
    \draw[->] (0,-1)--(0,5)node[right]{$y$};
    \draw[very thick] (-1.5,-.5)--(4,5);
\draw[->](-.2,0) arc (0:45:.8)node[below]{$\theta$};
\draw(2,3)node[left]{$P_1$}--(2,0)node[below]{$x_1$};
\draw(3,4)node[left]{$P$}--(3,0)node[below]{$x$};
\draw[|<->|](3.2,4)--node[right]{$f(x)-f(x_1)$}(3.2,3);
\draw[|<->|](2,2.8)--node[below]{$x-x_1$}(3,2.8);
\draw(2,3)--(3,3);\node at (.25,-.25){$O$};
\end{tikzpicture}
    \caption{}
\end{figure}

设$x_1$, $f(x_1)$和$x$, $f(x)$分别是直线上点$P_1$和$P$的坐标，为了反映函数变化快慢的问题，无论$x<x_1$还是$x>x_1$，自然地考虑在点$x_1$邻近，函数与自变量的相应的改变量的比：
\[\frac{f(x)-f(x_1)}{x-x_1}=\frac{(kx+b)-(kx_1+b)}{x-x_1}=k=\tan\theta\]
上面的表达式称为函数的\textbf{差商}，它表示函数在区间$[x_1,x]$或$[x,x_1]$上对于自变量$x$的\textbf{平均变化率}．由于$k$是不随变量$x$变动而变动的常数，因此，一次函数在自变量的任何一个区间内的平均变化率都是常数．

如果让自变量的变化区间的长度无限地缩短，也就是让$x$无限地接近于$x_1$时，平均变化率所趋向的极限
\[\lim_{x\to x_1} \frac{f (x) -f (x_1)}{x-x_1} =k\]

\subsection{平滑曲线的切线与变率}
一般的函数$y=f(x)$的图象通常不是直线，由于函数
和它的图象的多样性，为了讨论的方便起见，我们先把讨论的范围限制在“平滑”的曲线上，常用的函数$y=f(x)$的图象往往是“平滑的”，平滑性的直观内涵是：用愈高倍的显微镜去观察曲线的微段，就愈象直线段，比较明确的几何说法是：一条曲线在$P$点的平滑性就是存在唯一的一条过$P$点的切线，它无限地逼近曲线在$P$点邻近的微段，于是，当$y=f(x)$的图象$C$在$P$点存在唯一的一条切线时，我们就说曲线$C$在$P$点平滑，而$P$点叫做曲线的平滑点，一条在每点都平滑的曲线叫做平滑的曲线，一个函数的图象是平滑曲线时，我们就称这种函数为平滑函数．

同学可能会问这样一个问题：在曲线的点$P$存在唯一的一条切线的含义是什么？因为迄今我们对于一般的曲线的切
线还未下过定义呢！

我们从图2.2和2.3注意到不能把切线定义为与曲线只有一个交点的直线．这样的定义限制得既太紧同时又太松．因
为，照此定义，图2.2所示的直线就不是过曲线上$P$点的切线了，实际上，尽管图2.2的直线与曲线还有一个交点$Q$, 但它在曲
线$P$点邻近却与曲线密合，故仍应该是过曲线$P$点的切线；又图2.3表明过抛物线上任何一点$P$与$y$轴平行的直线虽然与抛物线只有一个交点，但它的其余部分却远离$P$点邻近的弧，故它不应该是抛物线的切线，定义切线的可行途径是从割线开始，并应用极限的概念．

\begin{figure}[htp]
    \centering
    \begin{minipage}[t]{0.48\textwidth}
    \centering
    \begin{tikzpicture}[>=latex, scale=.7]
\draw[->](-4,0)--(4,0)node[right]{$x$};
\draw[->](0,-2)--(0,3.5)node[right]{$y$};
\draw[domain=-3.2:3, very thick, samples=100]plot(\x, {.06*\x^3+.5})node[above]{$y=f(x)$};
\node at (0,0)[below left]{$O$};
\node at (1.5,.7025)[above]{$P$};
\draw[domain=-4:3.5, samples=10, thick]plot(\x, {.405*\x+.095});
\node at (-3,-1.12)[below right]{$Q$};
    \end{tikzpicture}
    \caption{}
    \end{minipage}
    \begin{minipage}[t]{0.48\textwidth}
    \centering
    \begin{tikzpicture}[>=latex, scale=.8]
  \draw[->](-3,0)--(3,0)node[right]{$x$};
\draw[->](0,-1)--(0,4)node[right]{$y$};
\draw[domain=-2:2, very thick, samples=100]plot(\x, {\x*\x});
\node at (-1.2,1)[left] {$f(x)=x^2$};
\draw(1,0.2)--(1,2);
\node at (1,1)[right]{$P$};
\node at (0,0)[below left]{$O$};
\draw[domain=.2:2, samples=10, thick]plot(\x, {2*\x-1});


    \end{tikzpicture}
    \caption{}
    \end{minipage}
  \end{figure}



如图\ref{fig:PQ_tangent}所示，取曲线$y=f(x)$上点$P$附近的另一点$Q$, 通过这两点画一条直线，这直线叫做过曲线上$P$点的割线，让$Q$点沿曲线向点$P$移动，这条割线将达到极限位置，此极限位置与$Q$点从哪一侧趋向于$P$是无关的，我们称这个割线的极限位置为过曲线上$P$点的切线．

\begin{figure}[htp]
    \centering
\begin{tikzpicture}[>=latex, scale=1.5]
    \draw[->](-.5,0)--(4,0)node[right]{$x$};
    \draw[->](0,-.5)--(0,3)node[right]{$y$};
    \node at (0,0)[below left]{$O$};
    \draw[domain=.5:3.2, samples=100, very thick]plot(\x, {-.5*(\x-2.25)^2+2})node[right]{$y=f(x)$};
    \tkzDefPoints{1/1.22/P, 3/1.72/Q, 3.5/1.22/P'}
    \tkzLabelPoints[above](P,Q)
    \tkzDrawSegments(P,P')
    \draw[domain=.2:2, samples=10, thick]plot(\x, {1.25*\x-0.03});
    \tkzDefPoints{3/3.72/Q'}
    \tkzDrawLines(P,Q)
    \tkzMarkAngles[mark=none, size=.7](P',P,Q)
    \tkzMarkAngles[mark=none, size=.3](P',P,Q')
    \tkzLabelAngle[pos=.5](P',P,Q'){$\alpha_1$}
    \tkzLabelAngle[pos=.9](P',P,Q){$\alpha$}
\end{tikzpicture}
    \caption{}\label{fig:PQ_tangent}
\end{figure}

割线的这种极限位置的存在性这一假设，与曲线在点$P$具有唯一的一条切线或确定的方向的假设是等价的．

现在我们要对曲线$y=f(x)$用解析式子把割线的这种极限位置存在的过程表示出来．
设$\alpha$是割线$PQ$同正$x$轴构成的夹角，$\alpha_1$是过点$P$点的切线同正$x$轴构成的夹角，于是
\[\lim_{Q\to P}\alpha=\alpha_1\]
设$x_1,y_1$和$x,y$分别是点$P$和$Q$的坐标，这时，我们立即得到
\[\tan\alpha=\frac{y-y_1}{x-x_1}=\frac{f(x)-f(x_1)}{x-x_1}\]
因此，上述求极限的过程（不考虑垂直切线$\alpha_1=\frac{\pi}{2}$的情况）可由下式来表示：
\[\lim_{x\to x_1}\frac{f(x)-f(x_1)}{x-x_1}=\lim_{\alpha\to \alpha_1}\tan\alpha=\tan\alpha_1\]
这就是说过曲线$y=f(x)$上$P(x_1,y_1)$点的切线的斜率等于$y=f(x)$的差商当$x\to x_1$时的极限．

\begin{example}
    求抛物线$y=ax^2+bx+c$在$x_0$处的切线的斜率．
\end{example}

\begin{solution}
    解依题意$(x_0,f(x_0)=ax_0^2+bx_0+c)$在抛物线上，并设$(x_0+h,f(x_0+h))$是抛物线上点$(x_0,f(x_0))$的附近的一点，我们有
\begin{align*}
&\qquad \lim_{h\to 0}\frac{f(x_0+h)-f(x_0)}{(x_0+h)-x_0}\\
&=\lim_{h\to 0}\frac{[a(x_0+h)^2+b(x_0+h)+c]-[ax^2_0+bx_0+c]}{h}    \\
&=\lim_{h\to 0}\frac{(2ax_0+b)h+ah^2}{h}\\
&=\lim_{h\to 0}[(2ax_0+b)+ah]=2ax_0+b
\end{align*}
所以抛物线$y=ax^2+bx+c$在$x_0$处的切线的斜率是$2ax_0+b$.
\end{solution}

\begin{example}
    一质点沿一直线在$t$秒内移动的距离是$s=s(t)=t^2+4t$．
    
    求：质点的初速度；在两秒末的速度；前两秒内的平均速度．
\end{example}

\begin{analyze}
    如果质点从起点开始所走的距离$s$是时间$t$的线性函数，则由2.1知道该质点在每一时刻的速度都是常数，它的大小由平均速度来确定，即等于一次函数的斜率，此时我们说该质点作匀速运动，但是，如果运动不再是匀速的，即质点的速度每时每刻都是变的，那么我们将时刻$t$的速度（也叫做瞬时速度）理解成什么呢？为了回答这个问题，我们考察差商
\[\frac{\Delta s}{\Delta t}=\frac{s(t)-s(t_0)}{t-t_0}\]
或者写成\[\frac{s(t_0+\Delta t)-s(t_0)}{\Delta t}\]
这个差商称为在$t_0$和$t_0+\Delta t$之间的这段时间间隔上的质点的平均速度，对照着$s=s(t)=t^2+4t$的图象来看，这个平均速度也就是过曲线上的$P(t_0,s(t_0))$
点及它邻近一点$Q(t_0+\Delta t,s(t_0+\Delta t))$的割线的斜率（图2.5）．当$\Delta t$ 很小时，可以认为，从时刻$t_0$到$t_0+\Delta t$这段时间内，速度来不及有很大变化，可以近似地看成匀速运动，因而这段时间内的平均速度就可以看成时刻$t_0$的瞬时速度的近似值．

\begin{figure}[htp]
    \centering
    \begin{tikzpicture}[>=latex, scale=.5]
\draw[->](-5.5,0)--(3,0)node[right]{$t$};
\draw[->](0,-5)--(0,7)node[right]{$s$};

\draw[domain=-5.25:0, samples=300, dashed]plot(\x, {\x*\x+4*\x});
\draw[domain=0:1.25, samples=100, thick]plot(\x, {\x*\x+4*\x});
\foreach \x in {-4,-3,-2,-1,1,2}
{
    \draw(\x,0)--(\x,.2);
}
\node at (-2,0)[below]{$-2$};
\node at (1,0)[below]{$1$};
\node at (0,0)[below left]{$O$};
    \end{tikzpicture}

    \caption{}
\end{figure}



显然，从时刻$t_0$到时刻$t_0+\Delta t$, 质点走过的路程为
\begin{align*}
    \Delta s&=s(t_0+\Delta t)-s(t_0)\\
    &=(t_0+\Delta t)^2+4(t_0+\Delta t)-(t^2_0+4t_0)\\
    &=(2t_0+4)\Delta t+(\Delta t)^2
\end{align*}
所以这段时间内的平均速度为
\begin{align*}
    \frac{\Delta s}{\Delta t}&=\frac{s(t_0+\Delta t)-s(t_0)}{\Delta t}\\
    &=\frac{(2t_0+4)\Delta t+(\Delta t)^2}{\Delta t}\\
    &=(2t_0+4)+\Delta t
\end{align*}
$\Delta t$越小，这个平均速度就越接近时刻$t_0$的瞬时速度$v_0$, 我们自然令$\Delta t\to 0$, 求差商的极限值，得到
\begin{align*}
    \lim_{\Delta t\to 0}\frac{s(t_0+\Delta t)-s(t_0)}{\Delta t}&= \lim_{\Delta t\to 0}[(2t_0+4)+\Delta t]\\
    &=2t_0+4
\end{align*}

这样平均速度$\frac{s(t_0+\Delta t)-s(t_0)}{\Delta t}$，当$\Delta t\to 0$时的极限值
就表达了质点在时刻$t_0$的瞬时速度，把它记作
\[s'(t_0)=   \lim_{\Delta t\to 0}\frac{s(t_0+\Delta t)-s(t_0)}{\Delta t}\]
\end{analyze}

\begin{solution}
\begin{enumerate}
\item 质点的初速度
\[s'(0)=2\x0+4=4(\ms)\]
\item 质点在两秒末的速度
\[s'(2)=2\x2+4=8(\ms)\]
\item \[\text{质点在前两秒内的平均速度}=\frac{\text{在前两秒内所走距离}}{\text{时间}}=\frac{2^2+4\x2}{2}
=6(\ms)\]
\end{enumerate}
\end{solution}

从这个问题可以看出质点在$t_0$时刻的瞬时速度$s'(t_0)$的几何意义就是曲线$s=s(t)$在$P(t_0,s(t_0))$点的切线的斜率，所以函数在某点的变率有确定值与函数的图象在该点有唯一的一条切线是两个密切相关的概念．

\subsection{微商（导数）的定义}
从上面所举的两个例子来看，问题来自不同的领域：
\begin{enumerate}
    \item 求过曲线上一点的切线，
    \item 求函数在某点的变率.
\end{enumerate}
但解决的方法却完全一样，就是计算函数的差商的极限，这种极限反映了自然界中很多不同现象在量方面的共性，因此有必要从这些具体问题中把它抽象出来加以研究，再反过来去解决这类具体问题．

\begin{blk}{定义}
    设$y=f(x)$是定义在闭区间$[a,b]$上的一个函数，$x_0\in (a,b)$, 如果极限
\[\lim_{\Delta x\to 0}\frac{f (x_0+\Delta x) -f (x_0)}{\Delta x}\]
存在，我们就说$f(x)$在点$x_0$处\textbf{可微}，并称这极限为函数$f(x)$在$x_0$点的\textbf{微商}（或导数），记为$f'(x_0)$或
$\frac{\dd y}{\dd x}\Big|_{x=x_0}$.
\end{blk}

显然，$f'(x_0)$的值与点$x_0$有关，当点$x_0$在开区间$(a,b)$内变化时，$f'(x_0)$也将跟着变化，因此，如果函数$f(x)$在开区间$(a,b)$内每点都可微（即存在有限导数），那么$f'(x)$便是一个新的函数，称为$f(x)$的\textbf{导函数}．

求已知函数的导函数$f'(x)$的运算，称为微商 运算，计算过程如下：
\begin{enumerate}
\item 设$\Delta x$为自变量某个值$x$的改变量：
\item 计算$f(x)$在点$x$的相应改变量
\[\Delta y=f (x+\Delta x) -f (x) \]
\item 计算$f(x)$在点$x$的差商
\[\frac{\Delta y}{\Delta x}=\frac{f (x+\Delta x) -f (x)}{\Delta x}\]
\item 计算
\[\lim_{\Delta x\to 0}\frac{f (x+\Delta x) -f (x)}{\Delta x}=f'(x)=\frac{\dd y}{\dd x}\]
\end{enumerate}

应当注意，这里$\frac{\dd y}{\dd x}$是一个独立的记号，它表示函数$f(x)$
在点$x$的导数，不能把它当成一个分数来看待，必须把它看成一个整体．

从微商的定义可以看出：
\begin{enumerate}
\item 曲线在一点的切线的斜率，就是函数在这一点的变率（微商或导数）．    
\item 微商所涉及的是函数的“局部”性质，也就是说，函数$y=f(x)$在一点$x_0$处是否可微只与函数$y=f(x)$在$x=x_0$处及其近旁的性质有关，而与其它地方无关．    
\item 如果$f(x)$在点$x$可微，按照极限存在的条件，必须且只须
\[\lim_{\Delta x\to 0^+}\frac{f (x+\Delta x) -f (x)}{\Delta x},\qquad \lim_{\Delta x\to 0^-}\frac{f (x+\Delta x) -f (x)}{\Delta x}\]
同时存在而且相等．

上面两式分别称为$f(x)$在点$x$的\textbf{右导数}和\textbf{左导数}，记为$f'_+(x)$和$f'_-(x)$.
\end{enumerate}

\begin{example}
今有一个正在膨胀的肥皂泡，
\begin{enumerate}
\item 求肥皂泡的体积对于半径的增大率，
\item 如果肥皂泡的半径每秒增大0.1cm，问当半径为2cm时，体积的增大率是多少？
\end{enumerate}
\end{example}

\begin{solution}
    肥皂泡的体积$V$与半径$r$的函数关系是．
\[ V(r) =\frac{4}{3}\pi r^3\]
体积$V(r)$对于半径$r$的增大率，依导函数定义，就是$V(r)$对于$r$的导函数，故
\begin{align*}
    V'(r)&=\lim_{\Delta r\to 0}\frac{V(r+\Delta r)-V(r)}{\Delta r}\\
    &=\lim_{\Delta r\to 0}\frac{\frac{4\pi }{3}[(r+\Delta r)^3-r^3]}{\Delta r}\\
    &=\frac{4\pi }{3}\lim_{\Delta r\to 0}[3r^2+3r(\Delta r)+(\Delta r)^2]\\
    &=4\pi r^2
\end{align*}

因此，肥皂泡的体积$V$对于半径$r$的增大率是$4\pi r^2$. 当$r=2$时，\[V'(2)=4\pi \cdot 2^2=16\pi \]
上式表示在$r=2$时，肥皂泡体积对于半径的增大率是$16\pi$, 它
的含义也可以这样理解，体积增大的快慢在$r=2$时，是$16\pi$倍于半径增大的快慢．

由于半径的增大率是每秒0.1cm，故体积在半径等于2时的增大率是
\[16\pi \x0.1=1.6\pi ({\rm cm^3/s})\]
\end{solution}

\begin{ex}
\begin{enumerate}
    \item 设$y=\frac{1}{x},\quad (x\ne 0)$，求
\begin{enumerate}
    \item 当$x$取改变量$\Delta x$后，函数改变量$\Delta y$的表达式；
    \item 当$x=3$, $\Delta x=-1$时，$\Delta y$的值；
    \item 当$x=3$时，$\frac{\Delta y}{\Delta x}$的表达式和$\frac{\dd y}{\dd x}\Big|_{x=3}$的值．
\end{enumerate}

    \item 设$\phi(x)=\frac{2}{x^2},\quad (x\ne0)$, 求
\begin{enumerate}
\item 当$x=1$, $\Delta x=-0.5$时，$\Delta \phi(x)$的值；
\item 对于任何$x\; (x\ne 0)$, 在$\Delta x$的间隔内，$\phi(x)$的平均变率$\frac{\Delta \phi}{\Delta x}$的表达式；
\item $\phi(x)$在任何点$x\; (x\ne 0)$处的瞬时变率$\phi'(x)$.
\end{enumerate}

    \item 平均变化率$\frac{\Delta y}{\Delta x}=\frac{f (x+\Delta x) -f (x)}{\Delta x}$依赖于哪两个变量？在平均变化率取极限求瞬时变化率的过程中，$x$是变量还是常量？$\Delta x$是变量还是常量？
    \item 点$P(2, 8)$在抛物线$y=x^2+2x$上，点$Q$为抛物线上任何一点．
\begin{enumerate}
    \item 求割线$PQ$的斜率的表达式；
  \item 当$Q$点横坐标为$2.1, 1.9, 2.002, 1.998, 2+h$时，求割线斜率的值；
\item 求在$P(2, 8)$点处，抛物线$y=x^2+2x$的切线方程和法线方程．
\end{enumerate}    

\item 求圆面积对于它的半径的变率，对于它的直径的变率．
\item 圆的半径的变率为2cm/s, 求圆面积在半径等于4cm
时的变率．
\end{enumerate}
\end{ex}

\subsection{函数的可微性与连续性的关系}

\begin{blk}
    {定理} 如果函数$f(x)$在点$x_0$可微，那么$f(x)$在点$x_0$处连续．
\end{blk}

\begin{proof}
设$f(x)$在点$x_0$处是可微的，也就是$f(x)$在$x_0$处有导数，则
\[f'(x_0)=\lim_{x\to x_0}\frac{f(x)-f(x_0)}{x-x_0}\]
所以
\begin{align*}
    \lim_{x\to x_0}[f(x)-f(x_0)]&=\lim_{x\to x_0}\frac{f(x)-f(x_0)}{x-x_0}(x-x_0)\\
    &=\lim_{x\to x_0}\frac{f(x)-f(x_0)}{x-x_0}\cdot \lim_{x\to x_0}(x-x_0)\\
&=f'(x_0)\cdot 0=0
\end{align*}
\end{proof}

上面的定理说明函数在其有导数的点一定连续，在其不
连续的点一定没有导数．但是它的逆命题不一定成立，即连续的函数不一定有导数．前面我们曾指出一个函数可微的充分必要条件是它的左导数和右导数都存在而且相等，下面给出在某些点不可微的函数的例子．

\begin{example}
    求函数$f(x)=|x^2-1|$的导函数的定义域．
\end{example}

\begin{figure}[htp]
    \centering
\begin{tikzpicture}[>=latex, scale=1]
    \draw[->](-2.5,0)--(2.5,0)node[right]{$x$};
    \draw[->](0,-.5)--(0,5)node[right]{$y$};
    \node at (0,0)[below left]{$O$};
    \draw[domain=-2.3:-1, samples=100, very thick]plot(\x, {\x*\x-1});
    \draw[domain=1:2.3, samples=100, very thick]plot(\x, {\x*\x-1});
    \draw[domain=-1:1, samples=100, very thick]plot(\x, {-\x*\x+1});
    \node at (0,0)[below left]{$O$};
    \node at (-1,0)[below left]{$-1$};
    \node at (1,0)[below right]{$1$};
    
    \draw[domain=-2:-.5, samples=10, thick]plot(\x, {-2*\x-2});
    \draw[domain=-1.5:.5, samples=10, thick]plot(\x, {2*\x+2});
    \draw[domain=-.5:1.5, samples=10, thick]plot(\x, {-2*\x+2});
    \draw[domain=.5:2, samples=10, thick]plot(\x, {2*\x-2});

\end{tikzpicture}
    \caption{}\label{fig:f(x)=|x^2-1|}
\end{figure}
    
\begin{solution}
函数$f(x)=|x^2-1|$是一个到处连续的函数，它的图象如图\ref{fig:f(x)=|x^2-1|}所示．

去掉函数式中的绝对值的符号，$f(x)$便可以写成分段函数式：
\[f(x)=\begin{cases}
    x^2-1,& x\le -1\quad \text{或}\quad x\ge 1\\
    -(x^2-1), & -1\le x\le 1
\end{cases}\]
所以
\[f'(x)=\begin{cases}
    2x ,&  x< -1\quad \text{或}\quad x> 1\\
    -2x,& -1<x<1
\end{cases}\]
但是当$x=-1$和1时，函数$f(x)$的左导数和右导数不相等，即此时导数不存在．因为当$x=-1$时，函数$f(x)$在$-1$点左邻域的表达式是
\[f (x) =x^2-1\]
而在$-1$点右邻域的表达式是
\[f (x) =- (x^2-1)\]
所以$f(x)$在$x=-1$点的左导数是
\begin{align*}
    f'_-(-1)&=\lim_{\Delta x\to 0^-}\frac{f(-1+\Delta x)-f(-1)}{\Delta x}\\
    &=\lim_{\Delta x\to 0^-}\frac{[(-1+\Delta x)^2-1]-0}{\Delta x}\\
    &=\lim_{\Delta x\to 0^-}(-2+\Delta x)=-2
\end{align*}
$f(x)$在$x=-1$点的右导数是
\begin{align*}
    f'_+(-1)&=\lim_{\Delta x\to 0^+}\frac{f(-1+\Delta x)-f(-1)}{\Delta x}\\
    &=\lim_{\Delta x\to 0^+}\frac{-[(-1+\Delta x)^2-1]-0}{\Delta x}\\
    &=\lim_{\Delta x\to 0^+}(2-\Delta x)=2 
\end{align*}
由于$f'_-(-1)\ne f'_+(-1)$, 所以我们说$f(x)$在$x=-1$处的导数不存在，同样可得
\[-2=f'_- (1)\ne f'_+(1)=2\]
故$f(x)$在$x=1$处也不可微．

因此，$f(x)=|x^2-1|$的导函数的定义域是$x\ne\pm 1$的点的集合，它的函数值如下：
\[f'(x)=\begin{cases}
    2x,& |x|>1\\
    -2x,& |x|<1
\end{cases}\]
\end{solution}

对照图\ref{fig:f(x)=|x^2-1|}来看，尽管函数$f(x)=|x^2-1|$处处连续，但是在$x=-1$和1这两点，曲线的特征是当动点由左侧趋于定
点1(或$-1$)时，割线的极限位置存在；当动点由右侧趋于定点1(或$-1$)时，割线的极限位置也存在，这两个极限位置分别叫1(或$-1$)点的\textbf{左、右切线}，并且它们之间的夹角不为\textbf{平角}，点$-1$和1分别是曲线的一种\textbf{角点}．
    
图\ref{fig:vertical_tangent}的情形是曲线$y=f(x)$在$P(x_0,y_0)$点连续并且在此点有平行于$y$轴的切线．现在我们来讨论图\ref{fig:vertical_tangent}中所示各函数在点$x_0$处的导数：

\begin{figure}[htp]
    \centering
\begin{tikzpicture}[>=latex]
\begin{scope}
\draw[->](-.5,0)--(4,0)node[right]{$x$};
\draw[->](0,-.5)--(0,4)node[right]{$y$};
\node at (2,-1){(a)};
\node at (0,0)[below left]{$O$};
\draw(2,0)node[below]{$x_0$}--(2,3.5);
\node at (2,2)[right]{$P(x_0,y_0)$};
\tkzDefPoints{1.5/.2/A, 2/2/B, 2.5/3.8/C}
\draw[very thick](A) to [bend right=18] (B) to [bend left=18](C);

\end{scope}
\begin{scope}[xshift=6cm]
    \draw[->](-.5,0)--(4,0)node[right]{$x$};
\draw[->](0,-.5)--(0,4)node[right]{$y$};
\node at (2,-1){(b)};
\node at (0,0)[below left]{$O$};
\draw(2,0)node[below]{$x_0$}--(2,3.5);
\node at (2,2)[right]{$P(x_0,y_0)$};
\tkzDefPoints{2.5/.2/A, 2/2/B, 1.5/3.8/C}
\draw[very thick](A) to [bend right=-18] (B) to [bend left=-18](C);

\end{scope}
\begin{scope}[yshift=-6cm]
    \draw[->](-.5,0)--(4,0)node[right]{$x$};
\draw[->](0,-.5)--(0,4)node[right]{$y$};
\node at (2,-1){(c)};
\node at (0,0)[below left]{$O$};
\draw(2,0)node[below]{$x_0$}--(2,3.5);
\node at (2,2)[right]{$P(x_0,y_0)$};
\tkzDefPoints{1.5/.2/A, 2/2/B, 2.5/.2/C}
\draw[very thick](A) to [bend right=18] (B) to [bend left=-18](C);


\end{scope}

\begin{scope}[xshift=6cm, yshift=-6cm]
    \draw[->](-.5,0)--(4,0)node[right]{$x$};
\draw[->](0,-.5)--(0,4)node[right]{$y$};
\node at (2,-1){(d)};
\node at (0,0)[below left]{$O$};
\draw(2,0)node[below]{$x_0$}--(2,3.5);
\node at (2,2)[right]{$P(x_0,y_0)$};
\tkzDefPoints{1.5/3.8/A, 2/2/B, 2.5/3.8/C}
\draw[very thick](A) to [bend right=-18] (B) to [bend left=18](C);



\end{scope}
\end{tikzpicture}
    \caption{}\label{fig:vertical_tangent}
\end{figure}

图(a)表示的函数在点$x_0$的邻近递增，因此过曲线上$P(x_0,y_0)$和$Q(x_0+\Delta x,y_0+\Delta y)$两点的割线，无论$Q$点在$P$点的哪一侧，$\Delta x$与$\Delta y$都有相同正、负号，于是
\[\frac{\Delta y}{\Delta x}=\frac{f(x_0+\Delta x)-f(x_0)}{\Delta x}>0\]
而且当$\Delta x\to 0$时，割线的倾斜角$\varphi$（$0^{\circ}\le \varphi\le 180^{\circ}$）便以$\pi/2$为极限，从而
\[\lim_{\Delta x\to 0}\frac{f(x_0+\Delta x)-f(x_0)}{\Delta x}=\lim_{\varphi\to \tfrac{\pi^-}{2}}\tan\varphi=+\infty\]
可见$f(x)$在点$x_0$不可微（不存在有限导数），但是我们常常把这种情况：当$\Delta x\to 0$时，$\frac{\Delta y}{\Delta x}\to\infty$简略地叙述为函
数$f(x)$在$x_0$的导数为正无穷大，记作$f'(x_0)=+\infty$．

图(b)表示的函数$f$在点$x_0$的邻近递减，因此过曲线上$P(x_0,y_0)$和$Q(x_0+\Delta x,y_0+\Delta y)$两点的割线，无论$Q$点在$P$点的哪一侧，$\Delta x$与$\Delta y$都有相反的正、负号，于是
\[\frac{\Delta y}{\Delta x}=\frac{f(x_0+\Delta x)-f(x_0)}{\Delta x}<0\]
而且当$\Delta x\to 0$时，$\varphi\to \frac{\pi^+}{2}$，从而
\[\lim_{\Delta x\to 0}\frac{f(x_0+\Delta x)-f(x_0)}{\Delta x}=\lim_{\varphi\to \tfrac{\pi^+}{2}}\tan\varphi=-\infty\]
因此$f(x)$在点$x_0$处不可微，但是我们常常把这种情况叙述为函数$f(x)$在$x_0$处的导数为负无穷大，记作$f'(x_0)=-\infty$．

图\ref{fig:vertical_tangent}(c)表示的函数在点$x_0$处连续，让$x$由$x_0$变动到
$x_0+\Delta x$，于是由$\Delta x$所引起的相应的函数的改变量$\Delta y$的情形是：
\begin{itemize}
    \item 当$\Delta x<0$时，有$\Delta y<0$，从而$\frac{\Delta y}{\Delta x}>0$；
    \item 当$\Delta x>0$时，有$\Delta y<  0$，从而$\frac{\Delta y}{\Delta x}<0$．
\end{itemize}
于是：
\begin{itemize}
    \item 当$\Delta x\to 0^-$时，$\varphi\to \frac{\pi^-}{2}$，从而$\frac{\Delta y}{\Delta x}=\tan\varphi\to +\infty$；
    \item 当$\Delta x\to 0^+$时，$\varphi\to \frac{\pi^+}{2}$，从而$\frac{\Delta y}{\Delta x}=\tan\varphi\to -\infty$．
\end{itemize}
因此$f(x)$在点$x_0$不可微，但是我们常把这种情况叙述为$f(x)$在点$x_0$的左导数$f'_-(x_0)=+\infty$, 而它的右导数为$f'_+(x_0)=-\infty$, 曲线在$x_0$处由上升转为下降有一个尖点，它的切线垂直于$x$轴．

图(d)表示的函数在点$x_0$处连续，但是
\begin{itemize}
    \item 当$\Delta x\to 0^-$时，$\frac{\Delta y}{\Delta x}=\tan\varphi\to -\infty$；
    \item 当$\Delta x\to 0^+$时，$\frac{\Delta y}{\Delta x}=\tan\varphi\to +\infty$．
\end{itemize}
因此，$f(x)$在点$x_0$不可微，我们常常把这种情况叙述为$f'_-(x_0)=-\infty$和$f'_+(x_0)=+\infty$, 点$x_0$是曲线的一个尖点，它的切线垂直于$x$轴．

\begin{example}
设$f(x)=\begin{cases}
    0,& x=0\\
    x\sin\frac{1}{x},& x\ne 0
\end{cases}$

讨论它在0点的导数．
\end{example}

\begin{solution}
在第一章，我们已经说明了
\[\lim_{x\to 0}x\sin\frac{1}{x}=0=f(0)\]
因此，$f(x)$在点$x=0$连续，很明显$f(x)$在其它各点处也都连续，所以说$f(x)$是一个到处连续的函数，我们将说明它在原点不存在导数．

因为
\[\frac{f(0+\Delta x)-f(0)}{\Delta x}=\frac{\Delta x\sin\frac{1}{\Delta x}}{\Delta x}=\sin\frac{1}{\Delta x}\]
而当$\Delta x\to 0$时，$\sin\frac{1}{\Delta x}$
没有极限，故$f'(0)$不存在，就
其几何意义来说，当动点沿着曲线$y=x\sin\frac{1}{x}$趋于原点$O$
时，割线$OQ$不断地在$-\frac{\pi}{4}\le \theta\le \frac{\pi}{4}$这个幅度之内摆动，而不趋于任何极限位置，即切线不存在（图\ref{fig:x*sin(1/x)}）．
\end{solution}

\begin{figure}[htp]
    \centering
\begin{tikzpicture}[>=latex, scale=3]
\draw[->](-1.8,0)--(1.8,0)node[right]{$x$};
\draw[->](0,-.5)--(0,1.5)node[right]{$y$};
\draw[domain=.02:1.8, very thick, samples=300]plot(\x, {\x*sin(1/\x r)});
\draw[domain=.02:1.8, very thick, samples=300]plot(-\x, {\x*sin(1/\x r)});
\draw[thick](-.4,-.4)--(1.3,1.3);
\draw[thick](.4,-.4)--(-1.3,1.3);
\draw[thick](15:1.5)--(-165:1);
\draw[dashed](-1.8,1)--(1.8,1);

\foreach \x in {1,2,3,4}
{
   \draw(\x/pi,0)node[below]{$\frac{\x}{\pi}$}--(\x/pi,.04);
}

\node at (0,0)[above]{$O$};
\end{tikzpicture}
    \caption{}\label{fig:x*sin(1/x)}
\end{figure}

通过以上的例子，我们知道函数的连续点未必是它的可以微分点．几种常见的不可微分点如例2.4的角点，平行于$y$轴的切线的切点，特别是曲线上的尖点和在例2.5的曲线上的不存在割线的极限位置的点．

假如区间$[a,b]$上的函数$f(x)$, 在开区间$(a,b)$中存在着导数$f'(x)$, 并且$f'_+(a)$和$f'_-(b)$都存在，则称$f(x)$在闭区间上可以微分（可导）．

\begin{ex}
\begin{enumerate}
    \item 说明下面函数在点$x=0$不可微：
\begin{multicols}{2}
\begin{enumerate}
    \item $f(x)=|x|$
    \item $f(x)=\begin{cases}
        x^2,& x\le 0\\
        x,& x\ge 0
    \end{cases}$
    \item $f(x)=\sqrt[3]{x}$
    \item $f(x)=\sqrt{|x|}$
\end{enumerate}
\end{multicols}
\item 说明函数$f(x)=\begin{cases}
    x^2,& x\ge 0\\
    -x^2,& x<0
\end{cases}$
    处处可微．
\item \begin{enumerate}
\item 假设$g(x)=f(x+c)$. 从定义出发证明$g'(x)=f'(x+c)$, 并绘图说明．    
\item 设$f$可微并有周期$\varphi$, 证明$f'$也是有周期的．
\end{enumerate}

\end{enumerate}
\end{ex}

\section{微商运算的基本法则}
\subsection{几个基本函数的微商（导数）}
\subsubsection{常数函数的导数恒等于0}

\begin{proof}
设$f(x)=c$ （$c$是常数），则根据常数函数的性质，由自变数$x$的改变量$\Delta x$引起的相应的函数的改变量为：
\[\Delta y=f (x+\Delta x) -f (x)=c-c=0\]
显然差商$\frac{\Delta y}{\Delta x}=\frac{0}{\Delta x}=0$, 故
\[f' (x) =\lim_{\Delta x\to 0}\frac{\Delta y}{\Delta x}=
\lim_{\Delta x\to 0} 0=0\]
这个事实是明显的，因为当自变数$x$有增、减时，函数值并无增、减，因此，它的变化率为零．
\end{proof}

\subsubsection{一次函数（直线函数）的导数是常数}

若$f(x)=kx+b$, 则$f'(x)=k$. 

\subsubsection{若$f(x)=x^n$ （$n$是正整数），则$f'(x)=nx^{n-1}$}

\begin{proof}
  由自变量的改变量$\Delta x$引起的函数改变量为
\[\Delta y=f (x+\Delta x) -f (x) = (x+\Delta x)^n-x^n\] 
按牛顿二项式定理展开，得  
\begin{align*}
\Delta y&= x^n+nx^{n-1}\Delta x+\frac{n(n-1)}{1\cdot 2}x^{n-2}(\Delta x)^2+\cdots+nx(\Delta x)^{n-1}+(\Delta x)^n-x^n\\
&=nx^{n-1}\Delta x+\frac{n(n-1)}{1\cdot 2}x^{n-2}(\Delta x)^2+\cdots+nx(\Delta x)^{n-1}+(\Delta x)^n
\end{align*}
其中差商
\[\frac{\Delta y}{\Delta x}=nx^{n-1}+\frac{n(n-1)}{1\cdot 2}x^{n-2}\Delta x+\cdots+nx(\Delta x)^{n-2}+(\Delta x)^{n-1}\]
式中除第一项外，其余各项随$\Delta x\to 0$而趋于0，因此：
\[f'(x)=\lim_{\Delta x\to 0}\frac{\Delta y}{\Delta x}=nx^{n-1}\]
\end{proof}

\subsubsection{若$f(x)=x^{\mu}$ (其中$\mu$是任意实数，函数的定义域依赖于$\mu$), 则$f'(x)=\mu x^{\mu-1}$}

\begin{proof}
    \[\frac{\Delta y}{\Delta x}=\frac{(x+\Delta x)^{\mu}-x^{\mu}}{\Delta x}=x^{\mu-1}\cdot \frac{\left(1+\frac{\Delta x}{x}\right)^{\mu}-1}{\frac{\Delta x}{x}}\]
利用第一章中已经算出的极限
\[\lim_{x\to 0}\frac{(1+x)^{\mu}-1}{x}=\mu\]
就得到
\[f'(x)=\lim_{\Delta x\to 0}\frac{\Delta y}{\Delta x}=\mu x^{\mu-1}\]
\end{proof}

特殊情形：
\begin{enumerate}
    \item 若$f(x)=\frac{1}{x}=x^{-1},\; (x\ne 0)$，则$f'(x)=(-1)x^{-2}=-\frac{1}{x^2}$
    \item 若$f(x)=\sqrt{x}=x^{\tfrac{1}{2}},\; (x\ge 0)$，则$f'(x)=\frac{1}{2}x^{-\tfrac{1}{2}}=\frac{1}{2\sqrt{x}}$
\end{enumerate}

\subsubsection{若$f(x)=\log_a x\; (0<a\ne 1,\; 0<x<+\infty)$，则$f'(x)=\frac{\log_a e}{x}$}

\begin{proof}
\begin{align*}
    \frac{\Delta y}{\Delta x}&=\frac{\log_a(x+\Delta x)-\log_a x}{\Delta x}\\
    &=\frac{1}{x}\cdot \frac{\log_a\left(1+\frac{\Delta x}{x}\right)}{\frac{\Delta x}{x}}\\
    &=\frac{1}{x}\cdot \log_a\left(1+\frac{\Delta x}{x}\right)^{\tfrac{1}{\Delta x/x}}
\end{align*}  

利用极限$\Lim_{t\to 0}(1+t)^{1/t}=e$，以及对数函数是连续函数，得到
\begin{align*}
    f'(x)&=\lim_{\Delta x\to 0}\frac{\Delta y}{\Delta x}=\frac{1}{x}\lim_{\tfrac{\Delta x}{x}\to 0}\log_a\left(1+\frac{\Delta x}{x}\right)^{\tfrac{1}{\Delta x/x}}\\
    &=\frac{1}{x}\log_a\left[\lim_{\tfrac{\Delta x}{x}\to 0}\left(1+\frac{\Delta x}{x}\right)^{\tfrac{1}{\Delta x/x}}\right]\\
    &=\frac{1}{x}\log_a e=\frac{\log_a e}{x}
\end{align*}

特别：当$f(x)=\ln x$时，$f'(x)=\frac{1}{x}$．
\end{proof}

上面的结果表明对数函数（$a>1$时）的增大速度是与自
变数的值成反比的，当自变数无限增大时，增大速度就保持正值而趋向于0. 由于自然对数的导数比较简单，所以在理论研究中常常采用自然对数．

\subsubsection{若$f(x)=a^x\; (0<a\ne 1,\; -\infty<x<+\infty)$，则$f'(x)=a^x\cdot \ln a$}

\begin{proof}
    \[\frac{\Delta y}{\Delta x}=\frac{a^{x+\Delta x}-a^x}{\Delta x}=a^x\cdot \frac{a^{\Delta x}-1}{\Delta x}\]
利用$\Lim_{\Delta x\to 0}a^{\Delta x}=1$，则可设$a^{\Delta x}=1+\alpha$，且当$\Delta x\to 0$时，有$\alpha\to 0$. 在此等式的两边取对数，得：
\[\Delta x=\log_a(1+\alpha)\]
于是
\[\frac{\Delta y}{\Delta x}=a^x\cdot \frac{\alpha}{\log_a(1+\alpha)}=a^x\cdot \frac{1}{\log_a(1+\alpha)^{1/\alpha}}\]
因此
\begin{align*}
    f'(x)&=\lim_{\Delta x\to 0}\frac{\Delta y}{\Delta x}=a^x\cdot \frac{1}{\Lim_{\alpha\to 0}\log_a(1+\alpha)^{1/\alpha}}\\
    &=a^x\cdot \frac{1}{\log_a\left[\Lim_{\alpha\to 0}(1+\alpha)^{1/\alpha}\right]}\\
    &=a^x\cdot \frac{1}{\log_a e}
\end{align*}
又因为$\ln a=\frac{1}{\log_a e}$，所以
\[f'(x)=a^x\cdot \ln a\]

特别，若$f(x)=e^x$，则$f'(x)=e^x$．
\end{proof}

这个事实表明指数函数（当$a>1$时）的增大速度与函数值成正比例，当底数为$e$时，导数的结果特别简单．


\subsubsection{若$f(x)=\sin x$, 则$f'(x)=\cos x$}
\begin{proof}
    \begin{align*}
\frac{\Delta y}{\Delta x}&=\frac{\sin(x+\Delta x)-\sin x}{\Delta x}\\
&=\frac{2\cos\left(x+\frac{\Delta x}{2}\right)\sin\frac{\Delta x}{2}}{\Delta x}\\
&=\cos\left(x+\frac{\Delta x}{2}\right)\cdot \frac{\sin\frac{\Delta x}{2}}{\frac{\Delta x}{2}}
    \end{align*}
因此
\begin{align*}
    \lim_{\Delta x\to 0}\frac{\Delta y}{\Delta x}&=\lim_{\Delta x\to 0}\left[\cos\left(x+\frac{\Delta x}{2}\right)\cdot \frac{\sin\frac{\Delta x}{2}}{\frac{\Delta x}{2}}\right]\\
    &=\lim_{\Delta x\to 0}\cos\left(x+\frac{\Delta x}{2}\right)\cdot \lim_{\Delta x/2\to 0}\frac{\sin\frac{\Delta x}{2}}{\frac{\Delta x}{2}}
\end{align*}
根据$\cos x$的连续性与极限$\Lim_{t\to 0}\frac{\sin t}{t}=1$, 得到
$f' (x) =\cos x$.
\end{proof}

\subsubsection{若$f(x)=\cos x$, 则$f'(x)=-\sin x$}
\begin{proof}
\begin{align*}
\frac{\Delta y}{\Delta x}&=\frac{\cos(x+\Delta x)-\cos x}{\Delta x}\\
&=\frac{-2\sin\left(x+\frac{\Delta x}{2}\right)\cdot \sin\frac{\Delta x}{2}}{\Delta x}\\
&=-\sin\left(x+\frac{\Delta x}{2}\right)\cdot \frac{\sin\frac{\Delta x}{2}}{\frac{\Delta x}{2}}
    \end{align*}
因此
\begin{align*}
    \lim_{\Delta x\to 0}\frac{\Delta y}{\Delta x}&=-\lim_{\Delta x\to 0}\left[\sin\left(x+\frac{\Delta x}{2}\right)\cdot \frac{\sin\frac{\Delta x}{2}}{\frac{\Delta x}{2}}\right]\\
    &=-\lim_{\Delta x\to 0}\sin\left(x+\frac{\Delta x}{2}\right)\cdot \lim_{\Delta x/2\to 0}\frac{\sin\frac{\Delta x}{2}}{\frac{\Delta x}{2}}\\
    &=-\sin  x
\end{align*}
\end{proof}

现在将上面的导数公式列成下表：
\begin{enumerate}
    \item $\frac{\dd c}{\dd x}=0$
    \item $\frac{\dd(kx+b)}{\dd x}=k$
    \item $\frac{\dd x^n}{\dd x}=nx^{n-1}$\quad （$n$为自然数）
    \item $\frac{\dd x^\mu}{\dd x}=\mu x^{\mu-1}$\quad （$\mu$为任意实数）
    \item $\frac{\dd\log_a x}{\dd x}=\frac{\log_a e}{x},\qquad \frac{\dd \ln x}{\dd x}=\frac{1}{x}$
    \item $\frac{\dd a^x}{\dd x}=a^x\cdot \ln a,\qquad \frac{\dd e^x}{\dd x}=e^x$
    \item $\frac{\dd\sin x}{\dd x}=\cos x,\qquad \frac{\dd \cos x}{\dd x}=-\sin x$
\end{enumerate}

\subsection{求导的基本法则}

现在要建立一些求导数的公式，用这些公式求函数的导数．

\begin{blk}
{定理1} 函数和的导数等于各函数的导数的和，即：
\[[f (x) +g (x) ] '=f' (x) +g' (x)\] 
\end{blk}

\begin{proof}
    设$u(x)=f(x)+g(x)$, 任意固定$x$, 任取非零数$\Delta x$，有 
\begin{align*}
    \frac{\Delta u(x)}{\Delta x}&=\frac{[f(x+\Delta x)+g(x+\Delta x)]-[f(x)+g(x)]}{\Delta x}\\
    &=\frac{f(x+\Delta x)-f(x)}{\Delta x}+\frac{g(x+\Delta x)-g(x)}{\Delta x}\\
    &=\frac{\Delta f(x)}{\Delta x}+\frac{\Delta g(x)}{\Delta x}
\end{align*}
令$\Delta x\to 0$，即得：
\begin{align*}
    [f(x)+g(x)]'&=\lim_{\Delta x\to 0}\frac{\Delta u(x)}{\Delta x}\\
    &=\lim_{\Delta x\to 0}\frac{\Delta f(x)}{\Delta x}+\lim_{\Delta x\to 0}\frac{\Delta g(x)}{\Delta x}\\
    &=f'(x)+g'(x)
\end{align*}
\end{proof}

这个法则可以推广到任意有限个函数．

\begin{blk}
{定理2} $[A\varphi (x)]'=A\varphi'(x)$, 其中$A$是常数，即常因数不因微分法而改变．    
\end{blk}

\begin{proof}
设$f(x)=A\varphi(x)$， 
\begin{align*}
    [A\varphi(x)]'&=f'(x)=\lim_{\Delta x\to 0}\left[A\frac{\varphi(x+\Delta x)-\varphi(x)}{\Delta x}\right]\\
    &=A\lim_{\Delta x\to 0}\frac{\varphi(x+\Delta x)-\varphi(x)}{\Delta x}\\
    &=A\varphi'(x)
\end{align*}
\end{proof}

\begin{example}
    求$3x^2-7x-1$的导数．
\end{example}

\begin{solution}
\[  (3x^2-7x-1)'=(3x^2)'+(-7x)'+(-1)'=6x-7\]
\end{solution}    

\begin{blk}{推论}
设$f(x)=a_nx^n+a_{n-1}x^{n-1}+\cdots+a_1x+a_0\; (a_n\ne 0)$，则
\[f'(x)=na_nx^{n-1}+(n-1)a_{n-1}x^{n-2}+\cdots+a_1\]
\end{blk}

\begin{example}
求过抛物线$y=3x^2+6$外一点$P(2,-9)$所引抛物线的两条切线的方程．
\end{example}

\begin{solution}
在点$x$处的抛物线的切线的斜率为
\[k=y'_x=6x\]
故过抛物线上一点$T(x_1,y_1)$的切线方程为
\[y-y_1=6x_1 (x-x_1)\] 

$\because\quad y_1=3x^2_1+6$

$\therefore\quad $切线方程可化简为
\[    y=6x_1x-3x_1^2+6\]
设切线通过$P(2,-9)$点，于是
\[-9=12x_1-3x_1^2+6\]
即：$x_1^2-4x_1-5=0$

$\therefore\quad x_1=-1$或$5$，从而$y_1=9$或81.

故所求切线方程为：
\[y-9=-6(x+1)\quad\text{和}\quad y-81=30(x-5)\]
即：
\[y+6x-3=0\quad\text{和}\quad  y-30x+69=0\]
\end{solution}

\begin{example}
求三次函数$f(x)=ax^3+3bx^2+3cx+d$的图象与$x$轴相切的条件．
\end{example}

    
\begin{solution}
由函数$f(x)=ax^3+3bx^2+3cx+d$求得它的图象任意一点$x=x_0$处的切线斜率为
\[f' (x_0) =3ax^2_0+6bx_0+3c=3 (ax^2_0+2bx_0+c) \]
    
要此曲线与$x$轴相切于$(x_0,0)$点，$x_0$必须满足条件
\begin{numcases}{}
    ax^3_0+3bx^2_0+3cx_0+d=0\\
    ax^2_0+2bx_0+c=0
\end{numcases}
将(2.1)改写成
\begin{equation}
    x_0(ax^2_0+2bx_0+c)+bx_0^2+2cx_0+d=0
\end{equation}
(2.2)代入得
\begin{equation}
    bx_0^2+2cx_0+d=0
\end{equation}
于是解(2.1)和(2.2)即解(2.2)和(2.4)，要方程(2.2)和(2.4)有公共解，必须
\[D=\begin{vmatrix}
    a&2b\\b&2c
\end{vmatrix}=2(ac-b^2)\ne 0\]
于是
\[x^2_0=\frac{c^2-bd}{b^2-ac},\qquad x_0=\frac{ad-bc}{2(b^2-ac)}\; (b^2-ac\ne 0)\]
所以，$a,b,c,d$必须适合条件
\[\begin{cases}
    b^2-ac\ne 0\\
    \frac{c^2-bd}{b^2-ac}=\frac{(ad-bc)^2}{4(b^2-ac)^2}
\end{cases}\]
即
\[\begin{cases}
    b^2-ac\ne 0\\
   (ad-bc)^2 =4(b^2-ac)(c^2-bd)
\end{cases}\]
\end{solution}

\begin{blk}
    {定理3} 二函数之积的导数为两项之和，其中第一项是第一个因子的导数与第二个因子的乘积，而第二项是第二个因子的导数与第一个因子的乘积，即
\[[f (x) \cdot g (x) ] '=f' (x) \cdot g (x) +f (x) \cdot g' (x) \]
\end{blk}

\begin{proof}
    设 $h(x)=f(x)\cdot g(x)$, 任取非零数$\Delta x$, 则
\begin{align*}
    [f(x)\cdot g(x)]'&=\lim_{\Delta x\to 0}\frac{\Delta h(x)}{\Delta x}\\
    &=\lim_{\Delta x\to 0}\frac{f(x+\Delta x)g(x+\Delta x)-f(x)g(x)}{\Delta x}\\
    &=\lim_{\Delta x\to 0}\frac{1}{\Delta x}\Bigl[f(x+\Delta x)g(x+\Delta x)-f(x)g(x+\Delta x)\\
    &\qquad\qquad +f(x)g(x+\Delta x)-f(x)g(x) \Bigr]\\
    &=\lim_{\Delta x\to 0}\frac{f(x+\Delta x)-f(x)}{\Delta x}\cdot g(x+\Delta x)\\
    &\qquad +\lim_{\Delta x\to 0}f(x)\cdot \frac{g(x+\Delta x)-g(x)}{\Delta x}\\
    &=f' (x) \cdot g (x) +f (x) \cdot g' (x) 
\end{align*}
\end{proof}

重复应用这个定理，我们便可以求两个以上的函数之积的导数，例如
\[\frac{\dd }{\dd x}(uvw)=(vw)\frac{\dd u}{\dd x}+u\frac{\dd }{\dd x}(vw)\]
但是
\[\frac{\dd }{\dd x}(vw)=w\frac{\dd v}{\dd x}+v\frac{\dd w}{\dd x}\]
代入上式便得到
\[\frac{\dd }{\dd x}(uvw)=vw\frac{\dd u}{\dd x}+uw\frac{\dd v}{\dd x}+uv\frac{\dd w}{\dd x}\]


\begin{example}
    设$f(x)=\sqrt{x}(x^2-3)$，求$f'(x)$
\end{example}


\begin{solution}
    \begin{align*}
f'(x)&=\left[\sqrt{x}(x^2-3)\right]'\\
&=\left(\sqrt{x}\right)'(x^2-3)+\sqrt{x}(x^2-3)'\\
&=\frac{1}{2\sqrt{x}}(x^2-3)+\sqrt{x}\cdot 2x = \frac{5x^2-3}{2\sqrt{x}}
    \end{align*}
\end{solution}

\begin{example}
    设$f(x)=\frac{1}{2}\sin^2 x+x\sin x+\frac{7}{x^2}$，求$f'(x)$
\end{example}

\begin{solution}
\begin{align*}
    f'(x)&=\left(\frac{1}{2}\sin^2 x\right)'+(x\sin x)'+\left(\frac{7}{x^2}\right)'\\
    &=\frac{1}{2}\left[(\sin x)'\sin x+\sin x(\sin x)'\right]+(\sin x+x\cos x)-\frac{14}{x^3}\\
    &=\sin x\cos x+\sin x+x\cos x-\frac{14}{x^3}
\end{align*}
\end{solution}

\begin{blk}
    {定理4} 两个函数的商的导数，等于分子的导数与分母的积，减去分母的导数与分子的积，再除以分母的平方，即
\[\left(\frac{u(x)}{v(x)}\right)'=\frac{v(x)\cdot u'(x)-u(x)\cdot v'(x)}{[v(x)]^2}\]
\end{blk}

\begin{proof}
    设$f(x)=\frac{u(x)}{v(x)}$，任取非零数$\Delta x$，则
\begin{align*}
    \left(\frac{u(x)}{v(x)}\right)'&=\lim_{\Delta x\to 0}\frac{\Delta f(x)}{\Delta x}=\lim_{\Delta x\to 0}\frac{\frac{u(x+\Delta x)}{v(x+\Delta x)}-\frac{u(x)}{v(x)}}{\Delta x}\\
    &=\lim_{\Delta x\to 0}\frac{1}{\Delta x}\left\{\frac{u(x+\Delta x)\cdot v(x)-u(x)\cdot v(x+\Delta x)}{v(x+\Delta x)\cdot v(x)}\right\}\\
    &=\lim_{\Delta x\to 0}\frac{1}{\Delta x}\left\{\frac{v(x)[u(x+\Delta x)-u(x)]}{v(x)v(x+\Delta x)}-\frac{u(x)[v(x+\Delta x)-v(x)]}{v(x)v(x+\Delta x)}\right\}\\
    &=\frac{v(x)\Lim_{\Delta x\to 0}\frac{u(x+\Delta x)-u(x)}{\Delta x}}{v(x)\Lim_{\Delta x\to 0}v(x+\Delta x)}-\frac{u(x)\Lim_{\Delta x\to 0}\frac{v(x+\Delta x)-v(x)}{\Delta x}}{v(x)\Lim_{\Delta x\to 0}v(x+\Delta x)}\\
    &=\frac{v(x) u'(x)-u(x) v'(x)}{v^2(x)}
\end{align*}
\end{proof}


\begin{example}
求证：$(\tan x)'=\sec^2 x,\qquad (\cot x)'=-csc^2 x$
\end{example}

\begin{proof}    
\begin{align*}
    (\tan x)'&=\left(\frac{\sin x}{\cos x}\right)'\\
    &=\frac{\cos x(\sin x)'-\sin x(\cos x)'}{\cos^2 x}\\
    &=\frac{\cos^2 x+\sin^2 x}{\cos^2x}\\
    &=\frac{1}{\cos^2 x}=\sec^2 x
\end{align*}
\[(\cot x)'=\left(\frac{\cos x}{\sin x}\right)'=\frac{-\sin^2 x-\cos^2 x}{\sin^2 x}=-\csc^2 x\]    
\end{proof}

\begin{example}
    求$\left(\frac{2x-1}{x^2+1}\right)'$
\end{example}

\begin{solution}
\begin{align*}
    \left(\frac{2x-1}{x^2+1}\right)'&=\frac{(x^2+1)(2x-1)'-(2x-1)(x^2+1)'}{(x^2+1)^2}\\
&=\frac{(x^2+1)\cdot 2-(2x-1)(2x)}{(x^2+1)^2}\\
&=\frac{2(x^2+1)-(4x^2-2x)}{(x^2+1)^2}\\
&=\frac{-2(x^2-x-1)}{(x^2+1)^2}
\end{align*}
\end{solution}

\begin{ex}
\begin{enumerate}
    \item  求下面各函数的导函数：
\begin{multicols}{2}
\begin{enumerate}
    \item $3 x^{3}-1$
    \item  $x+2 x^{2}+3 x^{3}$
    \item  $1-\frac{1}{x^{2}}$
    \item $\frac{x^{2}+1}{x}$
    \item  $1+\frac{1}{x}-\frac{1}{x^{2}}$
    \item  $\frac{2 x^{5}-7 x^{3}-3 x^{2}}{6 x^{2}}$
    \item  $(x+2)\left(x^{2}+1\right)$
    \item  $x+\sqrt{x}$
    \item $x-\frac{1}{\sqrt{x}}$
    \item  $\frac{x+1}{\sqrt{x}}$
    \item  $\frac{2 x \sqrt{x}-x^{\tfrac{5}{2}}+3 x^{\tfrac{1}{2}}}{-\sqrt{x}}$ \item  $e^{x}-e^{-x}$
    \item $\frac{x^{n}-x^{-n}}{x} $
    \item $\sin x+\cos x$
    \item $\sin x-\cos x$
\end{enumerate}
\end{multicols}
\item  求下面各函数的导函数：
\begin{multicols}{2}
\begin{enumerate}
    \item $\left(x^{2}+1\right)^{3} \sqrt{x}$
    \item $\frac{x+1}{\sqrt{x}}$
    \item  $x \sin x$
    \item $x \cot  x$
    \item  $\frac{\tan  x}{x}$;
    \item  $x \log _{2} x+x^{2} \tan  x$
    \item  $\frac{x^{2}}{x^{3}+c^{3}}$
    \item  $x^{3} e^{-x}$;
    \item  $a^{-x} \sin x$ 
    \item $a^{x}\{(x+1) \cos x+\log_2 x\}$
    \item $\frac{\sqrt{x}+1}{2 x+1}$ 
    \item $\frac{2 x-1}{x^{2}+1}$
    \item $\frac{(x-3)(x-4)}{(x-5)}$
    \item  $\frac{x^{2}\left(x^{2}-1\right)\left(x^{3}-1\right)}{x+1}$
    \item  $\frac{x}{x+1}-\frac{1}{x-1}$
    \item  $x \sin x \cos x$
    \item $\sec x $ 
    \item $\csc x$
    \item $\frac{\sin x}{1+\tan  x}$
\end{enumerate}
\end{multicols}

\item 已知抛物线$y=x^2-x$上一点的切线平行于直线$y=x$．求此点的坐标．
\item 求曲线$y=x^3-x^2$上这样的点，使过该点的切线与$x$轴平行．
\item 若抛物线$y=ax^2+bx+c$通过原点，且过该点的切线的斜率等于2, 又抛物线过$(1, 1)$点，求$a$、$b$、$c$.
\item 对于函数$f(x)=x^3+ax^2+bx+c$, 方程$f(x)+\frac{1}{2}x=0$
有不相等的三根$2,\alpha,\beta$, 若$f'(\alpha)=f'(\beta)$, $f'(1)=0$, 求$a,b,c$.
\item 求三次函数$y=f(x)$, 使它同时满足下面的条件：
\begin{enumerate}
    \item 用$x+1$去除$f(x)$与用$(x-1)(x-2)$去除，$f(x)$所得余式相同；
    \item 过曲线$y=f(x)$上的点$(1,f(1))$的切线方程是
$y=-2x+1$
\end{enumerate}
\item 
过$(2, 0)$点求与曲线$y=\frac{1}{x}$相切的直线方程．
\item 
\begin{enumerate}
\item 在抛物线$y=x^2$上求这样一点，使过该点的切线与过$(a,a^2)$点的切线垂直；
\item 求抛物线两条正交切线交点的轨迹方程．
\end{enumerate}

\item 设抛物线方程为$y=x^2+ax+b$, 试问点$(x_0,y_0)$位于
何处时，可以从点$(x_0,y_0)$对此抛物线作出两条切线或一条切线，或作不出切线？
\item 求抛物线$y=x^2+ax$, $y=x^2+bx\; (a\ne b)$的公切线的方
程．
\end{enumerate} 
\end{ex}

\subsection{复合函数的求导法则}
\begin{blk}
    {定理} 假设函数$y=g(x)$在点$x$可导，而函数$z=f(y)$在点$g(x)$可导，那么复合函数$z=\varphi(x)=f(g(x))$在点$x$可导，并且
\[[f(g(x))]'=f'(g(x))\cdot g'(x)\]
简写成$[f(g(x))]'=f'_g\cdot g'_x$，或者写成
\[\frac{\dd z}{\dd x}=\frac{\dd z}{\dd y}\cdot \frac{\dd y}{\dd x}\]
\end{blk}

\begin{proof}
    设$D_g$, $D_f$分别是函数$g(x),f(y)$的定义域，令$x\in D_g$, $g(x)\in D_f$, 任取非零数$\Delta x$, 且使$g(x+\Delta x)\in D_f$，于是$\Delta y=g(x+\Delta x)-g(x)$, 又根据$g(x)$在点$x$连续，故当$\Delta x\to 0$时，$\Delta y\to 0$. 因为$g(x)$在点$x$可导，可以设
\begin{equation}
    \alpha=\frac{\Delta y}{\Delta x}-g'(x)
\end{equation}
而且由上式知道
\[\lim_{\Delta x\to 0}\alpha=\lim_{\Delta x\to 0}\left(\frac{\Delta y}{\Delta x}-g'(x)\right)=g'(x)-g'(x)=0\]
把(2.5)改写成
\begin{equation}
    \Delta y=g'(x)\Delta x+\alpha\Delta x
\end{equation}
其中$\alpha$随$\Delta x$一同趋向于零，且$\alpha \Delta x$是比$\Delta x$更小的量．

\begin{enumerate}
    \item 若$g'(x)\ne 0$，只要$\Delta x$够小，由(2.6)知$\Delta y\ne 0$，那么
\begin{align*}
    \frac{\varphi(x+\Delta x)-\varphi(x)}{\Delta x}&=\frac{f(g(x+\Delta x))-f(g(x))}{g(x+\Delta x)-g(x)}\cdot \frac{g(x+\Delta x)-g(x)}{\Delta x}\\
    &=\frac{f(g(x)+\Delta y)-f(g(x))}{\Delta y}\cdot \frac{\Delta y}{\Delta x}
\end{align*}
令$\Delta x\to 0$，由上式得
\[\varphi'(x)=f'(g(x))\cdot g'(x)\]
即
\[[f(g(x))]'=f'(g(x))\cdot g'(x)\]

\item 如果$g'(x)=0$，则当$\Delta x\to 0$时，有两种情形：
\begin{enumerate}
    \item 若$\Delta y=g(x+\Delta x)-g(x)=0$，则
\begin{align*}
    \frac{\varphi(x+\Delta x)-\varphi(x)}{\Delta x}&=\frac{f(g(x+\Delta x))-f(g(x))}{\Delta x}\\
    &=\frac{f(g(x))-f(g(x))}{\Delta x}=0
\end{align*}
\item 若$\Delta y\ne 0$，则
\[\frac{\varphi(x+\Delta x)-\varphi(x)}{\Delta x}=\frac{f(g(x)+\Delta y)-f(g(x))}{\Delta y}\cdot \frac{\Delta y}{\Delta x}\]
令$\Delta x\to 0$，得到
\[\varphi'(x)=f'(g(x))\cdot 0=0\]
\end{enumerate}
合并上述两种情形，都有
\[\varphi'(x)=[f(g(x))]'=0\]
于是定理得证．
\end{enumerate}
\end{proof}

这个定理是说：复合函数对自变量的导数，等于已知函数对中间变量的导数，乘以中间变量对自变量的导数．


\begin{example}
    设$\varphi (x)=2^{\sin x}$, 求$\varphi'(x)$．
\end{example}

    
\begin{solution}
    把$\varphi (x)$看成$u=\sin x$和$f(u)=2^u$的复合函数，于是$\varphi  (x) =2^{\sin x}=f (\sin x)$.

依上述定理，得
\begin{align*}
  \varphi' (x) &= [f (u) ]'\cdot u'_x= (2^u)'\cdot (\sin x)'\\
&=2^u\ln2 \cdot \cos x\\
&=2^{\sin x}\cdot \cos x \cdot \ln2  
\end{align*}
\end{solution}

\begin{example}
    设$y=\sin(x^2+x+1)\cdot \cos(x^3+2x^2)$, 求$y'_x$.
\end{example}

\begin{solution}
    设$u=x^2+x+1$, $v=x^3+2x^2$, 则
\begin{align*}
  y'&= (\sin u)'\cdot u'_x\cdot \cos v+\sin u\cdot (\cos v)'\cdot v'_x\\
&=\cos u\cdot u_x'\cdot \cos v-\sin u\cdot \sin v\cdot  v'_x\\
&= (2x+1) \cos (x^2+x+1) \cos (x^3+2x^2)\\
&\qquad - (3x^2+4x) \sin (x^2+x+1) \sin (x^3+2x^2)  
\end{align*}
\end{solution}
    
\begin{example}
设$f(x)=\frac{x}{\sqrt{9+x^2}}$，求$f'(x)$.
\end{example}


\begin{solution}
\begin{align*}
    f'(x)&=\frac{\sqrt{9+x^2}(x)'-x\left(\sqrt{9+x^2}\right)'}{9+x^2}\\
    &=\frac{\sqrt{9+x^2}-x\cdot \frac{1}{2}(9+x^2)^{-\tfrac{1}{2}}(9+x^2)'}{9+x^2}\\
    &=\frac{\sqrt{9+x^2}-\frac{x^2}{\sqrt{9+x^2}}}{9+x^2}=\frac{9}{\left(9+x^2\right)^{3/2}}
\end{align*}
\end{solution}
    
\begin{blk}
{推论1} 如果函数有两个以上的中间变量，上面的求导法则可以推广使用．
\end{blk}

例如，设$y=f(u)$, $u=g(v)$, $v=h(x)$, 那么对于复合函数．$y=\varphi (x)=f\{g[h(x)]\}$的导数可以这样来求：
\[\frac{\dd y}{\dd x}=\frac{\dd y}{\dd u}\cdot \frac{\dd u}{\dd v}\cdot \frac{\dd v}{\dd x}\]
即：$y'_x=y'_u\cdot u'_v\cdot v'_x$

事实上，依前定理有
\[\frac{\dd y}{\dd x}=\frac{\dd y}{\dd u}\cdot \frac{\dd u}{\dd x}\]
这里$u=g[h(x)]$．再用同一定理
\[\frac{\dd u}{\dd x}=\frac{\dd u}{\dd v}\cdot \frac{\dd v}{\dd x}\]
所以
\[\frac{\dd y}{\dd x}=\frac{\dd y}{\dd u}\cdot \frac{\dd u}{\dd v}\cdot \frac{\dd v}{\dd x}\]


\begin{example}
    设$y=\ln\sqrt{5-2x+3x^4},\quad (5-2x+3x^4>0)$，求$y'_x$.
\end{example}

    
\begin{solution}
    设$y=\ln u$，$u=\sqrt{v}$，$v=5-2x+3x^4$，于是：
\begin{align*}
    \frac{\dd y}{\dd x}&=\frac{\dd y}{\dd u}\cdot \frac{\dd u}{\dd v}\cdot \frac{\dd v}{\dd x}\\
    &=\frac{1}{u}\cdot \left(\frac{1}{2}v^{-\tfrac{1}{2}}\right)\cdot (-2+12x^3)\\
    &=\frac{1}{\sqrt{5-2x+3x^4}}\cdot \frac{1}{2}\frac{1}{\sqrt{5-2x+3x^4}}\cdot (-2+12x^3)\\
    &=\frac{6x^3-1}{5-2x+3x^4}
\end{align*}
\end{solution}

求复合函数的导数，关键在于分析清楚函数的复合关系，适当选定中间变量，根据复合函数求导法则由外向里逐层求导，直到最后一个中间变量对自变量求导为止，每次求导时，必须明确是哪个变量对哪个变量求导，所有这些导数的乘积就是复合函数的导数．

\begin{blk}{推论}
    \[(\ln|g(x)|)'=\frac{g'(x)}{g(x)}\quad (g(x)\ne 0)\]
\end{blk}

\begin{proof}
  \[\ln|g(x)|\text{有意义}\quad \Longleftrightarrow\quad g(x)\ne 0\]  
又
\[|g(x)|'\text{存在}\quad \Longleftrightarrow\quad g(x)\ne 0\]
所以由
\[|g(x)|'=\begin{cases}
    g'(x) ,  & g(x)>0\\
    [-g(x)]'=-g'(x), & g(x)<0
\end{cases}\]
和复合函数求导法则，便得
\[[\ln|g(x)|]'=\frac{1}{|g(x)|}\cdot |g(x)|'=\frac{g'(x)}{g(x)}\quad (g(x)\ne 0)\]
\end{proof}

\begin{example}
    求$(\ln|x^2-5x+4|)'$
\end{example}

\begin{solution}
\[(\ln|x^2-5x+4|)'=\frac{(x^2-5x+4)'}{x^2-5x+4}=\frac{2x-5}{x^2-5x+4}\]
其中：$x^2-5x+4\ne 0$，即$x\ne 1$和$x\ne 4$.
\end{solution}



\begin{example}
求$\left[\frac{(x^2+1)^{\tfrac{1}{2}}(x-1)^2}{(x+1)^{\tfrac{3}{2}}}\right]'$
\end{example}

    
\begin{solution}
设$y=\frac{(x^2+1)^{\tfrac{1}{2}}(x-1)^2}{(x+1)^{\tfrac{3}{2}}}$，则
\begin{align*}
    \ln|y|&=\frac{1}{2}\ln(x^2+1)+2\ln|x-1|-\frac{3}{2}\ln|x+1|\\
    \frac{1}{y}\cdot y'&=\frac{1}{2}\frac{2x}{x^2+1}+\frac{2}{x-1}-\frac{3}{2(x+1)}\\
    &=\frac{3x^3+7x^2-x+7}{2(x^2+1)(x-1)(x+1)}
\end{align*}    
\begin{align*}
    y'&=\frac{3x^3+7x^2-x+7}{2(x^2+1)(x-1)(x+1)}\cdot \frac{(x^2+1)^{\tfrac{1}{2}}(x-1)^2}{(x+1)^{\tfrac{3}{2}}}\\
    &=\frac{(3x^3+7x^2-x+7)(x-1)}{2(x^2+1)^{\tfrac{1}{2}}(x+1)^{\tfrac{5}{2}}}
\end{align*}
    
\end{solution}

\begin{example}
设$f(x)=u(x)^{v(x)}\quad (u(x)>0)$，求$f'(x)$.
\end{example}

\begin{solution}
先求此函数的对数，有
\[\ln f(x)=v(x)\ln u(x)\]
两边对$x$求导，即得
\[\frac{1}{f(x)}\cdot f'(x)=v'(x)\ln u(x)+v(x)\cdot\frac{1}{u(x)}\cdot u'(x)\]
    所以
\[f'(x)=u(x)^{v(x)}\left[v'(x)\ln u(x)+\frac{v(x)u'(x)}{u(x)}\right]\]
\end{solution}

\begin{example}
一人在放风筝，且知风筝在80米的高度飘行，放线速度为2米/秒，求当线长为100米时，风筝和此人之间的水平距离的变率．
\end{example}


\begin{solution}
    设在某时刻$t$, 线长为$\ell(t)$米，水平距离为$x(t)$米，根据勾股定理得到绳长和水平距离的关系式
\[\ell^2 (t) =x^2 (t) +6400\]
两边对时间$t$求导数，得到
\[\frac{\dd \ell^2}{\dd \ell}\cdot \frac{\dd \ell}{\dd t}=\frac{\dd(x^2+6400)}{\dd x}\cdot \frac{\dd x}{\dd t}\]
即
\[\ell\cdot \frac{\dd \ell}{\dd t}=x\cdot \frac{\dd x}{\dd t}\]
但是已知$\frac{\dd\ell}{\dd t}=2({\rm m/s})$，又当$\ell=100$时，
\[x=\sqrt{100^2-6400}=\sqrt{3600}=60\]
于是
\[60\frac{\dd x}{\dd t}=100\times 2\]

$\therefore\quad \frac{\dd x}{\dd t}=\frac{10}{3}=3\frac{1}{3}({\rm m/s})$

答：当$\ell=100$米时，风筝和人之间的水平距离的变率为$3\frac{1}{3}$（米／秒）．
\end{solution}

\begin{ex}
\begin{enumerate}
    \item 求下面函数对于 $x$ 的导数:
\begin{multicols}{2}
\begin{enumerate}
    \item $(x+1)^{2} $ 
\item $(1-3 x)^{3} $
\item  $\sqrt{x+1}$
\item  $\frac{1}{\sqrt{x+1}}$
\item  $\sqrt{x^{2}+1}$
\item  $\frac{1}{\sqrt{x^{2}+1}}$
\item  $\sqrt{2 x^{2}+x-1}$
\item  $\frac{1}{\sqrt{\tan x}}$
\item  $\frac{1}{\sin ^{3} x} $
\item  $\cos (a x+b)$
\item  $\cos ^{3}(a x+b) $
\item  $\cot \sqrt{x} $
\item  $\sec \left(\frac{1}{x}\right)$
\item  $x \sqrt{\sin x}$
\item  $x^{2} \sqrt{\sec 2x}$
\item  $\ln \left(1+x^{2}\right)$
\item $\log _{a}(\sin x)$
\item $\log _{a}(\sec x)$
\item  $\sqrt{\ln x}$
\item $e^{a x^{2}}$
\item $e^{1+x-x^{2}}$
\item  $a^{\sqrt{2x+1}}$
\end{enumerate}
\end{multicols}

 \item    对下面各函数先取对数再求导数:
 \begin{multicols}{2}
    \begin{enumerate}
        \item   $\frac{x-1}{(x+1)^{2}}$
        \item   $\frac{\sqrt{x-1}(x+1)}{x^{2}+1}$
        \item   $\frac{(x-3)(x-4)}{(x-5)}$
        \item   $\sqrt{\frac{x-1}{x+1}}$ 
        \item  $(x+2)^{2}(x+3)^{3}(x+4)^{4}$
\end{enumerate}
\end{multicols}
\item  求下列函数对于 $x$ 的微商:
\begin{enumerate}
\begin{multicols}{2}
  \item $e^{2 x}(2 \cos 3 x+3 \sin 3 x)$
\item $\ln \left(x+\sqrt{x^{2} \pm a^{2}}\right)$
\item $\frac{2 x+1}{\sqrt{3-4 x-4 x^{2}}}$  
\item $\cos[\ln(x^2+9)]$
\item $\ln(1-\tan^2 x)$
\item $x^x+(\ln x)^x$
\end{multicols}
\item $(1-x)^{\tfrac{1}{2}}(1+3x)^{\tfrac{3}{20}}(1-2x)^{-\tfrac{2}{5}}$
\end{enumerate}

\item 若$y=e^{ax}\cos(bx+c)$，求证：
\[\frac{1}{\sqrt{a^2+b^2}}\frac{\dd y}{\dd x}=e^{ax}\cos(bx+c+r)\]
其中：$\tan r=b/a$.

\item 设$f(x)=\begin{cases}
    x^2\sin\frac{1}{x},& x\ne 0\\
    0,& x=0
\end{cases}$，求$f(x)$的导函数．

（提示：对于$x\ne 0$，可用本节的方法得出，而$f'(0)$需直接由定义得出）．

\item 有一个长度为5米的梯子贴靠在铅直的墙上，假设其下端
沿地板以3米/秒的速率离开墙脚而滑动，则
\begin{enumerate}
    \item 当其下端离开墙脚1.4米时，梯子上端下滑之速率为多少？
    \item 何时梯子的上下端能以相同的速率移动？
    \item 何时其上端下滑之速度为4米/秒？
\end{enumerate}

\item 求曲线$y=\cos^2\sqrt{x}$在点$x=0$处的切线方程．
\item 在$\triangle ABC$中，$a=7$cm, $b=5$cm, $c=8$cm, $P$为$AB$边上
一动点，从$A$往$B$以每秒2厘米的速度运动．当$P$点在
$AB$边中点时，$CP$长度的变率是多少？
\item 求曲线$y=2\left(e^x+\frac{1}{5}e^{-3x}\right)$在横坐标$x_0=0$的点的切线方程．
\item 试由公式$1+x+x^2+\cdots+x^n=\frac{x^{n+1}-1}{x-1}$分别导出$1+2x+3x^2+\cdots+nx^{n-1}$与$x+2^2x^2+3^2x^3+\cdots +n^2x^n$的公式．
\end{enumerate}  
\end{ex}


\subsection{反函数求导法则}

\begin{blk}
    {定理} 若函数$y=f(x)$在$x_0$的某一邻域内连续，且严格单调，又函数$f(x)$在点$x_0=f^{-1}(y_0)$的导数$f'(x_0)$存在并且异于零，则反函数$x=f^{-1}(y)$在对应点$y_0=f(x_0)$
的导数存在，并且等于$\frac{1}{f'(x_0)}=\frac{1}{f'(f^{-1}(y_0))}$
，即
\[[f^{-1}(y_0)]'=\frac{1}{f'(x_0)}=\frac{1}{f'[f^{-1}(y_0)]}\]
\end{blk}

\begin{proof}
给数值$y=y_0$以任意增量$\Delta y$, 则反函数$x=f^{-1}(y)$也得到对应增量$\Delta x$, 当$\Delta y\ne 0$时，由于$y=f(x)$的严格单调性，从而$x=f^{-1}(y)$也是严格单调的，也有$\Delta x\ne 0$, 于是有
\begin{align*}
    \frac{f^{-1}(y_0+\Delta y)-f^{-1}(y_0)}{\Delta y}&=\frac{\Delta x}{f(x_0+\Delta x)-f(x_0)}\\
    &=\frac{1}{\frac{f(x_0+\Delta x)-f(x_0)}{\Delta x}}
\end{align*}
现在令$\Delta y\to 0$，则由于反函数$x=f^{-1}(y)$也是连续的，故增量$\Delta x\to 0$，于是
\[\lim_{\Delta y\to 0}\frac{f^{-1}(y_0+\Delta y)-f^{-1}(y_0)}{\Delta y}=\frac{1}{\Lim_{\Delta y\to 0}\frac{f(x_0+\Delta x)-f(x_0)}{\Delta x}}\]
上面等式右边的分母就趋于极限$f'(x_0)\ne 0$，而左边按定义就是$[f^{-1}(y_0)]'$，因此
\[[f^{-1}(y_0)]'=\frac{1}{f'(x_0)}\]
再将$x_0=f^{-1}(y_0)$代入上式，得到
\[[f^{-1}(y_0)]'=\frac{1}{f'[f^{-1}(y_0)]}\]
\end{proof}

\begin{blk}{推论}
\begin{align*}
    (\arcsin x)'&=\frac{1}{\sqrt{1-x^2}},\qquad -1<x<1\\
    (\arccos x)'&=\frac{-1}{\sqrt{1-x^2}},\qquad -1<x<1\\
    (\arctan x)'&=\frac{1}{{1+x^2}},\qquad -\infty<x<\infty\\
    (\arccot x)'&=\frac{-1}{{1+x^2}},\qquad -\infty<x<\infty\\
\end{align*}
\end{blk}

\begin{proof}
如果我们限制正弦函数$x=\sin y$的定义域为$\left(-\frac{\pi}{2},\frac{\pi}{2}\right)$，那么$x=\sin y$的反函数存在，即
\[    y=\arcsin x,\qquad x\in (-1, 1)\]
于是
\[    (\arcsin x)'=\frac{1}{(\sin y)'}=\frac{1}{\cos y}=\frac{1}{\cos(\arcsin x)}\]
由于$\arcsin x\in \left(-\frac{\pi}{2},\frac{\pi}{2}\right)$，
所以$\cos y>0$, 从而
\[    \cos (\arcsin x) =\sqrt{1-\sin^2 (\arcsin x)} =\sqrt{1-x^2}\]
    最后得到
\[  (\arcsin x)'=\frac{1}{\sqrt{1-x^2}},\qquad x\in (-1, 1) \]

我们除去了数值$x=\pm 1$, 因为在它的对应值$y=\pm\frac{\pi}{2}$
处导数$(\sin y)'=\cos y=0$.

在第四册， 我们已经证明$\arcsin x+\arccos x=\frac{\pi}{2}$，所以
\[(\arccos x)'=\left(\frac{\pi}{2}-\arcsin x\right)'=-\frac{1}{\sqrt{1-x^2}},\qquad -1<x<1\]
其余同理可证．
\end{proof}

\begin{example}
求$(\arcsec x)'$.
\end{example}


\begin{solution}
函数$x=\sec y$在区间$\left(0,\frac{\pi}{2}\right)$上，由1递增到$+\infty$，而在区间$\left(\frac{\pi}{2},\pi\right)$上，由$-\infty$递增到$-1$．因此，$x=\sec y$，$y\in\left(0,\frac{\pi}{2}\right)\cup \left(\frac{\pi}{2},\pi\right)$有反函数存在，即：
\[y=\arcsec x,\qquad x\in (-\infty,-1)\cup(1,+\infty)\]

依反函数求导法则，有
\begin{align*}
    (\arcsec x)'&=\frac{1}{(\sec y)'}=\frac{1}{\sec y\cdot \tan y}\\
    &=\frac{1}{\sec(\arcsec x)\cdot \tan(\arcsec x)}\\
    &=\begin{cases}
        \frac{1}{x\sqrt{\sec^2(\arcsec x)-1}}, & 0<\arcsec x<\frac{\pi}{2}\\
        \frac{1}{x\left(-\sqrt{\sec^2(\arcsec x)-1}\right)}, & \frac{\pi}{2}<\arcsec x<\pi \\
    \end{cases}\\
    &=\begin{cases}
        \frac{1}{x\sqrt{x^2-1}}, & x>1\\
       - \frac{1}{x\sqrt{x^2-1}}, & x<-1
    \end{cases}
\end{align*}
\end{solution}

\begin{example}
若已知$(\log_a x)'=\frac{\log_a e}{x},\quad x\in(0,+\infty)$，试利用反函数求导法则求$(a^x)'$    
\end{example}

\begin{solution}
    因为$y=a^x,\quad (a>0,\; a\ne 1,\; -\infty<x<+\infty)$是$x=\log_a y\; (y>0)$的反函数，所以
\begin{align*}
    (a^x)'&=\frac{1}{(\log_a y)'}=\frac{1}{\frac{\log_a e}{y}}=\frac{y}{\log_a e}\\
    &=\ln a\cdot (a^x)=a^x\ln a
\end{align*}
\end{solution}

\begin{example}
求函数$f(x)=2x\arctan 2x-\ln\sqrt{1+4x^2}$的导数．
\end{example}

\begin{solution}
\begin{align*}
    f'(x)&=\left[(2x)'\arctan 2x+2x\cdot \frac{1}{1+(2x)^2}\cdot (2x)'\right]\\
    &\qquad -\frac{1}{\sqrt{1+4x^2}}\cdot \frac{1}{2\sqrt{1+4x^2}}\cdot (1+4x^2)'\\
    &=2\arctan 2x+\frac{4x}{1+4x^2}-\frac{4x}{1+4x^2}\\
    &=2\arctan 2x
\end{align*}
\end{solution}


\begin{example}
求曲线$y=\arcsin x$和$y=\arccos x$的交角．
\end{example}

\begin{solution}
设这两条曲线交于$T(x_0,y_0)$点，则
\[y_0=\arcsin x_0=\arccos x_0,\qquad 0<y_0<\frac{\pi}{2}\]
  由此得
\begin{equation}
    \sin(\arcsin x_0)=\sin(\arccos x_0)
\end{equation}  
即
\[x_0=\sqrt{1-x^2_0}\]
解得
\[x_0=\pm\frac{1}{\sqrt{2}}\]
因为$0<y_0<\frac{\pi}{2}$，所以$x_0>0$. 如此应取$x_0=\frac{1}{\sqrt{2}}$. 由于
\[(\arcsin x)'=\frac{1}{\sqrt{1-x^2}},\qquad (\arccos x)'=\frac{-1}{\sqrt{1-x^2}}\]
所以在点$x=\frac{1}{\sqrt{2}}$处，$y=\arcsin x$的切线的斜率为
\[\tan\alpha\Big|_{x=\tfrac{1}{\sqrt{2}}}=\frac{1}{\sqrt{1-\frac{1}{2}}}=\sqrt{2}\]
又$y=\arccos x$的切线的斜率是
\[\tan\beta\Big|_{x=\tfrac{1}{\sqrt{2}}}=\frac{-1}{\sqrt{1-\frac{1}{2}}}=-\sqrt{2}\]

设$\theta$是这两条曲线在交点$\left(\frac{1}{\sqrt{2}},\frac{\pi}{4}\right)$的夹角，也就是过$T$点的两条切线的交角，于是
\[\tan\theta=\frac{\tan\beta-\tan\alpha}{1+\tan\beta\tan\alpha}=\frac{-\sqrt{2}-\sqrt{2}}{1+\left(-\sqrt{2}\right)\sqrt{2}}=\frac{-2\sqrt{2}}{-1}=2\sqrt{2}\]
因此，这两条曲线的交角$\theta=\arctan 2\sqrt{2}$. 
\end{solution}

\begin{example}
一飞机在离地面2公里的高度，以每小时200公里的速度飞临某目标的上空，以便进行航空摄影，试求飞机飞至该目标正上方时摄影机转动的角速度（图2.9）．
\end{example}

\begin{solution}
我们把该目标取为坐标原点，设飞机和目标的水平距离为$x$, 显然它是时间$t$的函数$x=x(t)$. 现在要求出飞机在目标正上方时，$\frac{\dd \theta}{\dd t}$之值，
为此，先找出$\theta$和$x$的关系，从图2.9中可以看出
\[\tan\theta=\frac{2}{x},\quad  0<\theta\le \frac{\pi}{2}\]

$\therefore\quad \theta=\arctan \frac{2}{x}$
于是
\[\frac{\dd \theta}{\dd t}=\frac{-2}{x^2+4}\frac{\dd x}{\dd t}\]
但根据题意知道$\frac{\dd x}{\dd t}=-200$, 这里负号表示$x$在减小，故得
\[\frac{\dd \theta}{\dd t}=\frac{400}{x^2+4}\]
当飞机在目标正上方时，即$x=0$时，
$\frac{\dd \theta}{\dd t}=100\text{弧度/时}$，将弧度化为角度即为$\frac{\dd \theta}{\dd t}=\frac{5}{\pi}\text{度/秒}$
\end{solution}

\begin{figure}[htp]
    \centering
    \begin{minipage}[t]{0.48\textwidth}
    \centering
\begin{tikzpicture}[>=latex, scale=.9]
\draw[->](-1,0)--(5,0)node[right]{$x$};
\draw[->](0,-1)--(0,4.5)node[right]{$y$};
\node at (0,0)[below left]{$O$};
\draw[dashed](0,3.5)--(4,3.5);
\draw[very thick](4,3.5)--(0,0);
\draw[<->](4,3.5)--node[fill=white]{2}(4,0);
\draw[<->](1,0) arc (0:41.2:1)node[below right]{$\theta$};
\draw[<->](3,3.5) arc (180:180+41.2:1)node[above left]{$\theta$};
\draw[ultra thick, <-](4,3.5)--(4.6,3.5);
\end{tikzpicture}
    \caption{}
    \end{minipage}
    \begin{minipage}[t]{0.48\textwidth}
    \centering
    \begin{tikzpicture}[>=latex, scale=1.5]
        \draw[very thick](-1,2.5)--(-.25,2.5)--(-.25,2)--(-1.5,0)--(1.5,0)--(.25,2)--(.25,2.5)--(1,2.5);
\draw[|<->|](1.75,0)--node[fill=white]{12cm}(1.75,2.5);
\draw[dashed](1,2.5)--(1.75,2.5);
\draw[|<->|](0,-.15)--node[below]{$R=5$cm}(1.5,-.15);
\draw(-1.75/2,1)--(1.75/2,1);
\draw[<->](0,0)--node[fill=white]{$y$}(0,1);
\draw[|<->|](0,.8)--node[fill=white]{$r$}(1.75/2,.8);
    \end{tikzpicture}
    \caption{}
    \end{minipage}
    \end{figure}

\begin{example}
有一底半径为5cm, 高为12cm的正圆锥形容器（图2.10）, 以每秒2$\rm cm^3$的速度往容器中倒水，试求容器水位等于锥高一半时，水面上升速度．
\end{example}


\begin{solution}
    设在某时刻$t$, 容器水高为$y$, 则此时水的体积$V$为
 \begin{align*}
     V&=\frac{\pi R^2h}{3}-\frac{\pi r^2}{3}(h-y)\\
     &=\frac{\pi}{3}\cdot 5^2\cdot 12-\frac{\pi r^2}{3}(12-y)\\
     &=100\pi-\frac{\pi r^2}{3}(12-y)
 \end{align*}
而
\[\frac{r}{R}=\frac{h-y}{h}\]
即
\[r=\frac{R}{h}(h-y)=\frac{5}{12}(12-y)\]
代入上式，得
\begin{align*}
    V&=100\pi-\frac{\pi}{3}\cdot \frac{5^2}{12^2}(12-y)^3\\
    &=100\pi-\frac{25\pi}{432}(12-y)^3
\end{align*}
又由题设知：$V=2t$，所以
\[2t=100\pi-\frac{25\pi}{432}(12-y)^3\]
即
\begin{equation}
    t=50\pi-\frac{25\pi}{864}(12-y)^3
\end{equation}
这就确定了水的高度与时间$t$的关系，水面上升速度就是
$y$关于时间$t$的变化率，也就是$\frac{\dd y}{\dd t}$．根据反函数求导法则，有
\[\frac{\dd y}{\dd t}=\frac{1}{\frac{\dd t}{\dd y}}\]
由(2.8)得
\[\frac{\dd t}{\dd y}=-\frac{25\pi}{864}\cdot 3(12-y)^2(12-y)'=\frac{25\pi}{288}(12-y)^2\]

$\therefore\quad \frac{\dd y}{\dd t}=\frac{288}{25\pi(12-y)^2}$

当$y=\frac{1}{2}h=6$时，
\[\left.\frac{\dd y}{\dd t}\right|_{y=6}=\frac{288}{25\x 36\pi}\approx 0.102({\rm cm/s})\]

答：水面上升速度为每秒0.102厘米．
\end{solution}

\begin{ex}
\begin{enumerate}
    \item 求下列函数的微商：
\begin{multicols}{2}
\begin{enumerate}
  \item $\arcsin \frac{x}{4}$
\item  $\arccos x^{2}$
\item  $\arctan \frac{1}{x}$
\item  $\arcsin \sqrt{x}$
\item  $(\arcsin x)^{2}$
\item  $\arctan (\tan  x)$
\item $\arctan \frac{2 x}{1-x^{2}}$
\item  $\arctan (2 \tan  x)$
\item $x^{2} \arccot \frac{x}{a}$
\item $\ln \sqrt{1+4 x^{2}} \cdot \arccot \frac{1-x}{1+x}$
\item  $\frac{1-\arccos 2 x}{\arctan x}$
\item  ${\rm arccsc}\, x$
\end{enumerate}
\end{multicols}

\item 证明
   \[
    \frac{\dd}{\dd x}\left\{\arctan \left(\frac{\sin x}{\cos 2 x}\right)\right\}=\frac{3 \cos x-\cos 3 x}{2(1-\sin x \sin 3 x)}
   \]
    \item 求下列函数的导数：
\begin{enumerate}
\item $\arcsin \frac{6 x^{2}}{\sqrt{4+81 x^{8}}}$
\item $\arctan \frac{1}{1-3 x^{2}}-\arctan \frac{1}{1+3 x^{2}}$
\item $\arccot \frac{4 \sqrt{x}}{1-4 x}$
\end{enumerate}

\item 一圆锥形容器，深10尺，上顶圆半径为4尺；求
\begin{enumerate}
    \item 灌入水时，水的体积$V$对水面高度$h$的变化率；
    \item 体积$V$对容器截面圆半径$R$的变化率．
\end{enumerate}

\item 一圆锥形容器底面朝上放着，它的顶角为$2\arctan\frac{3}{4}$, 
今向里面倒进某种液体，试求
\begin{enumerate}
\item 当液面半径$r$为3, 半径增加的速度$\frac{\dd r}{\dd t}$为$\frac{1}{4}$时，
体积增加的速度$\frac{\dd V}{\dd t}$；
\item 当液面半径为6, 体积增加的速度为24时，半径增加的速度．
\end{enumerate}

\item 水从高为18厘米，底半径为6厘米的圆锥形漏斗中流入半
径为5厘米的圆柱形筒内．已知此漏斗中水深为12厘米时，漏斗中水面的下降速度为1厘米/分，求此时圆筒中水面的上升速度．
\item 底数$a$为什么值时，直线$y=x$才能与对数曲线$y=\log_a x$相切？在何处相切？
\item 求曲线$y=e^{-x}$上的一点，使过该点的切线与直线$y=-ex$平行，并写出该点的法线方程．
\item 求曲线$y=\frac{1}{2}(1+2x^2\pm \sqrt{1+4x^2})$上横坐标$x=u$的点
处的切线方程．这切线还与曲线交于何处？$y$的反函数为何？
\end{enumerate}    
\end{ex}

\subsection{隐函数的导数}

到现在为止，我们所讨论的函数都是可以由自变量明显地表示出来的函数，例如，方程
$y=x^3+1$
就定义了一个明显的函数$f$, 这里$f(x)=x^3+1$. 这种函数叫做显函数．但是并不是所有的函数都可以用这种明显的表达式来定义的，例如由下面的$x,y$的方程
\[x^3-x=y^3-y^2+24\]
就不容易解出用$x$表示$y$的解，或用$y$表示$x$的解．但是，却有可能存在一个函数$f$, 使得方程
\[x^3-x=f^3(x)-f^2(x)+24\]
对于函数$f$的定义域中的一切$x$值都成立，这样一个函数，如果存在的话，就叫做被这个方程定义的隐函数．最简的情形是二元一次方程：
\[Ax+By+C=0\]
当$B\ne 0$时，就定义一个隐函数$f:\; \mathbb{R}\mapsto \mathbb{R}$, 这里$f(x)=-\frac{A}{B}x-\frac{C}{B}$. 

应该注意：
\begin{enumerate}
\item 方程$F(x,y)=0$的解有可能不存在，例如方程$x^2+y^2+1=0$没有实数解，因此，由它不能定义一个隐函数．
\item 由方程$F(x,y)=0$可能定义几个函数．例如，由方程
$x^2+y^2=16$
就定义了两个函数：
\[f_1 (x) =\sqrt{16-x^2},\qquad x\in [-4, 4] \]
和
\[f_2 (x) =-\sqrt{16-x^2},\qquad x\in  [-4, 4] \]
\end{enumerate}

\begin{blk}
    {定义} 方程
    \begin{equation}
 F (x,y) =0       
    \end{equation}
给出了变元$x,y$之间的一个关系．如果存在一个函数$y=
f(x)$, 或$x=g(y)$使得对于函数的定义域中的一切值，等式$F(x,f(x))=0$或$F(g(y),y)=0$恒成立，我们就说方程(2.9)定义了一个\textbf{隐函数}：
\[y=f(x)\qquad \text{或}\qquad x=g(y)\]
\end{blk}

要对方程(2.9)给出的隐函数求导数，通常的办法是直接对这个方程求导数，而不是先解出显函数，再来求导数．


\begin{example}
设$xy+x^2-1=0$，求$y'_x$.
\end{example}

\begin{solution}
先把$y$看成$x$的函数，即$y=f(x)$, 然后，在等式
$xf (x) +x^2-1=0$
的两边对$x$求导数，得到
\begin{align*}
    (x\cdot f(x))'+(x^2)'&=0\\
xf'(x)+f(x)+2x&=0
\end{align*}

$\therefore\quad f' (x) =-\frac{2x+f (x)}{x}=-\frac{2x+\frac{1-x^2}{x}}{x}=-\frac{1+x^2}{x^2}$    
\end{solution}

\begin{example}
    方程$x^2+y^2=16$定义了两个可微函数：
$y=\sqrt{16-x^2}$和$y=-\sqrt{16-x^2},\; x\in [-4, 4]$. 求每一个函数的导数．
\end{example}


\begin{solution}
    我们用$f(x)$代表上面两个函数中的任何一个，则
\[x^2+f^2 (x) =16,\qquad  x\in  [-4, 4] \]
对上面等式的两边微分，得到
\begin{align*}
    2x+2f (x) \cdot f' (x) &=0\\
    f'(x)&=-\frac{x}{f(x)}
\end{align*}
即
\begin{align*}
    \left(\sqrt{16-x^2}\right)'&=-\frac{x}{\sqrt{16-x^2}}\\
    \left(-\sqrt{16-x^2}\right)'&=\frac{x}{\sqrt{16-x^2}},\qquad x\in[-4,4]
\end{align*}
\end{solution}

\begin{rmk}
应用这个方法求隐函数的导数时，我们作了两点假设：
\begin{enumerate}
    \item 存在一个由方程$F(x,y)=0$定义的函数；
    \item 这样定义的函数是可微的．
\end{enumerate} 
至于在什么条件下，这种假设是合理的问题，已经超出了本书的范围，故不在此讨论．
\end{rmk}

\begin{example}
求证椭圆$\frac{x^2}{a^2}+\frac{y^2}{b^2}=1$在点$(x_0,y_0)$处切线的方程为$\frac{x_0x}{a^2}+\frac{y_0y}{b^2}=1$.
\end{example}

\begin{proof}
\begin{align*}
    \left(\frac{x^2}{a^2}+\frac{y^2}{b^2}\right)'_x&=(1)'_x\\
    \frac{2x}{a^2}+\frac{2y}{b^2}\cdot y'_x&=0
\end{align*}
若$y\ne 0$，则：
\[y'_x=-\frac{b^2x}{a^2y}\]
于是，在点$(x_0,y_0)$处（这里$y_0\ne 0$）的切线方程为
\[y-y_0=-\frac{b^2x_0}{a^2y_0}(x-x_0)\]
即：$b^2x_0x+a^2y_0y=b^2x^2_0+a^2y^2_0$

$\because\quad $点$(x_0,y_0)$在椭圆上，

$\therefore\quad b^2x_0^2+a^2y^2_0=a^2b^2$，因此，所求切线方程是
\[b^2x_0 x+a^2y_0 y=a^2b^2\]
也就是
\[\frac{x_0x}{a^2}+\frac{y_0y}{b^2}=1\]

若$y=0$则$x=\pm a$, 显然在点$(-a,0)$和点$(a,0)$处的切线方程分别是$x=-a$, $x=a$, 它们是方程$\frac{x_0x}{a^2}+\frac{y_0y}{b^2}=1$的特殊形式．
\end{proof}
    
\begin{example}
    如果将抛物线$y^2=2px$绕它的对称轴旋转成旋转抛物面，那么与对称轴平行的光线射到曲面上，经曲面反射后通过焦点，试证明之．
\end{example}


\begin{proof}
从物理学知道，光线经镜面反射后，要服从光的反射定律：
\begin{enumerate}
    \item 在均匀介质中，光沿直线传播．
    \item 当光线碰到平面镜时，光线要反射．入射线与镜面法线夹角（记作$i$）等于反射线与镜面法线的夹角（记作$r$），见图2.11, 角$i$称为入射角，角$r$称为反射角．
\item 光从曲面镜的反射也遵守同样的法则：
  \[  \text{入射角}i=\text{反射角}r\]
    此时镜面的法线定义为与镜面在入射点的切线相垂直的直线．
\end{enumerate}

\begin{figure}[htp]
    \centering
    \begin{tikzpicture}[>=latex, scale=1]
        \fill[pattern=north east lines](-2,-.25) rectangle (2,0);
        \draw(-2,0)node[left]{镜面}--(2,0);
        \draw[->, very thick](140:3)node[left]{入射线}--(140:2);
        \draw[->, very thick](140:2)--(0,0)--(40:3)node[right]{反射线};
        \draw(0,0)--(0,3)node[above]{法线};
        \draw[<->, thick](90:1.8) arc (90:40:1.8);
        \draw[<->, thick](90:1.8) arc (90:140:1.8);
        \node at (60:2){$r$};
        \node at (120:2){$i$};
            \end{tikzpicture}
    \caption{}
\end{figure}

假设光线经位于$xy$平面上的抛物线$y^2=2px$形的镜面反射．当入射线与x轴平行且与抛物线交于$P_1(x_1,y_1)$点时，我们要求反射线的方
程，并证明它通过焦点$F\left(\frac{p}{2},0\right)$. 如果过$P_1(x_1,y_1)$点的切线与$x$轴成$\theta$角，那么平行于$x$轴的入射光线与法线的夹角是$90^{\circ}-\theta$. 由光的反射定律，反射光线与法线的夹角也是$90^{\circ}-\theta$. 于是，从图2.12上看出，反射光线与$x$轴的夹角是$2\theta$.

现在，在方程$y^2=2px$的两边对$x$求导数，得到
\[2yy'=2p\quad \Rightarrow\quad y'=\frac{2p}{2y}=\frac{p}{y}\]
所以，由导数的几何意义得出，在$P_1(x_1,y_1)$点的切线的斜率是
\[\tan\theta=y'_{x_1}=\frac{p}{y_1}\]
于是反射光线的斜率是
\begin{align*}
    \tan2\theta=\frac{2\tan\theta}{1-\tan^2\theta}&=\frac{2\cdot \frac{p}{y_1}}{1-\frac{p^2}{y^2_1}}\\
    &=\frac{2py_1}{y^2_1-p^2}=\frac{2py_1}{2px_1-p^2}=\frac{2y_1}{2x_1-p}
\end{align*}
从而，得出反射线方程
\[y-y_1=\frac{2y_1}{2x_1-p}(x-x_1) \]
令$y=0$, 得到反射光线与$x$轴交点的横坐标，
\[0-y_1=\frac{2y_1}{2x_1-p}(x-x_1)\quad \Rightarrow\quad x=\frac{p}{2}\]
即反射线通过焦点$F\left(\frac{p}{2},0\right)$. 由于$P_1(x_1,y_1)$是抛物线上
任意一点，所以只要入射光线与$x$轴平行，反射光线必通过焦
点．如果把光源放在焦点$F\left(\frac{p}{2},0\right)$上，那么从焦点发出的
光线经抛物线镜面反射后的反射光线是与对称轴（$x$轴）平行的光线．我们常利用这个原理获得平行光．
\end{proof}

\begin{figure}[htp]\centering
\begin{tikzpicture}[>=latex, scale=1.2]
\draw[->](-1.5,0)--(5,0)node[right]{$x$};
\draw[->](0,-3)--(0,3)node[right]{$y$};
\draw[domain=2.5:-2.5, samples=100, very thick]plot({0.5*\x*\x}, \x)node[right]{$y^2=2px$};
\draw[domain=-1:3.5, samples=10, thick]plot(\x, {0.707*(\x+1)});
\draw[thick](.5,0)node[below right]{$F\left(\frac{p}{2},0\right)$}--(1,1.414)--(4,1.414);
\draw[->, thick](4,1.414)--(3,1.414);
\tkzDefPoints{-1/0/A, 1/1.414/P, .5/0/F, 4/1.414/B, 3/2.828/C, 5/0/D}
\tkzMarkAngles[mark=none, size=.45](D,A,P D,F,P A,P,F B,P,C)
\tkzLabelAngle[pos=.6](D,A,P){$\theta$}
\tkzLabelAngle[pos=.6](D,F,P){$2\theta$}
\tkzLabelAngle[pos=.6](A,P,F){$\theta$}
\tkzLabelAngle[pos=.6](B,P,C){$\theta$}

\draw[domain=.5:1.5, samples=10]plot(\x, {1.414*(-\x+2)});
\tkzDefPoints{1.5/0.707/E}
\tkzMarkAngles[mark=none, size=.6, <->](E,P,B)
\tkzLabelAngle[pos=1.2](E,P,B){$90^{\circ}-\theta$}
\node at (P) [left]{$P_1(x_1,y_1)$};
\node at (0,0)[below left]{$O$};
\end{tikzpicture}
\caption{}
\end{figure}

\begin{ex}
\begin{enumerate}
    \item 求下列隐函数在指定点的导数$\frac{\dd y}{\dd x}$:
\begin{enumerate}
    \item $y=\cos x+\frac{1}{2}\sin y$, 点$\left(\frac{\pi}{2},0\right)$
    \item $ye^x+\ln y=1$, 点$\left(0,1\right)$
    \item $x^2+xy+y^2=1$, 点$\left(1,-1\right)$
    \item $x^2y^2+x^2+y^2=9$, 点$\left(1,2\right)$
\end{enumerate}
    \item 求下列隐函数的导数$\frac{\dd y}{\dd x}$:
\begin{multicols}{2}
\begin{enumerate}
    \item $\frac{x^2}{a^2}+\frac{y^2}{b^2}=1$
    \item $y^2=8x$
    \item $x^2=2y$
    \item $\sqrt{x}+\sqrt{y}=1$
    \item $x^3+y^3=a^3$
    \item $x^{\tfrac{2}{3}}+y^{\tfrac{2}{3}}=a^{\tfrac{2}{3}}$
    \item $y-\cos(x+y)=0$
    \item $y=x+\arctan y$
    \item $y=1-\ln(x+y)+e^y$
    \item $\arctan \frac{y}{x}=\ln\sqrt{x^2+y^2}$
\end{enumerate}
\end{multicols}

\item 求下列方程在指定点的导数$\frac{\dd y}{\dd x}$:
\begin{enumerate}
    \item $x-\sqrt{xy}+y=2$，点$(x_1,y_1)$
    \item $(ax)^{\tfrac{2}{3}}-(by)^{\tfrac{2}{3}}=(a^2+b^2)^{\tfrac{2}{3}}$，点$(x_0,y_0)$
    \item $(x+y)^{\tfrac{2}{3}}+(x-y)^{\tfrac{2}{3}}=2a^{\tfrac{2}{3}}$在$x=0$的点；在$x=y$的点；在$x=-y$的点
    \item $x^2+xy+y^2=c^2$，在斜率等于$-1$的点
\end{enumerate}
\item 若$x\sqrt{y}+y\sqrt{x}=c$，求$\frac{\dd y}{\dd x}$和$\frac{\dd x}{\dd y}$.
\item 求椭圆$x^2+4y^2=9$的切线方程，使它和直线$2x+3y=0$
平行．
\item 曲线$y=(x-2)(x-3)(x-4)$和$x$轴交于$P(2, 0)$, $Q(3,0)$和$R(4, 0)$点．
\begin{enumerate}
    \item 求证在$P$点和$R$点处的切线互相平行；
    \item 求曲线在$Q$点处的法线与$y$轴的交点．
\end{enumerate}
\item 求曲线 $2x^3-xy+8=0$上的$P$点，使在该点的切线通过
原点，过$P$点的切线与曲线交于另一点$Q$的坐标．
\item \begin{enumerate}
    \item 求证直线$\ell x +my+n=0$与椭圆$\frac{x^2}{a^2}+\frac{y^2}{b^2}=1$, 双曲线$\frac{x^2}{a^2}-\frac{y^2}{b^2}=1$, 抛物线$y^2=2px$相切的条件分别是$a^2\ell^2+b^2m^2=n^2$, $a^2\ell^2-b^2m^2=n^2$, $pm^2=2n\ell$.
    \item 证明对于任意$k$, 直线$y=kx+\sqrt{a^2k^2+b^2}$是椭圆曲线$\frac{x^2}{a^2}+\frac{y^2}{b^2}=1$的切线．    
\end{enumerate}
\item \begin{enumerate}
    \item 求过$P(5, 3)$点与椭圆$\frac{x^2}{16}+{y^2}=1$相切的直线
    方程.
    \item 求过(a)中的两切点的直线方程．
\end{enumerate}

\item 椭圆$\frac{x^2}{16}+{y^2}=1$的切线与$x$轴成$60^{\circ}$角，求该切线的方程．
\item 通过双曲线$\frac{x^2}{a^2}-\frac{y^2}{b^2}=1$的焦点而与横轴垂直的弦的
两个端点分别记作$A$和$B$, 求证过$A$、$B$两点的切线交于
横轴上的点$Q$, 且$Q$点离中心的距离等于$\frac{a}{e}$.
\item 求证过双曲线上一点$P$的切线与过$P$点的焦点半径成等角．
\end{enumerate}
\end{ex}

\subsection{高阶导数}
对于给定的可微函数$f(x)$, 它的导数$f'(x)$仍是$x$的函数．如果$f'(x)$是可微的，它的导函数$(f'(x))'=f''(x)$或$f^{(2)}(x)$称为$f(x)$的\textbf{二阶导函数}，……如此类推，$f(x)$
的$n$阶导函数$f^{(n)}(x)$, 可以归纳地定义为：
\begin{align*}
  f ^{(1)} (x)&=f'(x)\\
  f ^{(n)} (x)&=\left(f^{(n-1)}(x)\right)',\qquad n=2,3,\ldots  
\end{align*}
因此，求高阶导函数，实际上就是连续地求一阶导函数．

高阶导函数有时也记作$\frac{\dd^n y}{\dd x^n}$或$\frac{\dd^n }{\dd x^n}f(x)$. 例如
\begin{align*}
    f''(x)&=\frac{\dd}{\dd x}f'(x)=\frac{\dd}{\dd x}\left(\frac{\dd y}{\dd x}\right)=\frac{\dd^2 y}{\dd x^2}\\
    f^{(3)}(x)&=\frac{\dd}{\dd x}\left(\frac{\dd^2 y}{\dd x^2}\right)=\frac{\dd^3 y}{\dd x^3}
\end{align*}

二阶导数有它的实际意义，例如动点的速度$v$等于它所
经过的路程$S$对于时间$t$的变率：$v=\frac{\dd S}{\dd t}$；而加速度$a$是速度$v$
对于时间$t$的变率：$a=\frac{\dd v}{\dd t}$，
这就是说加速度是路程$S$对于时间$t$的二阶导数：$a=\frac{\dd^2 S}{\dd t^2}$．
在以后的学习中，我们还会知道二阶导数有它的重要的几何意义．我们已经知道一次函数$f(x)=ax+b\; (a\ne 0)$的图象是一条直线，而它的二阶导数$f''(x)=(f'(x))'=(a)'=0$. 因此，若一函数的二阶导数$f''(x)\ne 0$, 那么它的图象就不是直线．

\begin{example}
设$y=\frac{x}{1-x}$，求$\frac{\dd^2 y}{\dd x^2}$, $\frac{\dd^3 y}{\dd x^3}$.
\end{example}

\begin{solution}
\begin{align*}
    y&=\frac{x}{1-x}=-1+\frac{1}{1-x}\\
\frac{\dd y}{\dd x}&=-\frac{(1-x)'}{(1-x)^2}=\frac{1}{(1-x)^2}
\end{align*}
\begin{align*}
    \frac{\dd^2 y}{\dd x^2}&=\frac{-[(1-x)^2]'}{(1-x)^4}=\frac{-2(1-x)(-1)}{(1-x)^4}=\frac{2}{(1-x)^3}\\
    \frac{\dd^3 y}{\dd x^3}&=\frac{-2[(1-x)^3]'}{(1-x)^6}=\frac{-2\cdot 3(1-x)^2(-1)}{(1-x)^6}=\frac{6}{(1-x)^4}  
\end{align*}
\end{solution}


\begin{example}
若$f(x)=\sqrt{a^2+x^2}$，求$f''(x)$
\end{example}


\begin{solution}
\textbf{解法1: }
\begin{align*}
    f'(x)&=\frac{1}{2}(a^2+x^2)^{-\tfrac{1}{2}}(a^2+x^2)'\\
    &=\frac{1}{2}(a^2+x^2)^{-\tfrac{1}{2}}(2x)=\frac{x}{\sqrt{a^2+x^2}}\\
\end{align*}
\begin{align*}
f''(x)&=\frac{\sqrt{a^2+x^2}\cdot x'-x\left(\sqrt{a^2+x^2}\right)'}{a^2+x^2}\\
&=\frac{\sqrt{a^2+x^2}-\frac{x^2}{\sqrt{a^2+x^2}}}{{a^2+x^2}}\\
&=\frac{a^2+x^2-x^2}{(a^2+x^2)^{\tfrac{3}{2}}}=\frac{a^2}{(a^2+x^2)^{\tfrac{3}{2}}}
\end{align*}

\textbf{解法2: } 将原来的函数变形为隐函数，用隐函数微分法求二阶导数．

设$y=\sqrt{a^2+x^2}$, 则$y^2=a^2+x^2$, 对$x$微分得：
\[y\cdot y'_x=x\]    
再对$x$微分得：
\[y\cdot y_x''+(y_x')^2=1\]
将$y'=\frac{x}{y}$代入，得
\[y\cdot y''_x+\frac{x^2}{y^2}=1\]
因此：
\[y''_x=\frac{1-\frac{x^2}{y^2}}{y}=\frac{1-\frac{x^2}{x^2+a^2}}{\sqrt{a^2+x^2}}=\frac{a^2}{(a^2+x^2)^{\tfrac{3}{2}}}\]
\end{solution}

\begin{example}
    若$y=\frac{x^2}{\sqrt{x^2+1}}$，求证：
\[(x^2+1)\frac{\dd^2 y}{\dd x^2}+3x\frac{\dd y}{\dd x}-y=2\sqrt{x^2+1}\]
\end{example}

\begin{proof}
将原式去分母
\[y\sqrt{x^2+1}=x^2\]
对$x$微分得
\[\sqrt{x^2+1}\cdot\frac{\dd y}{\dd x}+y\cdot \frac{x}{\sqrt{x^2+1}}=2x\]
两边乘以$\sqrt{x^2+1}$，得到
\[(x^2+1)\cdot \frac{\dd y}{\dd x}+y\cdot x=2x\sqrt{x^2+1}\]
再对$x$微分得：
\[(x^2+1)\frac{\dd^2 y}{\dd x^2}+(2x)\cdot \frac{\dd y}{\dd x}+x\frac{\dd y}{\dd x}+y=2\sqrt{x^2+1}+\frac{2x^2}{\sqrt{x^2+1}}\]
即：
\[(x^2+1)\frac{\dd^2 y}{\dd x^2}+3x\frac{\dd y}{\dd x}+y=2\sqrt{x^2+1}+2y\]
因此：
\[(x^2+1)\frac{\dd^2 y}{\dd x^2}+3x\frac{\dd y}{\dd x}-y=2\sqrt{x^2+1}\]
\end{proof}

有时需要求一些函数的$n$阶导数，但是在一般情形下，函数的$n$阶导数是如此的复杂，以致求函数的$n$阶导数的通式是办不到的，然而有一些函数，我们不难写出它的$n$阶导数的通式．

\begin{example}
  设$y=(ax+b)^m$, 求$y^{(n)}\; (m>n)$  
\end{example}

\begin{solution}
\begin{align*}
y'&=ma(ax+b)^{m-1}\\
y''&=m(m-1)a^2(ax+b)^{m-2}\\
y^{(3)}&=m(m-1)(m-2)a^3(ax+b)^{m-3}\\
\cdots\quad &\cdots\quad \cdots
\end{align*}
由此推知：
\[y^{(n)}=m(m-1)\cdots(m-n+1)a^n(ax+b)^{m-n}=\frac{m!}{(m-n)!}a^n(ax+b)^{m-n}\]
这个结果的正确性可用数学归纳法证明．
\end{solution}


\begin{example}
设$y=\frac{1}{x^2-4x+3}$，求$y^{(n)}$.
\end{example}


\begin{solution}
\[y=\frac{1}{x^2-4x+3}=\frac{1}{2}\left(\frac{1}{x-3}-\frac{1}{x-1}\right)\]
\begin{align*}
    y'&=\frac{1}{2}\left[\frac{-1}{(x-3)^2}-\frac{-1}{(x-1)^2}\right]=\frac{(-1)}{2}\left[\frac{1}{(x-3)^2}-\frac{1}{(x-1)^2}\right]\\
    y''&=\frac{(-1)}{2}\left[\frac{-2(x-3)}{(x-3)^4}-\frac{-2(x-1)}{(x-1)^4}\right]=\frac{(-1)^2 2!}{2}\left[\frac{1}{(x-3)^3}-\frac{1}{(x-1)^3}\right]\\
    \cdots\quad &\cdots\quad \cdots
\end{align*}
假设一般地有：
\[y^{(n)}=\frac{(-1)^n n!}{2}\left[\frac{1}{(x-3)^{n+1}}-\frac{1}{(x-1)^{n+1}}\right]\]
那么：
\begin{align*}
    y^{(n+1)}&=\frac{(-1)^n n!}{2}\left[\frac{-(n+1)(x-3)^n}{(x-3)^{2(n+1)}}-\frac{-(n+1)(x-1)^n}{(x-1)^{2(n+1)}}\right]\\
    &=\frac{(-1)^{n+1} (n+1)!}{2}\left[\frac{1}{(x-3)^{n+2}}-\frac{1}{(x-1)^{n+2}}\right]
\end{align*}
即仍有此规律．因此，对于任何自然数$n$，都有
\[y^{(n)}=\frac{(-1)^n n!}{2}\left[\frac{1}{(x-3)^{n+1}}-\frac{1}{(x-1)^{n+1}}\right]\]
成立．
\end{solution}

\begin{example}
    设$y=\sin(ax+b)$，求$y^{(n)}$.
\end{example}


\begin{solution}
    \[y'=a\cos(ax+b)=a\sin\left(ax+b+\frac{\pi}{2}\right)\]
因此，求$\sin(ax+b)$的导数是将函数乘以$a$，同时将幅角加上$\frac{\pi}{2}$.

不难推知：
\[y^{(n)}=a^n\sin\left(ax+b+\frac{n\pi}{2}\right)\]
请读者用数学归纳法证明这个猜想正确．
\end{solution}

\begin{example}
设$f(x)=a_0+a_1x+a_2x^2+\cdots+a_nx^n$，$h$为任取非零实数，求证：
\[f(a+h)=f(a)+f'(a)h+\frac{f''(a)}{2!}h^2+\cdots+\frac{f^{(n)}(a)}{n!}h^n\]
\end{example}


\begin{proof}
已给$f(x)=a_0+a_1x+\cdots+a_nx^n$，令$x=a+h$，并按$h$的升幂展开每一项，得到一个形式如
\begin{equation}
    f(a+h)=C_0+C_1 h+C_2 h^2+\cdots +C_n h^n
\end{equation}
的表达式．

我们将说明，系数$C_0,C_1,\ldots,C_n$由$f(x)$在点$a$的函数值$f(a)$以及直到$n$阶为止的各阶导数在$x=a$的值表示，即
\[C_0=f (a) ,\quad  C_i =\frac{1}{i!}f^{(i)} (a) ,\qquad i=1, 2,\ldots ,n\]

为了证明这个事实，(2.10)中把$h$看作自变量，首先令$h\to 0$, 立刻得到
\[f (a) =C_0\]

然后，应用复合函数求导法则，对方程(2.10), 关于$h$逐次求导，依次得到：
\begin{align*}
f'(a+h)&=C_1+2C_2 h+3C_3h^2+\cdots +nC_nh^{n-1}\\
f''(a+h)&=2C_2+3\x 2C_3h+\cdots+n(n-1)C_nh^{n-2}\\
f^{(3)}(a+h)&=3!C_3+\frac{4!}{1!}C_4h+\frac{5!}{2!}C_5h^2+\cdots+\frac{n!}{(n-3)!}C_nh^{n-3}\\
\cdots\quad &\cdots\quad \cdots\\
f^{(i)}(a+h)&=i! C_i+\frac{(i+1)!}{1!}C_{i+1}h+\cdots+\frac{n!}{(n-i)!}C_nh^{n-i}\\
\cdots\quad &\cdots\quad \cdots\\
f^{(n-1)}(a+h)&=(n-1)!C_{n-1}+n!C_nh\\
f^{(n)}(a+h)&=n!C_n
\end{align*}
把$h=0$代入上面一个方程中，便得到
\[C_1=f^{(1)}(a),C_2=\frac{f^{(2)}(a)}{2!},C_3=\frac{f^{(3)}(a)}{3!},\ldots,C_i=\frac{f^{(i)}(a)}{i!},\ldots,C_n=\frac{f^{(n)}(a)}{n!}\]

把上面$C_0=f(a),\; C_i=\frac{f^{(i)}(a)}{i!},\quad i=1,2,\ldots,n$代入(2.10)中，最后得到：
\begin{equation}
    f(a+h)=f(a)+f'(a)h+\frac{f^{(2)}(a)}{2!}h^2+\cdots+\frac{f^{(n)}(a)}{n!}h^n
\end{equation}
这个公式叫做多项式的泰勒公式．

在(2.11)中，令$a+h=x$, 则(2.11)可以改写成
\[    f(x)=f(a)+f'(a)(x-a)+\frac{f^{(2)}(a)}{2!}(x-a)^2+\cdots+\frac{f^{(n)}(a)}{n!}(x-a)^n\]
\end{proof}

\begin{ex}
\begin{enumerate}
    \item 求下列函数的$n$阶导数：
\begin{multicols}{2}
\begin{enumerate}
    \item $f(x)=\frac{1-x}{1+x}$
    \item $f(x)=\sin^2 x$
\end{enumerate}
\end{multicols}
    \item 若$y=\tan(m\cdot \arctan x)$，验证
\[(1+x^2)\frac{\dd^2 y}{\dd x^2}=2(my-x)\frac{\dd y}{\dd x}\]
\item 若$y=\frac{A\cos mx+B\sin mx}{x^n}$，其中$A,B,m,n$是常数，验证
\[\frac{\dd^2 y}{\dd x^2}+2\frac{n}{x}\frac{\dd y}{\dd x}+\left[\frac{n(n-1)}{x^2}+m^2\right]y=0\]

\item 若$h(x)=f(x)g(x)$，其中$f(x)$, $g(x)$和$h(x)$都是$n$阶可微，求证：
\begin{align*}
    h^{(n)}(x)&=f(x)\cdot g^{(n)}(x)+C^1_n f^{(1)}(x)g^{(n-1)}(x)+\cdots \\
    &\qquad +C^h_n f^{(h)}(x)g^{(n-h)}(x)+\cdots\\
    &\qquad  +C^{n-1}_n f^{(n-1)}(x)g^{(1)}(x)+C^n_n f^{(n)}(x)g(x)
\end{align*}
\item 求下列函数的$n$阶导数：
\begin{multicols}{2}
\begin{enumerate}
\item $f(x)=x\ln x$
\item $f(x)=xe^x$
\item $f(x)=e^{ax}\sin\beta x$
\end{enumerate}
\end{multicols}
\item 证明 $x^3-7x^2+15x-9=0$有一个二重根．
\item 检验方程 $x^3-6x^2+11x-6=0$有无重根？
\item 方程$x^4-6x^2-8x-3=0$有重根，解此方程．
\item 应用泰勒公式将多项式$(x+a)^n$按$x$的升幂展开．
\item 若$f(x)=x^3-3kx+2k+8$在实数范围内能分解成一个二重因式，试求$k$值和它的因式乘积．
\end{enumerate}   
\end{ex}

\subsection{参数方程表示的函数的导数}
用方程$y=f(x)$来表示曲线在几何上受很大的限制：这样表示的曲线与平行于$y$轴的任意直线相交不能多于一点．通常，把曲线分成可以表为$y=f(x)$的若干部分来克服这个限制，例如，以原点为中心，$r$为半径的圆，可以由定义在闭区间$[-r,r]$上的两个函数$y=\sqrt{r^2-x^2}$和$y=-\sqrt{r^2-x^2}$给出．

曲线最直接和最灵活的描述法是参数表示．我们不把直角坐标$x$或$y$中的一个看成另一个的函数，而把两个坐标$x$和$y$都看成为第三个自变量$t$的函数，$t$叫做\textbf{参数}或\textbf{参变量}．当$t$在$t$轴上或$t$轴上的一个区间$[a,b]$变化时，由$t$轴上的原象点，通过参数方程$x=\varphi(t)$, $y=\psi(t)$映射到曲线$C$上的一点$(x,y)=(\varphi(t),\psi(t))$, 使得对应于$t$的某个区间$[a,b]$或$(-\infty, \infty)$, 一对函数值$x=\varphi(t)$, $y=\psi(t)$定义的点集
$$\Big\{\big(\varphi(t), \psi(t)\big)\Big|t\in [a,b]\Big\}$$
就是曲线$C$上的全部点，而且没有其它点．

在解析几何中，我们已经知道：

圆的参数方程是
\[x=a\cos t,\qquad y=a\sin t\]
这里$t$表示圆心角．

椭圆的参数方程是
\[x=a\cos t,\qquad y=b\sin t\quad (a>b>0)\]
这里$t$是偏心角．（图2.13）

\begin{figure}[htp]
    \centering
\begin{tikzpicture}[>=latex]
    \draw[<->|](0,-.25)--node[fill=white]{$a$}(2,-.25);
    \draw[<->|](.25,0)--node[fill=white]{$b$}(.25,1.2);
\draw[->](-2.8,0)--(2.8,0)node[right]{$x$};
\draw[->](0,-2.8)--(0,2.8)node[right]{$y$};
\draw[thick](0,0) circle (2);
\draw[thick](0,0) circle (1.2);
\draw[very thick](0,0) ellipse [x radius=2, y radius=1.2];
\node at (0,0)[below left]{$O$};
\draw(0,0)--(30:3);
\draw(30:1.2)node[above]{$B$}--(1.04,0)node[above right]{$N$};
\draw(30:2)node[above]{$A$}--(1.73,0)node[above right]{$M$};
\draw(.5,0) arc (0:30:.5)node[right]{$t$};
\draw(30:1.2)--(1.9,.6)node[right]{$P$};

\end{tikzpicture}
    \caption{}
\end{figure}

曲线$C$可以有许多参数表示法，就好象同一条路可以有许多不同走法，例如抛物线$y^2=2px$可以用在其上的点$(x,y)$的切线的斜率作参数，由于
$2yy'=2p$, 所以抛物线$y^2=2px$的参数方程是
\[x=\frac{p}{2(y')^2},\qquad y=\frac{p}{y'}\]
其中$y'\in (-\infty,0)\cup (0, \infty)$.

如果用$\frac{1}{y'}=\cot\theta$作参数，那么它的参数方程是
\[x=\frac{p}{2}\cot^2 \theta,\qquad y=p\cot\theta\]

我们也可以用$\frac{y}{2p}=\frac{x}{y}=\lambda,\quad \lambda\in (-\infty, \infty)$作参数，于
是又得一参数方程
\[x=2p\lambda^2,\qquad y=2p\lambda\]


\begin{example}
讨论由参数方程
\[x=\sin t,\qquad y=\cos^2 t,\quad t\in[0,2\pi]\]
所定义的曲线$C$.
\end{example}

\begin{solution}
因为$\sin^2 t+\cos^2 t=1 $，所以曲线$C$在抛物线
\begin{equation}
    x^2+y=1
\end{equation}
上．反过来，抛物线(2.12)并不与曲线$C$一致，我们要说明限制在长方形区域$R=\{-1\le x\le 1,\; 0\le y\le 1\}$内的抛物线(2.12)上的部分是曲线$C$（图2.14）．

\begin{figure}[htp]
    \centering
\begin{tikzpicture}[>=latex, scale=1.5]
    \draw[->](-2,0)--(2,0)node[right]{$x$};
    \draw[->](0,-1.5)--(0,2)node[right]{$y$};
\node at (0,0)[below left]{$O$};
\draw(-1,0) rectangle (1,1);
\draw[domain=-1.5:1.5, samples=100, very thick, dashed]plot(\x, {1-\x*\x});
\draw[domain=-1:1, samples=100, very thick]plot(\x, {1-\x*\x});
\node at (-1,0)[below left]{$(-1,0)$};
\node at (1,0)[below right]{$(1,0)$};
\node at (0,1)[above right]{$(0,1)$};
\end{tikzpicture}
    \caption{}
\end{figure}

事实上，$x=\sin t$在闭区间$[0,2\pi]$上处处连续，且在闭区间$\left[\frac{\pi}{2},\frac{3\pi}{2}\right]$上由1单调递减到$-1$，因此根据反函数存在定理，对于每一个$x\in[-1,1]$，存在唯一的一个$t$，使得$\sin t=x$成立，从而得到
\[y=1-x^2=1-\sin^2 t=\cos^2 t\in [0,1]\]
这就说明了，点$(x,1-x^2)\; (x\in [-1, 1])$在曲线$C$上，而且只有这些点在曲线$C$上，换言之，在(2.12)上的除去这些点的其它任何一点不在曲线$C$上．
\end{solution}

在参数方程$x=\varphi(t)$, $y=\psi(t)$中，由于$t$轴是有序的，我们可以用明显的方式把曲线$C$的点规定一个次序或“指向”，即如果$t_1<t_2$，那么由$t_1$映射的点在由$t_2$映射的点的前面．例如例2.38的曲线上的动点$P$随着$t$增加的方向这样运动：$P$点在$t=0$时，由点$(0,1)$开始移动到在$t=\frac{\pi}{2}$时的$(1,0)$点，然后沿原路线返回移动到$t=\pi$时的$(0,1)$点，再继续移动到在$t=\frac{3\pi}{2}$时的$(-1, 0)$点，最后又返回到在$t=2\pi$时的$(0, 1)$点．

如果函数$\varphi(t)$, $\psi(t)$的定义域是$t$轴上的闭区间$[a,b]$, 而且在这区间上的$t$的不同值对应曲线$C$上不同的点，
那么我们称$C$为\textbf{简单弧}，例2.38中的抛物线，由上面知道不是简单弧，而抛物线$x=t$, $y=t^2,\; t\in [0, 1]$是简单弧的一个例子．

如果曲线$C$上的始点$P(\varphi(a),\psi(a))$与终点$P(\varphi(b),\psi(b))$相同，并且除$t=a$和$t=b$外，在$[a,b]$上没有另外一对不同值对应曲线上同一点，那么曲线$C$叫做\textbf{简单闭曲线}．$x=a\cos t$, $y=a\sin t,\; t\in [-\pi,\pi]$就是简单的闭曲线．

有时我们使用$y=f(x)$表示曲线$C$或$C$的一部分是方便的，如果$C$用参数方程$x=\varphi(t)$, $y=\psi(t)$给出，而对于曲线的一部分$t_2\le t\le t_1$, 函数$\varphi(t),\psi(t)$中的一个，譬如说$x=\varphi(t)$是单调的和连续的，那么这样的表示法总是可能的，办法是先由$x=\varphi(t)$解出它的反函数$t=\varphi^{-1}(x)$, 再代入$y=\psi(t)$中，便得到
\[y=\psi\left(\varphi^{-1}(x)\right)=f(x)\]
即$y=f(x)$是由$y=\psi(t)$和$t=\varphi^{-1}(x)$复合而成的复合函数．

要求曲线$C$的斜率，我们不必重新建立$y$对$x$的函数关系$y=f(x)$, 再去求导，只须由方程组
\begin{equation}
    \begin{cases}
        x=\varphi(t)\\ y=\psi(t)
    \end{cases}
\end{equation}
当$\frac{\dd\varphi}{\dd t}\ne 0$时，按照复合函数求导法则可以直接得到$f'(x)$.

事实上，当$\varphi'(t)\ne 0$时，$t=\varphi^{-1}(x)$的导数存在，并且是$\frac{\dd t}{\dd 
x}=\frac{1}{\varphi'(t)}$.

按照复合函数求导法则，得到
\[f'(x)=\frac{\dd y}{\dd t}\cdot \frac{\dd t}{\dd x}=\frac{\psi'(t)}{\varphi'(t)}\]
于是，当$\frac{\dd\varphi}{\dd t}\ne 0$时，就得到曲线$C$的斜率：
\begin{equation}
    \frac{\dd y}{\dd x}=\frac{\psi'(t)}{\varphi'(t)}
\end{equation}
当$\frac{\dd\psi}{\dd t}\ne 0$时，就得到
\begin{equation}
    \frac{\dd x}{\dd y}=\frac{\varphi'(t)}{\psi'(t)}
\end{equation}
因此，只要$[\varphi'(t)]^2+[\psi'(t)]^2\ne 0$, 切线总是存在的：如果$\psi'(t)=0$, 切线是水平的；如果$\varphi'(t)=0$, 切线是垂直的．

当曲线由参数方程组(2.13)给出，那么曲线上点$(x_0,y_0)$处的切线方程为
\begin{equation}
    (y-y_0)\frac{\dd x}{\dd t}-(x-x_0)\frac{\dd y}{\dd t}=0
\end{equation}

法线的斜率是$-\frac{1}{y'_x}$，因此法线方程为
\begin{equation}
    (x-x_0)\frac{\dd x}{\dd t}+(y-y_0)\frac{\dd y}{\dd t}=0
\end{equation}

\begin{example}
求椭圆$\begin{cases}
    x=a\cos t\\ y=b\sin t
\end{cases}$在点$t$的切线方程和法线方程．
\end{example}



\begin{solution}
椭圆在点t的切线方程，按照(2.16)得到
\[(y-b\sin t) (-a\sin t )-(x-a\cos t)(b\cos t)=0\]
    化简得
 \[   y ( - a\sin t) - x(b\cos t)+ab=0\]
    两边除以$-ab$, 得到在点$t$的切线方程
\[\frac{x}{a}\cos t+\frac{y}{b}\sin t=1\]

在点$t$的法线方程，按照(2.17)得到
\[(x-a\cos t) (-a\sin t)+(y-b\sin t)(b\cos t)=0\]
    化简得
\[   -ax\sin t+by\cos t+ (a^2-b^2) \sin t\cdot \cos t=0\]
    或
    \[xa\sin t-yb\cos t-\frac{c^2}{2}\sin 2t=0\]
\end{solution}    

\begin{example}
    求曲线$\begin{cases}
        x=t^2\\y=t^3
    \end{cases}$
    在点$t$的切线方程，并求此切线与曲线的另一交点的坐标.
\end{example}

\begin{solution}
在点$t$的切线方程是
\[(y-t^3)(2t)-(x-t^2)(3t^2)=0\]

若$t\ne 0$, 则得切线方程
  \[  2y-3tx+t^3=0\]
    设此切线与曲线交于$(T^2,T^3)$点，于是数$T$满足等式
   \[ 2T^3-3tT^2+t^3=0\]
    这是$T$的三次方程，它表示切线与曲线相交于三个点，由于切点的$t$值是它的二重根，根据根和系数的关系，得
\[2t+T=-\frac{-3t}{2}\quad \Rightarrow\quad T=-\frac{1}{2}t\]
故切线与曲线的另一交点的坐标是$\left(\frac{1}{4}t^2,-\frac{1}{8}t^3\right)$

如果再假定$\varphi''(t)$，$\psi''(t)$存在，则$y$对$x$的二阶导数可按照复合函数求导法及商的导数法来求：
\begin{align*}
    y''=\frac{\dd y'}{\dd x}&=\frac{\dd y'}{\dd t}\cdot \frac{\dd t}{\dd x}\\
    &=\frac{\dd }{\dd t}\left(\frac{\psi'(t)}{\varphi'(t)}\right)\cdot \frac{1}{\varphi'(t)}\\
    &=\frac{\varphi'(t)\psi''(t)-\psi'(t)\varphi''(t)}{\big\{\varphi'(t)\big\}^3}
\end{align*}
\end{solution}

\begin{ex}
\begin{enumerate}
    \item 设$x=\frac{1}{2}p\lambda^2$, $y=p\lambda$是抛物线以$\lambda$为参数的方程，求在
    点$\lambda$的切线方程和法线方程．
    \item 过点$(x_0,y_0)$作抛物线$y^2=2px$的两条切线，设切点的参数为$\lambda_1$和$\lambda_2$, 求证
    \[\begin{cases}
        p(\lambda_1+\lambda_2)=2y_0\\
        p\lambda_1\lambda_2=2x_0
    \end{cases}\]
    \item 过抛物线$y^2=2px$的定长为$c$的弦的两端作切线，求切线交点的轨迹方程．
    
    \item 求曲线
    $\begin{cases}
        x=t (1-t^2) \\ y=1+t^2
    \end{cases}$在点$t$的切线方程．

\item 曲线$y^2=x^3$上点$P(t^2,t^3)$的切线交$x$轴于$T$点，过$T$
点与$PT$垂直的直线交$y$轴于$N$点，求证$\frac{ON^2}{OT}$是一个与
$t$无关的常数，这里$O$是原点．
\item 如果曲线$ay^2=x^3$上$P(at^2,at^3)$点的切线交曲线于另一点$Q$, 交$y$轴于$R$点，
\begin{enumerate}
    \item 求$Q$点坐标；
    \item 如果
由点$P$引$x$轴的垂线的垂足为$N$点，求证：$OQ$和$RN$的斜率的绝对值相等．
\end{enumerate}
\item \begin{enumerate}
    \item 求以$a$为参数的抛物线系
\[y+\frac{35}{16}a=\frac{a^3}{3}\left(x+\frac{3}{4a}\right)^2\]
的顶点的轨迹方程．
\item 证明对应于参数$a$的抛物线与双曲线$xy=-\frac{7}{6}$ 
相切于点$\left(\frac{1}{a},-\frac{7a}{6}\right)$，又与双曲线$xy=\frac{10}{3}$, 相切
于点$\left(-\frac{2}{a},-\frac{5}{3}a\right)$.
\end{enumerate} 

\item 试证曲线
$x=a(\cos t+t\sin t)$, $y=a(\sin t-t\cos t)$的每
条法线都与圆$x^2+y^2=a^2$相切

\item 若$x=\sec\theta$ 和$y=2\tan\theta$,
\begin{enumerate}
    \item 求$\frac{\dd y}{\dd x}$和$\frac{\dd^2 y}{\dd x^2}$以$\theta$表示之；
\item 求证：$y^3\frac{\dd^2 y}{\dd x^2}+16=0$
\end{enumerate}

\item 求证过椭圆$x=a\cos\varphi$, $y=b\sin\varphi$上两点$\varphi_1$、$\varphi_2$的切线
的交点的坐标为
$\left(a\frac{\cos\frac{\varphi_1+\varphi_2}{2}}{\cos\frac{\varphi_1-\varphi_2}{2}}, b\frac{\sin\frac{\varphi_1+\varphi_2}{2}}{\cos\frac{\varphi_1-\varphi_2}{2}}\right)$
\end{enumerate}
\end{ex}

\subsection{微分与变率（微商）}

\subsubsection{微分的概念}

当函数$y=f(x)$的自变量$x$在点$x_0$处得到一个微小的改变量$\Delta x=x-x_0$时，函数$y$的值得到相应的改变量，即
\begin{equation}
    \Delta y=f (x_0+\Delta x) -f (x_0) 
\end{equation} 
函数的改变量$\Delta y$最能表现函数的特征，无论函数变化的大小，或者是函数在某一点的递增或递减的变化趋势，都是由它来刻画的．

从数的观点来看，一个变量就是一个可以取很多不同的值的符号，所以我们也可以把$\Delta x$看成一个新变量，这样$\Delta y$就是$x_0$和$\Delta x$这两个变量的函数．如

\begin{example}
    设$y=x^3$，那么
\[\Delta y=(x_0+\Delta x)^3-x^3_0 =3x_0^2\Delta x+3x_0(\Delta x)^2+(\Delta x)^3\]
\end{example}

\begin{example}
    设$f(x)=x^2-\frac{1}{x}$，那么
\begin{align*}
    \Delta y&=\left[(x+\Delta x)^2-\frac{1}{x+\Delta x}\right]-\left[x^2-\frac{1}{x}\right]\\
    &=2x\Delta x+(\Delta x)^2+\frac{\Delta x}{x(x+\Delta x)}
\end{align*} 
\end{example}

从上面的例子看出，通常$\Delta y=f(x_0+\Delta x)-f(x_0)$是一个很复杂的关系，但是如果函数$f(x)$在点$x_0$存在变率$f'(x_0)$, 我们可以把$\Delta y$分解成一个线性部分和一个比$\Delta x$变小得更快的微量部分．事实上
\begin{equation}
 f'(x_0)=\lim_{\Delta x\to 0}\frac{f(x_0+\Delta x)-f(x_0)}{\Delta x}
\end{equation}
我们用$\varepsilon(x_0,\Delta x)$表示$\frac{\Delta 
y}{\Delta x}$与$f'(x_0)$的差，那么，我们可以写出
\begin{equation}
\frac{f(x_0+\Delta x)-f(x_0)}{\Delta x}=f'(x_0)+\varepsilon(x_0,\Delta x)
\end{equation}
这里因为(2.19)，而有$\Lim_{\Delta x\to 0}\Delta x \cdot \varepsilon(x_0,\Delta x)=0$.

将(2.20)的两边乘以$\Delta x$，我们得到
\begin{equation}
    f(x_0+\Delta x)-f(x_0)=f'(x_0)\Delta x+\Delta x\cdot \varepsilon(x_0,\Delta x)
\end{equation}
这个等式右边的第一项是与$\Delta x$成比例的量，其比例常数是$f(x)$在点$x_0$的变率$f'(x_0)$, 而第二项$\varepsilon(x_0,\Delta x)\Delta x$和$\Delta x$比较起来，当$|\Delta x|$很小时，$|\varepsilon(x_0,\Delta x)\cdot \Delta x|$就更小，从极限的观点来看，就是
\[\lim_{\Delta x\to 0} \frac{\varepsilon(x_0,\Delta x)\cdot \Delta x}{\Delta x} =\lim_{\Delta x\to 0}\varepsilon(x_0,\Delta x)   =0\]
这也就是说，$\varepsilon(x_0,\Delta x)\cdot \Delta x$比$\Delta x$变小得更快．因此，若$f'(x_0)\ne 0$, (2.21)中的线性部分$f'(x_0)\Delta x$, 当$|\Delta x|$很小时就是函数改变量$f(x_0+\Delta x)-f(x_0)$的很好的近似值，它的误差等于$\varepsilon(x_0,\Delta x)\cdot \Delta x$, 由于这个与$\Delta x$成比例的部分是$\Delta y$的\textbf{主要部分}，我们把$f'(x_0)\Delta x$叫做函数$f(x)$在$x=x_0$处\textbf{微分}．

\begin{blk}
    {定义} 设函数$y=f(x)$是可微的，我们称$f'(x)\Delta x$为函数$y=f(x)$在点$x$处与自变量的改变量$\Delta x$相对应的微分，记作$\dd y$, 于是
\begin{equation}
   \dd y=f' (x) \Delta x 
\end{equation}
\end{blk}

因为函数$f(x)=x$的导数等于1, 所以对于任何$x$来说，微分$\dd f(x)$便等于$\dd x=(x')\Delta x=\Delta x$. 因此，习惯上，常以$\dd x$代替$\Delta x$, 于是(2.22)就可以写成
\begin{equation}
    \dd y=f'(x)\dd x
\end{equation}
从而
\[f' (x) =\frac{\dd y}{\dd x}\]
即变率等于微分的商，这也就是称变率（导数）为微商的根据．

下面我们用前面的例2.41和例2.42将$\dd y$和$\Delta y$的值作个比较．


\begin{example}
\[y=x^3,\qquad \dd y=3x^2\dd x,\qquad \Delta y=3x^2\dd x+(3x+\dd x)(\dd x)^2\]

当$x=1$时，$\dd y=3\dd x$, $\Delta y=3\dd x+(3+\dd x)(\dd x)^2$.

\begin{center}
\begin{tabular}{cccccc}
    \hline
$\dd x$ & 10& 1& 0.1 & 0.01& $\cdots$\\
    \hline
$\dd y$& 30&3&0.3& 0.03&$\cdots$\\
\hline
$\Delta y$& 1330& 7&0.331&0.030301&$\cdots$\\
    \hline
\end{tabular}
\end{center}
我们注意到$|\dd x|$愈小，$\dd y$与$\Delta y$的值愈靠近，而且 $\Delta y$与$\dd y$的差的绝对值比$\dd x$的绝对值变小得更快．
\end{example}

\begin{example}
\[f(x)=x^2-\frac{1}{x},\qquad \dd f(x)=f'(x)\dd x=\left(2x+\frac{1}{x^2}\right)\dd x\]
\[\Delta y=2x\dd x+(\dd x)^2+\frac{\dd x}{x(x+\dd x)}\]
令$x=2$，则$\dd y=4.25\dd x$, $\Delta y=\left[4+\dd x+\frac{1}{2(2+\dd x)}\right]\dd x$
\begin{center}
    \begin{tabular}{cccccc}
        \hline
    $\dd x$ & 1& 0.5& 0.1 & 0.01& $\cdots$\\
        \hline
    $\dd y$& 4.25&2.13&0.425& 0.0425&$\cdots$\\
    \hline
    $\Delta y$& 5.17& 2.35&0.434&0.0426&$\cdots$\\
        \hline
    \end{tabular}
    \end{center}
上面的例子验证了我们在前面讨论的结论：当自变量的改变量$\Delta x$很小时，可以用微分$\dd y$来近似地代替函数改变量
$\Delta y$.
\end{example}

\begin{example}
    半径为$r$的金属圆片加热膨胀，求当半径从$r$变到$r+\Delta r$时，圆面积$S$变化大小和此时圆面积$S$与$\Delta r$相应的微分$\dd S$, 并试用几何图形解释所得的结果．
\end{example}

\begin{solution}
    已知圆面积$S=\pi r^2$, 设$\Delta S$是半径从$r$变到$r+\Delta r$时圆面积改变量，即圆环的面积（图2.15），则
\begin{align*}
    \Delta S=\pi(r+\Delta r)^2-\pi r^2&=\pi\left[2r\Delta r+(\Delta r)^2\right]\\
    &=\frac{1}{2}\left[2\pi r+2\pi(r+\Delta r)\right]\Delta r
\end{align*}
此时，$\dd S=2\pi r\Delta r$.

\begin{figure}[htp]\centering
	\begin{minipage}[t]{0.48\textwidth}
		\centering
        \begin{tikzpicture}[>=latex, scale=.8]
\draw[thick](0,0) circle(2);
\draw[thick](0,0) circle(2.8);
\draw(0,0)--(2.8,0);
\draw[|<->|](0,-.2)--node[below]{$r$}(2,-.2);
\draw[|<->|](2.8,.2)--node[above]{$\Delta r$}(2,.2);
    \end{tikzpicture}
    \caption{}
\end{minipage}
\begin{minipage}[t]{0.48\textwidth}
    \centering
\begin{tikzpicture}[>=latex, scale=.8]
\tkzDefPoints{-1.5/0/A, 1.5/0/B, 2.7/1.2/C, -2.7/1.2/D}
\tkzDrawPolygon(A,B,C,D)
\tkzDefPoints{1.5/1.2/E, -1.5/1.2/F}
\tkzDrawPolygon[pattern=north east lines](A,D,F)
\tkzDrawPolygon[pattern=north west lines](B,C,E)
\node at (0,1.2)[above]{$2\pi(r+\Delta r)$};
\node at (0,0)[below]{$2\pi r$};
\draw[|<->|](3,0)--node[fill=white]{$\Delta r$}(3,1.2);
\end{tikzpicture}
\caption{}
\end{minipage}
\end{figure}

把圆沿半径剖开，把圆周拉直，圆环就“变成”一个梯形，它的两个底边长分别是$2\pi r$, $2\pi(r+\Delta r)$, 高为$\Delta r$（图2.16），其中的矩形面积就是微分$\dd S=2\pi r\cdot \Delta r$, 它表示圆面积$S$在从$r$到$r+\Delta r$的整段变化过程中都以同一瞬时变率$(\pi r^2)'=2\pi r$作匀速变化所产生的面积增量．图2.16中的两个小直角三角形的面积表示$\Delta S$与$\dd S$的误差，即
\[\Delta S-\dd S=\pi  (\Delta r)^2\]
它比$|\Delta r|$更快地变小，即
\[\lim_{\Delta r\to 0}\frac{\Delta S-\dd S}{\Delta r}=\lim_{\Delta r\to 0}\pi(\Delta r)=0\]
\end{solution}

\subsubsection{微分的几何意义}

从几何的观点来说，变率的根本想法就是以直代曲，即以在点$x_0$处的切线$y=f(x_0)＋f'(x_0)(x-x_0)$来逼近曲线$y=f(x)$的微段．因此，$f'(x_0)\Delta x$的几何意义可以从图2.17看出．

\begin{figure}[htp]
    \centering
    \begin{tikzpicture}[>=latex, xscale=3]
\draw[->](-.25,0)--(2,0)node[right]{$x$};
\draw[->](0,-1)--(0,5)node[right]{$y$};
\node at (0,0)[below left]{$O$};
\draw[domain=.1:1.6, samples=100, very thick]plot(\x, {exp(\x)});
\draw(.5,0)node[below]{$M(x_0)$}--(.5,1.65)node[above]{$P$};
\draw(1.4,0)node[below]{$M'(x_0+\Delta x)$}--(1.4,4.055)node[above]{$P'$};
\draw(.1,1.65)--(1.8,1.65);
\draw[thick](1.4,4.055)--(.5,1.65);
\draw[domain=.1:1.8, samples=10, thick]plot(\x, {1.65*\x+0.825});
\draw[<->](.5,1.3)--node[fill=white]{$\Delta x$}(1.4,1.3);
\draw[<->|](1.45,1.65)node[below]{$N$}--node[right]{$f'(x_0)\Delta x$}(1.45,3.135)node[above]{$T$};
\tkzDefPoints{.5/1.65/P, 1.4/1.65/N, 1.4/3.135/P'}
%\tkzMarkAngle[mark=none, size=.3](N,P,P')
\tkzLabelAngle[pos=.2](N,P,P'){$\alpha$}
    \end{tikzpicture}
    \caption{}
\end{figure}

在直角三角形$TPN$中，$PN=\Delta x$, $\tan\alpha=f'(x_0)$, 故
\[NT=PN\cdot \tan\alpha=\Delta x\cdot f'(x_0)=f'(x_0)\Delta x\]
也就
是切线函数$y=f (x_0) +f' (x_0) (x-x_0)$的相应改变量．至于函数的改变量则是曲线$y=f(x)$本身
的纵坐标在这一段的改变量，在图2.17中用线段$NP'$来代表．所以$f(x)$相应于这一段的微分就是以曲线的切线改变量代替函数的改变量．从图2.17可以看出，当$|\Delta x|$减小时，$TP'$跟着减小，但$TP'$比$|\Delta x|$要变小得更快．

回顾上面的讨论，微分的意义是，当$|\Delta x|$充分小时，曲线$y=f(x)$在点$x_0$的切线
\[y=f (x_0) +f' (x_0) (x-x_0)\]
的改变量$\dd y=f'(x_0)\Delta x=f'(x_0)(x-x_0)$, 就是相应于函数改变量$\Delta y=f(x_0+\Delta x)-f(x_0)$的微分，它可以近似地表达函数改变量，而差是一个比$|\Delta x|=|x-x_0|$变小得更快的微量．因此我们得到
\[\Delta y\approx \dd y\qquad \text{或}\qquad f(x_0+\Delta x)-f(x_0)\approx f'(x_0)\Delta x\]
上式可改写为
\[f (x_0+\Delta x) \approx f (x_0) +f' (x_0) \Delta x\]
可见$f(x)$在$x_0+\Delta x$点的函数值，可以通过它在$x_0$的函数值$f(x_0)$加上一个与$\Delta x$成正比例的量$f'(x_0)\Delta x$来表达，比例系数就是$f$在点$x_0$的微商．上面的公式常用来作近似计算．

\begin{example}
    计算$\sqrt{101}$
\end{example}


\begin{solution}
设$y=\sqrt{x}$，令$x_0=100$, $\Delta x=1$，则
\[\dd y=\frac{\Delta x}{2\sqrt{x_0}}=\frac{1}{2\sqrt{100}}=\frac{1}{20}=0.05\]
因此：
\[\sqrt{101}\approx \sqrt{100}+\dd y=10+0.05=10.05\]
\end{solution}



\begin{example}
一个立方体形状的盒子，它的设计边长等于4(dm), 但允许误差为0.05(dm), 求盒子体积的可能误差．
\end{example}


\begin{solution}
设盒子的边长由$x$改为$x+\dd x$, 那么盒子体积的改变量是
\[\Delta V= (x+\dd x)^3-x^3\]
但是
\[\Delta V\approx \dd V=3x^2\dd x\]
令$x=4$, $\dd x=\pm 0.05$, 那么
\[\dd V=3\x16\x (\pm 0. 05) =\pm 2.4\]
因此，体积的可能误差近似地等于$\pm 2.4({\rm dm})^3$.
\end{solution}

\begin{example}
    证明当$|x|$很小时，$\ln(1+x)\approx x$.
\end{example}

\begin{solution}
    设$f(x)=\ln(1+x)$, 令 $x_0=0$, $\Delta x=x$，则
\[\dd f (x) =f' (x_0) \Delta x=\frac{1}{1+x_0}\Delta x=x\]
因此：
\[f (x_0+\Delta x) =\ln (1+x) \approx \ln 1+x=x\]
即：$\ln(1+x)\approx x$
\end{solution}


\begin{example}
    当$|x|$很小时，导出近似公式：$\tan x\approx x$.
\end{example}


\begin{solution}
    设$f(x)=\tan x$, 令 $x_0=0$, $\Delta x=x$, 则
\[\dd f (x) =\sec^2x_0 \dd x=\sec^2 0\cdot x=1\cdot x=x\]
因此：
\[f (x_0+\Delta x) =\tan x\approx \tan 0+x=x\]
即：$\tan x\approx x$.
\end{solution}

\begin{example}
如图2.18, 加工锥形工件时，已知工件两头直径分别为$d_1,d_2$, 长度为$\ell$, 当斜角$\alpha$很小时，导出近似公式：
\[\alpha=28.6^{\circ}\x\frac{d_1-d_2}{\ell}\]
\end{example}

\begin{figure}[htp]
    \centering
\begin{tikzpicture}[>=latex]
    \draw[|<->|](-0.5,.9)--node[fill=white]{$d_2$}(-.5,-.9);
\draw[|<->|](4.8,1.5)--node[fill=white]{$d_1$}(4.8,-1.5);
\draw[very thick](4,0) ellipse [x radius=.5, y radius=1.5];
\draw[dashed](0,0) ellipse [x radius=.3, y radius=.9];
\draw[dashed](0,-.9) rectangle (4,.9);
\draw[very thick](0,-.9)--(4,-1.5);\draw[very thick](0,.9)--(4,1.5);
\draw[dashdotted](0,0)--(4,0);
\draw[dashed](4,-1.5)--(4,1.5);
\draw[|<->|](0,1.8)--node[fill=white]{$\ell$}(4,1.8);
\draw[very thick] (0,.9) arc [start angle =90, end angle =270, x radius=.3, y radius=.9] ;
\draw (1.5,.9) arc (0:9:1.5)node[below right]{$\alpha$};
\end{tikzpicture}  
    \caption{}
\end{figure}

\begin{solution}
由图：$\tan\alpha=\frac{d_1-d_2}{2\ell}$. 当$\alpha$很小时，根据$\tan\alpha\approx \alpha$，得到：
\[\alpha\approx \frac{d_1-d_2}{2\ell}\]
又：$1\text{弧度}=\left(\frac{180}{\pi}\right)^{\circ}\approx 57.3^{\circ}$，因此：
\[\alpha\approx 57.3^{\circ}\x \frac{d_1-d_2}{2\ell}\approx 28.6^{\circ}\x \frac{d_1-d_2}{\ell}\]
\end{solution}

\subsubsection{微分的运算}

如果函数$f(x)$可微，即导数$f'(x)$存在，求函数的微分只要求出导数再乘以自变量的改变量$\Delta x$就行了，求已知函数的微分或求导数的方法都叫做微分法．

我们已经知道当$x$是自变量时，$\dd x=\Delta x$, 原来的微分表达式$\dd y=f'(x)\Delta x$就可以改写为$\dd y=f'(x)\dd x$, 例如$\dd(\sin x)=\cos x \dd x$. 所以我们可以根据函数的和、积、商的求导数法则得到函数的和、积、商的求微分法则，即
\begin{enumerate}
    \item $\dd(cu)=c\dd u$
    \item $\dd(u\pm v)=\dd u\pm \dd v$
    \item $\dd(uv)=v\dd u+u\dd v$
    \item $\dd\left(\frac{u}{v}\right)=\frac{v\dd u-u\dd v}{v^2}$
\end{enumerate}
这里只证明最后一个法则，其余请读者自证．

\begin{proof}
\begin{align*}
\dd\left(\frac{u}{v}\right)=\left(\frac{u}{v}\right)'\dd x&=\frac{u'v-uv'}{v^2}\dd x\\
&=\frac{v(u'\dd x)-u(v'\dd x)}{v^2}\\
&=\frac{v\dd u-u\dd v}{v^2}
\end{align*}
\end{proof}


现在我们来研究求复合函数的微分法则．

如果$y=f(x)$, $x=\varphi (t)$, 那么$y$是$t$的复合函数$y=f(\varphi (t))$, 这时$x$和$y$的微分根据定义都要以$t$为自变量来确定．在这情况下，须注意$\dd x\ne \Delta x$, 比如若$x=t^2$, 则$\dd x=2t\dd t$, 而$\Delta x=(t+\Delta t)^2-t^2=2t\Delta t+(\Delta t)^2$, 显然$\dd x\ne \Delta x$, 对于复合函数$y=f(\varphi (t))$的微分，我们有下面重要定理．

\begin{blk}
    {定理} 设$y=f(x)$, $x=\varphi (t)$都是可微函数，$x$是中间变量，那么关系式
$\dd y=f' (x) \dd x$
仍成立．
\end{blk}

\begin{proof}
   \[ \dd y=[f(\varphi (t))]'\dd t\]
根据复合函数求导数法则有
\[[f (\varphi  (t) ) ] '=f' (\varphi  (t) ) \cdot \varphi ' (t)\]
因此
\[\dd y=f' (\varphi  (t) ) \cdot \varphi' (t) \dd t\]

因为$x=\varphi  (t)$, $\dd x=\varphi ' (t) \dd t$，最后得到
\[\dd y=f' (x) \dd x\]
\end{proof}

这个定理是说不论$x$是自变量或是中间变量，函数$f(x)$
的微分形式不变，只须注意，若$x$是中间变量，$\dd x$不是任意增量$\Delta x$, 而是$x$的微分，$\dd y$不能写成$f'(x)\Delta x$即$\dd y\ne f'(x)\Delta x$, 这个性质叫做\textbf{微分不变性}．有了微分不变性，我们在作微分运算时，就无需指明对哪一个变量的微分，而在求函数的变率时，总要指明对中间变量求导数还是对自变量求导数．另外，求复合函数的导数法则用微分符号写出，就是
\[\frac{\dd y}{\dd t}=\frac{\dd y}{\dd x}\cdot \frac{\dd x}{\dd t}\]
它就成为简单的代数恒等式了．

\begin{example}
    求由方程$y^3=x^2+xy+y^2$确定的隐函数的导数$\frac{\dd y}{\dd x}$.
\end{example}

\begin{solution}
    对等式两边求微分得到
\[3y^2\dd y=2x\dd x+ (x\dd y+y\dd x) +2y\dd y\]
移项
\[(3y^2-x-2y)\dd y=(2x+y)\dd x\]
最后得到
\[\frac{\dd y}{\dd x}=\frac{2x+y}{3y^2-x-2y}\]
\end{solution}

\begin{example}
    求$\dd\sin [\tan(x^2-x+1)]$.
\end{example}


\begin{solution}
\begin{align*}
    \dd\sin \left[\tan(x^2-x+1)\right]&=\cos\left[\tan(x^2-x+1)\right]\dd \tan(x^2-x+1)\\
    &=\cos\left[\tan(x^2-x+1)\right]\cdot \sec^2 (x^2-x+1)\cdot \dd (x^2-x+1)\\
    &=\cos\left[\tan(x^2-x+1)\right]\sec^2(x^2-x+1)\cdot (2x-1)\dd x 
\end{align*}
\end{solution}
（例2.52的计算中，将微分符号“$\dd$”逐步顺着复合函数
的层次从$\sin$到$\tan$到$x^2-x+1$, 由外到里，犹如剥笋．）


\begin{ex}
\begin{enumerate}
    \item 边长4cm的正方形铁皮，在加热中边长增加了0.001cm, 
求此时刻面积的可能误差$\Delta S$, 和可能误差的近似值$\dd s$, 又正方形的面积近似值是多少？
\item 半径为10cm的球，在冷却中$R$缩短了0.001cm, 求此时
刻体积的微分$\dd V$和体积的近似值．
\item 当$|x|$很小时，导出下列近似公式：
\begin{multicols}{2}
\begin{enumerate}
    \item  $\sqrt{1+x} \approx 1+\frac{x}{2}$
\item $e^{x} \approx 1+x$
\item  $(1+x)^{n} \approx 1+n x$
\item  $\sqrt[n]{1+x} \approx 1+\frac{1}{n} x$
\item  $\sin x \approx x$
\item  $\arcsin x \approx x$
\item  $\ln (1+\sin x) \approx x$
\end{enumerate}
\end{multicols}

\item  计算下列各式的值：
\begin{multicols}{2}
\begin{enumerate}
\item  $\sqrt[3]{1.02}$
\item  $\sqrt[3]{0.998}$
\item   $e^{-0.1}$
\item   $\tan 0.01$ 
\item   $\sin 0.1^{\circ}$
\item   $\ln 1.0021$   
\end{enumerate}
\end{multicols}

\item 求下列函数的微分:
\begin{multicols}{2}
\begin{enumerate}
    \item $y=\frac{x+5}{x^{2}-2}$
    \item  $y=\left(\sqrt{1+x^{2}}\right)^{n}$
    \item  $y=\ln (\sin \sqrt{x})$
    \item  $y=\sin ^{3}(\sqrt{x})$
    \item  $y=\arcsin \sqrt{2 x}$
    \item  $y=\arctan (\cos x)$ 
    \item $y=\sin (\ln x)$
    \item $y=x^{n} \ln (n x)$
    \item $y=e^{x} \ln x$
    \item $y=\left(e^{x}+e^{-x}\right)^{2}$
\end{enumerate}
\end{multicols}

\item 一直圆锥内接于半径等于$R$的球中，球心在直圆锥的内部
且到圆锥底面的距离等于$p$, 若$p$增加一个微量$x$而球半径$R$不变，求证圆锥体积的增量约等于$\frac{\pi}{3}(R+p)(R-3p) x$.
\item 直圆锥的底面半径为$r$ cm, 高为$h$ cm, 若半径$r$和高$h$分别
有一个增量$\rho$ cm和$\lambda$ cm, 求证圆锥体积的增量近似地等于$\left(\frac{2\rho}{r}+\frac{\lambda}{h}\right)$.
\end{enumerate}
\end{ex}

\subsection*{习题2.2}
\begin{enumerate}
    \item 求下列函数的导函数：
	\begin{enumerate}
    	\begin{multicols}{2}
			\item $f(x)=2 x^{3}-3 x^{3}+4 x-5 $
			\item $f(x)=x^{5}-6 x^{2}+13$
			\item $f(x)=\frac{x^{3}-4 x}{x^{2}+1} $
			\item $f(x)=\cos m x \sin n x$       
    	\end{multicols}

		\item $f(x)=\left[\sin x+\cos x+\left(x^{2}+x+1\right)\right]^{2}$
		\item $f(x)=\tan\left[\sin x+\cos x+\left(x^{2}+x+1\right)\right]$
		\begin{multicols}{2}
			\item $f(x)=\cot \left(x^{2}+1\right)$
			\item $f(x)=\sqrt{a^{2}+b^{2} x^{2}}$
			\item $f(x)=\sqrt[3]{x^{4}+x^{2}+1}$
			\item $f(x)=\cos \left(\sqrt[3]{3 x^{4}+5 x^{2}+6}\right)$
			\item $f(x)=(\sqrt{x}+1)\left(\frac{1}{\sqrt{x}}-1\right)$
			\item $f(x)=x \sin x+\cos x$
			\item $f(x)=2^{\sin ^{2}\tfrac{1}{x}} $
			\item $f(x)=x^{\sin x}$       
		\end{multicols}
		\begin{multicols}{2}
		\item $f(x)=e^{-x^{2}} \cos \left(e^{-x^{2}}\right) $
		\item  $f(x)=\arctan\sqrt{x^{2}-2 x}$
		\item  $f(x)=x e^{1-\cos x}$
		\item  $f(x)=x \arccos x-\sqrt{1-x^{2}}$
		\item  $f(x)=(\sin x)^{\cos x}$
		\item  $f(x)=\sqrt[5]{\frac{x-5}{\sqrt[5]{x^{2}+2}}}$
		\end{multicols}
		\item  $f(x)=x-\ln \left(2 e^{x}+1+\sqrt{e^{2 x}+4 e^{x}+1}\right)$
		\begin{multicols}{2}
		\item $f(x)=\sqrt{x \sqrt{x \sqrt{x}}}$
		\item  $f(x)=\log _{a} e ^{\sin ^{2} x+1}$
		\end{multicols}
		\item  $f(x)=\ln \tan  \frac{x}{2}-\cot x \cdot \ln (1+\sin x)-x$
	\end{enumerate}

\item 验证函数$f(x)=e^x \sin x$满足方程
\[f'' (x) -2f' (x) +2f (x) =0.\]

\item 验证函数$f(x)=\frac{x-3}{x+4}$满足方程
\[2 [f' (x) ]^2= [f (x) -1] f'' (x).\]

\item $P$点在直线上运动，在时刻$t$, 点$P$与直线上的定点$O$的距离由方程
$s^2=a^2+V^2t^2$
给出，其中$a$和$V$是常量，试用$s$表示$P$点在时刻$t$的速度和加速度．

\item 在$xy$平面上，求通过点$P(-3, 6)$且与曲线$y=x^3-5x^2+x+5$
相切，并且切点的横坐标满足$x\ge 0$的切线方程．

\item 求证方程$16x^4-24x^2+16x-3=0$在$(0, 1)$内有
一个三重根，又在$(-2,-1)$内有一个单根．

\item 设三次函数$y=x^3-3ax^2+bx+c$ （其中$a,b,c$是
常数）．
\begin{enumerate}
    \item 该曲线与直线$y=2x-1$切于点$(2, 3)$, 试用$a$的式子表示$b,c$;
    \item 另外，该曲线还与$y=2x+3$相切，试确定$a,b,c$的值．
\end{enumerate}
\item 在中午十二点整，甲船以6公里/小时的速率向东行，乙船在甲船之北16公里，以8公里/小时的速率向南行，在下午一点整两船相近之速率为多少？在下午两点整两船相离之速率是多少？
\item 已知圆心为$O$, 半径为1的圆与直线$\ell$相切于点$A$, 一动点$P$自切点$A$沿直线$\ell$向右移动时，取弧$\wideparen{AC}$的长为$\frac{2}{3}AP$，
直线$PC$与直线$AO$交于点$M$(图2.19), 又知当$AP=\frac{3\pi}{4}$时，点$P$的速率为$v$, 求这时$M$点的速率．

\begin{figure}[htp]
    \centering
\begin{tikzpicture}
\draw(-2,0)--(2,0)node[right]{$\ell$};
\draw[very thick](0,1) circle (1);
\draw(0,0)node[below]{$A$}--(0,2.5);
\draw(0,1.7)node[left]{$M$}--(1.7,0)node[below]{$P$};
\node at (30:1)[right]{$C$};
\tkzDefPoints{0/1/O}
\tkzLabelPoints[left](O)
\tkzDrawPoints(O)

\end{tikzpicture}    
    \caption{}
\end{figure}

\item 溶液自深18cm, 顶直
径12cm的正圆锥形漏斗中漏入一直径为10cm的圆柱形筒中，开始时漏斗中盛满了水，已知当溶液在漏斗中深为12cm时，其水平面下落速率为1cm/s, 问此时圆柱形筒中水平面上升之速率为多少？
\item 求曲线$y=\arctan\ln x$与$x$轴的交角．
\item 设质点$P$在时刻$t$的位移方程是
\[x=\frac{58}{240}+A\cos\left(\sqrt{24}gt+\alpha\right)\]
其中$A,\alpha$是待定常数，$g\approx 9.8{\rm m/s}$，且知当$t=0$时，$x=0$, $\frac{\dd x}{\dd t}=0$，求$A$和$\alpha$.


\begin{figure}[htp]
    \centering
\begin{tikzpicture}[>=latex, scale=1.7]
    \draw[dashdotted] (-3.5,0)--(1.2,0);
    \draw (-3.5,.5) rectangle (-1.5,-.5);
    \draw (0,0) circle (.75);
    \node at (-.75,0)[below left]{$K_0$};
    \node at (.75,0)[above right]{$K_1$};
\draw [fill=gray!20](-2.5,.5) rectangle (-2.6,-.5);
\draw[very thick] (0,0)node[below]{$O$}--node[above]{$r$}(90+30:.75)node[above]{$K$}--node[above right]{$\ell$}(-2.5,0)node[below]{$M$};
\draw [dashed](-3.3,.5)node[below]{$M_0$}--(-3.3,-.5);
\draw [dashed](-1.7,.5)node[below]{$M_1$}--(-1.7,-.5);
\draw (-.25,0) arc (180:180-60:.25)node[below left]{$\varphi$};
\draw[dashed] (90+30:.75)--(-.75/2,0);
\draw[|<->|] (-3.3,-.7)--node[fill=white]{$x$}(-2.5,-.7);

\draw (-2,0) arc (0:15:.5)node[below right]{$\alpha$};


\end{tikzpicture}
    \caption{}
\end{figure}



\item 图2.20表示曲柄机构装置的图样，连接杆$MK$
的长$\ell=125$cm, 曲柄$OK$的长$r=25$cm, 连杆与汽缸的轴$OM$成$\alpha$角，曲柄与同一轴成$\varphi$角．
\begin{enumerate}
    \item 求用$\varphi$的式子表示活塞$M$的位移$x$的公式；
    \item 如果$\frac{\dd \varphi}{\dd t}=\omega$ （$\omega$
    是常量），求活塞$M$的速率．
\end{enumerate}
\item 在抛物线$y=ax^2+bx+c\; (a\ne 0)$上两点作切线，
它们在何处相交？
\item 求抛物线$y^2=4ax$在点$(a\lambda^2, 2a\lambda)$的切线方程，并求它的两条互相垂直的切线的交点的轨迹方程．
\item 
\begin{enumerate}
    \item 求曲线$4y^3=27x^2$ 在点$(2t^3, 3t^2)$的切
线方程；
\item 求互相垂直的两条切线的交点的轨迹方程．
\end{enumerate}

\item 计算$\sqrt{\frac{2.037^2-1}{2.037^2+1}}$的近似值．
\item 计算$\cos151^{\circ}$的近似值．
\item 设$f(x)=e^{0.1x(1-x)}$, 试计算$f(1.05)$的近
似值．
\item 若球体积的相对误差$\frac{\Delta V}{V}=0.01$, 试计算球半径的
相对误差$\frac{\Delta R}{R}$约等于多少？
\end{enumerate}


\section{微商运算的初步应用{}---{}函数的增减与极值}

\subsection{局部极值和中值定理}
我们在第一章中曾说明在闭区间$[a,b]$上连续的函数必有最大值和最小值．函数$f$在闭区间$M=[a,b]$上的最大值的大小涉及到两个因素，既依赖于$f$，也依赖于$M$.

如图2.21，若$N=[c,d]\subset M=[a,b]$，则
\[
\max_N f \le \max_M f.
\]
相应地，
\[
\min_N f \ge \min_M f.
\]
在这两种情形里都可能出现严格不等式（即不带等号的不等式）．

\begin{figure}[htp]
    \centering
    \begin{tikzpicture}[>=latex]
        \draw[->](-.5,0)--(5.5,0)node[right]{$x$};
        \draw[->](0,-.5)--(0,4.5)node[right]{$y$};
        \node at (0,0)[below left]{$O$};
        \draw[rounded corners=10pt, very thick](.5,1)--(1.2,2)--(2,1.6)--(2.3,2.2)--(2.9,3.6)--(3.6,3.4)--(4.3,2.5)--(4.8,4);
        \node at (3,4){$y=f(x)$};
        \draw(.5,1)--(.5,0)node[below]{$a$};
        \draw(1.2,1.85)--(1.2,0)node[below]{$c$};
        \draw(4.3,2.65)--(4.3,0)node[below]{$d$};
        \draw(4.8,4)--(4.8,0)node[below]{$b$};
        \end{tikzpicture}
    \caption{}
\end{figure}

现在我们将阐明一种只决定于$f$的性质的局部极大与极小的概念．

\begin{blk}
    {定义} 设函数$f$在点$x$的某个$\delta>0$的邻域，即$(x-\delta,x+\delta)$内有定义，如果对于充分小的$|\Delta x|<\delta$, 能使$x+\Delta x\in (x-\delta,x+\delta)$, 而且对于一切$x'=x+\Delta x\in (x-\delta,x+\delta)$, 都有
\[f(x+\Delta x)\le f(x)\quad \text{或}\quad f(x+\Delta x)\ge f(x)\]
那么函数$f$在点$x$达到局部极大值（或局部极小值）．
\end{blk}

\begin{blk}
    {费马定理} 如果函数$f$在区间$(a,b)$上有定义，并且在这个区间内的一点$c$取局部极大值（或极小值），又若$f$在$c$点可微，那么
    \[f' (c) =0\]
\end{blk}

\begin{proof}
    为确定起见，设$f(x)$在点$c\in (a,b)$取最大值、于是对于一切点$x\in (a,b)$, 就有
\[
f (x) \le f (c). 
\]
由导数定义：
\[\lim_{x\to c}\frac{f (x) -f (c)}{x-c} =f' (c) \]
并且极限$f'(c)$的存在与$x$是从右边还是从左边趋向$c$无关．

\begin{enumerate}
    \item 如果$x>c$，我们有$\frac{f (x) -f (c)}{x-c}\le 0$，从而
\[\lim_{x\to c^+}\frac{f (x) -f (c)}{x-c}\le 0\]
\item 如果$x<c$，我们有$\frac{f (x) -f (c)}{x-c}\ge 0$，从而
\[\lim_{x\to c^-}\frac{f (x) -f (c)}{x-c}\ge 0\]
\end{enumerate}
根据假设$f$在$c$点可微，因此这两个极限必相等，且应等于$f'(c)$, 这表示$f'(c)\ge 0$和$f'(c)\le 0$, 由此得：$f' (c) =0$.

$f$在$c$点有局部极小值的情形留给读者自证．
\end{proof}

\begin{blk}
  {罗尔定理} 如果函数$f$在闭区间$[a,b]$上连续，在开区间$(a,b)$内可微，且在两端点的函数值相等，即$f(a)=f(b)$, 那么至少存在一个点$c$, $c\in (a,b)$, 使$f'(c)=0$.  
\end{blk}

\begin{proof}
如果$f$在$[a,b]$上恒为常数，无疑对于任何
$c\in (a,b)$有$f'(c)=0$. 因此，可设$f$在$[a,b]$中不为常数，此时$f$在$[a,b]$上的最大值$M$和最小值$m$至少有一个不等于$f(a)=f(b)$, 为确定起见，不妨设$M\ne f(a)=f(b)$．故$f$在$[a,b]$上的最大值$M$不会在$a$点或$b$点达到，这时根据闭区间上连续函数的性质，在$(a,b)$内至少有一点$c$, 使$f(c)=M$, 于是，由费马定理，即得$f'(c)=0$. 

本定理的几何意义如图2.22, 曲线$y=f(x)$在$a,b$处有相同高度，如果在点$c$处、曲线达到最高度，并有切线存在，那么，曲线在该处的切线与$x$轴平行．同样，在曲线的最低点处的切线也与$x$轴平行．
\end{proof}

\begin{figure}[htp]
    \centering
    \begin{minipage}[t]{0.48\textwidth}
    \centering
\begin{tikzpicture}[>=latex, scale=1]
    \draw[->](-.5,0)--(5,0)node[right]{$x$};
    \draw[->](0,-.5)--(0,2.8)node[right]{$y$};
    \node at (0,0)[below left]{$O$};
    \draw[domain=.5:4.5, samples=100, thick]plot(\x, {sin((\x-.5)*pi/2 r)+.5});
    \draw[dashed](.5,0)node[below]{$a$}--(.5,.5)--(4.5,.5)--(4.5,0)node[below]{$b$};
    \draw(1.5,0)node[below]{$c$}--(1.5,1.5);
\end{tikzpicture}
    \caption{}
    \end{minipage}
    \begin{minipage}[t]{0.48\textwidth}
    \centering
    \begin{tikzpicture}[>=latex, scale=1.5]
        \draw[->](-1.5,0)--(1.5,0)node[right]{$x$};
        \draw[->](0,-.25)--(0,2)node[right]{$y$};
        \draw[thick](-1,0)node[below]{$(-1,0)$}--(0,1)node[above right]{$(0,1)$}--node[right]{$g(x)=1-|x|$}(1,0)node[below]{$(1,0)$};
        \node at (0,0)[below left]{$O$};
    \end{tikzpicture}
    \caption{}
    \end{minipage}
    \end{figure}


\begin{rmk}
    注意 如果罗尔定理中的条件不被满足，结论就不成立．例如$g(x)=1-|x|$在$[-1, 1]$上连续，且$g(-1)=g(1)=0$, 但在$x=0$处，$g'(x)$不存在，对于这个$g$, 在$(-1,1)$中就没有$c$能使$g'(c)=0$. 它的图象如图2.23所示，
\end{rmk}

\begin{blk}
    {拉格朗日中值定理} 如果函数$f$在闭区间$[a,b]$上连续，在开区间$(a,b)$内可导，那么在$(a,b)$内至少有一点$c$, 使得
\[f'(c)=\frac{f(b)-f(a)}{b-a}\]
\end{blk}

罗尔定理是拉格朗日定理的特殊情形，我们借助于罗尔定理来证明这个重要的定理．

\begin{proof}
 我们作一个辅助函数$\Phi$使它具有函数f的性质：即它在$[a,b]$上连续且在$(a,b)$内可微，同时满足$\Phi (a)=\Phi (b)$. 为此，设$\Phi  (x) =f (x) +k (x-a)$, 并由条件$\Phi (a)=\Phi (b)$来决定参数$k$. 

因为当$x=a$时，有$\Phi (x)=f(a)$; 当$x=b$时，有$\Phi (b)=f(b)+k(b-a)$, 所以
\[k=-\frac{f (b) -f (a)}{b-a}\]
这样得到的函数
\[\Phi  (x) =f (x) -\frac{f (b) -f (a) }{b-a}(x-a)\]
就满足罗尔定理中的所有的条件，由罗尔定理，至少存在一点$c\in (a,b)$使得$\Phi '(c)=0$,
但
\[\Phi ' (x) =f' (x) -\frac{f (b) -f (a)}{b-a}\]
所以
\[\Phi ' (c) =f' (c) -\frac{f (b) -f (a)}{b-a} =0\]
即
\[f' (c) =\frac{f (b) -f (a)}{b-a}\]
\end{proof}

为了便于应用，拉格朗日定理的结论，通常写成如下形式：
\begin{equation}
   f (b) -f (a) =f' (c) (b-a),\qquad a<c<b 
\end{equation}
在上式中，令$b=a+h$，$c=a+\theta h$，$0<\theta<1$，得到中值定理的另一形式：
\begin{equation}
    f (a+h) -f (a) =hf' (a+\theta h) ,\qquad 0<\theta<1
\end{equation}

中值定理的几何意义如图2.24所示．在闭区间$[a,b]$上有一条连续的曲线，且过曲线$y=f(x)$上每一点（端点除外），都可以作出一条切线，那么在曲线上至少有一点$M(c, f( c)),\; c\in (a,b)$, 使得过点$M$的切线与弦$AB$平行．

\begin{figure}[htp]
    \centering
\begin{tikzpicture}[>=latex, scale=1.5]
\draw[->](-.5,1)--(4,1)node[right]{$x$};
\node at (0,1)[below left]{$O$};
\draw[->](0,.5)--(0,3)node[right]{$y$};
\draw[domain=.8:3.2, samples=100, very thick]plot(\x, {(\x-1)*(\x-2)*(\x-3)+2});
\draw(.8,1)node[below]{$a$}--(.8,1.472)node[left]{$A$};
\draw(3.2,1)node[below]{$b$}--(3.2,2.528)node[right]{$B$};
\draw[thick](.8,1.472)--(3.2,2.528);
\draw[dashed](.8,1.472)--(3.2,1.472)node[right]{$C$};
\draw[dashed](1.31,2.3603)node[above left]{$M(c,f(c))$}--(1.31,1)node[below]{$c$};

\draw[domain=.5:2.5, samples=10, thick]plot(\x, {0.44*\x+1.784});

\end{tikzpicture}
    \caption{}
\end{figure}



\begin{example}
求$f(x)=1+\frac{4}{x}$在$(1, 4)$内满足中值定理$f(b)-f(a)=f'(c)\cdot (b-a)$的点$c$的值．
\end{example}


\begin{solution}
    显然$f$在$[1, 4]$上连续，又$f'(x)=-\frac{4}{x^2}$, 故$f$在$(1, 4)$内可微，根据中值定理，存在点$c\in (1, 4)$, 使得
\[    f (4) -f (1) = (4-1) f' (c)\]
成立．

$\because\quad f (4) =2,\quad f (1) =5$,

$\therefore\quad 2-5=3 \left(-\frac{4}{c^2}\right),$
即：$c^2=4$，得$c=\pm 2$.

又$\because\quad 2\in (1, 4),\quad -2\notin (1, 4)$,

$\therefore\quad c=2$. 于是
\[    f (4) -f (1) = (4-1) f' (2) \]
\end{solution}

\begin{example}
    设$f(x)=a_0+a_1x+a_2x^2+\cdots+a_nx^n$, 求证在$f'(x)=0$的相邻二根之间，至多有$f(x)=0$的一个根．   
\end{example}

\begin{proof}
因为多项式函数处处连续和可微，所以若在$f'(x)=0$的两相邻实根$x_1,x_2$之间有$f(x)=0$的二实根$\alpha,\beta\; (\alpha<\beta)$, 那么根据罗尔定理，至少存在$f'(x)=0$的一个根$c$, 使得$\alpha<c<\beta$．但是$(\alpha,\beta)\subset (x_1,x_2)$，于是$x_1<c<x_2$，由此得出矛盾，所以在$f'(x)=0$的相邻两根之间，至多有$f(x)=0$的一个根．
\end{proof}

\begin{example}
    求证不等式$\frac{h}{1+h}<\ln(1+h)<h$，
    对于$h\in (-1, 0)\cup (0,+\infty)$成立．
\end{example}


\begin{proof}
    设$f(x)=\ln x$, 于是$(\ln x)'=\frac{1}{x}$. 由于$\ln x$在区间$[a,b],\; (0<a<b)$上连续且在$(a,b)$内可微，根据中值定理$(a,b)$内存在$c$，使得：
\begin{equation}
    \ln b-\ln a=\frac{1}{c}(b-a),\quad a<c<b
\end{equation}

若$h>0$, 令$a=1$, $b=1+h$, 于是$c=1+\theta h$, 这样(2.26)可写成
\[\ln (1+h) =\frac{h}{1+\theta h},\quad 0<\theta<1\]
若把上式中的$\theta$看作变数，由0变到1, 那么$\ln(1+h)=\frac{h}{1+\theta h}=\frac{1}{\frac{1}{h}+\theta}$
便由$h$递减到$\frac{h}{1+h}$，所以
\[h>\ln(1+h)>\frac{h}{1+h},\quad h>0\]

若$-1<h<0$，令$a=1+h$，$b=1$，于是$c=1+\theta h$，这样
(2.26)可以写成
\[\ln1-\ln (1+h) =\frac{-h}{1+\theta h}\]
仍得
\[\ln (1+h) =\frac{h}{1+\theta h}\]
同样可证
\[h> \ln (1+h) >\frac{h}{1+h}\]
因此，当$h>-1$且$h\ne 0$时，不等式
\[\frac{h}{1+h}<\ln(1+h)<h\]
成立．
\end{proof}



\begin{ex}
\begin{enumerate}
    \item 曲线$y=x^2$上的某点的切线平行于连接横坐标为$a,b$两点的弦，求出此点．
    \item 曲线$y=\frac{1}{x}$上的某点的切线平行于连接横坐标为$a,b\;(0<a<b)$两点的弦，求出此点．
    \item 若$\varphi(x)=(x-a)^m(x-b)^n$，验证罗尔定理成立．
    \item 用罗尔定理证明，无论$m$为何值，多项式函数$f_m(x)=x^3-3x+m$在$[0, 1]$之间至多有一个根．
    \item 证明下列不等式：
\begin{enumerate}
    \item $|\sin x-\sin y|\le |x-y|$
    \item $\frac{a-b}{a}<\ln \frac{a}{b}<\frac{a-b}{b},\quad 0<b<a$
    \item $\frac{h}{1+h^2}<\arctan h<h,\quad h>0$
    \item 若$0<m<1$，且$x>1$，则：$\frac{m(x-1)}{x^{1-m}}<x^m-1<m(x-1)$
    
    （提示：利用中值定理形式1，取$b=x,\; a=1$）
    \item 若$h>-1$，$h\ne 0$，$n>1$，则：$(1+h)^{\tfrac{1}{n}}<1+\frac{h}{n}$
    \item 若$h>-1$，$h\ne 0$，则：$1-\frac{h}{2}<\frac{1}{\sqrt{1+h}}<1-\frac{h}{2\sqrt{(1+h)^3}}$
\end{enumerate}
\item 说明方程$3x^5-25x^3+60x-20=0$有三个实数根，给出根所在的区间，使每个区间只含一个根．
\end{enumerate}
\end{ex}

\subsection{函数的增减与极值——中值定理的应用}

\subsubsection{函数的增减性}

我们已经学过增函数与减函数的概念，现在要阐明导数的正负与函数的增减的关系．

读者已经知道常量的导数等于零，现在反过来问：如果
函数$f$在给定区间上的导数永远是零，那么这个函数在这个区间上就是一个常量？由拉格朗日定理可得到肯定的回答．

\begin{blk}
  {定理1} 如果函数$f$在开区间$(a,b)$上的每一点的导数等于零，即$f'(x)=0$, 那么函数$f$在$(a,b)$上是一个常数．  
\end{blk}

\begin{proof}
  我们从区间$(a,b)$中固定某一点$x_0$, 再取其它的任何一点$x\in (a,b)$. 对于区间$[x_0,x]$或$[x,x_0]$, 函数$f$满足拉格朗日定理条件，因此，我们可以写出
  \[f (x) -f (x_0) =f' (c) (x-x_0),\qquad c\in (x_0,x)\text{ 或 }(x,x_0)\]
  也就是$c\in (a,b)$, 依已知条件$f'(c)=0$, 因而对任何一点$x\in (a,b)$, $f(x)=f(x_0)=$常数．  
\end{proof}



\begin{example}
    求证：如果$x>1$，那么$2\arctan x+\arcsin\frac{2x}{1+x^2}=\pi$成立．
\end{example}

\begin{proof}
设$f(x)=2\arctan x+\arcsin\frac{2x}{1+x^2}$，因为当$x>1$时，$\left|\frac{2x}{1+x^2}\right|<1$，故对于所指定的$x$值的范围，$f(x)$有定义，又因为：
\begin{align*}
    f'(x)&=2(\arctan x)'+\left(\arcsin\frac{2x}{1+x^2}\right)'\\
    &=\frac{2}{1+x^2}+\frac{1}{\sqrt{1-\left(\frac{2x}{1+x^2}\right)^2}}\cdot \left(\frac{2x}{1+x^2}\right)'\\
    &=\frac{2}{1+x^2}+\frac{1}{\sqrt{1-\left(\frac{2x}{1+x^2}\right)^2}}\cdot \frac{(1+x^2)\cdot 2-2x(2x)}{(1+x^2)^2}\\
    &=\frac{2}{1+x^2}+\frac{2(1-x^2)}{(x^2-1)(1+x^2)}\\
    &=2\left(\frac{1}{1+x^2}-\frac{1}{1+x^2}\right)=0
\end{align*}
所以：
\[2\arctan x+\arcsin\frac{2x}{1+x^2}=c\; \text{（常数）}\]
另一方面，因为
\[f\left(\sqrt{3}\right)=2\arctan\sqrt{3}+\arcsin\frac{\sqrt{3}}{2}=\frac{2\pi}{3}+\frac{\pi}{3}=\pi\]
因此：
\[2\arctan x+\arcsin\frac{2x}{1+x^2}=\pi,\quad (x>1)\]
\end{proof}


\begin{blk}
    {定理2} 若$f(x)$在$[a,b]$连续，在$(a,b)$内$f'(x)＞0$ (或$f'(x)<0$)，那么$f(x)$在$[a,b]$上严格递增（或严格递减）．
\end{blk}


\begin{proof}
    设$a\le x_1<x_2\le b$，由于在闭区间$[x_1,x_2]$上满足拉格朗日定理的条件，因此，在开区间$(x_1,x_2)$上可以找到点$c$，使得
\[f(x_2)-f(x_1)=(x_2-x_1)f'(c)\quad (x_1<c<x_2)\]
根据条件，如果在$(a,b)$内，$f'(x)>0$ (或$f'(x)<0$), 那么$f'(c)>0$ (或$f'(c)<0$), 从而
\[f(x_2)-f(x_1)>0\quad (\text{或}<0)\]
因为对于任意$x_1$和$x_2$ (其中$a\le x_1<x_2\le b$)，上面不等式成立，所以$f(x)$在$[a,b]$上严格递增（或严格递减）．
\end{proof}

定理2的几何意义是明显的，$f'(x)$是函数图象上点$x$处的切线的斜率$\tan\theta$, 这里$0\le\theta \le \pi$. 这斜率的符号指出切线是向上或向下倾斜的，而曲线本身也就随着它向上升或向下降．

如果$f'(x)$在个别的$x$值等于零，其他条件不变，则定理
2仍然成立，这一点很容易证实，就是把区间$[a,b]$用使$f'(x)=0$的点来分成若干个小区间，再在每个小区间上分别应用定理2. 图2.25和2.26表示定理2的几何意义．

\begin{figure}[htp]
    \centering
    \begin{minipage}[t]{0.48\textwidth}
    \centering
    \begin{tikzpicture}[>=latex, scale=.7]
  \draw[->](-2,0)--(5,0)node[right]{$x$};
\draw[->](0,-1)--(0,5)node[right]{$y$};
\node at (0,0)[below left]{$O$};
\draw[domain=-.2:4, samples=100, very thick]plot(\x, {.2*(\x-2)^3+2.5});
\draw[->, domain=-1.5:3, thick, samples=10]plot(\x, {1.2*\x+1.22});
\draw[->, domain=.5:5, thick, samples=10]plot(\x, {1.2*\x-1.03});
\draw(.2,2.5)--(3.5,2.5);
\draw[->, thick](.86+.5,0) arc (0:50.2:.5)node[right]{$\theta'$};
\draw[->, thick](-1.02+.5,0) arc (0:50.2:.5)node[right]{$\theta$};
    \end{tikzpicture}
    \caption{}
    \end{minipage}
    \begin{minipage}[t]{0.48\textwidth}
    \centering
    \begin{tikzpicture}[>=latex, scale=.7]
  \draw[->](-1,0)--(7.5,0)node[right]{$x$};
\draw[->](0,-1)--(0,5)node[right]{$y$};
\node at (0,0)[below left]{$O$};
\draw[domain=.7:5, samples=100, very thick]plot(\x, {-.2*(\x-3)^3+2.5});
\draw[->, domain=2.5:6.5, thick, samples=10]plot(\x, {-1.2*\x+7.23});
\draw[->, domain=.5:5, thick, samples=10]plot(\x, {-1.2*\x+5});
\draw(.2,2.5)--(5.5,2.5);
\draw[->, thick](6.025+.5,0) arc (0:129.8:.5)node[above right]{$\theta'$};
\draw[->, thick](4.15+.5,0) arc (0:129.8:.5)node[above right]{$\theta$};


    \end{tikzpicture}
    \caption{}
    \end{minipage}
  \end{figure}



\begin{example}
证明$\tan x>x,\quad 0<x<\frac{\pi}{2}$
\end{example}


\begin{proof}
我们只须证明$\tan x-x>0,\; 0<x<\frac{\pi}{2}$．为此，设$f(x)=\tan x-x,\; x\in\left(0,\frac{\pi}{2}\right)$，因为$f(x)$在$\left[0,\frac{\pi}{2}\right)$
上连续，且在$\left(0,\frac{\pi}{2}\right)$内$f'(x)=\sec^2x-1=\tan^2 x>0$，依定理2, $f(x)$在$\left[0,\frac{\pi}{2}\right)$上递增，所以有
\[f (x) > f (0) =\tan 0-0=0,\quad x\in \left(0,\frac{\pi}{2}\right)\]
即：$\tan x-x>0$，
也就是$\tan x>x,\quad 0<x<\frac{\pi}{2}$.
\end{proof}

\begin{example}
        求证当$x>0$时，有
$e^x> 1+\ln (1+x) $.
\end{example}


\begin{proof}
设$f(x)=e^x-1-\ln(1+x)\; (x>0)$, 则
\[f' (x) =e^x-\frac{1}{1+x},\qquad 
f''(x) =e^x +\frac{1}{(1+x)^2}\] 
    显然，当$x>0$时，$f''(x)>0$, 又$f'(x)$在$[0,\infty)$上连续，由定理2, $f'(x)$在$[0,\infty)$上递增，所以
\[f' (x) > f' (0) =e^0-\frac{1}{1+0}=0\quad  (x> 0) \]
又$f(x)$在$[0, \infty)$连续，同理$f(x)$在$[0, \infty)$上递增，所以当$x>0$时，有
\[f (x) > f (0) =e^0-1-\ln (1+0) =0\]
即：
\[e^x-1-\ln(1+x)>0\quad \text{或}\quad e^x>1+\ln(1+x)\]
\end{proof}


\begin{example}
    确定$f(x)=2x^3-9x^2+12x+5$的递增区间和递减区间，并画出这个函数图象的草图．
\end{example}


\begin{solution}
  \[  f'(x)=6x^2-18x+12=6(x-1)(x-2)\]

  当$x<1$或$x>2$时，$f'(x)>0$, 所以$f(x)$在区间$(-\infty,1)$和$(2,+\infty)$内递增；

  当$1<x<2$时，$f'(x)<0$，所以$f(x)$在区间$(1,2)$内递减．
  
  图象的$y$截距$f(0)=5$，为确定图象的$x$截距，根据例2.54和连续函数中间值定理，由计算的结果
  \[f(-1)=-18,\quad f(1)=10,\quad f(2)=9,\quad f(+\infty)>0\]
  得知$f(x)$在$(1, 2)$和$(2,+\infty)$内没有零点，即图象在$x>1$的范围内与$x$轴不相交：由于$f(x)$在$x<1$的区间内递增，并且$f(-1)\cdot f(0)<0$, 故$f(x)$的图象在
  $x<1$的范围内与$x$轴只有一个交点，它的横坐标在区间$(-1,0)$内．

  把$f(x)$的递增、递减情况列成下表，就更加醒目．

  \begin{center}
\begin{tabular}{c|ccccc}
\hline
$x$ & $(-\infty,1)$& 1& $(1,2)$ & 2& $(2,+\infty)$\\
\hline
$f'(x)$ & $+$ &0&$-$&0&$+$\\
$f(x)$ & $\nearrow$ &10&$\searrow$ &9&$\nearrow$ \\
\hline
\end{tabular}
  \end{center}
图象的草图如图2.27.

\begin{figure}[htp]
    \centering
\begin{tikzpicture}[>=latex, yscale=.3]
\draw[->](-2,0)--(4,0)node[right]{$x$};
\draw[->](0,-2)--(0,15)node[right]{$y$};
\draw[domain=-1:3, samples=300, very thick]plot(\x, {2*\x^3-9*\x^2+12*\x+5});
\draw[dashed](1,0)node[below]{1}--(1,10);
\draw[dashed](2,0)node[below]{2}--(2,9);
\draw(-1,0)node[below]{$-1$}--(-1,.2);
\end{tikzpicture}
    \caption{}
\end{figure}
\end{solution}

\begin{ex}
\begin{enumerate}
    \item     若$|f(x)-f(y)|\le (x-y)^n,\; n>1$，
   求证$f$是一常数．
    \item 若$f'(x)=g'(x)$, $x\in(-\infty,+\infty)$, 求证$f(x)=g (x) +c$.
    \item 证明在任何区间内，函数$x-\sin x$是增函数，又$\alpha$为何
    值，函数$\alpha x-\sin x$是增函数，$\alpha x-\sin x$为减函数？
    \item 证明$\frac{x}{\sin x}$在$\left(0,\frac{\pi}{2}\right]$内是增函数．
    \item 若$x>0$, 求证
\begin{multicols}{2}
  \begin{enumerate}
\item  $\cos x-1+\frac{1}{2}x^2> 0$
\item $\sin x-x+\frac{1}{6}x^3> 0$
\item  $x-\frac{x^2}{2}<\ln (1+x)<x$
\end{enumerate}  
\end{multicols}

\item 证明下列不等式当$0<x<\frac{\pi}{2}$时成立：
\begin{multicols}{2}
    \begin{enumerate}
        \item $\tan x> x+\frac{x^3}{3}$
        \item $\frac{\sin x}{x}> \frac{2}{\pi}$
    \end{enumerate}
\end{multicols}

\item 已知$x+e^x=y+e^y$, 问$\sin x=\sin y$是否为真？
\item 若$a>b$, 求证$a^3-3a^2+3a>b^3-3b^2+3b$.
\item 设$f(0)=0$, 且$f'$在$[0,+\infty)$内递增，证明
$ g(x) =\frac{f (x)}{x}$在$(0,+\infty)$内递增．
\item 
\begin{enumerate}
    \item 已知$a,b$为实数，并且$e<a<b$，其中$e$是自然对数的底，证明：$a^b>b^a$.
    \item 如果正实数$a,b$满足$a^b=b^a$, 且$a<1$, 则$a=b$.
\end{enumerate}

\item 证明：对于所有$x\in\left(0,\frac{\pi}{2}\right]$, 有不等式
$(\sin x)^{-2}\le x^{-2}+1-\frac{4}{\pi^2}$
成立．
\end{enumerate}
\end{ex}


\subsubsection{函数的极大值与极小值}

从例2.59我们看到$x=1$和$x=2$分别是函数$f(x)=2x^3-9x^2+12x+5$取到极大值$f(1)=10$和极小值$f(2)=9$的点，我们称$x=1$为极大点，$x=2$为极小点，以后，极大值与极小值统称为\textbf{极值}，极大点与极小点统称为\textbf{极值点}．

下面我们来讨论如何确定一个函数的极值．
\begin{enumerate}
    \item 若$x_0$是$f(x)$的极值点，且$f(x)$在点$x_0$可微，那么总存在$x_0$的一个邻域，使$f(x)$在点$x_0$的邻域中满足费马条件，因而必有$f'(x_0)=0$. 例如例2.59, $f'(1)=f' (2) =0$.
    \item 若$f(x)$在$x_0$不可微，这时$x_0$也可能是极值点，例如，$f(x)=|x|$, 它在$x_0=0$这一点不可微，但从图2.28即可看出$x_0=0$是它的极小点．
\end{enumerate}

\begin{figure}[htp]
    \centering
    \begin{minipage}[t]{0.48\textwidth}
    \centering
\begin{tikzpicture}[>=latex, scale=1.3]
    \draw[->](-2,0)--(2,0)node[right]{$x$};
\draw[->](0,-.5)--(0,2)node[right]{$y$};
\draw[very thick](-1.8,1.8)--(0,0)node[below left]{$O$}--(1.8,1.8)node[right]{$y=|x|$};
\end{tikzpicture}
    \caption{}
    \end{minipage}
    \begin{minipage}[t]{0.48\textwidth}
    \centering
    \begin{tikzpicture}[>=latex, scale=.35]
        \draw[->](-5,0)--(5,0)node[right]{$x$};
\draw[->](0,-8)--(0,9)node[right]{$y$};
\draw[domain=-2:2, samples=100, very thick]plot(\x, {\x^3});
\node at (0,0)[below right]{$O$};
\foreach \x in {2,4,6,8}
{
    \draw(0,\x)node[left]{$\x$}--(.3,\x);
}
\draw(2,0)node[below]{2}--(2,.3);
    \end{tikzpicture}
    \caption{}
    \end{minipage}
    \end{figure}


这就告诉我们连续函数的极值点应该从它的导函数$f'(x)$的零点和$f(x)$的不可微点中去找，但是，$f'(x)$的零点和$f(x)$的不可微点只是$f(x)$可能达到极值的点，并不一定是极值点，例如$f(x)=x^3$, $x=0$是$f'(x)=3x^2$的零点，然而在整个区间$(-\infty,+\infty)$上，$f'(x)\ge 0$, 而且只在$x=0$时，$f'(x)=0$, 即函数$f(x)=x^3$在$(-\infty,+\infty)$上严格递增，因而$x=0$不是极值点，它的图象如图2.29, 在$y$轴的左侧，图象位于切线$y=0$的下方；而在$y$轴的右侧，图象位于切线$y=0$的上方，我们把方程$f'(x)=0$的根称为函数$f(x)$的\textbf{驻点}．

又如函数
\[g(x)=\begin{cases}
    \frac{1}{2}(x-1), & x\in[1,2]\\
    \frac{3}{2}(x-2)+ \frac{1}{2}(2-1), & x\in[2,3]\\
 (x-3) +\frac{1}{2}(2-1) + \frac{3}{2}(3-2), & x\in[3,4]
\end{cases}\]
的图象是折线，如图2.30, 易知
\[g' (x) > 0,\quad x\in (1, 2) \cup (2, 3) \cup (3, 4) \]
但$g'(2)$和$g'(3)$不存在，由于$g(x)$在$[1, 4]$上严格递增，所以$x=2$, $x=3$都不是极值点．

\begin{figure}[htp]
    \centering
\begin{tikzpicture}[>=latex]
    \draw[->](-.5,0)--(5,0)node[right]{$x$};
    \draw[->](0,-.5)--(0,4)node[right]{$y$};
\draw[domain=1:2, samples=5, very thick]plot(\x, {.5*(\x-1)});
\draw[domain=2:3, samples=5, very thick]plot(\x, {1.5*(\x-2)+.5});
\draw[domain=3:4, samples=5, very thick]plot(\x, {(\x-3)+2});
\foreach \x in {1,2,3,4}
{
    \node at (\x,0)[below]{$\x$};
}
\draw[dashed](2,0)--(2,.5);
\draw[dashed](3,0)--(3,2)node[above left]{$y=g(x)$};
\draw[dashed](4,0)--(4,3);
\node at (0,0)[below left]{$O$};
\end{tikzpicture}
    \caption{}
\end{figure}

综上讨论，我们得到函数$f$取极值的必要条件：

\begin{blk}
    {定理1（极值的必要条件）}若$f(x)$在$[a,b]$上连续，$x_0$是$f(x)$的极值点，那么$x_0$只可能是$f'(x)$的零点或$f(x)$的不可微点．
\end{blk}

下面给出极值点的两个充分性判别法．

\begin{blk}
  {定理2} 如果函数$f$在$x_0$的邻域内连续，并且在$(x_0-\delta,x_0)$和$(x_0,x_0+\delta)\; (\delta>0)$可微：
\begin{itemize}
    \item 在$(x_0-\delta,x_0)$内，$f'(x)<0\; (>0)$;
    \item 在$(x_0,x_0+\delta)$内，$f'(x)>0\; (<0)$;
\end{itemize}
那么，$x_0$是$f$的局部极小点（极大点）．
\end{blk}

\begin{proof}
按照前一节中定理2, $f$在$(x_0-\delta,x_0)$内严格递减（递增），而在$(x_0,x_0+\delta)$内严格递增（递减），故$f(x_0)$必为极小值（极大值）．
\end{proof}

\begin{rmk}
    这个定理并不要求函数$f$在点$x_0$处存在导数，如果导数存在，那么根据费马定理，它等于零．
\end{rmk}


\begin{example}
    讨论$y=\sqrt[3]{(x-1)(x+2)^2}$的极值，并绘出图象的草图．
\end{example}

\begin{solution}
将$y=\sqrt[3]{(x-1)(x+2)^2}$的两边立方，得
\begin{align*}
    y^3&= (x-1) (x+2)^2\\
    3y^2\cdot y'&= (x+2)^2+ (x-1)\cdot 2\cdot  (x+2)\\
    &= (x+2) [x+2+2(x-1)]\\
    &=3x (x+2) 
\end{align*}

$\therefore\quad 
y'=\frac{x (x+2)}{y^2}=\frac{x (x+2)}{\sqrt[3]{(x-1)^2(x+2)^4}}$
即
\[y'=\frac{x}{\sqrt[3]{(x-1)^2(x+2)}}\quad 
(x\ne 1,\; x\ne -2)\]
\begin{itemize}
    \item 当$x=0$时，$y'=0$, 函数图象在点$x=0$处有水平切线；
    \item 当$x=-2$和1时，$y'$不存在，但是由$\Lim_{x\to 1}y'=+\infty$和$\Lim_{x\to -2^+}y'=-\infty$, 而$\Lim_{x\to -2^-}y'=+\infty$得知函数图象在点$x=1$和$x=-2$处有平行$y$轴的切线，所以函数的可能的极值点是$x=-2, 0, 1$这三个点．
\end{itemize}

为确定曲线的$y$截距，令$x=0$, 得$y=-\sqrt[3]{4}\approx -1.59$; 为确定它的$x$截距，令$y=0$, 得$x=-2$和$x=1$, 现列表讨论它的增减性与极值如下：
\begin{center}
\begin{tabular}{c|ccccccc}
\hline
$x$ &$(-\infty,-2)$ & $-2$ & $(-2,0)$ & 0&$(0,1)$& 1& $(1,+\infty)$\\
\hline
$y'$ & $+$ & 不存在 & $-$ &0&$+$&不存在&$+$\\
$y$ & $\nearrow$ &极大值&$\searrow$ &极小值&$\nearrow$ &0&$\nearrow$\\ 
&&0&&$-\sqrt[3]{4}\approx -1.59$\\
\hline
\end{tabular}
\end{center}

函数$y=\sqrt[3]{(x-1)(x+2)^2}$的图象如图2.31.
\begin{figure}[htp]
    \centering
\begin{tikzpicture}[>=latex]
    \draw[->](-4,0)--(3,0)node[right]{$x$};
    \draw[->](0,-3)--(0,3)node[right]{$y$};
\draw[domain=-4:-2.001, samples=100, very thick]plot(\x, {-exp(1/3*ln(1-\x)+2/3*ln(-\x-2))});

\draw[domain=-1.999:.999, samples=100, very thick]plot(\x, {-exp(1/3*ln(1-\x)+2/3*ln(\x+2))});
\draw[domain=1.001:2.25, samples=100, very thick]plot(\x, {exp(1/3*ln(\x-1)+2/3*ln(\x+2))});
\draw[very thick](1,.4)--(1,-.4);

\node at (-2,0) [above]{$-2$};
\node at (1,0) [below right]{$1$};
\draw(2,0)node[below]{$2$}--(2,.1);
\draw(-3,0)node[below]{$-3$}--(-3,.1);
\draw(-1,0)node[below]{$-1$}--(-1,.1);
\draw[dashed](-2,-3)--(-2,3);
\draw[dashed](1,-3)--(1,3);
\end{tikzpicture}
    \caption{}
\end{figure}
\end{solution}

\begin{blk}
{定理3} 设函数$f(x)$在区间$(a,b)$中有二阶连续的导函数，又$c\in (a,b)$, 使得$f' (c) =0$,
\begin{enumerate}
    \item 如果$f''(c)<0$, 则$c$是$f(x)$的一个极大点．
    \item 如果$f''(c)>0$, 则$c$是$f(x)$的一个极小点．
\end{enumerate}
\end{blk}

\begin{proof}
设$f''(c)<0$, 由于$f''(x)$的连续性，故存在点$c$的一个$\delta>0$的邻域$(c-\delta,c+\delta)\subset (a,b)$, 使得$f''(x)＜0$, $x\in (c-\delta,c+\delta)$. 按前节定理2, $f'(x)$在$(c-\delta,c+\delta)$中是递减的，又由于$f'(c)=0$, 因此，
$f'(x)$在点$c$处由正变为负，由前面的定理2, $c$是$f(x)$的极大点．

当$f''(c)>0$时，同理可证$c$是$f(x)$的一个极小点．

为说明定理3的应用，我们求例2.60的二阶导函数：
\begin{align*}
   y''&=\frac{y^2 (2x+2) -x (x+2) \cdot 2y\cdot y'}{y^4}\\
&=\frac{2 [ (x+1) y-x (x+2) y']}{y^3} 
\end{align*}

从例2.60知，当$x=0$时，$y'=0$. 又$f(0)=-\sqrt[3]{4}$. 故在$x=0$处
\[y''_{x=0}=\frac{2\left[(0+1)\left(-\sqrt[3]{4}\right)-0\right]}{-4}=\frac{\sqrt[3]{4}}{2}>0\]
由定理3, $x=0$是函数$y=\sqrt[3]{(x-1)(x+2)^2}$的极小点．
\end{proof}




\begin{example}
    讨论$y=x^3+ax^2+bx+c$的极值．
\end{example}


\begin{solution}
\[y=x^3+ax^2+bx+c,\qquad 
    y'=3x^2+2ax+b\]

$ \Delta=a^2-3b$是导函数的判别式．
    我们可以分$\Delta>0$, $\Delta=0$和$\Delta<0$的三种情形讨论$y'$的正负区间．
\begin{enumerate}
    \item 若$\Delta>0$, 则
\[    y'=3x^2+2ax+b=3 (x-x_1) (x-x_2)\]
    这里$x_1=\frac{1}{3}\left(-a-\sqrt{a^2-3b}\right)$和$x_2=\frac{1}{3}\left(-a+\sqrt{a^2-3b}\right)$
    是导函数$y'$的二个实根．
\begin{enumerate}
    \item 当$x\in (-\infty,x_1)\cup (x_2,+\infty)$时，$y'>0$;
    \item 当$x\in (x_1,x_2)$时，$y'<0$
\end{enumerate}
    所以，$x_1$是$y$的一个极大点，$x_2$是$y$的一个极小点．
    
\item 若$\Delta=0$, 则
    \[y'=3x^2+2ax+\frac{a^2}{3}=3\left(x+\frac{a}{3}\right)^2\]
这里，$y'$虽然在$-\frac{a}{3}$处等于0, 但$y'$在$-\frac{a}{3}$的左、右两侧都取正值，即$y=x^3+ax^2+bx+c$是一个递增函数．因此，$-\frac{a}{3}$不是$y$的极值点．
\item 若$\Delta<0$, 则对于任意$x\in\mathbb{R}$, 恒有\[y'=3x^2+2ax+b=3 \left(x+\frac{a}{3}\right)^2+\frac{3b-a^2}{3}> 0\]
因此，$y=x^3+ax^2+bx+c$是一个增函数，也是没有极值点．
\end{enumerate}    
\end{solution}

\begin{ex}
\begin{enumerate}
    \item 先求出$f'(x)$, 然后由$f'(x)$的正负去讨论$f(x)$的增减与极值，并画出函数图象的草图：
\begin{multicols}{2}
    \begin{enumerate}
        \item $f (x) =2x^3+9x^2-24x+5$
        \item $f (x) =x^4-6x^2-8x+15$
        \item $f (x) =-x^5-5x^4$
        \item $f (x) = (x-1)^3$
     \item  $f (x) =|x^2-1|$
     \item $f (x) = (x^2-1)^3+1$
    \end{enumerate}
\end{multicols}
\item 讨论$f(x)=\frac{x^2-2x-3}{x^2+1}$的增减和极值，并画函数图象的草图．
    \item 求$f(\theta)=\sin\theta-\frac{1}{2}\sin2\theta$在$(-\pi,\pi)$内的极大值和极
    小值．
    \item 
\begin{enumerate}
    \item 求两条曲线$y=\frac{4e^x+2e^{-x}}{3}$和$y=e^{-x}$的交点的坐标；
    \item 求函数$y=\frac{4e^x+2e^{-x}}{3}$的极值点，又它在此点取极大值，还是极小值？
\end{enumerate}

\item 讨论函数$y=x^2e^{-x^2}$的增减与极值．
\item 求$f(x)=x^x\; (x>0)$的极值，并说明此极值是极大值
还是极小值．
\item 求函数$\frac{\ln x}{x}$的极大值，

\item 函数$f(x)=\frac{ax+b}{(x-1)(x-4)}$在点$(2,-1)$有一极值，
求$a,b$并证明此极值是极大值．
\item 证明$f(x)=\left(\alpha-\frac{1}{\alpha}-x\right)(4-3x^2)$仅有一个极大值
和一个极小值，且它们的差是$\frac{4}{9}\left(\alpha+\frac{1}{\alpha}\right)^3$, 又此差值的绝对值的最小值是多少？
\item 函数$f(x)=x(x-1)(x-a)$有绝对值相等，符号相反的
极大值与极小值，试确定常数$a$的值．
\item 求曲线$y=e^{-x}\sin x$上的极值点，并证明极值点上的纵坐
标成等比数列，又该等比数列的公比是多少？  
\end{enumerate} 
\end{ex}


\subsubsection{函数的最大值和最小值}
在实践中有一类“最大”、“最小”、“最省”的问题，例如生产搪瓷杯时，就要考虑在一定容积下，杯子的高和直径取多大时，用料最省．又如把一根圆木锯成矩形截面横梁时，怎样选取矩形的长和宽才能使横梁强度最大？这类“最大”、“最小”、“最省”等问题，在数学上叫做最大值、最小值问题．

我们曾在第一章中指出过“在闭区间$[a,b]$上连续的函数$f(x)$总是具有最大值和最小值”，但是，微积分并未提供确定函数$f(x)$的最大值、最小值的直接方法．如果我们讨论的函数$f(x)$在$(a,b)$上有连续的一阶导数，那么只要求出$f(x)$在$(a,b)$内的驻点，即$f'(x)=0$的根，就可以求出$f(x)$的局部极值．显然，$f(x)$在$[a,b]$上的最大值和最小值一定出现在局部极大值、局部极小值和端点值之中，从几何上来说，最大值、最小值如果不是位于定义区间的端点上，则分别是曲线的波峰和波谷．从图2.33可以看到局部极大值可能小于局部极小值．

\begin{figure}[htp]
    \centering
\begin{tikzpicture}[>=latex]
        \draw[->](-.5,0)--(5,0)node[right]{$x$};
        \draw[->](0,-.5)--(0,4.5)node[right]{$y$};
        
        \node at (0,0) {};
        \draw plot[smooth] coordinates{(.5,4)(1,2.5)(1.9, 3.5)(3.2, .8)(4,1.5)(4.4,1)};
%        \pscurve[linewidth=2pt](.5,4)(1,2.5)(1.9, 3.5)(3.2, .8)(4,1.5)(4.4,1)
%        };
        \draw[dashed](.5,4)--(.5,0)node[below]{$a$};
        \draw[dashed](4.4,1)--(4.4,0)node[below]{$b$};
\end{tikzpicture}
    \caption{}
\end{figure}

下面给出求函数$f(x)$在$[a,b]$上最大值和最小值的法则，这个法则适用于在$[a,b]$上连续，在$(a,b)$内可
微的函数$f(x)$,并且$f'(x)$最多在有限个点上等于零．
\begin{enumerate}
    \item 求$f(x)$在$(a,b)$内的驻点，即解$f'(x)=0$;
    \item 计算$f(x)$在驻点和端点的函数值，并把这些值加以比较，其中最大的一个为最大值，最小的一个为最小值．
\end{enumerate}

假如$f'(x)$在$(a,b)$内有一个驻点$c$, 且在整个$(a,b)$
上，$f''(x)<0\; (>0)$, 那么点$c$是最大值（最小值）点．

事实上，如果$f''(x)<0,\; x\in (a,b)$, 那么函数$f'(x)$在$[a,b]$上递减，因而$c$是其唯一的驻点，由中值定理知：

当$a\le x<c$时，$f'(x)-f'(c)=(x-c)\cdot f''(\xi)$，这里$x<\xi<c$．

$\because\quad f' (c) =0,\quad f'' (\xi)<0,\quad x-c<0$,

$\therefore\quad f'(x)>0$

同理说明当$c<x\le b$时，$f'(x)<0$．再根据$f(x)$在$[a,c]$上递增而在$[c,b]$上递减．故当$x\ne c$时有$f(x)<f(c)$．由此可知$c$是最大值点．当$f''(x)>0$时，也可以进行同样的论证．

如果$f(x)$在$[a,b]$内有不可微的点，而且这种点只有有限个，那么求最大值和最小值的法则改为求出三种点的对应函数值再加以比较，即$f(x)$的驻点的函数值，$f(x)$的不可微点的函数值，最后还有端点的函数值$f(a)$和$f(b)$, 这些值中最大的就是最大值，最小的就是最小值．

在实际问题中，如果在$(a,b)$内部，$f(x)$只有一个驻点$c$, 而且从实际问题本身又可以知道在$(a,b)$内必定有最大值或最小值，那么$f(c)$就是所要求的最大值或最小值，不需要算出$f(a)$, $f(b)$了．

\begin{example}
求函数$f(x)=x^{\tfrac{2}{3}}-(x^2-1)^{\tfrac{1}{3}}$在区间$[-2,2]$上的最大值和最小值．
\end{example}


\begin{solution}
\begin{align*}
f'(x)&=\frac{2}{3}x^{-\tfrac{1}{3}}-\frac{1}{3}(x^2-1)^{-\tfrac{2}{3}}\cdot 2x\\
&=\frac{2}{3} \frac{(x^2-1)^{\tfrac{2}{3}}-x^{\tfrac{4}{3}}}{x^{\tfrac{1}{3}}(x^2-1)^{\tfrac{2}{3}}}    
\end{align*}
    故$f(x)$的驻点是$\pm\frac{1}{\sqrt{2}}$, 不可微点是$x=0$, $x=\pm 1$, 且这些点都在$[-2, 2]$内，因此$f(x)$的极值点可能是
\[-1,\quad \frac{-1}{\sqrt{2}},\quad  0,\quad \frac{1}{\sqrt{2}},\quad  1\]
    由于$f(x)$是偶函数，故仅需计算
\[    f (0) =1,\qquad f\left(\frac{1}{\sqrt{2}}\right) =\sqrt[3]{4},\qquad f (1) =1\]
    又在端点$x=2$处的函数值$f(2)=\sqrt[3]{4}-\sqrt[3]{3}$, 把上面这些函数值相互比较，可见$f(x)$在$[-2, 2]$中的最大值为$\sqrt[3]{4}\approx 1.59$, 最小值为$\sqrt[3]{4}-\sqrt[3]{3}\approx 1.59-1. 44\approx 0.15$.
\end{solution}

\begin{example}
    设$a>0$, 证明
$    f (x) =\frac{1}{1+|x|}+\frac{1}{1+|x-a|}$
    的最大值是$\frac{2+a}{1+a}$.
\end{example}    

\begin{proof}
把$f(x)$写成分段函数
\[f(x)=\begin{cases}
    \frac{1}{1-x}+\frac{1}{1+a-x}, & x<0\\
    \frac{1}{1+x}+\frac{1}{1+a-x}, & 0\le x\le a\\
    \frac{1}{1+x}+\frac{1}{1+x-a}, & x>a
\end{cases}\]
于是
\[f'(x)=\begin{cases}
    \frac{1}{(1-x)^2}+\frac{1}{(1+a-x)^2}, & x<0\\
    \frac{-1}{(1+x)^2}+\frac{1}{(1+a-x)^2}, & 0\le x\le a\\
    \frac{-1}{(1+x)^2}-\frac{1}{(1+x-a)^2}, & x>a
\end{cases}\]
因而$f(x)$在$(-\infty, 0]$上递增，而在$[a,+\infty)$上递减，由此得知$f(x)$在$[0,a]$上的最大值也就是在$\mathbb{R}$上的最大值．

设对于$(0,a)$内的$x$有$f'(x)=0$, 则
\[(1+x)^2- (1+a-x)^2=0\]
上式唯一解是$x=\frac{a}{2}$. 
因为$f\left(\frac{a}{2}\right)=\frac{4}{2+a}$, $f(0)=f(a)=\frac{2+a}{1+a}$，
并且
\[\frac{2+a}{1+a}-\frac{4}{2+a}= \frac{a^2}{(1+a)(2+a)}>0\quad  (a>0)\]
所以$f(x)$的最大值是$\frac{2+a}{1+a}$.
\end{proof}

\begin{example}
    用铁皮做体积为$V$的圆柱形封闭容器，如何使材料最省？
\end{example}


\begin{solution}
\textbf{解法1:} 设圆柱形底面半径为$r$, 高为$h$, 于是需用材料即圆柱表面积
\[S=2\pi rh+2\pi r^2\]
但$V=\pi r^2h$为已知，因此，$h=\frac{V}{\pi r^2}$，
从而问题变为求函数$S(r)=\frac{2V}{r}+2\pi r^2\; (0<r<+\infty)$的最小值．

根据经验，当$r$很小时，$S$很大，又当$r$很大时，$S$也很大，因此在$(0,+\infty)$内必有$r$的一个值使$S$最小，即它是$S'(r)=0$的一个根，由
\[S' (r) =-\frac{2V}{r^2}+4\pi r=0\]
得唯一正根$r=\sqrt[3]{\frac{V}{2\pi}}$．
从这个实际问题知道，当$r=\sqrt[3]{\frac{V}{2\pi}}$
时，$S$为最小，这时相应的高为
\[h=\frac{V}{\pi r^2}=\frac{2\pi r^3}{\pi r^2}=2r=2\sqrt[3]{\frac{V}{2\pi}}\]
这说明圆柱的高和直径相等时，用料最省．

\textbf{解法2:}  我们可以不先将下面两个式子中的$h$消去，再求$\frac{\dd S}{\dd r}=0$的根
\begin{align}
    S&=2\pi r h+2\pi r^2\\
    \pi r^2 h&=V
\end{align}
而是保留辅助变量$h$，将
(2.27), (2.28)对$r$求导数，只须记着$V$是常量，$h$是$r$的函数，于是由$\frac{\dd S}{\dd r}=0$，得
\begin{equation}
    \frac{\dd S}{\dd r}=2\pi\left(h+r\frac{\dd h}{\dd r}+2r\right)=0
\end{equation}
由(2.28)得
\begin{equation}
    \pi\left(2rh+r^2\frac{\dd h}{\dd r}\right)=0
\end{equation}
由此得：
\[\frac{\dd h}{\dd r}=-\frac{2h}{r}\]
将它代入(2.29)，得
\[2\pi\left[h+r\left(-\frac{2h}{r}\right)+2r\right]=0\]
所以：$2r=h$. 

同样地说明了圆柱的高和直径相等时，用料最省．
\end{solution}


\begin{example}
由半径为$\ell$的圆铁皮，剪去一个扇形，把剩下的部分围成一个圆锥形容器，为了使所得圆锥容器有最大的容积，问剪去扇形的中心角应该多大？
\end{example}

\begin{figure}[htp]
    \centering
\begin{tikzpicture}[>=latex]
\begin{scope}[scale=1.5]
\draw[pattern=north east lines](0,0)--(-135:1.5) arc (-135:135:1.5)--(0,0);
\draw[<->](-135:.2) arc (-135:135:.2);
\node at (.5,0)[fill=white]{$x$};
\node at (1,1)[above right]{$x\ell=2\pi r$};
\tkzDefPoint(-30:1.5){A}
\tkzDefPoint(-30:2){B}
\tkzDrawCircle(B,A)
\draw[->](B)--node[above]{$r$}+(10:.5);
\draw[decorate,decoration={brace,mirror,raise=2pt}](135:1.5)--node[below left]{$\ell$}(0,0);


\end{scope}
\begin{scope}[xshift=7cm, scale=.7]
\draw[dashed](0,-2) ellipse [x radius=1.5cm, y radius=.5cm];
\draw[thick](-1.5,-2)--(0,3)--node[above right]{$\ell$}(1.5,-2);
\draw[dashed](0,3)--node[left]{$h$}(0,-2)--node[below]{$r$}(1.5,-2);
\draw[thick](1.5,-2) arc [start angle =0, end angle=-180, x radius=1.5cm, y radius=.5cm];
\end{scope}


\end{tikzpicture}
    \caption{}
\end{figure}


\begin{solution}
    设剪去扇形后，所剩材料的圆心角为$x$（弧度），扇形的弧长为$x\ell$, 由它围成一个圆锥，其底面半径为$r$, 母线长为$\ell$, 高为$h$, 如图2.34. 于是
\[r=\frac{\ell x}{2\pi}\]
圆锥高为
\[h=\sqrt{\ell^2-r^2}=\sqrt{\ell^2-\left(\frac{\ell x}{2\pi}\right)^2}=\frac{\ell }{2\pi}\sqrt{4\pi^2-x^2}\]
因此圆锥的体积为
\begin{align*}
    V&=\frac{1}{3}\pi r^2 h\\
    &=\frac{1}{3}\pi\left(\frac{\ell x}{2\pi}\right)^2\cdot \frac{\ell}{2\pi}\sqrt{4\pi^2-x^2}\\
    &=\frac{\ell^3 x^2}{24\pi^2}\sqrt{4\pi^2-x^2}\quad (0\le x\le 2\pi)
\end{align*}
由于$V(0)=0$, $V(2\pi)=0$, 故在$(0, 2\pi)$内必有一个值$x_0$使$V$的值最大，因而使$V'(x_0)=0$. 由于
\begin{align*}
    V'(x)&=\frac{\ell^3}{24\pi^2}\left[2x\sqrt{4\pi^2-x^2}+\frac{1}{2}x^2(4\pi^2-x^2)^{-\tfrac{1}{2}}\cdot (-2x)\right]\\
    &=\frac{\ell^3}{24\pi^2}\left(\frac{8\pi^2 x-3x^3}{\sqrt{4\pi^2-x^2}}\right)
\end{align*}
令$V'(x)=0$, 即$8\pi^2 x-3x^3=0$, 解得
\[x_1=0,\quad x_2=-2\pi\sqrt{\frac{2}{3}}, \quad x_3=2\pi\sqrt{\frac{2}{3}}\]
这里只有$x_3=2\pi\sqrt{\frac{2}{3}}\in (0, 2\pi)$, 故知$x=2\pi\sqrt{\frac{2}{3}}$是$V(x)$在$[0, 2\pi]$内的唯一的极大点，因而也是最大点．于是最后得到结论，剪去扇形的圆心角为
\[2\pi-2\pi\sqrt{\frac{2}{3}}=\frac{2}{3}\pi\left(3-\sqrt{6}\right)\; \text{（弧度）}\]
时，才能使所围成的容器有最大的容积．
\end{solution}

\begin{example}
    求证：若$p>q>1$，又$x\ge 0$且$x\ne 1$，则$\frac{x^p-1}{p}>\frac{x^q-1}{q}$.
\end{example}

\begin{proof}
设$f(x)=\frac{x^p-1}{p}-\frac{x^q-1}{q}\; (x\ge 0)$，则：
\[f'(x)=x^{p-1}-x^{q-1}=x^{q-1}\left(x^{p-q}-1\right)\quad (p>q>1)\]
    函数的增、减如下表所示：
\begin{center}
\begin{tabular}{c|cccc}
\hline
$x$&0&$(0,1)$&1& $(1,+\infty)$\\
\hline
$f'(x)$&0&$-$&0&$+$\\
$f(x)$&$\frac{1}{q}-\frac{1}{p}$&$\searrow$&极小值0& $\nearrow$\\ 
\hline
    \end{tabular}
\end{center}

$\because\quad f(0)=\frac{1}{q}-\frac{1}{p}>0=f(1)$

$\therefore\quad f(1)=0$不仅是极小值也是最小值．于是：
\[f(x)=\frac{x^p-1}{p}-\frac{x^q-1}{q}\ge 0\quad (x\ge 0)\]
因为$x\ne 1$，所以
\[\frac{x^p-1}{p}>\frac{x^q-1}{q}\quad (x\ge 0,\; x\ne 1)\]
\end{proof}

\begin{example}
 （光的折射定律）设有两种介质，$X'X$是它们的交界线，光由一种介质中的某点$A$走向另一种介质中的某点$B$. 又设当光走过交界线$X'X$时，它的速度由$v_1$变成$v_2$. 依据物理学中的一个基本原理“光由$A$到达$B$时，所采取的途径要使用的时间为最小”，如果两种介质是均匀的（例如一为空气，一为水），那么$v_1,v_2$都是常数，从而光在这两种介质中所走的路径是两段直线$AP$, $PB$, $P$表示光达到交界线时与这条线的交点，如果$AB$不与$X'X$垂直，一般$APB$不是直线，我们来定出$AP$, $BP$与$X'X$的法线的夹角$\theta_1$和$\theta_2$的关系．
\end{example}

    \begin{figure}[htp]
    \centering
\begin{tikzpicture}[>=latex, scale=.7]
\draw[->](-1,0)node[left]{$X'$}--(6,0)node[right]{$X$};
\draw[->](0,-4)node[right]{$S'$}--(0,4)node[right]{$S$};
\node at (0,0)[below left]{$O$};
\draw[->, very thick](0,3)node[left]{$A(0,a)$}--(2,0)node[above left]{$P$};
\draw[->, very thick](2,0)--(4.2,-4)node[right]{$B(c,b)$};
\draw[dashed](2,-2)--(2,2);
\draw[<->](0,-.4)--node[fill=white]{$x$}(2,-.4);
\draw[|<->](4.2,.4)--node[fill=white]{$c-x$}(2,.4);
\draw[dashed](4.2,-4)--(4.2,0);
\draw(2,.75) arc (90:90+33.7:.75)node[above]{$\theta_1$};
\draw(2,-.75) arc (-90:-90+28.8:.75)node[below]{$\theta_2$};


\end{tikzpicture}
    \caption{}
\end{figure}

\begin{solution}
如图2.35建立坐标系，$A$点坐标是$(0,a)$, $B$点的坐标是$(c,b)$, 设$P$点的坐标是$(x,0)$, 光由$A$到$P$所需时间是$\frac{\sqrt{a^2+x^2}}{v_1}$，从$P$到$B$所需时间是$\frac{\sqrt{b^2+(c-x)^2}}{v_2}$, 
则整个时间$t(x)$是
\[t (x) =\frac{\sqrt{a^2+x^2}}{v_1}+\frac{\sqrt{b^2+(c-x)^2}}{v_2}\]

为了使$t(x)$极小，应使$t'(x)=0,\; x\in (0,c)$. 由$t(x)$得
\begin{align*}
    t' (x) &=\frac{x}{v_1\sqrt{a^2+x^2}}-\frac{c-x}{v_2\sqrt{b^2+(c-x)^2}}\\
    t'' (x) &=\frac{-x^2}{v_1(a^2+x^2)^{\tfrac{3}{2}}}+\frac{1}{v_1\sqrt{a^2+x^2}}+\frac{1}{v_2\sqrt{b^2+(c-x)^2}}-\frac{(c-x)^2}{v_2\left[b^2+(c-x)^2\right]^{\tfrac{3}{2}}}\\
    &=\frac{a^2}{v_1(a^2+x^2)^{\tfrac{3}{2}}}+\frac{b^2}{v_2\left[b^2+(c-x)^2\right]^{\tfrac{3}{2}}}>0.
\end{align*}
既然$t'(0)<0$, $t'(c)>0$, 及$t'$是严格递增的，故在$(0,c)$内有唯一的$x_0$, 使得
\[t'(x_0)=0\]
从而当$x=x_0$时，$t(x_0)$是极小值也是最小值，于是由$t'(x_0)=0$得到
\[\frac{x_0}{v_1\sqrt{a^2+x^2_0}}=\frac{c-x_0}{v_2\sqrt{b^2+(c-x_0)^2}}\]
即
\[\frac{\sin\theta_1}{v_1}=\frac{\sin\theta_2}{v_2}\]
这就是光的折射定律．
\end{solution}

\begin{ex}
\begin{enumerate}
    \item 由$y=0$, $x=8$, $y=x^2$围成一曲边三角形$OAB$, 在曲边$\wideparen{OB}$上，求一点使得过此点所作$y=x^2$的切线与$OA$、$AB$所围成的三角形面积最大．
    \item 椭圆$\frac{x^2}{a^2}+\frac{y^2}{b^2}=1$的内接矩形中怎样的矩形面积最大？又此最大的矩形面积比椭圆的内接正方形的面积大多少？
    \item 设$\triangle ABC$是定圆的内接三角形，其中$A$、$B$两点固定，$C$点在圆周上变动，试求其面积最大时的位置．
    \item 要做一个底面为长方形的带盖箱子，其体积为72立方厘米，其底边成$1:2$的关系，问各边的长怎样，才能使表面积最小？
    \item 设有底为等边三角形的直棱柱，体积为$V$, 要使其总面积为最小，问底边长应为多少？
    \item 设正三棱锥的侧面积为$S$, 问当侧面与底面所成角$\alpha$为何值时，其体积最大？
    \item 平面上通过一个已知点$P(1, 4)$引一条直线，要使它在两个坐标轴上的截距都为正，且它们的和为最小，求这条直线的方程．
    \item 直圆柱内接于半径为$R$的球中，当圆柱的高为何值时，其
    体积最大．
    \item 三个点$A$, $B$和$C$不在同一直线上，$\angle ABC=60^{\circ}$, 汽车以80公里/小时的速度由$A$向$B$行驶，同时火车以50公里/小时的速度由$B$向$C$行驶．若$AB=200$公里，问运动开始几小时后汽车与火车的距离最小？
    \item 船航行一昼夜的耗费由两部分组成：一为固定部分，设为$a$元；另一为变动部分，设它与速度立方成比例，试问船应以怎样的速度$v$行驶为最经济？
    \item 在某种产品的制造过程中，次品率$y$依赖于日产量$x$, 即$y=y(x)$, 已知
\[y(x)=\begin{cases}
    \frac{1}{101-x},& 0\le x\le 100,\\
    1,& x>100
\end{cases}\]
其中$x$为正整数．又该厂每生产一件产品可盈利$A$元，但每生产出一件次品就要损失$\frac{A}{3}$元．问为了获得最大盈利，该厂的日产量应定为多少？
\item 抛物线$y=x^2$上的哪一点到直线$x-y-2=0$的距离最短？\item 一弹在空中飞行（不计空气阻力），其弹道方程为
\[y=mx-\frac{(m^2+1)x^2}{800}\]
这里原点取在此弹发射之点，其中$m$为曲线在原点处的切线斜率．
\begin{enumerate}
    \item 若要此弹击中同一水平面上最远距离的目标；
    \item 若要此弹击中300米远处一直立墙壁上的最大高度，问$m$之值各为多少？
\end{enumerate}
\item 若直角三角形的周长$2p$一定，求它的面积的最大值．
\end{enumerate}  
\end{ex}

\subsection{函数的凸性与拐点}
上面已经讨论了函数的增减性与极值，这对于我们了解函数的性质与图象形状有很大的帮助．现在我们还要进一步考察函数在什么区间是向上凸的，什么区间是向下凸的，以及上凸函数和下凸函数的一些性质．

\begin{blk}
    {定义1} 函数$f$在区间$(a,b)$内，称为向下（上）凸的，指对于任意$c\in (a,b)$, 曲线位于在点$(c,f(c))$处的切线上（下）面（见图2.37和2.38）.
\end{blk}    


\begin{figure}[htp]
    \centering
    \begin{minipage}[t]{0.48\textwidth}
    \centering
\begin{tikzpicture}[>=latex, scale=.9]
\draw[->](-.5,0)--(5,0)node[right]{$x$};
\draw[->](0,-.5)--(0,4)node[right]{$y$};
\draw[domain=1:4, samples=100, very thick]plot(\x,{0.5*(\x-2)^2+1})node[right]{$y=f(x)$};
\draw[dashed](1,1.5)--(1,0)node[below]{$a$};
\draw[dashed](4,3)--(4,0)node[below]{$b$};
\draw[domain=.5:5, samples=10,thick]plot(\x,{0.5*\x-0.125});
\node at (0,0)[below left]{$O$};
\node at (2.5,1.125)[below]{$(c,f(c))$};
\end{tikzpicture}
    \caption{}
    \end{minipage}
    \begin{minipage}[t]{0.48\textwidth}
    \centering
    \begin{tikzpicture}[>=latex, scale=.9]
\draw[->](-.5,0)--(5,0)node[right]{$x$};
\draw[->](0,-.5)--(0,4)node[right]{$y$};
\draw[domain=1.5:4, samples=100, very thick]plot(\x,{-0.5*(\x-3)^2+2})node[right]{$y=f(x)$};
\draw[dashed](1.5,.875)--(1.5,0)node[below]{$a$};
\draw[dashed](4,1.5)--(4,0)node[below]{$b$};
\node at (2.5,1.875)[above]{$(c,f(c))$};
\draw[domain=.5:5, samples=10,thick]plot(\x,{0.5*\x+0.625});
\node at (0,0)[below left]{$O$};
    \end{tikzpicture}
    \caption{}
    \end{minipage}
    \end{figure}


\begin{blk}
    {定义2} 函数$f$的图象若经过某点$x=c$时，它的图象
$y=f(x)$由向上凸变为向下凸，或者由向下凸变为向上凸，则点$x=c$称为函数$f$的\textbf{拐点}，如图2.39和2.40.
\end{blk}    
    

\begin{figure}[htp]
    \centering
    \begin{minipage}[t]{0.48\textwidth}
    \centering
\begin{tikzpicture}[>=latex, yscale=.8]
    \draw[->](0,.5)--(5,.5)node[right]{$x$};
    \draw[->](1,0)--(1,4)node[right]{$y$};
    \draw[domain=.5*pi:1.5*pi, samples=100, very thick]plot(\x,{sin(\x r)+2});
    \node at (1,.5)[below left]{$O$};
\node at (3.1416, 2)[above right]{$(c,f(c))$};
\draw[domain=2:4.4, samples=10, thick]plot(\x,{-\x+5.1416});
\draw[dashed](3.1416,2)--(3.1416,.5)node[below]{$c$};
\end{tikzpicture}
    \caption{}
    \end{minipage}
    \begin{minipage}[t]{0.48\textwidth}
    \centering
    \begin{tikzpicture}[>=latex, yscale=.8]
   \draw[->](-3,.5)--(2,.5)node[right]{$x$};
    \draw[->](-2,0)--(-2,4)node[right]{$y$};
    \draw[domain=-.5*pi:.5*pi, samples=100, very thick]plot(\x,{sin(\x r)+2});
    \node at (-2,.5)[below left]{$O$};
\node at (0,2)[above left]{$(c,f(c))$};
\draw[domain=-1.25:1.25, samples=10, thick]plot(\x,{\x+2});   
\draw[dashed](0,2)--(0,.5)node[below]{$c$};

    \end{tikzpicture}
    \caption{}
    \end{minipage}
    \end{figure}


这时过$x=c$的点的切线，把曲线分成两部分，换言之，存在充分小的邻域$\delta>0$, 使得对于一切$x\in (c-\delta,c)$, 曲线位于在点$c$处的切线的一方，而对于$x\in (c,c+\delta)$, 曲线在该切线的另一方．

\begin{blk}
  {凸性判定定理} 设函数$f$在区间$(a,b)$内有二阶导数，则
\begin{enumerate}
    \item 函数f在区间$(a,b)$内为向下凸的充分必要条件是，对于每个$x\in (a,b)$, 有$f''(x)\ge 0$;
    \item 函数f在区间$(a,b)$内为向上凸的充分必要条件是，对于每个$x\in (a,b)$, 有$f''(x)\le 0$.
\end{enumerate}
\end{blk}

\begin{proof}
    假设曲线$y=f(x)$在$(a,b)$内是向下凸的，因为函数$f$在$(a,b)$内有二阶导数，则对于任意一点$c\in(a,b)$, $f'(c)$存在，曲线$y=f(x)$在点$x=c$处的切线方程是
\begin{equation}
    y=f (c) +f' (c) (x-c)
\end{equation}
    对于每个点$x\in (a,b)$, 曲线$y=f(x)$的对应点的纵坐
标与切线上的对应点的纵坐标之差用$g(x)$表示，如图2.41, 于是
\begin{equation}
   g (x) =f (x) -[f (c) +f' (c) (x-c)]=[f(x)-f (c) ] - (x-c) f' (c)  
\end{equation}
根据中值定理，得到
\[g (x) = (x-c) f' (d) - (x-c) f'(c)\]
这里$d\in (x,c)$或$(c,x)$. 又$g(x)=(x-c)\cdot [f'(d)-f'(c)]$, 再根据中值定理，得到
\begin{equation}
    g (x) = (x-c) (d-c)f''(e)
\end{equation}
这里$e\in (d,c)$或$(c,d)$.

于是对于一切$x\in (a,b)$, 点$(x,f(x))$不在切线(2.31)的下方的充要条件是$g(x)\ge 0$, 显然只有点$c$使$g(c)=0$. 根据等式(2.33), 并且注意到
\[(x-c)(d-c)>0\]
便知道$g(x)\ge 0$的充分必要条件是$f''(e)\ge 0$. 这里$e\in(a,b)$且随着$c$和$x$的不同值而改变．因此，依曲线向下凸的定义，曲线向下凸的充分必要条件是对于任意$x\in (a,b)$有$f''(x)\ge 0$.

用同样的方法便可说明曲线$y=f(x)$向上凸的充分必要条件是$f''(x)\le 0$.
\end{proof}

\begin{figure}[htp]
    \centering
    \begin{minipage}[t]{0.48\textwidth}
    \centering
\begin{tikzpicture}[>=latex, scale=1]
    \draw[->](-.5,0)--(5,0)node[right]{$x$};
\draw[->](0,-.5)--(0,4)node[right]{$y$};
\draw[domain=.75:3.1, samples=100, very thick]plot(\x, {-1/(\x-3.5)+1})node[right]{$y=f(x)$};
\draw[domain=.75:3, samples=10, thick]plot(\x,{\x-0.5});
\draw[dashed] (2.5,2)node[left]{$(c,f(c))$}--(2.5,0)node[below]{$c$};
\draw[dashed] (1,1.4)--(1,0)node[below]{$x$};
\draw[|<->|](.8,1.4)--node[left]{$g(x)$}(.8,.5);
\node at (0,0) [below left]{$O$};
\end{tikzpicture}
    \caption{}
    \end{minipage}
    \begin{minipage}[t]{0.48\textwidth}
    \centering
    \begin{tikzpicture}[>=latex, scale=1]
       \draw[->](-.5,0)--(5,0)node[right]{$x$};
\draw[->](0,-.5)--(0,4)node[right]{$y$};
\draw[domain=.25:4, samples=100, very thick]plot(\x,{0.4*(\x-1.5)^2+1})node[right]{$y=f(x)$};
\draw[domain=.25:4, samples=10, thick]plot(\x,{.6*\x+.5})node[right]{$y=\ell(x)$};

\draw[dashed] (1,1.1)--(1,0)node[below]{$a$};
\draw[dashed] (3.5,2.6)--(3.5,0)node[below]{$b$};
\draw(2,0)node[below]{$x$}--(2,1.7)node[above left]{$(x,\ell(x))$};
\draw[<->|](2.2,0)--node[right]{$(x,f(x))$}(2.2,1.1);
\node at (0,0)[below left]{$O$};
    \end{tikzpicture}
    \caption{}
    \end{minipage}
    \end{figure}

向下凸或向上凸的函数的重要性质如下面的定理所述：

\begin{blk}
    {定理} 设$f$是$[a,b]$上二次可微函数，并且$f''(x)>0$, $x\in (a,b)$, 那么对于每一个满足$a<x<b$的$x$，有
\[f(x)<f(b)\cdot \frac{x-a}{b-a}+f(a)\cdot \frac{b-x}{b-a}\]
\end{blk}

这一定理有一个明显的几何意义（图2.42），用$\ell(x)$表示右端的一次函数：
\[\ell(x)=f(b)\cdot \frac{x-a}{b-a}+f(a)\cdot \frac{b-x}{b-a}\]
    它在$x=a$和$x=b$的值是$\ell(a)=f(a)$和$\ell(b)=f(b)$, 即与$f(x)$在这些点的值一致．因为$y=\ell(x)$的图象是$y=f(x)$在$[a,b]$上的割线，因此$[a,b]$上的向下凸函数$f$的图象$y=f(x)$位于该区间的割线的下面．





\begin{proof}
设$F(x)=f(x)-f(b)\frac{x-a}{b-a}-f(a)\frac{b-x}{b-a}$，
这里$x\in [a,b]$. 我们需要证明在$[a,b]$上，$F(x)\le 0$. 因为$F(x)$在$[a,b]$上连续，所以$F(x)$在$[a,b]$的某点达到它的最大值，用$c$表示最大值点，我们证明$c$是两端点之一．

假设$a<c<b$，那么$c$也是$F(x)$的极大值点．因此$F'(c)=0$. 又
\begin{align*}
    F'(x)&=f'(x)-\frac{f(b)}{b-a}+\frac{f(a)}{b-a}\\
    F''(x)&=f''(x)>0,\qquad x\in(a,b)
\end{align*}
因此我们得出结论，$F(c)$是$[a,b]$上的最小值，即对于
任意$x\in [a,b]$且$x\ne c$, 有
\[F (x) > F (c) =f (c) -\ell (c)\]
这与$c$是极大值点矛盾，因此，唯一的可能是$c$是一个端点$a$或$b$, 但是$f(a)=\ell(a)$, $f(b)=\ell(b)$, 这表明$F(x)=f(x)-\ell(x)$的最大值是0, 因此在$[a,b]$上除端点以外的所有其它点$x$, 有
\[F (x) =f (x) -f (b) \frac{x-a}{b-a}-f (a) \frac{b-x}{b-a}<0\]
即
\[f (x)<f (b) \frac{x-a}{b-a}+f (a) \frac{b-x}{b-a}\]
这就完成了向下凸函数定理的证明．
\end{proof}

对于向上凸函数有类似的定理，即：区间$[a,b]$上向上凸的函数的图象位于$[a,b]$的割线上面，用式子表示就是
\[f (x)>f (b) \frac{x-a}{b-a}+f (a) \frac{b-x}{b-a},\quad x\in(a,b)\]

\begin{example}
    求证 $\cos\pi x>1-2x,\quad x\in\left(0, \frac{1}{2}\right)$.
\end{example}


\begin{proof}
    设$f(x)=\cos\pi x$, $\ell (x)=1-2x$, 因为
\[f (0) =1=\ell (0) ,\qquad f \left(\frac{1}{2}\right) =0=\ell\left(\frac{1}{2}\right)\]
所以：$y=1-2x$是曲线$y=\cos\pi x$在$\left[0,\frac{1}{2}\right]$上的割线．

又$f'(x)=-\pi \sin\pi x$, $f''(x)=-\pi^2\cos\pi x$, 且当$x\in \left(0,\frac{1}{2}\right)$时，有$f''(x)<0$, 即$f(x)$在$\left(0,\frac{1}{2}\right)$内是向上凸
的．因此，$y=\cos\pi x$在$\left[0,\frac{1}{2}\right]$上的图象是在该区间的割
线$y=1-2x$的上面，从而
\[\cos\pi x>1-2x,\qquad x\in \left(0,\frac{1}{2}\right)\]
\end{proof}

下面给出拐点满足的必要条件．

\begin{blk}
    {定理} 如果$x=c$是曲线$y=f(x)$的拐点，并且$f''(c)$存在，则$f''(c)=0$.
\end{blk}

\begin{proof}
    因为$x=c$是曲线的拐点，即它是曲线$y=f(x)$向上凸和向下凸的分界点，依凸性判定定理，我们可以取$a<c<b$, 使得$f''(x)$在$(a,c)$内的函数值的符号与在$(c,b)$内的函数值的符号相反．

设$g(x)=f'(x),\; x\in (a,b)$, 于是$g'(x)=f''(x)$经过$c$点由负变正或由正变负，因此点$c$是$g(x)$在$(a,b)$内的唯一的极值点，又$g'(c)=f''(c)$存在，根据费马定理
\[g' (c) =f'' (c) =0\]
\end{proof}


\begin{example}
    讨论函数$f(x)=x^3$和$f(x)=x^4$有无拐点．
\end{example}


\begin{solution}
    $f(x)=x^3$有导数$f'(x)=3x^2$, $f''(x)=6x$, 显然$f'(0)=f''(0)=0$, 而且：
当$x>0$时，$f''(x)>0$; 当$x<0$时，$f''(x)<0$, 即曲线$y=x^3$在$(-\infty, 0)$内是向上凸的，而在$(0,+\infty)$区间内是向下凸的．因此，$x=0$是曲线的拐点而不是极值点．

又$f(x)=x^4$有导数$f'(x)=4x^3$, $f''(x)=12x^2$, 显然$f'(0)=f''(0)=0$, 并且对于每一个$x\ne 0$都有$f''(x)>0$, 因此这个曲线$y=x^4$是处处向下凸的，$x=0$不是曲线$y=x^4$的拐点而是一个极小点．
\end{solution}


下面给出判定拐点的充分条件．
    
\begin{blk}
   {定理（拐点的充分条件）} 设$f$在点$c$处连续，在点$c$附近（可以不包括$c$点）具有二阶导数，如果自变量$x$通过$c$时，二阶导数变号，则$x=c$是曲线$y=f(x)$的拐点．
\end{blk}

\begin{proof}
    我们可以取$c$点的一个邻域$(a,b)$使$f''(x)$在
$(a,c)$上的函数值的符号和在$(c,b)$上的相反，依凸性判定定理，$f$经过$c$时将由向上凸变为向下凸或者由向下凸变为向上凸，即$x=c$是曲线$y=f(x)$的拐点．
\end{proof}


\begin{example}
说明$x=0$是曲线$y=x^{\tfrac{1}{3}}$的拐点．
\end{example}

\begin{proof}
    设$f(x)=x^{\tfrac{1}{3}}$，则
\[f'(x)=\frac{1}{3x^{\tfrac{2}{3}}},\qquad f''(x)=\frac{-2}{9x^{\tfrac{5}{3}}}\]
显然$f'(0)$和$f''(0)$都不存在，但是当$x<0$时，$f''(x)>0$, 即曲线$y=x^{\tfrac{1}{3}}$向下凸；当$x>0$时，$f''(x)<0$, 即曲线$y=x^{\tfrac{1}{3}}$向上凸，因此点$x=0$是曲线的拐点．

由$\Lim_{x\to 0}f'(x)=\Lim_{x\to 0}\frac{1}{3x^{\tfrac{2}{3}}}=+\infty$知，曲线在点$x=0$处的切线就是$y$轴本身（见图2.43）．
\end{proof}

\begin{figure}[htp]
    \centering
    \begin{minipage}[t]{0.48\textwidth}
    \centering
\begin{tikzpicture}[>=latex, scale=1]
    \draw[->](-2.5,0)--(2.5,0)node[right]{$x$};
\draw[->](0,-2)--(0,2)node[right]{$y$};
\draw[domain=-1.3:1.3, samples=100, very thick]plot({\x^3},\x)node[above]{$y=x^{\tfrac{1}{3}}$};
\node at (0,0) [above left]{$O$};
\end{tikzpicture}
    \caption{}
    \end{minipage}
    \begin{minipage}[t]{0.48\textwidth}
    \centering
    \begin{tikzpicture}[>=latex, xscale=.6]
       \draw[->](-4.5,0)--(5.5,0)node[right]{$x$};
\draw[->](0,-4)--(0,2)node[right]{$y$};
\draw[->](-4.5,-3)--(5.5,-3)node[right]{$x$};
\node at (0,0)[below left]{$O$};
\draw[domain=-4:5, samples=100, very thick]plot(\x, {sin((\x+pi/3)*180/pi)});
\draw[domain=-4:5, samples=100, very thick]plot(\x, {cos((\x+pi/3)*180/pi)-3});

\foreach \x/\xtext in {1.5*pi/m_3, -.5*pi/m_1, 0/I_1, pi/I_2}
{
   \draw[dashed](\x-pi/3, 0)node[above]{$\xtext$}--(\x-pi/3, -3);
}
\draw[dashed](.5*pi-pi/3, 1)--(.5*pi-pi/3, -3);
\node at (.5*pi-pi/3,0)[above right]{$m_2$};

\draw[domain=-2.5:.5, samples=10, thick]plot(\x, {\x+pi/3});
\draw[domain=1:3.2, samples=10, thick]plot(\x, {-\x+2*pi/3});
\node at (2,1)[right]{$y=f(x)$};
\node at (3,-4)[right]{$y=f'(x)$};
    \end{tikzpicture}
    \caption{}
    \end{minipage}
    \end{figure}

在图2.44中，函数$y=f(x)$的图象上的极值点$m_1$, $m_2$和$m_3$, 对应于函数$y=f'(x)$图象上的零点；$y=f(x)$的图象上的拐点对应于$y=f'(x)$的图象上的极值点．

\begin{ex}
\begin{enumerate}
    \item 讨论下列函数的极值、凸性及拐点，并出函数图象的
    草图；
\begin{multicols}{2}
\begin{enumerate}
    \item $f(x)=\frac{1}{1+x^2}$
    \item $y=(x-1)\sqrt[3]{x^2}$
\end{enumerate}
\end{multicols}
\item 如果函数$f(x)$在区间$(a,b)$内有二阶导数，并且向上凸，求证：对于满足条件$a\le x_1<x_2\le b$的任意两点$x_1,x_2$，恒有
\[f (\lambda x_1+ (1-\lambda) x_2) \ge \lambda f (x_1) + (1-\lambda)f (x_2) ,\quad  0\le\lambda\le 1\]
如果上面的函数在区间$(a,b)$内向下凸，则应有怎样的结论？
\item 求证函数
\begin{enumerate}
    \item $y=\log_ax \; (a>1)$在区间$(0,+\infty)$内是向上凸的；
    \item $f(x)=\sin x$在$[0,\pi]$内是向上凸的；
    \item $\log_a\frac{x_1+x_2}{2}\ge \frac{\log_a x_1+\log_a x_2}{2}\quad (x_1> 0,\; x_2>0)$
    \item $\sin\frac{x_1+x_2}{2}\ge \frac{\sin x_1+\sin x_2}{2}\quad  (0\le x_1,x_2\le \pi)$
\end{enumerate}

\item 如果函数$f(x)$在$(a,b)$内是向上凸的，求证：对于任意$a\le x_1<x_2<x_3\le b$，$\sum^3_{i=1}\lambda_i=1$，有
\[f(\lambda_1x_1+\lambda_2x_2+\lambda_3x_3)\ge \lambda_1f(x_1)+\lambda_2f(x_2)+\lambda_3f(x_3)\]

\item  求证对于任意正数$x_1,x_2,x_3$有
\[\ln\left(\frac{1}{6}x_1+\frac{1}{3}x_2+\frac{1}{2}x_3\right) \ge \ln \left(x_1^{\tfrac{1}{6}}x_2^{\tfrac{1}{3}}x_3^{\tfrac{1}{2}} \right)\]
\item  若$A,B,C$是任意三角形的三个内角，求证
\[\sin A+\sin B+\sin C\le \frac{3\sqrt{3}}{2}\]
又当$A,B,C$为何值时，上面等式成立？
\item  对于$a>0$, $b>0$, $0<\alpha<1$, $0<\beta<1$，$\alpha+\beta=1$，求证：
$a^{\alpha}b^{\beta}\le \alpha a+\beta b$成立．
（提示：利用对数函数是上凸函数）
\item  $f(x)=a_4x^4+a_3x^3+a_2x^2+a_1x+a_0$为具有正系数
的非零多项式，若为偶函数，求证$f(x)$在$(-\infty,+\infty)$内向下凸，且仅有一个极值．
\end{enumerate}
\end{ex}

\subsection{作函数的图象}
给出函数$y=f(x)$, 常常希望画出它的图象来，这对于许多问题是有用的，在初中我们所介绍的用列表法来描绘函数的图象，有很大的不精确性和盲目性，由于所取的点是有穷多个，很可能在所要画出的某两点之间的曲线有比较剧烈的变化，如果我们用比较平滑的曲线弧去连结，于是所得
曲线与实际相差很大．以前各节所介绍的函数性态可以帮助我们断定函数图象的主要特征，因此在取点之前，应先对函数的性态作一番讨论，兹将描绘曲线的一般步骤列出如下：

\begin{enumerate}
\item 确定函数的定义域;
\item 确定曲线关于坐标轴的对称性；
\item 确定曲线与坐标轴的交点（如果不易确定，不必强求）；
\item 确定函数的增减性，极大值与极小值；
\item 确定函数的凸性与拐点；
\item 确定曲线的渐近线；
\item 需要时，还得由曲线方程计算出一些适当的点的坐标；
\item 把上面所得的结果，按自变量大小顺序列入一个表格内，以观察图形的大概形态，然后描绘成图．
\end{enumerate}

在用例子说明如何绘图之前，我们需要补充说明如何求曲线的渐近线．

\begin{blk}
    {定义}假设曲线$y=f(x)$不限制在有限平面内，即其上的点$P$可以趋于无穷远，如果曲线的点$P$到直线$\ell:\; ax+by+c=0$的距离$d$, 当$P$沿曲线趋于无穷远时，
    $d\to 0$，
    那么我们称直线$\ell$为所给曲线的\textbf{渐近线}．
\end{blk}

\subsubsection{水平渐近线及其求法}

设$\ell_1:\; y=\alpha$是曲线$y=f(x)$的一条渐近线，曲线上的$P(x,f(x))$点到$\ell_1$的距离为
\[d=|\alpha-f(x)|\]
因$\ell_1$为渐近线，所以，当$x\to +\infty$或$x\to -\infty$时，应有$d\to 0$, 即
\[\lim_{x\to+\infty}f(x)=\alpha,\qquad \text{或}\qquad \lim_{x\to -\infty}f(x)=\alpha\]
反之，若上式成立，则$y=\alpha$必为$y=f(x)$的一条渐近线．

\subsubsection{铅直渐近线及其求法}

设$\ell_2:\; x=\beta$是曲线$y=f(x)$的一条渐近线，又$P(x,f(x))$为曲线上的任意一点，它到$x=\beta$的距离是
\[d=|\beta-x|\]
因为$\ell_2$是$y=f(x)$的渐近线，那么
\[\lim_{x\to\pm\beta} f (x) =\pm\infty\]
反之，若上式成立，则$\ell_2$为$y=f(x)$的一条渐近线．

例如，对于$y=\frac{1}{x}$, 当$x\to 0$, $y\to \pm\infty$，故$x=0$是$y=\frac{1}{x}$的一条渐近线．

\subsubsection{任意渐近线及其求法}

设直线$\ell_3:\; y=kx+b$是曲线$y=f(x)$的一条渐近线，曲线上任意一点$P(x,f(x))$到直线$\ell_3$的距离是
\[d=\left|\frac{kx+b-f(x)}{\sqrt{1+k^2}}\right|\]
因为$\ell_3$是渐近线，所以当$x\to\pm\infty$时，$d\to 0$, 即
\[kx+b-f (x) \to 0\]
我们必须用上式求出$k,b$, 才能确定$\ell_3$. 由于当$x\to\pm\infty$时，$\frac{1}{x}\to 0$, 所以
\begin{align*}
\lim_{x\to\pm\infty}\left[k+\frac{b}{x}-\frac{f(x)}{x}\right]&=\lim_{x\to\pm\infty}\frac{1}{x}[kx+b-f(x)]\\
&=\lim_{x\to\pm\infty}\frac{1}{x}\cdot \lim_{x\to\pm\infty}[kx+b-f(x)]=0
\end{align*}
由于$b$是常数，故有
\begin{align}
    k&=\lim_{x\to\pm\infty}\frac{f(x)}{x}\qquad \text{(渐近线方向)}\\
b&=\lim_{x\to\pm\infty}[f(x)-kx]
\end{align}

\begin{example}
求函数$f(x)=-3x^5+5x^3$的极值、拐点，并画出它的图象的草图．
\end{example}

\begin{solution}
\begin{enumerate}
\item 曲线的范围与截距：这函数$f$的定义域与值域都是无限的，它的曲线伸展在整个平面内，与$y$轴的交点是原点，而与$x$轴有三个交点，其零点是$0,\pm\sqrt{\frac{5}{3}}$.
\item 对称性：因为$f(-x)=-f(x)$, 所以曲线关于原点对称．
\item 函数的增、减与极值：
\[f' (x) =-15x^4+15x^2=-15x^2 (x-1) (x+1)\]
由此知$f(x)$的驻点是$-1, 0, 1$. $f'(x)$的符号变化的情
况是：
\begin{itemize}
    \item 当$x<-1$时，$f'(x)<0$, 这时$f(x)$递减；
    \item 当$-1<x<1$时, $f'(x)>0$，这时$f(x)$递增；
    \item 当$x>1$时，$f'(x)<0$, 这时$f(x)$递减．
\end{itemize}
因此，$x=-1$是极小点，极小值$f(-1)=-2$; $x=1$是极大点，极大值$f(1)=2$.
\item 函数的凸性与拐点：
\[f''(x)=-60x^3+30x=-60x\left(x+\frac{\sqrt{2}}{2}\right)\left(x-\frac{\sqrt{2}}{2}\right)\]
拐点的可能值是$-\frac{\sqrt{2}}{2},0,\frac{\sqrt{2}}{2}$. 又$f''(x)$的符号变化情况是
\begin{itemize}
    \item 当$x<-\frac{\sqrt{2}}{2}$时，$f''(x)>0$, 此时曲线向下凸；
    \item 当$-\frac{\sqrt{2}}{2}<x<0$时，$f''(x)<0$, 此时曲线向上凸；
    \item 当$0<x<\frac{\sqrt{2}}{2}$时，$f''(x)>0$, 此时曲线向下凸；
    \item 当$x>\frac{\sqrt{2}}{2}$时，$f''(x)<0$, 此时曲线向上凸；
\end{itemize}
因此：$x=-\frac{\sqrt{2}}{2},0,\frac{\sqrt{2}}{2}$的点都是拐点．

\item 列表：今将讨论情况列成下表，并作图（如图2.45）
\end{enumerate}    
\begin{center}\small
    \begin{tabular}{cccccccccccc}
    \hline
    $x$ & &$-1$&&$-\frac{\sqrt{2}}{2}$&&0&&$\frac{\sqrt{2}}{2}$&&1\\
    \hline
    $f'(x)$& $-$ & 0&$+$&$+$&$+$&0&$+$&$+$&$+$&0&$-$\\
    $f''(x)$&$+$&$+$&$+$&0&$-$&0&$+$&0&$-$&$-$&$-$\\
    $f(x)$&下凸&极小值& $\nearrow$  & $-1.24$&上凸 & 0& 下凸 &1.24& 上凸 & 极大值 & $\searrow$  \\
    & $\searrow $ &$-2$&&&$\nearrow $&&$\nearrow $&&$\nearrow $&2\\
    \hline
    \end{tabular}
    \end{center}
    
\begin{figure}[htp]
    \centering
\begin{tikzpicture}[>=latex, xscale=1.5]
\draw[->](-2.5,0)--(2.5,0)node[right]{$x$};
\draw[->](0,-3)--(0,3)node[right]{$y$}; 
\draw[domain=-1.4:1.4, samples=200, very thick]plot(\x, {-3*\x^5+5*\x^3});
\foreach \x in {-1,-2,1,2}
{
    \draw(\x,.1)--(\x,0)node[below]{$\x$};
}
\node at (0,0)[below right]{$O$};
\draw[dashed](-.707,-1.24)--(-.707,0)node[above]{$-\tfrac{\sqrt{2}}{2}$};
\draw[dashed](.707,1.24)--(.707,0)node[below]{$\tfrac{\sqrt{2}}{2}$};
\node at (1.29,0)[above]{$\sqrt{\tfrac{5}{3}}$};

\end{tikzpicture}
    \caption{}
\end{figure}
\end{solution}


\begin{example}    
描绘函数$y=\frac{(x-3)^2}{4(x-1)}$的图象．
\end{example}

\begin{solution}
    这个函数在$x=1$这点没有定义，并且，当$x\to 1^-$时，$y\to -\infty$; 当$x\to 1^+$时，$y\to +\infty$, 所以$x=1$是曲线的铅直渐近线，曲线在$x=1$这点断开．
    
求导数，得
\[y'= \frac{(x+1) (x-3)}{4 (x-1)^2}\]

当$x<-1$时，$y'>0$；当$-1<x<3$（$x\ne1$）时，$y'<0$；当$x>3$时，$y'>0$．因此，$x=-1$是极大点；$x=3$是极小点．又极大值$f(-1)=-2$; 极小值$f(3)=0$

求二阶导数，得
\[y''=\frac{2}{(x-1)^3}\]
\begin{itemize}
    \item 当$x<1$时，$y''<0$, 曲线向上凸；
    \item 当$x>1$时，$y''>0$, 曲线向下凸，但是曲线在$x=1$这一
    点断开，因此$x=1$不是曲线的拐点．
\end{itemize}

又因为
\[y=\frac{(x-3)^2}{4(x-1)}=\frac{1}{4}(x-5)+\frac{1}{x-1}\]
\[\lim_{x\to \infty}\left[\frac{(x-3)^2}{4(x-1)}-\frac{1}{4}(x-5)\right]=\lim_{x\to\infty}\frac{1}{x-1}=0\]
这就是说曲线以$y=\frac{1}{4}x-\frac{5}{4}$为渐近线．

综合起来得到下面结果：
\begin{center}
    \begin{tabular}{cccccccccccc}
    \hline
    $x$ & &$-1$&&0&&1&&3&\\
    \hline
    $y'$& $+$ & 0&$-$&$-$&$-$&$\x$&$-$&0&$+$\\
    $y''$&$-$&$-$&$-$&$-$&$-$&$\x$&$+$&$+$&$+$\\
    $y$&$\nearrow$ &极大值& $\searrow$  & $\searrow$&$\searrow$ & $\x$& $\searrow$ & 极小值 & $\nearrow$  \\
    &  &$-2$&&$-\frac{9}{4}$&$-\infty $&&$+\infty$&0 &\\
    \hline
    \end{tabular}
    \end{center}

画出具体图形如图2.46.
\begin{figure}[htp]
    \centering
\begin{tikzpicture}[>=latex]
    \draw[->](-3.5,0)--(6.5,0)node[right]{$x$};
    \draw[->](0,-5)--(0,3)node[right]{$y$}; 
\draw[domain=-3.5:6.5, samples=10, thick]plot(\x, {0.25*\x -1.25});
\draw[domain=-3:.7, samples=100, very thick]plot(\x, {(\x-3)^2/(4*(\x-1))});
\draw[domain=1.3:6, samples=100, very thick]plot(\x, {(\x-3)^2/(4*(\x-1))});
\node at (0,0)[below left]{$O$};
\draw[dashed](1,-5)--(1,3);
\foreach \x in {1,2,3,4}
{
    \draw(\x,0)node[below]{$\x$}--(\x,.1);
}
\foreach \x in {-1,-2}
{
    \draw(0,\x)node[left]{$\x$}--(.1,\x);
}

\node at (2,-1)[ right]{$y=\frac{1}{4}x-\frac{5}{4}$};
\node at (3,1)[ right]{$y=\frac{(x-3)^2}{4(x-1)}$};

\end{tikzpicture} 
    \caption{}
\end{figure}
\end{solution}



\begin{ex}
\begin{enumerate}
    \item 讨论下列函数的极值、凸性和拐点, 并画出函数图象.
\begin{multicols}{2}
\begin{enumerate}
\item  $y=\frac{6(x-1)}{x^{2}+3}$
\item  $y=(x-1)^{3}-3(x-1)$
\item  $y=\frac{x^{4}}{2}-x^{2}$
\item  $y=1+4 x^{3}-3 x^{4}$
\item  $y=\sqrt{(x+2)^{2}\left(16-x^{2}\right)}$
\item  $y=x^{x}\quad (x>0.1)$
\item  $y=x e^{-x}$
\end{enumerate}
\end{multicols}

    \item 作下列函数的图象, 并画出渐近线及在拐点处的切线:
    \begin{multicols}{2}
        \begin{enumerate}
            \item $y=x+\frac{1}{x}$
            \item $y=\frac{x^{2}}{1+x^{3}}$
            \item $y=\frac{2 x}{1-x^{2}}$
            \item $y=\frac{x^{3}}{x^{2}-4}$
            \item $y=\frac{(x-5)(x+3)}{(x+2)^{2}}$
\end{enumerate}
\end{multicols}
\end{enumerate}
\end{ex}

\subsection*{习题2.3}
\begin{enumerate}
    \item 求函数$f(x)=x^3+ax^2-3bx+c$的图象关于原点对称时，$f(x)$在区间$-2\le x\le 2$上的最小值，其中$b>0$.
    
    \item 函数$f(x)=3x^2-ax^3$在区间$0\le x\le 2$上的最小值
    是$-4$, 求
\begin{enumerate}
    \item $a$的值
    \item $f(x)$在区间$0\le x\le 2$的最大值$M$
\end{enumerate}
\item 设$a,b$为实数，求函数$y=x^3-3ax^2+bx+1$在区间$0\le x\le 1$上单调增加时，点$(a,b)$存在范围，并用图表示出来．

\item 设$f(x)=x^3-5x^2+3x+a$, 在怎样的情况下对于
$x$的正值，总有$f(x)>0$成立？
\item 设$a$是常数，对于函数$f(x)=4x^3-3x+a$
\begin{enumerate}
    \item 求$y=f(x)$的极大值和极小值；
    \item 确定当方程$f(x)=0$有三个不相等的实根时，$a$值的范围．
\end{enumerate}

\item 讨论函数$y=\frac{x^3}{2 (x+2)}$的极值，凸向，拐点，渐近线，
并作出它的图象．

\item 求证$\pi<\frac{\sin\pi x}{x(1-x)}\le 4$.
\item 对于已知半径为$R$的球，求具有最小体积的外切圆锥
的体积．
\item 一长方体的对角线长$\sqrt{11}$cm，全面积为$14{\rm cm}^2$, 求它的体积的最大值，并求体积为最大时三边的长．
\item 一浮标由三部分组成，一个圆筒与两个相等的圆
锥，其中每一个圆锥的高等于圆筒的高，问当表面积一定时，什么样的形状会有最大的体积．
\end{enumerate}     



\chapter{求函数从$a$到$b$的和与积分}
函数有两个基本性质——变率与和，在前一章，我们研究了函数的瞬时变率的概念，即
\[f'(x_0) =\lim_{h\to 0} \frac{f (x_0+h) -f (x_0)}{h}\]
以及它的应用．在本章我们将研究对于一个给定函数$f(x)$“求从$a$到$b$的和”这个概念．在给出定义之前，我们先举几个实例来看一看．

\paragraph{速率与距离} 
一列行驶中的火车，它的行进速率$v$是时间$t$的函数，即
$v=f(t)$, 如图3.1, 我们可以从速率表读得当时的速率，很自然地，我们想知道火车在从$t=a$到$t=b$这一段时间间隔内一共走了多少路程．从函数的观点看，所谓求在时间$[a,b]$内所走过的距离就是求速率函数$v=f(t)$由$t=a$到$t=b$的和．

\begin{figure}[htp]
    \centering
\begin{tikzpicture}[>=latex]
\draw[->](-.5,0)--(6,0)node[right]{$t$};
\draw[->](0,-.5)--(0,3)node[right]{$v=f(t)$};
\draw[very thick](0,0)--(1,1.5)--(1.6,1.5);
\draw[dashed](1,1.5)--(1,0)node[below]{$a$};
\node at (0,0)[below left]{$O$};
\draw plot[smooth] coordinates{
(1.6,1.5)(2.5,2.5)(3.5,2.2)(4,1.2)(4.3,.8)(4.8,.6)(5,0)
};
\draw[dashed](3.5,2.2)--(3.5,0)node[below]{$b$};

\end{tikzpicture}
    \caption{}
\end{figure}

让我们用$D(f,[a,b])$
表示所走过的距离，这个记号强调$D$依赖于$f$和区间$[a,b]$．

\paragraph{变力所作的功}

假定某物体在一个平行于$OX$轴的力$P$的作 用下沿直线
$OX$运动，力$P$的方向与物体运动的方向一致，并且力的大小随离开$O$点的距离而改变，即变力$P$是所在点的横坐标$x$的函数$P=P(x)$, 如图3.2．假定物体在这个变力$P$作用之下，从直线$OX$的一点$a$移到另一点$b$, 那么力$P$所作的功就是变力函数$P(x)$由$x=a$到$x=b$的和，我们用
$W (P, [a,b])$
表示力$P$所作的功，它表明$W$依赖于$P(x)$和$[a,b]$．

\begin{figure}[htp]
    \centering
    \begin{minipage}[t]{0.48\textwidth}
    \centering
    \begin{tikzpicture}[>=latex, scale=1]
\draw[->](-.5,0)--(4,0)node[right]{$x$};
\draw[->](0,-.5)--(0,3)node[right]{$P(x)$};
\draw[domain=0:3.5, samples=100, very thick]plot(\x, {3/(\x+1)-.5});
\draw(1,0)node[below]{$a$}--(1,1);
\draw(3,0)node[below]{$b$}--(3,.25);
\node at (0,0)[below left]{$O$};
    \end{tikzpicture}
    \caption{}
    \end{minipage}
    \begin{minipage}[t]{0.48\textwidth}
    \centering
    \begin{tikzpicture}[>=latex, scale=1]
\draw[->](-.5,0)--(4,0)node[right]{$x$};
\draw[->](0,-.5)--(0,3)node[right]{$y$};
\draw[domain=0:3.5, samples=100, very thick]plot(\x, {.3*(\x-2.2)^2+.5})node[above right]{$y=f(x)$};
\node at (0,0)[below left]{$O$};
\draw(.5,0)node[below]{$a$}--(.5,1.367)node[above right]{$M_0$};
\draw(3,0)node[below]{$b$}--(3,.692)node[above]{$M$};
\fill[pattern=north east lines, domain=.5:3, samples=100, very thick]plot(\x, {.3*(\x-2.2)^2+.5})--(3,.692)--(3,0)--(.5,0)--(.5,1.367);


    \end{tikzpicture}
    \caption{}
    \end{minipage}
\end{figure}

\paragraph{曲边梯形的面积}

令$y=f(x)$是函数$f$的图象，表示一条曲线，如图3.3. 我们要求曲线上的一段弧$M_0M$ 与其两端的纵坐标线及$x$轴上的线段$[a,b]$所围成的图形的面积．这样的图形（它有三条边是直线，其中两条互相平行，第三条与前两条
互相垂直，而第四条边是曲线）叫做曲边梯形．显然曲边梯形的面积$A$依赖于$y=f(x)$和$x$轴上的线段$[a,b]$, 记这一面积为$A (f, [a,b])$．

在上一章，我们利用函数$f(x)$和它的图象之间的对应关系，就可以把函数的变率和它的图象的切线的斜率相对应地一并分析讨论，这样，一方面可以把函数的变率这种“数量”的概念用切线来形象化，便于想象；而另一方面又可以把“切线”这种几何概念数量化，便于计算，在这一章，我们也要利用函数$f(x)$和它的图象之间的对应关系，把函数$f(x)$由$x=a$到$x=b$的“和”与曲线$y=f(x)$的曲边梯形的“面积”相对应地一并分析讨论，并且说明函数的“求从$a$到$b$的和”恰好对应于求曲边梯形的面积．这也就是为什么把函数“求从$a$到$b$的和”这种基本运算叫做积分的道理．

\section{{``}和{''}与{``}面积{''}}

对于任意曲线围成的图形，我们还没有规定它的“面积”的意义，和它密切相关的“函数从$a$到$b$的和”的概念至今也没有明确的解析的定义．在这一节我们要把这两个概念由“直观的定性理解”推进到“数量化的定量定义”．唯有确立了它们的“解析的定义”，它们才真正地成为能算好用的量．

\subsection{{``}和{''}与{``}面积{''}的基本性质}

现在让我们先从“函数的和”与“曲线形的面积”的直观内涵来分析一下它们分别所应有的基本性质．

\subsubsection{{``}曲线形的面积{''}的基本性质}

从面积的直观内涵容易看出下列两点：
\begin{enumerate}
\item (单调性）设区域$R_1$包含在$R_2$之内，即$R_1\subseteq R_2$, 则：
$R_1$的面积$\le R_2$的面积．（图3.4）
\item (可加性）设区域$R$可以用一条曲线分割成两块区
域$R_1+R_2$, 则有：
$R$的面积$=R_1$的面积$+R_2$的面积．（图3.5）
\end{enumerate}

\begin{figure}[htp]
    \centering
    \begin{minipage}[t]{0.48\textwidth}
    \centering
    \begin{tikzpicture}[>=latex, scale=1]
\draw[very thick](0,0) ellipse [x radius=2, y radius =1];
\draw[very thick](-.5,0) ellipse [x radius=1, y radius =.6];
\node at (0,1.3){$R$};
\node at (-.5,0){$R_1$};
\node at (1.3,0){$R_2$};
    \end{tikzpicture}
    \caption{}
    \end{minipage}
    \begin{minipage}[t]{0.48\textwidth}
    \centering
    \begin{tikzpicture}[>=latex, scale=1]
\draw[very thick](0,0) ellipse [x radius=2, y radius =1];
\node at (0,1.3){$R$};
\node at (-1,0){$R_1$};
\node at (1,0){$R_2$};
\draw[very thick](0,1) to [bend left](.3,.6)to [bend left](0,-.3)to [bend right](0,-1);
    \end{tikzpicture}
    \caption{}
    \end{minipage}
\end{figure}

\subsubsection{{``}函数的和{''}的基本性质}

同样地，从“函数从$a$到$b$的和”的直观内涵容易看出下列两个基本性质，即

\begin{blk}{性质1：单调性}
设函数$f(t)$和$g(t)$在$a\le t\le b$上，$f(t)\le g(t)$恒成立，那么：

“$f(t)$由$t=a$到$t=b$的和”$\le $“$g(t)$由$t=a$到$t=b$
的和”．
\end{blk}

例如，两部车子在$t=a$到$t=b$的时间内，甲车的速率$f(t)\le $乙车的速率$g(t)$, 则甲车在上述时间内所经过的里程$\le $乙车在上述时间内所经过的里程．

\begin{blk}{性质2：可加性}
设$a<b<c$，那么，有：

“$f(t)$由$t=a$到$t=b$的和”$+$“$f(t)$由$t=b$到$t=c$的和”$=$“$f(t)$由$t=a$到$t=c$的和”．
\end{blk}

上述基本性质是一目了然的，让我们先用来说明一些简单的基本事实．

\begin{example}
    设函数$y=f(t)=k$（常数），则由和的直观内涵，显然应有：
\[\text{常数函数从$a$到$b$的和}=k(b-a)\]

例如，以等速率每小时$k$公里行驶的火车从$t=a$小时到$t=b$小时内所走的总里程应该是$k(b-a)$公里．

现在让我们利用函数的图象来观察上述数值$k(b-a)$的几何意义．

\begin{figure}[htp]
    \centering
    \begin{minipage}[t]{0.48\textwidth}
    \centering
    \begin{tikzpicture}[>=latex, scale=.9]
\draw[->](-.5,0)--(5,0)node[right]{$t$};
\draw[->](0,-.5)--(0,3)node[right]{$y$};
\fill[pattern=north east lines](1,0) rectangle (4,1.5);
\draw[dashed](0,1.5)--(1,1.5)--(1,0)node[below]{$t=a$};
\draw[dashed](4,1.5)--(4,0)node[below]{$t=b$};
\node at (0,0)[below left]{$O$};
\draw[very thick](1,1.5)--(4,1.5);
\node at (0,.75)[left]{$k>0$};
\node at (2.5,.75){\Large $+$};

    \end{tikzpicture}
    \caption{}
    \end{minipage}
    \begin{minipage}[t]{0.48\textwidth}
    \centering
    \begin{tikzpicture}[>=latex, scale=.9]
\draw[->](-.5,0)--(5,0)node[right]{$t$};
\draw[->](0,-2)--(0,1.5)node[right]{$y$};
\fill[pattern=north east lines](1,0) rectangle (4,-1.5);
\draw[dashed](0,-1.5)--(1,-1.5)--(1,0)node[above]{$t=a$};
\draw[dashed](4,-1.5)--(4,0)node[above]{$t=b$};
\node at (0,0)[below left]{$O$};
\draw[very thick](1,-1.5)--(4,-1.5);
\node at (0,-.9)[left]{$k<0$};
\node at (2.5,-.75){\Large $-$};
    \end{tikzpicture}
    \caption{}
    \end{minipage}
\end{figure}

\begin{enumerate}
    \item 当$k>0$时，函数$f(t)=k$的图象在$t$轴上方，$k(b-a)$对应于一个矩形的面积，如图3.6中阴影部分的面积．
    \item 当$k<0$时，函数$f(t)=k$的图象在$t$轴下方，$k(b-a)$对应于一个矩形的面积的负值，如图3.7中阴影部分的面积的负值．
\end{enumerate}

假如我们把$t$轴之上的面积定义为正的，而把$t$轴之下的面积定义为负的，则当$y=k$（常数）时，常数函数从$a$到$b$的和$=k(b-a)=$上述的有号面积．
\end{example}

\begin{example}
设$y=f(t)$是一个阶梯函数，即分段地是常数函数：
设
\[a=t_0<t_1<t_2<\cdots<t_{i-1}<t_i<\cdots <t_n=b\]
当$t_{i-1}\le t<t_i$时，$f(t)=k_i,\; i=1,2,\ldots,n$．由性质2，我们可以分段地用常数函数求和，就可以得到阶梯函数$f(t)$从$a$到$b$的和．

阶梯函数$f(t)$从$a$到$b$的和
\[=k_1 (t_1-t_0) +k_2 (t_2-t_1) +\cdots +k_n (t_n-t_{n-1})\]

作出阶梯函数的图象（图3.8），上述函数和就等于下列逐段由矩形并起来的区域的“有号面积”．
\begin{figure}[htp]
    \centering
\begin{tikzpicture}[>=latex, yscale=.7]
\draw[->](-5,0)--(7,0)node[right]{$t$};
\draw[->](0,-2)--(0,4)node[right]{$y=f(t)$};
\node at (0,0)[above right]{$O$};

\foreach \x/\xtext in {-3/t_1, -1/t_2, 1.5/t_3, 2.5/t_4, 3.5/t_5, 4.3/t_6, 5.5/t_{n-1},-4.5/{t_0=a},6.5/{t_n=b}}
{
    \node at (\x,0)[below]{$\xtext$};
}

\draw[pattern=north east lines, thick](-4.5,0) rectangle(-3,1.5);
\draw[pattern=north east lines, thick](-3,1) rectangle(-1,0);
\draw[pattern=north east lines, thick](-1,0) rectangle(1.5,-1.5);
\draw[pattern=north east lines, thick](1.5,0) rectangle(2.5,1);
\draw[pattern=north east lines, thick](2.5,0) rectangle(3.5,2);
\draw[pattern=north east lines, thick](3.5,0) rectangle(4.3,-1);
\draw[pattern=north east lines, thick](5.5,0) rectangle(6.5,1);
\foreach \x in {-3.75, -2, 2, 3, 6}
{
    \node at (\x,.5){$+$};
}
\node at (.75,-.9){$-$};
\node at (3.9,-.65){$-$};
\node at (4.8,.5){…};\node at (4.8,-.5){…};

\end{tikzpicture}
    \caption{}
\end{figure}
\end{example}

\subsection{逼近法求和}

对于一个比较一般的函数，例如$y=kt+c$, $y=t^2$, $y=\sin t$等等，我们必须给函数的和下一个适当的定义，同时要提供计算这个和的方法，这里我们将要再一次地用逼近法的观点去解答上述问题．回顾上一章讨论变率时，我们的基本想法是用折线函数去逼近一般的平滑函数，所以折线函数是讨论函数变率的简单好用的基本函数，因为它们的变率十分简单，分段地是个常数而且无限逼近函数在某一点的变率．由上一段的讨论；我们知道阶梯函数的“从$a$到$b$的和”是十分简明的，它们在“函数求从$a$到$b$的和”这方面是不是也扮演着
这种既简单又基本的角色呢？这就得看一看能否用阶梯函数的从$a$到$b$的和去无限逼近一般性的“函数从$a$到$b$的和”了！

现在让我们使用几个简单的实例来试试看．

\begin{example}
假设物体以初速度$u$, 加速度$a>0$在直线上运动，于是物体在任何时刻$t$的速度是$v=f (t) =u+at$, 求物体从$t=0$到$t=T$时所经过的距离，即求速度函数$f(t)$在时间间隔$[0,T]$上的和．
\end{example}

\begin{analyze}
    一个自然的作法是用两个阶梯函数$g_n(t)$和$G_n(t)$把上述函数$f(t)=u+at$夹逼在中间，即
    \[g_n(t)\le f(t)\le G_n(t)\]
对于任何$0\le t\le T$都成立，那么由性质1，$g_n(t)$从0到$T$的和$\le f(t)$从$0$到$T$的和$\le G_n(t)$从0到$T$的和．

如果能使$g_n(t)$从0到$T$的和与$G_n(t)$从0到$T$的和，当$n\to\infty$时的极限相同，则由上述夹逼关系就可以看出$f(t)$从0到$T$的和必须等于这个共同极限．
\end{analyze}

\begin{figure}[htp]
    \centering
\begin{tikzpicture}[>=latex]
\draw[->](-.5,1)--(8,1)node[right]{$t$};
\draw[|<->|](-.2,1)--node[fill=white]{$u$}(-.2,2);
\draw[->](0,0)--(0,6)node[right]{$v=f(t)=u+at,\; (0\le t\le T)$};
\node at (0,1)[below left]{$O$};
\draw[ultra thick](0,2)node[above left]{$A$}--(7,5.5)node[above]{$B$};
\draw[<->](7.3,1)node[below]{$C$}--node[fill=white]{$u$}(7.3,2);
\draw[|<->|](7.3,5.5)--node[fill=white]{$aT$}(7.3,2);
\foreach \x in {1,2,3,...,6,7}
{
    \ifnum\x<7 \node at (\x,1)[below]{$t_{\x}$};\fi
    \draw[dashed](\x,1)--(\x,2+\x*.5);
    \draw[very thick](\x-1, 1.5+\x*0.5)--(\x, 1.5+ \x*0.5)--(\x, 2+\x*0.5);
    \draw[thick](\x-1, 1.5+\x*0.5)--(\x-1, 2+\x*0.5)--(\x, 2+ \x*0.5);
}
\draw[dashed](0,2)--(7,2);
\draw(7,0)--(7,1);
\draw[<->](0,.2)--node[fill=white]{$T$}(7,.2);
\end{tikzpicture}
    \caption{}\label{fig:chapter3_sec1_line_segmentation}
\end{figure}

下面就是把上述想法付诸实践的具体做法之一．

\begin{solution}
\begin{enumerate}
    \item 把时间间隔$n$等分，取分点$t_i=\frac{i}{n}T,\; i=0,1, 2,\ldots,n$, 于是在每个分点$t_i=\frac{i}{n}T$处的速度分别是
\[f (t_0) =f (0) =u,\;  f(t_1) =u+at_1,\; f (t_2) =u+at_2,\ldots,f (t_n) =f (T) =u+aT\]
\item 用阶梯函数近似代替$v=f(t)=u+at$. 

假设物体在时间的每个小区间$[t_i,t_{i+1}],\; i=0,1, 2,\ldots,(n-1)$内，以小区间起点处的速度作匀速运动，我们得到一个速度的阶梯函数
\[g_n (t)=f(t_i)=u+at_i,\quad t_i\le t<t_{i+1},\quad i=0,1,2,\ldots,(n-1)\]
因为$f(t)=u+at$是严格递增的，所以$g_n(t)$有以下性质：
\[g_n (t)=f(t_i)\le f(t),\quad t_i\le t<t_{i+1},\quad i=0,1,2,\ldots,(n-1)\]
从而，对于任何$0\le t\le T$，有$g_n(t)\le f(t)$.

假设物体在时间的每个小区间$[t_i,t_{i+1}],\; i=0,1, 2,\ldots,(n-1)$内，以小区间终点处的速度作匀速运动，我们得到另一个速度的阶梯函数
\[G_n(t)=f(t_{i+1})=u+at_{i+1},\quad t_i< t\le t_{i+1},\quad i=0,1,2,\ldots,(n-1)\]

由于$f(t)=u+at$是递增的，所以$G_n(t)$有以下性质
\[G_n (t)=f(t_{i+1})\ge f(t),\quad t_i< t\le t_{i+1},\quad i=0,1,2,\ldots,(n-1)\]
从而对于任何$0\le t\le T$，有
\[G_n(t)\ge f(t)\]
这样，我们就得到了$v=f(t)$的夹逼阶梯函数：
\[ g_n(t) \le f (t) \le G_n (t),\quad 0\le t\le T\]

（在图\ref{fig:chapter3_sec1_line_segmentation}中，绘出了$f(t)=u+at$的图象和当$n=7$时阶梯函数$g_7(t)$和$G_7(t)$的图象.）

\item 求阶梯函数的总和

应用基本性质2, 我们得到$g_n(t)$从0到$T$的和
\[f(t_0)(t_1-t_0)+f(t_1)(t_2-t_1)+\cdots+f (t_{n-1}) (t_n-t_{n-1})\]

$\because\quad t_1-t_0=t_2-t_1=\cdots=t_n-t_{n-1}=\frac{T}{n}$

因此：
\begin{align*}
\text{$g_n(t)$从0到$T$的和}&=\left[f(t_0)+f(t_1)+\cdots+f(t_{n-1})\right]\cdot \frac{T}{n}\\
&=\left[u+(u+at_1)+(u+at_2)+\cdots +(u+at_{n-1})\right]\cdot \frac{T}{n}\\
&=\left[nu+a(t_1+t_2+\cdots+t_{n-1})\right]\cdot \frac{T}{n}\\
&=\left[nu+a\left(\frac{T}{n}+\frac{2T}{n}+\cdots+\frac{(n-1)T}{n}\right)\right]\cdot \frac{T}{n}\\
&=uT+a\cdot [1+2+\cdots+(n-1)]\left(\frac{T}{n}\right)^2\\
&=uT+a\cdot \frac{n(n-1)}{2}\cdot \frac{T^2}{n^2}\\
&=uT+\frac{1}{2}aT^2\left(1-\frac{1}{n}\right)
\end{align*}

\begin{align*}
    \text{$G_n(t)$从0到$T$的和}&=f(t_1)(t_1-t_0)+f(t_2)(t_2-t_1)+\cdots+f(t_n)(t_n-t_{n-1})  \\
   &= \left[f(t_1)+f(t_2)+\cdots+f(t_{n})\right]\cdot \frac{T}{n}\\
    &=\left[(u+at_1)+(u+at_2)+\cdots +(u+at_{n})\right]\cdot \frac{T}{n}\\
    &=\left[nu+a\left(\frac{T}{n}+\frac{2T}{n}+\cdots+\frac{nT}{n}\right)\right]\cdot \frac{T}{n}\\
    &=uT+a\cdot (1+2+\cdots+n)\cdot \left(\frac{T}{n}\right)^2\\
    &=uT+\frac{1}{2}aT^2\left(1+\frac{1}{n}\right)
    \end{align*}
综合上述计算和基本性质1, 即得：$g_n(t)$从0到$T$的和$\le f(t)$从0到$T$的和$\le G_n(t)$从0到$T$的
和，即
\[uT+\frac{1}{2}aT^2\left(1-\frac{1}{n}\right)\le f(t)
\text{从0到$T$的和}\le uT+\frac{1}{2}aT^2\left(1+\frac{1}{n}\right)\]

\item 求阶梯函数和式的极限．

因为
\begin{align*}
    \lim_{n\to\infty}\left\{uT+\frac{1}{2}aT^2\left(1-\frac{1}{n}\right)\right\}&=\lim_{n\to\infty}\left\{uT+\frac{1}{2}aT^2\left(1+\frac{1}{n}\right)\right\}\\
    &=uT+\frac{1}{2}aT^2
\end{align*}
所以，这个共同的极限值$uT+\frac{1}{2}aT^2$就是物体在变速$v=f(t)=u+at$运动下，从$t=0$到$t=T$所走的距离，也就是函数$v=f(t)=u+at$从$t=0$到$t=T$的和．
\end{enumerate}

从几何上看，上述极限值$uT+\frac{1}{2}aT^2=\frac{[u+(u+aT)]T}{2}$就是在图\ref{fig:chapter3_sec1_line_segmentation}中的梯形$OABC$的面积．
\end{solution}    

\begin{example}
    设函数$f(x)=x^2$, $a=0$, $b>0$, 求$f(x)$由0到$b$的和，相应地，求抛物线$y=x^2$在线段$[0,b]$上所盖的曲边三角形$OBC$的面积（图\ref{fig:chapter3_sec1_y=x^2}）.
\end{example}

\begin{figure}[htp]
    \centering
\begin{tikzpicture}[>=latex, scale=.5]
\draw[->](-5,0)--(5,0)node[right]{$x$};
\draw[->](0,-1.5)--(0,11)node[right]{$y$};
\draw[domain=-3.25:3.25, samples=100, very thick]plot(\x, {\x*\x});
\node at (0,0)[below left]{$O$};
\foreach \x in {.5,1,...,3}
{
    \draw[thick](\x-.5,0) rectangle (\x, \x*\x);
}
\foreach \x in {.5,1,...,2.5}
{
    \draw(\x,\x*\x) -- (\x+.5,\x*\x);
}
\node at (3,9)[right]{$C$};
\node at (3,0)[below]{$B$};
\node at (-2.5,6)[left]{$y=x^2$};
\end{tikzpicture}
    \caption{}\label{fig:chapter3_sec1_y=x^2}
\end{figure}

\begin{solution}
\begin{enumerate}
    \item 把线段$[0,b]$ $n$等分，等分点的坐标是
   $ x_i=\frac{b}{n}i\;  (i=0, 1, 2,\ldots,n)$ 
   \item 作阶梯函数$g_n(x)$和
    $G_n (x)$.

    根据$f(x)=x^2$在$x>0$时递增，定义
\begin{align*}
g_n(x)&=f(x_i)=x^2_i=\left(\frac{b}{n}\right)^2i^2,\quad x_i\le x<x_{i+1}\\
G_n(x)&=f(x_{i+1})=x^2_{i+1}=\left(\frac{b}{n}\right)^2(i+1)^2,\quad x_i<x\le x_{i+1}\\
\end{align*}
其中，$i=0,1,2,\ldots,(n-1)$．于是，$g_n(x)\le f(x)\le G_n(x),\quad 0\le x\le b$

\item 求阶梯函数的和
\begin{align*}
s_n&=g_n(x)\text{从0到$b$的和}\\
&=f\left(x_{0}\right) \cdot \frac{b}{n}+f\left(x_{1}\right) \cdot \frac{b}{n}+\cdots+f\left(x_{n-1}\right) \cdot \frac{b}{n}\\
&=\left\{f\left(x_{0}\right)+f\left(x_{1}\right)+\cdots+f\left(x_{n-1}\right)\right\} \cdot \frac{b}{n} \\
&=\left\{0^{2}+1^{2}+2^{2}+\cdots+(n-1)^{2}\right\} \cdot\left(\frac{b}{n}\right)^3\\
&=\frac{1}{6}n(n-1)(2 n-1) \cdot\left(\frac{b}{n}\right)^{3} =\frac{b^3}{6}\left(1-\frac{1}{n}\right)\left(2-\frac{1}{n}\right) 
\end{align*}
\begin{align*}
   S_{n}= G_{n}(x) \text { 从0到$b$的和}
   &= \left\{f\left(x_{1}\right)+f\left(x_{2}\right)+\cdots+f\left(x_{n}\right)\right\} \cdot \frac{b}{n} \\
   &=\left\{1^2+2^{2}+\cdots+n^{2}\right\} \cdot\left(\frac{b}{n}\right)^{3} \\
   &=\frac{1}{6}n(n+1)(2 n+1) \cdot\left(\frac{b}{n}\right)^{3} \\
   &=\frac{b^3}{6}\left(1+\frac{1}{n}\right)\left(2+\frac{1}{n}\right) 
\end{align*}

综合上述计算和性质1，有：
\begin{align*}
\frac{b^3}{6}\left(1-\frac{1}{n}\right)\left(2-\frac{1}{n}\right)&=s_n\le \begin{bmatrix}
\text{函数$f(x)=x^2$从0到$b$的和，}\\
\text{相应的曲边三角形$OBC$的面积}
\end{bmatrix}\\
&\le S_n=\frac{b^3}{6}\left(1+\frac{1}{n}\right)\left(2+\frac{1}{n}\right)
\end{align*}

\item 令$n\to\infty$，求阶梯函数和式的极限．因为
\begin{align*}
    \frac{b^3}{6} \left(1-\frac{1}{n}\right)\left(2-\frac{1}{n}\right)\to \frac{b^3}{3}\qquad (n\to\infty)\\
    \frac{b^3}{6} \left(1+\frac{1}{n}\right)\left(2+\frac{1}{n}\right)\to \frac{b^3}{3}\qquad (n\to\infty)\\
\end{align*}
所以$\frac{1}{3}b^3$是能被$s_n$、$S_n$左、右夹逼的唯一实数，而所求的“$f(x)=x^2$从0到$b$的和”以及“曲边三角形$OBC$的面积，”这两个值都是夹逼在$s_n$和$S_n$中间的实数，故它们都等于$\frac{1}{3}b^3$.
\end{enumerate}
\end{solution}

\begin{blk}
    {推论} $f(x)=x^2$从$a$到$b$的和$=\frac{1}{3}b^3-\frac{1}{3}a^3\quad  (b> a> 0)$.
\end{blk}


\begin{example}
设$y=f(x)$是一个定义在$[a,b]$上的递增连续函数，试说明它从$a$到$b$的和是可以确定的．
\end{example}

\begin{figure}[htp]
    \centering
\begin{tikzpicture}[>=latex]
\draw[->](-.5,0)--(5,0)node[right]{$x$};
\draw[->](0,-.5)--(0,4)node[right]{$y$};
\draw[domain=.5:4, samples=50, very thick]plot(\x, {.2*\x*\x+.5});
\draw[dashed](1,.7)--(1,0)node[below]{$a$};
\draw[dashed](4,3.7)--(4,0)node[below]{$b$};
\draw[dashed](2,1.3)--(2,0);
\draw[dashed](3,2.3)node[left]{$y=f(x)$}--(3,0);
\node at (0,0)[below left]{$O$};
\end{tikzpicture}
    \caption{}\label{fig:chapter3_sec1_monotone f}
\end{figure}

\begin{solution}
\begin{enumerate}
    \item     取$[a,b]$之间的$n$等分点，将它分成$n$段，即有
\[x_0=a<x_1<x_2<\cdots<x_n=b\]
它的第$i$段是$[x_{i-1},x_i]$，且$x_i-x_{i-1}=\frac{b-a}{n}$，$x_i=a+\frac{i}{n}(b-a)$

\item 定义$f(x)$的上、下夹逼阶梯函数如下：
\[
\begin{split}
    g_n(x)&=f(x_{i-1}),\qquad x_{i-1}\le x\le x_i \\
    G_n(x)&=f(x_i),\qquad x_{i-1}\le x\le x_i 
\end{split} \quad (i=1,2,3,\ldots,n)
\]
由$f(x)$的递增性，对于任何一个$x\in[a,b]$，有
\[g_1(x)\le g_2(x)\le\cdots \le g_n(x)\le\cdots \le f(x)\le \cdots\le G_n(x)\le\cdots\le G_2(x)\le G_1(x)\]
简写成
\[g_n(x)\le f(x)\le G_n(x),\qquad a\le x\le b\]
\item 求相应的阶梯函数从$a$到$b$的和
\begin{equation}
\begin{split}
    s_n=g_n(x)\text{从$a$到$b$的和}
    &=\frac{b-a}{n}\left[f(x_0)+f(x_1)+\cdots+f(x_{n-1})\right]\\
    &=\frac{b-a}{n}\sum^n_{i=1}f(x_{i-1})
\end{split}
\end{equation}
\begin{equation}
    \begin{split}
    S_n=G_n(x)\text{从$a$到$b$的和}
    &=\frac{b-a}{n}\left[f(x_1)+f(x_2)+\cdots+f(x_{n})\right]\\
    &=\frac{b-a}{n}\sum^n_{i=1}f(x_{i})   
    \end{split}
    \end{equation}
\end{enumerate}
由性质1, 有：
\[s_1<s_2<\cdots<s_n<\cdots<f(x)\text{从$a$到$b$的和}<\cdots<S_n<\cdots<S_2<S_1\]
简写成
\begin{equation}
    s_n<f(x)\text{从$a$到$b$的和}<S_n
\end{equation}
在例3.3, 例3.4中，$f(x)$有明确的解析式，我们可以用求和公式直接求出$s_n$和$S_n$的表达式，从而可以得出它们的共同极限值，在这里，我们只知道$f(x)$的递增性，当然无法将和式$s_n$和$S_n$进一步化简，但是我们可以说明，当$n\to\infty$时，$S_n-s_n\to 0$.

事实上，
\begin{align*}
    S_n-s_n&=\frac{b-a}{n}\sum^n_{i=1}f(x_{i})   -\frac{b-a}{n}\sum^n_{i=1}f(x_{i-1}) \\
    &=\frac{b-a}{n}\sum^n_{i=1}\left[f(x_{i})-f(x_{i-1})\right]   \\
    &=\frac{b-a}{n}\left[f(x_{n})-f(x_{0})\right]  =\frac{b-a}{n}\left[f(b)-f(a)\right]   
\end{align*}
因此，当$n\to\infty$时，$S_n-s_n\to 0$.

又$f(x)$从$a$到$b$的和是介于$s_n$和$S_n$中间的唯一实数，
因此
\[\lim_{n\to\infty} s_n=\lim_{n\to\infty} S_n=f(x)\text{ 从$a$到$b$的和}\]

以上简单明了的分析，说明了下面两点互相关联的事实：
\begin{enumerate}
    \item 当函数$f(x)$在$a\le x\le b$上递增时，则它从$a$到$b$的和可以用上述两个夹逼阶梯函数从$a$到$b$的和去无限逼近.
    \item 由于$\Lim_{n\to\infty}s_n$与$\Lim_{n\to\infty}S_n$存在，且$\Lim_{n\to\infty}s_n=\Lim_{n\to\infty}S_n$, 我们可以用$f(x)$的上、下夹逼阶梯函数序列，即：$g_1 (x)\le g_2(x)\le \cdots\le g_n(x)\le \cdots\le f(x)\le \cdots\le G_n(x)\le \cdots\le G_2(x)\le G_1(x),\; (a\le x\le b)$的从$a$到$b$的和的极限$\Lim_{n\to\infty}s_n=\Lim_{n\to\infty}S_n$作为上述$f(x)$从$a$到$b$的和的数量化定义．
\end{enumerate}
\end{solution}

\begin{blk}
    {定义1}若$f(x)$是$[a,b]$上的递增连续函数，将$[a,b]$等分成$n$段，分点坐标$x_i=a+\frac{i}{n}(b-a),\; i=0,1, 2,\ldots,n$，则存在两个阶梯函数：
\[    g_n (x) =f (x_{i-1}) ,\quad x_{i-1}\le  x< x_i\]
    和
\[    G_n (x)=f(x_i),\quad x_{i-1}<x\le x_i,\quad (i=1,2,\ldots,n)\]
满足下面的性质：
\begin{enumerate}
    \item 对于任何$x\in[a,b]$，$g_n(x)\le f(x)\le G_n(x)$
    \item 相应的阶梯函数从$a$到$b$的和
\[s_n=\frac{b-a}{n}\sum^n_{i=1}f(x_{i-1}),\qquad S_n=\frac{b-a}{n}\sum^n_{i=1}f(x_{i})  \]
\end{enumerate}
当$n\to\infty$时，$S_n-s_n\to 0$，那么$f(x)$从$a$到$b$的和是$s_n$与$S_n$的共同极限，即：
\begin{align*}
f(x)\text{从$a$到$b$的和}&=\lim_{n\to\infty}\frac{b-a}{n}\sum^n_{i=1}f(x_{i-1})\\
&=\lim_{n\to\infty}\frac{b-a}{n}\sum^n_{i=1}f(x_{i})
\end{align*}

同样地，如果$f(x)$是$[a,b]$上的递减连续函数．那么将$[a,b]$ $n$等分，分点坐标$x_i=a+\frac{i}{n}(b-a)$, 这时有两个阶梯函数
\[    g_n (x) =f (x_{i}) ,\quad x_{i-1}<  x\le x_i\]
    和
\[    G_n (x)=f(x_{i-1}),\quad x_{i-1}\le x< x_i,\quad (i=1,2,\ldots,n)\]
满足下面的性质：
\begin{enumerate}
    \item 对于任何$x\in[a,b]$，$g_n(x)\le f(x)\le G_n(x)$
    \item 相应的阶梯函数$g_n(x)$与$G_n(x)$相应的和
\begin{align*}
    s_n&=g_n\text{从$a$到$b$的和}=\frac{b-a}{n}\sum^n_{i=1}f(x_{i})\\
    S_n&=G_n\text{从$a$到$b$的和}=\frac{b-a}{n}\sum^n_{i=1}f(x_{i-1})  
\end{align*}
\end{enumerate}
当$n\to\infty$时，$S_n-s_n=\frac{b-a}{n}\left[f(a)-f(b)\right]\to 0$，即：
\[\Lim_{n\to\infty}s_n=\Lim_{n\to\infty}S_n\]
\end{blk}

\begin{blk}
    {定义2} 递减连续函数从$a$到$b$的和为上述夹逼阶梯函数$g_n(x)$和$G_n(x)$从$a$到$b$的和的共同极限，即  
\begin{align*}
    f(x)\text{从$a$到$b$的和}&=\lim_{n\to\infty} \frac{b-a}{n}\sum^n_{i=1}f(x_i)\\
    &=\lim_{n\to\infty} \frac{b-a}{n}\sum^n_{i=1}f(x_{i-1})\\
\end{align*}
\end{blk}

通常我们把递增或递减的函数合称为单调函数，常见的
函数$f(x)$都是分段单调连续的，例如，$y=\sin x$本身虽然不是单调的，但是它在$\left[-\frac{\pi}{2}+2k\pi,\frac{\pi}{2}+2k\pi\right]\; (k\in\mathbb{Z})$这些区间的每一段是递增的；而在$\left[\frac{\pi}{2}+2k\pi,\frac{3\pi}{2}+2k\pi\right]\; (k\in\mathbb{Z})$这些区间的每一段是递减的．因此，对于一般在$[a,b]$上连续的函数，如果存在有限个分点使得$f(x)$在每个分段上都是单调的，我们可以逐段取上、下夹逼阶梯函数，合起来作为定义在$[a,b]$上的函数$f(x)$的阶梯函数，于是得出
\[g_n (x)\le f(x)\le G_n(x),\qquad x\in [a,b]\]
又因为有限个在分段上趋于0的量的和仍趋于0, 所以
\[\text{“$G_n(x)$从$a$到$b$的和”}-\text{“$g_n(x)$从$a$到$b$的和”}\to 0\]

总结以上讨论，我们叙述为存在定理和定义如下：

\begin{blk}
    {定理} 设$y=f(x)$是一个定义在$[a,b]$上分段单调函数，则存在满足下列性质的两系列阶梯函数：
    \[g_n(x)\le f(x)\le G_n(x),\qquad x\in[a,b]\]
而且当$n\to\infty$时，“$G_n(x)$从$a$到$b$的和，”与“$g_n(x)$从$a$到$b$的和”趋于共同极限．
\end{blk}


\begin{blk}
    {定义3} 设$f(x)$是$[a,b]$上的连续函数，如果 存在有限个分点：
\[a=a_0<a_1<\cdots<a_k<\cdots<a_\ell=b\]
使得$f(x)$在每个分段$[a_{k-1},a_k]$都是单调的，再将每个分段
    $[a_{k-1},a_k]$都$n$等分，则当$n\to\infty$时，有
\begin{align*}
f(x)\text{从$a$到$b$的和}&=\lim_{n\to\infty}\frac{a_1-a_0}{n}\sum^n_{i=1}f(x_i)+\lim_{n\to\infty}\frac{a_2-a_1}{n}\sum^{2n}_{i=n+1}f(x_i)+\\
&\cdots +\lim_{n\to\infty}\frac{a_\ell-a_{\ell-1}}{n}\sum^{\ell n}_{i=(\ell-1)n+1}f(x_i)
\end{align*}
\end{blk}

为了说明上述定义是合理的，我们就得证明上述$f(x)$从$a$到$b$的和与夹逼的阶梯函数列$\{G_n(x)\}$和$\{g_n(x)\}$的选取无关，其证明如下：

\begin{proof}
    设$\{\bar G_m(x)\}$和$\{\bar g_m(x)\}$是另外一组满足存在定理的上、下夹逼函数列，则由下述不等式
\[g_n(x)\le f(x)\le G_n(x),\quad \bar g_m(x)\le f(x)\le \bar G_m(x),\qquad a\le x\le b\]
即有
\[g_n(x)\le \bar G_m(x),\qquad \bar g_m(x)\le G_n(x)\]
所以由基本性质1，有
\begin{align*}
    s_n&=g_n(x)\text{从$a$到$b$的和}\le \bar G_m(x)\text{从$a$到$b$的和}=\bar S_m\\
    \bar s_m&=\bar g_m(x)\text{从$a$到$b$的和}\le G_n(x)\text{从$a$到$b$的和}=S_n\\
\end{align*}
所以
\[\lim_{n\to\infty}s_n\le\lim_{m\to\infty}\bar S_m,\qquad \lim_{m\to\infty}\bar s_m\le \lim_{n\to\infty}S_n\]
但是
\[\lim_{n\to\infty}s_n\le \lim_{m\to\infty}\bar S_m=\lim_{m\to\infty}\bar s_m\le \lim_{n\to\infty} S_n=\lim_{n\to\infty}s_n\]
所以上述极限必须相等，即：
\[\lim_{n\to\infty}s_n= \lim_{m\to\infty}\bar S_m=\lim_{m\to\infty}\bar s_m=\lim_{n\to\infty} S_n\]
\end{proof}

\begin{example}
求$f(x)=2x-x^2$从0到2的和．
\end{example}

\begin{solution}
    由$f'(x)=2-2x$, $f''(x)=-2$知$f(1)=2-1=1$是
$f(x)$的极大值，并且$y=2x-x^2$在$[0, 1]$上递增，在$[1,2]$上递减，于是，我们把区间$[0, 1]$和$[1, 2]$都$n$等分，设分点坐标$x_i=\frac{i}{n},\; i=0, 1, 2,\ldots,2n$, 即有
\[0=x_0<x_1<x_2<\cdots <x_n=1<x_{n+1}<x_{n+2}<\cdots <x_{2n}=2\]
且$\Delta x_i=x_i-x_{i-1}=\frac{1}{n}$，由定义3得到：
\begin{align*}
    f(x)&=2x-x^2\text{从0到2的和}\\
    &=\lim_{n\to\infty}\frac{1}{n}\sum^n_{i=1}\left[2\left(\frac{i}{n}\right)-\left(\frac{i}{n}\right)^2\right]+\lim_{n\to\infty}\frac{1}{n}\sum^{2n}_{i=n+1}\left[2\left(\frac{i}{n}\right)-\left(\frac{i}{n}\right)^2\right] \\
&=\lim_{n\to\infty}\frac{1}{n}\sum^{2n}_{i=1}2\left(\frac{i}{n}\right)-\lim_{n\to\infty}\frac{1}{n}\sum^{2n}_{i=1}\left(\frac{i}{n}\right)^2\\
&=2\lim_{n\to\infty}\frac{1}{n^2}\sum^{2n}_{i=1}i-\lim_{n\to\infty}\frac{1}{n^3}\sum^{2n}_{i=1}i^2\\
&=2\lim_{n\to\infty}\frac{1}{n^2}\cdot \frac{(2n+1)2n}{2}-\lim_{n\to\infty}\frac{1}{n^3}\cdot \frac{2n(2n+1)(4n+1)}{6}\\
&=2\lim_{n\to\infty}\left(2+\frac{1}{n}\right)-\lim_{n\to\infty}\frac{1}{3}\left(2+\frac{1}{n}\right)\left(4+\frac{1}{n}\right)\\
&=4-\frac{8}{3}=\frac{4}{3}
\end{align*}
\end{solution}

\begin{ex}
\begin{enumerate}
    \item 已知质点的运动速度$v=t+4$, 试求质点在前10秒内所走的路程．
    \item 求$f(x)=x^3$在$1\le x\le 2$上的和．
\end{enumerate}
\end{ex}


\section{定积分的定义和基本性质}
\subsection{定积分定义}

设函数$f(x)$在$[a,b]$上连续，如果存在有限个分点
\[a=a_0<a_1<a_2<\cdots<a_{\ell-1}<a_{\ell}=b\]
使得$f(x)$在每个分段$[a_{k-1},a_k]\; (k=1, 2,\ldots,\ell)$上都是单调的，我们把$f(x)$从$a$到$b$的和叫做$f(x)$从$a$到$b$的\textbf{定积分}，并记作
\[\int^b_a f(x) \dd x\]

这里，积分符号所使用的是长$S$形的求和号的变形，而从部分区间长$\Delta x_i=\frac{a_k-a_{k-1}}{n}$过渡到极限，则通过字母$\dd$来表示．

我们把数$a$与$b$称为\textbf{积分限}（$a$称为\textbf{下限}，$b$称为\textbf{上限}）．区间$[a,b]$称为\textbf{积分区间}，函数$f(x)$称为\textbf{被积函数}，乘积$f(x)\dd x$称为\textbf{被积表达式}，定积分符号下出现的字母$x$叫做\textbf{积分变量}．

上述定积分定义用定积分符号表示就是：
\[\int^b_a f(x) \dd x=\sum^{\ell}_{k=1}\int^{a_k}_{a_{k-1}}f(x)\dd x\]
其中$\Int^{a_k}_{a_{k-1}}f(x)\dd x$是$f(x)$为单调的第$k$个分段$[a_{k-1},a_k]$上的夹逼阶梯函数$g_n(x)$和$G_n(x)$从$a_{k-1}$到$a_k$的和的共同极限，即：
\begin{align*}
    \int^{a_k}_{a_{k-1}}f(x)\dd x&=\lim_{n\to \infty}\frac{a_k-a_{k-1}}{n}\sum^n_{i=1}f(x_{i-1})\\
    &=\lim_{n\to \infty}\frac{a_k-a_{k-1}}{n}\sum^n_{i=1}f(x_{i})
\end{align*}

定积分定义的另一种表述形式是：设函数$f(x)$在$a\le x\le b$上分段单调连续，如果存在两系列上、下夹逼阶梯函数$\{g_n(x)\}$, $\{G_n(x)\}$, 使得
\[g_n (x)\le f(x)\le G_n(x),\qquad a\le x\le b\]
并且$g_n(x)$从$a$到$b$的和$s_n$与$G_n(x)$从$a$到$b$的和$S_n$具有相同的
极限，这个极限叫做$f(x)$从$a$到$b$的定积分$\Int^b_a f(x) \dd x$.

综上所述，在积分符号中，我们只对$a<b$, 即积分下限小于积分
上限的情形给出了定义，若$a>b$，我们定义
\[\int^a_b f(x) \dd x=-\int^b_a f(x) \dd x\]
此外，由于定积分可以解释为曲边梯形的面积，自然可以规定
\[\int^a_a f(x) \dd x=0\]
作了这样的规定之后，不论$a<b$, $a>b$或$a=b$, 定积分都有意义了．


\begin{ex}
\begin{enumerate}
    \item 求证 $\Int_{a}^{b} k x\dd x=\frac{k b^{2}}{2}-\frac{k a^{2}}{2} \quad(a<b)$
    \item  求 $\Int_{8}^{0}\left(x^{2}-4 x\right) \dd x$
    \item  求 $\Int_{-1}^{2}\left(x^{3}-3 x\right) \dd x$
\end{enumerate}
\end{ex}

\subsection{逼近法求曲线形的面积}

任意一条曲线围成的图形（图\ref{fig:chapter3_sec2_arbitrary_shape}）常常可以用两组互相垂直的直线把它分成若干部分，每一部分都是一个曲边梯形（图\ref{fig:chapter3_sec2_good_piece}）, 在这里并不排除下述情形（图\ref{fig:chapter3_sec2_degenerated_piece}）：两条平行的边中有一条缩成了一点，因而曲边梯形变成了曲边三角形，这样一来，我们的问题就化成了求曲边梯形面积的问题．

\begin{figure}[htp]
    \centering
    \begin{minipage}[t]{0.45\textwidth}
        \centering
        \begin{tikzpicture}[>=latex, scale=1]
    \draw(-2,3)--(2,3);
    \draw(-2,2)--(2,2); 
    \draw(-.5,0)--(-.5,4);
    \draw(.5,0)--(.5,4); 
    \node at (0,0){};
    \draw[thick] plot[smooth cycle] coordinates{
 (-0.4,3.75)(0.4,3.85)(1.0,3.6)(1.15,3.0)(0.95,2.1)(0.55,1.1)(0.0,0.62)(-0.8,0.62)(-1.35,0.95)(-1.3,1.7)(-1.0,2.6)
};
        \end{tikzpicture}
        \caption{}\label{fig:chapter3_sec2_arbitrary_shape}
        \end{minipage}
    \begin{minipage}[t]{0.25\textwidth}
    \centering
    \begin{tikzpicture}[>=latex, scale=1]
    \draw(0,1.3)--(0,0)--(2.5,0)--(2.5,1.8);
    \node at (0,0) {};
    \draw plot[smooth] coordinates{(0,1.3)(1.4,2)(2.5,1.8)};
    \end{tikzpicture}
    \caption{}\label{fig:chapter3_sec2_good_piece}
    \end{minipage}
    \begin{minipage}[t]{0.25\textwidth}
    \centering
    \begin{tikzpicture}[>=latex, scale=1]
    \draw(0,2)--(0,0)--(1.5,0);
    \node at (0,0){};
    \draw plot[smooth] coordinates{(0,2)(0.5,1.8)(1,1.3)(1.5,0)};
    \end{tikzpicture}
    \caption{}\label{fig:chapter3_sec2_degenerated_piece}
    \end{minipage}
  \end{figure}

\begin{blk}
  {命题} 设$y=f(x)$是定义在$a\le x\le b$的分段单调连续函数，而且$f(x)\ge 0$, 则区域$R=\{a\le x\le b,\; 0\le y\le f(x)\}$的
面积等于
\[\int^b_a f (x) \dd x \]  
\end{blk}

\begin{figure}[htp]
    \centering
\begin{tikzpicture}[>=latex]
    \draw[->, thick](-.5,0)--(7,0)node[right]{$x$};
    \draw[->, thick](0,-.5)--(0,5)node[right]{$y$};
\draw[domain=.5:6, samples=100, very thick]plot(\x, {1.5*sin(1.2*(\x-1) r)+3});
\draw(.5,0)node[below]{$a$} rectangle  (1.1, .18+3);
\draw(1.1,0) rectangle  (1.8, 1.23+3);
\draw(1.8,0) rectangle  (2.5, 1.46+3);
\draw[pattern=north east lines](2.5,0) rectangle  (3.5, .212+3);
\draw[pattern=north east lines](3.5,0) rectangle  (4.2, -.965+3);
\draw[pattern=north east lines](4.2,0) rectangle  (5, -1.49+3);
\draw(5,0) rectangle  (5.5, -1.16+3);
\draw(5.5,0) rectangle (6, -.419+3);

\draw[pattern=north east lines](.5,0) rectangle  (1.1, -.847+3);
\draw[pattern=north east lines](1.1,0) rectangle  (1.8, .18+3);
\draw[pattern=north east lines](1.8,0) rectangle  (2.5, 1.23+3);
\draw(2.5,0) rectangle  (3.5, 1.46+3);
\draw(3.5,0) rectangle  (4.2, .212+3);
\draw(4.2,0) rectangle  (5, -.965+3);
\draw[pattern=north east lines](5,0) rectangle  (5.5,-1.49+3);
\draw[pattern=north east lines](5.5,0) rectangle (6, -1.16+3);

\node at (6,0)[below]{$b$};
\node at (0,0)[below left]{$O$};
\node at (4.5,3.2)[above]{$y=f(x)$};

\end{tikzpicture}
    \caption{}\label{fig:chapter3_sec2_riemann_sum}
\end{figure}


\begin{proof}
    由$\Int^b_a f (x) \dd x$的定义得知，存在两系列阶梯函数
$\{g_n(x)\}$和$\{G_n(x)\}$, 满足下面的性质：
\[g_n(x)\le f (x) \le G_n(x)\]
而且$G_n(x)$从$a$到$b$的和与$g_n(x)$从$a$到$b$的和趋于共同的极限$\Int^b_a f (x) \dd x$. 令
\begin{align*}
    R_n&=\{a\le x\le b,\; 0\le y\le g_n(x)\}\\
    R'_n&=\{a\le x\le b,\; 0\le y\le G_n(x)\}
\end{align*}
它们都是由高高低低的狭长方形所组成的区域（图\ref{fig:chapter3_sec2_riemann_sum}），则：
\[R_n\subset \text{曲边梯形区域 }R \subset R'_n\]
而且
\[R_n\text{的面积}=g_n(x)\text{从$a$到$b$的和}s_n,\qquad R'_n\text{的面积}=G_n(x)\text{从$a$到$b$的和}S_n\]
即有
\[s_n<\text{曲边梯形区域$R$的面积}A<S_n\]
由定积分定义有：
\[\lim_{n\to\infty}s_n=\lim_{n\to\infty}S_n=\Int^b_a f (x) \dd x\]
又曲边梯形面积$A(f,\; [a,b])$也是夹逼在$s_n$和$S_n$中间的实数，所以
\[A(f,\; [a,b])=\Int^b_a f (x) \dd x\]
\end{proof}

\begin{ex}
\begin{enumerate}
    \item 求$\Int^1_{-2}e^x\dd x$. 

提示：$\Lim_{\Delta x\to 0}\frac{e^{\Delta x}-1}{\Delta x}=1$.    

\item 求在$x=0$与$x=\pi$间正弦曲线$y=\sin x$与$Ox$轴所包的面积．
\end{enumerate}
\end{ex}    

\subsection{定积分的基本性质}

为简单起见，我们约定以下被积函数在积分区间上连续并且可以分段单调，即$f(x)$是在$[a,b]$上连续并且只有有限个极大值和极小值的函数．

\begin{blk}{性质1}
若$a<b<c$，那么，
\[\int^b_a f(x)\dd x+\int^c_b f(x)\dd x=\int^c_a f(x)\dd x\]
\end{blk}

\begin{proof}
因为$\Int^b_a f(x)\dd x$存在，故对于在$[a,b]$上的
$f(x)$, 存在一组上、下夹逼函数列$\{g_n(x)\}$和$\{G_n(x)\}$, 使得
\[g_n (x) \le f (x) \le G_n (x),\qquad x\in [a,b]\]
且相应的阶梯函数在$[a,b]$上的和$s_n$与$S_n$适合
\begin{equation}
    s_n<\int^b_a f(x)\dd x<S_n,\qquad S_n-s_n\to 0
\end{equation}
对于在$[b,c]$上的$f(x)$, 同样得到
\[\bar g_n (x) \le f (x) \le \bar G_n (x),\qquad x\in [b,c]\]
且
\begin{equation}
    \bar  s_n<\int^c_b f(x)\dd x<\bar S_n,\qquad \bar S_n-\bar s_n\to 0
\end{equation}
$(3.4)+(3.5)$得到
\begin{equation}
    s_n+ \bar  s_n<\int^b_a f(x)\dd x+\int^c_b f(x)\dd x<S_n+\bar S_n
\end{equation}

另一方面，对于函数$f(x)\; (a\le x\le c)$存在下面一组阶梯函数列：
\[\tilde g_n(x)=\begin{cases}
    g_n(x),  &  a\le x<b\\
    \max\left(g_n(b), \bar g_n(b)\right),  &  x=b\\
    \bar g_n(x), & b<x\le c
\end{cases}\]
\[\tilde G_n(x)=\begin{cases}
    G_n(x),  &  a\le x<b\\
    \min\left(G_n(b), \bar G_n(b)\right),  &  x=b\\
    \bar G_n(x), & b<x\le c
\end{cases}\]

由性质2，阶梯函数列$\tilde g_n(x)$与$\tilde G_n(x)$的从$a$到$c$的和分别是$s_n+\bar s_n$和$S_n+\bar S_n$．因为$\tilde g_n(x)$与$\tilde G_n(x)$满足条件：
\begin{enumerate}
    \item $\tilde g_n(x)\le f(x)\le \tilde G_n(x),\quad a\le x\le c$
    \item 
    \begin{align*}
    \lim_{n\to\infty}\left[(S_n+\bar S_n)-(s_n+\bar s_n)\right]&=\lim_{n\to\infty}\left[(S_n-s_n)+(\bar S_n-\bar s_n)\right]\\
    &=\lim_{n\to\infty}(S_n-s_n)+\lim_{n\to\infty}(\bar S_n-\bar s_n)=0
    \end{align*}
\end{enumerate}
所以根据定积分定义得到
\[\lim_{n\to\infty}(S_n+\bar S_n)=\lim_{n\to\infty}(s_n+\bar s_n)=\int^c_a f(x)\dd x \]

因为$\Int^c_a f(x)\dd x $与$\Int^b_a f(x)\dd x +\Int^c_b f(x)\dd x $都是被$s_n+\bar s_n$与$S_n+\bar S_n$所夹逼的唯一实数，所以
\[\Int^c_a f(x)\dd x=\Int^b_a f(x)\dd x +\Int^c_b f(x)\dd x\]
\end{proof}

\begin{blk}{性质2}
若$f(x)=f_1(x)+f_2(x)$，则：
\[\Int^b_a f(x)\dd x=\Int^b_a f_1(x)\dd x +\Int^b_a f_2(x)\dd x\]
\end{blk}

\begin{proof}
设$\{g_n(x)\}$与$\{G_n(x)\}$是$f_1(x)$在$[a,b]$上的上、下夹逼阶梯函数列，又$\{\bar g_n(x)\}$与$\{\bar G_n(x)\}$是$f_2(x)$在$[a,b]$上的上、下夹逼阶梯函数列．$s_n, S_n, \bar s_n, \bar S_n$是相应的阶梯函数从$a$到$b$的和．于是由
$\Int^b_a f_1(x)\dd x$和$\Int^b_a f_2(x)\dd x$的存在，得：
\begin{align}
g_n(x)\le f_1(x)\le G_n(x),\qquad a\le x\le b\\
\bar g_n(x)\le f_2(x)\le \bar G_n(x),\qquad a\le x\le b
\end{align}
并且
\begin{align}
s_n<\Int^b_a f_1(x)\dd x<S_n,\quad \text{且 } S_n-s_n\to 0\\
\bar s_n<\Int^b_a f_2(x)\dd x<\bar S_n,\quad \text{且 } \bar S_n-\bar s_n\to 0
\end{align}
由$(3.9)+(3.10)$，得到
\begin{equation}
   s_n+ \bar s_n <\Int^b_a f_1(x)\dd x+\Int^b_a f_2(x)\dd x<S_n+\bar S_n
\end{equation}
另一方面，由$(3.7)+(3.8)$，得到
\[g_n(x)+\bar g_n(x)\le f_1(x)+ f_2(x)\le G_n(x)+\bar G_n(x),\qquad a\le x\le b\]
而且，$g_n(x)+\bar g_n(x)$与$G_n(x)+\bar G_n(x)$也是在$[a,b]$上的阶梯函数，我们要说明它们是$f_1(x)+ f_2(x)$在$[a,b]$上的一组夹逼阶梯函数列．

因为：
\begin{align*}
    g_n(x)+\bar g_n(x)\text{从$a$到$b$的和}&=\left[ g_n(x)\text{从$a$到$b$的和}\right]+\left[\bar g_n(x)\text{从$a$到$b$的和}\right]\\
    &=s_n+\bar s_n\\
    G_n(x)+\bar G_n(x)\text{从$a$到$b$的和}&=\left[ G_n(x)\text{从$a$到$b$的和}\right]+\left[\bar G_n(x)\text{从$a$到$b$的和}\right]\\
    &=S_n+\bar S_n
\end{align*}
而且
\begin{equation}
    \left(S_n+\bar S_n\right)-\left(s_n+\bar s_n\right)=\left(S_n-s_n\right)+\left(\bar S_n-\bar s_n\right)\to 0
\end{equation}
所以，由$\Int^b_a \bigl(f_1(x)+ f_2(x)\bigr)\dd x$的定义，得到
\begin{equation}
    \lim_{n\to\infty}\left(S_n+\bar S_n\right)=\lim_{n\to\infty}\left(s_n+\bar s_n\right)=\Int^b_a \bigl(f_1(x)+ f_2(x)\bigr)\dd x
\end{equation}
由(3.11)和(3.13)，$\Int^b_a \bigl(f_1(x)+ f_2(x)\bigr)\dd x$与$\Int^b_a f_1(x)\dd x + \Int^b_a  f_2(x)\dd x$
是被$s_n+\bar s_n$与$S_n+\bar S_n$所夹逼的唯一实数，所以
\[\Int^b_a \bigl(f_1(x)+ f_2(x)\bigr)\dd x=\Int^b_a f_1(x)\dd x + \Int^b_a  f_2(x)\dd x\]
\end{proof}

\begin{blk}{性质3}
\[\Int^b_a kf(x)\dd x =k\Int^b_a f(x)\dd x \]
\end{blk}

此法则的证明大致和性质2的证明相同，即当$\{g_n(x)\}$和$\{G_n(x)\}$上、下夹逼$f(x)$时，那么在$k>0$的情形，$\{kg_n(x)\}$和$\{kG_n(x)\}$上、下夹逼$kf(x)$; 在$k<0$的情形，$\{kG_n(x)\}$和$\{kg_n(x)\}$上、下夹逼$kf(x)$. 我们把证明的过程留给读者去写．

\begin{blk}{性质4}
    若$m\le f(x)\le M$, $a\le x\le b$，那么
    \[m(b-a)\le \Int^b_a f(x)\dd x \le M(b-a) \]
    \end{blk}

\begin{proof}
    设$\{g_n(x)\}$与$\{G_n(x)\}$是$f(x)$在$[a,b]$上的夹逼阶梯函数列，$s_n$与$S_n$是相应的阶梯函数在$[a,b]$上的和，根据积分的定义，有
\[\lim_{n\to\infty}s_n=\lim_{n\to\infty}S_n=\int^b_a f(x)\dd x\]
换言之，任给一个正数$\varepsilon$，存在自然数$N$，使得当$n>N$时，有
\begin{equation}
    S_n<\int^b_a f(x)\dd x+\varepsilon
\end{equation}
成立．

又由于$G_n(x)\ge f(x)\ge m,\quad x\in[a,b]$，根据性质1，得到：
\begin{equation}
    S_n\ge m(b-a)
\end{equation}
所以，当$n>N$时，由(3.14)和(3.15)，有
\[m(b-a)\le S_n<\int^b_a f(x)\dd x+\varepsilon\]
即
\begin{equation}
    m(b-a)<\int^b_a f(x)\dd x+\varepsilon
\end{equation}
成立．

因为这个不等式(3.16)对于每个正数$\varepsilon$都成立，所以
\[m(b-a)\le \int^b_a f(x)\dd x\]
同样证明，得到
\[M(b-a)\ge \int^b_a f(x)\dd x\]
因此
\[m(b-a)\le \int^b_a f(x)\dd x\le M(b-a)\]
\end{proof}

\begin{figure}[htp]
    \centering
\begin{tikzpicture}[>=latex]
\draw[->, thick](-1.5,0)--(3,0)node[right]{$x$};
\draw[->, thick](-1,-.5)--(-1,4)node[right]{$y$};
\node at (-1,0)[below left]{$O$};
\draw[domain=-.1:2.3, samples=50, very thick]plot(\x, {\x*(\x-1)*(\x-2)+2});
\draw(-.1,0)node[below]{$a$} rectangle (2.3,2.9);
\node at (2.3,0)[below]{$b$};
\draw(1.577,0)node[below]{$x_0$}--(1.577,1.615);
\draw(-.1,1.615)--(2.3,1.615);
\draw[dashed](-1,1.615)node[left]{$m$}--(-.1,1.615);
\draw[dashed](-1,2.9)node[left]{$M$}--(-.1,2.9);

\node at (1.1,2.5){$y=f(x)$};
\end{tikzpicture}
    \caption{}\label{fig:chapter3_sec2_y=f(x)_bounded}
\end{figure}

从几何图形来看，以曲线$y=f(x)$为一边而以线段$[a,b]$为底边的曲边梯形界于以$[a,b]$为底边，高分别等于$m$和$M$的两个矩形之内，
故曲边梯形的面积$\Int^b_a f(x)\dd x$
在两个矩形面积$m(b-a)$和$M(b-a)$之间（图\ref{fig:chapter3_sec2_y=f(x)_bounded}）．

这个性质是说分段单调的连续函数$f(x)$在区间$[a,b]$上的定积分是有界的．

\begin{blk}{性质5}
如果$f(x)\le \varphi(x),\; a\le x\le b$，那么：
\[\int^b_a f(x)\dd x\le \int^b_a \varphi(x)\dd x\]
\end{blk}

\begin{proof}
    设$\{g_n(x)\}$与$\{G_n(x)\}$是$f(x)$在$[a,b]$上的夹逼阶梯函数列，又$\{\bar g_n(x)\}$与$\{\bar G_n(x)\}$是$\varphi(x)$在$[a,b]$上的夹逼阶梯函数列，
又$s_n, S_n, \bar s_n, \bar S_n$是相应的阶梯函数从$a$到$b$的和．于是根据定积分的定义，有：
\[s_n<\int^b_a f(x)\dd x<S_n,\qquad \bar s_n<\int^b_a \varphi(x)\dd x<\bar S_n\]
当$n\to\infty$时，得到：
\begin{align*}
\lim_{n\to\infty}S_n =\lim_{n\to\infty}s_n =\int^b_a f(x)\dd x\\
\lim_{n\to\infty}\bar S_n =\lim_{n\to\infty}\bar s_n =\int^b_a \varphi(x)\dd x\\
\end{align*}
换言之，任给正数$\varepsilon$，存在$N_1$，使得当$n>N_1$时，有
\[s_n>\int^b_a f(x)\dd x-\frac{\varepsilon}{2}\]

当$n>N_2$时，有
\[\bar S_n<\int^b_a \varphi(x)\dd x+\frac{\varepsilon}{2}\]

因为$\bar G_n(x)>\varphi(x)>f(x)>g_n(x),\quad x\in[a,b]$，根据性质1，得到：$\bar S_n\ge s_n$.

所以，当$n>\max(N_1,N_2)$时，有
\[\int^b_a f(x)\dd x<s_n+\frac{\varepsilon}{2}<\bar S_n+\frac{\varepsilon}{2}<\int^b_a \varphi(x)\dd x+\varepsilon\]
因为这个不等式对于每个正数$\varepsilon$都成立，所以必然有
\[\int^b_a f(x)\dd x\le \int^b_a \varphi(x)\dd x\]
\end{proof}

\begin{blk}{性质6：定积分中值定理}
   设函数$f(x)$在闭区间$[a,b]$上分段单调连续，又
$m=f(c)$, $M=f(d)$分别是$f(x)$在$[a,b]$上的最小值和最大值，则在$[a,b]$上至少存在一点$\xi$, 使得下面的等式成立：
\[\int^b_a f (x) \dd x=f(\xi) (b-a)\] 
\end{blk}

\begin{proof}
    因为$m\le f (x) \le M,\quad x\in [a,b]$, 
由性质4, 即得
\[m (b-a)\le \int^b_a f (x) \dd x\le M(b-a)\]
上面不等式的两端各除以$(b-a)$, 得
\[f(c)=m\le \frac{\int^b_a f (x) \dd x}{b-a}\le M=f(d)\]
因为$f(x)$在$[a,b]$上连续，再由连续函数的中间值定理，必存在一个$\xi\in(c,d)$, 使得
\[f(\xi)=\frac{1}{b-a}\int^b_a f (x) \dd x\]
两边再乘以$(b-a)$，得
\[\int^b_a f (x) \dd x=f(\xi)(b-a)\]
这就是我们所要证明的．
\end{proof}

这个公式的几何意义是：以线段$[a,b]$为底边，以曲线$y=f(x)$为曲边的曲边梯形，它的面积等于同一底边
而高为$f(\xi)$的一个矩形的面积（图\ref{fig:chapter3_sec2_integral_mean_value}）．因此，$f(\xi)$称为曲边梯形的平均高度．

我们也称
\[f(\xi)=\frac{1}{b-a}\int^b_a f (x) \dd x\]
为$f(x)$在$[a,b]$上的平均值．

\begin{figure}[htp]
    \centering
    \begin{minipage}[t]{0.48\textwidth}
    \centering
    \begin{tikzpicture}[>=latex, xscale=.8]
\draw[->, thick](-.5,0)--(5,0)node[right]{$x$};
\draw[->, thick](0,-.5)--(0,4)node[right]{$y$};
\node at (0,0)[below left]{$O$};
\draw[domain=.7:4, very thick, samples=100]plot(\x, {ln(\x)+1.5})node[above]{$y=f(x)$};
\draw[dashed](.7,0)node[below]{$a$} rectangle (4,2.3);
\draw[dashed](2.23,0)node[below]{$\xi$}--(2.23,2.3)node[below right]{$f(\xi)$};
\draw[dashed](4,2.89)--(4,2.3);
\node at (4,0)[below]{$b$};

    \end{tikzpicture}
    \caption{}\label{fig:chapter3_sec2_integral_mean_value}
    \end{minipage}
    \begin{minipage}[t]{0.48\textwidth}
    \centering
    \begin{tikzpicture}[>=latex, scale=.5]
  \draw[->, thick](-2,0)--(6,0)node[right]{$x$};
\draw[->, thick](0,-8.5)--(0,2)node[right]{$y$};
\node at (0,0)[below left]{$O$};
\draw[domain=-1:4, very thick, samples=100]plot(\x, {2*\x-\x*\x});
\fill[pattern=north east lines, domain=2:4, very thick, samples=100]plot(\x, {2*\x-\x*\x})--(4,0)node[above]{4}
--(2,0)node[above]{2};
\draw(1,0)node[below]{1}--(1,.2);
    \end{tikzpicture}
    \caption{}\label{fig:chapter3_sec2_example3.7}
    \end{minipage}
  \end{figure}

\begin{example}
求$\Int^4_2 (2x-x^2)\dd x$.
\end{example}


\begin{solution}
    根据本节的定积分计算法则得：
\begin{align*}
    \Int^4_2 (2x-x^2)\dd x&=\Int^4_2 2x\dd x+\Int^4_2 (-x^2)\dd x\\
    &=2\Int^4_2 x\dd x-\Int^4_2 x^2\dd x
\end{align*}
利用例3.4的计算结果，得
\[\Int^4_2 x\dd x=\frac{4^2-2^2}{2}=6,\qquad \Int^4_2 x^2\dd x=\frac{4^3-2^3}{3}=18\frac{2}{3}\]
因此
\[\Int^4_2 (2x-x^2)\dd x=2\x 6-18\frac{2}{3}=-6\frac{2}{3}\]
它的几何意义是图\ref{fig:chapter3_sec2_example3.7}中阴影区域的有号面积．
\end{solution}

\begin{example}
  设$f(x)=\begin{cases}
      x^2,& 0\le x\le 1\\
      1,& 1<x\le 2
  \end{cases}$,  求$f(x)$在$[0,2]$上的平均值.
\end{example}

\begin{solution}
由平均值定义，再利用性质1，有    
\begin{align*}
    f(\xi)&=\frac{1}{b-a}\int^b_a f(x)\dd x=\frac{1}{2}\int^2_0 f(x)\dd x\\
    &=\frac{1}{2}\left[\int^1_0 f(x)\dd x+\int^2_1 f(x)\dd x\right]\\
    &=\frac{1}{2}\left[\int^1_0 x^2\dd x+\int^2_1 1\dd x\right]\\
    &=\frac{1}{2}\left[\frac{1}{3}+1\right]=\frac{2}{3}
\end{align*}
\end{solution}

\begin{ex}
\begin{enumerate}
    \item 利用以前定积分的结果和定积分的性质，计算下列定积分：
\begin{enumerate}
    \item $\Int^3_{-1}(2x^2-4x)\dd x$
    \item 若$f(x)=\begin{cases}
        x,& 0\le x<1\\
        x-2,& 1\le x\le 2
    \end{cases}$，则$\Int^2_0 f(x)\dd x=?$
    \item 若$f(x)=\begin{cases}
        x,& 0\le x\le 1\\
        x-2,& 1< x\le 2
    \end{cases}$，则$\Int^2_0 f(x)\dd x=?$
\item 若$g(t)=2t^2+|t|-1,\; -1\le t\le 1$，则$\Int^1_{-1} g(t)\dd t=?$
\end{enumerate}

\item 将图中阴影部分的面积用定积分表示．
\item 设$\ell(t)=mt+b$, （$m,b$是常数），$t\in [c,d]$，

求证：$\Int^d_c (mt+b)\dd t=\frac{\ell(c)+\ell(d)}{2}(d-c)$
\item 证明：
\begin{enumerate}
    \item 若$0<x<10$，则$\frac{1}{1016}\le\frac{1}{x^3+16}\le \frac{1}{16}$
    \item $\frac{5}{508}\le \Int^{10}_0\frac{1}{x^3+16}\dd x\le \frac{5}{8}$
\end{enumerate}
\end{enumerate}
\end{ex}

\begin{figure}[htp]
    \centering
\begin{tikzpicture}[>=latex]
    
\begin{scope}[scale=.6]
\draw[->](-2,0)--(4,0)node[right]{$x$};
\draw[->](0,-3)--(0,7)node[right]{$y$};
\fill[domain=-1:3, samples=50, pattern=north east lines]plot(\x,{2*(\x-1)^2-2})--(3,0)--(-1,0)--(-1,6);
\draw[domain=-1:3, samples=50, very thick]plot(\x,{2*(\x-1)^2-2})node[above]{$y=2x^2-4x$};
\node at (0,0)[above right]{$O$};
\node at (-1,0)[below]{$-1$};
\node at (3,0)[below]{$3$};
\node at (-1,3)[left]{$x=-1$};
\node at (3,3)[right]{$x=3$};
\node at (1,-2)[below]{$(1,-2)$};
\node at (1,-3.8){(i)};
\end{scope}
\begin{scope}[xshift=5cm, yshift=-.8cm, scale=.9]
\draw[->](-1,0)--(4,0)node[right]{$x$};
\draw[->](0,-1)--(0,5.5)node[right]{$y$};
\node at (0,0)[below left]{$O$};
\draw[domain=0:2.2, samples=50, very thick]plot(\x,{\x^2})node[above]{$y=x^2$};
\draw[thick](0,0)--(3.5,3.5)node[right]{$y=x$};
\node at (1,0)[below]{$1$};
\node at (2,0)[below]{$2$};
\draw[dashed](0,1)node[left]{1}--(1,1)--(1,0);
\fill[domain=0:1, samples=50, pattern=north east lines]
plot(\x,{\x^2})--(2,2)--(2,0)--(0,0);
\node at (2,-1.5){(ii)};
\end{scope}

\end{tikzpicture}
    \caption*{第2题}
\end{figure}



\subsection*{习题3.2}
\begin{enumerate}
    \item 计算下列定积分：
\begin{multicols}{3}
\begin{enumerate}
    \item $\Int^1_4 |x|\dd x$
    \item $\Int^2_4 |x|\dd x$
    \item $\Int^4_1 |x|\dd x$
\end{enumerate}
\end{multicols}
\item 求$\Int^b_a |x|\dd x$.

提示：分$a<b<0$, $a<0<b$, $0<a<b$三种情况讨论．
\item 证明：
\begin{enumerate}
    \item 若$\frac{\pi}{4}\le x\le\frac{\pi}{2}$，则$\frac{2}{\pi}\le \frac{\sin x}{x}\le \frac{2\sqrt{2}}{\pi}$
    \item $\frac{1}{2}\le \Int^{\tfrac{\pi}{2}}_{\tfrac{\pi}{4}}\frac{\sin x}{x}\dd x\le \frac{\sqrt{2}}{2}$
\end{enumerate}
\item 已知作用在作直线运动的质点上的力是$F=s^2+1$, 试求从距离1到10之间所作的功．
\item 求$y=\begin{cases}
    A\sin\frac{2\pi t}{T}, & 0\le t\le \frac{T}{2}\\
    0,& \frac{T}{2}\le t\le T
\end{cases}$ 在$[0,T]$上的平均值．
\end{enumerate}
\chapter{微积分学基本定理}
在前两章中，我们分别引入了函数的变率（导数），函数的和（定积分）的基本概念，本章将研究函数的导函数与函数的求和函数这两者之间的互逆关系，并说明我们可以用求导函数的逆运算方法来计算定积分．

\section{微积分学基本定理}
\subsection{导函数与求和函数}
函数$f(x)$在点$x_0$处的变化率（导数）的定义是
\[f' (x_0) =\lim_{h\to 0}\frac{f (x_0+h) -f (x_0)}{h}\]
显然， $f'(x_0)$的值与$f(x)$在点$x_0$的值以及在点$x_0$的邻近的函数值有关，当点$x_0$在$(a,b)$内变化时，$f'(x_0)$也跟着变化，那么$f'(x_0)$便是一个新函数称为$f(x)$的导函数．计算一个函数的导函数是一件比较简便的事情．

定积分$\int^b_a f(x)\dif x$的定义是把区间$[a,b]$无限细分
而得到上下夹逼阶梯函数的和的共同极限，其几何意义是曲线$y=f(x)$和直线$x=a$, $x=b$, 及$y=0$所围成的区域的有号面积．

假如我们考虑$f(x)$在一个变动的区间$[a,x]$上的和，
即让区间的左端点固定，右端点变动，则
\[S_f(x)=\text{函数$f$从$a$到$x$的和}=\int^x_a f (x) \dif x\]
可以看作上限变量$x$的函数，在这里，积分符号中的$x$既表示积分变量，又表示积分上限，容易混淆，因此，为了区别起见，我们用字母$t$来代表积分变量，这样上式就写成
\[S_f (x)=\int^x_a f(t)\dif t\]

和函数$S_f(x)$在$x=x_0$处的值$S_f(x_0)$的几何意义就是曲线$y=f(t)$, 直线$t=a$, $t=x_0$, $y=0$所围成的区域的有号面积，它是随区域的变动界线$t=x_0$的变动而变动的（图4.1）．

\begin{figure}[htp]
    \centering
    \begin{tikzpicture}[>=latex]
        \draw[->](-3,0)--(5,0)node[right]{$t$};
        \draw[->](0,-1.5)--(0,2)node[right]{$y$};
    \node at (0,0) [below left]{$O$};
    \node at (0,0){};
        \draw[line width=1.5pt] plot[smooth] coordinates{ (-2.75,-.5)(-2,-1)(-.8,0)(1,1.25)(2.5,0)(3.5,-1.1)(4,-1.3) 
    };
    \draw(-2,-1)node[below]{$M'$}--node[left]{$t=x_0'$}(-2,0)node[above]{$P'$};
    \draw(1,1.25)node[above]{$B$}--node[right]{$t=a$}(1,0)node[below]{$A$};
    \draw(3.5,-1.1)node[below]{$M$}--node[right]{$t=x_0$}(3.5,0)node[above]{$P$};
    \node at (-.8,0)[above]{$D$};
    \node at (2.5,0)[above]{$C$};
    \node at (2,1)[right]{$y=f(t)$};
    
    \end{tikzpicture}
    \caption{}
\end{figure}

例如，当变动界限（积分上限）在图中的$PM$位置时，则
\[S_f(x_0)=\int^{x_0}_a f(t)\dif t=ACB\text{的面积}-CPM\text{的面积}\]

当变动界限在图中的$P'M'$位置时，则
\[\begin{split}
    S_f(x'_0)&=\int^{x_0'}_a f(t)\dif t=-\int^a_{x_0'} f(t)\dif t\\
    &=-\left\{-P'DM'\text{的面积}+DAB\text{的面积}  \right\}\\
    &=P'DM'\text{的面积}-DAB\text{的面积}
\end{split}\]

例如，折线函数
\[g(x)=\begin{cases}
    \frac{1}{2}(x-1) & x\in[1,2]\\
    \frac{3}{2}(x-2)+\frac{1}{2}(2-1) & x\in [2,3]\\
    (x-3)+\frac{1}{2}(2-1)+\frac{3}{2}(3-2) & x\in [3,4]
\end{cases}\]
的导函数（除去在折线段的那些交接点处不作定义外）是阶梯函数
\[g'(x)=\begin{cases}
    \frac{1}{2} & x\in [1,2)\\
    \frac{3}{2} & x\in (2,3)\\
    1& x\in (3,4]
\end{cases}\]
反过来，该阶梯函数的和函数，是上述折线函数$g(x)$, 我们有如下的图解关系：
\begin{center}
    \begin{tikzpicture}[>=latex]
        \node (A) at (0,0) {\{折线函数\}};
        \node (B) at (5,0) {\{阶梯函数\}};
\draw[->](A) to [bend left=30]node[above]{求导} (B);
\draw[->](B) to [bend left=30]node[above]{求和} (A);

    \end{tikzpicture}
\end{center}

上述简明的例子表明“微分”与“积分”（或求函数由$a$到$b$的和）之间的运算关系应该是互逆的．

\subsection{微积分学基本定理}

\begin{blk}
  {定理1} 设$f(t)$是在$[a,b]$上的分段单调连续函数，又它的和函数是
\[S_f(x)=\int^x_a f (t) \dif t\]
那么
\[\frac{\dif }{\dif x}S_f(x)=\frac{\dif }{\dif x}\int^x_a f (t) \dif t=f(x),\qquad a\le x\le b\]
\end{blk}

\begin{figure}[htp]
    \centering
    \begin{tikzpicture}[>=latex, xscale=3]
        \draw[->](.4,0)--(3,0)node[right]{$t$};
        \draw[->](.6,-.5)--(.6,4.5)node[right]{$y$};
    \draw[domain=.9:2.5, very thick, samples=100]plot(\x, {2.5*(\x-1)*(\x-2)*(\x-3)+2})node[right]{$y=f(t)$};
    \draw(.9,0)node[below]{$a$}--(.9,1.4225);
    \draw[pattern=north east lines, domain=1.4:1.7, samples=50]plot(\x, {2.5*(\x-1)*(\x-2)*(\x-3)+2})--(1.7,2.6825)--(1.7,0)node[below]{$x$}--(1.4,0)node[below]{$x+\bar h$}--(1.4,2.96);
    \draw[pattern=north west lines, domain=1.7:2, samples=50]plot(\x, {2.5*(\x-1)*(\x-2)*(\x-3)+2})--(2,2)--(2,0)node[below]{$x+h$}--(1.7,0)--(1.7,2.6825);
    \draw[<-](1.55,2)--(1.55,3.5)node[right]{$S_f(x)-S_f(x+\bar h),\; (\bar h<0)$};
    \draw[<-](1.85,.5)--(2.1,.5)node[right]{$S_f(x+h)-S_f(x),\; (h>0)$};
    \node at (.6,0)[below left]{$O$};
    \end{tikzpicture}
    \caption{}
\end{figure}

\begin{proof}
  如图4.2，设$h>2$，则
\[\begin{split}
    \frac{S_f(x+h)-S_f(x)}{h}&=\frac{\Int^{x+h}_a f(t)\dd t-\Int^{x}_a f(t)\dd t}{h}\\
    &=\frac{\Int^{x+h}_x f(t)\dd t}{h}\qquad \text{（定积分性质1）}\\
    &=\frac{hf(\xi)}{h}\quad (x<\xi<x+h)\\
    &=f(\xi)\qquad \text{（定积分中值定理）}
\end{split}\]
于是
\begin{equation}
    \lim_{h\to 0^+}\frac{S_f(x+h)-S_f(x)}{h}=\lim_{\xi\to x}f(\xi)=f(x)
\end{equation}

设$\bar h<0$，则
\[\begin{split}
    \frac{S_f(x)-S_f(x+\bar h)}{\bar h}&=\frac{\Int^{x}_a f(t)\dd t-\Int^{x+\bar h}_a f(t)\dd t}{\bar h}\\
    &=\frac{\Int^{x}_{x+\bar h} f(t)\dd t}{\bar h}\qquad \text{（定积分性质1）}\\
    &=\frac{-\bar h f(\bar\xi)}{\bar h}=-f(\bar\xi)\quad (x+\bar h<\bar \xi<x)\\
\end{split}\]
即：
\[\frac{S_f(x+\bar h)-S_f(x)}{\bar h}=f(\bar\xi),\qquad (x+\bar h<\bar \xi<x)\]
于是
\begin{equation}
    \lim_{h\to 0^-}\frac{S_f(x+\bar h)-S_f(x)}{\bar h}=\lim_{\bar\xi\to x}f(\bar\xi)=f(x)
\end{equation}
由(4.1)和(4.2)，得
\[\frac{\dd }{\dd x}S_f(x)=\frac{\dd }{\dd x}\int^x_a f(t)\dd t=f(x)\]
\end{proof}

\begin{blk}
    {引理} 如果两个函数$F(x)$和$G(x)$满足条件
\[F' (x) =G' (x) ,\qquad a\le x\le b\]
那么
\[G(x)=F(x)+C,\qquad (C\text{是常数})\]
\end{blk}

\begin{proof}
    由$G'(x)-F'(x)=0$, 得
\[[G (x) -F (x) ]'=0\]
根据中值定理（参看第三章）知，对于每一个$x\in (a,b)$, 恒有
\[G(x)-F(x)=C,\qquad (C\text{是常数})\]
所以
\[G (x) =F (x) +C\]
\end{proof}

\begin{blk}
 {定义} 设$f(x)$是在区间$[a,b]$上给定的函数，如果可微函数$G(x)$满足条件
\[G'(x)=f(x)\qquad (\text{或}\quad \dd G(x)=f(x)\dd x)\]
那么就称$G(x)$为$f(x)$在$[a,b]$上的一个\textbf{原函数}．
\end{blk}

例如，设$G'(x)=\cos x$, 那么$G(x)=\sin x$就称为$\cos x$的一个原函数．

\begin{blk}
    {定理2} 设$f(x)$是给定的分段单调连续函数，如果一个可微函数$G(x)$满足条件
\begin{equation}
    G'(x)=f(x)\qquad \text{（或$\dd G(x)=f(x)\dd x$）},\quad a\le x\le b
\end{equation}    
那么
\[S_f (x)=\int^x_a  f (t) \dd t=G (x) -G (a) \]
\end{blk}

\begin{proof}
    由定理1知，$f(x)$的一个原函数存在，就是它的和函数，即
\[\frac{\dd }{\dd x}\int^x_a f(t)\dd t=f(x)\]
题设$G(x)$是$f(x)$的另一个原函数，故由引理得
\begin{equation}
    \int^x_a    f (t) dt=G (x) +C 
\end{equation}
这里的常数要由条件$\Int^a_a f(t)\dd t=0$来确定．

在(4.4)中，令$x=a$，得
\[0=G(a)+C\]
所以：$C=-G(a)$，代入(4.4)，得
\[\int^x_a f(t)\dd t=G(x)-G(a)\]
\end{proof}

在上面等式中，令$x=b$, 便得到下面著名的莱布尼兹-牛顿公式：
\[\int^b_a f (t) \dd t= G (b) -G (a) \]
这个公式告诉我们，求分段单调连续函数$f(x)$在$[a,b]$上的定积分可以化简为去求$f(x)$的一个原函数$G(x)$, 而函数值的差$G(b)-G(a)$就是定积分的值．

\begin{example}
    设$f(x)=k$（常数函数），则$G(x)=kx$是一个满足$G'(x)=k=f(x)$的原函数，所以
\[\int^b_a k\dd x=kb-ka=k(b-a)\]
以后我们引用一个符号：
\[G(x)\Bigg|^b_a=G(b)-G(a)\]
\end{example}

\begin{example}
    设$f(x)=kx^n,\; (n\in\mathbb{N})$，则
\[G(x)=k\frac{x^{n+1}}{n+1}\]
是满足$G'(x)=kx^n$的一个原函数，所以
\[\int^b_akx^n\dd x=k\frac{x^{n+1}}{n+1}\Bigg|^b_a=\frac{k}{n+1}(b^{n+1}-a^{n+1})\]
\end{example}

\begin{example}
设$f(x)=x^{-n}=\frac{1}{x^n}\; (x\ne 0,\; n\in\mathbb{N})$，则
\[G(x)=\frac{x^{-n+1}}{-n+1}\]
是满足$G'(x)=x^{-n}$的一个原函数．要特别注意定理2中$f(x)$是连续的条件，在任何包含0的区间上，函数$f(x)=\frac{1}{x^n}$无界，即有第二类间断点，如果$a,b$异号，则定积分$\Int^b_a x^{-n}\dd x$无意义，但若$a,b$同时为正或同时为负，则
\[\int^b_ax^{-n}\dd x=\frac{x^{-n+1}}{-n+1}\Bigg|^b_a=\frac{b^{-n+1}}{-n+1}-\frac{a^{-n+1}}{-n+1}\]
当然，又有$n\ne -1$时，上式才能成立．

如果$n=-1$，依复合函数求导法则，得
\[(\ln|x|)'=\frac{1}{|x|}\cdot |x|'\]
若$x>0$，则
\[(\ln|x|)'=\frac{1}{x}\cdot 1=\frac{1}{x}\]
若$x<0$，则
\[(\ln|x|)'=\frac{1}{-x}\cdot (-1)=\frac{1}{x}\]
总之
\[(\ln|x|)'=\frac{1}{x},\quad (x\ne 0)\]
因此
\[\begin{split}
\int^b_a \frac{1}{x}\dd x=\ln |x|\Bigg|^b_a &=\ln|b|-\ln |a|\\
&=\ln\left|\frac{b}{a}\right|=\ln\frac{b}{a}\quad (x\ne 0,\; ab>0)
\end{split}\]
\end{example}

\begin{example}
    求$\Int^3_1 |x-2|\dd x$
\end{example}

\begin{solution}
\[\begin{split}
    \Int^3_1 |x-2|\dd x&=\Int^2_1 -(x-2)\dd x+\Int^3_2 (x-2)\dd x\\
&=\left(-\frac{x^2}{2}+2x\right)\Bigg|^2_1+\left(\frac{x^2}{2}-2x\right)\Bigg|^3_2\\
&=\frac{1}{2}+\frac{1}{2}=1
\end{split}\]
\end{solution}

\begin{example}
求$\Lim_{n\to\infty}\frac{\sum^n_{i=1}i^p}{n^{p+1}}\quad (p\in \mathbb{N})$
\end{example}


\begin{solution}
\[\begin{split}
    \Lim_{n\to\infty}\frac{\sum^n_{i=1}i^p}{n^{p+1}}&=\Lim_{n\to\infty}\sum^n_{i=1}\left(\frac{i}{n}\right)^p\cdot \frac{1}{n}\\
    &=\int^1_0 x^p\dd x=\left.\frac{x^{p+1}}{p+1}\right|^1_0=\frac{1}{p+1}
\end{split}\]
\end{solution}

\begin{ex}
\begin{enumerate}
\item 用求原函数的方法, 计算下列定积分:
\begin{multicols}{2}
\begin{enumerate}
 \item $\Int_{2}^{3}\left(x^{3}+6 x^{2}+4 x+1\right) \dd x$
\item $\Int_{1}^{5} \frac{v^{2}+2}{v^{4}} \dd v$
\item $\Int_{-3}^{-2}\left(x+\frac{2}{x}\right)^{2} \dd x$
\item $\Int_{1}^{10}\left(x+\frac{1}{x}\right) \dd x$
\item $\Int_{a}^{b}|x-1| \dd x$
\end{enumerate}
\end{multicols}
\item 求曲线 $y=2 x-x^{2}$ 与 $x$ 轴之间的面积.
\item 设 $y=x(x-1)(x-2)$，
求 $\Int_{0}^{1} y \dd x,\; \Int_{1}^{2} y\dd x,\; \Int_{0}^{2} y \dd x$,
并说明所得结果的几何意义．

\item 求直线$y=2x$和曲线$y=x^2$之间的面积．
\item 求由曲线$y=\frac{1}{x^2}$和直线$x=1$, $x=4$所围成的曲边梯形的
面积．
\item 求抛物线方程使它具有以下的性质：
\begin{enumerate}
\item 抛物线通过$(0, 0)$, $(1, 2)$两点，
\item 它的对称轴平行$y$轴，且向上凸，
\item 它与$x$轴所围的面积最小．
\end{enumerate}

\item 求正数$a$, 使$\Int^1_0 |x^2-a^2|\dd x$最小．
\item 求下列函数的导数：
\begin{multicols}{2}
    \begin{enumerate}
     \item $F(x)=\Int_{a}^{x}\sin^2 t \dd t$
    \item $F(x)=\Int_{a}^{x^2} \sin^2 t \dd t$
    \end{enumerate}
    \end{multicols}
\end{enumerate}

\end{ex}

\section{不定积分}
由上节的微积分基本定理得知，在$[a,b]$上的两个可
微函数$G_1(x)$和$G_2(x)$均以$f(x)$为其导数，则$G_1(x)$与$G_2(x)$在$[a,b]$上只差一个常数$C$. 由此可见，如果$G'(x)=f(x)$, 且$C$为任意常数，那么$[G(x)+C]'=G'(x)=f(x)$. 这就是说：如果 函数$f(x)$有一个原函数$G(x)$, 那么就有无穷多个原函数，并且所有原函数刚好组成函数族
\[\{G(x)+C\;|\; C\text{是任意常数}\}\]

\begin{blk}
   {定义} 函数$f(x)$在区间$[a,b]$上的原函数全体叫做$f(x)$的不定积分，记作
   \[\int f(x)\dd x\]
   设$G(x)$是$f(x)$的一个原函数，那么由不定积分的定义得到
\[\int f (x) \dd x=G (x) +C\]
   其中$C$是任意常数．
\end{blk}   

关于不定积分，易见有如下性质：

\begin{blk}{性质1}
   积分与微分互为逆运算，即：
\[\begin{cases}
    \dd\Int f(x)\dd x=f(x)\dd x\\
    \Int \dd f(x)=f(x)+C
\end{cases}\]
\end{blk}

事实上，因为$\frac{\dd }{\dd x}\Int^x_a f(t)\dd t=f(x)$，所以
\[\int f(x)\dd x=\int^x_a f(t)\dd t+C\qquad \text{（基本定理1和不定积分定义）}\]
于是
\[\left(\int f(x)\dd x\right)'=\left(\int^x_a f(x)\dd x\right)'=f(x)\]
从而
\[\dd\int f(x)\dd x=f(x)\dd x\]
又
\[\begin{split}
    \int \dd f(x)&=\int f'(x)\dd x=\int^x_a f'(t)\dd t+C_1\\
    &=f(x)-f(a)+C_1\quad \text{（基本定理）}\\
    &=f(x)+C
\end{split}\]
这里$C=-f(a)+C_1$. 

从性质1看出，若不要不定积分等式中的任意常数项，则当符号$\Int$与$\dd$紧接着时，无论哪个在前，哪个在后，都可以消掉．

因为微分，积分是两种互逆的运算，所以它们的运算法则是互相对应的，现将它们相对应的法则叙述如下：

\begin{blk}{性质2}
 \[\begin{cases}
     {\rm d} [f(x)+g(x)]=\dd f(x)+\dd g(x)\\
     \Int [f(x)+g(x)] \dd x=\Int f(x)\dd x+\Int g(x)\dd x+C
 \end{cases}\]
 \end{blk}

 上面不定积分公式的正确性不难证明，只要求出两边的微商，由于所得到的微商恒等，就可以知道它们只差一个常数项，因为，由性质1
\[\begin{split}
    \left(\int[f(x)+g(x)]\dd x\right)'&=f(x)+g(x)\\
    \left(\int f(x)\dd x+\int g(x)\dd x\right)'&=\left(\int f(x)\dd x\right)'+\left(\int g(x)\dd x\right)'=f(x)+g(x)
\end{split}\]
所以$f(x)+g(x)$的两个原函数相差一个常数项，即：
\[\int[f(x)+g(x)]\dd x=\int f(x)\dd x+\int g(x)\dd x+C\]
同理可得：

\begin{blk}{性质3}
\[\begin{cases}
    {\rm d} [a f(x)]=a\dd f(x)\\
    \Int af(x)\dd x=a\Int f(x)\dd x
\end{cases}\qquad (a\ne 0)\]
\end{blk}

\subsection{基本积分表}

由基本微商公式表，倒转顺序，就可以得到下面这个表：
\begin{itemize}
    \item $\frac{\dd}{\dd x}C=0,\qquad \Int 0\dd x=C$
    \item $\frac{\dd }{\dd x}x^{\mu}=\mu x^{\mu-1},\qquad \Int x^{\mu}\dd x=\frac{x^{\mu+1}}{\mu+1}+C\quad \text{（$\mu$是不等于$-1$的实数）}$
    \item $\frac{\dd }{\dd x}\ln|x|=\frac{1}{x},\qquad \Int\frac{\dd x}{x}=\ln|x|+C\quad (x\ne 0)$
    \item $\frac{\dd }{\dd x}e^x=e^x,\qquad \Int e^x\dd x=e^x+C$
\end{itemize}
还因为$a^x=e^{x\ln a}$，由复合函数求导法则，得：
\[(a^x)'=e^{x\ln a}(x\ln a)'=\ln a\cdot e^{x\ln a}=\ln a\cdot a^x\]
所以：
\begin{itemize}
    \item $\frac{\dd}{\dd x}a^x=\ln a\cdot a^x,\qquad \Int a^x\dd x=\frac{a^x}{\ln a}+C$
    \item $\frac{\dd}{\dd x} \sin x=\cos x, \qquad \Int \cos x \dd x=\sin x+C$.
    \item $\frac{\dd}{\dd x} \cos x=-\sin x, \qquad \Int \sin x \dd x=-\cos x+C$.
    \item  $\frac{\dd}{\dd x}\tan  x=\sec ^{2} x, \qquad \Int \sec^{2} x \dd x=\tan x+C$
    \item $\frac{\dd}{\dd x}\cot x=-\csc ^{2} x, \qquad \Int \csc^{2} x \dd x=-\cot x+C$
    \item $\frac{\dd}{\dd x} \arcsin x=\frac{1}{\sqrt{1-x^{2}}} ,\qquad \Int \frac{\dd x}{\sqrt{1-x^{2}}}=\arcsin x+C$ 
    
    $\frac{\dd}{\dd x} \arccos x=\frac{-1}{\sqrt{1-x^{2}}},\qquad \Int\frac{-1}{\sqrt{1-x^{2}}}\dd x=-\arccos x+C'$
    \item  $\frac{\dd}{\dd x}\arctan x=\frac{1}{x^{2}+1}, \qquad \Int \frac{\dd x}{x^{2}+1}=\arctan x+C$ 
    
    $\frac{\dd}{\dd x} \arccot x=\frac{-1}{x^{2}+1}, \qquad \Int \frac{-1}{x^{2}+1}\dd x=-\arccot  x+C'$
\end{itemize}

现在来看不定积分的几何意义．

由不定积分的定义知道，求函数$f(x)=3x^2$的原函数全体，就是求不定积分$\Int 3x^2\dd x$, 也就是求满足微分方程
\begin{equation}
    \frac{\dd y}{\dd x}=3x^2
\end{equation}
的所有的原函数．

显然，$y=x^3$就是微分方程(4.5)的一个原函数，它的图象上各点切线的斜率由(4.5)确定．我们称曲线$y=x^3$为
微分方程(4.5)的积分曲线，把它沿$y$轴上下平移任意一段距离，便得到曲线族
\begin{equation}
    y=x^3+C
\end{equation}
并且(4.6)中任意一条曲线与曲线$y=x^3$在横坐标相同的点的切线平行，其斜率都等于$3x^2$, 因此，曲线族(4.6)是满足微分方程(4.5)的所有积分曲线，也就是不定积分$\Int 3x^2\dd x=x^3+C$的几何表示（图4.3）．

为要确定(4.5)的某条积分曲线的位置，应当给出这条积分曲线必须经过的一点，例如它经过点$P(1, 3)$. 我们说这条件也是微分方程(4.5)的一个相应的原函数要满足的初值条件
$y\big|_{x=1}=3$
把它们代入(4.6), 得到
\[3=1^3+C\quad \Rightarrow\quad C=2\]
因此，满足微分方程$\frac{\dd y}{\dd x}=3x^2$和初值条$y\big|_{x=1}=3$的原函
数或相应的积分曲线是
\[y=x^3+2\]

通常把满足微分方程
\begin{equation}
\frac{\dd y}{\dd x}=f(x)
\end{equation}
的所有解叫做给定函数$f(x)$的不定积分，即
\[y=\int f(x)\dd x\]
把函数$f(x)$的一个原函数$y=G(x)$的图象叫做微分方程(4.7)的一条\textbf{积分曲线}，这样，不定积分
\[y=\int f(x)\dd x=G(x)+C \quad \text{（$C$为任意常数）}\]
的几何意义就表示与$y=G(x)$在各点有相同斜率$\frac{\dd y}{\dd x}=f(x)$的积分曲线族（图4.3）.
\begin{figure}[htp]
    \centering
\begin{tikzpicture}[>=latex, xscale=3, yscale=.8]
\draw[->](-1,0)--(1.25,0)node[right]{$x$};
\draw[->](0,-4)--(0,4)node[right]{$y$};
\node at (0,0)[below left]{$O$};
\foreach \y in {1,2,3}
{
    \draw[domain=-.8:.8, samples=100, very thick]plot(\x, {\x*\x*\x+\y})node[right]{$y=x^3+$\y};
    \draw[domain=-.8:.8, samples=100, very thick]plot(\x, {\x*\x*\x-\y})node[right]{$y=x^3-$\y};
}
\draw[domain=-.8:.8, samples=100, very thick]plot(\x, {\x*\x*\x})node[right]{$y=x^3$};
\end{tikzpicture}
    \caption{}
\end{figure}

由微积分基本定理2知道，如果$y=G(x)$满足下面两个条件：
\[G'(x)=f(x)\quad \text{且}\quad G(a)=0\]
那么，$y=G(x)$就是$f(x)$的求和函数$G(x)=\Int^x_a f(x)\dd x$.

\begin{example}
求$\Int \sin (kx)\dd x,\quad \Int \cos(kx)\dd x,\quad \Int e^{kx}\dd x,\quad \Int \frac{\dd x}{x^2+a^2}$
\end{example}

\begin{solution}
    因为
\[\begin{split}
    (\cos kx)'&=-k\sin kx\\
    (\sin kx)'&=k\cos kx\\
    (e^{kx})'&=ke^{kx}\\
    \left(\arctan\frac{x}{a}\right)'&=\frac{1}{\left(\frac{x}{a}\right)^2+1}\cdot\left(\frac{x}{a}\right)'=\frac{a}{x^2+a^2}
\end{split}\]
所以
\[\begin{split}
    \int \sin (kx)\dd x&=-\frac{1}{k}\cos kx+C\\
    \Int \cos(kx)\dd x&=\frac{1}{k}\sin kx+C\\
    \Int e^{kx}\dd x&= \frac{1}{k}e^{kx}+C\\
    \Int \frac{\dd x}{x^2+a^2}&=\frac{1}{a}\arctan\frac{x}{a}+C \\
\end{split}  \]
\end{solution}

\begin{example}
    证明：若$\Int f(t)\dd t=F(t)+C$，则
\[\int f(kx)\dd x=\frac{1}{k}F(kx)+C\]
\end{example}

\begin{proof}
只须证明$f(kx)$的一个原函数是$\frac{1}{k}F(kx)$即可．
我们求$\frac{1}{k}F(kx)$对$x$的导数，于是
\[\frac{\dd}{\dd x}\left[\frac{1}{k}F(kx)\right]=\frac{1}{k}F'(kx)\cdot (kx)'=F'(kx)\]
由题设知$F' (t) =f (t)$, 所以
\[\frac{\dd}{\dd x}\left[\frac{1}{k}F(kx)\right]=F'(kx)=f(kx)\]
因此
\[\int f (kx) \dd x=\frac{1}{k}F (kx) +C\]
\end{proof}

\begin{example}
求$\Int \sin kx\sin\ell x\dd x$
\end{example}

\begin{solution}
\[\begin{split}
    \Int \sin kx\sin\ell x\dd x&=\int \frac{1}{2}[\cos(k-\ell)x-\cos(k+\ell)x]\dd x\\
    &=\frac{1}{2}\int \cos (k-\ell)x\dd x-\frac{1}{2}\int \cos (k+\ell)x\dd x\\
&=\frac{\sin(k-\ell)x}{2(k-\ell)}-\frac{\sin(k+\ell)x}{2(k+\ell)}+C
\end{split}\]
    这里$k\ne \pm \ell$．（如果$k=\pm\ell$，则另当考虑）
\end{solution}


\begin{example}
求$\Int \left(\frac{2a}{\sqrt{x}}-\frac{b}{x^2}+2xc^{\tfrac{2}{3}}\right)\dd x$
\end{example}


\begin{solution}
    \[\begin{split}
        \Int \left(\frac{2a}{\sqrt{x}}-\frac{b}{x^2}+2xc^{\tfrac{2}{3}}\right)\dd x&=2a\int x^{-\tfrac{1}{2}}\dd x-b\int x^{-2}\dd x+3c\int x^{\tfrac{2}{3}}\dd x\\
        &=4a\sqrt{x}+\frac{b}{x}+\frac{9}{5}cx^{\tfrac{5}{3}}+C
    \end{split}\]
\end{solution}

\begin{example}
    求$\Int \frac{3x^2}{1+x^2}\dd x$
\end{example}

\begin{solution}
\[\begin{split}
    \Int \frac{3x^2}{1+x^2}\dd x&=3\int \left(1-\frac{1}{x^2+1}\right)\dd x\\
    &=3\int \dd x-3\int \frac{\dd x}{x^2+1}\\
    &=3x-3\arctan x+C
\end{split}\]
\end{solution}

以上各例都是利用不定积分的性质2、3, 同时配合使用三角恒等式或设法将被积函数拆成几项，从而把一个比较复杂的积分化成若干个可以查基本积分表的积分．

\begin{ex}
\begin{enumerate}
    \item 计算下列不定积分:
\begin{multicols}{2}
\begin{enumerate}
\item  $\Int\left(x^{3}+3 x+5\right) \dd x$
\item  $\Int(x+2)^{4} \dd x$
\item  $\Int\left(x^{-2}+x^{-1}+x^{\tfrac{1}{3}}+x^{\tfrac{1}{2}}\right) \dd x$
\item  $\Int\left(a^{\tfrac{2}{3}}-x^{\tfrac{2}{3}}\right)^{3} \dd x$
\item  $\Int\left(\frac{1}{x}+e^{x}\right) \dd x$
\item  $\Int \frac{2 \cdot 3^{x}-5 \cdot 2^{x}}{3^{x}} \dd x$
\item  $\Int \left(\frac{3}{1+x^{2}}-\frac{2}{\sqrt{1-x^{2}}}\right) \dd x$
\item  $\Int \frac{2+x^{2}}{1+x^{2}} \dd x$
\item  $\Int \frac{\dd x}{x^{2}\left(1+x^{2}\right)}$
\item  $\Int \frac{\dd x}{x^{2}\left(a^{2}+x^{2}\right)}$
\item  $\Int 3^{x} e^{x} \dd x$
\end{enumerate}
\end{multicols}

\item  利用三角恒等式, 计算下列不定积分:
    \begin{multicols}{2}
        \begin{enumerate}
    \item  $\Int \tan^{2} x \dd x$
    \item  $\Int \frac{\cos 2 x}{\cos x-\sin x} \dd x$
    \item  $\Int \frac{2+\sin ^{2} x}{\cos ^{2} x} \dd x$
    \item  $\Int \sin ^{2} x \dd x$
    \item  $\Int \cos ^{2} x \dd x$
    \item  $\Int \cos k x \cos \ell x \dd x$
    \item  $\Int \sin 2 x \cos 5 x \dd x$
    \item  $\Int \frac{\dd t}{\sin ^{2} t \cos ^{2} t}$
    \item  $\Int \sin ^{3} x \dd x$
    \item  $\Int \cos ^{3} x \dd x$
\end{enumerate}
\end{multicols}    

\item 若$\Int f(t)\dd t=F(t)+C$, 求证
\[\int f (ax+b) \dd x=\frac{1}{a}F (ax+b) +C\]
\item  若$\frac{\dd y}{\dd x}=4x-6$, 且当$x=0$时，$y=9$, 求$y$的极小值．
\item  若$\frac{\dd y}{\dd x}=ax+b$, 求原函数$y=f(x)$使其满足下列条件：
\begin{enumerate}
    \item 当$x=0$时，$y=5$;
    \item 积分曲线向下凸且曲线$y=f(x)$在$x$轴上截得的弦长等于4;
    \item 曲线$y=f(x)$上纵坐标的最小值等于$-4$.
\end{enumerate}
\item 试求一个次数最低的多项式，使之当$x=1$时，有最大值6，而当$x=3$时，有最小值2.
\item  求下列各式的极限：
\begin{enumerate}
    \item $\Lim_{n\to\infty}\frac{\sqrt[n]{e}+\sqrt[n]{e^2}+\cdots+\sqrt[n]{e^n}}{n}$
    \item $\Lim_{n\to\infty}\frac{\sqrt[n]{e}+\sqrt[n]{e^4}+\cdots+\sqrt[n]{e^{2n}}}{n}$
\end{enumerate}
\end{enumerate}

\end{ex}


\section{不定积分的计算方法}
从上一节看到，虽然利用积分运算法则及基本积分表可以求出不少函数的原函数，但是实际遇到的积分仅凭这一些
方法还不能完全解决，本节将介绍几种典型的方法，利用这些方法，我们就可以计算更多的不定积分，此外，我们还制作了较为详细的积分表，可供使用时查阅．

\subsection{第一换元法}

有一些不定积分，将积分变量进行某种变映后就能由基本积分公式求出，例如求$\Int e^{2x}\dd x$, 在基本积分公式中只有$\Int e^x\dd x$, 没有直接求出它的公式，我们看到所求积分的被积函数$e^{2x}$是$y=e^u$和$u=2x$的复合函数，如果在积分表达式$e^{2x}\dd x$中，把被积函数看成中间变量$u$的函数，即设$u=2x$则$e^{2x}=e^u$, 从而$\dd u=2\dd x$, 由此可见只要再给微分$\dd x$凑上一个常数因子2, 就得到中间变量$u$的微分$\dd u$, 于是
\[\int e^{2x}\dd x=\int e^{2x}\cdot \frac{1}{2}\dd (2x) =\frac{1}{2}\int e^u\dd u=\frac{1}{2}e^u+C\]
然后再代回原来的积分变量，得
\[\int e^{2x}\dd x=\frac{1}{2}e^{2x}+C\]

\begin{example}
    求$\Int \frac{\dd x}{\sqrt{a^2-x^2}}\quad (a>0)$
\end{example}

\begin{solution}
    因为被积函数$\frac{1}{\sqrt{a^2-x^2}}=\frac{1}{a\sqrt{1-\left(\frac{x}{a}\right)^2}}$
是$y=\frac{1}{a\sqrt{1-u^2}}$和$u=\frac{x}{a}$的复合函数，且$\dd u=\frac{1}{a}\dd x$，于是
\[\begin{split}
    \int \frac{\dd x}{\sqrt{a^2-x^2}}=\int \frac{\dd x}{a\sqrt{1-\left(\frac{x}{a}\right)^2}}&=\int \frac{\frac{1}{a}\dd x}{\sqrt{1-\left(\frac{x}{a}\right)^2}}\\
    \text{令}u=\frac{x}{a}\to &= \int \frac{\dd u}{\sqrt{1-u^2}}\\
    \text{查表}\to &=\arcsin u+C\\
    u=\frac{x}{a} \to &=\arcsin \frac{x}{a}+C
\end{split}\]
\end{solution}

\begin{example}
    求$\Int \tan x\dd x$
\end{example}


\begin{solution}
\[\begin{split}
    \Int \tan x\dd x=\Int  \frac{\sin x}{\cos x}\dd x&=-\int \frac{\dd \cos x}{\cos x}\\
    u=\cos x\to &=-\int \frac{\dd u}{u}\\
    &=-\ln|u|+C=-\ln |\cos x|+C=\ln |\sec x|+C
\end{split}\]
\end{solution}

从上述两例看到，在求不定积分时，首先要与已知的基本积分公式相对比，并利用简单的变量代换，把要求的积分凑成公式所具有的形式，求出以后，再把原来的变量代回．以上变量代换中关键的一步是把原来的被积函数$f(x)$的某部分$\varphi(x)$换成中间变量$u$, 那么原来的被积表达式$f(x)\dd x$就拆成了$g[\varphi(x)]=g(u)$和$\varphi'(x)\dd x=\dd u$这两部分的乘积，即
\[f (x) \dd x=g [\varphi (x) ] \varphi' (x) \dd x=g (u) \dd u\]
于是：若$\Int g(x)\dd x=F(x)+C$，那么
\[\int g[\varphi(x)]\varphi'(x)\dd x=\int g(u)\dd u=F(u)+C=F[\varphi(x)]+C\]
上述换元法叫做第一换元法．

\begin{example}
求$\Int \sin^{3}x\dd x$
\end{example}

\begin{solution}
设$u=\cos x$，$\dd u=-\sin x\dd x$，于是
\[\begin{split}
    \Int \sin^{3}x\dd x&=\int (1-\cos^2 x)\cdot \sin x\dd x\\
  u=\cos x\to  &=-\int (1-u^2)\dd u\\
  &=-u+\frac{u^3}{3}+C=\frac{\cos^3 x}{3}-\cos x+C
\end{split}\]
\end{solution}

对换元积分法较熟练以后，所设换元变量$u$可以不写出．

\begin{example}
求$\Int \frac{\dd x}{\sqrt{ax+b}}$
\end{example}

\begin{solution}
    \[\Int \frac{\dd x}{\sqrt{ax+b}}=\frac{1}{a}\int \frac{\dd(ax+b)}{\sqrt{ax+b}}=\frac{2}{a}\sqrt{ax+b}+C\]
\end{solution}

\begin{example}
    求$\Int \frac{\dd x}{a^2-x^2}$
\end{example}

\begin{solution}
\[\begin{split}
    \Int \frac{\dd x}{a^2-x^2}&=\frac{1}{2a}\int\left(\frac{1}{a-x}+\frac{1}{a+x}\right)\dd x\\
    &=\frac{1}{2a}\left[-\int\frac{\dd(a-x)}{a-x}+\int\frac{\dd(a+x)}{a+x}\right]\\
    &=\frac{1}{2a}\left[-\ln|a-x|+\ln|a+x|\right]=\frac{1}{2a}\ln\left|\frac{a+x}{a-x}\right|
\end{split}\]
\end{solution}

\begin{example}
    求$\Int \sec x\dd x$
\end{example}

\begin{solution}
    \[\Int \sec x\dd x=\int \frac{\cos x}{\cos^2 x}\dd x=\int \frac{\dd(\sin x)}{1-\sin^2 x}\]
利用上题结果，得：
\[\begin{split}
    \int \sec x\dd x&=\frac{1}{2}\ln\left|\frac{1+\sin x}{1-\sin x}\right|+C\\
    &=\frac{1}{2}\ln\frac{1+\sin x}{1-\sin x}+C\\
    &=\frac{1}{2}\ln\frac{(1+\sin x)^2}{\cos^2 x}+C\\
    &=\ln|\sec x+\tan x|+C
\end{split}\]

\end{solution}

\begin{example}
    求$\Int \frac{1}{\sin^2x +3\cos^2x}\dd x$
\end{example}

\begin{solution}
    \[\begin{split}
        \Int \frac{1}{\sin^2x +3\cos^2x}\dd x&=\Int \frac{1}{(\tan^2x +3)\cos^2x}\dd x\\
        &=\int \frac{\dd (\tan x)}{\tan^2 x+\left(\sqrt{3}\right)^2}=\frac{1}{\sqrt{3}}\int \frac{\dd\left(\frac{\tan x}{\sqrt{3}}\right)}{\left(\frac{\tan x}{\sqrt{3}}\right)^2+1}\\
        &=\frac{1}{\sqrt{3}}\arctan\left(\frac{\tan x}{\sqrt{3}}\right)+C      
    \end{split}\]
\end{solution}

\begin{ex}
求下列不定积分：
\begin{enumerate}
    \begin{multicols}{2}
\item  $\Int(3 x+5)^{100} \dd x$
\item  $\Int\pi e^{-\tfrac{\pi}{4} x} \dd x$
\item  $\Int \frac{\ln x}{x} \dd x$
\item  $\Int\frac{1}{x^{2}} \cdot e^{\tfrac{1}{x}} \dd x$
\item  $\Int\frac{\cos \sqrt{x}}{\sqrt{x}} \dd x$
\item  $\Int \frac{1}{(\arcsin x)^{2} \sqrt{1-x^{2}}} \dd x$
\item  $\Int\frac{\dd x}{x^{2}+a^{2}}$
\item  $\Int \frac{\dd x}{x^{2}-a^{2}}$
\item  $\Int \frac{1}{\sqrt{x}} e^{\sqrt{x}} \dd x$
\item  $\Int \frac{2 a x+b}{2 \sqrt{a x^{2}+b x+c}} \dd x$
\item  $\Int\frac{1}{1+e^{x}} \dd x$        
    \end{multicols}
    \begin{multicols}{2}
\item  $\Int \frac{e^{x}}{e^{2 x}+1} \dd x$    
\item $\Int \frac{\cos x}{e^{\sin x}} \dd x$
\item $\Int \frac{\dd x}{\cos ^{2} x \sqrt{\tan x}}$
\item $\Int  \frac{\cos a x}{b-\sin a x} \dd x\quad (a \neq 0) $
\item $\Int x \sqrt{a x^{2}+c}\; \dd x$
\item $\Int \frac{\sin ^{2} x \cos x \dd x}{1+\sin x}$
\item $\Int \frac{\sqrt{\arctan x}}{1+x^{2}} \dd x$
\item $\Int \frac{1}{1+\sin x} \dd x$        
    \end{multicols}
    \begin{multicols}{2}
\item $\Int \frac{\sqrt{1+\cos x}}{\sin x} \dd x$
\item $\Int \frac{\dd x}{x^{2}+x-6}$
\item $\Int \frac{x \dd x}{3 x^{2}+x+2}$
\item $\Int  \frac{\dd x}{3 x^{2}+x+2}$
\item $\Int  \cot x \dd x$
\item $\Int  \frac{\dd x}{\sin x}$
\item $\Int \frac{\arcsin x}{\sqrt{1-x^{2}}} \dd x$        
    \end{multicols}

\end{enumerate}
\end{ex}

\subsection{第二换元积分法}

有些积分不能很容易地凑出微分，而是一开始就要作代换，把要求的积分化简，然后再求积分．

\begin{example}
    求$\Int \frac{1}{1+\sqrt{x}}\dd x$.
\end{example}

\begin{solution}
    这个积分的表达式不容易拆成中间变量$u=\varphi(x)$
的函数$g(\varphi(x))=g(u)$和中间变量的微分$\varphi'(x)\dd x=\dd u$的乘积，因此前面的第一换元法对此题用起来不便．这时我们设法作一个代换，以便把被积函数中的根号去掉，化成新变量的有理式．

令$x=u^2$, 即作代换$\sqrt{x}=u$, 于是
\[\begin{split}
    \frac{1}{1+\sqrt{x}}&=\frac{1}{1+u}\\
    \dd x&=2u\dd u
\end{split}\]
从而
\[\begin{split}
    \Int \frac{1}{1+\sqrt{x}}\dd x&=\int \frac{1}{1+u}2u\dd u\\
    &=2\int \frac{(1+u)-1}{1+u}\dd u\\
    &=2\left[\int \dd u-\int \frac{1}{1+u}\dd u\right]\\
    &=2\left[u-\ln |1+u|\right]+C\\
    &=2\left[\sqrt{x}-\ln\left(1+\sqrt{x}\right)\right]+C
\end{split}\]
\end{solution}

这里由于把积分变量$x$换成了$u$的函数$x=u^2$, 因此在最后结果中，用变量$u$表出的函数，还要还原成$x$的函数，这就要求从所作的代换$x=u^2$中解出反函数$u=\sqrt{x}$来．

我们把第二换元积分法用定理形式叙述如下：

\begin{blk}{定理}
设$x=\varphi(u)$是可微函数，且$\varphi'(u)\ne 0$，若
\begin{equation}
    \int f[\varphi(u)]\varphi'(u)\dd u=F(u)+C
\end{equation}
那么
\begin{equation}
    \int f(x)\dd x=F\left[\varphi^{-1}(x)\right]+C
\end{equation}
\end{blk}

\begin{proof}
只须证明$\frac{\dd}{\dd x}F[\varphi^{-1}(x)]=f(x)$即可．

因为$\varphi'(u)\ne 0$，故$x=\varphi(u)$的反函数$u=\varphi^{-1}(x)$存在，于是
\[\begin{split}
    \frac{\dd}{\dd x}F[\varphi^{-1}(x)]&=\frac{\dd}{\dd u}F[\varphi^{-1}(x)]\cdot \frac{\dd }{\dd x}\varphi^{-1}(x)\\
    &=\frac{\dd}{\dd u}F(u)\cdot \frac{1}{\frac{\dd}{\dd u}\varphi(u)}\\
    &=F'(u)\cdot \frac{1}{\varphi'(u)}\\
    &=\left\{f[\varphi(u)]\varphi'(u)\right\}\cdot \frac{1}{\varphi'(u)}\\
    &=f[\varphi(u)]\\
    x=\varphi(u)\to &=f(x)
\end{split}\]
\end{proof}


\begin{example}
    求$\Int\frac{\dd x}{\sqrt{x}+\sqrt[3]{x}}$.
\end{example}

\begin{solution}
    \[\Int\frac{\dd x}{\sqrt{x}+\sqrt[3]{x}}=\int \frac{\dd x}{\left(\sqrt[6]{x}\right)^3+\left(\sqrt[6]{x}\right)^2}\]
设$x=u^6$，$u>0$，则可化去被积函数中的根号，于是
\[\frac{1}{\sqrt{x}+\sqrt[3]{x}}=\frac{1}{u^3+u^2}\quad \Rightarrow\quad \dd x=6u^5\dd u\]
从而
\[\begin{split}
    \Int\frac{\dd x}{\sqrt{x}+\sqrt[3]{x}}&=\int \frac{6u^5}{u^3+u^2}\dd u=\int\frac{6u^3}{u+1}\dd u\\
    &=6\int \left(u^2-u+1-\frac{1}{u+1}\right)\dd u\\
    &=6\left\{\frac{u^3}{3}-\frac{u^2}{2}+u-\ln|u+1|\right\}+C\\
    &=2\sqrt{x}-3\sqrt[3]{x}+6\sqrt[6]{x}-\ln\left(\sqrt[6]{x}+1\right)+C
\end{split}
    \]
\end{solution}

\begin{example}
    求$\Int \frac{x+1}{\sqrt[3]{3x+1}}\dd x$.
\end{example}

\begin{solution}
    设$3x+1=t^3$，即$x=\frac{1}{2}(t^3-1)$，于是：$\dd x=t^2\dd t$
\[\begin{split}
    \Int \frac{x+1}{\sqrt[3]{3x+1}}\dd x&=\int \frac{\frac{1}{3}(t^3-1)+1}{t}t^2\dd t\\
&=\frac{1}{3}\int (t^4+2t)\dd t\\
&=\frac{1}{3}\left(\frac{1}{5}t^5+t^2\right)+C=\frac{1}{15}(3x+1)^{\tfrac{5}{3}}+\frac{1}{3}(3x+1)^{\tfrac{2}{3}}+C\\
&=\frac{1}{5}(x+2)(3x+1)^{\tfrac{2}{3}}+C
\end{split}\]
\end{solution}

上面所介绍的两种变量代换法的基本思想是一致的．第一种是把被积函数中的一个小部分看作一个变量，即“化繁为简”；第二种则把积分变量代以一个新变量的函数，表面上看来是“化简为繁”，但实际上克服了求积分的困难．应当注意，在使用第二种变量代换法时，要保证变换是可微的和一对一的，否则会出现错误.

下面介绍当被积函数含有二次根式$\sqrt{a^2-x^2}$, $\sqrt{a^2+x^2}$, $\sqrt{x^2-a^2}$时，怎样作代换．

\begin{example}
    求$\Int \sqrt{a^2-x^2}\dd x\quad (a>0)$.
\end{example}

\begin{solution}
因为\[\sqrt{a^2-x^2}=\sqrt{a^2\left(1-\frac{x^2}{a^2}\right)}=a\sqrt{1-\left(\frac{x}{a}\right)^2}\]
所以可作正弦代换
\[\frac{x}{a}=\sin t\quad \left(-\frac{\pi}{2}\le t\le \frac{\pi}{2}\right)\]
$\sin t$在$\left[-\frac{\pi}{2},\frac{\pi}{2}\right]$上连续且递增，有反函数$t=\arcsin\frac{x}{a},\; -1\le x\le 1$. 于是
\[\begin{split}
    \dd x&=\dd (a\sin t)=a\cos t\dd t\\
    \sqrt{a^2-x^2}&=a\sqrt{1-\left(\frac{x}{a}\right)^2}=a\sqrt{1-\sin^2 t}=a\cos t
\end{split}\]
从而：
\[\begin{split}
    \Int \sqrt{a^2-x^2}\dd x&=\int a\cos t\cdot a\cos t\dd t=a^2\int \cos^2 t\dd t\\ 
&=a^2\int \frac{1+\cos 2t}{2}\dd t=\frac{a^2}{2}\left[\int \dd t+\frac{1}{2}\int \cos 2t\dd (2t)\right]\\
&=\frac{a^2}{2}\left(t+\frac{1}{2}\sin 2t\right)+C=\frac{a^2}{2}\left(t+\sin t\cos t\right)+C\\
&=\frac{a^2}{2}\left[\arcsin\frac{x}{a}+\frac{x}{a}\sqrt{1-\left(\frac{x}{a}\right)^2}+C\right]\\
&=\frac{a^2}{2}\arcsin\frac{x}{a}+\frac{x}{2}\sqrt{a^2-x^2}+C
\end{split}\]
\end{solution}

\begin{example}
    求$\Int \frac{\dd x}{\sqrt{a^2+x^2}}\quad (a>0)$.
\end{example}

\begin{solution}
因为$\sqrt{a^2+x^2}=a\sqrt{1+\left(\frac{x}{a}\right)^2}$，设$\frac{x}{a}=\tan t\; \left(-\frac{\pi}{2}<t<\frac{\pi}{2}\right)$，则存在反函数
\[t=\arctan\frac{x}{a}\quad \Rightarrow\quad \dd x=a\sec^2 t\dd t\]
又$\sqrt{a^2+x^2}=a\sqrt{1+\tan^2 t}=a\sec t$
，从而：
\[\begin{split}
    \Int \frac{\dd x}{\sqrt{a^2+x^2}}&=\int\frac{a\sec^2 t\dd t}{a\sec t}=\int\sec t\dd t\\
    &=\ln|\sec t+\tan t|+C_1\\
    &=\ln\left|\frac{1}{a}\sqrt{a^2+x^2}+\frac{x}{a}\right|+C_1\\
    &=\ln\left|x+\sqrt{a^2+x^2}\right|-\ln a+C_1\\
    &=\ln\left|x+\sqrt{a^2+x^2}\right|+C\qquad (C=C_1-\ln a)
\end{split}\]
\end{solution}

\begin{example}
    求$\Int \frac{\dd x}{\sqrt{x^2-a^2}}\quad (a>0)$
\end{example}

\begin{solution}

    \begin{center}
\begin{tikzpicture}[scale=.7]
\draw[very thick](0,0)--node[below]{$a$}(3,0)--node[right]{$\sqrt{x^2-a^2}$}(3,4)--node[left]{$x$}(0,0);
\draw(.75,0) arc (0:53.13:.75)node[right]{$t$};
\end{tikzpicture}
    \end{center}

    作正割代换．设$x=a\sec t\; \left(0<t<\frac{\pi}{2}\right)$，此函数在$t=\frac{\pi}{2}$点间断，在区间$\left[0,\frac{\pi}{2}\right)$上由$a$递增到$+\infty$，而在区间$\left(\frac{\pi}{2},\pi\right]$上由$-\infty$递增到$-a$．因此，$x=a\sec t$在$\left(0,\frac{\pi}{2}\right]$上的反函数是 
\[t=\arcsec\frac{x}{a}\quad \left(1<\frac{x}{a}<\infty,\; \text{即}x>a\right)\]
于是，当$x>a$时，有
\[\begin{split}
    \sqrt{x^2-a^2}&=\sqrt{a^2\sec^2 t-a^2}=\sqrt{a^2(\sec^2 t-1)}\\
    &=a\sqrt{\tan^2 t}=a|\tan t|\\
    &=a\tan t \quad \left(0\le t<\frac{\pi}{2}\right)
\end{split}\]
\[\dd x=\dd(a\sec t)=a\frac{\sin t}{\cos^2 t}\dd t\]
从而：
\[\begin{split}
    \Int \frac{1}{\sqrt{x^2-a^2}}\dd x&=\int\frac{1}{a\tan t}\cdot \frac{a\sin t}{\cos^2 t}\dd t\\
    &=\int \frac{1}{\cos t}\dd t=\ln|\sec t+\tan t|+C_1\\
    &=\ln\left|\frac{x}{a}+\frac{\sqrt{x^2-a^2}}{a}\right|+C_1\\
    &=\ln\left|x+\sqrt{x^2-a^2}\right|+C_1-\ln a=\ln\left|x+\sqrt{x^2-a^2}\right|+C
\end{split}\]
其中$C=C_1-\ln a$.

又函数$x=a\sec t$在$\left(\frac{\pi}{2},\pi\right]$上的反函数是
\[t=\arcsec\frac{x}{a}\quad \left(-\infty<\frac{x}{a}<-1,\; \text{即}x<-a\right)\]
于是，当$x<-a$时，有
\[\sqrt{x^2-a^2}=a\sqrt{\tan^2 t}=a|\tan t|=-a\tan t\quad \left(\frac{\pi}{2}<t\le \pi\right)\]
\[\dd x=\dd(a\sec t)=a\frac{\sin t}{\cos^2 t}\cdot \dd t\]
\[\begin{split}
    \Int \frac{1}{\sqrt{x^2-a^2}}\dd x&=\int-\frac{1}{a\tan t}\cdot \frac{a\sin t}{\cos^2 t}\dd t\\
    &=-\int \frac{1}{\cos t}\dd t=-\ln|\sec t+\tan t|+C_2\\
    &=-\ln\left|\frac{x}{a}-\frac{\sqrt{x^2-a^2}}{a}\right|+C_2\\
    &=-\ln\left|x-\sqrt{x^2-a^2}\right|+C_2+\ln a\\
    &=\ln\left|\frac{1}{x-\sqrt{x^2-a^2}}\right|+C'\qquad (C'=C+\ln a)\\
    &=\ln\left|\frac{x+\sqrt{x^2-a^2}}{a^2}\right|+C'=\ln\left|x+\sqrt{x^2-a^2}\right|+C'-\ln a^2\\
    &=\ln\left|x+\sqrt{x^2-a^2}\right|+C''\qquad (C''=C'-\ln a^2)
\end{split}\]
因此，不论$x>a$或$x<-a$，都有
\[\int\frac{\dd x}{\sqrt{x^2-a^2}}=\ln\left|x+\sqrt{x^2-a^2}\right|+C\]
\end{solution}

到现在为止，我们又得到四个重要的不定积分公式：

\begin{blk}{}
\begin{align}
\Int \frac{\dd x}{a^2+x^2}&=\frac{1}{a}\arctan\frac{x}{a}+C\\
\Int\frac{\dd x}{a^2-x^2}&=\frac{1}{2a}\ln\left|\frac{a+x}{a-x}\right|+C\\
 \Int\frac{\dd x}{\sqrt{a^2-x^2}}&=\arcsin\frac{x}{a}+C\\
\Int \frac{\dd x}{\sqrt{x^2\pm a^2}}&=\ln\left|x+\sqrt{x^2\pm a^2}\right|+C
\end{align}
\end{blk}

利用这些公式可以进一步求出许多不定积分．

\begin{blk}
 {类型A}
\[\int \frac{\dd x}{ax^2+bx+c}\]
\end{blk}

如果被积函数的分母可以分解因式，则把被积函数分解成部分分式，就很容易分项求出不定积分．

如分母不能分解因式，那么将它配方得
\[ax^2+bx+c=a\left(x+\frac{b}{2a}\right)^2+\frac{4ac-b^2}{4a}\]
于是
\begin{itemize}
    \item 当$4ac>b^2$时，上面不定积分归结到(4.10);
    \item 当$4ac<b^2$时，上面不定积分归结到(4.11).
\end{itemize}

\begin{example}
    求$\Int \frac{\dd x}{3x^2+x-1}$.
\end{example}

\begin{solution}
\[\begin{split}
    \Int \frac{\dd x}{3x^2+x-1}&=\frac{1}{3}\int \frac{\dd x}{\left(x+\frac{1}{6}\right)^2-\frac{13}{36}}\\
&=\frac{1}{3}\frac{1}{2\sqrt{\frac{13}{36}}}\ln\left|\frac{x+\frac{1}{6}-\frac{\sqrt{13}}{6}}{x+\frac{1}{6}+\frac{\sqrt{13}}{6}}\right|+C\\
&=\frac{1}{\sqrt{13}}\ln\left|\frac{6x+1-\sqrt{13}}{6x+1+\sqrt{13}}\right|+C
\end{split}\]    
\end{solution}

\begin{blk}
    {类型B}
   \[\int \frac{px+q}{ax^2+bx+c}\dd x\]
   \end{blk}

注意到$\frac{\dd}{\dd x}(ax^2+bx+c)=2ax+b$，被积函数的分子可以写成
\[px+q=\frac{p}{2a}(2ax+b)+q-\frac{pb}{2a}\]
于是
\[\int\frac{px+q}{ax^2+bx+c}\dd x=\frac{p}{2a}\int\frac{2ax+b}{ax^2+bx+c}\dd x+\left(q-\frac{bp}{2a}\right)\int\frac{\dd x}{ax^2+bx+c}\]
等式右端第一项是$\frac{p}{2a}\ln(ax^2+bx+c)$，第二项可归结到类型A.

\begin{example}
    求$\Int \frac{5x-1}{3x^2+x+2}\dd x$.
\end{example}

\begin{solution}
\[\begin{split}
    \Int \frac{5x-1}{3x^2+x+2}\dd x&=\frac{5}{6}\int \frac{6x+1}{3x^2+x+2}\dd x-\frac{11}{6}\int\frac{\dd x}{3x^2+x+2}\\
&=\frac{5}{6}\ln(3x^2+x+2)-\frac{11}{6}\cdot \frac{1}{3}\cdot \int \frac{\dd x}{\left(x+\frac{1}{6}\right)^2+\frac{23}{36}}\\
&=\frac{5}{6}\ln(3x^2+x+2)-\frac{11}{6}\cdot\frac{1}{3}\cdot\frac{6}{\sqrt{23}}\arctan\frac{6x+1}{\sqrt{23}}+C\\
&=\frac{5}{6}\ln(3x^2+x+2)-\frac{11}{3\sqrt{23}}\arctan\frac{6x+1}{\sqrt{23}}+C
\end{split}\]    
\end{solution}

\begin{blk}{类型C}
\[\int \frac{\dd x}{\sqrt{ax^2+bx+c}}\]
\end{blk}

\begin{itemize}
    \item 若$a>0$, 则被积函数的根底式经配方后，所求积分可归结到积分公式(4.13);
    \item  若$a<0$, 则此积分可归结到积分公式(4.12).
\end{itemize}

\begin{example}
    求$\Int \frac{\dd x}{\sqrt{2+x-3x^2}}$
\end{example}

\begin{solution}
    \[\begin{split}
        \Int \frac{\dd x}{\sqrt{2+x-3x^2}}&=\frac{1}{\sqrt{3}}\int \frac{\dd x}{\sqrt{\frac{2}{3}+\frac{1}{3}x-x^2}}\\
        &=\frac{1}{\sqrt{3}}\int \frac{\dd x}{\sqrt{\frac{25}{36}-\left(x-\frac{1}{6}\right)^2}}\\
        &=\frac{1}{\sqrt{3}}\arcsin\frac{x-\frac{1}{6}}{\frac{5}{6}}+C\\
        &=\frac{1}{\sqrt{3}}\arcsin\frac{6x-1}{5}+C       
    \end{split}\]
\end{solution}

\begin{blk}{类型D}
\[\int \frac{px+q}{\sqrt{ax^2+bx+c}}\dd x\]
\end{blk}

被积函数的分子可写成
\[px+q=\frac{p}{2a}(2ax+b)+q-\frac{pb}{2a}\]
所以
\[
    \int \frac{px+q}{\sqrt{ax^2+bx+c}}\dd x=\frac{p}{2a}\int \frac{2ax+b}{\sqrt{ax^2+bx+c}}\dd x+\left(q-\frac{pb}{2a}\right)\int\frac{\dd x}{\sqrt{ax^2+bx+c}}
\]
等式右端第一个积分等于被积函数的分母的2倍，即$2\sqrt{ax^2+bx+c}$, 第二个是类型C的积分．

\begin{example}
    求$\Int \frac{5x-1}{\sqrt{3x^2+x+2}}\dd x$.
\end{example}

\begin{solution}
    \[\begin{split}
        \Int \frac{5x-1}{\sqrt{3x^2+x+2}}\dd x&=\frac{5}{6}\int \frac{6x+1}{\sqrt{3x^2+x+2}}\dd x-\frac{11}{6}\int \frac{\dd x}{\sqrt{3x^2+x+2}}\\
&=\frac{5}{6}\cdot 2\sqrt{3x^2+x+2}\\
&\qquad -\frac{11}{6}\cdot\frac{1}{\sqrt{3}}\ln\left[x+\frac{1}{6}+\sqrt{\left(x+\frac{1}{6}\right)^2+\frac{23}{36}}\right]+C\\
&=\frac{5}{3}\sqrt{3x^2+x+2}-\frac{11}{6\sqrt{3}}\ln\left[x+\frac{1}{6}+\sqrt{\left(x+\frac{1}{6}\right)^2+\frac{23}{36}}\right]+C        
    \end{split}\]
\end{solution}

\begin{blk}{类型E}
    \[\int \frac{\dd x}{x\sqrt{ax^2+bx+c}}\]
\end{blk}

作变量代换$x=\frac{1}{y}$，所求积分可化为类型C．因为$\dd x=-\frac{1}{y^2}\dd y$，并且
\[\int \frac{\dd x}{x\sqrt{ax^2+bx+c}}=\int\frac{-\frac{1}{y^2}\dd y}{\frac{1}{y}\sqrt{\frac{a}{y^2}+\frac{b}{y}+c}}=-\int\frac{\dd y}{\sqrt{a+by+cy^2}}\]

\begin{example}
    求$\Int\frac{\dd x}{x\sqrt{x^2+x+1}}$
\end{example}

\begin{solution}
\[\begin{split}
    \Int\frac{\dd x}{x\sqrt{x^2+x+1}}&=-\int \frac{-\frac{1}{y^2}\dd y}{\frac{1}{y}\sqrt{\frac{1}{y^2}+\frac{1}{y}+1}}=-\int \frac{\dd y}{\sqrt{y^2+y+1}}\\
    &=-\int \frac{\dd y }{\sqrt{\left(y+\frac{1}{2}\right)^2+\frac{3}{4}}}  \\
    &=-\ln \left[y+\frac{1}{2}+\sqrt{\left(y+\frac{1}{2}\right)^2+\frac{3}{4}}\right]+C \\
    &=-\ln \left[y+\frac{1}{2}+\sqrt{y^2+y+1}\right]+C\\
    &=-\ln\left[\frac{1}{x}+\frac{1}{2}+\sqrt{\frac{1}{x^2}+\frac{1}{x}+1}\right]+C\\
    &=-\ln\frac{2+x+2\sqrt{x^2+x+1}}{2x}+C
\end{split}\]
\end{solution}
    
\begin{ex}
\begin{enumerate}
    \item 求下列不定积分：
\begin{multicols}{2}
\begin{enumerate}
    \item $\Int \frac{1}{1+\sqrt[3]{x}}\dd x$
    \item $\Int \frac{1}{(x+2) \sqrt{x+1}}\dd x$
    \item $\Int  \frac{x}{\sqrt[3]{1-x}}\dd x $
    \item $\Int \frac{\sqrt{x}}{\sqrt[3]{x^{2}}-\sqrt[4] {x}}\dd x$
    \item $\Int  \frac{\sqrt{x^{2}-a^{2}}}{x}\dd x $
    \item $\Int  \frac{\sqrt{x^{2}-a^{2}}}{x^{2}}\dd x$
    \item $\Int  \frac{\sqrt{a^{2}-x^{2}}}{x}\dd x$
    \item $\Int \frac{\sqrt{a^{2}-x^{2}}}{x^{2}}\dd x$
    \item $\Int \frac{x^{2}}{\sqrt{1-x^{2}}}\dd x$
    \item $\Int \frac{1}{x^{4} \sqrt{1+x^{2}}}\dd x$
\end{enumerate}
\end{multicols}

    \item 求下列不定积分：
\begin{multicols}{2}
\begin{enumerate}
\item $\Int \frac{1}{\sqrt{4 x^{2}+6 x}}\dd x $
\item $\Int \frac{x}{\sqrt{4 x^{2}+6 x}}\dd x$
\item $\Int  \frac{1}{\sqrt{1-x-4 x^{2}}}\dd x$
\item $\Int \frac{x+1}{\sqrt{1-x-4 x^{2}}}\dd x$
\item $\Int \frac{1}{x \sqrt{x^{2}+1}}\dd x$
\item $\Int \frac{1}{x \sqrt{1-x^{2}}}\dd x$
\item $\Int \frac{1}{x \sqrt{x^{2}-4 x}}\dd x$
\item $\Int \sqrt{\frac{x-1}{x+1}}\dd x$  
\end{enumerate}
\end{multicols}

\end{enumerate}    
\end{ex}

\subsection{有理函数的积分}
如果$P(x)$, $Q(x)$是两个多项式，则$\Int \frac{P(x)}{Q(x)}\dd x$可以化
归为下列类型的积分：
\begin{center}
\begin{tabular}{p{.3\textwidth}p{.3\textwidth}}
        $\Int \frac{A \dd x}{x-a}$ & $\Int \frac{A \dd x}{(x-a)^{n}}$ \\
        $\Int \frac{p x+q}{x^{2}+b x+c} \dd x$ & $\Int \frac{p x+q}{\left(x^{2}+b x+c\right)^{n}} \dd x$
\end{tabular}
\end{center}

对这些积分，我们是有办法的．下面我们通过举例来说明求这些积分的大概的方法．

如果有理分式的分子次数高于分母的次数时，则只要作一次除法就可以把它写成一个多项式与一个真分式的和，而前者是容易求积的．所以我们只须注意真分式的积分如何求．

\begin{example}
    求$\Int \frac{x^3+x+1}{x^2+1}\dd x$
\end{example}

\begin{solution}
    \[\begin{split}
        \Int \frac{x^3+x+1}{x^2+1}\dd x&=\int\left(x+\frac{1}{x^2+1}\right)\dd x\\
&=\int x\dd x+\int\frac{\dd x}{1+x^2}\\
&=\frac{x^2}{2}+\arctan x+C        
    \end{split}\]
\end{solution}

\begin{example}
    求$\Int \frac{\dd x}{x^4-1}$
\end{example}

\begin{solution}
$\because\quad x^4-1=(x-1)(x+1)(x^2+1)$

$\therefore\quad \frac{1}{x^4-1}=\frac{A}{x-1}+\frac{B}{x+1}+\frac{Cx+D}{x^2+1}$

把等式右端通分，再比较两端系数，于是等式
\[A (x^2+1) (x+1)+B(x^2+1)(x-1)+(Cx+D) (x^2-1)=1\]
的左边多项式的$x^3,x^2,x$的系数应为零，常数项应为1, 从而得到
\begin{numcases}{}
    A+B+C=0\\
    A-B+D=0\\
    A+B-C=0\\
    A-B-D=1
\end{numcases}
由$(4.14)+(4.16)$和$(4.15)+(4.17)$得到：
\[A+B=0,\qquad A-B=\frac{1}{2}\]
解得：
\[A=\frac{1}{4},\quad B=-\frac{1}{4},\quad C=0,\quad D=-\frac{1}{2}\]
即：
\[\frac{1}{x^4-1}=\frac{1}{4(x-1)}-\frac{1}{4(x+1)}-\frac{1}{2(x^2+1)}\]
从而
\[\begin{split}
    \int\frac{\dd x}{x^4-1}&=\frac{1}{4}\int \frac{\dd x}{x-1}-\frac{1}{4}\int\frac{\dd x}{x+1}-\frac{1}{2}\int\frac{\dd x}{1+x^2}\\
    &=\frac{1}{4}\ln\left|\frac{x-1}{x+1}\right|-\frac{1}{2}\arctan x+C
\end{split}\]
\end{solution}

\begin{example}
求$\Int \frac{3x^2+x-2}{(x-2)^2(1-2x)}\dd x$
\end{example}

\begin{solution}
设$\frac{3x^2+x-2}{(x-2)^2(1-2x)}=\frac{A}{1-2x}+\frac{B}{x-2}+\frac{C}{(x-2)^2}$，则
\[3x^2+x-2=A(x-2)^2+B(1-2x)(x-2)+C(1-2x)\]

令$1-2x=0$，则$A=-\frac{1}{3}$；令$x-2=0$，则$C=-4$. 比较等式两端$x^2$的系数，得
\[3=A-2B\quad\Rightarrow\quad B=-\frac{5}{3}\]
所以
\[\frac{3x^2+x-2}{(x-2)^2(1-2x)}=-\frac{1}{3(1-2x)}-\frac{5}{3(x-2)}-\frac{4}{(x-2)^2}\]
从而
\[\begin{split}
\int \frac{3x^2+x-2}{(x-2)^2(1-2x)}\dd x &=\frac{1}{6}\int \frac{\dd x}{x-\frac{1}{2}}-\frac{5}{3}\int \frac{\dd x}{x-2}-4\int \frac{\dd x}{(x-2)^2}\\
&=\frac{1}{6}\ln\left|x-\frac{1}{2}\right|-\frac{5}{3}\ln|x-2|+\frac{4}{x-2}+C
\end{split}\]
\end{solution}

三角函数有理式的积分
\[\int R (\sin x,\cos x) \dd x\]
恒可用“半角变换”法，即设$t=\tan\frac{x}{2}$或$x=2\arctan t$化为$t$的有理函数的积分．

事实上，由$t=\tan\frac{x}{2}$, 知
\[\begin{split}
\sin x&=\frac{2 \sin \frac{x}{2} \cos \frac{x}{2}}{\sin ^{2} \frac{x}{2}+\cos ^{2} \frac{x}{2}}=\frac{2 \tan \frac{x}{2}}{1+\tan^{2} \frac{x}{2}}=\frac{2 t}{1+t^{2}}\\
\cos x &=\frac{\cos ^{2}-\frac{x}{2}-\sin ^{2} \frac{x}{2}}{\sin ^{2} \frac{x}{2}+\cos ^{2} \frac{x}{2}}=\frac{1-\tan^{2} \frac{x}{2}}{1+\tan^{2} \frac{x}{2}} 
=\frac{1-t^{2}}{1+t^{2}}
\end{split}\]
又由 $x=2 \arctan t$, 知
\[\dd x=\frac{2 \dd t}{1+t^{2}}\]
于是
\[
\int R(\sin x,  \cos x) \dd x=\int R\left(\frac{2 t}{1+t^{2}}, \frac{1-t^{2}}{1+t^{2}}\right) \cdot \frac{2}{1+t^{2}} \dd x 
\]
上式右端是 $t$ 的有理函数.

\begin{example}
    求$\Int \frac{\dd\theta}{\sin\theta}$.
\end{example}

\begin{solution}
令$t=\tan\frac{\theta}{2}$，则：$\theta=2\arctan t,\quad \dd\theta=\frac{2\dd t}{1+t^2}$，从而：
\[\begin{split}
    \int\frac{\dd\theta}{\sin\theta}=\int\frac{1+t^2}{2t}\cdot \frac{2\dd t}{1+t^2}&=\int\frac{\dd t}{t}\\
    &=\ln|t|+C=\ln\left|\tan\frac{\theta}{2}\right|+C
\end{split}\]
又$\tan\frac{\theta}{2}=\frac{\sin\frac{\theta}{2}}{\cos\frac{\theta}{2}}=\frac{2\sin^2 \frac{\theta}{2}}{2\sin\frac{\theta}{2}\cos \frac{\theta}{2}}=\frac{1-\cos\theta}{\sin\theta}=\csc\theta-\cot\theta$，
所以
\[\int\frac{\dd\theta}{\sin\theta}=\ln\left|\tan\frac{\theta}{2}\right|+C=\ln|\csc\theta-\cot\theta|+C\]    
\end{solution}

\begin{example}
    求$\Int \frac{\dd\theta}{3+2\cos\theta}$
\end{example}

\begin{solution}
    令$t=\tan\frac{\theta}{2}$，则：$\dd\theta=\frac{2\dd t}{1+t^2}$，于是
\[\begin{split}
    \Int \frac{\dd\theta}{3+2\cos\theta}=\int\frac{\frac{2\dd t}{1+t^2}}{3+2\frac{1-t^2}{1+t^2}}&=\int \frac{2\dd t}{5+t^2}\\
    &=\frac{2}{\sqrt{5}} \arctan\frac{t}{\sqrt{5}}+C\\
    &=\frac{2}{\sqrt{5}}\arctan\left(\frac{1}{\sqrt{5}}\tan\frac{\theta}{2}\right)+C
\end{split}\]
\end{solution}

\begin{ex}
\begin{enumerate}
\item 求下列不定积分:
\begin{multicols}{2}
\begin{enumerate}
\item  $\Int \frac{1}{x^{2}-3 x+2}\dd x$
\item  $\Int  \frac{x}{x^{2}-3 x+2}\dd x$
\item  $\Int  \frac{x^{2}}{x^{2}-3 x+2}\dd x$
\item  $\Int \frac{x^{2}+1}{x+3}\dd x$
\item  $\Int \frac{1}{x^{2}-3 x+1}\dd x$
\item  $\Int \frac{x}{x^{2}-3 x+1}\dd x$
\item  $\Int \frac{x^{2}}{x^{2}-3 x+1}\dd x$
\end{enumerate}
\end{multicols}

\item 求下列不定积分:
\begin{multicols}{2}
    \begin{enumerate}
        \item  $\Int  \frac{\cos \theta \dd \theta}{1+\sin ^{2} \theta}$
        \item  $\Int \frac{\sin x+\cos x}{\cos ^{2} x}\dd x$
        \item  $\Int  \frac{\dd x}{4 \cos x+3 \sin x+5}$
        \item  $\Int \frac{\sin x}{\cos x+\sin x}\dd x$
    \end{enumerate}
\end{multicols}
\end{enumerate}
\end{ex}

\subsection{分部积分法}

分部积分法是从两个函数乘积的微分公式得到的求积分方法．

设$u=u(x)$, $v=v(x)$是两个可微函数，乘积$u\cdot v$的微分法则是$$\dd (u\cdot v) =u\dd v+v\dd u$$
移项，得$u\dd v=\dd (w\cdot v) -v\dd u$，
两边积分得
\begin{equation}
    \int u\dd v=u\cdot v-\int v\dd u
\end{equation}
公式(4.18)表示如果积分$\Int u\dd v=\Int u(x)v'(x)\dd x$不可能直接求出，就改为求积分$\Int v\dd u=\Int v(x)u'(x)\dd x$, 设$F(x)$是$v(x)\cdot u'(x)$的一个原函数，那么$u(x)\cdot v(x)-F(x)$是$u(x)v'(x)$的一个原函数，用公式(4.18)求不定积分的方法叫做\textbf{分部积分法}，公式(4.18)叫做\textbf{分部积分公式}．

\begin{example}
    求$\Int x\sin 3x\dd x$.
\end{example}

\begin{solution}
被积函数$x\sin3x$不能由前面的基本积分表求出它的原函数，我们把$x\sin3x$分成两部分，设$u(x)=x$, $v'(x)=\sin3x$, 于是
\[\int v' (x) \dd x= \int \sin3x\dd x=-\frac{\cos3x}{3}+C\]
即$v(x)=-\frac{\cos3x}{3}$是$v'(x)=\sin3x$的一个原函数，根据公
式(4.18)，有
\[\begin{split}
    \int x\sin 3x\dd x&=\int x\dd\left(-\frac{\cos 3x}{3}\right)\\
    &=x\left(-\frac{\cos 3x}{3}\right)-\int\left(-\frac{\cos3x}{3}\right)\dd x \\
&=-\frac{x\cos 3x}{3}+\frac{1}{3}\int \cos 3x\dd x\\
&=-\frac{x\cos 3x}{3}+\frac{\sin 3x}{9}+C
\end{split}\]
\end{solution}

运用公式(4.18)时，应首先将被积表达式分成函数$u(x)$和微分$\dd v=v'(x)\dd x$两部分，选取$u(x)$和$v'(x)\dd x$的原则是：
\begin{itemize}
    \item 充当微分部分的$v'(x)\dd x$, 其$v'(x)$的原函数能够比较容易地求出；
\item 
充当$u(x)$的部分经求微商后要化简，以便使积分$\Int v\dd u$较$\Int u\dd v$简单，容易求出结果．
\end{itemize}

在应用分部积分法求不定积分时，往往要连续使用这个
方法多次才能求出最后结果．

我们把能够应用分部积分法求积分的例题分为以下几类：

\begin{blk}
 {类型A} 被积函数是$x^ne^{ax}$, $x^n\sin ax$, $x^n\cos ax$, 其中$x$的方幂因子经求微商后降次，而非$x$的方幂因子的原函数是容易求得的．   
\end{blk}

\begin{example}
求$\Int x^2 e^{2x}\dd x$.
\end{example}

\begin{solution}
\[\begin{split}
    \Int x^2 e^{2x}\dd x=\int x^2 \dd\left(\frac{e^{2x}}{2}\right)&=\frac{x^2e^{2x}}{2}-\int \frac{e^{2x}}{2}\dd(x^2)\\
&=\frac{x^2e^{2x}}{2}-\int xe^{2x}\dd x=\frac{x^2e^{2x}}{2}-\left(\frac{xe^{2x}}{2}-\int \frac{e^{2x}}{2}\dd x\right)\\
&=\frac{x^2}{2}e^{2x}-\frac{x}{2}e^{2x}+\frac{e^{2x}}{4}+C\\    
&=e^{2x}\left(\frac{x^2}{2}-\frac{x}{2}+\frac{1}{4}\right)+C
\end{split}\]
\end{solution}

\begin{blk}
    {类型B} 被积函数是$x^n\sin^{-1}x$, $x^n\tan^{-1}x$, $x^n\ln x$, 其中乘积中的因子$\tan^{-1} x$, $\ln x$的导数是$x$的有理函数，$\sin^{-1}x$的
导数是$\frac{1}{\sqrt{1-x^2}}$，
它们的积分可求，而$x^n$的原函数是$\frac{x^{n+1}}{n+1}$.
\end{blk}


\begin{example}
求$\Int x\arctan x\dd x$.    
\end{example}

\begin{solution}
  \[\begin{split}
    \Int x\arctan x\dd x=\int \arctan x \dd\left(\frac{x^2}{2}\right)&=\frac{x^2\arctan x}{2}-\int \frac{x^2}{2}\dd(\arctan x)\\
    &=\frac{x^2\arctan x}{2}-\frac{1}{2}\int \frac{x^2}{1+x^2}\dd x\\
    &=\frac{x^2\arctan x}{2}-\frac{1}{2}\int \left(1-\frac{1}{1+x^2}\right)\dd x\\
    &=\frac{x^2\arctan x}{2}-\frac{1}{2}x+\frac{1}{2}\arctan x+C
  \end{split}\]  
\end{solution}


\begin{example}
    求$\Int x^3\ln x\dd x$.
\end{example}

\begin{solution}
\[\begin{split}
    \Int x^3\ln x\dd x=\int \ln x\dd\left(\frac{x^4}{4}\right)&=\frac{x^4}{4}\ln x-\int \frac{x^4}{4}\dd(\ln x)\\
    &= \frac{x^4}{4}\ln x-\int  \frac{x^4}{4}\cdot \frac{1}{x}\dd x\\
    &= \frac{x^4}{4}\ln x- \frac{x^4}{16}+C
\end{split}\]
\end{solution}

\begin{blk}
    {类型C} 被积函数是$\sin^{-1}x$, $\tan^{-1}x$, $\ln x$, 这里没有乘积的形式，不妨设$v'(x)=1$.
\end{blk}

\begin{example}
求$\Int \arcsin x\dd x$.
\end{example}

\begin{solution}
\[\begin{split}
    \Int \arcsin x\dd x=x\arcsin x-\int x\dd(\arcsin x)&=x\arcsin x-\int\frac{x}{\sqrt{1-x^2}}\dd x  \\
    &=x\arcsin x+\frac{1}{2}\int \frac{\dd(1-x^2)}{\sqrt{1-x^2}}\\
    &=x\arcsin x+\frac{1}{2}\cdot 2\sqrt{1-x^2}+C\\
    &=x\arcsin x+\sqrt{1-x^2}+C
\end{split}\]

另解：设$y=\arcsin x,\; -1\le x\le 1$，则$\sin y=x,\; -\frac{\pi}{2}\le y\le \frac{\pi}{2}$
\[\begin{split}
    \Int \arcsin x\dd x=x\arcsin x-\int x\dd(\arcsin x)&=y\sin y-\int \sin y \dd y\\
    &=y\sin y+\cos y+C\\
    &=x\arcsin x+\sqrt{1-\sin^2 y}+C\\
    &=x\arcsin x+\sqrt{1-x^2}+C
\end{split}\]
\end{solution}

\begin{example}
    求$\Int \ln^2 x\dd x$.
\end{example}

\begin{solution}
    \[\begin{split}
        \Int \ln^2 x\dd x=x(\ln x)^2-\int x\dd (\ln^2 x)&=x(\ln x)^2-\int x\cdot 2\ln x\cdot \frac{1}{x}\dd x\\
        &=x(\ln x)^2-2\int \ln x\dd x\\
        &=x(\ln x)^2-2\left[x\ln x-\int x\d(\ln x)\right]\\
        &=x(\ln x)^2-2x\ln x+2\int \dd x\\
        &=x(\ln x)^2-2x\ln x+2x+C   
    \end{split}\]
\end{solution}

\begin{blk}
    {类型D} 被积函数是$e^{ax}\sin bx$, $\sec^3 x$, $\sqrt{a^2+x^2}$. 对于这类例题，两次使用这个方法可以得到含有这个积分作未知元的一次方程．
\end{blk}

\begin{example}
    求$\Int e^{ax}\sin bx\dd x$
\end{example}

\begin{solution}
\[\begin{split}
    I=\Int e^{ax}\sin bx\dd x
    &=\int \sin bx \dd\left(\frac{e^{ax}}{a}\right)\\
&=\sin bx\cdot \frac{e^{ax}}{a}-\int \frac{e^{ax}}{a}\dd(\sin bx)\\
&=\sin bx\cdot \frac{e^{ax}}{a}-\int \frac{e^{ax}}{a}b\cos bx\dd x\\
&=\sin bx\cdot \frac{e^{ax}}{a}-\frac{b}{a}\int \cos bx\dd\left(\frac{e^{ax}}{a}\right)\\
&=\sin bx\cdot \frac{e^{ax}}{a}-\frac{b}{a}\left[\cos bx\cdot \frac{e^{ax}}{a}-\int \frac{e^{ax}}{a}\dd\left(\cos bx\right)\right]\\
&=\sin bx\cdot \frac{e^{ax}}{a}-\frac{b}{a^2}\cos bx\cdot e^{ax}-\frac{b^2}{a^2}\int e^{ax}\sin bx\dd x\\
&=\sin bx\cdot \frac{e^{ax}}{a}-\frac{b}{a^2}\cos bx\cdot e^{ax}-\frac{b^2}{a^2}I
\end{split}\]
所以：
\[\begin{split}
    I\left(1+\frac{b^2}{a^2}\right)&=\frac{1}{a}\sin bx\cdot e^{ax}-\frac{b}{a^2}\cos bx\cdot e^{ax}\\
    I&=\frac{e^{ax}(a\sin bx-b\cos bx)}{a^2+b^2}+C
\end{split}\]    
\end{solution}

\begin{example}
    求$\Int \sqrt{a^2+x^2}\dd x$.
\end{example}

\begin{solution}
\[    \begin{split}
    I =\int \sqrt{a^{2}+x^{2}} \dd x 
    &=x \sqrt{a^{2}+x^{2}}-\int x \dd\left(\sqrt{a^{2}+x^{2}}\right) \\
    &=x \sqrt{a^{2}+x^{2}}-\int \frac{x^{2}}{\sqrt{a^{2}+x^{2}}} \dd x \\
    &=x \sqrt{a^{2}+x^{2}}-\int \frac{x^{2}+a^{2}}{\sqrt{a^{2}+x^{2}}} \dd x+\int \frac{a^{2}}{\sqrt{a^{2}+x^{2}}} \dd x \\
    &=x \sqrt{a^{2}+x^{2}}-I+a^{2} \ln \left(x+\sqrt{a^{2}+x^{2}}\right)
    \end{split}\]
    所以
\[    I=\frac{1}{2} x \sqrt{a^{2}+x^{2}}+\frac{1}{2} a^{2} \ln \left(x+\sqrt{a^{2}+x^{2}}\right)\]
\end{solution}

应用分部积分公式(4.18)还能导出一些求积分的递归公式．所谓求积分的递归公式就是建立一个被积函数含有正整数$n$的积分与被积函数含有非负整数$n-1$或$n-2$的积分的关系式．每次要求$n=k$时的积分，重复利用这个关系式总能回到当$n=0, 1, 2$时的起始积分，于是所求积分总能用$n$的起始值$n=0, 1, 2$中的一个已知积分表示出来．这也就是递归的
意思．
    
\begin{example}
    求$\Int \sin^n x\dd x$.
\end{example}

\begin{solution}
当$n=0$时，$I_0=\Int (\sin x)^0\dd x=\Int \dd x=x+C$

又：
\[
\begin{split}
I_{n} =\int \sin ^{n} x \dd x&=\int \sin ^{n-1} x \dd(-\cos x) \\
&=-\cos x \sin ^{n-1} x+\int \cos x \dd\left(\sin ^{n-1} x\right) \\
&=-\cos x \sin ^{n-1} x+(n-1) \int \cos x \sin ^{n-2} x \cdot \cos x \dd x \\
&=-\cos x \sin ^{n-1} x+(n-1) \int \sin ^{n-2} x\left(1-\sin ^{2} x\right) \dd x
\end{split}
\]
所以
$$
I_{n}=-\cos x \sin ^{n-1} x+(n-1) I_{n-2}-(n-1) I_{n}
$$
即
\begin{equation}
   n I_{n}=-\cos x \sin ^{n-1} x+(n-1) I_{n-2} 
\end{equation}

现在我们用递归关系(4.19)来求 $\Int \sin ^{4} x \dd x$. 
设 $I_{4}=\Int \sin ^{4} x \dd x$, 于是
\begin{align}
4 I_{4}&=-\cos x \sin ^{3} x+3 I_{2} \\
2 I_{2}&=-\cos x \sin x+I_{0} 
\end{align}
将(4.21)代入(4.20)得：
\[\begin{split}
    4I_4&=-\cos x\sin^3 x+3\left(-\frac{\sin x\cos x}{2}+\frac{1}{2}I_0\right)\\
    &=-\cos x\sin^3 x-\frac{3}{2}\cos x\sin x+\frac{3}{2}x+C
\end{split}\]
所以
\[I_4=\int \sin^4 x\dd x=-\frac{1}{4}\cos x\sin^3 x-\frac{3}{8}\cos x\sin x+\frac{3}{8}x+C\]


\begin{ex}
\begin{enumerate}
    \item 求下列不定积分：
\begin{multicols}{2}
\begin{enumerate}
    \item $\Int x \sin x\dd x$
    \item $\Int x \cos 3 x\dd x$
    \item $\Int x \sin ^{2} \frac{x}{2}\dd x$
    \item $\Int x e^{-x}\dd x$
    \item $\Int (x-1) e^{x}\dd x$
    \item $\Int \left(x^{2}+x\right) e^{-x}\dd x$
    \item $\Int  x \sec ^{2} x\dd x$
    \item $\Int  x \arcsin x\dd x$
    \item $\Int  \frac{\arcsin x}{x^{2}}\dd x$
    \item $\Int  \frac{\arctan x}{x^{2}}\dd x$
    \item $\Int  x^{2} e^{3 x}\dd x$
    \item $\Int  e^{a x} \sin ^{2} b x\dd x$
    \item $\Int  \sin x \cdot \ln (\tan x)\dd x$
    \item $\Int  x^{2} \arctan x\dd x$
    \item $\Int  \frac{x^{2} e^{x}\dd x}{(x+2)^{2}}$
    \item $\Int \frac{x e^{x}}{(x+1)^{2}}\dd x$
\end{enumerate}
\end{multicols}
 \item   求下列不定积分:    
   \begin{multicols}{2}
        \begin{enumerate}
    \item $\Int \arctan \sqrt{x}\dd x$
    \item $\Int \ln ^{2} x \dd {x}$
    \item $\Int \ln \left(x+\sqrt{1+x^{2}}\right)\dd x$
    \item $\Int (\arcsin x)^{2}\dd x$
    \item $\Int \sin(\ln x)\dd x$
    \item $\Int \sqrt{4+x^2}\dd x$
\end{enumerate}
\end{multicols}
\end{enumerate}
\end{ex}
\end{solution}
    
\section{定积分的计算与应用}

\subsection{定积分的计算}

由微积分的基本定理得知，定积分$\Int^b_a f(x)\dd x$的计算归结为不定积分$\Int f(x)\dd x$的计算，由后者计算得到$f(x)$的原函数$F(x)$, 则
\[\Int^b_a f(x)\dd x=F (x)\Bigg|^b_a=F (b) -F (a)\] 
于是计算不定积分的各种方法也可移植过来求定积分．

\subsubsection{定积分的换元积分法}

第一换元法：设$\Int f(u)\dd u=F(u)$, $u=\varphi(x)$可微，
那么
\[\begin{split}
    \int^b_a f(\varphi(x))\varphi'(x)\dd x&=\int^{\varphi(b)}_{\varphi(a)}f(u)\dd u\\
    &=F(u)\Big|^{\varphi(b)}_{\varphi(a)}=F[{\varphi(b)}]-F[{\varphi(a)}]
\end{split}\]

\begin{example}
    求$\Int^{\tfrac{\pi}{2}}_0 \sin^3x \cos x\dd x$.
\end{example}

\begin{solution}
    因为
\[\int \sin^3x\cos x\dd x= \int \sin^3x\dd (\sin x) =\frac{1}{4}\sin^4x+C\]
所以
\[\begin{split}
    \Int^{\tfrac{\pi}{2}}_0 \sin^3x \cos x\dd x&=\frac{1}{4}\sin^4 x\Big|^{\tfrac{\pi}{2}}_0 \\
    &=\frac{1}{4}\left[\sin^4\frac{\pi}{4}-0\right]=\frac{1}{4}
\end{split}\]
\end{solution}

\begin{blk}
 {定理1（第二种变量代换法）}设$f(x)$在$[a,b]$上连续，作代换$x=\varphi (t)$, 其中$\varphi (t)$在某一闭区间$[\alpha,\beta]$上有连续导数$\varphi '(t)$, 当$\alpha\le t\le \beta$时，$a\le \varphi (t)\le b$, 且$\varphi (\alpha)=a$, $\varphi (\beta)=b$, 那么
\[\int^b_a f(x) \dd x= \int^{\beta}_{\alpha} f [\varphi  (t) ] \varphi ' (t) \dd t\]   
\end{blk}

\begin{proof}
    设$f(x)$的一个原函数是$G(x)$, 即有
\[\frac{\dd}{\dd x}G (x) =f (x)\]
于是由微积分基本定理得
\begin{equation}
    \int^b_a f (x)\dd x=G(x)\Big|^b_a=G(b) -G (a)
\end{equation}
根据复合函数求导法则，知
\[\begin{split}
    \frac{\dd}{ \dd t}G[\varphi  (t) ] &=\frac{\dd}{ \dd x}G (x)\cdot  \frac{\dd}{ \dd t}\varphi  (t)\\
    &=f(x)\cdot \varphi'(t)=f[\varphi(t)]\varphi'(t)
\end{split}\]
所以$G[\varphi (t)]$是$f[\varphi (t)]\varphi '(t)$的一个原函数，从而
\begin{equation}
    \begin{split}
\int^{\beta}_{\alpha} f [\varphi  (t) ] \varphi ' (t) \dd t&=G [\varphi  (t) ]\Big|^{\beta}_{\alpha}\\
&=       G [\varphi  (\beta) ] -G [\varphi  (\alpha) ]=G (b) -G (a)
    \end{split}
\end{equation}
因为(4.22)(4.23)两个积分相等，所以
\[\int^b_a f(x) \dd x= \int^{\beta}_{\alpha} f [\varphi  (t) ] \varphi ' (t) \dd t\]
\end{proof}

这个公式与不定积分的换元公式很类似，所不同的是，后者最后需要将变量还原，而现在只须把积分限作相应的改变．

\begin{example}
    求$\Int^a_0 \sqrt{a^2-x^2}\dd x$
\end{example}

\begin{solution}
作代换$x=a\sin t$，当$x$从0变到$a$时，相应地$t$自0变到$\frac{\pi}{2}$，于是
\[\begin{split}
    \Int^a_0 \sqrt{a^2-x^2}\dd x=a^2\int^{\tfrac{\pi}{2}}_0 \cos ^2 t\dd t&=a^2\int^{\tfrac{\pi}{2}}_0 \frac{1+\cos 2t}{2}\dd t\\
    &=\frac{a^2}{2}\left(t+\frac{1}{2}\sin 2t\right)\Bigg|^{\tfrac{\pi}{2}}_0=\frac{\pi a^2}{4} 
\end{split}\]
\end{solution}

\begin{example}
    求$\Int^4_0 \frac{1}{1+\sqrt{x}}\dd x$
\end{example}

\begin{solution}
    设$x=u^2$，于是
\[\frac{1}{1+\sqrt{x}}=\frac{1}{1+u},\qquad \dd x=2u\dd u\]
定限：当$x=0$时，从变换$\sqrt{x}=u$中解出$u=0$；当$x=4$时，解出$u=2$. 于是
\[\begin{split}
    \Int^4_0 \frac{1}{1+\sqrt{x}}\dd x=\int^2_0 \frac{1}{1+u}\cdot 2u\dd u&=2\int^2_0 \left(1-\frac{1}{1+u}\right)\dd u\\
    &=2\left[u-\ln(1+u)\right]\Big|^2_0=2(2-\ln 3)
\end{split}\]
\end{solution}

\begin{example}
求证：$\Int^a_0 f(x)\dd x=\Int^a_0 f(a-x)\dd x$
\end{example}

\begin{proof}
    设$x=a-t$，则$\dd x=-\dd t$，且当$x=0$时，$t=a$；当$x=a$时，$t=0$. 于是：
\[\begin{split}
    \Int^a_0 f(x)\dd x=\int^0_a f(a-t)(-\dd t)&=-\int^0_a f(a-t)\dd t\\
    &=\int^a_0 f(a-t)\dd t=\Int^a_0 f(x)\dd x
\end{split}\]
\end{proof}

我们在第三节曾指出，使用第二种变量代换时，要注意条件是否满足，否则会出现矛盾，这里举一反例．

求$\Int^{1}_{-1}\frac{\dd x}{1+x^2}$时，如作变换$x=\frac{1}{t}$，则
\[\Int^{1}_{-1}\frac{\dd x}{1+x^2}=-\Int^{1}_{-1}\frac{\dd t}{1+t^2}\]
于是$\Int^{1}_{-1}\frac{\dd x}{1+x^2}=0$，这当然不对，问题在于，代换$x=\frac{1}{t}$
在$[-1,1]$上并不是连续变化的．


\subsubsection{定积分的分部积分法}
\begin{blk}
    {定理2} 设$u=u(x)$, $v=v(x)$在$[a,b]$上有连续的导数$u'(x)$, $v'(x)$, 则有分部积分公式
\[\int^b_a u\dd v=u\cdot v\Big|^b_a- \int^b_a v\dd u\]
\end{blk}

\begin{proof}
    由牛顿-莱布尼兹公式，显然有
\[\begin{split}
    u (x) v (x)\Big|^b_a &=\int^b_a (u (x) v (x) ) '\dd x\\
&=\int^b_a u' (x) v (x) \dd x+\int^b_a u (x) v' (x) \dd x
\end{split}\]
移项，得
\[\int^b_a u (x) \dd [v (x) ] =u (b) v (b) -u (a) v (a) -\int^b_a v (x) \dd [u (x) ] \]
\end{proof}

这个公式与不定积分的分部积分公式很相似，但是，这里每一项都带着积分限．

\begin{example}
求$\Int^{\sqrt{3}}_{-\sqrt{3}}|\arctan x|\dd x$
\end{example}

\begin{solution}
因为$|\arctan x|$是偶函数，所以
\[I=\Int^{\sqrt{3}}_{-\sqrt{3}}|\arctan x|\dd x=2\Int^{\sqrt{3}}_{0}|\arctan x|\dd x=2\Int^{\sqrt{3}}_{0}\arctan x\dd x\]
设$u(x)=\arctan x$, $v'(x)=1$, 则$v(x)=x$，利用分部积分公式，得到：
\[\begin{split}
    I=2\Int^{\sqrt{3}}_{0}\arctan x\dd x &=2\left[x\arctan x\Big|^{\sqrt{3}}_0-\int^{\sqrt{3}}_0 x\dd (\arctan x)\right]\\
    &=2\left[\sqrt{3}\arctan\sqrt{3}-\int^{\sqrt{3}}_0 x\cdot \frac{1}{1+x^2}\dd x \right]\\
    &=2\sqrt{3}\cdot \frac{\pi}{3}-\int^{\sqrt{3}}_0 \frac{\dd(1+x^2)}{1+x^2}\\
    &=\frac{2\sqrt{3}}{3}\pi -\ln(1+x^2)\Big|^{\sqrt{3}}_0= \frac{2\sqrt{3}}{3}\pi -\ln 4
\end{split}\]
\end{solution}


\begin{example}
    求$\Int^{\tfrac{\pi}{2}}_0 \sin^n \theta\dd \theta$.
\end{example}

\begin{solution}
    设$I_n=\Int^{\tfrac{\pi}{2}}_0 \sin^n \theta\dd \theta$，利用例4.42的结果，得到：
\[nI_n=\left[-\cos\theta\sin^{n-1}\theta\right]\Big|^{\tfrac{\pi}{2}_0} +(n-1)I_{n-2}\]
因为方括号的式子对于积分上，下限的值等于零，所以上面等式化简为
\[I_n=\frac{n-1}{n}I_{n-2}\]
分别考虑$n$是偶数和奇数两种情况：
\begin{enumerate}
    \item 若$n$是偶数，那么
\[\begin{split}
    I_n=\frac{n-1}{n}I_{n-2}&=\frac{n-1}{n}\cdot \frac{n-3}{n-2}I_{n-4}\\
    &=\cdots=\frac{n-1}{n}\cdot \frac{n-3}{n-2}\cdot \frac{n-5}{n-4}\cdots \frac{1}{2}I_{0}
\end{split}
    \]
此处$I_0=\Int^{\tfrac{\pi}{2}}_0 1\cdot \dd \theta=\frac{\pi}{2}$，所以
\[I_n=\frac{(n-1)(n-3)\cdots 3\cdot 1}{n(n-2)\cdots 4\cdot 2}\cdot \frac{\pi}{2}\]

\item 若$n$是奇数，由于
\[I_1=\Int^{\tfrac{\pi}{2}}_0 \sin\theta\dd\theta =(-\cos\theta)\Big|^{\tfrac{\pi}{2}}_0=1\]
所以
\[I_n=\frac{(n-1)(n-3)\cdots 4\cdot 2}{n(n-2)\cdots 5\cdot 3}\]
\end{enumerate}
\end{solution}

\begin{ex}
\begin{enumerate}
    \item 求下列定积分：
\begin{multicols}{2}
\begin{enumerate}
    \item $\Int_{0}^{\tfrac{\pi}{4}}(1+\sin \theta)^{2} \dd \theta $ 
    \item $\Int_{0}^{\tfrac{\pi}{4}} \frac{\cos \theta \dd \theta}{1+\sin ^{2} \theta}$
    \item $\Int_{0}^{1} \frac{x^{2}}{\left(1+x^{2}\right)^{2}} \dd x$
    \item $\Int_{0}^{\tfrac{1}{\sqrt{2}}} \frac{x^{2}}{\sqrt{1-x^{2}}} \dd x$
    \item $\Int_{\tfrac{1}{\sqrt{2}}}^{1} \sqrt{1-x^{2}} \dd x$
    \item $\Int_{0}^{x} \frac{\cos ^{2} x \dd x}{a^{2} \cos ^{2} x+b^{2} \sin x}$
    \item $\Int_{0}^{2 n}[\sin \omega t -\sin (\omega t+\varphi)]^{2} \dd t$
    \item $\Int_{2}^{4} \sqrt{(4-x)(x-2)} \dd x$
    \item $\Int_{1}^{5} \frac{1}{\sqrt{2 x-1}} \dd x$ 
    \item $\Int_{0}^{1} \frac{x+1}{\sqrt{1 - x^{2}}} \dd x$
    \item $\Int_{\tfrac{\pi}{6}}^{\tfrac{\pi}{2}} \cot \theta \dd \theta$
    \item $\Int_{1}^{0} \frac{x^{2}}{ 1+x^{3}} \dd x$
    \item $\Int^{\tfrac{\pi}{2}}_0 \frac{\sin x}{\cos^2 x}\dd x$
    \item $\Int^{2a}_0 x^2\sqrt{2ax-x^2}\dd x$
\end{enumerate}
\end{multicols}

\item 证明：$\Int^{\tfrac{\pi}{2}}_0 \frac{\cos x}{\cos x+\sin x}\dd x=\Int^{\tfrac{\pi}{2}}_0 \frac{\sin x}{\cos x+\sin x}\dd x$

（提示：用变量代换$x=\frac{\pi}{2}-y$.）
\item \begin{enumerate}
    \item 求常数$A,B,C,D$使等式$$\frac{1+2x^2}{(1-x)(1+x)^2}=A+\frac{B}{1-x}+\frac{C}{(1+x)^2}+\frac{D}{1+x}$$成立，并计算
    \[
    \Int^{\tfrac{1}{2}}_0 \frac{1+2x^3}{(1-x)(1+x)^2}\dd x.
    \]
    \item 求常数$a$，使$\Int^1_0 \frac{x-a}{(x+1)(3x+1)}\dd x=0$.
\end{enumerate}
\item 求极限：$\Lim_{n\to\infty}\left(\frac{1}{n+1}+\frac{1}{n+2}+\cdots +\frac{1}{n+n}\right)$
\item 求一三次多项式$f(x)=ax^3+bx^2+cx+d$, 使它满足下
列条件：
\begin{enumerate}
    \item $f (-1) =f (1) =0$;
    \item $\Int^1_{-1} f(x)\dd x=6$.
\end{enumerate}
\item 
\begin{enumerate}
\item 在同一个坐标系内，作$y=(x-1)(x-2)$和$y=\Int^x_0 (t-1)(t-2)\dd t$的函数图象．
\item 求函数$\Int^x_0 (t-1)(t-2)\dd t$在闭区间$[0, 3]$上的最大值和最小值；
\item 求在曲线$y=(x-1)(x-2)$, $x$轴和纵坐标$x=0$, $x=3$之间的区域的面积．
\end{enumerate}
\item 
\begin{enumerate}
    \item 若$f(2a-x)=-f(x)$, 求证$\Int^{2a}_0 f (x) \dd x=0$；
\item 
若$f(2a-x)=f(x)$, 求证$\Int^{2a}_0 f (x) \dd x=2\int^a_0 f (x) \dd x$. 
\item 试计算$\int^{\pi}_0 \sin^6 x\cos^3x\dd x$和$\Int^{\pi}_0 \sin^3x \cos^2 x\dd x$.
\end{enumerate}

\item 求$\Int^{\tfrac{\pi}{2}}_0 \cos^n \theta\dd \theta$.
\item 求$\Int^{\tfrac{\pi}{4}}_0 x^3\cos2x\dd x$.
\item 求$\Int^{\tfrac{3\pi}{2}}_0 \sin^3\theta \cos^2 \theta\dd \theta$.
\end{enumerate} 
\end{ex}

\subsection{平面图形的面积}
我们已经证明过在曲线$y=f(x)$与$x$轴，$x=a$, $x=b$之间的曲边梯形的有号面积等于
\[\int^b_a y\dd x=\int^b_a f(x)\dd x\]
对于任意的平面图形．我们总可以取适当的直角坐标系，使得它的边界为两条曲线
\[y=f (x) ,\qquad  y=g (x) ,\quad x\in [a,b]\]
所描述，现在我们来求这样的区域$R$的面积（图4.4）．

\begin{blk}
   {定理} 若$f$和$g$是在区间$[a,b]$上的连续函数，而且$f(x)\ge g(x)$, 那么在曲线$y=f(x)$, $y=g(x)$和$x=a$, $x=b$之间的区域$R$的面积等于
\[A=\int^b_a [f (x) -g (x) ] \dd x\] 
\end{blk}

\begin{figure}[htp]\centering
	\begin{minipage}[t]{0.48\textwidth}
		\centering
        \begin{tikzpicture}[>=latex, scale=.8]
\draw[->](-2.5,0)--(3,0)node[right]{$x$};
\draw[->](0,-3)--(0,3)node[right]{$y$};
\draw[domain=-2:2, samples=100, very thick]plot(\x, {.15*(-\x*\x-4*\x +5) })node[right]{$y=f(x)$};
\draw[domain=-2:2, samples=100, very thick]plot(\x, {.1*(\x*\x-4*\x -15)})node[right]{$y=g(x)$};
\draw[domain=-2:2, samples=100, pattern=north east lines, dashed]plot(\x, {.15*(-\x*\x-4*\x +5) })--plot(-\x, {.1*(\x*\x+4*\x -15)});
\node at (0,0)[below left]{$O$};
\draw[dashed](-1.75,-2)--(-1.75,2)node[above]{$x=a$};
\draw[dashed](1.75,-2)--(1.75,2)node[above]{$x=b$};



    \end{tikzpicture}
    \caption{}
\end{minipage}
\begin{minipage}[t]{0.48\textwidth}
    \centering
\begin{tikzpicture}[>=latex, scale=.8]
\draw[->](-2.5,0)--(3,0)node[right]{$x$};
\draw[->](0,-3)--(0,3)node[right]{$y$};
\draw[domain=-2:2, samples=100, very thick]plot(\x, {.15*(-\x*\x-4*\x +5) })node[right]{$y=f(x)-k$};
\draw[domain=-2:2, samples=100, very thick]plot(\x, {.1*(\x*\x-4*\x -15)})node[right]{$y=g(x)-k$};
\draw[domain=-2:2, samples=100, pattern=north east lines, dashed]plot(\x, {.15*(-\x*\x-4*\x +5) })--plot(-\x, {.1*(\x*\x+4*\x -15)});
\node at (0,0)[below left]{$O$};

\node at (0,3)[left]{$y'$};
\draw[thick, ->](-2.5,-2.5)--(3,-2.5)node[right]{$x'$};
\draw[dashed](-1.75,-2.5)node[below]{$a$}--(-1.75,2);
\draw[dashed](1.75,-2.5)node[below]{$b$}--(1.75,2);
\node at (0,-2.5)[below right]{$O'$};
\end{tikzpicture}
\caption{}
\end{minipage}
\end{figure}

\begin{proof}
让我们选择一个常数$k$使其小于函数$g$在区间$[a,b]$上的最小值，则对于每一个$x\in [a,b]$, $g(x)
-k\ge 0$.

又因为$f(x)\ge g(x)$, 因此对于每一个$x\in [a,b]$, 
也有$f(x)-k\ge 0$. 将坐标系$x,y$平行于$y$轴平移到新坐标系$x',y'$, 使新原点在原来坐标系中的坐标为$(0,k)$ (图4.5), 这个坐标系的平移由变换
\[\begin{cases}
    x=x'\\
    y=y'+k
\end{cases}\]
实现，于是原来的函数图象在新坐标的方程是
\[f_1 (x) =f (x) -k,\qquad g_1 (x) =g (x) -k\]
函数$f_1$和$g_1$在区间$[a,b]$上仍连续且取非负值，显然在函数$f_1$和$g_1$之间的区域$R=\{a\le x\le b,\; g_1(x)\le y\le f_1(x)\}$就是原来的区域$R=\{a\le x\le b,\; g(x)\le y\le f(x)\}$, 但是，$R'$的区域的面积等于
\[A=\int^b_a f_1(x)\dd x-\int^b_a g_1(x)\dd x\]
根据定积分的性质进一步得到
\[\begin{split}
    A&=\int^b_a [f_1(x)-g_1(x)]\dd x\\
    &=\int^b_a \Big\{\big[f(x)-k\big]-\big[g(x)-k\big]\Big\}\dd x\\
    &=\int^b_a [f(x)-g(x)]\dd x\\
\end{split}\]
\end{proof}

\begin{example}
求直线$y=x-2$和抛物线$y=2x-x^2$所围成的图形的面积$A$.
\end{example}

\begin{solution}
直线与抛物线交于$(2, 0)$, $(-1,-3)$两点，作出它的草图（图4.6）. 依本节定理，得到
\[\begin{split}
    A=\int^2_1 [(2x-x^2)-(x-2)]\dd x
    &=\int^2_1 (-x^2+x+2)\dd x\\
    &=\left[-\frac{1}{3}x^3+\frac{1}{2}x^2+2x\right]\Bigg|^2_1=\frac{9}{2}
\end{split}\]
\end{solution}


\begin{figure}[htp]\centering
	\begin{minipage}[t]{0.48\textwidth}
		\centering
        \begin{tikzpicture}[>=latex, scale=.7]
\draw[->](-2,0)--(5,0)node[right]{$x$};
\draw[->](0,-5)--(0,2.5)node[right]{$y$};
\draw[domain=-1.25:3.25, samples=100, very thick]plot(\x, {2*\x-\x*\x});
\draw[domain=-2:4, samples=10, very thick]plot(\x, {\x-2});
\node at (1,1)[above]{$(1,1)$};
\node at (2,0)[below right]{$(2,0)$};
\node at (-1,-3)[left]{$(-1,-3)$};
\node at (0,0)[above left]{$O$};
\draw[domain=-1:2, samples=100, pattern=north west lines]plot(\x, {2*\x-\x*\x})--(-1,-3);

    \end{tikzpicture}
    \caption{}
\end{minipage}
\begin{minipage}[t]{0.48\textwidth}
    \centering
\begin{tikzpicture}[>=latex, scale=.8]
    \draw[->](-2,0)--(3.5,0)node[right]{$x$};
\draw[->](0,-3.5)--(0,3.5)node[right]{$y$};
\node at (0,0)[below left]{$O$};
\draw[domain=-2.5:2.5, samples=100, very thick]plot({2-0.5*\x*\x}, \x);
\draw[domain=-2.75:2.75, samples=100, very thick]plot({1-0.25*\x*\x}, \x);
\foreach \x in {2,-2}
{
    \node at (0,\x) [left]{$\x$};
}
\node at (1,0)[below left]{1};
\node at (2,0)[below right]{2};
\draw[domain=-2:2, samples=100, pattern=north east lines]plot({2-0.5*\x*\x}, \x)--plot({1-0.25*\x*\x}, -\x);


\end{tikzpicture}
\caption{}
\end{minipage}
\end{figure}


\begin{example}
求抛物线$y^2=-4(x-1)$, $y^2=-2(x-2)$所围成的图形的面积．
\end{example}

\begin{solution}
    先作出它们的草图（图4.7）, 且解得它的两个交点$(0, 2)$和$(0,-2)$. 于是
\[\begin{split}
    A=\int^2_{-2} \left[\left(2-\frac{1}{2}y^2\right)-\left(1-\frac{1}{4}y^2\right)\right]\dd y
    &=\frac{1}{4}\int^2_{-2} (4-y^2)\dd y\\
    &=\frac{1}{4}\left[4y-\frac{1}{3}y^3\right]\Bigg|^2_{-2}=\frac{8}{3}.
\end{split}\]
\end{solution}

\begin{example}
求半径等于$r$, 中心角等于$\alpha \; \left(0<\alpha\le \frac{\pi}{2}\right)$的扇形的面积．
\end{example}

\begin{figure}[htp]
    \centering
\begin{tikzpicture}[>=latex, scale=1.5]
\draw[->](-1.5,0)--(2,0)node[right]{$x$};
\draw[->](0,-1.5)--(0,1.5)node[right]{$y$};
\draw[very thick] (0,0) circle (1);
\draw[pattern=north west lines](0,0)--(1,0)node[below right]{$B$} arc (0:60:1)node[above right]{$P(r\cos\alpha, r\sin\alpha)$}--(0,0);
\node at (0,0)[below left]{$O$};
\draw(.5,0)node[below]{$C$}--(60:1);
\draw[<->](.2,0) arc (0:60:.2)node[right]{$\alpha$};

\end{tikzpicture}
    \caption{}
\end{figure}

\begin{solution}
    如图4.8, 扇形$OBP$的区域划分为$\triangle OCP$的区域与曲边三角形$CBP$的区域的和，因此
\[A (OBP) =A (OCP) +A (CBP) \]
因为$P$点坐标是$(r\cos\alpha, r\sin\alpha)$，所以
\[A(OCP)=\frac{1}{2}(r\cos\alpha)\cdot (r\sin\alpha)=\frac{r^2}{4}\sin 2\alpha\]
又扇形的圆方程是$x^2+y^2=r^2$，所以
\[A(CBP)=\int^r_{r\cos \alpha}\sqrt{r^2-x^2}\dd x\]
从而
\[A(OBP)=\frac{r^2}{4}\sin2\alpha+\int^r_{r\cos \alpha}\sqrt{r^2-x^2}\dd x\]

令$x=r\sin\theta,\; \theta\in\left[-\frac{\pi}{2},\frac{\pi}{2}\right]$，则
\[\theta=\arcsin\frac{x}{r},\qquad \dd x=r\cos\theta \dd\theta\]
且当$x=r\cos \alpha$时，
\[\theta=\arcsin(\cos\alpha)=\arcsin\left[\sin\left(\frac{\pi}{2}-\alpha\right)\right]=\frac{\pi}{2}-\alpha\]
当$x=r$时，$\theta=\arcsin 1=\frac{\pi}{2}$

所以
\[\begin{split}
    A(OBP)&=\frac{r^2}{4}\sin2\alpha+\int^{\tfrac{\pi}{2}}_{\tfrac{\pi}{2}-\alpha}(r\cos\theta)(r\cos\theta)\dd \theta\\
&=\frac{r^2}{4}\sin2\alpha+r^2\int^{\tfrac{\pi}{2}}_{\tfrac{\pi}{2}-\alpha}\cos^2\theta\dd\theta\\ 
&=\frac{r^2}{4}\sin2\alpha+r^2\left.\left[\frac{\theta}{2}+\frac{\sin 2\theta}{4}\right]\right|^{\tfrac{\pi}{2}}_{\tfrac{\pi}{2}-\alpha}\\
&=\frac{r^2}{2}\alpha
\end{split}\]

当$\alpha=2\pi$时，由上面公式就得到圆面积$\pi r^2$.
\end{solution}

\begin{example}
$P$点是位于曲线$y=x^2$上的动点，从原点$O$向点$A(1, 1)$移动，过二点$O,P$引直线，它与曲线$y=x^2$及直线$x=1$所围图形的面积为图4.9中阴影部分所示，求使这部分面积最小时$P$点的坐标．
\end{example}

\begin{figure}[htp]\centering
        \begin{tikzpicture}[>=latex, scale=3]
\draw[->](-.2,0)--(1.5,0)node[right]{$x$};
\draw[->](0,0)node[below]{$O$}--(0,1.5)node[right]{$y$};
\draw(0,1)--(1,1)node[right]{$A(1,1)$}--(1,0);
\draw[domain=0:1.15, samples=100, very thick]plot(\x, {\x*\x})node[above right]{$y=x^2$};
\draw[domain=0:1, samples=100, thick, pattern=north west lines]plot(\x, {\x*\x})--(1,.8)--(0,0);
\node at (.8,.64)[left]{$P(t,t^2)$};
        \end{tikzpicture}
        \caption{}
    \end{figure}

\begin{solution}
    设$P$点的坐标是$(t,t^2),\; (0\le t\le 1)$. 直线$OP$的方程是$y=tx$，设阴影部分的面积为$S$，则
\[\begin{split}
S&=\int^t_0 (tx-x^2)\dd x+\int^1_t (x^2-tx)\dd x\\
&=\left[\frac{t}{2}x^2-\frac{x^3}{3}\right]\Bigg|^t_0+\left[\frac{x^3}{3}-\frac{t}{2}x^2\right]\Bigg|^1_t\\
&=\frac{t^3}{3}-\frac{t}{2}+\frac{1}{3}
\end{split}\]
\[S'=t^2-\frac{1}{2}=\left(t+\frac{\sqrt{2}}{2}\right)\left(t-\frac{\sqrt{2}}{2}\right)\]
由$S'$符号的变化，当$t=\frac{\sqrt{2}}{2}$时，$S$极小也就是最小，故所求点的坐标是$\left(\frac{\sqrt{2}}{2},\frac{1}{2}\right)$.
\end{solution}

\begin{example}
    求椭圆$\frac{x^2}{a^2}+\frac{y^2}{b^2}=1\; (a>b>0)$的面积$A$（图4.10）.
\end{example}

\begin{figure}[htp]
    \centering
\begin{tikzpicture}[>=latex]
    \draw[->](-3,0)--(3,0)node[right]{$x$};
\draw[->](0,-2)--(0,2)node[right]{$y$};
\node at (0,0) [below left]{$O$};
\draw[very thick](0,0) ellipse [x radius=2.2, y radius=1.3];

\end{tikzpicture}
    \caption{}
\end{figure}

\begin{solution}
由于图形的对称性，只要求出椭圆在第一象限内的面积，再乘以4就是该椭圆的面积，所以
\[\frac{A}{4}=\int^a_0 y\dd x\]
这里将椭圆用参数方程
$\begin{cases}
    x=a\cos t\\ y=b\sin t
\end{cases}$
表出，当$x=a$时，$t=\arccos 1=0$; 当$x=0$时，$t=\arccos 0=\frac{\pi}{2}$, 于是
\[\begin{split}
    \frac{A}{4}=\int^0_{\tfrac{\pi}{2}}b\sin t \dd(a\cos t)&=\int^0_{\tfrac{\pi}{2}}(b\sin t )(-a\sin t)\dd t\\
&=ab \int^{\tfrac{\pi}{2}}_0 \sin^2 t\dd t=\frac{1}{4}ab\pi 
\end{split}\]
所以：$A=\pi ab$.
\end{solution}

\begin{ex}
\begin{enumerate}
    \item 求在曲线$y=\sin2x$, $y=0$, $x\in[0, 2\pi]$之间的面积．
    \item 求下列曲线所围图形的面积：
\begin{enumerate}
    \item $y=8-\frac{1}{2}x^2$和$y=3.2$ 
    \item $y=x^2-5x+7$和$y=-2x^2+10x-5$
    \item $y=\sin x$和$y=\frac{2}{\pi} x,\; x\ge 0$
\end{enumerate}
\item 求以曲线$y=\sin x\; \left(0\le x\le \frac{\pi}{4}\right)$；$y=\cos x\; \left(\frac{\pi}{4}\le x\le \frac{\pi}{2}\right)$和$y=0$为边界的区域的面积．

\item \begin{enumerate}
    \item 证明点$x=t^2-1$, $y=t(t^2-1)$在曲线$y^2=x^2(x+1)$上；
    \item 当$t$由$t=-2$变化到$t=2$时，作出曲线$\begin{cases}
        x=t^2-1\\x=t (t^2-1)
    \end{cases}$的图象；
\item 求由曲线的环路所围区域的面积．
\end{enumerate}

\item 在曲线$y=x^2-2x+2$与过点$(2, 3)$的直线所围的图形
中，求图形面积最小时的直线方程．
\item 求抛物线$y=x^2$及过此抛物线上点$(2, 4)$的切线和$x$轴所围图形的面积．
\item 已知曲线$C:\; y=x^2-7x+10\; (x>0)$与点$A(0, 1)$, 
回答下列问题：
\begin{enumerate}
    \item 求过点$A$所引曲线$C$的切线方程以及切点$T$的坐标．
    \item 设$S$是直线$x=1$与曲线$C$的交点．若曲线$C$上的动点$P$从$S$点运动到$T$点时，将线段$AP$所在的区域用阴影表示出来，并求出这个区域的面积．
\end{enumerate}

\item 求通过$(0, 0)$, $(1, 2)$两点的抛物线，要求它具有以
下的性质：
\begin{enumerate}
    \item 它的对称轴平行于$y$轴，且凸向上，
    \item 它与$x$轴所围区域的面积最小．
\end{enumerate}

\item 对称轴平行$y$轴的抛物线与圆$x^2+y^2=4$切于$A(0, 2)$点并通过$B(-2, 0)$点，计算抛物线与圆所围区域的面积．
\end{enumerate}
\end{ex}

\subsection{利用横断面算体积法}
对于空间的一个立体，我们取一条直线作为坐标轴，取直线上一点作为原点（图4.11），假定这立体紧夹于在$x=a$与$x=b$两点所作的$x$轴的垂直平面之间，又如果在离原点$x$处作这轴的一个垂直平面，并且该垂直平面截得立体的截面的面积是$x$的连续函数$A(x),\; a\le x\le b$, 那么立体的体积可由定积分来计算：
\begin{figure}[htp]
    \centering
\begin{tikzpicture}[>=latex]
    \draw[->](-.5,0)--(6,0)node[right]{$x$};
\tkzDefPoints{.8/0/a, 2/0/A, 3.3/0/x_i, 5/0/b, 0/0/O}
\tkzDrawPoints(a,A,x_i,b,O)
\tkzLabelPoints[below](a,x_i,b,O)
\tkzLabelPoint[below](A){$x_{i-1}$}
\node at (0,0){};
\draw[line width = 1pt] plot[smooth] coordinates{(.8,0)(1.2,.8)(2,1.3)(3.3,1.15)(5,0)(3.3,-1.15)(2,-1.3)(1.2,-.8)(.8,0)
};
\draw[dashed](A) ellipse [x radius=.35, y radius=1.3]; 
\draw[dashed](x_i) ellipse [x radius=.3, y radius=1.15]; 
\draw[thick](2,1.3) arc [start angle =90, end angle =270, x radius=.35, y radius=1.3];
\draw[thick](3.3,1.15) arc [start angle =90, end angle =270, x radius=.3, y radius=1.15];

\end{tikzpicture}
    \caption{}
\end{figure}

\begin{enumerate}
    \item 分$[a,b]$为$n$份，其分点
\[a=x_0<x_1<x_2<\cdots<x_n=b\]
    记$\|P_n\|=\max(x_i-x_{i-1})\quad i=1, 2,\ldots,n$.

    以平面$x=x_i,\; i=0, 1, 2,\ldots,n$截此立体为$n$个小单元，于是整个体积被看作这$n$个小单元的和．
\item  求阶梯函数的总和作为所求立体体积的近似和．考虑这样一个单元，它介于两个截面$x=x_{i-1},x=x_i$之间，于是以$(x_i-x_{i-1})$为高，以$A(x_{i-1})$作为底面积的直棱柱的体积是
  \[  A (x_{i-1}) (x_i-x_{i-1})\]
    于是，我们得到所求的体积$V$的近似表达式
    \[\sum^n_{i=1}A (x_{i-1}) (x_i-x_{i-1})\] 
    \item 对和取极限：当断面无限增加，$\|P_n\|\to 0$时，有
\[V=\lim_{n\to\infty}\sum^n_{i=1}A (x_{i-1}) (x_i-x_{i-1})=\int^b_a A(x)\dd x\]
\end{enumerate}

现在将上面的讨论总结为下面的定理．

\begin{blk}
   {定理} 如果已知一个给定的立体的垂直于一定方向所有的横断面的面积，取这些横断面的垂直方向作为$x$轴的方向，则这立体的体积由公式
\[V=\int^b_a A (x) \dd x\]
表达．其中$A(x)$是横坐标为$x$的横断面的面积，$a$、$b$为这立体的两端断面的横坐标． 
\end{blk}

\begin{example}
    求底边长为$a$, 高为$h$的正四棱锥的体积．
\end{example}

\begin{figure}[htp]
    \centering
\begin{tikzpicture}[>=latex]
    \tkzDefPoints{0/0/A, 3/0/B, 4/1/C}
    \tkzDefPointsBy[translation= from B to C](A){D}
    \tkzInterLL(A,C)(B,D)  \tkzGetPoint{O}
    \tkzDefPoints{0/3.5/E, 5/3.5/F}
    \tkzDefPointBy[projection = onto E--F](O) \tkzGetPoint{OO}
    \draw[|<->|](.7,1.4)--node[fill=white]{$y$}(2.7,1.4);
    \tkzDrawSegments(OO,A OO,B OO,C A,B B,C)
    \tkzDrawSegments[dashed](OO,D OO,O C,D D,A)
    
    \foreach \x in {A,B,C,D,O}
    {
        \tkzDefPointWith[linear, K=.66](OO,\x)  \tkzGetPoint{\x'}
        \tkzDefPointWith[linear, K=.33](OO,\x)  \tkzGetPoint{\x''}
    }
    \tkzDrawSegments(A',B' B',C' A'',B'' B'',C'')
    \tkzDrawSegments[dashed](C',D' D',A' C'',D'' D'',A'')
\draw[->](OO)--(2,4)node[right]{$x$};
\draw(-1,3.5)--(5,3.5);
\draw[dashed](.2,.5)--(O);
\draw(-1,0.5)--(.2,.5);
\draw[<->](-.8,0.5)--node[fill=white]{$h$}(-.8,3.5);
\draw[<->](4.8,0.5)--node[fill=white]{$x$}(4.8,1.5);
\draw[<->](4.8,1.5)--node[fill=white]{$h-x$}(4.8,3.5);
\draw[<->](3.8,1.5)--node[fill=white]{$\Delta x$}(3.8,2.5);
\tkzDefMidPoint(B,C) \tkzGetPoint{G}
\tkzDefMidPoint(B',C') \tkzGetPoint{G'}
\tkzDefMidPoint(B'',C'') \tkzGetPoint{G''}
\tkzDrawSegments[dashed](O,G O',G' O'',G'')
\tkzDefPointWith[linear, K=2](O,G) \tkzGetPoint{H}
\tkzDefPointWith[linear, K=3](O',G') \tkzGetPoint{H'}
\tkzDefPointWith[linear, K=4.5](O'',G'') \tkzGetPoint{H''}
\tkzDrawSegments(G,H G',H' G'',H'')
\tkzLabelPoints[below](O)

\draw[|<->|](0,-0.25)--node[fill=white]{$a$}(3,-.25);


\end{tikzpicture}
    \caption{}
\end{figure}

\begin{solution}
    设$x$为从底面到截面的距离，$y$为截面正方形的边长（图4.12），于是
\[\frac{h-x}{h}=\frac{y}{a}\]
即：$y=\frac{a}{h}(h-x)$.

截面面积
\[A(x)=\frac{a^2}{h^2}(h-x)^2\]
所以
\[\begin{split}
    V&=\int^h_0 \left(\frac{a}{h}\right)^2(h-x)^2 \dd x\\
    &=-\left(\frac{a}{h}\right)^2\frac{(h-x)^3}{3}\Bigg|^h_0=\frac{1}{3}a^2 h
\end{split}\]
\end{solution}

如果立体是一个旋转体，即先给出一条平面曲线
$y=f(x)$, 绕$x$轴旋转所得出的立体，我们的横断面就是以$y$为半径的圆，所以这立体的体积就等于
\[V=\pi \int^b_a y^2\dd x\]
如果$c\le y\le d$，$x=g(y)$绕$y$轴旋转，则
\[V=\pi \int^d_c x^2\dd y\]


\begin{example}
若球的半径是$a$, 球缺的高是$h$, 求球缺的体积．
\end{example}

\begin{figure}[htp]
    \centering
\begin{tikzpicture}[>=latex, scale=1.5]
    \draw[->](-1.5,0)--(2,0)node[right]{$x$};
    \draw[->](0,-1.5)--(0,1.5)node[right]{$y$};
    \draw[very thick](0,0)circle(1);
\draw(.4,1)--(.4,-1.5);
\draw(1,0)--(1,-1.5);
\draw[<->](0,-1.3)--node[below]{$a-h$}(.4,-1.3);
\draw[<->](0.4,-1.1)--node[fill=white]{$h$}(1,-1.1);
\draw[<->](-.2,0)--node[fill=white]{$a$}(-.2,1);

\draw[pattern=north east lines](45:1) arc (45:50:1)--(.64,0)--(1/1.414,0)--(45:1);
\draw[<->](0,1.3)--node[below]{$x$}(.64,1.3);
\draw(.64,1.5)--(.64,0);
\end{tikzpicture}
    \caption{}
\end{figure}


\begin{solution}
球是由圆$x^2+y^2=a^2$绕$x$轴旋转所得到的立体，且球心到球缺底面的距离为$a-h$, 设球心到截面的距离为$x$, 则截面圆的面积$A (x) =\pi y^2=\pi  (a^2-x^2)$.（图4.13）

所以
\[\begin{split}
    V=\pi\int^a_{a-h}(a^2-x^2)\dd x&=\pi\left[a^2 x-\frac{1}{3}x^3\right]\Bigg|^a_{a-h}\\
    &=\frac{\pi h^2}{3}(3a-h)
\end{split} \]
\end{solution}

\begin{example}
    计算底面半径分别为$r$和$R$, 高为$h$的圆台体积$V$.
\end{example}

\begin{solution}
此圆台可看作$xy$平面上的直线段
\[y=\frac{R-r}{h}x+r,\; (0\le x\le h),\quad x=0,\quad x=h,\quad y=0\]
所围成的直角梯形绕$x$轴旋转所得到的立体，从而
\[\begin{split}
    V&=\pi\int^h_0 \left(\frac{R-r}{h}x+r\right)^2 \dd x\\
    &=\pi\left[\left(\frac{R-r}{h}\right)^2\frac{x^3}{3}+2r\cdot \frac{R-r}{h}\cdot \frac{x^2}{2}+r^2 x\right]\Bigg|^h_0\\
    &=\frac{\pi h}{3}(R^2+Rr+r^2)
\end{split}\]
\end{solution}

\begin{example}
在第一象限中，由$y=\frac{x}{1-x}\; (x\ne 1)$, $y=1$, $y=0$和$x=1$所围成的图形绕着直线$x=1$旋转，求所生成的旋转体的体积（图4.14）．
\end{example}

\begin{figure}[htp]
    \centering
\begin{tikzpicture}[>=latex, scale=2.5]
    \draw[->](-.25,0)--(1.8,0)node[right]{$x$};
    \draw[->](0,-.25)--(0,1.8)node[right]{$y$};
\draw[domain=0:.6, samples=100, very thick]plot(\x, {\x/(1-\x)});
\draw(1,0)node[below]{1}--(1,1.5)node[above]{$x=1$};
\draw[dashed](0,1)node[left]{1}--(1,1);
\draw[domain=0:.5, samples=100, pattern=north east lines]plot(\x, {\x/(1-\x)})--(1,1)--(1,0)--(0,0);

\draw[<->|](1.1,0)--node[fill=white]{$y$}(1.1,.4);
\draw[domain=0.286:.333, samples=10, pattern=north west lines]plot(\x, {\x/(1-\x)})--(1,.5)--(1,0.4)--(0.286,0.4);
\draw[|<->|](0.286,0.25)--node[fill=white]{$1-x$}(1,.25);

\end{tikzpicture}
    \caption{}
\end{figure}

\begin{solution}
$y$轴垂直于旋转体的横断面，设原点到横断面的距离为$y$, 则该横断面的面积
\[A (y) =\pi  (1-x)^2\]
由于函数$y=\frac{x}{1-x}$是单射的，故反函数存在，由它解出
\[x=\frac{y}{1+y},\qquad 1-x=\frac{1}{1+y}\]
于是
\[A(y)=\frac{\pi}{(1+y)^2}\]
\[\begin{split}
    V=\int^1_0 A(y)\dd y&=\pi\int^1_0 \frac{\dd y}{(1+y)^2}\\
&=-\pi\left(\frac{1}{1+y}\right)\Bigg|^1_0=-\pi\left(\frac{1}{2}-1\right)=\frac{\pi}{2}
\end{split}\]
\end{solution}

\begin{figure}[htp]
    \centering
    \begin{tikzpicture}[>=latex, scale=1.5]
\draw[dashed](0,0) ellipse [x radius=1.5, y radius=.5];
\draw[dashed](0,3) ellipse [x radius=1.5, y radius=.5];
\draw[dashed](-1.5,0)--(-1.5,3);
\draw[dashed](1.5,0)--(1.5,3);
\draw(.5,-.45)node[below right]{$B$}--(-.5,.45)node[above left]{$A$};
\draw[thick] (-.5,.45) .. controls (1.8,2.8) and (2.1,1.6) .. (.5,-.45);
\draw(0,0)--(1.4,1.7);
\draw[dashed](0,0)--(1.4,.2)--(1.4,1.7);
\tkzDefPoints{0/0/O, 1.4/.2/B1, 1.4/1.7/C}
\tkzMarkAngle[mark=none, size=.4, ->](B1,O,C)
\tkzLabelAngle[pos=.6](B1,O,C){$60^{\circ}$}
\tkzLabelPoints[left](O)
\draw(1.5,0) arc [start angle = 0, end angle =-70, x radius=1.5, y radius=.5];
\draw(1.5,0)--(1.5,1.5);
    \end{tikzpicture}    
    \caption{}
\end{figure}

\begin{ex}
\begin{enumerate}
    \item 底面半径为$a$的直圆柱，用通过底面一条直径并与底面成$60^{\circ}$角的平面去
    截，得到图4.15的立体，求其体积．
    \item 由抛物线$y^2=4ax$和直线$x=h$所围成
    的区域绕$x$轴旋转一周，求所生成旋转体的体积．
    \item 求底面半径为$r$, 高为$h$的圆锥体的
    体积．
    \item 设由椭圆$\frac{x^2}{a^2}+\frac{y^2}{b^2}=1$所围成的区域绕$x$轴旋转，求此旋转
    椭圆体的体积．
    \item 由曲线$y=\frac{\ln x}{x}$, $x$轴和通过曲线的最高点的纵坐标所围
    成的区域绕$x$轴旋转，求该旋转体的体积．
    \item 两条抛物线有公共的顶点和公共的对称轴，但位于两个
    互相垂直的平面内，一个运动着的椭圆，其中心在两抛物线的公共轴上，它的平面与公共轴垂直，它的顶点位于两条抛物线上，如果椭圆从公共顶点开始，移动一段距离$h$, 求由此椭圆所生成的体积．
    \item 求$y=4-x^2$与$x$轴之间的区域绕直线$y=-1$一周所生成的体积．
    \item 求圆$x^2+(y-b)^2=a^2\; (0<a<b)$绕$x$轴旋转所生成的旋转体的体积．
\end{enumerate}
\end{ex}

\section{简单初值问题{}---{}不定积分的简单应用}

在本章第2节，我们曾说过，一个函数的原函数有无穷多个，它们彼此之间相差一个常数．但是，在许多问题中，我们感兴趣的往往并不是这无穷多个原函数的全体（不定积分），而是某一个特定的原函数，在几何上，就是一条特定的积分曲线．为了求得这条积分曲线，必须知道它所经过的一个点，也就是说，要知道初始条件，在实际问题中，求满足一定初始条件的原函数问题是很多的，下面举一些简单的实例．

\begin{example}
    已知一质点$m$作具有匀加速度$a$的直线运动，设它的初速度为$u$, 求质点在$t$秒后所经过的距离．
\end{example}

\begin{solution}
由于$\frac{\dd v}{\dd t}=a$，所以
\begin{equation}
    v=at+C
\end{equation}
因为，当$t=0$时，$v=u$，代入(4.24)得：$C=u$，所以
\begin{equation}
    v=at+u
\end{equation}
由于$\frac{\dd s}{\dd t}=u+at$，分离变量，把它写成
\[\dd s=(u+at)\dd t\]
两边求不定积分，得
\[\int \dd s=\int (u+at)\dd t\]
\begin{equation}
    s=ut+\frac{1}{2}at^2+C'
\end{equation}
由初始条件：当$t=0$时，$s=0$，代入(4.26)，得：$C'=0$，所以
\begin{equation}
    s=ut+\frac{1}{2}at^2
\end{equation}
为所求．(4.25)和(4.27)就是我们常见的运动学公式.
\end{solution}


\begin{example}
假设在原点处斜抛一质点$m$, 它的初速度$v_0$与水平$x$轴成$\alpha$角，又作用在质点$m$上的力只有向下的重力，求质点$m$的轨迹方程．
\end{example}

\begin{solution}
    设质点$m$的坐标是$(x,y)$，它的质量为$M$克，于是沿着$x$轴和$y$轴的速度分量分别是
    \[v_x=\frac{\dd x}{\dd t},\qquad v_y=\frac{\dd y}{\dd t}.\]
    而沿着$x$轴和$y$轴的加速度分量分别是
    \[\frac{\dd v_x}{\dd t}=\frac{\dd^2 x}{\dd t^2},\qquad \frac{\dd v_y}{\dd t}=\frac{\dd^2 y}{\dd t^2}.\]
    根据牛顿第二定律，得到
\[
\begin{cases}
    M\frac{\dd v_x}{\dd t}=0\\
    M\frac{\dd v_y}{\dd t}=-Mg
\end{cases}
\]
即：
\[\begin{cases}
    \dd v_x=0\\ 
    \dd v_y=-g\dd t
\end{cases}\]
所以
%\begin{align}
\[
\begin{cases}
    v_x=C_1\\ 
    v_y=-gt+C_2
\end{cases}
\] %\end{align}
常数$C_1$和$C_2$由初始条件决定，所以
\[ %\begin{align}
\begin{cases}
v_x=\frac{\dd x}{\dd t}=v_0\cos\alpha\\
v_y=\frac{\dd y}{\dd t}=-gt+v_0\sin\alpha
\end{cases}
\] 
即：
\[\begin{cases}
    \dd x=v_0\cos\alpha \dd t\\ \dd y=(-g t+v_0\sin\alpha)\dd t
\end{cases}\]
再求不定积分得
\[\begin{cases}
    x=(v_0\cos\alpha)t+C_3\\
    y=-\frac{gt^2}{2}+(v_0\sin\alpha)t+C_4
\end{cases}\]
由于，当$t=0$时，$x=y=0$，因此：
\[C_3=C_4=0\]
于是得到质点$m$以时间$t$为参变数的轨迹方程
\begin{align}
\begin{cases}
    x=(v_0\cos\alpha)t\\
    y=-\frac{gt^2}{2}+(v_0\sin\alpha)t
\end{cases}
\end{align}
消去$t$，得：
\[y=\frac{-gx^2}{2v^2_0 \cos^2\alpha}+x\tan\alpha\]
化为标准式，得
\[\left(x-\frac{v^2_0\sin2\alpha}{2g}\right)^2=-\frac{2v^2_0\cos^2\alpha}{g}\left(y-\frac{v^2_0\sin^2\alpha}{2g}\right)\]
这是顶点在$\left(\frac{v^2_0\sin2\alpha}{2g},\frac{v^2_0\sin^2\alpha}{2g}\right)$处的抛物线方程（如图4.16），且从顶点到焦点的方向是$y$轴的负方向，又焦点到准线的距离是$\frac{v^2_0 \cos^2\alpha}{g}$，因此，从顶点到它的上方的准线距离等于$\frac{v^2_0 \cos^2\alpha}{2g}$，而准线到$x$轴的距离就等于
\[\frac{v^2_0 \cos^2\alpha}{2g}+\frac{v^2_0 \sin^2\alpha}{2g}=\frac{v^2_0}{2g}\]
故准线的位置仅依赖于初速度的大小，而与抛射角的大小无关．

\begin{figure}[htp]
    \centering
\begin{tikzpicture}[>=latex]
\draw[->](-.5,0)--(5,0)node[right]{$x$};
\draw[->](0,-.5)--(0,4)node[right]{$y$};
\draw[domain=0:4.2, samples=100, very thick]plot(\x, {0.5*(-\x*\x+4*\x)});
\draw(-.5,2.75)--(4.5,2.75)node[right]{$y=\frac{v^2_0}{2g}$};
\node at (3.5,1)[right]{$y=\frac{-gx^2}{2v^2_0 \cos^2\alpha}+x\tan\alpha$};
\node at (0,0)[below left]{$O$};
\draw(2,2.75)--(2,0);
\tkzDefPoints{2/2/A, 2/1.25/S}
\tkzDrawPoints(A,S)
\tkzLabelPoints[above right](A,S)
\draw[->, thick](0,0)--(1,2)node[right]{$v_0$};


\end{tikzpicture}
    \caption{}
\end{figure}
\end{solution}

下面我们来讨论人口的增长与放射性元素蜕变的问题．

在很多单纯的自然现象中，一个随着时间而变化的量$y(t)$的变化率常常和$y(t)$当时的值成正比例，用数学式表达，即存在一个随该现象而定的常数$k$（与$t$无关），使得
\[y' (t) =ky (t) \]
或
\begin{equation}
    \frac{\dd }{\dd t}y(t)=ky(t)
\end{equation}
人口的增长与放射性元素的蜕变就是这样的两个例子．

假如人口按某一固定的“增长率”$k$增长，譬如时间单位是“年”，增长率2\%, 则$k=2\%$. 设人口随时间变化的函数是$y(t)$, 本来$y(t)$是整数，但由于数量很大，我们可以把它看作一个平滑函数，$y(t)$满足下面的条件：
\[y(t+\Delta t)-y(t)=ky(t)\cdot \Delta t\]
换言之，人口增加数量与这段时间开始的人口数，以及这段时间间隔成比例，也即
\[\frac{y(t+\Delta t)-y(t)}{\Delta t}=ky(t)\]
令$\Delta t\to 0$, 由上式得到
\[\frac{\dd }{\dd t}y(t)=ky(t)\]
也就是说，人口在时刻$t$的瞬时变率与当时的值$y(x)$成正比例．为求出函数$y(t)$, 我们来解具有初始条件的微分方程
\begin{numcases}{}
    \frac{\dd }{\dd t}y(t)=ky(t)\\
    y(0)=A
\end{numcases}
将(4.33)改写成
\[\frac{\dd y(t)}{y(t)}=k\dd t\]
即：$\ln[y(t)]=kt+C$，由初始条件确定常数$C$的值，得：
\[\ln A=C\]
所以：$\ln[y(t)]=kt+\ln A$，
移项，得：
\[\ln\frac{y(t)}{A}=kt\]
由此得：
\[\frac{y(t)}{A}=e^{kt}\]
所以：$y(t)=Ae^{kt}$为所求．



总结上面的讨论，得到以下命题：

\begin{blk}{命题}
    设$y'(t)=ky(t)$，$y(0)=A$，则：
    \[y(t)=Ae^{kt}\]
\end{blk}




\begin{example}
    若某城市现有人口10万，按照每年3\%的比率增长，问几年后人口是现有人口的2倍？
\end{example}

\begin{solution}
    设$t_1$年后人口加倍，于是依题意有
\[10e^{3\%t_1}=2\x10\]
即：$e^{0.03t_1}=2$, 两边取常用对数，得
\[\begin{split}
    0.03t_1\lg e&=\lg2\\
t_1&= \frac{100\lg2}{3\lg e}\approx \frac{100}{3}\cdot\frac{0.3010}{0. 4343}\\
&\approx \frac{100}{3}\cdot 0.6931\approx 23.
\end{split} \]
答：约23年后，人口会加倍．
\end{solution}

从问题的解法看出，人口的加倍期仅与增长率有关，即$t_1=\frac{\ln2}{k}$, 下面是二者之间的对照表．

\begin{center}
\begin{tabular}{c|ccccc}
\hline
    增长率$k$ & 0.05&0.04&0.03&0.02&0.01\\
\hline
加倍期$t_1$ & 约14年& 约17年& 约23年& 约34.5年& 约39年\\
\hline
        \end{tabular}
\end{center}

同样地可把上述命题应用于放射性元素蜕变问题．

\begin{example}
    放射性元素的数量随时间$t$变化的函数是$N(t)$, $A=N(0)$是起算时的数量，$k$是蜕变率，且知放射性元素经过$\Delta t$这段时间的蜕变，它减少的数量$\Delta N(t)=N(t+\Delta t)-N(t)$与在时刻$t$的数量$N(t)$及这段蜕变时间$\Delta t$成比
例，即
\[\Delta N (t) =-kN (t) \cdot \Delta t\]
因为放射性元素不断减少，所以$\Delta N$是负数，问经过多少年放射性元素剩下了一半？
\end{example}

\begin{solution}
    依题意，有
\[\frac{\dd N}{\dd t}=-kN (t) \]
根据本节命题，得到
\[N (t) =Ae^{-kt}\]

设经过$t_1$年，放射性元素减少了一半，通常称$t_1$为放射性元素的半衰期，于是
\[t_1=\frac{\ln 2}{k}\approx \frac{0. 6931}{k}\]
\end{solution}

下面是一些放射性元素的半衰期的简表．

\begin{center}
    \begin{tabular}{c|ccc}
    \hline
        元素 & 氩Ar$^{39}$ &碳C$^{14}$&氢H$^3$  \\
    半衰期 & 269年& 5730年& 12.3年 \\
    \hline
    元素 & 碘I$^{129}$ &铀U$^{238}$& 铀U$^{235}$ \\
    半衰期 &  $16\x10^6$年& $4.5\x 10^9$年& $0.71\x 10^9$年\\
    \hline
            \end{tabular}
    \end{center}

我们可以把岩石或化石中所含的各种微量的放射性元素，当作天然的计时仪表，因为各种放射性元素的微量都可以测量十分准确，所以我们可以通过对残留微量的测定，
用半衰期去推算岩石或化石的“年龄”，甚至于我们所居住的地球的“年龄”．

\begin{ex}
\begin{enumerate}
    \item 求下列积分曲线：
\begin{multicols}{2}
\begin{enumerate}
    \item $\begin{cases}
        y'=2x\\
        y|_{x=x_0}=y_0
    \end{cases}$
    \item $\begin{cases}
        y'=2\sqrt{y}\\
        y|_{x=0}=0
    \end{cases}$
    \item $\begin{cases}
        y'=\frac{2}{3\sqrt[3]{x}}\\
        y|_{x=0}=4
    \end{cases}$
\end{enumerate}
\end{multicols}
\item 在$20.2^{\circ}{\rm C}$, 细菌受到5\%的消毒溶液消毒，每小时细菌死
    亡的比率为11\%, 假如开始时有1000个细菌，那么在24小时后还剩下多少？
    \item 化学反应速度$v$和温度$t$的增加是按照有机体生长的规律
    进行的，当$t=0$时，$v=30$; 当$t=9.90$时，$v=60$, 求反应率以及当$t=20$时，$v$的值．
\item 若函数$y=f(x)$的变率$\frac{\dd y}{\dd x}$与当时的函数值的平方成
    正比例，且当$x=0$时，$y=f(0)=2$; 当$x=1$时，$y=f(1)=4$, 求函数$y=f(x)$.
\end{enumerate} 
\end{ex}

\subsection*{习题4.5}
\begin{enumerate}
    \item 计算下列不定积分：
\begin{multicols}{2}
\begin{enumerate}
    \item $\displaystyle\int \frac{\dd x}{1-3x^2}$
    \item $\displaystyle\int \frac{1-\sin^{-3}x}{1-\cos^2 x}\dd x$
    \item $\displaystyle\int \frac{\dd x}{\sin x\cos x}$
    \item $\displaystyle\int x\sec^2 x\dd x$
    \item $\displaystyle\int xe^{2x+1}\dd x$
    \item $\displaystyle\int \frac{x^2+2}{x(x-1)}\dd x$
\end{enumerate}
\end{multicols}

\item 求下列定积分：
\begin{multicols}{2}
    \begin{enumerate}
        \item  $\displaystyle\int_{-1}^{3} \frac{\dd x}{x^{2}+7 x+10} $
        \item  $\displaystyle\int_{\sqrt{3}}^{0} \frac{\dd x}{2+x^{2}}$
        \item  $\displaystyle\int_{0}^{1} x \sqrt{\left(2-x^{2}\right)^{3}} \dd x$
        \item  $\displaystyle\int_{0}^{\ln 3} \frac{\dd x}{e^{x}+5+6 e^{-x}}$
        \item  $\displaystyle\int_{1}^{\tfrac{7}{5}} \frac{x^{\tfrac{3}{4}}+x^{\tfrac{1}{4}}}{x^{\tfrac{9}{4}}-x^{\tfrac{5}{4}}+x^{\tfrac{1}{4}}} \dd x$
        \item  $\displaystyle\int_{1}^{2} \frac{2-x+x^{2}}{(1+x)(1-x)^{2}} \dd x$
        \item  $\displaystyle\int_{1}^{2} x \ln x \dd x$
        \item  $\displaystyle\int_{-\pi}^{\pi} \sin m x \sin n x \dd x$
        (此处 $m, n$ 是整数且 $m \pm n\ne 0$)
\item  $\displaystyle\int_{\tfrac{a}{\sqrt{3}}}^{a} \frac{\dd x}{\left(a^{2}+x^{2}\right)^{3 / 2}}$
\end{enumerate}
\end{multicols}
\item  求下列极限：
    \begin{enumerate}
\item  $\displaystyle\lim _{n \to  \infty} \frac{1}{n}\left\{\left(1+
\frac{1}{n}\right)^2+\left(1+\frac{2}{n}\right)^2+\cdots+
\left(1+\frac{n}{n}\right)^{2}\right\}$
\item $\displaystyle\lim _{n \to  \infty} \frac{1^{4}+2^{4}+\cdots+n^{4}}{n^{5}}$
\item $\Lim_{n\to\infty}\frac{1}{n}\left(\sin\frac{\pi}{n}+\sin\frac{2\pi}{n}+\cdots+\sin\frac{n\pi}{n}\right)$
\item $\Lim_{n\to\infty}\left(\frac{1}{\sqrt{n+1}}+\frac{1}{\sqrt{n+2}}+\cdots+\frac{1}{\sqrt{n+n}}\right)\cdot \frac{1}{\sqrt{n}}$
\end{enumerate}

\item 利用求$y=\frac{1}{x}$的定积分，证明
\[\ln(n+1)<1+\frac{1}{2}+\frac{1}{3}+\cdots+\frac{1}{n}<1+\ln n.\]
\item 设$0<a\le 1$，求证：
\[\int^a_0\frac{\dd x}{x^3+1}<\int^a_0\frac{\dd x}{x^4+1}. \]
\item 设$0\le x\le \frac{\pi}{4}$，利用不等式$0\le \sin x\le x$证明：
\[\frac{\pi}{4}<\int^{\pi/4}_0\frac{\dd x}{\sqrt{1-\sin x}}<2-\sqrt{4-\pi}.\]

\item 若$m,n$是正整数，求证：
\[\int^1_0 x^m (1-x)^n\dd x=\frac{n}{m+1}\int^1_0x^{m+1}(1-x)^{n-1}\dd x. \]
\item 求证：$\displaystyle\int^\pi_0 xf(\sin x)\dd x=\frac{\pi}{2}\int^\pi_0f(\sin x)\dd x$.
\item 若$I_n=\displaystyle\int^{\pi/4}_0 \tan^n \theta \dd \theta\quad (n\ge 2)$. 试求$I_n$和$I_{n-2}$的关系式．
\item 若$I_n=\displaystyle\int \frac{\dd x}{(x^2+a^2)^n}$，求证：
\[2na^2I_{n+1}=\frac{x}{(x^2+a^2)^n}+(2n-1)I_n\]
\item 设$S_n=\frac{1}{\sqrt{1}}+\frac{1}{\sqrt{2}}+\cdots+\frac{1}{\sqrt{n}}$，$n$为正整数．
\[T_n=\frac{1}{\sqrt{1+\frac{1}{2}}}+\frac{1}{\sqrt{2+\frac{1}{2}}}+\cdots+\frac{1}{\sqrt{n+\frac{1}{2}}}.\]
求证：
\begin{multicols}{2}
\begin{enumerate}
    \item $S_n\ge \displaystyle\int^{n+1}_1 \frac{\dd x}{\sqrt{x}}$
    \item $\Lim_{n\to\infty}\frac{1}{S_n}=0$
    \item $S_{n+1}-1\le T_n\le S_n$
    \item $\Lim_{n\to\infty}\frac{T_n}{S_n}=1$
\end{enumerate}
\end{multicols}

\item 设$f(x)=(ax+b)e^{x/2}$, 求$a,b$的值，使得等式
\[f (x) =e^{x/2}-1+\frac{1}{2}\int^x_0 f (t) \dd t\]
对于任何$x$的值都成立．
\item 求由曲线$y=\frac{4}{x^2}$, $y=x-1$, $x=1$所围区域的面积．
\item 求在曲线$x^2+y^2=9$和$y=\frac{1}{3}(-x^2+1)$之间位于
第二象限那部分的面积．
\item 求积分$I(a)=\displaystyle\int^1_{-1}|x-a|e^x\dd x$在$|a|\le 1$时的最大值．
\item 设$0<t<\pi$，考察$x$的函数$f(x)=2\sin x\sin(t-x)$
\begin{enumerate}
    \item 画出$f(x)$在区间$0\le x\le \pi$上的图象，再设
    \[F (t) =\int^\pi_0 |f (x)|\dd x.\]
根据所作的图象说明$F(t)$所表示的是什么？
\item 用关于$t$的式子写出$F(t)$, 并求出使$F(t)$取最小值时$t$的值．
\end{enumerate}
\item  设$T_1$是由抛物线$y=4x^2$和直线$x=a$, $x=1$, $y=0$所
围成的区域，$T_2$是抛物线$y=4x^2$和二直线$x=a$, $y=0$围成的区域，但是$0<a<1$，试求：
\begin{enumerate}
\item $T_1$绕$x$轴旋转而成的旋转体体积$V_1$;    \item $T_2$绕$y$轴旋转而成的旋转体体积$V_2$;    \item 使$V_1+V_2$为最大时$a$的值．
\end{enumerate}
\item 设$0\le t\le \frac{\pi}{2}$，曲线$y=\sin x$与三条直线$x=t$, $x=2t$, 
$y=0$围成部分绕$x$轴旋转所成的旋转体的体积为$V(t)$．当$t=\alpha$时$V(t)$最大，试求$\cos\alpha$的值．
\item 设曲线$y=f(x)$通过点$P(3, 4)$, 并满足微分方程
$\frac{\dd y}{\dd x}=\frac{y}{x}$，
求此积分曲线．

\item 若$\frac{\dd y}{\dd x}=\cos(x+y)$，求不定积分．（提示：设$z=x+y$）
\end{enumerate}

\fi

\end{document}
